\documentclass[UTF8]{ctexart}
\usepackage[paper=a4paper,left=2.5cm,right=2.5cm,top=3cm,bottom=3cm]{geometry}
\usepackage{hyperref}
\usepackage{amsmath}
\usepackage{amsthm}
\usepackage{amssymb}
\usepackage{mathrsfs}
\usepackage{graphicx}
\usepackage{xcolor}
\usepackage{tikz}
\usetikzlibrary{calc,positioning,shapes.geometric}
\usepackage{tikz-cd}
\usepackage{longtable}
\usepackage{enumitem}
\providecommand{\xleftrightarrow}[2][]{\mathrel{\overset{#2}{\underset{#1}{\longleftrightarrow}}}}
\DeclareMathOperator{\cosupp}{cosupp}

\definecolor{coverNavy}{RGB}{20,35,65}
\definecolor{coverBlue}{RGB}{48,103,158}
\definecolor{coverGold}{RGB}{218,170,70}
\definecolor{coverPaper}{RGB}{247,244,235}
\newcommand{\statementlabel}[1]{\par\medskip\noindent\textbf{#1}\quad}

\theoremstyle{plain}
\newtheorem{definition}{定义}[subsection]
\newtheorem{lemma}{引理}[subsection]
\newtheorem{theorem}{定理}[subsection]
\newtheorem{corollary}{推论}[subsection]
\newtheorem{remark}{注记}[subsection]
\newtheorem{example}{例}[subsection]
\newtheorem{proposition}{命题}[subsection]

\makeatletter
\renewcommand{\proofname}{\textbf{证明}}
\renewenvironment{proof}[1][\proofname]{\par
  \pushQED{\qed}%
  \normalfont
  \topsep6\p@\@plus6\p@\relax
  \trivlist
  \item[\hskip\labelsep
        \normalfont\bfseries
    #1\quad]\ignorespaces
}{%
  \popQED\endtrivlist\@endpefalse
}
\makeatother
\renewcommand{\qedsymbol}{\ensuremath{\square}}

\ctexset{
    section = {
        name = {第,章},
        number = \arabic{section},
        format = \Large\bfseries\centering,
        aftername = \quad
    },
    subsection = {
        format = \large\bfseries
    }
}

\begin{document}

\begin{titlepage}
\thispagestyle{empty}
\begin{tikzpicture}[remember picture,overlay]

\fill[coverPaper] (current page.south west)
  rectangle (current page.north east);

\fill[coverNavy] (current page.north west)
  rectangle ($(current page.north east)+(0,-5.7cm)$);

\foreach \r/\op in {2.35/0.18,1.72/0.30,1.08/0.46} {
  \draw[coverGold,opacity=\op,line width=1.25pt]
    ($(current page.north east)+(-3.05cm,-2.85cm)$)
    circle (\r cm);
}

\node[
  anchor=west,
  text=white,
  font=\bfseries\fontsize{35}{42}\selectfont
] at ($(current page.north west)+(1.45cm,-2.0cm)$)
{流形上的层};

\node[
  anchor=west,
  text=coverGold,
  font=\fontsize{15}{20}\selectfont
] at ($(current page.north west)+(1.5cm,-3.15cm)$)
{层论与微局部分析讲义};

\node[
  text=coverNavy,
  font=\fontsize{23}{30}\selectfont
] at ($(current page.center)+(0,1.15cm)$)
{$\mathrm{SS}(F)\subset T^{*}X$};

\node[
  text=coverBlue,
  font=\large
] at ($(current page.center)+(0,0.25cm)$)
{导出范畴 · 六种运算 · 微支撑 · 反常层};

\node[
  draw=coverNavy,
  fill=white,
  rounded corners=4pt,
  minimum width=4.7cm,
  minimum height=1.35cm,
  font=\bfseries\Large,
  text=coverNavy
] at ($(current page.center)+(-3.1cm,-1.8cm)$)
{导出范畴};

\node[
  draw=coverNavy,
  fill=white,
  rounded corners=4pt,
  minimum width=4.7cm,
  minimum height=1.35cm,
  font=\bfseries\Large,
  text=coverNavy
] at ($(current page.center)+(3.1cm,-1.8cm)$)
{六种运算};

\node[
  draw=coverBlue,
  fill=white,
  rounded corners=4pt,
  minimum width=4.7cm,
  minimum height=1.35cm,
  font=\bfseries\Large,
  text=coverBlue
] at ($(current page.center)+(-3.1cm,-3.65cm)$)
{微局部化};

\node[
  draw=coverBlue,
  fill=white,
  rounded corners=4pt,
  minimum width=4.7cm,
  minimum height=1.35cm,
  font=\bfseries\Large,
  text=coverBlue
] at ($(current page.center)+(3.1cm,-3.65cm)$)
{反常层};

\draw[coverGold, line width=1pt, opacity=.85]
  ($(current page.center)+(-.75cm,-1.8cm)$) --
  ($(current page.center)+(.75cm,-1.8cm)$);
\draw[coverGold, line width=1pt, opacity=.85]
  ($(current page.center)+(-.75cm,-3.65cm)$) --
  ($(current page.center)+(.75cm,-3.65cm)$);
\draw[coverGold, line width=1pt, opacity=.85]
  ($(current page.center)+(-3.1cm,-2.5cm)$) --
  ($(current page.center)+(-3.1cm,-2.95cm)$);
\draw[coverGold, line width=1pt, opacity=.85]
  ($(current page.center)+(3.1cm,-2.5cm)$) --
  ($(current page.center)+(3.1cm,-2.95cm)$);

\coordinate (base) at ($(current page.south west)+(3.45cm,3.85cm)$);
\coordinate (uone) at ($(base)+(90:1.25cm)$);
\coordinate (utwo) at ($(base)+(18:1.25cm)$);
\coordinate (uthree) at ($(base)+(306:1.25cm)$);
\coordinate (ufour) at ($(base)+(234:1.25cm)$);
\coordinate (ufive) at ($(base)+(162:1.25cm)$);

\draw[coverNavy,line width=1pt]
  (uone) -- (utwo) -- (uthree) -- (ufour) -- (ufive) -- cycle;
\draw[coverGold,line width=.9pt,opacity=.85]
  (base) -- (uone)
  (base) -- (utwo)
  (base) -- (uthree)
  (base) -- (ufour)
  (base) -- (ufive);

\foreach \point in {uone,utwo,uthree,ufour,ufive} {
  \fill[coverNavy] (\point) circle (3pt);
}
\fill[coverGold] (base) circle (3pt);

\node[coverNavy,font=\small] at ($(uone)+(0,0.35cm)$) {$U_1$};
\node[coverNavy,font=\small] at ($(utwo)+(0.42cm,0)$) {$U_2$};
\node[coverNavy,font=\small] at ($(uthree)+(0.42cm,0)$) {$U_3$};
\node[coverNavy,font=\small] at ($(ufour)+(-0.42cm,0)$) {$U_4$};
\node[coverNavy,font=\small] at ($(ufive)+(-0.42cm,0)$) {$U_5$};

\node[
  anchor=south east,
  text=coverNavy!78,
  font=\large,
  align=right
] at ($(current page.south east)+(-1.35cm,1.6cm)$) {
  $Rf_*,\ f^{-1},\ Rf_!,\ f^!$\\[5pt]
  $F^\wedge=\mathrm{Fourier}\text{-}\mathrm{Sato}(F)$\\[5pt]
  $\mu_M(F)$\\[5pt]
  ${}^{p}\mathsf{D}^{\leq 0}(X)$
};

\draw[coverGold,line width=1.1pt]
  ($(current page.south west)+(8.1cm,2.0cm)$) --
  ($(current page.south west)+(13.7cm,2.0cm)$);

\node[
  anchor=west,
  text=coverNavy,
  font=\large
] at ($(current page.south west)+(8.1cm,2.35cm)$)
{整理：Yi Wang};

\node[
  anchor=west,
  text=coverNavy,
  font=\large
] at ($(current page.south west)+(8.1cm,1.35cm)$)
{日期：\today};

\end{tikzpicture}
\null
\end{titlepage}

\newpage
\hypersetup{hidelinks}
\tableofcontents
\newpage

本书的目的是对层理论进行独立的阐述。

层是由让·勒雷 (Jean Leray) 在上次世界大战期间制造的，当时他是德国战俘营的一名战俘。

层理论的目的非常普遍：它是从局部信息中获取全局信息，或者定义“障碍”，其特征是局部属性不再全局成立：例如，流形并不总是可定向的，或者微分方程可以局部可解，但不能全局求解。因此，层理论是代数拓扑（例如奇异同源）的一部分的广泛概括，其对应于恒定层，或更一般地，对应于局部常值层。层有许多自然的例子，例如定向层、可微分或全纯函数的层、微分方程组解的层、作为直像获得的可构造层等。

然而，一开始并不清楚这样的一般理论是否有任何应用，直到它被成功地应用于几个复变函数的理论，并且可以想象，如果没有嘉当、塞尔和后来的格罗滕迪克在五十年代发展的实质性工作，勒雷的原始著作将很难被理解（参见胡泽尔的《简史》，见下文）。

当与同调代数工具相结合时，层理论发挥出其全部力量。事实上，勒雷还引入了谱序列的概念，它与嘉当-艾伦伯格的派生函子一起，自然地引出了格罗腾迪克的导出范畴理论。在长期保留给一些代数几何专家之后，导出范畴理论开始被充分认可为数学的基本工具。特别是，如果没有它，就不可能像六十年代的 Verdier 那样对庞加莱的对偶性给出如此美丽的概括（在格洛腾迪克在 étale 上同调框架中的相关工作之后），或者系统地处理格洛腾迪克所谓的层“六种运算”，即函子  $ Rf_*, f^{-1}, Rf_1, f^1, \otimes^L $  和  $ R\mathcal{H}om $  。正如我们将看到的，这种形式主义导致了深刻而有力的公式，这些公式在一般背景下解释了经典结果。我们已经提到了庞加莱对偶性，但还有许多其他主题，例如莱夫谢茨不动点公式或欧拉-庞加莱指数。当然，这些函子是函数经典运算的抽象版本：用于积分的正像、用于合成的逆像、张量

产品为产品。 （这三个运算通过“对偶性”给出了六种运算，这种运算在函数中没有对应项，但我们在这里将针对可构造函数引入它。）

在这个阶段，我们已经简单地解释了人们可以粗略地称为“经典层理论”的东西，但是佐藤干夫在1969年出现了一个新的、基本的想法，它为层理论提供了一个新的视角，即“微局部观点”，而在这个新框架内发展层理论也是本书的目的之一。佐藤的主要兴趣是研究线性微分方程组解的解析奇异性。早在 1959 年，他就使用局部上同调来定义超函数束，并将它们解释为全纯函数的边界值之和。十年后，他引入了一系列微函数，以识别“边界值来自哪个方向”。为了实现这一点，佐藤引入了特化函子 $ v_M $ （沿着流形 $ X $ 的子流形 $ M $ ），及其傅立叶变换微局部化的函子 $ \mu_M $ 。这些函子将  $ X $  上的层导出范畴分别发送到  $ X $  中的  $ M $  的法向束和共形束上的层导出范畴，它们使我们能够精确分析  $ M $  邻域上的层范畴，同时考虑到  $ M $  的所有法向（或共形）方向。

当试图将佐藤的理论应用于微双曲系统的研究时，作者逐渐意识到他们使用的唯一信息是系统的特征变化（在复流形X的余切丛中），并且有可能忘记X的复结构，甚至忘记主题是偏微分方程的事实。他们所做的只不过是在“非禁止”方向上（即穿过非特征真实超曲面）（系统的全纯解层复合体）的非特征变形。请注意，这些非特征变形技术以前已经出现过，但这里作者处理的是微微分系统，并且确实需要微局部几何。特别是，他们引入了 $ \gamma $ 拓扑，一种微局域截止。后来，在 1982 年，通过抽象他们之前关于微双曲系统的工作，他们引入了束微支撑的概念：粗略地说，如果 F 不具有由余常态接近 p 的半空间支持的上同调，则实流形 X 的余切丛  $ T^*X $  的点 p 不属于 SS(F)，即束 F 的微支撑。这个新定义使我们能够“微观局部”研究层，特别是使接触变换（余切束上的自然变换）作用于层，类似于对 Sato-Kawai-Kashiwara [1] 中的微函数或 Hormander [2] 中的傅里叶分布进行量化接触变换。

微支撑的思想与莫尔斯理论密切相关。众所周知，如果  $\phi$  是紧实流形  $X$  上的实函数，则只要  $t$  不满足  $\phi$  的临界值，空间  $\{x \in X; \phi(x) < t\}$  的上同调群就是同构的。如果考虑  $X$  上的束  $F$  的上同调群的类似问题，则可以在帮助下获得相应的结果

F 的微支撑。对于复流形，可构造层的微支撑也可以使用 Grothendieck-Deligne 的消失循环函子来描述。请注意，消失循环函子是佐藤微局部化函子的特例。

微支撑具有深刻的几何意义，我们要证明这个集合是对合的。这特别给出了微分方程组特征簇的经典相应结果的纯实数和几何证明。

正如我们将在本书中看到的那样，关于层的微局域观点加深了层理论并带来了许多应用。我们在此仅讨论几个例子。

(a) 层论中的许多态射在满足某些微局域条件的情况下变成同构。例如，考虑流形  $ f: Y \to X $  的态射，并令 F 是 X 上的一个层（或者更好，一个层的复数）。那么存在一个规范态射

\[
\omega_{Y/X}\otimes f^{-1}F\to f^!F\ ,
\]

 $ \left(\omega_{Y/X}\right. $  是相对对偶复形）。众所周知，如果 f 是光滑的，这个态射就是同构，但实际上还有一个更强的结果：只要“f 对于 F 是非特征的”，这个态射就是同构； （如果 f 是闭嵌入，这意味着 SS(F) 与 X 中 Y 的共法丛的交集包含在零部分中）。

(b) 令 $X$  为复流形， $\mathcal{M}$  为 $X$  上的线性微分方程组，即左相干 $\mathcal{D}_X$  模，其中 $\mathcal{D}_X$  表示 $X$  上的全纯微分算子环层。让

\[
\mathrm{F}=\mathrm{R}\mathcal{H o m}_{\mathfrak{L}_{x}}(\mathcal{M},\mathcal{O}_{X})
\]

是系统全纯解的层复形。具有解析系数的线性偏微分方程的几何情况可以用下面的公式简洁地表达，其证明基本上依赖于柯西-科瓦列夫斯基定理：

\[
\mathrm{SS}(\boldsymbol{F})=\operatorname{char}(\mathcal{M})~.
\]

这里 $ \operatorname{char}(\mathcal{M}) $ 表示系统的特征簇。

如果将（a）的结果应用于 Y 是实解析流形的情况，则我们获得了经典分析结果的纯层理论证明。

(c) 如果存在 X 的亚解析分层，使得所有上同调群在层上局部常值，则实解析流形上的层 F 的复形是弱可构造的。这一条件将被证明等价于微局部条件，即 F 的微支撑是一个子

 $ T^*X $  的解析拉格朗日子集。此外，如果满足一定的有限性条件，可以将 $ \mathrm{SS}(F) $ 支持的拉格朗日环与 $ F $ 关联起来，并得到 $ F $ 在 $ X $ 上的全局欧拉-庞加莱指数，作为该环与自然关联到零部分的环的交数。通过这个结果，我们有了一个有效的微局部工具来计算索引。

在本书中，我们希望让读者相信这一观点至关重要。从头开始（导出范畴），我们将介绍经典理论（“六操作”）和微局域理论（微局域化、微支撑、接触变换）。然后我们将应用该机制来研究实流形上的可构造层，最后简要讨论其在线性偏微分方程中的应用。

更详细地说，本书的内容如下。

第一章包含了关于同调代数的基本事实，这些知识对于本书的其余部分，特别是导出范畴的理论来说是必要的（t 结构的概念除外，推迟到第十章）。

第二章和第三章包含导出范畴语言中关于层的“经典”概念，包括六种运算，以及傅立叶-佐藤变换，该变换交换矢量丛上的层（更准确地说，层导出范畴的对象）和对偶向量丛上的层。

第四章专门讨论微局部化。在回顾了流形  $ X $  中子流形  $ M $  的法向变形的几何构造之后，我们定义了特化函子  $ \nu_M $  ，它将  $ X $  上的层发送到法向束  $ T_M X $  上的层，及其傅里叶-佐藤变换，即微局部化函子  $ \mu_M $  。我们还定义了  $ \mu_M $  、函子  $ \mu hom $  的自然泛化，并研究所有这些函子的函子性质。

第五章介绍了层的微支撑。在根据其微支撑的几何形状证明层的全局扩展定理后，我们利用 $ \gamma $ 拓扑来“微局部”切断层。然后我们研究非特征情况（对于逆像）或适当情况（对于直像）的函子运算下微支撑的行为。作为应用，获得了层的莫尔斯不等式。

在第六章中，我们使用微支撑来定位层 $ \mathbf{D}^b(X) $ 的导出范畴相对于 $ T^*X $ 的子集 $ \Omega $ （获得新的三角范畴 $ \mathbf{D}^b(X;\Omega) $ ），并定义“微局部”逆像或直像。接下来，我们通过研究微支撑相对于函子运算的行为，将第五章的结果扩展到一般情况。特别是，我们证明 $ \mu_M(F) $ 的微支撑包含在 $ \mathrm{SS}(F) $ 沿 $ T_M^*X $ 的法向锥体中。这个不等式是微双曲系统定理的层理论版本，并将在整本书中使用。本章还包含一个重要的结果，即微支撑的对合性定理。

在第七章中，我们对层进行接触变换。如果 $\chi$ 是 $T^*X$ 和 $T^*Y$ 的两个开子集 $\Omega_X$ 和 $\Omega_Y$ 之间的接触变换，那么在适当的条件下，可以构造一个同构 $\mathbf{D}^b(X; \Omega_X) \simeq\mathbf{D}^{b}(Y;\Omega_{Y})$  ，并且这种同构与  $\mu\text{hom}$  兼容。当计算  $X$  的子流形  $M$  上的常值层  $A_{M}$  的图像时，人们自然会想到沿光滑拉格朗日流形的纯束的概念。对于某些  $A$  -模  $L$  和某些移位  $d$  ，纯束“一般”（和微局部）同构于某些束  $L_{M}\left[d\right]$  ，但移位的计算需要拉格朗日平面三重态的惯性指数的整个机制。在本章中，我们特别计算核的微局部组成的变化。

第八章详细研究了实流形上的可构造层。我们引入 $\mu$ 分层的（微局部）概念，然后证明当且仅当其微支撑是亚解析和各向同性的（因此是拉格朗日）时，层是弱可构造的。然后我们可以应用前面的结果来研究可构造层的函子运算。在复流形上，我们证明如果层在实底层流形上是可构造的，而且如果它的微支撑在  $C^{*}$  的作用下是不变的。由此，我们推导出非真直像的定理。最后，我们证明邻近循环和消失循环函子是特化和微局部化函子的特例。

第九章在对偶复形的帮助下引入了亚解析链和循环的概念。然后，使用函子  $\mu\text{hom}$  ，我们将其特征循环与可构造层  $F$  相关联，并证明该拉格朗日循环与  $T^*X$  的零部分的交集数给出  $F$  在  $X$  上的全局欧拉-庞加莱指数。我们还计算局部欧拉-庞加莱指数，并将拉格朗日循环与  $X$  上的可构造函数联系起来，从而获得这些函数的新微积分。在本章中，我们还给出了可构造层的 Lefschetz 不动点公式。

第十章发展了反常层理论。在回顾了 t 结构的概念之后，我们定义了实流形上的反常层（即反常复合体），并表明它们形成了一个阿贝尔范畴。然后我们研究复流形上的反常层，给出反常性的微观局部特征，并证明反常性通过各种函子运算得以保留。

在第十一章中，我们简要介绍了如何将层理论应用到线性偏微分方程组的研究中。在对 $ \mathcal{O}_X $ 和 $ \mathcal{D}_X $ 模的理论进行简短回顾之后，我们证明了(i.3)中的一个包含，并推断出完整 $ \mathcal{D}_X $ 模的全纯解的复形是反常的。我们还介绍了超函数和微函数的层，并从（i.3）推导出椭圆或双曲系统微函数解的一些基本结果。在本章中，我们还对层  $ \mathcal{O}_X $  进行量化接触变换。这个结果有很多重要的应用，这里不予讨论。

我们以一个简短的附录来结束这本书，其中我们收集了辛几何，特别是惯性指数所需的所有结果（带有一些证明）。

每章开头都有一个简短的介绍，并包括不同难度的练习。其中一些练习（尤其是第一章）是本书课程中使用的辅助结果。一般来说，此类练习的证明很简单，否则会给出提示。

我们尽可能避免在文本中提供参考文献。相反，我们选择以一些历史评论来结束每一章。原因是，大多数时候，一个定理有着漫长而复杂的历史，每次都引用对结果做出贡献的每个人是很乏味的。另一方面，仅引用发起该主题或将其付诸最终形式的人似乎是不合适的。

层理论的起源和开端相当复杂，本书受益于 Christian Houzel 的历史著作，他同意对这部分历史进行详细描述。

\section*{简史：层理论的开端}
\addcontentsline{toc}{section}{简史：层理论的开端}
作者：克里斯蒂安·乌泽尔

\subsection*{勒雷的课程（1945）}
让·勒雷 (Jean Leray) 在奥特里什 (Autriche) 的第十七军团战俘监狱中，为囚禁大学的代数拓扑学课程提供了组织者的帮助。这是 1934 年在 J. Schauder 的儿子文章中找到的，关于应用程序度的无限概念和 Brouwer 固定点的理论[33]。 Leray 需要这样一个函数空间中的定理，以获得流体力学中出现的非线性方程解的存在性（这些解未必正则，也未必唯一）。

Leray 的课程于 1945 年战争结束时发表在《Journal de Liouville》上 [29]。代数拓扑是在新的基础上发展起来的，避免了局部可定向性或线性的假设以及细分或简单逼近的方法。重点是上同调，它只是在战前 [47] 才与同调明确区分开来，特别是在 de Rham 的工作之后 [37]；环中系数空间的上同调总是具有乘法结构，而我们在紧可定向簇的情况下使用的同调中的乘法结构是通过庞加莱对偶性从上同调推导出来的。当谈到同调时，勒雷将上同调重新命名为“同调”，并谈到“贝蒂群”。为了定义具有环 A 中的系数的拓扑空间 E 的上同调，它受到 Čech 过程 [9] 的启发，但它用更适合代数拓扑的概念替代了覆盖集的概念：覆盖的概念。为了定义一个覆盖，我们首先给自己一个“抽象复形”，一个对应于各种维度 p 的有限类型的自由交换群序列，并提供基数 (  $ X^{pa} $  )ₐ，以及维  $ p + 1 $  群的协同边运算符  $ X^{pa} \mapsto \dot{X}^{pa} $  元素的数据（我们通过线性扩展到该群的其他元素），遵循协同边的协同边为零的公理。我们通过将每个 $ X^{pa} $ 与支持 $ |X^{pa}| $ （E的非空部分）相关联来使复杂的“具体”；我们强加这样的公理：如果 $ X^{q\beta} $  遵循 $ X^{pa} $ （也就是说，通过基本元素的有限序列与其连接，每个基本元素都介入前一个元素的共边界）， $ X^{q\beta} $  的支持包含在 $ X^{pa} $  的支持中。如果支撑封闭，则具体复形 K 是一个覆盖，对于 E 的任何点 x，其支撑包含 x 的元素形成的子复形是单纯形

（它的上同调是微不足道的）并且维度 0 的元素的和  $ K^{0} $  是一个余循环（称为单位（余）循环）。  $ E $  与 $ A $  中的系数的上同调类是 $ E $  的任意覆盖 $ K $  的形式 $ L^{p} $  的上同调类（ $ K $  的基本元素 $ p $  中的 $ A $  中的系数的线性组合），同意在任何时候将 $ L^{p} $  与“交集” $ L^{p}K^{0} $  一起识别 $ K' $  是另一个封面（我们使用复张量积定义交集； $ X^{pa} \otimes X^{q\beta} $  的支持度是 $ X^{pa} $  和 $ X^{q\beta} $  的支持度的交集，我们通过取消空支持元素来传递商）。在正规空间的情况下，我们可以通过仅取稳定族的交集的覆盖来计算上同调，并且对于  $ E $  的任何有限开覆盖​​  $ \rho $  ，包含其支撑为“ $ \rho $  -small”的覆盖；在紧凑空间的情况下，如果其支撑是“简单的”（也就是说上同调平凡的），我们可以对单个覆盖感到满意。 Leray 将 Hopf [23] 在某些紧致可定向流形上的结果扩展到紧致拓扑空间的情况。他发展了对偶理论，使我们能够恢复贝蒂群。他的课程的第一部分以介绍  $ E $  本身的连续应用的 Lefschetz 数结束，在  $ E $  是紧连通的并且允许有限“凸”覆盖的情况下（也就是说，通过交集为空或简单的简单闭包）。

在第二部分中，Leray 将  $\xi:E\to E$  的 Lefschetz 数与  $\xi$  限制为  $\xi$  封闭的稳定数进行比较。它引入了“伪循环”（ $E$  的紧凑  $B$  部分的上同调的射影极限的元素）以扩展到先前在紧凑情况下演示的非紧凑情况结果。他从可微流形的细胞分解中定义了覆盖（“瓦片”），并在可定向和无边的情况下建立了庞加莱对偶性。第三部分是Leray在代数拓扑方面的工作的起源，定义了 $E$ 的开 $O$ 中方程 $x=\xi(x)$ 的解的总索引 $i(O)$ （其中 $\xi:F\to E$ 是在包含 $\overline{O}$ 的闭 $F$ 中定义的连续映射）；我们假设 $E$ “凸形”（紧凑连接并允许简单封闭对象的重叠，其有限交集为空或简单，其内部构成拓扑的基础），并且我们考虑具有整数系数的上同调。如果所考虑的方程在  $O$  的边界上无解，则定义总索引  $i(O)$  ，并且它仅取决于  $\xi$  到  $\overline{O}$  的限制；它是同伦不变的（在  $\xi$  上），并且在  $O=E$  的情况下等于 Lefschetz 数。如果  $O$  中方程的所有解都属于  $O$  中包含的不相交开放空间  $O_\alpha$  的并集，则  $i(O)$  是  $i(O_\alpha)$  的总和。对于孤立解 $x$ ，我们将其索引定义为总索引 $i(V)$ ，其中 $V$ 是 $x$ 的一个相当小的邻域，如果 $O$ 仅包含孤立解，则 $i(O)$ 是这些解的索引之和。通过确定方程的任何解都是孤立的并且指数为 1，在高电流电阻尾流理论中获得了唯一性定理。Leray 仍然为稍微更一般形式的方程定义了指数，并将他的理论应用于 Fredholm 方程。它还考虑 $x=\xi(x,x')$ 形式的方程的情况，其中 $\xi:F\to E$ 是一个应用程序

在 $E \times E'$  的闭合 $F$  中定义的连续（ $E$  假设简单凸面）；索引被  $Z^p$  的  $E'$  上的投影操作替换。   $O$  、 $Z^pE \times E'$  的伪循环。

\subsection*{层理论和谱序列}
勒雷在被囚禁期间对代数拓扑的研究使他走得更远。对纤维空间的同调性与其基底和纤维的同调性之间的关系的研究使得有必要引入新的工具：局部系数的上同调，从一点到另一点都不同；通过一系列近似计算上同调。更一般地说，在连续  $ \xi: X \to Y $  应用程序的情况下，这些工具将用于研究 X 与 Y 的上同调以及  $ \xi $  的纤维； Picard 使用这种类型的研究来确定复杂代数曲面的同源性 [36]，而 Lefschetz 已将这种方法扩展到更高的维度 [28]。局部系数的想法也以局部系数系统的特定形式独立出现于 N. Steenrod [41]。

在[31]中，再现了他 1947-48 年和 1949-50 年在法兰西学院的课程，Leray 在局部紧拓扑空间上调用了层  $ \mathcal{B}(F_1) $  。它强制  $ \mathcal{B}(\emptyset) = 0 $  和节运算的传递性：  $ F_2(F_1 b) = F_2 b $  if  $ F_2 \subset F_1 \subset F $  。如果  $ \mathcal{B}(W) $ （ $ \infty $  的 W 闭邻域）的归纳极限为零，并且如果对于所有闭 F， $ \mathcal{B}(F) $  是  $ \mathcal{B}(V) $  的归纳极限，其中 V 穿过  $ F \cup \infty $  的闭邻域；如果满足这两个条件中的第一个条件，第二个条件满足非紧 F，并且如果  $ \mathcal{B}(K) $  是  $ \mathcal{B}(V) $  的归纳极限，则称它是干净的，V 是紧 K 的 K 的闭邻域。除了层之外，Leray 还引入了复合体和覆盖，就像在他 1945 年的课程中一样；但是，根据 Henri Cartan [5] 的建议，复数不再需要有基，而是提供乘法：它们是微分环，具有与每个元素 k 关联的 X 的闭合 S(k) 定律，称为元素的支持度。给定 X 上的复数  $ \mathcal{H} $  和 F 的闭合 F，我们因此获得层  $ \mathcal{B}: F \mapsto F\mathcal{H} $ ，据说与  $ \mathcal{H} $  相关。复数  $ \mathcal{H} $  是超上同调的覆盖，其中  $ \mathcal{B} $  提供了使其成为微分束的导数，Leray 使用精细覆盖而不是像 1945 年课程中那样使用所有覆盖；如果对于紧凑化  $ X \cup \infty $  的任何有限开覆盖​​（ $ V_v $ ），我们可以分解，那么复数  $ \mathcal{H} $  就可以了

将  $ \mathcal{H} $  的相同自同构转化为线性映射  $ \lambda_v $ :  $ \mathcal{H} \to \mathcal{H} $  的总和，使得  $ S(\lambda_v k) \subset \overline{V}_v \cap S(k) $  对于  $ \mathcal{H} $  的任何元素  $ k $  和任何索引  $ v $  。  $ \mathcal{C} $  éch 或 Alexander [1] 的构造保证了有限维局部紧空间 X 上精细覆盖的存在；如果 X 的维度为  $ n $  ，我们可以选择大于 n 的零精细覆盖。如果  $ \mathcal{X} $  是 X 的精细覆盖， $ \mathcal{H} $  是单元支撑上的适当微分层  $ \sum k_\mu \otimes b_\mu $  是点 x 的集合，使得  $ x \sum k_\mu \otimes b_\mu $  不为零（如果  $ x \in S(k) $  ，我们同意  $ x(k \otimes b) = xk \otimes xb $  ，否则我们将其设置为零）。复数 $ \mathcal{X} \circ \mathcal{H} $  是 $ \mathcal{X} \otimes \mathcal{H} $  与具有空支持的元素的理想的商（具有相同的支持定律）。 Leray 表明 $ \mathcal{X} \circ \mathcal{H} $  的上同调并不取决于精细覆盖 $ \mathcal{X} $  的选择，并且他注意到 $ \mathcal{H}(X \circ \mathcal{H}) $  这个上同调。

为了确定正常空间的上同调是通过交集稳定族的覆盖来计算的，并且其中对于每个开覆盖 ρ ，都有一个带有 ρ 小支撑的覆盖，Leray 在他 1945 年的课程中使用了一个基本引理，根据该引理  $K^{*}$  。如果  $K^{*}$  是覆盖且  $C'$  是复数，例如  $K^{*}$  ，则 C' 和 C' 具有相同的上同调。 e 是  $C$  的任何支持 e 的单纯形。通过分析这个引理的证明，他提出了谱环的概念，并引入了滤波环（整数滤波递减），他首先将其称为低估（术语“滤波”来自 H. Cartan [3]）。滤波后的差分环 A（参见 Koszul [26]）与分级差分环的序列 (  $\mathcal{H}_{r}A_{r}$  ) 相关联，其中  $\mathcal{H}_{r}A$  的差分度为  $r$  。它的定义如下：我们记下 $\mathcal{A}^{(p)}$ 过滤 $\geqslant p$ 的 $A$ 的元素集合，并放入 $\mathcal{C}^{p} = \mathcal{C} \cap \mathcal{A}^{(p)}$ 和 $\mathcal{D}^{p} = \mathcal{D} \cap \mathcal{A}^{(p)}$ ，其中 $\mathcal{C}$ 是 $A$ 和 $\mathcal{D}$ 的cobord的单元集合。然后，我们用 $\mathcal{C}_{r}^{p}$  指定 $a \in \mathcal{A}^{(p)}$  的集合，其共边界是过滤 $\geqslant p + r$ ，并用 $\mathcal{D}_{r}^{p}$  指定 $a \in \mathcal{C}_{r}^{p-r}$  的共边界 $\delta a$  的集合。这些组分别显示为 $\mathcal{C}^{p}$  和 $\mathcal{D}^{p}$  的阶近似 $r$ ；我们看到  $\mathcal{C}_{r}^{p}$  （或  $\mathcal{D}_{r}^{p}$  ）是  $r$  的减（或增）函数，并且它包含  $\mathcal{C}^{p}$  （或包含在  $\mathcal{D}^{p}$  中）。在度  $p$  中，  $\mathcal{H}_{r}A$  被定义为商  $\mathcal{C}_{r}^{p}/(\mathcal{C}_{r-1}^{p+1} + \mathcal{D}_{r-1}^{p})$  ，并且根据定义，该商中  $c_{r}^{p} \in \mathcal{C}_{r}^{p}$  的类  $h_{r}^{p}$  的共边界  $\delta r, h_{r}^{p}$  是  $\delta c_{r}^{p}$  mod 的类。   $(\mathcal{C}_{r-1}^{p+1} + \mathcal{D}_{r-1}^{p})$ 。我们验证 $\mathcal{H}_{r}A$ 的同源性与 $\mathcal{H}_{r+1}A$ 规范一致；因此，对于  $s \geqslant r$  ，每个  $\mathcal{H}_{r}A$  都标识为  $\mathcal{H}_{r}A$  的子商。对于度数  $p$  的刻度环  $\mathcal{H}_{\infty}A$  值  $\mathcal{C}^{p}/(\mathcal{C}^{p+1} + \mathcal{D}^{p})$  也是如此，它与与上同调环  $\mathcal{H}A$  相关的刻度环同构（提供来自  $A$  的过滤）。光谱环 (  $\mathcal{H}_{r}A$  ) 将该毕业标定为  $\mathcal{H}_{r}A$  的“归纳极限”的子环，并且当 $A$  的过滤有界更高时，确定为该极限本身；此外，如果  $\mathcal{H}_{r}A$  对于  $r > l$  的微分消失，则序列  $\mathcal{H}_{r}A$  对于  $r$  的这些值是平稳的，并且所研究的毕业生被识别为  $\mathcal{H}_{r}A$  。因此，我们看到，如果  $A$  的过滤在上面有界，并且  $\mathcal{H}_{r+1}A$  集中在  $0$  度上，则  $r > l$  的  $\mathcal{H}A$  与  $\mathcal{H}_{r}A$  同构。光谱环在功能上取决于  $A$  ；如果具有上限过滤的过滤环的同构  $\lambda: A' \to A$  会导致  $\mathcal{H}_{r+1}A$  的同构。

在  $ \mathcal{H}_{l+1}\mathcal{A} $  上，它引发了关于  $ \mathcal{H}\mathcal{A}^{\prime} $  对  $ \mathcal{H}\mathcal{A} $  过滤的同构。

环A仍然是滤波微分，Leray仍然考虑分级微分环 $ \mathcal{H} $ （具有1次微分）和滤波张量积 $ \mathcal{H}^l \otimes A $ ，其中 $ \mathcal{H} $ 的次 $ p $ 的元素被视为过滤 $ lp $ （ $ l $ 给定整数）。当  $ \mathcal{H} $  无扭转时  $ \mathcal{H}_{l+1}(\mathcal{H}^l \otimes A) $  与  $ \mathcal{H}^l \otimes \mathcal{H}_l A $  的上同构同构（规范地）；我们注意到，将 A 的微分替换为 0 并不会改变  $ \mathcal{H}_l(\mathcal{H}^l \otimes A) $ ；这允许使用嘉当的结果，在适当的假设下建立 Künneth 公式。

如果 $\mathcal{B}$ 是空间 $X$ 上的适当滤波差分层， $\mathcal{X}$ 是 $X$ 和 $l$ 的精细覆盖，是一个整数，则谱环 $(\mathcal{H}_r(X^l \circ \mathcal{B}))$ 和上同调 $\mathcal{H}(X^l \circ \mathcal{B})$ 不依赖于 $\mathcal{X}$ 的选择，并且Leray分别注意到它们 $(\mathcal{H}_r(X^l \circ \mathcal{B}))$ 和 $\mathcal{H}(X^l \circ \mathcal{B});\mathcal{H}_{l+1}(X^l \circ \mathcal{B})$  与  $X^l \circ \mathcal{F}_l(\mathcal{B})$  的上同调同构，其中  $(\mathcal{F}_r(\mathcal{B}))$  是与  $\mathcal{B}$  相关的谱束。当  $\mathcal{H}$  是一个真分级微分，具有下界度，具有度差  $>0$  ，并且对于所有  $x \in X$  ，  $\mathcal{H}(x)$  都集中在度  $0$  上，我们通过取  $l=1$  ，推导出  $\mathcal{H}(X \circ \mathcal{B})=\mathcal{H}(X \circ \mathcal{F}(\mathcal{B}))$  ，其中  $\mathcal{F}(\mathcal{B})$  是 $\mathcal{B}$ 。如果  $(F_\mu)$  是有限闭覆盖，并且如果  $\mathcal{H}^*$  是关联的（自由）Cech 复形，则假设  $\mathcal{H}(F \circ \mathcal{B})=\mathcal{H}(\mathcal{B}(F))$  对于  $F_\mu$  的所有  $F$  非空交集意味着  $\mathcal{H}(\mathcal{H}^* \otimes \mathcal{B})=\mathcal{H}(X \circ \mathcal{B})$  。

为了在局部紧凑空间之间连续应用 $\xi: X \to Y$ ，对于 $X$ 上的任何干净滤波差分层 $\mathcal{B}$ 以及任何整数对 $(l, m)$ 与 $l < m$ ，谱环 $\mathcal{H}_r(\xi^{-1}Y^m \circ X^l \circ \mathcal{B})$ ，通过精细覆盖 $X$ 的 $\mathcal{X}$ 和 $Y$ 的 $\mathcal{Y}$ 来定义；覆盖 $\xi^{-1}\mathcal{Y}$ 是 $\mathcal{Y}$ 与元素 $y$ 的理想的商，使得 $\xi^{-1}(S(y))$ 为空，并且 $\xi^{-1}\mathcal{Y}^m \circ \mathcal{X}^l$ 是 $X$ 的良好覆盖。因此， $\mathcal{H}(\xi^{-1}\mathcal{Y}^m \circ \mathcal{X}^l \circ \mathcal{B})$ 是 $\mathcal{H}(X \circ \mathcal{B})$ ，具有独立于 $\mathcal{X}$ 和 $\mathcal{Y}$ 的选择的某种过滤，并且其毕业生是 $\mathcal{H}_r(\xi^{-1}Y^m \circ X^l \circ \mathcal{B})$ 的归纳极限的子环，并且在有限维的适当假设下与该极限一致；我们有 $\mathcal{H}_{m+1}(\xi^{-1}Y^m \circ X^l \circ \mathcal{B}) = \mathcal{H}(Y^m \circ \xi \mathcal{F}_m(X^l \circ \mathcal{B}))$ ，其中 $\xi \mathcal{F}_r(X^l \circ \mathcal{B})$ 是与 $\mathcal{H} \circ \mathcal{B}$ 相关的光谱层的 $\xi$ 的图像（ $\xi$ 在 $X$ 上的层 $\mathcal{F}$ 的图像 $\xi \mathcal{F}$ 由 $\xi \mathcal{F}(F) = \mathcal{F}(\xi^{-1}F)$ 定义为 $Y$ 的所有闭合 $F$  ）。

文章 [30] 的结尾专门讨论了常值层（“与环相同的层”）和 Steenrod [41] 意义上的局部系数系统（“与环局部同构的束”）的特殊情况。文章 [31] 的结果在 1946 年的注释 [29] 中公布，内容于 1950 年在法兰西学院发表，该文章将一般理论应用于 Σ 是纤维 F 的纤维化的情况，并且我们考虑 X 与系数在（常数）环 A 中的上同调； Y 上的图像层 $ \mathcal{B} = \xi \mathcal{F}(X \circ A) $  与 $ \mathcal{H}(F \circ A) $  局部同构。 Leray 推导的部分结果是由 G. Hirsch 独立获得的[22]。当 F 与球体具有相同的同源性时，我们特别发现 Gysin [20] 的结果得到了 Chern 和 Spanier [10] 的补充；当Y与球体具有相同的同源性时，我们找到Wang的结果[44]。

\subsection*{嘉当研讨会和韦伊对德拉姆定理的证明}
Leray 的想法启发了 A. Weil 和 H. Cartan 致力于代数拓扑的革新。第一个在 1945 年与勒雷进行了一次谈话，并从中汲取了德拉姆定理现在经典证明的想法；这个演示的第一个版本是在 1947 年的一封信中传达给嘉当的[45]，完整的版本于 1952 年发表在 Commentarii Mathematici Helvetici 上[46]。 Weil 将与仿紧可微流形 V 的“简单覆盖” $ \mathcal{U} = (U_i) $  相关的两个复数双数放入对偶性中；我们将  $ \mathcal{U} $  强制为局部有限，并在  $ U_i $  的每个非空交点  $ U_j $  （J 属于  $ \mathcal{U} $  的神经）处进行可微收缩。第一个双复形由“辅元素” $ \Omega = (\omega_H) $  组成，以双度 (m, p) 表示，其中，对于所有  $ H = (i_0, \ldots, i_p) $  ，使得  $ U_{|H|} = \bigcap_v U_i, \neq \emptyset $  、 $ \omega_H $  是  $ U_{|H|} $  中 m 次的微分形式； 1 次微分 d 和  $ \partial $  是由外微分和限制的交替和定义的，并且有两个相应的同伦运算符，分别使用  $ U_{|H|} $  在一点处的收缩和  $ \mathcal{U} $  下属单元的分区来定义。这些算子使得在 m 阶 de Rham 群（在 V 模精确形式上闭合的 m 形式）和  $ \mathcal{U} $  神经的第 m 上同调群之间建立同构成为可能。第二个复数由“元素”  $ T = (t_H) $  以二度 (m, p) 形式形成，其中，对于上述所有 $ H = (i_0, \ldots, i_p) $ ， $ t_H $  是维度为 m 的有限奇异链，并有  $ U_{|H|} $  的支持；两个微分（-1 次）分别对应于  $ \mathcal{U} $  的奇异同源性和神经同源性，并且仍然有两个同伦算子给出了 V 的奇异同源性和神经同源性之间的同构。两个复合体之间的耦合由奇异链上微分形式的积分给出（(T,  $ \Omega $  ) 仅当 T 或  $ \Omega $  有限时才定义）。

1948 年至 1951 年间，高等师范学院的嘉当研讨会致力于代数拓扑。此外，嘉当在 1948 年的 CNRS 代数拓扑会议上提出了某些结果。1948-49 年的研讨会包含层理论的第一个版本（呈示 12 至 17），该版本已不再流通，并在 1950-51 年由新的呈示取代。研讨会。丛概念采用的形式源自 M. Lazard：拓扑（正则）空间上的 K 模丛  $F$ （ $K$  交换环）  $\mathcal{X}$  是一个扩展空间（来自 Godement 的术语） $p: F \to \mathcal{X}$ ，其中每根纤维  $p^{-1}(x) = F_x$  具有  $K$  -模的结构，因此元素的加法和乘法 $K$  （离散）对于  $F$  的拓扑是连续的。对于 $\mathcal{X}$ 的每个开 $X$ ，我们将 $X$ 中 $F$ 的 $s: X \to F$ 部分的模 $\Gamma(F, X)$ 关联起来（以 $p \circ s = id_x$ 为特征），并且每个纤维 $F_x$ 是 $X$ 邻域的 $\Gamma(F, X)$ 的感应极限 $x$ 。我们可以从与任何开放 $X$ 模 $F_x$ 关联的数据以及与开放的同态的任何包含 $X \subset Y$ 定义一个层 $F$ 

限制  $f_{XY}: F_Y \to F_X$  ，具有限制的传递性；我们将 $x$ 的 $X$ 邻域的 $F_X$ 作为纤维 $F_x$ ，并用合适的拓扑提供它们的不相交和。请注意，规范态射  $F_X \to \Gamma(F, X)$  通常既不是单射也不是满射，尽管我们自然以这种方式定义的各种函数束都是这种情况。例如， $n$  度的 Alexander-Spanier cochanes 的丛  $C^n$  因此是根据  $X^{n+1}$  在  $K$  中的应用的似乎的数据  $C_X^n$  定义的，并且，对于  $n \geq 1$  ，  $C_X^n \to \Gamma(C^n, X)$  不是单射的，但如果  $X$  是仿紧的，则它是满射的。对于奇异 cochanes 的束  $S^n$  来说也是如此。

支持族，因此我们将 $\mathcal{X}$  的闭仿紧紧凑族称为 $\Phi$ ，它是遗传的、稳定的而不是有限联合，并且 $\Phi$  的任何元素都有属于 $\Phi$  的邻域。 Leray 只考虑了紧族的情况，在  $\mathcal{X}$  局部紧的情况下。我们注意到 $\Gamma(F)$ 是 $F$ 的全局部分的模， $\Gamma_{\phi}(F)$ 是其支持（ $x$ 的集合，例如 $s(x) \neq 0$ ；它总是封闭的）属于 $\Phi$ 的部分的子模。当 $f: F \to G$ 是层的满射同态时，相应的同态 $\Gamma_{\phi}(F) \to \Gamma_{\phi}(G)$ 一般不是满射的，但如果我们假设 $f$ 的核 $F'$ 是好的，也就是说，对于 $\mathcal{X}$ 的每个局部有限开覆盖（ $\mathcal{U}^i$ ），  $F'$  是自同态之和，除了包含在具有相同索引的开中的闭之外，每个自同态均为零。 Alexander-Spanier 联链束或单个联链束都是细层；对于可微流形上的外微分形式束也是如此。嘉当以公理化的方式定义了 $\mathcal{X}$ 与束中的系数和族 $\Phi$ 中的支持的上同调，要求基环 $K$ 为主环。它是一系列函子 $F \mapsto H_{\phi}(\mathcal{X}, F)$ ，其值为 $K$  -模， $q<0$  为零， $q=0$  与 $\Gamma_{\phi}(F)$  一致；此外，层的每个精确序列 (1)  $0 \to F' \to F \to F'' \to 0$  都与该序列的同态  $\delta_q: H_{\phi}(\mathcal{X}, F'') \to H_{\phi}^{q+1}(\mathcal{X}, F')$  函子相关联。我们对这些数据施加以下公理：如果  $F$  很好并且  $q>0$  则  $H_{\phi}(\mathcal{X}, F)$  为零；与 (1) 相关的序列  $\cdots \to H_{\phi}(\mathcal{X}, F') \to H_{\phi}(\mathcal{X}, F) \to H_{\phi}(\mathcal{X}, F'') \to H_{\phi}^{q+1}(\mathcal{X}, F') \to \cdots$  是精确的。在建立了这种理论的唯一性之后，嘉当通过证明如果  $0 \to K \to C_0 \to C_1 \to \cdots \to C_n \to \cdots$  是一个精确序列，其中  $C_n$  是细层并且没有扭转（“恒定层  $K”)$  的精细分辨率，我们通过对任何层  $F$  提出  $H_{\phi}(\mathcal{X}, F) = H^q(\Gamma_{\phi}(C \circ F))$  来获得上同调理论，其中基本层 $C$  是复数  $C_0 \to C_1 \to \cdots$  ，而 o 表示  $\mathcal{X}$  上层的张量积。但是 Alexander-Spanier 协链给出了  $K$  的精细分辨率，而对于奇异协链也是如此，我们称空间为  $F$  并使用  $\Phi$  来支持  $\Phi$  。  $F$  的分辨率  $0 \to F \to A_0 \to A_1 \to \cdots$  我们通过 $F$  进行张量化获得 $\Phi$  的 $F$  分辨率。

 $ 0 \to K \to C_0 \to C_1 \to \cdots $  其中  $ C_n $  没有挠场且  $ \Phi $  -内射，也就是说，对于任何外同态  $ F' \to F'' $  ，  $ \Gamma_\phi(C \circ F') \to \Gamma_\phi(C \circ F'') $  是满射的；然后我们说  $ C $  复合体是  $ \Phi $  基本层。

Leray 的封面与 Cartan 演示中的贝壳相对应。  $\mathcal{X}$  上的外壳是为  $\mathcal{X}$  的每个点  $x$  提供的  $K$  模  $A$  ，以及商模  $\varphi_x: A \to A_x$  ；  $A$  的元素 $\alpha$  的支持 $\sigma(\alpha)$  是 $x$  的集合，例如 $\varphi(\alpha) \neq 0$  ，并且我们强制它是闭集且不为空，除非 $\alpha = 0$  。我们还可以按照Leray的方式，通过支撑的数据来定义一个壳；在这种情况下，  $A_x$  是  $A$  与  $\alpha$  的子模（例如  $x \notin \sigma(\alpha)$  ）的商。如果我们具有拓扑的统一性，其中开集基由 $\{\varphi_x(\alpha)|x \in X\}$ 组成，其中 $\alpha \in A$ 和 $X$ 是 $\mathcal{X}$ 的开集，则 $A_x$  ( $x \in \mathcal{X}$ )的不相交和 $F = \mathcal{F}(A)$ 成为束；从  $A$  到  $\Gamma(\mathcal{F}(A))$  的自然态射是单射的，但一般不是满射的。

为了以接近 Leray [30] 最初提出的方式直接获得同调光纤空间的谱序列，S. Eilenberg（演讲 8）将该理论置于有序集同调的公理框架中。它考虑一个（部分）有序的元素集 $A, B, \ldots$ ，其中最小元素为0，最大元素为1，并将 $(A, B)$ 与 $B < A$ 中的任意一对 $(A, B)$ 关联为 $(A, B)$ 中的交换群 $H(A, B)$ 函子（态射由顺序给出），并与 $(A, B, C)$ 中的任何三元组 $(A, B, C)$ 关联为同态 $d: H(A, B) \to H(B, C)$ 函子 $(A, B, C)$  ，因此周期序列 $\cdots \to H(B, C) \to H(A, C) \to H(A, B) \to H(B, C) \to \cdots$  是精确的。我们写  $H(A)$  而不是  $H(A, 0)$  ，与三元组  $(A, A, 0)$  相关的确切序列表明  $H(A, A) = 0$  。如果 $(A_p)_{p \in \mathbb{Z}}$ 是一个以 $A$ 为界的严格递增的元素序列，则对于所有 $p$ ，我们将 $B_p$ 和 $C_p$ 分别定义为 $H(A_p)$ 在 $H(A)$ 和 $H(A_p, A_{p-1})$ 中的图像，将 $D_p$ 定义为 $H(A_p, A_{p-1}) \to H(A, A_{p-1})$ 的内核；对于 $k \geq 1$ ，我们仍然用 $C_p^k$ 指定 $H(A_p, A_{p-1})$ 中 $H(A_p, A_{p-k})$ 的图像，用 $D_p^k$ 指定 $H(A_p, A_{p-1}) \to H(A_{p+k-1}, A_{p-1})$ 的内核。我们有 $D_p^1 = 0$ ， $C_p^1 = H(A_p, A_{p-1})$ ；  $C_p^k$  （分别为  $D_p^k$  ）随  $k$  减少（分别增加）并包含  $C_p^p = C_p$  （分别包含在  $D_p^p = D_p$  中）。最后我们放置 $E_p = C_p/D_p$ 和 $E_p^k = C_p^k/D_p^k$ （因此 $E_p^1 = H(A_p, A_{p-1})$ ）；我们证明 $E_p \simeq B_p/B_{p-1}$ 。使用三元组 $(A_p, A_{p-1}, A_{p-k-1})$  和 $(A_p, A_{p-k}, A_{p-k-1})$  的态射 $d$ ，我们构造了一个微分 $d_p^k: E_p^k \to E_{p-k}^k$ ，使 $(E_p^k)_p$  成为一个复合体，并且我们证明了该复合体的同源性通过 $E^{k+1}$  进行了识别。还有一种逆变理论，其中箭头的方向相反（上同调）。我们通过将  $A$  固定微分模  $A$ （ $A$  基环）的微分子模作为有序集，找到 Koszul 代数意义上的谱序列；如果  $A' \supset B'$  是两个这样的子模，我们将  $H(A', B') = H(A'/B')$  与它们相关联，并且三元组  $A' \supset B' \supset C'$  的  $d$  由精确序列  $0 \to B'/C' \to A'/C' \to A'/B' \to 0$  定义。光纤 $F$  的光纤空间 $\tau: X \to B$  的同源谱序列，其基 $B$  是有限单纯复形，通过子空间 $X_p = \tau^{-1}(B^p)$  过滤 $X$  获得，其中 $B^p$  是 $B$  的 $p$  骨架；有序集是  $X$  子空间的有序集，如果  $X' \supset Y'$  是两个这样的子空间，则我们设置  $H(X', Y') = H(X', Y'; G)$  （相对于群  $G$  中的系数的同源性）。因此，我们有一个谱序列， $E_{p+q, p}^{1} = H_{p+q}(X_p, X_{p-1}; G) $  （使用三元组边运算符  $ (X_p, X_{p-1}, X_{p-2}) $  ）和  $ E_{p+q,p}^2 \simeq H_p(B; H_q(F; G)) $  基数与局部系统  $ H_q(F; G) $  中系数的同源性。作为谱序列的另一个应用，嘉当研究（第 11 和 12 讲）空间与算子群的同调和上同调。

在层理论中，嘉当定义了由  $\mathcal{X}$  上层  $F$  的复合体定义的两个超上同调谱序列。如果  $A$  是一个具有 1 次边界的分级壳，并且其支持属于一个族  $\Phi$  ，则对于所有  $p$  ，我们定义一个（功能）同态  $H^{p+k}(A) \to H_{\Phi}^{p}(\mathcal{X}, H^{k}(\mathcal{F}(A)))$  如果满足以下条件： a)  $\mathcal{X}$  具有  $\Phi$  有限维（即存在一个整数  $n$  使得 $H_{\Phi}^{q}(\mathcal{X}, F) = 0$  对于  $q > n$  和任何层  $F$  ) 或  $A$  的度是下界的； b)  $H^{q}(\mathcal{F}(A)) = 0$  为  $q > k$  。如果我们假设  $H^{q}(\mathcal{F}(A)) = 0$  对应  $q \neq k$  ，  $A \simeq \Gamma_{\Phi}(\mathcal{F}(A))$  和  $A$  同伦  $\Phi$  - 则它甚至是同构。可微流形 $\mathcal{X}$ 的庞加莱对偶性是通过考虑奇异链的束 $S$ 来推导的，在维度 $p$ 中表示为 $S^{-p}$ （以便具有1次微分）；该层是同伦薄的，除了  $p = -n$  之外， $H^{p}(S)$  为零，其中  $n$  是  $\mathcal{X}$  的维度。层  $T = H^{-n}(S)$  与  $K$  局部同构，并由  $\mathcal{X}$  的局部方向定义；我们有  $H^{p-n}(\Gamma_{\Phi}(S \circ F)) \simeq H_{\Phi}^{p}(\mathcal{X}, T \circ F)$  ，其中第一个成员再次表示为  $H_{n-p}^{\Phi}(\mathcal{X}, F)$  ，是  $\Phi$  的  $\mathcal{X}$  的奇异同源性。

J-P 在 $\Phi$  上同调的框架中公开了连续映射 $f:E\to B$  的谱序列。温室（谈话 21）；我们分别在 $E$ 和 $B$ 上为自己提供支持家庭 $\Psi$ 和 $\Phi$ ，并且我们假设它们在适当的意义上是“适应”的。如果  $G$  是  $E$  上的一个层，我们计算它的上同调，就像甲壳  $A^0=\Gamma_{\Psi}(C\circ G)$  的上同调一样，其中  $C$  是  $E$  上的基本层；我们通过将所有 $x\in A^0$ 与支持 $\overline{f}(\sigma(x))$ 关联起来，在 $B$ 上获得一个外壳 $A$ （其中 $\sigma(x)$ 是 $E$ 中 $x$ 的支持）。然后存在一个频谱序列，使得 $E_{p+q,p}^2=H_{\Phi}^p(B,H^q(\mathcal{F}(A)))$ 和 $E$ 是与 $H_{\Psi}(E,G)$ 相关的经过适当滤波的毕业生；  $B$  上的束  $H^q(\mathcal{F}(A))$  的纤维是  $H_{\Psi}^q(F_b,G)$  (  $b\in B$  )。

K. Oka 在 1950 年和 1951 年发表的文章 [34] 和 [35] 引入了一个与多变量分析函数理论中的束非常接近的概念：不确定域的理想（我们现在称之为理想束）。对奥卡这项工作的反思无疑启发了嘉当对层理论的阐述，并且使他能够将奥卡的结果重新表述为层上同调定理（著名的定理 A 和 B）； 1951-52 年的嘉当研讨会在连贯的分析框架内专门讨论了这些问题。 Cartan 和 Serre 于 1953 年建立了具有相干束系数的紧解析空间的上同调的有限性 [8]1；紧凑复流形上局部自由解析层的对偶性归因于 J-P。 1953年的温室[39]。关于分析空间理论的模型，Serre [40]

1954 年 les 基础的代数几何和 faisceaux 代数相干的同同学的重述。

\subsection*{成熟期：1955-58年}
这一时期的标志是同调代数的系统发展及其在层上同调中的应用。 Cartan 和 Eilenberg 于 1956 年出版的《同调代数》[7] 一书至今仍是经典。他使用 Eilenberg 和 MacLane [12] 于 1942 年引入的函子语言，并将上同调理论与卫星函子和派生函子的概念联系起来。在模范畴上定义的协变加性函子  $ T $  的第一个卫星  $ S_1T $  和  $ S^1T $  构造如下：如果  $ A $  是一个模，我们考虑使用  $ P $  直和  $ Q $  单射的精确序列  $ 0 \to M \to P \to A \to 0 $  和  $ 0 \to A \to Q \to N \to 0 $ ，并且设置  $ S_1T(A) = \text{Ker}(T(M) \to T(P)) $  、  $ S^1T(A) = \text{Coker}(T(Q) \to T(N)) $  （我们表明，这并不依赖于精确序列的选择，直至唯一的同构）。然后我们输入  $ S_{n+1}T = S_1(S_nT) $  和  $ S^{n+1}T = S^1(S^nT) $  ，并且我们还为  $ S_nT $  和  $ S_0T = S^0T = T $  标注  $ S^{-n}T $  。对于模的任何精确序列 $ 0 \to A'\to A \to A''\to 0 $ ，我们有连接同态 $ \theta_1: S_1T(A'') \to T(A') \in \theta^1: T(A'') \to S^1T(A') $ ；在无限序列  $ \cdots \to S^nT(A') \to S^nT(A) \to S^nT(A'') \to S^{n+1}T(A') \to \cdots $  中，两个连续态射的复合为零，如果  $ T $  是半精确的，则整个序列是精确的。加法逆变函子有一个对偶理论。从协变加法函子  $ T $  派生的函子是通过设置  $ R^nT(A) = H^n(T(X)) $  同时定义的（不再通过归纳），其中  $ 0 \to A \to X^0 \to X^1 \to \cdots $  是模  $ A $  的单射解析；对于逆变函子，我们将采用投影解析  $ \cdots \to Y^1 \to Y^0 \to A \to 0 $  。相反，通过在  $ T $  是协变的情况下采用射影解析，在  $ T $  是逆变的情况下采用单射解析，我们定义了左导函子  $ L_nT $  。当  $ T $  是协变（或逆变）时，模  $ 0 \to A'\to A \to A'' \to 0 $  的任何精确序列都定义连接同态  $ R^nT(A'') \to R^{n+1}T(A') $  和  $ L_nT(A'') \to L_{n-1}T(A') $  （分别为  $ R^nT(A') \to R^{n+1}T(A'') $  和  $ L_nT(A') \to L_{n-1}T(A'' $  ）），并且派生函子的长序列是精确的。我们通过对每个参数进行适当的解析，以相同的方式定义从多个参数派生的函子。我们有函子态射  $ \tau^0: T \to R^0T $  和  $ \sigma_0: L_0T \to T $  ；为了使第一个（或第二个）成为同构， $ T $  精确位于左侧（或右侧）是必要且充分的。当这个条件得到验证时，对于  $ T $  来说，准确地说， $ R^1T = 0 $  （或  $ L_1T = 0 $  ）是必要且充分的。从态射 $ \sigma_0 $ 和 $ \tau^0 $ ，我们定义态射 $ \sigma_n: L_nT \to S_nT $ 和 $ \tau^n: S^nT \to R^nT $ ；如果 $ T $  是右精确的，则第一个是同构，如果 $ T $  是左精确的，则第二个是同构。导出函子理论是定义函子Tor（由张量积导出）和Ext（由函子Hom导出）； Tor 函子使得能够以令人满意的方式表达 Künneth [27] 获得的关系。

两个拓扑空间的乘积与因子的同源性之间的关系，正是为了理解这些关系，嘉当和艾伦伯格开展了同调代数的发展。由此构建的理论还使得统一各种上同调理论成为可能：代数上同调（Hochschild [24]）、有限群（Tate）、李代数（Chevalley-Eilenberg [11]）。本书第十五章介绍了谱序列的理论，随后的章节介绍了它的应用；它本质上是复合函子的谱序列，其中术语 $ E_{2}^{pq} $ 由派生函子组成，而结果是与复合函子的导数相关的毕业生（用于适当的过滤）。本书的最后一章致力于超同调：如果  $ T $  是模的加性函子， $ A $  是模的复合体，我们证明双复形  $ T(X) $  的上同调，其中  $ X $  是  $ A $  （它是双复形）的“Cartan-Eilenberg 解析”，不依赖于  $ X $  的选择。该上同调具有过滤，使得相关的梯度是两个谱序列的结果，这两个谱序列的项 $ E_{2}^{pq} $ 分别是 $ H^{p}(R^{q}T(A)) $ 和 $ (R^{q}T)(H^{p}(A)) $ 。

D. Buchsbaum [2] 为同源代数提出了一个更通用的抽象框架：精确范畴的框架，格罗滕迪克以阿贝尔范畴的名义引入了该框架。嘉当和艾伦伯格在书的附录中给出了确切范畴的公理。 Grothendieck [14] 的工作旨在为具有束中系数的拓扑空间的上同调以及从模函子导出的函子理论找到一个通用框架，其类比是显而易见的；它于 1955 年春天在堪萨斯大学展出。格洛腾迪克将阿贝尔范畴定义为加法范畴，其中每个态射都有一个核心和一个副核，图像中共像的自然态射是同构。 Si  $(A_i)$  是一个对象 $A$  的子对象的递增过滤器族，对于所有 $B \subset A$ ，我们有 $(\sum A_i) \cap B = \sum (A_i \cap B)$ ，我们说 $C$  有足够的单射。

格洛腾迪克将  $C$  上的  $\partial$  -函子称为加法函子序列  $(T^i)$ ，例如协变体，对于  $C$  中的任何精确序列  $0 \to A' \to A \to A'' \to 0$  、连接同态序列  $\partial: T^i(A'') \to T^{i+1}(A')$  、精确序列中的函子，例如长序列的两个连续态射的复合 $\cdots \to T^i(A') \to T^i(A) \to T^i(A'') \to T^{i+1}(A') \to \cdots$  为零。它将协变函子  $F$  的右卫星  $S^i F$  的序列定义为通用  $\partial$  -函子，因为任何态射  $f^0: F \to T^0$ （其中  $(T^i)$  是  $\partial$  -函子）都扩展到  $\partial$  -函子的唯一态射  $(S^i F) \to (T^i)$ ；这采用了嘉当和艾伦伯格对卫星的公理化描述。通过考虑  $\partial^*$  函子，以相同的方式定义左侧卫星。如果任何对象嵌入到单射对象中，则对于任何协变加法函子  $F$  都存在正确的卫星  $S^i F$  (  $i \geq 0$  )，并且它们形成一个  $\partial$  -精确函子：

例如， $F$  是否恰好位于左侧或右侧。派生函子在 Cartan-Eilenberg 书中通过单射或射影解析来定义，必须假设其存在。格洛腾迪克将谱序列和超上同调理论扩展到阿贝尔范畴的情况。特别是，它在此框架中给出了复合函子的谱序列，对于协变加法函子 $F:C\to C'\to C'et G:C'\to C''$ ；我们假设 $C$ 和 $C'$ 有足够的单射， $G$ 是精确的， $R^qG(F(I))=0$ 对于所有 $q>0$ 如果 $I$ 是 $C$ 的单射对象，并且我们获得一个通向 $R(GF)$ （过滤的）的谱函子，并且对于所有对象 $E_2^{pq}(A)=R^pG(R^qF(A))A$ 来自 $C$ 。

格洛腾迪克在拓扑空间上调用集合的预层，层对他来说是特定的预层，但满足以下条件：如果 (  $ U_i $  ) 是一个由非空开口端覆盖的开放 U 形，并且如果 (  $ f_i $  ) 是一个元素族  $ f_i \in F(U_i) $  ，使得  $ f_i $  和  $ f_j $  到  $ U_i \cap U_j $  的限制在 $ F(U_i \cap U_j) $  无论索引对 (i, j) 是多少，都存在一个唯一的  $ f \in F(U) $ ，对于任何索引 i，其对  $ U_i $  的限制是  $ f_i $ 。这个条件表示自然映射 $ F(U) \to \Gamma(\widetilde{F}, U) $ 是双射的，通过 $ \widetilde{F} $ 指定与预层F相关的嘉当意义上的层（扩展空间）。我们类似地定义具有范畴中的值的层，例如群的层，环的层等。如果O是环层，则它承认一个生成器（层的直接和）  $ O_U $  与 U 中的 O 以及 X - U 中的 0 一致，其中 U 是

层中系数的上同调是由 Grothendieck 提出的，没有对空间做出任何限制性假设，这使得该理论可以应用于我们在代数几何中遇到的非分离空间（Zariski 拓扑；Serre [40]，对于相干层，使用 Cech 的仿射开放空间覆盖的上同调，在确定这些空间对于相干系数来说是上同调平凡的）。格洛腾迪克考虑的支持族比嘉当的支持族更普遍：它们是非空的 $ \varPhi $ 族，具有递增遗传和过滤封闭部分。函子  $ F \to \varGamma_{\varPhi}(F) $  （在  $ \varPhi $  中支持交换群 F 的全局部分，即交换群的一束）是左精确的，上同调  $ H_{\varPhi}^{q}(X, F) $  被定义为该函子的右导函子序列；它可以通过 F 的单射分辨率来计算。连续映射  $ f: Y \to X $  的谱序列表现为复合函子谱序列的特例；如果 $ \varPhi $  和 $ \varPsi $  是 $ f_{\varPsi} $  上的支持族，则 $ f_{\varPsi} $  是在 $ \varPsi $  中具有支持的直像函子， $ R^{q}f_{\varPsi} $  是其q 次右导数，而 $ \varPsi' $  是 $ \varPsi $  的适当修改（如果 $ \varPsi $  包括所有

Y 的闭集，它是 Y 的闭 A 的集合，例如  $ \overline{f(A)} \in \Phi $  。高阶直像层 $ R^q f_{\psi}(F) $  与预层 $ U \mapsto H_{\psi(U)}^q(f^{-1}(U), F) $  相关联（ $ \Psi(U) $  是 $ \Psi $  在X 的开U 上方的位置）。同样，将开覆盖  $ \mathcal{U} $  的 Čech 上同调与真上同调联系起来的谱序列是复合函子谱序列的应用；它的术语 $ E_2^q $  值 $ H^p(\mathcal{U}, H^q(F)) $ ，其中 $ H^q(F) $  是与预层 $ V \mapsto H^q(V, F) $  关联的层（它是预层范畴中层范畴的包含函子的q 次右导数）。

文章最后内容是关于 Ext 的 faisceaux de 模的全部新理论，以及与 Ext 的 locaux 和 Ext globaux 相关的套件光谱，以及与 faisceaux 或 faisceaux 操作组的系数空间的同源关系等价的。格罗滕迪克的 Ext de faisceaux 理论服务于 variétés algébriques 投影的二元理论的公式和演示（参见 [15] 和 [16]）。

第一本系统阐述层理论的书由 R. Godement 于 1958 年出版 [13]。所采用的观点与格洛腾迪克的观点接近；主要创新包括引入某些范畴的极其有用的非循环层：松弛层和软层，它们比内射层和精细层更容易使用来构建分辨率。如果层 F 在开放空间中的任意截面延伸到整个空间，我们就说层 F 是松弛层；内射层是松弛的。当阿贝尔层的外同态  $ \mathcal{L} \to \mathcal{L}'' $  的核是松弛的时，同态  $ \mathcal{L}(U) \to \mathcal{L}''(U) $  对于任何开 U 来说都是满射的。每个层 F 都规范地浸入板束  $ \mathcal{C}^0(\mathcal{F})\colon U \mapsto \prod_{x\in U} \mathcal{F}(x) $  中，并且我们推导出一个规范的松弛分辨率  $ 0 \to \mathcal{F} \to \mathcal{C}^0 \to \mathcal{C}^1 \to \cdots $ ，这使得可以用  $ \mathcal{F} $  中的系数计算上同调，因为松弛对于任何支撑系列 $ \Phi $ ，层都是非循环的 $ \Phi $ 。如果  $ \mathcal{F} $  的任何部分在闭合的一个之上，我们说空间 X 上的层  $ \mathcal{F} $  是软的，我们说它是  $ \Phi $  -软（ $ \Phi $  平行紧致族），如果对于所有  $ S \in \Phi $  ，限制  $ \mathcal{F}|S $  是一个软束。如果 X 是仿紧致的，则任何松弛层都是软层。如果外同态 $ \mathcal{L} \to \mathcal{L}'' $ 的核是 $ \Phi $  -mou，则同态 $ \Gamma_{\Phi}(\mathcal{L}) \to \Gamma_{\Phi}(\mathcal{L}'') $ 是满射的；由此可见， $ \Phi $  -软层是 $ \Phi $  -无环的，并且我们可以使用 $ \Phi $  -软分辨率来计算 $ \Phi $  -上同调。如果环层 $ \mathcal{A} $ 是 $ \Phi $  -软，则任何 $ \mathcal{A} $  -模也是如此；  $ \Phi $  -细阿贝尔丛的特征在于它们的局部自同态丛（在 $ \mathbb{Z} $  上）是 $ \Phi $  -软的。

\subsection*{层的导出范畴和操作}
为了代数几何的目的，格洛滕迪克在 1957-1965 年显着发展了同调代数和层理论，引入了导出范畴的概念和形式主义

层运算在导出范畴的框架内进行。此外，为了定义具有常数（“离散”）系数的方案的良好上同调，有必要通过用站点和拓扑空间替换拓扑空间的概念来扩大拓扑空间的概念。 1960 年至 1968 年由格洛腾迪克主持的 IHS 代数几何研讨会就带有这些理论革新的标志。 1961-62 年涉及相干代数层的局部对偶理论；在那里，我们发现了局部封闭子空间中支持的上同调、对偶模的概念以及图上准相干层的 Ext 的研究。为了获得令人满意的形式的对偶定理，需要有导出范畴的语言，J-L. Verdier 根据 Grothendieck 的思想（参见 [42]）在 1963 年的论文中对其进行了格式化。在这种语言中，我们不是研究上同调不变量，而是直接研究复形，通过充分丰富态射集，使得导致上同调同构的复形态射变成同构；这避免了移动到子商的危险操作。导出范畴意义上的总派生函子是在适当的假设下组成的，就像我们派生的函子一样，谱序列仅在我们想要计算上同调时出现。 R. Hartshorne 于 1963-64 年在哈佛大学举办了一次研讨会，根据格洛腾迪克的“前言”讨论相干代数层的对偶性理论；他于 1966 年发表了这个研讨会（[21]）。对偶定理是在相对情况下用诺特（前）模式的 $ f: X \to Y $  态射来表述的；紧邻 $ \mathcal{O}_Y $  -模的层的通常逆像 $ f^* $ 及其全左导数 $ L_f^* $ （在导出范畴的意义上），我们在适当的假设下定义一个逆像函子 $ f^1 $ （在导出范畴中），提供迹态射 $ \mathrm{Tr}_f: Rf_s f^1 \to id $ ，以便对于F具有对偶同构 $ R\mathrm{Hom}_{\mathcal{O}_X}(F, f'G) \simeq R\mathrm{Hom}_{\mathcal{O}_Y}(Rf_s F, G) $ 从  $ \mathcal{O}_X $-模 派生的范畴和从  $ \mathcal{O}_Y $-模 派生的范畴中的 G。当f是相对维数 $ n, f'(G) = f^*(G) \otimes \omega[n] $ 的光滑时，其中 $ \omega = \Omega_{X/Y}^n $ 是n次相对微分形式的束；当 f 完成时， $ f'(G) = \mathcal{H}\mathrm{Hom}_{\mathcal{O}_Y}(f_s \mathcal{O}_X, G) $  。总体构造极其复杂。格洛腾迪克 (Grothendieck) 在 1964-65 年的研讨会上对等上同调进行了类似的构造，他在会上给出了 l-adic 上同调的对偶性理论；它是庞加莱对偶性的代数和相对类比。 Verdier [43] 于 1965 年处理了普通拓扑中的相应理论，从而给出了庞加莱对偶性的相对概括，在独立拓扑空间的连续  $ f: X \to Y $  应用的情况下，Y 局部仿紧，f 在以下意义上是“局部欧几里得”：局部在 X 上，它由闭浸入和随后的投影  $ V \times \mathbb{R}^n \to V $ （Y 的开）组成；附加假设为 f 的有限上同调维数和 X 上  $ \mathbb{Z} $  的有限拓扑维数。

Grothendieck 的方法采用了 M. Sato 和 M. Kashiwara ([38] 和 [25]) 的方法，包括部分驱动方程系统研究和微区域分析。这是重要的发展结果，是 70 年后的选择的回归，$D$-模 的技术问题是几何代数方面的问题。

\subsection*{参考书目}
1 Alexander, J.W.：在抽象空间的连接环上。安.数学。 37, 698–708 (1936)

2 Buchsbaum, D.A.：精确范畴和对偶性。跨。 A.M.S. 80, 1–34 (1955)

3 Cartan, H.：论群运行空间的上同调。 C.R.阿卡德。巴黎 226, 148–150, 303–305 (1948)

4 Cartan, H.：“代数拓扑”研讨会，第一年（1948-49）

5 Cartan, H.：关于代数拓扑甲壳的概念。代数拓扑学座谈会，CNRS 1-2 (1947)

6 Cartan, H.：研讨会“群的上同调、谱序列、层”。第三年（1950–51）

7 Cartan, H. 和 Eilenberg, S.：同调代数。普林斯顿大学出版社 (1956)

8 Cartan, H. 和 Serre, J-P.：关于紧解析簇的有限性定理。 C.R.阿卡德。巴黎 237, 128–130 (1953)

9 Čech, E.：复数乘法。安.数学。 37, 681–697 (1936)

10 Chern, S.S. 和 Spanier, E.：纤维束的同源结构。过程。纳特。阿卡德。科学。美国 36, 248–255 (1950)

11 Chevalley, C. 和 Eilenberg, S.：李群和李代数的上同调理论。跨。 A.M.S. 63, 85–124 (1948)

12 Eilenberg, S. 和 Maclane, S.：群扩展和同源性。安.数学。 43, 758–831 (1942)

13 Godement, R.：Théorie des faisceaux。赫尔曼，巴黎 (1958)

14 Grothendieck, A.：关于同调代数的某些点。东北数学. J. 9, 119–221 (1957)

15 Grothendieck, A.：关于 faisceaux algébriques cohérents 的二元理论。塞姆。布尔巴基 149 (1957)

16 Grothendieck, A.：抽象代数簇的上同调。实习生。数学大会。 1958 年爱丁堡，剑桥 103–118 (1960)

17 Grothendieck, A.：Cohomologie locale des faisceaux cohérents et théorèmes de Lefschetz locaux et globaux。塞姆。 de Géométrie algébrique 1962 (SGA2)，北荷兰省，阿姆斯特丹 (1968)

18 Grothendieck, A.：L-adic 上同调和函数 L. Sém。代数几何 1964–65 (SGA5)。莱克特。数学笔记。 589. 施普林格，柏林 海德堡 纽约 (1977)

19 Grothendieck, A.：Récoltes et semailles。蒙彼利埃大学预出版

20 Gysin, W.：关于流形映射和纤维的同源理论。评论。数学。赫尔夫。 14, 61–122 (1941)

21. Hartshorne, R.：残差和对偶性。莱克特。数学笔记。 20.施普林格，柏林海德堡纽约（1966）

22 Hirsch, G.：《纤维空间同源组》。公牛社会。数学。比利时 23-33 (1947-48)

23 Hopf, H.：连续表示理论中的一些问题。恩斯。数学。 35, 334–347 (1936)

24 Hochschild，G.：关于结合代数的上同调群，Ann。数学。 46, 58–67 (1945)

25 Kashiwara, M.：偏微分方程组的代数研究。论文，大学。东京 (1970)

26 Koszul, J-L.：李代数的同调和上同调。公牛。苏克。数学。法国 78, 65–127 (1950)

27 Künneth, H.：关于产品歧管的扭转数。数学。安. 91, 65–85 (1923)

28 Lefschetz, S.：关于代数簇的某些数值不变量及其在阿贝尔簇上的应用。跨。 A.M.S. 22, 327–382 (1921)

29 Leray, J.：关于拓扑空间的形式和表示的不动点。 J.马斯。 Pures et Appl.，第 9 系列 24, 95-167 (1945)；关于拓扑空间的闭点集的位置，同上。 169-199；关于方程和变换，同上。 201-248

30 Leray, J.：纤维空间在其基底上的投影的同源环的性质。 C.R.阿卡德。巴黎 223, 395–397 (1946)

31 Leray, J.：局部紧空间和连续应用的谱环和滤波同调环。 J.马斯。 Pures et Appl.，9<sup>e</sup> 系列 29, 1-139 (1950)

32 Leray, J.：纤维相连的纤维空间的同源性，同上。 169–213

33 Leray, J. 和 Schauder, J.：拓扑和函数方程。安. ENS 51, 45–78 (1934)

34 Oka, K.：关于一些算术概念。公牛。苏克。数学。法国 78, 1–27 (1950)

35 Oka, K.：基本引理。数学杂志。苏克。日本 3, 204–214 和 259–278 (1951)

36 Picard, E.：关于两个自变量的代数函数的记忆。 J.马斯。纯粹和应用。 5、135–319（1889）

37 de Rham, G.：关于 n 维簇的分析位置。 J.马斯。 Pures et Appl.，9<sup>e</sup> 系列 10, 115–200 (1931)

38 Sato, M.：超函数和偏微分方程。过程。实习生。泛函分析及相关主题会议东京 1969, 91–94 Univ。东京出版社 (1969)

分析和相关主题东京 1969 年，91-94。大学。东京出版社 (1969)

39 Serre, J-P.：对偶定理。通讯。数学。赫尔夫。 29, 9–26 (1955)

40 Serre, J-P.：相干代数层。安.数学。 61, 197–278 (1955)

41 Steenrod, N.：与局部系数的同源性。安.数学。 44, 610–627 (1943)

42 Verdier, J-L.：导出范畴（状态 0）。阅读。数学笔记。 569, 262–312。施普林格，柏林海德堡纽约 (1977)

43 Verdier, J-L.：局部紧空间上同调的对偶性。 SEM。布尔巴基 300 (1965–66)

44 Wang, H.C.：球体上纤维束的同源群。杜克数学。 J. 16, 33–38 (1949)

45 Weil, A.：H. Cartan 的信。作品 2, 44。柏林 (1985)

46 Weil, A.：《拉姆的理论》。通讯。数学。赫尔夫。 26 (1952)

47 Whitney, H.：关于综合体中的产品。安.数学。 39, 397–432 (1938)

48 Dieudonné, J.：1900-1960 年代数和微分拓扑史。伯克豪瑟，波士顿 (1989)

\section{同调代数}
\subsection{概要}
本章包含理解本书其余部分所必需的同调代数基础：范畴和函子、三角范畴、定位、导出范畴、ind-objects 和 pro-objects、Mittag-Leffler 条件。

由于不可能在一个章节中介绍整个理论的所有细节，因此我们省略了一些辅助结果作为练习，并且我们将 t 结构理论推迟到第十章，直到第十章才使用。

当然，读者也可以查阅有关该主题的（经典）书籍和论文，例如 Bourbaki [1]、Cartan-Eilenberg [1]、Freyd [1]、Gabriel-Zisman [1]、Godement [1]、Grothendieck [1]、Hilton-Stammbach [1]、Iversen [1]、MacLane [1]、Mitchell [1]、Northcott [1]，特别是关于导出范畴，Deligne [1]、Gelfand-Manin [1]、Hartshorne [1] 和 Verdier [2]。

注意。 1.0。众所周知，操纵集合的范畴可能是危险的。避免这种危险的一种方法是留在给定的宇宙中，这就是我们所假设的。这对我们的目的没有影响，我们不会进一步讨论这一点。

\subsection{范畴与函子}
\statementlabel{定义 1.1.1}范畴 $ \mathcal{C} $  由以下数据组成。

(i) 一个族 Ob(C)，其成员称为 C 的对象，

(ii) 对于  $\mathrm{Ob}(\mathcal{C})$  的所有对  $(X, Y)$  ，有一个集合  $\mathrm{Hom}_{\mathcal{C}}(X, Y)$  ，其元素称为从  $X$  到  $Y$  的态射，

(iii) 对于  $\mathrm{Ob}(\mathcal{C})$  的任何三元组  $(X, Y, Z)$  ，从  $\mathrm{Hom}_{\mathcal{C}}(X, Y) \times \mathrm{Hom}_{\mathcal{C}}(Y, Z)$  到  $\mathrm{Hom}_{\mathcal{C}}(X, Z)$  的映射，称为合成映射，并表示为  $(f, g) \mapsto g \circ f$  。

这些数据满足：

(1.1.1) 态射的复合是结合律，

(1.1.2) 对于任何  $ X \in \mathrm{Ob}(\mathcal{C}) $  都存在  $ \mathrm{id}_{X} \in \mathrm{Hom}_{\mathcal{C}}(X, X) $  ，使得对于任何  $ f \in \mathrm{Hom}_{\mathcal{C}}(X, Y) $  和任何  $ g \in \mathrm{Hom}_{\mathcal{C}}(Y, X) $  存在  $ f \circ \mathrm{id}_{X} = f $  和  $ \mathrm{id}_{X} \circ g = g $  。

请注意，对于每个  $X \in \mathrm{Ob}(\mathcal{C})$  ，  $\mathrm{id}_{X}$  都是唯一的。

我们简写 $f: X \to Y$ 来表示态射 $f \in \mathrm{Hom}_{\mathcal{C}}(X, Y)$ 。如果存在  $g: Y \to X$  使得  $f \circ g = \mathrm{id}_{Y}$  和  $g \circ f = \mathrm{id}_{X}$  ，态射  $f: X \to Y$  称为  $\mathrm{isomorphism}$  。如果  $f$  是同构，我们写为  $f: X \xrightarrow{\sim} Y$  或  $X \simeq Y$  。

示例 1.1.2。我们用 $ \mathcal{S} $ 表示集合的范畴和集合的映射。

 $\mathcal{C}$  的子范畴  $\mathcal{C}'$  是范畴  $\mathcal{C}'$  ，使得  $\mathrm{Ob}(\mathcal{C}') \subset \mathrm{Ob}(\mathcal{C})$  和对于  $\mathrm{Ob}(\mathcal{C}')$  、  $\mathrm{Hom}_{\mathcal{C}'}(X, Y) \subset \mathrm{Hom}_{\mathcal{C}}(X, Y)$  的任何对  $(X, Y)$  ，具有导出的组合定律，以及  $\mathrm{id}_X \in \mathrm{Hom}_{\mathcal{C}'}(X, X)$  。

如果此外  $ \mathrm{Hom}_{\mathcal{C}^{\prime}}(X, Y) = \mathrm{Hom}_{\mathcal{C}}(X, Y) $  ，则  $ \mathcal{C}^{\prime} $  称为  $ \mathcal{C} $  的满子范畴。

示例 1.1.3。我们用 Top 表示拓扑空间和连续映射的范畴。那么Top是Set的一个子类，但不是完整的子类。

让 $\mathcal{C}$  成为一个范畴。相反的范畴，表示为  $\mathcal{C}^{\circ}$  ，定义如下：

\[
\left\{\begin{aligned}{}&{{}\operatorname{O b}(\mathcal{C}^{\circ})=\operatorname{O b}(\mathcal{C})\enspace,}\\ {}&{{}\operatorname{H o m}_{\mathcal{C}^{\circ}}(X,Y)=\operatorname{H o m}_{\mathcal{C}}(Y,X)\quad\mathrm{f o r~a n y~p a i r}\left(X,Y\right)\mathrm{o f}\mathcal{C}\enspace,}\\ {}&{{}\mathrm{w i t h~t h e~o b v i o u s~c o m p o s i t i o n~l a w}.}\\ \end{aligned}\right.
\]

令  $f: X \to Y$  为  $\mathcal{C}$  中的态射。如果对于任何  $W \in \text{Ob}(\mathcal{C})$  和  $\text{Hom}_{\mathcal{C}}(W, X)$  的任意一对  $(g, g')$  使得  $f \circ g = f \circ g',$  一个有  $g = g'.$  ，则有人说  $f$  是  $\text{monomorphism}$  如果  $f$  是  $\mathcal{C}^\circ$  中的  $\text{monomorphism}$  ，则说  $f$  是  $\text{epimorphism}$  ，即，对于任何 $Z \in \text{Ob}(\mathcal{C})$ 和 $\text{Hom}_{\mathcal{C}}(Y, Z)$ 的任何一对 $(h, h')$ ，使得 $h \circ f = h' \circ f$ ，其中一个有 $h = h'.$ 

在范畴  $\mathcal{C}$  中，如果  $\operatorname{Hom}_{\mathcal{C}}(P, Y)$  对于任何  $Y \in \operatorname{Ob}(\mathcal{C})$  恰好有一个元素，则对象  $P$  称为初始对象。类似地，如果  $\operatorname{Hom}_{\mathcal{C}}(X, Q)$  只有一个元素，即  $Q$  是  $\mathcal{C}^\circ$  中的初始对象，则对象  $Q$  称为  $\operatorname{final}$  。请注意，两个初始（或最终）对象自然是同构的。

\statementlabel{定义 1.1.4}让 $ \mathcal{C} $  和 $ \mathcal{C}' $  成为两个范畴。从 $ \mathcal{C} $ 到 $ \mathcal{C}' $ 的函子F由以下数据和规则组成。

(i) 映射 $ F: \mathrm{Ob}(\mathcal{C}) \to \mathrm{Ob}(\mathcal{C}' $ ，

(ii) 对于  $ \mathrm{Ob}(\mathcal{C}) $  的任何对  $ (X, Y) $  ，映射  $ F: \mathrm{Hom}_{\mathcal{C}}(X, Y) \to \mathrm{Hom}_{\mathcal{C}'}(F(X), F(Y)) $  。

这些数据满足：

\[
F(\mathrm{i d}_{X})=\mathrm{i d}_{F(X)}\mathrm{~,~}
\]

\[
F(f\circ g)=F(f)\circ F(g)\enspace.
\]

有时有人说  $F$  是从  $\mathcal{C}$  到  $\mathcal{C}'$  的协变函子，而从  $\mathcal{C}^\circ$  到  $\mathcal{C}'$  的函子被称为从  $\mathcal{C}$  到  $\mathcal{C}'$  的逆变函子。

例如让  $X \in \mathrm{Ob}(\mathcal{C})$  。那么 $\mathrm{Hom}_{\mathcal{C}}(X, \cdot) : Z \mapsto \mathrm{Hom}_{\mathcal{C}}(X, Z)$  是从 $\mathcal{C}$  到 $\mathcal{S}$  集合的（协变）函子， $\mathrm{Hom}_{\mathcal{C}}(\cdot, X) : Z \mapsto \mathrm{Hom}_{\mathcal{C}}(Z, X)$  是从 $\mathcal{C}$  到 $\mathcal{S}$  集合的逆变函子。

\statementlabel{定义 1.1.5}令  $ F_1 $  和  $ F_2 $  为从  $ \mathcal{C} $  到  $ \mathcal{C}' $  的两个函子。从 $ F_1 $  到 $ F_2 $  的态射 $ \theta $  由以下数据组成。

对于任何  $ X \in \mathrm{Ob}(\mathcal{C}) $  ，一个元素  $ \theta(X) \in \mathrm{Hom}_{\mathcal{C}}(F_1(X), F_2(X)) $  。

这些数据满足：

(1.1.6) 对于任何 $f \in \mathrm{Hom}_{\mathcal{C}}(X, Y)$ ，下图可交换：

\[
\begin{array}{l l l}{F_{1}(X)\xrightarrow[]{\theta(X)}F_{2}(X)}\\ {\displaystyle\bigvee F_{1}(f)}&{\displaystyle\bigvee F_{2}(f)}\\ {F_{1}(Y)\xrightarrow[]{\theta(Y)}F_{2}(Y).}\\ \end{array}
\]

请注意，如果  $\mathcal{C}$  和  $\mathcal{C}'$  是两个范畴，则得到一个新范畴，其对象是从  $\mathcal{C}$  到  $\mathcal{C}'$  的函子，态射是此类函子的态射。

\statementlabel{定义 1.1.6}如果存在  $X \in \mathrm{Ob}(\mathcal{C})$  使得  $F$  与函子  $\mathrm{Hom}_{\mathcal{C}}(X, \cdot)$  同构，则从  $\mathcal{C}$  到  $\mathcal{Set}$  的函子  $F$  称为可表示的。

在这种情况下，X 在同构上是唯一的，被称为 F 的代表。

逆变函子有一个类似的定义。

\statementlabel{定义 1.1.7}让 $ F: \mathcal{C} \to \mathcal{C}' $  成为一个函子。

(i) 如果满足以下条件，则称 F 完全忠实：

(1.1.7) 对于 $\mathrm{Ob}(\mathcal{C})$  的任何对 $(X, Y)$ ，映射 $\mathrm{Hom}_{\mathcal{C}}(X, Y) \to \mathrm{Hom}_{\mathcal{C}^{\prime}}(F(X), F(Y))$  是双射的。

(ii) 如果满足 (1.1.7) 并且满足以下条件，则称 F 是范畴的等价：

(1.1.8) 对于任何 $ X' \in \text{Ob}(\mathcal{C}') $  都存在 $ X \in \text{Ob}(\mathcal{C}) $  和同构 $ X' \to F(X) $  。

请注意， $F$  是范畴的等价物，当且仅当存在函子  $F': \mathcal{C}' \to \mathcal{C}$  以及函子  $F \circ F' \simeq \mathrm{id}_{\mathcal{C}}$  和  $F' \circ F \simeq \mathrm{id}_{\mathcal{C}}$  的同构。在这种情况下，有人说  $F'$  是  $F$  的准逆。

我们可以将范畴  $\mathcal{C}$  嵌入到从  $\mathcal{C}$  到  $\mathcal{Set}$  的所有逆变函子的范畴  $\mathcal{C}^{\vee}$  中。

让  $ h: \mathcal{C} \to \mathcal{C}^{\vee} $  成为函子  $ X \mapsto \mathrm{Hom}_{\mathcal{C}}(\cdot, X) $  。

\statementlabel{命题 1.1.8}(i) 对于任何  $ X \in \mathrm{Ob}(\mathcal{C}) $  和  $ F \in \mathrm{Ob}(\mathcal{C}^{\vee}) $  ，我们有：

\[
\operatorname{H o m}_{\mathcal{C}^{\vee}}(h(X),F)\simeq F(X)\enspace.
\]

(ii) 函子 h 是完全忠实的。

\statementlabel{证明}(i) 对于 $f \in \mathrm{Hom}_{\mathcal{G}^\vee}(h(X), F)$ ，我们将 $\phi(f) \in F(X)$  关联如下。由于  $f$  是从  $h(X)$  到  $F$  函子的态射， $f(X)$  给出了从  $h(X)(X) = \mathrm{Hom}_{\mathcal{G}}(X, X)$  到  $F(X)$  的映射。然后将  $\phi(f)$  定义为  $\mathrm{id}_X$  的图像。相反，对于 $s \in F(X)$ ，我们可以如下关联 $\psi(s) \in \mathrm{Hom}_{\mathcal{G}^\vee}(h(X), F)$ 。对于  $Y \in \mathrm{Ob}(\mathcal{C})$  ，考虑映射：

\[
\begin{aligned}{h(X)(Y)}&{{}=\operatorname{H o m}_{\mathcal{C}}(Y,X)}\\ {}&{{}\to\operatorname{H o m}_{\mathfrak{S e l}}(F(X),F(Y))}\\ {}&{{}\to F(Y),}\\ \end{aligned}
\]

其中最后一个箭头由 s 定义。很容易检查 $ \phi $  和 $ \psi $  彼此相反。

(ii) 将 (i) 应用于 $ h(Y) $ ，我们得到：

\[
\begin{aligned}{\operatorname{H o m}_{\mathcal{G}^{\vee}}(h(X),h(Y))}&{{}\simeq h(Y)(X)}\\ {}&{{}}\\ {}&{{}=\operatorname{H o m}_{\mathcal{G}}(X,Y)~.~}&{}&{{}\Box}\\ \end{aligned}
\]

请注意，逆变函子  $F: \mathcal{C} \to \mathcal{S}$  是可表示的，当且仅当对于某些  $X \in \mathrm{Ob}(\mathcal{C})$  而言，它与  $h(X)$  同构。

同样，我们可以考虑  $ \mathcal{C}^{\vee} $  的相反范畴  $ \mathcal{C}^{\wedge} $  ，即从  $ \mathcal{C} $  到  $ \mathcal{Set} $  的（协变）函子范畴。范畴之间存在天然的等价性：

\[
\mathcal{C}^{\wedge}\simeq\mathcal{C}^{\circ\vee\circ}
\]

并且函子  $ h': \mathcal{C} \to \mathcal{C}^\wedge $  、  $ X \mapsto \mathrm{Hom}_{\mathcal{C}}(X,\cdot) $  完全忠实。

\subsection{阿贝尔范畴}
\statementlabel{定义 1.2.1}加性范畴 $ \mathcal{C} $  是范畴 $ \mathcal{C} $ ，使得：

(i) 对于 $\mathrm{Ob}(\mathcal{C})$  中的任何对 $(X,Y)$  ， $\mathrm{Hom}_{\mathcal{C}}(X,Y)$  具有加性（即阿贝尔）群的结构，且复合律是双线性的，

(ii) 存在一个对象 0 使得  $ \mathrm{Hom}_{\mathcal{C}}(0,0)=0 $  ，

(iii) 对于  $\mathrm{Ob}(\mathcal{C})$  的任何对  $(X, Y)$  函子

\[
W\mapsto\operatorname{H o m}_{\mathcal{C}}(X,W)\times\operatorname{H o m}_{\mathcal{C}}(Y,W){~i s~r e p r e s e n t a b l e~},
\]

(iv) 对于  $\mathrm{Ob}(\mathcal{C})$  的任何对  $(X, Y)$  ，函子

\[
W{\longmapsto}\operatorname{H o m}_{\mathcal{C}}(W,X)\times\operatorname{H o m}_{\mathcal{C}}(W,Y){\mathrm{~i s~r e p r e s e r t a b l e}}.
\]

注意，(iii)和(iv)在条件(i)和(ii)下是等价的。此外，(iii) 和 (iv) 中给出的代表是同构的（参见练习 I.3）。我们用 $ X \oplus Y $ 来表示这样的代表，并将其称为X和Y的直和。同构：

\[
\operatorname{H o m}_{\mathfrak{e}}(X\oplus Y,X\oplus Y)\simeq\operatorname{H o m}_{\mathfrak{e}}(X,X\oplus Y)\oplus\operatorname{H o m}_{\mathfrak{e}}(Y,X\oplus Y)
\]

定义两个态射，与  $\mathrm{id}_{X\oplus Y}$  、  $i_1: X \to X \oplus Y$  和  $i_2: Y \to X \oplus Y$  关联。这些态射满足以下普遍性质。对于任意一对态射  $f: X \to W$  和  $g: Y \to W$  ，都存在  $h: X \oplus Y \to W$  使得下图可交换：

\begin{center}
\includegraphics[width=0.31\textwidth]{流行上的层_Chn-assets/img_in_image_box_405_829_785_1032.png}
\end{center}

通过反转箭头也有类似的评论。

如果  $F$  是从  $\mathcal{C}$  到  $\mathcal{C}'$  的函子，两个加法范畴，如果对于  $\mathrm{Ob}(\mathcal{C})$  的任何一对  $(X, Y)$  ，从  $\mathrm{Hom}_{\mathcal{C}}(X, Y)$  到  $\mathrm{Hom}_{\mathcal{C}}(F(X), F(Y))$  的映射  $F$  是群同态，则称  $F$  是加法。

现在我们假设 $\mathcal{C}$  是可加的。让 $Z \in \mathrm{Ob}(\mathcal{C})$  。函子  $\mathrm{Hom}_{\mathcal{C}}(Z, \cdot)$  将态射  $f: X \to Y$  与群同态相关联：

\[
\operatorname{H o m}_{\mathcal{C}}(Z,f)\colon\operatorname{H o m}_{\mathcal{C}}(Z,X)\to\operatorname{H o m}_{\mathcal{C}}(Z,Y)\enspace.
\]

有一个类似的定义：

\[
\operatorname{H o m}_{\mathcal{C}}(f,Z)\colon\operatorname{H o m}_{\mathcal{C}}(Y,Z)\to\operatorname{H o m}_{\mathcal{C}}(X,Z)\enspace.
\]

\statementlabel{定义 1.2.2}让 $ f \in \mathrm{Hom}_{\mathcal{C}}(X, Y) $  。

(i) 如果函子：

\[
\operatorname{Ker}(\operatorname{Hom}_{\mathcal{C}}(\cdot,f)):Z\mapsto\operatorname{Ker}(\operatorname{Hom}_{\mathcal{C}}(Z,f))=\left\{u\in\operatorname{Hom}_{\mathcal{C}}(Z,X);f\circ u=0\right\}
\]

是可表示的，其代表称为f的核，记为Ker f。 (ii) 同样，如果函子：

\[
\operatorname{Ker}(\operatorname{Hom}_{\mathcal{C}}(f,\cdot)):Z\mapsto\operatorname{Ker}(\operatorname{Hom}_{\mathcal{C}}(f,Z))=\left\{u\in\operatorname{Hom}_{\mathcal{C}}(Y,Z);u\circ f=0\right\}
\]

是可表示的，其代表称为f的cokernel，记作Coker f。

请注意，等式 Ker  $ f = 0 $ （或 Coker f = 0）相当于说 f 是单同态（或泛同态）。

假设 f 有一个内核。存在函子的态射：

\[
\beta:\mathrm{H o m}_{\mathcal{C}}(\cdot,\mathrm{K e r}\,f)\to\mathrm{H o m}_{\mathcal{C}}(\cdot,X)
\]

 $ \beta(\mathrm{id}_{\mathrm{Kerf}}) $  定义态射：

\[
\alpha:\mathbf{K e r}f\to X,
\]

这样 $ \beta = \mathrm{Hom}_{\mathcal{C}}(\cdot, \alpha) $  。

那么  $\alpha$  满足以下通用性质：对于任何态射  $e: W \to X$  使得  $f \circ e = 0$  、  $e$  通过  $\alpha$  因式分解，也就是说，下面的虚线箭头唯一存在（使下图可交换）：

\begin{center}
\includegraphics[width=0.26\textwidth]{流行上的层_Chn-assets/img_in_image_box_434_816_747_990.png}
\end{center}

类似地，如果  $f$  有一个 cokernel，则存在函子的态射：

\[
\delta:\mathrm{H o m}_{\mathcal{C}}(\mathrm{C o k e r}f,\cdot)\to\mathrm{H o m}_{\mathcal{C}}(Y,\cdot)~,
\]

其中定义：

\[
\gamma:Y\to\mathrm{Coker}f
\]

 $ \gamma $  是下图所示的通用问题的解决方案（其中  $ g \circ f = 0 $  ）：

\[
\begin{array}{c}X\xrightarrow[f]{}\begin{array}{c}Y\xrightarrow[g]{}Z\\\left.\vphantom{\sum}\right|_{\gamma}\\\downarrow\quad\end{array}\\ \mathrm{Coker}f~.\end{array}
\]

假设  $\alpha: \mathrm{Ker} f \to X$  有一个 cokernel。在这种情况下，用  $\mathcal{Coim} f$  表示它，并将其称为  $f$  的余像。类似地，假设  $\gamma: Y \to \mathrm{Coker} f$  有一个内核：用  $\mathrm{Im} f$  表示它，并将其称为  $f$  的图像。根据核和辅助核的通用性质，如果 $\mathrm{Coim} f$  和 $\mathrm{Im} f$  存在，则存在自然态射：

\[
\operatorname{Coim}f\to\operatorname{Im}f.
\]

\statementlabel{定义 1.2.3}如果满足以下两个条件，加法范畴 $ \mathcal{C} $  称为阿贝尔范畴。

(i) 对于任何态射  $ f: X \to Y $  ， Ker f 和 Coker f 都存在。

(ii) 规范态射  $ \operatorname{Coim} f \to \operatorname{Im} f $  是同构。

示例1.2.4。 (a) 让 $\mathcal{B}\text{an}(\mathbb{C})$  表示 $\mathbb{C}$  和线性连续映射上的Banach 向量空间的范畴。这是一个附加范畴。此外，如果  $f: X \to Y$  是一个态射， $f$  承认一个核和一个 cokernel（将  $Y/\overline{\mathrm{Im}f}$  作为一个 cokernel，其中  $\overline{\mathrm{Im}f}$  表示  $Y$  中  $\mathrm{Im}f$  的闭包）。请注意，可能会发生  $\mathrm{Ker}f = \mathrm{Coker}f = 0$  ，但  $f$  不是同构。因此范畴 $\mathcal{B}\text{an}(\mathbb{C})$  不是阿贝尔的。

(b) 设 $A$  为带单位的环。那么  $\mathfrak{Mod}(A)$  ，左  $A$  -模和  $A$  -线性映射的范畴是阿贝尔范畴。用 $A^{op}$  表示 $A$  的相反环，阿贝尔范畴 $\mathfrak{Mod}(A^{op})$  等于右 $A$  模的范畴。

事实上，对于某个环 A，任何阿贝尔范畴都等价于 Mod(A) 的满子范畴（参见 Mitchell [1]）。

(c) 非交换范畴的加性范畴的一个有趣的例子是过滤环上的过滤左模的范畴（参见第十一章）。

从现在开始，我们假设 $ \mathcal{C} $ 是一个阿贝尔范畴。

\statementlabel{定义 1.2.5}态射序列：

\[
X\xrightarrow[f]{}Y\xrightarrow[g]{}Z
\]

称为精确序列，如果：

（一） $ g \circ f = 0 $ ，

(ii) 自然态射 $ \mathrm{Im}f \to \mathrm{Ker}g $  是同构。

更一般地，如果任何连续的箭头对是精确的，则态射序列被称为精确的。

因此，如果  $f: X \to Y$  是态射，我们就得到精确的序列：

\[
0\to\mathbf{K e r}\;f\to X\to\mathrm{I m}\;f\to0,
\]

\[
0\to\operatorname{I m}f\to Y\to\operatorname{C o k e r}f\to0~.
\]

请注意，序列  $ 0 \to X \to Y $  （或  $ X \to Y \to 0 $  ）是精确的，当且仅当 f 是单同态（或同态）。

\statementlabel{定义 1.2.6}令 $\mathcal{C}$  和 $\mathcal{C}'$  为两个阿贝尔范畴。如果对于  $\mathcal{C}$  中的任何精确序列，从  $\mathcal{C}$  到  $\mathcal{C}'$  的加法函子  $F$  被称为左（或右）精确：

\[
0\to X^{\prime}\to X\to X^{\prime \prime}
\]

（分别： $ X' \to X \to X'' \to 0 $ ）序列：

\[
0\to F(X^{\prime})\to F(X)\to F(X^{\prime \prime})
\]

（分别： $ F(X') \to F(X) \to F(X'') \to 0 $ ）是准确的。

如果 F 既是左精确的又是右精确的，则 F 称为精确的。

从  $\mathcal{C}$  到  $\mathcal{C}'$  的逆变函子  $F$  称为左精确（或右精确、或精确），如果是这样，则  $F$  被视为从  $\mathcal{C}^\circ$  到  $\mathcal{C}'$  的函子。

示例 1.2.7。让 $ X \in \mathrm{Ob}(\mathcal{C}) $  。那么  $ \mathrm{Hom}_{\mathcal{C}}(X, \cdot) $  和  $ \mathrm{Hom}_{\mathcal{C}}(\cdot, X) $  都是从  $ \mathcal{C} $  到  $ \mathfrak{Mod}(\mathbb{Z}) $  的左精确函子。

\statementlabel{定义 1.2.8}让 $ X \in \mathrm{Ob}(\mathcal{C}) $  。如果函子  $ \mathrm{Hom}_{\mathcal{C}}(\cdot, X) $  （或  $ \mathrm{Hom}_{\mathcal{C}}(X,\cdot) $  ）是精确的，则可以说 X 是单射的（或射影的）。

请注意， $ Z \in \mathrm{Ob}(\mathcal{C}) $  是单射的，当且仅当，对于如下所有图表，其中行是精确的，虚线箭头存在（使图表可交换）：

\begin{center}
\includegraphics[width=0.25\textwidth]{流行上的层_Chn-assets/img_in_image_box_429_616_727_784.png}
\end{center}

事实上  $Z$  是单射的 iff 对于所有单态  $f: X \to Y$  ，  $\mathrm{Hom}_{\mathcal{G}}(f, Z)$  是满射的。

从这个评论中我们得知，如果  $ 0 \to X \to Y \to Z \to 0 $  是精确序列，并且如果  $ X $  是单射序列，则序列分裂（参见练习 I.5）。

此外，如果 $0 \to X' \to X \to X'' \to 0$  是 $\mathcal{C}$  中的精确序列并且 $X'$  是单射的，则 $X$  是单射的当且仅当 $X''$  是单射的。

通过反转箭头，对于射影物体也有类似的陈述。

术语 1.2.9。在阿贝尔范畴  $\mathcal{C}$  中，如果  $0 \to X \to Y \to Z \to 0$  是精确序列，有时会说  $X$  是  $Y$  的子对象， $Z$  是  $Y$  的商。有时甚至会写 $Z = Y/X$ 。

符号 1.2.10。我们将无差别地写 $ \operatorname{Mod}(\mathbb{Z}) $ 或 $ \mathbb{U} $ 来表示阿贝尔群的范畴。

如果  $A$  是一个环，而  $X$  是一个  $A$  模，则可以写  $\mathrm{Hom}_{A}(X,\cdot)$  而不是  $\mathrm{Hom}_{\mathrm{Rob}(A)}(X,\cdot)$  ， $\mathrm{Hom}_{A}(\cdot,X)$  也类似。

如果 $A$  是可交换的，则这些函子在 $\text{Mod}(A)$  中取它们的值，并且保留相同的符号来表示如此获得的从 $\text{Mod}(A)$  到 $\text{Mod}(A)$  的函子。对函子  $X\otimes_A \cdot \text{or} \cdot \otimes_A X$  进行类似的处理。

\subsection{复形范畴}
令 $C$  为加法范畴。

\statementlabel{定义 1.3.1}$\mathcal{C}$  中的复数  $X$  由数据  $\{X^n, \mathrm{d}_X^n\}_{n\in\mathbb{Z}}$  组成，因此对于任何  $n\in\mathbb{Z}$  ：

\[
X^{n}\in\mathrm{O b}({\mathcal C})\quad,\quad\mathtt{d}_{X}^{n}\in\mathrm{H o m}_{{\mathcal C}}(X^{n},X^{n+1})\quad{a n d}\quad\mathtt{d}_{X}^{n+1}\circ\mathtt{d}_{X}^{n}=0
\]

从复数  $X$  到复数  $Y$  的态射  $f$  是态射  $f^n:X^n\to Y^n$  的序列  $\{f^n\}_{n\in\mathbb{Z}}$  ，这样对于任何  $n$  ：

\[
\mathrm{d}_{Y}^{n}\circ f^{n}=f^{n+1}\circ\mathrm{d}_{X}^{n}.
\]

我们用 $ \mathbf{C}(\mathcal{C}) $ 表示由此获得的 $ \mathcal{C} $ 的复合体范畴。这也是一个加性范畴，如果 $ \mathcal{C} $  是阿贝尔范畴，那么 $ \mathbf{C}(\mathcal{C}) $  也是阿贝尔范畴。

人们经常把一个复数写成一个序列：

\[
X^{n-1}\xrightarrow{\mathrm{d}_{X}^{n-1}}X^{n}\xrightarrow{\mathrm{d}_{X}^{n}}X^{n+1}\longrightarrow\cdots.
\]

 $ \mathbf{d}_X = \{\mathbf{d}_X^n\}_{n} $  族被称为复数  $ X $  的微分。  如果  $ X^n = 0 $  对于  $ |n| \gg 0 $  （分别为  $ n \ll 0 $  和  $ n \gg 0 $  ），则  $ A $  复数  $ X $  被认为是有界的（分别为下界、上界）。  $ \mathbf{C}(\mathcal{C}) $  的完整子类由有界复合体（分别为下界复合体和上界复合体）组成，表示为  $ \mathbf{C}^b(\mathcal{C}) $ （分别为  $ \mathbf{C}^+(C) $  和  $ \mathbf{C}^-(C) $  ）。

我们将  $\mathcal{C}$  识别为  $\mathbf{C}(\mathcal{C})$  的满子范畴，由复合体  $X$  组成，使得  $X^n = 0$  对应  $n \neq 0$  。

\statementlabel{定义 1.3.2}令  $k$  为整数，并令  $X \in \mathrm{Ob}(\mathbf{C}(\mathcal{C}))$  。通过设置定义一个新的复数 $X[k]$ ：

\[
\left\{\begin{aligned}{X[k]^{n}}&{{}=X^{n+k},}\\ {\mathbf{d}_{X[k]}^{n}}&{{}=(-1)^{k}\mathbf{d}_{X}^{n+k}.}\\ \end{aligned}\right.
\]

对于  $ \mathbf{C}(\mathcal{C}) $  中的态射  $ f: X \to Y $  ，可以通过设置定义  $ f[k]: X[k] \to Y[k] $ ：

\[
f[k]^{n}=f^{n+k}.
\]

从 $ \mathbf{C}(\mathcal{C}) $ 到 $ \mathbf{C}(\mathcal{C}) $ 的函子[k]称为k次移位函子。

\statementlabel{定义 1.3.3}如果 $\mathcal{C}$  中存在态射 $s^n: X^n \to Y^{n-1}$ ，那么对于任何 $n$  ， $\mathbf{C}(\mathcal{C})$  中的态射 $f: X \to Y$  被称为同伦于零：

\[
f^{n}=s^{n+1}\circ\mathrm{d}_{X}^{n}+\mathrm{d}_{Y}^{n-1}\circ s^{n}.
\]

如果 $f - g$  是 $\text{homotopic}$  到零，则可以说 $f$  是 $\text{homotopic}$  到 $g$ 。我们用 $\text{Ht}(X, Y)$  表示由 $\text{homotopic}$  到零的态射组成的 $\text{Hom}_{\mathbf{C}(\mathcal{G})}(X, Y)$  子群。人们很容易看出，合成图 $\text{Hom}_{\mathbf{C}(\mathcal{G})}(X, Y) \times \text{Hom}_{\mathbf{C}(\mathcal{G})}(Y, Z) \to \text{Hom}_{\mathbf{C}(\mathcal{G})}(X, Z)$ 将 $\text{Ht}(X, Y) \times \text{Hom}_{\mathbf{C}(\mathcal{G})}(Y, Z)$ 和 $\text{Hom}_{\mathbf{C}(\mathcal{G})}(X, Y) \times \text{Ht}(Y, Z)$ 发送到 $\text{Ht}(X, Z)$ 。这允许定义一个新范畴 $\mathbf{K}(\mathcal{G})$ ，如下所示。

\statementlabel{定义 1.3.4}范畴 $ \mathbf{K}(\mathcal{C}) $  定义为

\[
\left\{\begin{aligned}{}&{{}\operatorname{O b}(\mathtt{K}(\mathcal{C}))=\operatorname{O b}(\mathtt{C}(\mathcal{C})),}\\ {}&{{}\operatorname{H o m}_{\mathtt{K}(\mathcal{C})}(X,Y)=\operatorname{H o m}_{\mathtt{C}(\mathcal{C})}(X,Y)/\operatorname{H t}(X,Y).}\\ \end{aligned}\right.
\]

类似地定义了范畴  $ \mathbf{K}^{b}(\mathcal{C}) $  、  $ \mathbf{K}^{+}(\mathcal{C}) $  和  $ \mathbf{K}^{-}(\mathcal{C}) $  。它们是  $ \mathbf{K}(\mathcal{C}) $  的满子范畴。

从现在开始，直到本节结束，我们假设 $ \mathcal{C} $  是阿贝尔算子。

\statementlabel{定义 1.3.5}对于  $ X \in \mathrm{Ob}(\mathbb{C}(\mathcal{C})) $  ，一组。

\[
Z^{k}(X)=\operatorname{K e r}\mathsf{d}_{X}^{k},\qquad B^{k}(X)=\operatorname{I m}\mathsf{d}_{X}^{k-1},
\]

\[
H^{k}(X)=\operatorname{C o k e r}(B^{k}(X)\to Z^{k}(X))~.
\]

我们称  $ H^{k}(X) $  为复数 X 的 k 次上同调。

换句话说：

\[
\begin{array}{r}{H^{k}(X)=\mathrm{K e r}\;{\bf d}_{X}^{k}/\mathrm{I m}\;{\bf d}_{X}^{k-1}}\end{array}.
\]

请注意，  $ H^k(\cdot) $  是从  $ \mathbf{C}(\mathcal{C}) $  到  $ \mathcal{C} $  的加法函子，并且：

\[
H^{k}(X)=H^{0}(X[k])~.
\]

如果 $f: X \to Y$  与零同伦，则 $H^k(f): H^k(X) \to H^k(Y)$  是零态射。因此  $H^k(\cdot)$  是从  $\mathbb{K}(\mathcal{C})$  到  $\mathcal{C}$  的明确定义的函子。

有精确的顺序：

\[
X^{k-1}\longrightarrow Z^{k}(X)\longrightarrow H^{k}(X)\longrightarrow0,
\]

\[
0\longrightarrow H^{k}(X)\longrightarrow\operatorname{C o k e r}(\mathbf{d}_{X}^{k-1})\longrightarrow X^{k+1},
\]

\[
0\longrightarrow Z^{k-1}(X)\longrightarrow X^{k-1}\longrightarrow B^{k}(X)\longrightarrow0,
\]

\[
0\longrightarrow B^{k}(X)\longrightarrow X^{k}\longrightarrow\operatorname{C o k e r}(\mathbf{d}_{X}^{k-1})\longrightarrow0,
\]

\[
0\to H^{k}(X)\to\operatorname{C o k e r}\mathbf{d}_{X}^{k-1}\xrightarrow[\mathbf{d_{v}^{k}}]{}Z^{k+1}(X)\xrightarrow{\quad}\:H^{k+1}(X)\xrightarrow{\quad}0\enspace.
\]

\statementlabel{命题 1.3.6}令  $ 0 \to X \to Y \to Z \to 0 $  为  $ \mathbf{C}(\mathcal{C}) $  中的精确序列。那么  $ \mathcal{C} $  中存在一个规范的长精确序列：

\[
\cdots\longrightarrow H^{n}(X)\longrightarrow H^{n}(Y)\longrightarrow H^{n}(Z)\underset{\delta}{\longrightarrow}H^{n+1}(X)\longrightarrow\cdots,
\]

更准确地说，如果

\begin{center}
\includegraphics[width=0.48\textwidth]{流行上的层_Chn-assets/img_in_image_box_288_1476_869_1649.png}
\end{center}

是  $\mathbb{C}(\mathcal{C})$  中精确序列的交换图，那么所有图：

\[
\begin{array}{c} H^{n}(Z) \longrightarrow H^{n+1}(X) \\ \downarrow \quad \downarrow \\  \downarrow \quad \downarrow \end{array}
\]

\[
H^{n}(Z^{\prime}) \quad \longrightarrow \quad H^{n+1}(X^{\prime})
\]

通勤。

\statementlabel{证明}考虑具有精确行的交换图：

\textbackslash{}[\textbackslash{}begin\{array\}\{c\}\{\textbackslash{}mathrm\{C o k e r\}(\textbackslash{}mathsf\{d\}\_\{X\}\textasciicircum{}\{n-1\})\textbackslash{}xrightarrow\{\textbackslash{}quad\}\textbackslash{}mathrm\{C o k e r\}(\textbackslash{}mathsf\{d\}\_\{Y\}\textasciicircum{}\{n-1\})\textbackslash{}xrightarrow\{\textbackslash{}quad\}\textbackslash{}mathrm\{C o k e r\}(\textbackslash{}mathsf\{d\}\_\{Z\}\textasciicircum{}\{n-1\})\textbackslash{}xrightarrow\{\textbackslash{}quad\}\textbackslash{}mathbb\{0\}\}\textbackslash{}\textbackslash{} \{\textbackslash{}begin\{array\}\{c\}\textbackslash{}left\textbackslash{}downarrow\textbackslash{}right\textbackslash{}downarrow\textbackslash{}mathsf\{d\}\_\{X\}\textasciicircum{}\{n\}\textbackslash{}\textbackslash{}\textbackslash{}end\{array\}\textbackslash{}quad\textbackslash{}left\textbackslash{}downarrow\textbackslash{}right\textbackslash{}downarrow\textbackslash{}mathsf\{d\}\_\{Y\}\textasciicircum{}\{n\}\textbackslash{}quad\textbackslash{}left\textbackslash{}downarrow\textbackslash{}right\textbackslash{}downarrow\textbackslash{}mathsf\{d\}\_\{Z\}\textasciicircum{}\{n\}\}\textbackslash{}\textbackslash{} \{\textbackslash{}begin\{array\}\{c c c c c\}\textbackslash{}mathbb\{0\}\&\textbackslash{}mathbb\{Z\}\textasciicircum{}\{n+1\}(X)\&\textbackslash{}mathbb\{X\}\textasciicircum{}\{n+1\}\textbackslash{}mathbb\{Y\}\&\textbackslash{}mathbb\{X\}\textasciicircum{}\{n+1\}(Z)。\&\textbackslash{}mathbb\{X\}\textasciicircum{}\{n+1\}(Z)\textbackslash{}mathbb\{Y\}\textbackslash{}mathbb\{X\}\textasciicircum{}\{

结果来自 (1.3.9) 和练习 I.9。构造的功能留给读者。 ⑨

让 $ X \in \mathrm{Ob}(\mathbb{C}(\mathcal{C})) $  。我们通过以下方式定义截断的复合体 $ \tau^{\leq n}(X) $ 和 $ \tau^{\geq n}(X) $ ：

\[
\tau^{\leqslant n}(X)\colon\cdots\to X^{n-2}\to X^{n-1}\to\operatorname{K e r}\mathrm{d}_{X}^{n}\to0\to\cdots,
\]

\[
\tau^{\geq n}(X)\colon\cdots\to0\to\mathrm{Coker}\,d_{X}^{n-1}\to X^{n+1}\to X^{n+2}\to\cdots.
\]

然后我们在  $ \mathbf{C}(\mathcal{C}) $  中有态射：

\[
\tau^{\leq n}(X)\to X\quad,\qquad X\to\tau^{\geq n}(X)
\]

对于  $ n' \leq n $  ：

\[
\tau^{\leqslant n^{\prime}}(X)\to\tau^{\leqslant n}(X),\qquad\tau^{\geqslant n^{\prime}}(X)\to\tau^{\geqslant n}(X).
\]

此外：

\statementlabel{命题 1.3.7}(i) 对于 k > n，自然态射  $ H^k(\tau^{\leq n}(X)) \to H^k(X) $  是  $ k \leq n $  和  $ H^k(\tau^{\leq n}(X)) = 0 $  的同构。

(ii) 对于 k < n，自然态射  $ H^k(X) \to H^k(\tau^{\geq n}(X)) $  是  $ k \geq n $  和  $ H^k(\tau^{\geq n}(X)) = 0 $  的同构。

\statementlabel{证明}很简单。

\statementlabel{备注 1.3.8}令  $X$  和  $Y$  为  $\mathbf{C}(\mathcal{C})$  的两个对象。有时有人会说，如果  $X$  和  $Y$  在  $\mathbf{K}(\mathcal{C})$  中同构，则它们是同伦等价的，也就是说，如果  $f \in \mathrm{Hom}_{\mathbf{C}(\mathcal{C})}(X, Y)$  存在  $\mathbf{K}(\mathcal{C})$  中的同构。这样的  $f$  称为同伦等价。

注释 1.3.9。 (i) 考虑序列 $\{X_n, \mathbf{d}_n^X\}_{n \in \mathbb{Z}}$ ，其中 $X_n \in \mathrm{Ob}(\mathcal{C})$ 、 $\mathbf{d}_n^X \in \mathrm{Hom}_{\mathcal{C}}(X_n, X_{n-1})$  和 $\mathbf{d}_n^X \circ \mathbf{d}_{n+1}^X = 0$ 。那么我们仍然会说这个序列是

事实上，设置  $ X^n = X_{-n} $  、  $ \mathrm{d}_X^n = \mathrm{d}_{-n}^X $  ，序列  $ \{X^n, \mathrm{d}_X^n\} $  就是我们之前意义上的复数。

(ii) 我们有时用 $ X' $ （分别 $ X $ ）表示复数 $ \{X^n, \mathrm{d}_X^n\} $ （分别 $ \{X_n, \mathrm{d}_n^X\} $ ）。对象 Ker  $ \mathrm{d}_{n-1}^X / \mathrm{Im} \mathrm{d}_n^X $  称为  $ X $  的  $ n $  全息群，并用  $ H_n(X) $  表示。

符号 1.3.10。在续集中，我们将设置：  $ \tau^{<n}(X) = \tau^{\leq(n-1)}(X) $  和  $ \tau^{>n}(X) = \tau^{\geq(n+1)}(X) $  。

\subsection{映射锥}
令  $\mathcal{C}$  为加法范畴，让  $f: X \to Y$  为  $\mathbf{C}(\mathcal{C})$  中的态射。

\statementlabel{定义 1.4.1}$f$  的映射锥，用  $M(f)$  表示，是  $\mathbf{C}(\mathcal{C})$  的对象，定义如下：

\[
\left\{\begin{array}{l}{M(f)^{n}=X^{n+1}\oplus Y^{n},}\\ {\mathbf{d}_{M(f)}^{n}=\left(\begin{matrix}{\mathbf{d}_{X[1]}^{n}}&{0}\\ {f^{n+1}}&{\mathbf{d}_{Y}^{n}}\\ \end{matrix}\right).}\\ \end{array}\right.
\]

回想一下  $ d_{X[1]}^{n} = -d_{X}^{n+1} $  。

我们通过以下方式定义态射 $ \alpha(f)\colon Y\to M(f) $ 和 $ \beta(f)\colon M(f)\to X[1] $ ：

\[
\alpha(f)^{n}=\left(\begin{array}{c}{0}\\ {\mathrm{i}\mathrm{d}_{\gamma^{n}}}\end{array}\right),
\]

\[
\beta(f)^{n}=(\mathrm{id}_{x^{n+1}},0).
\]

\statementlabel{引理 1.4.2}对于  $\mathbf{C}(\mathcal{C})$  中的任何  $f: X \to Y$  ，存在  $\phi: X[1] \to M(\alpha(f))$  使得：

\[
\phi{\mathrm{~i s~a n~i s o m o r p h i s m~i n~}}\mathbf{K}(\mathcal{C})\enspace.
\]

下图在  $ \mathbf{K}(\mathcal{C}) $  中进行通勤：

\[
\begin{array}{l l l l}{Y\xrightarrow{\alpha(f)}M(f)\xrightarrow{\beta(f)}X[1]}&{\xrightarrow{-f[1]}Y[1]}\\ {\left\downarrow\right\downarrow\operatorname{i d}_{Y}}&{\left\downarrow\right\downarrow\operatorname{i d}_{M(f)}}&{\left\downarrow\right\downarrow\phi}&{\left\downarrow\right\downarrow\operatorname{i d}_{Y[1]}}\\ {Y\xrightarrow{\alpha(f)}M(f)\xrightarrow{\alpha(\alpha(f))}M(\alpha(f))}&{\frac{}{}\beta(\alpha(f))}&{Y[1]\enspace.}\\ \end{array}
\]

请注意，这样的结果在  $ \mathbf{C}(\mathcal{C}) $  中不成立。进一步注意，即使在  $ \mathbf{K}(\mathcal{C}) $  中，  $ \phi $  也不是唯一的。这是许多尚未完全阐明的问题的根源。

\statementlabel{证明}我们有：

\[
M(\alpha(f))^{n}=Y^{n+1}\oplus M(f)^{n}=Y^{n+1}\oplus X^{n+1}\oplus Y^{n}~.
\]

我们通过以下方式定义 $ \phi^n: X[1]^n \to M(\alpha(f))^n $ 和 $ \psi^n: M(\alpha(f))^n \to X[1]^n $ ：

\[
\phi^{n}=\left(\begin{array}{c}{{-f^{n+1}}}\\ {{\mathrm{i d}_{X^{n+1}}}}\\ {{0}}\end{array}\right),\qquad\psi^{n}=(0,\mathrm{i d}_{X^{n+1}},0).
\]

然后引理由以下观察得出。

(a)  $ \phi = (\phi^n)_n $  和  $ \psi = (\psi^n)_n $  是复形的态射，

（b） $ \psi \circ \phi = id_{x_{[1]}} $ ，

(c)  $ \phi \circ \psi $  与  $ id_{M(\alpha(f))} $  同伦，

（d） $ \psi \circ \alpha(\alpha(f)) = \beta(f) $ ，

（e） $ \beta(\alpha(f)) \circ \phi = -f[1] $ 。

除 (c) 外，所有这些属性都可以直接检查。为了得到(c)，我们通过以下方式定义 $ s^n: M(\alpha(f))^n \to M(\alpha(f))^{n-1} $ ：

\[
s^{n}=\left(\begin{array}{l l l}{0}&{0}&{\mathsf{i d}_{\mathsf{Y}^{n}}}\\ {0}&{0}&{0}\\ {0}&{0}&{0}\end{array}\right).
\]

然后有人验证：

\[
\mathrm{i d}_{M(\alpha(f))^{n}}-\phi^{n}\circ\psi^{n}=s^{n+1}\circ\mathrm{d}_{M(\alpha(f))}^{n}+\mathrm{d}_{M(\alpha(f))}^{n-1}\circ s^{n}~.\quad\Box
\]

将  $ \mathbf{K}(\mathcal{C}) $  中的三角形定义为态射  $ X \to Y \to Z \to X[1] $  的序列，将三角形的态射定义为  $ \mathbf{K}(\mathcal{C}) $  中的交换图：

\[
\begin{array}{c c c c c c c}{X}&{\longrightarrow}&{Y}&{\longrightarrow}&{Z\quad\longrightarrow\quad}&{X[1]}\\ {\phi\Big\downarrow}&{}&{\Big\downarrow}&{}&{\Big\downarrow}&{\Big\downarrow}\\ {\downarrow}&{}&{\downarrow}&{}&{\downarrow}&{\downarrow}\\ {X^{\prime}}&{\longrightarrow}&{Y^{\prime}}&{\longrightarrow}&{Z^{\prime}}&{\longrightarrow}&{X^{\prime}[1]~.}\\ \end{array}
\]

\statementlabel{定义 1.4.3}对于  $\mathbf{C}(\mathcal{C})$  中的某些  $f$  来说， $\mathbf{K}(\mathcal{C})$  中的三角形  $X \to Y \to Z \to X$  [1] 如果与三角形  $X^{\prime} \xrightarrow{f} Y^{\prime} \xrightarrow{\alpha(f)} M(f) \xrightarrow{\beta(f)} X^{\prime}[1]$  同构，则称为杰出三角形。

\statementlabel{命题 1.4.4}$ \mathbf{K}(\mathcal{C}) $  中的杰出三角形的集合满足以下属性：(TR 0)–(TR 5)。

(TR 0) 区分与已区分三角形同构的三角形。

(TR 1) 对于任何 $ X \in \text{Ob}(\mathbb{K}(\mathcal{C})) $  ， $ X \xrightarrow{\text{idx}} X \longrightarrow 0 \longrightarrow X[1] $  是一个杰出的三角形。

(TR 2)  $\mathbb{K}(\mathcal{C})$  中的任何  $f: X \to Y$  都可以嵌入到一个独特的三角形  $X \xrightarrow{f} Y \to Z \to X[1]$  中。

(TR 3)  $ X \xrightarrow{f} Y \xrightarrow{g} Z \xrightarrow{h} X $  [1] 是一个可区分三角形当且仅当  $ Y \xrightarrow{g} Z \xrightarrow{h}  X $  [1]  $ \xrightarrow{-f[1]}  Y $  [1] 是一个特殊的三角形。

(TR 4) 给定两个不同的三角形 $ X \xrightarrow{f} Y \to Z \to X $  [1] 和  $ X' \xrightarrow{f'} Y' \to Z' \to X' $  [1]，交换图

\[
\begin{array}{l}X\xrightarrow{f}Y\\\left|\begin{array}{l}u\\\end{array}\right.\\\left|\begin{array}{ccc}v&\\&\\v&\\&\\X^{\prime}\xrightarrow{f^{\prime}}&Y^{\prime}\end{array}\right.\\\end{array}
\]

可以嵌入三角形的态射中（不一定是唯一的）。 (TR 5)（八面体公理）。假设给定不同的三角形：

\[
x\xrightarrow{\quad f\quad}Y\xrightarrow{\quad}\longrightarrow Z^{\prime}\xrightarrow{\quad}\longrightarrow X[1]~,
\]

\[
Y\xrightarrow{\quad g\quad}Z\xrightarrow{\quad}\cdots\xrightarrow{}X^{\prime}\xrightarrow{\quad}\cdots\longrightarrow Y[1]~,
\]

\[
{\cal X}\xrightarrow{g\circ f}{\cal Z}\xrightarrow{\quad}\;{\cal Y}^{\prime}\xrightarrow{\quad}\;{\cal X}[1]\;,
\]

那么存在一个显着的三角形

\[
Z^{\prime}\to Y^{\prime}\to X^{\prime}\to Z^{\prime}[1]
\]

使得下图是可交换的：

\begin{center}
\includegraphics[width=0.45\textwidth]{流行上的层_Chn-assets/img_in_image_box_350_1075_895_1545.png}
\end{center}

\statementlabel{证明}性质(TR 0) 和(TR 2) 是显而易见的，并且(TR 3) 由引理1.4.2 得出。

由于 $f:0\to X$ 的映射锥是 $X$ ，所以三角形 $0\longrightarrow X\xrightarrow{\mathrm{id}x}X\longrightarrow0[1]$ 被区分。应用 (TR 3) 我们得到 (TR 1)。让我们证明（TR 4）。我们可以假设 $X\xrightarrow{f}Y\to Z\to X[1]$ 和 $X'\xrightarrow{f'}Y'\to Z'\to X'[1]$ 是 $X\xrightarrow{f}Y\xrightarrow{\alpha(f)}$ 

分别为 $ M(f) \xrightarrow{\beta(f)} X[1] $  和 $ X' \xrightarrow{f'} Y' \xrightarrow{\alpha(f')} M(f') \xrightarrow{\beta(f')} X'[1] $  。我们将构造一个态射 $ w: M(f) \to M(f') $ ，使得：

\[
\left\{\begin{aligned}{}&{{}w\circ\alpha(f)=\alpha(f^{\prime})\circ v,}\\ {}&{{}u[1]\circ\beta(f)=\beta(f^{\prime})\circ w.}\\ \end{aligned}\right.
\]

根据  $\mathbf{K}(\mathcal{C})$  的定义，存在  $s^n: X^n \to Y^{n-1}$  使得  $v^n \circ f^n - f'^n \circ u^n = s^{n+1} \circ \mathrm{d}_X^n + \mathrm{d}_{Y'}^{n-1} \circ s^n$  。我们通过以下方式定义 $w^n: M(f)^n = X^{n+1} \oplus Y^n \to M(f')^n = X^{n+1} \oplus Y'^n$ ：

\[
w^{n}=\left(\begin{array}{l l}{{u^{n+1}}}&{{0}}\\ {{s^{n+1}}}&{{v^{n}}}\end{array}\right).
\]

然后直接计算可知w是复形的态射并且满足(1.4.4)。

让我们证明（TR 5）。我们可以假设  $ Z' = M(f) $  、  $ X' = M(g) $  和  $ Y' = M(g \circ f) $  。让我们通过以下方式定义 $ u: Z' \to Y' $ 和 $ v: Y' \to X' $ ：

\[
u^{n}\colon X^{n+1}\oplus Y^{n}\to X^{n+1}\oplus Z^{n},\qquad u=\left(\begin{matrix}{{\mathrm{i d}}_{X^{n+1}}}&{0}\\ {0}&{g^{n}}\\ \end{matrix}\right),
\]

\[
v^{n}\colon X^{n+1}\oplus Z^{n}\to Y^{n+1}\oplus Z^{n},\qquad v^{n}=\left(\begin{matrix}f^{n+1}&0\\ 0&\mathsf{i d}_{Z^{n}}\end{matrix}\right).
\]

我们将  $w: X' \to Z' [1]$  定义为复合  $X' \to Y [1] \to Z' [1]$  。那么(TR 5)中的图是可交换的，足以证明 $Z' \xrightarrow{u} Y' \xrightarrow{v} X' \xrightarrow{w} Z' [1]$ 是一个杰出的三角形。为此，我们将构造一个同构  $\phi: M(u) \to X'$  及其逆  $\psi: X' \to M(u)$  ，使得  $\phi \circ \alpha(u) = v$  和  $\beta(u) \circ \psi = w$  。我们有：

\[
M(u)^{n}=M(f)^{n+1}\oplus M(g\circ f)^{n}=X^{n+2}\oplus Y^{n+1}\oplus X^{n+1}\oplus Z^{n}
\]

和 $ X'' = M(g)^n = Y^{n+1} \oplus Z^n $ 。我们通过以下方式定义 $ \phi $ 和 $ \psi $ ：

\[
\phi^{n}=\left(\begin{matrix}{0}&{\mathsf{i d}_{{Y}^{n+1}}}&{f^{n+1}}&{0}\\ {0}&{0}&{0}&{\mathsf{i d}_{Z^{n}}}\\ \end{matrix}\right)\;,\qquad\psi^{n}=\left(\begin{matrix}{0}&{0}\\ {\mathsf{i d}_{{Y}^{n+1}}}&{0}\\ {0}&{0}\\ {0}&{\mathsf{i d}_{X^{n+1}}}\\ \end{matrix}\right)\;.
\]

然后可以轻松检查  $\phi$  和  $\psi$  是复数和  $\phi \circ \alpha(u) = v, \beta(u) \circ \psi = w$  的态射。

我们有 $\phi \circ \psi = id_{X^{\prime}}$ 。如果我们定义：

\[
s^{n}\colon M(u)^{n}\to M(u)^{n-1}\quad,\qquad s^{n}=\left(\begin{matrix}{0}&{0}&{\mathsf{i d}_{X^{n+1}}}&{0}\\ {0}&{0}&{0}&{0}\\ {0}&{0}&{0}&{0}\\ {0}&{0}&{0}&{0}\\ \end{matrix}\right)
\]

那么：

\[
\begin{array}{r}{(\mathrm{i d}_{M(u)}-\psi\circ\phi)^{n}=s^{n+1}\circ\mathrm{d}_{M(u)}^{n}+\mathrm{d}_{M(u)}^{n-1}\circ s^{n}.}\end{array}
\]

因此  $ \psi \circ \phi $  等于  $ \mathbf{K}(\mathcal{C}) $  中的  $ id_{M(u)} $  。 □

\statementlabel{备注 1.4.5}属性 (TR 5) 可以通过以下八面体图来可视化：

\begin{center}
\includegraphics[width=0.28\textwidth]{流行上的层_Chn-assets/img_in_image_box_394_441_737_738.png}
\end{center}

\begin{center}图1.4.1\end{center}

\subsection{三角范畴}
我们通过抽象  $ \mathbf{K}(\mathcal{C}) $  的属性获得三角范畴的概念。

令  $\mathcal{C}$  为加法范畴，并具有自同构  $T:\mathcal{C}\to\mathcal{C}$  。我们有时为  $T$  编写 [1] ，为  $T^k$  编写  $[k]$  ，（即  $X[1]$  为  $T(X)$  ，或  $f[1]$  为  $T(f)$  ）。

 $ \mathcal{C} $  中的三角形是态射序列

\[
X\to Y\to Z\to T(X)~.
\]

\statementlabel{定义 1.5.1}三角范畴 $ \mathcal{C} $  由以下数据和规则组成。

(1.5.1) 加法范畴  $ \mathcal{C} $  和自同构  $ T: \mathcal{C} \to \mathcal{C} $  ，

（1.5.2）三角形族，称为杰出三角形。

当设置 $ X[1] = T(X) $ 时，这些数据满足命题1.4.4的公理(TR 0)-(TR 5)。

令 $(\mathcal{C}, T)$  和 $(\mathcal{C}^{\prime}, T^{\prime})$  为两个三角范畴。我们说，如果  $F \circ T \simeq T^{\prime} \circ F$  和  $F$  将  $\mathcal{C}$  的可区分三角形发送到  $\mathcal{C}^{\prime}$  的可区分三角形，则从  $\mathcal{C}$  到  $\mathcal{C}^{\prime}$  的加法函子  $F$  是三角范畴的函子。

显然，对于加法范畴  $\mathcal{C}$  ，  $\mathbf{K}(\mathcal{C})$  是一个三角范畴。现在让 $\mathcal{C}$  成为三角范畴，让 $\mathcal{A}$  成为阿贝尔范畴。

\statementlabel{定义 1.5.2}如果对于任何可区分的三角形  $X \to Y \to Z \to T(X)$  ，序列  $F(X) \to F(Y) \to F(Z)$  是精确的，则加法函子  $F: \mathcal{C} \to \mathcal{A}$  称为上同调函子。

对于上同调函子  $F$  ，我们将  $F \circ T^k$  写为  $F^k$  。然后对于任何杰出的三角形 $X \to Y \to Z \to T(X)$ ，我们获得一个长的精确序列：

\[
\cdots\to F^{k-1}(Z)\to F^{k}(X)\to F^{k}(Y)\to F^{k}(Z)\to F^{k+1}(X)\to\cdots.
\]

\statementlabel{命题 1.5.3}(i) 如果  $ X \xrightarrow{J} Y \xrightarrow{g} Z \to T(X) $  是一个杰出的三角形，则  $ g \circ f = 0 $  。

(ii) 对于任何 $W \in \mathrm{Ob}(\mathcal{C})$  ， $\mathrm{Hom}_{\mathcal{C}}(W, \cdot)$  和 $\mathrm{Hom}_{\mathcal{C}}(\cdot, W)$  是上同调函子。

\statementlabel{证明}(i) 根据(TR 1)， $ X \xrightarrow{\mathrm{id}_{X}} X \longrightarrow 0 \longrightarrow T(X) $  是一个显着的三角形。因此，通过 (TR 4) 存在态射 $ \phi: 0 \to Z $ ，它使得下图可交换：

\[
\begin{array}{c}X\xrightarrow{\quad}X\quad\xrightarrow{\quad}0\quad\xrightarrow{\quad}T(X)\\\mathrm{id}_{x}\bigg\downarrow\quad\begin{array}{c}f\\\bigg\downarrow\\\quad\scriptstyle X\xrightarrow{\quad f\quad}Y\quad\begin{array}{c}g\\\quad\scriptstyle Y\xrightarrow{\quad g\quad}Z\quad\begin{array}{c}f\\\quad\scriptstyle Z\xrightarrow{\quad}T(X)\quad\end{array}\end{array}\end{array}\end{array}
\]

因此 $ g \circ f = \phi \circ 0 = 0 $  。

(ii) 令 $X \xrightarrow{f} Y \xrightarrow{\theta} Z \to T(X)$  为一个独特的三角形。为了证明  $\operatorname{Hom}_{\mathcal{C}}(W, \cdot)$  是上同调函子，只需证明对于任何  $\phi \in \operatorname{Hom}_{\mathcal{C}}(W, Y)$  和  $g \circ \phi = 0$  ，我们都可以找到  $\psi \in \operatorname{Hom}_{\mathcal{C}}(W, X)$  和  $\phi = f \circ \psi$  。这是从 (TR 1)、(TR 3) 和 (TR 4) 得出的，这意味着下面的虚线箭头可以完成：

\begin{center}
\includegraphics[width=0.42\textwidth]{流行上的层_Chn-assets/img_in_image_box_325_1299_827_1486.png}
\end{center}

 $ \operatorname{Hom}_{\mathcal{C}}(\cdot,W) $  是上同调函子的证明是类似的。 ⑨

\statementlabel{备注 1.5.4}令 $\mathcal{C}$  为加法范畴， $f: X \to Y$  为 $\mathbf{C}(\mathcal{C})$  中的态射。一般来说，复合  $X \xrightarrow{f} Y \xrightarrow{\alpha(f)} M(f)$  在  $\mathbf{C}(\mathcal{C})$  中不为零，而仅在  $\mathbf{K}(\mathcal{C})$  中。

\statementlabel{推论 1.5.5}让

\begin{center}
\includegraphics[width=0.42\textwidth]{流行上的层_Chn-assets/img_in_image_box_335_201_839_377.png}
\end{center}

是可区分三角形的态射。如果  $\phi$  和  $\psi$  是同构，那么  $\theta$  也是同构。

\statementlabel{证明}对于任何  $W \in \mathrm{Ob}(\mathcal{C})$  ，让我们将函子  $\mathrm{Hom}_{\mathcal{C}}(W, \cdot)$  应用于上图。我们得到一个交换图，其行是精确的。由于  $\mathrm{Hom}_{\mathcal{C}}(W, \phi)$  和  $\mathrm{Hom}_{\mathcal{C}}(W, \psi)$  是同构， $\mathrm{Hom}_{\mathcal{C}}(W, T(\phi))$  和  $\mathrm{Hom}_{\mathcal{C}}(W, T(\psi))$  也是同构，我们通过练习 1.8 得出  $\mathrm{Hom}_{\mathcal{C}}(W, \theta)$  是同构。根据命题 1.1.8，这就完成了证明。   $\square$ 

\statementlabel{命题 1.5.6}令 $\mathcal{C}$  为阿贝尔范畴。那么函子 $H^0(\cdot): \mathbf{K}(\mathcal{C}) \to \mathcal{C}$  是一个上同调函子。

\statementlabel{证明}足以证明如果  $f: X \to Y$  是  $\mathbf{C}(\mathcal{C})$  中的态射，则序列

\[
H^{0}(Y)\to H^{0}(M(f))\to H^{0}(X[1])
\]

是准确的。

由于  $0 \to Y \to M(f) \to X[1] \to 0$  是  $\mathbf{C}(\mathcal{C})$  中的精确序列，因此结果来自命题 1.3.6。   $\square$ 

\statementlabel{定义 1.5.7}令  $\mathcal{C}$  为阿贝尔范畴，让  $f: X \to Y$  为  $\mathbf{K}(\mathcal{C})$  中的态射。如果  $H^n(f)$  是每个  $n$  的同构，则可以说  $f$  是准同构（简称 qis）。

因此  $f$  是一个  $q$  当且仅当每个  $n$  都有  $H^n(M(f)) = 0$  。如果  $f$  是  $q$  是，则简写为  $X \xrightarrow[qis]{} Y$  。

注释 1.5.8。让 $\mathcal{C}$  成为一个三角范畴。在后续部分中，我们将经常写 $X \longrightarrow Y \longrightarrow Z \xrightarrow[+1]{}$ 而不是 $X \to Y \to Z \to T(X)$ ，以表示一个特殊的三角形。

\statementlabel{定义 1.5.9}让 $\mathcal{C}$  成为一个三角范畴。  如果三角形 $X \xrightarrow{\alpha} Y \xrightarrow{\beta} Z \xrightarrow{-\gamma} T(X)$  是可区分的， $\mathcal{C}$  中的 $A$  三角形 $X \xrightarrow{\alpha} Y \xrightarrow{\beta} Z \xrightarrow{\gamma}T(X)$  称为反可区分的。

请注意，范畴  $\mathcal{C}$  具有反可区分三角形族，是一个三角范畴。我们用 $\mathcal{C}^{a}$  表示它。

\subsection{范畴的局部化}
令 $C$  为一个范畴，并让 $S$  为 $C$  中的态射族。

\statementlabel{定义 1.6.1}如果满足下面的(S 1)-(S 4)，则称S 是乘法系统。

(S 1) 对于任何 $ X \in \text{Ob}(\mathcal{C}) $  、 $ \text{id}_X \in S $  。

(S2) 对于 S 的任何对  $ (f,g) $ ，使得组合  $ g\circ f $  存在， $ g\circ f\in S $  。

(S 3) 任何图表：

\begin{center}
\includegraphics[width=0.15\textwidth]{流行上的层_Chn-assets/img_in_image_box_526_557_704_731.png}
\end{center}

使用  $g \in S$  ，可以完成交换图：

\begin{center}
\includegraphics[width=0.15\textwidth]{流行上的层_Chn-assets/img_in_image_box_526_842_704_1017.png}
\end{center}

与 $ h \in S $  。同上，所有箭头都相反。

(S4) 如果 f 和 g 属于  $ \mathrm{Hom}_{\mathcal{C}}(X, Y) $  ，则以下条件等效：

(i) 存在  $t: Y \to Y', t \in S$  ，使得  $t \circ f = t \circ g$  ，

(ii) 存在  $ s: X' \to X $  、  $ s \in S $  ，使得  $ f \circ s = g \circ s $  。

\statementlabel{定义 1.6.2}令 $\mathcal{C}$  为范畴， $S$  为乘法系统。范畴  $\mathcal{C}_S$  称为  $S$  的  $\mathcal{C}$  本地化，定义如下：

\[
\mathsf{O b}(\mathcal{C}_{S})=\mathsf{O b}(\mathcal{C})\enspace,
\]

(1.6.2) 对于  $ \mathrm{Ob}(\mathcal{C}) $  中的任何对  $ (X, Y) $  ，

\[
\operatorname{H o m}_{\mathcal{G}_{S}}(X,Y)=\left\{(X^{\prime},s,f);X^{\prime}\in\operatorname{O b}(\mathcal{C}),s:X^{\prime}\to X,f\colon X^{\prime}\to Y,s\in S\right\}/\mathcal{R}
\]

其中 R 是以下等价关系：

\[
(X^{\prime},s,f)\mathcal{R}(X^{\prime \prime},t,g)
\]

iff 存在交换图

\begin{center}
\includegraphics[width=0.28\textwidth]{流行上的层_Chn-assets/img_in_image_box_424_162_761_465.png}
\end{center}

与 $ u \in S $  。

 $(X', s, f) \in \mathrm{Hom}_{\mathcal{C}_s}(X, Y)$ 和 $(Y', t, g) \in \mathrm{Hom}_{\mathcal{C}_s}(Y, Z)$ 的组成定义如下。我们用(S3)求交换图：

\begin{center}
\includegraphics[width=0.49\textwidth]{流行上的层_Chn-assets/img_in_image_box_296_628_888_850.png}
\end{center}

使用  $ t' \in S $  ，我们设置：

\[
(Y^{\prime},t,g)\circ(X^{\prime},s,f)=(X^{\prime \prime},s\circ t^{\prime},g\circ h)\enspace.
\]

使用公理 (S 1)–(S 4)，我们很容易看出  $\mathcal{C}_{S}$  是一个范畴。

我们用 Q 表示函子：

\[
Q:\mathcal{C}\to\mathcal{C}_{S}
\]

由  $ Q(X) = X \text{ for } X \in \text{Ob}(\mathcal{C}), \text{and } Q(f) = (X, \mathrm{id}_{X}, f) \text{for } f \in \mathrm{Hom}_{\mathcal{C}}(X, Y) $  定义。

\statementlabel{命题 1.6.3}(i) 对于  $ s \in S $  ，  $ Q(s) $  是  $ \mathcal{C}_S $  的同构。

(ii) 令  $\mathcal{C}'$  为另一个范畴， $F:\mathcal{C} \to \mathcal{C}'$  为函子，使得  $F(s)$  是所有  $s \in S$  的同构。然后  $F$  通过  $Q$  进行唯一因子分解。

\statementlabel{证明}很简单。

\statementlabel{备注 1.6.4}它源自命题 1.6.3。那：

\[
({\mathcal C}^{\circ})_{s}\simeq({\mathcal C}_{S})^{\circ}~.
\]

因此，通过将条件（1.6.2）替换为以下内容，我们得到了与  $ \mathcal{C}_{s} $  等效的范畴：

 $ (1.6.2)' \ \mathrm{Hom}_{\mathcal{G}_5}(X, Y) = \{(Y', t, g); Y' \in \mathrm{Ob}(\mathcal{C}), t: Y \to Y', g: X \to Y', t \in S\} / \mathcal{R}' $ 

其中  $ \mathcal{R}' $  的定义与  $ \mathcal{R} $  类似。

\statementlabel{命题 1.6.5}让 $\mathcal{C}$  成为一个范畴， $\mathcal{C}'$  成为一个完整的子范畴。令 $S$  为 $\mathcal{C}$  中的乘法系统，并令 $S'$  为属于 $S$  的 $\mathcal{C}'$  的态射族。假设  $S'$  是  $\mathcal{C}'$  中的乘法系统，并假设以下条件之一成立：

(i) 只要  $f: X \to Y$  是  $S$  中的态射，对于  $Y \in \mathrm{Ob}(\mathcal{C}')$  ，就存在  $g: W \to X$  ，对于  $W \in \mathrm{Ob}(\mathcal{C}')$  和  $f \circ g \in S$  ，

(ii) 与 (i) 相同，但箭头相反。

那么本地化  $\mathcal{C}_{s}^{\prime}$  是  $\mathcal{C}_{s}$  的满子范畴。

当  $ \mathcal{C} $  是三角范畴时，可以相对于  $ \mathcal{C} $  的一系列对象来本地化  $ \mathcal{C} $  。

\statementlabel{定义 1.6.6}令 $\mathcal{C}$  为三角范畴，并令 $\mathcal{N}$  为 $\mathrm{Ob}(\mathcal{C})$  的子族。如果 $\mathcal{N}$  满足下面的(N 1)–(N3)，则称 $\mathcal{N}$  为空系统。

（N1） $ 0 \in \mathcal{N} $ ，

(N2)  $X \in \mathcal{N}$  当且仅当  $X[1] \in \mathcal{N}$  ，

(N 3) 如果  $X \to Y \to Z \to X$  [1] 是一个杰出的三角形，且  $X \in \mathcal{N}$  、  $Y \in \mathcal{N}$  ，则  $Z \in \mathcal{N}$  。

现在我们设置：

\[
\begin{cases}{S(\mathcal{N})=\{f:X\to Y;f\mathrm{~i s~e m b e d d e d~i n t o~a~d i s t i n g u i s h e d}}\\ {\operatorname{t r i a n g l e}~X\xrightarrow{f}Y\to Z\to X[1],\mathrm{w i t h}~Z\in\mathcal{N}\}.}\\ \end{cases}
\]

\statementlabel{命题 1.6.7}假设 $\mathcal{N}$  是一个空系统。那么 $S(\mathcal{N})$  是一个乘法系统。

\statementlabel{证明}性质(S 1) 是从(N 1) 和(TR 1) 推导出来的。让我们证明（S 2）。令  $ X \xrightarrow{f} Y \to Z' \to X[1] $  和  $ Y \xrightarrow{g} Z \to X' \to Y[1] $  为两个不同的三角形，其中  $ X' \in \mathcal{N} $  、  $ Z' \in \mathcal{N} $  。通过(TR 2)，存在一个显着的三角形 $ X \xrightarrow{\theta \circ f} Z \longrightarrow Y' \longrightarrow X[1] $ ，并且通过(TR 5)，存在一个显着的三角形 $ Z' \to Y' \to X' \to Z'[1] $ 。通过 (N 2)、(N 3) 和 (TR 3)，我们有：  $ Y' \in \mathcal{N} $  。因此 $ g \circ f \in S(\mathcal{N}) $  。

为了证明 (S 3)，请考虑一个杰出的三角形  $ Z \xrightarrow{g} Y \xrightarrow{k} X' \to Z[1] $  和  $ X' \in \mathcal{N} $  ，并让  $ f: X \to Y $  。存在一个显着的三角形

\[
W\longrightarrow X\xrightarrow{k\circ f}X^{\prime}\longrightarrow W[1]~.
\]

然后通过 (TR 4) 和 (TR 3)，我们得到了可区分三角形的态射：

\begin{center}
\includegraphics[width=0.41\textwidth]{流行上的层_Chn-assets/img_in_image_box_334_1500_834_1692.png}
\end{center}

由于  $ X' \in \mathcal{N} $  ， h 属于  $ S(\mathcal{N}) $  。

通过反转箭头也可以得到类似的证明。

最后我们证明（S4）。让  $f: X \to Y$  和  $t: Y \to Y'$  ，以及  $t \in S(\mathcal{N})$  和  $t \circ f = 0$  。我们将证明存在 $s: X' \to X$ 、 $s \in S(\mathcal{N})$ 、 $f \circ s = 0$ 。令  $Z \xrightarrow{g} Y \xrightarrow{t} Y' \to Z[1]$  为一个杰出的三角形，其中  $Z \in \mathcal{N}$  。根据 (TR 1)、(TR 3)、(TR 4)，存在  $h: X \to Z$  使得  $f = g \circ h$  。如果我们将 $h$ 嵌入到一个独特的三角形 $X' \xrightarrow{s} X \xrightarrow{h} Z \to X' [1]$ 中，那么 $s$ 将满足所需的属性。逆蕴涵的证明类似。   $\square$ 

符号 1.6.8。令  $\mathcal{C}$  为三角范畴， $\mathcal{N}$  为  $\mathcal{C}$  中的空系统。我们写 $\mathcal{C}/\mathcal{N}$ 而不是 $\mathcal{C}_{\mathcal{S}(\mathcal{N})}$ 。

\statementlabel{命题 1.6.9}令 $ \mathcal{C} $  为三角范畴， $ \mathcal{N} $  为空系统。

(i) 通过将那些与 C 中的可区分三角形的图像同构的三角形作为可区分的三角形，C/N 成为三角范畴。

(ii) 用 $ Q $  表示自然函子 $ \mathcal{C} \to \mathcal{C}/N $  。我们有  $ Q(X) \simeq 0 $  对应  $ X \in \mathcal{N} $  。

(iii) 三角范畴的任何函子  $F: \mathcal{C} \to \mathcal{C}'$  ，使得所有  $X \in \mathcal{N}$  的  $F(X) \simeq 0$  都通过  $Q$  进行唯一因子分解。

命题 1.6.3 的证明是显而易见的。

\statementlabel{命题 1.6.10}令 $\mathcal{C}$  为三角范畴， $\mathcal{N}$  为 $\mathcal{C}$  中的空系统， $\mathcal{C}'$  为 $\mathcal{C}$  的完整三角化子范畴，使得 $\mathcal{C}$  中的任何可区分三角形 $X \to Y \to Z \to X$  [1] 与 $X \in \mathrm{Ob}(\mathcal{C}')$  、 $Y \in \mathrm{Ob}(\mathcal{C}')$  都是 $\mathcal{C}'$  中的可区分三角形。让 $\mathcal{N}' = \mathcal{N} \cap \mathrm{Ob}(\mathcal{C}')$  。然后：

(i)  $\mathcal{N}^{\prime}$  是  $\mathcal{C}^{\prime}$  中的空系统。

(ii) 此外，假设  $\mathcal{C}$  中的任何态射  $Y \to Z$  与  $Y \in \mathrm{Ob}(\mathcal{C}')$  、  $Z \in \mathcal{N}$  通过对象  $\mathcal{N} \cap \mathrm{Ob}(\mathcal{C}')$  进行因式分解。那么  $\mathcal{C}'/\mathcal{N}'$  是  $\mathcal{C}/\mathcal{N}$  的满子范畴。

\statementlabel{证明}（一）明确。

(ii) 我们将验证命题 1.6.5 的条件 (i)。

令  $X \to Y \to Z \to X[1]$  与  $Y \in \mathrm{Ob}(\mathcal{C}')$  、  $Z \in \mathcal{N}$  成为一个可区分的三角形。根据假设，态射  $Y \to Z$  通过  $Y \to Z'\to Z$  和  $Z' \in \mathcal{N} \cap \mathrm{Ob}(\mathcal{C}')$  进行因式分解。将 Axiom TR 5 应用于态射  $Y \to Z'$  和  $Z' \to Z$  我们发现一个显着的三角形  $Y \to Z' \to W \to Y[1]$  使得  $W[-1] \to Y$  因式分解为  $W[-1] \to X \to Y$  。这样就完成了证明。

\statementlabel{备注 1.6.11}令  $\mathcal{N}$  为  $\mathcal{C}$  中的空系统， $Q$  函子  $\mathcal{C} \to \mathcal{C}/\mathcal{N}$  。那么  $X \in \mathrm{Ob}(\mathcal{C})$  满足  $Q(X) \simeq 0$  当且仅当存在  $Y \in \mathrm{Ob}(\mathcal{C})$  使得  $X \oplus Y \in \mathrm{Ob}(\mathcal{N})$  。这再次相当于 $X \oplus X[1] \in \mathrm{Ob}(\mathcal{N})$ 。证明很简单。

\subsection{导出范畴}
在本节中， $ \mathcal{C} $  将表示阿贝尔范畴。

我们将把前面的构造应用于三角范畴 $ \mathbf{K}(\mathcal{C}) $ 。很明显：

\[
\mathcal{N}=\left\{X\in\mathbf{O b}(\mathbb{K}(\mathcal{C}));H^{n}(X)=0\mathrm{f o r~a n y}n\right\}
\]

是一个零系统。请注意，根据命题 1.5.6， $ S(\mathcal{N}) $  由  $ \mathbf{K}(\mathcal{C}) $  的准同构组成。

\statementlabel{定义 1.7.1}我们设置  $ \mathbf{D}(\mathcal{C}) = \mathbf{K}(\mathcal{C}) / \mathcal{N} $  并将  $ \mathbf{D}(\mathcal{C}) $  称为  $ \mathcal{C} $  的导出范畴。

通过将 $ \mathbf{K}(\mathcal{C}) $ 替换为 $ \mathbf{K}^{b}(\mathcal{C}) $ （分别为 $ \mathbf{K}^{+}(\mathcal{C}) $ ，分别为 $ \mathbf{K}^{-}(\mathcal{C}) $ ），我们类似地定义了导出范畴 $ \mathbf{D}^{b}(\mathcal{C}) $ （分别为 $ \mathbf{D}^{+}(\mathcal{C}) $ ，分别为 $ \mathbf{D}^{-}(\mathcal{C}) $ ）。根据命题 1.6.3，函子  $ H^{n}(\cdot): \mathbf{K}(\mathcal{C}) \to \mathcal{C} $  因式分解  $ \mathbf{D}(\mathcal{C}) $  。我们仍然用 $ H^{n}(\cdot) $ 来表示这样获得的从 $ \mathbf{D}(\mathcal{C}) $ 到 $ \mathcal{C} $ 的函子。

\statementlabel{命题 1.7.2}(i)  $ \mathbf{D}^{b}(\mathcal{C}) $ （分别为 $ \mathbf{D}^{+}(\mathcal{C}) $ 、 $ \mathbf{D}^{-}(\mathcal{C}) $ ）等价于由对象 $ X $ 组成的 $ \mathbf{D}(\mathcal{C}) $ 的满子范畴，使得 $ H^{n}(X) = 0 $ 为 $ |n| \gg 0 $ （分别为 $ n \ll 0 $ 、分别为 $ n \gg 0 $ ）。

(ii) 通过函子  $\mathcal{C} \to \mathbf{K}(\mathcal{C}) \to \mathbf{D}(\mathcal{C})$  的组合， $\mathcal{C}$  等价于由对象  $X$  组成的  $\mathbf{D}(\mathcal{C})$  的满子范畴，使得  $H^n(X) = 0$  代表  $n \neq 0$  。

\statementlabel{证明}该命题紧接着命题 1.6.10 和以下事实：对于  $X \in \mathrm{Ob}(\mathbb{K}(\mathcal{C}))$  ， $H^j(X) = 0$  对于  $j < n$  （分别  $H^j(X) = 0$  对于  $j > n$  ）， $X \to \tau^{\geq n}(X)$  （分别  $\tau^{\leq n}(X) \to X$  ），是准同构。   $\square$ 

\statementlabel{备注 1.7.3}令  $ X \in \text{Ob}(\mathbb{K}(\mathcal{C})) $  、  $ Q(X) $  其图像在  $ \mathbb{D}(\mathcal{C}) $  中。那么  $ Q(X) = 0 $  当且仅当  $ X $  在  $ \mathbb{K}(\mathcal{C}) $  中与零准同构：这直接从备注 1.6.11 和函子  $ H^n(\cdot) $  的可加性得出。现在让  $ f: X \to Y $  成为  $ \mathbb{C}(\mathcal{C}) $  中的态射。根据定义， $ \mathbb{D}(\mathcal{C}) $  中的  $ f $  为 0，当且仅当存在准同构  $ g: X' \to X $  ，使得  $ f \circ g $  与 0 同伦，或者，当且仅当存在准同构  $ h: Y \to Y' $  使得  $ h \circ f $  与 0 同伦。注意，一般来说，不存在准同构 $ g: X' \to X $  （分别为  $ h: Y \to Y' $  ），使得  $ f \circ g = 0 $  （分别为  $ h \circ f = 0 $  ）位于  $ \mathbb{C}(\mathcal{C}) $  中。

示例 1.7.4。令  $ \mathcal{C} = \mathfrak{M}\mathfrak{o}\mathfrak{d}(\mathbb{Z}) $  和  $ f: X \to Y $  由下式给出：

\textbackslash{}[\textbackslash{}begin\{array\}\{ccc\}X\&\textbackslash{}cdots\textbackslash{}longrightarrow\&0\textbackslash{}longrightarrow\&0\textbackslash{}longrightarrow\textbackslash{}mathbb\{Z\}\textbackslash{}longrightarrow\textbackslash{}cdots\textbackslash{}longrightarrow0\textbackslash{}longrightarrow\textbackslash{}cdots\textbackslash{}\textbackslash{} \textbackslash{}left\textbackslash{}downarrow f:\&\textbackslash{}quad\textbackslash{}quad\textbackslash{}quad\&\textbackslash{}quad\textbackslash{}quad\textbackslash{}quad\textbackslash{}

那么  $f$  与  $0$  同伦，但不存在  $g:X'\to X$  使得  $g$  是准同构和  $f\circ g=0$  。 （事实上，这样的 $g$ 必须是 $0$ ，因为 $f$ “是单射的，这是荒谬的。）

\statementlabel{命题 1.7.5}令  $\mathcal{C}$  为阿贝尔范畴，让  $0 \to X \xrightarrow{f} Y \xrightarrow{g} Z \to 0$  为  $\mathbf{C}(\mathcal{C})$  中的精确序列。令  $M(f)$  为  $f$  的映射锥，并令  $\phi^n: M(f)^n = X^{n+1} \oplus Y^n \to Z^n$  为态射  $(0, g^n)$  。那么  $\{\phi^n\}_{n}: M(f) \to Z$  是复形的态射，  $\phi \circ \alpha(f) = g$  ，而  $\phi$  是拟同构。

\statementlabel{证明}很容易看出 $\phi$  是复形的态射。此外，我们有一个精确的序列：

\[
0\to M(\operatorname{id}_{X})\xrightarrow{\gamma}M(f)\to Z\to0
\]

其中  $ \gamma $  与态射  $ id_{x} \to f $  相关联。最后一个态射由交换图描述：

\[
\begin{array}{l}X\xrightarrow{\mathrm{i d}}X\\\left|\mathrm{i d}\right.\\\downarrow X\xrightarrow{\quad f\quad}Y\quad.\end{array}
\]

根据命题 1.3.6，检查所有  $ n \in \mathbb{Z} $  的  $ H^n(M(\mathrm{id}_X)) = 0 $  就足够了。由于  $ M(\mathrm{id}_X) $  在  $ \mathbf{K}(\mathcal{C}) $  中为零，所以这是显而易见的。   $ \square $ 

在命题 1.7.5 的情况下，区分三角形  $ X \to Y \to Z^{\frac{n}{2}}  X[1] $  被称为与精确序列  $ 0 \to X \to Y \to Z \to 0 $  相关联的区分三角形。这里 $ h = \beta(f) \circ \phi^{-1} $ 。

请注意，上面区分的三角形产生了一个长的精确序列：

\[
H^{n}(X)\longrightarrow H^{n}(Y)\longrightarrow H^{n}(Z)\xrightarrow[H^{n}(h)]{H^{n+1}}(X)\longrightarrow H^{n+1}(X)
\]

和  $ H^{n}(h) = -\delta $  、 $ \delta $  在命题 1.3.6 中定义。

另请注意，如果  $X, Y, Z$  集中在零度（即：是  $\mathcal{C}$  的对象），则态射  $h: Z \to X[1]$  在  $\mathbf{D}^{+}(\mathcal{C})$  中为零，当且仅当精确序列分裂时。然而，对于所有 $n \in \mathbb{Z}$  来说，总是有 $H^n(h) = 0$  。

\statementlabel{备注 1.7.6}我们在§3中定义了 $\mathbf{C}(\mathcal{C})$ 上的函子 $\tau^{\geq n}(\cdot)$ 和 $\tau^{\leq n}(\cdot)$ 。很简单，它们将零态射同伦变换为零态射同伦，并且根据命题 1.3.7，它们将拟同构变换为拟同构。这样就得到了函子  $\tau^{\geq n}:\mathbf{D}(\mathcal{C}) \to \mathbf{D}^+(C)$  和  $\tau^{\leq n}:\mathbf{D}(C) \to \mathbf{D}^-(C)$  。

应用命题 1.7.5，我们得到  $ \mathbf{D}(\mathcal{C}) $  中的区别三角形：

\[
\tau^{\leq n}(X)\longrightarrow X\longrightarrow\tau^{\geq n+1}(X)\xrightarrow[+1]{}
\]

\[
\tau^{\leqslant n-1}(X)\longrightarrow\tau^{\leqslant n}(X)\longrightarrow H^{n}(X)[-n]\xrightarrow[+1]{}
\]

\[
H^{n}(X)[-n]\longrightarrow\tau^{\geq n}(X)\longrightarrow\tau^{\geq n+1}(X)\xrightarrow[+1]{}.
\]

事实上，对于  $X \in \mathrm{Ob}(\mathbf{C}(\mathcal{C}))$  ，  $\tau^{\geq n+1}(X)$  与  $\mathrm{Coker}(\tau^{\leq n}(X) \to X)$  准同构，  $\tau^{\leq n-1}(X)$  与  $\mathrm{Ker}(\tau^{\leq n}(X) \to H^n(X)[-n])$  准同构，  $\tau^{\geq n+1}(X)$  与  $\mathrm{Coker}(H^n(X)[-n] \to \tau^{\geq n}(X))$  准同构。

\statementlabel{命题 1.7.7}令 I 为 C 的完全可加子范畴，使得：

(1.7.5) 对于任何  $ X \in \mathrm{Ob}(\mathcal{C}) $  ，都存在  $ X' \in \mathrm{Ob}(\mathcal{I}) $  和精确序列  $ 0 \to X \to X' $  。

然后：

(i) 对于任何  $ X \in \text{Ob}(\mathbb{K}^{+}(\mathcal{C})) $  ，存在  $ X' \in \text{Ob}(\mathbb{K}^{+}(\mathcal{I})) $  和准同构  $ f: X \to X' $  ，

(ii) 令  $\mathcal{N}$  由 (1.7.1) 给出，并令  $\mathcal{N}' = \mathcal{N} \cap \mathrm{Ob}(\mathbb{K}^{+}(\mathcal{I}))$  。然后是规范函子：

\[
\mathsf{K}^{+}(\mathcal{I})/\mathcal{N}^{\prime}\to\mathsf{D}^{+}(\mathcal{C})
\]

是范畴的等价。

\statementlabel{证明}根据命题 1.6.5，(ii) 由 (i) 得出。现在让 $ X \in \mathrm{Ob}(\mathbb{K}^{+}(\mathcal{C})) $  。我们将通过归纳构造一个复数 $ X'_{i} \to X'_{p-1} \to X'_{p} \to 0 \to \cdots $ 和复数 $ X \to X'_{i} \to p $ 的态射，使得 $ X^{j,p} $ 对于所有j属于 $ \mathcal{F} $ ，对于j<p $ H^{j}(X) \simeq H^{j}(X'_{i} \to p) $ 并且 $ H^{p}(X) \to \mathrm{Coker}\, d_{X'_{i}p}^{p-1} $ 是单态。

这对于  $p \ll 0$  是可能的。假设我们已经构建了  $X'_{p}$  。设置 $Z^{p+1} = \mathrm{Coker}\,d_{X'_{p}}^{p-1} \oplus_{\mathrm{Coker}\,d_{X}^{p-1}} X^{p+1}$ （参见练习I.6）并选择 $X^{p+1} \in \mathrm{Ob}(\mathcal{P})$ ，使得存在单态 $Z^{p+1} \to X^{p+1}$ 。

让我们将练习 I.6 的结果应用到图表中

\[
\begin{array}{c}\text{Coker}\mathrm{d}_{X}^{p-1}\longrightarrow X^{p+1}\\\downarrow\quad\downarrow\quad\downarrow\\\text{Coker}\mathrm{d}_{X^{\prime}\leq p}^{p-1}\longrightarrow X^{\prime p+1}\end{array}
\]

（态射 $X^{p} \to X^{p+1}$  和 $X^{p+1} \to X^{p+1}$  以明显的方式定义）。那么 $H^{p}(X) \simeq H^{p}(X_{\leq p+1}^{\prime})$  就可以得出 $H^{p}(X) \to \mathrm{Coker}\, \mathrm{d}_{X_{\leq p}}^{p-1}$  是单态的事实，而且 $H^{p+1}(X) \to \mathrm{Coker}\, \mathrm{d}_{X_{\leq p+1}^{\prime}}$  是单态的事实。

\statementlabel{推论 1.7.8}在命题 1.7.7 的情况下，假设 (1.7.5) 并且：

(1.7.6) 存在一个整数  $d \geqslant 0$  ，使得对于  $\mathcal{C}$  、  $X^0 \to X^1 \to \cdots \to X^d \to 0$  中的任何精确序列以及  $j < d$  的  $X^j \in \mathrm{Ob}(\mathcal{I})$  ，我们有  $X^d \in \mathrm{Ob}(\mathcal{I})$  。

那么对于任何  $ X \in \mathrm{Ob}(\mathbb{K}^{b}(\mathcal{C})) $  ，都存在  $ X' \in \mathrm{Ob}(\mathbb{K}^{b}(\mathcal{I})) $  和准同构  $ X \to X' $  。

\statementlabel{证明}通过前面的命题，我们可以找到 $ X' \in \mathrm{Ob}(\mathbb{K}^+(\mathcal{I})) $  和一个拟同构 $ X \to X' $  。

如果我们假设  $H^j(X)=0$  为  $j>n_\circ$  ，那么我们得到  $H^j(X')=0$  为  $j>n_\circ$  。因此 $\tau^{\leq n_\circ+d}(X')\to X'$  是准同构。现在  $(\tau^{\leq n_\circ+d}(X'))^k$  属于  $k<n_\circ+d$  的  $\mathrm{Ob}(\mathcal{I})$  ，并且  $k>n_\circ+d$  属于  $0$  。由于  $H^k(\tau^{\leq n_\circ+d}(X'))=0$  对于  $k>n_\circ$  ，条件 (1.7.6) 意味着  $(\tau^{\leq n_\circ+d}(X'))^{n_\circ+d}\in\mathrm{Ob}(\mathcal{I})$  。最后  $X$  有界，准同构  $X\to X'$  定义了准同构  $X\to\tau^{\leq k}X$  ，对于  $k\gg0$  。   $\square$ 

最后一个命题在以下情况下尤其重要。

\statementlabel{定义 1.7.9}有人说 $\mathcal{C}$ 有足够的单射，如果对于任何 $X \in \mathrm{Ob}(\mathcal{C})$ ，在 $\mathcal{C}$ 中存在单射对象 $X'$ 和单态 $X \to X'$ 。

换句话说，如果单射对象的子范畴满足（1.7.5），则 \& 具有足够的单射。

\statementlabel{命题 1.7.10}假设 $\mathcal{C}$ 有足够的单射，并让 $\mathcal{I}$ 表示单射对象的 $\mathcal{C}$ 的满子范畴。那么从 $\mathbf{K}^{+}(\mathcal{I})$ 到 $\mathbf{D}^{+}(\mathcal{C})$ 的自然函子就是范畴的等价。

\statementlabel{证明}由命题1.7.7足以证明：

\[
\mathcal{N}\cap\mathrm{O b}(\mathbb{K}^{+}(\mathcal{I}))=0.
\]

也就是说，任何  $ X \in \mathrm{Ob}(\mathbb{C}^{+}(\mathcal{I})) $  使得  $ H^n(X) = 0 $  对于任何  $ n $  都是同伦于零。设置 $ Z^n = \mathrm{Ker}\, d_X^n $ 。那么我们就有了精确的序列：

\[
0\longrightarrow Z^{n}\xrightarrow{i^{n}}X^{n}\xrightarrow{j^{n}}Z^{n+1}\longrightarrow0.
\]

通过 $n$  的归纳，我们得到所有 $Z$  都是单射的。因此序列 (1.7.8) 分裂，并且存在态射：

\[
k^{n}\colon X^{n}\to\mathbb{Z}^{n}~,\qquad t^{n}\colon\mathbb{Z}^{n+1}\to X^{n}
\]

这样  $ k^n \circ i^n = \mathrm{id}_{Z^n} $  、  $ j^n \circ t^n = \mathrm{id}_{Z^{n+1}} $  、  $ k^n \circ t^n = 0 $  和  $ \mathrm{id}_{X^n} = i^n \circ k^n + t^n \circ j^n $  。然后  $ s^n = t^{n-1} \circ k^n: X^n \to X^{n-1} $  给出同伦，即：  $ \mathrm{id}_{X^n} = \mathrm{d}_X^{n-1} \circ s^n + s^{n+1} \circ \mathrm{d}_X^n $  。   $ \square $ 

现在让 $\mathcal{C}$  是一个阿贝尔范畴， $\mathcal{C}'$  是一个完整的阿贝尔子范畴。用  $\mathbf{D}_{\mathcal{C}}^{+}(\mathcal{C})$  表示  $\mathbf{D}^{+}(\mathcal{C})$  的完整三角子类，由其上同调对象属于  $\mathcal{C}'$  的复形组成。存在一个自然函子：

\[
\mathbf{D}^{+}(\mathcal{C}^{\prime})\xrightarrow[\delta]\longrightarrow\mathbf{D}_{\mathcal{C}^{\prime}}^{+}(\mathcal{C})~.
\]

我们将给出一个有用的标准，确保  $\delta$  是等价的。首先我们说  $\mathcal{C}'$  是  $\mathcal{C}$  的厚子类，如果对于任何精确序列  $\mathcal{C}$  中的  $Y \to Y' \to X \to Z \to Z'$  和  $\mathcal{C}'$  中的  $Y, Y', Z, Z'$  ，  $X$  属于  $\mathcal{C}'$  。

\statementlabel{命题 1.7.11}令 $\mathcal{C}$  为阿贝尔范畴， $\mathcal{C}'$  为厚完整阿贝尔子范畴。假设对于任何单态  $f: X' \to X$  和  $X' \in \mathrm{Ob}(\mathcal{C}')$  ，都存在态射  $g: X \to Y$  和  $Y \in \mathrm{Ob}(\mathcal{C}')$  ，使得  $g \circ f$  是单态。那么(1.7.9)中的函子 $\delta$ 就是范畴的等价。

\statementlabel{证明}根据命题1.6.5，足以证明：

\[
\begin{cases}{{f o r~}a n y\;X\in\operatorname{O b}(\mathbf{D}_{\mathcal{C}^{\prime}}^{+}(\mathcal{C})){~t h e r e~e x i s t s~}X^{\prime}\in\operatorname{O b}(\mathbf{K}^{+}(\mathcal{C}^{\prime})){}}\\ {{a n d~}a{q u a s i-i s o m o r p h i s m}X\xrightarrow{\sim}X^{\prime}{~.}}\\ \end{cases}
\]

 $X'$  的构造与命题 1.7.7 的证明类似。定义了复数  $X'_{≤p}$  :  $\cdots \to X^{p-1} \to X^{p} \to 0 \to \cdots$  和态射  $X \to X'_{≤p}$  ，使得  $X^{j}$  属于  $\mathcal{C}'$  ，  $H^j(X) \simeq H^j(X'_{≤p})$  属于  $j < p$  和  $H^p(X) \to$  Coker  $\mathrm{d}_{X^{<p}}^{p-1}$  是单态，我们构造  $X^{p+1}$  如下。

让 $ M = \text{Coker } \text{d}_{X' \subset P}^{p-1} \oplus \text{Coker } \text{d}_{X}^{p-1} \text{Ker } \text{d}_{X}^{p+1} \text{ and } N = \text{Coker } \text{d}_{X' \subset P}^{p-1} \oplus \text{Coker } \text{d}_{X}^{p-1} X^{p+1} $  。我们有一个精确的序列（参见练习 I.6 和 (1.3.9)）：

\[
0\to H^{p}(X)\to\mathbf{C o k e r d}_{X\subsetneq p}^{p-1}\to M\to H^{p+1}(X)\to0\enspace.
\]

因此  $M$  属于  $\mathcal{C}'$  。将假设应用于单态  $i: M \to N$  ，我们发现态射  $g: N \to X'^{p+1}$  与  $X'^{p+1} \in \mathrm{Ob}(\mathcal{C}')$  和  $g \circ i$  是单态。一种以自然方式定义态射 $X'^{p} \to X'^{p+1}$  和 $X^{p+1} \to X'^{p+1}$ ，并按照命题1.7.7 的证明检查复数 $X'_{\leq p+1}$  是否具有所需的属性。   $\square$ 

\statementlabel{备注 1.7.12}利用与命题 1.7.11 中相同的假设， $ \delta $  导出等价性：

\[
\begin{array}{r}{\mathbf{D}^{b}(\mathcal{C}^{\prime})\simeq\mathbf{D}_{\mathcal{C}^{\prime}}^{b}(\mathcal{C})\;.}\end{array}
\]

（使用函子  $ \tau^{\leq n} $  表示  $ n \gg 0 $  。）

评论 1.7.13。让我们总结一下  $\mathbf{D}(\mathcal{C})$  的构造。我们从阿贝尔范畴  $\mathcal{C}$  开始，并考虑  $\mathcal{C}$  的复形的范畴  $\mathbf{C}(\mathcal{C})$  。然后我们确定  $\mathbf{C}(\mathcal{C})$  同伦于零的态射是零态射：我们得到范畴  $\mathbf{K}(\mathcal{C})$  。  $\mathbf{K}(\mathcal{C})$  在  $\mathbf{C}(\mathcal{C})$  上的优点是，许多在  $\mathbf{C}(\mathcal{C})$  中不交换的图表，实际上在  $\mathbf{K}(\mathcal{C})$  中这样做，使  $\mathbf{K}(\mathcal{C})$  成为三角剖分

范畴。然后我们想要  $ \mathbf{K}(\mathcal{C}) $  中的态射，从而导致上同调的同构是可逆的。为此，我们“本地化” $ \mathbf{K}(\mathcal{C}) $ ，并得到 $ \mathbf{D}(\mathcal{C}) $ 。

符号 1.7.14。让 $A$  成为一个环。如果没有混淆的风险，我们将写  $\mathbf{D}(A)$  而不是  $\mathbf{D}(\mathfrak{Mod}(A))$  。

\subsection{导出函子}
在本节中， $ \mathcal{C} $  和  $ \mathcal{C}' $  将表示两个阿贝尔范畴， $ F: \mathcal{C} \to \mathcal{C}' $  表示一个加性函子。

我们将用  $ Q $  表示自然函子  $ \mathbf{K}^{+}(\mathcal{C}) \to \mathbf{D}^{+}(\mathcal{C}) $  或  $ \mathbf{K}^{+}(\mathcal{C}') \to \mathbf{D}^{+}(\mathcal{C}') $  。

\statementlabel{定义 1.8.1}令 $ T: \mathbb{D}^{+}(\mathcal{C}) \to \mathbb{D}^{+}(\mathcal{C}') $  为三角范畴的函子，并令 $ s $  为函子的态射：

\[
\begin{array}{r}{s\colon Q\circ\mathbb{K}^{+}(F)\to T\circ Q}\end{array},
\]

其中  $ \mathbf{K}^{+}(\mathbf{F}): \mathbf{K}^{+}(\mathcal{C}) \to \mathbf{K}^{+}(\mathcal{C}') $  是与  $ F $  自然关联的函子。假设对于三角范畴  $ G: \mathbf{D}^{+}(\mathcal{C}) \to \mathbf{D}^{+}(\mathcal{C}') $  的任何函子，态射：

\[
\mathrm{H o m}(T,G)\xrightarrow[\mathfrak{s}]{\quad\mathrm{H o m}(Q\circ\mathbb{K}^{+}(F),G\circ Q)\quad}
\]

是同构。

那么(T,s)，它在同构上是唯一的，被称为F的右导函子，记为RF。函子  $H^{n}\circ RF$  也表示为  $R^{n}F$  ，称为 F 的第  $n$  导出函子。

让我们给出一个有用的标准来确保 RF 的存在。从现在开始直到命题 1.8.7，我们假设 F 保持精确。

\statementlabel{定义 1.8.2}$\mathcal{C}$  的满加性子范畴  $\mathcal{I}$  被称为相对于  $F$  的单射（或简称  $F$  -单射），如果：

(i) 满足条件（1.7.5），

(ii) 如果 $ 0 \to X' \to X \to X'' \to 0 $ 

Ob(J)，然后 $ X'': $ 

(iii) 如果  $0 \to X' \to X \to X'' \to 0$  是  $\mathcal{C}$  中的精确序列，并且如果  $X', X, X''$  位于  $\mathrm{Ob}(\mathcal{I})$  中，则序列  $0 \to F(X') \to F(X) \to F(X'') \to 0$  是精确的。

请注意，在条件 (i) 和 (ii) 下，条件 (iii) 等效于仅假设  $ X' \in \mathrm{Ob}(\mathcal{I}) $  的类似条件，因为假设 F 保持精确。

设 $\mathcal{I}$  为 $F$  -内射。然后我们可以很容易地检查 $F$ 将准同构为零的 $\mathbf{K}^{+}(\mathcal{I})$ 对象转换为满足相同属性的 $\mathbf{K}^{+}(\mathcal{C}^{\prime})$ 对象。因此函子的组成

\[
{\bf K}^{+}({\mathcal{I}})\xrightarrow{{\bf K}^{+}({\cal F})}{\bf K}^{+}({\mathcal{C}}^{\prime})\xrightarrow{\quad\quad}\mathbf{D}^{+}({\mathcal{C}}^{\prime})
\]

通过  $ \mathbf{K}^{+}(\mathcal{I})/\mathcal{N} \cap \operatorname{Ob}(\mathbf{K}^{+}(\mathcal{I})) $  进行因子分解，其中  $ \mathcal{N} $  由 (1.7.1) 给出。由于 $ \mathbf{K}^{+}(\mathcal{I})/\mathcal{N} \cap \operatorname{Ob}(\mathbf{K}^{+}(\mathcal{I})) $  根据命题1.7.7 等价于 $ \mathbf{D}^{+}(\mathcal{C}) $ ，我们得到：

\statementlabel{命题 1.8.3}假设存在  $\mathcal{C}$  的 F 内射子范畴  $\mathcal{I}$  。那么上面构造的从  $\mathbf{K}^{+}(\mathcal{I})/\mathcal{N} \cap \mathrm{Ob}(\mathbf{K}^{+}(\mathcal{I}))$  到  $\mathbf{D}^{+}(\mathcal{C}')$  的函子就是  $F$  的右派生函子。

\statementlabel{备注 1.8.4}根据 RF 的通用属性，前面的构造不依赖于  $\mathcal{I}$  。

\statementlabel{备注 1.8.5}令 $\mathcal{I}$  为 $\mathcal{C}$  的单射对象的满子范畴，并假设 $\mathcal{C}$  有足够的单射，（即：满足（1.7.5））。那么  $\mathcal{I}$  是  $F$  相对于任何左精确函子  $F$  的内射，因为  $\mathcal{I}$  中的任何序列都会分裂（参见练习 I.5）。特别是 $RF$  在这种情况下总是存在。

\statementlabel{备注 1.8.6}假设存在  $\mathcal{C}$  的  $F$  内射子范畴  $\mathcal{I}$  ，令  $n$  为整数，并令  $X \in \mathrm{Ob}(\mathbf{K}^{+}(\mathcal{C}))$  使得  $H^{k}(X) = 0$  对于  $k < n$  。然后  $R^{k}F(X) = 0$  为  $k < n$  和  $R^{n}F(X) = F(H^{n}(X))$  。事实上，对于这样的  $X$  ，我们可以找到  $X^{\prime} \in \mathrm{Ob}(\mathbf{K}^{+}(\mathcal{I}))$  和准同构  $X \to X^{\prime}$  ，使得  $X^{\prime k} = 0$  对应  $k < n$  。当然，如果  $X \in \mathrm{Ob}(\mathbf{K}^{+}(\mathcal{I}))$  ，则  $R^{k}F(X) = F(H^{k}(X))$  对于所有  $k$  。特别是如果  $X \in \mathrm{Ob}(\mathcal{I})$  、  $R^{k}F(X) = 0$  对于  $k \neq 0$  。 （ $\mathcal{C}$  的对象  $X$  使得  $k \neq 0$  的  $R^{k}F(X) = 0$  被称为  $F$  -非循环，参见练习 I.19。）

\statementlabel{命题 1.8.7}令 $\mathcal{C}, \mathcal{C}^{\prime}, \mathcal{C}^{\prime\prime}$  为三个阿贝尔范畴，并令 $F : \mathcal{C} \to \mathcal{C}^{\prime}$  、 $F^{\prime} : \mathcal{C}^{\prime} \to \mathcal{C}^{\prime\prime}$  为两个左精确函子。假设存在  $\mathcal{C}$  的完整加性子范畴  $\mathcal{I}$  （分别为  $\mathcal{C}^{\prime}$  的  $\mathcal{I}^{\prime}$  ），即  $F$  -内射（分别为  $F^{\prime}$  -内射），并且  $F(\mathrm{Ob}(\mathcal{I})) \subset \mathrm{Ob}(\mathcal{I}^{\prime})$  。那么  $\mathcal{I}$  是  $(F^{\prime} \circ F)$  -内射，我们有：

\[
R(F^{\prime}\circ F)=R F^{\prime}\circ R F\enspace.
\]

\statementlabel{证明}很简单。

请注意，如果我们只假设  $RF, RF', R(F' \circ F)$  存在，我们就会发现函子的规范态射：

\[
R(F^{\prime}\circ F)\to R F^{\prime}\circ R F\enspace.
\]

事实上我们有：

\[
\mathrm{H o m}(R(F^{\prime}\circ F),R F^{\prime}\circ R F)=\mathrm{H o m}(Q\circ\mathbb{K}^{+}(F^{\prime})\circ\mathbb{K}^{+}(F),R F^{\prime}\circ R F\circ Q)
\]

和态射：

\[
Q\circ\mathbb{K}^{+}(F)\stackrel{\alpha}{\longrightarrow}R F\circ Q
\]

\[
Q\circ\mathbb{K}^{+}(F^{\prime})\xrightarrow{\beta}R F^{\prime}\circ Q
\]

给出态射：

\[
Q\circ\mathbb{K}^{+}(F^{\prime})\circ\mathbb{K}^{+}(F)\xrightarrow{\beta\circ\mathbb{K}^{+}(F)}R F^{\prime}\circ Q\circ\mathbb{K}^{+}(F)\xrightarrow{R F^{\prime}\circ\alpha}R F^{\prime}\circ R F\circ Q\qquad.
\]

\statementlabel{命题 1.8.8}令 $\mathcal{C}$ 和 $\mathcal{C}'$ 为两个阿贝尔范畴，令 $F', F, F''$ 为从 $\mathcal{C}$ 到 $\mathcal{C}'$ 的三个左精确函子，令 $\lambda: F' \to F$ 和 $\mu: F \to F''$ 为函子的态射。假设存在  $\mathcal{C}$  的完整加性子范畴  $\mathcal{I}$  ，它相对于  $F', F$  和  $F'$  是单射的。进一步假设：

(1.8.3) 对于任何  $ X \in \mathrm{Ob}(\mathcal{I}) $  ，序列

\[
0\to F^{\prime}(X)\to F(X)\to F^{\prime \prime}(X)\to0{~i s~e x a c t~}.
\]

那么自然存在函子  $ \nu: RF'' \to RF' [1] $  的态射，这样对于任何  $ X \in \text{Ob}(\mathbf{D}^+(\mathcal{C})) $  ，序列：

\[
RF^{\prime}(X)\xrightarrow{R\lambda(X)}RF(X)\xrightarrow{R\mu(X)}RF^{\prime \prime}(X)\xrightarrow{\nu(X)}RF^{\prime}(X)[1]
\]

是  $ \mathbf{D}^{+}(\mathcal{C}^{\prime}) $  中的一个显着三角形。

\statementlabel{证明}对于任何  $ X \in \mathrm{Ob}(\mathbb{K}^{+}(\mathcal{I})) $  ，我们都有一个精确的序列：

\[
0\xrightarrow{\quad\quad\quad\quad}F^{\prime}(X^{n})\xrightarrow{\quad\lambda(X^{n})\quad}F(X^{n})\xrightarrow{\quad\mu(X^{n})\quad}F^{\prime\prime}(X^{n})\xrightarrow{\quad\quad\quad\quad}0
\]

因此，根据备注 1.7.5， $ F''(X) $  与  $ \mathbf{D}^+(\mathcal{C}') $  中  $ \lambda(X) $  的映射锥  $ M(\lambda(X)) $  同构。

态射  $\alpha(\lambda(X)) \colon M(\lambda(X)) \to F'(X)[1]$  给出了  $\mathbf{D}^{+}(\mathcal{C}')$  中从  $F''(X)$  到  $F'(X)[1]$  的态射。通过商，我们得到：

\[
\nu:R F^{\prime \prime}\to R F^{\prime}[1]\enspace.
\]

声明的其余部分很简单。 ⑨

为了结束本节，让我们考虑右精确函子  $F$  的情况。然后，通过反转箭头，定义  $\mathcal{C}$  的  $F$  投影子范畴和  $F$  的左派生函子的概念，表示为  $LF$  。

更准确地说，如果满足以下条件， $ \mathcal{C} $  的满加性子范畴  $ \mathcal{P} $  称为 F-射影（F 右正合）：

(i) 对于任何  $ X \in \mathrm{Ob}(\mathcal{C}) $  ，都存在  $ X' \in \mathrm{Ob}(\mathcal{P}) $  和精确序列  $ X' \to X \to 0 $  ，

(ii) 如果  $ 0 \to X' \to X \to X'' \to 0 $  是  $ \mathcal{C} $  中的精确序列，并且如果  $ X'' $  和  $ X $  位于  $ \mathrm{Ob}(\mathcal{P}) $  中，则  $ X' $  也在  $ \mathrm{Ob}(\mathcal{P}) $  中，

(iii) 如果  $ 0 \to X' \to X \to X'' \to 0 $  是  $ \mathcal{C} $  中的精确序列，并且如果  $ X', X, X'' $  位于  $ \mathrm{Ob}(\mathcal{P}) $  中，则序列  $ 0 \to F(X') \to F(X) \to F(X'') \to 0 $  是精确的。

然后构造左派生函子：

\[
L F:\mathbf{D}^{-}(\mathcal{C})\to\mathbf{D}^{-}(\mathcal{C}^{\prime})
\]

与右派生函子类似。

示例 1.8.9。让 $A$  成为一个环。那么范畴 $\mathfrak{Mod}(A)$ 就有足够的单射和足够的射影（参见Cartan-Eilenberg [1]）。此外，让 $M$  是一个正确的 $A$  模。那么（左）平坦 $A$-模 的范畴相对于函子  $M\otimes_A^{-1}$  是投影的。

我们将在接下来的章节中遇到许多派生函子。

\statementlabel{备注 1.8.10}派生函子的构造可以扩展到仅在  $ \mathbf{K}^{+}(\mathcal{C}) $  上定义的函子。更准确地说，让  $ F $  成为从  $ \mathbf{K}^{+}(\mathcal{C}) $  到  $ \mathbf{K}^{+}(\mathcal{C}') $  的三角范畴的函子。假设给定 $ \mathbf{K}^{+}(\mathcal{C}) $ 的完整三角子范畴 $ \mathcal{I} $ ，使得满足下面的(1.8.4)和(1.8.5)。

(1.8.4) 对于任何  $ X \in \text{Ob}(\mathbb{K}^+(\mathcal{C})) $  ，存在  $ X \to X' $  与  $ X' \in \text{Ob}(\mathcal{I}) $  的准同构。

(1.8.5) 如果 $ X \in \mathrm{Ob}(\mathcal{I}) $  与 0 拟同构，则  $ F(X) $  与 0 拟同构。

那么我们可以像我们在命题1.8.3中所做的那样定义 $ RF: \mathbf{D}^+(\mathcal{C}) \to \mathbf{D}^+(\mathcal{C}') $ ，并且RF将满足定义1.8.1的通用性质。

\statementlabel{备注 1.8.11}令 $F: \mathcal{C} \to \mathcal{C}'$  为逆变函子。然后它定义了一个逆变函子 $\mathbf{K}(F): \mathbf{K}(\mathcal{C}) \to \mathbf{K}(\mathcal{C}')$ ，如下所示。如果 $X = (X^n)_{n \in \mathbb{Z}}$  属于 $\mathbf{C}(\mathcal{C})$  则：

\[
\left\{\begin{aligned}{}&{{}(K(F)(X))^{n}=F(X^{-n})}\\ {}&{{}\mathsf{d}_{\mathsf{K}(F)(X)}^{n}=(-1)^{n+1}F(\mathsf{d}_{X}^{-n-1}).}\\ \end{aligned}\right.
\]

将  $F$  视为规范逆变函子  $\mathcal{C}^\circ \to \mathcal{C}$  ，我们有  $\mathbf{K}(\mathcal{C}^\circ) \simeq (\mathbf{K}(\mathcal{C}))^\circ$  ，如果  $\mathcal{C}$  是阿贝尔算子，则  $\mathbf{D}(\mathcal{C}^\circ) \simeq (\mathbf{D}(\mathcal{C}))^\circ$  。我们还有 $\mathbf{K}^{\pm}(\mathcal{C}^\circ) \simeq (\mathbf{K}^{\mp}(\mathcal{C}))^\circ$  和 $\mathbf{D}^{\pm}(\mathcal{C}^\circ) \simeq (\mathbf{D}^{\mp}(\mathcal{C}))^\circ$  。

\subsection{双复形}
设 C 为加法范畴。

\statementlabel{定义 1.9.1}$ \mathcal{C} $  中的双复数  $ (X, \mathbf{d}_X) $  由数据  $ \{X^{n,m}, \mathbf{d}_X^{n,m}, \mathbf{d}_X^{n,m}\} $  组成，其中  $ X^{n,m} \in \mathrm{Ob}(\mathcal{C}) $  、  $ \mathbf{d}_X^{n,m}: X^{n,m} \to X^{n+1,m} $  、  $ \mathbf{d}_X^{n,m}: X^{n,m} \to X^{n,m+1} $  对于任何对  $ (n,m) $  以及：

\[
\mathbf{d}_{X}^{\prime2}=0~,\qquad\mathbf{d}_{X}^{\prime\prime2}=0~,\qquad\mathbf{d}_{X}^{\prime}\circ\mathbf{d}_{X}^{\prime\prime}=\mathbf{d}_{X}^{\prime\prime}\circ\mathbf{d}_{X}^{\prime}~.
\]

(1.9.1) 的含义如下：  $ \mathbf{d}_{X}^{\prime n+1,m} \circ \mathbf{d}_{x}^{\prime n,m} = 0 $  ，其他关系也类似。

我们有时将“简单复杂”称为定义 1.3.1 意义上的复杂。

令 $X$  和 $Y$  为两个双复形。从 $X$  到 $Y$  的态射 $f$  以显而易见的方式定义。然后我们得到 $\mathcal{C}$ 上双复形的范畴 $\mathbf{C}^{2}(\mathcal{C})$ 。

令 $X$  为双重复数。对于给定的  $n \in \mathbb{Z}$  ，让  $X_I^n$  表示简单复数：

\[
X_{I}^{n}=\left\{X^{n,m},\mathsf{d}_{X}^{\prime\prime n,m}\right\}_{m\in{\mathbb{Z}}}.
\]

态射族  $ \{d_{x}^{\prime n,m}\}_{m\in\mathbb{Z}} $  定义了一个态射：

\[
\mathbf{d}_{I}^{n}\colon X_{I}^{n}\to X_{I}^{n+1}
\]

显然 $ \mathrm{d}_{I}^{n+1} \circ \mathrm{d}_{I}^{n} = 0 $ 。因此我们构造了一个函子：

\[
F_{I}\colon\mathbf{C}^{2}(\mathcal{C})\to\mathbf{C}(\mathbf{C}(\mathcal{C}))\enspace,
\]

\[
X\mapsto\{X_{I}^{n},\mathsf{d}_{I}^{n}\}.
\]

这个函子显然是范畴的等价物。

通过反转第一个和第二个索引（即  $ d' $  和  $ d'' $  ），我们得到另一个等价：

\[
F_{I I}\colon\mathbf{C}^{2}(\mathcal{C})\to\mathbf{C}(\mathbf{C}(\mathcal{C}))\enspace,
\]

\[
X\mapsto\left\{X_{I I}^{m},\mathsf{d}_{I I}^{m}\right\}\;,
\]

其中  $ X_{II}^{m}=\left\{X^{n,m},d_{X}^{\prime n,m}\right\}_{n\in\mathbb{Z}} $  和  $ d_{II}^{m}=\left\{d_{X}^{\prime\prime n,m}\right\}_{n\in\mathbb{Z}} $  。

假设 X 满足以下有限性：

(1.9.2) 对于任何 $ k \in \mathbb{Z} $  ，集合 $ \{(n, m) \in \mathbb{Z} \times \mathbb{Z} : n + m = k, X^{n,m} \neq 0\} $  是有限的。

然后就可以将  $X$  关联到一个简单的复杂  $s(X)$  。一套：

\[
\mathbf{s}(X)^{k}=\bigoplus_{k=n+m}X^{n,m}~.
\]

让 $ i_{n,m}: X^{n,m} \to \bigoplus_{k=n^{\prime}+m^{\prime}} X^{n^{\prime},m^{\prime}} $  和 $ p_{n,m}: \bigoplus_{k=n^{\prime}+m^{\prime}} X^{n^{\prime},m^{\prime}} \to X^{n,m} $  成为自然的

态射分别从  $X^{n,m}$  到  $\mathrm{s}(X)^{k}$  和从  $\mathrm{s}(X)^{k}$  到  $X^{n,m}$  。其中一个定义：

\[
\mathbf{d}_{\mathbf{s}(X)}^{k}:\mathbf{s}(X)^{k}\to\mathbf{s}(X)^{k+1}
\]

by:

\[
P_{n^{\prime},m^{\prime}}\circ\mathsf{d}_{\mathtt{s}(X)}^{k}\circ i_{n,m}=\left\{\begin{aligned}{}&{{}\mathsf{d}_{X}^{\prime n,m}}&{}&{{}\mathrm{i f}m=m^{\prime}}\\ {}&{{}(-1)^{n}\mathsf{d}_{X}^{\prime\prime n,m}}&{}&{{}\mathrm{i f}n=n^{\prime}}\\ {}&{{}0}&{}&{{}\mathrm{o t h e r w i s e}}\\ \end{aligned}\right.
\]

对于 $ n + m = k $ ， $ n' + m' = k + 1 $ 。

\statementlabel{命题 1.9.2}数据  $ \{\mathbf{s}(X)^k, \mathbf{d}_{\mathbf{s}(X)}^k\}_{k\in\mathbb{Z}} $  ，在  $ \mathcal{C} $  中定义了一个复数，即  $ \mathbf{d}_{\mathbf{s}(X)}^{k+1} \circ \mathbf{d}_{\mathbf{s}(X)}^k = 0 $  。

\statementlabel{证明}很简单。

我们称 $ s(X) $  为与X 相关的简单复形。

人们可以认为 $ \mathbf{C}^2(\mathcal{C}) $ 对象的满加性子范畴 $ \mathbf{C}_f^2(\mathcal{C}) $ 满足(1.9.2)，并且 $ s(\cdot) $ 成为加法函子：

\[
{\mathsf{s}}(\cdot):{\mathbb{C}}_{f}^{2}({\mathcal{C}})\to{\mathbb{C}}({\mathcal{C}})~.
\]

现在我们假设 $ \mathcal{C} $  是一个阿贝尔范畴。

我们使用  $ F_I $  或  $ F_{II} $  定义函子  $ \tau_I^{\leq n} $  、  $ \tau_I^{\geq n} $  、  $ \tau_{II}^{\leq m} $  、  $ \tau_{II}^{\geq m} $  。例如：

\[
\tau_{I}^{\leqslant n}=(F_{I})^{-1}\circ\tau^{\leqslant n}\circ F_{I}\enspace.
\]

换句话说， $ \tau_I^{\leq n}(X) $  是双重复数：

\[
\cdots\longrightarrow X^{n-1,\cdot}\xrightarrow{\mathrm{d}_{X}^{n-1,\cdot}}X^{n,\cdot}\longrightarrow\operatorname{Ker}\mathrm{d}_{X}^{n,\cdot}\longrightarrow0\longrightarrow\cdots
\]

或等效于  $ \mathbf{C}(\mathcal{C}) $  中的复合体：

\[
\cdots{\longrightarrow}X_{I}^{n-1}{\xrightarrow{\quad}}\underbrace{X_{I}^{n-1}}_{\mathrm{d}_{I}^{n-1}}\longrightarrow X_{I}^{n}{\longrightarrow}\mathbf{K e r}\mathrm{d}_{I}^{n}{\longrightarrow}0{\longrightarrow}\cdots.
\]

我们还介绍了简单的复杂的：

\[
H_{I}^{p}(X)=H^{p}(F_{I}(X))~,
\]

和双重复合体：

\[
H_{I}(X)\colon\cdots\longrightarrow H_{I}^{p}(X)\xrightarrow[0]{}\scriptstyle{H_{I}^{p+1}(X)}\scriptstyle{\longrightarrow}\cdots
\]

对于  $ H_{I I}^{\eta}(X) $  和  $ H_{I I}(X) $  也类似。

\statementlabel{定理 1.9.3}令 $f: X \to Y$  为双复形的态射，其中 $X$  和 $Y$  均满足(1.9.2)。假设  $f$  引发同构：

\[
f:H_{I}H_{I I}(X)\simeq H_{I}H_{I I}(Y)\enspace.
\]

那么 $ \mathbf{s}(f) \colon \mathbf{s}(X) \to \mathbf{s}(Y) $  是拟同构。

\statementlabel{证明}让 $ q \in \mathbb{Z} $  。我们有一个区分三角形的交换图：

\[
\begin{array}{c}\mathrm{s}\tau_{II}^{\leq q^{-1}}(X)\longrightarrow\mathrm{s}\tau_{II}^{\leq q}(X)\longrightarrow H_{II}^{q}(X)[-q]\xrightarrow[+1]{}\\\left|\mathrm{s}\tau_{II}^{\leq q^{-1}}(f)\quad\right|\left|\mathrm{s}\tau_{II}^{\leq q}(f)\quad\right|\quad\left|H_{II}^{q}(f)[-q]\right.\\\left|\mathrm{s}\tau_{II}^{\leq q^{-1}}(Y)\quad\right|\longrightarrow\mathrm{s}\tau_{II}^{\leq q}(Y)\longrightarrow H_{II}^{q}(Y)[-q]\xrightarrow[+1]{}\end{array}
\]

（参见练习 I.25）。

假设  $ X_{II}^{q} = 0 $  为  $ q \ll 0 $  。

然后，由于  $q \ll 0$  的  $\mathfrak{st}_{II}^{\leq q}(f)$  与  $0$  准同构，并且  $H_{II}^{q}(f)$  是所有  $q$  的准同构，根据假设，我们得到  $\mathfrak{st}_{II}^{\leq q}(f)$  是所有  $q$  的准同构。由于  $q \gg 0$  的  $H^{k}(\mathfrak{st}_{II}^{\leq q}(f)) = H^{k}(\mathfrak{s}(f))$  和  $k$  被修复（再次参见练习 I.25），我们在这种情况下得出结论。为了处理一般情况，我们将前面的结果应用于 $\tau_{II}^{\geq q}(f)$ 。这给出了结果，因为对于固定的  $k$  ，  $H^{k}(\mathfrak{s}(f)) = H^{k}(\mathfrak{st}_{II}^{\geq q}(f))$  对于  $q \ll 0$  。   $\square$ 

最后，我们注意到 s 将与 0 的态射同伦转换为与 0 的态射同伦。更准确地说，如果存在  $ t_1^{n,m}: X^{n,m} \to Y^{n-1,m} $  和  $ t_2^{n,m}: X^{n,m} \to Y^{n,m-1} $  态射，则  $ \mathbf{C}_f^2(\mathcal{C}) $  中的态射  $ f: X \to Y $  被称为与 0 同伦：

\[
\mathrm{d}^{n n-1,m}\circ t_{1}^{n,m}=t_{1}^{n+1,m}\circ\mathrm{d}^{n n,m},
\]

\[
\mathrm{d}^{\prime n,m-1}\circ t_{2}^{n,m}=t_{2}^{n,m+1}\circ\mathrm{d}^{\prime n,m},
\]

以及：

\[
f^{n,m}=\mathrm{d}^{\prime n-1,m}\circ t_{1}^{n,m}+t_{1}^{n+1,m}\circ\mathrm{d}^{\prime n,m}+\mathrm{d}^{\prime\prime n,m-1}\circ t_{2}^{n,m}+t_{2}^{n,m+1}\circ\mathrm{d}^{\prime\prime n,m}~.
\]

那么 $ s(f):s(X)\to s(Y) $  与零同伦。特别是，如果  $ F_{I}(f):F_{I}(X)\to F_{I}(Y) $  与零同伦，则  $ s(f) $  与零同伦， $ F_{II}(f) $  也类似。

\subsection{双函子}
令  $C$  、  $C'$  和  $C''$  为三个范畴。

\statementlabel{定义 1.10.1}从 $ \mathcal{C} \times \mathcal{C}' $  到 $ \mathcal{C}'' $  的双函子 $ F $  由以下数据组成：

(i) 映射 $F: \mathrm{Ob}(\mathcal{C}) \times \mathrm{Ob}(\mathcal{C}') \to \mathrm{Ob}(\mathcal{C}')$ ，

(ii) 对于任何对  $(X, Y) \in \mathrm{Ob}(\mathcal{C})$  和任何对  $(X', Y') \in \mathrm{Ob}(\mathcal{C}')$  ，映射  $F: \mathrm{Hom}_{\mathcal{C}}(X, Y) \times \mathrm{Hom}_{\mathcal{C}'}(X', Y') \to \mathrm{Hom}_{\mathcal{C}'}(F(X, X'), F(Y, Y')))$  使得对于任何  $X \in \mathrm{Ob}(\mathcal{C})$  （分别  $X' \in \mathrm{Ob}(\mathcal{C}')$  ）， $F(X, \cdot)$  （分别  $F(\cdot, X')$  ）是从  $\mathcal{C}$  到  $\mathcal{C}'$  （分别从  $\mathcal{C}'$  到  $\mathcal{C}'$  的函子 $\mathcal{C}'$  )，如果  $f \in \mathrm{Hom}_{\mathcal{C}}(X, Y)$  、  $g \in \mathrm{Hom}_{\mathcal{C}'}(X', Y')$  则  $F(f, Y') \circ F(X, g) = F(Y, g) \circ F(f, X')$  。

双函子被称为加法、左精确、右精确、精确、三角范畴的双函子、上同调双函子等，如果对于每个变量都是这样的话。

人们以一种自然的方式定义双函子的态射。

示例 1.10.2。  $ _{\mathcal{C}}(\cdot,\cdot) $  是 $ \mathcal{C}^{\circ} \times \mathcal{C} $  设置的双功能。

示例 1.10.3。让 $ A $  成为一个环。那么  $ \otimes_A $  是从  $ \mathfrak{Mod}(A^{op}) \times \mathfrak{Mod}(A) $  到  $ \mathfrak{Mod}(\mathbb{Z}) $  的双函子。

现在我们假设  $\mathcal{C}, \mathcal{C}'$  和  $\mathcal{C}''$  是三个阿贝尔范畴， $F$  是从  $\mathcal{C} \times \mathcal{C}'$  到  $\mathcal{C}''$  的加法左精确双函子。

对于  $\mathrm{Ob}(\mathbf{C}^{+}(\mathcal{C}))$  中的复数  $X$  和  $\mathrm{Ob}(\mathbf{C}^{+}(\mathcal{C}'))$  中的复数  $X'$  ，  $F(X,X')$  是  $\mathcal{C}''$  中的双复数。通过将 $F(X,X')$  与简单复数 $s(F(X,X'))$  相关联，可以定义一个双函子：

\[
\mathbb{C}^{+}(F)\colon\mathbb{C}^{+}(\mathcal{C})\times\mathbb{C}^{+}(\mathcal{C}^{\prime})\to\mathbb{C}^{+}(\mathcal{C}^{\prime \prime})
\]

并传递给商，一个双函子：

\[
\mathbb{K}^{+}(F)\colon\mathbb{K}^{+}(\mathcal{C})\times\mathbb{K}^{+}(\mathcal{C}^{\prime})\to\mathbb{K}^{+}(\mathcal{C}^{\prime\prime})\enspace.
\]

（由于第 9 节中的最后评论，这是可能的。）

为了定义双函子的导出函子，从三角范畴的双函子开始会很方便。因此，让 $ F: \mathbf{K}^+(\mathcal{C}) \times \mathbf{K}^+(\mathcal{C}') \to \mathbf{K}^+(\mathcal{C}') $  成为三角范畴的双函子。令  $ \mathcal{I} $  （分别为  $ \mathcal{I}' $  ）为  $ \mathbf{K}^+(\mathcal{C}) $  （分别为  $ \mathbf{K}^+(\mathcal{C}') $  ）的完整三角子范畴。

考虑以下两个条件。

\[
\left\{\begin{array}{l}\text{For any}X\in\text{Ob}(\mathbb{K}^{+}(\mathcal{C}))(\text{resp.Ob}(\mathbb{K}^{+}(\mathcal{C}^{\prime})))\,\text{there is a}\\ \text{quasi-isomorphism from}X\text{to an object of}\mathcal{I}(\text{resp.}\mathcal{I}^{\prime})\end{array}\right.
\]

对于任何 $ X \in \text{Ob}(\mathcal{I}) $  和任何 $ X' \in \text{Ob}(\mathcal{I}') $  ，如果 $ X $  或 $ X' $  与零准同构，则复数 $ F(X, X') $  与零准同构。

\statementlabel{命题 1.10.4}在假设 (1.10.1)–(1.10.2) 下，存在双函子或三角范畴  $ RF: \mathbf{D}^{+}(\mathcal{C}) \times \mathbf{D}^{+}(\mathcal{C}') \to \mathbf{D}^{+}(\mathcal{C}'') $  使得下图可交换：

\[
\begin{array}{c c}{\mathcal{I}\times\mathcal{I}^{\prime}}&{\stackrel{F}{\longrightarrow}\mathbf{K}^{+}(\mathcal{C}^{\prime\prime})}\\ {\bigg\downarrow\mathcal{Q}\times\mathcal{Q}^{\prime}}&{\bigg\downarrow\mathcal{Q}^{\prime\prime}}\\ {\mathbf{D}^{+}(\mathcal{C})\times\mathbf{D}^{+}(\mathcal{C}^{\prime})\stackrel{R F}{\longrightarrow}\mathbf{D}^{+}(\mathcal{C}^{\prime\prime})~,}\\ \end{array}
\]

其中  $Q$  、  $Q'$  、  $Q''$  分别是从  $\mathcal{I}$  、  $\mathcal{I}'$  、  $\mathcal{C}$  '' 到  $\mathbf{D}^{+}(\mathcal{C})$  、  $\mathbf{D}^{+}(\mathcal{C}')$  、  $\mathbf{D}^{+}(\mathcal{C}'')$  的本地化函子。

此外，RF 满足以下通用属性：对于每个双函子或三角范畴  $G: \mathbf{D}^{+}(\mathcal{C}) \times \mathbf{D}^{+}(\mathcal{C}') \to \mathbf{D}^{+}(\mathcal{C}')$  ，规范同态：

\[
\operatorname{Hom}(RF,G)\to\operatorname{Hom}(Q^{\prime \prime}\circ F,G\circ(Q\times Q^{\prime}))
\]

是同构。这里  $ Q'' \circ F $  和  $ G \circ (Q \times Q') $  是从  $ \mathbf{K}^+(\mathcal{C}) \times \mathbf{K}^+(\mathcal{C}') $  到  $ \mathbf{D}^+(\mathcal{C}') $  的双函子。

\statementlabel{证明}很简单。

\statementlabel{推论 1.10.5}令 $F: \mathbb{K}^+(\mathcal{C}) \times \mathbb{K}^+(\mathcal{C}') \to \mathbb{K}^+(\mathcal{C}'')$  为三角范畴的双函子。假设存在  $\mathbb{K}^+(\mathcal{C})$  的满子范畴  $\mathcal{I}$ ，使得：

(1.10.3) 对于任何  $ Y \in \text{Ob}(\mathbb{K}^+(\mathcal{C}')) $  ，对于函子  $ F(\cdot, Y) $  满足条件 (1.8.4) 和 (1.8.5) 。

(1.10.4) 对于任何 $ X \in \text{Ob}(\mathcal{I}) $  、 $ Y \in \text{Ob}(\mathbb{K}^{+}(\mathcal{C}')) $  、 $ F(X, Y) $  ，如果Y 与零准同构，则 $ F(X, Y) $  与零准同构。

然后  $F$  接纳一个派生函子，对于  $Y\in\mathrm{Ob}(\mathbb{K}^{+}(\mathcal{C}'))$  ，函子  $F(\cdot,Y)\colon\mathbb{K}^{+}(\mathcal{C})\to\mathbb{K}^{+}(\mathcal{C}'')$  接纳一个派生函子，表示为  $R_I F(\cdot,Y)$  。此外，对于  $X\in\mathrm{Ob}(\mathbb{K}^{+}(\mathcal{C}))$  来说，有一个自然的同构：

\[
R F(X,Y)\simeq R_{I}F(X,Y)\enspace.
\]

\statementlabel{证明}$R_I F(\cdot, Y)$  的存在源自备注 1.8.10。  $RF$  的存在性通过取  $\mathcal{S}' = \mathbf{K}^+(\mathcal{C}')$  由命题 1.10.4 得出。等式 $RF(X, Y) = R_I F(X, Y)$  由构造得出。   $\square$ 

让我们回到  $F$  是  $\mathcal{C} \times \mathcal{C}'\mathcal{C}'$  的左精确双函子的情况。令  $\mathcal{I}$  （分别为  $\mathcal{I}'$  ）为  $\mathcal{C}$  （分别为  $\mathcal{C}'$  ）的满加性子范畴。

\statementlabel{定义 1.10.6}如果对于任何  $X \in \mathrm{Ob}(\mathcal{I})$  和任何  $X' \in \mathrm{Ob}(\mathcal{I}')$  ，  $\mathcal{I}$  是  $F(\cdot, X')$  -内射并且  $\mathcal{I}'$  是  $F(X, \cdot)$  -内射，我们说  $(\mathcal{I}, \mathcal{I}')$  是  $F$  -内射。

\statementlabel{命题 1.10.7}假设 $(\mathcal{I}, \mathcal{I}')$  是 $F$  -内射。那么 $(\mathbf{K}^{+}(\mathcal{I}), \mathbf{K}^{+}(\mathcal{I}'))$ 满足(1.10.1)–(1.10.2)。

\statementlabel{证明}应用命题 1.7.7 得到 (1.10.1)。让我们验证一下（1.10.2）。首先我们要证明如果 $X \in \mathrm{Ob}(\mathcal{I})$ 、 $Y \in \mathrm{Ob}(\mathbf{K}^{+}(\mathcal{I}^{\prime}))$  和 $Y$  与零准同构，则 $F(X, Y)$  与零准同构。

在这种情况下，

\[
H_{I I}^{q}(F(X,Y))=0\qquad\mathrm{f o r~a l l~q~}.
\]

应用定理1.9.3，我们得到 $F(X,Y)$  与零准同构。通过颠倒  $X$  和  $Y$  的角色也可以得到类似的证明。 ⑨

\statementlabel{推论 1.10.8}令 $F: \mathcal{C} \times \mathcal{C}' \to \mathcal{C}''$  为阿贝尔范畴的左精确双函子。假设对于任何  $Y \in \mathrm{Ob}(\mathcal{C}')$  ，对于函子  $F(\cdot, Y)$  ，存在满足定义 1.8.2 的条件 (i)、(ii)、(iii) 的  $\mathcal{C}$  的完整加性子范畴  $\mathcal{I}$  。此外，假设函子  $F(X, \cdot)$  对于任何  $X \in \mathrm{Ob}(\mathcal{I})$  都是精确的。那么范畴 $\mathbf{K}^+(\mathcal{I})$ 满足(1.10.3)–(1.10.4)。

\statementlabel{证明}使用  $ \mathcal{I}' = \mathcal{C}' $  应用命题 1.10.7。 □

\statementlabel{命题 1.10.9}令 $F: \mathcal{C} \times \mathcal{C}' \to \mathcal{C}''$  为阿贝尔范畴的左精确双函子，并令 $G: \mathcal{C}'' \to \mathcal{C}'''$  为阿贝尔范畴的左精确函子。假设 $\mathcal{C}$ 、 $\mathcal{C}'$ 和 $\mathcal{C}''$ 分别存在完整的加法子范畴 $\mathcal{I}$ 、 $\mathcal{I}'$ 和 $\mathcal{I}''$ ，使得 $(\mathcal{I}, \mathcal{I}')$ 是 $F$  -单射， $\mathcal{I}'$ 是 $G$  -单射，并且：

\[
F(\mathrm{O b}(\mathcal{I}),\mathrm{O b}(\mathcal{I}^{\prime}))\subset\mathrm{O b}(\mathcal{I}^{\prime \prime})~.
\]

那么导出函子 $R(G\circ F)\colon\mathbf{D}^{+}(\mathcal{C})\times\mathbf{D}^{+}(\mathcal{C}')\to\mathbf{D}^{+}(\mathcal{C}''))$ 存在，并且有：

\[
R(G\circ F)\simeq R G\circ R F\enspace.
\]

\statementlabel{证明}很简单。

对于函子的其他组合也有类似的结果。

\statementlabel{备注 1.10.10}在推论 1.10.5 中，如果我们互换  $ \mathcal{C} $  和  $ \mathcal{C}' $  的角色，那么就可以定义  $ R_{II}F(X,\cdot) $  。当然，人们会得到：

\[
R_{I}F(X,Y)\simeq R_{I I}F(X,Y)\simeq R F(X,Y)\enspace.
\]

示例 1.10.11。令 $\mathcal{C}$  为阿贝尔范畴。那么  $\mathrm{Hom}_{\mathcal{C}}(\cdot,\cdot)$  是从  $\mathcal{C}^{\circ}\times\mathcal{C}$  到  $\mathrm{Mod}(\mathbb{Z})$  的加性双函子。如果  $\mathcal{C}$  承认足够的单射（或者足够的射影），那么  $\mathrm{Hom}_{\mathcal{C}}(\cdot,\cdot)$  承认一个右派生函子：

\[
R\mathbf{H}\mathbf{o}\mathbf{m}_{\mathcal{C}}(\cdot,\cdot):(\mathbf{D}^{-}(\mathcal{C}))^{\circ}\times\mathbf{D}^{+}(\mathcal{C})\to\mathbf{D}^{+}(\mathfrak{M o b}(\mathbb{Z}))~.
\]

函子  $ H^n(\cdot) \circ \mathrm{RHom}_{\mathcal{G}}(\cdot, \cdot) $  表示为  $ \mathrm{Ext}^n_{\mathcal{G}}(\cdot, \cdot) $  。

示例 1.10.12。让 $A$  成为一个环。那么  $\otimes_A \cdot$  是从  $\mathfrak{Mod}(A^{op}) \times \mathfrak{Mod}(A)$  到  $\mathfrak{Mod}(\mathbb{Z})$  的加性双函子。对于两个变量来说，它都是一个右正合函子。一个由  $\cdot \otimes_A^L \cdot$  表示来自  $\mathbf{D}^-(\mathfrak{Mod}(A^{op})) \times \mathbf{D}^-((\mathfrak{Mod}(A)) \to \mathbf{D}^-(\mathfrak{Mod}(\mathbb{Z})))$  的左派生函子，一组： $\operatorname{Tor}_n^A(N,M) = H^{-n}(N \otimes_A^L M)$  。

符号 1.10.13。令 $A$  为交换环。仍然用  $R\mathrm{Hom}_{A}(\cdot,\cdot)$  或  $\otimes_{A}^{L}$  表示  $\mathrm{Hom}_{A}(\cdot,\cdot)$  或  $\otimes_{A}(\cdot,\cdot)$  的派生函子，其值在  $\mathbf{D}(\mathrm{Mod}(A))$  中。

\statementlabel{备注 1.10.14}令  $F: \mathcal{C} \times \mathcal{C}' \to \mathcal{C}''$  为阿贝尔范畴的双函子，并令  $G: \mathcal{C}' \times \mathcal{C} \to \mathcal{C}''$  为  $G(Y, X) = F(X, Y)$  给出的关联双函子。我们定义同构  $r: \mathbf{K}^+ (F) \simeq \mathbf{K}^+ (G)$  ，如下所示。令  $X = (X^n)_{n \in \mathbb{Z}}$  、  $Y = (Y^n)_{n \in \mathbb{Z}}$  分别为  $\mathbf{C}(\mathcal{C})$  和  $\mathbf{C}(\mathcal{C}')$  的两个对象。那么 $r(X, Y)$  就是

 $ (-1)^{nm}\mathrm{id}\colon F(X^{n}, Y^{m}) \to G(Y^{m}, X^{n}) $  的直接和。检查 $ r(X, Y) $  是简单复形的态射。

\statementlabel{备注 1.10.15}令 $ F: \mathcal{C}' \times \mathcal{C}^\circ \times \mathcal{C} \to \mathcal{C}'' $  为三函数， $ G: \mathcal{C}' \to \mathcal{C}'' $  为函子。假设对所有  $ Y \in \mathrm{Ob}(\mathcal{C}') $  和  $ X \in \mathrm{Ob}(\mathcal{C}) $  给出态射

\[
\alpha(Y,X)\colon F(Y,X,X)\to G(Y)
\]

关于 Y 的函子，使得对于  $\mathcal{C}$  中的任何态射  $f: X \to X'$ ，该图：

\begin{center}
\includegraphics[width=0.44\textwidth]{流行上的层_Chn-assets/img_in_image_box_317_511_848_685.png}
\end{center}

通勤。然后对于 $ Y \in \text{Ob}(\mathbf{K}(\mathcal{C}')) $  和 $ X \in \text{Ob}(\mathbf{K}(\mathcal{C})) $  可以定义态射

\[
\mathbb{K}(\alpha)(Y,X)\colon K(F)(Y,X,X)\to K(G)(Y)
\]

作为态射的直和

\[
F(Y^{m},X^{n},X^{n})\to G(Y^{m})~.
\]

使用备注 1.8.11 中引入的约定，可以轻松检查  $ \mathbf{K}(\alpha)(X, Y) $  是否是复数的态射。

例如，如果  $A$  是具有有限全局维数的交换环（参见练习 I.28），我们可以在  $\mathbf{D}^{b}(\mathfrak{Mod}(A))$  中定义  $X$  和  $Y$  ：

\[
\mathbf{D}(\alpha)(Y,X)\colon R\mathbf{H o m}_{A}(X,Y)\overset{L}{\underset{A}{\otimes}}\;X\to Y\;.
\]

事实上，采用准同构  $ Y \to I $ ，其中 I 是有界且单射的，以及准同构  $ P \to X $ ，其中 P 是有界且射影的。

然后是态射：

\[
\mathbb{K}\left(\alpha\right)(I,P){:~}\operatorname{H o m}_{A}\left(P,I\right)\otimes_{A}P\to I
\]

由上述注释明确定义，并定义  $ \mathbf{D}(\alpha)(Y, X) $  。

\statementlabel{备注 1.10.16}让 $F$  如备注1.10.14 中所示。然后，我们通过身份识别 $F(X[p], Y) \simeq F(X, Y)[p]$ ：

\[
(F(X[p],Y))^{n}\simeq\prod_{k}F(X^{k+p},Y^{n-k})\simeq(F(X,Y)[p])^{n}
\]

和  $ F(X, Y[q]) \simeq F(X, Y)[q] $  的身份：

\[
(F(X,Y[q]))^{n}\simeq\prod_{k}F(X^{k},Y^{n+q-k})\xrightarrow{(-)^{k\cdot q}}\prod_{k}F(X^{k},Y^{n+q-k})\simeq(F(X,Y)[q])^{n},
\]

以便这成为复形的态射。

然后

\[
\begin{array}{c}F(X[p],Y[q])\xrightarrow{\quad}\quad F(X,Y[q])[p]\\\downarrow\quad\downarrow\quad\downarrow\quad\downarrow\\F(X[p],Y)[q]\xrightarrow{\quad}\quad F(X,Y)[p+q]\end{array}
\]

根据 pq 的奇偶性进行交换或反交换。

\statementlabel{备注 1.10.17}符号问题总是非常微妙的。 Deligne [1] 和 Berthelot-Breen-Messing [1] 的第一页对它们进行了非常仔细的处理，但最后作者采用的约定与我们不同。特别是，在 (1.4.3) 中，它们将  $ \beta $  替换为  $ -\beta $  。

\subsection{归纳对象与射影对象}
令  $ \mathcal{C}^\vee $  为从  $ \mathcal{C} $  到  $ \mathcal{Set} $  的逆变函子的范畴。让 $ h: \mathcal{C} \to \mathcal{C}^\vee $  成为函子： $ X \mapsto \mathrm{Hom}_{\mathcal{C}}(\cdot, X) $  。通过这个函子，我们将 $ \mathcal{C} $ 视为 $ \mathcal{C}^\vee $ 的满子范畴（参见§1）。回想一下，从  $ \mathcal{C} $  到  $ \mathcal{Set} $  的逆变函子  $ F $  是可表示的，当且仅当  $ F $  与  $ \mathcal{C}^\vee $  中的  $ \mathcal{C} $  对象同构。

让 $\mathcal{I}$  成为一个范畴。对于另一个范畴  $\mathcal{C}$  ，  $\mathcal{C}$  中由  $\mathcal{I}$  索引的归纳系统根据定义是从  $\mathcal{I}$  到  $\mathcal{C}$  的函子。类似地，射影系统是从  $\mathcal{I}^\circ$  到  $\mathcal{C}$  的函子。令  $F_\circ$  为  $\mathfrak{Set}$  中的射影系统，由  $\mathcal{I}$  索引，由  $F_\circ(i) = \{\mathrm{pt}\}$  给定任何  $i \in \mathrm{Ob}(\mathcal{I})$  。那么对于由  $\mathcal{I}$  索引的射影系统  $F : \mathcal{I}^\circ \to \mathfrak{Set}$  ，从  $F_\circ$  到  $F$  的态射集称为  $F$  的射影极限，并用  $\lim_{\overleftarrow{f}} F$  或  $\lim_{i \in \mathcal{I}} F(i)$  表示。即 $\lim_{\overleftarrow{f}} F$  由 $\{x(i)\}_{i \in \mathrm{Ob}(\mathcal{I})}$  组成，使得 $x(i) \in F(i)$  和 $F(h)x(j) = x(i)$  对于任何 $h : i \to j$  。

我们有时会写  $ i \in \mathcal{I} $  而不是  $ i \in \mathrm{Ob}(\mathcal{I}) $  和  $ \lim_{i \to \infty} \lim_{i \to \infty} \lim_{i \to \infty} \lim_{i \to \infty} f(f(f(f))) $  。

\statementlabel{定义 1.11.1}设 I 和 C 为两个范畴。

(a) 对于  $\mathcal{C}$  中由  $\mathcal{I}$  索引的归纳系统  $F$ ，一个由  $\lim\overrightarrow{f}F$  表示从  $\mathcal{C}$  到  $\mathcal{Set}$  的函子：

\[
X\mapsto\operatorname*{l i m}_{i\in{\mathcal{S}}}\operatorname{H o m}_{\mathcal{C}}(F(i),X)\enspace.
\]

(b) 对于  $\mathcal{C}$  中由  $\mathcal{I}$  索引的射影系统  $F$ ，其中  $\lim_{\overleftarrow{f}} F$  表示从  $\mathcal{C}^\circ$  到  $\text{Set}$  的函子：

\[
X\mapsto\operatorname*{l i m}_{i\in{\mathcal{S}}}\operatorname{H o m}_{\mathcal{C}}(X,F(i))\enspace.
\]

如果这些函子是可表示的，则在  $\mathcal{C}$  中保留相同的符号来表示它们的代表，并分别将它们称为归纳极限和射影极限。

如果 $F$  是 $\mathcal{S}$  中的射影系统，则定义1.11.1(b) 与前一个一致。

\statementlabel{定义 1.11.2}如果范畴 $\mathcal{I}$ 非空且满足以下条件，则称为过滤剂。

\[
\left\{\begin{aligned}{}&{{}{F o r}\;i,j\in\operatorname{O b}(\mathcal{I}),\;{t h e r e~e x i s t}\;k\in\operatorname{O b}(\mathcal{I})\;{a n d~m o r p h i s m s}}\\ {}&{{}\begin{aligned}{}&{{}i\to k\;{a n d}j\to k\enspace.}\\ \end{aligned}}\\ \end{aligned}\right.
\]

\[
\left\{\begin{aligned}{}&{{}{F o r~t w o~m o r p h i s m s}f,g\in\operatorname{H o m}_{\mathcal{S}}(i,j),{t h e r e~e x i s t s~a~m o r p h i s m}}\\ {}&{{}h:j\to k{~s u c h~t h a t}h\circ f=h\circ g.}\\ \end{aligned}\right.
\]

下面的结果很容易证明。

\statementlabel{命题 1.11.3}令 $F$  为 $\mathfrak{Set}$  中的归纳系统，由 $a$  过滤剂范畴 $\mathscr{I}$  索引。然后 $\lim_{\mathscr{F}} F$ 由 $(\bigsqcup_{i\in\text{Ob}(\mathscr{I})} F(i)) / \sim$ 表示，其中对于 $x \in F(i)$ 和 $y \in F(j)$ ， $x \sim y$ 当且仅当存在 $f: i \to k$ 和 $g: j \to k$ 使得 $F(f)x = F(g)y$ 。

\statementlabel{定义 1.11.4}令 $ \mathcal{I} $  和 $ \mathcal{C} $  为两个范畴， $ \mathcal{I} $  为过滤剂。

(a) 对于  $\mathcal{C}$  中由  $\mathcal{I}$  索引的归纳系统  $F$ ，用“  $\lim\overrightarrow{f}$  ”  $F$  表示从  $\mathcal{C}^{\circ}$  到  $\mathcal{Set}$  的函子：

\[
X\mapsto\varinjlim_{i\in\mathcal{I}}\operatorname{H o m}_{\mathcal{C}}(X,F(i))\enspace.
\]

(b) 对于  $\mathcal{C}$  中由  $\mathcal{I}$  索引的射影系统  $F$ ，其中一个用“  $\lim_{\mathcal{I}}$  ”  $F$  表示从  $\mathcal{C}$  到  $\mathcal{Set}$  的函子：

\[
X\mapsto\operatorname*{l i m}_{i\stackrel{\rightarrow}{\in}\mathcal{I}}\operatorname{H o m}_{\mathcal{C}}(F(i),X)\enspace.
\]

如果这些函子是可表示的，则可以使用相同的符号来表示它们在 C 中的代表。

从 $\mathcal{C}^\circ$ （或 $\mathcal{C}$ ）到 $\mathfrak{Set}$ （即 $\mathcal{C}^\vee$ （或 $\mathcal{C}^\wedge$ ）的对象）的函子如果与“ $\lim$ ” $F$ （或“ $\lim$ ”同构），则称为ind-object（或pro-object）  $F$  ) 用于  $\mathcal{C}$  中的感应（或射影）系统  $F$ ，由过滤剂范畴  $\mathcal{I}$  索引。

示例 11.1.5。令 $ (I, \leq) $  为有序集。我们将  $ I $  关联到一个范畴，表示为  $ \mathcal{I} $  ，如下所示。我们设定：

\[
\left\{\begin{aligned}{}&{{}\operatorname{O b}(\mathcal{I})=I~,}\\ {}&{{}\operatorname{H o m}_{\mathcal{S}}(i,j){c o n s i s t s~o f~a~s i n g l e~e l e m e n t~}s_{j i}{i f~}i\leqslant j{a n d}}\\ {}&{{}{i s~e m p t y~o t h e r w i s e~~.}}\\ \end{aligned}\right.
\]

那么对于范畴  $\mathcal{C}$  ，由  $I$  索引的  $\mathcal{C}$  中的归纳系统  $F$  根据定义是由  $\mathcal{I}$  索引的  $\mathcal{C}$  中的归纳系统。

人们通常写 $\{X_{i},\rho_{j,i}\}$ 而不是 $F$ ，其中 $X_{i}=F(i)$ ， $\rho_{i,j}=F(s_{ij})$ 。类似地，定义了一个由  $\mathcal{I}$  索引的射影系统。

如果  $I$  是有向有序集（即对于任何  $i \in I, j \in I$  都存在  $k \in I$  和  $k \geq i$  、  $k \geq j$  ），则范畴  $\mathcal{I}$  是过滤剂。

令 $F$  为 $\mathcal{C}$  中的归纳系统，由过滤剂范畴 $\mathcal{I}$  索引。假设 $\lim F(j)$ 是可表示的，并用 $K$ 表示其代表。我们得到：

\[
\nearrow
\]

\[
\operatorname{H o m}_{\mathcal{C}}(K,K)\simeq\varprojlim_{j}\operatorname{H o m}_{\mathcal{C}}(F(j),K)\enspace.
\]

这个同构定义了一系列态射， $\rho_{j}: F(j) \to K$  满足：

\[
\rho_{j}\circ F(s)=\rho_{i}~,\qquad\mathrm{f o r~a n y}s:i\to j
\]

和  $ \{K,\rho_{i}\}_{i} $  满足由图表示的通用属性：

\[
F(s)\overset{X_{j}\longrightarrow}{\overset{f_{j}}{\longrightarrow}}\overset{\rho_{j}\longrightarrow}{\underset{\ell\to\infty}{\rho_{i}}}\overset{K\to\infty}{\longrightarrow}X\enspace.
\]

也就是说，对于任何态射族  $f_i: X_i \to X$  使得  $f_j \circ F(s) = f_i$  对于任何  $s: i \to j$  ，虚线箭头唯一存在，使得图可交换。

现在假设“ $ \underline{\text{lim}} $ ” $ F(j) $ 可表示，并用L表示其代表。

对于任何  $ i \in \mathrm{Ob}(\mathcal{I}) $  ，同构：

\[
\operatorname{H o m}_{\mathcal{C}}(F(i),L)\simeq\varinjlim\operatorname{H o m}_{\mathcal{C}}(F(i),F(j))
\]

和  $\mathrm{id}_{F(i)}$  定义态射  $\rho_{i}\colon F(i)\to L$  ，它满足 (1.11.4)。此外，根据同构：

\[
\operatorname{H o m}_{\mathcal{C}}(L,L)\simeq\varinjlim\operatorname{H o m}(L,F(j))~,
\]

我们找到一个  $ i_o \in \mathrm{Ob}(\mathcal{S}) $  和一个态射  $ f: L \to F(i_o) $  这样：

\[
\rho_{i_{0}}\circ f=\operatorname{i d}_{L}\;.
\]

对于任何  $ X \in \mathrm{Ob}(\mathcal{C}) $  和任何  $ i \in \mathrm{Ob}(\mathcal{I}) $  ，态射的组成：

\[
\mathrm{H o m}_{\mathcal{C}}(X,F(i))\xrightarrow{\rho_{i}}\mathrm{H o m}_{\mathcal{C}}(X,L)\xrightarrow{f}\mathrm{H o m}_{\mathcal{C}}(X,F(i_{\circ}))\xrightarrow{\quad}\varinjlim_{j}\mathrm{H o m}_{\mathcal{C}}(X,F(j))
\]

与规范一致。因此：

(1.11.6) 对于任何  $i \in \mathrm{Ob}(\mathcal{I})$  都存在  $j \in \mathrm{Ob}(\mathcal{I})$  、  $s: i \to j$  和  $t: i_\circ \to j$  使得  $F(t) \circ f \circ \rho_i = F(s)$  。

事实上，这些属性确保了“ $ \lim_{x \to 0} $ ”在 $ \mathrm{Ob}(\mathcal{C}) $ 中的存在。

\statementlabel{命题 1.11.6}让 $L \in \mathrm{Ob}(\mathcal{C})$  。那么 $L$ 表示“ $\lim_{\overrightarrow{f}}$ ” $F(j)$ 当且仅当对于任何 $i \in I$ 、 $i_\circ \in \mathrm{Ob}(\mathcal{I})$ 和 $f: L \to F(i_\circ)$ 都存在 $\rho_i: F(i) \to L$ ，使得满足(1.11.4)、(1.11.5)和(1.11.6)。

\statementlabel{证明}我们已经看到了这些条件的必要性。为了证明相反的情况，请注意对于任何  $ X \in \mathrm{Ob}(\mathcal{C}) $  ，两个映射：

\[
\varinjlim_{\vec{j}}\operatorname{H o m}_{\mathcal{C}}(X,F(j))\xrightarrow{\{\rho_{j}\}}\operatorname{H o m}_{\mathcal{C}}(X,L)
\]

and

\[
\mathrm{H o m}_{\mathcal{C}}(X,L)\xrightarrow{f}\mathrm{H o m}_{\mathcal{C}}(X,F(i_{\circ}))\longrightarrow\varinjlim_{i}\mathrm{H o m}_{\mathcal{C}}(X,F(j))
\]

彼此相反。 ⑨

\statementlabel{推论 1.11.7}如果“  $\lim\limits_{j}\overrightarrow{F}(j)$  可由  $L$  表示，则  $\lim\limits_{j}\overrightarrow{F}(j)$  可由  $L$  表示。

\statementlabel{推论 1.11.8}令 $\mathcal{C}'$  为另一个范畴， $T$  为从 $\mathcal{C}$  到 $\mathcal{C}'$  的函子，让 $F$  为 $\mathcal{C}$  中的归纳系统。如果“ $\lim$ ” $F(j)$ 存在于 $\mathcal{C}$ 中，则“ $\lim$ ” $TF(j)$ 存在于 $\mathcal{C}'$ 中，并由 $T$ (“ $\lim$ ” $F(j)$ )表示。

\statementlabel{证明}命题 1.11.6 中  $ \rho_{i}, i_{\circ} $  、 f 满足 (1.11.4)–(1.11.6) 的存在性被 T 保留。 $ \Box $ 

在这里我们讨论了归纳极限，但同样的结果也适用于射影极限，因为  $\mathcal{C}$  中的射影系统是  $\mathcal{C}^{\circ}$  中的归纳系统。

\subsection{米塔--列夫勒条件}
在本节中，我们将主要研究阿贝尔群的射影系统，索引为  $\mathbb{N}$  。我们用 $\mathfrak{Ab}$  表示阿贝尔群的范畴。我们将有序集 $(\mathbb{N}, \leq)$ 视为一个范畴，并用 $\underline{\mathbb{N}}$ 表示它（参见§11）。因此，由  $\mathbb{N}$  索引的阿贝尔群射影系统是从  $\mathbb{N}^\circ$  到  $\mathfrak{Ab}$  的函子。让 $X$ 成为这样一个函子。我们还写  $\{X_n, \rho_{n,p}\}$  而不是  $X$  。因此  $\rho_{n,p}: X_p \to X_n$  是为  $p \geq n$  定义的阿贝尔群的态射，并且这些态射满足：

\[
\rho_{n,n}=\operatorname{id}_{X_{n}}~,\qquad\rho_{n,p}\circ\rho_{p,q}=\rho_{n,q}\qquad\mathrm{for}~n\leqslant p\leqslant q~.
\]

令Hom  $ (\mathbb{N}^\circ, \mathbb{Ab}) $  为从 $ \mathbb{N} $  到 $ \mathbb{Ab} $  的逆变函子的范畴（即，由 $ \mathbb{N} $  索引的阿贝尔群射影系统的范畴）。这显然是一个阿贝尔范畴。

由于范畴  $ \mathcal{Ab} $  承认射影极限（也承认归纳极限，但我们暂时不研究它们），因此存在一个明确定义的函子：

\[
\varprojlim:\mathrm{Hom}(\mathbb{N}^{\circ},\mathfrak{A}\mathfrak{b})\to\mathfrak{A}\mathfrak{b}
\]

它与射影系统及其射影极限相关联。如果 $ X = \{X_n, \rho_{n,p}\} $ 是一个射影系统，我们应该写 $ \lim X $ 以及 $ \lim_{n} X_n $ 来表示它的射影极限。

请注意，函子  $ \lim $  是精确的。我们将构造  $ \text{Hom}(\mathbb{N}^\circ, \mathbb{U})_\text{b}) $  的满子范畴，相对于该函子的单射。

\statementlabel{定义 1.12.1}令 $X = \{X_n, \rho_{n,p}\}$  为交换群射影系统。有人说 $X$  满足Mittag-Leffler 条件（简称M-L），如果对于任何 $n \in \mathbb{N}$  ， $X_n$  子群的递减序列 $\{\rho_{n,p}(X_p)\}_{p \geq n}$  是平稳的。

\statementlabel{命题 1.12.2}令 $ 0 \to X' \to X \to X'' \to 0 $  为交换群射影系统的精确序列。

(i) 假设 $X'$  和 $X''$  满足 $M$  -L 条件。那么 $X$  满足 $M$  -L 条件。

(ii) 假设X满足M-L条件。那么 $ X'' $ 满足M-L条件。

\statementlabel{证明}(i) 对于任何  $n, N_{\circ}, N_{1}, p$  和  $n \leqslant N_{\circ} \leqslant N_{1} \leqslant p$ ，我们有一个包含精确行的交换图：

\begin{center}
\includegraphics[width=0.53\textwidth]{流行上的层_Chn-assets/img_in_image_box_254_1140_886_1744.png}
\end{center}

对于固定 $n$ ，我们选择 $N_{\circ}$ ，然后选择 $N_{1} \geqslant N_{\circ}$ ，这样

\[
\mathrm{I m}(X_{p}^{\prime}\to X_{n}^{\prime})=\mathrm{I m}(X_{N_{\circ}}^{\prime}\to X_{n}^{\prime})\qquad\mathrm{f o r~a n y~}p\geqslant N_{\circ}\quad\mathrm{a n d}
\]

\[
\mathrm{I m}(X_{p}^{\prime\prime}\to X_{N}^{\prime\prime})=\mathrm{I m}(X_{N_{1}}^{\prime\prime}\to X_{N_{0}}^{\prime\prime})\qquad\mathrm{f o r~a n y~}p\geqslant N_{1}\enspace.
\]

自 $ g(\mathrm{Im}(X_{N_1} \to X_{N_0})) = g(\mathrm{Im}(X_p \to X_{N_0})) $ 以来，我们有：

\[
\operatorname{I m}(X_{N_{1}}\to X_{N_{\circ}})\subset\operatorname{I m}(X_{p}\to X_{N_{\circ}})+f(X_{N_{\circ}}^{\prime})~.
\]

因此：

\[
\begin{aligned}{\operatorname{I m}(X_{N_{1}}\to X_{n})}&{{}\subset\operatorname{I m}(X_{p}\to X_{n})+f(\operatorname{I m}(X_{N_{\circ}}^{\prime}\to X_{n}^{\prime}))}\\ {}&{{}\subset\operatorname{I m}(X_{p}\to X_{n})\enspace.}\\ \end{aligned}
\]

(ii) 是显而易见的。 □

\statementlabel{命题 1.12.3}令 $ 0 \to X' \to X \xrightarrow{J} X'' \to 0 $  为射影系统的精确序列。假设 $ X' $  满足M-L 条件。然后是顺序：

\[
0\to\varprojlim X^{\prime}\to\varprojlim X\to\varprojlim X^{\prime \prime}\to0
\]

是准确的。

我们选择一个

 $\operatorname{Im}(X_{v(n)}^{\prime}) \to X_{n}^{\prime})\operatorname{ai}_{v}g(x_{v(n)})=x_{v(n)}^{\prime\prime}$  ，那么对于  $\operatorname{each}_{x}$  存在  $x_{p}\in X_{p}$  ，它满足  $g(x_{p})=x_{p}^{\prime\prime}$  和  $\rho_{n,p}(x_{p})=\rho_{n,v(n)}(x_{v(n)})$  ，i.o 检查这一点让  $y_{p}\in X_{p}$  和  $g(y_{p})=x_{p}^{\prime\prime}$  。那么 $\rho_{v(n),p}(y_{p})-x_{v(n)}$  属于 $f(X_{v(n)})$ ，因此 $\rho_{n,p}(y_{p})-\rho_{n,v(n)}(x_{v(n)})\in\operatorname{Im}(f(X_{p}^{\prime})\to f(X_{n}^{\prime}))$ ，我们可以找到 $z_{p}\in f(X_{p}^{\prime})$ ，使得 $y_{p}-z_{p}$  具有所需的属性。现在我们构造一个序列  $x=\{x_{v(n)}\}_{n}$  使得  $x_{v(n)}^\prime$  和  $x_{v(n)}^\prime$  )。

假设我们已经为 $ i < n_o $ 构造了 $ x_{v(i)} \in X_{v(i)} $ ，那么我们可以通过前面的注释构造 $ x_{v(n_o)} $ 。

那么  $ x_{n} = \rho_{n,v(n)}(x_{v(n)}) $  满足  $ g(x_{n}) = x_{n}' $  和  $ \rho_{n-1,n}(x_{n}) = x_{n-1} $  。 □

现在我们考虑阿贝尔群复形的射影系统，或者等价地，阿贝尔群射影系统的复形。令 $ X' = \{X^k, \mathrm{d}^k\} $  为阿贝尔群射影系统的复形：对于每个 $ k $  ， $ X^k = \{X_n^k, \rho_{n,p}^k\} $  是阿贝尔群射影系统，并且态射 $ \mathrm{d}^k, \rho_{n,p}^k $  满足自然相容条件。

对于 $X^{\prime}$ ，我们可以将阿贝尔群的复形关联起来，表示为 $X_{\infty}^{\prime}$ ，通过：

\[
X_{\infty}^{\cdot}=\varprojlim X^{\cdot}=\left\{\varprojlim X^{k},\mathsf{d}^{k}\right\}\quad.
\]

对于每个  $ n \in \mathbb{N} $  ，自然态射  $ X_{\infty}^: \to X_n $  为每个  $ k \in \mathbb{Z} $  定义，

态射：

\[
\phi_{k}:H^{k}(X_{\infty}^{.})\to\varprojlim_{n}H^{k}(X_{n}^{.})\enspace.
\]

\statementlabel{命题 1.12.4}假设对于每个  $ k \in \mathbb{Z} $  ，系统  $ X^k $  满足 M-L 条件。然后：

(a) 对于每个  $k$  ，  $\phi_{k}$  是满射的，

(b) 此外，如果对于给定的  $i$  ，系统  $H^{i-1}(X^*)$  （即射影系统  $\{H^{i-1}(X_n^*, H^{i-1}(\rho_{n,p})\})$  ）满足  $M$  -L 条件，则  $\phi_i$  是双射的。

\statementlabel{证明}设置  $ Z_n^k = \mathrm{Ker}(d_n^k; X_n^k \to X_n^{k+1}) $  和  $ B_n^k = \mathrm{Im}(d_{n-1}^k; X_{n-1}^k \to X_n^k) $  。我们有精确的序列：

\[
0\to Z_{n}^{k}\to X_{n}^{k}\to B_{n}^{k+1}\to0\enspace,
\]

\[
0\to B_{n}^{k}\to Z_{n}^{k}\to H^{k}(X_{n}^{\cdot})\to0\quad,
\]

\[
0\to\operatorname*{l i m}_{\leftarrow\atop{n}}\;B_{n}^{k}\to\operatorname*{l i m}_{\leftarrow\atop{n}}Z_{n}^{k}\to\operatorname*{l i m}_{\leftarrow\atop{n}}H^{k}(X_{n}^{\cdot})\to0\enspace.
\]

事实上，(1.12.4) 由命题 1.12.3 得出，因为射影系统 \textbackslash{}\{B\_n\textasciicircum{}k\textbackslash{}\}\_\{n\} 满足命题 1.12.2 的 M-L 条件。此外， $ \lim(\cdot) $  是左精确函子，我们有：

\[
\begin{aligned}{\varprojlim_{n}Z_{n}^{k}}&{{}=\operatorname{K e r}\left(\varprojlim_{n}X_{n}^{k}\to\varprojlim_{n}X_{n}^{k+1}\right)}\\ {}&{{}=\operatorname{K e r}(X_{\infty}^{k}\to X_{\infty}^{k+1})\enspace.}\\ \end{aligned}
\]

考虑图 (1.12.6)，其行是精确的：

\begin{center}
\includegraphics[width=0.76\textwidth]{流行上的层_Chn-assets/img_in_image_box_139_1212_1052_1410.png}
\end{center}

从该图中我们得知 $\phi_{k}$  是满射的。

然后假设 $ \{H^{i-1}(X_n^i)\}_n $ 满足M-L条件。由(1.12.3)和命题1.12.2，我们发现射影系统 $ \{Z_n^{i-1}\}_n $ 满足M-L条件。将命题 1.12.3 应用于 (1.12.2)，我们得到精确的序列：

\[
0\to\varprojlim\ Z_{n}^{i-1}\to\varprojlim\ X_{n}^{i-1}\to\varprojlim\ B_{n}^{i}\to0\enspace.
\]

因此，在图(1.12.6)中， $ \psi_{i-1} $ 是满射的，而 $ \phi_i $ 是单射的。 □

\statementlabel{备注 1.12.5}让我们用  $ \mathrm{Hom}(\mathcal{I}, \mathcal{U}\mathrm{b}) $  表示从范畴  $ \mathcal{I} $  到  $ \mathcal{U}\mathrm{b} $  的函子范畴，即由  $ \mathcal{I} $  索引的阿贝尔群归纳系统的范畴。有一个明确定义的函子：

\[
\lim_{\rightarrow}:\mathrm{Hom}(\mathcal{I},\mathfrak{A}\mathfrak{b})\rightarrow\mathfrak{A}\mathfrak{b}
\]

它与归纳系统相关，即它的归纳极限。与函子  $ \lim $  相反，函子  $ \lim $  是精确的。

特别是，考虑阿贝尔群的归纳系统复体  $ X' = \{X^k, d^k\} $  ，其中每个  $ X^k = \{X_i^k, \rho_{i,j}^k\} $  都是一个归纳系统。然后我们对每个  $ k \in \mathbb{Z} $  都有：

\[
H^{k}(\varinjlim X^{\cdot})\simeq\varinjlim H^{k}(X^{\cdot})~.
\]

为了结束本节，我们将回顾一下由  $ \mathbb{R} $  索引的集合的归纳和射影极限的结果（参见 Kashiwara [5]）。

\statementlabel{命题 1.12.6}令  $\{X_s, \rho_{s,t}\}$  为由  $\mathbb{R}$  索引的集合的射影系统。假设对于每个  $s \in \mathbb{R}$  ，规范映射：

\[
\lambda_{s}\colon X_{s}\to\varprojlim_{r<s}X_{r}
\]

and

\[
\mu_{s}\colon\varinjlim_{t>s}X_{t}\to X_{s}
\]

都是单射的（分别是满射的）。那么所有映射  $ \rho_{s_0,s_1} $  (  $ s_0 \leqslant s_1 $  ) 都是单射的（或满射的）。

\statementlabel{证明}我们首先证明单射性。让  $x$  和  $y$  属于  $X_{s_1}$  ，这样  $\rho_{s_0,s_1}(x)=\rho_{s_0,s_1}(y)$  就属于某些  $s_0<s_1$  。让：

\[
I=\left\{s\in\mathbb{R};s\leqslant s_{1},\rho_{s,s_{1}}(x)=\rho_{s,s_{1}}(y)\right\}.
\]

那么  $I$  包含  $s_0$  ，并且  $s \in I$  ，  $r < s$  ，暗示  $r \in I$  。让 $s_2 = \sup I$  。由于  $\lambda_{s_2}$  是单射的，因此  $s_2$  属于  $I$  。如果  $s_2 = s_1$  ，则  $x = y$  。假设 $s_2 < s_1$  。由于  $\mu_{s_2}$  是单射的，因此存在  $s > s_2$  和  $s \in I$  。这是一个矛盾。

现在我们证明满射性。让 $ s_0 < s_1 $ ， $ x_0 \in X_{s_0} $ 。令  $ A $  为  $ (s, x) $  与  $ s_0 \leqslant s \leqslant s_1 $  、  $ x \in X_s $  和  $ \rho_{s,s_0}(x) = x_0 $  的对的集合。我们可以按如下方式订购 $ A $ ：

\[
(s,x)\leqslant(s^{\prime},x^{\prime})\quad iff\quad s\leqslant s^{\prime}\quad and\quad\rho_{s,s^{\prime}}(x^{\prime})=x\ .
\]

让我们证明  $A$  是归纳排序的。令  $B \subset A$  、  $B$  完全有序，并令：

\[
I=\left\{s\in\mathbb{R};s_{0}\leqslant s\leqslant s_{1},there exists x\in X_{s},with(s,x)\in B\right\}
\]

让 $s_2 = \sup I$  。如果  $s_2 \in I$  ，则  $B$  拥有最大元素。如果  $s_2$  不属于  $I$  ，则存在  $(s_2, x_2) \in A$  大于  $B$  的任意元素的满射性。这证明 $A$  是归纳有序的。

令  $(s,x)$  为  $A$  的最大元素。如果 $s=s_1$ ，证明就结束了。否则，存在  $s',s<s'\leq s_1$  和  $x'\in X_s$  使得  $\rho_{s,s'}(x')=x$  ，因为  $\mu_s$  是满射的。这是一个矛盾。   $\square$ 

\subsection{习题}
\statementlabel{练习 I.1}令 $\mathcal{C}$  为加法范畴。证明只有一种方法可以赋予每个  $\operatorname{Hom}_{\mathcal{C}}(X, Y)$  一个加法群的结构，使得组合律  $\operatorname{Hom}_{\mathcal{C}}(X, Y) \times \operatorname{Hom}_{\mathcal{C}}(Y, Z) \to \operatorname{Hom}_{\mathcal{C}}(X, Z)$  对于所有  $X, Y, Z \in \operatorname{Ob}(\mathcal{C})$  都是可加的。 （提示：先确定 $0 \in \operatorname{Hom}_{\mathcal{C}}(X, Y)$ ，然后利用乘积的存在性确定 $\operatorname{Hom}_{\mathcal{C}}(X, Y)$ 的加法。）

\statementlabel{练习 I.2}令 $ \mathcal{C} $  和 $ \mathcal{C}' $  为两个范畴，并令 $ F : \mathcal{C} \to \mathcal{C}' $  和 $ G : \mathcal{C}' \to \mathcal{C} $  为函子。

(i) 证明以下两个条件是等价的。

(a) 函子存在态射

\[
\alpha:F\circ G\to\operatorname{i d}_{\mathcal{C}^{\prime}},\qquad\beta:\operatorname{i d}_{\mathcal{C}}\to G\circ F
\]

这样，对于任何  $Y\in\mathrm{Ob}(\mathcal{C}')$  ，组合  $G(Y)\xrightarrow{\beta(G(Y))}G\circ F\circ G(Y)\xrightarrow{G(\alpha(Y))}G(Y)$  等于  $\mathrm{id}_{G(Y)}$  ，对于任何  $X\in\mathrm{Ob}(\mathcal{C})$  ，组合  $F(X)\xrightarrow{F(\beta(X))}F\circ G\circ F(X)\xrightarrow{\alpha(F(X))}F(X)$  等于  $\mathrm{id}_{F(X)}$  。

(b) 从  $ \mathcal{C}^\circ \times \mathcal{C}' $  到 Set:  $ \text{Hom}_{\mathcal{C}'}(F(X), Y) \simeq \text{Hom}_{\mathcal{C}}(X, G(Y)) $  存在双函子的同构。

在这种情况下，我们说  $G$  是  $F$  的右伴随函子， $F$  是  $G$  的左伴随函子。

(ii) 证明，对于函子 $F: \mathcal{C} \to \mathcal{C}'$ （或 $G: \mathcal{C}' \to \mathcal{C}$ ），其右（或左）伴随函子（如果存在）在同构上是唯一的。

(iii) 证明函子  $F: \mathcal{C} \to \mathcal{C}'$  （分别为  $G: \mathcal{C}' \to \mathcal{C}$  ）具有右（分别为左）伴随函子当且仅当对于任何  $Y \in \mathrm{Ob}(\mathcal{C}')$  ，  $X \mapsto \mathrm{Hom}_{\mathcal{C}'}(F(X), Y)$  可表示，（分别对于任何  $X \in \mathrm{Ob}(\mathcal{C})$  ，  $Y \mapsto \mathrm{Hom}_{\mathcal{C}}(X, G(Y))$  可表示）。

\statementlabel{练习 I.3}在定义 1.2.1 的条件 (i)、(ii)、(iii) 下，证明  $Z$  是函子  $W \mapsto \mathrm{Hom}_{\mathcal{C}}(X, W) \oplus \mathrm{Hom}_{\mathcal{C}}(Y, W)$  的代表，当且仅当存在态射  $i_1: X \to Z$  、  $i_2: Y \to Z$  、  $p_1: Z \to X$  、  $p_2: Z \to Y$  ，使得  $p_2 \circ i_1 = 0$  、  $p_1 \circ i_2 = 0$  、  $p_1 \circ i_1 = \mathrm{id}_X$ 、 $p_2 \circ i_2 = \mathrm{id}_Y$  和  $i_1 \circ p_1 + i_2 \circ p_2 = \mathrm{id}_Z$  。

\statementlabel{练习 I.4}令  $ X \xrightarrow{i_1} Z \xrightarrow{p_2} Y $  为加法范畴  $ \mathcal{C} $  中的态射序列，以及  $ p_2 \circ i_1 = 0 $  。证明下列条件是等价的。

(i) 对于任何  $ W \in \mathrm{Ob}(\mathcal{C}) $  ，序列：

\[
0\to\operatorname{H o m}_{\mathcal{C}}(W,X)\to\operatorname{H o m}_{\mathcal{C}}(W,Z)\to\operatorname{H o m}_{\mathcal{C}}(W,Y)\to0
\]

是准确的。

(ii) 对于任何  $ W \in \mathrm{Ob}(\mathcal{C}) $  ，序列：

\[
0\leftarrow\operatorname{H o m}_{\mathfrak{s}}(X,W)\leftarrow\operatorname{H o m}_{\mathfrak{s}}(Z,W)\leftarrow\operatorname{H o m}_{\mathfrak{s}}(Y,W)\leftarrow0
\]

是准确的。

(iii) 可以找到满足练习 I.3 条件的 $ i_2: Y \to Z $  和 $ p_1: Z \to X $ 。

如果满足这些条件，则称序列 $ 0 \to X \to Z \to Y \to 0 $  分裂。在这种情况下，  $ Z $  与  $ X \oplus Y $  同构，并且可以说  $ X $  是  $ Z $  的直接被加数。

(iv) 证明如果  $\mathcal{C}$  是阿贝尔范畴或三角范畴，如果存在  $i_1: X \to Z$  和  $p_1: Z \to X$  使得  $p_1 \circ i_1 = \mathrm{id}_X$  ，则  $X$  是  $Z$  的直加数。

\statementlabel{练习 I.5}令 $ 0 \to X \to Z \to Y \to 0 $  为阿贝尔范畴中的精确序列。如果 X 是单射的，或者 Y 是射影的，则序列会分裂。

\statementlabel{练习 I.6}令 \& 为阿贝尔范畴。

(i) 令  $f: X \to Z$  和  $g: Y \to Z$  为  $\mathcal{C}$  中的两个态射。证明  $\mathrm{Ker}(X \oplus Y \to Z)$  代表函子

\[
W\mapsto\operatorname{H o m}_{\mathcal{C}}(W,X)\mathop{\times}_{\operatorname{H o m}_{\mathcal{C}}(W,Z)}\operatorname{H o m}_{\mathcal{C}}(W,Y)\enspace.
\]

该函子的代表表示为  $ X \times_Z Y $  。

类似地，通过反转箭头，可以为两个态射  $f: Z \to X$  和  $g: Z \to Y$  定义  $X \oplus_Z Y$  ，作为函子  $W \to \mathrm{Hom}_{\mathcal{C}}(X, W) \times_{\mathrm{Hom}_{\mathcal{C}}(Z, W)} \mathrm{Hom}_{\mathcal{C}}(Y, W)$  的代表（即 Coker  $(Z \to X \oplus Y)$  ）。

(ii) 表明，在 (i) 的情况下，我们有一个交换图：

\begin{center}
\includegraphics[width=0.50\textwidth]{流行上的层_Chn-assets/img_in_image_box_300_1215_897_1664.png}
\end{center}

其中  $ f' $  和  $ g' $  以明显的方式定义。

（三）让

\begin{center}
\includegraphics[width=0.16\textwidth]{流行上的层_Chn-assets/img_in_image_box_515_200_712_386.png}
\end{center}

是一个交换图。那么下面的条件是等价的：

(a)  $ X' \to X \times_Y Y' $  是一个外态，

(b)  $ X \oplus_{X'} Y' \to Y $  是单态，

(c) 图中：

\begin{center}
\includegraphics[width=0.82\textwidth]{流行上的层_Chn-assets/img_in_image_box_96_559_1076_1179.png}
\end{center}

所有序列都是准确的。

(iv) 设  $ f: X \to Y $  为  $ \mathbf{C}(\mathcal{C}) $  中的态射。假设对于每个  $ n $  ，平方

\begin{center}
\includegraphics[width=0.24\textwidth]{流行上的层_Chn-assets/img_in_image_box_465_1284_752_1460.png}
\end{center}

满足(iii)中的等效条件。那么 $f$  是拟同构。 （提示：使用图表：

\begin{center}
\includegraphics[width=0.70\textwidth]{流行上的层_Chn-assets/img_in_image_box_180_1562_1022_1741.png}
\end{center}

\statementlabel{练习 I.7}令 \& 为阿贝尔范畴。

(i) 对于  $\mathcal{C}$  的对象  $Z$  ，设  $\mathcal{P}(Z)$  为其对象为表态  $f: Z' \to Z$  的范畴，态射  $(f: Z' \to Z) \to (f': Z'' \to Z)$  由  $h: Z' \to Z''$  和  $f' \circ h = f$  定义。证明 $\mathcal{P}(Z)$  是助滤剂，即 $\mathcal{P}(Z)^\circ$  是滤剂。

(ii) 对于 $ X \in \mathrm{Ob}(\mathcal{C}) $ ，设置 $ \tilde{h}_Z(X) = \lim_{\overrightarrow{\mathcal{P}}(\tilde{Z})} \mathrm{Hom}_{\mathcal{C}}(Z', X) $ ，其中 $ \mathrm{Hom}_{\mathcal{C}}(Z', X) $  是从 $ \mathcal{P}(Z) $  到 $ \mathfrak{U}\mathfrak{b} $  的函子，其将 $ \mathcal{P}(Z) $  中的 $ \mathrm{Hom}_{\mathcal{C}}(Z', X) $  关联到 $ f: Z' \to Z $ 。

证明  $k$  在  $\operatorname{Hom}_{c}(X, X')$  中满足  $...$  。还要证明  $C$  中的序列与  $k$  和  $Z \in \operatorname{Ob}(C)$  完全相同。

请注意，这个练习将使我们能够将阿贝尔范畴上的许多问题转化为 Ab 上的问题。

\statementlabel{练习 I.8（五个引理）}令  $ \mathcal{C} $  为阿贝尔范畴，并考虑  $ \mathcal{C} $  中具有精确行的交换图：

\begin{center}
\includegraphics[width=0.54\textwidth]{流行上的层_Chn-assets/img_in_image_box_261_782_910_958.png}
\end{center}

\statementlabel{证明}

(i) 如果 $f_{0}$  是同态， $f_{1}$  和 $f_{3}$  是单态，那么 $f_{2}$  是单态，

(ii) 如果 $f_{4}$  是单态， $f_{1}$  和 $f_{3}$  是外同态，则 $f_{2}$  是外同态。

\statementlabel{练习 I.9}令 $ \mathcal{C} $  为阿贝尔范畴。考虑  $ \mathcal{C} $  中具有精确行的交换图：

\begin{center}
\includegraphics[width=0.48\textwidth]{流行上的层_Chn-assets/img_in_image_box_289_1291_868_1478.png}
\end{center}

证明存在自然精确序列：

\[
\mathrm{Ker}\alpha\longrightarrow\mathrm{Ker}\beta\longrightarrow\mathrm{Ker}\gamma\xrightarrow{\varphi}\mathrm{Coker}\alpha\longrightarrow\mathrm{Coker}\beta\longrightarrow\mathrm{Coker}\gamma,
\]

以便下图可以通勤：

\begin{center}
\includegraphics[width=0.39\textwidth]{流行上的层_Chn-assets/img_in_image_box_356_158_828_479.png}
\end{center}

\statementlabel{练习 I.10}令 \& 为阿贝尔范畴。考虑精确序列图：

\begin{center}
\includegraphics[width=0.50\textwidth]{流行上的层_Chn-assets/img_in_image_box_293_607_890_784.png}
\end{center}

并假设 $ M_0 $  和 $ M'_0 $  是单射的。构造一个同构  $ M_0 \oplus M'_1 \simeq M'_0 \oplus M_1 $  。

\statementlabel{练习 I.11}令 $\mathcal{C}$  为阿贝尔范畴，并让 $X \in \mathrm{Ob}(\mathbb{C}(\mathcal{C}))$  为这样，对于任何 $Y \in \mathrm{Ob}(\mathcal{C})$  ，阿贝尔群 $\mathrm{Hom}_{\mathcal{C}}(Y, X)$  的复形是精确的。证明  $X$  中的  $\mathbb{K}(\mathcal{C})$  为 0（参见练习 I.4）。

\statementlabel{练习 I.12}令  $ \mathcal{C} $  为三角范畴，并考虑  $ \mathcal{C} $  中的交换图：

\begin{center}
\includegraphics[width=0.33\textwidth]{流行上的层_Chn-assets/img_in_image_box_383_1165_785_1331.png}
\end{center}

其中第一行是一个显着的三角形和  $ f[1] \circ h' = 0 $  。在以下假设之一下证明第二行也是一个杰出的三角形：

(i) 对于任何  $ P \in \mathrm{Ob}(\mathcal{C}) $  ，序列  $ \mathrm{Hom}(P, X) \to \mathrm{Hom}(P, Y) \to \mathrm{Hom}(P, Z') \to \mathrm{Hom}(P, X[1]) $  是精确的，

(ii) 对于任何  $ Q \in \mathrm{Ob}(\mathcal{C}) $  ，序列  $ \mathrm{Hom}(X[1], Q) \to \mathrm{Hom}(Z', Q) \to \mathrm{Hom}(Y, Q) \to \mathrm{Hom}(X, Q) $  是精确的。

\statementlabel{练习 I.13}令 $ X_i \longrightarrow Y_i \longrightarrow Z_i \xrightarrow{+1} (i = 1,2) $  为三角范畴中的两个三角形。证明它们是可区分三角形当且仅当它们的直和 $ X_1 \oplus X_2 \longrightarrow Y_1 \oplus Y_2 \longrightarrow Z_1 \oplus Z_2 \xrightarrow{+1} $  是可区分三角形。 （提示：

使用练习 I.12。为了证明它是足够的，考虑一个杰出的三角形 $ \delta: X_1 \longrightarrow Y_1 \longrightarrow U \xrightarrow{+1} \text{and prove that the composition } \delta_1 \to \delta_1 \oplus \delta_2 \to \delta $ 是一个同构，其中 $ \delta_1 $ 和 $ \delta_2 $ 是给定的三角形。）

\statementlabel{练习 1.14}令 $\mathcal{C}$  为范畴， $S$  为乘法系统。对于 $X \in \mathrm{Ob}(\mathcal{C})$ ，令 $S_X$  为范畴，其对象是态射 $s: X' \to X$  和 $s \in S$ ，态射定义如下。对于 $s: X' \to X$  和 $s': X'' \to X$  ，我们设置：

\[
\mathrm{H o m}_{S_{x}}(s,s^{\prime})=\left\{h\in\mathrm{H o m}_{\mathcal{C}}(X^{\prime\prime},X^{\prime});s^{\prime}=s\circ h\right\}\;.
\]

(i) 显示范畴 $(S_{X})^{\circ}$  是过滤剂。

(ii) 证明对于  $ X $  、 $ Y \in \text{Ob}(\mathcal{C}) $  ：

\[
\mathrm{H o m}_{\mathcal{C}_{S}}(X,Y)=\varinjlim_{S_{X}^{\circ}}\mathrm{H o m}_{\mathcal{C}}(X^{\prime},Y)\enspace.
\]

这里  $\mathrm{Hom}_{\mathcal{G}}(X', Y)$  是从  $S_X^\circ$  到  $\mathfrak{S}$  集的函子，它将  $\mathrm{Hom}_{\mathcal{G}}(X', Y)$  关联到  $s': X' \to X$  。

(iii) 通过反转箭头，定义范畴 $S_{Y}^{a}$ 并证明：

\[
\mathrm{H o m}_{\mathcal{C}_{S}}(X,Y)=\varinjlim_{S_{Y}^{a}}\mathrm{H o m}_{\mathcal{C}}(X,Y^{\prime})~.
\]

\statementlabel{练习 I.15}（参见 Deligne [1]。）让 $\mathcal{C}$  成为一个范畴。人们将范畴 $\mathrm{Ind}(\mathcal{C})$ 定义为 $\mathcal{C}^{\vee}$ 的满子范畴（参见§1），由与“ $\lim \overrightarrow{f}$ ” $F$ 同构的对象组成，用于 $\mathcal{C}$ 中的某些归纳系统 $F$ ，由过滤剂范畴 $\mathcal{I}$ 索引。

现在假设  $\mathcal{C}$  是阿贝尔的，对于  $X \in \mathrm{Ob}(\mathbf{K}^{+}(\mathcal{C}))$  来说， $S_X$  表示其对象由准同构  $u: X \to X',$  组成的范畴，态射  $h:(u:X \to X') \to (u': X \to X'')$  由  $v: X' \to X''$  和  $v \circ u = u'$  定义。

(i) 证明函子  $\sigma: \mathbf{D}^{+}(\mathcal{C}) \to \mathrm{Ind}(\mathbf{K}^{+}(\mathcal{C})), X \mapsto \varinjlim_{\overline{S}_{X}} X'$  定义明确且完全忠实。这里  $X'$  是从  $S_{X}$  到  $\mathbf{K}^{+}(\mathcal{C})$  的函子，它将  $X'$  关联到  $u: X \to X'$  。

(ii) 令 $F: \mathcal{C} \to \mathcal{C}'$  为阿贝尔范畴的左精确函子。通过  $T(X) = \lim_{\overrightarrow{S}_X} F(X')$  定义函子  $T: \mathbf{D}^+(\mathcal{C}) \to \mathrm{Ind}(\mathbf{K}^+(\mathcal{C}'))$  。假设如果存在  $Y \in \mathrm{Ob}(\mathbf{D}^+(\mathcal{C}'))$  使得  $T(X) \simeq \sigma(Y)$  ，则  $F$  可在  $X \in \mathrm{Ob}(\mathbf{D}^+(\mathcal{C}))$  处导出。证明在这种情况下  $Y$  是唯一的，并且如果  $F$  在每个  $X \in \mathrm{Ob}(\mathbf{D}^+(\mathcal{C}))$  处可导，则  $F$  承认右派生函子和  $\sigma \circ RF \simeq T$  。

\statementlabel{练习 I.16}设 C 为加法范畴。

(i) 对于 $ X \in \mathbf{C}^{-}(\mathcal{C}) $  和 $ Y \in \mathbf{C}^{+}(\mathcal{C}) $  证明：

\[
Z^{0}(\mathsf{s}(\operatorname{H o m}_{\mathcal{C}}(X,Y)))=\operatorname{H o m}_{\mathsf{C}(\mathcal{C})}(X,Y),
\]

\[
B^{0}(\mathrm{s}(\mathrm{Hom}_{\mathcal{C}}(X,Y)))=\mathrm{Ht}(X,Y),
\]

\[
H^{0}(\mathsf{s}(\mathsf{H o m}_{\mathcal{C}}(X,Y)))=\mathsf{H o m}_{\mathsf{K}(\mathcal{C})}(X,Y)\enspace.
\]

（s 在（1.9.4）中定义。）

(ii) 此外，假设  $\mathcal{C}$  是一个具有足够单射（或者足够射影）的阿贝尔范畴。证明对于 $X \in \mathrm{Ob}(\mathbf{D}^{-}(\mathcal{C}))$  和 $Y \in \mathrm{Ob}(\mathbf{D}^{+}(\mathcal{C}))$  有：

\[
H^{0}(R\mathrm{H o m}_{\mathcal{C}}(X,Y))=\mathrm{H o m}_{\mathsf{D}(\mathcal{C})}(X,Y).
\]

\statementlabel{练习 I.17}令 $\mathcal{C}$  为阿贝尔范畴。有人说  $\mathcal{C}$  具有同源维度  $\leq n$  ，其中  $n$  是非负整数，如果  $\text{Ext}^{j}(X, Y) = 0$  对于  $j > n$  和任何  $X, Y \in \text{Ob}(\mathcal{C})$  。这里 $\text{Ext}^{j}(X, Y) = \text{Hom}_{\mathcal{D}(\mathcal{C})}(X, Y[j])$ 。假设 $\mathcal{C}$  有足够的单射。证明下列条件是等价的。

(i)  $ \mathcal{C} $  具有同源维度  $ \leq n $  。

(ii) 对于任何  $ X \in \mathrm{Ob}(\mathcal{C}) $  ，存在长度为  $ \leq n $  的  $ X $  的单射解析，即精确序列：  $ 0 \to X \to I^0 \to \cdots \to I^n \to 0 $  与所有  $ I^j $  的单射。

满足这些等价条件的最小 $ n \in \mathbb{N} \cup \{\infty\} $ 称为 $ \mathcal{C} $ 的同调维数，并表示为 $ \mathrm{hd}(\mathcal{C}) $ 。

\statementlabel{练习 I.18}令 $\mathcal{C}$  成为带有 $\mathrm{hd}(\mathcal{C}) \leq 1$  的阿贝尔范畴。证明对于任何  $X \in \mathrm{Ob}(\mathbf{D}^{b}(\mathcal{C}))$  ，在  $\mathbf{D}^{b}(\mathcal{C})$  中具有同构：

\[
X\simeq\bigoplus_{k}H^{k}(X)[-k]~.
\]

（示例： $\mathcal{C} = \mathfrak{Mod}(A)$ ，其中  $A$  是主要理想域。特别是当  $A$  是一个字段，或者  $A = \mathbb{Z}$  时。）

\statementlabel{练习 I.19}令 $\mathcal{C}$  和 $\mathcal{C}'$  为两个阿贝尔范畴， $F: \mathcal{C} \to \mathcal{C}'$  为左精确函子，并令 $\mathcal{I}$  为 $\mathcal{C}$  的 $F$  内射子范畴。我们将  $\mathcal{C}$  和  $F$  的对象称为  $X$  -非循环 if  $R^k F(X) = 0$  for  $k \neq 0$  。令  $\mathcal{J}$  为  $\mathcal{C}$  的满子范畴，由  $F$  -非循环对象组成。

(i) 证明 $\mathcal{J}$  是 $F$  -内射。

(ii) 证明对于整数  $ n \geq 0 $  ，以下条件是等价的：

(a)  $ R^k F(X) = 0 $  对于任何  $ k > n $  和任何  $ X \in \mathrm{Ob}(\mathcal{C}) $  ，

(b) 对于任何  $ X \in \mathrm{Ob}(\mathcal{C}) $  ，都存在精确序列

\[
0\to X\to X^{0}\to\cdots\to X^{n}\to0
\]

用 $ X^j \in \mathrm{Ob}(\mathcal{J}) $  表示 $ 0 \leq j \leq n $  ，

(c) 如果  $ X^0 \to \cdots \to X^n \to 0 $  是精确序列，并且如果  $ X^j \in \mathrm{Ob}(\mathcal{J}) $  对应  $ j < n $  ，则  $ X^n $  属于  $ \mathcal{J} $  。

在这种情况下，有人说  $F$  具有上同调维度  $\leq n$  。

\statementlabel{练习 I.20}在命题1.8.7的情况下，假设 $F$ （分别 $F'$ ）具有上同调维数 $\leq r$ （分别 $\leq r'$ ），（参见练习I.19）。证明  $F' \circ F$  具有上同调维数  $\leq r + r'$  。

\statementlabel{练习 I.21}令 $\mathcal{C}$  和 $\mathcal{C}'$  为两个阿贝尔范畴， $F:\mathcal{C} \to \mathcal{C}'$  为左精确函子。假设存在  $\mathcal{C}$  的  $F$  内射子范畴  $\mathcal{I}$  。让

 $ X \in \mathrm{Ob}(\mathbb{D}^+\mathcal{C}) $  使得 $ R^i F(H^j(X)) = 0 $  对于所有 $ i > 0 $  、所有 $ j \leqslant j_\circ $  。证明同构：

\[
R^{j} F(X)\simeq F(H^{j}(X))\quad\mathrm{f o r~a l l}\quad j\leqslant j_{\circ}~.
\]

\statementlabel{练习 I.22}在命题1.8.7的情况下，设 $ X \in \mathrm{Ob}(\mathbb{D}^{+}(\mathcal{C})) $ ，并假设 $ R^j F(X) = 0 $ 为 $ j < n $ 。证明：

\[
R^{n}(F^{\prime}\circ F)(X)\cong F^{\prime}\circ R^{n}F(X)
\]

（提示：应用备注 1.8.6。）

\statementlabel{练习 I.23}令 $\mathcal{C}$  为阿贝尔范畴， $\mathcal{I}$  为满子范畴。假设 $\mathcal{I}$ 满足条件(1.7.5)、(1.7.6)并且：

如果  $0 \to X' \to X \to X'' \to 0$  是  $\mathcal{C}$  中的精确序列，并且如果  $X'$  属于  $\mathcal{I}$  ，则  $X''$  属于  $\mathcal{I}$  当且仅当  $X$  属于  $\mathcal{I}$  。

让 $ ^{\star} = \emptyset $  或 $ b $  或- 或+。

(a) 证明任何  $ X \in \text{Ob}(\mathbb{C}^*(\mathcal{C})) $  与  $ \mathbb{C}^*(\mathcal{I}) $  的对象  $ Y $  准同构。

(b) 令 $ \mathcal{C}' $  为另一个阿贝尔范畴， $ F : \mathcal{C} \to \mathcal{C}' $  为左精确函子。假设范畴 $ \mathcal{I} $  是 $ F $  -内射。证明  $ RF $  作为从  $ \mathbf{D}^*(\mathcal{C}) $  到  $ \mathbf{D}^*(\mathcal{C}') $  的函子存在。

(c) 令 $\mathcal{C}$  为第三阿贝尔范畴， $G: \mathcal{C} \times \mathcal{C}' \to \mathcal{C}$  为左精确双函子。假设对于每个  $X' \in \mathrm{Ob}(\mathcal{C}')$  ，范畴  $\mathcal{I}$  是  $G(\cdot, X')$  -内射。证明  $RG$  作为从  $\mathbf{D}^{-}(\mathcal{C}) \times \mathbf{D}^{-}(\mathcal{C}')$  到  $\mathbf{D}^{-}(\mathcal{C}')$  和  $\mathbf{D}^*(\mathcal{C}) \times \mathbf{D}^b(\mathcal{C}')$  到  $\mathbf{D}^*(\mathcal{C}')$  的双函子存在，（参见 Hartshorne [1, p. 42]）。

\statementlabel{练习 I.24}(i) 令 $F: \mathcal{C} \to \mathcal{C}'$  为阿贝尔范畴的左精确函子。让 $X \in \mathrm{Ob}(\mathcal{D}^+)(\mathcal{C})$  。构造自然态射： $H^j(RF(X)) \to F(H^j(X))$  。

(ii) 令 $\mathcal{C}, \mathcal{C}^{\prime}, \mathcal{C}^{\prime\prime}$  为阿贝尔范畴，并令 $F$  为从 $\mathcal{C} \times \mathcal{C}^{\prime}$  到 $\mathcal{C}^{\prime\prime}$  的双函子。让  $X \in \mathrm{Ob}(\mathbb{D}^*(\mathcal{C}))$  、  $Y \in \mathrm{Ob}(\mathbb{D}^*(\mathcal{C}^{\prime}))$  和  $* = + \mathrm{or} -$  。

(a) 假设  $F$  为左精确， $* = +$  （分别为  $F$  为右精确， $* = -$  ）。构造  $p$  、  $q \in \mathbb{Z}$  的自然态射：

\[
H^{p+q}(R F(X,Y))\to F(H^{p}(X),H^{q}(Y))
\]

（分别 $ F(H^{p}(X), H^{q}(Y)) \to H^{p+q}(LF(X, Y)) $ ）。）

(b) 假设 $F$  是准确的。证明  $n \in \mathbb{Z}$  的同构：

\[
H^{n}(F(X,Y))\simeq\bigoplus_{p+q=n}F(H^{p}(X),H^{q}(Y))\enspace.
\]

\statementlabel{练习 I.25}设C是阿贝尔范畴，X是C中的双复形，满足条件(1.9.2)。

(i) 证明以下三角形在  $ \mathbf{D}(\mathcal{C}) $  中是可区分的：

\[
\begin{aligned}{}&{{}\mathsf{s}\tau_{I I}^{\leqslant n-1}(X)\longrightarrow\mathsf{s}\tau_{I I}^{\leqslant n}(X)\longrightarrow H_{I I}^{n}(X)[-n]\stackrel{\longrightarrow}{+1},}\\ {}&{{}H_{I I}^{n}(X)[-n]\longrightarrow\mathsf{s}\tau_{I I}^{\geqslant n}(X)\longrightarrow\mathsf{s}\tau_{I I}^{\geqslant n+1}(X)\stackrel{\longrightarrow}{+1}.}\\ \end{aligned}
\]

(ii) 固定 $ k \in \mathbb{Z} $ 。证明自然态射：

\[
H^{k}(\mathsf{s}\tau_{I I}^{\leq n}(X))\to H^{k}(\mathsf{s}(X))
\]

（分别为  $ H^k(\mathfrak{s}(X)) \to H^k(\mathfrak{s}\tau_{II}^{\geq n}(X)) $  ），是  $ n \gg 0 $  的同构，（分别为  $ n \ll 0 $  ）。

(iii) 整数 $ k \in \mathbb{Z} $ 固定，证明：

\[
H^{k}(\mathfrak{s t}_{I I}^{\leq n}(X))=0\qquad\mathrm{f o r}\quad n\ll0
\]

and

\[
H^{k}({\sf s r}_{{I I}}^{\geq n}(X))=0\qquad{\mathrm{f o r}}\quad n\gg0~.
\]

（提示： $ H^k(s(X)) $  仅取决于  $ \{X^{n,m}; k-1 \leq n + m \leq k+1\} $  。）

\statementlabel{练习 I.26}（J.-P.施奈德斯）。在练习 I.25 的情况下，假设：

\[
H_{I I}^{q}(X)\mathop{\simeq}_{q i s}0\qquad{\mathrm{e x c e p t~f o r}}\quad q=q_{0},q_{1}
\]

其中 $ q_{0} < q_{1} $  。证明我们有一个独特的三角形：

\[
H_{I I}^{q_{0}}(X)[-q_{0}]\longrightarrow s(X)\longrightarrow H_{I I}^{q_{1}}(X)[-q_{1}]\stackrel{+1}{\longrightarrow}
\]

（提示：使用练习 I.25。）

\statementlabel{练习 I.27}令  $\mathcal{C}$  为阿贝尔（或三角）范畴。 1 表示  $\operatorname{K}(\mathcal{C})$  作为由  $\mathcal{C}$  的对象通过关系  $X = X' + X''$  生成的自由阿贝尔群的商而获得的阿贝尔群，如果  $\mathcal{C}$  中存在精确序列  $0 \to X' \to X \to X'' \to 0$  （分别是  $\mathcal{C}$  中的杰出三角形  $X' \longrightarrow X \longrightarrow X'' \longrightarrow _{+1}$  ）。有人称  $\operatorname{K}(\mathcal{C})$  为  $\mathcal{C}$  的 Grothendieck 小组。

现在让 $\mathcal{C}$  成为阿贝尔范畴。证明函子  $i: \mathcal{C} \to \mathbf{D}^b(\mathcal{C})$  、  $X \mapsto X$  诱导群同构  $\mathbf{K}(\mathcal{C}) \simeq \mathbf{K}(\mathbf{D}^b(\mathcal{C}))$  ，其逆由  $X \mapsto \sum_j (-1)^j [H^j(X)]$  给出，其中  $[Z]$  是  $\mathbf{K}(\mathcal{C})$  中  $Z$  的类。

\statementlabel{练习 I.28}设A为环。证明下列条件(i)-(iii)是等价的。

(i) Mod(A) 具有同调维度  $ \leq n $  。

(ii) 任何左模 M 都具有长度  $ \leq n $  的单射分辨率。

(iii) 任何左模 M 的投影分辨率长度为  $ \leq n $  ，

（即存在一个精确的序列  $ 0 \to P_n \to \cdots \to P_0 \to M \to 0 $  ，所有  $ P_j $  都是投影）。

一套 $\mathrm{gld}(A)=\sup(\mathrm{hd}(\mathfrak{Mod}(A),\mathrm{hd}(\mathfrak{Mod}(A^{op}))\mathrm{and}\mathrm{calls}\mathrm{gld}(A)\mathrm{the}\mathrm{global}\mathrm{homological}\mathrm{dimension}\mathrm{of}A.$ 

\statementlabel{练习 I.29}设A为环。

(i) 证明自由模是射影的。

(ii) 证明射影模是自由模的直加数。

(iii) 如果函子 $ \cdot \otimes_A M $  是精确的，则 $ A $  模 $ M $  被称为\textbackslash{}textit\{flat\}。证明射影模是平坦的。

(iv) 令 n 为非负整数。证明下列条件是等价的。

(a)  $ \mathrm{Tor}_{j}^{A}(N,M)=0 $  对于任何 j>n 以及对于任何右模 N 和左模 M。

(b) 任何左模  $M$  都具有长度为  $\leq n$  的平坦分辨率（即存在一个精确序列  $0 \to P^n \to \cdots \to P^0 \to M \to 0$  且所有  $P^{j_1}$  都是平坦的）。

(b)  $ ^{op} $  ：与 (b) 相同，将左侧替换为右侧。

将  $A$  的弱全局维度 wgld  $(A)$  定义为满足这些条件的最小  $n\in\mathbb{N}\cup{+\infty}$ 。

(v) 证明 $ \mathrm{wgld}(A) \leq \mathrm{gld}(A) $  。

\statementlabel{练习 I.30}令 $A$  为交换环。如果 $\mathbf{D}^{b}(\mathfrak{M}\mathbf{o b}(A))$  的对象 $X$  与有限生成的射影 $A$  模的有界复形同构，则称该对象是完美的。

(i) 证明如果 $X \longrightarrow Y \longrightarrow Z \xrightarrow[+1]{}$  是 $\mathbf{D}^{b}(\mathfrak{M}\partial(A))$  中的一个显着三角形，并且 $X$  、 $Y$  是完美的，那么 $Z$  也是完美的。

(ii) 证明完美对象的直接加法是完美的。 （提示：对于射影模  $ P' $  的有界复形以及  $ j > 0 $  的准同构  $ P' \to X' \oplus Y' $  与  $ P^j = X^j = Y^j = 0 $  ，让  $ \widetilde{P} $  、  $ \widetilde{X} $  、  $ \widetilde{Y} $  表示通过替换  $ P^0 $  、  $ P^{-1} $  、  $ X^{-1} $  获得的复形，  $ Y^{-1} $  和  $ 0 $  、  $ P^{-1} \oplus P^0 $  、  $ X^{-1} \oplus P^0 $  、  $ Y^{-1} \oplus P^0 $  分别，然后使用态射  $ P^0 \xrightarrow{h} \widetilde{X}^{-1} \oplus \widetilde{Y}^{-1} $  构造准同构  $ \widetilde{P} \to \widetilde{X} \oplus \widetilde{Y} $  ，使得  $ P^0 \to \widetilde{X}^{-1} \oplus \widetilde{Y}^{-1} \oplus X^0 \oplus Y^0 \oplus P^0 $  是  $ (0, 0, id_{p^0}) $  。

(iii) 令 $ M \in \text{Ob}(\mathbb{D}^b(\mathfrak{Mod}(A))) $  并假设 $ M $  是完美的。让 $ M^* = \text{RHom}(M, A) $  。证明  $ M^* $  是完美的，并且规范态射  $ M \to M^{**} $  是同构。

现在假设  $A$  是 Noetherian 且  $\mathrm{gld}(A) < \infty$  。

(iv) 让 $\mathfrak{Mod}^{f}(A)$  表示有限生成的 $A$  模的阿贝尔范畴。证明 $\mathbf{D}^{b}(\mathfrak{Mod}^{f}(A))$  的任何对象都是完美的。

(v) 用  $ \mathbf{D}_f^v(\mathfrak{Mod}(A)) $  表示  $ \mathbf{D}^v(\mathfrak{Mod}(A)) $  的完整子类，由上同调群属于  $ \mathfrak{Mod}^f(A) $  的对象组成。证明自然函子  $ \mathbf{D}^b(\mathfrak{Mod}^f(A)) \to \mathbf{D}_f^b(\mathfrak{Mod}(A)) $  是等价的。 （提示：使用命题 1.7.11）（参见 [SGA6] Exposé I.）

\statementlabel{练习 I.31}(i) 让 $ M \in \text{Ob}(\mathbb{D}^b(\mathfrak{W}_\text{od}(\mathbb{Z}))) $  并让 $ M^* = \text{RHom}(M, \mathbb{Z}) $  。证明  $ M^* = 0 $  蕴含  $ M = 0 $  。 （提示：可以假设  $ H^k(M) = 0 $  为  $ k > 0 $  ）。使用区分三角形 $ \tau^{\leq -1}(M) \longrightarrow M \longrightarrow H^0(M) \xrightarrow{+1} $ ，证明简化为以下内容：

“让  $ M $  成为  $ \mathbb{Z} $  模，使得  $ \text{Hom}(M, \mathbb{Z}) = \text{Ext}^1(M, \mathbb{Z}) = 0 $  。然后  $ M = 0 $  ”。

为了证明这个结果，需要证明这样的  $M$  是无扭转的，然后证明  $M$  是可整除的（与  $n \in \mathbb{Z} \setminus \{0\}$  的乘法是满射的），然后证明  $M \simeq \mathbb{Q} \otimes_{\mathbb{Z}} M$  。如果  $M$  不为零，则可以写  $M = \mathbb{Q} \oplus L$  ，这将意味着  $\mathrm{Ext}^1(\mathbb{Q}, \mathbb{Z}) = 0$  ，这是不正确的。）

(ii) 证明如果  $ M^* \in \mathrm{Ob}(\mathbb{D}^b(\mathfrak{M}\mathrm{od}^f(\mathbb{Z}))) $  则  $ M \in \mathrm{Ob}(\mathbb{D}^b(\mathfrak{M}\mathrm{od}^f(\mathbb{Z}))) $  。

\statementlabel{练习 I.32}令 $k$  为交换域， $X \in \text{Ob}(\mathbf{D}^b(\mathfrak{Mod}(k)))$  。一套  $X^* = \text{RHom}(X, k)$  。 （请记住，注释 1.8.11 中给出了差异。）

(i) 假设  $X$  属于  $\mathbf{D}_{f}^{b}(\mathfrak{M o b}(k))$  ，证明自然同构：

\[
X\simeq X^{**}\quad,\qquad X^{*}\otimes X\simeq R\mathrm{H o m}(X,X)\quad,
\]

并将态射  $ X^* \otimes X \to k $  构造为态射  $ (X^n)^* \otimes X^n \to k $  的直积。

(ii) 令 $ X \in \text{Ob}(\mathbb{D}_f^b(\mathfrak{M}_\text{od}(k))) $  并让 $ v \in \text{Hom}(X, X) $  。一套：

\[
\mathrm{t r}(v)=\sum_{j}(-1)^{j}\mathrm{t r}(H^{j}(v))\quad,
\]

其中  $ \operatorname{tr}(H^j(v)) $  是自同态  $ H^j(v): H^j(X) \to H^j(X) $  的迹。现在让  $ Y \in \mathrm{Ob}(\mathbb{K}^b(\mathfrak{Mod}^f(k))) $  和  $ v \in \mathrm{Hom}(Y, Y) $  。证明：

\[
\mathrm{t r}(v)=\sum_{j}(-1)^{j}\mathrm{t r}(v^{j})\quad.
\]

(iii) 考虑  $ \mathbf{D}_{f}^{b}(\mathfrak{Mod}(k)) $  中不同三角形的自同态：

\[
\begin{array}{c}X^{\prime}\longrightarrow X\longrightarrow X^{\prime\prime}\xrightarrow[+1]{}\\\left|\vphantom{\begin{array}{c}X^{\prime}\\ X^{\prime}\end{array}}v^{\prime}\right.\quad\left|\vphantom{\begin{array}{c}X^{\prime}\\ X^{\prime}\end{array}}v\quad\vphantom{\begin{array}{c}X^{\prime}\\ X^{\prime}\end{array}}\right.\\\quad\left|\vphantom{\begin{array}{c}X^{\prime}\\ X^{\prime}\end{array}}X\right.\quad\left|\vphantom{\begin{array}{c}X^{\prime}\\ X^{\prime}\end{array}}X\quad\vphantom{\begin{array}{c}X^{\prime}\\ X^{\prime}\end{array}}\right.\\\longrightarrow X\longrightarrow X^{\prime\prime}\xrightarrow[+1]{}.\\\end{array}
\]

证明 $ \mathrm{tr}(v) = \mathrm{tr}(v') + \mathrm{tr}(v'' $ 。

(iv) 在(ii)的情况下，通过态射证明 $ \operatorname{tr}(v) $ 是v的像：

\[
H^{0}(R\mathrm{H o m}(X,X))\simeq H^{0}(X^{*}\otimes X)\to k.
\]

通常一组，对于  $ X \in \text{Ob}(\mathbb{D}_f^b(\mathfrak{Mod}(k))) $  ：

\[
\chi(X)=\sum_{j}(-1)^{j}\mathrm{d i m}H^{j}(X)\enspace.
\]

当然， $ \chi(X) = \mathrm{tr}(\mathrm{id}_X) $  位于  $ k $  中（参见 [SGA 6] Exposé I）。

\statementlabel{练习 I.33}令 $k$  为交换域， $V$  为 $k$  向量空间， $u: V \to V$  为自同态。如果  $\dim(u^n(V)) < \infty$  对于某些  $n > 0$  ，我们说  $u$  位于跟踪类中，在这种情况下我们设置：

\[
\mathrm{t r}(u)=\mathrm{t r}(u|_{u^{n}(V)}).
\]

(i) 证明 $ \operatorname{tr}(u) $  的定义与n 的选择无关。

(ii) 令 $ V \to W \to v $  为k 向量空间的态射。证明  $ u \circ v $  位于跟踪类中当且仅当  $ v \circ u $  位于跟踪类中，并且在本例中为  $ \mathrm{tr}(u \circ v) = \mathrm{tr}(v \circ u) $  。

(iii) 考虑下面的交换图，其中的行是精确的：

\begin{center}
\includegraphics[width=0.48\textwidth]{流行上的层_Chn-assets/img_in_image_box_310_1561_889_1751.png}
\end{center}

当且仅当  $u'$  和  $u''$  位于跟踪类中时，证明  $u$  位于跟踪类中，并在这种情况下证明  $\mathrm{tr}(u) = \mathrm{tr}(u') + \mathrm{tr}(u'')$  。

\statementlabel{练习 I.34}令 $ k $  为交换域， $ X \in \text{Ob}(\mathbb{D}_f^b(\mathfrak{Mod}(k))) $  。一套：

\[
\mathsf{b}_{i}(X)=\mathsf{d i m}\:H^{i}(X)\;,\qquad\mathsf{b}_{i}^{*}(X)=(-1)^{i}\sum_{j\leqslant i}(-1)^{j}\mathsf{b}_{j}(X)\;.
\]

令  $ X' \longrightarrow X \longrightarrow X''\xrightarrow[+1]{} $  为  $ \mathbf{D}_{f}^{b}(\mathfrak{M}\mathfrak{o}\mathfrak{b}(k)) $  中的一个显着三角形。证明：

\[
\chi(X)=\chi(X^{\prime})+\chi(X^{\prime\prime})~,
\]

\[
\mathbf{b}_{i}^{*}(X)\leqslant\mathbf{b}_{i}^{*}(X^{\prime})+\mathbf{b}_{i}^{*}(X^{\prime \prime})~.
\]

（ $ \chi(X) $  在练习 I.32 中定义）。

\statementlabel{练习 I.35}(i) 令 $I$  为有向集， $\{X_i\}$  为由 $I$  在范畴 $\mathcal{C}$  中索引的归纳系统。证明“ $\lim$ ” $X_i$  是 $\mathcal{C}^\vee$  中的归纳极限。更精确地证明对于任何 $F \in \mathrm{Ob}(\mathcal{C}^\vee)$ 

\[
\mathrm{H o m}_{\mathcal{C}^{\vee}}(\mathrm{“}\varprojlim\mathrm{”}X_{i},F)\simeq\varprojlim F(X_{i})\enspace.
\]

(ii) 如果  $\{Y_{j}\}$  是由  $\mathcal{C}$  中的有向集  $J$  索引的归纳系统，则证明

\[
\operatorname{H o m}_{\mathcal{G}^{\vee}}\left(\underset{}{\substack{\mathbf{\updownarrow}i}}{\operatorname*{l i m}}X_{i},\underset{}{\operatorname*{l i m}}{\overrightarrow{j}}Y_{j}\right)\simeq\underset{}{\operatorname{\lhd}i}{\operatorname*{l i m}}\underset{}{\operatorname{\lhd}j}{\operatorname{H o m}}_{\mathcal{G}}(X_{i},Y_{j})~.
\]

\statementlabel{练习 I.36}令 $A$  为诺特环，并令 $\mathfrak{Mod}^{f}(A)$  为有限生成的 $A$  模的范畴。令  $\{X_{i},\rho_{i,j}\}$  为此范畴的归纳系统，由有向有序集  $I$  索引。证明如果 $\lim\limits_{\substack{\vec{j}\\}}\,X_{j}$ 存在于 $\mathfrak{Mod}^{f}(A)$ 中，那么它代表“lim” $X_{i}$ 。

Mod\textasciicircum{}\{f\}(A)，则表示“ $ \lim_{x \to a} x $ ”。

\statementlabel{练习 I.37}令 $ \mathcal{C} $  为加法范畴。我们用  $ \text{End}(\mathcal{C}) $  表示函子  $ \text{id}_{\mathcal{C}}: \mathcal{C} \to \mathcal{C} $  的自同态集合。

(i) 证明 End(C) 是交换环。

(ii) 证明对于环  $ A $  ，  $ \text{End}(\mathfrak{Mod}(A)) $  与  $ A $  的中心同构（即  $ \{a \in A; ax = xa $  对于任何  $ x \in A\} $  ）。

(iii) 如果 $A$  是交换环，并且如果给出环同态 $A \to \mathrm{End}(\mathcal{C})$ ，我们称 $\mathcal{C}$  为 $A$  之上的加法范畴。证明 $\mathrm{Hom}_{\mathcal{C}}(X, Y)$  具有 $A$  -模的结构，并证明态射的复合是 $A$  -双线性。

(iv) 假设  $ \mathcal{C} $  是交换诺特环 A 上的阿贝尔范畴。

(a) 证明对于 $M \in \text{Ob}(\mathfrak{Mod}^f(A))$  和 $X \in \text{Ob}(\mathcal{C})$  ，函子 $Y \mapsto \text{Hom}_A(M, \text{Hom}_{\mathcal{C}}(X, Y))$  是可表示的。我们用  $X \otimes_A M$  表示其代表。

(b) 证明  $ \otimes_A $  是从  $ \mathcal{C} \times \mathfrak{Mod}^f(A) $  到  $ \mathcal{C} $  的右精确双函子。

(c) 证明上面的双函子有一个左派生函子  $ \otimes_A^L: \mathbf{D}^-(\mathcal{C}) \times \mathbf{D}^-(\mathfrak{Mod}^f(A)) \to \mathbf{D}^-(\mathcal{C}) $  。

(d) 类似地讨论函子  $ \mathrm{Hom}_{A}(\cdot,\cdot):\mathfrak{Mod}^{f}(A)^\circ\times\mathcal{C}\to\mathcal{C} $  （即  $ \mathcal{C}^\circ $  中的函子  $ \otimes_{A} $  ）。

\statementlabel{练习 I.38}令 $\mathcal{I}$  和 $\mathcal{I}'$  为过滤范畴， $\phi: \mathcal{I} \to \mathcal{I}'$  为函子。如果我们有：

(a) 对于任何  $ i' \in \mathrm{Ob}(\mathcal{I}') $  ，都存在  $ i \in \mathrm{Ob}(\mathcal{I}) $  和态射  $ i' \to \phi(i) $  。

(b) 对于  $ i \in \text{Ob}(\mathcal{I}) $  、 $ i' \in \text{Ob}(\mathcal{I}') $  和态射  $ f: \phi(i) \to i' $  ， $ \mathcal{I} $  中存在态射  $ g: i \to i_1 $  ，使得  $ \phi(g): \phi(i) \to \phi(i_1) $  通过  $ f $  因式分解。令  $ \mathcal{C} $  为范畴，让  $ F $  为由  $ \mathcal{I}' $  索引的归纳（或射影）系统。证明：

\[
\varinjlim_{F}F\circ\phi\to\varinjlim_{F}F\qquad\mathrm{~a n d~}\qquad\varinjlim_{F}F\circ\phi\to\varinjlim_{F}F
\]

（分别地，$\varprojlim_{\mathcal I'} F \to \varprojlim_{\mathcal I} F\circ\phi$ 和 $\varprojlim_{\mathcal I'} F \to \varprojlim_{\mathcal I} F\circ\phi$）是同构。

\statementlabel{练习 I.39}令 $\mathcal{C}$  为阿贝尔范畴。对于 $X$ ， $Y \in \mathrm{Ob}(\mathbb{D}^{b}(\mathcal{C}))$ ，我们设置 $\mathrm{Ext}^{j}(X, Y) = \mathrm{Hom}_{\mathbb{D}(\mathcal{C})}(X, Y[j]).$ 

(i) 设  $ X, Y \in \mathrm{Ob}(\mathcal{C}) $  和  $ n \geqslant 1 $  。对于精确的序列

\[
E:0\to Y\to Z_{n}\to Z_{n-1}\to\cdots\to Z_{1}\to X\to0
\]

定义一个元素 C(E) ∈ Ext*(X, Y)。人们将这样的精确序列称为 X 乘 Y 的 n 扩展。

(ii) 证明 Ext  $ ^{n} $  (X, Y) 的任何元素对于 X 的某些扩展 E 通过 Y 都可以写成 C(E)。

(iii) 让 $E':0\to Y\to Z'_n\to\cdots\to Z'_1\to X\to0$  成为另一个扩展。证明  $C(E)=C(E')$  当且仅当存在扩展  $E":0\to Y\to Z''_n\to\cdots\to Z''_1\to X\to0$  和交换图：

\begin{center}
\includegraphics[width=0.49\textwidth]{流行上的层_Chn-assets/img_in_image_box_308_1161_900_1476.png}
\end{center}

\begin{center}（通常 Ext  $ ^{n} $  (X, Y) 称为米田扩展）。\end{center}

\subsection{注记}
我们在本书的开头参考了 C. Houzel 的《简史》，了解上同调理论的详细历史。

这里我们只提一下，人们普遍承认导出范畴的思想是由 Grothendieck [4] 提出的，并由 Verdier 在他的论文中发展起来。不幸的是，除了短论文 Verdier [2] 之外，这篇论文从未发表过。

在这个复杂的理论出现之前，人们使用了“派生函子”（参见Cartan-Eilenberg [1]和Buchsbaum的附录），以及Leray [1,2]在1945年引入的“谱序列”的基本技术。Grothendieck [1]在东北的著名论文通过澄清和统一该理论在当时产生了巨大的影响。

本书中没有使用光谱序列。事实上，我们在这里不需要这种完全通用的微积分，定理 1.9.3 足以满足我们的目的。

最后，让我们回顾一下 Grothendieck [2] 中引入的ind-objects 和 pro-objects，以及 Grothendieck [3] 中引入的 Mittag-Leffler 条件。

\section{层}
\subsection{概要}
在本章中，我们在拓扑空间上构造了层的阿贝尔范畴，以及常见的关联函子，例如逆像 $ f^{-1} $ 、直像 $ f_* $ 、真直像 $ f_i $ 、张量积 $ \otimes $ 和内宏 $ \mathcal{Hom} $ 。利用第一章的结果，定义了层的导出范畴 $ \mathbf{D}^b(X) $ ，以及前面的派生函子。在本章中，我们还介绍了内射层、平层、松弛层、 $ c $  -软层的概念，并给出了我们稍后将使用的层理论工具：非特征变形引理和上同调的同伦不变性。尽管我们并不真正需要它，但我们（简短地）提出了 $ \operatorname{\tilde{Cech}} $ 上同调。我们通过回顾实流形或复流形上的一些自然层来结束本章。

我们在这里解释的大多数结果都是经典的，我们参考 Bredon [1]、Godement [1]、Iversen [1] 进行进一步的发展。

\subsection{预层}
令 $X$  为拓扑空间。我们用  $\mathrm{OP}(X)$  表示  $X$  的开放子集集合，按包含关系排序，并通过设置将其与范畴  $\mathcal{OP}(X)$  关联起来（参见 I §1）：

\[
\left\{\begin{aligned}{\operatorname{O b}(\mathfrak{D P}(X))}&{{}=\operatorname{O P}(X)}\\ {\operatorname{H o m}_{\mathfrak{D P}(X)}(V,U)}&{{}=\left\{\operatorname{p t}\right\}\quad\mathrm{i f}\;V\subset U~,}\\ {}&{{}=\emptyset\quad\mathrm{o t h e r w i s e}~.}\\ \end{aligned}\right.
\]

（这里 V 和 U 是 X 的开子集，\{pt\} 是包含一个元素的集合。）

\statementlabel{定义 2.1.1}让 $\mathcal{C}$  成为一个范畴。  $X$  上的 presheaf  $F$  的值在  $\mathcal{C}$  中，是从  $\mathcal{P}(X)$  到  $\mathcal{C}$  的逆变函子，并且 presheaves 的态射是此类函子的态射。

换句话说， $X$  上的预层是由  $\mathfrak{D}\mathfrak{P}(X)$  索引的射影系统（参见 I §11）。

在本书中，我们将只考虑交换群的预层（和层）。因此，除非另有说明，前面的定义中我们取 $ \mathcal{C} = \mathfrak{Ab} $ ，预层表示交换群的预层。

因此，presheaf 为每个开集  $U \subset X$  分配一个阿贝尔群  $F(U)$  ，并为每个开包含  $V \subset U$  分配一个群同态  $\rho_{V,U}: F(U) \to F(V)$  ，称为限制态射，条件为：

\[
\begin{cases}{\rho_{U,U}=\operatorname{i d}_{F(U)},\rho_{\mathcal{W},U}=\rho_{\mathcal{W},V}\circ\rho_{V,U}\mathrm{f o r}}\\ {\mathrm{a n y~t r i p l e t~}W\subset V\subset U\mathrm{~o f~o p e n~s u b s e t s}}\\ \end{cases}.
\]

预滑态射  $\phi: F \to G$  是群同态族  $\phi_U: F(U) \to G(U)$  ，与限制态射兼容。也就是说，如果用相同的字母  $\rho_{V,U}$  表示  $F$  和  $G$  的限制态射，则对于开集的所有对  $U$  、  $V$  和  $V \subset U$  ，下图可交换：

\[
\begin{array}{c}F(U)\xrightarrow{\quad\phi_{U}\quad}G(U)\\\rho_{V,U}\bigg\downarrow\quad\bigg\downarrow\rho_{V,U}\\F(V)\xrightarrow{\quad\phi_{V}\quad}G(V)\quad\end{array}
\]

一种通过 $ U \mapsto 0 $  定义预层0，一种通过 $ U \mapsto F(U) \oplus G(U) $  定义两个预层 $ F $  和 $ G $  的直接和 $ F \oplus G $ 。因此， $ X $  上的预滑组范畴（交换群）和预滑组态射是一个加法范畴。我们用 $ \mathfrak{P}\mathfrak{S}\mathfrak{h}(X) $  表示它。

令 $\phi: F \to G$  为预滑的态射。对应关系：  $U \mapsto \mathrm{Ker} \phi_U$  （分别为  $U \mapsto \mathrm{Coker} \phi_U$  ）在  $X$  上定义了一个 presheaf，并且该 presheaf 是  $\mathfrak{P}\mathfrak{S}\mathfrak{h}(X)$  中  $\phi$  的内核（分别为 cokernel）。一个用 $\mathrm{Ker}\phi$ （或 $\mathrm{Coker}\phi$ ）表示。由于定义 1.2.3 的性质 (ii) 明显满足，我们发现  $\mathfrak{P}\mathfrak{S}\mathfrak{h}(X)$  是阿贝尔范畴。

术语 2.1.2。令  $F$  为  $X$  的预层， $U$  为  $X$  的开子集。元素  $s \in F(U)$  称为  $U$  上  $F$  的一部分。如果  $V$  是  $U$  的开子集，人们通常会写  $s|_{V}$  而不是  $\rho_{V,U}(s)$  ，并将其称为  $s$  到  $V$  的限制。

一组 $x \in X$ ：

\[
F_{x}=\varinjlim_{U}F(U),
\]

其中  $U$  贯穿  $x$  的开放邻域系列。群 $F_x$ 被称为 $x$ 处的 $F$ 的茎， $F_x$ 中的 $s \in F(U)$ （其中 $x \in U$ ）的像被称为 $x$ 处的 $s$ 的萌芽，并用 $s_x$ 表示。

其中一个用 $ F|_{U} $  表示由 $ U \supset V \mapsto F(V) $  定义的 $ U $  上的预层，并将其称为 $ F $  到 $ U $  的限制。

\subsection{层}
设X为拓扑空间。

\statementlabel{定义 2.2.1}X 上的预层（交换群）F 如果满足下面的条件 (S 1) – (S 2)，则称为束。

(S1) 对于任何开集  $U \subset X$  、任何开覆盖  $U = \bigcup_{i \in I} U_i$  、任何部分  $s \in F(U)$  、 $s|_{U_i} = 0$  对于所有  $i$  都意味着  $s = 0$  。

(S2) 对于任何开集  $U \subset X$  、任何开覆盖  $U = \bigcup_{i \in I} U_i$  、任何族  $s_i \in F(U_i)$  满足所有对  $(i, j)$  的  $s_i|_{U_i \cap U_j} = s_j|_{U_i \cap U_j}$  ，存在  $s \in F(U)$  使得所有  $i$  满足  $s|_{U_i} = s_i$  。

请注意，我们可以通过稍加修改来定义集轮的概念，但我们不会在本书中使用它。

此外请注意， (S 1) 和 (S 2) 相当于说，对于任何开子集  $U \subset X$  和任何开覆盖  $U = \bigcup_{i \in I} U_i$  ，通过有限交稳定，态射  $F(U) \to \lim_{\leftarrow} F(U_i)$  是同构。

如果  $F$  是一个层，则  $F(\varnothing) = 0$  。

如果  $F$  是  $X$  上的一个层，并且  $U$  在  $X$  中打开，则  $F|_{U}$  是  $U$  上的一个层。

将一束  $F$  的支持定义为  $\text{supp}(F)$  ，作为开集  $U \subset X$  并集的补集，使得  $F|_{U} = 0$  。类似地，我们将  $U$  上  $F$  的  $s$  部分的  $\text{support}$  定义为开集  $V \subset U$  的并集  $U$  中的补集，使得  $s|_{V} = 0$  。我们用 $\text{supp}(s)$  表示它。我们有 $\text{supp}(s) = \{x \in U; s_x \neq 0\}$ 。

人们将层的态射定义为底层预层的态射。那么我们用  $\mathfrak{Sh}(X)$  表示的层范畴显然是  $\mathfrak{PSh}(X)$  的满加性子范畴。其中一个用  $\Gamma(U;\cdot)$  表示从  $\mathfrak{Sh}(X)$  到  $\mathfrak{Ab}$  的函子  $F\mapsto F(U)$  。因此 $\Gamma(U;F)=F(U)$  。

\statementlabel{命题 2.2.2}令 $\phi: F \to G$  为层的态射。那么  $\phi$  是同构当且仅当对于任何  $x \in X$  ，导出的态射  $\phi_x: F_x \to G_x$  是同构。

\statementlabel{证明}其必要性是显而易见的。

相反，假设  $ \phi_x $  是所有  $ x \in X $  的同构，并让  $ U $  是  $ X $  的开子集。让我们证明 $ \phi_U: F(U) \to G(U) $  是单射的。让  $ s \in F(U) $  与  $ \phi_U(s) = 0 $  。然后是  $ (\phi_U(s))_x = \phi_x(s_x) = 0 $  和  $ s_x = 0 $  对于所有  $ x \in U $  。

根据公理 (S1)，这意味着  $s = 0$  。让我们证明  $\phi_U: F(U) \to G(U)$  的满射性。让 $t \in G(U)$  。然后根据假设，存在  $U$  和  $s_i \in F(U_i)$  的开覆盖  $\{U_i\}$  ，使得  $\phi(s_i) = t|_{U_i}$  。由于  $\phi(s_i)|_{U_i \cap U_j} = \phi(s_j)|_{U_i \cap U_j}$  ，  $\phi$  的单射性意味着  $s_i|_{U_i \cap U_j} = s_j|_{U_i \cap U_j}$  ，因此存在  $s \in F(U)$  使得  $s|_{U_i} = s_i$  。然后就可以轻松检查  $\phi(s) = t$  。   $\square$ 

\statementlabel{命题 2.2.3}给定  $X$  上的预层  $F$  ，存在一个层  $F^+$  和一个态射  $\theta: F \to F^+$  ，这样对于任何束  $G$  ，同态由  $\theta$  给出：

\[
\mathrm{H o m}_{\mathfrak{S b}(X)}(F^{+},G)\to\mathrm{H o m}_{\mathfrak{P S b}(X)}(F,G)
\]

是同构。换句话说， $F \mapsto F^+$  是包含函子 $\mathfrak{S h}(X) \to \mathfrak{P S h}(X)$  的左伴随函子（参见练习 I.2）。

此外 $(F^+,\theta)$ 在同构上是唯一的，并且对于任何 $x\in X,\theta_x:F_x\to F_x^+$ 都是同构。

\statementlabel{证明}对于任何开集 $U \subset X$ ，令 $F^+(U)$ 为从 $U$ 到 $\bigsqcup_{x \in U} F_x$ 的函数 $s$ 的集合，使得对于任何 $x \in U$ ， $s(x) \in F_x$ ，并且存在 $x$ ， $V \subset U$ 和 $t \in F(V)$ 的开邻域 $V$ ，使得 $t_y = s(y)$  对于所有  $y \in V$  。那么 $F^+$ 显然满足(S 1) 和(S 2)。

态射 $\theta: F \to F^+$ 定义如下：将 $x \mapsto s_x$ 给出的从 $U$ 到 $\bigsqcup_{x \in U} F_x$ 的映射分配给 $s \in F(U)$ 。那么很容易看出，对于任何  $x \in X$  ，  $\theta_x: F_x \to (F^+)_x$  都是同构。特别是，如果  $G$  是一个捆，那么  $\theta: G \to G^+$  是前面命题的同构。因此，对于预层 $F$ 和束 $G$ ，我们可以构造同态 $\text{Hom}_{\mathfrak{B}\mathfrak{S}(X)}(F, G) \to \text{Hom}_{\mathfrak{S}(X)}(F^+, G)$ 作为 $\text{Hom}_{\mathfrak{B}\mathfrak{S}(X)}(F, G) \to \text{Hom}_{\mathfrak{S}(X)}(F^+, G^+) \lesssim \text{Hom}_{\mathfrak{S}(X)}(F^+, G)$ 的组合。很容易检查得到的同态是  $\text{Hom}_{\mathfrak{S}(X)}(F^+, G) \to \text{Hom}_{\mathfrak{P}\mathfrak{S}(X)}(F, G)$  的逆。   $\square$ 

我们将关联的层称为  $ F^{+} $  预层 F。

令 $\phi: F \to G$  为层的态射。 presheaf  $U \mapsto \mathrm{Ker} \phi_U$  是一个层，并且是范畴  $\mathfrak{Sh}(X)$  中  $\phi$  的内核。相反，  $U \mapsto \mathrm{Coker} \phi_U$  并不总是一个层，但关联的层是  $\mathfrak{Sh}(X)$  中  $\phi$  的一个核心。因此，我们将用  $\mathrm{Ker} \phi$  表示层  $U \mapsto \mathrm{Ker} \phi_U$  ，用  $\mathrm{Coker} \phi$  表示与预层  $U \mapsto \mathrm{Coker} \phi_U$  相关的层。应注意  $\mathrm{Coker} \phi$  在  $\mathfrak{Sh}(X)$  和  $\mathfrak{Sh}(X)$  中的含义不同。由于我们始终使用层进行工作，因此不会有混淆的风险。

请注意，对于  $ x \in X $  ：

\[
\left(\mathrm{Ker}\phi\right)_{x}=\mathrm{Ker}\phi_{x},
\]

\[
\left(\mathrm{Coker}\phi\right)_{x}=\mathrm{Coker}\phi_{x},
\]

其中  $ \phi_x: F_x \to G_x $  是与族  $ \phi_U, x \in U $  相关的态射。

两个层 $F$  和 $G$  的直和被定义为底层预层的直和。这显然是一个层，是  $\mathfrak{S h}(X)$  中的直和。我们用 $F\oplus G$  表示它。

\statementlabel{命题 2.2.4}$X$  上交换群的束的范畴  $\mathfrak{Sh}(X)$  是交换范畴。

\statementlabel{证明}令 $\phi: F \to G$  为层的态射。令 $K = \operatorname{Coim}\phi, L = \operatorname{Im}\phi$  和 $\psi: K \to L$  为自然态射。对于每个  $x \in X$  ，  $\psi_x: K_x \to L_x$  都是同构，由 (2.2.1) 和 (2.2.2) 实现。那么结果由命题2.2.2得出。   $\Box$ 

\statementlabel{备注 2.2.5}通过应用命题 2.2.2，我们发现层复合体  $ F' \to F \to F'' $  是精确的当且仅当对于每个  $ x \in X $  组的序列  $ F_x' \to F_x \to F_x'' $  是精确的。特别是从  $ \mathfrak{S}h(X) $  到  $ \mathfrak{Ab} $  的函子  $ F \mapsto F_x $  是精确的。另一方面，从 $ \mathfrak{S}h(X) $  到 $ \mathfrak{Ab} $  的函子 $ \Gamma(X; \cdot) $  只保持精确。

人们自然遇到的大多数层都具有比仅仅阿贝尔群结构更丰富的结构。

让我们用 $ \mathbf{Ring} $ 来表示酉环的范畴和酉环的态射。值为  $ \mathbf{Ring} $  的预层称为环预层。如果这样的预层是一个层（其值在 $ \mathbf{Ab} $ 中），则它被称为环层。

\statementlabel{定义 2.2.6}让  $ \mathcal{R} $  成为  $ X $  上的环层。一个  $ \mathcal{R} $ -模  $ M $  ，（或  $ \mathcal{R} $  上的模层），是一个层  $ M $  ，使得对于每个开集  $ U \subset X $  ，  $ M(U) $  是左  $ \mathcal{R}(U) $  -模，并且对于任何开包含  $ V \subset U $  ，限制态射与模的结构兼容，即，  $ \rho_{V,U}(sm) = \rho_{V,U}(s) \cdot \rho_{V,U}(m) $  对于任何 $ s \in \mathcal{R}(U) $  、 $ m \in M(U) $  。

人们以一种明显的方式定义了  $ \mathscr{R} $  模的态射概念。

类似地定义了右  $ \mathcal{R} $  模的层概念。

其中一个用  $ \operatorname{Mod}(\mathcal{R}) $  表示（左） $ \mathcal{R} $  模的范畴。用 $ \mathcal{R}^{op} $ 表示与 $ \mathcal{R} $ 相对的环层（即 $ \mathcal{R}^{op}(U) = \mathcal{R}(U)^{op} $ ）， $ \operatorname{Mod}(\mathcal{R}^{op}) $ 相当于右 $ \mathcal{R} $ 模的范畴。

对于 $F$ ， $G \in \mathrm{Ob}(\mathfrak{Mod}(\mathscr{R}))$ ，我们应该写 $\mathrm{Hom}_{\mathscr{R}}(F, G)$ 而不是 $\mathrm{Hom}_{\mathfrak{Mod}(\mathscr{R})}(F, G)$ 。

因此  $ \text{Hom}_{\mathcal{R}}(\cdot, \cdot) $  是从  $ \text{Mod}(\mathcal{R})^\circ \times \text{Mod}(\mathcal{R}) $  到  $ \mathbb{U}b $  的双函子。

人们立即检查 $ \text{Mod}(\mathcal{R}) $  是阿贝尔范畴。此外，当且仅当  $ \text{Mod}(\mathcal{R}) $  中的序列被精确地视为  $ \text{Sh}(X) $  中的序列时，它才是精确的。

让我们用 $\mathbb{Z}_X$  来表示与预层 $U \mapsto \mathbb{Z}$  相关的 $X$  上的层（参见下面的定义2.2.11 以了解该结构的概括）。那么 $\mathbb{Z}_X$  是环层，其中一个有：

\[
\mathfrak{S h}(X)=\mathfrak{M o b}(\mathbb{Z}_{X})\enspace.
\]

我们将无差别地使用其中一种或另一种符号。此外，对于 $F$ ， $G \in \mathrm{Ob}(\mathfrak{S}h(X))$ ，我们将写 $\mathrm{Hom}(F, G)$ 而不是 $\mathrm{Hom}_{\mathcal{Z}_X}(F, G) = \mathrm{Hom}_{\mathfrak{S}h(X)}(F, G)$ 。

现在让 R 成为 X 上的环层，让 F 和 G 成为两个 R 模。考虑预层：

\[
U\mapsto\operatorname{H o m}_{\mathcal{A}|_{U}}(F|_{U},G|_{U})~.
\]

这个预层显然是一堆阿贝尔群（甚至是一堆 $\mathcal{R}$  -模，如果 $\mathcal{R}$  是可交换的）。

\statementlabel{定义 2.2.7}其中 $\mathcal{H}_{\mathcal{M}_{\mathcal{R}}}(F,G)$  表示(2.2.4) 中定义的层，并将其称为 $G$  中 $F$  的解层（在 $\mathcal{R}$  上）。

请注意，一般来说，自然态射：

\[
(\mathcal{H o m}_{\mathcal{R}}(F,G))_{x}\to\mathrm{H o m}_{\mathcal{R}_{x}}(F_{x},G_{x})
\]

既不是单射的，也不是满射的。通过构建，我们有：

\[
\Gamma(X;\mathcal{H o m}_{\mathcal{R}}(F,G))=\operatorname{H o m}_{\mathcal{R}}(F,G)\enspace.
\]

双函子 $ \text{Hom}_{\mathcal{R}}(\cdot,\cdot) $  对于它的每个参数都是精确的。现在让 $ F $  为右 $ \mathcal{R} $  模， $ G $  为左 $ \mathcal{R} $  模。

\statementlabel{定义 2.2.8}其中一个用  $F\otimes_{\mathcal{R}}G$  表示与预层  $U\mapsto F(U)\otimes_{\mathcal{R}(U)}G(U)$  关联的层，并将其称为  $F$  和  $G$  （在  $\mathcal{R}$  上）的张量积。

当 $ \mathscr{R} $  可交换时， $ F \otimes_{\mathscr{R}} G $  是一堆 $ \mathscr{R} $  模。当  $ \mathscr{R} = \mathbb{Z}_X $  时，只需写  $ F \otimes G $  而不是  $ F \otimes_{\mathbb{Z}_X} G $  。通过张量积的构造， $ x \in X $  可以得到：

\[
\left(F\otimes_{\mathcal{B}}G\right)_{x}=F_{x}\otimes_{\mathcal{B}_{x}}G_{x}.
\]

因此，函子  $\otimes_{\mathcal{R}}$  在其每个参数中都是右正合的。令 $\mathcal{S}$ 为另一个环层，并让 $H$ 为双( $\mathcal{R},\mathcal{S}$ )模(即 $H$ 被赋予左 $\mathcal{R}$ 模结构和右 $\mathcal{S}$ 模结构，使得两个动作可交换)。那么就有一个自然的同构：

\[
F\mathop{\otimes}_{\mathcal{R}}\left(H\mathop{\otimes}_{\mathcal{S}}G\right)\simeq\left(F\mathop{\otimes}_{\mathcal{R}}H\right)\mathop{\otimes}_{\mathcal{S}}G~.
\]

在这种情况下，我们只需写 $ F \otimes_{\mathscr{R}} H \otimes_{\mathscr{S}} G $ 即可。

\statementlabel{命题 2.2.9}令 $ \mathcal{R} $  为环层， $ \mathcal{S} $  为交换环层， $ \mathcal{S} \to \mathcal{R} $  为环层的态射，使得其图像包含在 $ \mathcal{R} $  的中心。令 F 和 G 为两个  $ \mathcal{R} $  模，H 为一个  $ \mathcal{S} $  模。那么就有规范的同构：

\[
\begin{aligned}{\mathcal{H}{o m}_{\mathcal{R}}\bigg(H\mathop{\otimes}_{\mathcal{S}}F,G\bigg)}&{{}\simeq\mathcal{H}{o m}_{\mathcal{R}}(F,\mathcal{H}{o m}_{\mathcal{S}}(H,G))}\\ {}&{{}\simeq\mathcal{H}{o m}_{\mathcal{S}}(H,\mathcal{H}{o m}_{\mathcal{R}}(F,G))~.}\\ \end{aligned}
\]

\statementlabel{证明}对于  $X$  的每个开集  $U$  都有自然同构：

\[
\begin{aligned}{\operatorname{H o m}_{\mathcal{H}(U)}\bigg(H(U)\mathop{\otimes}_{\mathcal{S}(U)}F(U),G(U)\bigg)}&{{}\simeq\operatorname{H o m}_{\mathcal{H}(U)}(F(U),\operatorname{H o m}_{\mathcal{S}(U)}(H(U),G(U)))}\\ {}&{{}\simeq\operatorname{H o m}_{\mathcal{S}(U)}(H(U),\operatorname{H o m}_{\mathcal{H}(U)}(F(U),G(U))).}\\ \end{aligned}
\]

让我们用  $ H \otimes_{\mathcal{S}}^{\vee} F $  表示 presheaf  $ U \mapsto H(U) \otimes_{\mathcal{S}(U)} F(U) $  并令  $ \mathrm{Hom}_{\mathcal{R}^{\vee}}(\cdot, \cdot) $ 

表示 $ \mathcal{R} $  -模的预滑范畴中的态射群。然后：

\[
\begin{aligned}{\operatorname{H o m}_{\mathcal{R}^{\vee}}\bigg(H\mathop{\otimes}_{\mathcal{G}}^{\vee}F,G\bigg)}&{{}\simeq\operatorname{H o m}_{\mathcal{R}}(F,\operatorname{H o m}_{\mathcal{S}}(H,G))}\\ {}&{{}\simeq\operatorname{H o m}_{\mathcal{S}}(H,\mathcal{H o m}_{\mathcal{R}}(F,G))}\\ \end{aligned}
\]

仍然需要应用命题 2.2.3。 ⑨

\statementlabel{推论 2.2.10}在命题 2.2.9 的情况下，有规范态射：

\[
\mathcal{H o m}_{\mathfrak{A}}(F,G)\mathop{\otimes}_{\mathcal{S}}F\to G\qquad{~i n~}\mathfrak{M o b}(\mathcal{R}),
\]

\[
\mathcal{H}{o m}_{\mathcal{R}}(F,G)\mathop{\otimes}_{\mathcal{S}}H\to\mathcal{H}{o m}_{\mathcal{R}}\bigg(F,G\mathop{\otimes}_{\mathcal{S}}H\bigg)\qquad{i n~}\mathfrak{W o d}(\mathcal{S}),
\]

\[
\mathcal{H}{o m}_{\mathcal{R}}(\mathrm{F},\mathrm{G})\to\mathcal{H}{o m}_{\mathcal{R}\circ\mathrm{p}}(\mathcal{H}{o m}_{\mathcal{R}}(\mathrm{G},\mathcal{R}),\mathcal{H}{o m}_{\mathcal{R}}(\mathrm{F},\mathcal{R}))\qquad{i n~}\mathfrak{M o d}(\mathcal{S})\enspace.
\]

\statementlabel{证明}(i) 到(2.2.9)我们有：

\[
\mathcal{H}{o m}_{\mathcal{R}}\bigg(\mathcal{H}{o m}_{\mathcal{R}}(F,G)\otimes_{\mathcal{S}}F,G\bigg)\simeq\mathcal{H}{o m}_{\mathcal{S}}(\mathcal{H}{o m}_{\mathcal{R}}(F,G),\mathcal{H}{o m}_{\mathcal{R}}(F,G))~.
\]

 $ \mathcal{Hom}_{\mathfrak{A}}(F, G) $  的恒等式定义 (2.2.10)。

(ii) 通过 (2.2.10) 中的 H 张量，我们有一个态射

\[
\mathcal{H o m}_{\mathfrak{A}}(F,G)\mathop{\otimes}_{\mathcal{S}}H\mathop{\otimes}_{\mathcal{S}}F\to G\mathop{\otimes}_{\mathcal{S}}H~.
\]

鉴于(2.2.9)，这定义了(2.2.11)。

(iii) 设置简称为 $ \mathbf{D}_{\mathcal{R}}F = \mathcal{H} \text{om}_{\mathcal{R}}(F, \mathcal{R}) $ 。那么  $ \mathbf{D}_{\mathcal{R}} $  是从  $ \mathfrak{Mod}(\mathcal{R}) $  到  $ \mathfrak{Mod}(\mathcal{R}^{op}) $  的函子。我们有态射：

\[
\begin{aligned}{\mathcal{H o m}_{\mathcal{R}^{\operatorname{o p}}}(\mathbf{D}_{\mathcal{R}}F,\mathbf{D}_{\mathcal{R}}F)}&{{}\to\mathcal{H o m}_{\mathcal{R}^{\operatorname{o p}}}\bigg(\mathcal{H o m}_{\mathcal{R}}\bigg(F,G\otimes_{\mathcal{S}}\mathbf{D}_{\mathcal{R}}G\bigg),\mathbf{D}_{\mathcal{R}}F\bigg)}\\ {}&{{}\to\mathcal{H o m}_{\mathcal{R}^{\operatorname{o p}}}\bigg(\mathcal{H o m}_{\mathcal{R}}(F,G)\otimes_{\mathcal{S}}\mathbf{D}_{\mathcal{R}}G,\mathbf{D}_{\mathcal{R}}F\bigg)}\\ {}&{{}\simeq\mathcal{H o m}_{\mathcal{S}}(\mathcal{H o m}_{\mathcal{R}}(F,G),\mathcal{H o m}_{\mathcal{R}^{\operatorname{o p}}}(\mathbf{D}_{\mathcal{R}}G,\mathbf{D}_{\mathcal{R}}F))}\\ \end{aligned}
\]

 $D_{\mathcal{R}}F$  的恒等式定义 (2.2.12)。 ⑨

现在我们将在  $ \mathfrak{Mod}(\mathscr{R}) $  范畴中构造直和归纳极限。

令  $\{F_i\}_{i \in I}$  为  $\mathcal{R}$  模的射影系统，由有序集  $I$  索引。预层  $U \mapsto \lim_{i \in I} F_i(U)$  显然是一个层，并且它是  $\{F_i\}$  的射影极限。

在 Mod(R) 中。因此我们用 $ \lim_{i} F_i $  表示它。换句话说，范畴 Mod(R) 中的射影极限存在，并且  $ \Gamma(U; \cdot) $  与  $ \lim_{i} $  交换，即

\[
\varGamma(U;\varliminf_{\leftarrow}F_{i})\simeq\varliminf_{\leftarrow}\varGamma(U;F_{i})~.
\]

请注意，如果  $x \in X$  ，则自然态射  $(\lim_{\leftarrow} F_i)_x \to \lim_{\leftarrow} (F_i)_x$  不一定是同构。

类似地，如果  $\{F_j\}_{j\in J}$  是  $\mathcal{R}$  模的归纳系统，由有序集  $J$  索引，我们用  $\lim_{\overrightarrow{j}} F_j$  表示与预层  $U\mapsto\lim_{j\in J}F_j(U)$  关联的层。请注意，对于  $x\in X$  ，有一个同构：

\[
\left(\varinjlim_{j}F_{j}\right)_{x}\simeq\varinjlim_{j}(F_{j})_{x}.
\]

然而，规范同态  $ \lim \Gamma(U; F_j) \to \Gamma(U; \lim F_j) $  不一定是同构。

作为前述构造的一个特例（即：当  $I$  或  $J$  上的阶为平凡阶时），我们获得了层族的乘积  $\prod_{i\in I} F_i$  或直和  $\bigoplus_{j\in J} F_j$  的概念。

在结束本节之前，让我们定义 X 上可能遇到的“最简单”层。

\statementlabel{定义 2.2.11}(i) 令 $M$  为阿贝尔群。其中一个用  $M_{X}$  表示与预层  $U \mapsto M$  相关联的层（U 在  $X$  中打开），并表示  $M_{X}$  是  $X$  上带有茎  $M$  的常量层。

(ii) 设 $F$  为 $X$  上的一个层。有人说 $F$ 在 $X$ 上是局部常值的，如果存在一个开覆盖 $X = \bigcup_{i} U_{i}$ ，使得对于每个 $i$ ， $F|_{U_{i}}$ 是常值层。

请注意，如果  $ M_x $  是具有茎  $ M $  的常值层，则  $ (M_x)_x = M $  对于所有  $ x \in X $  。还要注意，对于开集 $ U $ ， $ \Gamma(U; M_x) $  与从 $ U $  到赋予离散拓扑的拓扑空间 $ M $  的连续函数集是同构的。我们将在第 9 节中给出局部常值层的例子。

\statementlabel{备注 2.2.12}令 $\{pt\}$  表示具有单个元素的集合，并令 $A$  为环。然后  $A$  在空间  $\{pt\}$  上定义环层，有时将  $\mathfrak{Mod}(A)$  识别为  $\mathfrak{Mod}(A_{\{pt\}})$  很有用，即将  $A$  模视为一束  $A_{\{pt\}}$  模。

\subsection{层上的运算}
令 $X$  和 $Y$  为两个拓扑空间， $f: Y \to X$  为连续映射。

\statementlabel{定义 2.3.1}(i) 令 G 为 Y 上的层。G 由 f 的直像，表示为  $ f_{*} $  G，是 X 上的层，定义为：

\[
U\mapsto f_{*}G(U)=G(f^{-1}(U))~,\qquad U{~o p e n~i n~X~}.
\]

(ii) 令 F 为 X 上的层。F 由 f 的逆像，表示为  $f^{-1}F$  ，是与前层相关的 Y 上的层：

 $ V \mapsto \lim_{\rightarrow} F(U) $  ，V 在 Y 中开路，其中 U 穿过 X 中  $ f(V) $  的开邻域族。

人们以一种明显的方式定义了 Y 上（或 X 上）层态射的正像（或逆像）。因此我们得到两个函子：

\[
\begin{aligned}{}&{{}f_{*}\colon\textcircled{5}\mathfrak{h}(Y)\to\textcircled{5}\mathfrak{h}(X),}\\ {}&{{}f^{-1}\colon\textcircled{5}\mathfrak{h}(X)\to\textcircled{5}\mathfrak{h}(Y).}\\ \end{aligned}
\]

如果 $ \mathcal{R} $ （分别 $ \mathcal{S} $ ）是 $ X $ （分别 $ Y $ ）上的环层，那么 $ f^{-1}\mathcal{R} $ （分别 $ f_*\mathcal{S} $ ）是 $ Y $ （分别在 $ X $ ）上的环层，并且 $ f_* $ 和 $ f^{-1} $ 归纳函子，仍表示为  $ f_* $  和  $ f^{-1} $  ：

\[
f_{*}\colon\mathfrak{M o d}(\mathcal{S})\to\mathfrak{M o d}(f_{*}\mathcal{S})\enspace,
\]

\[
f^{-1}:\mathfrak{M o d}(\mathcal{R})\to\mathfrak{M o d}(f^{-1}\mathcal{R})\enspace.
\]

示例 2.3.2。让 $ a_X: X \to \{\text{pt}\} $  。 （回想一下 $ \{\text{pt}\} $  是只有一个元素的集合。）令 $ M $  为阿贝尔群。然后：

\[
\begin{array}{r}{M_{X}=a_{X}^{-1}M_{\{\mathrm{p t}\}}.}\end{array}
\]

设 F 为 X 上的一个层，则：

\[
\begin{array}{r}{\varGamma(X;F)=\varGamma(\{\mathsf{p t}\};a_{X*}{F})\simeq a_{X*}{F}\enspace.}\end{array}
\]

让 $ y \in Y $  和 $ F $  成为 $ X $  上的一个层。然后：

\[
(f^{-1}F)_{y}=F_{f(y)}\enspace.
\]

通过这个公式，我们看到 $ f^{-1} $  是一个精确函子。另一方面，函子  $ f_{*} $  仅左精确，如下所示：函子  $ \Gamma(U;\cdot) $  的左精确。

函子存在自然态射。

\[
f^{-1}\circ f_{*}\to\operatorname{i d}\qquad\mathrm{~i n~}\mathfrak{M o b}(f^{-1}\mathcal{R}),
\]

\[
\mathrm{i d}\to f_{*}\circ f^{-1}\qquad\mathrm{~i n~\mathfrak{M o b}(\mathcal{R})~}.
\]

\statementlabel{命题 2.3.3}让  $ \mathcal{R} $  成为  $ X $  上的环层。让  $ F \in \text{Ob}(\text{Mod}(\mathcal{R})) $  和  $ G \in \text{Ob}(\text{Mod}(f^{-1}\mathcal{R})) $  。然后：

\[
\operatorname{H o m}_{\mathcal{R}}(F,f_{*}G)\simeq\operatorname{H o m}_{f^{-1}\mathcal{R}}(f^{-1}F,G)\enspace.
\]

换句话说，  $ f^{-1} $  是  $ f_{*} $  的左伴随，  $ f_{*} $  是  $ f^{-1} $  的右伴随。

\statementlabel{证明}我们有同态

\[
\mathrm{H o m}_{\mathcal{R}}(F,f_{*}G)\xrightarrow[\alpha]{}\mathrm{H o m}_{f^{-1}\mathcal{R}}(f^{-1}F,f^{-1}f_{*}G)\xrightarrow[\beta]{}\mathrm{H o m}_{f^{-1}\mathcal{R}}(f^{-1}F,G)
\]

and

\[
\mathrm{H o m}_{f^{-1}\mathcal{R}}(f^{-1}F,G)\underset{\gamma}{\to}\mathrm{H o m}_{f*f^{-1}\mathcal{R}}(f_{*}f^{-1}F,f_{*}G)\underset{\delta}{\to}\mathrm{H o m}_{\mathcal{R}}(F,f_{*}G)\enspace.
\]

同态 $ \alpha $ 和 $ \gamma $ 以显而易见的方式定义， $ \beta $ 和 $ \delta $ 分别从(2.3.4)和(2.3.5)推导出来。然后，我们可以轻松检查 $ \beta \circ \alpha $  和 $ \delta \circ \gamma $  是否彼此相反。   $ \Box $ 

\statementlabel{推论 2.3.4}在命题2.3.3的情况下，有：

\[
\mathcal{H o m}_{\mathcal{R}}(\mathbb{F},f_{*}G)\simeq f_{*}\mathcal{H o m}_{f^{-1}\mathcal{R}}(f^{-1}\mathbb{F},G)\enspace.
\]

\statementlabel{证明}设 U 是 X 的开子集，则：

\[
\begin{aligned}{\varGamma(U;f_{*}\mathcal{H}o m_{f^{-1}\mathcal{R}}(f^{-1}F,G))}&{{}=\operatorname{H o m}_{f^{-1}\mathcal{R}|_{f^{-1}(U)}}(f^{-1}F|_{f^{-1}(U)},G|_{f^{-1}(U)})}\\ {}&{{}=\operatorname{H o m}_{\mathcal{R}|_{U}}(F|_{U},f_{*}G|_{U})}\\ {}&{{}=\varGamma(U;\mathcal{H}o m_{\mathcal{R}}(F,f_{*}G))~.\quad\Box}\\ \end{aligned}
\]

令 $Z$  为另一个拓扑空间， $g: Z \to Y$  为连续映射。通过构建正像或逆像，我们立即看到：

\[
(f\circ g)_{*}=f_{*}\circ g_{*},
\]

\[
(f\circ g)^{-1}=g^{-1}\circ f^{-1}.
\]

逆像转换为张量积。更准确地说：

\statementlabel{命题 2.3.5}令  $ F_{1} $  （分别为  $ F_{2} $  ）为右（分别为左）R 模。然后有一个规范的同构：

\[
f^{-1}F_{1}\mathop{\otimes}_{f^{-1}\mathcal{R}}f^{-1}F_{2}\simeq f^{-1}\biggl(F_{1}\mathop{\otimes}_{\mathcal{R}}F_{2}\biggr).
\]

\statementlabel{证明}态射 (2.3.10) 由以下因素引发：

\[
F_{1}(U)\mathop{{\otimes}}_{\mathcal{R}(U)}F_{2}(U)\to\biggl(F_{1}\mathop{{\otimes}}_{\mathcal{R}}F_{2}\biggr)(U)~,\qquad{U\mathrm{~o p e n~i n}\;X}~.
\]

为了证明这是一个同构，取  $ y \in Y $  并设置  $ x = f(y) $  。然后：

\[
\begin{aligned}{\left(f^{-1}F_{1}\mathop{\otimes}_{f^{-1}\mathcal{R}}f^{-1}F_{2}\right)_{y}}&{{}\simeq(f^{-1}F_{1})_{y}\mathop{\otimes}_{(f^{-1}\mathcal{R})_{y}}(f^{-1}F_{2})_{y}}\\ {}&{{}\simeq(F_{1})_{x}\mathop{\otimes}_{\mathcal{R}_{x}}(F_{2})_{x}}\\ {}&{{}\simeq\left(F_{1}\mathop{\otimes}_{\mathcal{R}}F_{2}\right)_{x}\;.\quad\Box}\\ \end{aligned}
\]

我们现在将构造与  $X$  的子集  $Z$  关联的函子。我们赋予 $Z$ 诱导拓扑，并用 $j: Z \hookrightarrow X$ 表示包含。我们设置 $F \in \mathrm{Ob}(\mathfrak{S} \mathfrak{h}(X))$ ：

\[
\begin{array}{r}{F|_{z}=j^{-1}F}\end{array},
\]

\[
\Gamma(\mathbf{Z};F)=\Gamma(\mathbf{Z};j^{-1}F).
\]

请注意，(2.3.12) 与之前为 Z open 引入的符号一致。

存在自然态射  $ \Gamma(X; F) \to \Gamma(Z; F|_{z}) $  。如果 $ s \in \Gamma(X; F) $ 我们用 $ s|_{z} $ 表示它在 $ \Gamma(Z; F|_{z}) $ 中的图像，并将其称为s到Z的限制。

假设 Z 在 X 中是封闭的。在这种情况下，我们设置：

\[
F_{Z}=j_{*}j^{-1}F\quad.
\]

因此我们有一个自然态射 $ F \to F_Z $ 。此外：

\[
\left\{\begin{aligned}{}&{{}F_{Z}|_{Z}=F|_{Z},}\\ {}&{{}F_{Z}|_{X\setminus Z}=0.}\\ \end{aligned}\right.
\]

特别是， $ (F_Z)_x = F_x $  如果 $ x \in Z $  ， $ (F_Z)_x = 0 $  如果 $ x \notin Z $  。当  $ Z $  是  $ X $  的局部闭子集时，仍然可以在  $ X $  上构造满足 (2.3.13) 的束  $ F_Z $  。如果  $ Z $  在  $ X $  中打开，则一组：

\[
F_{Z}=\operatorname{K e r}(F\to F_{X\setminus Z})\enspace.
\]

在一般情况下，让我们写成 $ Z = U \cap A $ ，其中U在X中开，A在X中闭。我们设置：

\[
F_{Z}=\left(F_{U}\right)_{A}.
\]

这个定义不依赖于U和A的选择（参见下面命题中的（i））。

让我们列出函子  $(\cdot)_{z}: F \mapsto F_{z}$  的主要属性。

\statementlabel{命题 2.3.6}设 Z 为 X 的局部闭子集，设 F 为 X 上的一个层。

(i) 该层 $F_{z}$  满足(2.3.13)。此外，X 上满足 (2.3.13) 的任何层都同构于  $F_{z}$  。

(ii) 函子 $ (\cdot)_{z}: F \mapsto F_{z} $  是精确的。

(iii) 令  $ Z' $  为 X 的另一个局部闭子集。则：

\[
(F_{2})_{z^{\prime}}=F_{z\circ z^{\prime}}.
\]

(iv) 假设 $Z$  在 $X$  中是封闭的，并让 $j: Z \subsetneq X$  是 $Z$  在 $X$  中的包含。然后：

\[
F_{Z}=j_{*}j^{-1}F\quad.
\]

(v) 令  $Z'$  为  $Z$  的闭子集（Z 为局部闭）。那么下面的顺序是准确的：

\[
0\to F_{Z\setminus Z^{\prime}}\to F_{Z}\to F_{Z^{\prime}}\to0~.
\]

(vi) 令 $ Z_{1} $  和 $ Z_{2} $  为X 的两个闭子集。则序列：

\[
0\longrightarrow F_{Z_{1}\cup Z_{2}}\mathop{\longrightarrow}_{\alpha}F_{Z_{1}}\oplus F_{Z_{2}}\mathop{\longrightarrow}_{\beta}F_{Z_{1}\cap Z_{2}}\longrightarrow0
\]

是准确的。这里  $\alpha=(\alpha_{1},\alpha_{2})$  和  $\beta=(\beta_{1},-\beta_{2})$  分别由自然态射  $F_{z_{1}\cup z_{2}}\to F_{z_{i}}$  和  $F_{z_{i}}\to F_{z_{1}\cap z_{2}}(i=1,2)$  诱导。

(vii) 令 $ U_{1} $  和 $ U_{2} $  为X 的两个开子集。则序列：

\[
0\longrightarrow F_{U_{1}\cap U_{2}}\xrightarrow[\gamma]{}F_{U_{1}}\oplus F_{U_{2}}\xrightarrow[\delta]{}F_{U_{1}\cup U_{2}}\longrightarrow0
\]

是准确的。这里  $\gamma=(\gamma_{1},\gamma_{2})$  和  $\delta=(\delta_{1},-\delta_{2})$  分别由自然态射  $F_{U_{1}\cap U_{2}}\to F_{U_{i}}$  和  $F_{U_{i}}\to F_{U_{1}\cap U_{2}}(i=1,2)$  诱导。

\statementlabel{证明}很容易验证（i）并且其他断言随之而来。 ⑨

\statementlabel{推论 2.3.7}令  $Z_{1}$  和  $Z_{2}$  为  $X$  的两个闭子集。然后是顺序：

\[
0\to\varGamma(\mathbb{Z}_{1}\cup\mathbb{Z}_{2};F)\to\varGamma(\mathbb{Z}_{1};F)\oplus\varGamma(\mathbb{Z}_{2};F)\to\varGamma(\mathbb{Z}_{1}\cap\mathbb{Z}_{2};F)
\]

是准确的。

\statementlabel{证明}将左精确函子  $ \Gamma(X;\cdot) $  应用于命题 2.3.6 (vi) 的精确序列，并注意  $ \Gamma(X;F_z)\simeq\Gamma(Z;F) $  ，对于  $ Z $  是  $ X $  的闭子集。   $ \square $ 

还有另一个层可以在功能上与 X 的局部闭子集 Z 相关联。设 U 是 X 的开子集，Z 是 U 的闭子集。一组：

\[
\Gamma_{Z}(U;F)=\operatorname{K e r}(F(U)\to F(U\backslash Z))~.
\]

因此 $ \Gamma_{Z}(U; F) $ 是 $ \Gamma(U; F) $ 的子群，由其支持包含在Z中的部分组成。

令  $V$  为包含  $Z$  的  $U$  的开子集。规范态射  $\Gamma_z(U; F) \to \Gamma_z(V; F)$  是同构。因此对于  $X$  的局部闭子集  $Z$ 

我们可以将  $ \varGamma_Z(X; F) $  定义为  $ \varGamma_Z(U; F) $  ，其中  $ U $  是  $ X $  的任何开子集，其中包含  $ Z $  作为闭子集。请注意，预层 $ U \mapsto \varGamma_{Z \cap U}(U; F) $  是一个层。

\statementlabel{定义 2.3.8}其中  $ \Gamma_Z(F) $  表示层  $ U \mapsto \Gamma_{Z \cap U}(U; F) $  ，并将其称为由 Z 支持的 F 部分的层。

\statementlabel{命题 2.3.9}设 Z 为 X 的局部闭子集，并设 F 为 X 上的层。

(i) 从 $\mathfrak{Sh}(X)$  到 $\mathfrak{Ab}$  的函子 $\Gamma_z(X;\cdot):F \mapsto \Gamma_z(X;F)$  和从 $\mathfrak{Sh}(X)$  到 $\mathfrak{Sh}(X)$  的函子 $\Gamma_z(\cdot):F \mapsto \Gamma_z(F)$  保持精确。此外：

\[
\Gamma_{z}(X;\cdot)=\Gamma(X;\cdot)\circ\Gamma_{z}(\cdot).
\]

(ii) 令  $ Z' $  为 X 的另一个局部闭子集。然后：

\[
\varGamma_{Z^{\prime}}(\cdot)\circ\varGamma_{Z}(\cdot)=\varGamma_{Z\cap Z^{\prime}}(\cdot).
\]

(iii) 假设 $Z$  在 $X$  中是开放的，并让 $i:Z \hookrightarrow X$  是 $Z$  在 $X$  中的注入。然后：

\[
\varGamma_{Z}(\cdot)=i_{*}\circ i^{-1}\quad.
\]

(iv) 令  $ Z' $  为 X 的局部闭子集 Z 的闭子集。那么下面的序列是精确的：

\[
0\to\varGamma_{Z^{\prime}}(F)\to\varGamma_{Z}(F)\to\varGamma_{Z\setminus Z^{\prime}}(F)\enspace.
\]

(v) 令  $U_{1}$  和  $U_{2}$  为  $X$  的两个开子集。然后是顺序：

\[
0\longrightarrow\varGamma_{U_{1}\cup U_{2}}(F)\xrightarrow[\alpha]{}\varGamma_{U_{1}}(F)\oplus\varGamma_{U_{2}}(F)\xrightarrow[\beta]{}\varGamma_{U_{1}\cap U_{2}}(F)
\]

是准确的。这里  $\alpha=(\alpha_{1},\alpha_{2})$  和  $\beta=(\beta_{1},-\beta_{2})$  分别由态射  $\Gamma_{U_{1}\cup U_{2}}(F)\to\Gamma_{U_{i}}(F)$  和  $\Gamma_{U_{i}}(F)\to\Gamma_{U_{1}\cap U_{2}}(F)$  (  $i=1,2$  ) 导出。

(vi) 令 $ Z_{1} $  和 $ Z_{2} $  为X 的两个闭子集。则序列：

\[
0\longrightarrow\varGamma_{\mathbb{Z}_{1}\cap\mathbb{Z}_{2}}(F)\xrightarrow[\gamma]{}\varGamma_{\mathbb{Z}_{1}}(F)\oplus\varGamma_{\mathbb{Z}_{2}}(F)\xrightarrow[\delta]{}\varGamma_{\mathbb{Z}_{1}\cup\mathbb{Z}_{2}}(F)
\]

是准确的。这里  $\gamma=(\gamma_{1},\gamma_{2})$  和  $\delta=(\delta_{1},-\delta_{2})$  分别由态射  $\Gamma_{z_{1}\cap z_{2}}(F)\xrightarrow[\gamma_{i}]{\vphantom{\gamma_{i}}\vphantom{\gamma_{i}}}\Gamma_{z_{i}}(F)$  和  $\Gamma_{z_{i}}(F)\xrightarrow[\delta_{i}]{}\Gamma_{z_{1}\cup z_{2}}(F)$  (  $i=1,2$  ) 导出。

\statementlabel{证明}很简单。

我们在范畴  $\mathcal{Sh}(X)$  上定义了几个函子：函子  $\mathcal{Hom}(\cdot,\cdot),\;\dot{\otimes}\cdot,f_{*},f^{-1},(\cdot)_{z},\,\Gamma_{Z}(\cdot),\,\Gamma(X;\cdot)$  。现在我们要研究它们的一些关系。

\statementlabel{命题 2.3.10}令  $ \mathcal{R} $  为  $ X $  上的环层，并令  $ F \in \text{Ob}(\mathfrak{Mod}(\mathcal{R})) $  。令  $ Z $  为  $ X $  的局部闭子集。那么我们就有自然同构：

\[
\mathcal{R}_{2}\mathop{{\otimes}}_{\mathcal{A}}F\simeq F_{2}~,
\]

\[
\mathcal{H o m}_{\mathfrak{A}}(\mathcal{R}_{Z},\mathrm{F})\simeq\varGamma_{Z}(F)\enspace.
\]

\statementlabel{证明}第一个语句 (2.3.15) 源自

\[
\left(\mathcal{R}_{Z}\mathop{\otimes}_{\mathcal{R}}F\right)\bigg|_{Z}\simeq\left(\mathcal{R}_{Z|Z}\right)\mathop{\otimes}_{\mathcal{R}|z}(F|_{Z})\simeq\left(\mathcal{R}|_{Z}\right)\mathop{\otimes}_{\mathcal{R}|z}(F|_{Z})\simeq F|_{Z}
\]

和  $ (\mathcal{R}_Z \otimes_{\mathcal{R}} F)|_{X \setminus Z} = 0 $  ，与命题 2.3.6 (i) 结合。

为了证明（2.3.16），假设首先  $Z$  在  $X$  中打开。对于 $X$  的任何开子集 $U$  ， $\Gamma(U; \mathcal{R}_Z)$  元素的支持在 $U \cap Z$  中是封闭的。从这句话中我们推断出自然态射：

\[
\operatorname{H o m}_{\mathcal{R}}(\mathcal{R}_{Z},F)\to\operatorname{H o m}_{\mathcal{R}|_{Z}}(\mathcal{R}_{Z}|_{Z},F|_{Z})
\]

是同构。右侧的术语是  $ \Gamma(Z; F) $  。然后通过用  $ U $  替换  $ X $  ，对于  $ U $  open in  $ X $  ，在这个同构中，我们得到  $ Z $  open 的结果。  $ Z $  闭合的情况如下：将函子  $ \mathcal{Hom}_{\mathcal{R}}(\cdot, F) $  应用于精确序列  $ 0 \to \mathcal{R}_{X \setminus Z} \to \mathcal{R} \to \mathcal{R}_Z \to 0 $  并与精确序列  $ 0 \to \Gamma_Z(F) \to F \to \Gamma_{X \setminus Z}(F) $  进行比较。

最后假设  $ Z = A \cap U $  ，其中  $ A $  是封闭的，而  $ U $  在  $ X $  中是开放的。我们得到：

\[
\begin{aligned}{\mathcal{H}{{{o}}m}_{\mathcal{R}}(\mathcal{R}_{Z},F)}&{{}=\mathcal{H}{{{o}}m}_{\mathcal{R}}(\mathcal{R}_{A\cap U},F)}\\ {}&{{}=\mathcal{H}{{{o}}m}_{\mathcal{R}}(\mathcal{R}_{A},\varGamma_{U}(F))}\\ {}&{{}=\varGamma_{A}(\varGamma_{U}(F))}\\ {}&{{}=\varGamma_{U\cap A}(F)\enspace.\enspace\Box}\\ \end{aligned}
\]

\statementlabel{备注 2.3.11}通过使用命题2.2.9、2.3.3、2.3.4、2.3.5、2.3.9、2.3.10，我们可以获得许多其他的态射或同构。例如，令  $ f: Y \to X $  为连续映射， $ \mathscr{R} $  为  $ X $  上的环层， $ Z $  为  $ X $  的局部闭子集，并令  $ F $  、  $ F_1 $  、  $ F_2 $  （分别为  $ G $  、  $ G_1 $  、  $ G_2 $  ）是 $ \mathscr{R} $  -模（分别为 $ f^{-1} \mathscr{R} $  -模）的层。 （考虑张量积时，  $ F_1 $  、  $ G $  和  $ G_1 $  将是右模。）然后有自然态射或同构：

\[
\left(F_{1}\mathop{{\otimes}}_{\mathcal{R}}F_{2}\right)_{Z}\simeq F_{1}\mathop{{\otimes}}_{\mathcal{R}}(F_{2})_{Z}\simeq(F_{1})_{Z}\mathop{{\otimes}}_{\mathcal{R}}F_{2},
\]

\[
\mathcal{H}{o m}_{\mathcal{R}}((F_{1})_{Z},F_{2})\simeq\mathcal{H}{o m}_{\mathcal{R}}(F_{1},\varGamma_{Z}(F_{2}))\simeq\varGamma_{Z}\mathcal{H}{o m}_{\mathcal{R}}(F_{1},F_{2})~,
\]

\[
f^{-1}F_{z}\simeq(f^{-1}F)_{f^{-1}(z)},
\]

\[
\begin{array}{r}{\varGamma_{Z}f_{*}G\simeq f_{*}\varGamma_{f^{-1}(Z)}(G),}\end{array}
\]

\[
f_{*}G\mathop{\otimes}_{\mathcal{R}}F\to f_{*}\left(G\mathop{\otimes}_{f^{-1}\mathcal{R}}f^{-1}F\right)\;,
\]

\[
f_{*}G_{1}\mathop{{\otimes}_{\mathcal{R}}}f_{*}G_{2}\to f_{*}\biggl(G_{1}\mathop{{\otimes}_{f^{-1}\mathcal{R}}}G_{2}\biggr),
\]

\[
\begin{aligned}{}&{{}f_{*}\mathcal{H}{{}o m}_{f^{-1}\mathcal{R}}(G_{1},G_{2})\to\mathcal{H}{{}o m}_{\mathcal{R}}(f_{*}G_{1},f_{*}G_{2})\enspace,}\\ {}&{{}\quad f^{-1}\mathcal{H}{{}o m}_{\mathcal{R}}(F_{1},F_{2})\to\mathcal{H}{{}o m}_{f^{-1}\mathcal{R}}(f^{-1}F_{1},f^{-1}F_{2})\enspace.}\\ \end{aligned}
\]

例如，让我们构造态射 (2.3.21)。考虑态射链（我们不会简写为  $ \mathscr{R} $  ）：

\[
\begin{aligned}{\operatorname{H o m}(G\otimes f^{-1}F,G\otimes f^{-1}F)}&{{}\to\operatorname{H o m}(f^{-1}f_{*}G\otimes f^{-1}F,G\otimes f^{-1}F)}\\ {}&{{}\simeq\operatorname{H o m}(f^{-1}(f_{*}G\otimes F),G\otimes f^{-1}F)}\\ {}&{{}\simeq\operatorname{H o m}(f_{*}G\otimes F,f_{*}(G\otimes f^{-1}F))~.}\\ \end{aligned}
\]

 $ G \otimes f^{-1} F $  的恒等态射图像给出了所需的态射。

符号2.3.12。令 $p_X: X \to S$  和 $p_Y: Y \to S$  为两个连续映射，并令 $X \times_S Y = \{(x, y) \in X \times Y; p_X(x) = p_Y(y)\}$  为 $X$  和 $Y$  在 $S$  上的纤维积。我们分别用 $q_1$ 和 $q_2$ 表示从 $X \times_S Y$ 到 $X$ 和 $Y$ 的投影，用 $p$ 表示 $X \times_S Y \to S$ 的投影（参见图2.3.25）。

\begin{center}
\includegraphics[width=0.28\textwidth]{流行上的层_Chn-assets/img_in_image_box_431_871_769_1315.png}
\end{center}

令 $ \mathcal{R} $  为S 上的环层，令 $ F $  （分别为 $ G $  ）为一束 $ p_X^{-1}(\mathcal{R}^{op}) $  模（分别为 $ p_Y^{-1}\mathcal{R} $  模）。一套：

\[
F\boxtimes_{\mathcal{R}}G=q_{1}^{-1}F\:\otimes\:q_{2}^{-1}G\:\:.
\]

如果没有混淆的风险，我们只需在这些公式中写 $ F \boxtimes_S G $ ，如果 $ S = \{pt\} $ ，我们就不写 $ S $ 。请注意，如果  $ S = \{pt\} $  、  $ X = Y $  且 if one 由  $ \delta $  表示对角线嵌入，则：

\[
F\otimes G\simeq\delta^{-1}(F\boxtimes G)\enspace.
\]

束 $ F \boxtimes_S G $  称为 $ F $  和 $ G $  的外张量积（在 $ S $  上）。

\subsection{内射层、柔软层与平坦层}
令 $X$  为拓扑空间， $\mathcal{R}$  为 $X$  上的环层， $F$  为 $\mathfrak{Mod}(\mathcal{R})$  的对象。如果  $F$  在范畴  $\mathfrak{Mod}(\mathcal{R})$  中是单射的，我们可以说  $F$  是  $\mathcal{R}$  -单射的。

\statementlabel{命题 2.4.1}(i) 假设 $F$  是 $\mathcal{R}$  -内射。令  $U$  为  $X$  的开子集。那么 $F|_{U}$  是 $\mathcal{R}|_{U}$  -内射。

(ii) 令 $f: Y \to X$  为连续映射，并令 $G$  为 $Y$  上的 $f^{-1} \mathcal{R}$  内层。那么 $f_* G$  是 $\mathcal{R}$  -内射。

\statementlabel{证明}(i) 令 $ i: U \hookrightarrow X $  为包含图。然后对于  $ \mathcal{R}|_{U} $  -模  $ G $  我们得到 (2.3.18)：

\[
\begin{aligned}{\operatorname{H o m}_{\mathcal{R}|U}(G,F|_{U})}&{{}=\operatorname{H o m}_{\mathcal{R}|U}((i_{*}G)|_{U},F|_{U})}\\ {}&{{}=\operatorname{H o m}_{\mathcal{R}}((i_{*}G)_{U},F)\enspace.}\\ \end{aligned}
\]

由于函子  $G \mapsto (i_{*}G)_{U}$  是精确的，所以函子  $G \mapsto \mathrm{Hom}_{\mathcal{A}|U}(G, F|_{U})$  也是精确的。

(ii) 应用命题 2.3.3。 ⑨

\statementlabel{推论 2.4.2}假设 $F$  是 $\mathcal{R}$  -内射。那么函子 $\mathcal{Hom}_{\mathcal{R}}(\cdot,F)$  是精确的。

\statementlabel{命题 2.4.3}令 $ \mathcal{R} $  为X 上的环层。那么范畴 $ \mathfrak{Mod}(\mathscr{R}) $  就有足够的单射。

\statementlabel{证明}令 $\hat{X}$  为具有离散拓扑的空间 $X$ ，并令 $f: \hat{X} \to X$  为自然映射。让  $F \in \mathrm{Ob}(\mathfrak{Mod}(\mathscr{R}))$  并假设我们已经找到了  $f^{-1} \mathcal{R}$  -内层  $I$  和单态  $f^{-1}F \to I$  。应用左精确函子  $f_*$  并使用态射 (2.3.5)，我们得到精确序列  $0 \to F \to f_* I$  。由于 $f_* I$  是 $\mathscr{R}$  命题2.4.1 的内射，因此足以证明 $\hat{X}$  上的结果。

令  $F \in \mathrm{Ob}(\mathfrak{Mod}(f^{-1}\mathcal{R}))$  ，对于每个  $x \in X$  ，令  $0 \to F_x \to I_x$  为  $\mathfrak{Mod}(\mathcal{R}_x)$  中的精确序列，并带有  $I_x$  单射（参见例 1.8.9）。然后  $\prod_{x \in X} I_x$  在  $\widehat{X}$  上定义了一个单层  $I$  ，并且  $0 \to F \to I$  是精确的。   $\square$ 

\statementlabel{备注 2.4.4}一般来说，范畴  $ \mathfrak{Mod}(\mathscr{R}) $  没有足够的射影（参见练习 II.23）。

\statementlabel{定义 2.4.5}如果对于 $X$  的任何开子集 $U$ ，限制态射 $\Gamma(X; F) \to \Gamma(U; F)$  是满射的，则 $X$  上的束 $F$  是松弛的。

\statementlabel{命题 2.4.6}设 F 是 X 上的松弛捆。

(i) 对于任何开子集 $ U \subset X $  ，U 上的束 $ F|_{U} $  是松弛的。

(ii) 令 $f: X \to Y$  为连续映射。然后层 $f_{*}$ 就松弛了。

(iii) 令  $ Z $  为  $ X $  的局部闭子集。然后层 $ \Gamma_Z(F) $ 就松弛了。

(iv) 设 Z 为 X 的局部闭子集， $ Z' $  Z 的闭子集。则以下序列是精确的：

\[
0\to\varGamma_{Z^{\prime}}(F)\to\varGamma_{Z}(F)\to\varGamma_{Z\setminus Z^{\prime}}(F)\to0~.
\]

(v) 令  $U_{1}$  和  $U_{2}$  为  $X$  的两个开子集。然后是顺序：

\[
0\longrightarrow\varGamma_{U_{1}\cup U_{2}}(F)\xrightarrow[\alpha]{}\varGamma_{U_{1}}(F)\oplus\varGamma_{U_{2}}(F)\xrightarrow[\beta]{}\varGamma_{U_{1}\cap U_{2}}(F)\longrightarrow0
\]

是精确的（ $ \alpha $  和  $ \beta $  在命题 2.3.9 中定义）。

(vi) 令 $ Z_{1} $  和 $ Z_{2} $  为X 的两个闭子集。则序列：

\[
0\longrightarrow\varGamma_{\mathbb{Z}_{1}\cap\mathbb{Z}_{2}}(F)\xrightarrow[\gamma]{}\varGamma_{\mathbb{Z}_{1}}(F)\oplus\varGamma_{\mathbb{Z}_{2}}(F)\xrightarrow[\delta]{}\varGamma_{\mathbb{Z}_{1}\cup\mathbb{Z}_{2}}(F)\longrightarrow0
\]

是真映射，$\gamma$ 和 $\delta$ 在命题 2.3.9 中定义。

(vii) 令  $ \mathcal{R} $  为 X 上的环层，G 为  $ \mathcal{R} $  模，H 为  $ \mathcal{R} $  内射模。然后 $ \mathcal{Hom}_{\mathcal{R}}(G,H) $ 就松弛了。特别是  $ \mathcal{R} $  -内射模是不稳定的。

\statementlabel{证明}(i)和(ii)是显而易见的。

(iii) 用 $U$  替换 $X$ ，其中 $U$  是 $X$  的开子集，包含 $Z$  作为闭子集，我们可以从一开始就假设 $Z$  是闭子集。令  $U$  为  $X$  的开子集。我们必须证明限制态射： $\Gamma_Z(X; F) \to \Gamma_{Z \cap U}(U; F)$  是满射的。让 $s \in \Gamma_{Z \cap U}(U; F)$  。我们首先在  $X \setminus Z$  上通过  $0$  扩展  $s$  。让 $s'$  成为这个扩展名。然后我们利用  $F$  的松弛性将  $s'$  扩展到  $X$  。

(iv) 对于任何开集  $U$  ， $\Gamma(U; \Gamma_Z(F)) \to \Gamma(U; \Gamma_{Z\setminus Z'}(F)) \simeq \Gamma(U\setminus Z'; \Gamma_Z(F))$  是 (iii) 的满射。

(v) 根据命题 2.3.9，仍需证明  $ \beta $  是一个外同态，但限制态射  $ \Gamma_{U_1}(F) \to \Gamma_{U_1 \cap U_2}(F) $  本身是一个外同态。

(vi) 让 $s \in \Gamma_{Z_1 \cup Z_2}(X; F)$  。我们可以找到  $s_i \in \Gamma_{Z_i}(X \setminus Z_1 \cap Z_2; F)$  (i = 1,2) 使得  $s = s_1 - s_2$  在  $X \setminus Z_1 \cap Z_2$  上。我们将 $s_1$  和 $s_2$  扩展到 $X$  上。让 $s'_1$  和 $s'_2$  成为这些扩展。然后 $s'_1 - s'_2 = s + s', s' \in \Gamma_{Z_1 \cap Z_2}(X; F)$  和 $(s'_1 - s') - s'_2 = s$  。

(vii) 设  $U$  为  $X$  的开子集。将精确函子  $\mathrm{Hom}_{\mathcal{R}}(\cdot,H)$  应用于精确序列  $0\to G_U\to G\to G_{X\setminus U}\to0$  ，我们根据推论 2.4.2 得到结果。 □

\statementlabel{命题 2.4.7}令  $0 \to F' \to F \to F'' \to 0$  为  $\mathfrak{S h}(X)$  中的精确序列。假设 $F'$  是松弛的。那么序列 $0 \to \Gamma(X; F') \to \Gamma(X; F) \to \Gamma(X; F'') \to 0$  是精确的。

\statementlabel{证明}令  $s'' \in \Gamma(X; F'')$  ，并令  $\mathfrak{S}$  为  $(U, s)$  的对的集合，使得  $U$  在  $X$  、  $s \in \Gamma(U; F)$  中打开，并且  $s$  被发送到  $s''|_{U}$  。我们通过设置  $(U, s) \leq (V, t)$  if  $U \subset V$  和  $t|_{U} = s$  来排序  $\mathfrak{S}$  。那么 $\mathfrak{S}$  显然是归纳排序的。令  $(U, s)$  为最大元素并假设  $U \neq X$  。让 $x \in X \setminus U$  。存在 $x$  的开放邻域 $V$  和 $t \in \Gamma(V; F)$  部分，使得 $t$  被发送到 $s''|_{V}$  。在 $U \cap V$  上， $s - t$  属于 $\Gamma(U \cap V; F')$  。让  $r \in \Gamma(X; F')$  在  $U \cap V$  上扩展  $s - t$  。将  $t$  替换为

 $t-r$  ，我们可以在 $U\cap V$  上假设 $t=s$  。因此  $s$  可以扩展到  $U\cup V$  ，这是一个矛盾。   $\square$ 

\statementlabel{推论 2.4.8}在命题2.4.7的情况下，设Z是X的局部闭子集，则序列：

\[
0\to\varGamma_{Z}(X;F^{\prime})\to\varGamma_{Z}(X;F)\to\varGamma_{Z}(X;F^{\prime\prime})\to0
\]

and

\[
0\to\varGamma_{Z}(F^{\prime})\to\varGamma_{Z}(F)\to\varGamma_{Z}(F^{\prime \prime})\to0
\]

是准确的。

\statementlabel{证明}对于开集 $U$ ，使得 $U \cap Z$  在 $U$  中闭集，我们有一个交换图：

\begin{center}
\includegraphics[width=0.76\textwidth]{流行上的层_Chn-assets/img_in_image_box_106_624_1019_1212.png}
\end{center}

根据命题 2.4.7，这里第二行和第三行以及所有列都是精确的。因此顶行是准确的（参见练习 I.8）。   $ \square $ 

\statementlabel{推论 2.4.9}令  $ 0 \to F' \to F \to F'' \to 0 $  为  $ \mathfrak{Sh}(X) $  中的精确序列。假设 $ F' $  和 $ F $  是松弛的。然后 $ F'' $ 就松弛了。

\statementlabel{证明}令  $U$  为  $X$  的开子集。在下图中， $\alpha$  和  $\beta$  是满射：

\[
\begin{array}{c}\Gamma(X;F)\longrightarrow\Gamma(X;F^{\prime \prime})\\\left|\begin{array}{cc}\alpha&\\&\end{array}\right.\quad\left|\begin{array}{c}\gamma\\ \\\end{array}\right.\\\Gamma(U;F)\xrightarrow{\quad\beta\quad}\Gamma(U;F^{\prime \prime})\quad.\end{array}
\]

因此 $ \gamma $  是满射的。 ⑨

根据定义 1.8.2，我们发现由松弛层组成的  $\mathfrak{Sh}(X)$  的完整子类相对于函子  $\Gamma(X;\cdot),\Gamma_{Z}(\cdot),f_{*}$  是单射的。

单射或松弛的属性是局部属性。

\statementlabel{命题 2.4.10}让 $X = \bigcup_{i \in I} U_{i}$ 成为一个开放的覆盖物。

(i) 让 $ F \in \mathrm{Ob}(\mathfrak{S h}(X)) $  。如果  $ F|_{U_i} $  对于所有  $ i $  都是松弛的，那么  $ F $  也是松弛的。

(ii) 令 $ \mathcal{R} $  为 $ X $  上的环层，并令 $ F \in \mathrm{Ob}(\mathrm{Mod}(\mathcal{R})) $  。如果对于所有 $ i $ ， $ F|_{U_i} $  是 $ \mathcal{R}|_{U_i} $  -内射， $ F $  是 $ \mathcal{R} $  -内射。

\statementlabel{证明}为了证明(i)，只需证明

(2.4.1) 层 F 是松弛的当且仅当对于任何开集  $U$  ，  $F \to \Gamma_U(F)$  是外同态。

由于松弛意味着 $F \mapsto \Gamma_U(F)$  是一个外态，所以我们将展示相反的情况。

令  $s_{\circ}$  成为  $F$  上  $X$  的开放子集  $V_{\circ}$  的一部分。为了证明 $s_{\circ}$ 扩展到全局部分，让 $\mathfrak{S}$ 成为开子集 $V$ 和 $s \in \Gamma(V; F)$ 的对 $(s, V)$ 族。我们通过设置  $(s, V) \leq (s', V')$  if  $V \subset V'$  和  $s'|_{V} = s$  来排序  $\mathfrak{S}$  。然后 $\mathfrak{S}$  被归纳排序。因此，存在最大元素  $(s, V)$  使得  $(s, V) \geq (s_{\circ}, V_{\circ})$  。我们将展示 $V = X$ 。否则就会有 $x \in X \setminus V$ 。由于  $F_x \to \Gamma_V(F)_x$  是满射的，因此将存在  $x$  和  $t \in \Gamma(W; F)$  的开邻域  $W$  ，使得  $t|_{V \cap W} = s|_{V \cap W}$  。那么就会存在  $s' \in \Gamma(W \cup V; F)$  这样  $s'|_{W} = t$  和  $s'|_{V} = s$  。这与  $(s, V)$  的最大值相矛盾。因此 $V = X$ ，我们得到(2.4.1)。

(ii) 函子  $\mathrm{Hom}_{\mathscr{R}}(\cdot,F)$  是精确的，并且对于  $\mathscr{R}$ -模  $G$  的任何层， $\mathrm{Hom}_{\mathscr{R}}(G,F)$  层是松弛的。因此函子  $\mathrm{Hom}_{\mathscr{R}}(\cdot,F)=\Gamma(X;\cdot)\circ\mathrm{Hom}_{\mathscr{R}}(\cdot,F)$  是精确的。   $\square$ 

\statementlabel{定义 2.4.11}让  $ \mathcal{R} $  成为  $ X $  、  $ F \in \text{Ob}(\mathfrak{Mod}(\mathcal{R})) $  上的环层。如果函子  $ \cdot \otimes_{\mathcal{R}} F $  （从右  $ \mathcal{R} $ -模 到  $ \mathfrak{Sh}(X) $  的范畴）是精确的，则可以说  $ F $  是  $ \mathcal{R} $  -flat。

当不存在混淆风险时，我们将说“平坦”而不是“ $\mathscr{R}$  -平坦”。根据 (2.2.7)， $F$  是  $\mathscr{R}$  -flat 当且仅当  $F_x$  是所有  $x \in X$  的  $\mathscr{R}_x$  -flat 模。

\statementlabel{命题 2.4.12}让 $ F \in \text{Ob}(\mathfrak{Mod}(\mathscr{R})) $  。然后存在一个外同态  $ P \to F $  ，其中  $ P $  是  $ \mathscr{R} $  -flat。

\statementlabel{证明}让  $\mathfrak{S}$  表示对族  $(U,s)$  ，其中  $U$  在  $X$  和  $s\in\Gamma(U;F)$  中开。对于  $(U,s)\in\mathfrak{S}$  ，我们用  $\mathcal{R}(U,s)$  表示由  $(U,s)$  索引的层  $\mathcal{R}_{U}$  。然后我们设置：

\[
P=\bigoplus_{(U,s)\in\mathfrak{S}}\mathcal{R}(U,s)\enspace.
\]

态射链  $ \mathscr{R}_{U} \to F_{U} \to F $  ，其中  $ 1 \in \Gamma(U; \mathscr{R}_{U}) $  部分被发送到

 $s \in \Gamma(U; F_U)$  ，定义了表态  $P \to F$  ，并且  $P$  是平坦的，因为对于每个  $x \in X$  ，  $P_x$  是一个免费的  $\mathscr{R}_x$  模。 □

\statementlabel{命题 2.4.13}令  $ 0 \to F' \to F \to F'' \to 0 $  为  $ \mathfrak{Mod}(\mathscr{R}) $  中的精确序列。如果 $ F $  和 $ F'' $  是 $ \mathscr{R} $  -平坦，那么 $ F' $  也是平坦的

\statementlabel{证明}对于所有  $ x \in X $  ，这紧随  $ \mathscr{R}_x $  -模的相应属性而来。   $ \square $ 

请注意，根据前面的两个结果，对于任何右  $\mathscr{R}$  -模  $G$  ，范畴  $\operatorname{Mod}(\mathscr{R})$  相对于函子  $G\otimes_{\mathscr{R}}\cdot$  具有足够的射影对象。

\subsection{局部紧空间上的层}
当拓扑空间 X 具有一些有限性属性时，例如局部紧空间，新的层类和新函子是令人感兴趣的。在本节中，除非另有说明（例如在命题 2.5.1 和备注 2.5.3 中），所有空间都应该是局部紧的，特别是 Hausdorff。

\statementlabel{命题 2.5.1}设 X 是一个（不一定是局部紧致的）拓扑空间，Z 是一个子空间，F 是 X 上的层。考虑规范态射：

\[
\psi:\varinjlim_{U}\varGamma(U;F)\to\varGamma(Z;F)\enspace,
\]

其中 U 的范围穿过 X 中 Z 的开邻域族。

(i) 态射 $ \psi $  是单射的。

(ii) 假设 X 是 Hausdorff 并且 Z 是紧的。则 ψ 是同构。

(iii) 假设 X 是仿紧且 Z 是闭的。那么 $\psi$  是同构。

在进入证明之前，回想一下，仿紧空间  $X$  是一个 Hausdorff 空间，使得对于  $X$  的每个开覆盖  $(U_i)_{i \in I}$  ，都存在  $X$  的开覆盖  $(V_j)_{j \in J}$  ，比  $(U_i)_{i \in I}$  更精细（即对于每个  $j \in J$  都存在  $i \in I$  使得  $V_j \subset U_i$  ）并且局部有限。另请记住，如果  $X$  是仿紧且  $(U_i)_{i \in I}$  是开局部有限覆盖，则存在一个开覆盖  $(V_i)_{i \in I}$  使得  $\bar{V}_i \subset U_i$  对于所有  $i$  。仿紧空间的闭子空间是仿紧的，无穷大可数的局部紧空间以及度量空间是仿紧的。

命题证明 2.5.1。 (i) 如果  $ \Gamma(Z; F) $  中的  $ s \in \Gamma(U; F) $  为零，则所有  $ x \in Z $  中的  $ s_x = 0 $  为零。因此  $ s $  在  $ Z $  的开邻域上为零。

(ii) 和 (iii) 让  $ s \in \Gamma(Z; F) $  。那么存在一系列开集  $ \{U_i\}_{i \in I} $  和  $ s_i \in \Gamma(U_i; F) $  使得  $ s|_{U_i \cap Z} = s_i|_{U_i \cap Z} $  和  $ \bigcup U_i \supset Z $  。在情况 (ii) 中，我们可以假设  $ I $  是有限的，在情况 (iii) 中，我们可以假设  $ \{U_i\} $  是局部有限覆盖

 $X$ 。在任何一种情况下，我们都可以找到一系列开子集 $\{V_i\}_{i\in I}$ ，使得 $\bar{V}_i\subset U_i$ 、 $\{\bar{V}_i\}$ 是局部有限的并且 $\bigcup V_i\supset Z$ 。对于 $x\in X$  设置 $I(x)=\{i\in I;x\in\bar{V}_i\}$  并定义

\[
W=\left\{x\in\bigcup V_{i};s_{i x}=s_{j x}\mathrm{f o r a n y}i,j\in I(x)\right\}\;.
\]

那么  $I(x)$  是一个有限集，每个  $x$  都有一个邻域  $W_x$  ，使得  $I(y) \subset I(x)$  对于任何  $y \in W_x$  。因此 $W$  是开放的， $W$  通过其构造包含 $Z$ 。由于  $s_i|_{W \cap V_i \cap V_j} = s_j|_{W \cap V_i \cap V_j}$  ，存在  $\tilde{s} \in \Gamma(W; F)$  使得  $\tilde{s}|_{W \cap V_i} = s_i|_{W \cap V_i}$  。那么  $\tilde{s}$  满足  $\psi(\tilde{s}) = s$  。   $\square$ 

令 $f: Y \to X$  为连续映射（X 和Y 不一定是局部紧凑的）。回想一下，如果  $f$  是封闭的（即：  $Y$  的任何封闭子集的图像在  $X$  中是封闭的），并且它的纤维是相对豪斯多夫（纤维中的两个不同点在  $Y$  中具有不相交的邻域）和紧凑的，那么  $f$  是正确的。如果 $X$  和 $Y$  是局部紧致的，则 $f$  当且仅当 $X$  的任何紧致子集的逆像是紧致的。让 $G$  成为 $Y$  上的一个层。通过设置 $U$  open in  $X$  来定义 $f_*G$  的子层 $f_!G$  ：

\[
\varGamma(U;f,G)=\{s\in\varGamma(f^{-1}(U);G);f:\operatorname{s u p p}(s)\to U\mathrm{~i s~p r o p e r}\}\enspace.
\]

由于真性在  $X$  上是局部的，因此显然由 (2.5.1) 定义的预层是  $f_*(G)$  的子束。该束称为具有  $G$  真支撑直像。其中一个用  $f_!$  表示从  $\mathfrak{Sh}(Y)$  到  $\mathfrak{Sh}(X)$  的函子  $G \mapsto f_!G$  。如果  $\mathcal{R}$  是  $X$  上的环层，则该函子会导出一个函子，仍用  $f_!,$  从  $\mathfrak{Mod}(f^{-1}\mathcal{R})$  到  $\mathfrak{Mod}(\mathcal{R})$  表示。函子  $f_!$  显然是精确的。还规定：

\[
\varGamma_{c}(X;F)=\{s\in\varGamma(X;F);\mathrm{s u p p}(s)\mathrm{i s c o m p a c t a n d H a u s d o r f f}\}~.
\]

如果  $ a_X $  是映射：  $ X \to \{\text{pt}\} $  ，我们找到  $ \Gamma_c(X; F) \simeq a_X! F $  。

让 $ g: Z \to Y $  成为另一个连续映射。我们有：

\[
f_{!}\circ g_{!}=(f\circ g)_{!}.
\]

特别是，我们有：

\[
\Gamma_{c}(X;f_{!}G)\simeq\Gamma_{c}(Y;G)\enspace.
\]

\statementlabel{命题 2.5.2}令 X 和 Y 为局部紧空间（特别是豪斯多夫空间）， $f: Y \to X$  为连续映射， $G$  为  $Y$  上的层。然后对于  $x \in X$  ，规范态射：

\[
\alpha:(f_{!}G)_{x}\rightarrow\Gamma_{c}(f^{-1}(x);G|_{f^{-1}(x)})
\]

是同构。

\statementlabel{证明}让我们首先证明 $\alpha$  是单射的。令  $V$  成为  $x\in X$  的开放邻域，并令  $t\in\Gamma(V;f_!G)$  。然后  $t$  由节  $s\in\Gamma(f^{-1}(V);G)$  定义

这样  $\mathrm{supp}(s) \to V$  是正确的。如果  $\alpha(t) = 0$  ，则  $\mathrm{supp}(s) \cap f^{-1}(x) = \emptyset$  和  $x \notin f(\mathrm{supp}(s))$  。最后一组是封闭的，存在  $x$  的开邻域，其上有  $t = 0$  。

接下来我们证明 $\alpha$  是满射的。让  $s \in \Gamma_c(f^{-1}(x); G|_{f^{-1}(x)})$  ，并放置  $K = \text{supp}(s)$  。根据命题 2.5.1， $X$  和  $t \in \Gamma(U; G)$  中存在  $K$  的开邻域  $U$  ，使得  $t|_{\kappa} = s|_{\kappa}$  。通过缩小 $U$ ，我们可以假设 $t|_{U \cap f^{-1}(x)} = s|_{U \cap f^{-1}(x)}$ 。让  $V$  成为  $K$  的相对紧凑的开放邻域，以及  $\bar{V} \subset U$  。由于  $x$  不属于  $f(\bar{V} \cap \text{supp}(t) \setminus V)$  ，因此存在  $x$  的开邻域  $W$  ，使得  $f^{-1}(W) \cap \bar{V} \cap \text{supp}(t) \subset V$  。通过设置定义 $\tilde{s} \in \Gamma(f^{-1}(W); G)$ ：

\[
\begin{array}{r}{\left.\tilde{s}\right|_{f^{-1}(\bar{W})\setminus(\mathsf{s u p p}(t)\cap\bar{V})}=0\mathrm{~,~}}\end{array}
\]

\[
\tilde{s}|_{f^{-1}(W)\cap V}=t|_{f^{-1}(W)\cap V}.
\]

由于  $ \operatorname{supp}(\tilde{s}) $  包含在  $ f^{-1}(W) \cap \operatorname{supp}(t) \cap \bar{V} $  中，因此  $ f $  在此集合上是真映射。此外，我们还有 $ \tilde{s}|_{f^{-1}(x)} = s $ 。   $ \Box $ 

\statementlabel{备注 2.5.3}令 $f: Y \to X$  为连续映射， $Y$  和 $X$  不一定是局部紧凑的，并令 $G$  为 $Y$  上的一束。假设  $f$  在  $supp(G)$  上是真映射。通过与命题 2.5.2 的证明类似的论证，我们发现对于  $x \in X$  ，自然态射：

\[
(f_{*}{\cal G})_{x}\to\varGamma(f^{-1}(x);{\cal G}|_{f^{-1}(x)})
\]

是同构。

在下面的命题 2.5.4 中，我们不假设 X 是局部紧的。

\statementlabel{命题 2.5.4}令  $Z$  为  $X$  的局部闭子集， $i: Z \hookrightarrow X$  包含图。

(i) 函子 $i_{!}$  是精确的。

(ii) 让 $ F \in \mathrm{Ob}(\mathfrak{S}\mathfrak{h}(X)) $  。然后：

\[
F_{Z}\simeq i_{!}\circ i^{-1}(F)\enspace.
\]

\statementlabel{证明}(i)  $ (i_!F)_x \simeq F_x $  或 0 根据  $ x \in Z $  或  $ x \notin Z $  。

(ii) 我们有  $ (i_! \circ i^{-1}F)|_{X \setminus Z} = 0 $  和  $ (i_! \circ i^{-1}F)|_Z = i^{-1}F $  。然后应用命题 2.3.6 (i)。   $ \Box $ 

在本节结束之前，所有空间将再次被假定为局部紧空间。

\statementlabel{定义 2.5.5}让 $ F \in \mathrm{Ob}(\mathfrak{S}\mathfrak{h}(X)) $  。如果对于 X 的任何紧子集 K，限制态射  $ \Gamma(X; F) \to \Gamma(K; F) $  是满射的，则可以说 F 是 c-soft。

根据命题 2.5.1，松弛层，特别是射射层，是 c-软的。在第 9 节中，我们将给出许多 c-软层的例子。

\statementlabel{命题 2.5.6}让 $F \in \mathrm{Ob}(\mathfrak{S}\mathfrak{h}(X))$  。那么  $F$  是  $c$  -soft 当且仅当对于  $X$  的任何闭子集  $Z$  ，限制态射  $\Gamma_c(X; F) \to \Gamma_c(Z; F|_Z)$  是满射的。

\statementlabel{证明}如果  $K$  是紧凑的，则  $\Gamma(K; F) = \Gamma_c(K; F|_K)$  。这就证明了条件的充分性。

相反，假设  $F$  是  $c$  -soft 并让  $s \in \Gamma(Z; F|_{Z})$  具有紧凑支持  $K$  。令 $U$  为 $X$  中 $K$  的相对紧凑的开邻域。通过设置  $\tilde{s}|_{Z \cap \overline{U}} = s$  、  $\tilde{s}|_{\partial U} = 0$  来定义  $\tilde{s} \in \Gamma(\partial U \cup (Z \cap \overline{U}); F)$  ，并将  $\tilde{s}$  扩展到节  $t \in \Gamma(X; F)$  。由于  $t = 0$  在  $\partial U$  的邻域上，我们可以假设  $t$  由  $\overline{U}$  支持。   $\square$ 

\statementlabel{命题 2.5.7}设 F 为 X 上的 c-软束。

(i) 令  $Z$  为  $X$  的局部闭子集。那么 $F|_{Z}$  就是c-soft。

(ii) 令 $f: X \to Y$  为连续映射。那么 $f_! F$  就是 $c$  -soft。

(iii) 对于 (i) 中的 Z， $ F_{z} $  是 c-soft。

\statementlabel{证明}(i) 如果 Z 开路，则结果明确。如果 Z 是闭合的，则应用命题 2.5.6。

(ii) 如果  $K$  是  $Y$  的紧凑子集，则  $\Gamma(K; f_1F) = \Gamma_c(f^{-1}(K); F)$  。自  $\Gamma_c(Y; f_1F) = \Gamma_c(X; F)$  以来，结果来自命题 2.5.6。

(iii) 从 (i) 和 (ii) 得出，因为  $ F_Z = f_1(F|_Z) $  、  $ f $  表示包含图  $ Z \hookrightarrow X $  。   $ \square $ 

\statementlabel{命题 2.5.8}令 $0 \to F' \to F \to F'' \to 0$  为 $X$  上层的精确序列，并假设 $F'$  为 $c$  -soft。令 $f: X \to Y$  为连续映射。那么序列 $0 \to f_1 F' \to f_1 F \to f_1 F'' \to 0$  是精确的。特别是序列  $0 \to \Gamma_c(X; F') \to \Gamma_c(X; F) \to \Gamma_c(X; F'') \to 0$  是精确的。

\statementlabel{证明}对于所有  $y \in Y$  ，束  $F'(f^{-1}(y))$  是  $c$  -soft on  $f^{-1}(y)$  。因此，根据命题 2.5.2 足以证明特定情况  $f: X \to \{\mathrm{pt}\}$  的结果。

令 $s'' \in \Gamma_c(X; F'')$  和 $U$  成为 $\mathrm{supp}(s'')$  的相对紧凑的开邻域。我们将证明  $s''$  位于  $\Gamma_c(X; F) \to \Gamma_c(X; F'')$  的图像中。通过用 $F'_U, F_U, F''_U$ 替换 $F', F$ 和 $F''$ ，然后用 $\bar{U}$ 替换 $X$ ，我们可以假设 $X$ 是紧凑的。对于  $s'' \in \Gamma(X; F''),$  令  $\{K_i\}_{i=1,\ldots,n}$  为  $X$  的紧凑子集的有限覆盖，使得存在  $s_i \in \Gamma(K_i; F)$  其图像为  $s''|_{\mathcal{K}_i}$  。让我们通过 $n$  进行归纳论证。对于  $n \geq 2$  ，  $K_1 \cap K_2, s_1 - s_2$  上定义了  $\Gamma(K_1 \cap K_2; F')$  的元素，因此扩展到  $s' \in \Gamma(X; F')$  。将  $s_2$  替换为  $s_2 + s',$  我们可以假设  $s_1|_{\mathcal{K}_1 \cap \mathcal{K}_2} = s_2|_{\mathcal{K}_1 \cap \mathcal{K}_2}$  。因此存在 $t \in \Gamma(K_1 \cup K_2; F)$ 使得 $t|_{\mathcal{K}_i} = s_i$  ( $i = 1,2$ )。归纳就这样进行。   $\square$ 

\statementlabel{推论 2.5.9}令  $ 0 \to F' \to F \to F'' \to 0 $  为  $ \mathfrak{Sh}(X) $  中的精确序列。假设 $ F' $  和 $ F $  是 $ c $  -soft。那么 $ F'' $  就是 $ c $  -soft。

\statementlabel{证明}与推论 2.4.9 类似。 □

通过前面的结果，我们发现  $X$  上的  $c$  -软层的范畴对于函子  $\mathcal{F}_c(X;\cdot)$  、  $f_!$  和  $\mathcal{F}(K;\cdot)$  是单射的（ $K$  紧凑）。当 $X$ 在无穷大时可数时，我们得到更好的结果。

\statementlabel{命题 2.5.10}假设  $X$  是局部紧致的并且在无穷大处可数。那么  $c$  -soft sheaves 的范畴相对于函子  $\Gamma(X;\cdot)$  是单射的。

\statementlabel{证明}令  $0 \to F' \to F \to F'' \to 0$  为  $\mathfrak{Sh}(X)$  中的精确序列，使用  $F'c$  -soft。令  $\{K_n\}_{n \in \mathbb{N}}$  为  $X$  的紧凑子集的递增序列，例如  $X = \bigcup_n K_n$  和  $K_n \subset \mathrm{Int}(K_{n+1})$  对于所有  $n$  。根据命题 2.5.7 和 2.5.8，所有序列  $0 \to \Gamma(K_n; F') \to \Gamma(K_n; F) \to \Gamma(K_n; F'') \to 0$  都是精确的。由于  $\Gamma(K_{n+1}; F') \to \Gamma(K_n; F')$  对于所有  $n$  都是满射，我们可以应用命题 1.12.3 和序列：

\[
0\to\varprojlim_{n}\varGamma(K_{n};F^{\prime})\to\varlimsup_{n}\varGamma(K_{n};F)\to\varlimsup_{n}\varGamma(K_{n};F^{\prime \prime})\to0
\]

是准确的。因为对于  $X$  、  $\Gamma(X;G) \simeq \varprojlim_n \Gamma(K_n;G)$  上的任何束  $G$  ，证明都是完整的。 □

让我们研究一下函子 $f_{1}$  和之前介绍的其他函子之间的一些关系。首先考虑局部紧空间的笛卡尔平方：

\[
\begin{array}{c}Y^{\prime}\xrightarrow{f^{\prime}}X^{\prime}\\g^{\prime}\left|\begin{array}{c}\Box\\\hline Y\end{array}\right.\left|\begin{array}{c}g\end{array}.\right.\\f\xrightarrow{f}X\end{array}
\]

回想一下，这意味着该图是可交换的，并且  $ Y' $  作为拓扑空间与纤维积  $ Y \times_X X' = \{(y, x') \in Y \times X'; \, f(y) = g(x')\} $  同构。 （图中的正方形是为了强调该图是笛卡尔图。）

\statementlabel{命题 2.5.11}函子具有规范的同构：

\[
g^{-1}\circ f_{!}\stackrel{}{\sim}f_{!}^{\prime}\circ g^{\prime-1}\quad.
\]

\statementlabel{证明}首先我们将构造规范态射：

\[
f_{!}\circ g_{*}^{\prime}\to g_{*}\circ f_{!}^{\prime}~.
\]

令  $G \in \mathrm{Ob}(\mathfrak{S}\mathfrak{h}(Y'))$  和  $V$  成为  $X$  的开子集。节 $t \in \Gamma(V; f_! \circ g_*^*G)$  由节 $s \in \Gamma((f \circ g')^{-1}(V); G)$  定义，使得 $\mathrm{supp}(s) \subset g'^{-1}(Z)$  为 $f^{-1}(V)$  的子集 $Z$  在 $V$  上的正确子集。那么  $g'^{-1}(Z) \to g^{-1}(V)$  是真映射，  $s$  定义了  $g_* f'_G$  的一部分。这定义了(2.5.7)。

为了证明 (2.5.6)，请考虑  $ G \in \text{Ob}(\mathfrak{Sh}(Y)) $  。根据命题 2.3.3 我们有：

\[
\mathrm{H o m}(g^{-1}\circ f_{!}G,f_{!}^{\prime}\circ g^{\prime-1}G)=\mathrm{H o m}(f_{!}G,g_{*}\circ f_{!}^{\prime}\circ g^{\prime-1}G)\enspace.
\]

态射 $f_! \to f_! \circ g'_* \circ g'_{'-1} \to g_* \circ f'_! \circ g'_{'-1}$  引发态射(2.5.6)。

为了证明这是一个同构，让我们采用  $ x' \in X' $  。然后：

\[
\begin{aligned}{(g^{-1}\circ f_{!}G)_{x^{\prime}}}&{{}=(f_{!}G)_{g(x^{\prime})}}\\ {}&{{}=\varGamma_{c}(f^{-1}(g(x^{\prime}));G).}\\ \end{aligned}
\]

映射  $g'$  引发同构  $f^{\prime-1}(x^{\prime}) \simeq f^{-1}(g(x^{\prime}))$  和同构  $\Gamma_c(f^{-1}(g(x^{\prime})); G) \simeq \Gamma_c(f^{\prime-1}(x^{\prime}); g^{\prime-1}G) \simeq (f_!^{\prime} \circ g^{\prime-1}G)_x^v$  。

最后，我们将研究 $f_{i}$ 与张量积的关系。

\statementlabel{引理 2.5.12}设 A 为环，M 为扁平 A 模。令 F 为一​​束右 $ A_{x} $  模。那么就有一个自然的同构：

\[
\varGamma_{c}(X;F)\mathop{\otimes_{A}}\strut M\simeq\varGamma_{c}\biggl(X;F\mathop{\otimes_{A_{X}}}\strut M_{X}\biggr).
\]

特别是，如果 $F$  是c-soft，则 $F \otimes_{A_X} M_X$  是c-soft。

\statementlabel{证明}我们可以假设 $X$  是紧凑的。如果  $X = \bigcup_{j} K_{j}$  是  $X$  的紧凑子集的有限覆盖，我们有一个精确的序列：

\[
0\longrightarrow\varGamma(X;F)\xrightarrow{\lambda}\bigoplus_{j}\varGamma(K_{j};F)\xrightarrow{\mu}\bigoplus_{j,k}\varGamma(K_{j}\cap K_{k};F)\enspace.
\]

由于  $M$  是平坦的，因此在应用函子  $\otimes_A M$  后该序列仍然准确。因此我们有一个包含精确行的交换图：

\[
\begin{aligned}0&\xrightarrow{\quad}\quad\varGamma(X;F)\otimes_{A}M\xrightarrow{\quad\lambda\quad}\quad\bigoplus_{j}\varGamma(K_{j};F)\otimes_{A}M\xrightarrow{\quad\mu\quad}\quad\bigoplus_{j,k}\varGamma(K_{j}\cap K_{k};F)\otimes_{A}M\\&\quad\qquad\alpha\bigg\downarrow\quad\beta\bigg\downarrow\quad\gamma\bigg\downarrow\\0&\xrightarrow{\quad}\quad\varGamma\bigg(X;F\otimes_{A x}M_{\bar{X}}\bigg)\xrightarrow{\quad\lambda^{\prime}\quad}\bigoplus_{j}\varGamma\bigg(K_{j};F\otimes_{A x}M_{\bar{X}}\bigg)\xrightarrow{\mu^{\prime}}\bigoplus_{j,k}\varGamma\bigg(K_{j}\cap K_{k};F\otimes_{A x}M_{\bar{X}}\bigg)\end{aligned}
\]

让我们首先证明 $ \alpha $  是单射的。我们有  $ x \in X $  的同构：

\[
\operatorname*{l i m}_{\stackrel{\rightarrow}{U}}\left(\varGamma(U;F)\mathop{\otimes}_{\mathcal{A}}M\right)\stackrel{\sim}{\to}\operatorname*{l i m}_{\stackrel{\rightarrow}{U}}\varGamma\left(U;F\mathop{\otimes}_{\mathcal{A}_{X}}M_{X}\right),
\]

其中  $U$  范围贯穿  $x$  的开放邻域族：事实上 (2.5.9) 的两边都与  $F_x \otimes_A M$  同构。因此，如果  $s \in \Gamma(X; F) \otimes_A M$  满足  $\alpha(s) = 0$  ，我们可以找到一个有限覆盖  $X = \bigcup_j K_j$  使得  $\lambda(s) = 0$  。然后 $s = 0$ 。如果我们将此结果应用于  $K_j$  和  $K_j \cap K_k$  而不是  $X$  ，我们会发现  $\beta$  和  $\gamma$  是单射的。

为了证明  $\alpha$  是满射的，让我们取  $t \in \Gamma(X; F \otimes_{A_x} M_x)$  。到 (2.5.9) 为止，存在有限覆盖  $X = \bigcup_j K_j$  ，使得  $\lambda'(t)$  位于  $\beta$  的图像中。那么  $y$  的单射性意味着  $t$  在  $\alpha$  的图像中。   $\square$ 

\statementlabel{命题 2.5.13}令 $f: Y \to X$  为连续映射，让 $\mathcal{R}$  为 $X$  上的环层，让 $F \in \mathrm{Ob}(\mathfrak{Mod}(\mathcal{R}))$  并让 $G \in \mathrm{Ob}(\mathfrak{Mod}(f^{-1}\mathcal{R}^{op}))$  。

(i) 存在自然态射：

\[
f_{!}G\mathop{{\otimes}_{\mathring{\mathcal{R}}}}F\to f_{!}\biggl(G\mathop{{\otimes}_{{f^{-1}}\mathring{\mathcal{R}}}}f^{-1}F\biggr).
\]

(ii) 假设 F 是一个平坦  $ \mathcal{R} $  模。那么(2.5.10)是同构。

\statementlabel{证明}首先请注意，态射 (2.5.10) 是由态射 (2.3.21) 引发的。现在假设 F 是平坦的。为了证明 (2.5.10) 是同构，让我们选择  $ x \in X $  ，并应用命题 2.5.2 和引理 2.5.12。我们得到：

\[
\begin{aligned}{\left(f_{!}\bigg(G\mathop{\otimes}_{f^{-1}\mathcal{R}}f^{-1}F\bigg)\right)_{x}}&{{}\simeq\varGamma_{c}\bigg(f^{-1}(x);G\mathop{\otimes}_{f^{-1}\mathcal{R}}f^{-1}F\bigg)}\\ {}&{{}\simeq\varGamma_{c}\bigg(f^{-1}(x);G\mathop{\otimes}_{(\mathcal{R}_{x})_{Y}}(F_{x})_{Y}\bigg)}\\ {}&{{}\simeq\varGamma_{c}(f^{-1}(x);G)\mathop{\otimes}_{\mathcal{R}_{x}}F_{x}}\\ {}&{{}\simeq(f_{!}G)_{x}\mathop{\otimes}_{\mathcal{R}_{x}}F_{x}}\\ {}&{{}\simeq\bigg(f_{!}G\mathop{\otimes}_{\mathcal{R}}F\bigg)_{x}\enspace.\enspace\Box}\\ \end{aligned}
\]

当然，利用命题2.5.13，我们可以获得新的自然态射。例如，如果 $ G_1 $ 和 $ G_2 $ 表示两个 $ f^{-1}R $  -模，我们就有态射（当考虑张量积时， $ G_1 $ 将是一个正确的模）：

\[
\begin{aligned}f_{!}G_{1}\otimes_{\mathcal{R}}f_{*}G_{2}&\xrightarrow{\quad}f_{!}\bigg(G_{1}\otimes_{f^{-1}\mathcal{R}}G_{2}\bigg)\;,\\f_{!}\mathcal{H}om_{f^{-1}\mathcal{R}}(G_{1},G_{2})&\xrightarrow{\quad}&\mathcal{H}om_{\mathcal{R}}(f_{*}G_{1},f_{!}G_{2})\\\downarrow&\downarrow&\\f_{*}\mathcal{H}om_{f^{-1}\mathcal{R}}(G_{1},G_{2})&\xrightarrow{\quad}&\mathcal{H}om_{\mathcal{R}}(f_{!}G_{1},f_{!}G_{2})\end{aligned}
\]

公式(2.5.11)紧随(2.5.10)和(2.3.4)得出。让我们解释一下（2.5.12）中的顶部箭头。设置 $ H = \mathcal{H}_{om_{f^{-1}}\mathcal{R}}(G_1, G_2) $ 。然后还有自然态射：

\[
f_{!}H\otimes f_{*}G_{1}\to f_{!}(H\otimes G_{1})\to f_{!}G_{2}
\]

命题 2.2.9 给出了顶部箭头。 (2.5.12)中的其他箭头也有类似的定义。请注意，如果  $ \mathcal{R} $  是可交换的，则这些态射是  $ \mathcal{R} $  -线性的。

\subsection{层上同调}
我们将把第一章的结果应用到层范畴中。令 $X$  为拓扑空间， $\mathcal{R}$  为 $X$  上的环层。  $X$  上（左） $\mathcal{R}$  -模的层的范畴 $\mathfrak{Mod}(\mathcal{R})$  是阿贝尔范畴。因此，我们可以考虑导出范畴​​  $\mathbf{D}(\mathfrak{Mod}(\mathcal{R}))$  及其满三角子范畴  $\mathbf{D}^*(\mathfrak{Mod}(\mathcal{R}))$  ，其中  $* = +, -, b$  。我们简单写一下：

\[
\mathbf{D}^{\ast}(\mathcal{R})=\mathbf{D}^{\ast}(\mathfrak{M o b}(\mathcal{R}))~,\qquad^{\ast}=\mathcal{O},+,-,b~.
\]

特别是，如果  $A$  是环，则范畴  $\mathbf{D}(A_X)$  是  $X$  上  $A$  -模的层范畴的导出范畴。例如 $\mathbf{D}(\mathbb{Z}_X) = \mathbf{D}(\mathfrak{S}\mathfrak{h}(X))$ 。

为了定义前面介绍的函子的导出函子，我们考虑以下情况： $Z$ 是 $X$ 的局部闭子集， $f: Y \to X$ 是连续映射， $g: W \to Y$ 是连续映射， $\mathcal{S} \to \mathcal{R}$ 是 $X$ 上环的层态射，其图像包含在 $\mathcal{R}$ 的中心，  $\mathcal{S}$  是可交换的。当考虑函子  $f_!$  （或  $g_!$  等）时，我们将始终隐式假设所有空间都是局部紧凑的。

由于范畴 $ \operatorname{Mod}(\mathscr{R}) $ 有足够的单射，我们可以导出所有左精确函子。我们得到函子：

\[
\begin{array}{l l}{R\varGamma_{Z}(X;\cdot)}&{:\mathbf{D}^{+}(\mathcal{R})\to\mathbf{D}^{+}(\mathfrak{U b})\;,}\\ {}&{}\\ {R\varGamma_{Z}(\cdot)}&{:\mathbf{D}^{+}(\mathcal{R})\to\mathbf{D}^{+}(\mathcal{R})\;,}\\ {}&{}\\ {R\varGamma(Z;\cdot)}&{:\mathbf{D}^{+}(\mathcal{R})\to\mathbf{D}^{+}(\mathfrak{U b})\;,}\\ {}&{}\\ {R f_{*}}&{:\mathbf{D}^{+}(f^{-1}\mathcal{R})\to\mathbf{D}^{+}(\mathcal{R})\;,}\\ {}&{}\\ {R\mathcal{H}o m_{\mathcal{R}}(\cdot,\cdot):\mathbf{D}^{-}(\mathcal{R})^{\circ}\times\mathbf{D}^{+}(\mathcal{R})\to\mathbf{D}^{+}(\mathcal{S})\;,}\\ {}&{}\\ {R\mathcal{H}o m_{\mathcal{S}}(\cdot,\cdot):\mathbf{D}^{-}(\mathcal{S})^{\circ}\times\mathbf{D}^{+}(\mathcal{R})\to\mathbf{D}^{+}(\mathcal{R})\;,}\\ {}&{}\\ {R\varGamma_{c}(X;\cdot)}&{:\mathbf{D}^{+}(\mathcal{R})\to\mathbf{D}^{+}(\mathfrak{U b})\;,}\\ {}&{}\\ {R f_{!}}&{:\mathbf{D}^{+}(f^{-1}\mathcal{R})\to\mathbf{D}^{+}(\mathcal{R})\;.}\\ \end{array}
\]

此外，函子  $(\cdot)_Z: F \to F_Z$  和  $f^{-1}: F \to f^{-1}F$  准确地说，它们自然地扩展到导出范畴。我们得到 $* = \emptyset$ 、+、-、b 的函子：

\[
\begin{aligned}{}&{{}(\cdot)_{Z}\colon\mathsf{D}^{\ast}(\mathscr{R})\to\mathsf{D}^{\ast}(\mathscr{R})\enspace,}\\ {}&{{}f^{-1}\colon\mathsf{D}^{\ast}(\mathscr{R})\to\mathsf{D}^{\ast}(f^{-1}\mathscr{R})\enspace.}\\ \end{aligned}
\]

回想一下，为了计算这些派生函子，需要用与  $F$  准同构的内射层  $I$  复形替换层复形  $F$  ，并将这些函子应用于  $I$  。此外，对于给定的函子，在与该函子相关的子类单射中选择  $I^j$  就足够了。

示例 2.6.1。让 $ F \in \mathrm{Ob}(\mathcal{S}\mathfrak{h}(X)) $  。考虑一个精确序列  $ 0 \to F \to F^0 \to \cdots $ ，其中  $ F^{j'} $  是松弛的。将 F 识别为复数  $ \cdots \to 0 \to F \to 0 \to \cdots $  ，其中 F 处于零度，F 与复数准同构：

\[
F:\cdots\to0\to F^{0}\to F^{1}\to\cdots
\]

（其中  $ F^0 $  代表零度）。然后  $ R\Gamma_Z(X;F) $  由复数  $ \Gamma_Z(X;F^-) $  表示。

为了定义左派生函子  $ \cdot \otimes_R^L $  ，我们需要对  $ \mathscr{R} $  做出假设。

\statementlabel{定义 2.6.2}令 R 为 X 上的环层。一组（参见练习 I.29）：

\[
\mathsf{w g l d}(\mathcal{R})=\operatorname*{s u p}_{x\in X}\left(\mathsf{w g l d}(\mathcal{R}_{x})\right)\;.
\]

我们将 wgld(  $ \mathcal{R} $  ) 称为  $ \mathcal{R} $  的弱全局维度。

当考虑环层 R 上的张量积时，我们始终假设 R 具有有限的弱全局维数：

\[
\mathsf{w g l d}(\mathcal{R})<\infty.
\]

应用命题 2.4.12、推论 1.7.8 和练习 I.23，我们发现如果  $ F \in \text{Ob}(\mathbb{D}^b(\mathcal{R})) $  （或  $ \text{Ob}(\mathbb{D}^+(R)) $  ），则  $ F $  与平坦  $ \mathcal{R} $  模的有界复形（或从下面有界的复形）准同构。因此我们可以定义左派生函子：

\[
\begin{aligned}{}&{{}\overset{L}{\cdot}\otimes_{\mathcal{R}}\cdot\colon\mathsf{D}^{\ast}(\mathcal{R}^{{o p}})\times\mathsf{D}^{\ast}(\mathcal{R})\to\mathsf{D}^{\ast}(\mathcal{S})\enspace,}\\ {}&{{}\overset{L}{\cdot}\otimes_{\mathcal{S}}\cdot\colon\mathsf{D}^{\ast}(\mathcal{R})\times\mathsf{D}^{\ast}(\mathcal{S})\to\mathsf{D}^{\ast}(\mathcal{R})}\\ \end{aligned}
\]

与 $ * = -,\ +, b $  。

特别要考虑符号 2.3.12 的情况。我们得到  $ * = -,\ +,\ b $  的函子：

\[
\begin{array}{c}{^{L}\times}\\ {\boxminus_{\mathcal{A}}\cdot\colon\mathbf{D}^{*}(p_{X}^{-1}(\mathcal{R}^{o p}))\times\mathbf{D}^{*}(p_{Y}^{-1}\mathcal{R})\to\mathbf{D}^{*}(p^{-1}\mathcal{S})\enspace.}\\ {^{S}\times}\\ \end{array}
\]

（这里 R 和 S 是 S 上的环层。）

现在我们将系统地使用命题 1.8.7 和 1.10.9，简要描述所有这些函子之间存在的众多关系中的一些。

让 $ F \in \mathrm{Ob}(\mathbb{D}^{+}(\mathcal{R})) $  。然后：

\[
R\varGamma_{Z}(X;F)\simeq R\varGamma(X;R\varGamma_{Z}(F))\enspace.
\]

事实上，如果 $F$  是单射的， $\Gamma_{Z}(F)$  也是单射的。

让  $ F \in \mathrm{Ob}(\mathbb{D}^b(\mathcal{R})) $  并让  $ G \in \mathrm{Ob}(\mathbb{D}^+(R)) $  。然后：

\[
R\mathrm{H o m}_{\mathcal{R}}(F,G)\simeq R\varGamma(X;R\mathcal{H}\mathrm{o m}_{\mathcal{R}}(F,G))\enspace.
\]

事实上，松弛层的范畴相对于函子  $ \Gamma(X;\cdot) $  是单射的，并且如果 F 和 G 是具有 F 单射的  $ \mathcal{R} $  模的层，则  $ \mathcal{H}om_{\mathcal{H}}(G,F) $  是松弛的。

让 $ F \in \mathrm{Ob}(\mathbb{D}^{+}(g^{-1}f^{-1}\mathcal{R})) $  。然后：

\[
R(f\circ g)_{*}F\simeq R f_{*}R g_{*}F\enspace.
\]

事实上，如果 $F$ 是松弛的， $g_{*}F$ 是松弛的，并且松弛层形成了关于直像的单射范畴。类似地，将“flabby”替换为“c-soft”，可以得到：

\[
R(f\circ g)_{!}F\simeq R f_{!}R g_{!}F\enspace.
\]

\statementlabel{命题 2.6.3}令  $\phi: \mathbf{D}^+(\mathcal{R}) \to \mathbf{D}^+(\mathcal{S})$  为从健忘函子  $\mathfrak{Mod}(\mathcal{R}) \to \mathfrak{Mod}(\mathcal{S})$  获得的函子。让  $F \in \mathrm{Ob}(\mathbf{D}^-)(\mathcal{R})$  、  $G \in \mathrm{Ob}(\mathbf{D}^+)(\mathcal{R})$  和  $H \in \mathrm{Ob}(\mathbf{D}^-)(\mathcal{S})$  ，（和  $\mathrm{wgld}(\mathcal{S}) < \infty$  ）。然后：

\[
\phi(R\mathcal{H}o m_{\mathcal{S}}(H,G))\simeq R\mathcal{H}o m_{\mathcal{S}}(H,\phi(G)).
\]

(ii) 我们有：

\[
\begin{aligned}{{R}\mathcal{H}{o m}_{\mathcal{R}}\bigg(F\mathop{\otimes}_{\mathcal{S}}^{L}H,G\bigg)}&{{}\simeq{R}\mathcal{H}{o m}_{\mathcal{R}}(F,{R}\mathcal{H}{o m}_{\mathcal{S}}(H,G))}\\ {}&{{}\simeq{R}\mathcal{H}{o m}_{\mathcal{S}}(H,{R}\mathcal{H}{o m}_{\mathcal{R}}(F,G))\qquad{i n}\mathbb{D}^{+}(\mathcal{S}).}\\ \end{aligned}
\]

\statementlabel{证明}根据命题 2.2.9，如果  $H$  平铺于  $\mathcal{S}$  且  $G$  是单射的，则  $\mathcal{Hom}_{\mathcal{S}}(H,G)$  是单射  $\mathcal{R}$  模的复合体。因此  $R\mathcal{Hom}_{\mathcal{R}}(F\otimes^L_{\mathcal{S}}H,G)=\mathcal{Hom}_{\mathcal{R}}(F\otimes_{\mathcal{S}}H,G)$  和  $R\mathcal{Hom}_{\mathcal{R}}(F,R\mathcal{Hom}_{\mathcal{S}}(H,G))\simeq\mathcal{Hom}_{\mathcal{R}}(F,\mathcal{Hom}_{\mathcal{S}}(H,G))$  。那么(ii)中的第一个同构由命题2.2.9得出。为了证明 (i) 和 (ii) 中的第二个同构，我们将用（上面有界的）复数  $H$  替换  $H$  ，使得  $H$  是由命题 2.4.12 中的构造得到的  $\mathcal{S}_U$  的直和。然后对于任何松弛的 $\mathcal{S}$  -模 $K$ ， $\operatorname{Ext}^{j}_{\mathcal{S}}(H^n,K)=0$ 对于 $j\neq0$ 。因此，对于一个松散的  $\mathcal{S}$  -模  $K$  、  $R\mathcal{Hom}_{\mathcal{S}}(H,K)\simeq\operatorname{Hom}_{\mathcal{S}}(H,K)$  的复合体。现在假设 $G$ 单射。然后 (i) 从以下事实得出结论： $\phi(G)$  是一个由松散的  $\mathcal{S}$  模组成的复合体。类似地，由于  $\mathcal{Hom}_{\mathcal{R}}(F,G)$  是由命题 2.4.6 (vii) 得出的松弛  $\mathcal{S}$  模的复合体， $R\mathcal{Hom}_{\mathcal{S}}(H,R\mathcal{Hom}_{\mathcal{R}}(F,G))\simeq$ 

\[
\begin{aligned}{\operatorname{H o m}_{\mathsf{D}^{+}(\mathcal{R})}\biggl(F\mathop{\otimes}_{\mathcal{S}}^{L}H,G\biggr)}&{{}\simeq\operatorname{H o m}_{\mathsf{D}^{+}(\mathcal{S})}(H,R\mathcal{H o m}_{\mathcal{R}}(F,G))}\\ {}&{{}\simeq\operatorname{H o m}_{\mathsf{D}^{+}(\mathcal{R})}(F,R\mathcal{H o m}_{\mathcal{S}}(H,G)).}\\ \end{aligned}
\]

在(2.6.7)中选择 $ H = \mathcal{S}_Z $ ，我们得到：

\[
\begin{aligned}{R\mathcal{H}o m_{\mathcal{R}}(F_{Z},G)}&{{}\simeq R\varGamma_{Z}R\mathcal{H}o m_{\mathcal{R}}(F,G)}\\ {}&{{}\simeq R\mathcal{H}o m_{\mathcal{R}}(F,R\varGamma_{Z}(G))\qquad\mathrm{i n~}\mathsf{D}^{+}(\mathcal{S})\enspace.}\\ \end{aligned}
\]

此外，按照推论 2.2.10 进行论证，可以得到自然态射：

\[
R\mathcal{H}{_{O m}}_{\mathcal{R}}(F,G)\mathop{{\otimes}}_{\mathcal{S}}^{L}F\to G\qquad\mathrm{i n}\mathrm{}\mathbb{D}^{+}(\mathcal{R}),
\]

\[
R\mathcal{H}o m_{\mathcal{R}}(F,G)\overset{L}{\otimes}_{\dot{\mathcal{S}}}H\to R\mathcal{H}o m_{\mathcal{R}}\bigg(F,G\overset{L}{\otimes}_{\mathcal{S}}H\bigg)\qquad\mathrm{i n}\mathbb{D}^{+}(\mathcal{S}),
\]

\[
R\mathcal{H}o m_{\mathcal{R}}(F,G)\to R\mathcal{H}o m_{\mathcal{R}}(R\mathcal{H}o m_{\mathcal{R}}(G,\mathcal{R}),R\mathcal{H}o m_{\mathcal{R}}(F,\mathcal{R}))\quad\mathrm{i n}\;\mathsf{D}^{+}(\mathcal{R})
\]

在(2.6.12)中，我们假设 $ \mathcal{R} $ 是可交换的， $ \operatorname{wgld}(\mathcal{R}) < \infty $ 和 $ F $ 、 $ G $ 、 $ R \operatorname{Hom}_{\mathcal{R}}(G, \mathcal{R}) $ 是有界的。态射(2.6.12)定义如下。

到 (2.6.11) 我们有：

\[
R\mathcal{H}o m_{\mathcal{R}}(F,G)\overset{L}{\otimes}_{\mathcal{R}}R\mathcal{H}o m_{\mathcal{R}}(G,\mathcal{R})\to R\mathcal{H}o m_{\mathcal{R}}\bigg(F,G\overset{L}{\otimes}_{\mathcal{R}}R\mathcal{H}o m_{\mathcal{R}}(G,\mathcal{R})\bigg).
\]

另一方面，(2.6.10) 给出 $ G \otimes_R^LR \mathrm{Hom}_{\mathcal{R}}(G,\mathcal{R}) \to \mathcal{R} $  。将它们结合起来，我们得到：

\[
R\mathcal{H}{o m}_{\mathcal{R}}(F,G)\overset{L}{\otimes}_{\mathcal{R}}R\mathcal{H}{o m}_{\mathcal{R}}(G,\mathcal{R})\to R\mathcal{H}{o m}_{\mathcal{R}}(F,\mathcal{R})\enspace.
\]

然后由（2.6.8）得出（2.6.12）。

\statementlabel{命题 2.6.4}(i) 设  $ G \in \text{Ob}(\mathbb{D}^+(f^{-1}\mathcal{R})) $  和  $ F \in \text{Ob}(\mathbb{D}^-(\mathcal{R})) $  。然后：

\[
R\mathrm{H o m}_{\mathcal{R}}(F,R f_{*}G)\simeq R\mathrm{H o m}_{f^{-1}\mathcal{R}}(f^{-1}F,G)\enspace.
\]

(ii) 函子 $f^{-1} : \mathbf{D}^{+}(\mathcal{R}) \to \mathbf{D}^{+}(f^{-1}\mathcal{R})$  和 $Rf_* : \mathbf{D}^{+}(f^{-1}\mathcal{R}) \to \mathbf{D}^{+}(\mathcal{R})$  彼此伴随，即：

\[
\mathrm{H o m}_{\mathsf{D}^{+}(\mathcal{R})}(F,R f_{*}G)\simeq\mathrm{H o m}_{\mathsf{D}^{+}(f^{-1}\mathcal{R})}(f^{-1}F,G)\enspace.
\]

\statementlabel{证明}(i) 如果 $G$  是单射的，则 $f_*G$  是单射的。因此  $R\mathrm{Hom}_{\mathcal{A}}(F, Rf_*)(\cdot)$  是  $\mathrm{Hom}_{\mathcal{A}}(F, f_*)(\cdot)$  的派生函子。由于  $R\mathrm{Hom}_{f^{-1}\mathcal{A}}(f^{-1}F, \cdot)$  是  $\mathrm{Hom}_{f^{-1}\mathcal{A}}(f^{-1}F, \cdot)$  的派生函子，因此结果由命题 2.3.3 得出。

(ii) 证明类似。 ⑨

同样的论证给出：

\[
\begin{array}{r}{\mathrm{R}\mathcal{H}o m_{\mathcal{R}}(F,\mathrm{R}f_{*}\mathrm{G})\simeq\mathrm{R}f_{*}\mathrm{R}\mathcal{H}o m_{f^{-1}\mathcal{R}}(f^{-1}F,\mathrm{G})\qquad\mathrm{i n~}\mathbb{D}^{+}(\mathcal{S})\enspace.}\end{array}
\]

到(2.6.14)，我们得到自然态射：

\[
\mathrm{i d}\to R f_{*}\circ f^{-1}\qquad\mathrm{i n~\mathbb{D}^{+}(\mathcal{R})~},
\]

\[
f^{-1}\circ R f_{*}\to\mathtt{i d}\quad\mathrm{i n~\mathbb{D}^{+}(f^{-1}\mathcal{R})~}.
\]

\statementlabel{命题 2.6.5}让 $ F_1 \in \text{Ob}(\mathbf{D}^+(\mathcal{R}^{op})) $ 和 $ F_2 \in \text{Ob}(\mathbf{D}^+(\mathcal{R})) $ （和 $ \text{wgld}(\mathcal{R}) < \infty $ ）。然后：

\[
f^{-1}F_{1}\mathop{\bigotimes_{f^{-1}\mathcal{R}}}^{L}f^{-1}F_{2}\xrightarrow{\quad}f^{-1}\left(F_{1}\mathop{\bigotimes_{\mathcal{R}}}^{L}F_{2}\right)\quad{i n}\quad\mathbb{D}^{+}(f^{-1}\mathcal{S})\enspace.
\]

\statementlabel{证明}如果  $F$  在  $\mathscr{R}$  上平坦，则  $f^{-1}F$  在  $f^{-1}\mathscr{R}$  上平坦。因此  $f^{-1}(\cdot)\otimes_{f^{-1}\mathscr{R}}^{L}f^{-1}(\cdot)$  是  $f^{-1}(\cdot)\otimes_{f^{-1}\mathscr{R}}f^{-1}(\cdot)$  的派生函子，并且仍然需要应用命题 2.3.5。   $\square$ 

\statementlabel{命题 2.6.6}让  $ G \in \text{Ob}(\mathbb{D}^+(f^{-1}\mathcal{R}^{op})) $  并让  $ F \in \text{Ob}(\mathbb{D}^+(\mathcal{R})) $  和  $ \text{wgld}(\mathcal{R}) < \infty $  。然后：

\[
R f_{!}G\mathop{{\otimes}_{\mathcal{R}}}^{L}F\xrightarrow{\quad}R f_{!}\left(G\mathop{{\otimes}_{f^{-1}\mathcal{R}}}^{L}f^{-1}F\right)\qquad{i n~}\mathbb{D}^{+}(\mathcal{S})\enspace.
\]

（根据本节开头的约定，X 和 Y 是局部紧空间。）

\statementlabel{证明}首先假设 $F$  是一个平捆。根据引理 2.5.12，可以看出函子  $\cdot \otimes_{f^{-1}\mathcal{R}} f^{-1}F$  将  $c$  -软层发送到  $f_!$  -内射层。因此  $R f_!(\cdot \otimes_{f^{-1}\mathcal{R}} f^{-1}F)$  是  $f_!(\cdot \otimes_{f^{-1}\mathcal{R}} f^{-1}F)$  的派生函子。由于  $R f_!(\cdot) \otimes_{\mathcal{R}} F$  是  $f_!(\cdot) \otimes_{\mathcal{R}} F$  的派生函子，因此在这种情况下结果来自命题 2.5.13。为了处理一般情况，我们注意到，如果  $F \in \mathrm{Ob}(\mathbf{D}^{+}(\mathcal{R}))$  ，则  $F$  与从平层下方界定的复数准同构。   $\square$ 

\statementlabel{命题 2.6.7}考虑卡特平方 (2.5.5) 并设  $ G \in \mathrm{Ob}(\mathbb{D}^{+}(f^{-1}\mathcal{R})) $  。然后：

\[
g^{-1}\circ R f_{!}G\simeq R f_{!}^{\prime}\circ g^{\prime-1}G\qquad{~i n~}\mathbf{D}^{+}(g^{-1}\mathcal{R})\enspace.
\]

\statementlabel{证明}由于  $ g^{-1} \circ Rf_1 $  是  $ g^{-1} \circ f_1 $  的派生函子，因此根据命题 2.5.11 足以证明  $ Rf_1' \circ g^{' - 1} $  是  $ f_1' \circ g^{' - 1} $  的派生函子。

用 $\mathrm{I}_{Y}$  表示由层 $G$  组成的 $\mathrm{Mod}(f^{-1}\mathcal{R})$  的子范畴，使得对于所有 $x\in X$  ， $G|_{f^{-1}(x)}$  是 $c$  -soft，并在 $\mathrm{Mod}(g^{\prime -1}f^{-1}\mathcal{R})$  中类似地定义 $\mathrm{I}_{Y}$  。然后  $\mathrm{I}_{Y}$  相对于  $g^{\prime -1}$  是单射的，  $g^{\prime -1}$  将  $\mathrm{I}_{Y}$  发送到  $\mathrm{I}_{Y}$  。  $\mathrm{I}_{Y}$  相对于  $f^{\prime}$  是单射的。这样就完成了证明。

还有许多其他的自然态射或同构可以通过组合上面的公式来获得。让我们只提及其中的一些，将证明留给读者。在这些公式中，当  $ \otimes_{\mathcal{R}} $  出现时，我们假设  $ \operatorname{wgld}(\mathcal{R}) < \infty $  。让 $ F \in \operatorname{Ob}(\mathbf{D}^+((\mathcal{R})) $ ， $ G \in \operatorname{Ob}(\mathbf{D}^+(f^{-1}\mathcal{R}^{op})) $ 。我们得到：

\[
R f_{*}G\mathop{{\otimes}_{\mathcal{R}}}^{L}\underset{}{F\to R f_{*}}{\left(G\mathop{{\otimes}_{f^{-1}\mathcal{R}}}^{L}f^{-1}F\right)}\qquad\mathrm{i n}\mathbf{D}^{+}(\mathcal{S})\enspace.
\]

让 $ G_1 \in \mathrm{Ob}(\mathbb{D}^+(f^{-1}\mathcal{R}^{op})) $ ， $ G_2 \in \mathrm{Ob}(\mathbb{D}^+(f^{-1}\mathcal{R})) $ 。我们得到：

\[
R f_{*}G_{1}\mathop{{\otimes}_{\mathcal{R}}}^{L}R f_{*}G_{2}\to R f_{*}\left(G_{1}\mathop{{\otimes}_{{f^{-1}\mathcal{R}}}^{L}}G_{2}\right)\qquad\mathrm{i n~}\mathbb{D}^{+}(\mathcal{S})\enspace,
\]

\[
R f_{!}G_{1}\mathop{\otimes}_{\mathcal{R}}^{L}R f_{*}G_{2}\to R f_{!}\left(\begin{matrix}{G_{1}\mathop{\bigotimes}_{f^{-1}\mathcal{R}}^{L}G_{2}}\\ \end{matrix}\right)\qquad\mathrm{i n}\mathbf{D}^{+}(\mathcal{S})\enspace.
\]

令  $ G_1 \in \mathrm{Ob}(\mathbb{D}^b(f^{-1}\mathcal{R})) $  、  $ G_2 \in \mathrm{Ob}(\mathbb{D}^+(f^{-1}\mathcal{R})) $  并假设  $ f_* $  和  $ f_! $  具有有限上同调维数（参见练习 I.19）。我们得到：

\[
R f_{*}R\mathcal{H o m}_{f^{-1}\mathcal{R}}(G_{1},G_{2})\to R\mathcal{H o m}_{\mathcal{R}}(R f_{*}G_{1},R f_{*}G_{2})\qquad\mathrm{i n~\mathcal{D}^{+}(\mathcal{S})~},\qquad
\]

\[
R f_{*}R\mathcal{H o m}_{f^{-1}\mathcal{R}}(G_{1},G_{2})\to R\mathcal{H o m}_{\mathcal{R}}(R f_{!}G_{1},R f_{!}G_{2})\qquad\mathrm{~i n~}\mathsf{D}^{+}(\mathcal{S}),
\]

\[
R f_{!}R\mathcal{H o m}_{f^{-1}\mathcal{R}}(G_{1},G_{2})\to R\mathcal{H o m}_{\mathcal{R}}(R f_{*}G_{1},R f_{!}G_{2})\qquad\mathrm{~i n~}\mathsf{D}^{+}(\mathcal{S})~.
\]

让  $ F_1 \in \mathrm{Ob}(\mathbb{D}^b(\mathcal{R})) $  和  $ F_2 \in \mathrm{Ob}(\mathbb{D}^+(R)) $  。我们得到：

\[
f^{-1}R\mathcal{H}o m_{\mathcal{R}}(F_{1},F_{2})\to R\mathcal{H}o m_{f^{-1}\mathcal{R}}(f^{-1}F_{1},f^{-1}F_{2})\qquad\mathrm{i n~}\mathbb{D}^{+}(f^{-1}\mathcal{S})\enspace.
\]

最后我们将应用命题1.8.8、2.3.6、2.3.9和2.4.6。设 $ \dot{U}_1 $ 和 $ U_2 $ （分别为 $ Z_1 $ 和 $ Z_2 $ ）为X的两个开（或闭）子集，令Z为X的局部闭子集， $ Z' $ 为Z的闭子集。

对于  $ F \in \mathrm{Ob}(\mathbb{D}^{+}(\mathcal{R})) $  ，我们得到了区分的三角形：

\[
R\varGamma_{U_{1}\cup U_{2}}(F)\longrightarrow R\varGamma_{U_{1}}(F)\oplus R\varGamma_{U_{2}}(F)\longrightarrow R\varGamma_{U_{1}\cap U_{2}}(F)\underset{+1}{\longrightarrow},
\]

\[
R\varGamma_{Z_{1}\cap Z_{2}}(F)\longrightarrow R\varGamma_{Z_{1}}(F)\oplus R\varGamma_{Z_{2}}(F)\longrightarrow R\varGamma_{Z_{1}\cup Z_{2}}(F)\xrightarrow[+1]{},
\]

\[
F_{U_{1}\cap U_{2}}\longrightarrow F_{U_{1}}\oplus F_{U_{2}}\longrightarrow F_{U_{1}\cup U_{2}}\xrightarrow[+1]{},
\]

\[
F_{Z_{1}\cup Z_{2}}\longrightarrow F_{Z_{1}}\oplus F_{Z_{2}}\longrightarrow F_{Z_{1}\cap Z_{2}}\xrightarrow[+1]{},
\]

\[
R\varGamma_{Z^{\prime}}(F)\longrightarrow R\varGamma_{Z}(F)\longrightarrow R\varGamma_{Z\setminus Z^{\prime}}(F)\xrightarrow[+1]{}~,
\]

\[
F_{Z\setminus Z^{\prime}}\longrightarrow F_{Z}\longrightarrow F_{Z^{\prime}}\xrightarrow[+1]{}\cdot
\]

注释 2.6.8。令  $ \mathcal{R} $  为  $ X $  上的环层，  $ Z $  为  $ X $  、  $ F \in \mathrm{Ob}(\mathbf{D}^+(R)) $  、  $ K \in \mathrm{Ob}(\mathbf{D}^b(R)) $  、  $ G \in \mathrm{Ob}(\mathbf{D}^+(R^{op})) $  的局部闭子集。一套：

\[
\begin{array}{r}{H_{Z}^{j}(F)=H^{j}(R\varGamma_{Z}(F))\enspace,}\end{array}
\]

\[
\begin{array}{r}{H^{j}(X;F)=H^{j}(R\varGamma(X;F)),}\end{array}
\]

\[
H_{2}^{j}(X;F)=H^{j}(R\varGamma_{Z}(X;F))\enspace,
\]

\[
H^{j}(\mathbf{Z};F)=H^{j}(R\Gamma(\mathbf{Z};F))\enspace,
\]

\[
\begin{array}{r}{H_{c}^{j}(X;F)=H^{j}(R\varGamma_{c}(X;F))\enspace,}\end{array}
\]

\[
\begin{array}{r}{\mathcal{E}x t_{\mathcal{R}}^{\mathrm{j}}(\mathbb{K},\mathbb{F})=H^{\mathrm{j}}(\mathbb{R}\mathcal{H o m}_{\mathcal{R}}(\mathbb{K},\mathbb{F}))~,}\end{array}
\]

\[
\mathrm{E x t}_{\mathcal{A}}^{j}(K,F)=H^{j}(R\mathrm{H o m}_{\mathcal{A}}(K,F))~,
\]

\[
\mathcal{T}o r_{j}^{\mathcal{R}}(G,F)=H^{-j}\left(\begin{matrix}{\scriptstyle{L}}\\ {\scriptstyle\bigotimes\scriptstyle\mathcal{R}}\\ \end{matrix}\right).
\]

请注意，如果  $X$  是局部紧致空间，则  $H_c^j(X;F)\simeq\lim_{\overrightarrow{k}}H_k^j(X;F)$  ，其中  $K$  的范围贯穿  $X$  的紧致子集族。

\statementlabel{备注 2.6.9}对于  $X$  的不一定局部闭子集  $Z$  ，我们定义了函子  $F \mapsto \Gamma(Z; F) = \Gamma(Z; F|_{z})$  。因此可以考虑其派生函子  $F \mapsto R\Gamma(Z; F)$  。应注意  $R\Gamma(Z; F)$  可能与应用于  $F|_{z}$  的函子  $R\Gamma(Z; F)$  不同。然而，在以下情况之一中存在自然同构 $R\Gamma(Z; F) \simeq R\Gamma(Z; F|_{z})$ ：

(i) Z 开路，

(ii) X 是豪斯多夫且 Z 是紧的，

(iii) Z 在 X 的仿紧开子集 U 中是封闭的。

在所有这些情况下，人们还具有：

\[
H^{j}(Z;F)\xleftarrow{\simeq}\varinjlim H^{j}(U;F)
\]

其中 U 的范围穿过 X 中 Z 的开邻域族。

\statementlabel{备注 2.6.10}当将函子  $R\Gamma(X;\cdot)$  应用于可区分的三角形 (2.6.28)–(2.6.32) 时，我们得到了我们不写的新的可区分三角形。当将函子  $H^{0}(\cdot)$  应用于这些三角形时，我们得到长的精确序列。从(2.6.28)到(2.6.31)推导的长精确序列被称为Mayer-Victoris序列。例如，有一个精确的序列：

\[
\cdots\to H^{k-1}(U_{1}\cap U_{2};F)\to H^{k}(U_{1}\cup U_{2};F)\to H^{k}(U_{1};F)\oplus H^{k}(U_{2};F)\to\cdots
\]

以及通过将  $ R\Gamma_c(X;\cdot) $  与  $ Z = X $  、  $ Z' = K $  紧凑应用到 (2.6.33) 获得的精确序列：

\[
\cdots\to H^{k-1}(K,F)\to H_{c}^{k}(X\setminus K;F)\to H_{c}^{k}(X;F)\to\cdots.
\]

注释 2.6.11。 (i) 让 $ F \in \mathrm{Ob}(\mathbb{D}^+)(\mathbb{Z}_X)) $  。它的支持表示为  $ \mathrm{supp}(F) $  ，是  $ X $  的闭子集：

\[
\mathrm{supp}(F)=\overline{\bigcup_{j\in\mathbb{Z}}\mathrm{supp}H^{j}(F)}~.
\]

(ii) 令 $ A $  为环，并令 $ M \in \text{Ob}(\mathbb{D}^+)(\mathfrak{Mod}(A))) $  。一套：

\[
M_{x}=a_{x}^{-1}M,
\]

其中  $a_X$  是映射  $X \to \{\mathrm{pt}\}$  。

(iii) 令 $A$  为交换环。当在 $\mathbf{D}^{+}(A_{X})$ 范畴中工作时，我们将 $A$ 称为“基环”。如果没有混淆的风险，我们将写 Hom,  $\otimes$  ,  $\boxtimes$  ，而不是 Hom  $_{A}$  ,  $\otimes_{A}$  ,  $\boxtimes_{A}$  。从第三章开始，我们经常会写 $\mathbf{D}^{+}(X)$ 而不是 $\mathbf{D}^{+}(A_{X})$ 。

\subsection{若干消没定理}
在本节中，我们收集了一些关于层上同调群消失的有用结果，或者等效地，在某些特殊情况下关于上同调群同构的结果。这些结果可以通过“非特征变形”（如命题 2.7.2 中所示）获得，也可以通过同伦方法（借助 Mittag-Leffler 程序）获得，我们现在针对层进行说明。

\statementlabel{命题 2.7.1}令  $X$  为拓扑空间，并令  $F \in \mathrm{Ob}(\mathbf{D}^+)(\mathbb{Z}_X)$  。令 $\{U_n\}_{n \in \mathbb{N}}$  为 $X$  的开子集的递增序列， $\{Z_n\}_{n \in \mathbb{N}}$  为 $X$  的闭子集的递减序列。设置 $U = \bigcup_n U_n$ ， $Z = \bigcap_n Z_n$ 。

(i) 对于任何 $j$  ，自然映射 $\varphi_j: H_Z^j(U; F) \to \lim_{\leftarrow} H_{Z_n}^j(U_n; F)$  是满射的。

(ii) 假设对于给定的  $j$  ，射影系统  $\{H_{Z_n}^{j-1}(U_n;F)\}_n$  满足  $M\text{-}L$  条件。那么 $\varphi_i$  是双射的。

(iii) 现在让 $\{X_n\}_{n\in\mathbb{N}}$ 成为 $X$ 的子集的递增族，对于所有 $n$ 满足 $X=\bigcup_n X_n$ 和 $X_n\subset\mathrm{Int}(X_{n+1})$ 。假设对于给定的  $j$  ，射影系统  $\{H^{j-1}(X_n;F)\}_{n}$  满足  $M\cdot L$  条件。那么自然映射 $H^j(X;F)\to\lim_{}\quad H^j(X_n;F)$  是双射的。

\statementlabel{证明}我们可以假设 F 是松弛层的复合体。

为了证明 (i) 和 (ii) 用  $ E_{n} $  表示与双复形相关的简单复形：

\[
\begin{array}{c}\cdots\longrightarrow\quad\varGamma(U_{n};F^{i-1})\quad\longrightarrow\quad\varGamma(U_{n};F^{i})\quad\longrightarrow\cdots\\\quad\downarrow\quad\downarrow\quad\downarrow\\\longrightarrow\quad\varGamma(U_{n}\backslash Z_{n};F^{i-1})\quad\longrightarrow\quad\varGamma(U_{n}\backslash Z_{n};F^{i})\longrightarrow\cdots.\end{array}
\]

然后 $H_{Z_{n}}^{j}(U_{n}:F)$ 自系统 $\{E_{n}^{i}\}_{n}$ 讽刺1.12 d

(iii) 从  $ \{H^{j-1}(\text{Int } X_n; F)\}_n $  满足  $ M $  -L 条件和  $ \lim_{F \to \mathcal{C}} H^j(X_n; F) \simeq \lim_{F \to \mathcal{C}} H^j(\text{Int } X_n; F) $  的事实得出。   $ \square $ 

\statementlabel{命题 2.7.2（非特征变形引理）}令  $X$  为豪斯多夫空间  $F \in \mathrm{Ob}(\mathbb{D}^+)(\mathbb{Z}_X)$  ，并令  $\{U_t\}_{t \in \mathbb{R}}$  为  $X$  的开子集族。我们假设以下条件：

(i)  $ U_t = \bigcup_{s<t} U_s $  ，对于所有 $ t \in \mathbb{R} $  。

(ii) 对于 $ (s,t) $  和 $ s \leq t $  的所有对，集合 $ \overline{U}_{t}\backslash U_{s} \cap \operatorname{supp}(F) $  是紧的。

(iii) 设置  $ Z_s = \bigcap_{t > s} (U_t \setminus U_s) $  ，我们对所有  $ (s, t) $  与  $ s \leq t $  和所有  $ x \in Z_s \setminus U_t $  配对：

\[
({\cal R}\Gamma_{(X\setminus U_{t})}(F))_{x}=0\enspace.
\]

那么我们就得到了所有  $t \in \mathbb{R}$  的同构：

\[
R\varGamma\left(\bigcup_{s}U_{s};F\right)\rightsquigarrow R\varGamma(U_{t};F)\enspace.
\]

\statementlabel{证明}考虑以下断言：

\[
(a)_{k}^{s}\colon\ \operatorname*{l i m}_{\overrightarrow{t>s}}H^{k}(U_{t};F)\xrightarrow{\sim}H^{k}(U_{s};F)\enspace,
\]

\[
(b)_{k}^{t}\colon\operatorname*{l i m}_{\overleftarrow{s<t}}H^{k}(U_{s};F)\stackrel{}{\sim}H^{k}(U_{t};F)\enspace.
\]

假设  $(a)_k^s$  对于所有  $s \in \mathbb{R}$  和所有  $k \in \mathbb{Z}$  得到证明，并且  $(b)_k^t$  对于所有  $t \in \mathbb{R}$  和所有  $k < k_\circ$  得到证明。由命题1.12.6我们得到：

\[
H^{k}(U_{t};F)\xrightarrow{\sim}H^{k}(U_{s};F)
\]

对于所有  $ k < k_o $  ，所有  $ (s, t) $  与  $ s \leq t $  配对。那么对于 $ t $ 固定，序列 $ \{H^{k_o-1}(U_{t-1/n}; F)\}_n $ 满足 $ M $  -  $ L $ 条件， $ (b)_k^t $ 满足命题1.12.4。通过对  $ k $  进行归纳论证，我们发现  $ (b)_k^t $  对于所有  $ t \in \mathbb{R} $  所有  $ k \in \mathbb{Z} $  都将得到满足。再次将命题 2.7.1 应用于系统  $ \{H^k(U_n; F)\}_{n \in \mathbb{N}} $  ，我们通过对  $ k $  进行归纳得到结果。

现在我们要证明 $(a)_{k}^{s}$ 。将  $X$  替换为  $\operatorname{supp}(F)$  ，我们可以从一开始就假设  $\overline{U_{t}\backslash U_{s}}$  对于所有  $t\geqslant s$  都是紧凑的。考虑一下不同的三角形：

\[
R\varGamma_{(X\setminus U_{t})}(F)|_{Z_{s}}\longrightarrow R\varGamma_{(X\setminus U_{s})}(F)|_{Z_{s}}\longrightarrow R\varGamma_{(U_{t}\setminus U_{s})}(F)|_{Z_{s}}\xrightarrow[~+1~]{}
\]

由于根据假设，前两项为零，因此我们得到  $ R\Gamma_{(U_t \setminus U_s)}(F)|_{Z_s} = 0 $  。因此，对于所有  $ k \in \mathbb{Z} $  和所有  $ t \geq s $  我们有（参见符号 2.6.8）：

\[
\begin{aligned}{0}&{{}=H^{k}(Z_{s},R\varGamma_{(U_{t}\setminus U_{s})}(F))}\\ {}&{{}=\varliminf_{U\to\overrightarrow{Z}_{s}}H^{k}(U\cap U_{t};R\varGamma_{(X\setminus U_{s})}(F)),}\\ \end{aligned}
\]

其中  $U$  范围涵盖  $Z_s$  的开放邻域系列。由于对于任何这样的  $U$  都存在  $t'$  和  $t \geq t' > s$  这样  $U \cap U_t \supset U_{t'} \setminus U_s$  ，我们得到：

\[
0=\varinjlim_{t>s}H^{k}(U_{t};R\varGamma_{(X\setminus U_{s})}(F))\enspace.
\]

这意味着  $ (a)_{k}^{s} $  。 ⑨

我们将在第三章中系统地研究欧几里得空间上的层，但现在我们需要一个特殊的结果。

\statementlabel{引理 2.7.3}令 I 为  $ \mathbb{R} $  的区间 [0,1] 。设 F 为 I 上的一个层，则：

(i)  $ H^j(I; F) = 0 $  对于所有  $ j > 1 $  。

(ii) 假设 $ F(I) \to F_t $  对于所有 $ t \in I $  都是满射。然后 $ H^j(I; F) = 0 $  对于所有 $ j \geq 1 $  。

(iii) 如果  $F$  是  $I$  上的常值层  $M_I$  ，则态射  $M \to R\Gamma(I; F) \to F_t$  是任何  $t \in I$  的同构。

\statementlabel{证明}让 $ j \geqslant 1 $  ，并让 $ s \in H^j(I; F) $  。让 $ f_{t_1, t_2}(t_1 \leqslant t_2) $ 成为自然映射：

\[
f_{t_{1},t_{2}}\colon H^{j}(I;F)\to H^{j}([t_{1},t_{2}];F)\enspace.
\]

设置 $ J = \{t \in [0,1]; f_{0,t}(s) = 0\} $ 。然后  $ 0 \in J $  和  $ 0 \leqslant t' \leqslant t $  、  $ t \in J $  意味着  $ t' \in J $  。此外， $ J $  是开放的，因为我们对所有 $ t_0 < 1 $  都有：

\[
H^{j}([0,t_{\circ}];F)=\varinjlim_{t>t_{\circ}}H^{j}([0,t];F)\enspace,
\]

对于某些  $ t>t_{0} $  ，  $ f_{0,t_{0}}(s)=0 $  意味着  $ f_{0,t}(s)=0 $  。

考虑与分解  $ [0, t_{\circ}] = [0, t] \cup [t, t_{\circ}] $  相关的 Mayer-Vietoris 序列，其中  $ 0 \leqslant t \leqslant t_{\circ} \leqslant 1 $  ：

\[
\cdots\to H^{j}([0,t_{\circ}];F)\to H^{j}([0,t];F)\oplus H^{j}([t,t_{\circ}];F)\to H^{j}(\{t\};F)\to\cdots.
\]

对于  $j > 1$  或  $j = 1$  假设  $F(I) \to F_{t}$  是满射的，我们得到：

\[
H^{j}([0,t_{\circ}];F)\simeq H^{j}([0,t];F)\oplus H^{j}([t,t_{\circ}];F)\enspace.
\]

然后选择 $ t_{\circ} = \sup J $ 。由于  $ \lim_{t \to t_{\circ}} H^{j}([t, t_{\circ}]; F) = 0 $  ，存在  $ t < t_{\circ} $  与

 $ f_{t,t\circ}(s) = 0 $ 。另一方面 $ f_{0,t}(s) = 0 $ 。因此  $ f_{0,t\circ}(s) = 0 $  (2.7.2) 和  $ J = [0,1] $  。

(iii) 自  $H^j(I;F)=0$  至  $j>0$  至 (ii)、 $R\Gamma(I;F)\simeq\Gamma(I;F)$  。此外，组合 $M\to\Gamma(I;F)\to F_t\simeq M$  是恒等式。因此，足以证明如果  $s\in\Gamma(I;F)$  和  $s_t=0$  ，则  $s$  消失。这是因为  $\mathrm{supp}(s)$  是开放和封闭的。   $\square$ 

现在考虑三个拓扑空间  $S$  、  $X$  、  $Y$  和连续映射的交换图：

\begin{center}
\includegraphics[width=0.15\textwidth]{流行上的层_Chn-assets/img_in_image_box_490_497_672_679.png}
\end{center}

在这种情况下，可以说  $f: Y \to X$  是  $S$  上的连续映射。

让 R 表示 S 上的环层。

让 $ F \in \mathrm{Ob}(\mathbb{D}^{+}(\mathcal{R})) $  。态射 (2.3.5) 引发态射：

\[
R p_{x*}\circ p_{x}^{-1}F\to R p_{x*}\circ R f_{*}\circ f^{-1}\circ p_{x}^{-1}F\simeq R p_{Y*}\circ p_{Y}^{-1}F~,
\]

如果  $f$  是真映射，则函子的态射

\[
R p_{X!}\circ p_{X}^{-1}F\to R p_{X!}\circ R f_{*}\circ f^{-1}\circ p_{X}^{-1}F\simeq R p_{Y!}\circ p_{Y}^{-1}F\enspace.
\]

我们用 $ f^{\#} $ （分别是 $ f_{c}^{\#} $ ）表示态射：

\[
R p_{X*}\circ p_{X}^{-1}\to R p_{Y*}\circ p_{Y}^{-1}\qquad\mathrm{(r e s p.~}R p_{X!}\circ p_{X}^{-1}\to R p_{Y!}\circ p_{Y}^{-1}{\mathrm{)}}
\]

由 (2.7.4)（分别为 2.7.5）给出。

\statementlabel{定义 2.7.4}(i) 令 $ f_0: Y \to X $  和 $ f_1: Y \to X $  为S 上的两个连续映射。如果在S 上存在 $ h: Y \times I \to X $  、 $ (I = [0,1]) $  上的连续映射，则可以说 $ f_0 $  和 $ f_1 $  在S 上同伦，这样，用 $ j_t $  (  $ t \in I $  ) 表示从Y 到 $ y \mapsto (y,t) $  的映射 $ Y \times I $ ，一个有 $ f_i = h \circ j_i $  ( $ i = 0,1 $ )。此外，如果 h 是真映射，则可以说  $ f_0 $  和  $ f_1 $  是正确的同伦。

(ii) 令  $f: Y \to X$  为  $S$  上的连续映射。如果 $S$ 上存在连续映射 $g: X \to Y$ ，使得 $f \circ g$ 和 $g \circ f$ 分别与 $\mathrm{id}_{X}$ 和 $\mathrm{id}_{Y}$ 同伦，则我们说 $f$ 是 $S$ 上的同伦同构。

\statementlabel{命题 2.7.5}(i) 令  $ f_0: Y \to X $  和  $ f_1: Y \to X $  为 S 上的两个同伦映射。然后  $ f_0^\# = f_1^\# $  。

(ii) 如果  $f_0$  和  $f_1$  是正确同伦的，则  $f_{0c}^{\#} = f_{1c}^{\#}$  。

\statementlabel{证明}(i) 令 $h: Y \times I \to X$  为 $S$  之上 $f_0$  和 $f_1$  之间的同伦。我们有一个交换图：

\begin{center}
\includegraphics[width=0.19\textwidth]{流行上的层_Chn-assets/img_in_image_box_469_156_706_487.png}
\end{center}

其中  $p: Y \times I \to Y$  表示投影。

根据备注 2.5.3 和引理 2.7.3，我们有从  $ \mathbf{D}^{+}(\mathcal{R}) $  到  $ \mathbf{D}^{+}(p_{Y}^{-1}\mathcal{R}) $  的函子同构：

\[
p_{Y}^{-1}\xrightarrow{\sim}R p_{*}\circ p^{-1}\circ p_{Y}^{-1}\xrightarrow{\sim}j_{i}^{-1}\circ p^{-1}\circ p_{Y}^{-1}\simeq p_{Y}^{-1}
\]

所有这些同构的组合就是  $ p_{Y}^{-1} $  上的恒等式。

我们得到从  $ \mathbf{D}^{+}(\mathcal{R}) $  到  $ \mathbf{D}^{+}(\mathcal{R}) $  函子态射的交换图：

\begin{center}
\includegraphics[width=0.49\textwidth]{流行上的层_Chn-assets/img_in_image_box_282_805_873_1127.png}
\end{center}

这里态射  $ \alpha_{i} $  是由以下推导出来的：

\[
R p_{*}\circ p^{-1}\circ p_{\Upsilon}^{-1}\to j_{i}^{-1}\circ p^{-1}\circ p_{\Upsilon}^{-1}=p_{\Upsilon}^{-1}.
\]

因此  $ \alpha_{0}=\alpha_{1} $  (2.7.7) 和  $ f_{0}^{\#}=f_{1}^{\#} $  。

(ii) 证明类似。 ⑨

\statementlabel{备注 2.7.6}令  $ X' $  （分别为  $ Y' $  ）为  $ X $  （分别为  $ Y $  ）的闭子集，并让  $ f: Y \to X $  为  $ S $  上的映射，使得  $ f^{-1}(X') \subset Y' $  。让 $ F \in \text{Ob}(\mathbf{D}^+\left(\mathcal{R}\right)) $  。我们有自然态射：

\[
\begin{aligned}{R p_{X*}\circ R\varGamma_{X^{\prime}}\circ p_{X}^{-1}F}&{{}\to R p_{X*}\circ R f_{*}\circ f^{-1}\circ R\varGamma_{X^{\prime}}\circ p_{X}^{-1}F}\\ {}&{{}\to R p_{X*}\circ R f_{*}\circ R\varGamma_{f^{-1}(X^{\prime})}\circ f^{-1}\circ p_{X}^{-1}F}\\ {}&{{}\to R p_{Y*}\circ R\varGamma_{Y^{\prime}}\circ p_{Y}^{-1}F\enspace.}\\ \end{aligned}
\]

让我们用  $ f_{Y^{\prime},X^{\prime}}^{\#} $  表示函子的态射：

\[
f_{Y^{\prime},X^{\prime}}^{\#}\colon R p_{X\ast}\circ R\varGamma_{X^{\prime}}\circ p_{X}^{-1}\to R p_{Y\ast}\circ R\varGamma_{Y^{\prime}}\circ p_{Y}^{-1}\enspace.
\]

如果 $ Y' = f^{-1}(X') $ ，我们写 $ f_X' $ 而不是 $ f_Y' $ 。现在假设定义 2.7.4 中的  $ f_0 $  、  $ f_1 $  、 h 满足：

\[
\mathrm{f o r~e a c h~}\lambda\in I,h_{\lambda}^{-1}(X^{\prime})\subset Y^{\prime},\mathrm{~w h e r e~}h_{\lambda}(\cdot)=h(\cdot,\lambda)\enspace.
\]

那么通过修改命题2.7.5的证明，我们得到：

\[
f_{0Y^{\prime},X^{\prime}}^{\#}=f_{1Y^{\prime},X^{\prime}}^{\#}.
\]

\statementlabel{推论 2.7.7}(i) 设  $f: Y \to X$  为  $S$  的同伦同构。那么，对于任何  $F \in \mathrm{Ob}(\mathbb{D}^{+}(\mathcal{R}))$  ，  $f^{*}: R p_{X^*} p_X^{-1} F \to R p_{Y^*} p_Y^{-1} F$  都是同构。

(ii) 令 $f: Y \to X$  为连续映射， $s: X \to Y$  为 $f$  的连续部分。假设  $s \circ f$  与  $X$  上的恒等式同伦。那么对于任何  $F \in \mathrm{Ob}(\mathbf{D}^{+}(A_X))$  、  $F \to R f_*f^{-1}F$  和组合  $R f_*f^{-1}F \to R f_* R s_*s^{-1}f^{-1}F \simeq F$  都是同构。

(iii) 特别是，如果  $X$  是可收缩拓扑空间，则  $M \simeq R\Gamma(X; M_X)$  对于任何  $M \in \mathrm{Ob}(\mathbb{D}^+(\mathfrak{Mod}(A)))$  。

(iv) 设 $f: Y \to X$  为连续映射。假设  $f$  对于可收缩纤维是正确的（因此是满射的）。那么对于  $F \in \mathrm{Ob}(\mathbf{D}^{+}(A_X))$  ，  $F \to R f_* f^{-1} F$  是同构。

\statementlabel{证明}(i) 令  $g: X \to Y$  为定义 2.7.4 (ii) 中的映射。然后，根据命题 2.7.5， $f^\# \circ g^\# = \mathrm{id}$  和  $g^\# \circ f^\# = \mathrm{id}$  。

(ii) 是(i) 的特例。

(iii) 是(ii) 的特例。

(iv) 让 $ F \in \mathrm{Ob}(\mathbb{D}^{+}(A_X)) $  并让 $ x \in X $  。根据备注 2.5.3，我们有：

\[
\begin{aligned}(Rf_{*}\circ f^{-1}F)_{x}&\simeq R\varGamma(f^{-1}(x);f^{-1}F)\\&\simeq R\varGamma(f^{-1}(x);(f|_{f^{-1}(x)})^{-1}F_{x})\\&\simeq F_{x}~,\end{aligned}
\]

由（三）。 ⑨

请注意，推论 2.7.7 (iv) 通常称为 Vietoris-Begle 定理。我们将使其更加精确，并将其扩展到不适当的情况。

令 $f: Y \to X$  为连续映射。让我们用 $\mathfrak{Mod}(A_Y/f)$ 表示由层 $G$ 组成的 $\mathfrak{Mod}(A_Y)$ 的满子范畴，使得对于每个 $x \in X$ ， $\operatorname{sheaf} G|_{f^{-1}(x)}$ 是局部常值的，让我们用 $\mathbf{D}_f^+(A_Y)$ 表示由对象 $G$ 组成的 $\mathbf{D}^+(A_Y)$ 的满子范畴，使得 $H^j(G) \in \mathfrak{Mod}(A_Y/f)$ 对于所有 $j$ 。请注意， $\mathfrak{Mod}(A_Y/f)$  是  $\mathfrak{Mod}(A_Y)$  的厚子范畴（参见 I §7）， $\mathbf{D}_f^+(A_Y)$  是三角范畴。函子  $f_*$  和  $f^{-1}$  导出函子，我们仍然

用  $ f_{*} $  和  $ f^{-1} $  表示：

\[
\mathfrak{M o d}(A_{X})\xrightarrow[f_{*}]{\overset{f^{-1}}{\longrightarrow}}\mathfrak{M o d}(A_{Y}/f)\enspace.
\]

“包含”函子  $ \text{Mod}(A_Y/f) \to \text{Mod}(A_Y) $  引出一个函子： $ \mathbf{D}^+(\text{Mod}(A_Y/f)) \to \mathbf{D}_f^+(A_Y) $  。因此，我们获得函子（我们仍然用  $ R f_* $  和  $ f^{-1} $  表示）：

\[
\mathbf{D}^{+}(A_{X})\xrightarrow[R f_{*}]{\quad f^{-1}\quad}\mathbf{D}_{f}^{+}(A_{Y})\enspace.
\]

我们将给出一个充分条件，以便这些函子彼此互逆。假设给定一个 Y 闭子集族 $ \{Y_n\}_{n\in\mathbb{N}} $ ，满足：

\[
\left\{\begin{aligned}{Y}&{{}=\bigcup_{n}Y_{n},Y_{n}\subset\operatorname{I n t}(Y_{n+1})\qquad\mathrm{f o r~a l l}n,}\\ {f|_{Y_{n}}}&{{}\qquad\mathrm{i s~p r o p e r~w i t h~c o n t r a c t i b l e~f i b e r s~f o r~a l l}n.}\\ \end{aligned}\right.
\]

\statementlabel{命题 2.7.8}在假设（2.7.13）下，（2.7.11）（或（2.7.12））中的函子 $ f^{-1} $ 和 $ f_{*} $ （分别为 $ f^{-1} $ 和 $ Rf_{*} $ ）彼此互逆。

\statementlabel{证明}我们只给出式(2.7.12)情况下的证明，其他情况类似且更简单。

首先注意可收缩空间非空， $ f(Y_{n}) = X $  对于所有 n。

此外请注意，可收缩空间上的局部常值层是常值层（参见练习 II.4）。我们首先证明 f 正确时的结果。

让 $G \in \mathrm{Ob}(\mathbf{D}_f^+(A_Y))$  ，并让 $y \in Y$  。那么  $H^j(G)|_{f^{-1}f(y)}$  是一个常值层，因此  $H^i(f^{-1}f(y); H^j(G)) = 0$  对于所有  $j \in \mathbb{Z}$  和所有  $i > 0$  。我们得到：

\[
\begin{aligned}{H^{j}(f^{-1}\circ R f_{*}^{}G)_{y}}&{{}\simeq H^{j}(R f_{*}^{}G)_{f(y)}}\\ {}&{{}\simeq H^{j}(f^{-1}f(y);G)}\\ {}&{{}\simeq\varGamma(f^{-1}f(y);H^{j}(G))}\\ {}&{{}\simeq H^{j}(G)_{y}}\\ \end{aligned}
\]

因此  $f^{-1} \circ Rf_* \simeq \mathrm{id}_{\mathbf{D}_f^+(A_\gamma)}$  ，我们通过推论 2.7.7 知道  $Rf_* \circ f^{-1} \simeq \mathrm{id}_{\mathbf{D}^+(A_\chi)}$  。

现在我们处理一般情况。令  $ F \in \mathrm{Ob}(\mathbb{D}^+) (A_X) $  和  $ V $  成为  $ X $  的开子集。对于每个  $ j $  ，族  $ \{H^j (f^{-1}(V) \cap Y_n; f^{-1}F)\}_n $  满足  $ M $  -  $ L $  条件。应用引理 2.7.1，我们得到：

\[
\begin{aligned}{H^{j}(V;F)}&{{}\simeq H^{j}(f^{-1}(V)\cap Y_{n};f^{-1}F)}\\ {}&{{}\simeq H^{j}(f^{-1}(V);f^{-1}F).}\\ \end{aligned}
\]

因此 $F \mapsto Rf_* \circ f^{-1}F$  是同构。

最后让 $G \in \mathrm{Ob}(\mathbf{D}_f^+(A_Y))$ 、 $y \in Y$ 和 $V$ 成为 $X$ 中 $f(y)$ 的开邻域。由于对于所有  $j$  ，族  $\{H^j(f^{-1}(V) \cap Y_n; G)\}_n$  满足  $M\text{-}L$  条件，我们得到：

\[
H^{j}({f}^{-1}(V);G)\simeq H^{j}({f}^{-1}(V)\cap Y_{n};G).
\]

因此。

\[
\begin{aligned}{H^{j}(f^{-1}\circ R f_{*}^{}G)_{y}}&{{}\simeq H^{j}(R f_{*}^{}G)_{f(y)}}\\ {}&{{}\simeq\varinjlim_{V\ni f(y)}H^{j}(f^{-1}(V);G)}\\ {}&{{}\simeq\varinjlim_{V\ni f(y)}H^{j}(f^{-1}(V)\cap Y_{n};G)}\\ {}&{{}\simeq H^{j}(G)_{y}\;,}\\ \end{aligned}
\]

和  $G \mapsto f^{-1} \circ Rf_{*}$  G 是同构。 □

\subsection{覆盖的上同调}
令  $X$  为拓扑空间，并令  $\mathcal{U} = \{U_j\}_{j \in J}$  为  $X$  的开子集族。

对于整数  $ p \geqslant 0 $  和  $ \alpha = (\alpha_0, \ldots, \alpha_p) \in J^{p+1} $  ，我们设置：

\[
U_{\alpha}=\bigcap_{\nu=0}^{p}U_{\alpha_{\nu}}.
\]

如果  $\sigma$  是  $\{0,\ldots,p\}$  的排列，我们用  $\mathrm{sgn}\sigma$  表示它的签名。如果  $\alpha=(\alpha_{0},\ldots,\alpha_{p})$  我们设置  $\alpha^{\sigma}=(\alpha_{\sigma(0)},\ldots,\alpha_{\sigma(p)})$  。

让 $F$  成为 $X$  上的一个层。对于  $X$  的两个开子集  $U$  和  $V$  ，以及  $U \subset V$  ，我们用  $\rho_{V,U}$  表示规范态射  $F_{U} \to F_{V}$  （参见§3）。

对于 U 和 F，我们将 X 上的层复合体关联起来（参见符号 1.3.9），表示为  $ \mathcal{C}_{.}(\mathcal{U}; F) $  ，如下所示。

对于  $ p \in \mathbb{Z} $  、  $ p < 0 $  、  $ \mathcal{C}_p(\mathcal{U}; F) = 0 $  。

对于  $p \geq 0$  ，  $\mathcal{C}_p(\mathcal{U}; F)$  是  $\bigoplus_{\alpha \in J^{p+1}} F_{U_\alpha}$  的子层，由交替的部分组成：  $(s_\alpha)_{\alpha \in J^{p+1}}$  定义  $\mathcal{C}_p(\mathcal{U}; F)$  的一个部分，当且仅当满足以下两个条件时。

\[
\left\{\begin{aligned}{}&{{}\mathtt{F o r a n y p e r m u t a t i o n}\sigma\operatorname{o f}\{0,\ldots,p\},}\\ {}&{{}s_{\alpha}=(\mathtt{s g n}\sigma)s_{\alpha^{\sigma}}.}\\ \end{aligned}\right.
\]

\[
\begin{cases}\operatorname{If}\alpha=(\alpha_{0},\ldots,\alpha_{p})\operatorname{and}\alpha_{v}=\alpha_{v^{\prime}},\operatorname{for a pair of distinct}v\operatorname{and}v^{\prime}\\\operatorname{in}\{0,\ldots,p\},\operatorname{then}s_{\alpha}=0\end{cases}
\]

微分 $ d_p: \mathcal{C}_p(\mathcal{U}; F) \to \mathcal{C}_{p-1}(\mathcal{U}; F) $  由下式给出：

\[
\left\{\begin{aligned}{}&{{}\mathbf{f o r}\;\beta\in J^{p}\mathbf{a n d}s=(s_{\alpha})_{\alpha}\in\mathcal{C}_{p}(\mathcal{U};F),}\\ {}&{{}(\mathbf{d}_{p}(s))_{\beta}=\sum_{j\in J}\rho_{U_{\beta},U_{(j,\beta)}}(s_{(j,\beta)}).}\\ \end{aligned}\right.
\]

在这里， $ (j, \beta) = (j, \beta_{0}, \ldots, \beta_{p-1}) \in J^{p+1} $ 

人们立即检查：

\[
\mathbf{d}_{p-1}\circ\mathbf{d}_{p}=0.
\]

示例 2.8.1。 (i) 如果  $ J = \{0, 1\} $  ，则  $ \mathcal{C}_{.}(\mathcal{U}; F) $  是复数：

\[
0\to F_{U_{1}\cap U_{2}}\to F_{U_{1}}\oplus F_{U_{2}}\to0\enspace.
\]

(ii) 如果  $ J = \{0, 1, 2\} $  ，则  $ \mathcal{C}_{.}(\mathcal{U}; F) $  是复数：

\[
0\to F_{U_{1}\cap U_{2}\cap U_{3}}\to F_{U_{2}\cap U_{3}}\oplus F_{U_{3}\cap U_{1}}\oplus F_{U_{1}\cap U_{2}}\to F_{U_{1}}\oplus F_{U_{2}}\oplus F_{U_{3}}\to0~.
\]

设置 $U = \bigcup_j U_j$ 。通过设置 $\delta|_{\mathcal{F}_{U_i}} = \rho_{U,U_i}$ 来定义“增强图” $\delta: \mathcal{C}.(\mathcal{U}; F) \to F_U$ 。然后立即验证：

\[
\delta\circ\mathrm{d}_{1}=0.
\]

通过这种关系，  $ \delta $  引发了复合体  $ \mathscr{C}_(\mathscr{U}; F) \to F_U $  的态射，其中  $ F_U $  被识别为复合体  $ \cdots \to 0 \to F_U \to 0 \to \cdots $  。

\statementlabel{引理 2.8.2}态射 $ \mathcal{C}_{.}(\mathcal{U}; F) \to F_{U} $  是准同构。

\statementlabel{证明}让  $\widetilde{\mathcal{C}}.(\mathcal{U};F)$  表示增广复数  $\mathcal{C}.(\mathcal{U};F)\to F$  （即复数  $\cdots\xrightarrow{\mathfrak{d}_{1}}\mathcal{C}_{0}(\mathcal{U};F)\xrightarrow{\mathfrak{d}}F\longrightarrow0\longrightarrow\cdots$  ）。让 $j_{\circ}\in J$ 并让 $J^{\prime}=J\backslash\{j_{\circ}\}$ ， $\mathcal{U}^{\prime}=\{U_{j}\}_{j\in J^{\prime}}$ 。态射 $\rho_{X,U_{j_{\circ}}}$  引发态射：

\[
\alpha_{j_{*}}:\widetilde{\mathcal{C}}_{.}(\mathcal{U}^{\prime};F_{U_{j_{*}}})\to\widetilde{\mathcal{C}}_{.}(\mathcal{U}^{\prime};F)\enspace.
\]

首先让我们注意  $ \widetilde{\mathcal{C}}(\mathcal{U}; F) $  是  $ \alpha_{j_s} $  的映射锥。事实上，让我们用 $ \bigoplus_{\alpha \in J^{p+1}} F_{U_s} $  表示 $ \bigoplus_{\alpha \in J^{p+1}} F_{U_s} $  的子层，该子层由满足(2.8.1) 和(2.8.2) 的部分组成。然后：

\[
\begin{aligned}{\widetilde{\mathcal{C}}_{p}(\mathcal{U};F)}&{{}=\bigoplus_{\alpha\in J^{p+1}}^{\prime}F_{U_{\alpha}}}\\ {}&{{}=\left(\bigoplus_{\alpha=(j_{\circ},\beta),\beta\in J^{p}}^{\prime}F_{U_{j_{\circ}}\cap U_{\beta}}\right)\oplus\left(\bigoplus_{\alpha\in J^{\prime p+1}}^{\prime}F_{U_{\alpha}}\right)}\\ {}&{{}=\widetilde{\mathcal{C}}_{p-1}(\mathcal{U}^{\prime};F_{U_{j_{\circ}}})\oplus\widetilde{\mathcal{C}}_{p}(\mathcal{U}^{\prime};F)}\\ \end{aligned}
\]

检查  $ \widetilde{\mathcal{C}}_p(\mathcal{U}; F) $  的微分是否是  $ \alpha_{j_0} $  映射锥的微分。由于 $ \alpha_{j_0}|_{U_{j_0}} $  是同构，因此复数 $ \widetilde{\mathcal{C}}.(\mathcal{U}; F)|_{U_{j_0}} $  与零同伦。因此  $ \mathcal{C}.(\mathcal{U}; F) \to F_U $  是  $ U $  的准同构，并且在  $ U $  之外两项均为零，这就完成了证明。   $ \square $ 

\statementlabel{备注 2.8.3}事实上，我们已经证明，在  $U$  上局部，  $\mathcal{C}_{.}(\mathcal{U};F)$  与  $F$  是同伦的。

现在我们假设 U 是 X 的开覆盖。我们设置：

\[
\mathcal{C}(\mathcal{U};\mathbb{F})=\mathcal{H o m}_{\mathbb{I}_{X}}(\mathcal{C}.(\mathcal{U};\mathbb{Z}_{X}),\mathbb{F})\enspace,
\]

\[
C^{\prime}(\mathcal{U};F)=\operatorname{H o m}_{\mathbb{Z}_{X}}(\mathcal{C}.(\mathcal{U};\mathbb{Z}_{X}),F)\enspace.
\]

然后 $ \Gamma(X;\mathcal{C}(\mathcal{U};F)) = C^{\prime}(\mathcal{U};F) $ 。此外，由于 $ \operatorname{Hom}_{\mathcal{I}_X}(\mathbb{Z}_U;F) = \Gamma(U;F) $ ， $ C^p(\mathcal{U};F) $ 是 $ \prod_{\alpha\in J^{p+1}} F(U_\alpha) $ 的子模，由满足(2.8.1)–(2.8.2)的部分组成。

\statementlabel{命题 2.8.4}增强图  $\delta$  引起层的准同构： $F \xrightarrow{\sim} \mathcal{C}^{\prime}(\mathcal{U}; F)$  。

\statementlabel{证明}将备注 2.8.3 应用于捆  $\mathbb{Z}_{X}$  。 □

\statementlabel{命题 2.8.5（Leray 非循环覆盖定理）}令 $ \mathcal{U} = \{U_j\}_{j \in J} $  为 $ X $  的开放覆盖，让 $ F $  为 $ X $  上的一个层。假设对于所有  $ \alpha \in J^p $  (  $ p \geq 0 $  )、所有  $ k > 0 $  、  $ H^k(U_\alpha; F) = 0 $  。那么  $ R\Gamma(X; F) $  与  $ C^\cdot(\mathcal{U}; F) $  准同构。

\statementlabel{证明}令  $I^\prime$  为与  $F$  准同构的松弛层复合体。然后 $R\Gamma(X;F) \simeq \Gamma(X;I^\prime)$ 。根据命题 2.8.4 和定理 1.9.3， $\Gamma(X;I^\prime)$  与  $\mathfrak{s}(\Gamma(X;\mathcal{C}^\prime(\mathcal{U};I^\prime)))$  准同构，其中  $\mathfrak{s}(\cdot)$  表示与双复形相关的简单复形。从  $\Gamma(X;\mathcal{C}^\prime(\mathcal{U};I^\prime)) \simeq C^\prime(\mathcal{U};I^\prime)$  开始，需要注意的是，根据定理 1.9.3 和假设， $C^\prime(\mathcal{U};F)$  与  $\mathfrak{s}(C^\prime(\mathcal{U};I^\prime))$  准同构。   $\square$ 

\statementlabel{备注 2.8.6}覆盖  $ H^n(C^r(\mathcal{U}; F)) $  的上同调通常称为 Čech 上同调。

\subsection{实流形与复流形上层的例子}
我们将在这里解释一些例子，其中大部分我们将在第十一章中详细讨论。

2.9.1.层 $ \mathcal{C}_X^0 $ 。在拓扑空间  $ X $  上，与复值连续函数的空间  $ \mathcal{C}^0(U) $  的  $ X $  的开子集  $ U $  相关联的预层，以及通常的函数限制操作，显然是一个层。我们用 $ \mathcal{C}_X^0 $  表示它。请注意，常值层  $ \mathbb{Z}_X $  可以被识别为具有  $ \mathbb{Z} $  中的值的函数的  $ \mathcal{C}_X^0 $  的子束。

2.9.2.层 $ \mathcal{L}_{loc,dx}^1 $ 。令  $ U $  为欧几里得空间  $ \mathbb{R}^n $  的开子集，并令  $ L^1(U; dx) $  表示  $ U $  上的可积函数空间，关于

dx， $ \mathbb{R}^n $  上的勒贝格测度。预层 $ U \mapsto L^1(U; dx) $  不是捆。  $ \mathbb{R}^n $  上的关联层表示为  $ \mathcal{L}_{loc,dx}^1 $  。

2.9.3。环化空间。环化空间  $ (X, \mathcal{A}_X) $  是一个拓扑空间  $ X $  ，具有环层  $ \mathcal{A}_X $  。环化空间的态射，  $ f: (Y, \mathcal{A}_Y) \to (X, \mathcal{A}_X) $  是一个连续映射  $ f: Y \to X $  以及环层的态射  $ f^{-1} \mathcal{A}_X \to \mathcal{A}_Y $  。如果  $ A $  是一个环，而  $ \mathcal{A}_X $  是一堆  $ A $  -代数（即我们有一个态射  $ A_X \to \mathcal{A}_X $  ），那么  $ (X, \mathcal{A}_X) $  被称为  $ A $  -环化空间。

2.9.4.   $ C^*-manifolds $ 。令  $ \alpha $  为整数  $ (0 \leq \alpha < \infty) $  或  $ \alpha = \infty $  或  $ \alpha = \omega $  。我们用 $ \mathcal{C}_R^n $  表示 $ C^\alpha $  类复值函数 $ \mathbb{R}^n $  上的层（ $ C^\omega $  表示实解析）。维数  $ n $  的类  $ C^\alpha $  的实流形  $ C^\alpha $  是一个无限远可数的局部紧空间  $ M $ ，并赋予环层  $ \mathcal{C}_M^\alpha $ ，使得  $ (M, \mathcal{C}_M^\alpha) $  与  $ (\mathbb{R}^n, \mathcal{C}_R^n) $  局部同构为  $ \mathbb{C} $  环化空间。

其中一个表示  $\dim X$  （或  $\dim_{\mathbb{R}}X$  ），实流形  $X$  的维数。请注意，在文献中，通常用  $\mathcal{A}_{M}$  表示层  $\mathcal{C}_{M}^{\omega}$  。

我们参考 Guillemin – Pollack [1] 的微分几何基础课程。

2.9.5。方向、形式和密度。在 $C^0$  -流形 $M$  上，我们还必须考虑方向束 $\sigma_M$  。该层与  $\mathbb{Z}_M$  局部同构，并且  $M$  上方向的选择（如果存在）相当于选择同构  $\sigma_M \simeq \mathbb{Z}_M$  。我们将在下一章仔细学习 $\sigma_M$ 。

现在假设  $\alpha = \infty$  或  $\alpha = \omega$  。令 $p$  为整数。我们用 $\mathcal{C}_M^{\alpha,(p)}$  表示度 $p$  的微分形式束，系数在 $\mathcal{C}_M^\alpha$  中。我们用 $\mathrm{d} : \mathcal{C}_M^{\alpha,(p)} \to \mathcal{C}_M^{\alpha,(p+1)}$  表示外导数。

回想一下，如果  $(x_{1}, \ldots, x_{n})$  是  $M$  上的局部坐标系，则可以唯一地写出  $p$  形式  $f$ ：

\[
f=\sum_{|I|=p}f_{I}\mathrm{d}x_{I},
\]

其中  $ I = \{i_1, \ldots, i_p\} \subset \{1, \ldots, n\} $  、  $ (i_1 < i_2 < \cdots < i_p) $  、  $ \mathrm{d}x_I = dx_{i_1} \wedge \cdots \wedge dx_{i_p} $  和  $ f_I $  是  $ \mathcal{C}_M^a $  的一部分。然后：

\[
\mathrm{d}f=\sum_{i=1}^{n}\sum_{|I|=p}\frac{\partial f_{I}}{\partial x_{i}}dx_{i}\wedge dx_{I}.
\]

我们还介绍一下捆：

\[
\mathcal{V}_{M}^{\alpha}=\mathcal{C}_{M}^{\alpha,(n)}\otimes\operatorname{o r}_{M}
\]

 $ (\alpha = \infty \text{ or } \alpha = \omega) $  ，并将其称为 M 上的  $ C^{\alpha} $  密度束。

可以将  $C^{\infty}$  密度与紧凑型支持集成。我们用 $\int_{M}$ 表示积分图：

\[
\textstyle\int_{\cal M}\cdot:\Gamma_{c}(M;\mathcal{V}_{M}^{\infty})\to\mathbb{C}\enspace.
\]

层 $ \mathcal{C}_{M}^{\alpha,(p)} $  和 $ \mathcal{V}_{M}^{\alpha} $  是 $ \mathcal{C}_{M}^{\alpha} $  模的层。

根据“统一分区”的存在性，可以得出层 $ \mathcal{C}_M^\alpha, \mathcal{C}_M^{\alpha,(p)}, \mathcal{V}_M^\alpha $  是 $ c $  -软 $ \alpha \neq \omega $  。层  $ \mathcal{C}_M^\omega, \mathcal{C}_M^{\omega,(p)}, \mathcal{V}_M^\omega $  相对于函子  $ \Gamma(M; \cdot) $  是非循环的（即  $ H^j(M; \mathcal{C}_M^\omega) = 0 $  对应  $ j > 0 $  。参见 Grauert [1]）。

2.9.6。分布和超功能。在 $C^\infty$  -流形 $M$  上，我们还自然地定义了Schwartz 分布的束 $\mathcal{D}\ell_M$ （参见Schwartz [2]、de Rham [1]）。回想一下，  $\mathcal{D}\ell_M$  是  $c$  -软束， $\Gamma_c(M; \mathcal{D}\ell_M)$  是  $\Gamma(M; \mathcal{V}_M^\infty)$  的拓扑对偶，最后一个空间被赋予了 Fréchet 空间的自然拓扑。

在  $C^\omega$  -流形  $M$  上，我们可以类似地定义 Sato 超函数的束  $\mathcal{B}_M$ （参见 Sato [1]）。这是一个松弛的束， $\Gamma_c(M;\mathcal{B}_M)$  是  $\Gamma(M;\mathcal{V}_M^\omega)$  的拓扑对偶，最后一个空间被赋予了 DFS 空间的自然拓扑（参见 Martineau [1] 或 Schapira [1] 的详细说明）。然而，Sato 的构造是纯粹的上同调，我们将在下面的 2.9.13 小节中回顾它。

积分图 (2.9.1) 定义了一个配对：

\[
\left\{\begin{aligned}{\varGamma(M;\mathcal{C}_{M}^{\infty})\times\varGamma_{c}(M;\mathcal{V}_{M}^{\infty})\to\mathbb{C}}\\ {(f,g)\mapsto\smallint_{M}f g}\\ \end{aligned}\right..
\]

通过这种配对，我们得到了从  $ \mathcal{C}_M^\infty $  到  $ \mathcal{D}\ell_M $  的层态射，并且证明了这种态射是单射的。此外，在实解析流形  $ M $  上，注入  $ \Gamma(M; \mathcal{V}_M^\omega) \to \Gamma(M; \mathcal{V}_M^\infty) $  会引发态射  $ \mathcal{D}\ell_M \to \mathcal{B}_M $  ，它也是单射的。

人们还可以引入具有分布（或超函数）作为系数的 p 形式的层 $ \mathcal{D}\ell_M^{(p)} = \mathcal{C}_M^\infty, (p) \otimes \mathcal{C}_M^\infty \mathcal{D}\ell_M $ （或 $ \mathcal{P}_M^{(p)} = \mathcal{C}_M^\omega, (p) \otimes \mathcal{C}_M^\omega \mathcal{P}_M $ ）。这些是 c-软（或松弛）层。

2.9.7。德拉姆复合体。令 $M$  成为 $C^{\infty}$  流形。根据庞加莱引理，序列：

\[
0\to\mathbb{C}_{M}\to\mathcal{C}_{M}^{\infty,(0)}\xrightarrow[\mathtt{d}]{}\cdots\to\mathcal{C}_{M}^{\infty,(n)}\to0
\]

是准确的。因此，束 $ \mathbb{C}_M $  与 $ c $  软层的复合体准同构：

\[
\mathbb{C}_{M}\xrightarrow[{q i s}]{}{\mathrm{q i s}}(0\longrightarrow\mathcal{C}_{M}^{\infty,(0)}\xrightarrow[{d}]{}\cdots\longrightarrow\mathcal{C}_{M}^{\infty,(n)}\longrightarrow0).
\]

这允许显式计算上同调群  $ H^j(M;\mathbb{C}_M) $  或  $ H_c^j(M;\mathbb{C}_M) $  。例如，通过将  $ R\Gamma(M;\cdot) $  应用于 (2.9.4)，可以得到：

\[
R\varGamma(M;\mathbb{C}_{M})\simeq(0\longrightarrow\varGamma(M;\mathcal{C}_{M}^{\infty,(0)})\xrightarrow[\scriptstyle d]{}\cdots\longrightarrow\varGamma(M;\mathcal{C}_{M}^{\infty,(n)})\longrightarrow0)\enspace.
\]

将  $ \mathcal{C}_M^\infty $  替换为  $ \mathcal{D}\ell_M $  时，结果相同。如果  $ M $  是实解析的，则将  $ \mathcal{C}_M^\infty $  替换为  $ \mathcal{C}_M^\omega $  或  $ \mathcal{B}_M $  时，结果相同，但应注意  $ \mathcal{C}_M^\omega $  不再是  $ c $  -soft，而只是  $ \Gamma(M; \cdot) $  -非循环。另一方面， $ \mathcal{B}_M $  不仅是  $ c $  -软的，而且是松弛的，这允许计算  $ M $  中局部闭集  $ Z $  的上同调群  $ H_Z^j(M;\mathbb{C}_M) $  。

复合体 (2.9.3) 被称为 M 上的 de Rham 复合体。

2.9.8。复流形。让  $ \mathcal{O}_{C^n} $  表示  $ \mathbb{C}^n $  上的全纯函数束。维度  $ n $  的复流形  $ X $  是  $ \mathbb{C} $  环化空间  $ (X, \mathcal{O}_X) $  与  $ (\mathbb{C}^n, \mathcal{O}_{C^n}) $  局部同构。

其中  $\dim_{c}X$  表示复流形  $X$  的维数。我们参考 Wells [1] 来了解复微分几何的基本概念，并推荐 Banica-Stanasila [1] 的书来了解解析几何的进一步发展。

我们用  $\mathcal{O}_{X}^{(p)}$  表示  $X$  上的全纯 p 型束，并用  $\partial$  表示全纯微分。有时会设置：

\[
\begin{array}{r}{\Omega_{x}=\mathcal{O}_{X}^{(n)}\otimes\mathcal{o r}_{X},}\end{array}
\]

其中 $or_{X}$  是 $X$  上的定向层（参见第三章）。庞加莱引理对于全纯系数有效，束  $C_{X}$  与复形准同构：

\[
0\longrightarrow\mathcal{O}_{X}^{(0)}\xrightarrow[\partial]{}\cdots\longrightarrow\mathcal{O}_{X}^{(n)}\longrightarrow0~.
\]

2.9.9。多尔博复合体。令 $(X, \mathcal{O}_X)$  为复流形。我们用  $(\overline{X}, \mathcal{O}_{\overline{X}})$  表示具有  $X$  上反全纯函数束  $\mathcal{O}_{\overline{X}}$  的拓扑空间  $X$  。 （回想一下，如果  $\mathbb{C}$  上的复共轭映射与  $f: X \to \mathbb{C}$  的复合是全纯的，则  $f: X \to \mathbb{C}$  是反全纯的。）因此  $(\overline{X}, \mathcal{O}_{\overline{X}})$  是另一个复流形。

我们用  $ X^{\mathbb{R}} $  表示底层实解析流形。通过将  $ X^{\mathbb{R}} $  识别到  $ X \times \overline{X} $  的对角线，我们发现  $ X \times \overline{X} $  是  $ X^{\mathbb{R}} $  的复化。事实上：

\[
\mathcal{O}_{X\times\bar{X}}|_{X^{{{\mathsf{R}}}}}\simeq\mathcal{C}_{X^{{{\mathsf{R}}}}}^{\omega}\;.
\]

在 $X \times \bar{X}$ 上，我们可以考虑 $X$ 的全纯微分 $\partial$ 和 $\bar{X}$ 的 $\bar{\partial}$ 。因此  $X \times \bar{X}$  上的微分  $\partial$  被分解为  $\mathrm{d} = \partial + \bar{\partial}$  ，这会引起层  $\mathcal{C}_{X^{\mathbb{R}}}^{\alpha, (r)}(\alpha = \infty$  或  $\alpha = \omega)$  的分裂，如下所示

\[
\mathcal{C}_{X^{\mathtt{R}}}^{\alpha,(r)}=\bigoplus_{p+q=r}\mathcal{C}_{X}^{\alpha,(p,q)}~,
\]

其中  $ \mathcal{C}_{X}^{\alpha,(p,q)} $  是  $ X $  上的  $ (p,q) $  -形式的捆。在 $ X $  上的局部全纯坐标 $ (z_{1},\ldots,z_{n}) $  系统中， $ \mathcal{C}_{X}^{\alpha,(p,q)} $  的 $ f $  部分唯一地写为：

\[
{f}=\sum_{\substack{|I|=p,|J|=q}}f_{I,J}\mathbf{d}z_{I}\wedge\mathbf{d}\overline{{z}}_{J}~,
\]

其中 $ \mathrm{d}z_I = \mathrm{d}z_{i_1} \wedge \cdots \wedge \mathrm{d}z_{i_p} $ ， $ \mathrm{d}\overline{z}_J = \mathrm{d}\overline{z}_{j_1} \wedge \cdots \wedge \mathrm{d}\overline{z}_{j_q} $ ，类似于§9.5的符号。特别是：

\[
\partial f=\sum_{I}\sum_{J}\sum_{i=1}^{n}\frac{\partial f_{I,J}}{\partial z_{i}}\mathrm{d}z_{i}\wedge\mathrm{d}z_{I}\wedge\mathrm{d}\overline{z}_{J}.
\]

Dolbeault 引理断言复形：

\[
0\longrightarrow\mathcal{O}_{X}^{(p)}\longrightarrow\mathcal{C}_{X}^{\infty,(p,0)}\xrightarrow[\overline{{\partial}}]{\quad}\mathcal{C}_{X}^{\infty,(p,1)}\longrightarrow\cdots\longrightarrow\mathcal{C}_{X}^{\infty,(p,n)}\longrightarrow0
\]

是准确的，并且将  $ \mathcal{C}_X^{\infty,(p,q)} $  替换为  $ \mathcal{C}_X^{\omega,(p,q)} $  或  $ \mathcal{D}\ell_X^{(p,q)} $  或  $ \mathcal{B}_X^{(p,q)} $  会得到类似的结果。特别是，我们发现  $ \mathcal{O}_X^{(p)} $  与松弛层的复合体准同构（参见 Komatsu [1] 或 Schapira [1]）：

\[
0\longrightarrow\mathcal{B}_{X}^{(p,0)}\xrightarrow[\overline{{\partial}}]{\quad\cdots\quad}\cdots\longrightarrow\mathcal{B}_{X}^{(p,n)}\longrightarrow0\quad.
\]

正如 Golovin [1] 所示，这是一个单射  $ \mathcal{O}_{X} $  模的复合体。

2.9.10.对  $\mathcal{O}_X$  的操作。令 $f: Y \to X$  为复流形的态射。根据环化空间态射的定义，  $f$  诱导态射  $f^{-1}\mathcal{O}_X \to \mathcal{O}_Y$  。还有另一个态射可以在导出范畴  $\mathbf{D}^+(\mathbf{C}_\lambda)$  中定义：

\[
R f_{!}\varOmega_{Y}[\dim_{\mathbb{C}}Y]\to\varOmega_{X}[\dim_{\mathbb{C}}X]~.
\]

这种态射可以描述如下。

让 $ n = \dim_{\mathbb{C}} X $ ， $ m = \dim_{\mathbb{C}} Y $ ， $ l = m - n $ 。态射： $ f^{-1} \mathcal{C}_X^{\infty}, (m - p, m - q) \to \mathcal{C}_Y^{\infty}, (m - p, m - q) $  通过对偶性定义态射：

\[
f_{!}\mathcal{D}\ell_{Y}^{(p,q)}\otimes o\mathfrak{r}_{Y}\to\mathcal{D}\ell_{X}^{(p-l,q-l)}\otimes o\mathfrak{r}_{X}\enspace.
\]

然后从 (2.9.11) 以及  $ \Omega_{Y} $  和  $ \Omega_{X} $  的 Dolbeault 决议推导出 (2.9.10) 。

2.9.11.  $\mathcal{O}_X$  的上同调。我们参考 Hörmander [1] 对 $\mathcal{O}_X$  上同调的详细研究。让我们回想一下，如果  $\Omega$  在  $\mathbb{C}^n$  中开，那么如果  $H^j(\Omega; \mathcal{O}_X) = 0$  对于所有  $j > 0$  ，则可以说  $\Omega$  是伪凸的。例如，凸域是伪凸的，如果  $n = 1$  ，则所有域都是伪凸的。最后的结果概括如下：

\[
\left\{\begin{aligned}{}&{{}{i f}\;\Omega\;{i s~o p e n~i n}\;\mathbb{C}^{n},\quad}&{}&{{}{t h e n}}\\ {}&{{}H^{j}(\Omega;\mathcal{O}_{\mathbb{C}^{n}})=0\quad}&{}&{{}{f o r\quad a l l}j\geqslant n\enspace,}\\ \end{aligned}\right.
\]

（参见 Malgrange [1]，他使用 Dolbeault 分辨率获得了 (2.9.12)，并且方程  $ \left(\sum_{i=1}^{n}\frac{\partial}{\partial z_{i}}\frac{\partial}{\partial\overline{z}_{i}}\right)f = g $  在  $ \varGamma(\Omega;\mathcal{C}_{\mathbb{R}^{2n}}^{\infty}) $  中始终可解。

令 $X$  为维度 $n$  的复流形， $Z$  为 $X$  的局部闭子集， $x \in Z\setminus\mathrm{Int}Z$  。然后：

\[
H_{Z}^{j}(\mathcal{O}_{X})_{x}=0\qquad\text{对 }j\notin[1,n]\enspace.
\]

事实上，对于 $j=0$ 来说，这只不过是“解析连续原理”，而对于 $j>n$ 来说，这是从(2.9.12)或(2.9.9)得出的(即： $\mathcal{O}_{X}$ 具有松弛维度 $n$ ，参见练习II.9)。

由于 Martineau [2] 和 Kashiwara（参见 Sato-Kawai-Kashiwara [1]）， $ H_Z^j(\mathcal{O}_X) $  的消失有一个有用的标准。令  $ X = \mathbb{C}^n $  和  $ Z $  为  $ X $  和  $ x \in Z $  的闭凸子集。

\[
\begin{aligned}&\text{若不存在过 }x\text{ 的 }d\text{ 维仿射复空间 }L\text{，使得 }L\cap Z\text{ 是 }L\text{ 中 }x\text{ 的邻域，}\\&\text{则 }H_Z^j(\mathcal{O}_X)_x=0\quad\text{对 }j\leq n-d\text{。}\end{aligned}
\]

2.9.12.全纯函数的边界值。令  $\Omega$  为  $\mathbb{C}^n$  的严格伪凸开子集，具有  $C^2$  -边界。回想一下，在  $\partial\Omega$  上，本地存在坐标的全纯变化，它将  $\Omega$  与  $\mathbb{C}^n$  的严格凸开子集互换。

令  $j$  为嵌入  $\Omega \hookrightarrow \bar{\Omega}$  。我们在  $\bar{\Omega}$  上有一个显着的三角形：

\[
\mathcal{O}_{X}|_{\bar{\Omega}}\to R j_{*}\mathcal{O}_{\Omega}\longrightarrow R\Gamma_{\partial\Omega}(\mathcal{O}_{X}|_{\bar{\Omega}})\stackrel{\longrightarrow}{+1}\enspace.
\]

由于  $ H_{\partial\Omega}^{0}(\mathcal{O}_{X}|_{\bar{\Omega}})=0 $  ，预层  $ U\mapsto H_{U\cap\partial\Omega}^{1}(U\cap\bar{\Omega};\mathcal{O}_{X}|_{\bar{\Omega}}) $  等于束  $ H_{\partial\Omega}^{1}(\mathcal{O}_{X}|_{\bar{\Omega}}) $  ，（参见练习 II.13）。此外， $ H_{\partial\Omega}^{k}(\mathcal{O}_{X}|_{\bar{\Omega}})=0 $  代表 $ k>1 $  ，因为 $ R^{k}j_{*}\mathcal{O}_{\Omega}=0 $  代表 $ k>0 $  。我们还有：

\[
{t h e~s h e a f~}H_{\partial\Omega}^{1}(\mathcal{O}_{X}|_{\bar{\Omega}}){~i s~f l a b b y~}.
\]

为了证明(2.9.16)，我们可以根据命题2.4.10假设 $\Omega$ 是严格凸的。令  $U$  为  $\mathbb{C}^n$  的凸开子集。通过将函子  $R\Gamma(U;\cdot)$  应用于三角形 (2.9.15)，我们发现：

\[
\begin{array}{r}{\varGamma(U\cap\bar{\Omega};H_{\partial\Omega}^{1}(\mathcal{O}_{X}|_{\bar{\Omega}}))\simeq\mathcal{O}_{X}(\Omega\cap U)/\mathcal{O}_{X}(\bar{\Omega}\cap U).}\end{array}
\]

事实上， $ H^k(U \cap \bar{\Omega}; \mathcal{O}_X) = 0 $  代表 $ k > 0 $  ，因为 $ U \cap \bar{\Omega} $  承认 $ U $  中的凸开邻域的基本系统。

令  $\omega$  为  $\partial\Omega$  的开子集。存在  $\mathbb{C}^n$  的凸开子集  $U$  ，使得  $U\cap\overline{\Omega}=\omega$  和  $U\cup\Omega$  是凸的。然后是 Mayer-Victoris 序列：

\[
0\to\mathcal{O}_{X}(U\cup\Omega)\to\mathcal{O}_{X}(U)\oplus\mathcal{O}_{X}(\Omega)\to\mathcal{O}_{X}(U\cap\Omega)\to0
\]

是精确的，并且映射 $ \mathcal{O}_X(\Omega)/\mathcal{O}_X(\overline{\Omega}) \to \mathcal{O}_X(\Omega \cap U)/\mathcal{O}_X(\overline{\Omega} \cap U) $ 是满射的，这证明了(2.9.16)。

2.9.13.佐藤功能亢进。令  $M$  为维度  $n$  的实解析流形，并令  $X$  为  $M$  的复化。 （回想一下， $X$  在 $M$  的邻域中唯一定义）佐藤超函数的束 $\mathcal{B}_{M}$  定义为：

\[
\mathcal{B}_{M}=H_{M}^{n}(\mathcal{O}_{X})\otimes\mathfrak{o r}_{M/X}\;,
\]

其中 $ \sigma t_{M/X} = \sigma t_{M} \otimes \sigma t_{X} $ （参见第三章）。

请注意，根据（2.9.14），复数 $ R\Gamma_M(\mathcal{O}_X)[n] $ 集中在0度，并且我们有：

\[
\mathcal{B}_{M}\simeq R\varGamma_{M}(\mathcal{O}_{X})[n]\otimes o r_{M/X}.
\]

由于  $ j < n $  的束  $ H_M^j(\mathcal{O}_X) $  为零，我们发现（练习 II.13）预层  $ U \mapsto H_{U \cap M}^n(U; \mathcal{O}_X) $  是一个层，并且等于  $ \mathcal{B}_M $  。 （我们经常会识别  $ X $  上的层  $ \mathcal{B}_M $  及其对  $ M $  的限制。）此外，从 (2.9.12) 可以得出，该层是松弛的。该层与第 2.9.6 节中描述的层一致（更多详细信息，参见第 XI 章）。

2.9.14.局部常值层的示例。取  $X = \mathbb{C}$  ，并让  $z$  为  $X$  上的全纯坐标。令 $\alpha$  为复数，并令 $P$  为全纯微分运算符 $z \frac{\partial}{\partial z} - \alpha$  。考虑  $X$  上的层复合体：

\[
F:=0\longrightarrow\mathcal{O}_{X}\xrightarrow[P]{}\mathcal{O}_{X}\longrightarrow0\enspace.
\]

束  $ H^0(F)|_{X \setminus \{0\}} \simeq \mathrm{Ker}(P)|_{X \setminus \{0\}} $  是局部常值层，因为在  $ X \setminus \{0\} $  的每个开连通子集和单连通子集  $ U $  上，  $ H^0(F)|_{U} $  与由  $ z^\alpha $  （的分支）生成的常值层  $ \mathbb{C}_U $  同构。然而，如果  $ \alpha \notin \mathbb{Z} $  、  $ \Gamma(X \setminus \{0\}); H^0(F)) = 0 $  ，由于  $ X \setminus \{0\} $  上不存在非零全纯函数  $ f $  ，所以  $ Pf = 0 $  的解。其中一个为  $ \alpha \notin \mathbb{Z} $  ：

 $ H^0(F)|_{X \setminus \{0\}} $ ：等级一的局部常值层。

\[
H^{0}(F)|_{\{\mathbf{0}\}}=0,
\]

\[
H^{1}(F)=0.
\]

当  $ \alpha = 0, 1, 2, \ldots $  时，我们有：

\[
H^{0}(F)\simeq\mathbb{C}_{x},
\]

\[
H^{1}(F)\simeq\mathbb{C}_{\{0\}}~.
\]

当  $ \alpha = -1, -2, \ldots $  时，我们有：

\[
H^{0}(F)\simeq\mathbb{C}_{X\setminus\{0\}}\qquad\mathrm{a n d}\qquad H^{1}(F)=0~.
\]

复数 F 是所谓“反常层”的一个简单例子。此类配合物将在第八章和第十章中进行研究。

\subsection{习题}
\statementlabel{练习 II.1}令  $\mathbb{N}$  为具有拓扑的集合  $\mathbb{N}$  ，其开子集为集合  $\{0,1,\ldots,n\}$  、  $(n\geq-1)$  和  $\mathbb{N}$  。在识别  $\mathbb{N}$  上的预层和  $\mathbb{N}$  上的群射影系统后，证明预层  $F$  是一个层 iff  $\Gamma(\mathbb{N};F)=\lim F$  。也证明

(a)  $ H^j(\underline{\mathbb{N}}; F) = 0 $  为  $ j \neq 0 $  , 1。

（b） $ H^1(\mathbb{N}; F) = \{\{x_n\}_{n \in \mathbb{N}}; x_n \in F_n\}/\{\{x_n\}_{n \in \mathbb{N}}; x_n \in F_n \text{ and there are } y_n \in F_n \text{ with } x_n = y_n - y_{n+1}\} $ 。

\statementlabel{练习 二.2}令  $ A $  和  $ B $  为  $ X $  和  $ A \cup B = X $  的两个闭子集。对于  $ F \in \mathrm{Ob}(\mathbf{D}^+(X)) $  ，证明自然同构  $ (R\Gamma_B(F))_A \simeq R\Gamma_B(F_A) $  。

\statementlabel{练习 二.3}(i) 令 $U$  为 $X$  、 $x \in \bar{U} \setminus U$  的开子集。通过考虑层  $F = Z_U$  ，证明公式  $(\mathcal{H}om(F, G))_x \simeq \mathrm{Hom}(F_x, G_x)$  一般而言是假的。

(ii) 给出  $X$  上的束  $F$  的示例， $X$  的闭子集  $Z$  ，  $X$  的开子集  $U$  ，使得  $Z \cap U = \varnothing$  但  $R\varGamma_Z(F_U) \neq 0$  。观察  $\varGamma_Z(F_U) = 0$  ，因此在这种情况下，组合的派生函子不是派生函子的组合。

\statementlabel{练习 二.4}证明可收缩空间 X 上的局部常值层 F 是常值层。

\statementlabel{练习 二.5}令 $X$  为仿紧空间。如果对于  $X$  的所有闭子集  $Z$  ，自然映射  $\Gamma(X; F) \to \Gamma(Z; F)$  是满射的，则有人说  $X$  上的束  $F$  是软的。证明如果  $F$  是软的，则  $H^i(X; F) = 0$  对于  $i > 0$  。

\statementlabel{练习 二.6}设 X 为局部紧空间，F 为 X 上的层。

(a) 证明 $F$  是 $c$  -软当且仅当 $H_{c}^{i}(U; F) = 0$  对于所有 $i > 0$  、 $X$  的所有开子集 $U$  ； （有人写 $H_{c}^{i}(U; F)$ 而不是 $H_{c}^{i}(U; F|_{U}$ ）。

(b) 此外假设 X 在无穷大处是可数的。证明如果 F 是 c-soft，则 F 是软的。

(c) 证明 c-soft 的性质在 X 上是局部的。

\statementlabel{练习 二.7}令 R 为 c-软环层。证明任何 R 模束都是 c 软的。

\statementlabel{练习 二.8}设 $X$  为无穷大可数的局部紧致空间。如果对于 $X$  的开子集 $U$  的任何一对闭子集 $Z_1$  和 $Z_2$ ， $X$  上的束 $F$  被称为 $\text{supple}$ ， $\Gamma_{Z_1}(U; F) \oplus \Gamma_{Z_2}(U; F) \to \Gamma_{Z_1 \cup Z_2}(U; F)$  是满射的（参见Bengel-Schapira [1]）。

(a) 证明松弛的层是柔软的。

(b) 证明如果层  $F$  是柔韧的，则对于任何闭子集  $Z$  来说， $\Gamma_{Z}(F)$  是  $c$  -软的。

(c) 证明柔顺特性是 X 上的局部特性。

（注：在实解析流形  $M$  上，束  $\mathcal{D}\ell_M$  是  $c$  -软但不柔软。另一方面，商束  $\mathcal{D}\ell_M/\mathcal{A}_M$  是柔软但不松弛 - 束  $\mathcal{D}\ell_M$  和  $\mathcal{A}_M$  是第 9 节中定义的。）

\statementlabel{练习 二.9}设X为拓扑空间。

(a) 对于一个层  $F$  和一个非负整数  $n$  ，证明以下条件是等价的。

(i) 存在一个精确的序列  $ 0 \to F \to F^0 \to \cdots \to F^n \to 0 $ ，其中  $ F^j $  是松弛的。

(ii) 如果  $ 0 \to F \to F^0 \to \cdots \to F^n \to 0 $  是精确的，并且如果  $ 0 \leq j < n $  的  $ F^j $  是松弛的，那么  $ F^n $  是松弛的。

(iii) 对于任何闭集 S， $ H_{S}^{k}(X; F) = 0 $  对于 k > n。

(iv) 对于任何局部闭集 S， $ H_S^k(X; F) = 0 $  对于 k > n。

(v) 对于任何闭集 S， $ H_{S}^{k}(F) = 0 $  对于 k > n。

(vi) 对于任何局部闭集 S， $ H_{S}^{k}(F) = 0 $  对于 k > n。

满足这些等效条件的最小  $n$  称为  $F$  的松弛维数。  $X$  上所有层 $F$  的松弛尺寸的最大值称为 $X$  的松弛尺寸。

(b) 对于局部紧空间 X，定义 X 上层 F 的 c 软维数和 X 的 c 软维数，并给出这种情况下相应的等效条件 (i)-(iv)。

(c) 证明，对于层 $F$  和局部紧空间 $X$ ， $F \leq$  的 $c$  -软维数， $F \leq 1$  的松弛维数+ $F$  的 $c$  -软维数。

\statementlabel{练习 二.10}令  $ \mathcal{R} $  为  $ X $  上的环层，并令  $ M \in \mathrm{Ob}(\mathfrak{Mod}(\mathcal{R})) $  。

(a) 证明 $M$ 是单射的当且仅当对于 $\mathcal{R}$ 的任何子 $\mathcal{R}$ 模 $\mathcal{I}$ （我们说 $\mathcal{I}$ 是 $\mathcal{R}$ 的理想），自然同态：

\[
\begin{array}{r}{\varGamma(X;M)\simeq\mathrm{H o m}_{\mathcal{R}}(\mathcal{R},M)\to\mathrm{H o m}_{\mathcal{R}}(\mathcal{I},M)}\end{array}
\]

是主观的。

(b) 令 $A$  为一个字段。证明  $A_{x}$  的任何理想都与层  $A_{U}$  同构，其中  $U$  在  $X$  中开放。由此推断，当且仅当束 $M$  是松弛的时， $A_{x}$  -模 $M$  是单射的。

\statementlabel{练习 二.11}令 $f: Y \to X$  为局部紧空间的连续映射。让 $G$  成为 $Y$  上的一个层。证明下面两个条件等价。

(a) 对于任何  $ x \in X $  ， $ G|_{f^{-1}(x)} $  都是 c-soft。

(b) 对于 Y 的任何开子集 U，任何 j > 0， $ R^{j}f_{!}G_{U} = 0 $  。

\statementlabel{练习 二.12}设X为拓扑空间。

(a) 令  $(F_\lambda)_{\lambda \in \mathcal{A}}$  为  $X$  上的层归纳系统，由有向有序集  $\mathcal{A}$  索引。假设  $X$  是紧致且 Hausdorff，证明所有  $k \in \mathbb{N}$  的同构  $\lim_{\overrightarrow{\lambda}} H^k(X; F_\lambda) \simeq H^k(X; \varliminf_{\overrightarrow{\lambda}} F_\lambda)$  。

(b) 令 $ (F_n)_{n \in \mathbb{N}} $  为 $ X $  上层的射影系统，并令 $ Z $  为 $ X $  的局部闭子集。假设  $ \{H_Z^{k-1}(X; F_n)\}_n $  满足  $ M $  -  $ L $  条件，证明同构  $ H_Z^k(X; \lim_{\leftarrow} F_n) \cong \lim_{\leftarrow} H_Z^k(X; F_n) $  。

\statementlabel{练习 二.13}令  $F$  为  $X$  上的一个层，并让  $Z$  为  $X$  的局部封闭子集。假设  $R^j\varGamma_Z(F) = 0$  为  $j < n$  。证明预层  $U \mapsto H_Z^n(U; F)$  是一个层，并且等于  $R^n\varGamma_Z(F)$  ，（参见练习 I.22）。

\statementlabel{练习 二.14}让  $ X = \bigcup_{i \in I} U_i $  成为  $ X $  的开放覆盖。

对于每个  $i \in I$  ，让  $F_i$  成为  $U_i$  上的一束，对于每对  $(i,j)$  ，让  $\phi_{ij}$  成为同构  $F_j|_{U_i \cap U_j} \xrightarrow{\sim} F_i|_{U_i \cap U_j}$  。假设  $\phi_{ii} = \mathrm{id}_{F_i}$  和  $\phi_{ij} \circ \phi_{jk} = \phi_{ik}$  上

 $ U_i \cap U_j \cap U_k $ 。然后证明  $ X $  上存在一个层  $ F $  ，并且对于每个  $ i $  ，同构  $ \phi_i: F|_{U_i} \simeq F_i $  ，使得  $ \phi_{ij} = \phi_i \circ \phi_j^{-1} $  位于  $ U_i \cap U_j $  上，对于所有对  $ (i, j) $  。证明  $ F $  在同构方面是唯一的。

\statementlabel{练习 二.15}(i) 令 F 为 X 上层的下界复形。构造自然态射：

\[
H^{j}(\varGamma(X;F^{\cdot}))\to H^{j}(R\varGamma(X;F^{\cdot}))\enspace.
\]

(ii) 令 $ \mathcal{U} = \{U_i\}_{i} $  为 $ X $  的开放覆盖物，并令 $ F $  为 $ X $  上的一个层。构造规范态射：

\[
H^{j}(C^{\cdot}(\mathcal{U};F))\to H^{j}(X;F)\enspace.
\]

\statementlabel{练习 二.16}令 $A$  为交换环， $A^\times$  为 $A$  的可逆元素群。令  $X$  为拓扑空间， $\mathcal{U} = \{U_i\}_{i \in I}$  为  $X$  的开放覆盖，并令  $c \in C^2(\mathcal{U}; A_X^\times)$  与  $\delta c = 0$  （参见第 8 节）。令 $c'$  为 $H^2(C^\prime(\mathcal{U}; A_X^\times))$  中 $c$  的类， $c''$  为 $H^2(X; A_X^\times)$  中 $c'$  的图像（参见练习II.15）。定义一个范畴 $\mathfrak{Sh}(X; c)$ ，其对象是系列 $\{F_i, \rho_{ij}\}$ ， $F_i$ 是 $A_{U_i}$ 模， $\rho_{ij} : F_j|_{U_i \cap U_j} \simeq F_i|_{U_i \cap U_j}$ 是同构，这样：

\[
\rho_{i j}\rho_{j k}\rho_{k i}=\mathbf{c}_{i j k}\cdot\mathbf{i d}_{F_{i}|U_{i}\cap U_{j}\cap U_{k}}\qquad\mathrm{f o r~a l l~}i,j,k,
\]

并且这个范畴的态射是用显而易见的方式定义的。

(i) 证明 $\mathfrak{Sh}(X;c)$  是阿贝尔范畴。

(ii) 令  $ \tilde{c} $  为  $ C^2(\mathcal{U}; A_X^\times) $  的另一个元素并假设  $ \tilde{c}'' = c'' $  。证明  $ \mathfrak{Sh}(X; c) $  和  $ \mathfrak{Sh}(X; \tilde{c}) $  之间存在范畴的等价性。

\statementlabel{练习 二.17}令 $X$  为局部紧空间， $\mathcal{R}$  为 $X$  与 $\mathrm{wgld}(\mathcal{R}) < \infty$  上的交换环层，并令 $Z_1$  和 $Z_2$  为 $X$  的两个局部闭子集。

(i) 对于  $\mathbf{D}^{+}(\mathcal{R})$  中的  $F_1$  和  $F_2$  ，构造自然态射：

\[
R\varGamma_{Z_{1}}(F_{1})\overset{L}{\otimes}_{\mathcal{R}}R\varGamma_{Z_{2}}(F_{2})\to R\varGamma_{Z_{1}\cap Z_{2}}\bigg(\overset{L}{F_{1}}\overset{\bigotimes}{\otimes}_{\mathcal{R}}F_{2}\bigg)\;.
\]

(ii) 假设 $ \mathscr{R} = A_X $ （其中 $ A $  是交换环）。构造态射：

\[
R\varGamma_{\mathbb{Z}_{1}}(X;F_{1})\overset{L}{\otimes}_{\cal A}R\varGamma_{\mathbb{Z}_{2}}(X;F_{2})\to R\varGamma_{\mathbb{Z}_{1}\cap\mathbb{Z}_{2}}\biggl(\overset{L}{X};F_{1}\overset{}{\otimes}_{\cal A}F_{2}\biggr)
\]

和态射  $ (p, q \in \mathbb{Z}) $  ：

\[
H_{Z_{1}}^{p}(X;F_{1})\mathop{\otimes}_{\cal A}H_{Z_{2}}^{q}(X;F_{2})\to H_{Z_{1}\cap Z_{2}}^{p+q}\left(\begin{matrix}{L}\\ {X;F_{1}\mathop{\otimes}_{\cal A}F_{2}}\\ \end{matrix}\right)\;.
\]

最后的态射通常称为杯积。 （提示：使用练习 I.24。）

\statementlabel{练习 二.18}令  $f_i: Y_i \to X_i$  、  $i = 1, 2$  为空间  $S$  上局部紧空间的连续映射（参见 (2.7.3)）。让  $p_{Y_i}$  表示映射  $Y_i \to S$  并设置  $f = f_1 \times_S f_2: Y_1 \times_S Y_2 \to X_1 \times_S X_2$  。令  $\mathcal{R}$  为  $S$  上的交换环层，与  $\mathrm{wgld}(\mathcal{R}) < \infty$  ，并令  $G_i \in \mathrm{Ob}(\mathbf{D}^+) (p_{Y_i}^{-1} \mathcal{R})$  。

(i) 证明同构

\[
R f_{1!}G_{1}\mathop{\boxtimes}_{S}^{\cal L}{\cal R}R f_{2!}G_{2}\xrightarrow{\sim}R f_{!}\left(G_{1}\mathop{\boxtimes}_{S}^{\cal L}{\cal R}G_{2}\right)\;.
\]

这种同构被称为 Künneth 公式。 （提示：将问题简化为  $ S = X_1 = X_2 = \{\text{pt}\} $  的情况。然后使用同构（2.6.19），类似于下面命题 3.1.15 的证明。）

(ii) 假设 $ S = X_1 = X_2 = \{\text{pt}\} $  和 $ \mathcal{R} $  是一个字段。证明：

\[
H_{c}^{n}(Y_{1}\times Y_{2};G_{1}\boxtimes G_{2})\simeq\bigoplus_{p+q=n}(H_{c}^{p}(Y_{1};G_{1})\otimes H_{c}^{q}(Y_{2};G_{2})).
\]

（提示：使用练习 I.24。）

\statementlabel{练习 二.19}假设  $X$  是局部紧致的，并且  $A$  与  $\text{wgld}(A) < \infty$  是可交换的。令  $F \in \text{Ob}(\mathbf{D}^{+}(A_X))$  ，令  $\Omega$  （或  $Z$  ）为  $X$  的开（或闭）子集，并用  $a_X$  表示映射  $X \to \{\text{pt}\}$  。证明同构：

\[
R\varGamma(\varOmega;F)\simeq R\mathfrak{a}_{X\ast}R\mathcal{H}\mathfrak{o}m(A_{\varOmega},F)\enspace,
\]

\[
R\varGamma_{c}(\varOmega;F)\simeq R\mathsf{a}_{X!}\biggl(A_{\varOmega}\overset{L}{\otimes}F\biggr),
\]

\[
R\Gamma_{Z}(X;F)\simeq R a_{X*}R\mathcal{H}o m(A_{Z},F)\enspace,
\]

\[
R\varGamma_{c}(Z;F)\simeq R\mathsf{a}_{X!}\biggl(A_{Z}\overset{L}{\otimes}F\biggr).
\]

\statementlabel{练习 二.20}（参见练习 IX.10）。令 $A$  为 $\mathrm{wgld}(A)$  有限的交换环，并令 $E$  为实数有限维向量空间。在  $\mathrm{Ob}(\mathbb{D}^{+}(A_E))$  上，通过设置  $F*G = Rs_1(F \boxtimes^L G)$  来定义卷积运算，其中  $s$  是  $\mathrm{map}E \times E \to E$  、 $(x, y) \mapsto x + y$  。

(a) 证明对于  $\mathbf{D}^{+}(A_E)$  中的  $F$  、  $G$  、  $H$  ，我们有  $F*G \simeq G*F$  、  $F*(G*H) \simeq (F*G)*H$  、  $A_{\{0\}}*F \simeq F$  。

(b) 令 $ Z_1 $  和 $ Z_2 $  为E 的两个紧凸子集。证明 $ A_{Z_1} * A_{Z_2} \simeq A_{Z_1 + Z_2} $  。

(c) 证明  $ A_{\gamma} * A_{\text{int}\gamma} = 0 $  ，其中  $ \gamma $  是真闭凸锥。

(d) 假设  $ E = \mathbb{R}^n $  。让 $ Z_1 = [-1,1]^n $ ， $ Z_2 =] -1,1[^n $ 。证明 $ A_{Z_1} * A_{Z_2} \simeq A_{\{0\}}[-n] $ 。

\statementlabel{练习 二.21}令 $X$  为拓扑空间， $\{X_n\}_{n\in\mathbb{Z}}$  为 $n \ll 0$  的闭子集的递减序列，其中 $\bigcap_n X_n = \emptyset$  和 $X_n = X$  。让 $F \in \mathrm{Ob}(\mathbf{D}^+) (X)$ 

令人满意： $ H_{X_n \setminus X_{n+1}}^k(F) = 0 $  对于 $ k \neq n $  。使用区分的三角形 $ R\Gamma_{X_{n+1} \setminus X_{n+2}}(F) \longrightarrow R\Gamma_{X_n \setminus X_{n+2}}(F) \longrightarrow R\Gamma_{X_n \setminus X_{n+1}}(F) \xrightarrow{+1} $ ，定义 $ d^n $ ： $ H_{X_n \setminus X_{n+1}}^n(F) \longrightarrow H_{X_{n+1} \setminus X_{n+2}}^{n+1}(F) $ ，一组定义 $ K^n = H_{X_n \setminus X_{n+1}}^n(F) $ 。

(i) 证明 $(K',d')$  是 $X$  上层的复合体。

(ii) 证明 $ H_{X_n}^k(F) = 0 $  对于 $ k < n $  并证明同构 $ H_{X_n}^n, (F) \simeq H^n(F) $  。

(iii) 设  $ \tilde{G}^n = \Gamma_{X_n}(F^n) \cap (\mathrm{d}_F^n)^{-1}(\Gamma_{X_{n+1}}(F^{n+1})) $  ，其中  $ F $  是层  $ \{F^n, \mathrm{d}_F^n\}_{n \in \mathbb{Z}} $  的复合体。构造态射  $ \mathrm{d}_G^n: G^n \to G^{n+1} $  ，证明  $ G = (G^*, \mathrm{d}_G^-) $  是复形，构造态射  $ G \to K $  和  $ G \to F $  并证明如果所有  $ F^n $  都是松弛的，则这些是拟同构。得出  $ F \simeq K $  中的  $ \mathbf{D}^+(X) $  的结论。 （参见哈特肖恩 [1]。）

\statementlabel{练习 二.22}令  $ X = U_1 \cup U_2 $  为  $ X $  的开或闭覆盖，并令  $ F_i \in \text{Ob}(\mathbf{D}^+) (U_i) $  、  $ i = 1, 2 $  。令  $ \phi $  成为  $ \mathbf{D}^+(U_1 \cap U_2) $  中的同构  $ F_1|_{U_1 \cap U_2} \sim F_2|_{U_1 \cap U_2} $  。构造  $ F \in \text{Ob}(\mathbf{D}^+) (X) $  和同构  $ \psi_i: F|_{U_i} \simeq F_i $  使得  $ \psi_2|_{U_1 \cap U_2} = \phi \circ \psi_1|_{U_1 \cap U_2} $  。

\statementlabel{练习 二.23}令 $ X = \mathbb{R}^n $  与 $ n \geq 1 $  并让 $ k $  成为交换域。证明如果  $ P \in \text{Ob}(\mathfrak{Mod}(k_X)) $  是射影的，则  $ P = 0 $  。 （提示：假设  $ P \neq 0 $  ，证明存在  $ X $  的非空开子集  $ U $  ，使得  $ k_U $  是  $ P $  的直和，并推导出矛盾。）

\statementlabel{练习 二.24}令 $k$  为交换环。在本练习中，“环”是指环  $A$  与态射  $k \to A$  一起，使得  $k$  的图像位于  $A$  的中心。令 $X$  为拓扑空间，让 $A, B, C, D$  为 $X$  上的环层。我们假设  $A, B, C, D$  比  $k$  平坦。

(i) 证明平坦（或单射） $(\mathcal{A} \otimes_k \mathcal{B})$  -模在  $\mathcal{A}$  上平坦（或单射）。

(ii) 证明函子 $ \otimes_{\mathcal{G}} \colon $ 

\[
\mathfrak{M o d}\left(\mathcal{A}\mathop{\otimes}_{k}\mathcal{B}^{o p}\right)\times\mathfrak{M o d}\left(\mathcal{B}\mathop{\otimes}_{k}\mathcal{C}^{o p}\right)\to\mathfrak{M o d}\left(\mathcal{A}\mathop{\otimes}_{k}\mathcal{C}^{o p}\right)
\]

有一个左派生函子 $ \otimes^L_{\mathcal{B}} \colon $ 

\[
\mathbb{D}^{-}\left(\mathcal{A}\mathop{\otimes}_{k}\mathcal{B}^{o p}\right)\times\mathbb{D}^{-}\left(\mathcal{B}\mathop{\otimes}_{k}\mathcal{C}^{o p}\right)\to\mathbb{D}^{-}\left(\mathcal{A}\mathop{\otimes}_{k}\mathcal{C}^{o p}\right)
\]

并证明  $ \mathbf{D}^{-}(\mathcal{A} \otimes_k \mathcal{D}^{op}) $  中的同构：

\[
\left(K\mathop{\otimes}_{\mathcal{B}}^{L}M\right)\mathop{\otimes}_{\mathcal{C}}^{L}N\simeq K\mathop{\otimes}_{\mathcal{B}}^{L}\left(M\mathop{\otimes}_{\mathcal{C}}^{L}N\right)
\]

对于  $ K \in \mathrm{Ob}(\mathbb{D}^-(\mathcal{A} \otimes_k \mathcal{B}^{op})) $  、  $ M \in \mathrm{Ob}(\mathbb{D}^-(\mathcal{B} \otimes_k \mathcal{C}^{op})) $  、  $ N \in \mathrm{Ob}(\mathbb{D}^-(\mathcal{C} \otimes_k \mathcal{D}^{op})) $  。

(iii) 证明函子  $ \mathrm{Hom}_{\mathcal{A}}(\cdot,\cdot) $  ：

\[
\mathfrak{M o d}\left(\mathcal{A}\mathop{\otimes}_{k}\mathcal{B}^{o p}\right)^{\circ}\times\mathfrak{M o d}\left(\mathcal{A}\mathop{\otimes}_{k}\mathcal{C}^{o p}\right)\to\mathfrak{M o d}\left(\mathcal{B}\mathop{\otimes}_{k}\mathcal{C}^{o p}\right)
\]

有一个右派生函子  $R\mathcal{H}_{om_{A}}(\cdot,\cdot)$  ：

\[
\mathbf{D}^{-}\left(\mathcal{A}\mathop{\bigotimes}_{k}\mathcal{D}^{o p}\right)^{\circ}\times\mathbf{D}^{+}\left(\mathcal{A}\mathop{\bigotimes}_{k}\mathcal{C}^{o p}\right)\to\mathbf{D}^{+}\left(\mathcal{B}\mathop{\bigotimes}_{k}\mathcal{C}^{o p}\right),
\]

并证明 $ \mathbf{D}^{+}(\mathcal{C} \otimes_k \mathcal{D}^{op}) $ 中的同构

\[
{~R~}\mathcal{H}{_{O m}}_{\mathcal{A}}\bigg(M\mathop{\mathop{\otimes}_{\mathcal{B}}}^{L}N,K\bigg)\simeq{R}\mathcal{H}{_{O m}}_{\mathcal{B}}(N,{R}\mathcal{H}{_{O m}}_{\mathcal{A}}(M,K))
\]

对于  $ M \in \mathrm{Ob}(\mathbb{D}^{-}(\mathcal{A} \otimes_k \mathcal{B}^{op})) $  、  $ N \in \mathrm{Ob}(\mathbb{D}^{-}(\mathcal{B} \otimes_k \mathcal{C}^{op})) $  和  $ K \in \mathrm{Ob}(\mathbb{D}^{+}(\mathcal{A} \otimes_k \mathcal{D}^{op})) $  。

(iv) 令 $f: Y \to X$  为连续映射， $\mathcal{B}$  和 $\mathcal{A}$  分别为 $Y$  和 $X$  上的环层，均平坦于 $k$  上。假设存在  $n \in \mathbb{N}$  使得  $\mathrm{wgyd}(\mathcal{A}_x) \leq n$  对于所有  $x \in X$  ，为  $K \in \mathrm{Ob}(\mathbf{D}^b(\mathcal{B} \otimes_k f^{-1} \mathcal{A}^{op}))$  构造一个关于  $M \in \mathrm{Ob}(\mathbf{D}^-(\mathcal{A}))$  和  $N \in \mathrm{Ob}(\mathbf{D}^+(\mathcal{A}))$  的态射、函子：

\[
\mathfrak{f}^{-1}R\mathcal{H o m}_{\mathcal{A}}(M,N)\to R\mathcal{H o m}_{\mathcal{B}}\Biggl(\underset{}{K}{\overset{L}{\bigotimes_{\mathfrak{f}^{-1}\mathcal{A}}}f^{-1}}M,K\underset{}{_{\mathfrak{f}^{-1}\mathcal{A}}}{\overset{L}{\bigotimes}}f^{-1}N\Biggr)
\]

在 $ \mathbf{D}^{+}(k_{Y}) $  中。

\statementlabel{练习 二.25}令 $X$  为拓扑空间， $S$  为闭子集，并令 $\{S_k\}_{k \in \mathbb{Z}}$  为 $S$  的闭子集的递减序列，使得 $S_k = S$  对于 $k \ll 0$  和 $S_k = \varnothing$  对于 $k \gg 0$  。让 $F \in \mathrm{Ob}(\mathbb{D}^+) (X)$  并让 $r \in \mathbb{Z}$  。假设  $H_{S_k}^j(F)|_{S_k \setminus S_{k+1}} = 0$  对于  $j < r$  和所有  $k$  。证明  $H_S^j(F) = 0$  对于任何  $j < r$  。

（提示：通过对  $k$  进行归纳论证，假设  $H_S^j(F)|_{X\setminus S_{k-1}}=0$  为  $j<r$  并考虑不同的三角形：

\[
R\Gamma_{(S_{k-1}\setminus S_{k})}(F)|_{X\setminus S_{k}}\longrightarrow R\Gamma_{S}(F)|_{X\setminus S_{k}}\longrightarrow R i_{*}i^{-1}R\Gamma_{S}(F)\xrightarrow[+1]{}
\]

其中  $i$  表示开放嵌入  $X\backslash S_{k-1}\to X\backslash S_k$  。

\statementlabel{练习 二.26}令  $ F \in \mathrm{Ob}(\mathbb{D}^b(\mathbb{Z}_X)) $  和  $ \mathcal{U} = \{U_j\}_{j=1,\ldots,r} $  为  $ X $  的有限开覆盖。构造自然态射：

\[
H^{0}\left(\bigcap_{j}U_{j};F\right)\to H^{r-1}(X;F).
\]

（提示：应用练习 II.15。）

\subsection{注记}
我们在本书的开头参考 C. Houzel 的《简史》，了解层理论的详细历史。

让我们简单回顾一下，Leray [1,2] 于 1945 年左右引入了束，以及谱序列，这是当时研究束上同调的基本工具。该理论在嘉当 [1] 的推动下发挥了其全部重要性，也许可以将现代形式主义的层上同调与东北大学的格罗滕迪克 [1] 的论文进行比较。 Godement [1] 的书也极大地促进了该理论的普及，Grothendieck 强调了导出范畴框架中层函子运算的重要性，首先是在相干情况下（Grothendieck [4]），然后是在离散系数情况下（[SGA 4]，[SGA 5]）。

在第二章中，读者将会遇到“Čech上同调”、“Künneth公式”、“Mayer-Vietoris序列”、“Vietoris-Begle定理”、“庞加莱引理”、“de Rham复形”等。在层理论的其他重要贡献中，我们要提到Godement对松弛层和软层的介绍，以及局部上同调（函子）的介绍。  $ \varGamma_{z}(\cdot) $  及其派生函子）作者：Grothendieck [1]。还应该提到的是，将在第五章中发挥关键作用的“非特征变形引理”（命题2.7.2）是由Kashiwara提出的[3,5]。它是莫尔斯理论的一种变体，基于所谓的格罗腾迪克的米塔格-莱弗勒定理 [3]。

\section{庞加莱--韦尔迪耶对偶与傅里叶--佐藤变换}
\subsection{概要}
在本章中，我们介绍了研究流形层的两个基本工具。首先，按照 Verdier [1]，我们构造函子  $ Rf_1 $  的右伴随  $ f^1 $  。如果 $ f: Y \to X $ 是满足适当条件的局部紧空间的连续映射，并且如果 $ F $ （或 $ G $ ）属于 $ \mathbf{D}^+(A_X) $ （或 $ \mathbf{D}^+(A_Y) $ ），则得到公式：

\[
\operatorname{H o m}(R f_{!}G,F)=\operatorname{H o m}(G,f^{!}F)\enspace.
\]

这概括了经典的庞加莱对偶性。事实上，如果  $X$  简化为单点，如果  $F = A$  和  $Y$  是拓扑流形，则可以计算  $f^!A$  （在  $Y$  上称为“对偶复形”），并表明它与  $Y$  上的方向束同构，并移动了维度。有了函子 $f^!\in$ ，就可以推导出许多漂亮的层理论新公式。

其次，我们详细研究了傅立叶-佐藤变换，这是一种交换矢量丛上的圆锥层（在导出范畴中）和对偶矢量丛上的圆锥层的操作。

在本章中，我们研究拓扑流形上的层：定向层、松弛维数、上同调构造层，并且我们还介绍 $ \gamma $ 拓扑，作为第五章的准备。

对于§1-§4 内容的另一种方法，人们还可以参考 Borel 等人 [1]、Gelfand-Manin [1]、Iversen [1] 和 [SHS]。

本章的许多材料可以被认为是经典的，尽管它们事先没有经过系统的处理。

公约3.0。在这一章以及除第 XI 章之外的所有后续章节中，为了使演示不那么复杂，我们将在拓扑空间  $X$  上使用  $A_{x}$  模的层，其中假定基环  $A$  与有限全局维度可交换。 （回想一下，参见练习 I.29， $\text{wgld}(A) \leq \text{gld}(A)$  。）

如果没有混淆的风险，我们分别写 $ \mathbf{D}^{+}(X) $ 或 $ \mathbf{D}^{b}(X) $ 而不是 $ \mathbf{D}^{+}(A_{X}) $ 或 $ \mathbf{D}^{b}(A_{X}) $ 。我们还写  $ F \otimes^{L}G $  和  $ R\mathcal{H}om(F,G) $  而不是  $ F \otimes_{Ax}^{L}G $  和  $ R\mathcal{H}om_{Ax}(F,G) $  。

所有流形都是有限维且无穷大可数的。子流形始终是局部封闭的。

在本章中，除了与  $ \gamma $  拓扑相关的拓扑空间外，所有拓扑空间都假设为局部紧致的。

当组合两个函子时，例如 $ f^!\circ Rf_t $ ，我们不应该总是写符号 $ \circ $ 。

\subsection{庞加莱--韦尔迪耶对偶}
令 $f: Y \to X$  为局部紧空间的连续映射。本节的目的是构造函子  $Rf_!: \mathbf{D}^+ (A_Y) \to \mathbf{D}^+ (A_X)$  的右伴随  $f^!$  。因此函子  $f^!$  将满足：

\[
\mathrm{H o m}_{\texttt{D}^{+}(A_{X})}(R f_{!}G,F)=\mathrm{H o m}_{\texttt{D}^{+}(A_{Y})}(G,f^{!}F)
\]

对于 $ F \in \mathrm{Ob}(\mathbb{D}^{+}(A_X)) $ ， $ G \in \mathrm{Ob}(\mathbb{D}^{+}(A_Y)) $ 。

如果考虑  $Y$  是  $n$  维定向流形  $X = \{\mathrm{pt}\}$  、  $A = \mathbb{Q}$  、  $G = \mathbb{Q}_Y$  、  $F = \mathbb{Q}_{\{\mathrm{pt}\}}$  的特殊情况，我们将看到  $f^1 \mathbb{Q}_{\{\mathrm{pt}\}} \simeq \mathbb{Q}_Y[n]$  ，因此：

\[
\operatorname{H o m}(R\varGamma_{c}(Y;\mathbb{Q}_{Y})[n],\mathbb{Q})\simeq R\varGamma(Y;\mathbb{Q}_{Y})\enspace.
\]

取第 j 个上同调群，我们得到：

\[
(H_{c}^{n-j}(Y;\mathbb{Q}_{Y}))^{*}\simeq H^{j}(Y;\mathbb{Q}_{Y}),
\]

其中 $ ^{*} $ 表示Q上的对偶向量空间。这是经典的庞加莱对偶性。

在开始构建 $f^!,$ 之前，让我们启发式地解释一下。对于  $F \in \mathrm{Ob}(\mathbb{D}^+ (A_X))$  和  $VY$  的开放子集，我们有：

\[
\begin{aligned}R\Gamma(V;f^{!}F)&=R\mathrm{Hom}(A_{V},f^{!}F)\\&=R\mathrm{Hom}(Rf_{!}A_{V},F)\end{aligned}
\]

最后一个复数可以通过采用  $ A_x $  的 c-soft 解析  $ K $  来计算。然后  $ Rf_1A_V = f_1K_V $  ，如果  $ F $  是内射层的复形，我们得到：

\[
R\Gamma(V;f^{!}F)=\operatorname{H o m}(f_{!}K_{V},F).
\]

为了执行这个构造，我们需要对 f 进行一些假设。

\statementlabel{定义 3.1.1}$Y$  上的束  $G$  称为  $f$  -soft 如果对于任何  $x \in X$  ，束  $G|_{f^{-1}(x)}$  是  $c$  -soft。

鉴于练习 II.6，G 是 f-soft 当且仅当对于 Y 的任何开子集 V 和任何  $ j \neq 0 $  、  $ R^j f_t G_V = 0 $  。

为了构造函子 $f^!$ ，我们假设：

(3.1.3) 函子 $f_!: \mathfrak{Mod}(\mathbb{Z}_Y) \to \mathfrak{Mod}(\mathbb{Z}_X)$  具有有限上同调维数。

这意味着存在一个整数  $ r \geq 0 $  使得  $ R^{j}f_i = 0 $  对应  $ j > r $  。这也相当于以下条件之一（参见练习 I.19）：

\[
\begin{cases}{{f o r~}a n y\;G\in\mathsf{O b}(\mathfrak{S h}(Y)),\;{t h e r e~e x i s t s~a n~e x a c t}}\\ {{s e q u e n c e}\;0\to G\to G^{0}\to\cdots\to G^{r}\to0,\;{w h e r e~t h e~G^{j,s}~}}\\ {{a r e~f-s o f t~}\;,}\\ \end{cases}
\]

\[
\begin{cases}{{f o r~a n y~e x a c t~s e q u e n c e~i n~\mathfrak{S h}(Y):}}\\ {G^{0}\to\cdots\to G^{r}\to0,{w i t h~G^{j}~f-s o f t~f o r~}j<r,}\\ {{t h e~s h e a f~G^{r}~i s~f-s o f t~}.}\\ \end{cases}
\]

请注意， $f_!$  具有上同调维度  $\leq r$  当且仅当对于任何  $x\in X$  ，  $\Gamma_c(f^{-1}(x);\cdot)$  具有上同调维度  $\leq r$  。

令 $K$  为 $\mathbb{Z}_Y$  模， $F$  为 $A_X$  模。让我们通过以下方式定义  $A$  模的 presheaf  $f_K^!F$ ：

\[
(f_{\mathbf{K}}^{!}F)(V)=\operatorname{H o m}_{A_{X}}\biggl(f_{!}\biggl(A_{Y}\otimes_{\mathbf{Z}_{Y}}K_{V}\biggr),F\biggr).
\]

（回想一下，对于两个  $A_{X}$  -模  $F$  和  $G$  ，我们写  $\mathrm{Hom}_{Ax}(G,F)$  而不是  $\mathrm{Hom}_{\mathfrak{Rod}(Ax)}(G,F).$  ）

对于  $V' \subset V$  ，  $f_!(A_Y \otimes_{\mathbb{Z}_Y} K_{V'}) \to f_!(A_Y \otimes_{\mathbb{Z}_Y} K_V)$  定义限制图  $(f_K^! F)(V) \to (f_K^! F)(V')$  。

\statementlabel{引理 3.1.2}令 K 为平坦且 f-soft  $ \mathbb{Z}_{Y} $  模。

(i) 对于 Y 上的任何层 G， $ G \otimes_{Z_Y} K $  是 f-soft。

(ii)  $ G \mapsto f_!(G \otimes_{\mathbb{Z}_Y} K) $  是从  $ \mathfrak{Mod}(\mathbb{Z}_Y) $  到  $ \mathfrak{Mod}(\mathbb{Z}_X) $  的精确函子。

\statementlabel{证明}(i) Y 上的任何层 G 都有一个分辨率：

\[
{\rightarrow}G^{-r}{\rightarrow}\cdots{\rightarrow}G^{0}{\rightarrow}G{\rightarrow}0,
\]

其中每个  $G^j$  都是在  $Y$  中打开的层  $\mathbb{Z}_V$  、  $V$  的直积，通过命题 2.4.12 证明中的构造。因此 $G^j\otimes_{\mathbb{Z}_V}K$  是一个 $f$  -软捆。自从

\[
\to G^{-r}\otimes\underset{\mathbb{Z}_{Y}}{\bigotimes}K\to\cdots\to G^{0}\otimes\underset{\mathbb{Z}_{Y}}{\bigotimes}K\to G\otimes\underset{\mathbb{Z}_{Y}}{\bigotimes}K\to0
\]

是精确的，(3.1.4)' 给出了期望的结果，使 r 足够大。

(ii) 紧接着(i)。 ⑨

\statementlabel{引理 3.1.3}令 K 为平坦且 f 软的  $ \mathbb{Z}_{Y} $  模，并令 F 为单射  $ A_{X} $  模。

(i) presheaf  $f_{K}^{!}$  是一个层，这个层是一个单射  $A_{Y}$  模。

(ii) 我们有关于  $ G \in \text{Ob}(\mathfrak{M}_\partial(A_Y)) $  的规范同构、函子：

\[
\operatorname{H o m}_{A_{X}}\biggl(f_{!}\biggl(G\mathop{\otimes}_{\mathbb{Z}_{\mathbb{Y}}}K\biggr),F\biggr)\xrightarrow{\sim}\operatorname{H o m}_{A_{Y}}(G,f_{K}^{!}F).
\]

\statementlabel{证明}在证明中，我们将写 $ \cdot \otimes K $ 而不是 $ \cdot \otimes_{\mathbb{Z}_{Y}} K $ 。

我们首先证明 $ f_K^!F $ 是一个层。令  $ V = \bigcup V_j $  为  $ Y $  的开子集  $ V $  的开覆盖。那么有一个确切的顺序：

\[
\bigoplus_{j,k}A_{V_{j}\cap V_{k}}\to\bigoplus_{j}A_{V_{j}}\to A_{V}\to0\enspace.
\]

应用引理 3.1.2 (ii)，我们得到精确的序列：

\[
f_{!}\bigg(\bigoplus_{j,k}A_{V_{j}\cap V_{k}}\otimes K\bigg)\to f_{!}\bigg(\bigoplus_{j}A_{V_{j}}\otimes K\bigg)\to f_{!}(A_{V}\otimes K)\to0~.
\]

由于 F 是单射的，因此序列：

\[
\begin{aligned}0&\to\mathrm{Hom}_{A_{X}}(f_{!}(A_{V}\otimes K),F)\to\mathrm{Hom}_{A_{X}}\bigg(f_{!}\bigg(\bigoplus_{j} A_{V_{j}}\otimes K\bigg),F\bigg)\\&\to\mathrm{Hom}_{A_{X}}\bigg(f_{!}\bigg(\bigoplus_{j,k}A_{V_{j}\cap V_{k}}\otimes K\bigg),F\bigg)\end{aligned}
\]

是精确的，并且最后一个序列与以下序列同构：

\[
0\to(f_{K}^{!}F)(V)\to\prod_{j}(f_{K}^{!}F)(V_{j})\to\prod_{j,k}(f_{K}^{!}F)(V_{j}\cap V_{k})\enspace.
\]

这表明 $ f_{k}^{\prime}F $  是一个层。

现在让我们定义一个同态：

\[
\alpha(G)\colon\mathrm{H o m}_{A_{X}}(f_{!}(G\otimes K),F)\to\mathrm{H o m}_{A_{Y}}(G,f_{K}^{!}F)\enspace.
\]

让 $ \phi \in \mathrm{Hom}_{A_x}(f_!(G \otimes K), F) $  。那么对于  $ Y $  的任何开集  $ V $  ，我们有  $ A_x $  -模的态射链：

\[
\begin{aligned}{G(V)\mathop{\otimes}_{\mathcal{A}}f_{!}(A_{Y}\otimes K_{V})}&{{}\longrightarrow f_{!}(G\otimes K_{V})}\\ {}&{{}\longrightarrow f_{!}(G\otimes K)}\\ {}&{{}\stackrel{\phi}{\longrightarrow}F.}\\ \end{aligned}
\]

这给出了从  $G(V)$  到  $(f_{\mathbb{K}}^1F)(V) = \mathrm{Hom}(f_!(A_Y \otimes K_V), F)$  的态射，并且该态射相对于  $V$  在  $Y$  中开放，我们获得了一个元素  $\alpha(G)(\phi) \in \mathrm{Hom}_{A_Y}(G, f_{\mathbb{K}}^1F)$  。

我们将通过(a)、(b)、(c)三个步骤证明 $ \alpha(G) $ 是同构。

(a) 当 $ G = A_{V} $  ，V 在Y 中开路时， $ \alpha(G) $  是同构。事实上，

\[
\begin{aligned}{\operatorname{H o m}_{A_{X}}(f_{!}(G\otimes K),F)}&{{}=\operatorname{H o m}_{A_{X}}(f_{!}(A_{V}\otimes K),F)}\\ {}&{{}\simeq(f_{K}^{!}F)(V)}\\ {}&{{}\simeq\operatorname{H o m}_{A_{Y}}(G,f_{K}^{!}F).}\\ \end{aligned}
\]

(b) 当 $ G = \oplus_j A_{V_j} $  时，对于 $ Y, \alpha(G) $  的开子集族 $ V_j $  是同构。事实上，这是从  $ \alpha(G) = \prod_j \alpha(A_{V_j}) $  以来的 (a) 得出的。

(c) 现在让  $G$  为任意  $A_Y$  模。然后根据命题 2.4.12，有一个精确的序列  $0 \to G'' \to G' \to G \to 0$  ，其中  $G' \simeq \oplus_j A_{V_j}$  。因此 $\alpha(G')$  是同构。考虑交换图：

\[
\begin{array}{ccc}0\longrightarrow\mathrm{Hom}(f_{!}(G\otimes K),F)\longrightarrow\mathrm{Hom}(f_{!}(G^{\prime}\otimes K),F)\longrightarrow\mathrm{Hom}(f_{!}(G^{\prime \prime}\otimes K),F)\\\downarrow\Bigg|_{\substack{\alpha(G)}}&\Bigg|_{\substack{\alpha(G^{\prime})}}&\Bigg|_{\substack{\alpha(G^{\prime \prime})}}\\0\longrightarrow\quad\mathrm{Hom}(G,f_{K}^{!}F)&\longrightarrow&\mathrm{Hom}(G^{\prime},f_{K}^{!}F)\quad\longrightarrow\quad\mathrm{Hom}(G^{\prime \prime},f_{K}^{!}F)\\\end{array}
\]

两行均符合引理 3.1.2。由于 $\alpha(G^{\prime})$  是双射的， $\alpha(G)$  是单射的。将此结果应用于层 $G^{\prime}$ ，我们发现 $\alpha(G^{\prime})$  是单射的。因此 $\alpha(G)$  是双射的。

由于这个事实，我们得到  $ \mathrm{Hom}_{A_Y}(\cdot, f_K^!F) $  是  $ \mathrm{Mod}(A_Y) $  上的精确函子。因此 $ f_K^!F $  是单射的。   $ \square $ 

让  $f_{1}$  具有上同调维度  $\leq r$  。

\statementlabel{引理 3.1.4}层 $\mathbb{Z}_Y$  承认一个分辨率 $0 \to \mathbb{Z}_Y \to K^0 \to \cdots \to K^r \to 0$ ，使得所有 $K^j$  都是平坦的并且 $f$  -软 $\mathbb{Z}_Y$  -模。

\statementlabel{证明}我们将通过与命题 2.4.3 的证明相同的方法构造 $\mathbb{Z}_Y$  的分辨率。令 $p: \hat{Y} \to Y$  为从 $\hat{Y}$ （赋予离散拓扑的集合 $Y$ ）到 $Y$  的自然映射。然后我们定义： $K^0 = p_* p^{-1} \mathbb{Z}_Y$ ， $K^1 = p_* p^{-1}(K^0/\mathbb{Z}_Y)$ ，...， $K^j = p_* p^{-1}(\mathrm{Coker}(K^{j-2} \to K^{j-1}))$ ，对于 $1 < j < r$ 和 $K^r = \mathrm{Coker}(K^{r-2} \to K^{r-1})$ 。

然后我们有精确的序列  $ 0 \to \mathbb{Z}_Y \to K^0 \to \cdots \to K^r \to 0 $  ，  $ K^j $  对于  $ 0 \leq j < r $  来说是松弛的，因此  $ f $  -soft，并且从 (3.1.4)' 可以得出  $ K^r $  也是  $ f $  -soft。

因此，仍然需要证明  $K^j$  是  $\mathbb{Z}_Y$  的平坦  $0 \leqslant j \leqslant r$  模。为了证明这一点，我们将证明如果 $G$  是平坦的，那么 $p_* p^{-1}G$  和 $p_* p^{-1}G/G$  也是平坦的。因为对于 $y \in Y$ ，

\[
(p_{*}p^{-1}G)_{y}=\varinjlim_{y\in U}\prod_{y^{\prime}\in U}G_{y^{\prime}},
\]

and

\[
(p_{*}p^{-1}G/G)_{y}=\varinjlim_{y\in U}\prod_{y^{\prime}\in U\setminus\{y\}}G_{y^{\prime}}
\]

（U 的范围通过 y 的邻域系统），这些 Z 模是无扭转的，因此是平坦的。 □

让 $ \mathcal{I}(X) $  表示由单射对象组成的 $ \mathfrak{Mod}(A_X) $  的满子范畴。那么 $ \mathbf{K}^+(\mathcal{I}(X)) \to \mathbf{D}^+(A_X) $  是三角范畴的等价物。令  $ K $  为复形，如引理 3.1.4 中所示。对于  $ F \in \mathbf{K}^+(\mathcal{I}(X)) $  ，令  $ f_K^! F $  为与双复数  $ (f_{K-q}^! (F^p))^{p,q} $  关联的简单复数。那么 $ f_K^! F $  是单射 $ A_Y $  模的复合体。很容易检查 $ f_K^! $ 发送态射同伦

到零到一个态射同伦到零。传递到商，我们获得三角范畴的函子：

\[
{f}_{\mathbf{K}}^{!}\colon\mathbf{K}^{+}(\mathcal{I}(X))\to\mathbf{K}^{+}(\mathcal{I}(Y))\enspace.
\]

\statementlabel{定理 3.1.5}令 $f: Y \to X$  为局部紧空间的连续映射，使得 $f_!$  具有有限上同调维数。那么存在一个三角范畴函子  $f^!: \mathbf{D}^+ (A_X) \to \mathbf{D}^+ (A_Y)$  和  $\mathbf{D}^+ (A_Y)^\circ \times \mathbf{D}^+ (A_X)$  上的双函子同构：

\[
\mathrm{H o m}_{\mathsf{D}^{\cdot}(A_{X})}(R f_{!}(\cdot),\cdot)\simeq\mathrm{H o m}_{\mathsf{D}^{+}(A_{Y})}(\cdot,f^{!}(\cdot))\enspace.
\]

换句话说，  $f'$  是  $Rf_{!}$  的右伴随。

（请注意， $f'$  是唯一的，直至同构；参见练习 I.2。）

\statementlabel{证明}我们将证明  $f_{k}^{\prime}$  具有所需的属性。

对于  $ F \in \text{Ob}(\mathbb{K}^{+}(\mathcal{I}(X))) $  和  $ G \in \text{Ob}(\mathbb{K}^{+}(\mathcal{I}(Y))) $  ，引理 3.1.3 意味着：

\[
\operatorname{H o m}_{\mathsf{K}^{+}(\mathfrak{R o b}(A_{X}))}\biggl(f_{!}\biggl(G\mathop{\otimes}_{\mathbf{Z}_{Y}}K\biggr),F\biggr)\simeq\operatorname{H o m}_{\mathsf{K}^{+}(\mathfrak{R o b}(A_{Y}))}(G,f_{K}^{!}F)\enspace.
\]

由于  $G \simeq G \otimes_{\mathbb{Z}_Y} \mathbb{Z}_Y \to G \otimes_{\mathbb{Z}_Y} K$  是准同构，并且  $G \otimes_{\mathbb{Z}_Y} K$  是引理 3.1.2 中的  $f$  -软层的复形， $\mathbf{D}^+ (A_X)$  中的  $R f_! G \simeq f_!(G \otimes_{\mathbb{Z}_Y} K)$  。因此 $\mathrm{Hom}_{K^+(\mathfrak{M}_\mathrm{rob}(A_X))}(f_!(G \otimes_{\mathbb{Z}_Y} K), F) \simeq \mathrm{Hom}_{\mathbf{D}^+ (A_X)}(R f_! G, F)$  。另一方面，由于  $f_K^! F$  是单射  $A_Y$  模的复合体，

\[
\mathrm{H o m}_{\mathsf{K}^{+}(\mathfrak{M o d}(A_{Y}))}(G,f_{K}^{!}F)\simeq\mathrm{H o m}_{\mathsf{D}^{+}(A_{Y})}(G,f_{K}^{!}F)\enspace.
\]

这样就完成了证明。 ⑨

\statementlabel{备注 3.1.6}(i) 定理 3.1.5 可以概括如下。令 $f: Y \to X$ 如上，令 $\mathcal{R}$ （或 $\mathcal{S}$ ）为 $X$ （或 $Y$ ）上的环层，并假设被赋予环层 $f^{-1}\mathcal{R} \to \mathcal{S}$ 的态射。然后存在三角范畴  $f^!: \mathbf{D}^+(\mathcal{R}) \to \mathbf{D}^+(\mathcal{S})$  的函子，它与  $R f_!: \mathbf{D}^+(\mathcal{S}) \to \mathbf{D}^+(\mathcal{R})$  右邻。通过将  $(f_K^! F)(V)$  定义为  $\operatorname{Hom}_{\mathfrak{Mod}(\mathcal{R})}(f_!(\mathcal{S} \otimes_{Z_V} K_V), F)$  ，证明类似。

(ii) 定理 3.1.5 中的同构与移位相容。也就是说，对于  $ F \in \mathrm{Ob}(\mathbf{D}^{+}(A_X)) $  和  $ G \in \mathrm{Ob}(\mathbf{D}^{+}(A_Y)) $  ，下图是可交换的：

\[
\begin{array}{cccc}&\mathrm{Hom}_{\mathbb{D}^{+}(A\gamma)}(Rf_{!}^{\prime}G,F)&\xrightarrow{\quad}&\mathrm{Hom}_{\mathbb{D}^{+}(A\gamma)}(G,f^{!}F)\\&\bigg\downarrow\mathcal{E}&\bigg\downarrow\mathcal{E}&\\ \mathrm{Hom}_{\mathbb{D}^{+}(A\gamma)}((Rf_{!}^{\prime}G)[1],F[1])&\mathrm{Hom}_{\mathbb{D}^{+}(A\gamma)}(G[1],(f^{!}F)[1])&\\ &\bigg\downarrow\bigg|&\bigg\vert\\\mathrm{Hom}_{\mathbb{D}^{+}(A\gamma)}(Rf_{!}(G[1]),F[1])&\xrightarrow{\quad}&\mathrm{Hom}_{\mathbb{D}^{+}(A\gamma)}(G[1],f^{!}(F[1]))\end{array}
\]

为了强调基环  $A$  ，我们写一段时间  $f_A^!$  而不是  $f^!.$  因此  $f_A^!$  是从  $\mathbf{D}^+ (A_X)$  到  $\mathbf{D}^+ (A_Y)$  的函子。让 $\phi_X : \mathbf{D}^+ (A_X) \to \mathbf{D}^+ (Z_X)$  和 $\phi_Y : \mathbf{D}^+ (A_Y) \to \mathbf{D}^+ (Z_Y)$  成为健忘函子。

\statementlabel{命题 3.1.7}其中一个有：

\[
\phi_{Y}\circ f_{A}^{!}\simeq f_{\bar{Z}}^{!}\circ\phi_{X}\enspace.
\]

\statementlabel{证明}让 $ G \in \mathrm{Ob}(\mathbb{D}^{+}(\mathbb{Z}_{Y})) $ ， $ F \in \mathrm{Ob}(\mathbb{D}^{+}(A_{X})) $ 。我们有：

\[
\begin{aligned}{\operatorname{H o m}_{\mathsf{D}^{+}(\mathbb{Z}_{Y})}(G,f_{\mathbb{Z}}^{!}\circ\phi_{X}F)}&{{}\simeq\operatorname{H o m}_{\mathsf{D}^{+}(\mathbb{Z}_{X})}(R f_{!}G,\phi_{X}F)}\\ {}&{{}\simeq\operatorname{H o m}_{\mathsf{D}^{+}(A_{X})}\biggl(A_{X}\mathop{{\otimes}}_{\mathbb{Z}_{X}}^{L}R f_{!}G,F\biggr)}\\ {}&{{}\simeq\operatorname{H o m}_{\mathsf{D}^{+}(A_{X})}\biggl(R f_{!}\biggl(A_{Y}\mathop{{\otimes}}_{\mathbb{Z}_{Y}}^{L}G\biggr),F\biggr)}\\ {}&{{}\simeq\operatorname{H o m}_{\mathsf{D}^{+}(A_{Y})}\biggl(A_{Y}\mathop{{\otimes}}_{\mathbb{Z}_{Y}}^{L}G,f_{A}^{!}F\biggr)}\\ {}&{{}\simeq\operatorname{H o m}_{\mathsf{D}^{+}(\mathbb{Z}_{Y})}(G,\phi_{Y}\circ f_{A}^{!}F).}\\ \end{aligned}
\]

因此 $ f_{Z}^! \circ \phi_X F \simeq \phi_Y \circ f_A^! F $  。 □

现在让  $ g: Z \to Y $  成为局部紧空间的另一个连续映射。

\statementlabel{命题 3.1.8}假设 $f_{1}$  和 $g_{1}$  具有有限上同调维数。那么 $(f\circ g)_{1}$  具有有限上同调维数，并且有：

\[
(f\circ g)^{\prime}\simeq g^{!}\circ f^{!}
\]

作为从  $ \mathbf{D}^{+}(A_{X}) $  到  $ \mathbf{D}^{+}(A_{Z}) $  的函子。

\statementlabel{证明}它紧随公式： $ (f \circ g)_{!} \simeq f_{!} \circ g_{!} $  。 □

\statementlabel{命题 3.1.9}考虑连续映射和局部紧空间的笛卡尔平方：

\[
\begin{array}{c}Y^{\prime}\xrightarrow{\quad f^{\prime}\quad}X^{\prime}\\g^{\prime}\bigg\downarrow\quad\Box\quad\bigg\downarrow g\\Y\xrightarrow{\quad f\quad}X\quad.\end{array}
\]

假设  $f_{1}$  具有有限上同调维数。然后：

(i)  $f_{1}^{\prime}$  具有有限的上同调维数。

(ii) 函子具有从  $ \mathbf{D}^{+}(A_{X}) $  到  $ \mathbf{D}^{+}(A_{Y}) $  的规范同构：

\[
f^{!}\circ R g_{*}\simeq R g_{*}^{\prime}\circ f^{\prime!}.
\]

(iii) 我们有从  $ \mathbf{D}^{+}(A_X) $  到  $ \mathbf{D}^{+}(A_Y) $  的函子态射：

\[
g^{\prime^{-1}f!}\to f^{\prime!}g^{-1}\enspace.
\]

\statementlabel{证明}(i) 让 $ x' \in X' $  。那么  $ f^{-1}(x') $  与  $ f^{-1}(g(x)) $  同胚。因此  $ f_1' $  的上同调维数受  $ f_1 $  的上同调维数限制。

(ii) 对于 $ F \in \mathrm{Ob}(\mathbf{D}^{+}(A_{X^{-}})), G \in \mathrm{Ob}(\mathbf{D}^{+}(A_{Y})) $ ，我们在功能上有：

\[
\begin{aligned}{\operatorname{H o m}_{\mathsf{D}^{+}(A\gamma)}(G,f^{!}R g_{*}F)}&{{}\simeq\operatorname{H o m}_{\mathsf{D}^{+}(A x)}(R f_{!}G,R g_{*}F)}\\ {}&{{}\simeq\operatorname{H o m}_{\mathsf{D}^{+}(A x)}(g^{-1}R f_{!}G,F)}\\ {}&{{}\simeq\operatorname{H o m}_{\mathsf{D}^{+}(A x)}(R f_{!}^{\prime}\circ g^{\prime-1}G,F)}\\ {}&{{}\simeq\operatorname{H o m}_{\mathsf{D}^{+}(A\gamma)}(g^{\prime-1}G,f^{\prime!}F)}\\ {}&{{}\simeq\operatorname{H o m}_{\mathsf{D}^{+}(A\gamma)}(G,R g_{*}^{\prime}\circ f^{\prime!}F).}\\ \end{aligned}
\]

这意味着期望的结果。

(iii) 根据命题 2.5.11，对于  $F \in \mathrm{Ob}(\mathbb{D}^+) (A_X)$  ，我们有  $R f_! g'^{-1} f^! F \simeq g^{-1} R f_! f^! F \to g^{-1} F$  。这给出了所需的同态  $g'^{-1} f^! F \to f'^! g^{-1} F$  。

\statementlabel{命题 3.1.10}假设  $f_!$  具有有限上同调维数。然后对于 $F \in \mathrm{Ob}(\mathbf{D}^+(A_X))$ ， $G \in \mathrm{Ob}(\mathbf{D}^b(A_Y))$ ，我们有：

\[
R\mathrm{H o m}(R f_{!}G,F)\simeq R\mathrm{H o m}(G,f^{!}F)\enspace,
\]

\[
R\mathcal{H}o m(R f_{!}G,F)\simeq R f_{*}R\mathcal{H}o m(G,f^{!}F)\enspace.
\]

\statementlabel{证明}第一个同构源自第二个同构，通过应用函子  $ R\Gamma(X;\cdot) $  。

我们有规范态射：

\[
R f_{*}R\mathcal{H}o m(G,f^{!}F)\to R\mathcal{H}o m(R f_{!}G,R f_{!}f^{!}F)\enspace.
\]

通过与态射  $ Rf_1 \circ f^!F \to F $  组合，我们得到态射：

\[
R f_{*}R\mathcal{H}o m(G,f^{!}F)\to R\mathcal{H}o m(R f_{!}G,F)\enspace.
\]

设 V 为 X 的开子集。则：

\[
\begin{aligned}{H^{j}(R\varGamma(V;R f_{*}{R\mathcal{H}o m}(G,f^{!}F)))}&{{}\simeq\operatorname{H o m}_{\mathsf{D}^{+}(A_{f^{-1}(V)})}(G|_{f^{-1}(V)},f^{!}F[j]|_{f^{-1}(V)})}\\ {}&{{}\simeq\operatorname{H o m}_{\mathsf{D}^{+}(A_{V})}((R f_{!}G)|_{V},F[j]|_{V})}\\ {}&{{}\simeq H^{j}(R\varGamma(V;R\mathcal{H}o m(R f_{!}G,F))).}\\ \end{aligned}
\]

这样就完成了证明。 ⑨

\statementlabel{命题 3.1.11}假设  $f_1$  具有有限上同调维数。那么存在从  $\mathbf{D}^+ (A_X) \times \mathbf{D}^+ (A_X)$  到  $\mathbf{D}^+ (A_Y)$  函子的自然态射：

\[
f^{!}(\cdot)\overset{L}{\otimes}_{\mathcal{A}_{Y}}f^{-1}(\cdot)\to f^{!}\left(\begin{array}{c}{L}\\ {\cdot\bigotimes_{\mathcal{A}_{X}}\cdot}\\ \end{array}\right).
\]

\statementlabel{证明}让  $F$  、  $F_1$  和  $F_2$  属于  $\mathbf{D}^+ (A_X)$  。我们有：

\[
\begin{aligned}{\operatorname{H o m}_{\mathbb{D}^{+}(A_{Y})}\biggl(f^{!}F_{1}\mathop{\otimes}_{A_{Y}}^{L}f^{-1}F_{2},f^{!}F\biggr)}&{{}\simeq\operatorname{H o m}_{\mathbb{D}^{+}(A_{X})}\biggl(R f_{!}\biggl(f^{!}F_{1}\mathop{\otimes}_{A_{Y}}^{L}f^{-1}F_{2}\biggr),F\biggr)}\\ {}&{{}\simeq\operatorname{H o m}_{\mathbb{D}^{+}(A_{X})}\biggl(R f_{!}f^{!}F_{1}\mathop{\otimes}_{A_{X}}^{L}F_{2},F\biggr)~.}\\ \end{aligned}
\]

采取 $ F = F_1 \otimes_{A_X}^L F_2 $ 。然后我们得到所需的态射作为  $ R f_! f^! F_1 \otimes_{A_X}^L F_2 \to F_1 \otimes_{A_X}^L F_2 $  的图像。   $ \Box $ 

我们将在第 3 节中看到，如果 f 是拓扑淹没，那么命题 3.1.11 的态射就变成同构。

有一种情况，函子  $f^!$  是很熟悉的。

\statementlabel{命题 3.1.12}假设  $f: Y \to X$  是从 Y 到  $X$  的局部闭子集的同胚。然后：

\[
f^{!}(\cdot)\simeq f^{-1}\circ R\varGamma_{f(Y)}(\cdot).
\]

\statementlabel{证明}设置 $ Z = f(Y) $ 。让 $ F \in \mathrm{Ob}(\mathbb{D}^{+}(A_X)) $ ， $ G \in \mathrm{Ob}(\mathbb{D}^{+}(A_Y)) $ 。然后：

\[
\begin{aligned}{\operatorname{H o m}_{\mathsf{D}^{+}(A x)}(f_{!}G,F)}&{{}\simeq\operatorname{H o m}_{\mathsf{D}^{+}(A x)}(f_{!}G,R\varGamma_{Z}(F))}\\ {}&{{}\simeq\operatorname{H o m}_{\mathsf{D}^{+}(A_{Y})}(f^{-1}\circ f_{!}G,f^{-1}R\varGamma_{Z}(F))}\\ {}&{{}\simeq\operatorname{H o m}_{\mathsf{D}^{+}(A_{Y})}(G,f^{-1}R\varGamma_{Z}(F))~.}\\ \end{aligned}
\]

因此 $ f^!F \simeq f^{-1} \circ R\Gamma_Z(F) $  。 □

\statementlabel{命题 3.1.13}假设  $f_1$  具有有限上同调维数。然后对于 $F_1 \in \mathrm{Ob}(\mathbb{D}^b(A_X))$ ， $F_2 \in \mathrm{Ob}(\mathbb{D}^+ (A_X))$ ，我们有：

\[
f^{!}R{\mathcal{H}}o m(F_{1},F_{2})\simeq R{\mathcal{H}}o m(f^{-1}F_{1},f^{!}F_{2})~.
\]

\statementlabel{证明}让 $ G \in \mathrm{Ob}(\mathbb{D}^{+}(Y)) $  。然后：

\[
\begin{aligned}{\operatorname{H o m}_{\mathsf{D}^{+}(A\gamma)}(G,f^{!}R\mathcal{H}o m(F_{1},F_{2}))}&{{}\simeq\operatorname{H o m}_{\mathsf{D}^{+}(A\chi)}(R f_{!}G,R\mathcal{H}o m(F_{1},F_{2}))}\\ {}&{{}\simeq\operatorname{H o m}_{\mathsf{D}^{+}(A\chi)}\bigg(R f_{!}G\otimes\stackrel{L}{F_{1},F_{2}}\bigg)}\\ {}&{{}\simeq\operatorname{H o m}_{\mathsf{D}^{+}(A\chi)}\bigg(R f_{!}\bigg(G\otimes\stackrel{L}{f^{-1}F_{1}}\bigg),F_{2}\bigg)}\\ {}&{{}\simeq\operatorname{H o m}_{\mathsf{D}^{+}(A\chi)}\bigg(G\otimes\stackrel{L}{f^{-1}F_{1}},f^{!}F_{2}\bigg)}\\ {}&{{}\simeq\operatorname{H o m}_{\mathsf{D}^{+}(A\gamma)}(G,R\mathcal{H}o m(f^{-1}F_{1},f^{!}F_{2}))}\\ \end{aligned}
\]

由于这对于任何 G 都是成立的，因此证明如下。   $ \Box $ 

\statementlabel{命题 3.1.14}假设  $X$  具有有限的  $c$  -软维数（参见练习 11.9）。让 $F \in \mathrm{Ob}(\mathbf{D}^{+}(A_X))$ ， $G \in \mathrm{Ob}(\mathbf{D}^b(A_X))$ 。让  $q_1$  和  $q_2$  表示  $X \times X$  上的第一和第二个投影， $\Delta$  表示对角线。然后：

\[
R{\mathcal{H}}o m(G,F)\simeq R q_{1*}R\Gamma_{\varsigma}R{\mathcal{H}}o m(q_{2}^{-1}G,q_{1}^{!}F)\enspace.
\]

\statementlabel{证明}用 $\delta$  表示对角线嵌入。然后：

\[
\begin{aligned}{{R q}_{1*}{R}{\varGamma}_{\varDelta}{R}\mathcal{H}{o m}(q_{2}^{-1}G,q_{1}^{!}F)}&{{}\simeq\delta^{!}{R}\mathcal{H}{o m}(q_{2}^{-1}G,q_{1}^{!}F)}\\ {}&{{}\simeq{R}\mathcal{H}{o m}(\delta^{-1}q_{2}^{-1}G,\delta^{!}q_{1}^{!}F)}\\ {}&{{}\simeq{R}\mathcal{H}{o m}(G,F)~.\quad\Box}\\ \end{aligned}
\]

最后考虑两个局部紧空间X和Y，并分别用 $ q_{1} $ 和 $ q_{2} $ 表示从 $ X \times Y $ 到X和Y的投影。

\statementlabel{命题 3.1.15}假设 Y 具有有限的 c-软维数。让 $ F \in \mathrm{Ob}(\mathbb{D}^{+}(A_X)), G \in \mathrm{Ob}(\mathbb{D}^b(A_Y)) $  。然后有一个规范的同构：

\[
R\varGamma(X\times Y;R\mathcal{H}o m(q_{2}^{-1}G,q_{1}^{!}F))\simeq R\mathrm{H o m}(R\varGamma_{c}(Y;G),R\varGamma(X;F))\enspace.
\]

\statementlabel{证明}用 $ a_{X} $ （分别 $ a_{Y} $ ）表示投影 $ X \to \{\text{pt}\} $ （分别 $ Y \to \{\text{pt}\} $ ）。然后：

\[
\begin{aligned}{R\varGamma(X\times Y;R\mathcal{H}o m(q_{2}^{-1}G,q_{1}^{\prime}F))}&{{}\simeq R\mathsf{a}_{X\ast}R q_{1\ast}R\mathcal{H}o m(q_{2}^{-1}G,q_{1}^{\prime}F))}\\ {}&{{}\simeq R\mathsf{a}_{X\ast}R\mathcal{H}o m(R q_{1\ast}q_{2}^{-1}G,F)}\\ {}&{{}\simeq R\mathsf{a}_{X\ast}R\mathcal{H}o m(\mathsf{a}_{X}^{-1}R\mathsf{a}_{Y\ast}G,F)}\\ {}&{{}\simeq R\mathsf{H}o m(R\mathsf{a}_{Y\ast}G,R\mathsf{a}_{X\ast}F)~.\quad\Box}\\ \end{aligned}
\]

\statementlabel{定义 3.1.16}(i) 令 $f: Y \to X$  为连续映射，并假设 $f_1$  具有有限上同调维数。一套：

\[
\omega_{Y/X}=f^!A_X
\]

并将  $ \omega_{Y/X} $  称为相对对偶复形。如果 $ X = \{\mathrm{pt}\} $ ，则设置 $ \omega_{Y} = \omega_{Y/X} $ 并调用Y上的对偶复数 $ \omega_{Y} $ 。

(ii) 假设  $X$  具有有限的  $c$  -软维数，并设  $F \in \mathrm{Ob}(\mathbb{D}^b(X))$  。一套：

\[
\mathrm{D}_{x}F=R\mathcal{H}o m(F,\omega_{x})~,\qquad\mathrm{D}_{x}^{\prime}F=R\mathcal{H}o m(F,A_{x})~.
\]

有人称  $ D_{x}F $  为 F 的对偶。

为了避免混淆，我们写 DF 而不是  $ D_{x}F $  。

根据命题 3.1.11，函子存在自然态射：

\[
{f}^{-1}(\cdot)\otimes\omega_{Y/X}\to{f}^{1}(\cdot)\enspace.
\]

根据定理 3.1.5 和命题 3.1.10，对于  $ F \in \mathrm{Ob}(\mathbb{D}^{b}(A_{x})) $  ，我们有：

\[
\operatorname{H o m}_{\mathsf{D}^{\flat}(A x)}(F,\omega_{X})\simeq\operatorname{H o m}_{\mathsf{D}^{\flat}(\mathfrak{R o b}(A))}(R\varGamma_{c}(X;F),A)\qquad\mathrm{a n d}
\]

\[
R\mathrm{H o m}(F,\omega_{X})\simeq R\mathrm{H o m}(R\varGamma_{c}(X;F),A)\enspace.
\]

特别是，采用  $ F = A_z $  ，我们有  $ R\Gamma_z(X; \omega_X) \simeq \mathrm{RHom}(R\Gamma_c(Z; A_z), A) $  作为局部封闭子集  $ Z $  。

\subsection{流形上的消没定理}
令  $V$  为维度为  $n$  的实向量空间。

\statementlabel{引理 3.2.1}对于 V 上的任何层 F，我们有：

\[
H_{c}^{j}(V;F)=0\quad{f o r\quad~j>n~}.
\]

\statementlabel{证明}首先假设 $n=1$ ，并设置 $\widetilde{F}=i_{1}F$ ，其中 $i:V\to[0,1]$ 是从 $V$ 到区间 $]0,1[$ 的同胚。然后我们有：

\[
H_{c}^{j}(V;F)\xrightarrow{\sim}H^{j}([0,1];\widetilde{F})\enspace.
\]

因此，结果来自引理 2.7.3。在一般情况下，令  $V'$  为  $(n-1)$  维向量空间，并令  $f$  为满射线性映射  $V \to V'$  。根据前面的结果， $R^j f_i F = 0$  对于 $j \neq 0$  ， 1。由于 $R^0 f_i F \simeq \tau^{\leq 0} R f_i F$  和 $\tau^{\geq 1} R f_i F \simeq (R^1 f_i F) [-1]$  ，我们得到一个显着的三角形：

\[
R^{0}f_{!}F\longrightarrow R f_{!}F\longrightarrow(R^{1}f_{!}F)[-1]\xrightarrow[+1]{}.
\]

让我们将函子  $\Gamma_c(V'; \cdot)$  应用于这个三角形。我们得到长精确序列：

\[
\cdots\to H_{c}^{j}(V^{\prime};R^{0}f_{!}F)\to H_{c}^{j}(V;F)\to H_{c}^{j-1}(V^{\prime};R^{1}f_{!}F)\to\cdots
\]

结果是对 n 进行归纳。 ⑨

\statementlabel{命题 3.2.2}令 X 为 n 维  $ C^{0} $  流形，并令 F 为 X 上的一个层。然后：

(i) F 承认 c-软层的长度最多为 n，

(ii) F 承认松弛层的长度至多为 $ n + 1 $ ，

(iii)  $ H_c^j(X; F) = 0 $  为  $ j > n $  ，

(iv)  $ H^j(X; F) = 0 $  为  $ j > n $  。

(v) 令 Z 为 X 的局部闭子集。然后  $ H_{Z}^{j}(X; F) = 0 $  对于 j > n + 1。

回想一下，根据其定义， $ C^{0} $  流形在无穷大处是可数的。

\statementlabel{证明}令  $ 0 \to F \to F_0 \to F_1 \to \cdots $  为  $ F $  的松弛分辨率（参见 II §4）。考虑确切的顺序：

\[
0\to F\to F_{0}\to F_{1}\to\cdots\to F_{n-1}\to G_{n}\to0,
\]

其中 $G_{n}=\mathrm{Im}(F_{n-1}\to F_{n})$  。我们将证明 $G_{n}$  是 $c$  -soft。该属性在  $X$  上是局部的（参见练习 II.6），我们可以假设  $X$  在维度  $n$  的某个实向量空间  $V$  中是开的。令  $U$  为  $X$  的开子集。我们有：

\[
\boldsymbol{H}_{c}^{j}(\boldsymbol{U};\boldsymbol{F})=\boldsymbol{H}_{c}^{j}(\boldsymbol{V};\boldsymbol{F}_{U})
\]

根据引理 3.2.1，该组对于  $j > n$  为零。因此 $j > 0$  和 $G_n$  的 $H_c^j(U; G_n) = 0$  是 $c$  -soft，（练习II.6）。

现在 (iii) 由 (i) 和 (iv) 得出，因为如果 G 是 c-soft，则  $ H^j(X; G) = 0 $  对于 j > 0（命题 2.5.10）。

通过与 $Z \hookrightarrow X$ 相关的长精确序列，我们从(iv)获得(v)，并且从(v)得出(ii)。 □

\statementlabel{命题 3.2.3}令  $V$  为维度为  $n$  的实向量空间。让 $M \in \mathrm{Ob}(D^+(\mathfrak{Mod}(A)))$  ，让 $x \in V$  。

(i) 自然态射 $ \mathrm{R}\Gamma(V; M_{V}) \to (M_{V})_{x} \simeq M $  是同构。

(ii) 自然态射 $ \mathrm{R}\Gamma_{\{x\}}(V; M_{V}) \to \mathrm{R}\Gamma_{c}(V; M_{V}) $  是同构。

(iii) 存在同构：

\[
R\varGamma_{c}(\mathbb{R}^{n};M_{\mathbb{R}^{n}})\simeq M[-n]~.
\]

(iv) 令 $\phi$  为 $V$  的线性自同构，并令 $\phi_{c}^{*}$  为 $R\Gamma_{c}(V;M_{V})$  的相关自同构。然后是  $\phi_{c}^{*}=\mathrm{sgn}(\phi)\cdot\mathrm{id}_{R\Gamma_{c}(V;M_{V})}$  ，其中  $\mathrm{sgn}(\phi)$  是  $\phi$  行列式的符号。

(v)  $V$  的每个方向都与与 (iii) 中的同构兼容的同构  $R\Gamma_{c}(V;M_{V})\simeq M[-n]$  规范相关。

\statementlabel{证明}(i) 应用推论 2.7.7。

(ii) 我们可以假设  $x = 0$  。设置 $F = M_V$ 。如 (i) 中所示，  $R\Gamma(\{y; |y| > 0\}; F) \to R\Gamma(\{y; |y| > a\}; F)$  是  $a > 0$  的同构。因此 $R\Gamma_{\{0\}}(V; F) \to R\Gamma_{\{y; |y| \leq a\}}(V; F)$  是同构。从 $H_c^j(V; F) = \lim_{a} H_{\{y; |y| \leq a\}}^j(V; F)$ 开始，我们得到了想要的结果。

(iii), (iv) 首先假设 n = 1。用  $ U_1 $  和  $ U_2 $  表示  $ V \setminus \{0\} $  的连通分量并考虑区分三角形：

\[
R\Gamma_{\{0\}}(V;M_{V})\xrightarrow{\alpha}R\Gamma(V;M_{V})\xrightarrow{\beta}R\Gamma(U_{1};M_{V})\oplus R\Gamma(U_{2};M_{V})\xrightarrow[+1]{\gamma}.
\]

应用 (i)， $ \beta $  承认左逆，因此  $ \alpha = 0 $  并且我们有一个外同态：

\[
R\varGamma(U_{1};M_{V})\oplus R\varGamma(U_{2};M_{V})\xrightarrow{\gamma}R\varGamma_{\{0\}}(V;M_{V})[1],
\]

其中  $\beta = (\beta_1, \beta_2)$  和  $\beta_i$  是同构  $\Gamma(V; M_V) \simeq \Gamma(U_i; M_V)$  。然后，在这种情况下，(iii) 以及(iv) 都遵循，因为鉴于(3.2.1)， $\phi_c^\#[1] \circ \gamma = \pm 1$  根据 $\phi$  是否互换 $U_1$  和 $U_2$  或不互换。为了处理一般情况，我们通过  $\dim V$  进行归纳论证，并选择线性分解  $V \simeq L \times V$  和  $\dim L = 1$  。根据 Künneth 公式（练习 II.18），我们得到：

\[
R\varGamma_{c}(V;M_{V})\simeq R\varGamma_{c}(L;M_{L})\stackrel{L}{\otimes}R\varGamma_{c}(V^{\prime};A_{V^{\prime}}).
\]

然后是(iii)。为了证明 (iv)，我们注意到  $\phi_c^\#$  相对于  $\phi \in \mathrm{GL}(V)$  是局部常数（命题 2.7.5），并且  $\mathrm{GL}(V)$  有两个连通分量。因此，用  $\phi_c^\# = -1$  找到  $\phi$  就足够了。我们采用  $\phi = \psi \otimes \mathrm{id}_{V^*}$  ，其中  $\psi$  是  $L$  上的对映图，并得出 (3.2.2) 的结论。

(v) 由(iii) 和(iv) 得出。 ⑨

在§5中，我们需要以下结果。

\statementlabel{命题 3.2.4}令 $V$  为实数有限维向量空间， $X$  为 $V$  的开子集， $F$  为 $X$  上的一束。假设对于  $X$  的任何凸紧子集  $K$  ，限制图  $\Gamma(X; F) \to \Gamma(K; F)$  是满射的。然后对于  $X$  的任何凸开子集  $U$  ，  $H^j(U; F) = 0$  对于所有  $j > 0$  。

\statementlabel{证明}我们可以假设  $X = U$  是凸的。通过考虑  $X$  的紧凸子集  $\{K_n\}$  的递增序列并应用命题 2.7.1，足以证明  $H^j(K; F) = 0$  对于  $j > 0$  ，对于  $X$  的任何紧凸子集  $K$  。

我们对 dim  $V$  进行归纳，并注意如果  $\dim V = 1$  ，结果是引理 2.7.3 的特殊情况。令 $V = V' \oplus L$  为 $V$  与 $\dim L = 1$  的线性分解，并令 $f$  为投影 $V \to V'$  。令 $f_{\kappa}$  为 $f$  到 $K$  的限制。如果  $y \in V$  ，则层  $F|_{f^{-1}(y) \cap K}$  满足命题的假设。因此  $R^j f_{\kappa*}(F|_{\kappa}) = 0$  为  $j > 0$  。此外，束 $f_{\kappa*}(F|_{\kappa})$ 也满足命题的假设。因此，根据归纳假设，我们得到：

\[
\begin{array}{r}{H^{j}(K;F)=H^{j}(f(K);f_{K*}F)=0}\end{array}
\]

对于 $j>0$ 。 ⑨

\subsection{定向与对偶}
令 $f: Y \to X$  为局部紧空间的连续映射。

\statementlabel{定义 3.3.1}我们可以说  $f$  是具有纤维维数  $l$  的拓扑淹没，如果对于任何  $y \in Y$  存在  $y$  的开邻域  $V$  使得  $U = f(V)$  在  $X$  中开，并且存在交换图：

\begin{center}
\includegraphics[width=0.19\textwidth]{流行上的层_Chn-assets/img_in_image_box_510_149_738_336.png}
\end{center}

其中 p 是投影，h 是同胚。

请注意，如果 X 和 Y 是  $C^{1}$  -流形且  $f$  是  $C^{1}$  -淹没，则  $f$  是拓扑淹没。如果 $f$  是拓扑淹没，则根据命题3.2.2， $f_{1}$  具有有限上同调维数。

\statementlabel{命题 3.3.2}假设 f 是纤维尺寸为 l 的拓扑浸没。然后：

(i)  $ k\neq -l $  和  $ H^{-1}(f^{!}A_{X}) $  的  $ H^{k}(f^{!}A_{X})=0 $  与  $ A_{Y} $  局部同构。

(ii) 函子 $ f^!A_X \otimes_{A_Y} f^{-1}(\cdot) \to f^!(\cdot) $  的态射是同构。

\statementlabel{证明}我们首先要证明 (i) 当  $Y = \mathbb{R}^l$  和  $X = \{\text{pt}\}$  时。由(3.1.8)，对于 $Y$  、 $R\Gamma(U; f^!A_X) \simeq \text{RHom}(R\Gamma_c(U; A_Y), A)$  的任意开集 $U$  。此外，如果根据命题 3.2.3， $U$  同态于  $\mathbb{R}^l$  、  $R\Gamma_c(U; A_Y) \simeq A[-l]$  ，因此我们对于  $j \neq -l$  和  $\Gamma(U; H^{-l}(f^!A_X)) \simeq \text{Hom}(H^l_c(U; A_Y), A)$  获得  $H^j(U; f^!A_X) = 0$  。

由于 $ H_c^l(Y; A_Y) \leftarrow H_c^l(U; A_Y) $  是命题3.2.3 (ii) 的同构，所以 $ A \simeq \Gamma(Y; H^{-1}(f^1A_X)) \to \Gamma(U; H^{-1}(f^1A_X)) $  是同构。这表明  $ H^{-1}(f^1A_X) $  是与  $ A_Y $  同构的常值层。现在，让我们证明一般情况。问题是局部的，我们可以假设  $ Y = \mathbb{R}^l \times X $  和  $ f $  是投影。令 $ p $  为从 $ Y $  到 $ \mathbb{R}^l $  的投影，让 $ a_X $  、 $ a_{R^l} $  分别为从 $ X $  和 $ \mathbb{R}^l $  到 $ \{\text{pt}\} $  的投影。根据命题 3.1.9，我们有态射  $ p^{-1} \omega_{R^l} \to f^1 A_X $  。因此，对于任何  $ F \in \text{Ob}(\mathbf{D}^+(A_X)) $  ，我们都有态射链：

\[
p^{-1}\omega_{{\mathbb R}^{l}}\otimes f^{-1}F\to f^{!}A_{X}\otimes f^{-1}F\to f^{!}F\enspace.
\]

足以证明该组合是同构的。对于  $\mathbb{R}^l$  的开集  $U$  和  $X$  的开集  $V$  ，如果  $U$  同构于  $\mathbb{R}^l$  ，那么我们有一个同构链：

\[
\begin{aligned}{R\varGamma(U\times V;f^{!}F)}&{{}\simeq R\mathrm{H o m}(A_{U\times V},f^{!}F)}\\ {}&{{}\simeq R\mathrm{H o m}(R f_{!}A_{U\times V},F)}\\ {}&{{}\simeq R\mathrm{H o m}\bigg(R\varGamma_{c}(U;A_{U})\stackrel{L}{\otimes}A_{V},F\bigg)}\\ {}&{{}\simeq R\mathrm{H o m}(R\varGamma_{c}(U;A_{U}),A)\stackrel{L}{\otimes}R\mathrm{H o m}(A_{V},F)}\\ {}&{{}\simeq R\varGamma(U;\omega_{\mathbb{R}^{l}})\stackrel{L}{\otimes}R\varGamma(V;F)\enspace.\quad\Box}\\ \end{aligned}
\]

\statementlabel{定义 3.3.3}令  $f: Y \to X$  为具有纤维维度  $l$  的拓扑浸没。一套：

\[
{\mathit{o r}}_{\mathbf{Y}/\mathbf{X}}=H^{-1}(\omega_{\mathbf{Y}/\mathbf{X}})
\]

其中 $\omega_{Y/X}$  是相对对偶复形（定义3.1.16），并将 $or_{Y/X}$  称为相对方向束。如果 $X = \{\mathrm{pt}\}$ （在这种情况下Y是拓扑流形），则写 $or_{Y}$ 而不是 $or_{Y/X}$ ，并将 $or_{Y}$ 称为Y上的方向束。

根据命题 3.3.2 我们有：

\[
\omega_{{\mathbf{Y}}/X}\simeq o r_{{\mathbf{Y}}/X}[l]\enspace.
\]

\statementlabel{命题 3.3.4}令  $f: Y \to X$  为具有纤维维度  $l$  的拓扑浸没。

(i) 令 $\omega_{Y/X}^{Z}$  和 $or_{Y/X}^{Z}$  分别为相对于基环 $\mathbb{Z}$  的相对对偶复合体和相对定向层。然后我们有：

\[
\omega_{Y/X}\simeq A_{Y}\mathop{{\otimes}_{\mathbb{Z}_{Y}}}\omega_{Y/X}^{\mathbb{Z}}\qquad{~a n d~}\qquad\mathcal{o r}_{Y/X}\simeq A_{Y}\mathop{{\otimes}_{\mathbb{Z}_{Y}}}\mathcal{o r}_{Y/X}^{\mathbb{Z}}.
\]

(ii) 具有规范同构：

\[
o r_{\mathbf{Y}/X}\otimes o r_{\mathbf{Y}/X}\simeq A_{Y}~,
\]

\[
\mathcal{H o m}({o r}_{\Upsilon/X},A_{\Upsilon})\simeq{o r}_{\Upsilon/X}.
\]

(iii) 令 $g: Z \to Y$  为连续映射，并假设 $f \circ g$  是具有纤维维度 $m$  的拓扑浸没。让 $F \in \mathrm{Ob}(\mathbb{D}^+) (A_X)$  。然后：

\[
g^{1}\circ f^{-1}F\simeq(f\circ g)^{-1}F\otimes\mathfrak{o r}_{Z/X}\otimes g^{-1}\mathfrak{o r}_{Y/X}[m-l]~.
\]

\statementlabel{证明}(i) 由命题 3.3.2 得出。

(ii) 由练习 III.3 得出。

(iii) 根据命题 3.3.2 和前面的结果，我们有：

\[
(f\circ g)^{!}F\simeq(f\circ g)^{-1}F\otimes\mathfrak{o r}_{Z/X}[m]~,
\]

\[
f^{!}F\simeq f^{-1}F\otimes o{\mathfrak r}_{Y/X}[l]\enspace.
\]

由于  $or_{Y/X}$  与  $A_{Y}$  局部同构，结果如下：

\[
g^{!}\circ f^{-1}F\otimes o\mathcal{r}_{Y/X}[l]\simeq g^{!}\circ f^{!}F\simeq(f\circ g)^{-1}F\otimes o\mathcal{r}_{Z/X}[m]~.\quad\Box
\]

\statementlabel{备注 3.3.5}令 $ Y \xrightarrow{f} X $  为 $ p_Y \xrightarrow{S} p_X $  连续映射的交换图

局部紧凑的空间。假设 $p_X$ 和 $p_Y$ 分别是具有纤维尺寸 $m$ 和 $n$ 的拓扑浸没。然后 $f^!A_X \simeq or_{Y/S} \otimes f^{-1} or_{X/S}[n - m]$ ，它是

自然设置：

\[
o r_{\Upsilon/X}=o r_{\Upsilon/S}\otimes f^{-1}o r_{X/S}.
\]

请注意，如果  $g: Z \to Y$  是另一个连续映射，使得  $p_Y \circ g$  是拓扑淹没，则：

\[
\omega_{\mathbf{Z}/X}\simeq\omega_{\mathbf{Z}/\mathbf{Y}}\otimes g^{-1}\omega_{\mathbf{Y}/X}\qquad\mathrm{a n d}\qquad\mathfrak{o r}_{\mathbf{Z}/X}\simeq\mathfrak{o r}_{\mathbf{Z}/\mathbf{Y}}\otimes g^{-1}\mathfrak{o r}_{\mathbf{Y}/X}.
\]

\statementlabel{命题 3.3.6}令 X 为 n 维  $ C^{0} $  流形。

(i) or\_\{X\} 是与预层  $ U \mapsto \text{Hom}(H_{c}^{n}(U; A_{X}), A) $  关联的束。

(ii) 对于  $ x \in X $  ，存在规范同构  $ \sigma\iota_{X,x} \simeq \mathrm{Hom}(H^n_{\{x\}}(X; A_X), A) \simeq H^n_{\{x\}}(X; A_X) $  。

(iii) 假设 $X$  可微且有向。然后存在同构  $or_{X} \simeq A_{X}$  ，并且当反转  $X$  的方向时，这种同构会改变符号。

\statementlabel{证明}(i) 设  $U$  为与  $\mathbb{R}^n$  同胚的  $X$  的开子集。然后：

\[
\begin{aligned}{R\varGamma(U;\omega_{X})}&{{}\simeq R\mathrm{H o m}(A_{U};\omega_{X})}\\ {}&{{}\simeq R\mathrm{H o m}(R\varGamma_{c}(X;A_{U}),A)}\\ {}&{{}\simeq\mathrm{H o m}(H_{c}^{n}(U;A_{U}),A)[n],}\\ \end{aligned}
\]

鉴于(3.1.7)。

(ii) 由前面的讨论和命题 3.2.3 和 3.3.4 (ii) 得出。

(iii) 令  $ X = \bigcup_i U_i $  为  $ X $  的开覆盖，其开子集微分同胚于  $ \mathbb{R}^n $  ，并且方向兼容。那么同构 $ \phi_{U_i} : \sigma \circ_X |_{U_i} \simeq A_{U_i} $ 将通过下面的引理3.3.7粘合在一起，并且如果颠倒方向， $ \phi_{U_i} $ 将被 $ -\phi_{U_i} $ 替换，根据命题3.2.3。   $ \square $ 

\statementlabel{引理 3.3.7}令  $E$  为欧几里得空间  $\mathbb{R}^n$  ，并修复同构  $\sigma t_E \simeq A_E$  。令 $U$  和 $V$  为 $E$  和 $f: U \to V$  微分同胚的两个开子集。假设  $f$  的雅可比行列式在  $U$  的每个点上都是正值。那么下图是可交换的：

\[
\begin{array}{c c c c c}{}&{\sim}&{}&{\sim}&{}\\ {\mathcal{o r}_{U}}&{\phantom{...}}&{\mathcal{o r}_{E/U}}&{\phantom{...}}&{A_{U}}\\ {\bigg\downarrow\bigg\{}&{}&{}&{}&{\bigg\downarrow\bigg\}}\\ {\mathcal{f}_{o i}^{\#}}&{}&{}&{}&{\mathcal{f}_{A}^{\#}}\\ {\quad}\\ {f^{-1}(\mathcal{o r}_{V})}&{\sim}&{f^{-1}(\mathcal{o r}_{E/V})}&{\sim}&{f^{-1}(A_{V})\enspace.}\\ \end{array}
\]

（这里态射 $f_o^\#$ 和 $f_A^\#$ 定义如下。令 $\mathrm{a}_U$ （分别 $\mathrm{a}_V$ ）为 $U$ （分别 $V$ ）在 $\{\mathrm{pt}\}$ 上的投影。然后 $f_A^\#$ 是同构 $f^{-1} \circ \mathrm{a}_V^{-1} \simeq \mathrm{a}_U^{-1}$ 和 $f_o^\#$ 同构 $f^{-1} \circ \mathrm{a}_V^![-n] \simeq f^! \circ \mathrm{a}_V^![-n] \simeq \mathrm{a}_U^![-n]$ 。请注意 $f^{-1} \simeq f^!$ 。）

\statementlabel{证明}足以证明对于每个  $x_{\circ} \in U$  ，下图可交换：

\begin{center}
\includegraphics[width=0.40\textwidth]{流行上的层_Chn-assets/img_in_image_box_352_259_835_437.png}
\end{center}

我们可以假设  $ x_{\circ} = f(x_{\circ}) = 0 $  。设置 $ u = f'(0) $ 、 $ f_{\lambda}(x) = u(x) + \lambda(f(x) - u(x)) $ 、 $ 0 \leq \lambda \leq 1 $ 。我们可能会找到 0 的开邻域  $ U' $  和  $ V' $  使得  $ f_{\lambda}(U') \subset V' $  和  $ f_{\lambda}^{-1}(\{0\}) \cap U' \subset \{0\} $  。由于  $ u^{\#} $  作为  $ H_{\{0\}}^{n}(E; A_{E}) $  上的恒等式运行（命题 3.2.3），因此备注 2.7.6 中的  $ f^{\#} $  也是如此。   $ \Box $ 

我们可以从众所周知的函子中恢复函子 $f^!$ 。事实上，让 $f: Y \to X$  成为 $C^0$  流形的映射。我们可以将  $f$  分解为封闭嵌入和淹没的复合体， $f = p \circ j$  ：

\[
f\colon Y\subset_{\substack{\textsf{c j}}\to Y}\times X\xrightarrow{\quad}p\to X\enspace,
\]

其中  $p$  是投影， $j$  是图形映射， $j(y) = (y, f(y))$  。应用命题 3.1.8、3.1.12 和 3.3.2，我们得到  $F \in \mathrm{Ob}(\mathbb{D}^+) (A_X)$  ：

\[
f^{!}F\simeq j^{-1}R\varGamma_{j(Y)}(p^{-1}F)\otimes\mathcal{o r}_{Y}[\operatorname{d i m}Y]~.
\]

选择  $F = A_{x}$  ，我们得到：

\[
o r_{Y/X}\simeq j^{-1}(H_{j(Y)}^{\dim X}(A_{Y\times X}))\otimes o r_{Y}.
\]

对于  $ f = id_{x} $  ，我们得到：

\[
{o r}_{X}\simeq H_{X}^{\dim X}(A_{X\times X})|_{X}
\]

（其中 X 用  $ X \times X $  的对角线标识）。）

符号 3.3.8。在本书中，除了 X§3 之外，我们用  $\operatorname{dim} X$  表示实流形  $X$  的维数。如果  $f: Y \to X$  是  $C^0$  流形的态射，我们设置：

\[
\dim Y/X=\dim Y-\dim X.
\]

如果 Y 是子流形，我们还可以写：

\[
\operatorname{c o d i m}_{X}Y=-\operatorname{d i m}Y/X,
\]

如果不存在混淆的风险，我们会写 codim Y 而不是 codim  $ _{X} $  Y。请注意：

\[
\omega_{Y/X}\simeq\mathcal{o r}_{Y/X}[\dim Y/X]\enspace.
\]

那么很自然地设置：

\[
\begin{aligned}{\omega_{Y/X}^{\otimes-1}}&{{}=R\mathcal{H}{}_{}{o m}(\omega_{Y/X},A_{Y})\enspace,}\\ {}&{{}\simeq{o r}_{Y/X}[-\operatorname{d i m}Y/X]\enspace.}\\ \end{aligned}
\]

\statementlabel{命题 3.3.9}令 $f: Y \to X$  为局部紧空间的连续映射。假设：

(i) f 是拓扑淹没。

(ii)  $ Rf_1f^!\mathbb{Z}_X \to \mathbb{Z}_X $  是同构。

那么对于  $ F \in \mathrm{Ob}(\mathbb{D}^{+}(\mathbb{Z}_X)) $  ，态射  $ F \mapsto R f_{\star} f^{-1} F $  是同构。

\statementlabel{证明}令 $l$  为 $f$  的纤维尺寸。我们有  $f^!F \simeq f^{-1}F \otimes f^!\mathbb{Z}_X$  ，  $f^!\mathbb{Z}_X$  与  $\mathbb{Z}_Y[l]$  局部同构。所以：

\[
f^{-1}F\simeq R\mathcal{H}o m(f^{!}\mathbb{Z}_{x},f^{!}F)\enspace,
\]

\[
\begin{aligned}{R f_{\ast}f^{-1}F}&{{}\simeq R f_{\ast}R\mathcal{H}{o m}(f^{!}\mathbb{Z}_{X},f^{!}F)}\\ {}&{{}\simeq R\mathcal{H}{o m}(R f_{!}f^{!}\mathbb{Z}_{X},F)}\\ {}&{{}\simeq F\enspace.\enspace\Box}\\ \end{aligned}
\]

\statementlabel{备注 3.3.10}令  $f: Y \to X$  为局部紧空间的连续映射，并假设  $f$  是具有纤维维度  $l$  的拓扑浸没。那么命题 3.3.9 的条件 (ii) 将被满足当且仅当对于任何  $x \in X$  我们有同构：

\[
R\varGamma_{c}(f^{-1}(x);\omega_{f^{-1}(x)})\xrightarrow{\sim}\mathbb{Z}\enspace.
\]

这种同构等价于以下同构：

\[
\mathbb{Z}\xrightarrow{\sim}R\Gamma(f^{-1}(x);\mathbb{Z}_{f^{-1}(x)}).
\]

事实上，设置  $ M = R\Gamma_c(f^{-1}(x);\omega_{f^{-1}(x)}) $  和  $ M^* = R\mathrm{Hom}(M;\mathbb{Z}) $  。那么  $ M^* \simeq R\Gamma(f^{-1}(x);\mathbb{Z}_{f^{-1}(x)}) $  ，同构  $ M \simeq \mathbb{Z} $  （在  $ \mathbb{D}^b(\mathfrak{Mod}(\mathbb{Z})) $  中）等价于同构  $ \mathbb{Z} \simeq M^* $  ，鉴于练习 1.31。

现在让  $X$  为维度  $n$  的  $C^0$  流形，并让  $a_X$  表示映射  $X \to \{\mathrm{pt}\}$  。态射 $R a_X^! a_X^! A_{\{\mathrm{pt}\}} \to A_{\{\mathrm{pt}\}}$  定义态射：

\[
R\mathbf{a}_{X!}\omega_{X}\to A\enspace.
\]

取第0次上同调，我们得到“积分态射”，我们用 $ \int_{X} $ 表示

\[
\int_{X}:H_{c}^{n}(X;{\mathcal{o t}}_{X})\to A.
\]

另一方面，当 $A = \mathbb{C}$  和 $X$  是 $C^{\infty}$  流形时，有一个很好的-

已知态射  $ H_c''(X; \sigma_X) \to \mathbb{C} $  ，获得如下。束  $ \sigma_X $  与 de Rham 复合体准同构（参见 II §9）：

\[
0\longrightarrow\mathcal{C}_{X}^{\infty,(0)}\otimes o\mathfrak{r}_{X}\xrightarrow[d]{}\cdots\longrightarrow\mathcal{C}_{X}^{\infty,(n)}\otimes o\mathfrak{r}_{X}\longrightarrow0~.
\]

由于 $ \mathcal{C}_X^{\infty,(\tilde{j})}\otimes\sigma r_X $  是c-soft，所以我们有

\[
H_{c}^{n}(X;\mathcal{o r}_{X})\simeq\Gamma_{c}(X;\mathcal{C}^{\infty,(n)}\otimes\mathcal{o r}_{X})/\mathrm{d}\Gamma_{c}(X;\mathcal{C}_{X}^{\infty,(n-1)}\otimes\mathcal{o r}_{X})~.
\]

如果  $\phi$  是紧支持密度，即  $\Gamma_c(X;\mathcal{C}_X^{\infty,(n)}\otimes\sigma\iota_X)$  的元素，则  $\int_X\phi$  有意义，并且如果  $\phi = \mathrm{d}\psi$  为零，则根据斯托克斯定理，对于某些  $\psi \in \Gamma_c(X;\mathcal{C}_X^{\infty,(n-1)}\otimes\sigma\iota_X)$  。因此  $\int_X\mathrm{d}\rho$  定义了一个态射：

\[
\int_{X}:\varGamma_{c}(X;\mathcal{C}_{X}^{\infty,(n)}\otimes\mathcal{o r}_{X})/\mathsf{d}\varGamma_{c}(X;\mathcal{C}_{X}^{\infty,(n-1)}\otimes\mathcal{o r}_{X})\to\mathbb{C}.
\]

此态射 (3.3.16) 与 (3.3.15) 一致，直到符号为止。

我们不会在这里证明它并将其作为练习（参见练习 III.20）。我们只注意到，如果  $X$  连接，则  $H_c^n(X; o\tau_X) \simeq \mathrm{Hom}(H^0(X;\mathbb{C}_X);\mathbb{C}) \simeq \mathbb{C}$  ，这意味着 (3.3.15) 和 (3.3.16) 等于非零常数。

我们将在第九章中再次讨论积分态射。为了结束本节，我们将给  $ \mathrm{hd}(\mathfrak{Mod}(A_x)) $  一个绑定。

\statementlabel{命题 3.3.11}令 $X$  为维度 $n$  的 $C^{0}$  流形，并令 $A$  为环。那么范畴  $\mathfrak{Mod}(A_{x})$  的同调维数（参见练习 1.17）以  $3n + \mathrm{gld}(A) + 1$  为界。

\statementlabel{证明}让  $F$  和  $G$  属于  $\mathfrak{Mod}(A_x)$  。我们想证明：

\[
\operatorname{H o m}_{\mathsf{D}(A x)}(G,F[j])=0\qquad\mathrm{f o r}\qquad j>3n+\mathsf{g l d}(A)+1.
\]

让  $ \delta_X $  表示对角线嵌入  $ X \hookrightarrow X \times X $  。根据命题 3.1.14，我们有：

\[
R\mathcal{H}o m(G,F)\simeq\delta_{X}^{!}R\mathcal{H}o m(q_{2}^{-1}G,q_{1}^{!}F)\enspace.
\]

因此：

\[
\begin{aligned}{\operatorname{H o m}_{\mathsf{D}(A x)}(G,F[j])}&{{}\simeq H^{j}(R\varGamma(X;R\mathcal{H o m}(G,F)))}\\ {}&{{}\simeq H^{j}(R\varGamma_{\varDelta}(X\times X;R\mathcal{H o m}(q_{2}^{-1}G,q_{1}^{!}F)))~.}\\ \end{aligned}
\]

根据命题 3.1.15，我们有

\[
R\Gamma(U\times V;R\mathcal{H}o m(q_{2}^{-1}G,q_{1}^{!}F))\simeq R\mathrm{H o m}(R\Gamma_{c}(U;G),R\Gamma(U;F))~.
\]

因此 $H^j(R\mathcal{H}om(q_2^{-1}G, q_1^!F)) = 0$  对应 $j > n + \mathrm{gld}(A)$  ，由命题3.2.2 (iv) 得出，结果由命题3.2.2 (v) 得出。 □

请注意，上面的估计远非最佳。

\statementlabel{推论 3.3.12}令 $X$  成为 $C^0$  流形。那么  $R\mathcal{H}om(\cdot,\cdot)$  是从  $\mathbf{D}^b(A_X)^\circ \times \mathbf{D}^b(A_X)$  到  $\mathbf{D}^b(A_X)$  的明确定义的函子。

（回想一下，我们假设了  $ \mathrm{gld}(A) < \infty $  。）

\subsection{上同调构造层}
令 $X$  为具有有限 $c$  软维数的局部紧空间。回想一下（参见练习 I.30），如果  $\mathbf{D}^{b}(\mathfrak{Mod}(A))$  的对象  $M$  与有限生成的射影  $A$  模的有界复形准同构，则该对象被称为“完美”。

\statementlabel{定义 3.4.1}如果对于所有  $x \in X$  ，满足以下条件，则  $\mathbf{D}^{b}(A_{X})$  的对象  $F$  称为上同调可构造的。

(i) $\varinjlim_{x\in U}R\Gamma(U;F)$ 和 $\varprojlim_{x\in U}R\Gamma_c(U;F)$ 是可表示的（$U$ 遍历 $x$ 的开邻域）。

(ii) $\varinjlim_{x\in U}R\Gamma(U;F)\to F_x$ 和 $R\Gamma_{\{x\}}(X;F)\to\varprojlim_{x\in U}R\Gamma_c(U;F)$ 是同构。

(iii) 复合体 $F_{x}$  和 $R\Gamma_{\{x\}}(X;F)$  是完美的。

对于“lim”和“lim”，请参见 I §11。

\statementlabel{备注 3.4.2}请注意，(ii) 由 (i) 得出。事实上第一个同构是显而易见的。为了证明第二个，让我们采用紧子集的递减序列  $ \{K_n\}_{n} $  ，它形成  $ x $  的邻域系统。然后“ $ \lim $ ” $ R\Gamma_c(X; F) \simeq \text{“}\underbrace{\lim} \text{”} R\Gamma_{K_n}(X; F) $ 。因此，对于每个 $ k \in \mathbb{Z} $ ，“ $ \lim $ ” $ H_{K_n}^k(X; F) $  是可表示的，并且射影系统 $ \{H_{K_n}^k(X; F)\}_n $  满足 $ M $  -  $ L $  条件。根据命题 2.7.1，我们得到  $ \lim_{k \to \infty} H_{K_n}^k(X; F) \simeq H_{\{x\}}^k(X; F) $  。

\statementlabel{命题 3.4.3}假设 F 是上同调可构造的。然后：

(i) DF 是上同调可构造的（参见定义 3.1.16），

(ii)  $ F \rightarrow DDF $  是同构，

(iii) 对于任何  $ x \in X $  、  $ R\Gamma_{\{x\}}(X; \mathrm{DF}) \simeq R\mathrm{Hom}(F_x, A) $  和  $ (\mathrm{DF})_x \simeq R\mathrm{Hom}(R\Gamma_{\{x\}}(X; F), A) $  。

\statementlabel{证明}(i) 和 (iii) 根据 (3.1.8)，我们有：

\[
R\varGamma(U;\mathbf{D}F)\simeq R\mathrm{H o m}(R\varGamma_{c}(U;F),A),
\]

应用函子“ $ \underrightarrow{\lim} $ ”我们得到：

\[
\begin{aligned}{\varprojlim_{x\in U}R\varGamma(U;\mathbf{D}F)}&{{}\simeq R\mathbf{H o m}\left(\varprojlim_{x\in U}R\varGamma_{c}(U;F),A\right)}\\ {}&{{}\simeq R\mathbf{H o m}(R\varGamma_{\{x\}}(X;F),A)~,}\\ \end{aligned}
\]

这显示了(iii)中的第二个同构。

因此“  $\lim_{x\to U}R\Gamma(U;DF)$  是可表示的和完美的。让 $K$  成为 $x$  的紧邻域， $\dot{K}$  其内部。我们有：

\[
\begin{aligned}{R\varGamma_{K}(X;\mathbf{D}F)}&{{}\simeq R\mathbf{H o m}(A_{K},\mathbf{D}F)}\\ {}&{{}\simeq R\mathbf{H o m}(F_{K},\omega_{X})}\\ {}&{{}\simeq R\mathbf{H o m}(R\varGamma(X;F_{K}),A).}\\ \end{aligned}
\]

应用函子“l  $ \underline{i} $  m”，我们得到：

\[
\begin{aligned}{\varprojlim_{x\in U}R\varGamma_{c}(U;\mathbf{D}F)}&{{}\simeq\varprojlim_{x\in\mathring{K}}R\varGamma_{K}(X;\mathbf{D}F)}\\ {}&{{}\simeq R\mathbf{H o m}\bigg(\varprojlim_{x\in\mathring{K}}R\varGamma(X;F_{K}),A\bigg)}\\ {}&{{}\simeq R\mathbf{H o m}\bigg(\varprojlim_{x\in U}R\varGamma(U;F),A\bigg)}\\ {}&{{}\simeq R\mathbf{H o m}(F_{x},A)~,}\\ \end{aligned}
\]

并且这个复形是完美的，这表明了(iii)中的第一个同构。

(ii) 通过 (i)，DF 和 DDF 都是可构造的。对于 $ x \in X $ ，我们通过 (iii) 得到：

\[
\begin{aligned}{(\mathrm{D D F})_{x}}&{{}\simeq R\mathrm{H o m}(R\varGamma_{\{x\}}(X;\mathrm{D F}),A)}\\ {}&{{}\simeq R\mathrm{H o m}(R\mathrm{H o m}(F_{x},A),A)\simeq F_{x}~.}\\ \end{aligned}
\]

因此 $F \to DDF$  是同构。 ⑨

\statementlabel{命题 3.4.4}令 $X$ 和 $Y$ 为两个具有有限 $c$ 软维数的局部紧空间，并让 $q_1$ 和 $q_2$ 分别表示从 $X \times Y$ 到 $X$ 和 $Y$ 的第一和第二投影。令 $F \in \mathrm{Ob}(\mathbf{D}^b(A_X))$  和 $G \in \mathrm{Ob}(\mathbf{D}^+(A_Y))$  并假设 $F$  是上同调可构造的。然后：

\[
\mathbf{D}\overset{L}{\underset{}{\boldsymbol{F}}{\underset{\boxtimes}{\boxtimes}}}\boldsymbol{G}\to\boldsymbol{R}\mathcal{H}\mathcal{m}(\boldsymbol{q}_{1}^{-1}\boldsymbol{F},\boldsymbol{q}_{2}^{!}\boldsymbol{G})
\]

是同构。如果 $X$  是 $C^{0}$  流形，那么

\[
\mathrm{D}^{\prime}F\mathop{\boxtimes}^{L}G\to R\mathcal{H}o m(q_{1}^{-1}F,q_{2}^{-1}G)
\]

是同构。

\statementlabel{证明}证明第一个同构就足够了。

设U和V分别是X和Y的开子集。根据命题 3.1.15，我们有：

\[
R\Gamma(U\times V;R\mathcal{H}o m(q_{1}^{-1}F,q_{2}^{!}G))\simeq R\mathrm{H o m}(R\Gamma_{c}(U;F),R\Gamma(V;G))\enspace.
\]

应用函子“ $ \lim_{x\in U} $ ”我们得到：

\[
\begin{aligned}\varprojlim_{x\in U}R\varGamma(U\times V;R\mathcal{H}om(q_{1}^{-1}F,q_{2}^{\prime}G))&\simeq R\mathrm{Hom}\left(\varprojlim_{x\in U}R\varGamma_{c}(U;F),R\varGamma(V;G)\right)\\&\simeq R\mathrm{Hom}(R\varGamma_{\{x\}}(X;F),R\varGamma(V;G))\\&\simeq R\mathrm{Hom}(R\varGamma_{\{x\}}(X;F),A)\otimes^{L}R\varGamma(V;G)\\&\simeq(\mathrm{DF})_{x}\otimes^{L}R\varGamma(V;G)\ .\end{aligned}
\]

因此：

\[
\begin{aligned}{{R}\mathcal{H}{_{{o m}}}(q_{1}^{-1}F,q_{2}^{!}G)}&{{}\simeq q_{1}^{-1}\mathbf{D}F\overset{L}{\otimes}q_{2}^{-1}G}\\ {}&{{}\simeq q_{1}^{-1}\mathbf{D}^{\prime}F\overset{L}{\otimes}q_{2}^{!}G~.\quad\Box}\\ \end{aligned}
\]

示例 3.4.5。 (i) 令 X 为  $C^0$  流形，Y 为余维  $p$  的闭子流形。那么  $A_Y$  和  $R\Gamma_Y(A_X)$  在  $X$  上可同调构造，而且：

\[
\mathbf{D}_{X}(A_{Y})\simeq R\varGamma_{Y}(\omega_{X})\simeq\omega_{Y}~.
\]

(ii) 假设  $X$  是实数有限维向量空间，并令  $Z$  为  $X$  的闭（或开）凸子集。那么 $A_Z$  和 $R\Gamma_Z(A_X)$  在 $X$  上可同调构造（参见练习III.4）。

\statementlabel{命题 3.4.6}令 F 和 G 为  $ \mathbf{D}^{b}(X) $  的两个上同调可构造对象。然后：

\[
\begin{aligned}{R\mathcal{H}o m(G,F)}&{{}\simeq R\mathcal{H}o m(\mathbf{D}F,\mathbf{D}G)}\\ {}&{{}\simeq\mathbf{D}\bigg(\mathbf{D}(F)\overset{L}{\otimes}G\bigg).}\\ \end{aligned}
\]

\statementlabel{证明}证明规范态射：

\[
R\mathcal{H}o m(G,F)\to R\mathcal{H}o m(\operatorname{D}F,\operatorname{D}G)
\]

是同构，检查就足够了（参见命题 3.1.14）：

\[
q_{2}^{-1}{\bf D}{\cal G}\stackrel{L}{\otimes}q_{1}^{!}F\simeq q_{1}^{-1}({\bf D}{\bf D}{\cal F})\stackrel{L}{\otimes}q_{2}^{!}{\bf D}{\cal G},
\]

这是显而易见的。

另一方面我们有：

\[
\begin{aligned}{{R}\mathcal{H}{}_{}{o m}\bigg(\operatorname{D F}\otimes{}^{{L}}{G},\omega_{X}\bigg)}&{{}\simeq{R}\mathcal{H}{}_{}{o m}({G},{R}\mathcal{H}{}_{}{o m}(\operatorname{D F},\omega_{X}))}\\ {}&{{}\simeq{R}\mathcal{H}{}_{}{o m}({G},\operatorname{D D F})}\\ {}&{{}\simeq{R}\mathcal{H}{}_{}{o m}({G},F)~.\quad\Box}\\ \end{aligned}
\]

我们将在本书中始终遇到可构造的层。

\subsection{$\gamma$-拓扑}
令  $V$  为实数有限维向量空间，并令  $\gamma$  为  $V$  中的封闭凸锥（顶点位于  $0$  ）。人们可以在  $V$  上定义一个新的拓扑，与  $\gamma$  相关联。

\statementlabel{定义 3.5.1}$ V $  上的 $ \gamma $  拓扑是开集 $ \Omega $  满足的拓扑：

(i)  $ \Omega $  对通常的拓扑开放，

（二） $ \Omega + \gamma = \Omega $ 。

让  $X$  成为  $V$  的子集。一个由  $X_{v}$  表示具有诱导的  $\gamma$  拓扑的集合  $X$ ，并且由

\[
\phi_{\gamma}\colon X\to X_{\gamma}
\]

从  $X$  （具有通常的拓扑）到  $X_{\gamma}$  的自然连续映射。也可以写  $\phi_{\gamma}^{X}$  而不是  $\phi_{\gamma}$  以强调空格  $X$  。

请注意， $\{0\}$  拓扑是通常的拓扑，如果 $\gamma_1 \subset \gamma_2$  是两个封闭的凸锥，则自然映射 $X_{\gamma_1} \to X_{\gamma_2}$  是连续的。

示例 3.5.2。采取 $ X = \mathbb{R} $ ， $ \gamma = [0, +\infty] $ 。  $ \gamma $  开集是间隔  $ ]c, +\infty $  ；   $ -\infty \leq c \leq \infty $ 。

在下文中，我们将用  $ \gamma^{a} $  表示与  $ \gamma $  相反的圆锥体，即  $ \gamma^{a} = -\gamma $  。

如果  $\Omega$  是  $X$  的子集，我们可以说  $\Omega$  是  $\gamma$  -开（分别为  $\gamma$  - 闭，分别为  $\gamma$  -  $x$  的邻域），如果  $\Omega$  是开（分别为闭，分别为  $x$  的邻域）  $\gamma$  -拓扑。

\statementlabel{命题 3.5.3}令 $X$  为 $V$  的 $\gamma$  开子集，并让 $F$  为 $X_{\gamma}$  上的一个层。

(i) 令  $U$  为  $X$  的凸开子集。那么自然态射 $R\Gamma(U + \gamma; F) \to R\Gamma(U; \phi_{\gamma}^{-1}F)$  就是同构。

(ii) 令  $K$  为  $X$  的凸紧子集。那么自然态射 $R\Gamma(K + \gamma; \phi_{\gamma}^{-1}F) \to R\Gamma(K; \phi_{\gamma}^{-1}F)$  就是同构。

(iii) 自然态射 $ F \to R\phi_{y*} \phi_{y}^{-1} F $  是同构。

\statementlabel{证明}我们将证明分解为几个步骤。

(a) 令 $U$  为 $X$  的凸子集，并令 $\psi$  为自然映射： $\Gamma(U + \gamma; F) \to \Gamma(U; \phi_{\gamma}^{-1}F)$  。让我们证明 $\psi$  是单射的。令  $s$  成为  $F$  上  $U + \gamma$  的一部分，这样对于所有  $x \in U$  ，  $(\phi_{\gamma}^{-1}F)_x$  中的  $s_x = 0$  。对于这样的  $x$  ，存在一个  $\gamma$  开集  $W$  ，其中包含  $x$  ，使得  $s_y = 0$  对于所有  $y \in W$  。因此 $s_{x+v} = 0$  对于所有 $x \in U$  、 $v \in \gamma$  。

现在让我们证明 $\psi$  是满射的。  $U$  上的  $\phi_\gamma^{-1}F$  的部分  $s$  是由开放覆盖  $U = \bigcup_{i \in I} U_i$  定义的， $F$  的  $s_i$  部分是在  $U_i + \gamma$  上定义的， $s_x = s_{i,x}$  则用于任何  $x \in U_i$  。我们将证明  $x \in (U_i + \gamma) \cap (U_j + \gamma)$  意味着  $s_{i,x} = s_{j,x}$  。让 $x_i \in U_i \cap (x + \gamma^a)$ ， $x_j \in U_j \cap (x + \gamma^a)$ ， $x_t = tx_i + (1 - t)x_j$ ，( $t \in [0,1]$ )。然后 $x_t \in U \cap (x + \gamma^a)$ 。让：

\[
A=\left\{t\in[0,1];\mathrm{f o r~a n y}k\in I\mathrm{~s u c h~t h a t}x_{t}\in U_{k},\mathrm{~w e~h a v e~}s_{i,x}=s_{k,x}\right\}
\]

集合A也可以定义为：

\[
A=\left\{t\in[0,1];\mathrm{t h e r e~e x i s t s~}k\in I\mathrm{~s u c h~t h a t~i f~}x_{t}\in U_{k}\mathrm{~w e~h a v e~}s_{i,x}=s_{k,x}\right\}.
\]

因此  $A$  是开闭的，并且包含 1。因此  $A=[0,1]$  。因此 $s_{i}|_{(U_{i}+\gamma)\cap(U_{j}+\gamma)}=s_{j}|_{(U_{i}+\gamma)\cap(U_{j}+\gamma)}$  。因此存在  $\tilde{s}\in\Gamma(U+\gamma;F)$  这样  $\tilde{s}|_{U_{i}+\gamma}=s_{i}$  。这显示了  $\psi$  的满射性。因此我们得到了同构：

\[
\Gamma(U+\gamma;F)\simeq\Gamma(U;\phi_{\gamma}^{-1}F)
\]

（特别是 $ \Gamma(U + \gamma; F) \simeq \Gamma(U + \gamma; \phi_{\gamma}^{-1}F) $ ）。

(b) 令 K 为 X 的凸紧子集。通过 (a)，

\[
\operatorname*{l i m}_{\stackrel{\rightarrow}{U}}\varGamma(U;\phi_{\gamma}^{-1}F)\to\varGamma(K;\phi_{\gamma}^{-1}F)
\]

是一个同构，其中  $U$  范围贯穿  $\gamma$  系列 -  $K$  的开邻域。对于任何凸紧集  $L$  和  $K \subset L \subset K + \gamma$  ，任何  $\gamma$  的  $K$  开邻域都包含  $L$  ，因此

\[
\Gamma(L;\phi_{\gamma}^{-1}F)\stackrel{}{\sim}\Gamma(K;\phi_{\gamma}^{-1}F)\enspace.
\]

从  $ \varGamma(K + \gamma; \phi_{\gamma}^{-1}F) \simeq \lim_{\leftarrow} \varGamma(L; \phi_{\gamma}^{-1}F) $  开始，我们得到同构：

\[
\Gamma(K+\gamma;\phi_{\gamma}^{-1}F)\stackrel{}{\sim}\Gamma(K;\phi_{\gamma}^{-1}F)\enspace.
\]

(c) 自然态射 $F \to \phi_{\gamma*}\phi_{\gamma}^{-1}F$  是同构。事实上，如果  $U$  是  $X$  的凸  $\gamma$  开子集，那么：

\[
\begin{aligned}{\varGamma(U;\phi_{\gamma\ast}\phi_{\gamma}^{-1}F)}&{{}=\varGamma(U;\phi_{\gamma}^{-1}F)}\\ {}&{{}\simeq\varGamma(U;F).}\\ \end{aligned}
\]

(d) 假设  $F$  是松弛的，并让  $\Omega$  是  $X$  的凸开子集， $K$  是  $\Omega$  的凸紧子集。限制态射  $\Gamma(\Omega; \phi_{\gamma}^{-1}F) \to \Gamma(K; \phi_{\gamma}^{-1}F)$  是 (3.5.1) 的满射。应用命题 3.2.4 我们得到：

\[
R^{j}\varGamma(\varOmega;\phi_{\gamma}^{-1}F)=0\quad\mathrm{f o r}\quad j\neq0~.
\]

这意味着  $ R^j\Gamma(K;\phi_y^{-1}F) $  和  $ R^j\phi_{y*}\phi_y^{-1}F $  对于  $ j \neq 0 $  为零。因此，我们在 $ F $  是松弛的附加假设下证明了该命题。

(e) 令 $F$  为 $X_\gamma$  上的一个层，并让 $F^\prime$  为从松弛层下方界定的复形， $F^\prime$  与 $F$  准同构。到 (3.5.3) 我们就有了  $R\varGamma(U+\gamma;F)\simeq\varGamma(U+\gamma;F^\prime)\simeq\varGamma(U;\phi_\gamma^{-1}F^\prime)\simeq R\varGamma(U;\phi_\gamma^{-1}F)$  。这就证明了(i)，(ii)和(iii)的证明类似。   $\Box$ 

我们现在将描述另一种获取  $ \phi_{\gamma}^{-1}R\phi_{\gamma*}F $  的方法。假设  $ X = V $  ，并设置：

\[
Z(\gamma)=\left\{(x,y)\in X\times X;y-x\in\gamma\right\}.
\]

让我们用  $ q_{1} $  和  $ q_{2} $  表示  $ X \times X $  的第一个和第二个投影。

\statementlabel{命题 3.5.4}让 $ F \in \mathrm{Ob}(\mathbb{D}^{+}(X)) $  。存在自然的同构：

\[
R q_{1*}((q_{2}^{-1}F)_{Z(\gamma)})\simeq\phi_{\gamma}^{-1}R\phi_{\gamma*}F\enspace.
\]

\statementlabel{证明}令  $ \tilde{q}_j: Z(\gamma) \to X $  为  $ q_j $  (  $ j = 1, 2 $  ) 的限制。那么对于任何  $ \gamma $  - $ X $  的开子集  $ \Omega $  ：

\[
\tilde{q}_{1}^{-1}\Omega=\left\{(x,y)\in X\times X;x\in\Omega,y\in x+\gamma\right\}\subset\tilde{q}_{2}^{-1}\Omega.
\]

因此，对于  $Z(\gamma)$  上的任何束  $G$  ，我们有一个态射  $(\phi_\gamma \circ \tilde{q}_2)_* G \to (\phi_\gamma \circ \tilde{q}_1)_* G$  ，并且我们获得从  $\mathbf{D}^+ (Z(\gamma))$  到  $\mathbf{D}^+ (X_\gamma): R(\phi_\gamma \circ \tilde{q}_2)_* \to R(\phi_\gamma \circ \tilde{q}_1)_*$  的函子的态射。

因此，对于任何  $ F \in \mathrm{Ob}(\mathbb{D}^{+}(X)) $  ，我们获得态射：

\[
R\phi_{\gamma*}F\to R\phi_{\gamma*}R\tilde{q}_{2*}\tilde{q}_{2}^{-1}F\to R\phi_{\gamma*}R\tilde{q}_{1*}\tilde{q}_{2}^{-1}F\quad,
\]

这给出：

\[
\phi_{\gamma}^{-1}R\phi_{\gamma\ast}F\to R\tilde{q}_{1\ast}\tilde{q}_{2}^{-1}F\quad.
\]

让我们证明这是一个同构。

对于任何凸紧集 $K$ ，投影 $\tilde{q}_{1}^{-1}K \to \tilde{q}_{2}\tilde{q}_{1}^{-1}K = K + \gamma$  具有可收缩且适当的纤维。因此，根据推论 2.7.7 (iv)，

\[
R\varGamma(\tilde{q}_{1}^{-1}K;\tilde{q}_{2}^{-1}F)\simeq R\varGamma(K+\gamma;F).
\]

因此，对于任何  $ x \in X $  和整数  $ j $  ，

\[
\begin{aligned}{H^{j}(R\tilde{q}_{1*}\tilde{q}_{2}^{-1}F)_{x}}&{{}\simeq\varinjlim_{\vec{K}}H^{j}(R\varGamma(K;R\tilde{q}_{1*}\tilde{q}_{2}^{-1}F))}\\ {}&{{}\simeq\varinjlim_{\vec{K}}H^{j}(R\varGamma(\tilde{q}_{1}^{-1}K;\tilde{q}_{2}^{-1}F))}\\ {}&{{}\simeq\varinjlim_{\vec{K}}H^{j}(R\varGamma(K+\gamma;F))}\\ {}&{{}\simeq H^{j}(\phi_{\gamma}^{-1}R\phi_{\gamma*}F)_{x}~,}\\ \end{aligned}
\]

其中 K 的范围涵盖 x 的凸紧邻域族。这表明(3.5.5)是同构。 ⑨

\subsection{核}
令 $X$  和 $Y$  为两个具有有限 $c$  软维数的局部紧空间，并令 $q_1$  和 $q_2$  分别表示 $X \times Y$  上的第一和第二投影。让 $K \in \mathrm{Ob}(\mathbf{D}^b(X \times Y))$  。

\statementlabel{定义 3.6.1}我们通过以下方式定义函子 $ \Phi_K: \mathbf{D}^+ (Y) \to \mathbf{D}^+ (X) $ 和 $ \Psi_K: \mathbf{D}^+ (X) \to \mathbf{D}^+ (Y) $ ：

\[
\varPhi_{K}(G)=R q_{1!}\bigg(\overset{L}{K}\otimes\underset{}{q_{2}^{-1}G}\bigg),
\]

\[
\Psi_{K}(F)=R q_{2*}R\mathcal{H o m}(K,q_{1}^{!}F)\enspace.
\]

\statementlabel{命题 3.6.2}$ \mathbf{D}^{+}(Y)^{\circ} \times \mathbf{D}^{+}(X) $  上存在双函子同构：

\[
\mathrm{H o m}_{\mathsf{D}^{+}(x)}(\varPhi_{K}(\cdot),\cdot)\simeq\mathrm{H o m}_{\mathsf{D}^{+}(y)}(\cdot,\varPsi_{K}(\cdot))\enspace.
\]

换句话说， $ \Psi_{K} $  和 $ \Phi_{K} $  是伴随函子。

\statementlabel{证明}让 $ F \in \mathrm{Ob}(\mathbb{D}^{+}(X)) $ ， $ G \in \mathrm{Ob}(\mathbb{D}^{+}(Y)) $ 。然后：

\[
\begin{aligned}{\operatorname{H o m}_{\mathsf{D}^{+}(X)}(\varPhi_{\mathbb{K}}(G),F)}&{{}=\operatorname{H o m}_{\mathsf{D}^{+}(X)}\bigg(R q_{1!}\bigg(\overset{L}{K}\otimes\overset{L}{q_{2}^{-1}G}\bigg),F\bigg)}\\ {}&{{}\simeq\operatorname{H o m}_{\mathsf{D}^{+}(X\times Y)}\bigg(\overset{L}{K}\otimes\overset{L}{q_{2}^{-1}G},\overset{!}{q_{1}^{\vphantom{1}}}F\bigg)}\\ {}&{{}\simeq\operatorname{H o m}_{\mathsf{D}^{+}(X\times Y)}(q_{2}^{-1}G,R\mathcal{H}o m(K,\overset{!}{q_{1}^{\vphantom{1}}}F))}\\ {}&{{}\simeq\operatorname{H o m}_{\mathsf{D}^{+}(Y)}(G,R q_{2*}R\mathcal{H}o m(K,\overset{!}{q_{1}^{\vphantom{1}}}F))\enspace.\quad\Box}\\ \end{aligned}
\]

\statementlabel{命题 3.6.3}假设  $X = Y$  和  $K = A_\Delta$  ，其中  $\Delta$  是  $X \times X$  的对角线。那么  $\Phi_K$  和  $\Psi_K$  与  $\mathbf{D}^+ (X)$  中的恒等式同构。

\statementlabel{证明}根据命题 3.1.14，我们有：

\[
\begin{aligned}{F}&{{}\simeq R q_{2*}{R}\varGamma_{\varDelta}{R}\mathcal{H}{o m}(q_{2}^{-1}A_{X},q_{1}^{!}F)}\\ {}&{{}\simeq R q_{2*}{R}\mathcal{H}{o m}(A_{\varDelta},q_{1}^{!}F)}\\ {}&{{}\simeq\varPsi_{K}(F).}\\ \end{aligned}
\]

从  $ A_{\Delta} \otimes q_{2}^{-1}G \simeq (q_{2}^{-1}G)_{\Delta} $  开始，我们也有类似的  $ \varPhi_{K}(G) \simeq G $  。 □

令 $Z$  为具有有限 $c$  软维数的第三个局部紧空间。我们用 $q'_1$ 和 $q'_2$ （分别 $q''_1$ 和 $q''_2$ ）表示 $X \times Z$ （分别在 $Y \times Z$ ）上的第一和第二投影，并引入符号 $q_{ij}$ 来表示 $X \times Y \times Z$ 上的 $(i,j)$ 投影。例如  $q_{13}$  是到  $X \times Z$  的投影。

\begin{center}
\includegraphics[width=0.43\textwidth]{流行上的层_Chn-assets/img_in_image_box_320_632_841_943.png}
\end{center}

\statementlabel{命题 3.6.4}让 $ K_1 \in \mathrm{Ob}(\mathbb{D}^b(X \times Y)) $ ， $ K_2 \in \mathrm{Ob}(\mathbb{D}^b(Y \times Z)) $ 。设定：

\[
K=R q_{13!}\left(q_{12}^{-1}K_{1}\stackrel{L}{\otimes}q_{23}^{-1}K_{2}\right).
\]

然后 $ \Psi_{K_2} \circ \Psi_{K_1} \simeq \Psi_K $  和 $ \Phi_{K_1} \circ \Phi_{K_2} \simeq \Phi_K $  。

\statementlabel{证明}让 $ F \in \mathrm{Ob}(\mathbb{D}^{+}(A_X)) $  。然后：

\[
\begin{aligned}{\mathcal{W}_{K_{2}}\circ\mathcal{W}_{K_{1}}(F)}&{{}=R q_{2*}^{{^{\prime \prime}}}R\mathcal{H}{o m}(K_{2},q_{1}^{{^{\prime \prime}1}}R q_{2*}R\mathcal{H}{o m}(K_{1},q_{1}^{!}F))}\\ {}&{{}\simeq R q_{2*}^{{^{\prime \prime}}}R\mathcal{H}{o m}(K_{2},R q_{23*}q_{12}^{!}R\mathcal{H}{o m}(K_{1},q_{1}^{!}F))}\\ {}&{{}\simeq R q_{2*}^{{^{\prime \prime}}}R q_{23*}R\mathcal{H}{o m}(q_{23}^{-1}K_{2},R\mathcal{H}{o m}(q_{12}^{-1}K_{1},q_{12}^{!}q_{1}^{!}F))}\\ {}&{{}\simeq R q_{2*}^{^{\prime}}R q_{13*}R\mathcal{H}{o m}\bigg(q_{12}^{-1}K_{1}\otimes\stackrel{L}{q_{23}^{-1}}K_{2},q_{13}^{!}q_{1}^{{^{\prime}!}}F\bigg)}\\ {}&{{}\simeq R q_{2*}^{^{\prime}}R\mathcal{H}{o m}\bigg(R q_{13!}\bigg(q_{12}^{-1}K_{1}\otimes\stackrel{L}{q_{23}^{-1}}K_{2}\bigg),q_{1}^{{^{\prime}!}}F\bigg)}\\ {}&{{}=\mathcal{W}_{K}(F).}\\ \end{aligned}
\]

 $ \Phi_{K} $  的证明是类似的。 ⑨

我们设定：

\[
K_{1}\circ K_{2}=R q_{13!}\bigg(q_{12}^{-1}K_{1}\overset{L}{\otimes}q_{23}^{-1}K_{2}\bigg).
\]

\statementlabel{推论 3.6.5}假设  $Z = X$  、  $K_{2} \circ K_{1} \simeq A_{\Delta x}[l]$  、  $K_{1} \circ K_{2} \simeq A_{\Delta y}[l']$  对于某些班次  $l$  和  $l'$  。   $\Delta_{X}$  和 $\Delta_{Y}$  分别表示 $X \times X$  和 $Y \times Y$  的对角线）。那么  $\Phi_{K_{1}}$  、  $\Phi_{K_{2}}$  、  $\Psi_{K_{1}}$  和  $\Psi_{K_{2}}$  是范畴的等价。

请注意，在这种情况下，如果  $X \neq \emptyset$  ，则  $l$  和  $l'$  必须相等。

示例 3.6.6。令  $\tau: E \to X$  为局部紧致空间  $X$  上的实向量丛，纤维维数为  $n$  ，并令  $\pi: E^* \to X$  为对偶向量丛。令  $\dot{E}$  （分别为  $\dot{E}^*$  ）为删除了零部分的空间  $E$  （分别为  $E^*$  ），然后输入：

\[
S=\dot{E}/\mathbb{R}^{+},\qquad S^{*}=\dot{E}^{*}/\mathbb{R}^{+}.
\]

从 S 和  $ S^{*} $  到 X 的投影是纤维维数为 n - 1 的拓扑浸没。我们引入以下集合：

\[
\left\{\begin{aligned}D&=\left\{(x,y)\in S\times S^{*};\langle x,y\rangle\geqslant0\right\},\\ I&=\left\{(y,x)\in S^{*}\times S;\langle y,x\rangle>0\right\},\end{aligned}\right.
\]

以及  $ \mathbf{D}^b(S \times S^*) $  和  $ \mathbf{D}^b(S^* \times S) $  的对象：

\[
\left\{\begin{aligned}{K_{1}}&{{}=A_{D}\otimes\omega_{S^{\bullet}/X},}\\ {K_{2}}&{{}=A_{I}.}\\ \end{aligned}\right.
\]

\statementlabel{命题 3.6.7}其中一个有：

\[
K_{1}\circ K_{2}\simeq A_{\varDelta_{s}}\qquad{~i n~}\mathbf{D}^{b}(S\times S),
\]

\[
K_{2}\circ K_{1}\simeq A_{\varDelta_{S}*}\qquad{i n}\mathbf{D}^{b}(S^{*}\times S^{*})~.
\]

\statementlabel{证明}我们用  $p_{ij}$  表示  $S\times S^{*}\times S$  上的  $(i,j)$  次投影。考虑一下图表：

\begin{center}
\includegraphics[width=0.46\textwidth]{流行上的层_Chn-assets/img_in_image_box_305_1516_859_1689.png}
\end{center}

设置 $ L = p_{12}^{-1}(D) \cap p_{23}^{-1}(I) $ 。然后 $ K_{1} \circ K_{2} \simeq Rp_{13!}(A_{L} \otimes \omega_{S^*/X}) $ 。

让 $(x,x') = z \in S \times S$  。如果  $z \notin S \times_X S$  ，则  $p_{13}^{-1}(z) \cap L = \emptyset$  。如果  $z \in S \times_X S$  ，但  $z \notin S \times_S S$  ，则  $p_{13}^{-1}(z) \cap L$  同胚于  $\mathbb{R}^{n-1}$  的闭半空间，否则为空。因此在这种情况下  $(K_1 \circ K_2)_z = 0$  。

让 $ U = p_{13}^{-1}(\varDelta_S) \cap L $  。由于  $ (R p_{13!} A_L)_{\varDelta_S} \simeq R p_{13!}(A_U) $  和  $ U \to \varDelta_S $  与映射  $ \pi: I \to S $  同构，我们得到

\[
K_{1}\circ K_{2}|_{\varDelta_{S}}\simeq R\pi_{!}(A_{I}\otimes\omega_{S^{\bullet}/X})\simeq R\pi_{!}\pi_{!}^{!}A_{S}\enspace.
\]

映射  $ \pi: I \to S $  与  $ S $  上的本地投影  $ S \times \mathbb{R}^{n-1} \to S $  同构。因此 $ R\pi_i\pi_i^i A_S \to A_S $  是同构。

 $ K_{2} \circ K_{1} $  的证明是类似的。 ⑨

根据推论 3.6.5，我们看到  $\Phi_{K_{1}}, \Phi_{K_{2}}, \Psi_{K_{1}}$  和  $\Psi_{K_{2}}$  是范畴的等价物。

\subsection{傅里叶--佐藤变换}
让我们首先定义圆锥层。令  $ \mathbb{R}^{+} $  表示严格正数的乘法群，并令 X 为局部紧空间，赋予  $ \mathbb{R}^{+} $  的作用。换句话说，我们有一个连续的映射：

\[
\mu:X\times\mathbb{R}^{+}\to X\ ,
\]

满足每个  $x \in X$  ,  $t_1, t_2 \in \mathbb{R}^+$  。

\[
\left\{\begin{aligned}{\mu(x,t_{1}t_{\bar{2}})}&{{}=\mu(\mu(x,t_{1}),t_{2})}\\ {\mu(x,1)}&{{}=x}\\ \end{aligned}\right..
\]

\statementlabel{定义 3.7.1}(i) 我们用 $ \mathfrak{Mod}_{\mathbb{R}^+}(A_X) $ 表示由层 $ F $ 组成的 $ \mathfrak{Mod}(A_X) $ 的完整子类，这样对于 $ X $ 中 $ \mathbb{R}^+ $ 的任何轨道 $ b $ ， $ F|_{b} $ 是局部常值的层。

(ii) 我们用  $\mathbf{D}_{\mathbf{R}^{+}}^{+}(A_{X})$  （或简单地  $\mathbf{D}_{\mathbf{R}^{+}}^{+}(X)$  ）表示  $\mathbf{D}^{+}(A_{X})$  的满子范畴，由对象  $F$  组成，这样对于所有  $j \in \mathbb{Z}$  ，  $H^{j}(F) \in \mathfrak{Mod}_{\mathbf{R}^{+}}(A_{X})$  。

(iii) 我们将 $ \mathfrak{Mod}_{\mathbb{R}^{+}}(A_{X}) $ （或 $ \mathbb{D}_{R}^{+}(X) $ ）的对象称为圆锥对象。

考虑一下映射：

\[
X\subset_{\xrightarrow{j}}X\times\mathbb{R}^{+}\xrightarrow[\substack{\mu\to p}]{}
\]

其中  $ j(x) = (x; 1) $  ， p 是投影。一个具有自然态射：

\[
\mu^{-1}F\xleftarrow[\alpha]{}\rho^{-1}R p_{*}\mu^{-1}F\xrightarrow[\beta]{}\rho^{-1}F,
\]

其中  $\beta$  由 1 处求值的态射定义：

\[
R p_{*}\mu^{-1}F\to R p_{*}R j_{*}j^{-1}\mu^{-1}F\simeq F\enspace.
\]

\statementlabel{命题 3.7.2}让 $F \in \mathrm{Ob}(\mathbb{D}^{+}(X))$  。那么下面的条件是等价的。

（一） $ F \in \mathrm{Ob}(\mathbb{D}_{R^+}^{+}(X)) $ 。

(ii) (3.7.3) 中的态射 $ \alpha $  和 $ \beta $  是同构的。

(iii) 对于所有  $ j \in \mathbb{Z} $  ，  $ H^{j}(\mu^{-1}F) $  在 p 的纤维上是局部常数。

（四） $ \mu^{-1}F \simeq p^{-1}F $ 。

（五） $ \mu^!F \simeq p^!F $ 。

\statementlabel{证明}(iv)  $\Leftrightarrow$  (v) 和 (i)  $\Leftrightarrow$  (iii) 很清楚，以及 (ii)  $\Rightarrow$  (iv) 和 (iv)  $\Rightarrow$  (iii)。最后，(iii)  $\Rightarrow$  (ii) 由推论 2.7.7 得出。   $\square$ 

\statementlabel{推论 3.7.3}令  $U$  为  $X$  的开子集。假设对于 $X$  中 $\mathbb{R}^{+}$  的任何轨道 $b$  来说， $b \cap U$  是可收缩的（特别是非空的）。那么对于  $F \in \mathrm{Ob}(\mathbb{D}_{R^{+}}^{+}(X))$  ，限制态射  $R\Gamma(X; F) \to R\Gamma(U; F)$  是同构。

\statementlabel{证明}令  $ \mu': U \times \mathbb{R}^+\to X $  为  $ \mu $  的限制。那么 $ \mu' $ 就有可收缩纤维。

应用命题 3.3.9，我们得到：

\[
{\cal F}\simeq R\mu_{\ast}^{\prime}\mu^{\prime^{-1}}{\cal F}\quad,
\]

因此：

\[
\begin{aligned}{R\varGamma(X;F)}&{{}\simeq R\varGamma(U\times\mathbb{R}^{+};\mu^{\prime}{}^{-1}F)}\\ {}&{{}\simeq R\varGamma(U\times\mathbb{R}^{+};p^{-1}F)}\\ {}&{{}\simeq R\varGamma(U;F)~.~\Box}\\ \end{aligned}
\]

令  $X$  和  $Y$  为两个具有  $\mathbb{R}^+$  动作的空间。然后，空间 $X \times Y$ 被赋予 $\mathbb{R}^+ \times \mathbb{R}^+$ 的作用，我们将范畴 $\mathbf{D}_{\mathbf{R}^+ \times \mathbf{R}^+}^+(X \times Y)$ 的定义和双锥束的概念留给读者。

如果  $F \in \mathrm{Ob}(\mathbf{D}_{\mathbf{R}^+}^+) (X)$  ) 和  $G \in \mathrm{Ob}(\mathbf{D}_{\mathbf{R}^+}^+) (Y)$  ) 则  $F \boxtimes^L G \in \mathrm{Ob}(\mathbf{D}_{\mathbf{R}^+} \times \mathbf{R}^+) (X \times Y)$  。请注意， $X \times Y$  上的双圆锥束尤其是圆锥束，因为  $\mathbb{R}^+$  通过对角线嵌入  $\mathbb{R}^+ \hookrightarrow \mathbb{R}^+ \times \mathbb{R}^+$  对  $X \times Y$  进行对角作用。

现在让  $f: Y \to X$  是一个连续映射，并假设  $f_!$  具有有限的上同调维数，并且  $f$  通过  $\mathbb{R}^+$  的作用进行交换。我们得到一个交换图：

\textbackslash{}[\textbackslash{}begin\{array\}\{c\}Y\textbackslash{}xleftarrow\{\textbackslash{}quad\textbackslash{}mu\textbackslash{}quad\}Y\textbackslash{}times\textbackslash{}mathbb\{R\}\textasciicircum{}\{+\}\textbackslash{}xrightarrow\{\textbackslash{}quad p\_\{Y\}\textbackslash{}quad\}Y\textbackslash{}\textbackslash{}\textbackslash{}left.\textbackslash{}quad f\textbackslash{}right|\_\{\textbackslash{}quad\textbackslash{}quad\textbackslash{}quad\textbackslash{}quad\}\textbackslash{}quad\textbackslash{}bar\{f\}\textbackslash{}quad\textbackslash{}quad\textbackslash{}quad

\statementlabel{命题 3.7.4}(i) 让 $ F \in \mathrm{Ob}(\mathbb{D}_{R^+}^{+}(X)) $  。那么 $ f^{-1}F $  和 $ f^!F $  是圆锥曲线。

(ii) 让 $ G \in \text{Ob}(\mathbb{D}_{\mathbb{R}^+}^+)(\mathcal{Y})) $  。那么 $ R f_G $  和 $ R f_G $  是圆锥曲线。

(iii) 令 $F_1$  和 $F_2$  属于 $\mathrm{Ob}(\mathbb{D}_{\mathbb{R}^+}^+)(\mathcal{X})$  。那么 $F_1\otimes^L F_2$  是圆锥曲线。如果  $F_1\in\mathrm{Ob}(\mathbb{D}_{\mathbb{R}^+}^+)(\mathcal{X})$  ，则  $R\mathcal{Hom}(F_1,F_2)$  是圆锥曲线。

\statementlabel{证明}我们将应用命题3.7.2。

(i) 我们有同构：

\[
\mu^{-1}f^{-1}F\simeq\tilde{f}^{-1}\mu^{-1}F\simeq\tilde{f}^{-1}p_{X}^{-1}F\simeq p_{Y}^{-1}f^{-1}F.
\]

因此 $f^{-1}F$  是圆锥曲线。  $f^!$  的证明是类似的。

(ii) 我们有同构：

\[
\mu^{\downarrow}R f_{*}G\simeq R\tilde{f}_{*}\mu^{\downarrow}G\simeq R\tilde{f}_{*}p_{Y}^{\downarrow}G\simeq p_{X}^{\downarrow}R f_{*}G\enspace.
\]

因此 $ Rf_{*} $  G 是圆锥曲线。  $ Rf_{1} $  的证明是类似的。

（三）我们有：

\[
\begin{aligned}{\mu^{-1}\bigg(F_{1}\otimes\stackrel{L}{F_{2}}\bigg)}&{{}\simeq\mu^{-1}F_{1}\otimes\stackrel{L}{\mu^{-1}F_{2}}}\\ {}&{{}\simeq p^{-1}F_{1}\otimes\stackrel{L}{p^{-1}F_{2}}}\\ {}&{{}\simeq p^{-1}\bigg(F_{1}\otimes\stackrel{L}{F_{2}}\bigg)~,}\\ \end{aligned}
\]

\[
\begin{aligned}{\mu^{!}R\mathcal{H}o m(F_{1},F_{2})}&{{}\simeq R\mathcal{H}o m(\mu^{-1}F_{1},\mu^{!}F_{2})}\\ {}&{{}\simeq R\mathcal{H}o m(p^{-1}F_{1},p^{!}F_{2})}\\ {}&{{}\simeq p^{!}R\mathcal{H}o m(F_{1},F_{2})\enspace.\quad\Box}\\ \end{aligned}
\]

令  $ \tau: E \to Z $  为局部紧空间 Z 上纤维维数为 n 的实数向量丛。

我们将  $Z$  与  $E$  的零部分识别，并用  $i: Z \hookrightarrow E$  表示此嵌入。我们设定：

\[
\dot{E}=E\backslash Z~,\qquad\dot{\tau}=\tau|_{\dot{E}}~.
\]

我们用“a”表示 $ E $ 上的对映图 $ x \mapsto -x $ ，如果 $ A $ 是 $ E $ 的子集，我们用 $ A^a $ 表示其图像 $ a $ 。如果对于任何  $ z \in Z $  ，集合  $ \tau^{-1}(z) \cap A $  是凸的（分别是圆锥曲线，分别是真圆锥），我们就说  $ E $  的子集  $ A $  是凸的（分别是圆锥曲线，分别是真圆锥）。回想一下，如果圆锥体不包含直线，那么它就是正确的。

\statementlabel{命题 3.7.5}让 $ F \in \mathrm{Ob}(\mathbf{D}_{R^+}^{+}(E)) $  。然后：

（一） $ R\tau_{*}F \simeq i^{-1}F $ 。

（二） $ R\tau_!F \simeq i^!F $ 。

\statementlabel{证明}(i) 态射  $ \tau^{-1}R\tau_*, F \to F $  定义态射  $ R\tau_*F \simeq i^{-1}\tau^{-1}R\tau_*F \to i^{-1}F $  。为了检查这是一个同构，让我们采用  $ z \in Z $  ，并让

U 是  $ i(z) $  的凸开邻域。然后由推论3.7.3：

\[
\begin{aligned}{R\varGamma(U;F)}&{{}\simeq R\varGamma(\mathbb{R}^{+}U;F)}\\ {}&{{}\simeq R\varGamma(\tau(U);R\tau_{*}F).}\\ \end{aligned}
\]

取两边的上同调，以及 U 穿过  $i(z)$  开凸邻域族时的归纳极限，我们得到结果。

(ii) 态射 $ i_1\, i^1F \to F $  定义：

\[
i^{!}F\simeq R\tau_{!}i_{!}i^{!}F\to R\tau_{!}F\enspace.
\]

为了证明这是同构，我们如(i)中那样论证。令 $U$  为 $Z$  的开子集 $V$  的凸开邻域，使得 $\tau^{-1}V\cap\overline{U}\to V$  是正确的。然后由推论3.7.3：

\[
R\varGamma_{z}(\tau^{-1}V;F)\simeq R\varGamma_{\overline{{U}}}(\tau^{-1}V;F)\enspace.
\]

取两边的上同调，以及  $U$  范围穿过满足上述性质的  $\tau^{-1}V$  子集族时的归纳极限，我们得到  $R\Gamma(V;i^{\dagger}F)\simeq R\Gamma(V;R\tau_{\dagger}F)$  ，它显示  $i^{\dagger}F\simeq R\tau_{\dagger}F$  。   $\square$ 

令 $ \pi: E^* \to Z $  为对偶向量丛。我们类似地定义  $ i: Z \hookrightarrow E^* $  、  $ \dot{E}^* $  和  $ \dot{\pi} $  。

如果  $A$  是  $E$  的子集，则极集  $A^{\circ}$  定义为：

(3.7.6) $ A^{\circ} = \{y \in E^*; \pi(y) \in \tau(A) \text{ and } \langle x, y \rangle \geq 0 \text{ for all } x \in \tau^{-1} \pi(y) \cap A\} $ 

我们用  $ p_1 $  和  $ p_2 $  表示  $ E \times_Z E^* $  的第一个和第二个投影：

\begin{center}
\includegraphics[width=0.31\textwidth]{流行上的层_Chn-assets/img_in_image_box_386_1109_767_1417.png}
\end{center}

我们介绍：

\[
P=\left\{(x,y)\in E\underset{z}{\times}E^{*};\langle x,y\rangle\geqslant0\right\},
\]

\[
P^{\prime}=\left\{(x,y)\in E\underset{z}{\times}E^{*};\langle x,y\rangle\leqslant0\right\}
\]

和函子（参见 III \S 6）：

\[
\left\{\begin{array}{l}{\tilde{\Psi}_{P^{\prime}}=R p_{1*}\circ R\varGamma_{P^{\prime}}\circ p_{2}^{1}~,}\\ {}\\ {\tilde{\Phi}_{P^{\prime}}=R p_{2!}\circ(\cdot)_{P^{\prime}}\circ p_{1}^{-1}~,}\\ {}\\ {\tilde{\Psi}_{P}=R p_{2*}\circ R\varGamma_{P}\circ p_{1}^{-1}~,}\\ {}\\ {\tilde{\Phi}_{P}=R p_{1!}\circ(\cdot)_{P}\circ p_{2}^{!}~.}\end{array}\right.
\]

从命题 3.7.4 可以看出，这些函子从  $ \mathbf{D}_{R^+}^{+}(E) $  到  $ \mathbf{D}_{R^+}^{+}(E^*) $  或从  $ \mathbf{D}_{R^+}^{+}(E^*) $  到  $ \mathbf{D}_{R^+}^{+}(E) $  是明确定义的。

令  $ \sigma_{E/Z} $  为  $ E $  相对于  $ Z $  的相对方向束。由于该束在  $ \tau $  的纤维上是恒定的，因此我们有时会识别  $ \sigma_{E/Z} $  及其对  $ Z $  的限制。请注意  $ \sigma_{E/Z}|_{Z} \simeq \sigma_{Z/E} $ （其中  $ Z \to E $  由零嵌入定义）。此外， $ \sigma_{Z/E} $  和  $ \sigma_{Z/E^*} $  自然是同构的，因为向量空间上的方向定义了对偶空间上的方向（参见下面的备注 3.7.11）。

为了研究（3.7.7）中的函子，我们需要一个引理。

\statementlabel{引理 3.7.6}让 $ F \in \mathrm{Ob}(\mathbb{D}_{\mathbb{R}^+}^+)(\mathcal{E}) $  。然后  $ \mathrm{supp}((R\Gamma_P(p_1^{-1}F))_{p^r}) $  包含在  $ Z \times_Z E^* \subset E \times_Z E^* $  中。

\statementlabel{证明}设置 $U = E \times_Z E^* \setminus Z \times_Z E^*$ 。在  $U$  上，投影  $p_1$  与投影  $\mathbb{R}^n \times E \to E$  同构，并且集合  $P$  和  $P'$  分别与集合  $\{x_1 \geqslant 0\}$  和  $\{x_1 \leqslant 0\}$  局部同构，其中  $x_1$  表示  $\mathbb{R}^n$  上的第一个坐标。因此，为了证明  $(R\Gamma_P(p_1^{-1}F))_{P'}$  在  $U$  上为零，只需证明如果  $X$  是局部紧空间，如果  $p$  表示投影  $X \times \mathbb{R} \to X$  并且  $t$  是  $\mathbb{R}$  上的坐标，则：

\[
(R\varGamma_{\{\imath\geqslant0\}}(p^{-1}G))_{\{\imath=0\}}=0,
\]

对于任何 $ G \in \text{Ob}(\mathbb{D}^+(X)) $  。由于  $ p^{-1}G $  在  $ X $  上的向量束  $ X \times \mathbb{R} $  上是圆锥曲线，因此我们有：

\[
(R\varGamma_{\{t\geqslant0\}}(p^{-1}G))|_{t=0}\simeq R p_{\ast}R\varGamma_{\{t\geqslant0\}}(p^{-1}G)\enspace.
\]

为了看到它消失，只需显示相同的语句，并将  $ p^{-1}G $  替换为  $ p^l G $  即可。然后 $ R p_* R \Gamma_{\{t \geq 0\}} (p^l G) \simeq R \mathcal{H} o m(R p_! A_{\{t \geq 0\}}, G) $  和 $ R p_! A_{\{t \geq 0\}} = 0 $  。   $ \square $ 

\statementlabel{定理 3.7.7}从  $ \mathbf{D}_{\mathbf{R}^{+}}^{+}(E) $  到  $ \mathbf{D}_{\mathbf{R}^{+}}^{+}(E^{*}) $  、 $ \widetilde{\Phi}_{P'} $  和  $ \widetilde{\Psi}_{P} $  的两个函子自然是同构的。

\statementlabel{证明}让 $ F \in \mathrm{Ob}(\mathbf{D}_{R}^{+}(E)) $  。我们有同构：

\[
\begin{aligned}{\tilde{\varPhi}_{p^{\prime}}(F)}&{{}=R p_{2!}(p_{1}^{-1}F)_{p^{\prime}}}\\ {}&{{}\simeq R p_{2!}R\varGamma_{p}((p_{1}^{-1}F)_{p^{\prime}})}\\ {}&{{}\simeq R p_{2}!((R\varGamma_{p}(p_{1}^{-1}F))_{p^{\prime}})}\\ {}&{{}\simeq R p_{2*}((R\varGamma_{p}(p_{1}^{-1}F))_{p^{\prime}})}\\ {}&{{}\simeq R p_{2*}R\varGamma_{p}(p_{1}^{-1}F).}\\ \end{aligned}
\]

事实上，第一个同构来自命题 3.7.5 (ii)，第二个同构来自练习 II.2，第三个同构来自引理 3.7.6，最后一个同构来自命题 3.7.5 (i)。   $ \square $ 

这意味着从  $ \mathbf{D}_{\mathbf{R}^{+}}^{+}(E^{*}) $  到  $ \mathbf{D}_{\mathbf{R}^{+}}^{+}(E) $  、 $ \widetilde{\Phi}_{P} $  和  $ \widetilde{\Psi}_{P} $  的两个函子也是同构的。

\statementlabel{定义 3.7.8}让 $ F \in \mathrm{Ob}(\mathbf{D}_{R^+}^{+}(E)) $  。我们设定：

\[
\begin{aligned}{}&{{}F^{\wedge}=\widetilde{\Phi}_{P^{\prime}}(F)(=R p_{2!}(p_{1}^{-1}F)_{P^{\prime}})}\\ {}&{{}(\simeq\widetilde{\Psi}_{P}(F)=R p_{2*}R\varGamma_{P}(p_{1}^{-1}F))}\\ \end{aligned}
\]

并调用  $ F^{\wedge} $  F 的 Fourier-Sato 变换。

让 $ G \in \mathrm{Ob}(\mathbf{D}_{\mathbf{R}^+}^{+}(E^*)) $  。我们设定：

\[
\begin{aligned}{G^{\vee}}&{{}=\tilde{\Psi}_{P^{\prime}}(G)(=R p_{1*}{R}\varGamma_{P^{\prime}}(p_{2}^{!}G))}\\ {}&{{}\left(\simeq\tilde{\Phi}_{P}(G)=R p_{1!}(p_{2}^{!}G)_{P}\right)}\\ \end{aligned}
\]

并调用  $ G^{\vee} $  G 的逆傅立叶-佐藤变换。

当然，可以使用相同的公式通过交换 E 和  $ E^{*} $  来定义  $ F^{\vee} $  和  $ G^{\wedge} $  。请注意：

\[
F^{\vee}\simeq(F^{\wedge})^{a}\otimes\mathfrak{o t}_{E^{*}/Z}[n]\simeq(F^{\wedge})^{a}\otimes\omega_{E^{*}/Z}~,
\]

其中  $ (F^\wedge)^a $  是  $ E^* $  上对映图的  $ F^\wedge $  的反像。

事实上，(3.7.9) 是从  $ p_1^! F \simeq p_1^{-1} F \otimes \omega_{E \times_Z E^* / E} $  以来的定义得出的。

\statementlabel{定理 3.7.9}函子  $ ^ from \mathbf{D}_{R^+}^+(E) $  到  $ \mathbf{D}_{R^+}^+(E^*) $  和  $ ^ from \mathbf{D}_{R^+}^+(E^*) $  到  $ \mathbf{D}_{R^+}^+(E) $  是范畴的等价，彼此相反。特别是如果  $ F $  和  $ F' $  属于  $ \mathbf{D}_{R^+}^+(E) $  ，则：

\[
\mathrm{H o m}_{{\tt D}_{{\tt R}^{+}}(E)}(F^{\prime},F)\simeq\mathrm{H o m}_{{\tt D}_{{\tt R}^{+}}(E^{*})}(F^{\prime}{}^{\wedge},F^{\wedge})\enspace.
\]

\statementlabel{证明}根据命题 3.6.2 和定理 3.7.7，我们知道在选择同构  $ o\tau_{Z/E} \simeq o\tau_{Z/E} $  后，  $ \tilde{\Phi}_P $  和  $ \tilde{\Psi}_P $  是伴随函子， $ \tilde{\Phi}_P $  和  $ \tilde{\Psi}_P $  是伴随函子。因此，如果  $ F \in \mathrm{Ob}(\mathbf{D}_{\mathbf{R}^+}^{+}(E)) $  ，我们得到态射：

\[
F\to F^{\wedge\vee}~.
\]

为了证明 (3.7.10) 是同构，只需证明对于 E 的每个开凸子集 U，它都会产生同构：

\[
H^{j}(U;F)\simeq H^{j}(U;F^{\wedge v})~.
\]

应用推论 3.7.3，我们可以假设 U 是凸开锥体。在这种情况下我们有：

\[
\begin{aligned}{H^{j}(U;F^{\wedge\vee})}&{{}=\operatorname{H o m}_{\mathsf{D}_{\mathbb{R}^{+}}(E)}(A_{U},\widetilde{\Psi}_{P^{\prime}}\widetilde{\Phi}_{P^{\prime}}(F)[j])}\\ {}&{{}\simeq\operatorname{H o m}_{\mathsf{D}_{\mathbb{R}^{+}}(E^{*})}(\widetilde{\Phi}_{P^{\prime}}(A_{U}),\widetilde{\Phi}_{P^{\prime}}(F)[j])}\\ {}&{{}\simeq\operatorname{H o m}_{\mathsf{D}_{\mathbb{R}^{+}}(E^{*})}(\widetilde{\Phi}_{P^{\prime}}(A_{U}),\widetilde{\Psi}_{P}(F)[j])}\\ {}&{{}\simeq\operatorname{H o m}_{\mathsf{D}_{\mathbb{R}^{+}}(E)}(\widetilde{\Phi}_{P}\widetilde{\Phi}_{P^{\prime}}(A_{U}),F[j])}\\ {}&{{}\xleftarrow{\alpha}\operatorname{H o m}_{\mathsf{D}_{R^{+}}(E)}(A_{U},F[j]),}\\ \end{aligned}
\]

和  $\alpha$  来自态射  $\tilde{\Phi}_P\tilde{\Phi}_{P'}(A_U) \simeq \tilde{\Phi}_P\tilde{\Psi}_P(A_U) \to A_U$  。因此，足以证明  $\tilde{\Phi}_P\tilde{\Psi}_P(A_U)$  与  $A_U$  同构。这是在下面的引理 3.7.10 中完成的。类似地，我们可以证明对于  $G \in \mathrm{Ob}(\mathbf{D}_{\mathbf{R}^+}^+)(\mathcal{E}^*)$  ，  $G^{\vee^\wedge} \to G$  是同构。   $\square$ 

\statementlabel{引理 3.7.10}(i) 令 $ \gamma $  为包含零截面的E 的真闭凸锥。然后：

\[
(A_{\gamma})^{\wedge}\simeq A_{\mathtt{I n t}\gamma^{\circ}}.
\]

(ii) 令 U 为 E 的凸开锥体。则：

\[
(A_{U})^{\wedge}\simeq A_{U^{\circ}\alpha}\otimes\mathfrak{o r}_{E^{\bullet}/Z}[-n]~.
\]

\statementlabel{证明}(i) 让 $ y \in E^* $  。然后：

\[
\begin{aligned}{((A_{\gamma})^{\wedge})_{y}}&{{}\simeq R\varGamma_{c}(p_{2}^{-1}(y);A_{(\gamma\times_{Z}E^{*})\cap P^{\prime}})}\\ {}&{{}\simeq R\varGamma_{c}(p_{1}(p_{2}^{-1}(y)\cap P^{\prime})\cap\gamma;A_{E}).}\\ \end{aligned}
\]

设置 $\gamma_y = p_1(p_2^{-1}(y) \cap P') \cap \gamma$ 。如果 $y \notin \mathrm{Int} \gamma^\circ$  ， $\gamma_y$  是一个包含半线的闭凸真圆锥。那么 $\tau^{-1}(\pi(y))$  和 $\tau^{-1}(\pi(y)) \setminus \gamma_y$  是同胚的，在这种情况下我们发现 $((\mathcal{A}_y)^\wedge)_y = 0$ 。如果  $y \in \mathrm{Int} \gamma^\circ$  ，则  $\gamma_y = \{0\}$  和映射  $((A_y)^\wedge)_y \to ((\mathcal{A}_Z)^\wedge)_y$  是同构。这表明  $(A_y)^\wedge \simeq ((A_Z)^\wedge)_{\mathrm{Int}\gamma^\circ}$  和 (i) 自  $(A_Z)^\wedge \simeq A_{E^*}$  开始。

(ii) 证明类似。我们有：

\[
\begin{aligned}{((A_{U})^{\wedge})_{y}}&{{}\simeq R\varGamma_{c}(p_{2}^{-1}(y);A_{(U\times E^{*})\cap P^{\prime}})}\\ {}&{{}\simeq R\varGamma_{c}(p_{1}(p_{2}^{-1}(y)\cap P^{\prime})\cap U;A_{E}).}\\ \end{aligned}
\]

设置 $ \gamma_y = p_1(p_2^{-1}(y) \cap P') \cap U $ 。如果  $ y \notin U^{\circ a} $  则  $ \gamma_y = \emptyset $  和  $ ((A_U)^\wedge)_y = 0 $  。如果  $ y \in U^{\circ a} $  ，  $ ((A_U)^\wedge)_y \to (Rp_{2!}A_{E \times_2E^*})_y $  是同构。

因此 $ (A_U)^\wedge \simeq (R p_{2!} A_{E \times_Z E^*})_{U^\circ a} $  。然后结果来自同构  $ R p_{2!} A_{E \times_Z E^*} \simeq \sigma i_{E^*/Z}[-n] $  。   $ \square $ 

\statementlabel{备注 3.7.11}由于  $ \tilde{\Phi}_{P} $  和  $ \tilde{\psi}_{P} $  以及  $ \tilde{\Phi}_{P} $  和  $ \tilde{\psi}_{P} $  彼此伴随，因此我们有态射：

\[
\alpha^{\prime}(F)\colon F\to\widetilde{\Psi}_{P^{\prime}}\widetilde{\Phi}_{P^{\prime}}(F)\enspace,
\]

\[
\beta^{\prime}(G):\widetilde{\Phi}_{P^{\prime}}\widetilde{\Psi}_{P^{\prime}}(G)\to G\enspace,
\]

\[
\alpha(G):G\to\widetilde{\Psi}_{P}\widetilde{\Phi}_{P}(G)\enspace,
\]

\[
\beta(F)\colon\widetilde{\Phi}_{P}\widetilde{\Psi}_{P}(F)\to F\enspace.
\]

在这里，为了定义 $ \alpha(G) $ 和 $ \beta(F) $ ，我们需要识别 $ R_{\tau_*} \omega_{E/Z} $ 和 $ R\pi_* \omega_{E^*/Z} $ 。我们采取以下识别（注意问题是本地的）。

选择  $E$  上的任何负定对称形式。相关的同构  $E \simeq E^*$  产生  $R\tau_* \omega_{E/Z}$  和  $R\pi_* \omega_{E^*/Z}$  之间的同构。

然后，在识别了 $ \tilde{\Phi}_P $ 、 $ \tilde{\Psi}_P $ 以及 $ \tilde{\Phi}_P $ 和 $ \tilde{\Psi}_P $ （定理3.7.7）之后，我们可以证明 $ \alpha'(F) $ （分别 $ \alpha(G) $ ）和 $ \beta(F) $ （分别 $ \beta'(G) $ ）彼此相反。但我们不会在这里给出证据。

让我们总结一下傅立叶-佐藤变换的一些性质。

\statementlabel{命题 3.7.12}让 $ F \in \mathrm{Ob}(\mathbf{D}_{R^+}^{+}(E)) $  。

（一） $ F^{\wedge\wedge} \simeq F^a \otimes \sigma r_{E/Z}[-n] $ 。

(ii) 令  $U$  为  $E^{*}$  的凸开子集。然后：

\[
R\varGamma(U;F^{\wedge})\simeq R\varGamma_{U^{\circ}}(\tau^{-1}\pi(U);F)\simeq R\varGamma_{U^{\circ}}(E;F)\enspace.
\]

(iii) 令 $ \gamma $  为包含零截面的 $ E^{*} $  的闭凸真圆锥。然后：

\[
R\varGamma_{\gamma}(E^{*};F^{\wedge})\simeq R\varGamma(\mathtt{I n t}\gamma^{\circ a};F)\otimes\mathfrak{o r}_{E/\mathbb{Z}}[-n]~.
\]

（四）我们有：

\[
({\bf D}^{\prime}F)^{\vee}\simeq{\bf D}^{\prime}(F^{\wedge}),\qquad({\bf D}F)^{\vee}\simeq{\bf D}(F^{\wedge}).
\]

\statementlabel{证明}(i)、(ii)、(iii) 紧接着前面的结果。

（四）我们有：

\[
\begin{aligned}{{R\mathcal{H}o m}(F^{\wedge},A_{E^{*}})}&{{}={R\mathcal{H}o m}(R p_{2}!(p_{1}^{-1}F)_{P^{\prime}},A_{E^{*}})}\\ {}&{{}\simeq R p_{2*}R\mathcal{H}o m((p_{1}^{-1}F)_{P^{\prime}},p_{2}^{!}A_{E^{*}})}\\ {}&{{}\simeq R p_{2*}R\mathcal{H}o m((p_{1}^{-1}F)_{P^{\prime}},p_{1}^{-1}\omega_{E/Z})}\\ {}&{{}\simeq R p_{2*}R\varGamma_{P^{\prime}}(p_{1}^{!}(\mathsf{D}^{\prime}F))}\\ {}&{{}\simeq(\mathsf{D}^{\prime}F)^{\vee}.}\\ \end{aligned}
\]

D 的证明类似。 ⑨

现在我们将研究 Fourier-Sato 变换的一些函数性质。

令 $ Z' $  为另一个局部紧空间， $ f: Z' \to Z $  为连续映射。我们设置 $ E' = Z' \times_Z E $ 并用 $ f_\tau $ （分别 $ f_\pi $ ）表示与 $ f $ 关联的从 $ E' $ 到 $ E $ （分别 $ E'^* $ 到 $ E^* $ ）的映射。

\statementlabel{命题 3.7.13}(i) 让 $ F \in \mathrm{Ob}(\mathbf{D}_{R^+}^{+}(E)) $  。然后：

\[
(f_{\tau}^{!}F)^{\wedge}\simeq f_{\pi}^{!}(F^{\wedge}),
\]

\[
(f_{\tau}^{-1}F)^{\wedge}\simeq f_{\pi}^{-1}(F^{\wedge})
\]

(ii) 让 $ G \in \mathrm{Ob}(\mathbf{D}_{\mathbf{R}^+}^{+}(E)) $  。然后：

\[
\begin{aligned}{({R f}_{\tau*}{G})^{\wedge}}&{{}\simeq{R f}_{\pi*}({G}^{\wedge})~,}\\ {({R f}_{\tau!}{G})^{\wedge}}&{{}\simeq{R f}_{\pi!}({G}^{\wedge})~.}\\ \end{aligned}
\]

\statementlabel{证明}考虑下图，其中正方形是笛卡尔坐标：

\[
\begin{array}{c c c c}{E^{\prime*}}&{\xrightarrow{f_{*}}}&{E^{*}}&{}\\ {q_{2}^{\prime}\Bigg\downarrow}&{\Box}&{\Bigg\uparrow_{q_{2}}}&{}\\ {Z^{\prime}\times P}&{\xrightarrow{\tilde{f}}}&{P}&{}\\ {q_{1}^{\prime}\Bigg\downarrow}&{\Box}&{\Bigg\downarrow_{q_{1}}}&{}\\ {E^{\prime}}&{\xrightarrow{f_{*}}}&{E}&{.}\end{array}
\]

然后：

\[
\begin{aligned}{f_{\pi}^{!}F^{\wedge}\otimes\omega_{E^{*}/Z}=}&{{}f_{\pi}^{!}R q_{2*}q_{1}^{!}F}\\ {\simeq}&{{}R q_{2*}^{\prime}\tilde{f}^{!}q_{1}^{!}F}\\ {\simeq}&{{}R q_{2*}^{\prime}q_{1}^{\prime!}f_{\tau}^{!}F}\\ {\simeq}&{{}(f_{\tau}^{!}F)^{\wedge}\otimes\omega_{E^{*}/Z}~.}\\ \end{aligned}
\]

这显示了公式  $ f_{\pi}^!(F^\wedge) \simeq (f_\tau^! F)^\wedge $  。其他公式的证明类似。   $ \square $ 

现在让  $ E_1 $  和  $ E_2 $  成为  $ Z $  上的两个向量丛， $ f: E_1 \to E_2 $  是此类丛的态射。我们用 $ t'f: E_2^* \to E_1^* $  表示对偶态射。然后 $ \omega_{E_1/E_2} \simeq f^!A_{E_2} $  和 $ \omega_{E_2^t/E_1^t} \simeq t'f^!A_{E_1^*} $  。

\statementlabel{命题 3.7.14}(i) 让 $ F \in \mathrm{Ob}(\mathbb{D}_{R^+}^{+}(E_1)) $  。然后：

\[
\begin{aligned}{}&{{}{}^{t}f^{-1}(F^{\wedge})\simeq(R f_{!}F)^{\wedge},}\\ {}&{{}\quad{}^{t}f^{!}(F^{\vee})\simeq(R f_{*}F)^{\vee},}\\ {}&{{}\quad{}^{t}f^{!}(F^{\wedge})\simeq(R f_{*}F)^{\wedge}\otimes\omega_{E_{2}^{\pm}/E_{1}^{\pm}},}\\ {}&{{}{}^{t}f^{-1}(F^{\vee})\simeq(R f_{!}F)^{\vee}\otimes\omega_{E_{2}^{\pm}/E_{1}^{\pm}}.}\\ \end{aligned}
\]

(ii) 让 $ G \in \mathrm{Ob}(\mathbf{D}_{\mathbf{R}^{+}}^{+}(E_{2})) $  。然后：

\[
(f^{-1}G)^{\vee}\simeq R^{\mathfrak{r}}f_{!}(G^{\vee}),
\]

\[
(f^{!}G)^{\wedge}\simeq R^{\mathfrak{t}}f_{*}(G^{\wedge}),
\]

\[
(\omega_{E_{1}/E_{2}}\otimes f^!G)^{\vee}\simeq R^{t}f_{*}(G^{\vee}),
\]

\[
(\omega_{E_{1}/E_{2}}\otimes f^{-1}G)^{\wedge}\simeq R^{\mathfrak{t}}f_{!}(G^{\wedge})\enspace.
\]

\statementlabel{证明}设置  $ \widetilde{P}' = E_1 \times_{E_2} P'_2 = P'_1 \times_{E_1^*} E_2^* = \{(x, y) \in E_1 \times_Z E_2^*; \langle x, !f(y) \rangle = \langle f(x), y \rangle \leqslant 0\} $  ，并考虑下图，其中两个正方形是笛卡尔坐标：

\begin{center}
\includegraphics[width=0.30\textwidth]{流行上的层_Chn-assets/img_in_image_box_407_578_769_923.png}
\end{center}

然后：

\[
\begin{aligned}{{}^{t}f^{-1}F^{\wedge}}&{{}\simeq{}^{t}f^{-1}R q_{2!}q_{1}^{-1}F}\\ {}&{{}\simeq R q_{2!}^{\prime}R\tilde{q}_{2!}\tilde{q}_{1}^{-1}q_{1}^{-1}F}\\ {}&{{}\simeq R q_{2!}^{\prime}q_{1}^{\prime-1}R f_{!}F}\\ {}&{{}\simeq(R f_{!}F)^{\wedge}.}\\ \end{aligned}
\]

同样：

\[
\begin{aligned}{{}^{t}f^{!}F^{\vee}}&{{}\simeq{}^{t}f^{!}R q_{2*}q_{1}^{!}F}\\ {}&{{}\simeq R q_{2*}R\tilde{q}_{2*}\tilde{q}_{1}^{!}q_{1}^{!}F}\\ {}&{{}\simeq R q_{2*}q_{1}^{^{\prime}!}R f_{*}F}\\ {}&{{}\simeq(R f_{*}F)^{\vee}.}\\ \end{aligned}
\]

其他公式后面设置  $ F = G^{\vee} $  或  $ F = G^{\wedge} $  。 □

最后我们研究外张量积。令  $ E_1 $  和  $ E_2 $  为 Z 上的两个向量丛。我们用相同的符号  $ ^, $  表示  $ E_i $  (  $ i = 1,2 $  ) 或  $ E_1 \times_Z E_2 $  上的傅里叶佐藤变换。

\statementlabel{命题 3.7.15}让 $ F_i \in \mathrm{Ob}(\mathbf{D}_R^+(E_i)) $ ， $ i = 1, 2 $ 。然后：

\[
F_{1}^{\wedge}\mathop{\boxtimes}_{Z}^{L}F_{2}^{\wedge}\simeq\left(F_{1}\mathop{\boxtimes}_{Z}^{L}F_{2}\right)^{\wedge}~.
\]

\statementlabel{证明}对于 $i=1,2$ ， $p_{i}^{j}(j=1,2)$  和 $p_{i}$  分别表示在 $E_{j}\times_{Z}E_{j}^{*}(j=1,2)$  和 $(E_{1}\times_{Z}E_{2})\times_{Z}(E_{1}^{*}\times_{Z}E_{2}^{*})$  上定义的 $i$  投影。令  $P_{j}^{\prime}(j=1,2)$  和  $P^{\prime}$  表示  $E_{j}\times_{Z}E_{j}^{*}$  (  $j=1,2$  ) 和  $(E_{1}\times_{Z}E_{2})\times_{Z}(E_{1}^{*}\times_{Z}E_{2}^{*})$  的闭子集  $\{\langle x,y\rangle\leqslant0\}$  。我们有（参见练习 II.18）：

\[
F_{1}^{\wedge}\mathop{\boxtimes}_{Z}^{L}F_{2}^{\wedge}\simeq R p_{2!}\bigg((p_{1}^{1})^{-1}F_{1}\mathop{\boxtimes}_{Z}^{L}(p_{1}^{2})^{-1}F_{2}\bigg)_{P_{1}^{\prime}\times_{Z}P_{2}^{\prime}},
\]

\[
\bigg(\begin{matrix}{L}\\ {F_{1}\boxtimes}\\ {z}\\ \end{matrix}F_{2}\bigg)^{\wedge}\simeq R p_{2!}\bigg(\big(p_{1}^{1}\big)^{-1}F_{1}\boxtimes\begin{matrix}{L}\\ {z}\\ \end{matrix}\big(p_{1}^{2}\big)^{-1}F_{2}\bigg)_{P^{\prime}}~.
\]

我们设置 $ G = (p_1^1)^{-1} F_1 \boxtimes_Z^L (p_1^2)^{-1} F_2 $ 。映射 $ p_2 $  分解如下：

\[
\left(E_{1}\;{\times}\;E_{2}\right)\;{\times}\left(E_{1}^{*}\;{\times}\;E_{2}^{*}\right)\underset{\beta}{\longrightarrow}E_{1}^{*}\;{\times}\;E_{2}^{*}\;{\times}\;\mathbb{R}\;{\times}\;\mathbb{R}\underset{\alpha}{\longrightarrow}E_{1}^{*}\;{\times}\;E_{2}^{*}\;,
\]

其中  $ \beta(x_1, x_2, y_1, y_2) = (y_1, y_2, \langle x_1, y_1 \rangle, \langle x_2, y_2 \rangle) $  和  $ \alpha(y_1, y_2, t_1, t_2) = (y_1, y_2) $  。自  $ R\beta_1 G_{P'} \simeq (R\beta_1 G)_{\{t_1 + t_2 \leq 0\}} $  和  $ R\beta_1 G_{P_1 \times_2 P_2} \simeq (R\beta_1 G)_{\{t_1 \leq 0, t_2 \leq 0\}} $  以来，仍有待证明

\[
R\alpha_{!}((R\beta_{!}G)_{\{t_{1}+t_{2}\leq0\}})\to R\alpha_{!}((R\beta_{!}G)_{\{t_{1}\leq0,t_{2}\leq0\}})
\]

是同构。

我们将在  $ E_1^* \times_Z E_2^* $  的每一点证明这一点。因此，我们必须证明，如果  $ H $  是  $ \mathbb{D}^+ $  (  $ \mathbb{R} \times \mathbb{R} $  ) 的双圆锥对象，那么：

\[
R\varGamma_{c}(\mathbb{R}\times\mathbb{R};H_{\{t_{1}+t_{2}\leqslant0\}})\xrightarrow{\sim}R\varGamma_{c}(\mathbb{R}\times\mathbb{R};H_{\{t_{1}\leqslant0,t_{2}\leqslant0\}})~,
\]

或等价：

\[
R\varGamma_{c}(\mathbb{R}\times\mathbb{R};H_{\{0<t_{1}\leqslant t_{2}\}\cup\{0<t_{2}\leqslant-t_{1}\}})=0\enspace.
\]

那么这是因为  $H$  的上同调轮在开集  $\{t_1 t_2 \neq 0\}$  的每个连通分量上都是恒定的。   $\Box$ 

\statementlabel{备注 3.7.16}在本书的过程中，我们忽略了符号的问题。例如，标识 $ \sigma t_{Y/X} $ 和 $ \sigma t_{Y} \otimes f^{-1} \sigma t_{X} $ （对于态射 $ f: Y \to X $ ），或者矢量丛 $ E \to Z $ 的标识 $ \sigma t_{Z/E} \simeq \sigma t_{Z/E^*} $ （以及，作为特殊情况，标识 $ \sigma t_{T^*X} \simeq A_{T^*X} $ ）一般不会被描述。

\subsection{习题}
\statementlabel{练习 III.1}我们采用 $ A = \mathbb{Q} $ 。设置 $ X = \mathbb{R} $ 、 $ B = \{1/n; n \in \mathbb{N} \backslash \{0\}\} $ 、 $ Z = B \cup \{0\} $ 。

(i) 让 $ F = \mathbb{Q}_B $  。计算  $ D'F $  和  $ D'D'F $  。显示 $ F \neq D'D'F $  。

(ii) 让 $ G = \mathbb{Q}_Z $  。证明  $ G $  是软的，并且  $ H_{\{0\}}^1(X; G) $  在  $ \mathbb{Q} $  上是无限维的。

（提示：表明  $ H^{0}(X; F) $  （或  $ H^{0}(X; G) $  ）与 Q 中的值的所有序列（或所有平稳序列）的空间同构。）

(iii) 证明 $ H_{\{0\}}^{2}(X;\mathbb{Q}_{(X\backslash Z)}) \neq 0 $  。

\statementlabel{练习 III.2}证明  $ \mathbb{R}^n $  的软维数是  $ n $  ， $ \mathbb{R}^n $  的松弛维数是  $ n + 1 $  。 （提示：使用练习 III.1。）

\statementlabel{练习 三.3}令 X 为拓扑空间，F 为 X 上的层，与  $ Z_{X} $  局部同构。证明存在规范同构：

\[
F\otimes F\simeq\mathbb{Z}_{x}~,\qquad\operatorname{D}^{\prime}F\simeq F~.
\]

\statementlabel{练习 三.4}令 $X$  成为 $C^0$  流形。我们可以说  $\Omega$  在  $X$  （简称 l.c.t.）中是局部上同调平凡的，如果对于任何  $x \in \tilde{\Omega} \backslash \Omega$  ，有：

\[
(R\varGamma_{\bar{\Omega}}(A_{X}))_{x}=0\quad,\qquad(R\varGamma_{\Omega}(A_{X}))_{x}\simeq A\quad,
\]

（参见 Schapira [3]）。

(i) 证明  $\Omega$  是  $X$  中的  $l.c.t.$  当且仅当  $D'(A_\Omega) \simeq A_{\bar{\Omega}}$  和  $D'(A_{\bar{\Omega}}) \simeq A_\Omega$  。

(ii) 证明如果 $\Omega$  是 $X$  中的 $l.c.t.$ ，则 $\Omega = \mathrm{Int}(\bar{\Omega})$  。

(iii) 证明如果 $\Omega$  在 $\mathbb{R}^n$  中是凸的，那么 $\Omega$  是 $\mathbb{R}^n$  中的 $l.c.t.$ ，并且 $A_\Omega$  和 $A_{\bar{\Omega}}$  是上同调可构造的。

\statementlabel{练习 三.5}令  $V$  为实数  $n$  维向量空间，并令  $q$  为  $V$  上的二次形式。对于 $a \in \mathbb{R}$ ，我们设置 $Z_a = \{x \in V; q(x) \leq a\}$ 。令  $\varepsilon^{-}$  为  $q$  的非正特征值的数量。证明：

（一） $ RF_{c}(V;A_{Z_{a}})[\varepsilon^{-}]\simeq\begin{cases}A&a\geq0\\0&a<0\end{cases} $ ，

(ii) 使用命题 3.1.10，从 (i) 推出：

\[
R\varGamma_{Z_{a}}(V;A_{V})[n-\varepsilon^{-}]\simeq\begin{cases}{A}&{a\geqslant0}\\ {0}&{a<0}\\ \end{cases},
\]

\statementlabel{练习 三.6}令  $f: Y \to X$  为纤维尺寸为  $n \geqslant 1$  的球丛。证明杰出三角形的存在性：

\[
R^{0}f_{*}A_{Y}\longrightarrow R f_{*}A_{Y}\longrightarrow R^{n}f_{*}A_{Y}[-n]\stackrel{\longrightarrow}{+1}~,
\]

这导致：

\[
A_{X}\longrightarrow R f_{*}A_{Y}\longrightarrow f_{*}o r_{Y/X}[-n]\stackrel{\longrightarrow}{+1}.
\]

（提示：使用练习 I.26。）

\statementlabel{练习 三.7}令  $ \tau: E \to Z $  为具有纤维维度  $ n $  的向量丛， $ i: Z \hookrightarrow E $  为零截面嵌入。我们假设  $ E $  具有相对方向（即我们假设被赋予同构  $ A_E \simeq \sigma t_{E/Z} $  ）。

(i) 构造同构  $ \Gamma(Z; A_Z) \simeq H_Z^n(E; A_E) $  。 （1 的图像称为 Thom 类。）

(ii) 构造交换图：

\[
\begin{array}{ccc}R\tau_{!}A_{E}&\xrightarrow{\quad}R\tau_{*}A_{E}\\\left|\imath\right.&&\left|\imath\right.\\&\left|\vphantom{\imath}\right.\\A_{Z}[-n]&\xrightarrow{\quad}A_{Z}\end{array}
\]

和态射  $ \Gamma(Z; A_Z) \to H^n(Z; A_Z) $  。 （1的图像称为欧拉类。）

(iii) 通过映射证明 Euler 类是 Thom 类的镜像： $ H_{Z}^{n}(E; A_{E}) \to H^{n}(E; A_{E}) \simeq H^{n}(Z; A_{Z}) $  。

(iv) 证明如果丛允许无处消失的连续部分，则欧拉类为零（参见练习 V.11）。

\statementlabel{练习 三.8}令 $X$  为拓扑流形， $Y$  为余维 $l$  和 $j: Y \hookrightarrow X$  包含的闭子流形。假设被赋予同构  $o\iota_{Y/X} \simeq A_Y$  。然后构造长精确序列：

\[
H^{k}(Y;A_{Y})\xrightarrow{\alpha}H^{k+l}(X;A_{X})\longrightarrow H^{k+l}(X\setminus Y;A_{X})\longrightarrow H^{k+1}(Y;A_{Y}).
\]

（映射 $\alpha$  称为吉辛映射。 $1 \in H^0(Y; A_Y)$  的图像再次称为欧拉类。）

\statementlabel{练习 三.9}我们保留命题3.1.9的假设。

(i) 证明函子和交换图存在自然态射：

\begin{center}
\includegraphics[width=0.60\textwidth]{流行上的层_Chn-assets/img_in_image_box_312_288_1036_1057.png}
\end{center}

(ii) 对于 $ G \in \mathrm{Ob}(\mathbb{D}^{+}(A_{Y})) $  、 $ F \in \mathrm{Ob}(\mathbb{D}^{+}(A_{X})) $  ，构造一个规范交换图：

\begin{center}
\includegraphics[width=0.45\textwidth]{流行上的层_Chn-assets/img_in_image_box_325_1164_867_1344.png}
\end{center}

(iii) 对于 $\mathbf{D}^b(A_x)$  中的 $F_1$  和 $F_2$ ，构造交换图：

\begin{center}
\includegraphics[width=0.79\textwidth]{流行上的层_Chn-assets/img_in_image_box_102_1428_1045_1745.png}
\end{center}

\statementlabel{练习 三.10}利用命题 3.3.9 的假设证明，对于任意  $ F \in \mathrm{Ob}(\mathbf{D}^{+}(A_X)) $  ，态射  $ R f_! f^! F \to F $  是同构。

\statementlabel{练习 三.11}令 k 为交换域，X 为紧致 n 维  $ C^0 $  流形，F 为  $ \mathbf{D}^b(k_X) $  的上同调构造对象，DF 为对偶。

(i) 证明 $ H^j(X; F) $  和 $ H^{-j}(X; DF) $  是有限维k 向量空间，互为对偶。 （请注意  $ H^{-j}(X; DF) = H^{n-j}(X; D'F \otimes o t_X) $  ）。 （提示：通过 j 的归纳证明，对于任何紧子集 K 和任何开子集  $ U \supset K $  ，  $ \operatorname{Im}(H^j(U; F) \to H^j(K; F)) $  是有限维的。）

(ii) 现在假设  $ k = \mathbb{R} $  ，  $ X $  是紧凑的，面向类  $ C^\infty $  。使用带有  $ C^\infty $  函数和分布作为系数的 de Rham 复形来描述  $ H^j(X; \mathbb{R}_X) $  和  $ H^{n-j}(X; \mathbb{R}_X) $  之间的对偶性（参见 II §9）。

\statementlabel{练习 三.12}令  $E$  为有限维实向量空间， $\gamma$  为顶点位于  $0, \gamma^{a} = -\gamma$  的封闭凸锥。

(i) 证明同构：

\[
A_{\gamma^{a}}\stackrel{\sim}{\to}\phi_{\gamma}^{-1}R\phi_{\gamma*}(A_{\gamma^{a}})\enspace.
\]

(ii) 令  $\Omega$  为  $E$  的  $\gamma$  开子集。证明同构：

\[
A_{\Omega}\simeq\phi_{\gamma}^{-1}R\phi_{\gamma*}A_{\Omega}\enspace.
\]

\statementlabel{练习 三.13}让  $ X = \mathbb{R}^2 $  坐标为  $ (x, y) $  、  $ A = \{(x, y); y = 0\} $  、  $ B = \left\{(x, y); x \neq 0, y = x \sin \frac{1}{x}\right\} $  。让 $ F = \mathbb{Q}_A $ ， $ G = \mathbb{Q}_{\widetilde{B}} $ 。证明  $ F $  和  $ G $  是上同调可构造的（在  $ \mathbb{Q} $  上），但  $ F \otimes G $  和  $ R \mathcal{H} o m(F, G) $  不是。

\statementlabel{练习 三.14}令  $ S^n $  为欧几里得空间  $ \mathbb{R}^n $  的单位球面，令  $ a \in \mathbb{R} $  和  $ -1 \leq a < 1 $  ，并设置：

\[
\Omega=\left\{(x,y)\in\mathbf{S}^{n}\times\mathbf{S}^{n};\left<x,y\right>>a\right\}\;,\qquad K=A_{\Omega}\;.
\]

证明函子  $ \psi_{\kappa} $  和  $ \Phi_{\kappa} $  （参见 III §6）是  $ \mathbf{D}^{+}(A_{\mathbf{S}^{n}}) $  上范畴的等价。

\statementlabel{练习 三.15}令 $V$  为赋予其方向的向量空间 $\mathbb{R}^n$ ，并令 $m$  为整数 $1 \leq m \leq n$  。设定：

 $ X_m = \{\lambda; \lambda \text{ is an oriented } m\text{-dimensional linear subspace of } V\} $ 。

 $ \Omega = \{(\lambda, \mu) \in X_m \times X_{n-m}; \lambda \text{ and } \mu \text{ are transversal and the orientation of } \lambda \oplus \mu \text{ is that of } V\}. $ 

\[
K=A_{\Omega}\;.
\]

证明  $ \Psi_K $  和  $ \Phi_K $  定义了  $ \mathbf{D}^+\left(A_{X_m}\right) $  和  $ \mathbf{D}^+\left(A_{X_{n-m}}\right) $  之间范畴的等价性。

\statementlabel{练习 三.16}令  $ (t,x)=(t,x_1,\ldots,x_n) $  为  $ \mathbb{R}^{1+n} $  上的坐标，并令  $ \gamma_i  (i=1,2,3) $  为封闭圆锥体：

\[
\gamma_{1}=\left\{(t,x);t^{2}\geqslant\sum_{i=1}^{n}x_{i}^{2}\right\},
\]

\[
\gamma_{2}=\left\{(t,x);t^{2}\leqslant\sum_{i=1}^{n}x_{i}^{2}\right\},
\]

\[
\gamma_{3}=\left\{(t,x);t^{2}=\sum_{i=1}^{n}x_{i}^{2}\right\}.
\]

设置 $ \gamma_i^+ = \gamma_i \cap \{(x, t); t \geq 0\} $ ，以及： $ G_i = A_{\gamma_i} $ ， $ G_i^+ = A_{\gamma_i^+} $ 。计算所有这些层的傅立叶-佐藤变换。

\statementlabel{练习 三.17}令  $ \tau: E \to Z $  为局部紧致空间  $ Z $  上的向量丛。我们使用与例 3.6.6 中相同的符号，并用相同的字母  $ \gamma $  表示投影  $ \dot{E} \to S $  和  $ \dot{E}^* \to S^* $  。

令 $j: \dot{E} \subsetneq E$  为包含图。证明图的交换律，其中  $\tilde{\psi}_D$  和  $\tilde{\Phi}_D$  的定义类似于 (3.7.7)：

\[
\begin{array}{c c c c c c c c c}{\mathbf{D}^{+}(S)}&{\xrightarrow{\quad}}&{\widetilde{\Psi}_{D}}&{\quad}&{\mathbf{D}^{+}(S^{*})}&{\xrightarrow{\quad}}&{\widetilde{\Phi}_{D}}&{\quad}&{\mathbf{D}^{+}(S)}\\ {\downarrow}&{}&{}&{\quad}&{\quad}&{}&{\quad}&{\quad}\\ {\gamma^{-1}}&{}&{}&{\quad}&{\quad}&{\quad}&{\quad}&{\quad}\\ {}&{}&{}&{\quad}&{\quad}&{\quad}&{\quad}&{\quad}\\ {}&{}&{}&{\quad}&{\quad}&{\quad}&{\quad}&{\quad}\\ \end{array}\qquad\begin{array}{c c c c c c c c}{\mathbf{D}^{+}(S^{*})}&{\xrightarrow{\quad}}&{}&{}&{\mathbf{D}^{+}(S)}\\ {R\gamma_{*}}&{}&{}&{\downarrow}&{\gamma^{-1}}&{}&{}&{}\\ {}&{}&{}&{\downarrow}&{}&{}&{}&{}\\ \end{array}\qquad\begin{array}{c c c c c}{\quad}\\ {\quad}\\ \end{array}R\gamma_{*}
\]

\[
\begin{array}{l l l}{{\mathbf{D}}_{\mathbf{R}^{+}}^{+}(\dot{E})}&{{\mathbf{D}}_{\mathbf{R}^{+}}^{+}(\dot{E}^{*})}&{{\mathbf{D}}_{\mathbf{R}^{+}}^{+}(\dot{E}^{*})}\\ {\displaystyle\left|\vphantom{\int}R j_{*}\right.}&{\left.j^{-1}\right|}&{\displaystyle\left|\vphantom{\int}R j_{!}\right.}\\ {{\mathbf{D}}_{\mathbf{R}^{+}}^{+}(E)\xrightarrow[\wedge]{}\mathbf{D}_{\mathbf{R}^{+}}^{+}(E^{*})}&{{\mathbf{D}}_{\mathbf{R}^{+}}^{+}(E^{*})\xrightarrow[\vee]{}\mathbf{D}_{\mathbf{R}^{+}}^{+}(E)\enspace.}\\ \end{array}
\]

\statementlabel{练习 三.18}令 $X$  为维度 $n$  的 $C^0$  流形，令 $x \in X$  且令 $U$  为 $x$  的开邻域。定义  $x$  处的残基态射，将  $\text{Res}(x; \cdot)$  表示为复合：

\[
H^{n-1}(U\backslash\{x\};{o r}_{X})\to H_{\{x\}}^{n}(X;{o r}_{X})\xrightarrow{\sim}A\enspace.
\]

(i) 假设 $X$  是紧致的，并让 $Z$  是 $X$  的有限子集。让 $u \in H^{n-1}(X \setminus Z; \sigma t_X)$  。证明：

\[
\sum_{x\in Z}\operatorname{Res}(x;u)=0.
\]

（提示：将 $ H_{Z}^{n}(X; \sigma_{X}) $ 发送到 $ H^{n}(X; \sigma_{X}) $ 。）

(ii) 人们不再假设 $X$  紧凑。相反，人们假设  $H^{n-1}(X; \sigma t_X) = H^n(X; \sigma t_X) = 0$  。令 $Y$  为维度 $n - 1$  的紧凑 $C^0$  流形，并令 $\gamma$  为连续映射 $\gamma: Y \to X$  。假设被赋予同构  $\gamma^{-1} \sigma t_X \simeq \sigma t_Y$  。如果  $Z$  是  $X$  的闭子集，它不

与  $ \gamma(Y) $  相交，其中 1 表示  $ \gamma^{-1} $  从  $ \gamma^{-1} or_X \xrightarrow{\sim} or_Y $  推导出的映射  $ H^{n-1}(X \setminus Z; or_X) \to H^{n-1}(Y; or_Y) $  。

让 $ x \in X \backslash \gamma(Y) $  。由于  $ H^{n-1}(X \backslash \{x\}, o \iota_X) \simeq A $  ，只有一个  $ u \in H^{n-1}(X \backslash \{x\}, o \iota_X) $  使得  $ \text{Res}(x; u) = 1 $  。一套：

\[
\mathrm{Ind}(x;\gamma)=\int_{\gamma}\gamma^{-1}(u)
\]

其中  $\int_{Y}$  表示映射  $H^{n-1}(Y; \sigma_{\gamma}) \to A$  。

现在让  $Z$  成为  $X\backslash\gamma(Y)$  的有限子集，并让  $v\in H^{n-1}(X\backslash Z; o\iota_{X})$  。证明柯西留数公式：

\[
\int\limits_{Y}\gamma^{-1}(v)=\sum_{x\in\mathbb{Z}}\operatorname{I n d}(x;\gamma)\operatorname{R e s}(x;v)\enspace.
\]

（此练习摘自 Iversen [2]。）

\statementlabel{练习 三.19}我们保留命题 2.7.5 和图 (2.7.5) 的符号，并且假设所有空间都是局部紧的，并且函子  $h_!$  、  $p_{Y!}$  、  $p_X!$  具有有限的上同调维数。此外，我们假设 $h$ 是正确的。然后使用  $R f_j! \circ f_j^! \to \mathrm{id}_X$  定义  $j=0,1$  函子  $f_{j\neq} : R p_{Y!} \circ p_{Y}^! \to R p_{X!} \circ p_X^!$  的态射。证明 $f_{0\neq} = f_1\neq$ 。

\statementlabel{练习 三.20}在下面的 (i) 和 (ii) 中，我们设置  $ X = \mathbb{R} $  。我们选择  $ a, b \in X $  和  $ a < b $  并设置  $ I_b^- =] -\infty, b] $  、  $ I_a^+ = [a, +\infty[ $  。

(i) 证明 $ Z_x $  与复形准同构：

\[
K\colon0\to\mathbb{Z}_{I_{\bar{b}}}\oplus Z_{I_{a}^{+}}\to\mathbb{Z}_{I_{\bar{b}}\cap I_{a}^{+}}\to0~.
\]

让  $ \mathcal{C}_X $  表示 de Rham 复数  $ 0 \to \mathcal{C}_X^\infty \to \mathcal{C}_X^\infty $  、  $ (1) \to 0 $  （参见 II §9）。构造复形的态射： $ K \to \mathcal{C}_X^\phi $ ，使得组合

\[
\mathbb{Z}_{X}\simeq H^{0}(K)\to H^{0}(\mathcal{C}_{X}^{{:}})\simeq\mathbb{C}_{X}
\]

与与  $ \mathbb{Z} \subset \mathbb{C} $  相关的态射  $ \mathbb{Z}_X \to \mathbb{C}_X $  一致。

(ii) 证明 $ H_c^1(X; \mathbb{Z}_X) \simeq H^1(\Gamma_c(X; K)) $  ，并证明下图可换算为符号； （参见 (3.3.16) 了解  $ \int_X $  的定义。）

\[
\begin{array}{l c l}{H_{c}^{1}(X;\mathbb{Z}_{X})\xrightarrow{\sim}H^{1}(\varGamma_{c}(X;K))\xrightarrow{\phi}H^{1}(\varGamma_{c}(X;\mathcal{C}_{X}^{c}))}\\ {\quad\bigg|\bigg|_{\mathbb{Z}}}&{\xrightarrow{\quad}}&{\mathbb{C}}\\ \end{array}
\]

(iii) 证明当  $A = \mathbb{C}$  时，同态 (3.3.15) 和 (3.3.16) 在符号上相等。

\statementlabel{练习 三.21}令 $f: Y \to X$  为局部紧空间的拓扑淹没，并令 $Z_{1}, Z_{2}$  为 $Y$  的两个闭子集，使得 $Z_{1} \subset Z_{2}$  。假设

 $ H_c^j((Z_2\backslash Z_1)\cap f^{-1}(x);\mathbb{Z}_Y)=0 $  对于任何  $ j $  。证明对于任何  $ F\in\mathrm{Ob}(\mathbb{D}^b(X)) $  ，  $ Rf_*R\Gamma_{Z_1}(f^{-1}F)\to Rf_*R\Gamma_{Z_2}(f^{-1}F) $  是同构。

\subsection{注记}
我们在本书开头参考了 C. Houzel 的《简史》，了解层对偶性和函子  $ f' $  的详细历史。

让我们回想一下，函子  $f'$  是由 Grothendieck [4] 于 1963 年在他的局部诺特方案相干层的对偶性理论中首次引入的。该理论在其他环境中充当了对偶理论的模型：局部紧空间（Verdier [1]）、étale 上同调（[SGA 4]、[SGA 5]）以及最近的  $\mathcal{D}$ -模（参见 XI §2）。我们在此提出的理论（第 1 节和第 4 节）是 Verdier 的理论，在海德堡-斯特拉斯堡研讨会 [SHS] 中有详细介绍。请注意，§2 和 §3 的许多结果早已为人所知（例如，参见 Borel [1]、Borel-Moore [1]）。

Fourier-Sato 变换在数学的多个部分中起着重要作用（参见 Kashiwara-Hotta [1]、Malgrange [2]、Brylinski [2]），并且由于 Deligne 在代数案例中存在类似的构造（参见 Illusie [1]、Katz-Laumon [1]、Laumon [1]）。这种变换首先由 Sato（参见 Sato-Kawai-Kashiwara [1]）在 1970 年左右针对球丛引入，直到 1975 年之后（尤其是 Kashiwara [4] 的论文之后），人们才系统地研究这些问题中的向量丛（即包括零截面）。除了在 Brylinski-Malgrange-Verdier [1] 中执行的零部分的引入之外，§7 的所有结果基本上都包含在 Sato-Kawai-Kashiwara [1] 和 Kashiwara-Kawai [2] 中。

\section{特化与微局部化}
\subsection{概要}
令 $X$  为流形， $M$  为闭子流形。我们首先构造一个新的流形 $\widetilde{X}_M$ ，即 $X$ 中 $M$ 的法向变形。该流形的维度比  $X$  的维度大一，并且被赋予映射  $(p, t): \widetilde{X}_M \to X \times \mathbb{R}$  ，使得  $t^{-1}(c)$  同构于  $c \neq 0$  的  $X$  ，并且  $t^{-1}(0)$  同构于  $T_M X$  ，即  $X$  中  $M$  的正规丛。

我们使用这个流形将  $X$  上的一个层  $F$  关联起来（或者更一般地关联到  $F \in \mathrm{Ob}(\mathbf{D}^{b}(X))$  ），  $\mathbf{D}^{b}(T_{M}X)$  的对象  $\nu_{M}(F)$  称为  $F$  沿  $M$  的特化。其傅立叶-佐藤变换  $\mu_{M}(F)$  是  $F$  沿  $M$  的微局部化。

定义了函子  $v_{M}$  和  $\mu_{M}$  并研究了它们的函子性质后，我们接下来继续研究函子  $\mu\text{hom}$  。

这个新函子概括了微局部化函子，并将在本书的其余部分中发挥核心作用。 §2和§3的结果源自Sato-Kawai-Kashiwara [1]。

公约4.0。我们保持惯例3.0。此外，所有流形和流形的态射都被假定为  $ C^{\infty} $  类或实解析类。

从现在开始，我们将主要研究层有界复合体的范畴（参见推论 3.3.12）。当然，许多结果仍然适用于较弱的假设。

\subsection{法变形与法锥}
令 $X$  为 $n$  维流形。令 $M$  为余维 $l$  的闭子流形， $T_M X$  为 $X$  中 $M$  的正规丛，即 $M$  上由精确序列定义的向量丛：

\[
0\to T M\to M\underset{x}{\times}T X\to T_{M}X\to0.
\]

我们将构造一个新的流形 $ \widetilde{X}_{M} $ 和两个映射：

\[
\left\{\begin{aligned}{}&{{}p\colon\widetilde{X}_{M}\to X~,}\\ {}&{{}t\colon\widetilde{X}_{M}\to\mathbb{R}~,\quad}\\ \end{aligned}\right.
\]

这样

\[
\left\{\begin{aligned}{}&{{}p^{-1}(X\backslash M)}&{}&{{}\mathrm{i s~i s o m o r p h i c~t o~}(X\backslash M)\times(\mathbb{R}\backslash\{0\})\;,}\\ {}&{{}t^{-1}(\mathbb{R}\backslash\{0\})}&{}&{{}\mathrm{i s~i s o m o r p h i c~t o~}X\times(\mathbb{R}\backslash\{0\})\;,}\\ {}&{{}t^{-1}(0)}&{}&{{}\mathrm{i s~i s o m o r p h i c~t o~}T_{M}X\;.}\\ \end{aligned}\right.
\]

流形 $ \widetilde{X}_{M} $ 将被称为M在X中的法向变形。

为此，请考虑开放覆盖  $ X = \bigcup_i U_i $  和开放嵌入  $ \phi_i: U_i \hookrightarrow \mathbb{R}^n $  ，例如  $ U_i \cap M = \phi_i^{-1}(\{0\}^l \times \mathbb{R}^{n-l}) $  。设置  $ x = (x', x'') \in \mathbb{R}^l \times \mathbb{R}^{n-l} $  ，并定义：

\[
V_{i}=\left\{(x,t)\in\mathbb{R}^{n}\times\mathbb{R};(t x^{\prime},x^{\prime\prime})\in\phi_{i}(U_{i})\right\}.
\]

我们用 $ t_{V_i}\colon V_i \to \mathbb{R} $  表示投影 $ (x,t) \mapsto t $  ，用 $ p_{V_i}\colon V_i \to U_i $  表示映射： $ (x,t) \mapsto \phi_i^{-1}(tx',x'') $  。一个定义了映射

\[
\psi_{j i}\colon V_{i}\mathop{\times}_{U_{i}}(U_{i}\cap U_{j})\to\mathbb{R}^{n}
\]

通过设置 $ \psi_{ji}(x,t)=(\psi_{ji}^{\prime}(x,t),\psi_{ji}^{\prime\prime}(x,t)) $ ：

\[
(t\psi_{j i}^{\prime}(x,t),\psi_{j i}^{\prime\prime}(x,t))=\phi_{j}\phi_{i}^{-1}(t x^{\prime},x^{\prime\prime})~.
\]

这是可能的，因为  $\phi_{i}\phi_{i}^{-1}(tx',x'')$  的第一个  $l$  组件在  $t=0$  时消失。

如果 $t_{i} = t_{j}$  和 $x_{j} = \psi_{ji}(x_{i}, t_{i})$  则令 $\mathcal{R}$  为标识 $(x_{i}, t_{i}) \in V_{i}$  和 $(x_{j}, t_{j}) \in V_{j}$  的等价关系。我们设定：

\[
\widetilde{X}_{M}=\left(\bigsqcup_{i}V_{i}\right)\bigg/\mathcal{R}.
\]

可以直接检查 $ \widetilde{X}_M $ 实际上是一个流形，并且分别由 $ p|_{V_i} = p_{V_i} $ 和 $ t|_{V_i} = t_{V_i} $ 定义的映射 $ p: \widetilde{X}_M \to X $ 和 $ t: \widetilde{X}_M \to \mathbb{R} $ 是明确定义的，并且满足(4.1.3)中的前两个陈述。我们说 $ \widetilde{X}_M $ 是 $ X $ 中 $ M $ 的法线变形。

请注意，乘法群  $ \mathbb{R} \setminus \{0\} $  通过以下方式作用于  $ \widetilde{X}_M $ ：

\[
(c,(x^{\prime},x^{\prime \prime},t))\mapsto(c x^{\prime},x^{\prime \prime},c^{-1}t).
\]

让我们证明  $ \tilde{X}_M $  的超曲面  $ t^{-1}(0) $  可以与法线束  $ T_M X $  规范地等同。如上所述设置 $ \psi_{ji}(x,t) = (\psi_{ji}^{\prime}, \psi_{ji}^{\prime\prime}) $ ，并设置 $ \phi_j \circ \phi_i^{-1} = \phi_{ji} = (\phi_{ji}^{\prime}, \phi_{ji}^{\prime\prime}) $ 。然后：

\[
\left\{\begin{aligned}{\psi_{j i}^{\prime}(x,0)}&{{}=\sum_{k=1}^{l}x_{k}\frac{\partial}{\partial x_{k}}\phi_{j i}^{\prime}(0,x^{\prime\prime})~,}\\ {\psi_{j i}^{\prime\prime}(x,0)}&{{}=\phi_{j i}^{\prime\prime}(0,x^{\prime\prime})~.}\\ \end{aligned}\right.
\]

因此，族  $\{V_i \cap t^{-1}(0), \psi_{ji}(x,0)\}$  定义了局部图表对  $T_M X$  的覆盖，并且  $x = (x', x'') \in V_i \cap t^{-1}(0)$  与  $(0, x'') \in M$  处的法线向量  $x'$  进行了识别。

我们用 $\Omega$ 表示作为 $\mathbb{R}^+$ 的逆像获得的 $X_M$ 的开子集，通过映射 $t$ ，通过 $j$ 嵌入 $\Omega \hookrightarrow \widetilde{X}_M$ ，通过 $\tilde{p}$ 映射 $p \circ j$ ，通过 $s$ 浸入 $T_M X \hookrightarrow \widetilde{X}_M$ 。我们用  $\tau$  表示投影  $T_M X \to M$  ，用  $i$  表示包含  $M \hookrightarrow X$  。当不担心混淆时，我们有时会写 $\tau$ 而不是 $i \circ \tau$ 。

\[
\begin{array}{ccc}T_{M}X&\subset\underbrace{\longrightarrow}_{s}&\widetilde{X}_{M}\xleftarrow{\quad j\quad}\Omega\\\left.\begin{array}{c}\tau\\ \downarrow\\ \end{array}\right\}&&\left|\begin{array}{c}p\\ \tilde{p}\\ \end{array}\right.\\M&\underbrace{\longleftarrow}_{i}&\widetilde{X}\end{array}
\]

请注意， $ \tilde{p} $  是平滑的， $ \Omega $  通过映射  $ (\tilde{p}, t) $  与  $ X \times \mathbb{R}^+ $  同构。请注意， $ p^{-1}(M) $  是 $ T_M X $  和 $ M \times \mathbb{R} $  的并集， $ T_M X \cap (M \times \mathbb{R}) = M \times \{0\} $  与 $ T_M X $  的零部分一致。

\begin{center}
\includegraphics[width=0.40\textwidth]{流行上的层_Chn-assets/img_in_image_box_346_723_823_1205.png}
\end{center}

\begin{center}图4.1.1\end{center}

\statementlabel{定义 4.1.1}(i) 令 S 为 X 的子集。“沿 M 的 S 的法向圆锥”，表示为  $ C_{M}(S) $  ，是一个集合：

\[
C_{M}(S)=T_{M}X\cap\widetilde{p}^{-1}(S).
\]

(ii) 令 $S_{1}$  和 $S_{2}$  为 $X$  的两个子集。法锥 $C(S_{1},S_{2})$  是 $TX$  的集合 $C_{\Delta x}(S_{1}\times S_{2})$ ，其中 $\Delta_{x}$  是 $X\times X$  的对角线， $TX$  通过第一个投影与 $T_{\Delta x}(X\times X)$  进行识别。

通过其构造，  $C_M(S)$  是  $T_M X$  的闭圆锥子集，并且它在  $M$  上的投影是集合  $M \cap \overline{S}$  。

类似地，  $C(S_1, S_2)$  是  $TX$  的闭圆锥曲线子集。请注意，如果 $M$  是闭子流形，则 $C(S, M)$  是 $C_M(S)$  通过投影 $M \times_X TX \to T_M X$  的逆像，如以下命题所示。

\statementlabel{命题 4.1.2}(i) 令 $(x)$  为 $X$  上的局部坐标系，并令 $S_1$  和 $S_2$  为 $X$  的两个子集。让 $(x_o; v_o) \in TX$  。然后：

\[
\begin{array}{c}{\underset{}{\left(x_{\circ};v_{\circ}\right)\in C\left(S_{1},S_{2}\right)}{\Leftrightarrow}{~t h e r e~e x i s t s~a~s e q u e n c e~}\left\{\left(x_{n},y_{n},c_{n}\right)\right\}}\\ {{i n~}S_{1}\times S_{2}\times\mathbb{R}^{+}{~s u c h~t h a t:}}\\ \end{array}
\]

\[
x_{n}\xrightarrow[n]{}x_{\circ},\qquad y_{n}\xrightarrow[n]{}x_{\circ},\qquad c_{n}(x_{n}-y_{n})\xrightarrow[n]{}v_{\circ}.
\]

(ii) 令 $(x) = (x', x'')$  为 $X$  、 $M = \{x; x' = 0\}$  上的局部坐标系，并令 $S$  为 $X$  的子集。让 $x_{\circ} = (0, x_{\circ}') \in M$ ， $(x_{\circ}; v_{\circ}) \in T_{M}X$ 。然后：

\[
\begin{aligned}{(x_{\circ};v_{\circ})\in C_{M}(S)\Leftrightarrow}&{{}{t h e r e~e x i s t s~a~s e q u e n c e~}\left\{(x_{n},c_{n})\right\}\quad{i n~}S\times\mathbb{R}^{+}{~s u c h~t h a t,~}}\\ {}&{{}{s e t t i n g~}x_{n}=(x_{n}^{\prime},x_{n}^{\prime\prime}):x_{n}\xrightarrow[n]{}x_{\circ}~,~c_{n}x_{n}^{\prime}\xrightarrow[n]{}v_{\circ}~.}\\ \end{aligned}
\]

\statementlabel{证明}证明(ii)就足够了。令 $(x', x'', t)$  为 $\tilde{X}_M$  、 $p$  上的坐标： $\tilde{X}_M \to X$  为映射 $(x', x'', t) \mapsto (tx', x'')$  。

(a) 假设  $ (x_{\circ}, v_{\circ}) $  属于  $ C_M(S) $  。那么  $ \tilde{p}^{-1}(S) $  中存在序列  $ \{(x_n', x_n', t_n)\} $  使得  $ (x_n', x_n', t_n) \to (v_{\circ}, x_{\circ}》，0) $  。  $ S \times \mathbb{R}^+ $  中的序列  $ \{(t_nx_n', x_n》，t_n^{-1})\} $  具有所需的属性。

(b) 相反，令 $\{(x_n, c_n)\}$  为 $S \times \mathbb{R}^+$  中的序列，使得 $x_n \to (0, x_o^-)$  、 $c_n x_n' \to v_o$  。如果序列  $\{c_n\}$  无界，我们可以通过提取子序列来假设  $c_n \to +\infty$  。然后  $\tilde{p}^{-1}(S)$  中的序列  $\{(c_n x_n', x_n''', c_n^{-1})\}$  收敛到  $(v_o, x_o'', 0)$  。如果序列  $\{c_n\}$  有界，则  $v_o = 0$  。我们可以找到一个正数序列 $\{\varepsilon_n\}$ ，这样 $\varepsilon_n \to 0$ ， $\varepsilon_n^{-1} x_n' \to 0$ 。那么  $\{(\varepsilon_n^{-1} x_n', x_n''', \varepsilon_n)\}$  是  $\tilde{p}^{-1}(S)$  中的一个序列，它收敛到  $(0, x_o'', 0)$  。   $\square$ 

\statementlabel{命题 4.1.3}令 V 为  $ T_{M}X $  的圆锥开子集。

(i) 设 $W$  为 $\tilde{X}_M$  中 $V$  的开邻域，并设 $U = \tilde{p}(W \cap \Omega)$  。然后 $V \cap C_M(X \setminus U) = \varnothing$ 。

(ii) 相反，令  $U$  为  $X$  的开子集，使得  $V \cap C_M(X \setminus U) = \varnothing$  。那么 $\tilde{p}^{-1}(U) \cup V$  就是 $\bar{\Omega} = \Omega \cup T_M X$  中 $V$  的开邻域。

\statementlabel{证明}(i) 从  $ W \cap \tilde{p}^{-1}(X \setminus U) = \varnothing $  开始，我们就有了  $ V \cap \tilde{p}^{-1}(X \setminus U) = \varnothing $  。

(ii) 根据  $C_M(X \setminus U), C_M(X \setminus U) \cup \tilde{p}^{-1}(X \setminus U)$  的定义， $\Omega$  是封闭的。由于  $T_M X \setminus V$  是封闭的，并且  $C_M(X \setminus U)$  包含在  $T_M X \setminus V, (T_M X \setminus V) \cup \tilde{p}^{-1}(X \setminus U)$  中， $\Omega$  中是封闭的。它的互补集是  $\tilde{p}^{-1}(U) \cup V$  。

接下来的结果对于下面定理4.2.3的证明至关重要。

\statementlabel{命题 4.1.4}令  $V$  为  $T_M X$  的圆锥开子集。然后 $\widetilde{X}_M$ 中 $V$ 的开放邻域 $W$ 族，使得映射 $p: W \cap \Omega \to X$ 的所有纤维都连接起来，形成 $V$ 的邻域系统。

\statementlabel{证明}令 $W$  成为 $\widetilde{X}_{M}$  中 $V$  的开放邻域。设定：

\[
X^{\prime}=\tilde{X}_{M}~,\qquad\mathring{X}^{\prime}=\dot{T}_{M}X\cup p^{-1}(X\backslash M)=X^{\prime}\backslash(M\times\mathbb{R})~,\qquad S=\mathring{X}^{\prime}/\mathbb{R}^{+}~.
\]

（回想一下  $ \dot{T}_M X = T_M X \setminus M $  和  $ p^{-1}(M) = T_M X \cup (M \times \mathbb{R}) $  。）

那么 $S$  是流形， $\alpha: X' \to S$  是 $\mathbb{R}^+-$  丛。此外，  $S$  包含  $\dot{T}_M X/\mathbb{R}^+$  ，作为超曲面，而  $S \to X$  是正确的贴图。如有必要，缩小 $X$ ，让我们通过扩展 $\dot{T}_M X \to \dot{T}_M X/\mathbb{R}^+$ 的一部分来选择 $\dot{X}' \to S$ 的 $\sigma$ 部分。设定：

 $ W' = \bigcup_{x \in \sigma^{-1}(W)} \{\text{the connected component of } \alpha^{-1}(x) \cap W \text{ containing } \sigma(x)\} $ 。

通过构建，映射 $ W' \to S $ 的纤维被连接起来，并且 $ W' $ 是 $ V \cap \dot{T}_M X $ 的开邻域。

然后我们定义 $ W'' = W' \cup V \cup (W \cap t^{-1} \mathbb{R}^-) $ 。这是  $ V, W'' \subset W $  的开邻域， $ p: W'' \cap \Omega = W' \to X $  的纤维是相连的。   $ \square $ 

令  $f: Y \to X$  为流形的态射，让  $N$  为余维  $k$  的  $Y$  的闭子流形，并假设  $f(N) \subset M$  。我们用  $f_{|N}$  表示由  $f$  诱导的从  $N$  到  $M$  的映射。

我们用 $f'$  表示与 $f$  关联的从 $TY$  到 $Y \times_X TX$  的映射，用 $f_\tau$  表示基础变化 $Y \times_X TX \to TX$  。那么切线贴图 $Tf$ 就是复合贴图 $f_\tau \circ f':$ 

\[
Tf:TY\xrightarrow[f^{\prime}]{}{\mathrm{}}Y\xrightarrow[x]{}TX\xrightarrow[f_{*}]{}TX\enspace.
\]

我们用 $T_{N}f$ 表示与 $Tf$ 相关的从 $T_{N}Y$ 到 $T_{M}X$ 的映射，用 $f_{N}^{\prime}$ 表示从 $T_{N}Y$ 到 $N\times_{M}T_{M}X$ 的映射，以及 $f_{N\tau}$ 基本变化 $N\times_{M}T_{M}X\to T_{M}X$ 。

\[
T_{N}f\colon T_{N}Y{\xrightarrow{}}\!f^{\prime}_{N}\to N\nonumber\times_{M}T_{M}X\xrightarrow[f_{N\tau}]{}T_{M}X\quad.
\]

如果不怕混淆，我们仍然写 $f_{\tau}$ 而不是 $f_{N\tau}$ ， $f^{\prime}$ 而不是 $f_{N}^{\prime}$ 。

让我们用  $p_X$  、  $j_X$  、  $t_X$  、  $s_X$  、  $\tilde{p}_X$  表示与  $X$  中  $M$  的法线变形相关的映射，并让  $\Omega_X = t_X^{-1}(\mathbb{R}^+)$  ，同样对于  $p_Y$  、  $j_Y$  、  $t_Y$  ，  $s_Y$ ， $\tilde{p}_Y$ ， $\Omega_Y$ 。映射  $f$  定义了从  $\tilde{Y}_N$  到  $\tilde{X}_M$  的映射  $\tilde{f}'$  。如果  $y = (y', y'')$  （分别为  $x = (x', x'')$  ）是  $Y$  （分别为  $X$  ）上的局部坐标系，使得  $N = \{y; y' = 0\}$  （分别为  $M = \{x; x' = 0\}$  ）和  $f = (f_1, f_2)$  ，则：

\[
\left\{\begin{aligned}{}&{{}\tilde{f}^{\prime}(y^{\prime},y^{\prime\prime},t)=\left(\frac{1}{t}f_{1}(t y^{\prime},y^{\prime\prime}),f_{2}(t y^{\prime},y^{\prime\prime}),t\right)\quad\mathrm{f o r}\quad t\neq0,}\\ {}&{{}\tilde{f}^{\prime}(y^{\prime},y^{\prime\prime},0)=(\mathsf{d}_{y^{\prime}}f_{1}(0,y^{\prime\prime})\cdot y^{\prime},f_{2}(0,y^{\prime\prime}),0).}\\ \end{aligned}\right.
\]

因此我们得到下面的交换图，其中所有的平方都是笛卡尔的。

\[
\begin{array}{c c c c c|c c c c c}{T_{N}Y}&{\xrightarrow{s_{Y}}}&{\widetilde{Y}_{N}}&{\xleftrightarrow{j_{Y}}\quad}&{\Omega_{Y}}&{\xrightarrow{\quad}}&{\widetilde{p}_{Y}}&{Y}\\ {T_{N}f\Bigg\downarrow}&{}&{\Box}&{\widetilde{f}^{\prime}\Bigg\downarrow}&{}&{\Box}&{\widetilde{f}\Bigg\downarrow}&{}&{\Box}&{f\Bigg\downarrow}\\ {T_{M}X}&{\xleftrightarrow{s_{X}}}&{\widetilde{X}_{M}}&{\xleftrightarrow{j_{X}}\quad}&{\Omega_{X}}&{\xleftrightarrow{\quad}}&{\widetilde{p}_{X}}&{X.}\\ \end{array}
\]

请注意，  $ \tilde{p}_X $  和  $ \tilde{p}_Y $  是平滑的，而  $ t_Y = t_X \circ \tilde{f}' $  是平滑的。然而图表：

\[
\begin{array}{c c c c}{}&{\widetilde{Y}_{N}}&{\xrightarrow{\quad p_{Y}\quad}}&{Y}\\ {\left.\widetilde{f}^{\prime}\right|}&{}&{}&{\left|\begin{array}{l}{}\\ {}\\ {}\\ {}\end{array}\right|f}\\ {}&{\widetilde{X}_{M}}&{\xrightarrow{~-\quad\longrightarrow\quad~}}&{X}\\ \end{array}
\]

一般来说不是笛卡尔坐标系，即使 $f$ 是真映射， $\tilde{f}'$ 也可能是不正确的。如果 $f$ 对于 $M$ 来说是“干净的”或者更好的“横向的”，那么这种病态的情况就不会出现。为了回忆这些定义，使用余切丛和共角丛会更方便（参见下面的§3）。

令 $f: Y \to X$  为流形的态射， $M$  为 $X$  的闭子流形。

\statementlabel{定义 4.1.5}(i) 如果  $f^{-1}(M)$  是  $Y$  的子流形  $N$  ，并且映射  $f_N': N \times_M T_M^* X \to T_N^* Y$  是满射的，则可以说  $f$  相对于  $M$  是干净的。

(ii) 如果映射 $f'|_{Y\times_X T_M^*X}$  :  $Y\times_X T_M^*X \to T^*Y$  是单射的，则可以说 $f$  是 $M$  的横向。

如果  $f$  是封闭嵌入，则也可以说  $Y$  是干净的（或横向的），而不是说  $f$  是干净的（或横向的）。如果  $f$  横向于  $M$  ，则  $f$  相对于  $M$  是干净的，但反之亦然，如示例所示，其中  $Y$  是  $X$  的子流形， $M$  是  $Y$  的子流形。如果  $f$  相对于  $M$  和  $f^{-1}(M)=N$  是干净的，则映射  $p\times\tilde{f}^{\prime}:\tilde{Y}_{N}\to Y\times\tilde{X}_{M}$  是闭合嵌入。这是因为  $f_{N}^{\prime}:T_{N}Y\to N\times_{M}T_{M}X$  是一个封闭嵌入， $\tilde{Y}_{N}\subsetneq Y\times\tilde{X}_{M}$  、  $t_{Y}=t_{X}\circ\tilde{f}^{\prime}$  、  $T_{N}Y=t_{Y}^{-1}(0)$  和  $T_{M}X=t_{X}^{-1}(0)$  的单射性。

如果  $f$  横向于  $M$  和  $f^{-1}(M)=N$  ，则  $f': T_N Y \to Y \times_X T_M X$  是同构，并且平方 (4.1.12) 是笛卡尔矩阵。

在第七章中，我们需要对定义 4.1.5 进行以下概括。

\statementlabel{定义 4.1.6}令 $f_{1}: Y_{1} \to X$  和 $f_{2}: Y_{2} \to X$  为流形的两个态射。如果映射  $(f_{1}, f_{2}): Y_{1} \times Y_{2} \to X \times X$  相对于  $X \times X$  的对角线是横向的（或干净的），则可以说  $f_{1}$  和  $f_{2}$  是横向的（或干净的）。

请注意，如果  $f_{2}$  是将子流形  $Y_{2}$  嵌入到  $X$  中，那么这相当于说  $f_{1}$  相对于  $Y_{2}$  是横向的（或干净的）。

\subsection{特化}
令 M 为 X 的闭子流形，如§1 所示。在本书中，我们将保留第 1 节中介绍的符号，特别是 (4.1.5) 中的符号。

让 $ F \in \mathrm{Ob}(\mathbb{D}^{b}(X)) $  。

\statementlabel{引理 4.2.1}存在自然的同构：

\[
s^{-1}R j_{*}\tilde{p}^{-1}F\simeq s^{!}j_{!}\tilde{p}^{!}F\quad.
\]

\statementlabel{证明}考虑一下著名的三角形：

\[
(p^{-1}F)_{\Omega}\longrightarrow Rj_{*}\tilde{p}^{-1}F\longrightarrow R\Gamma_{\{t=0\}}((p^{-1}F)_{\Omega})[1]\longrightarrow+1
\]

应用 $ s^{-1} $  我们得到：

\[
\begin{aligned}{s^{-1}R j_{*}\tilde{p}^{-1}F}&{{}\simeq s^{\dagger}(p^{-1}F)_{\Omega}[1]}\\ {}&{{}\simeq s^{\dagger}j_{!}\tilde{p}^{\dagger}F\enspace,}\\ \end{aligned}
\]

自 $ \tilde{p}^! \simeq \tilde{p}^{-1}[1] $ 以来。 □

\statementlabel{定义 4.2.2}让 $ F \in \mathrm{Ob}(\mathbb{D}^{b}(X)) $  。一套：

\[
\nu_{M}(F)=s^{-1}R j_{*}\tilde{p}^{-1}F\simeq s^{!}j_{!}\tilde{p}^{!}F\quad,
\]

并表示  $ \nu_{M}(F) $  是 F 沿 M 的特化。

\statementlabel{定理 4.2.3}让 $ F \in \mathrm{Ob}(\mathbf{D}^{b}(X)) $  。然后：

(i)  $ \nu_M(F) \in \mathrm{Ob}(\mathbb{D}^b_{\mathbb{R}^+}(T_M X)) $  和 $ \mathrm{supp}(\nu_M(F)) \subset C_M(\mathrm{supp}(F)) $  。

(ii) 令  $V$  为  $T_M X$  的圆锥开子集。然后：

\[
H^{j}(V;\nu_{M}(F))=\varinjlim U^{j}(U;F)\enspace,
\]

其中  $U$  的范围贯穿  $X$  的开子集族，使得  $C_M(X \setminus U) \cap V = \emptyset$  ； （参见图 4.2.a）。

特别是如果  $ v \in T_M X $  ，则：

\[
H^{j}(\nu_{M}(F))_{v}=\varinjlim_{U}H^{j}(U;F)\enspace,
\]

其中  $U$  的范围贯穿  $X$  的开放子集族，使得  $v \notin C_{u}(X \setminus U)$  。

(iii) 令 A 为  $ T_{M}X $  的闭圆锥曲线子集。然后：

\[
H_{A}^{j}(T_{M}X;\nu_{M}(F))=\varinjlim_{Z,U}H_{Z\cap U}^{j}(U;F)\enspace,
\]

其中  $U$  的范围是  $X$  中  $M$  的开邻域族， $Z$  的范围是  $X$  的闭子集族，这样  $C_{M}(Z) \subset A$  ； （参见图 4.2.b）。

(iv) 具有同构：

\[
\nu_{M}(F)|_{M}\simeq R\tau_{*}(\nu_{M}(F))\simeq F|_{M},
\]

\[
(R\varGamma_{M}(\nu_{M}(F)))|_{M}\simeq R\tau_{!}(\nu_{M}(F))\simeq R\varGamma_{M}(F)|_{M}.
\]

（五） $ R\dot{\tau}_{*}(v_{M}(F)|_{\dot{\tau}_{M}X})\simeq R\Gamma_{X\backslash M}(F)|_{M}. $ 

\begin{center}
\includegraphics[width=0.24\textwidth]{流行上的层_Chn-assets/img_in_image_box_117_160_411_470.png}
\end{center}

\begin{center}图。 4.2a\end{center}

\begin{center}
\includegraphics[width=0.35\textwidth]{流行上的层_Chn-assets/img_in_image_box_629_159_1055_480.png}
\end{center}

\begin{center}图。 4.2b\end{center}

\statementlabel{证明}(i) 由于  $ \tilde{p}^{-1}F $  相对于  $ \mathbb{R}^{+} $  的动作是局部常数，因此相同的属性将适用于  $ Rj_{*}\tilde{p}^{-1}F $  ，然后适用于  $ s^{-1}Rj_{*}\tilde{p}^{-1}F $  。最后一句话是显而易见的。

(ii) 令  $U$  为  $X$  的开子集，使得  $V \cap C_M(X \setminus U) = \varnothing$  。我们有态射链：

\[
\left\{\begin{aligned}{R\varGamma(U;F)}&{{}\to R\varGamma(p^{-1}(U);p^{-1}F)}\\ {}&{{}\to R\varGamma(p^{-1}(U)\cap\varOmega;p^{-1}F)}\\ {}&{{}\to R\varGamma(\tilde{p}^{-1}(U)\cup V;R j_{*}j^{-1}p^{-1}F)}\\ {}&{{}\to R\varGamma(V;\nu_{M}(F))~,}\\ \end{aligned}\right.
\]

其中第三个箭头存在，因为 $ p^{-1}(U) \cup V $  是 $ \Omega $  中 $ V $  的邻域（命题4.1.3）。

由此我们得到态射：

\[
\operatorname*{l i m}_{\stackrel{\rightarrow}{U}}H^{k}(U;F)\to H^{k}(V;\nu_{M}(F))\enspace.
\]

让我们证明这是一个同构。我们有：

\[
\begin{aligned}{H^{k}(V;\nu_{M}(F))}&{{}\simeq\varinjlim_{\overrightarrow{w}}H^{k}(W;R j_{*}j^{-1}p^{-1}F)}\\ {}&{{}\simeq\varinjlim_{\overrightarrow{w}}H^{k}(W\cap\varOmega;p^{-1}F),}\\ \end{aligned}
\]

其中  $W$  的范围贯穿  $V$  的邻域系统。因此，我们可以通过命题 4.1.4 假设  $p: W \cap \Omega \to p(W \cap \Omega)$  具有连接纤维，即与  $\mathbb{R}$  同胚的纤维。然后我们应用命题3.3.9得到：

\[
H^{k}(W\cap\varOmega;p^{-1}F)\simeq H^{k}(p(W\cap\varOmega);F).
\]

由于 $ p(W \cap \Omega) $  的范围贯穿 $ X $  的开子集 $ U $  族，使得 $ C_M(X \setminus U) \cap V = \varnothing $  （命题4.1.3），我们得到(ii)。

(iii) 令 U 和 Z 如声明中所示。我们有态射链：

\[
\begin{aligned}{R\varGamma_{\mathcal{Z}\cap U}(U;F)}&{{}\to R\varGamma_{p^{-1}(\mathcal{Z}\cap U)}(p^{-1}(U);p^{-1}F)}\\ {}&{{}\to R\varGamma_{p^{-1}(\mathcal{Z}\cap U)\cap\varOmega}(p^{-1}(U)\cap\varOmega;p^{-1}F)}\\ {}&{{}\to R\varGamma_{(p^{-1}(\mathcal{Z}\cap U)\cap\varOmega)\cup A}(p^{-1}(U);R j_{*}j^{-1}p^{-1}F)}\\ {}&{{}\to R\varGamma_{A}(T_{M}X;\nu_{M}(F))~.}\\ \end{aligned}
\]

在这里，我们使用了  $ (p^{-1}(Z \cap U) \cap \Omega) \cup A $  在  $ p^{-1}(U) $  中闭合的事实。由此我们得到交换图：

\[
\begin{array}{c c c c c c}{\cdots\longrightarrow}&{\displaystyle\operatorname*{l i m}_{\overrightarrow{U}}H^{k-1}(U\setminus Z;F)}&{\longrightarrow}&{\displaystyle\operatorname*{l i m}_{\overrightarrow{U}}H_{Z\cap U}^{k}(U;F)\longrightarrow}&{\displaystyle\operatorname*{l i m}_{\overrightarrow{U}}H^{k}(U;F)}&{\longrightarrow\cdots}\\ {}&{\bigg\downarrow\vphantom{\overrightarrow{U}}\gamma_{k-1}}&{}&{\bigg\downarrow\vphantom{\overrightarrow{U}}\alpha_{k}}&{\bigg\downarrow\vphantom{\overrightarrow{U}}\beta_{k}}&{}\\ {}&{\cdots\longrightarrow}&{H^{k-1}(T_{M}X\setminus A;\nu_{M}(F))}&{\longrightarrow}&{H_{A}^{k}(T_{M}X;\nu_{M}(F))\longrightarrow}&{H^{k}(T_{M}X;\nu_{M}(F))\longrightarrow\cdots.}\\ \end{array}
\]

由于所有行都是精确的，并且所有  $ \gamma_{k} $  和  $ \beta_{k} $  都是 (ii) 的同构，因此所有  $ \alpha_{k} $  都是同构。

(iv) 考虑图 (4.1.5) 并用 k 表示 M 通过零部分嵌入  $ T_{M}X $  中。我们有态射：

\[
\begin{aligned}{F|_{{\cal M}}}&{{}\simeq k^{-1}s^{-1}p^{-1}F}\\ {}&{{}\to k^{-1}s^{-1}R j_{*}j^{-1}p^{-1}F}\\ {}&{{}\simeq\nu_{{\cal M}}(F)|_{{\cal M}}~,}\\ \end{aligned}
\]

and

\[
\begin{aligned}{k^{!}\nu_{M}(F)}&{{}\simeq k^{!}s^{!}j_{!}j^{!}p^{!}F}\\ {}&{{}\to k^{!}s^{!}p^{!}F}\\ {}&{{}\simeq i^{!}F\enspace,}\\ \end{aligned}
\]

这些态射是 (ii) 和 (iii) 的同构。 (iv) 中的其他同构来自命题 3.7.5。

(v) 我们有一个可区分三角形的态射：

\[
\begin{array}{c c c c c c}{R\varGamma_{M}(F)|_{M}}&{\xrightarrow{\quad}}&{F|_{M}}&{\xrightarrow{\quad}}&{R\varGamma_{X\setminus M}(F)|_{M}}&{\xrightarrow{\quad}+1}\\ {\quad\bigg\downarrow}&{\quad}&{\bigg\downarrow}&{\quad\bigg\downarrow}&{\quad\bigg\downarrow}&{\quad}\\ {R\varGamma_{M}(\nu_{M}(F))|_{M}}&{\xrightarrow{\quad}}&{R\tau_{*}\nu_{M}(F)}&{\xrightarrow{\quad}}&{R\dot{\tau}_{*}(\nu_{M}(F)|\dot{\tau}_{M X})}&{\xrightarrow{\quad}+1}\end{array}.
\]

由于左边和中间的垂直箭头是同构的（iv），所以右边的箭头也是同构的。 ⑨

令  $f: Y \to X$  为流形的态射，  $N$  （分别为  $M$  ）为  $Y$  （分别为  $X$  ）的闭子流形，并假设  $f(N) \subset M$  。我们遵循 IV §1 的符号（参见 (4.1.8), (4.1.9), (4.1.11)）。

\statementlabel{命题 4.2.4}让 $ G \in \text{Ob}(\mathbf{D}^b(Y)) $  。

(i) 存在规范态射的交换图：

\[
\begin{array}{c c c c}{R(T_{N}f)_{!}\nu_{N}(G)}&{\longrightarrow}&{\nu_{M}(R f,G)}\\ {\bigg\downarrow}&{}&{\bigg\downarrow}\\ {R(T_{N}f)_{*}\nu_{N}(G)}&{\longleftarrow}&{\nu_{M}(R f_{*}G)}\\ \end{array}.
\]

(ii) 此外，如果  $\mathrm{supp}(G) \to X$  和  $C_N(\mathrm{supp}(G)) \to T_M X$  是真映射，并且如果  $\mathrm{supp}(G) \cap f^{-1}(M) \subset N$  ，则所有这些态射都是同构。

特别是，如果 $f^{-1}(M)=N$  和 $f$  相对于 $M$  是干净的并且在supp(G) 上是真映射，那么所有这些态射都是同构。

\statementlabel{证明}(i) 我们有态射链（参见练习 III.9）：

\[
\begin{aligned}{{R}(T_{N}f)_{!}\nu_{N}(G)}&{{}={R}(T_{N}f)_{!}s_{Y}^{-1}{R}j_{Y*}\tilde{p}_{Y}^{-1}G}\\ {}&{{}\simeq s_{X}^{-1}{R}\tilde{f}_{!}^{\prime}{R}j_{Y*}\tilde{p}_{Y}^{-1}G}\\ {}&{{}\to s_{X}^{-1}{R}j_{X*}{R}\tilde{f}_{!}\tilde{p}_{Y}^{-1}G}\\ {}&{{}\simeq s_{X}^{-1}{R}j_{X*}\tilde{p}_{X}^{-1}{R}f_{!}G}\\ {}&{{}=\nu_{M}({R}f_{!}G)~,}\\ \end{aligned}
\]

and

\[
\begin{aligned}{\nu_{M}(R f_{*}{G})}&{{}=s_{X}^{-1}R j_{X*}\tilde{p}_{X}^{-1}R f_{*}{G}}\\ {}&{{}\simeq s_{X}^{-1}R j_{X*}R\tilde{f}_{*}\tilde{p}_{Y}^{-1}{G}}\\ {}&{{}\simeq s_{X}^{-1}R\tilde{f}_{*}^{\prime}R j_{Y*}\tilde{p}_{Y}^{-1}{G}}\\ {}&{{}\to R(T_{N}f)_{*}\nu_{Y}^{-1}R j_{Y*}\tilde{p}_{Y}^{-1}{G}}\\ {}&{{}=R(T_{N}f)_{*}\nu_{N}({G})\enspace.}\\ \end{aligned}
\]

交换律很容易通过构造得到。

(ii) 如果  $ \tilde{p}_{Y}^{-1}(\operatorname{supp}(G)) $  对  $ \tilde{X}_{M} $  是真映射，则所有态射都是同构，因为可以分别用  $ R\tilde{f}_{!}^{\prime} $  和  $ R(T_{N}f)_{!} $  替换  $ R\tilde{f}_{*}^{\prime} $  和  $ R(T_{N}f)_{*} $  。因此，足以证明对于 Y 的闭子集 Z， $ \tilde{p}_{Y}^{-1}(Z) $ 

如果  $ Z $  对  $ X $  正确，如果  $ C_N(Z) $  对  $ T_M X $  正确，并且如果  $ Z \cap f^{-1}(M) \subset N $  正确。由于 $ \widetilde{p}_Y^{-1}(Z) \to \widetilde{X}_M $  的纤维是紧凑的，因此足以表明这是一个闭合映射。令 $ \{u_n\}_{n} $  为 $ \widetilde{p}_Y^{-1}(Z) $  中的序列，使得 $ \{\widetilde{f}'(u_n)\}_n $  收敛。我们必须证明  $ \{u_n\}_{n} $  的子序列收敛。我们可以假设 $ \{\widetilde{p}_Y(u_n)\}_n $ 收敛。由于  $ \widetilde{p}_Y^{-1}(Z) \setminus T_N Y \to \widetilde{X}_M \setminus T_M X $  是真映射，我们可以假设  $ \{\widetilde{f}'(u_n)\}_n $  收敛到  $ T_M X $  点。那么  $ \{\widetilde{p}_Y(u_n)\} $  的极限点包含在  $ Z \cap f^{-1}(M) $  中，因此包含在  $ N $  中。

按照(4.1.10)中的 $X$ 和 $Y$ 的局部坐标系，设置 $u_n = (y'_n, y''_n, t_n)$  。然后， $t_n \to 0$ ， $t_n y'_n \to 0$ 。我们可以进一步假设 $t_n > 0$ （即 $u_n \in \tilde{p}_Y^{-1}(z)$ ）。由于 $\{y''_n\}_{n}$  收敛，因此足以证明 $\{|y'_n|\}_{n}$  有界。现在假设 $|y'_n| \to \infty$ ，我们将推断出一个矛盾。通过提取子序列， $\{y'_n/|y'_n|\}_{n}$  收敛到非零向量 $v$  。那么 $\{(y'_n/|y'_n|, y''_n, t_n|y'_n|\}_{n}$ 属于 $\tilde{p}_Y^{-1}(Z)$ 并收敛到不属于零部分的点 $p \in T_N Y$ 。另一方面， $\tilde{f}'(u_n) = \left( \frac{1}{t_n} f_1(t_n y'_n, y''_n), f_2(t_n y'_n, y''_n) \right)$  收敛，因此 $\left\{ \frac{1}{t_n |y'_n|} f_1(t_n y'_n, y''_n) \right\}_{n}$  收敛到零。这意味着  $T_N f(p)$  属于  $T_M X$  的零部分。因此  $C_N(Z) \cap (T_N f)^{-1}(T_N f(p)) \supset \mathbb{R}_{\geq 0} p$  ，这与  $C_N(Z) \to T_M X$  正确的事实相矛盾。   $\square$ 

为了研究逆像，首先注意如果  $ \tau_{X} $  表示映射  $ T_{M}X \to X $  ，那么：

\[
\tau_{x}^{!}A_{x}\simeq\tau_{x}^{-1}A_{x}~,
\]

这意味着  $ \tau_{Y}^{-1}f^!A_{X} \simeq (T_{N}f)^!A_{T_M X} $  ，或等效的  $ \tau^{-1}\omega_{Y/X} \simeq \omega_{T_M Y/T_M X} $  。然后：

\[
\left\{\begin{aligned}{\omega_{T_{N}Y/N\times_{M}T_{M}X}}&{{}\simeq\tau^{-1}\omega_{Y/X}\otimes\tau^{-1}\omega_{N/M}^{\otimes-1},}\\ {\omega_{N\times_{M}T_{M}^{*}X/T_{N}^{*}Y}}&{{}\simeq\pi^{-1}\omega_{N/M}\otimes\pi^{-1}\omega_{Y/X}^{\otimes-1}.}\\ \end{aligned}\right.
\]

为简单起见，我们通常会写  $\omega_{Y/X}$  而不是  $\tau^{-1}\omega_{Y/X}$  或  $\pi^{-1}\omega_{Y/X}$  ，对于  $\omega_{N/M}$  、  $\sigma\iota_{Y/X}$  等也是如此。我们还写  $\omega_{Y/X}\otimes\cdot$  而不是  $\omega_{Y/X}\otimes^{L}\cdot$  等。

\statementlabel{命题 4.2.5}让 $ F \in \mathrm{Ob}(\mathbb{D}^{b}(X)) $  。然后是规范态射：

\[
\alpha:(T_{N}f)^{-1}v_{M}(F)\to v_{N}(f^{-1}F)\enspace,
\]

\[
\beta:\nu_{N}(f^{!}F)\to(T_{N}f)^{\dagger}\nu_{M}(F)\;,
\]

如下图所示：

\[
\begin{array}{l c l}{\omega_{T_{N}Y/T_{M}X}\otimes(T_{N}f)^{-1}\nu_{M}(F)\xrightarrow[\omega_{Y/X}\otimes\alpha]{}\nu_{N}(\omega_{Y/X}\otimes f^{-1}F)}\\ {\Bigg\downarrow\begin{array}{c c c}{\gamma}\\ {\longleftarrow}\\ \end{array}}\\ {(T_{N}f)^{\dagger}\nu_{M}(F)}&{\xleftarrow[\beta]{}\nu_{N}(f^{\dagger}F)}\\ \end{array}
\]

这里垂直箭头是从（3.1.6）推导出来的。

所有这些态射在开集上都是同构，其中  $ T_{N}f: T_{N}Y \to T_{M}X $  是光滑的。

特别是，如果  $f: Y \to X$  和  $f|_{N}: N \to M$  是光滑的，则所有这些态射都是同构。

\statementlabel{证明}映射  $ \tilde{p}_{Y} $  和  $ \tilde{p}_{X} $  是平滑的，我们有一个同构  $ \tilde{p}_{Y}^{-1}f^{!} \simeq \tilde{f}^{!} \tilde{p}_{X}^{-1} $  。我们得到态射链：

\[
\begin{aligned}{(T_{N}f)^{-1}\nu_{M}(F)}&{{}=(T_{N}f)^{-1}s_{X}^{-1}R j_{X*}\tilde{p}_{X}^{-1}F}\\ {}&{{}\simeq s_{Y}^{-1}\tilde{f}^{\prime-1}R j_{X*}\tilde{p}_{X}^{-1}F}\\ {}&{{}\to s_{Y}^{-1}R j_{Y*}\tilde{f}^{-1}\tilde{p}_{X}^{-1}F}\\ {}&{{}\simeq s_{Y}^{-1}R j_{Y*}\tilde{p}_{Y}^{-1}f^{-1}F}\\ {}&{{}=\nu_{N}(f^{-1}F),}\\ \end{aligned}
\]

and

\[
\begin{aligned}{\nu_{N}(f^{!}F)}&{{}=s_{Y}^{-1}R j_{Y*}\tilde{p}_{Y}^{-1}f^{!}F}\\ {}&{{}\simeq s_{Y}^{-1}R j_{Y*}\tilde{f}^{!}\tilde{p}_{X}^{-1}F}\\ {}&{{}\simeq s_{Y}^{-1}\tilde{f}^{\prime!}R j_{X*}\tilde{p}_{X}^{-1}F}\\ {}&{{}\to(T_{N}f)^{!}s_{X}^{-1}R j_{X*}\tilde{p}_{X}^{-1}F}\\ {}&{{}=(T_{N}f)^{\prime}\nu_{M}(F).}\\ \end{aligned}
\]

交换律由证明得出。

在 $ T_{N}f $  光滑的点上， $ \tilde{f}' $  光滑， $ \alpha $  、 $ \beta $  和 $ \gamma $  是同构的。   $ \Box $ 

最后让我们研究一下张量积。

令X和Y为两个流形，并令M和N分别为X和Y的两个闭子流形。

\statementlabel{命题 4.2.6}让  $ F \in \mathrm{Ob}(\mathbf{D}^b(X)) $  和  $ G \in \mathrm{Ob}(\mathbf{D}^b(Y)) $  。那么  $ \mathbf{D}^b(T_{M \times N}(X \times Y)) $  中有一个自然态射：

\[
\nu_{M}(F)\mathop{\boxtimes}^{L}\nu_{N}(G)\to\nu_{M\times N}\left(\begin{matrix}{F\boxtimes^{L}G}\\ \end{matrix}\right).
\]

\statementlabel{证明}令  $p_X, j_X, \tilde{p}_X, s_X$  为与  $X$  中  $M$  的法线变形相关的自然贴图， $p_Y, \ldots, s_Y$  也类似。我们简称为  $p, j, \tilde{p}, s$  而不是  $p_{X \times Y}, \ldots, s_{X \times Y}$  ，并用  $p', j', \tilde{p}'$  、  $s'$  表示映射  $p_X \times p_Y, j_X \times j_Y$  、  $\tilde{p}_X \times \tilde{p}_Y, s_X \times s_Y$  。

从  $(X \times Y)_{M \times N}$  到  $\widetilde{X}_M \times \widetilde{Y}_N$  存在一个自然的封闭嵌入  $k$ ，它定义了交换图：

\begin{center}
\includegraphics[width=0.88\textwidth]{流行上的层_Chn-assets/diagram_4_2_6.png}
\end{center}

我们有态射：

\[
\begin{aligned}{\nu_{M}(F)\mathop{\boxtimes}^{L}\nu_{N}(G)}&{{}=s_{X}^{-1}R j_{X*}\tilde{p}_{X}^{-1}F\mathop{\boxtimes}^{L}s_{Y}^{-1}R j_{Y*}\tilde{p}_{Y}^{-1}G}\\ {}&{{}\to s^{{^{\prime}}-1}R j_{*}^{\prime}\tilde{p}^{\prime-1}\Bigg(F\mathop{\boxtimes}^{L}G\Bigg)}\\ {}&{{}\simeq s^{-1}k^{-1}R j_{*}^{\prime}\tilde{p}^{\prime-1}\Bigg(F\mathop{\boxtimes}^{L}G\Bigg)}\\ {}&{{}\to s^{-1}R j_{*}\tilde{k}^{-1}\tilde{p}^{\prime-1}\Bigg(F\mathop{\boxtimes}^{L}G\Bigg)}\\ {}&{{}\simeq s^{-1}R j_{*}\tilde{p}^{-1}\Bigg(F\mathop{\boxtimes}^{L}G\Bigg)\;.}&{}&{{}\Box}\\ \end{aligned}
\]

\statementlabel{推论 4.2.7}让  $F$  和  $G$  属于  $\mathrm{Ob}(\mathbb{D}^{b}(X))$  。那么  $\mathbb{D}^{b}(T_{M}X)$  中存在自然态射：

\[
\nu_{M}(F)\overset{L}{\otimes}\nu_{M}(G)\to\nu_{M}\left(F\overset{L}{\otimes}G\right)\;.
\]

\statementlabel{证明}令  $\varDelta_{X}$  为  $X \times X$  的对角线，  $\delta_{X}$  为嵌入  $\varDelta_{X} \hookrightarrow X \times X$  ，并类似地定义  $\varDelta_{TX}$  ，  $\varDelta_{T_{M}X}$  ，  $\delta_{TX}$  ，  $\delta_{T_{M}X}$  。请注意，使用符号 (4.1.9)  $\delta_{T_{M}X} = T_{M} \delta_{X}$  。然后应用命题4.2.5和4.2.6，我们得到态射：

\[
\begin{aligned}{\nu_{M}(F)\otimes\nu_{M}(G)}&{{}\simeq\delta_{T_{M}X}^{-1}\bigg(\nu_{M}(F)\boxtimes\nu_{M}(G)\bigg)}\\ {}&{{}\to\delta_{T_{M}X}^{-1}\bigg(\nu_{M\times M}\bigg(F\boxtimes G\bigg)\bigg)}\\ {}&{{}\to\nu_{M}\bigg(\delta_{X}^{-1}\bigg(F\boxtimes G\bigg)\bigg)}\\ {}&{{}\simeq\nu_{M}\bigg(F\otimes G\bigg)\;.\quad\Box}\\ \end{aligned}
\]

\subsection{微局部化}
令 $M$  为余维 $l$  的 $X$  的闭子流形，如§1 中所示。我们用  $T_M^*X$  表示  $X$  中  $M$  的共形丛，即映射  $M \times_X T^*X \to T^*M$  的内核。我们用 $\pi$ 表示投影 $T^*X \to X$ 或其限制 $T_M^*X \to M$ ，有时我们写 $\pi$ 而不是 $i \circ \pi$ ，其中 $i: M \to X$ 。

我们用  $ \dot{\pi} $  表示  $ \pi $  到  $ \dot{T}^{*}X = T^{*}X\backslash X $  的限制。

\statementlabel{定义 4.3.1}让 $ F \in \mathrm{Ob}(\mathbb{D}^b(X)) $  。  $ F $  沿  $ M $  的微局部化，表示为  $ \mu_M(F) $  ，是  $ \nu_M(F) $  的 Fourier-Sato 变换：

\[
\mu_{M}(F)=\nu_{M}(F)^{\wedge}\quad.
\]

应用定理 4.2.3 和 III §7 的结果。我们得到：

\statementlabel{定理 4.3.2}让 $ F \in \mathrm{Ob}(\mathbb{D}^{b}(X)) $  。然后：

（一） $ \mu_M(F) \in \mathrm{Ob}(\mathbb{D}_{\mathbb{R}^+}^b(T_M^*X)) $ 。

(ii) 令 $V$  为 $T_{M}^{*}$  X 的凸开锥体。则：

\[
H^{j}(V;\mu_{M}(F))=\operatorname*{l i m}_{\stackrel{\overrightarrow{U,Z}}{}}H_{Z\cap U}^{j}(U;F)\enspace,
\]

其中 U 的范围是 X 的开子集族，这样  $U \cap M = \pi(V)$  ，而 Z 的范围是闭子集族，这样  $C_{M}(Z) \subset V^{\circ}$  。特别是让  $p \in T_{M}^{*}X$  。然后：

\[
H^{j}(\mu_{M}(F))_{p}=\varinjlim_{Z}H_{Z}^{j}(F)_{\pi(p)}\;,
\]

其中  $Z$  的范围贯穿  $X$  的闭子集族，使得  $C_M(Z)_{\pi(p)} \subset \{v \in (T_M X)_{\pi(p)}; \langle v, p \rangle > 0\} \cup \{0\}$  。

(iii) 令 Z 为包含零截面 M 的  $ T_{M}^{*}X $  的真闭凸锥。则：

\[
H_{Z}^{j}(T_{M}^{\ast}X;\mu_{M}(F)\otimes\mathfrak{o r}_{M/X})=\varinjlim_{U}H^{j-l}(U;F)\enspace,
\]

其中  $U$  的范围贯穿  $X$  的开放子集族，使得  $C_{M}(X \setminus U) \cap \mathrm{Int} Z^{\circ a} = \emptyset$  。

(iv)

\[
\mu_{M}(F)|_{M}\simeq R\pi_{*}\mu_{M}(F)\simeq R\varGamma_{M}(F)|_{M}\simeq i^{!}F,
\]

\[
R\pi_{!}\mu_{_{M}}(F)\simeq R\varGamma_{_{M}}(\mu_{_{M}}(F))\simeq i^{-1}F\otimes\omega_{_{M/X}}~.
\]

请注意：

\[
i^{-1}F\otimes\omega_{M/X}=F|_{\cal M}\otimes o\tau_{M/X}[-l].
\]

将最后的结果应用于区分的三角形：

\[
R\varGamma_{M}(\mu_{M}(F))\longrightarrow R\pi_{*}\mu_{M}(F)\longrightarrow R\dot{\pi}_{*}\mu_{M}(F)\xrightarrow[+1]{}
\]

我们得到了显着的三角形：

\[
F|_{M}\otimes\omega_{M/X}\longrightarrow R\varGamma_{M}(F)|_{M}\longrightarrow R\dot{\pi}_{*}\mu_{M}(F)\xrightarrow[+1]{}~.
\]

现在让  $f: Y \to X$  为流形的态射，  $N$  为  $Y$  的余维  $k$  的闭子流形，以及  $f(N) \subset M$  。映射 $Tf$ （参见4.1.8）定义了映射：

\[
T^{*}Y\xleftarrow[{}^{t f^{\prime}}]{}Y\underset{x}{\times}T^{*}X\xrightarrow[{f_{\pi}}]{}T^{*}X~,
\]

从而诱发

\[
T_{N}^{*}Y\xleftarrow[{}^{t}f_{N}^{\prime}]{N}\times_{M}T_{M}^{*}X\xrightarrow[{f_{N\pi}}]{}T_{M}^{*}X\enspace.
\]

如果不怕混淆，我们会写  $ f_{\pi} $  而不是  $ f_{N\pi} $  ，也写 'f' 而不是 'f'\_\{\textbackslash{}pi\}。我们设定：

\[
T_{Y}^{*}X=\mathop{\bf K e r}({}^{t}f^{\prime}\colon Y\times\mathop{\bf T}^{*}X\to T^{*}Y)={}^{t}f^{\prime}{}^{-1}(T_{Y}^{*}Y)\enspace.
\]

\statementlabel{备注 4.3.3}在文献中，人们经常写 $ \rho_f $ 和 $ \bar{\omega}_f $ ，或者简单地分别写成 $ \rho $ 和 $ \bar{\omega} $ ，而不是 $ f' $ 和 $ f_\pi $ 。我们不会在这里使用这些符号。

我们将把 Fourier-Sato 函子应用于命题 4.2.4 和 4.2.5 的态射。

\statementlabel{命题 4.3.4}让 $ G \in \text{Ob}(\mathbb{D}^b(Y)) $  。那么，存在规范态射的交换图：

\[
\begin{array}{ccc}R f_{N\pi!}f_{N}^{\prime-1}\mu_{N}(G)&\xrightarrow{\quad}&\mu_{M}(R f_{!}G)\\\downarrow&&\downarrow\\R f_{N\pi*}({}^{t}f_{N}^{\prime}\mu_{N}(G)\otimes\omega_{Y/X}\otimes\omega_{N/M}^{\otimes-1})&\xleftarrow{\quad}&\mu_{M}(R f_{*}G)\end{array}
\]

如果  $\operatorname{supp}(G) \to X$  和  $C_N(\operatorname{supp}(G)) \to T_M X$  是真映射，并且如果  $f^{-1}(M) \cap \operatorname{supp}(G) \subset N$  ，那么这些态射是同构。

特别是如果  $f^{-1}(M)=N$  和  $f$  相对于  $M$  是干净的并且在supp(G)上是真映射，那么所有这些态射都是同构的。

\statementlabel{证明}让 $ H = \nu_N(G) $  。将命题 3.7.13 和 3.7.14 应用于  $ f' = f_\tau \circ f_N' $  ，我们得到：

\[
\begin{aligned}{(R(T_{N}f)_{!}H)^{\wedge}}&{{}\simeq R f_{\pi!}(R f_{N!}^{\prime}H)^{\wedge}}\\ {}&{{}\simeq R f_{\pi!}{}^{t}f_{N}^{\prime-1}(H^{\wedge})}\\ \end{aligned}\;,
\]

\[
\begin{aligned}{(R(T_{N}f)_{*}{\cal H})^{\wedge}}&{{}\simeq R f_{\pi*}(R f_{N*}^{\prime}{\cal H})^{\wedge}}\\ {}&{{}\simeq R f_{\pi*}({}^{[\imath}f_{N}^{\prime}{}^{[\imath}{\cal H}^{\wedge}\otimes\omega_{N\times_{M}T_{M}^{\ast}X/T_{N}^{\ast}Y}^{\otimes-1})~.}\\ \end{aligned}
\]

则由命题4.2.4和公式(4.2.4)得出结果。   $ \Box $ 

\statementlabel{命题 4.3.5}让 $ F \in \mathrm{Ob}(\mathbb{D}^{b}(X)) $  。

(i) 存在规范态射的交换图：

\[
\begin{array}{c c c}{R^{t}f_{N!}^{\prime}(\omega_{N/M}\otimes f_{N\pi}^{-1}\mu_{M}(F))}&{\longrightarrow}&{\mu_{N}(\omega_{Y/X}\otimes f^{-1}F)}\\ {}&{}&{}\\ {\downarrow}&{}&{\downarrow}\\ {R^{t}f_{N*}^{\prime}f_{N\pi}^{!}\mu_{M}(F)}&{\longleftarrow}&{\mu_{N}(f^{!}F)}\\ \end{array}.
\]

这里垂直箭头是从（3.1.6）推导出来的。

(ii)  $ I_{ff}: Y \to X $  和 $ f|_{N}: N \to M $  是光滑的，所有这些态射都是同构。

(iii) 如果 f 横向于 M 和  $ f^{-1}(M)=N $  ，则存在自然态射：

\[
R^{\mathfrak{t}}f_{N*}^{\prime}f_{N\pi}^{-1}\mu_{M}(F)\to\mu_{N}(f^{-1}F)\enspace.
\]

\statementlabel{证明}设置 $H = \nu_{M}(F)$ ，然后应用命题3.7.13和3.7.14。我们得到：

\[
\begin{aligned}{(T_{N}f)^{\prime}A_{T_{M}X}\otimes((T_{N}f)^{-1}H)^{\wedge}}&{{}\simeq(f_{N}^{\prime!}A_{N\times_{M}T_{M}X}\otimes f_{N}^{\prime-1}(f_{N\tau}^{\prime}A_{T_{M}X}\otimes f_{N\tau}^{-1}H))^{\wedge}}\\ {}&{{}\simeq R^{\prime}f_{N\uparrow}^{\prime}(f_{N\tau}^{\uparrow}A_{T_{M}X}\otimes f_{N\tau}^{-1}H)^{\wedge}}\\ {}&{{}\simeq R^{\prime}f_{N\uparrow}^{\prime}(f_{N\pi}^{\uparrow}A_{T_{M}^{\ast}X}\otimes f_{N\pi}^{-1}(H^{\wedge})),}\\ \end{aligned}
\]

\[
\begin{aligned}{((T_{N}f)^{\dagger}H)^{\wedge}}&{{}\simeq(f_{N}^{\prime\dagger}f_{N\tau}^{\dagger}H)^{\wedge}}\\ {}&{{}\simeq R^{\dagger}f_{N\ast}^{\prime}f_{N\pi}^{\dagger}(H^{\wedge})\enspace.}\\ \end{aligned}
\]

因此，(i) 由命题 4.2.5 得出，以及 (ii)。

如果  $f$  横向于  $M$  和  $f^{-1}(M)=N$  ，则  $f_{N}^{\prime}:T_{N}^{*}Y\to N\times_{M}T_{M}^{*}X$  是同构，我们有：  $\omega_{N/M}\simeq(\omega_{Y/X})|_{N}$  。

请注意，如果  $f$  和  $f|_{N}$  是平滑的，则  $Y \times_X T^*X$  是  $T^*Y$  的子丛，并且  $f_\pi$  会导致同构  $(Y \times_X T^*X) \cap T_N^*Y \simeq Y \times_X T_M^*X$  。

我们将回到第五章和第六章命题 4.3.4 和 4.3.5 中定义的态射，在层的微支撑的帮助下，我们将给出新的条件，以便它们是同构的。

最后我们研究张量积。

\statementlabel{命题 4.3.6}在命题 4.2.6 的情况下，存在自然态射：

\[
\mu_{M}(F){\overset{L}{\boxtimes}}\mu_{N}(G)\to\mu_{M\times N}\left(F{\overset{L}{\boxtimes}}G\right)\;.
\]

\statementlabel{证明}应用命题 4.2.6 和命题 3.7.15。 □

\statementlabel{命题 4.3.7}令 $M$  为 $X$  的子流形，并令 $\gamma: T_M^*X \times_M T_M^*X \to T_M^*X$  为加法给出的态射。那么，对于任何  $F$  、  $G \in \mathrm{Ob}(\mathbf{D}^b(X))$  ，都存在一个自然态射：

\[
R\gamma_{!}\left(\begin{matrix}{L}\\ {\mu_{M}(F)\mathop{\boxtimes}_{M}\mu_{M}(G)}\\ \end{matrix}\right)\to\mu_{M}\left(\begin{matrix}{L}\\ {F\otimes G}\\ \end{matrix}\right)\otimes\omega_{M/X}.
\]

\statementlabel{证明}令  $ \delta: T_M X \hookrightarrow T_M X \times_M T_M X $  为对角线嵌入。然后，我们有：

\[
\nu_{M}(F)\overset{L}{\otimes}\nu_{M}(G)\simeq\delta^{-1}\bigg(\nu_{M}(F)_{\underset{M}{\boxplus}}\nu_{M}(G)\bigg).
\]

注意到 $ \delta = \gamma $ 并应用命题3.7.13和3.7.15，我们得到：

\[
\begin{aligned}{\bigg(\omega_{T_{M}X/T_{M}X\times_{M}T_{M}X}\otimes\nu_{M}^{L}(F)\otimes\nu_{M}(G)\bigg)^{\wedge}}&{{}\simeq R^{i}\delta_{!}\bigg(\nu_{M}(F)\mathop{\boxtimes}_{M}^{L}\nu_{M}(G)\bigg)^{\wedge}}\\ {}&{{}\simeq R\gamma_{!}\bigg(\mu_{M}(F)\mathop{\boxtimes}_{M}^{L}\mu_{M}(G)\bigg).}\\ \end{aligned}
\]

然后，我们通过推论 4.2.7 得到所需的同态。 ⑨

\statementlabel{备注 4.3.8}命题4.3.4和4.3.5中的态射关系如下。 （证明留给读者。）

(a) 对于  $ \mathbf{D}^b(Y) $  中的任何态射  $ \varphi: G \to f^!F $  ，下图可交换：

\[
\begin{array}{c c c c}{R f_{N\pi^{!}}f_{N}^{\prime-1}\mu_{N}(G)}&{\xrightarrow[]{\quad\alpha\quad}}&{\mu_{M}(R f_{!}G)}\\ {\bigg\downarrow\varphi}&{}&{\bigg\downarrow\varphi}\\ {R f_{N\pi^{!}}f_{N}^{\prime-1}\mu_{N}(f^{!}F)}&{\xrightarrow[]{\quad\beta\quad}}&{\mu_{M}(F)}\end{array},
\]

其中  $ \alpha $  是命题 4.3.4 中第一个水平箭头给出的态射， $ \beta $  是命题 4.3.5 中第二个水平箭头给出的态射。

(b) 同样，对于  $ \mathbf{D}^b(X) $  中的任何态射  $ \psi: F \to Rf_*G $  ，下图可交换：

\[
\begin{array}{c c}{\mu_{M}(F)}&{\xrightarrow[\alpha^{\prime}]{}\;R f_{N\pi\ast}({}^{t}f_{N}^{\prime!}\mu_{N}(f^{-1}F)\otimes\omega_{Y/X}\otimes\omega_{N/M}^{\otimes-1})}\\ {\bigg\downarrow\psi}&{\bigg\downarrow\psi}\\ {\mu_{M}(R f_{*}{}^{}G)}&{\xrightarrow[\beta^{\prime}]{}\;R f_{N\pi\ast}({}^{t}f_{N}^{\prime!}\mu_{N}(G)\otimes\omega_{Y/X}\otimes\omega_{N/M}^{\otimes-1})\quad,}\end{array}
\]

其中  $ \alpha' $  由命题 4.3.5 中的第一个水平箭头给出， $ \beta' $  由命题 4.3.4 中的第二个水平箭头给出。

\subsection{微局部同态函子}
让  $f$  成为从  $Y$  到  $X$  的态射。令  $\mathcal{A}_f \subset X \times Y$  为  $f$  的图。我们用 $\mathcal{A}_X$ （或 $\mathcal{A}_Y$ ）表示 $X \times X$ （或 $Y \times Y$ ）的对角线。我们将确定

 $ Y \times_X T^* X $  和  $ T_{\mathcal{A}_f}^*(X \times Y) $  ，通过投影  $ T^*(X \times Y) \to (T^*X) \times Y $  ，类似地，我们通过第一个投影识别  $ T^*X $  （分别为  $ T^*Y $  ）和  $ T_{\mathcal{A}_X}^*(X \times X) $  （分别为  $ T_{\mathcal{A}_Y}^*(Y \times Y) $  ）。如果不怕混淆，我们将写 $ \mathcal{A} $ 而不是 $ \mathcal{A}_f $ 。这样我们就有了映射（参见（4.3.3））：

\[
\begin{array}{c c c c c}{T_{\mathcal{A}_{Y}}^{*}(Y\times Y)}&{\xleftarrow{\quad}}&{T_{\mathcal{A}_{f}}^{*}(X\times Y)}&{\xrightarrow{\quad}}&{T_{\mathcal{A}_{X}}^{*}(X\times X)}\\ {\bigg\downarrow\bigg\downarrow\bigg\downarrow}&{}&{\bigg\downarrow\bigg\downarrow}&{}&{\bigg\downarrow\bigg\downarrow}\\ {T^{*}Y}&{\xleftarrow{{\tiny~\scriptstyle{t f}~}}}&{Y\times\mathop{{T}^{*}}X}&{\xrightarrow{f_{*}}}&{T^{*}X}\\ \end{array}.
\]

请注意有用的公式：

\[
\omega_{Y\times_{X}T^{*}X/T^{*}Y}\otimes\omega_{Y\times_{X}T^{*}X/T^{*}X}\simeq A_{Y\times_{X}T^{*}X}\enspace.
\]

我们用  $ f_1: Y \times Y \to X \times Y $  表示映射  $ (f, \mathrm{id}_Y) $  并用  $ f_2: X \times Y \to X \times X $  表示映射  $ (id_X, f) $  。因此，我们有交换图，其中右侧的正方形是笛卡尔坐标， $ f_2 $  是  $ \mathcal{A}_X $  的横截面：

\[
\begin{array}{c c c c c}{Y\times Y}&{\xrightarrow[f_{1}]{}}&{X\times Y}&{\xrightarrow[f_{2}]{}}&{X\times X}\\ {\quad\quad\quad}&{}&{\quad\quad}&{}&{}\\ {\quad\cup}&{}&{\quad\quad\quad}&{}&{\quad\quad}\\ {\quad\quad\quad}&{\sim}&{}&{\quad\quad}&{}\\ {\varDelta_{Y}}&{\xrightarrow{\quad}}&{}&{\varDelta}&{\quad\varDelta_{X}\quad.}\\ \end{array}
\]

我们用 $q_{j}$ （分别 $\tilde{q}_{j}$ ， $q_{j}^{\prime}$ ）表示在 $X \times X$ （分别 $X \times Y$ ， $Y \times Y$ ）上定义的 $j$ -th投影（ $j = 1,2$ ）。

\statementlabel{定义 4.4.1}让 $ G \in \mathrm{Ob}(\mathbb{D}^{b}(Y)) $ ， $ F \in \mathrm{Ob}(\mathbb{D}^{b}(X)) $ 。我们设定：

\[
\mu\mathcal{h}o m(G\to F)=\mu_{\varDelta}R\mathcal{H}o m(\tilde{q}_{2}^{-1}G,\tilde{q}_{1}^{!}F),
\]

\[
\mu\mathrm{h o m}(F\leftarrow G)=(\mu_{\Delta}R\mathcal{H}\mathrm{o m}(\tilde{q}_{1}^{-1}F,\tilde{q}_{2}^{!}G))^{a}.
\]

(iii) 当 Y = X 且 f 为恒等式时，我们设置：

\[
\mu\mathcal{h}o m(G,F)=\mu\mathcal{h}o m(G\to F)=\mu_{\varDelta_{X}}(R\mathcal{H}o m(q_{2}^{-1}G,q_{1}^{!}F))\enspace.
\]

式(ii)中， $ (\cdot)^{a} $ 表示 $ Y \times_{X} T^{*}X $ 上的对映图的反像。

让  $ \pi $  表示从  $ T_{\mathcal{A}}^{*}(X \times Y) $  到  $ \mathcal{A} \simeq Y $  的投影。

\statementlabel{命题 4.4.2}我们有规范的同构：

\[
R\pi_{*}\mu\mathcal{h}o m(G\to F)\simeq R\mathcal{H}\mathcal{h}o m(G,f^{!}F),
\]

\[
R\pi_{*}\mu\mathcal{h}o m(F\leftarrow G)\simeq R\mathcal{H}o m(f^{-1}F,G).
\]

特别是当 Y = X 且 f 为恒等式时：

\[
R\pi_{*}\mu\mathrm{h o m}(G,F)\simeq R\mathcal{H o m}(G,F)\enspace.
\]

(ii) 假设 G（或 F）是上同调可构造的。然后：

\[
\begin{aligned}{R\pi_{1}{\mu}h o m(G\to F)}&{{}\simeq R\mathcal{H}o m(G,A_{Y})\overset{L}{\otimes}f^{-1}F\otimes\omega_{Y/X}}\\ {\left(\operatorname{r e s p.}\;R\pi_{1}{\mu}h o m(F\leftarrow G)\right)}&{{}\simeq f^{-1}R\mathcal{H}o m(F,A_{X})\overset{L}{\otimes}G}\\ \end{aligned}.
\]

特别是当 Y = X 且 f 为恒等式时

\[
R\pi_{{\mathbf{\mathit{\Pi}}}}\mu{\mathbf{\mathit{\Pi}}}_{{\mathbf{\mathit{\Pi}}}}({\mathbf{\mathit{\Pi}}}G,F)\simeq R{\mathbf{\mathit{\Pi}}}{\mathbf{\mathit{\Pi}}}_{{\mathbf{\mathit{\Pi}}}}({\mathbf{\mathit{\Pi}}})^{L}\otimes F.
\]

\statementlabel{证明}(i) 根据定理 4.3.2 我们有同构：

\[
\begin{aligned}{{R\pi_{*}\mu_{_{\Delta}}R\mathcal{H}o m(\tilde{q}_{2}^{-1}G,\tilde{q}_{1}^{-1}F)}}&{{}\simeq{R\tilde{q}_{2*}}{R\varGamma_{\varDelta}R\mathcal{H}o m(\tilde{q}_{2}^{-1}G,\tilde{q}_{1}^{!}F)}}\\ {}&{{}\simeq{R\tilde{q}_{2*}}{R\mathcal{H}o m(\tilde{q}_{2}^{-1}G,{R\varGamma_{\varDelta}\tilde{q}_{1}^{!}F})}}\\ {}&{{}\simeq{R\mathcal{H}o m(G,{R\tilde{q}_{2*}}{R\varGamma_{\varDelta}\tilde{q}_{1}^{!}F)}}}\\ {}&{{}\simeq{R\mathcal{H}o m(G,f^{!}F)}.}\\ \end{aligned}
\]

同样：

\[
\begin{aligned}{{R\pi_{*}\mu_{\varDelta}R\mathcal{H}o m(\tilde{q}_{1}^{-1}F,\tilde{q}_{2}^{!}G)}}&{{}\simeq{R\tilde{q}_{2*}}{R}\varGamma_{\varDelta}{R\mathcal{H}o m}(\tilde{q}_{1}^{-1}F,\tilde{q}_{2}^{!}G)}\\ {}&{{}\simeq{R\tilde{q}_{2*}}{R\mathcal{H}o m}((\tilde{q}_{1}^{-1}F)_{\varDelta},\tilde{q}_{2}^{!}G)}\\ {}&{{}\simeq{R\mathcal{H}o m}({R\tilde{q}_{2!}(\tilde{q}_{1}^{-1}F)}_{\varDelta},G)}\\ {}&{{}\simeq{R\mathcal{H}o m}(f^{-1}F,G).}\\ \end{aligned}
\]

(ii) 根据命题 3.4.4，证明是类似的。 ⑨

函子  $ \mu_{hom} $  概括了 Sato 微局部化的函子。

\statementlabel{命题 4.4.3}令 Y 为 X 的闭子流形，j 为嵌入  $ T_Y^*X \to T^*X $  并令  $ F \in \text{Ob}(\mathbb{D}^b(X)) $  。则有同构：

\[
\mu\mathrm{h o m}(A_{Y},F)\simeq j_{*}\mu_{Y}(F)\enspace.
\]

\statementlabel{证明}令  $f$  为沉浸感  $Y \hookrightarrow X$  。我们有：

\[
\begin{aligned}{\mu_{\varDelta_{x}}R\mathcal{H}o m(q_{2}^{-1}A_{Y},q_{1}^{!}F)}&{{}\simeq\mu_{\varDelta_{x}}(R f_{2*}f_{2}^{!}q_{1}^{!}F)}\\ {}&{{}\simeq\mu_{\varDelta_{x}}(R f_{2*}\tilde{q}_{1}^{!}F).}\\ \end{aligned}
\]

让我们将  $ T_{\Delta}^{*}(X \times Y) $  与  $ Y \times_X T_{\Delta_x}^{*}(X \times X) $  进行识别。

应用命题 4.3.4 我们得到：

\[
\mu_{\varDelta_{x}}(R f_{2*}\tilde{q}_{1}^{!}F)\simeq\mu_{\varDelta}(\tilde{q}_{1}^{!}F)\enspace.
\]

那么根据命题 4.3.5 我们有：

\[
\mu_{\varDelta}(\tilde{q}_{1}^{!}F)\simeq j_{*}\mu_{\Upsilon}(F)\enspace.\quad\Box
\]

人们可以使用 $ \gamma $ 拓扑来描述 $ \mu hom(G, F) $ 的茎（参见III5）。

\statementlabel{命题 4.4.4}假设 $X$  是向量空间。令 $F$  和 $G$  属于 $\mathrm{Ob}(\mathbf{D}^{b}(X))$  ，并令 $(x_{\circ};\xi_{\circ}) \in T^*X$  。然后：

\[
H^{j}(\mu{\mathcal{h}}o m(G,F))_{({x_{\circ}};\xi_{\circ})}=\operatorname*{l i m}_{\overrightarrow{U,\gamma}}H^{j}(R\Gamma(U;R\mathcal{H}o m(\phi_{\gamma}^{-1}R\phi_{\gamma\ast}G_{U},F))~,
\]

其中  $U$  范围穿过  $x_o$  的开邻域族，而  $\gamma$  范围穿过  $X$  的闭凸真锥族，使得  $\gamma \subset \{v \in X; \langle v, \xi_o \rangle < 0\} \cup \{0\}$  。

回想一下， $\phi_{\gamma}$  表示映射 $X \to X_{\gamma}$  ，其中 $X_{\gamma}$  是赋予 $\gamma$  拓扑的空间 $X$ 。

\statementlabel{证明}设 y 为 X 的闭真圆锥。设：

\[
Z_{\gamma}=\left\{(x,x^{\prime})\in X\times X;x^{\prime}-x\in\gamma\right\}.
\]

然后：

\[
H^{j}(\mu\mathcal{h}o m(G,F))_{({x}\circ;\xi\circ)}=\operatorname*{l i m}_{U\stackrel{\rightharpoonup}{,V,y}}H^{j}(R\varGamma_{Z_{\gamma}}(U\times V;R\mathcal{H}o m(q_{2}^{-1}G,q_{1}^{!}F))~,
\]

其中 U 和 V 的范围穿过  $x_{\circ}$  的开邻域族， $\gamma$  的范围穿过满足命题条件的闭凸真锥族。然后我们有：

\[
\begin{aligned}{R\varGamma_{Z_{\gamma}}(U\times V;R\mathcal{H}o m(q_{2}^{-1}G,q_{1}^{!}F))}&{{}\simeq R\varGamma(U\times X;R\mathcal{H}o m((q_{2}^{-1}G_{V})_{Z_{\gamma}},q_{1}^{!}F)}\\ {}&{{}\simeq R\varGamma(U;R\mathcal{H}o m(R q_{1}!(q_{2}^{-1}G_{V})_{Z_{\gamma}},F)\enspace.}\\ \end{aligned}
\]

那么结果由命题3.5.4得出。 ⑨

让我们继续描述  $ \mu h o m $  的一些功能属性。

\statementlabel{命题 4.4.5}在定义 4.4.1 的情况下，有规范态射的交换图：

\[
\begin{array}{c c c}{R^{\prime}f_{!}^{\prime}{\mu}h o m(G\to F)}&{\xrightarrow{\quad}}&{{\mu}h o m(G,f^{-1}F\otimes\omega_{Y/X})}\\ {\quad\bigg\downarrow}&{}&{\bigg\downarrow}\\ {R^{\prime}f_{*}^{\prime}{\mu}h o m(G\to F)}&{\xleftarrow{\quad}}&{{\mu}h o m(G,f^{\prime}F)}\\ \end{array}
\]

\begin{center}
\includegraphics[width=0.50\textwidth]{流行上的层_Chn-assets/img_in_image_box_266_148_871_724.png}
\end{center}

如果 f 是光滑的，则 (i) 和 (ii) 中的所有态射都是同构。如果 f 在 \textbackslash{}textit\{supp\}(G) 上是真映射，则 (iii) 和 (iv) 中的所有态射都是同构。

在进入证明之前，让我们回顾一下（定理 3.1.5 和公式（2.6.14）），态射  $ Rf_1G \to F $  定义态射  $ G \to f^1F $  ，反之亦然，态射  $ F \to Rf_*G $  定义态射  $ f^{-1}F \to G $  ，反之亦然。我们将系统地使用这些注释（用  $ Y \times_X T^*X $  替换 Y，用  $ T^*Y $  或  $ T^*X $  替换 X）。我们还将系统地使用  $ \mathrm{II}\S6 $  中构造的态射。

\statementlabel{证明}我们保留本节开头介绍的符号。然后我们有以下自然态射，通过应用命题 4.3.4 或 4.3.5 获得。

\[
\begin{aligned}{\mu h o m(G\to F)}&{{}=\mu_{\varDelta}R\mathcal{H}o m(\tilde{q}_{2}^{-1}G,\tilde{q}_{1}^{!}F)}\\ {}&{{}\to\mu_{\varDelta}R\mathcal{H}o m(\tilde{q}_{2}^{-1}G,R f_{1*}f_{1}^{-1}\tilde{q}_{1}^{!}F)}\\ {}&{{}\simeq\mu_{\varDelta}R f_{1*}R\mathcal{H}o m(q_{2}^{\prime-1}G,q_{1}^{\prime!}f^{-1}F)}\\ {}&{{}\to{^{\prime}}f^{\prime}{^{\prime}}\mu_{\varDelta_{Y}}R\mathcal{H}o m(q_{2}^{\prime-1}G,q_{1}^{\prime}{^{\prime}}f^{-1}F\otimes\omega_{Y/X}).}\\ \end{aligned}
\]

\[
\begin{aligned}{\mu\mathcal{h}o m(G,f^{!}F)}&{{}=\mu_{\varDelta_{\mathcal{Y}}}R\mathcal{H}o m(q_{2}^{\prime-1}G,q_{1}^{\prime}f^{!}F)}\\ {}&{{}\simeq\mu_{\varDelta_{\mathcal{Y}}}f_{1}^{!}R\mathcal{H}o m(\tilde{q}_{2}^{-1}G,\tilde{q}_{1}^{!}F)}\\ {}&{{}\to R^{t}f_{*}^{\prime}\mu_{\varDelta}R\mathcal{H}o m(\tilde{q}_{2}^{-1}G,\tilde{q}_{1}^{!}F).}\\ \end{aligned}
\]

如果 f 是光滑的，则最后一个态射是同构。

\[
\begin{aligned}{\mu h o m(F\leftarrow G)}&{{}=\mu_{\varDelta}R\mathcal{H}o m(\tilde{q}_{1}^{-1}F,\tilde{q}_{2}^{!}G)^{a}}\\ {}&{{}\to\mu_{\varDelta}R\mathcal{H}o m(R f_{1!}f_{1}^{!}\tilde{q}_{1}^{-1}F,\tilde{q}_{2}^{!}G)^{a}}\\ {}&{{}\simeq\mu_{\varDelta}R f_{1*}R\mathcal{H}o m(q_{1}^{^{\prime}-1}f^{!}F,q_{2}^{^{\prime}!}G)^{a}}\\ {}&{{}\to{^{\prime}}f^{!}{^{\prime}}\mu_{\varDelta_{Y}^{!}}R\mathcal{H}o m(q_{1}^{^{\prime}-1}f^{!}F,q_{2}^{^{\prime}!}G\otimes\omega_{Y/X})^{a}.}\\ \end{aligned}
\]

\[
\begin{aligned}{\mu\mathcal{h}o m(f^{-1}F,G)}&{{}\simeq\mu_{\Delta_{\Upsilon}}R\mathcal{H}o m(q_{1}^{\prime-1}f^{-1}F,q_{2}^{\prime}G)^{a}}\\ {}&{{}\simeq\mu_{\Delta_{\Upsilon}}R\mathcal{H}o m(f_{1}^{-1}\tilde{q}_{1}^{-1}F,f_{1}^{\prime}\tilde{q}_{2}^{\prime}G)^{a}}\\ {}&{{}\simeq\mu_{\Delta_{\Upsilon}}f_{1}^{\prime}R\mathcal{H}o m(\tilde{q}_{1}^{-1}F,\tilde{q}_{2}^{\prime}G)^{a}}\\ {}&{{}\to R^{\mathfrak{t}}f_{*}^{\prime}\mu_{\Delta}R\mathcal{H}o m(\tilde{q}_{1}^{-1}F,\tilde{q}_{2}^{\prime}G)^{a}.}\\ \end{aligned}
\]

如果  $f$  是光滑的，则最后一个态射是同构。

\[
\begin{aligned}{\mu{\mathcal{h}}o m(G\to F)}&{{}\to\mu_{\varDelta}R\mathcal{H}o m(f_{2}^{-1}R f_{2*}\tilde{q}_{2}^{-1}G,\tilde{q}_{1}^{!}F)}\\ {}&{{}\xrightarrow{\sim}\mu_{\varDelta}f_{2}^{!}R\mathcal{H}o m(q_{2}^{-1}R f_{*}G,q_{1}^{!}F)}\\ {}&{{}\to f_{\pi}^{!}\mu_{\varDelta_{X}}R\mathcal{H}o m(q_{2}^{-1}R f_{*}G,q_{1}^{!}F)\enspace.}\\ \end{aligned}
\]

\[
\begin{aligned}{\mu{\mathcal{h}}o m(R f_{!}G,F)}&{{}=\mu_{\varDelta_{X}}R\mathcal{H}o m(q_{2}^{-1}R f_{!}G,q_{1}^{!}F)}\\ {}&{{}\simeq\mu_{\varDelta_{X}}R f_{2*}R\mathcal{H}o m(\tilde{q}_{2}^{-1}G,\tilde{q}_{1}^{!}F)}\\ {}&{{}\to R f_{\pi*}\mu_{\varDelta}R\mathcal{H}o m(\tilde{q}_{2}^{-1}G,\tilde{q}_{1}^{!}F).}\\ \end{aligned}
\]

如果 $f$  在supp(G) 上是真映射，则最后一个态射是同构。

(iv,a)

\[
\begin{aligned}{\mu\mathcal{h}o m(F\leftarrow G)}&{{}\to\mu_{\varDelta}R\mathcal{H}o m(\tilde{q}_{1}^{-1}F,f_{2}^{!}R f_{2!}\tilde{q}_{2}^{!}G)^{a}}\\ {}&{{}\xrightarrow{\sim}\mu_{\varDelta}f_{2}^{!}R\mathcal{H}o m(q_{1}^{-1}F,q_{2}^{!}R f_{!}G)^{a}}\\ {}&{{}\to f_{\pi}^{!}\mu_{\varDelta_{x}}R\mathcal{H}o m(q_{1}^{-1}F,q_{2}^{!}R f_{!}G)^{a}.}\\ \end{aligned}
\]

\[
\begin{aligned}{\mu{\mathcal{h}}o m(F,R f_{*}G)}&{{}\simeq\mu_{\varDelta_{X}}R\mathcal{H}o m(q_{1}^{-1}F,q_{2}^{\prime}R f_{*}G)^{a}}\\ {}&{{}\simeq\mu_{\varDelta_{X}}R f_{2*}R\mathcal{H}o m(\tilde{q}_{1}^{-1}F,\tilde{q}_{2}^{\prime}G)^{a}}\\ {}&{{}\to R f_{\pi*}\mu_{\varDelta}R\mathcal{H}o m(\tilde{q}_{1}^{-1}F,\tilde{q}_{2}^{\prime}G)^{a}.}\\ \end{aligned}
\]

如果 $f$  在supp(G) 上是真映射，则最后一个态射是同构。

通过这些态射，我们得到了命题4.4.5中所述的所有水平箭头。垂直箭头是自然箭头，由  $ f_! \to f_* $  或  $ f^{-1} \otimes \omega_{Y/X} \to f^! $  定义。交换律留给读者。   $ \square $ 

\statementlabel{命题 4.4.6}在定义 4.4.1 的情况下，有规范态射：

\[
R^{t}f_{!}^{\prime}f_{\pi}^{-1}{\mu}h o m(R f_{!}G,F)\to{\mu}h o m(G,f^{-1}F\otimes\omega_{Y/X})
\]

\[
R^{t}f_{!}^{\prime}f_{\pi}^{-1}{\mu}{\overset{{\tiny~h o m}}{\sim}}(F,R f_{*}G)\to{\mu}{\overset{{\tiny~h o m}}{}}\left(f^{!}F,G\otimes\omega_{Y/X}\right)\;,
\]

\[
R f_{\pi!}{^{t} f^{^{\prime}-1}}\mu\mathcal{h}o m(G,f^{!}F)\to\mu\mathcal{h}o m(R f_{*}G,F)\enspace,
\]

\[
R f_{\pi!}f^{\prime-1}{\mu}h o m(f^{-1}F,G)\to{\mu}h o m(F,R f_{!}G)\quad.
\]

如果 f 在supp(G) 上是光滑且正确的，则(iii) 和(iv) 中的态射是同构。

\statementlabel{证明}我们有以下态射：

(i)

\[
\begin{aligned}{f_{\pi}^{-1}{\mu}h o m(R f_{!}G,F)}&{{}\to f_{\pi}^{-1}R f_{\pi*}{\mu}h o m(G\to F)}\\ {}&{{}\to{\mu}h o m(G\to F)}\\ {}&{{}\to{^{t} f^{\prime}}^{!}{\mu}h o m(G,f^{-1}F\otimes\omega_{Y/X})~,}\\ \end{aligned}
\]

(ii)

\[
\begin{aligned}{f_{\pi}^{-1}{\mu}h o m(F,R f_{*}{G})}&{{}\to f_{\pi}^{-1}R f_{\pi*}{\mu}h o m(F\leftarrow{G})}\\ {}&{{}\to{\mu}h o m(F\leftarrow{G})}\\ {}&{{}\to{^{t} f^{\prime}}^{!}{\mu}h o m(f^{!}{F},{G}\otimes\omega_{Y/X})\enspace,}\\ \end{aligned}
\]

(三)

\[
\begin{aligned}{^{t}f^{\prime-1}\mu h o m(G,f^{!}F)}&{{}\to{^{t}f^{\prime-1}R f_{*}^{\prime}\mu h o m(G\to F)}}\\ {}&{{}\to\mu h o m(G\to F)}\\ {}&{{}\to f_{\pi}^{!}\mu h o m(R f_{*}G,F)~,}\\ \end{aligned}
\]

(iv)

\[
\begin{aligned}{\mathfrak{t}f^{\prime-1}\mu h o m(f^{-1}F,G)}&{{}\to\mathfrak{t}f^{\prime-1}R\mathfrak{t}f_{*}^{\prime}\mu h o m(F\leftarrow G)}\\ {}&{{}\to\mu h o m(F\leftarrow G)}\\ {}&{{}\to f_{*}^{\prime}\mu h o m(F,R f_{!}G).}\\ \end{aligned}
\]

如果 $f$  是光滑的，而 $H \in \mathrm{Ob}(\mathbf{D}^b(Y \times_X T^*X))$  则 $H \to^t f'^{-1} R^t f_* H$  是同构的。应用命题 4.4.5，我们得到如果  $f$  在  $\mathrm{supp}(G)$  上是真映射，则 (iii) 和 (iv) 是同构。   $\square$ 

让  $ F_1 $  和  $ F_2 $  属于  $ \mathbf{D}^b(X) $  。存在规范态射：

(4.4.5)

\[
\begin{array}{c}\mu h o m(f^!F_2,f^{-1}F_1\otimes\omega_{Y/X})\stackrel{\mu h o m(f^!F_2,f^!F_1)}{\longleftrightarrow}\mu h o m(f^{-1}F_2\otimes\omega_{Y/X},f^!F_1)\\\mu h o m(f^{-1}F_2,f^{-1}F_1)\stackrel{\mu h o m(f^{-1}F_2\otimes f^!F_1)}{\longleftrightarrow}\end{array}
\]

同样，让  $ G_1 $  和  $ G_2 $  属于  $ \mathbf{D}^b(Y) $  。存在规范态射：

\[
\begin{array}{c c c c c}{\mu h o m(R f_{!}G_{2},R f_{!}G_{1})}&{}&{}&{}&{}\\ {\mu h o m(R f_{*}G_{2},R f_{!}G_{1})}&{\stackrel{\mu h o m(R f_{!}G_{2},R f_{!}G_{1})}{\searrow}}&{}&{\mu h o m(R f_{!}G_{2},R f_{*}G_{1})}&{.}\\ {}&{\searrow}&{}&{\nearrow}&{}\\ {\mu h o m(R f_{*}G_{2},R f_{*}G_{1})}&{}&{}&{}&{}\\ \end{array}
\]

\statementlabel{命题 4.4.7}令 $F_{1}, F_{2}, G_{1}, G_{2}$  如上。然后是规范态射的交换图：

\begin{center}
\includegraphics[width=0.69\textwidth]{流行上的层_Chn-assets/img_in_image_box_173_515_1000_903.png}
\end{center}

如果  $f$  是光滑的，则 (i) 中的态射都是同构。

如果 f 是闭嵌入，则 (ii) 中的态射都是同构。

\statementlabel{证明}(i) 根据命题 4.3.5，我们有一个交换图：

\[
\begin{array}{c c c}{f_{\pi}^{-1}\mu h o m(F_{2},F_{1})}&{\xrightarrow{\quad}}&{\mu_{\varDelta}f_{2}^{-1}R\mathcal{H}o m(q_{2}^{-1}F_{2},q_{1}^{!}F_{1})}\\ {}&{\downarrow}&{\downarrow}\\ {f_{\pi}^{!}\mu h o m(F_{2},F_{1})\otimes\omega_{Y/X}^{\otimes-1}}&{\longleftarrow\quad}&{\mu_{\varDelta}f_{2}^{!}R\mathcal{H}o m(q_{2}^{-1}F_{2}\otimes\omega_{Y/X},q_{1}^{!}F_{1})\enspace.}\\ \end{array}
\]

使用态射 $ f^{-1}R\mathcal{H}om(\cdot,\cdot) \to R\mathcal{H}om(f^!(\cdot),f^!(\cdot)) $ （参见练习 III.9），我们得到：

\[
\begin{array}{c c c}{f_{\pi}^{-1}{\mu}h o m(F_{2},F_{1})}&{\xrightarrow{\quad}}&{{\mu}h o m(f^{!}F_{2}\to F_{1})}\\ {\quad\bigg\downarrow}&{}&{\bigg\downarrow}\\ {f_{\pi}^{!}{\mu}h o m(F_{2},F_{1})\otimes\omega_{Y/X}^{\otimes-1}}&{\xleftarrow{\quad}}&{{\mu}h o m((f^{-1}F_{2}\otimes\omega_{Y/X})\to F_{1})~.}\\ \end{array}
\]

我们将  $ R^{\prime}f_{1}^{\prime} $  应用于 (4.4.7) 的第一个水平箭头，并将  $ R^{\prime}f_{*}^{\prime} $  应用于第二个水平箭头。那么结果由命题4.4.5(i)得出。

(ii) 根据命题 4.3.4，我们有一个交换图：

\[
\begin{array}{c c c}{\mathfrak{t}f^{\prime-1}\mu\mathfrak{h o m}(G_{2},G_{1})}&{\longrightarrow}&{\mu_{\varDelta}R f_{1^{!}}R\mathcal{H o m}(q_{1}^{\prime-1}G_{2},q_{2}^{\prime!}G_{1})^{\mathfrak{a}}}\\ {}&{\downarrow}&{\downarrow}\\ {\mathfrak{t}f^{\prime}\mu\mathfrak{h o m}(G_{2},G_{1})\otimes\omega_{Y/X}}&{\longleftarrow}&{\mu_{\varDelta}R f_{1*}\mathcal{R}\mathcal{H o m}(q_{1}^{\prime-1}G_{2},q_{2}^{\prime!}G_{1})^{\mathfrak{a}}}\\ \end{array},
\]

由此我们推断：

\[
\begin{array}{c c c c}{}&{\stackrel{t f^{\prime-1}}{\sim}\mu h o m(G_{2},G_{1})}&{\longrightarrow}&{\mu h o m(R f_{*}G_{2}\leftarrow G_{1})}\\ {}&{\downarrow}&{}&{\downarrow}\\ {\stackrel{t f^{\prime}!}{\sim}\mu h o m(G_{2},G_{1})\otimes\omega_{Y/X}}&{\longleftarrow}&{}&{\mu h o m(R f_{!}G_{2}\leftarrow G_{1})~.}\\ \end{array}
\]

现在我们将  $ Rf_{\pi} $  应用于 (4.4.8) 的第一个水平箭头，并将  $ Rf_{\pi*} $  应用于第二个水平箭头。那么结果由命题 4.4.5 (iv)得出。

假设 f 是光滑的。根据命题 4.4.5 (i)，我们有同构：

\[
R^{\mathfrak{t}}f_{!}^{\prime}{\mu}h o m(f^{!}F_{2}\to F_{1})\simeq{\mu}h o m(f^{!}F_{2},f^{-1}F_{1}\otimes\omega_{Y/X})\;.
\]

另一方面我们有：

\[
\begin{aligned}{\mu{\mathcal{h}}o m(f^{!}F_{2}\to F_{1})\otimes\omega_{Y/X}}&{{}\simeq\mu_{\varDelta}R\mathcal{H}o m(f_{2}^{-1}q_{2}^{-1}F_{2},f_{2}^{!}q_{1}^{!}F_{1})}\\ {}&{{}\simeq\mu_{\varDelta}f_{2}^{!}R\mathcal{H}o m(q_{2}^{-1}F_{2},q_{1}^{!}F_{1})}\\ {}&{{}\simeq f_{\pi}^{!}\mu_{\varDelta_{X}}R\mathcal{H}o m(q_{2}^{-1}F_{2},q_{1}^{!}F_{1})}\\ \end{aligned}
\]

根据命题 4.3.5。

这就证明了同构(i)。

最后假设 f 是封闭嵌入。由命题 4.4.5 (iii) 我们得到同构：

\[
R f_{\pi!\mu{\not}h o m}(G_{2}\to R f_{!}G_{1})\stackrel{}{\sim}\mu{\not}h o m(R f_{*}G_{2},R f_{!}G_{1})\quad.
\]

另一方面我们有：

\[
\begin{aligned}{\mu\mathcal{h}o m(G_{2}\to R f_{!}G_{1})}&{{}\simeq\mu_{\varDelta}R\mathcal{H}o m(\tilde{q}_{2}^{-1}G_{2},\tilde{q}_{1}^{!}R f_{!}G_{1})}\\ {}&{{}\simeq\mu_{\varDelta}R f_{1!}R\mathcal{H}o m(q_{2}^{^{\prime}-1}G_{2},q_{1}^{^{\prime}!}G_{1})}\\ {}&{{}\simeq{^{t}f^{\prime-1}}\mu_{\varDelta_{\mathcal{Y}}}R\mathcal{H}o m(q_{2}^{^{\prime}-1}G_{2},q_{1}^{^{\prime}!}G_{1})\enspace.}\\ \end{aligned}
\]

这就证明了同构(ii)。 □

现在我们将研究  $ \mu h o m $  的外张量积（参见符号 2.3.12）。

令 $p_X: X \to S$  和 $p_Y: Y \to S$  为流形的两个态射。我们用  $q_1$  和  $q_2$  表示在  $X \times Y$  或  $X \times_S Y$  上定义的第一个和第二个投影。我们假设：

\[
\begin{array}{c} X\times Y\text{is a submanifold of}X\times Y\enspace.\\ s\end{array}
\]

让  $j$  表示嵌入  $X \times_S Y \hookrightarrow X \times Y$  。从 $(X \times_S Y) \times_{X \times Y} T^*(X \times Y) \simeq T^*X \times_S T^*Y$ 开始，我们获得了两个映射：

\[
T^{*}\left(X\mathop{\bf\Phi}_{\mathcal{S}}\times Y\right)\xleftarrow{\quad t j^{\prime}\quad}T^{*}X\mathop{\bf\Phi}_{\mathcal{S}}\times T^{*}Y\xrightarrow{\quad j_{\pi}\quad}T^{*}X\times T^{*}Y~.
\]

\statementlabel{命题 4.4.8}令 $ F_1 $  和 $ F_2 $  属于 $ \mathbf{D}^b(X) $  并让 $ G_1 $  和 $ G_2 $  属于 $ \mathbf{D}^b(Y) $  。其中一个具有规范态射：

\[
\begin{aligned}{R^{t}j_{!}^{\prime}\bigg(\mu h o m(F_{2},F_{1})}&{{}\overset{L}{\underset{S}{\boxtimes}}\mu h o m(G_{2},G_{1})\bigg)}\\ {}&{{}\to\mu h o m\left(F_{2}\overset{L}{\underset{S}{\boxtimes}}G_{2},F_{1}\overset{L}{\underset{S}{\boxtimes}}G_{1}\right)\;,}\\ \end{aligned}
\]

\[
\begin{aligned}{R^{t}j_{!}^{\prime}\bigg(\mu}&{{}\mathcal{h o m}(F_{2},F_{1})^{a}\mathop{\boxtimes}_{s}^{L}\mu\mathcal{h o m}(G_{2},G_{1})\bigg)}\\ {}&{{}\to\mu\mathcal{h o m}(R\mathcal{H o m}(q_{1}^{-1}F_{1},q_{2}^{-1}G_{2}),R\mathcal{H o m}(q_{1}^{-1}F_{2},q_{2}^{-1}G_{1}))~.}\\ \end{aligned}
\]

\statementlabel{证明}(i) 让我们首先定义  $S = \{\mathrm{pt}\}$  时的态射。由命题 4.3.6 我们得到态射：

\[
\begin{aligned}{\mu_{\varDelta_{X}}(R\mathcal{H}o m(q_{2}^{-1}F_{2},q_{1}^{!}F_{1}))}&{{}^{\underline{{L}}}\boxtimes\mu_{\varDelta_{Y}}(R\mathcal{H}o m(q_{2}^{-1}G_{2},q_{1}^{!}G_{1}))}\\ {}&{{}\to\mu_{\varDelta_{X}\times\varDelta_{Y}}\bigg(R\mathcal{H}o m(q_{2}^{-1}F_{2},q_{1}^{!}F_{1})\boxtimes R\mathcal{H}o m(q_{2}^{-1}G_{2},q_{1}^{!}G_{1})\bigg)}\\ {}&{{}\to\mu_{\varDelta_{X}\times\varDelta_{Y}}\bigg(R\mathcal{H}o m\bigg(q_{2}^{-1}F_{2}\boxtimes q_{2}^{-1}G_{2},q_{1}^{!}F_{1}\boxtimes q_{1}^{!}G_{1}\bigg)\bigg)}\\ {}&{{}\simeq\mu\mathcal{h}o m\bigg(F_{2}\boxtimes G_{2},F_{1}\boxtimes G_{1}\bigg).}\\ \end{aligned}
\]

现在我们处理一般情况。应用前面的结果和命题 4.4.7(i)，我们得到：

\[
\begin{aligned}{R^{t}j_{!}^{\prime}\bigg(\mu h o m(F_{2},F_{1})}&{{}\mathop{\boxtimes}_{S}^{L}\mu h o m(G_{2},G_{1})\bigg)}\\ {}&{{}\simeq R^{t}j_{!}^{\prime}j_{\pi}^{-1}\bigg(\mu h o m(F_{2},F_{1})\mathop{\boxtimes}^{L}\mu h o m(G_{2},G_{1})\bigg)}\\ {}&{{}\to R^{t}j_{!}^{\prime}j_{\pi}^{-1}\mu h o m\bigg(F_{2}\mathop{\boxtimes}^{L}G_{2},F_{1}\mathop{\boxtimes}^{L}G_{1}\bigg)}\\ \end{aligned}
\]

\[
\begin{aligned}{}&{{}\to\mu{\mathcal{h}}o m{\left(j^{-1}{\left(\begin{matrix}{L}\\ {F_{2}\boxtimes G_{2}}\\ \end{matrix}\right)},j^{-1}{\left(\begin{matrix}{L}\\ {F_{1}\boxtimes G_{1}}\\ \end{matrix}\right)}\right)}}\\ {}&{{}\simeq\mu{\mathcal{h}}o m{\left(\begin{matrix}{L}\\ {F_{2}\boxtimes G_{2},F_{1}\boxtimes G_{1}}\\ {S}\\ \end{matrix}\right)}.}\\ \end{aligned}
\]

(ii) 让我们首先定义  $S = \{\mathrm{pt}\}$  时的态射。

我们用 $ q'_j $ （分别为 $ q''_j $ ，分别为 $ \tilde{q}_j $ ）表示在 $ X \times X $ （分别为 $ Y \times Y $ ，分别为 $ (X \times Y) \times (X \times Y) $ ）上定义的第j个投影。

\[
\begin{aligned}{}&{{}\mu h o m(F_{2},F_{1})^{a}\boxtimes\mu h o m(G_{2},G_{1})}\\ {}&{{}\quad\simeq\mu_{\varDelta_{X}}(R\mathcal{H}o m(q_{1}^{\prime-1}F_{2},q_{2}^{\prime!}F_{1}))\boxtimes\mu_{\varDelta_{Y}}(R\mathcal{H}o m(q_{2}^{\prime \prime-1}G_{2},q_{1}^{\prime \prime!}G_{1}))}\\ {}&{{}\quad\to\mu_{\varDelta_{X}\times\varDelta_{Y}}\Bigg(R\mathcal{H}o m(q_{1}^{\prime-1}F_{2},q_{2}^{\prime!}F_{1})\boxtimes R\mathcal{H}o m(q_{2}^{\prime \prime-1}G_{2},q_{1}^{\prime \prime!}G_{1})\Bigg)}\\ {}&{{}\quad\to\mu_{\varDelta_{X}\times\varDelta_{Y}}\Bigg(R\mathcal{H}o m(\tilde{q}_{1}^{-1}q_{1}^{-1}F_{2},\tilde{q}_{2}^{-1}q_{1}^{-1}F_{1})\otimes^{L}R\mathcal{H}o m(\tilde{q}_{2}^{-1}q_{2}^{-1}G_{2},\tilde{q}_{1}^{-1}q_{2}^{-1}G_{1})\Bigg)}\\ {}&{{}\quad\otimes\omega_{X\times Y}}\\ {}&{{}\quad\to\mu_{\varDelta_{X\times Y}}(R\mathcal{H}o m(R\mathcal{H}o m(\tilde{q}_{2}^{-1}q_{1}^{-1}F_{1},\tilde{q}_{2}^{-1}q_{2}^{-1}G_{2}),R\mathcal{H}o m(\tilde{q}_{1}^{-1}q_{1}^{-1}F_{2},\tilde{q}_{1}^{-1}q_{2}^{-1}G_{1})))}\\ {}&{{}\quad\otimes\omega_{X\times Y}}\\ {}&{{}\quad\simeq\mu h o m(R\mathcal{H}o m(q_{1}^{-1}F_{1},q_{2}^{-1}G_{2}),R\mathcal{H}o m(q_{1}^{-1}F_{2},q_{2}^{-1}G_{1}))~.}\end{aligned}
\]

现在我们使用前面的结果和命题 4.4.7 (i) 来处理一般情况。

\[
\begin{aligned}{}^{t}{j}_{!}^{\prime}\bigg(\mu{h o m}(F_{2},F_{1})^{a}&{{}\overset{L}{\boxtimes}\mu{h o m}(G_{2},G_{1})\bigg)}\\ {}&{{}\simeq{}^{t}{j}_{!}^{\prime}{j}_{\pi}^{-1}\bigg(\mu{h o m}(F_{2},F_{1})^{a}\overset{L}{\boxtimes}\mu{h o m}(G_{2},G_{1})\bigg)}\\ {}&{{}\to{}^{t}{j}_{!}^{\prime}{j}_{\pi}^{-1}\mu{h o m}(R\mathcal{H}{o m}(q_{1}^{-1}F_{1},q_{2}^{-1}G_{2}),R\mathcal{H}{o m}(q_{1}^{-1}F_{2},q_{2}^{-1}G_{1}))}\\ {}&{{}\to\mu{h o m}({j}^{!}R\mathcal{H}{o m}(q_{1}^{-1}F_{1},q_{2}^{-1}G_{2}),{j}^{!}R\mathcal{H}{o m}(q_{1}^{-1}F_{2},q_{2}^{-1}G_{1}))~.}\\ \end{aligned}
\]

在最后一项中，我们可以用  $ q_2^!G_2 $  和  $ q_2^!G_1 $  替换  $ q_2^{-1}G_2 $  和  $ q_2^{-1}G_1 $  。那么结果由命题3.1.13得出。   $ \Box $ 

现在，我们将组成函子  $ \mu h o m $  。

令  $g: Z \to Y$  和  $f: Y \to X$  为流形  $h = f \circ g$  的两个态射。我们分别定义  $q_i$  、  $q'_i$  、  $q''_i$  (i = 1,2)  $X \times Y$  、  $X \times Z$  、  $Y \times Z$  的  $i$  的第一个投影，以及  $q_{ij}X \times Y \times Z$  的  $(i,j)$  的第  $X \times Y \times Z$  的投影，（参见图表） (3.6.1))。我们使用符号  $\tilde{q}_{ij}$  来表示  $X \times Y \times Y \times Z$  上的  $(i,j)$  次投影，并且

我们用  $p_{ij}$  （分别为  $\tilde{p}_{ij}$  ）表示  $T^*X \times T^*Y \times T^*Z$  （分别为  $T^*X \times T^*Y \times T^*Y \times T^*Z$  ）的  $(i,j)$  次投影。我们用  $r_{ij}$  表示  $\tilde{p}_{ij}$  到  $\Delta_Y \times_{Y \times Y} T^*(X \times Y \times Y \times Z)$  的限制，最后我们用  $j$  表示对角线嵌入  $X \times Y \times Z \hookrightarrow X \times Y \times Y \times Z$  。请注意  $j'$  会导致同构：

\[
\varDelta_{Y\atop Y\times Y}\times T_{\varDelta_{f}\times\varDelta_{g}}^{*}(X\times Y\times Y\times Z)\xrightarrow[\scriptstyle i j^{\prime}]{\sim}T_{\varDelta_{f}\times Y\varDelta_{g}}^{*}(X\times Y\times Z),
\]

和  $ q_{13\pi} $  引发同构：

\[
\left(\varDelta_{f}\mathop{\times}_{\gamma}\varDelta_{g}\right)\mathop{\times}_{\varDelta_{h}}T_{\varDelta_{h}}^{*}(X\mathop{\times}Z)\xrightarrow{\sim}T_{\varDelta_{h}}^{*}(X\mathop{\times}Z)\enspace.
\]

因此我们有交换图：

\begin{center}
\includegraphics[width=0.90\textwidth]{流行上的层_Chn-assets/diagram_4_4_13.png}
\end{center}

\statementlabel{命题 4.4.9}让 $F$ （分别 $G$ ， $H$ ）属于 $\mathrm{Ob}(\mathbf{D}^{b}(X))$ ，（分别 $\mathrm{Ob}(\mathbf{D}^{b}(Y))$ ， $\mathrm{Ob}(\mathbf{D}^{b}(Z))$ ）。然后有一个规范态射：

\[
{}^{t f^{\prime}-1}\mu h o m(H\to G)\overset{L}{\otimes}g_{\pi}^{-1}\mu h o m(G\to F)\to\mu h o m(H\to F)~.
\]

\statementlabel{证明}设置  $ K_1 = R\mathcal{H}om(q_2^{-1}G, q_1^!F) $  、  $ K_2 = R\mathcal{H}om(q_2^{''-1}H, q_1''!G) $  和  $ L = g_{\pi}^{-1}\mu_{\Delta_f}(K_1) \otimes^L f^{''-1}\mu_{\Delta_g}(K_2) $  。

我们有态射链：

\[
\begin{aligned}{L\approx}&{{}\;{}^{t}q_{13}^{^{\prime}-1}{}^{t}j_{*}^{\prime}\Biggl((r_{12}^{-1}\mu_{\varDelta_{f}}(K_{1})\otimes r_{34}^{-1}\mu_{\varDelta_{g}}(K_{2})\Biggr)}\\ {\approx}&{{}\;{}^{t}q_{13}^{\prime-1}{}^{t}j_{*}^{\prime\prime}j_{\pi}^{-1}\Biggl(\tilde{p}_{12}^{-1}\mu_{\varDelta_{f}}(K_{1})\otimes\tilde{p}_{34}^{-1}\mu_{\varDelta_{g}}(K_{2})\Biggr)}\\ {\to}&{{}\;{}^{t}q_{13}^{\prime-1}{}^{t}j_{*}^{\prime\prime}j_{\pi}^{-1}\mu_{\varDelta_{f}\times\varDelta_{g}}\Biggl(\tilde{q}_{12}^{-1}K_{1}\otimes\tilde{q}_{34}^{-1}K_{2}\Biggr)}\\ {\to}&{{}\;{}^{t}q_{13}^{\prime-1}\mu_{\varDelta_{f}\times\gamma\varDelta_{g}}\Biggl(j^{-1}\Biggl(\tilde{q}_{12}^{-1}K_{1}\otimes\tilde{q}_{34}^{-1}K_{2}\Biggr)\Biggr)}\\ {\to}&{{}\;{}^{t}q_{13}^{\prime-1}\mu_{\varDelta_{f}\times\gamma\varDelta_{g}}(R\mathcal{H}o m(q_{13}^{-1}q_{2}^{\prime-1}H,q_{13}^{\prime}q_{1}^{\prime\prime}F))}\\ \end{aligned}
\]

\[
\begin{aligned}{}&{{}\to\mu_{\varDelta_{h}}({R q}_{13!}{R}\mathcal{H}{o m}(q_{13}^{-1}q_{2}^{\prime-1}H,q_{13}^{!}q_{1}^{\prime!}F))}\\ {}&{{}\to\mu_{\varDelta_{h}}({R}\mathcal{H}{o m}(q_{2}^{\prime-1}H,q_{1}^{\prime!}F))\;.}\\ \end{aligned}
\]

在这里，我们将命题 4.3.5 应用于映射  $j$  ，然后将命题 4.3.6 应用于映射  $q_{13}$  ，最后将命题 4.3.4 应用于映射  $q_{13}$  。 □

 $ \mu hom(\cdot \leftarrow \cdot) $  的组合也有类似的结果。

\statementlabel{推论 4.4.10}让  $ F_1 $  、  $ F_2 $  和  $ F_3 $  属于  $ \mathbf{D}^b(X) $  。然后我们有一个规范态射：

\[
\stackrel{L}{\mu h o m}(F_{1},F_{2})\otimes\mu h o m(F_{2},F_{3})\to\mu h o m(F_{1},F_{3})~.
\]

现在让  $X$  、  $Y$  、  $Z$  为三个流形。像往常一样，我们用  $q_{ij}$  表示在  $X\times Y\times Z$  上定义的  $(i,j)$  投影，用  $p_{ij}$  表示在  $T^*X\times T^*Y\times T^*Z$  上定义的 $(i,j)$  投影，并用  $p_{ij}^a$  表示  $p_{ij}$  的组成和  $j$  上的对映图。

\statementlabel{命题 4.4.11}对于属于  $\mathbf{D}^b(X \times Y)$  的  $K_1$  和  $F_1$  以及属于  $\mathbf{D}^b(Y \times Z)$  的  $K_2$  和  $F_2$  ，存在规范态射：

\[
\begin{aligned}{R p_{13}^{a}!}&{{}\bigg(p_{12}^{a-1}\mu{\mathcal{h}}o m(K_{1},F_{1})\overset{L}{\otimes}p_{23}^{a-1}\mu{\mathcal{h}}o m(K_{2},F_{2})\bigg)}\\ {}&{{}\to\mu{\mathcal{h}}o m(K_{1}\circ K_{2},F_{1}\circ F_{2})\enspace,}\\ \end{aligned}
\]

其中 $ K_{1} \circ K_{2} $  和 $ F_{1} \circ F_{2} $  由(3.6.2) 定义。

\statementlabel{证明}令  $j: X \times Y \times Z \to X \times Y \times Y \times Z$  为对角线嵌入，并考虑以下交换图：

\textbackslash{}[\textbackslash{}begin\{aligned\}T\textasciicircum{}\{*\}(X\textbackslash{}times Y)\&\textbackslash{}times T\textasciicircum{}\{*\}(Y\textbackslash{}times Z)\textbackslash{}xleftarrow\{p\_\{12\}\textasciicircum{}\{\textbackslash{}alpha\}\textbackslash{}times p\_\{23\}\textasciicircum{}\{\textbackslash{}alpha\}\}T\textasciicircum{}\{*\}X\textbackslash{}times T\textasciicircum{}\{*\}Y\textbackslash{}times T\textasciicircum{}\{*\}X\textbackslash{}xrightarrow\{\textbackslash{}quad\}\textbackslash{}\textbackslash{}\&\textbackslash{}quad\textbackslash{}quad\textbackslash{}quad\textbackslash{}quad\textbackslash{}quad\textbackslash{}quad\textbackslash{}quad\textbackslash{}quad\textbackslash{}quad j\_\{\textbackslash{}alpha\}\textbackslash{}

将命题 4.4.8 (i) 应用于映射  $X \times Y \to Y$  和  $Y \times Z \to Y$  ，我们得到态射：

\[
\begin{aligned}{R^{t}j_{!}^{\prime}j_{\pi}^{-1}\bigg(\mu{\mathcal{h}}o m(K_{1},F_{1})}&{{}^{\underline{{L}}}\boxtimes{\mu}{\mathcal{h}}o m(K_{2},F_{2})\bigg)}\\ {}&{{}\to{\mu}{\mathcal{h}}o m\bigg(q_{12}^{-1}K_{1}\otimes q_{23}^{-1}K_{2},q_{12}^{-1}F_{1}\otimes q_{23}^{-1}F_{2}\bigg).}\\ \end{aligned}
\]

将命题 4.4.6 (ii) 与公式 (4.4.6) 一起应用，我们得到：

\[
\begin{aligned}{}&{{}R q_{13\pi!}{}^{t}q_{13}^{\prime-1}\mu h o m\bigg(q_{12}^{-1}K_{1}\otimes q_{23}^{-1}K_{2},q_{12}^{-1}F_{1}\otimes q_{23}^{-1}F_{2}\bigg)}\\ {}&{{}\quad\to\mu h o m\bigg(R q_{13!}\bigg(q_{12}^{-1}K_{1}\otimes q_{23}^{-1}K_{2}\bigg),R q_{13!}\bigg(q_{12}^{-1}F_{1}\otimes q_{23}^{-1}F_{2}\bigg)\bigg)\;.}\\ \end{aligned}
\]

结合(4.4.16)和(4.4.17)，我们得到态射：

\[
\begin{aligned}{}&{{}R{q}_{13\pi!}{^{t} q}_{13}^{\prime-1}R^{\prime}{j}_{!}^{\prime}{j}_{\pi}^{-1}\bigg(\mu{\tilde{h}o m}(K_{1},F_{1})\overset{L}{\boxtimes}\mu{\tilde{h}o m}(K_{2},F_{2})\bigg)}\\ {}&{{}\quad\to\mu{\tilde{h}o m}(K_{1}\circ K_{2},F_{1}\circ F_{2})\enspace.}\\ \end{aligned}
\]

由于(4.4.18)的左侧与(4.4.14)的左侧同构，根据交换图(4.4.15)，结果如下。   $ \Box $ 

\subsection{习题}
\statementlabel{练习 IV.1}设 X 为 n 维流形，M 为余维 l 的闭子流形。

(i) 通过选择  $ \widetilde{X}_M $  上的局部坐标  $ (x', x''', t) $  和  $ X $  上的  $ (y', y'') $  使得  $ p(x', x'', t) = (y', y'') $  ，表明  $ dt/t^1 $  是相对形式  $ \Omega_{\widetilde{X}}^{(n+1)} \otimes_{p^{-1}} \mathcal{Q}_X^{(n) \otimes^{-1}} $  束的明确定义的部分，遍及  $ \widetilde{X}_M $  。

(ii) 表明  $ dt/t^1 $  本质上是在  $ \tilde{X}_M $  上定义的

(iii) 从 (i) 和 (ii) 推断相对方向束  $ \sigma_{\tilde{X}/X} $  与  $ A_{\tilde{X}} $  是规范同构的。

\statementlabel{练习 IV.2}令  $M$  为  $X$  的闭子流形， $S$  为  $X$  的闭子集。证明存在一个函子  $\nu: \mathbf{D}^{b}(X \setminus S) \to \mathbf{D}^{b}(T_{M}X \setminus C_{M}(S))$  使得下图可交换（直至同构）：

\[
\begin{array}{c c}{\mathbb{D}^{b}(X)}&{\xrightarrow{\quad\nu_{M}\quad}\quad\mathbb{D}^{b}(T_{M}X)}\\ {\quad\bigg\downarrow\quad}&{\quad\bigg\downarrow\quad}\\ {\mathbb{D}^{b}(X\setminus S)}&{\xrightarrow{\quad\nu\quad}\quad\mathbb{D}^{b}(T_{M}X\setminus C_{M}(S))~.}\end{array}
\]

\statementlabel{练习 IV.3}令  $M$  为  $X$  的闭子流形，并令  $S$  为  $X$  的局部闭子集。证明 $C_M(S) = \mathrm{supp} v_M(A_S)$ 。

\statementlabel{练习 四.4}令 $X$  为流形，令 $F$  和 $G$  属于 $\mathbf{D}^{b}(X)$  ，并假设 $F$  和 $G$  是上同调可构造的。证明：

\[
\mu h o m({\cal F},{\cal G})\simeq\mu h o m({\cal D}{\cal G},{\cal D}{\cal F})^{a}
\]

（参见命题 3.4.6）。

\statementlabel{练习 IV.5}令 $ E \to Z $  为向量丛， $ F \in \mathrm{Ob}(\mathbf{D}_{R^+}^b(E)) $  。证明  $ \nu_Z(F) \simeq F $  和  $ \mu_Z \simeq F^\wedge $  ，其中  $ Z $  与  $ E $  的零部分标识。

\statementlabel{练习 IV.6}让  $F_2$  和  $F_1$  属于  $\mathbf{D}^b(X)$  。

令 $(t;\tau)$  为 $T^*\mathbb{R}$  上的坐标。证明同构：

\[
\begin{aligned}{\mu h o m(F_{2},F_{1})}&{{}\simeq\mu h o m\bigg(\left.F_{2}\boxtimes A_{\{0\}},F_{1}\boxtimes A_{\{0\}}\right)\bigg|_{t=0,\tau=1},}\\ {\mu h o m(F_{2},F_{1})}&{{}\simeq\mu h o m\bigg(\left.F_{2}\boxtimes A_{\mathbb{R}},F_{1}\boxtimes A_{\mathbb{R}}\right)\bigg|_{t=0,\tau=0},}\\ \end{aligned}
\]

（提示：使用命题 4.4.7）。

\statementlabel{练习 IV.7}令  $M$  为  $X$  的闭子流形， $Z$  为  $X$  的闭子集，并令  $F$  和  $G$  属于  $\mathbf{D}^{b}(X)$  。构造以下自然态射并举例说明它们不是同构：

\[
A_{C_{M}(Z)}\to\nu_{M}(A_{Z})\enspace,
\]

\[
\nu_{M}R\mathcal{H}o m(G,F)\to R\mathcal{H}o m(\nu_{M}(G),\nu_{M}(F))\enspace,
\]

\[
\nu_{M}R\varGamma_{Z}(F)\to R\varGamma_{C_{M}(Z)}(\nu_{M}(F))\enspace.
\]

\subsection{注记}
炸毁奇点是几何学中常见的操作，并且“极坐标”的使用是经典的。然而，我们在 §1 中执行的子流形 M 的法向变形并不常见，因为我们获得的流形包含 M 的整个法向向量丛，而不仅仅是球体或射影丛。它是代数几何中最近引入的运算的真实情况的模拟（参见 Fulton [1]）。

我们在定义 4.1.1 中讨论的“正常视锥细胞”的概念与经典概念一致，并且为许多作者所熟悉。 Thom [1] 和 Whitney [1,2] 的名字自然地附在它上面。

专业化和微局部化函子  $\nu$  和  $\mu$  是由 Sato [2] 于 1969 年引入的（参见 Sato-Kawai-Kashiwara [1]）。我们的  $\nu$  构造与 Sato 的构造略有不同，因为我们使用  $M$  的正常变形。在代数情况下，Verdier [5] 进行了类似的构造，他将其与著名的“消失循环”函子联系起来（参见第八章第 6 章）。

当微局部化函子出现时，它首先被认为是定义微函数和解析波前集的工具。事实上，在 Sato 的工作之后不久，Hörmander [1] 引入了 $ C^{\infty} $  波前分布集和傅立叶积分算子，这是线性偏微分方程领域的激烈活动的开始。然后，观察余切丛中奇点的想法（当然这个想法起源于十九世纪）传播到数学的其他领域，例如费曼积分、解析几何、群表示或束理论，这构成了现在所谓的“微局域分析”。

§2 和 §3 的结果基本上包含在 Sato-Kawai-Kashiwara [1] 中。 Kashiwara-Schapira [3] 引入了§4 的函子 $ \mu h o m $  作为 $ \ell $  模的函子 $ \mathcal{H}o m $  的层理论版本（参见第XI 章§4）。

\section{层的微支撑}
\subsection{概要}
在流形  $X$  上，我们可以自然地将  $\mathbf{D}^{b}(X)$  的任何对象  $F$  关联到  $T^{*}X$  的闭圆锥子集， $F$  的微支撑，用  $\mathrm{SS}(F)$  表示。粗略地说， $\mathrm{SS}(F)$ 描述了 $X$ 的同向集合，其中 $F$ “不传播”，我们将在后续章节中证明 $\mathrm{SS}(F)$ 是 $T^{*}X$ 的对合子集。

这个概念还给出了层理论中各种函子的交换性条件；例如 $ f^1F $ 何时与 $ \omega_{Y/X} \otimes f^{-1}F $ 同构或 $ f^{-1}R\mathcal{H}om(F_1, F_2) $ 何时与 $ R\mathcal{H}om(f^{-1}F_1, f^{-1}F_2) $ 同构？

我们首先证明了SS(F)的三个定义的等价性，并使用微支撑给出了一个准则，即给定 $X$ 和 $\Omega_0 \subset \Omega_1$ 的两个开子集 $\Omega_0$ 和 $\Omega_1$ ，从 $R\Gamma(\Omega_1; F)$ 到 $R\Gamma(\Omega_0; F)$ 的限制态射是同构。第三章中介绍的 $\gamma$  与闭合凸真圆锥 $\gamma$  相关的拓扑是微支撑研究的自然工具。事实上，如果  $X$  是仿射的，而  $\phi_\gamma: X \to X_\gamma$  是削弱  $X$  拓扑的映射，则  $\phi_\gamma^{-1}R\phi_\gamma$  在以下意义上扮演截止函子的角色： SS( $\phi_\gamma^{-1}R\phi_{\gamma*}$  F) 包含在  $X \times \gamma^{o*}$  中，态射  $\phi_\gamma^{-1}R\phi_{\gamma*}$  F  $\to F$  是“一个 $X \times \operatorname{Int}\gamma^{o*}$  上的同构”。微局部同构的概念在这里定义，并将在下一章中展开。

在给出了一些微支撑的例子之后，我们研究了微支撑在层上的各种操作方面的行为：张量积和 Hom、正像或逆像、傅里叶佐藤变换。

在本章中，我们总是假设态射对于微支撑是正确的或“非特征的”。这一限制将在第六章中取消。我们在这里和下一章中获得的结果与带有解析系数的偏微分方程理论中的许多经典结果非常相似，我们将在第十一章中展示如何从微支撑理论中推导出其中一些经典结果。

我们经常非常严格地遵循 Kashiwara-Schapira [3] 的阐述，但有些证明实际上被简化了。

我们保持惯例4.0。

\subsection{微支撑的等价定义}
令 $X$  为流形。如果 $\alpha$  是整数或 $\alpha = +\infty$ ，我们将考虑类 $C^\alpha$  ( $1 \leqslant \alpha \leqslant \infty$ ) 的 $X$  上的函数。如果 $X$  是实解析函数，则 $C^\omega$  类的函数表示实解析函数。

像往常一样， $ \pi: T^*X \to X $  表示余切丛。

\statementlabel{命题 5.1.1}假设  $X$  在  $a$  向量空间  $E$  中是开集，令  $p = (x_o; \xi_o) \in T^*X$  并让  $F \in \mathrm{Ob}(\mathbf{D}^b(X))$  。那么以下条件是等效的，其中  $1 \leqslant \alpha \leqslant \infty$  或  $\alpha = \omega$  。

(1)ₐ 存在  $p$  的开邻域 U，使得对于在  $x_1$  邻域中定义的任何  $x_1 \in X$  和类  $C^\alpha$  的任何实函数  $\psi$  ，以及  $\psi(x_1) = 0$  、  $\mathrm{d}\psi(x_1) \in U$  ，我们有：

\[
({\cal R}\Gamma_{\{x;\psi(x)\geq0\}}(F))_{x_{1}}=0.
\]

(2)  $ E $  与  $ 0 \in \gamma $  和  $ F' \in \mathrm{Ob}(\mathbb{D}^b(E)) $  中存在真闭凸锥  $ \gamma $  使得：

（一） $ \gamma \setminus \{0\} \subset \{v; \langle v, \xi_o \rangle < 0\}, (i.e. \xi_o \in \mathrm{Int}(\gamma^{\circ a})), $ 

(b)  $ F'|_{v} \simeq F|_{v} $  为 X 中  $ x_{\circ} $  的邻域 U。

\[
\begin{array}{r}{R\phi_{y\ast}F^{\prime}=0\;(c f.\;\mathrm{I I I}\S5).}\end{array}
\]

(3) 存在 $x_{\circ}$  的邻域 $U$ 、 $\varepsilon > 0$  和一个真闭凸锥 $\gamma$ ，其中 $0 \in \gamma$  满足(2) 的条件(a)，因此如果我们设置：

\[
H=\left\{x;\left\langle x-x_{o},\xi_{o}\right\rangle\geqslant-\varepsilon\right\},
\]

\[
L=\left\{x;\left\langle x-x_{\circ},\xi_{\circ}\right\rangle=-\varepsilon\right\},
\]

然后 $ H \cap (U + \gamma) \subset X $ ，我们就有了自然的同构：

\[
R\varGamma(H\cap(x+\gamma);F)\xrightarrow{\simeq}R\varGamma(L\cap(x+\gamma);F)
\]

对于每个 $ x \in U $  。

\statementlabel{证明}如果  $ \xi_{\circ} = 0 $  ，这三个条件相当于  $ x_{\circ} $  邻域中的 F = 0。假设 $ \xi_{\circ} \neq 0 $ 。

(2) $\Rightarrow(1)_1$ 。我们可以假设  $X = E$  和  $F = F'$  。我们设置 $U = X \times \operatorname{Int}\gamma^{\circ a}$ 。对于任何具有  $\psi(x_1) = 0$  、  $\mathrm{d}\psi(x_1) \in U$  的函数  $\psi$  ，存在  $\gamma$  开集  $\Omega$  使得  $\Omega$  和  $\{x; \psi(x) < 0\}$  在  $x_1$  的邻域上重合。此外，我们可以假设，当 $\Omega'$ 在 $x_1$ 的 $\gamma$ 邻域系统上运行时， $\Omega' \setminus \Omega$ 在通常的拓扑中形成 $X \setminus \Omega$ 中 $x_1$ 的邻域系统。由此我们得到：

\[
\left(R\varGamma_{\{x;\;\psi(x)\geqslant0\}}(F)\right)_{x_{1}}=\left(R\varGamma_{(X_{\gamma}\setminus\varOmega_{\gamma})}(R\phi_{\gamma*}F)\right)_{x_{1}}=0{~.}
\]

(3) $\Rightarrow(2)$ 。我们可以假设  $X$  是  $\gamma$  -open 和  $U \subset H \setminus L$  。令  $\Omega_0$  和  $\Omega_1$  为两个  $\gamma$  开集，使得  $\Omega_0 \subset \Omega_1$  、  $x_0 \in \mathrm{Int}(\Omega_1 \setminus \Omega_0)$  、  $\Omega_1 \setminus \Omega_0 \subset \subset U$  。我们有：

\[
R\varGamma((x+\gamma)\cap H;F)\simeq R\varGamma(x+\gamma;F_{H})
\]

类似地，将  $H$  替换为  $L$  。从 (5.1.1) 中我们得到每个  $x \in U$  的  $R\Gamma(x + \gamma; F_{H\setminus L}) = 0$  ，因此：

\[
(\phi_{\gamma}^{-1}R\phi_{\gamma\ast}F_{H\setminus L})|_{U}=0.
\]

应用命题 3.5 3. 我们得到。

\[
\begin{aligned}{{R}\phi_{\gamma\ast}{R}\varGamma_{\varOmega_{1}\setminus\varOmega_{0}}(F_{H\setminus L})}&{{}\simeq{R}\varGamma_{\varOmega_{1\gamma}\setminus\varOmega_{0\gamma}}{R}\phi_{\gamma\ast}F_{H\setminus L}}\\ {}&{{}\simeq{R}\varGamma_{\varOmega_{1\gamma}\setminus\varOmega_{0\gamma}}{R}\phi_{\gamma\ast}\phi_{\gamma}^{-1}{R}\phi_{\gamma\ast}F_{H\setminus L}}\\ {}&{{}\simeq{R}\phi_{\gamma\ast}{R}\varGamma_{\varOmega_{1}\setminus\varOmega_{0}}\phi_{\gamma}^{-1}{R}\phi_{\gamma\ast}F_{H\setminus L}}\\ {}&{{}=0.}\\ \end{aligned}
\]

因此 $ R\Gamma_{\Omega_1\setminus\Omega_0}(F_{H\setminus L}) $  具有所需的属性。

(1) $ \omega \Rightarrow (3) $ 。我们可以假设  $ \xi_{\circ} = (1, 0, \ldots, 0) $  、  $ x_{\circ} = 0 $  和  $ F \in \mathrm{Ob}(\mathbb{D}^{b}(E)) $  。

设置 $ x = (x_1, x') $ ， $ x' = (x_2, \ldots, x_n) $ 。我们将  $ \varepsilon > 0 $  取得足够小，并按 (3) 中的方式定义 H 和 L。我们将  $ \delta > 0 $  取得足够小并设置：

\[
\gamma=\left\{x;x_{1}\leq-\delta|x^{\prime}|\right\}.
\]

我们选择  $\varepsilon$  、  $\delta$  和  $\gamma$  开子集  $\Omega_1$  ，这样  $0 \in \Omega_1$  和  $(\Omega_1 \cap H) \times \mathrm{Int} \gamma^{\circ a} \subset \mathbb{R}^+ \cdot U$  。然后我们将展示同构：

\[
R\varGamma(H\cap(x+\gamma);F)\simeq R\varGamma(L\cap(x+\gamma);F)\enspace,
\]

对于每个 $ x \in \Omega_1 \cap H $  。

对于每个  $a \in \Omega_1 \cap (H \setminus L)$  ，我们可以构造一个  $X$  的开放子集族  $\{\Omega_t(a)\}_{t \in \mathbb{R}^+}$  ，这样：

\[
\Omega_{t}(a)\subset a+\operatorname{I n t}\gamma~,
\]

\[
\Omega_{t}(a)\cap L=(a+\mathrm{I n t}\gamma)\cap L,
\]

\[
\Omega_{t}(a)=\bigcup_{r<t}\Omega_{r}(a),
\]

(iv)  $\Omega_{t}(a)$  具有实解析平滑边界  $\partial\Omega_{t}(a)$ 

(v)  $ Z_t(a) = \left( \bigcap_{s > t} \overline{\Omega_s(a) \setminus \Omega_t(a)} \right) \cap H $  包含在  $ \partial \Omega_t(a) $  中， $ \Omega_t(a) $  在  $ Z_t(a) $  处的外部共形包含在  $ U $  中，

\[
\left(\bigcup_{t>0}\Omega_{t}(a)\right)\cap H=(a+\operatorname{I n t}\gamma)\cap H,
\]

\[
\left(\bigcap_{t>0}\Omega_{t}(a)\right)\cap H=(a+\operatorname{I n t}\gamma)\cap L.
\]

例如，让：

\[
\psi_{t}(x^{\prime})=\delta^{2}|x^{\prime}-a^{\prime}|^{2}+\frac{(\delta^{2}|x^{\prime}-a^{\prime}|^{2}-(\varepsilon+a_{1})^{2})^{2}}{\sqrt{t+(\delta^{2}|x^{\prime}-a^{\prime}|^{2}-(\varepsilon+a_{1})^{2})^{2}}}\quad.
\]

那么我们可以采取：

\[
\Omega_{t}(a)=\left\{x;x_{1}<a_{1}\mathrm{a n d}(x_{1}-a_{1})^{2}>\psi_{t}(x^{\prime})\right\}
\]

\begin{center}
\includegraphics[width=0.29\textwidth]{流行上的层_Chn-assets/img_in_image_box_411_471_761_871.png}
\end{center}

\begin{center}图。 5.1\end{center}

现在我们设置：

\[
\tilde{\Omega}_{t}(a)=\Omega_{t}(a)\cup((a+\mathrm{I n t}\gamma)\backslash H),
\]

我们将证明同构：

\[
R\varGamma(a+\mathrm{I n t}\gamma;F)\xrightarrow{\sim}R\varGamma(\widetilde{\varOmega}_{t}(a);F)\enspace,
\]

对于每个 t > 0。

令  $j$  为嵌入  $a + \operatorname{Int}\gamma \hookrightarrow X$  ，并应用非特征变形引理（命题 2.7.2）。足以表明：

\[
\left(R\varGamma_{X\setminus\tilde{\Omega}_{t}(a)}(R j_{*}j^{-1}F)\right)_{y}=0
\]

对于任何 $y \in Z_s(a) \setminus \tilde{\Omega}_t(a)$  ，任何 $s \leq t$  。由于  $Z_s(a) \setminus \tilde{\Omega}_t(a) \subset Z_t(a)$  ，我们可以假设  $y \in Z_t(a) \setminus \tilde{\Omega}_t(a)$  。如果  $y$  属于  $H \setminus L$  则 (5.1.8) 由假设 (1)  $_\omega$  得出。现在假设  $y \in Z_t(a) \cap L = L \cap \partial(a + \mathrm{Int}\gamma)$  。由于  $a + \mathrm{Int}\gamma$  和  $\Omega_t(a)$  在  $y$  邻域内具有平滑的实解析边界，我们可以通过 (1)  $_\omega$  得到：

\[
\left(R\varGamma_{X\setminus\varOmega_{t}(a)}(F)\right)_{y}=\left(R\varGamma_{X\setminus(a+\mathtt{I n t}\setminus\gamma)}(F)\right)_{y}=0\mathrm{~,~}
\]

因此：

\[
\begin{array}{r}{(R\varGamma_{X\setminus\varOmega_{t}(a)}(R j_{*}j^{-1}F))_{y}=0\enspace.}\end{array}
\]

另一方面，  $ (a + \mathrm{Int}\gamma)\backslash\Omega_t(a) $  是  $ (a + \mathrm{Int}\gamma)\backslash\tilde{\Omega}_t(a) $  和  $ (a + \mathrm{Int}\gamma)\backslash(\Omega_t(a)\cup H) $  的不相交并集。因此 $ (R\Gamma_X\backslash\tilde{\Omega}_t(a)(Rj_*)j^{-1}F)) $  是 $ R\Gamma_X\backslash\Omega_t(a)(Rj_*)j^{-1}F) $  的直接加数，因此(5.1.10) 蕴含(5.1.8)。这显示了所有  $ t $  的 (5.1.7) 。

由(5.1.7)可得同构：

\[
\begin{aligned}{R\varGamma((a+\operatorname{l n t}\gamma)\cap H;F|_{\mathcal{H}})}&{{}\simeq R\varGamma(\tilde{\Omega}_{t}(a)\cap H;F|_{\mathcal{H}})}\\ {}&{{}\simeq R\varGamma(\Omega_{t}(a)\cap H;F|_{\mathcal{H}})~.}\\ \end{aligned}
\]

最后我们将展示任何  $ x \in \Omega_1 \cap H $  的同构：

\[
R\varGamma((x+\gamma)\cap H;F|_{\cal H})\simeq R\varGamma((x+\gamma)\cap L;F|_{\cal H})\enspace,
\]

由此立即推导出(5.1.4)。

让 $ v = (1, 0, \ldots, 0) $  。然后，族 $ \{(x + \rho v + \text{Int}\gamma) \cap H\}_{\rho > 0} $ 在 $ H $ 中形成 $ (x + \gamma) \cap H $ 邻域系统，族 $ \{\Omega_t(x + \rho v) \cap L\}_{\rho > 0, t > 0} $ 在 $ H $ 中形成 $ (x + \gamma) \cap L $ 邻域系统。因此，(5.1.12) 由(5.1.11) 得出。   $ \square $ 

\statementlabel{定义 5.1.2}设 X 为流形。

(i) 让 $ F \in \mathrm{Ob}(\mathbb{D}^b(X)) $  。  $ F $  的微支撑由  $ \mathrm{SS}(F) $  表示，是  $ T^*X $  的子集，定义如下：

命题5.1.1的 $ p \notin \mathrm{SS}(F) \Leftrightarrow \mathrm{condition}(1)_1 $ 得到满足。

(ii) 令 $u: F \to F'$  为 $\mathbf{D}^b(X)$  中的态射，并令 $A$  为 $T^*X$  的子集。如果  $u$  嵌入到与  $\mathrm{SS}(F'') \cap A = \emptyset$  不同的三角形  $F \xrightarrow[u]{} F' \longrightarrow F'' \xrightarrow[+1]{}$  中，我们就说  $u$  是  $A$  的同构。

请注意， $u$  是  $p$  的同构当且仅当存在  $p$  的邻域  $U$ ，使得对于任何  $x_{1}$  和任何函数  $\psi$ ，如命题 5.1.1 的 (1)  $_{\alpha}$  中，态射：

\[
\left(R\varGamma_{\{x;\;\psi(x)\geq0\}}(F)\right)_{x_{1}}\to\left(R\varGamma_{\{x;\;\psi(x)\geq0\}}(F^{\prime})\right)_{x_{1}}
\]

是同构。

从定义中立即推导出以下属性。

\statementlabel{命题 5.1.3}(i) 让 $ F \in \text{Ob}(\mathbb{D}^b(X)) $  。那么  $ \text{SS}(F) $  是  $ T^*X $  和  $ \text{SS}(F) \cap T_X^*X = \text{supp}(F) $  的闭圆锥曲线子集。

（二） $ \mathrm{SS}(F) = \mathrm{SS}(F[1]) $ 。

(iii) 令 $ F_1 \to F_2 \to F_3 \xrightarrow[+1]{} $  为 $ \mathbf{D}^b(X) $  中的一个显着三角形。然后对于 $ i,j,k\in{1,2,3} $ ：

\[
\begin{array}{r}{\mathrm{S S}(F_{i})\subset\mathrm{S S}(F_{j})\cup\mathrm{S S}(F_{k})\quad f o r\quad j\neq k,}\end{array}
\]

\[
(\mathbb{S}\mathbb{S}(F_{i})\backslash\mathbb{S}\mathbb{S}(F_{j}))\cup(\mathbb{S}\mathbb{S}(F_{j})\backslash\mathbb{S}\mathbb{S}(F_{i}))\subset\mathbb{S}\mathbb{S}(F_{k})\quad{f o r\quad}\;k\not=\mathbf{\Sigma}^{j}\mathbf{i},\mathbf{\Sigma}^{j}\mathbf{j}
\]

我们有时将 (iii) 中的性质称为微支撑的三角不等式。我们将在下面第六章证明微支撑总是 $ T^{*}X $ 的对合子集。

\statementlabel{备注 5.1.4}让 $ F \in \mathrm{Ob}(\mathbf{D}^b(X)) $  。包含 $ \mathrm{SS}(H^j(F)) \subset \mathrm{SS}(F) $ 一般来说是不正确的，但我们有 $ \mathrm{SS}(F) \subset \bigcup_j \mathrm{SS}(H^j(F)) $ （参见练习V.4和V.6）。

\statementlabel{备注 5.1.5}让  $ \phi_x $  成为健忘函子  $ \mathbf{D}^b(A_X) \to \mathbf{D}^b(\mathbb{Z}_X) $  并让  $ F \in \mathrm{Ob}(\mathbf{D}^b(A_X)) $  。然后：

\[
\mathrm{SS}(\boldsymbol{F})=\mathrm{SS}(\boldsymbol{\phi}_{X}(\boldsymbol{F}))~.
\]

用 $ \underline{\text{other}} $ 的话来说，SS(F)不依赖于基环A。

\statementlabel{备注 5.1.6}根据命题5.1.1， $F$ 的微支撑仅取决于 $X$ 的 $C^{1}$ 结构。

\subsection{传播}
令 $E$  为实数有限维向量空间， $X$  为 $E$  的开子集，并令 $F \in \mathrm{Ob}(\mathbb{D}^b(X))$  。

\statementlabel{命题 5.2.1}令  $U$  为  $X$  的开子集，并令  $\gamma$  为  $E$  与  $\gamma \ni 0$  中的闭真凸锥。令  $\Omega_0$  和  $\Omega_1$  为  $\gamma$  - $E$  的开子集。我们假设

\[
\mathrm{S S}(F)\cap(U\times\mathrm{I n t}\gamma^{\circ a})=\varnothing{,}
\]

\[
\Omega_{0}\subset\Omega_{1}\qquad{a n d}\qquad\Omega_{1}\backslash\Omega_{0}\subset U~,
\]

\[
{F o r~a n y~}x\in\Omega_{1},(x+\gamma)\backslash\Omega_{0}{~i s~c o m p a c t}.
\]

然后我们有

\[
R\varGamma(\varOmega_{1}\cap X;F)\to R\varGamma(\varOmega_{0}\cap X;F)\quad{i s~a n~i s o m o r p h i s m},
\]

and

\[
R\phi_{\gamma*}(R\varGamma_{X\setminus\varOmega_{0}}(F))|_{\varOmega_{1}}=0\mathrm{~,~}
\]

其中  $ \phi_{\gamma} $  是映射  $ X \to X_{\gamma} $  。

我们将从下面的特殊情况开始证明命题，然后推导出一般情况。

\statementlabel{引理 5.2.2}当  $\Omega_0$  对于某些  $\xi_0 \in \mathrm{Int} \gamma^{\circ a}$  具有  $\{x \in E; \langle x, \xi_0 \rangle < c\}$  形式时，命题 5.2.1 成立，并且

\[
\mathsf{S S}(F)\cap(U\times(\gamma^{\circ a}\backslash\{0\}))=\varnothing.
\]

\statementlabel{证明}我们赋予 $E$ 一个欧几里德结构，并让 $d(\cdot,\cdot)$ 表示距离函数。我们设定：

\[
\begin{aligned}W_{t}&=\left\{(s_{1},s_{2})\in\mathbb{R}^{2};s_{1}<t\right\}\cup\left\{(s_{1},s_{2})\in\mathbb{R}^{2};s_{2}<0\right\}\\&\quad\cup\left\{(s_{1},s_{2})\in\mathbb{R}^{2},s_{1}<2t,s_{2}<t,(s_{1}-2t)^{2}+(s_{2}-\mathfrak{t})^{2}>t^{2}\right\}\end{aligned}
\]

对于 0 < t。

我们设置  $ l(x) = \langle x, \xi_{\circ} \rangle - c $  、  $ H = \{x; l(x) \geq 0\} $  ，并定义 0 < t：

\[
U(t,x_{1})=\left\{x\in X;(d(x,x_{1}+\gamma),l(x))\in W_{t}\right\}\;.
\]

\begin{center}
\includegraphics[width=0.34\textwidth]{流行上的层_Chn-assets/img_in_image_box_116_570_528_979.png}
\end{center}

\begin{center}图。 5.2.a\end{center}

\begin{center}
\includegraphics[width=0.24\textwidth]{流行上的层_Chn-assets/img_in_image_box_742_570_1039_979.png}
\end{center}

\begin{center}图。 5.2.b\end{center}

那么  $U(t,x_1) \supset X \backslash H$  、  $\overline{U(t,x_1)} \subset \Omega_1$  if  $0 < t \ll 1$  和  $\{U(t,x_1) \cap H\}_{t>0}$  形成  $(x_1 + \gamma) \cap H$  的邻域系统。现在，请注意，函数 $d(x, x_1 + \gamma)$  属于 $X \backslash (x_1 + \gamma)$  上的 $C^1$  类，其微分属于该集合上的 $\gamma^{\circ a}$ 。因此  $U(t, x_1)$  有一个  $C^1$  边界，其外共法线属于  $\gamma^{\circ a}$  。因此  $(R\Gamma_{(X \setminus U(t, x_1)))}(F))_x = 0$  ，为  $x \in U \cap \partial U(t, x_1)$  。

我们选择足够小的 $\varepsilon$ ，以便 $\overline{U}(\varepsilon, x_1)$ 包含在 $\Omega_1$ 中。然后我们可以将非特征变形引理（命题 2.7.2）应用于  $\{U(t, x_1)\}_{0<t<\varepsilon}$  族。设置 $\widetilde{F} = R\Gamma_X \backslash \Omega_0(F)$ 。我们得到同构：

\[
R\varGamma(U(\varepsilon,x_{1});\widetilde F)\simeq R\varGamma(U(t,x_{1});\widetilde F)\enspace.
\]

通过取 (5.2.9) 项和  $ t > 0 $  的归纳极限的上同调，我们得到：

\[
R\varGamma(U(\varepsilon,x_{1});\widetilde{F})\simeq R\varGamma(x_{1}+\gamma;\widetilde{F})\enspace.
\]

我们将使用这种同构来证明对于每个  $ x \in \Omega_1 $  、  $ R\Gamma(x + \gamma; \tilde{F}) = 0 $  ，这将完成证明。选择  $ v \in \text{Int}\gamma $  并设置  $ x_t = x + tv $  。让：

\[
I=\left\{t\geqslant0;R\varGamma(x_{t}+\gamma;\widetilde{F})=0\right\}.
\]

那么  $I$  不为空，因为  $x_t \in \Omega_0$  对于  $t \gg 0$  。让 $t_\circ = \inf I$  。令 $\varepsilon > 0$  使得 $\overline{U(\varepsilon, x_{t_0})}$  包含在 $\Omega_1$  中，并令 $t > t_\circ$  使得 $t \in I$  且：

\[
x_{t_{0}}+\gamma\subset U(\varepsilon,x_{t})\subset U(\varepsilon,x_{t_{0}}).
\]

同构  $ \mathcal{R}\Gamma(U(\varepsilon,x_t,); \widetilde{F}) \simeq \mathcal{R}\Gamma(x_t, + \gamma; \widetilde{F}) $  通过为零的  $ \mathcal{R}\Gamma(U(\varepsilon,x_t); \widetilde{F}) $  因式分解。因此  $ t_0 $  属于  $ I $  。如果 $ t_0 $ 严格为正，我们发现 $ t $ 与 $ 0 < t < t_0 $ 和 $ \varepsilon > 0 $ 这样：

\[
x_{t}+\gamma\subset U(\varepsilon,x_{t_{0}})\subset U(\varepsilon,x_{t})\subset\overline{U(\varepsilon,x_{t})}\subset\Omega_{1}.
\]

由于同构  $ \mathcal{R}\Gamma(U(\varepsilon,x_t);\widetilde{F}) \to \mathcal{R}\Gamma(x_t + \gamma;\widetilde{F}) $  通过为零的  $ \mathcal{R}\Gamma(U(\varepsilon,x_{t_0});\widetilde{F}) $  因式分解，因此  $ t $  属于  $ I $  。这样 $ t_o = 0 $ 就完成了引理5.2.2的证明。 ⑨

命题 5.2.1 的证明结束。我们将命题 5.2.1 简化为引理 5.2.2。由于(5.2.4) 由(5.2.5) 得出，因此只要证明(5.2.5) 就足够了。当  $ \gamma = \{0\} $  、 $ F|_{U} = 0 $  时，命题是显而易见的。因此我们可以假设  $ \gamma \neq \{0\} $  。令  $ \gamma' $  为闭合真凸锥，使得  $ \operatorname{Int} \gamma' \supset \gamma \setminus \{0\} $  ，并令  $ \Omega $  为  $ \Omega_1 \setminus \Omega_0 $  的子集，使得  $ \Omega $  相对于  $ \Omega_1 \setminus \Omega_0 $  和  $ \Omega \subset \subset U $  上的  $ \gamma' $  拓扑是开的。然后这样的  $ \Omega $  形成  $ \Omega_1 \setminus \Omega_0 $  上的  $ \gamma $  拓扑的开集基础。因为对于这样的  $ \Omega $  和  $ \operatorname{Int} \gamma^{\circ a} \supset \gamma'^{\circ a} \setminus \{0\} $  来说，显示  $ R\Gamma(\Omega; R\Gamma_{X \setminus \Omega_0}(F)) = 0 $  就足够了，我们可以从一开始就假设 (5.2.6)。此外我们可以假设

\[
\Omega_{1}\backslash\Omega_{0}\subset\subset U~.
\]

由非特征变形引理（命题2.7.2），为了看到（5.2.5），足以证明

\[
R\varGamma_{H_{c}}R\varGamma_{\varOmega_{1}\setminus\varOmega_{0}}(F)_{\partial H_{c}}=0\qquad\mathrm{f o r~a n y~}c\enspace,
\]

我们设置了任意点  $ \xi_0 \in \operatorname{Int}\gamma^{\circ a} $  ：

\[
H_{c}=\left\{x;\left\langle x,\xi_{\circ}\right\rangle\geqslant c\right\}\qquad\mathrm{a n d}\qquad\partial H_{c}=\left\{x;\left\langle x,\xi_{\circ}\right\rangle=c\right\}\;.
\]

对于任何  $ x_o \in \partial H_c \cap U $  ，都存在  $ \gamma $  开子集  $ \Omega \ni x_o $  ，使得  $ \Omega \cap H_c \subset U $  。然后根据引理 5.2.2， $ R\phi_{\gamma*}(R\Gamma_{\Omega \cap H_c}(F)) = 0 $  。然后（5.2.12）从以下事实得出：当 $ W $ 穿过 $ \gamma $ 拓扑中的 $ x_o $ 邻域系统时， $ W \cap H_c $ 在 $ H_c $ 中形成 $ x_o $ 的邻域系统。这就完成了命题5.2.1的证明。   $ \square $ 

现在我们将考虑 $X = Y \times E$  的情况，其中 $E$  是有限维实向量空间， $Y$  是流形。令 $\gamma$  为 $E$  中的封闭凸锥， $0 \in \gamma$  ，（我们不要求 $\gamma$  是正确的）。我们设置  $X_\gamma = Y \times E_\gamma$  并用  $\phi_\gamma$  表示连续映射  $X \to X_\gamma$  。

我们将使用 $ \gamma $ 拓扑来切断层的微支撑。

让 $ F \in \mathrm{Ob}(\mathbb{D}^{b}(X)) $  。

\statementlabel{命题 5.2.3（微局部截止引理）}(i) 当且仅当态射 $ \phi_{\gamma}^{-1}R\phi_{\gamma^*}F \to F $  是同构时，微支撑SS(F) 包含在 $ T^*Y \times (E \times \gamma^{\circ a}) $  中。

(ii) 态射  $ \phi_{\gamma}^{-1}R\phi_{\gamma*}F \to F $  是  $ T^*Y \times (E \times \mathrm{Int}\gamma^{\circ a}) $  的同构。

\statementlabel{证明}我们可以假设 $Y$ 仿射。将 $\gamma$ 替换为 $\{0\} \times \gamma$ ，我们可以从一开始就假定 $X = E$ 。假设首先 $\phi_{\gamma}^{-1} R \phi_{\gamma^*} F \simeq F$ 。我们将证明包含 $\mathrm{SS}(F) \subset X \times \gamma^{\circ a}$ 。让 $\xi_0 \notin \gamma^{\circ a}$  。选择适当的凸闭锥  $\gamma'$  使得  $\gamma' \setminus \{0\} \subset \{v; \langle v, \xi_0 \rangle < 0\}$  和  $\gamma' + \gamma = E$  。那么对于 $X$ 的任何 $\gamma'$  -开非空凸子集 $\Omega$ ，我们有：

\[
R\varGamma(\varOmega;F)\simeq R\varGamma(\varOmega+\gamma;F)\simeq R\varGamma(X;F)\enspace.
\]

这意味着对于任何一对凸  $ \gamma' $  开子集  $ (\Omega_0, \Omega_1) $  和  $ \Omega_0 \subset \Omega_1 $  ，我们有：

\[
R\phi_{\gamma^{\prime}*}{\cal R}\varGamma_{\Omega_{1}\setminus\Omega_{0}}(F)=0\ ,
\]

因此 $ \mathrm{SS}(F) \cap (X \times \{\xi_o\}) = \emptyset $  。

相反，假设  $\mathrm{SS}(F) \subset X \times \gamma^{\circ a}$  。为了证明 $\phi_{\gamma}^{-1} R\phi_{\gamma*}F \simeq F$ 的同构，只要证明对于 $X$ 的任何相对紧开凸子集 $\Omega$ ，从 $R\Gamma(\Omega + \gamma; F)$ 到 $R\Gamma(\Omega; F)$ 的限制态射是同构即可。令 $\mathfrak{S}$  为 $(\Omega + \gamma)$  的开凸子集 $V$  的集合，使得 $V \supset \Omega$  和 $R\Gamma(V; F) \to R\Gamma(\Omega; F)$  是同构。根据命题 2.7.1， $\mathfrak{S}$  是归纳排序的。令  $V$  为  $\mathfrak{S}$  的最大元素。我们将通过矛盾来证明 $V = \Omega + \gamma$ 。如果  $V \neq \Omega + \gamma$  ，则存在  $x_{\circ} \in (\Omega + \gamma) \setminus \overline{V}$  。

\statementlabel{引理 5.2.4}令  $V$  为  $X$  的凸开子集，并令  $x_{\circ} \in X$  。令  $\lambda$  为由  $V$  生成的、顶点位于  $x_{\circ}$  的圆锥体，令  $\lambda'$  为以原点为顶点的闭圆锥体  $\overline{\lambda} - x_{\circ}$ ，令  $V_{1}$  为  $V \cup \{x_{\circ}\}$  凸包的内部。然后 $V_{1} \setminus V$  对于 $\lambda$  拓扑是局部封闭的。

引理 5.2.4 的证明。我们有：

\[
V_{1}=\left\{(1-t)x_{\circ}+t u;u\in V,0<t\leqslant1\right\}=\left\{(1-t)x_{\circ}+t u;u\in V,0<t<1\right\}~.
\]

设定：

\[
V_{2}=\left\{(1-t)x_{\circ}+t u;u\in V,t>0\right\}~,
\]

\[
V_{3}=\left\{(1-t)x_{\circ}+t u;u\in V,t>1\right\}.
\]

那么  $V_2$  和  $V_3$  是  $\lambda$  -open，并且足以证明  $V_1\backslash V = V_2\backslash V_3$  。由于我们有  $V \subset V_1 \subset V_2$  、  $V \subset V_3$  、  $V_2 \subset V_1 \cup V_3$  ，因此足以证明  $V_1 \cap V_3$  包含在  $V$  中。对于 $x \in V_1 \cap V_3$ ，我们可以写@@M​​ATH10@@， $0 < t \leqslant 1$ ， $1 < s$ 和 $u$ ， $v \in V$ 。然后我们有  $(s - t)x = (s - 1)tu + (1 - t)sv$  ，因此  $x \in V$  。   $\square$ 

命题 5.2.3 的证明结束。让 $V_1, \lambda$  和 $\lambda'$  如引理5.2.4 中所示。那么  $\lambda'$  是一个真凸锥，而  $V_1 \setminus V$  对于  $\lambda$  拓扑来说是局部封闭的，由前面的引理可知。另一方面，我们有一些  $v \in \gamma^a$  的 Int  $\lambda' \ni v$  。因此 $\lambda' \cap \gamma^a \subset \{0\}$  。根据命题 5.2.1，这意味着：

\[
(R\phi_{\lambda^{\prime}{*}}R\varGamma_{V_{1}\setminus V}(F))|_{V_{1}}=0\enspace.
\]

因此我们得到 $ V_1 \in \mathfrak{S} $ ，这是一个矛盾。

为了证明(ii)，设 $ \phi_{\gamma}^{-1}R\phi_{\gamma*}F \longrightarrow F \longrightarrow F' \xrightarrow{+1} be $ 为一个可区分的三角形。应用函子  $ R\phi_{\gamma*} $  我们得到  $ R\phi_{\gamma*}F' = 0 $  ，这意味着  $ \mathrm{SS}(F') \cap (X \times \mathrm{Int}\ \gamma^{\circ a}) = \emptyset $  ，根据命题 5.1.1。请注意，当  $ \gamma $  不正确时，  $ \mathrm{Int}\ \gamma^{\circ a} = \emptyset $  。

\statementlabel{备注 5.2.5}微支撑的切断将在VI§1中再次讨论。

\subsection{例子：与局部闭子集相关的微支撑}
回想一下，如果  $Z$  是  $X$  的局部闭子集，则  $X$  上的束  $A_Z$  是  $X\setminus Z$  上的零束，以及  $Z$  上带有茎  $A$  的常束。

令  $E$  为有限维实向量空间， $\gamma$  为顶点位于  $0$  的封闭凸锥。回想一下，我们为任何  $v \in \gamma\}$  设置了  $\gamma^a = -\gamma$  和  $\gamma^\circ = \{\xi \in E^*; \langle v, \xi \rangle \geq 0 \}$  。

\statementlabel{命题 5.3.1}其中一个有：

\[
\mathrm{SS}(A_{\gamma})\cap\pi^{-1}(0)=\gamma^{\circ}.
\]

\statementlabel{证明}$1\in\Gamma(E;A_E)$  节定义了  $1_\gamma\in\Gamma(E;A_\gamma)$  节，其支持是  $\gamma$  。因此  $\gamma^\circ$  包含在  $\pi^{-1}(0)\cap\mathrm{SS}(A_\gamma)$  中。另一方面，态射 $\phi_{\gamma^\alpha}^{-1}\circ R\phi_{\gamma^\alpha*}A_\gamma\to A_\gamma$  是同构。因此，结果来自微局域截止引理（命题 5.2.3）。   $\square$ 

\statementlabel{命题 5.3.2}设 X 为流形，M 为闭子流形。然后：

\[
\mathrm{SS}(A_{M})=T_{M}^{*}X\enspace.
\]

\statementlabel{证明}由于问题是  $M$  上的局部问题，因此我们可以假设  $X$  是向量空间， $M$  是向量子空间。那么这是前面结果的一个特例。   $\Box$ 

请特别注意  $ \mathrm{SS}(A_X) = T_X^* X $  ，  $ T^* X $  的零部分。

\statementlabel{命题 5.3.3}令  $ \varphi $  为类  $ C^1 $  的实函数，并假设  $ \mathrm{d}\varphi \neq 0 $  在集合  $ \{x; \varphi(x) = 0\} $  上。然后：

\[
\mathbf{S S}(A_{\{x;\phi(x)\geq0\}})=\left\{(x;\lambda\mathbf{d}\varphi(x));\lambda\varphi(x)=0,\lambda\geq0,\varphi(x)\geq0\right\}\;,
\]

\[
\mathbf{S S}(A_{\{x;\phi(x)>0\}})=\left\{(x;\lambda\mathbf{d}\phi(x));\lambda\varphi(x)=0,\lambda\leqslant0,\varphi(x)\geqslant0\right\}
\]

\statementlabel{证明}(i) 根据命题 5.3.1 选择一个局部坐标系，使得  $ \{x;\varphi(x)\geq0\} $  是一个封闭的半空间。

(ii) 通过考虑精确序列  $ 0 \to A_{\{x; \phi(x) > 0\}} \to A_X \to A_{\{x; \phi(x) \leq 0\}} \to 0 $  从 (i) 和三角不等式（命题 5.1.3）得出。   $ \square $ 

现在我们将讨论 $ \mathbb{R}^2 $  上的一些具体示例。我们用 $ (x_1, x_2) $  表示 $ \mathbb{R}^2 $  上的坐标，用 $ (\xi_1, \xi_2) $  表示对偶坐标。

示例 5.3.4。让 $ Z = \{x \in \mathbb{R}^2; x_1 > 0, -x_1^{3/2} \leq x_2 < x_1^{3/2}\} $  。然后：

\[
\left\{\begin{aligned}{\mathsf{S S}(A_{Z})=(T_{X}^{*}X\cap\bar{Z})\cup\{(x;\xi);\xi_{2}>0,x_{2}=-(2\xi_{1}/3\xi_{2})^{3},}\\ {x_{1}=(2\xi_{1}/3\xi_{2})^{2}\}.}\\ \end{aligned}\right.
\]

事实上，为  $x_{1}\geqslant0$  设置  $\varphi_{\pm}(x)=x_{2}\pm x_{1}^{3/2}$  ，为  $x_{1}\leqslant0$  、  $Z_{i}=\{x;\varphi_{i}(x)\geqslant0\},(i=\pm)$  设置  $\varphi_{\pm}(x)=x_{2}$  。使用命题 5.1.1 可以轻松检查，对于  $i=\pm$  ，  $\mathrm{SS}(A_{Z_{i}})$  是集合的闭包： $\{(x;\lambda\mathrm{d}\varphi_{i}(x));\lambda\varphi_{i}(x)=0,\lambda\geqslant0,\varphi_{i}(x)\geqslant0\}$  。那么 (5.3.1) 由命题 5.1.3 得出，因为  $Z=Z_{+}\backslash Z_{-}$  。

示例 5.3.5。让 $ Z = \{x \in \mathbb{R}^2; x_1 x_2 \geq 0\} $  。然后：

\[
\pi^{-1}(0)\cap\mathbf{S}\mathbf{S}(A_{Z})=\left\{\xi;\xi_{1}\xi_{2}\leqslant0\right\}\enspace.
\]

事实上，让  $Z^{\pm} = \{x \in Z; \pm x_1 \geq 0, \pm x_2 \geq 0\}$  。我们有一个精确的顺序；   $0 \to A_Z \to A_{Z^+} \oplus A_{Z^-} \to A_{\{0\}} \to 0$ 。根据命题 5.3.1 和三角不等式（命题 5.1.3）的结果，我们得到包含  $\pi^{-1}(0) \cap SS(A_Z) \supset \{\xi; \xi_1 \xi_2 \leq 0\}$  。为了证明逆包含，用  $\gamma$  表示圆锥  $Z^-$  。然后 $\phi_{\gamma}^{-1} R \phi_{\gamma^*} A_{Z^+} \simeq \phi_{\gamma}^{-1} R \phi_{\gamma^*} A_{\{0\}} \simeq A_{Z^+}$ ，因此 $\phi_{\gamma}^{-1} R \phi_{\gamma^*} A_Z \simeq \phi_{\gamma}^{-1} R \phi_{\gamma^*} A_{Z^-}$ 。从 $\operatorname{Int} \gamma^{\circ a} \cap SS(A_{Z^-}) = \emptyset$ 开始，我们通过命题5.2.3得到 $\operatorname{Int} \gamma^{\circ a} \cap SS(A_Z) = \emptyset$ 。将 $\gamma$  替换为 $\gamma^a$  我们得到(5.3.2)。

一般来说，很难精确计算 SS  $ (A_{z}) $  。为了给这个集合一个界限，我们引入：

\statementlabel{定义 5.3.6}令 S 为 X 的子集，其中 X 为流形。我们设定：

\[
N_{x}(S)=T_{x}X\backslash C_{x}(X\backslash S,S),
\]

\[
N_{x}^{*}(S)=(N_{x}(S))^{\circ},
\]

\[
N(S)=\bigcup_{x\in X}N_{x}(S),\qquad N^{*}(S)=\bigcup_{x\in X}N_{x}^{*}(S).
\]

我们将 $N(S)$  称为 $S$  的严格法锥体，将 $N^{*}(S)$  称为 $S$  的共形锥体。

（回想一下，正常圆锥  $C(X\backslash S,S)$  已在 IV §1 中定义。）

根据命题 4.1.2，非零向量  $ \theta \in T_x X $  属于  $ N_x(S) $  当且仅当在  $ x $  邻域的局部图中，存在包含  $ \theta $  的开锥体  $ \gamma $  和  $ x $  的邻域 U，使得：

\[
U\cap((S\cap U)+\gamma)\subset S.
\]

特别地，  $N(S)$  是  $TX$  的开凸锥体，它满足：

\[
\begin{aligned}{N_{x}(S)=T_{x}X}&{{}\Leftrightarrow C_{x}(X\backslash S,S)=\varnothing}\\ {}&{{}\Leftrightarrow x\notin\bar{S}\qquad\mathrm{o r}\qquad x\in\mathrm{I n t}S.}\\ \end{aligned}
\]

此外：

\[
N_{x}(S)=\emptyset\Leftrightarrow N_{x}^{*}(S)=T_{x}^{*}X\enspace,
\]

 $ N_x(S) \neq \emptyset \Leftrightarrow N_x^*(S) $  是一个真闭凸锥。

另请注意：

\[
N_{x}(X\backslash S)=N_{x}(S)^{a},
\]

\[
N_{x}^{*}(X\backslash S)=N_{x}^{*}(S)^{a}~.
\]

\statementlabel{命题 5.3.7}假设  $X$  是仿射的，并让  $\gamma$  是与  $\operatorname{Int}\gamma \neq \emptyset$  的真闭凸锥。令  $\Omega$  （分别为  $Z$  ）为  $X$  的开子集（分别为闭子集）。让 $x \in X$  。

(i) 如果 $ N_x^*(\Omega) \subset \text{Int} \gamma^\circ \cup \{0\} $ （或 $ N_x^*(Z)^a \subset \text{Int} \gamma^\circ \cup \{0\} $ ），则 $ \Omega $ （或 $ Z $ ）在 $ x $ 的邻域中是 $ \gamma $  -开（或 $ \gamma $  -闭）。

(ii) 相反，如果 $\Omega$ （或 $Z$ ）在 $x$ 的邻域中是 $\gamma$  -开（或 $\gamma$  -闭），则 $N_{x}^{*}(\Omega) \subset \gamma^{\circ}$ （或 $N_{x}^{*}(Z)^{a} \subset \gamma^{\circ}$ ）。

\statementlabel{证明}由于  $\Omega$  在  $x$  的邻域中是  $\gamma$  -开，当且仅当  $X\setminus\Omega$  在  $x$  和  $N^{*}(\Omega)=N^{*}(X\setminus\Omega)^{a}$  的邻域中是  $\gamma$  - 闭，因此足以证明  $\Omega$  的结果。

(i) 对于每个  $ v \in N_x(\Omega) $  、  $ v \neq 0 $  ，存在一个凸开锥  $ \gamma_v $  使得  $ \Omega $  是  $ \gamma_v $  -在  $ x $  的邻域内开（参见 (5.3.3)）。我们可以用有限数量的这样的锥体  $ \gamma_v $  覆盖  $ \gamma \setminus \{0\} $  ，结果如下。

(ii) 由(5.3.3)得出。 ⑨

\statementlabel{命题 5.3.8}设 X 为流形， $ \Omega $  为开子集，Z 为闭子集。然后：

\[
\mathsf{S S}(A_{\mathcal{Q}})\subset N^{*}(\mathcal{Q})^{a}~,
\]

\[
\mathsf{S S}(A_{Z})\subset N^{*}(\mathsf{Z})\enspace.
\]

\statementlabel{证明}我们可以假设  $X$  仿射，并且通过将三角不等式应用于  $0 \to A_{X \setminus Z} \to A_{X} \to A_{Z} \to 0$  足以证明  $\Omega$  的结果。

设  $ x \in X $  并假设  $ N_x^*(\Omega) \neq T_x^*X $  。令  $ \gamma $  为闭凸真圆锥，使得：

\[
N_{x}^{*}(\Omega)\subset\mathsf{Int}\gamma^{\circ}\cup\{0\}~.
\]

根据前面的命题， $\Omega$  在 $x$  的邻域中是 $\gamma$  -开，并且我们可以假设 $\Omega$  是 $\gamma$  -开。然后 $A_{\Omega} = \phi_{\gamma}^{-1} R \phi_{\gamma*} A_{\Omega}$ ，因此 $\mathrm{SS}(A_{\Omega}) \subset X \times \gamma^{\mathrm{oa}}(\mathrm{Propositiion} 5.2.3)$ 。由于该包含对于任何满足 (5.3.4) 的  $\gamma$  成立，因此我们得到了结果。   $\square$ 

\subsection{微支撑的函子性质}
在本节中，我们研究微支撑在几种操作下的行为，例如适当的正像、非特征情况下的逆像等。这里获得的一些结果将在后续章节中推广。

令 $X$  和 $Y$  为两个流形。我们像往常一样用 $q_1$ 和 $q_2$ （分别是 $p_1$ 和 $p_2$ ）表示在​​ $X \times Y$ （分别是 $T^*X \times T^*Y$ ）上定义的第一和第二投影。如果 $f$ 是从 $Y$ 到 $X$ 的映射，我们在第四章中用 $f'$ 和 $f_\pi$ 表示映射：

\[
T^{*}Y\xleftarrow[{}^{t f^{\prime}}]{}Y\underset{X}{\times}T^{*}X\xrightarrow[{f_{\pi}}]{}T^{*}X\enspace.
\]

\statementlabel{命题 5.4.1}让  $ F \in \mathrm{Ob}(\mathbb{D}^{b}(X)) $  和  $ G \in \mathrm{Ob}(\mathbb{D}^{b}(Y)) $  。然后：

\[
\mathrm{SS}\left(\stackrel{L}{F}\boxtimes G\right)\subset\mathrm{SS}(F)\times\mathrm{SS}(G)\enspace.
\]

\statementlabel{证明}我们可以假设 $X$ 和 $Y$ 是向量空间。让  $(x_o, y_o; \xi_o, \eta_o) \in T^*(X \times Y)$  ，并假设例如  $(x_o; \xi_o) \notin \mathrm{SS}(F)$  。取 $H, L, U, \gamma$ 满足命题5.1.1(3)中的条件作为 $X$ 上的层 $F$ 并设置 $\tilde{\gamma} = \gamma \times \{0\}$  。对于 $z = (x, y) \in U \times Y$ ，我们有：

\[
\begin{aligned}{R\varGamma\bigg(H\times Y\cap(z+\widetilde{\gamma});q_{1}^{-1}F\otimes{q_{2}^{-1}G}\bigg)}&{{}\simeq R\varGamma\bigg(H\cap(x+\gamma);F\otimes{G_{y}}\bigg)}\\ {}&{{}\simeq R\varGamma(H\cap(x+\gamma);F)\otimes{G_{y}}^{L}.}\\ \end{aligned}
\]

事实上， $H \cap (x + \gamma)$  是紧凑的，我们可以应用命题 2.6.6。

由于我们有一个与 (5.4.1) 类似的公式，其中  $H$  替换为  $L$  ，因此结果来自命题 5.1.1。 ⑨

\statementlabel{命题 5.4.2}让  $ F \in \mathrm{Ob}(\mathbb{D}^{b}(X)) $  和  $ G \in \mathrm{Ob}(\mathbb{D}^{b}(Y)) $  。然后：

\[
\mathrm{S S}(R\mathcal{H o m}(q_{2}^{-1}G,q_{1}^{-1}F))\subset\mathrm{S S}(F)\times\mathrm{S S}(G)^{a}.
\]

\statementlabel{证明}将  $ q_{1}^{-1}F $  替换为  $ q_{1}^{!}F $  就足以证明类似的说法。

我们可以假设 $X$ 和 $Y$ 是向量空间。让 $(x_o, y_o; \xi_o, -\eta_o) \notin \mathrm{SS}(F) \times \mathrm{SS}(G)^a$  。

首先假设  $ (x_o; \xi_o) \notin \mathrm{SS}(F) $  。这种情况下的证明与命题 5.4.1 类似。取 $ H, L, U, \gamma $ 满足命题5.1.1(3)中的条件作为 $ X $ 上的层 $ F $ 并设置 $ \tilde{\gamma} = \gamma \times \{0\} $  。对于 $ z = (x, y) \in U \times Y $ ，我们根据命题 3.1.15 得到：

\[
\begin{aligned}{H^{j}}&{{}(R\varGamma(H\times Y\cap(z+\tilde{\gamma});R\mathcal{H}o m(q_{2}^{-1}G,q_{1}^{!}F)))}\\ {}&{{}\simeq\varliminf_{\overrightarrow{w}}H^{j}(R\varGamma((H\cap(x+\gamma))\times W;R\mathcal{H}o m(q_{2}^{-1}G,q_{1}^{!}F)))}\\ {}&{{}\simeq\varliminf_{\overrightarrow{w}}H^{j}(R\mathbf{H o m}(R\varGamma_{c}(W,G),R\varGamma(H\cap(x+\gamma);F))),}\\ \end{aligned}
\]

其中 W 的范围穿过 y 的开邻域族。由于我们有一个类似的公式，其中 H 被 L 替换，因此在这种情况下，证明由命题 5.1.1 得出。

现在假设  $(y_{\circ};\eta_{\circ})\notin\mathrm{SS}(G)$  。对于 $Y$  上的束 $G$ ，取一个满足命题5.1.1 的条件(2) 的圆锥 $\gamma$ 。我们可以假设 $R\phi_{\gamma*}G=0$ 。

\statementlabel{引理 5.4.3}设 Y 为有限维实向量空间， $ \gamma $  为闭凸真锥， $ \gamma \ni 0 $  ， $ \Omega $  为 Y 的  $ \gamma^a $  开子集，使得对于 Y 的任何紧 K， $ \Omega \cap (K + \gamma) $  是相对紧的，并令  $ \Omega' $  为  $ \Omega $  的  $ \gamma^a $  开子集，使得  $ \Omega \setminus \Omega' $  相对紧在 Y 中，设  $ G \in \text{Ob}(\mathbb{D}^b(Y)) $  ，并假设  $ R\phi_{\gamma^*}G = 0 $  。然后：

（一） $ R\phi_{y*}G_{\Omega}=0 $ ，

（二） $ R\Gamma_c(\Omega'; G) \simeq R\Gamma_c(\Omega; G) $ 。

引理 5.4.3 的证明。令  $ \Omega_1 $  为  $ Y $  的  $ \gamma $  开子集并假设  $ \Omega_1 \cap \Omega \subset \subset Y $  。我们有：

\[
H^{j}(\Omega_{1};G_{\Omega})\simeq\varinjlim_{\vec{K}}H_{K}^{j}(\Omega_{1};G),
\]

其中 K 的范围涵盖  $\Omega$  中包含的  $\Omega_{1}$  的闭子集族。

对于这样的  $K, K' = \bar{K} + \gamma^{a}$  在  $Y, \bar{K}$  中闭合，表示  $K$  在  $Y$  中的闭合。然后我们有：

\[
K^{\prime}\cap\Omega_{1}\subset\Omega\enspace.
\]

事实上，如果  $ x = y + v \in \Omega_1 $  ，与  $ y \in \overline{K} $  ，  $ v \in \gamma^a $  ，那么  $ y = x - v $  属于  $ (\Omega_1 + \gamma) \cap \overline{K} = \Omega_1 \cap \overline{K} = K $  。因此 $ x \in K + \gamma^a \subset \Omega + \gamma^a = \Omega $  。

从  $ K \cap \Omega_1 \subset K' \cap \Omega_1 $  开始，我们得到 (5.4.3)：

\[
H^{j}(\Omega_{1};G_{\Omega})\simeq\varinjlim_{K^{\prime}}H_{K^{\prime}\cap\Omega_{1}}^{j}(\Omega_{1};G),
\]

其中  $ K' $  范围涵盖  $ \gamma $  闭子集系列，使得  $ K \cap \Omega_1 \subset \Omega $  。然后  $ R\Gamma_{K' \cap \Omega_1}(\Omega_1; G) \simeq R\Gamma_{K' \cap \Omega_1}(\Omega_1; R\phi_{\gamma^*}G) = 0 $  ，证明 (i)。

为了证明第二个断言，请注意  $ R\Gamma_c(\Omega, G) \simeq R\Gamma_c(Y, G_\Omega) $  和类似的公式，其中  $ \Omega $  替换为  $ \Omega' $  。那么就足以证明  $ R\Gamma_c(Y; G_{\Omega \setminus \Omega^*}) = 0 $  ，并且  $ \Omega \setminus \Omega' $  相对紧凑，足以证明  $ R\Gamma(Y; G_{\Omega \setminus \Omega^*}) = 0 $  。但是 $ R\Gamma(Y; G_\Omega) \simeq R\Gamma(Y; G_\Omega) = 0 $  (i)。   $ \square $ 

命题 5.4.2 的证明结束。通过前面的引理和命题3.1.15，我们发现对于 $X$ 中的任何开集 $W$ 以及 $Y$ 的任何一对 $\gamma^a$ 开子集 $\Omega$ 和 $\Omega'$ ，使得 $\Omega' \subset \Omega$ ， $\Omega \setminus \Omega' \subset \subset Y$ ，我们有：

\[
R\varGamma(W\times\varOmega;R\mathcal{H}o m(q_{2}^{-1}G,q_{1}^{!}F))\simeq R\varGamma(W\times\varOmega^{\prime};R\mathcal{H}o m(q_{2}^{-1}G,q_{1}^{!}F))\enspace.
\]

那么命题 5.1.1 的条件 (2) 满足层  $ R\mathcal{H}om(q_2^{-1}G, q_1^!F) $  和锥  $ \{0\} \times \gamma^a $  。   $ \square $ 

\statementlabel{命题 5.4.4}令  $f: Y \to X$  为流形  $G \in \mathrm{Ob}(\mathbb{D}^b(Y))$  的态射，并假设  $f$  在  $\mathrm{supp}(G)$  上是正确的。然后：

\[
\mathrm{S S}(R f_{*}G)\subset f_{\pi}({}^{t}f^{\prime}{}^{-1}(\mathrm{S S}(G))).
\]

此外，如果 f 是封闭嵌入，则上面的包含是等式。

\statementlabel{证明}令 $ x \in X $  和 $ \varphi $  成为 $ X $  上的真实 $ C^1 $  函数，使得 $ \varphi(x) = 0 $  和 $ \mathrm{d}(\varphi \circ f)(y) \notin \mathrm{SS}(G) $  对于所有 $ y \in f^{-1}(x) $  。因此我们有 $ R\Gamma_{\{\varphi \circ f \geq 0\}}(G)|_{f^{-1}(x)} = 0 $  和

\[
\begin{aligned}(R\varGamma_{\{\varphi\geq0\}}(Rf_{*}G))_{x}&\simeq(Rf_{*}R\varGamma_{\{\varphi\circ f\geq0\}}(G))_{x}\\&\simeq R\varGamma(f^{-1}(x);R\varGamma_{\{\varphi\circ f\geq0\}}(G))\\&=0.\end{aligned}
\]

（这里，  $\{\varphi \geq 0\}$  表示集合  $\{x \in X; \varphi(x) \geq 0\}$  ， $\{f \circ \varphi \geq 0\}$  也类似。\}）这证明了第一个断言。现在假设 $f$  是封闭式沉浸式体验。我们可以假设 $X$ 是一个向量空间， $Y$ 是一个向量子空间。假设 $p \notin SS(R f_* G)$  ，并令 $H, L, \gamma, U$  对于 $X$  上的层 $R f_* G$  满足命题5.1.1 的条件(3)。对于 $x \in U \cap Y$ ，我们有：

\[
R\varGamma(H\cap(x+\gamma);R f_{*}G)\xrightarrow{\sim}R\varGamma(L\cap(x+\gamma);R f_{*}G)\enspace.
\]

另一方面我们有：

\[
R\varGamma(H\cap(x+\gamma);R f_{*}G)\simeq R\varGamma((Y\cap H)\cap(x+(Y\cap\gamma));G)
\]

将H换成L，类似的公式成立。则由命题5.1.1得出结果。 ⑨

\statementlabel{命题 5.4.5}假设 $f: Y \to X$  是光滑的。

(i) 让 $ F \in \mathrm{Ob}(\mathbb{D}^{b}(X)) $  。然后：

\[
\mathsf{S S}(f^{-1}F)={}^{t}f^{\prime}(f_{\pi}^{-1}(\mathsf{S S}(F))).
\]

(ii) 让 $ G \in \text{Ob}(\mathbb{D}^b(Y)) $  。那么下面的条件是等价的：

(a) 所有  $ H^{j}(G) $  都是 f 纤维上的局部常值层。

(b) 在 Y 本地，存在  $ F \in \mathrm{Ob}(\mathbb{D}^b(X)) $  ，使得  $ G \simeq f^{-1} F $  ，

（c） $ \mathrm{SS}(G) \subset t f'(Y \times_X T^* X) $ 。

\statementlabel{证明}(i) 首先我们证明  $\mathrm{SS}(f^{-1}F)$  包含在  $f'(f_{\pi}^{-1}(\mathrm{SS}(F)))$  中。我们可以假设  $Y = \mathbb{R}^n \times \mathbb{R}^l$  、  $X = \mathbb{R}^n$  、  $f$  是投影  $(x, y) \mapsto x$  。让 $p = (x_o, y_o; \xi_o, \eta_o) \notin f'(f_{\pi}^{-1}(\mathrm{SS}(F)))$  。如果  $\eta_o \neq 0$  我们选择  $v \in \mathbb{R}^l$  这样  $\langle v, \eta_o \rangle < 0$  并设置  $\gamma = \{(0, tv); t \geq 0\}$  。对于任意  $\varepsilon > 0$  ，设置  $H = \{(x, y); \langle y, \eta_o \rangle \geq -\varepsilon\}$  、  $L = H \setminus \mathrm{Int} H$  。对于 $z \in \mathrm{Int} H$ ，我们有：

\[
R\varGamma(H\cap(z+\gamma);f^{-1}F)=R\varGamma(L\cap(z+\gamma);f^{-1}F),
\]

因为 $ f^{-1}F|_{(z+R(0,v))} $  是一个常值层。在这种情况下这证明了 $ p \notin \mathrm{SS}(f^{-1}F) $ 。

现在假设 $ \eta_{\circ} = 0 $ ， $ (x_{\circ}; \xi_{\circ}) \notin \mathrm{SS}(F) $ 。我们取  $ H, L, \gamma $  、  $ U $  满足命题 5.1.1 的条件（3），作为  $ X $  上的层  $ F $  ，并设置  $ \tilde{\gamma} = \gamma \times \{0\} $  。那么对于任何  $ z \in f^{-1}(U) $  我们有：

\[
R\varGamma(f^{-1}(H)\cap(z+\tilde{\gamma});f^{-1}F)\simeq R\varGamma(H\cap(f(z)+\gamma);F)
\]

和  $H$  替换为  $L$  的类似同构。因此 $p \notin \mathrm{SS}(f^{-1}F)$ 。

为了证明逆包含，我们使用命题 5.1.1 的条件 (1)，并注意如果  $ \varphi $  是  $ X $  和  $ x \in X $  上的实  $ C^1 $  函数，我们有：

\[
(R\Gamma_{\{\varphi\geqslant0\}}(F))_{x}=(R\Gamma_{\{\varphi\circ f\geqslant0\}}(f^{-1}F))_{y}
\]

对于任何 $ y \in f^{-1}(x) $  。

(ii) 我们可以假设  $ Y = X \times \mathbb{R}^l $  。然后 (a)  $ \Leftrightarrow $  (b) 通过命题 2.7.8，(b)  $ \Rightarrow $  (c) 通过 (i)，并且为了证明 (c)  $ \Rightarrow $  (a)，我们将命题 5.2.3 与  $ \gamma = \mathbb{R}^l $  一起应用。   $ \square $ 

\statementlabel{备注 5.4.6}在命题5.4.5的情况下，令 $ V = t f'(Y \times_X T^* X) $  。这是  $ T^* Y $  的平滑对合子流形。尽管包含  $ \mathrm{SS}(H^j(G)) \subset \mathrm{SS}(G) $  一般而言不成立（参见备注 5.1.4），但对于任何  $ j $  ，我们有等价物：  $ \mathrm{SS}(G) \subset V \Leftrightarrow \mathrm{SS}(H^j(G)) \subset V $  ，这是从命题 5.4.5 得出的。

为了描述层理论中各种操作下微支撑的行为，让我们为  $ T^{*}X $  中的两个锥体 A 和 B 引入符号  $ A + B $  ，通过

\[
(A+B)\cap\pi^{-1}(x)=\left\{a+b;a\in A\cap\pi^{-1}(x),b\in B\cap\pi^{-1}(x)\right\}\;.
\]

\statementlabel{引理 5.4.7}如果  $A$  和  $B$  是  $T^*X$  的闭锥体，并且如果  $A \cap B^a \subset T_X^*X$  ，则  $A + B$  也是  $T^*X$  中的闭锥体。

\statementlabel{证明}令 $(x)$  为 $X$  上的坐标系，并令 $(x; \xi)$  为 $T^*X$  上的关联坐标。该条件意味着  $\xi + \eta \neq 0$  上的  $\{(((x; \xi), (x; \eta)) \in A \times_X B; |\xi| + |\eta| = 1\}$  。因此，我们可以假设 $|\xi + \eta| \geq \varepsilon$ 对于这个集合上的一些 $\varepsilon > 0$ 。因此  $A \times_X B$  包含在  $C = \{(x; \xi, \eta); |\xi + \eta| \geq \varepsilon (|\xi| + |\eta|)\}$  中。由于  $(x; \xi, \eta) \mapsto (x; \xi + \eta)$  给出的  $\mu : C \to T^*X$  是一个真映射，因此  $A + B = \mu(A \times_X B)$  是闭合的。   $\square$ 

\statementlabel{命题 5.4.8}让 $ F \in \mathrm{Ob}(\mathbb{D}^{b}(X)) $  。

(a) 令  $ \Omega $  为  $ X $  的开放子集， $ j $  为开放嵌入  $ \Omega \hookrightarrow X $  。

(i) 假设  $ \mathrm{SS}(F) \cap N^*(\Omega)^a \subset T_X^* X $  。然后：

\[
\mathrm{S S}(R j_{*}j^{-1}F)\subset N^{*}(\varOmega)+\mathrm{S S}(F)\enspace.
\]

(ii) 假设  $ \mathrm{SS}(F) \cap N^*(\Omega) \subset T_X^* X $  。然后：

\[
\mathrm{S S}(R j,j^{-1}F)\subset N^{*}(\Omega)^{a}+\mathrm{S S}(F)\enspace.
\]

(b) 令 Z 为 X 的闭子集。

(i) 如果  $ \mathrm{SS}(F) \cap N^*(Z) \subset T_X^* X $  ，则  $ \mathrm{SS}(R\Gamma_Z(F)) \subset N^*(Z)^a + \mathrm{SS}(F) $  。

(ii) 如果  $ \mathrm{SS}(F) \cap N^*(Z)^a \subset T_X^* X $  ，则  $ \mathrm{SS}(F_Z) \subset N^*(Z) + \mathrm{SS}(F) $  。

\statementlabel{证明}(a) 我们可以假设 $X$  是一个向量空间。取 $x_{\circ} \in X$ ，我们将证明 $x_{\circ}$ 处的断言。

(i) 让  $ \xi_\circ \notin N_x^*(\Omega) + (\mathrm{SS}(F) \cap \pi^{-1}(x_\circ)) $  我们将显示  $ (x_\circ; \xi_\circ) \notin \mathrm{SS}(Rj_* j^{-1}F) $  。由假设：

\[
(N_{x_{0}}^{*}(\Omega)+\overline{{\mathbb{R}^{-}\xi_{0}}})\cap\mathrm{SS}(F)^{a}\subset\{0\}~,
\]

且  $ T_{x}^{*}X $  中存在闭凸真锥 K 使得：

\[
N_{x_{0}}^{*}(\Omega)+\overline{{\mathbb{R}^{-}\xi_{0}}}\subset\mathrm{Int}(K)\cup\{0\}~,
\]

\[
K^{a}\cap(\mathrm{S S}(F)\cap\pi^{-1}(x_{\circ}))\subset\{0\}~.
\]

令  $\gamma$  为  $K$  的极锥。然后  $\gamma \subset \{v; \langle v, \xi_0 \rangle < 0\} \cup \{0\}$  ，根据命题 5.3.7，我们可以假设  $\Omega$  是  $\gamma$  -开。令  $U$  成为  $x_0$  的相对紧凑的开邻域，使得  $(U \times \gamma^{\circ a}) \cap \mathrm{SS}(F) \subset T_X^*X$  。令  $\Omega_0$  和  $\Omega_1$  为两个  $\gamma$  开集，使得  $\Omega_0 \subset \Omega_1 \subset \Omega_0 \cup U$  。应用命题 5.2.1 我们得到：

\[
(R\phi_{\gamma\ast}R\varGamma_{X\setminus\varOmega_{0}}(F))|_{\varOmega_{1}}=0~.
\]

开集  $\Omega$  为  $\gamma$  -开集，  $R j_{*}\,j^{-1}$  与  $R\phi_{\gamma*}$  交换。因此，将 $F$  替换为 $R j_{*}\,j^{-1}F$  后，(5.4.6) 仍然成立，从而完成了命题5.1.1 的证明。

(ii) 证明类似。让 $\xi_\circ \notin N^*_\circ(\Omega)^a + (\mathrm{SS}(F) \cap \pi^{-1}(x_\circ))$  。我们可以找到一个闭凸真锥 $\gamma$ 和 $x_\circ$ 的邻域 $U$ ，使得 $\gamma \subset \{v; \langle v, \xi_\circ \rangle < 0\} \cup \{0\}$ 和 $(U \times \gamma^\circ a) \cap \mathrm{SS}(F) \subset T_X^*X$ ，并且我们可以假设 $\Omega$ 是 $\gamma^a$  -开。应用命题5.2.1，我们可以发现 $\gamma$  -开子集 $\Omega_0$ 和 $\Omega_1$ ，使得 $\Omega_0 \subset \Omega_1 \subset \Omega_0 \cup U$ 、 $\Omega_1 \setminus \Omega_0$ 是 $x_\circ$ 和 $R\phi_{\gamma^*} R \Gamma_{\Omega_1 \setminus \Omega_0}(F) = 0$ 的邻域。

设置 $F' = R\Gamma_{\Omega_1\setminus\Omega_0}(F)$ 。通过将  $\Omega$  替换为另一个与  $x_\circ$  邻域中的  $\Omega$  一致的  $\gamma^a$  开集，我们可以从一开始就假设  $\Omega \cap (K + \gamma)$  对于  $X$  的任何紧凑子集  $K$  在  $X$  中是相对紧凑的。然后我们通过引理 5.4.3 得到  $R\phi_{\gamma*}(F'_\Omega) = 0$  。这意味着命题 5.1.1 中的  $(x_\circ; \xi_\circ) \notin \mathrm{SS}(F'_\Omega)$ ，并且  $F'_\Omega$  在  $x_\circ$  的邻域中与  $F_\Omega = Rj_! j^{-1}F$  同构。

(b) 设置 $ \Omega = X \setminus Z $  ，通过将三角不等式应用于 $ R\Gamma_Z(F) \longrightarrow F \longrightarrow Rj_* j^{-1} F \xrightarrow[+1]{} $  和 $ Rj_! j^{-1} F \longrightarrow F \longrightarrow F_Z \xrightarrow[+1]{} $  从(a) 得出。   $ \square $ 

\statementlabel{推论 5.4.9}令  $Z$  为  $X$  、  $x \in Z$  的闭子集并假设  $N_x^*(Z) \neq T_x^*X$  。让 $F \in \mathrm{Ob}(\mathbb{D}^b(X))$  使得 $\mathrm{SS}(F) \cap N_x^*(Z) \subset \{0\}$  。然后 $(R\Gamma_Z(F))_x = 0$ 。

\statementlabel{证明}我们可以假设 $X$  是仿射的。设置  $F' = R\Gamma_{Z}(F)$  ，我们有：

\[
\mathrm{S S}(F^{\prime})\cap\pi^{-1}(x)\subset(\mathrm{S S}(F)\cap\pi^{-1}(x))+N_{x}^{*}(\mathrm{Z})^{a}\enspace.
\]

由于  $ ((SS(F) \cap \pi^{-1}(x)) + N_x^*(Z)^a) \cap N_x^*(Z) $  包含在  $ \{0\} $  中，因此存在一个闭凸真锥  $ \gamma $  使得  $ \operatorname{Int} \gamma^{\circ a} \cup \{0\} $  包含  $ N_x^*(Z) $  且

\[
({\mathrm{S S}}(F^{\prime})\cap\pi^{-1}(x))\cap\gamma^{\circ a}\subset\{0\}~.
\]

根据命题 5.3.7，我们可以假设 $Z$  是 $\gamma$  闭的。根据命题5.2.1，存在 $\gamma$  -开子集 $\Omega_1 \supset \Omega_0$ ，使得 $R\phi_{\gamma^*} R\Gamma_X \backslash \Omega_0(F')|_{\Omega_1} = 0$  和 $\Omega_1 \backslash \Omega_0$  是 $x$  的相对紧邻域。从 $N_x^*(Z) \subset \mathrm{Int} \gamma^{\circ a} \cup \{0\}$ 开始， $\Omega \cap Z$ 在 $Z$ 中形成 $x$ 的邻域系统，其中 $\Omega$ 贯穿 $\gamma$ 系列 -  $x$ 的开放邻域。因此

\[
\begin{aligned}{H^{j}(F^{\prime})_{x}}&{{}=\varinjlim_{\vec{\Omega}}H^{j}(\varOmega;F^{\prime})}\\ {}&{{}=\varinjlim_{\vec{\Omega}}H^{j}(\varOmega\cap\varOmega_{0};F^{\prime})~,}\\ \end{aligned}
\]

其中  $\Omega$  范围涵盖  $\gamma$  系列 -  $x$  的开放邻域。由于  $\Omega \cap \Omega_0 \cap Z = \varnothing$  对于  $\gamma$  - $x$  的开邻域  $\Omega$  使得  $\Omega \cap Z \subset \Omega_1 \setminus \Omega_0$  ，我们得到了结果。   $\square$ 

\statementlabel{推论 5.4.10}(i) 对于  $ X $  和  $ F \in \mathrm{Ob}(\mathbb{D}^b(X)) $  的闭合子流形  $ M $  ，

\[
\operatorname{s u p p}(\mu_{M}(F))\subset T_{M}^{*}X\cap\mathrm{S S}(F)\enspace.
\]

(ii) 让 $F$  和 $G$  属于 $\mathbf{D}^{b}(X)$  。然后：

\[
\operatorname{supp}\left(\mu\mathrm{hom}(G,F)\right)\subset\mathrm{SS}(G)\cap\mathrm{SS}(F).
\]

\statementlabel{证明}(i) 根据定理 4.3.2(ii)，对于  $p \in T_M^*X\backslash\mathrm{SS}(F)$  和  $H^j(\mu_M(F))_p \simeq \lim_{\overrightarrow{Z}} H_Z^j(F)_{\pi(p)}$  ，其中  $Z$  的范围贯穿此处描述的族。但我们可以进一步假设 $N_{\pi(p)}^*(Z) \subset T_M^*X\backslash\mathrm{SS}(F)$ 。那么推论 5.4.9 意味着  $H_Z^j(F)_{\pi(p)} = 0$  。

(ii) 由 (i) 和命题 5.4.2 得出。 □

我们将在下一章中给 SS( $ \mu_{hom}(G, F) $ ) 一个界限。

\statementlabel{推论 5.4.11}令  $M$  为  $X, F \in \mathrm{Ob}(\mathbb{D}^b(X))$  的闭子流形，并假设  $\mathrm{SS}(F) \cap T_M^* X \subset T_X^* X$  。然后：

（一） $ \mathrm{SS}(F_M) \subset \mathrm{SS}(F) + T_M^* X $ ，

(ii) 自然态射 $ F_M \otimes \omega_{M/X} \to R\Gamma_M(F) $  是同构。

回想一下， $ \omega_{M/X} = \sigma t_{M/X}[\dim M - \dim X] $  是相对对偶复形。

\statementlabel{证明}(i) 通过对 codim  $M$  的归纳论证，我们可以假设  $M$  是一个超曲面。由于结果是  $M$  上的本地结果，我们还可以假设  $M$  将  $X$  分隔成两个开放子集  $\Omega^+$  和  $\Omega^-$  ：  $X = \Omega^- \cup M \cup \Omega^+$  。令  $j_\pm$  为开放嵌入  $\Omega_\pm \hookrightarrow X$  。我们有区分三角形：

\[
R j_{-,!}j_{-,-}^{-1}F\oplus R j_{+,!}j_{+,}^{-1}F\longrightarrow F\longrightarrow F_{M}\mathop{\longrightarrow}_{+1}.
\]

那么结果由命题5.4.8得出。

(ii) 令 $ i: M \subsetneq T_M^* X $  为零部分， $ \dot{\pi}: \dot{T}_M^* X \to M $  为投影。然后我们就有了一个显着的三角形：

\[
i^{\downarrow}\mu_{M}(F)\longrightarrow R\pi_{*}\mu_{M}(F)\longrightarrow R\dot{\pi}_{*}(\mu_{M}(F)|_{\dot{T}_{M}^{*}X})\underset{+1}{\longrightarrow}.
\]

另一方面，根据定理 4.3.2， $ i^!\mu_M(F) = F_M \otimes \omega_{M/X} $  和  $ R\pi_* \mu_M(F) = R\Gamma_M(F) $  。根据推论 5.4.10， \textbackslash{}textit\{supp\}(  $ \mu_M(F) $  )  $ \subset  i(M) $  以及  $ R\dot{\pi}_*(\mu_M(F)|_{\dot{T}_M^*X}) = 0 $  。因此，我们得到(ii)。   $ \square $ 

\statementlabel{定义 5.4.12}令 $f: Y \to X$  为流形的态射， $A$  是 $T^*X$  的闭圆锥子集。如果满足以下条件，我们就说 $f$  不是 $A$  的特征：

\[
f_{\pi}^{-1}(A)\cap T_{Y}^{*}X\subset Y\times\mathop{T_{X}^{*}X}_{X}.
\]

如果  $F \in \mathrm{Ob}(\mathbb{D}^{b}(X))$  ，则表示  $f$  对于  $F$  是非特征的，如果  $f$  对于  $\mathrm{SS}(F)$  是非特征的。如果  $f$  是一种嵌入，那么也会说“Y 是非特征的”而不是“f 是非特征的”。

回想一下，  $ T_Y^*X $  表示  $ f' : Y \times_X T^*X \to T^*Y $  的内核。特别是，如果  $ f $  是平滑的，则  $ f $  对于  $ T^*X $  的任何圆锥子集来说都是非特征的。与引理 5.4.7 的证明类似的论证表明，如果  $ f $  对于  $ A $  来说是非特征的，那么  $ f'(f_\pi^{-1}(A)) $  也是一个闭圆锥曲线子集。

\statementlabel{命题 5.4.13}设  $ F \in \mathrm{Ob}(\mathbb{D}^b(X)) $  ，并假设  $ f: Y \to X $  是非特征的。然后：

（一） $ \mathrm{SS}(f^{-1}F) \subset \mathrm{f}'(f_{\pi}^{-1}(\mathrm{SS}(F))) $ ，

(ii) 自然态射 $ f^{-1}F \otimes \omega_{Y/X} \to f^!F $  是同构。

\statementlabel{证明}我们通过图map来分解f：

\[
Y{\xrightarrow{g}}Y\times X{\xrightarrow{h}}X,\quad f=h\circ g,
\]

其中 g(y) = (y, f(y)) 且 h 是 Y × X 上的第二个投影。那么分别证明 h 和 g 的结果就足够了，我们可以从一开始就假设 f 要么是平滑的，要么是闭合嵌入。那么，平滑情况下的结果由命题 5.4.5 得出，嵌入情况下的结果则由命题 5.4.4 和推论 5.4.11 得出  $ \square $ 

\statementlabel{命题 5.4.14}令 F 和 G 属于  $ \mathbf{D}^{b}(X) $  。

(i) 假设  $ \mathrm{SS}(F) \cap \mathrm{SS}(G)^a \subset T_X^* X $  。然后：

\[
\mathbb{S}\mathbb{S}\left(\overset{L}{F}\otimes G\right)\subset\mathbb{S}\mathbb{S}(F)+\mathbb{S}\mathbb{S}(G)\enspace.
\]

(ii) 假设  $ \mathrm{SS}(F) \cap \mathrm{SS}(G) \subset T_{X}^{*}X $  。然后：

\[
\mathrm{S S}(R{\mathcal{H}}o m(G,F))\subset\mathrm{S S}(F)+\mathrm{S S}(G)^{a}~.
\]

此外，如果 $G$  是上同调可构造的，则 $R\mathcal{H}om(G,A_X)\otimes^L F \to R\mathcal{H}om(G,F)$  是同构。

\statementlabel{证明}回想一下：

\[
{\cal F}\stackrel{L}{\otimes}{\cal G}\simeq\delta_{X}^{-1}\left({\cal F}\triangleleft{\cal G}\right),
\]

\[
R{\mathcal{H}}o m(G,F)\simeq\delta_{X}^{!}R{\mathcal{H}}o m(q_{2}^{-1}G,q_{1}^{!}F)\enspace,
\]

其中 $ \delta_X: X \to X \times X $  是对角线嵌入。因此，结果由命题 5.4.13、5.4.1、5.4.2 和命题 3.4.4 得出。   $ \square $ 

\statementlabel{备注 5.4.15}在命题5.4.8-5.4.13中，我们一直做出“非特征性”假设。这一假设将在第六章中被删除。

示例 5.4.16。 (i) 令 $t$  为 $\mathbb{R}$  、 $f$  上的坐标：  $\mathbb{R} \to \mathbb{R}$  为映射 $t \mapsto t^3$  ，并令 $G = A_R$  。然后  $R f_* G \simeq A_R$  但  $f_x^t f'^{-1}(T_R^* \mathbb{R})$  包含  $T_{\{0\}}^{*}\mathbb{R}$  。因此，命题 5.4.4 中的包含内容并不是一般意义上的等式。然而，如果  $f$  是全纯且有限的，并且如果  $G$  是  $\mathbb{C}$  可构造的（参见 VIII §5），我们就有等式（Kashiwara [5]）。

(ii) 让  $ X = \mathbb{R}^2 $  坐标为  $ (t, y) $  、  $ Y = \{(t, y); t = 0\} $  并让  $ Z = \{(t, y); t > 0, -t < y \leq t\} $  。然后：

\[
\mathrm{SS}(A_{2})\cap\pi^{-1}(0)=\left\{(\tau,\eta);\tau\leq-|\eta|\right\}
\]

其中 $(\tau, \eta)$  表示对偶坐标。因此  $Y$  对于束  $A_Z$  来说是非特征的。自  $(A_Z)_Y = R\varGamma_Y(A_Z) = 0$  以来，命题 5.4.13 中的内容一般来说不是等式。

现在我们将命题 5.4.4 扩展到非真的情况。

莫尔斯理论中众所周知的事实是，如果  $Y$  是紧流形， $\phi$  是  $Y$  上的实数  $C^1$  -函数，并且  $\phi$  在区间  $]t_1,t_2[$  上没有临界值，那么对于每个  $j$  ，上同调群  $H^j(\{x;\varphi(x)<t\};A_x)$  与  $t\in]t_1,t_2]$  都是同构的。在相对情况下，对于  $Y$  上的任何层  $G$  ， $\phi$  没有临界值的假设（即： $\mathrm{d}\varphi(x)\notin T_x^*X$  当  $\varphi(x)\in]t_1,t_2[]$  被  $\mathrm{d}\varphi(x)$  不属于  $\mathrm{SS}(G)$  的假设替换时，类似的结果成立。更准确地说，我们得到以下结果。

令 $f: Y \to X$  为流形的态射，并令 $\varphi$  为 $Y$  上的实 $C^1$  函数。我们设定：

\[
Y_{t}=\left\{y\in Y;\varphi(y)<t\right\}\mathrm{,~}\quad\widetilde{Y}_{t}=\left\{y\in Y;\varphi(y)\leq t\right\}
\]

我们用 $j_t$ （分别 $\tilde{j}_t$ ）表示嵌入 $Y_t \hookrightarrow Y$ （分别 $\tilde{Y}_t \hookrightarrow Y$ ），并用 $f_t$ （分别 $\tilde{f}_t$ ）表示映射 $f|_{Y_t} = f \circ j_t$ （分别 $f \circ \tilde{j}_t$ ）。让 $t_\circ \in \mathbb{R}$  。

\statementlabel{命题 5.4.17}设  $ G \in \text{Ob}(\mathbb{D}^b(Y)) $  ，并假设对于所有  $ t \in \mathbb{R} $  ，  $ \text{supp}(G) \cap \overline{Y}_t $  相对于  $ X $  是正确的。

(i) 假设对于任何  $ y \in Y \setminus \tilde{Y}_{t_0} $  ，我们有：

\[
\mathrm{d}\varphi(y)\notin\left(\mathbb{S S}(G)+{\mathfrak{t}}f^{\prime}\left(Y\times_{\boldsymbol{x}}T^{*}X\right)\right).
\]

然后：

\[
R f_{*}G\simeq R f_{t*}j_{t}^{-1}G\qquad\mathrm{f o r}\qquad t>t_{\circ}~,
\]

\[
R f_{*}{\cal G}\simeq R\tilde{f}_{t*}j_{t}^{-1}{\cal G}\qquad\mathrm{f o r}\qquad t\geq t_{\circ},
\]

\[
\mathbb{S S}(R f_{*}G)\subset f_{\pi}({}^{t}f^{\prime-1}(\mathbb{S S}(G)\cap\pi^{-1}(\widetilde{Y}_{t_{0}})))=f_{\pi}({}_{j}^{t f^{\prime}-1}(\mathbb{S S}(G))).
\]

(ii) 假设对于任何  $ y \in Y \setminus \tilde{Y}_{t_0} $  ，我们有：

\[
-\mathrm{d}\varphi(y)\notin\left(\mathrm{SS}(G)+{\mathrm{\scriptsize~}^{\prime}}f^{\prime}\left(Y\mathop{\times}_{X}T^{*}X\right)\right).
\]

然后：

\[
R f_{!}G\simeq R f_{t!}j_{t}^{-1}G\qquad\mathrm{f o r}\qquad t>t_{\circ}~,
\]

\[
R f_{!}G\simeq R\tilde{f}_{!*}j_{!}^{!}G\qquad{\mathrm{f o r}}\qquad t\geqslant t_{\circ},
\]

\[
\mathbb{S S}(R f_{!}G)\subset f_{\pi}({}^{t}f^{\prime-1}(\mathbb{S S}(G)\cap\pi^{-1}(\widetilde{Y}_{t_{\circ}})))\subset f_{\pi}({}^{t}f^{\prime-1}(\mathbb{S S}(G))).
\]

\statementlabel{证明}令 $ \tilde{f} $  为映射 $ (f, \varphi)\colon Y \to X \times \mathbb{R} $  ，并让 $ q $  和 $ \tilde{\varphi} $  分别表示在 $ X \times \mathbb{R} $  上定义的第一和第二投影。那么  $ \tilde{f} $  适用于  $ \operatorname{supp}(G) $  和  $ f = q \circ \tilde{f} $  。令  $ i_t $  （分别为  $ \tilde{i}_t $  ）表示嵌入  $ X \times \lceil -\infty, t \rceil \hookrightarrow X \times \mathbb{R} $  、（分别为  $ X \times \lceil -\infty, t \rceil \hookrightarrow X \times \mathbb{R} $  ），并设置  $ q_t = q \circ i_t $  、  $ \tilde{q}_t = q \circ \tilde{i}_t $  。然后，设置  $ \tilde{G} = R\tilde{f}_* G $  ，我们有：

\[
\begin{aligned}{R f_{t*}{j}_{t}^{-1}G}&{{}\simeq R q_{t*}{i}_{t}^{-1}\widetilde{G}~,\quad R\widetilde{f}_{t*}\widetilde{j}_{t}^{-1}G\simeq R q_{t*}\widetilde{i}_{t}^{-1}\widetilde{G}~,}\\ {R f_{t!}{j}_{t}^{-1}G}&{{}\simeq R q_{t!}{i}_{t}^{-1}\widetilde{G}~,\quad R\widetilde{f}_{t*}\widetilde{j}_{t}^{!}G\simeq R\widetilde{q}_{t*}\widetilde{i}_{t}^{!}\widetilde{G}~.}\\ \end{aligned}
\]

因此，证明  $ \tilde{G}, q, i, \tilde{\varphi} $  而不是  $ G, f, j, \varphi $  的结果就足够了。

(i) 该假设意味着：

\[
\mathsf{S S}(\widetilde{G})\cap(T^{*}X\times\{(t;\tau);t>t_{\circ}\})\subset T^{*}X\times\{(t;\tau);\tau\leqslant0\}\enspace.
\]

令 $U$  为局部图表中包含的 $X$  的开子集，并且在这样的图表中是凸的。根据命题 5.2.3 和推论 3.5.5，我们对所有  $t > t_0$  都有同构。

\[
R\varGamma(U\times\mathbb{R};\widetilde{G})\xrightarrow{\sim}R\varGamma(U\times]-\infty,t[;\widetilde{G})\enspace.
\]

为了看到这一点，只要显示同构  $ \mathcal{R}\Gamma(U \times \mathbf{]} t_{\circ}, \infty[\mathbf{;} \tilde{G}) \simeq \mathcal{R}\Gamma(U \times \mathbf{]} t_{\circ}, t[\mathbf{;} \tilde{G}) $  就足够了。通过同构 $ \mathbf{]} t_{\circ}, \infty[\mathbf{\simeq R} $ ，我们可以应用上面提到的命题。这就证明了(a)。

由于  $q$  在  $\mathrm{supp}(i_{t*}i_{t}^{-1}\tilde{G})$  上是真映射，我们得到命题 5.4.4：

\[
\mathrm{S S}(R q_{*}{R i}_{t*}{i}_{t}^{-1}\widetilde{G})\subset\{(x;\xi);(x,t^{\prime};\xi,0)\in\mathrm{S S}({R i}_{t*}{i}_{t}^{-1}\widetilde{G})\mathrm{f o r~s o m e}t^{\prime}\}
\]

由于  $ \mathrm{SS}(Ri_{t*}i_{t}^{-1}\tilde{G}) $  包含在  $ \mathrm{SS}(\tilde{G}) + (T_{X}^{*}X \times \{(t;\tau);\tau \leqslant 0\}) $  中，（命题 5.4.8）、（5.4.16）和（5.4.14）意味着  $ t > t_{0} $  ：

\[
\mathrm{S S}(R q_{*}{R i}_{t*}{i}_{t}^{-1}\tilde{G})\subset\{(x;\xi);(x,t^{\prime};\xi,0)\in\mathrm{S S}(\tilde{G})\mathrm{f o r~s o m e}t^{\prime}\}
\]

由此得出 (b)。

要获得  $ (a') $  ，只需注意，对于 X 的任何紧凑子集 K，

\[
H^{k}(K;R\tilde{q}_{t*}\tilde{i}_{t}^{-1}\tilde{G})\xrightarrow{\sim}\varinjlim_{t^{\prime}>t}H^{k}(K;R q_{t^{\prime}*}{i}_{t^{\prime}}^{-1}\tilde{G})\enspace.
\]

(ii) 证明类似，将(5.4.14)中的“ $ \tau \leq 0 $ ”替换为“ $ \tau \geq 0 $ ”，并将(5.4.15)替换为：

\[
R\varGamma_{U\times]-\infty,t[}(U\times\mathbb{R}:\tilde{G})\simeq R\varGamma_{U\times]-\infty,t^{\prime}[}(U\times\mathbb{R}:\tilde{G})
\]

对于任何  $ t, t' $  和  $ t_0 \leq t \leq t' $  。 ⑨

\statementlabel{备注 5.4.18}命题 5.4.17 中的条件意味着 $f$ 在以下意义下是相对非特征的。令 $T^*(Y/X)$ 为相对余切丛，即 $t'f':Y\times_X T^*X\to T^*Y$ 的余核，并令 $p:T^*Y\to T^*(Y/X)$ 为投影。于是命题 5.4.17 的条件可写为：对任意 $y\in Y\setminus Y_{t_0}$，有 $\pm p(d\varphi(y))\notin p(\mathrm{SS}(G))$。

作为命题 5.4.17 的一个特例，我们给出了一个有用的结果。

\statementlabel{推论 5.4.19（微局域莫尔斯引理）}令 $ F \in \text{Ob}(\mathbb{D}^b(X)) $  和 $ \phi: X \to \mathbb{R} $  成为 $ C^1 $  函数，使得 $ \phi: \text{supp}(F) \to \mathbb{R} $  是正确的。让  $ a, b \in \mathbb{R} $  与  $ a < b $  。

(i) 假设 $ \mathrm{d}\phi(x) \notin \mathrm{SS}(F) $  对于任何 $ x \in X $  使得 $ a \leqslant \phi(x) < b $  。然后是自然态射：

\[
\begin{aligned}R\varGamma(\phi^{-1}(]-\infty,b[);F)&\to R\varGamma(\phi^{-1}(]-\infty,a]);F)\\&\to R\varGamma(\phi^{-1}(]-\infty,a];F)\end{aligned}
\]

是同构。

(ii) 假设 $ -\mathrm{d}\phi(x) \notin \mathrm{SS}(F) $  对于任何 $ x \in X $ ，使得 $ a < \phi(x) \leq b $ （或 $ a \leq \phi(x) < b $ ）。然后是自然态射：

\[
R\Gamma_{\phi^{-1}(]-\infty,a])}(X;F)\to R\Gamma_{\phi^{-1}(]-\infty,b])}(X;F)
\]

（分别为  $ R\Gamma_c(\phi^{-1}(]-\infty,a[);F)\to R\Gamma_c(\phi^{-1}(]-\infty,b[);F) $  ），是同构。

作为上面获得的结果的简单应用，我们将证明层的莫尔斯不等式。为此，我们假设，直到第 4 节结束，基环  $A$  是一个域，我们用  $k$  表示它。如果  $V \in \text{Ob}(\mathbf{D}^b(\mathfrak{Mod}^f(k)))$  ，我们按照练习 I.34 进行设置：

\[
\mathbf{b}_{j}(V)=\dim H^{j}(V),
\]

\[
{\sf b}_{l}^{*}(V)=(-1)^{l}\sum_{j\leqslant l}(-1)^{j}{\sf b}_{j}(V)\enspace.
\]

现在让 $ F \in \mathrm{Ob}(\mathbb{D}^{b}(X)) $  和 $ \phi: X \to \mathbb{R} $  成为 $ C^\infty $  函数。我们设定：

\[
\varLambda_{\phi}=\left\{(x;\mathrm{d}\phi(x));x\in X\right\}.
\]

请注意，  $ A_{\phi} $  是  $ T^{*}X $  的（通常非圆锥）拉格朗日子流形。

\statementlabel{命题 5.4.20}我们提出以下假设（i）、（ii）、（iii）：

(i) 对于所有 $t \in \mathbb{R}$  ， $\{x \in \mathrm{supp}(F); \phi(x) \leq t\}$  是紧凑的，

(ii) 集合  $ A_\phi \cap \mathrm{SS}(F) $  是有限的，比如  $ \{p_1, \ldots, p_N\} $  ，

(iii) 设置 $ x_i = \pi(p_i) $  、 $ V_i = (R\Gamma_{\{\phi \geq \phi(x_i)\}}(F))_{x_i} $  属于所有 $ i = 1, \ldots, N $  的 $ \mathrm{Ob}(\mathbb{D}^b(\mathfrak{Mod}^f(k))) $  。

然后：

(a)  $ R\Gamma(X; F) $  属于  $ \mathbf{D}^b(\mathfrak{M} \circ \mathfrak{d}^f(k)) $  ，

(b) 设置 $ \mathbf{b}_j = \mathbf{b}_j(R\Gamma(X; F)), \mathbf{b}_j^* = \mathbf{b}_j^*(R\Gamma(X; F)), \mathbf{b}_j' = \sum_i \mathbf{b}_j(V_i) $  和 $ \mathbf{b}_j'^* = \sum_i \mathbf{b}_j^*(V_i) $  我们有：

\[
{\bf b}_{i}^{*}\leqslant{\bf b}_{i}^{\prime*}\qquad\text{for }~all~l~,
\]

\[
\sum_{j}(-1)^{j}\mathbf{b}_{j}=\sum_{j}(-1)^{j}\mathbf{b}_{j}^{\prime}\enspace.
\]

\statementlabel{证明}首先请注意，假设 (i) 意味着  $\phi$  在  $\operatorname{supp}(F)$  上是正确的。设置  $G = R\phi_* F$  并让  $t$  为  $\mathbb{R}$  上的坐标。让  $\phi(\{x_1, \ldots, x_N\}) = \{t_1, \ldots, t_L\}$  和  $t_i < t_{i+1}$  代表所有  $i$  。然后：

(i)′ 对于所有  $ t \in \mathbb{R} $  、 $ -\infty $  、 $ t[\cap \operatorname{supp}(G) $  都是紧凑的。

(ii)' $ \mathrm{SS}(G) \cap \{(t; \mathrm{d}t); t \in \mathbb{R}\} $  包含在  $ \bigcup_{i=1}^{L} \{(t_i; \mathrm{d}t)\} $  中（这源自命题 5.4.4）。

(iii)' $ (R\Gamma_{\{t\geq t_i\}}(G))_{t_i}\simeq\oplus_{\phi(x_j)=t_i}V_j $ 。

事实上，根据微支撑的定义，左侧与  $ (R\Gamma_{\{x;\phi(x)\geq t_i\}}(F))|_{\phi^{-1}(t_i)} $  同构，因此与  $ \oplus_{\phi(x_j)=t_i}(R\Gamma_{\{\phi(x)\geq t_i\}}(F))_{x_j} $  同构。由于  $ R\Gamma(X;F)=R\Gamma(\mathbb{R};G) $  和 G 满足与 F 相同的假设，我们可以从一开始就假设  $ X=\mathbb{R} $  和  $ \phi $  是恒等式。

现在设置  $ t_{o} = -\infty $  、  $ t_{L+1} = +\infty $  、  $ I_{t} =]-\infty $  、  $ t[, Z_{t} =]-\infty $  、  $ t] $  并简写为  $ I_{j} = I_{t_{j}} $  、  $ Z_{j} = Z_{t_{j}} $  。

根据非特征变形引理（命题 2.7.2），我们有同构：

\[
R\varGamma(I_{j+1};F)\xrightarrow{\sim}R\varGamma(I_{t};F)\;,\qquad t_{j}<t\leqslant t_{j+1}\enspace.
\]

取  $t > t_{j}$  的上同调和归纳极限，我们得到：

\[
H^{k}(I_{j+1};F)\simeq H^{k}(Z_{j};F)\enspace.
\]

考虑一下不同的三角形：

\[
(R\varGamma_{\{t\geq t_{j}\}}(F))_{t_{j}}\longrightarrow R\varGamma(Z_{j};F)\longrightarrow R\varGamma(I_{j};F)\xrightarrow[+1]{}.
\]

从 $ R\Gamma(I_1; F) = 0 $ 开始，根据(5.4.22)，我们通过归纳得到 $ R\Gamma(Z_j; F) $ 和 $ R\Gamma(I_j; F) $ 都属于 $ \mathbf{D}^b(\mathfrak{Mod}^f(k)) $ 。这就证明了(a)。

设置  $ \mathrm{b}_{j}(Z_{i}) = \mathrm{b}_{j}(R\Gamma(Z_{i}; F)) $  和  $ \mathrm{b}_{j}(I_{i}) = \mathrm{b}_{j}(R\Gamma(I_{i}; F)) $  ，简称。应用练习 I.34 的结果，我们得到 (5.4.23)：

\[
\mathsf{b}_{l}^{\prime*}\geqslant\sum_{i=1}^{L}\left(\mathsf{b}_{l}^{*}(Z_{i})-\mathsf{b}_{l}^{*}(I_{i})\right).
\]

自从

\[
\mathbf{b}_{j}=\sum_{i=1}^{L}\left(\mathbf{b}_{j}(Z_{i})-\mathbf{b}_{j}(I_{i})\right)
\]

由(5.4.22)，我们得到(5.4.20)。

最后，(5.4.21) 由 (5.4.20) 得出，因为  $ b_i^* = -b_{i-1}^* $  和  $ b_{i'}^* = -b_{i-1}' $  对于  $ l \gg 0 $  。   $ \square $ 

关于经典莫尔斯不等式，参见。米尔诺[1]。请注意， $\sum_{j}(-1)^{j}b_{j}$  是 $F$  在 $X$  上的欧拉-庞加莱索引。我们将在第九章中回到这个索引。

\subsection{锥层的微支撑}
在本节中，我们研究矢量束上层的微支撑。

令  $\tau: E \to Z$  为流形  $Z$  上的实向量丛（参见 III §7）。令 $\mu: \mathbb{R}^+ \times E \to E$  表示 $\mathbb{R}^+$  在 $E$  上的作用映射，并让 $e$  表示 $E$  上描述 $\mu$  的无穷小作用的向量场。因此，对于 $E$ 上的任何函数 $\varphi$ ，我们有 $(\mathbf{e}(\varphi))(x) = \frac{\partial}{\partial t} \varphi(\mu(t, x))|_{t=1}$ 。矢量场  $e$  通常称为欧拉矢量场。映射  $\mu$  定义了映射  $'\mu': \mathbb{R}^+ \times T^*E \to T^*\mathbb{R}^+ \times T^*E$  。当考虑  $'\mu'$  对  $1 \in \mathbb{R}^+$  的限制，并与投影  $T^*\mathbb{R}^+ \times T^*E \to T^*\mathbb{R}^+ \simeq \mathbb{R}^+ \times \mathbb{R} \to \mathbb{R}$  组合时，我们得到映射：

\[
\theta_{E}\colon T^{*}E\to\mathbb{R}~.
\]

这张映射只不过是 e 的主符号。

令 $(z, x)$  为坐标系（ $Z$  上的局部坐标系），这样 $(z)$  为 $Z$  的坐标， $(x)$  为线性坐标。让  $(z, x; \zeta, \xi)$  表示  $T^*E$  的关联坐标。因此：

\[
\left\{\begin{aligned}\mathbf{e}&=\left\langle x,\frac{\partial}{\partial x}\right\rangle=\sum_{j}x_{j}\frac{\partial}{\partial x_{j}}\ ,\\ \theta_{E}&=\left\langle x,\xi\right\rangle\ ,\\ \alpha_{E}&=\left\langle\zeta,\mathrm{d}z\right\rangle+\left\langle\xi,\mathrm{d}x\right\rangle\ .\end{aligned}\right.
\]

这里， $ \alpha_E $  是  $ T^*E $  上的基本 1-形式（参见 A§2）。

让我们回想一下， $ T^{*}(E/Z) $  表示 E 在 Z 上的相对余切丛。它由 E 上向量丛的精确序列定义：

\[
0\to E\times\mathop{T}{*}Z\to T{*}E\to T{*}(E/Z)\to0.
\]

对于任何  $ z \in Z $  ，  $ T^*(E/Z) $  在  $ z $  上的纤维  $ T^*(E/Z)_z $  可以用  $ z $  上  $ E $  的纤维  $ E_z $  的余切丛  $ T^*(E_z) $  来识别。另一方面，  $ T^*(E_z) $  与  $ E_z \times E_z^* $  同构。因此，对于每个  $ z $  ，我们获得从  $ T^*(E/Z)_z $  到  $ E_z^* $  的映射。这给出了一张映射

\[
T^{*}(E/Z)\to E^{*}\enspace.
\]

然后我们定义  $Z$  上的态射：

\[
\psi_{E}\colon T^{*}E\to E^{*}
\]

作为  $ T^*E \to T^*(E/Z) \to E^* $  的组合。

\statementlabel{命题 5.5.1}(i) 存在唯一映射 $ \Phi_E: T^*E \to T^*E^* $ ，使得 $ \alpha_E - \mathrm{d}\theta_E = \Phi_E^*(\alpha_E^*) $  且组合 $ T^*E \xrightarrow[\Phi_E]{} T^*E^* \xrightarrow[\alpha]{} E^* $  是 $ \psi_E $  。

（二） $ \theta_{E^{*}} \circ \Phi_{E} = -\theta_{E} $ 。

(iii)  $ \Phi_{E^{*}} \circ \Phi_{E} = a^{*} $  ，其中  $ a^{*} $  是 E 上反足图诱导的  $ T^{*}E $  的自同构。

\statementlabel{证明}令 $ (z, x) $  为 $ E $  上的局部坐标系，如上，并让 $ (z, y) $  表示 $ E^* $  上的对偶坐标。因此  $ E $  和  $ E^* $  在  $ Z $  上的规范配对由  $ E \times_Z E^* \ni (z, x, y) \mapsto \langle x, y \rangle = \sum_j x_j y_j $  给出。令  $ (z, x; \zeta, \xi) $  （分别为  $ (z, y; \zeta, \eta) $  ）为  $ T^*E $  （分别为  $ T^*E^* $  ）上的关联坐标。映射 $ \psi_E $  由：  $ (z, x; \zeta, \xi) \mapsto (z; \xi) $  给出，我们有：

\[
\alpha_{E}-\mathrm{d}\theta_{E}=\left\langle\zeta,\mathrm{d}z\right\rangle-\left\langle x,\mathrm{d}\xi\right\rangle.
\]

因此映射

\[
\varPhi_{E}:(z,x;\zeta,\xi)\mapsto(z,\xi;\zeta,-x)
\]

满足所需的条件。映射 $ \varPhi_E $  由这些条件唯一确定。事实上 $ \pi \circ \varPhi_E = \psi_E $ 意味着：

\[
\Phi_{E}(z,x,\zeta,\xi)=(z,\xi;\varphi_{0}(z,x,\zeta,\xi),\varphi_{1}(z,x,\zeta,\xi))\ ,
\]

另一个条件意味着：

\[
\langle\varphi_{0},\mathrm{d}z\rangle+\langle\varphi_{1},\mathrm{d}\xi\rangle=\langle\zeta,\mathrm{d}z\rangle-\langle x,\mathrm{d}\xi\rangle.
\]

(ii) 由公式(5.5.6)得出。

(iii) 从  $ \Phi_E^*(z, y, \zeta, \eta) = (z, \eta; \zeta, -y) $  开始，我们就有了  $ \Phi_{E^*} \circ \Phi_E(z, x, \zeta, \xi) = (z, -x, \zeta, -\xi) $  。   $ \Box $ 

\statementlabel{备注 5.5.2}映射  $ \varPhi_E $  是辛变换（即：  $ \varPhi_E $  保留  $ d\alpha_E $  ），但不是齐次辛变换（即  $ \varPhi_E $  不保留  $ \alpha_E $  ）。

我们用  $ S_{E} $  表示欧拉向量场的特征变化：

\[
S_{E}=\theta_{E}^{-1}(0).
\]

如果  $(z,x;\zeta,\xi)$  是  $T^*E$  上的坐标，则  $S_E$  由方程  $\langle x,\xi\rangle = 0$  定义。

最后注意  $ T^*E $  被赋予了  $ \mathbb{R}^+ $  的两个动作，一个来自  $ T^*E $  over  $ E $  的向量丛结构，另一个来自  $ E $  over  $ Z $  的向量丛结构。在上述局部坐标系中，这两个动作由以下方式描述：

\[
\begin{cases}(z,x;\zeta,\xi)\mapsto(z,x;\lambda\zeta,\lambda\xi),&\lambda\in\mathbb{R}^{+},\\ (z,x;\zeta,\xi)\mapsto(z,\mu x;\zeta,\mu^{-1}\xi),&\mu\in\mathbb{R}^{+}.\end{cases}
\]

如果  $ T^*E $  的子集在  $ \mathbb{R}^+ $  的两个操作下都不变，我们可以说  $ T^*E $  的子集是双圆锥的。

现在让我们在  $E$  上研究圆锥层。回想一下（III §7），范畴 $\mathbf{D}_{\mathbf{R}^r}^b(E)$ 表示 $\mathbf{D}^b(E)$ 的满子范畴，由对象 $F$ 组成，使得对于所有 $j$ ， $H^j(F)$ 在 $\mathbb{R}^+$ 的作用轨道上局部常值。

\statementlabel{命题 5.5.3}(i) 范畴  $ \mathbf{D}_{\mathbf{R}}^{b}(E) $  是  $ \mathbf{D}^{b}(E) $  的满子范畴，由对象  $ F $  组成，使得  $ \mathrm{SS}(F) \subset S_{E} $  。

(ii) 如果  $ F \in \mathrm{Ob}(\mathcal{D}_{\mathbb{R}^{+}}^{b}(E)) $  ，则  $ \mathrm{SS}(F) $  是双圆锥的。

\statementlabel{证明}(ii) 是显而易见的。让我们展示（i）。

设置  $ \dot{E} = E\backslash Z $  ，其中  $ Z $  用零部分标识。足以证明 $ \dot{E} $ 上的结果。然后从命题 5.4.5 得出，因为  $ \dot{E}/\mathbb{R}^{+} $  是流形，而  $ \dot{E} \times_{\dot{E}/\mathbb{R}^{+}} T^{*}(\dot{E}/\mathbb{R}^{+}) = \dot{E} \times_{E} S_{E} $  。   $ \square $ 

投影  $ \tau: E \to Z $  及其限制  $ \dot{\tau} $  到  $ \dot{E} $  定义了映射：

\textbackslash{}[\textbackslash{}begin\{array\}\{c\}T\textasciicircum{}\{*\}E\textbackslash{}xleftarrow\{\textbackslash{}iota\_\{\textbackslash{}tau\textasciicircum{}\{\textbackslash{}prime\}\}\}\textbackslash{}setminus E\textbackslash{}times\textbackslash{}underset\{z\}\{\textbackslash{}times\}T\textasciicircum{}\{*\}Z\textbackslash{}xrightarrow\{\textbackslash{}quad\textbackslash{}tau\_\{\textbackslash{}pi\}\textbackslash{}quad\}T\textasciicircum{}\{*\}Z\textbackslash{}\textbackslash{}\textbackslash{}quad\textbackslash{}

此外，E 的零部分允许我们定义沉浸：

\[
T^{*}Z\hookrightarrow\mathbb{E}\times_{z}T^{*}Z\hookrightarrow{}_{{}^{\prime}\tau^{\prime}}T^{*}E~.
\]

通过这种方式，我们有时将  $ T^*Z $  视为  $ T^*E $  的子集。请注意，如果  $ A $  是  $ T^*E $  的双圆锥闭子集，则：

\[
\tau_{\pi}{}^{t}\tau^{\prime-1}(A)=T^{*}Z\cap A.
\]

\statementlabel{命题 5.5.4}让 $ F \in \mathrm{Ob}(\mathbf{D}_{\mathbb{R}^n}^b(E)) $  。然后：

（一） $ \mathrm{SS}(R\tau_{*}(F)) \subset T^{*}Z \cap \mathrm{SS}(F) $ ，

（二） $ \mathrm{SS}(R\tau_!(F)) \subset T^*Z \cap \mathrm{SS}(F) $ ，

（三） $ \mathrm{SS}(R\dot{\tau}_{*}(F)) \subset \dot{\tau}_{\pi}^{t}\dot{\tau}^{\prime-1}(\mathrm{SS}(F)) $ ，

（四） $ \mathrm{SS}(R\dot{\tau}_{1}(F)) \subset \dot{\tau}_{\pi}^{t}\dot{\tau}^{\prime-1}(\mathrm{SS}(F)) $ 。

\statementlabel{证明}在  $Z$  本地，我们可以选择坐标系  $(z;x)$  ， $(x)$  表示纤维坐标。然后，通过在情况 (i) 或 (ii) 中设置  $Z = X$  、  $Y = E$  和  $Y_t = \{(z,x); |x|^2 < t\}$  在情况 (iii) 或 (iv) 中设置  $Y_t = \{(z,x); t^{-1} < |x|^2 < t\}$ ，得出命题 5.4.17 的结果。   $\square$ 

在第三章§7中，我们定义了圆锥束 $ F \in \mathrm{Ob}(\mathbf{D}_{\mathbf{R}^+}^b(E)) $ 的傅里叶-佐藤变换 $ F^\wedge $ 。回想一下：

\[
\begin{array}{r}{F^{\wedge}=R p_{2\ast}R\Gamma_{P}(p_{1}^{-1}F),}\end{array}
\]

其中  $ p_1 $  和  $ p_2 $  表示  $ E \times_Z E^* $  和  $ P = \{(z, x, y) \in E \times_Z E^*; \langle x, y \rangle \geqslant 0\} $  上定义的第一个和第二个投影。

\statementlabel{定理 5.5.5}让 $F \in \mathrm{Ob}(\mathbb{D}_{\mathbb{R}^+}^b(E))$  。然后，通过映射 $\Phi_E$ 识别 $T^*E$ 和 $T^*E^*$ （参见命题5.5.1），我们有：

\[
\mathrm{SS}(\boldsymbol{F})=\mathrm{SS}(\boldsymbol{F}^{\wedge})~.
\]

\statementlabel{证明}由定理3.7.9，足以证明包含 $ \mathrm{SS}(F^\wedge) \subset \mathrm{SS}(F) $ 。我们在  $ E $  上选择局部坐标  $ (z, x) $  ，并用  $ (z, y) $  表示  $ E^* $  上的双坐标，并用  $ (z, x; \zeta, \xi) $  和  $ (z, y; \zeta, \eta) $  分别表示  $ T^*E $  和  $ T^*E^* $  上的坐标。那么 $ \Phi_E $  由(5.5.6) 定义。让  $ (z_\circ, x_\circ; \zeta_\circ, \xi_\circ) \notin \mathrm{SS}(F) $  ，让我们证明  $ (z_\circ, \xi_\circ; \zeta_\circ, -x_\circ) \notin \mathrm{SS}(F^\wedge) $  。第一步，我们假设：

\[
x_{\circ}\neq0,\quad\xi_{\circ}\neq0,
\]

\[
R\Gamma_{Z}(F)=0.
\]

用  $j$  表示嵌入  $\dot{E} \hookrightarrow E$  ，用  $\tilde{j}$  表示嵌入  $\dot{E} \times_Z E^* \hookrightarrow E \times_Z E^*$  ，用  $\dot{p}_2$  表示投影  $\dot{E} \times_Z E^* \to E^*$  。假设（5.5.13），我们有：

\[
\begin{aligned}{{R}\varGamma_{P}(p_{1}^{-1}F)}&{{}\simeq{R}\varGamma_{P}(p_{1}^{-1}R j_{*}j^{-1}F)}\\ {}&{{}\simeq{R}\varGamma_{P}(R\tilde{j}_{*}j^{-1}p_{1}^{-1}F)\simeq{R}\tilde{j}_{*}\tilde{j}^{-1}{R}\varGamma_{P}(p_{1}^{-1}F),}\\ \end{aligned}
\]

因此：

\[
F^{\wedge}\simeq R\dot{p}_{2*}\tilde{j}^{-1}R\varGamma_{P}(p_{1}^{-1}F).
\]

假设 $ (z_{\circ}, \xi_{\circ}; \zeta_{\circ}, -x_{\circ}) \in \mathrm{SS}(F^{\wedge}) $  。将命题 5.5.4 应用于  $ p_2 $  ，我们发现一些  $ x_1 \in \dot{E} $  使得  $ (z_{\circ}, x_1, \xi_{\circ}; \zeta_{\circ}, 0, -x_{\circ}) \in \mathrm{SS}(R\Gamma_P(p_1^{-1}F)) $  。

如果  $ (z_o, x_1, \xi_o) \in \partial P $  ，则集合  $ P $  具有 (5.5.12) 的平滑边界，并且其此时的共法是协向量  $ \xi_o dx + x_1 dy $  。由于这一点不属于  $ \mathrm{SS}(p_1^{-1}F)^a $  ，我们可以应用命题 5.4.14 并找到  $ \xi_1 $  使得  $ (z_o, x_1, \xi_o; \zeta_o, \xi_1, 0) \in \mathrm{SS}(p_1^{-1}F) $  ，  $ (z_o, x_1, \xi_o; 0, \xi_1, x_o) \in \mathrm{SS}(A_P) $  。因此  $ x_o = kx_1 $  、  $ \xi_1 = k\xi_o $  对于某些  $ k \geq 0 $  和  $ (z_o, x_1; \zeta_0, \xi_1) \in \mathrm{SS}(F) $  。这是一个矛盾，因为 $ \mathrm{SS}(F) $  是双圆锥的。

要消除假设 (5.5.13)，请考虑已区分的三角形：

\[
R\varGamma_{Z}(F)\longrightarrow F\longrightarrow R j_{*}j^{-1}F\mathop{\longrightarrow}_{+1}.
\]

由于我们可以将前面的结果应用于  $ Rj_* j^{-1}F $  ，因此足以证明  $ (z_o, \xi_o, \zeta_o, -x_o) \notin \mathrm{SS}(R\Gamma_Z(F)^\wedge) $  。但这是根据命题 3.7.13 得出的。事实上，用  $ i $  表示嵌入  $ Z \hookrightarrow E $  。我们有 $ R\Gamma_Z(F)^\wedge \simeq (i_i i^! F)^\wedge \simeq \pi^{-1} i^! F $ ， $ \pi^{-1} i^! F $ 的微支撑包含在集合 $ \{(z, y; \zeta, \eta); \eta = 0\} $ 中。

最后我们处理一般情况。

考虑  $ \mathbb{R}^2 $  上的层  $ A_{R \times \{0\}} $  。应用命题 3.7.15 和 3.7.12 我们得到：

\[
(F\boxtimes A_{\mathbb{R}\times\{0\}})^{\wedge}\simeq F^{\wedge}\boxtimes A_{\{0\}\times\mathbb{R}}[-1]~.
\]

如果  $(z_{\circ}, x_{\circ}; \zeta_{\circ}, \xi_{\circ}) \notin \mathrm{SS}(F)$  ，则  $(z_{\circ}, x_{\circ}, 1, 0; \zeta_{\circ}, \xi_{\circ}, 0, 1) \notin \mathrm{SS}(F \boxtimes A_{\mathbb{R} \times \{0\}})$  ，因此  $(z_{\circ}, \xi_{\circ}, 0, 1; \zeta_{\circ}, -x_{\circ}, -1, 0) \notin \mathrm{SS}(F^{\wedge} \boxtimes A_{\{0\} \times \mathbb{R}})$  ，并且根据命题 5.4.4 和 5.4.5，这意味着  $(z_{\circ}, \xi_{\circ}; \zeta_{\circ}, -x_{\circ}) \notin \mathrm{SS}(F^{\wedge})$  。 □

\subsection{习题}
\statementlabel{练习 V.1}令 $X$  为流形， $V$  为 $T^*X$  的闭圆锥子集。其中一个表示  $\mathfrak{Sh}_V(X)$  （分别为  $\mathbf{D}_V^b(X)$  ）  $\mathfrak{Sh}(X)$  （分别为  $\mathbf{D}^b(X)$  ）的满子范畴，由对象  $F$  组成，使得  $\mathrm{SS}(F) \subset V$  。

(i) 现在假设 $X$  是一个赋有线性坐标 $(z) = (z_1, \ldots, z_n)$  的复向量空间，并令 $(z; \zeta)$  为复余切丛 $X \times X^*$  上的坐标。我们将  $X \times X^*$  与  $\langle \zeta, \mathrm{d}z \rangle + \langle \zeta, \mathrm{d}\overline{z} \rangle$  的实余切丛进行识别。让 $V = \{(z; \zeta); \sum_j z_j \zeta_j = 0\}$  。证明  $\mathfrak{S}h_V(X)$  由层组成，这些层在  $\mathbb{C}^*$  对  $X$  的作用轨道上局部常值。

(ii) 在情况(i) 中，讨论 $ \delta: \mathbf{D}^b(\mathfrak{S}\mathfrak{h}_\nu(X)) \to \mathbf{D}_\nu^b(X) $  是否是范畴的等价。

\statementlabel{练习 V.2}令  $F \in \mathrm{Ob}(\mathbf{D}^{b}(X))$  和  $\varphi$  成为  $X$  上的真实  $C^{1}$  函数。让 $x_{\circ} \in X$  与 $\varphi(x_{\circ}) = 0$  、 $\mathrm{d}\varphi(x_{\circ}) \notin \mathrm{SS}(F)$  。证明“ $\lim_{U} R\Gamma_{\{x; \varphi(x) \geq 0\}}(U; F) = 0$ ，其中 $U$ 穿过 $X$ 中 $x_{\circ}$ 的开邻域族。（对于ind-object“ $\lim_{U} \varphi(x_{\circ}) \notin \mathrm{SS}(F)$ ），参见I§11。）

\statementlabel{练习 V.3}令 $E$  为实向量空间， $F \in \mathrm{Ob}(\mathbf{D}_{\mathbf{R}^+}^b(E))$  。让 $\xi_\circ \in E^*$  。证明  $(0; \xi_\circ) \notin \mathrm{SS}(F)$  当且仅当存在  $\xi_\circ$  的凸圆锥开邻域  $\gamma$  的基本系统，使得  $R\Gamma_\gamma^* (E; F) = 0$  。 （提示：使用定理 5.5.5。）

\statementlabel{练习 V.4}令  $ C = \{(x, y, t) \in \mathbb{R}^3; x^2 + y^2 = t^2, t > 0\} $  和  $ j $  为嵌入  $ C \subsetneq \mathbb{R}^3 $  。设置 $ F = R j_* A_C $ 。

(i) 使用练习 V.3 的结果证明  $ \pi^{-1}(0) \cap SS(F) \neq T_{\{0\}}^{*} \mathbb{R}^{3} $  。

(ii) 证明 $ H^{1}(F) = A_{\{0\}} $  。

请注意，我们得到了  $F \in \mathrm{Ob}(\mathbf{D}^{b}(X))$  的示例，其中  $\mathrm{SS}(H^{j}(F))$  不包含在  $\mathrm{SS}(F)$  中。

\statementlabel{练习 V.5}设 X 为向量空间。

(i) 令  $Z$  为  $X$  和  $F = A_Z$  的闭凸子集。显示对于任何  $x \in Z$  、  $\mathrm{SS}(F) \cap \pi^{-1}(x) = \gamma^\circ$  、  $\gamma = C_x(Z)$  。

(ii) 令  $\Omega$  为  $X$  和  $F = A_\Omega$  的开凸子集。显示 $\mathrm{SS}(A_\Omega) = \mathrm{SS}(A_{\bar{\Omega}})^\alpha$  。

\statementlabel{练习 V.6}让 $ F \in \mathrm{Ob}(\mathbb{D}^{b}(X)) $  。证明包含 $ \mathrm{SS}(F) \subset \bigcup_{j} \mathrm{SS}(H^{j}(F)) $ 。

\statementlabel{练习 V.7}(i) 令 $ \{F_\lambda\}_{\lambda} $  为 $ X $  上层的归纳系统。证明包含 $ \mathrm{SS}(\lim F_\lambda) \subset \bigcup_{\lambda} \mathrm{SS}(F_\lambda) $ 。

(ii) 令 $\{F_n\}_{n\in\mathbb{N}}$  为 $X$  上层的射影系统。证明包含 $\mathrm{SS}(\lim\underbrace{F_n}_{n}) \subset \overline{\bigcup_n SS}(F_n)$ 。

\statementlabel{练习 V.8}令 $V$  为有限维实向量空间， $\Omega$  为 $V^*$  中的开锥体（不一定是凸锥体）。设置 $G = (A_\Omega)^\wedge \otimes \omega_V$ 。令  $\mu: V \times V \to V$  为映射  $(x, y) \mapsto y - x$  并令  $K = \mu^{-1}G$  。回想一下  $\mathrm{III}$  §6 的函子  $\phi_K$ 。

(i) 证明如果  $ F \in \mathrm{Ob}(\mathbb{D}^{b}(V)) $  ，则  $ \mathrm{SS}(\Phi_{K}(F)) \subset V \times \bar{\Omega} $  。

(ii) 证明存在一个自然态射  $ \Phi_{\kappa}(F) \to F $  ，它是  $ V \times \Omega $  上的同构。 （与命题 5.2.3 比较。）

\statementlabel{练习 V.9}让 $ F \in \mathrm{Ob}(\mathbb{D}^b(X \times \mathbb{R})) $  满足 $ \mathrm{SS}(F) \cap (T_X^* X \times T^* \mathbb{R}) \subset T_{X \times \mathbb{R}}^*(X \times \mathbb{R}) $  。令  $ \varphi: \mathbb{R} \to \mathbb{R} $  为  $ \varphi(s) = 0 $  对于  $ s \leq 0 $  和  $ \varphi(s) = s $  对于  $ s \geq 0 $  定义的连续映射。证明  $ \mathrm{SS}(\varphi_X^{-1} F) \cap (T_X^* X \times T^* \mathbb{R}) \subset T_{X \times \mathbb{R}}^{*}(X \times \mathbb{R}) $  其中  $ \varphi_X = \mathrm{id}_X \times \varphi $  。

（提示：让 $ S = \{(t, s) \in \mathbb{R}^2; \quad t = s \geq 0 \quad \text{or} \quad t = 0, \quad s \leq 0\} $ ）。然后  $ \varphi_X^{-1}F \simeq Rq_{2*}(q_1^{-1}F)_S $  ，其中  $ q_1 $  和  $ q_2 $  分别是从  $ X \times \mathbb{R}_t \times \mathbb{R}_s $  到  $ X \times \mathbb{R}_t $  和  $ X \times \mathbb{R}_s $  的投影。然后检查  $ \pm ds \notin (\mathrm{SS}(F) \times T_{\mathbb{R}_s}^* \mathbb{R}_s + \mathrm{SS}(A_{X \times S})) $  就足够了。）

\statementlabel{练习 V.10}令 $F_0$  和 $F_1$  属于 $\mathbf{D}^b(X)$  、 $F_0$  和 $F_1$  ，两者都具有紧支撑。假设  $F_0$  和  $F_1$  是  $\mathbf{h}omotopic$  如果存在  $H \in \mathrm{Ob}(\mathbf{D}^b(X \times \mathbb{R}))$  使得投影  $p: X \times \mathbb{R} \to \mathbb{R}$  在  $\mathrm{supp}(H)$  、  $\mathrm{SS}(H) \cap (T_X^* X \times T^* \mathbb{R}) \subset T_{X \times \mathbb{R}}^* (X \times \mathbb{R})$  和  $H|_{X \times \{i\}} \simeq F_i$  、  $i = 0, 1$  上正确。

(i) 证明如果  $H$  实现上述同伦，则存在  $\widetilde{H}$  满足与  $H$  相同的假设，而且  $\widetilde{H}|_{X\times[1,2]} \simeq F_1 \boxtimes A_{[1,2]}$  。 （提示：使用练习 V.9。）

(ii) 证明“同伦”是等价关系。 （提示：使用 (i) 和练习 II.22。）

(iii) 证明如果 $F_{0}$  和 $F_{1}$  同伦，则 $D_{x}F_{0}$  和 $D_{x}F_{1}$  同伦。

(iv) 证明如果 $F_0$  和 $F_1$  同伦，则对于任何 $k \in \mathbb{Z}$  ， $H^k(X; F_0)$  和 $H^k(X; F_1)$  同构。

(v) 对于 $ X = \mathbb{R} $  ，证明 $ A_{[0,1]} $  与 $ A_{\{0\}} $  同伦， $ A_{[0,1]} $  与零同伦， $ A_{[0,1]} $  与 $ A_{\{0\}}[-1] $  同伦。

\statementlabel{练习 V.11}(i) 令  $X$  为流形， $i: Y \hookrightarrow X$  为子流形  $Y$  的封闭嵌入。证明对于  $F \in \mathrm{Ob}(\mathbf{D}^b(X))$  ，当  $T_Y^* X \setminus \mathrm{SS}(F) \to Y$  具有连续部分时， $i^! F \to i^{-1} F$  是零态射。 （提示：考虑 $\mu_Y(i_* i^! F) \to \mu_Y(F) \to \mu_Y(i_* i^{-1} F).$ ）

(ii) 使用 (i) 再次证明练习 III.7。

(iii) 证明，在练习 III.8 的情况下，如果  $ \dot{T}_Y^* X \to Y $  具有连续部分，则组合  $ H^k(Y; A_Y) \to H^{k+l}(X; A_X) \to H^{k+l}(Y; A_Y) $  为零。

\statementlabel{练习 V.12}假设基环是交换域  $ k $  。令 X 为流形， $ \phi $  为 X 上的实数  $ C^\infty $  函数，并假设：

(i) 对于所有 $ t \in \mathbb{R} $ ；   $ \{x \in X; \phi(x) \leq t\} $  是紧凑的，

(ii) 集合  $\{x; \mathrm{d}\phi(x) = 0\}$  是有限的，例如  $\{x_1, \ldots, x_N\}$  ，并且在每个  $x_i$  处，Hessian  $H_\phi(x_i)$  是非简并的，具有  $\varepsilon_i$  负特征值。

证明  $ R\Gamma(X; k_x) $  属于  $ \mathbf{D}^b(\mathfrak{M}_\mathbf{od}^f(k)) $  ，并且：

\[
\chi(R\varGamma(X;k_{X}))=\sum_{i}(-1)^{\varepsilon_{i}},
\]

其中  $ \chi(\cdot) $  在练习 I.32 中定义。

（提示：使用命题 5.4.20，练习 III.5 和“莫尔斯引理”的本地版本。）

\statementlabel{练习 V.13}让 $ F \in \mathrm{Ob}(\mathbb{D}^{b}(X)) $  。证明公式

\[
SS(\mathbf{D}_{X} F)=SS(F)^{a}
\]

假设以下条件 (i) 或 (ii)。

(i) F 是上同调可构造的。

(ii) 基环 $ A $  满足： 如果 $ M \in \text{Ob}(\mathbb{D}^b(\mathfrak{M}_\text{ob}(A))) $  和 $ RHom_A(M, A) = 0 $  那么 $ M = 0 $  。

请注意，如果  $A$  是一个字段或  $\mathbb{Z}$ ，则 (ii) 为真（参见练习 I.31）。作者没有提供不满足此属性的环 $A$  的示例。

\statementlabel{练习 V.14}设置 $X = \mathbb{R}, B = \{0\} \cup \{1/n; n \in \mathbb{N} \backslash \{0\}\}$ 。让 $F = G = \mathbb{Q}_B$  。证明 $\nu_{\{0,0\}} (F \boxtimes G) \neq \nu_{\{0\}} (F) \boxtimes \nu_{\{0\}} (G)$ 。

\subsection{注记}
作者于 1982 年提出了实流形 X 上层复合体 F 的微支撑概念（Kashiwara-Schapira [2,3]）。它可以被认为是莫尔斯理论的概括（参见 Milnor [1]），因为它在推论 5.4.19 中清晰可见，甚至是 Goresky-MacPherson [2] 最近的“分层莫尔斯理论”（最后一个理论仅适用于可构造层，参见第八章）。

然而，我们的动机来自线性偏微分方程，特别是来自双曲系统的研究（参见 Kashiwara-Schapira [1]）。在研究此类系统时，人们会导致解的定义域发生“非特征变形”。这样看来，我们可能会忘记我们正在处理微分方程，我们可能会忘记 X 的复数结构，而唯一要记住的是系统在实余切丛中的特征变化，即系统复形解的微支撑（参见下面的定理 11.3.3）。正如我们将在

第十一章，该方法应用于微微分方程的研究，是非常富有成效的。

Kashiwara-Schapira [1] 中引入了  $ \gamma $  拓扑，以便使微微分算子在全纯函数上运行，并且命题 5.2.1 在该论文中得到了证明（当然是在更严格的背景下）。本章的结果是在 Kashiwara-Schapira [3] 中获得的，但命题 5.4.20 除外，该命题是 Schapira-Tose [1] 对 Kashiwara [7] 结果的变体。 （这个命题是经典莫尔斯不等式的推广，我们参考 Bott [1]。）

\section{微支撑与微局部化}
\subsection{概要}
令 $X$  为流形， $\Omega$  为 $T*X$  的子集。我们将三角范畴  $\mathbf{D}^{b}(X; \Omega)$  定义为  $\mathbf{D}^{b}(X)$  的本地化，其微支撑与  $\Omega$  不相交的对象的满子范畴。然后，在  $\Omega$  上“微局部”工作，并在  $X$  上使用一堆  $F$  就得到了精确的含义：它只是意味着将  $F$  视为  $\mathbf{D}^{b}(X; \Omega)$  的对象。有了这个新概念，我们引入了“微局部逆像”和“微局部正像”。这些是  $\mathbf{D}^{b}(X; p)$  范畴的pro-对象或内部对象， $\mathbf{D}^{b}(X)$  在  $p$  的本地化，但我们给出了确保其保留在  $\mathbf{D}^{b}(X; p)$  范畴中的条件。

 $ \mathbf{D}^{b}(X) $  的定位与函子  $ \mu_{hom} $  相关，公式如下：

\[
\operatorname{H o m}_{\mathsf{D}^{b}(X;p)}(G,F)=H^{0}(\mu\mathcal{h}o m(G,F))_{p}\enspace.
\]

该公式是定理 6.5.4 证明中的重要步骤，定理 6.5.4 断言 SS(F) 是  $ T^{*}X $  的对合子集。

在得到对合性定理之前，我们研究了经过各种操作（开放嵌入的直像、微局部化等）后层的微支撑，将前一章的结果扩展到特征情况或非真情况。特别地，我们得到公式：

\[
\mathrm{S S}(\mu\mathrm{h o m}(G,F))\subset C(\mathrm{S S}(F),\mathrm{S S}(G))~.
\]

该公式利用了我们在第 2 节中研究的余切束中的法锥。

接下来，我们描述“微局部”层的特征，其微支撑包含在渐合子流形中。特别是，我们表明，如果 SS(F) 包含在 X 的子流形 Y 的共正规丛中，则对于某些 A 模 L，F 与束  $ L_{Y} $  是微局部同构的。

最后我们研究了逆像函子和微局部化函子交换的情况，并获得了微双曲系统柯西问题结果的层理论版本。

\subsection{范畴 $\mathbf{D}^{b}(X;\Omega)$}
让  $V$  成为  $T^*X$  的子集。我们通过设置定义 $\mathbf{D}^b(X)$ 的满子范畴 $\mathbf{D}_V^b(X)$ ：

\[
\mathbf{O b}(\mathbf{D}_{V}^{b}(X))=\left\{F\in\mathbf{O b}(\mathbf{D}^{b}(X));\mathbf{S S}(F)\subset V\right\}.
\]

这是一个三角范畴，  $ \mathbf{D}_{\mathcal{V}}^{b}(X) $  中的区别三角形是  $ \mathbf{D}^{b}(X) $  的三角形，其对象属于  $ \mathbf{D}_{\mathcal{V}}^{b}(X) $  ，而  $ \mathrm{Ob}(\mathbf{D}_{\mathcal{V}}^{b}(X)) $  是  $ \mathbf{D}^{b}(X) $  中的空系统。因此，我们可以应用 I §6 的结果并针对  $ \mathrm{Ob}(\mathbf{D}_{\mathcal{V}}^{b}(X)) $  本地化  $ \mathbf{D}^{b}(X) $  。

\statementlabel{定义 6.1.1}让 $ \Omega = T^*X \setminus V $  。

(i) 我们设定：

\[
{\bf\Delta{\cal D}}^{b}(X;\Omega)={\bf\Delta{\cal D}}^{b}(X)/{\mathrm{O b}}({\bf\Delta{\cal D}}_{\nu}^{b}(X))\;.
\]

(ii) 让 $F$  和 $G$  属于 $\mathbf{D}^{b}(X)$  。如果它们在  $\mathbf{D}^{b}(X;\Omega)$  中同构，我们可以说  $F$  和  $G$  在  $\Omega$  上同构。

对于  $p \in T^*X$  ，我们写  $\mathbf{D}^b(X; p)$  而不是  $\mathbf{D}^b(X; \{p\})$  。回想一下， $\mathrm{Ob}(\mathbf{D}^b(X; \Omega)) = \mathrm{Ob}(\mathbf{D}^b(X))$  和 $\mathbf{D}^b(X; \Omega)$  中的态射 $u: G \to F$  由 $\mathbf{D}^b(X)$  中的两个态射 $v: G \to F'$  和 $w: F \to F'$  的数据给出，使得 $w$  是 $\Omega$  上的同构，或者由两个态射 $w': G' \to G$  和 $w': G' \to G$  的数据给出 $v': G' \to F$  中的  $\mathbf{D}^b(X)$  ，使得  $w'$  是  $\Omega$  的同构。更准确地说：

\[
\left\{\begin{aligned}{\operatorname{H o m}_{\mathsf{D}^{b}(X;\varOmega)}(G,F)}&{{}=\varinjlim\operatorname{H o m}_{\mathsf{D}^{b}(X)}(G,F^{\prime})}\\ {}&{{}=\varinjlim\operatorname{H o m}_{\mathsf{D}^{b}(X)}(G^{\prime},F)}\\ \end{aligned}\right.,
\]

其中 (6.1.2) 的第一行（或第二行）中的归纳极限被态射范畴  $ w : F \to F' $  （或  $ w : G' \to G $  ）取代，使得  $ w $  是  $ \Omega $  上的同构（参见练习 I.14）。请注意，如果  $ F \in \text{Ob}(\mathbb{D}^b(X; \Omega)) $  ，则  $ \text{SS}(F) $  在  $ \Omega $  中定义良好。

在 IV §4 中，我们引入了从  $ \mathbf{D}^b(X)^\circ \times \mathbf{D}^b(X) $  到  $ \mathbf{D}^b(T^*X) $  的双函子  $ \mu h o m $  。如果 $ F $  或 $ G $  的微支撑不与 $ \Omega $  相交，则根据推论5.4.10， $ \mu h o m(G, F)|_{\Omega} = 0 $  相交。因此  $ \mu h o m(\cdot, \cdot)|_{\Omega} $  是从  $ \mathbf{D}^b(X; \Omega)^\circ \times \mathbf{D}^b(X; \Omega) $  到  $ \mathbf{D}^b(\Omega) $  的明确定义的函子。我们仍然用  $ \mu h o m $  来表示这个函子。

\[
\mu\mathcal{h}o m\colon\mathbf{D}^{b}(X;\Omega)^{\circ}\times\mathbf{D}^{b}(X;\Omega)\to\mathbf{D}^{b}(\Omega)\enspace.
\]

同构  $ \text{Hom}_{\mathbf{D}^{b}(X)}(G, F) \simeq H^{0}(T^{*}X; \mu \text{hom}(G, F)) $  （命题 4.4.2）定义了态射：

\[
\mathrm{H o m}_{{\tt D}^{b}(X;\varOmega)}(G,F)\to H^{0}(\varOmega;\mu\mathcal{h}o m(G,F))\enspace.
\]

事实上，如果  $ w: F \to F' $  是  $ \mathbf{D}^b(X) $  中的态射，而  $ w $  是  $ \Omega $  上的同构，则  $ w $  会导致  $ \mu hom(G, F)|_{\Omega} \simeq \mu hom(G, F')|_{\Omega} $  同构。

一般来说，态射 (6.1.4) 不是同构（参见练习 VI 6）。然而：

\statementlabel{定理 6.1.2}令 $p \in T^*X$  并让 $F$  和 $G$  属于 $\mathbf{D}^b(X)$  。那么自然态射 $\mathrm{Hom}_{\mathbf{D}^b(X; p)}(G, F) \to H^0(\mu \mathrm{hom}(G, F))_p$  就是同构。

\statementlabel{证明}如果 $p \in T_{X}^{*}X$ ，则无需证明。

假设  $X$  是向量空间，而  $p = (x_o; \xi_o) \in T^*X$  。使用与命题 4.4.4 相同的符号，我们有：

\[
\begin{aligned}{H^{0}(\mu{\mathcal{h}}o m(G,F))_{p}}&{{}\simeq\varinjlim_{\overrightarrow{U},\gamma}H^{0}(R\varGamma(U;R\mathcal{H}o m(\phi_{\gamma}^{-1}R\phi_{\gamma*}G_{U},F)))}\\ {}&{{}\simeq\varinjlim_{\overrightarrow{U},\gamma}\operatorname{H o m}((\phi_{\gamma}^{-1}R\phi_{\gamma*}G_{U})_{U},F)\enspace.}\\ \end{aligned}
\]

令 $\alpha$  为定理6.1.2 的态射。首先我们证明 $\alpha$  是单射的。让  $u \in \mathrm{Hom}_{\mathbf{D}^{b}(X)}(G, F)$  和  $\alpha(u) = 0$  。存在 $U$  、 $\gamma$  使得复合态射 $(\phi_{\gamma}^{-1}R\phi_{\gamma*}G_U)_{U} \to G \to F$  为零。由于  $(\phi_{\gamma}^{-1}R\phi_{\gamma*}G_U)_{U} \to G$  是  $\mathbf{D}^{b}(X; p)$  中的同构（命题 5.2.3），因此我们在  $\mathbf{D}^{b}(X; p)$  中得到  $u = 0$  。让我们证明 $\alpha$  是满射的。如果  $v \in H^0(\mu \text{hom}(G, F))_p$  ，则存在  $U$  、  $\gamma$  和  $w \in \mathrm{Hom}_{\mathbf{D}^{b}(X)}((\phi_{\gamma}^{-1}R\phi_{\gamma*}G_U)_{U}, F)$  ，它们代表  $v$  。由于  $(\phi_{\gamma}^{-1}R\phi_{\gamma*}G_U)_{U} \to G$  是  $\mathbf{D}^{b}(X; p)$  中的同构，因此  $w$  定义了  $\mathrm{Hom}_{\mathbf{D}^{b}(X; p)}(G, F)$  的一个元素， $\alpha$  的图像是  $v$  。   $\square$ 

以下推论是证明微支撑对合性的重要一步。

让 $ F \in \mathrm{Ob}(\mathbb{D}^{b}(X)) $  。考虑自然态射：

\[
\begin{aligned}{\operatorname{H o m}(F,F)\xrightarrow{\sim}}&{{}H^{0}(R\varGamma(X;R\mathcal{H}o m(F,F)))}\\ {\xrightarrow{\sim}}&{{}H^{0}(R\varGamma(T^{*}X;\mu\mathcal{h}o m(F,F)))}\\ {\to}&{{}\varGamma(T^{*}X;H^{0}(\mu\mathcal{h}o m(F,F))).}\\ \end{aligned}
\]

用 s 表示  $\mathrm{id}_{F} \in \mathrm{Hom}(F, F)$  在  $\Gamma(T^{*}X; H^{0}(\mu \mathrm{hom}(F, F)))$  中的图像。

\statementlabel{推论 6.1.3}其中一个有：

\[
\mathrm{supp}(\mathrm{s})=\mathrm{supp}(\mu\mathrm{hom}(F,F))=\mathrm{SS}(F)\enspace.
\]

特别是，如果  $ p \in T^*X $  ，则有：

\[
p\notin\mathrm{SS}(F)\Leftrightarrow\mu\mathrm{hom}(F,F)_{p}=0\enspace.
\]

\statementlabel{证明}包含 $\operatorname{supp}(s) \subset \operatorname{supp}(\mu \operatorname{hom}(F, F))$  是明确的，并且包含 $\operatorname{supp}(\mu \operatorname{hom}(F, F)) \subset \mathrm{SS}(F)$  由推论5.4.10 得出。

让 $p \in T^*X$ ， $p \notin \mathrm{supp}(s)$ 。那么根据定理 6.1.2， $s_p = 0$  意味着  $\mathrm{id}_F \in \mathrm{Hom}_{\mathbf{D}^b(X; p)}(F, F)$  消失。因此  $\mathbf{D}^b(X; p)$  中的  $F = 0$  ，这意味着  $p \notin \mathrm{SS}(F)$  。

在下文中，我们将讨论层上的微局部操作，例如逆像和正像。为此，我们需要改进 V §2 的微局部截止引理。

\statementlabel{命题 6.1.4（精炼的微局部截止引理）}令 $ x_o $  为流形 $ X $  的一个点， $ K $  为 $ T_x^*X $  的真闭凸锥体， $ U \subset K $  为开锥体。令 $ F \in \text{Ob}(\mathbb{D}^b(X)) $  和 $ W $  成为 $ K \cap \text{SS}(F) \setminus \{0\} $  的圆锥邻域。则存在 $ F' \in \text{Ob}(\mathbb{D}^b(X)) $ 和满足以下条件的态射 $ u: F' \to F $ ：

u 是 U 的同构，

\[
\pi^{-1}(x_{o})\cap\mathrm{S S}(F^{\prime})\subset W\cup\{0\}~.
\]

\statementlabel{证明}通过加厚 $K$ ，我们可以假设 $\{0\} \cup \operatorname{Int} K \supset \bar{U}$ 。我们还可以假设  $X$  是向量空间，而  $x_o = 0$  。让我们采用一个闭合的真凸锥  $\gamma$  ，使得  $K^{\circ a} \subset \gamma \subset U^{\circ a}$  和  $\gamma \setminus \{0\}$  具有  $C^1$  -边界  $\delta\gamma$  。我们在 $X$ 上取坐标 $(x) = (x_1, \ldots, x_n)$ ，这样：

\[
\gamma\subset\{0\}\cup\{x;x_{1}<0\}~.
\]

令  $(x; \xi)$  为  $T^*X$  上的关联坐标。从  $\gamma^{\circ a} \subset K$  开始，存在一个 0 的开邻域  $\Omega$  ，使得：

(6.1.10) $ W \supset G = \{\xi \in \gamma^{0a} \setminus \{0\}; \text{there exists}(x; \xi) \in \pi^{-1}(\Omega) \cap \mathrm{SS}(F)\} $ 。

以 $ \varepsilon > 0 $  为例：

\[
\Omega\supset\left\{x\in\gamma;x_{1}\geqslant-\varepsilon\right\}.
\]

现在我们将取 X 的闭子集 Z，其中  $ C^{1} $  -边界  $ \delta Z $  满足以下条件（参见图 6.1.1）：

\[
0\in\mathrm{I n t}Z~,
\]

\[
Z\subset\left\{x;x_{1}\geq-\varepsilon\right\}\ ,
\]

\[
\left\{\begin{matrix}{\delta Z{a n d}\delta\gamma{a r e t a n g e n t a t t h e i r i n t e r s e c t i o n,}}\\ {{m o r e p r e c i s e l y f o r a n y}x\in\delta Z\cap\delta\gamma,N_{x}^{*}(Z)^{a}=N_{x}^{*}(\gamma).}\\ \end{matrix}\right.
\]

Z 的构造留给读者。

现在我们设置  $ F' = \phi_y^{-1}R\phi_{y*}R\varGamma_Z(F) $  并用  $ u $  表示规范态射  $ F' \to F $  。我们将证明 $ F' $  和 $ u $  满足所需的属性。根据微局域截止引理（命题 5.2.3）， $ u $  是  $ \mathrm{Int}Z \times $  的同构

\begin{center}
\includegraphics[width=0.24\textwidth]{流行上的层_Chn-assets/img_in_image_box_444_163_731_521.png}
\end{center}

\begin{center}图。 6.1.1\end{center}

整数 $ \gamma^{\circ a} $ 。由于 Int  $ \gamma^{\circ a} \supset U $  ，满足 (6.1.7)。同样的命题意味着 $ \mathrm{SS}(F') \subset X \times \gamma^{\circ a} $ ，因此仍有待证明：

\[
\mathrm{i f}\ \xi\in\delta(\gamma^{\circ a})\backslash\{0\}\quad\mathrm{a n d}\quad(0;\xi)\in\mathrm{S S}(F^{\prime})\qquad\mathrm{t h e n}\ \xi\in G.
\]

令  $ q_i: X \times X \to X $  表示  $ i $  第  $ (i = 1, 2) $  投影，并令  $ s: X \times X \to X $  为映射  $ s(x, x') = x - x' $  。根据命题 3.5.4，我们有：

\[
F^{\prime}\simeq R q_{2*}(s^{-1}A_{\gamma}\otimes q_{1}^{-1}R\varGamma_{Z}(F))\enspace.
\]

由条件(6.1.9)和(6.1.13)可知 $ q_2 : s^{-1}(\gamma) \cap q_1^{-1} Z \to X $ 是真映射，我们可以应用命题5.4.4、5.4.5和5.4.14来计算 $ \mathrm{SS}(F') $ 。我们得到：

\[
\begin{cases}{{i f}\left(0;\xi\right)\in\mathrm{S S}(F^{\prime}),{~t h e n~t h e r e~e x i s t s}}\\ {(x;\xi)\in\mathrm{S S}(A_{\gamma})^{a}\cap\mathrm{S S}(R\varGamma_{Z}(F))}\\ \end{cases}.
\]

然后根据(6.1.11)和(6.1.13)，x包含在 $ \Omega $ 中。因此（6.1.15）的证明简化为：

\[
\begin{cases}\mathrm{if}\xi\neq0\mathrm{and}(x;\xi)\in\mathrm{SS}(A_{\gamma})^{\alpha}\cap\mathrm{SS}(R\Gamma_{Z}(F))\\ \mathrm{then}(x;\xi)\in\mathrm{SS}(F).\end{cases}
\]

从  $\xi \neq 0$  和  $(x; \xi) \in \mathrm{SS}(A_\gamma)^a$  开始，我们就有了  $x \in \delta\gamma$  。如果  $x \in \mathrm{Int} Z$  ，则  $F \simeq R\Gamma_Z(F)$  位于  $x$  ，因此  $(x; \xi) \in \mathrm{SS}(F)$  。如果  $x \in \delta Z$  ，则通过 (6.1.14)， $N_x^*(Z) = N_x^*(\gamma)^a = \mathbb{R}_{\geq 0} \xi$  。假设 $(x; \xi) \notin \mathrm{SS}(F)$ 。那么命题 5.4.8 表明  $\xi \in \mathrm{SS}(R\Gamma_Z(F)) \cap \pi^{-1}(x) \subset -\mathbb{R}_{\geq 0} \xi + (\mathrm{SS}(F) \cap \pi^{-1}(x))$  ，这意味着  $(x; \xi) \in \mathrm{SS}(F)$  。这是一个矛盾，至此证明完毕。   $\square$ 

我们需要对前一个声明进行双重声明。为此，我们首先对命题 5.2.3 提出双重声明。

\statementlabel{引理 6.1.5（双微局域截止引理）}令  $ \gamma $  为向量空间  $ X $  中的闭凸真圆锥，其中  $ \operatorname{Int}\gamma \neq \emptyset $  。让 $ s: X \times X \to X $  成为映射

 $(x,x') \mapsto x - x'$  并让  $q_i: X \times X \to X$  为  $i$  次投影 (  $i = 1,2$  )。那么对于具有紧凑支持的 $F \in \mathrm{Ob}(\mathbf{D}^b(X))$ ，有一个自然的态射：

\[
u:F\to R q_{2!}R\Gamma_{s^{-1}(\gamma^{a})}(q_{1}^{!}F)\enspace,
\]

这样

（一） $ \mathrm{SS}(Rq_{2}R\Gamma_{s^{-1}(\gamma^{a})}(q_{1}^{!}F)) \subset X \times \gamma^{\circ a} $ ，

(ii)  $ u $  是  $ X \times \text{Int}\gamma^{\circ a} $  的同构。

\statementlabel{证明}我们有一个态射链：

\[
F\simeq R q_{2*}R\mathcal{H}o m(A_{s^{-1}(0)},q_{1}^{!}F)\to R q_{2*}R\mathcal{H}o m(s^{-1}A_{\gamma^{a}},F)\enspace.
\]

从  $ A_{\gamma a} \simeq \phi_{\gamma}^{-1} R \phi_{\gamma*} A_{\{0\}} $  开始，我们根据命题 5.2.3 得到：

\[
\mathrm{SS}(A_{\gamma^{a}})\subset X\times\gamma^{\circ a}~,
\]

 $ A_{\gamma^{a}} \to A_{\{0\}} $  是  $ X \times \mathrm{Int} \gamma^{\circ a} $  的同构。

那么命题 5.4.5、5.4.14 和 5.4.4 意味着  $ \mathrm{SS}(Rq_{2*}\mathrm{R}\mathcal{H}om(s^{-1}A_{\gamma^{a}},F)) \subset X \times \gamma^{\circ a} $  。这表明 (i)。

另一方面，(6.1.21) 意味着 $ \mathrm{SS}(A_{\gamma\setminus{0}})\subset X\times(X^*\setminus\mathrm{Int}\gamma^{\circ a}) $  。因此，再次根据命题 5.4.5、5.4.14 和 5.4.4，我们得到  $ \mathrm{SS}(Rq_{2*}\mathcal{R}\mathcal{H}om(s^{-1}A_{\gamma^a\setminus{0}},F)\subset X\times(X^*\setminus\mathrm{Int}\gamma^{\circ a}) $  。这表明 (ii)。   $ \square $ 

\statementlabel{命题 6.1.6（双重精化微局域截止引理）}令  $x_{\circ}$  、  $K$  、  $U$  、  $F$  和  $W$  如命题 6.1.4 中所示。则存在满足条件(6.1.7)和(6.1.8)的态射 $u: F \to F'$ 。

\statementlabel{证明}我们可以假设  $F$  具有紧凑的支持。证明与命题 6.1.4 类似，只是我们使用对偶截止函子  $Rq_{2!}R\Gamma_{s^{-1}(\gamma a)}(q_{1}^{!}F_{Z})$  而不是  $\phi_{\gamma}^{-1}R\phi_{\gamma*}(R\Gamma_{Z}(F)$  。更准确地说，我们在前面的证明中采用 $\gamma$ ， $\Omega$ ， $G$ ，并且我们采用 $\varepsilon>0$ ，这样

\[
\Omega\supset\{x\in\gamma^{\circ a};x_{1}\leq\varepsilon\}~.
\]

类似地，我们采用  $C^{1}$  的闭子集  $Z$  和  $C^{1}$  -边界  $\delta Z$  ，满足 (6.1.12) 和以下条件：

\[
Z\subset\left\{x;x_{1}\leq\varepsilon\right\},
\]

(6.1.14)'  $ \delta Z $  和  $ \delta\gamma^a $  在它们的交点处相切，更准确地说，对于任何  $ x \in \delta Z \cap \delta\gamma^a $  、  $ N_x^*(Z)^a = N_x^*(\gamma^a) $  。

我们设定：

\[
F^{\prime}=R q_{2!}R\Gamma_{s^{-1}(\gamma^{a})}(q_{1}^{!}F_{Z})\enspace.
\]

那么根据前面的引理，存在一个规范态射  $ F_Z \to F' $  。让  $ u $  成为组合  $ F \to F_Z \to F' $  。与前一个证明中相同的论点表明 $ u $  满足所需的属性。我们把细节留给读者。   $ \square $ 

现在我们准备介绍层上的“微局部操作”。令 $ f: Y \to X $  为流形的态射， $ p \in Y \times_X T^*X $  。我们设置 $ p_Y = t f'(p) $ ， $ p_X = f_\pi(p) $ 。我们将保留这个符号直到命题 6.1.10 结束。我们将利用ind-objects和pro-objects（参见I §11）。

\statementlabel{定义 6.1.7}(i) 让 $F \in \mathrm{Ob}(\mathbf{D}^{b}(X; p_{X}))$  。我们用 $f_{\mu}^{-1}F$ （或 $f_{\mu}^{!}F$ ）表示 $\mathbf{D}^{b}(Y; p_{Y})$ 的原对象（或ind对象）“ $\lim_{F^{\prime} \to F} f^{-1}F^{\prime}$ ”（或“ $\lim_{F \to F^{\prime}} f^{!}F^{\prime}$ ”）。这里  $F^{\prime} \to F$  （分别为  $F \to F^{\prime}$  ）范围涵盖  $\mathbf{D}^{b}(X)$  中的态射范畴，它们是  $p_{X}$  的同构。我们将 $f_{\mu}^{-1}F$  称为 $F$  在 $p$  处的微局部逆像。

(ii) 让 $G \in \mathrm{Ob}(\mathbb{D}^b(Y; p_Y))$  。我们用  $f_!^\mu G$  (resp.  $f_*^\mu G$  ) 表示  $\mathbb{D}^b(X; p_X)$  的原对象 (resp. ind-object) “  $\lim_{G^{\prime} \to G} R f_! G'$  (resp. “  $\lim_{G \to G'} R f_* G'$  )” 。这里  $G' \to G$  (resp.  $G \to G'$  ) 范围涵盖  $\mathbb{D}^b(Y)$  中的态射范畴，它们是同构的在 $p_Y$  处，我们将 $G$  处的 $G$  称为 $f_!^\mu G$  （分别为 $f_*^\mu G$  ）的微局部直像（分别为微局部直像）。

这四个操作的关系如下。

\statementlabel{命题 6.1.8}让  $ F \in \mathrm{Ob}(\mathbb{D}^{b}(X; p_{X})) $  和  $ G \in \mathrm{Ob}(\mathbb{D}^{b}(Y; p_{Y})) $  。

(i) 存在自然同构：

\[
\mathrm{H o m}_{\mathbb{D}^{b}(X;p_{X})}\hat{}(f_{!}^{\mu}G,F)\simeq\mathrm{H o m}_{\mathbb{D}^{b}(Y;p_{Y})}\hat{}(G,f_{\mu}^{!}F)\enspace,
\]

\[
\mathrm{H o m}_{{\tt D}^{b}(X;p_{X})^{\vee}}(F,f_{*}^{\mu}G)\simeq\mathrm{H o m}_{{\tt D}^{b}(Y;p_{Y})^{\wedge}}(f_{\mu}^{-1}F,G)\enspace.
\]

(ii) 存在规范态射：

\[
f_{!}^{\mu}G\to f_{*}^{\mu}G\quad,
\]

\[
\omega_{Y/X}\otimes f_{\mu}^{-1}F\to f_{\mu}^{!}F\enspace.
\]

\statementlabel{证明}证明很简单，我们只需证明（6.1.22）。我们有：

\[
\begin{aligned}{\operatorname{H o m}_{\mathsf{D}^{b}(X;p_{X})}\hat{}(f_{!}^{\mu}G,F)}&{{}\simeq\varinjlim_{G^{\prime}\to G}\operatorname{H o m}_{\mathsf{D}^{b}(X;p_{X})}(R f_{!}^{\prime}G^{\prime},F)}\\ {}&{{}\simeq\varinjlim_{G^{\prime}\to G}\varinjlim_{F\to F^{\prime}}\operatorname{H o m}_{\mathsf{D}^{b}(X)}(R f_{!}^{\prime}G^{\prime},F^{\prime})}\\ {}&{{}\simeq\varinjlim_{G^{\prime}\to G}\varinjlim_{F\to F^{\prime}}\operatorname{H o m}_{\mathsf{D}^{b}(Y)}(G^{\prime},f^{!}F^{\prime})}\\ {}&{{}\simeq\varinjlim_{F\to F^{\prime}}\operatorname{H o m}_{\mathsf{D}^{b}(Y;p_{Y})}(G,f^{!}F^{\prime})}\\ {}&{{}\simeq\operatorname{H o m}_{\mathsf{D}^{b}(Y;p_{Y})}\vee(G,f_{\mu}^{!}F)~.\quad\Box}\\ \end{aligned}
\]

现在我们将研究一些确保微局部逆像或微局部正像属于 $ \mathbf{D}^{b}(Y; p_{Y}) $  或 $ \mathbf{D}^{b}(X; p_{X}) $  的条件。

\statementlabel{命题 6.1.9}让 $ F \in \mathrm{Ob}(\mathbb{D}^{b}(X, p_{X})) $  。

(i) 如果  $f'^{-1}(p_Y) \cap f_\pi^{-1}(\mathrm{SS}(F)) \subset \{p\}$  在  $p$  的邻域上，则  $f_\mu^{-1}F$  和  $f_\mu^!F$  属于  $\mathbf{D}^b(Y, p_Y)$  并且态射  $\omega_{Y/X} \otimes f_\mu^{-1}F \to f_\mu^!$  是同构。此外，对于  $p$  的任何邻域  $W$  ，有：

(6.1.26)  $ \mathrm{SS}(f_{\mu}^{-1}F) \subset \iota f'(W \cap f_{\pi}^{-1}(\mathrm{SS}(F))) $  在  $ p_{Y} $  的邻域上。

(ii) 设  $ F \in \text{Ob}(\mathbb{D}^b(X)) $  并假设：

(a) f 对于 F 来说是非特征值，

\[
\mathrm{(b)}\mathcal{~t f~}^{\prime-1}(p_{\mathtt{Y}})\cap f_{\pi}^{-1}(\mathtt{S S}(F))\subset\{p\}.
\]

然后 $ f_{\mu}^{-1}F \simeq f^{-1}F $  和 $ f_{\mu}^{!}F \simeq f^{!}F $  。

\statementlabel{证明}如果 $ p_X \in T_X^*X $ ，则 $ f $ 对于 $ F $ 来说是非特征的，并且结果来自命题5.4.13。

假设 $p_X \in T^*X$ ，让我们首先证明关于 $f_\mu^{-1}$ 的结果。令  $F$  如 (i) 中所示，并设置  $x_\circ = \pi(p_X)$  、  $y_0 = \pi(p_Y)$  、  $V = \text{Ker}(T_x^*X \to T_y^*Y)$  。取一个适当的闭合凸锥 $K$ 和一个开放凸锥 $U$ ，使得 $p_X \in U \subset K$ ， $K \cap V \subset \{0\}$ 和 $f_x' f''^{-1}(p_Y) \cap \text{SS}(F) \cap K \subset \{p_X\}$ 。

通过精化的微局域截止引理，存在 $u: F' \to F$ ，使得 $u$  与 $p_X$  和 $F'$  处的同构满足条件(ii) (a) 和(ii) (b)。

此外，如果  $F'' \to F'$  是  $p_X$  和  $F'$  和  $F''$  处的同构满足 (ii) 中的条件，那么，通过将此态射嵌入到一个杰出的三角形  $F'' \longrightarrow F' \longrightarrow F_o \xrightarrow{+1}$  中，我们发现  $f$  对于  $F_o$  和  $f'f'^{-1}(p_Y) \cap f_{\pi}^{-1}(\mathrm{SS}(F_o)) = \emptyset$  来说是非特征的。因此，命题 5.4.13 提出了 $p_Y \notin \mathrm{SS}(f^{-1}F_o)$ 。这意味着  $f^{-1}F'' \to f^{-1}F'$  是  $\mathbf{D}^b(Y; p_Y)$  的同构。根据  $f_{\mu}^{-1}F$  的定义，“归纳极限”可以采用态射  $F' \to F$  的范畴，其中  $F'$  满足 (ii) 中的条件。那么所有  $f^{-1}F'$  在  $\mathbf{D}^b(Y; p_Y)$  中都是同构的，这证明了 (i) 和 (ii) 中关于  $f_{\mu}^{-1}$  的陈述。关于  $f_{\mu}^!$  的陈述也同样得到了证明，使用对偶精化微局域截止引理，并且同构  $\omega_{Y/X} \otimes f_{\mu}^{-1}F \simeq f_{\mu}^!F$  由 (ii) 和命题 5.4.13 得出。   $\square$ 

\statementlabel{命题 6.1.10}让 $ G \in \mathrm{Ob}(\mathbb{D}^{b}(Y; p_{Y})) $  。

(0) 我们有：

\[
f_{!}^{\mu}G\simeq\varprojlim_{V}R f_{!}G_{V}\simeq\varprojlim_{K}R f_{!}R\varGamma_{K}(G),
\]

\[
f_{*}^{\mu}G\simeq\operatorname*{l i m}_{\overrightarrow{V}}{R f}_{*}{R}\varGamma_{V}(G)\simeq\operatorname*{l i m}_{\overrightarrow{K}}{R f}_{*}(G_{K})\enspace,
\]

其中 $V$ （或 $K$ ）范围在 $y_{\circ}=\pi(p_{Y})$ 的开放（或封闭）邻域系统上。 （请注意，上面的同构是在范畴  $\mathbf{D}^{b}(X;p_{X})^{\wedge}$  和  $\mathbf{D}^{b}(X;p_{X})^{\vee}$  中定义的。）

(i) 如果  $ f_{\pi}^{-1}(p_X) \cap \mathfrak{t}f^{\prime-1}(\mathrm{SS}(G)) \subset \{p\} $  在  $ p $  的邻域上，则  $ f_!^\mu G $  和  $ f_*^\mu G $  属于  $ \mathrm{D}^b(X; p_X) $  并且态射  $ f_!^\mu G \to f_*^\mu G $  是同构。此外，对于  $ p $  的任何邻域  $ W $  ，有：

(6.1.27) p<sub>x</sub> 邻域上的 SS(f<sup>μ</sup>G) ⊂ f<sub>π</sub>(W ∩<sup>μ</sup>f<sup>γ</sup><sup>-1</sup>(SS(G)))

(ii) 如果  $\mathrm{supp}(G)$  对  $X$  是真映射，并且如果  $f_{\pi}^{-1}(p_{X})\cap{}^{t}f^{\prime-1}(\mathrm{SS}(G))\subset{p}$  则  $f_{!}^{\mu}G\simeq f_{*}^{\mu}G\simeq Rf_{*}^{*}G$  。

\statementlabel{证明}同构的“  $\lim_{\overleftarrow{V}}Rf_1G_V\simeq\lim_{\overleftarrow{K}}Rf_1R\Gamma_K(G)$  和”  $\lim_{\overrightarrow{V}}Rf_1R\Gamma_V(G)\simeq\lim_{\overrightarrow{K}}Rf_1G_K$  是显而易见的，因为存在自然态射  $G_V\to R\Gamma_K(G)$  、  $G_K\to R\Gamma_V(G)$  if  $V\subset K$  和  $R\Gamma_K(G)\to G_V$  、  $R\Gamma_V(G)\to G_K$  if  $K\subset V$  。

(a) 首先证明如果 G 满足 (i) 中的条件，则  $ y_{0} $  的开邻域集合 V 满足满足条件：

\[
\bar{V}\to X\mathrm{i s p r o p e r},
\]

\[
{}^{t}f^{\prime-1}(\mathbf{S}\mathbf{S}(G_{V}))\cap{f}_{\pi}^{-1}(p_{X})\subset\{p\}\mathrm{,}
\]

是  $y_{\circ}$  的邻域系统。

问题是局部的，我们可以假设 Y 和 X 是向量空间，  $ p_Y = (0; \eta_\circ) $  和  $ p_X = (0; \xi_\circ) $  。

那么存在  $ \varepsilon > 0 $  这样：

(6.1.30) $ \{y \in Y; f(y) = 0 \text{ and } (y; 'f'(y) \cdot \xi_0) \in \mathrm{SS}(G)\} \subset \{0\} \cup \{y; |y| > 2\varepsilon\} $ 。

因此存在  $ \varepsilon_{1} > 0 $  这样：

\[
(y;{}^{t}f^{\prime}(y)\cdot\xi_{\circ}+\eta)\notin\mathbf{S S}(G)\qquad\mathrm{i f}\qquad|y|=\varepsilon\qquad\mathrm{a n d}\qquad|\eta|\leqslant\varepsilon_{1}.
\]

现在，设置：

\[
V=\left\{y;\varepsilon_{1}|y|<\langle f(y),\xi_{\circ}\rangle+\varepsilon_{1}\varepsilon\right\}.
\]

那么只要证明V满足(6.1.29)就足够了。

通过（6.1.31）， $ \mathrm{SS}(G) \cap N^*(V) \cap \pi_{Y}^{-1}(f^{-1}(0)) \subset T_{Y}^{*}Y $ 。因此，根据命题 5.4.8，我们有：

\[
\mathrm{S S}(G_{V})\subset\mathrm{S S}(G)+N^{*}(V)^{a}\qquad\mathrm{o v e r}\qquad f^{-1}(0).
\]

由于  $ t f'^{-1}\mathrm{SS}(G_V) \cap f_\pi^{-1}(p_X) \cap \pi^{-1}(V) \subset \{p\} $  ，足以证明如果  $ y \in \overline{V} \setminus V $  ，则  $ (y; t f'(y) \cdot \xi_0) \notin \mathrm{SS}(G_V) $  。

到(6.1.33)，存在 $ k \geqslant 0 $  和 $ (y; \eta) \in \mathrm{SS}(G) $  使得 $ t f'(y) \cdot \xi_0 = \eta + k(-t f'(y) \cdot \xi_0 + \varepsilon_1 y/|y|) $  。

因此 $ \eta = (1 + k)^t f'(y) \cdot \xi_0 - k \varepsilon_1 y/|y| $  。

由于  $ |k\varepsilon_1y/(1+k)|y|\leq\varepsilon_1 $  ，这与 (6.1.31) 相矛盾。

(b) 让我们证明同构 $ f_!^\mu G \simeq \lim_{\leftarrow} R f_! G_V $  。令  $ G' \to G $  为  $ p_Y $  的同构。我们将其嵌入到一个独特的三角形 $ G' \longrightarrow G \longrightarrow G_0 \xrightarrow{+1} $ 中。

然后  $p_Y \notin \mathrm{SS}(G_0)$  ，并且通过 (a) 存在一个由  $V$  组成的  $y_0$  的开放邻域系统，使得  $\overline{V} \to X$  是正确的并且  $t'f'^{-1}(\mathrm{SS}(G_{0,V})) \cap f_\pi^{-1}(p_X) = \emptyset$  。对于这样的  $V$  ， $p_X \notin \mathrm{SS}(R f_t' G_{0,V})$  根据命题 5.4.4，因此  $R f_t' G_V' \to R f_t' G_V$  是  $\mathbf{D}^b(X; p_X)$  中的同构。对  $V$  和  $G$  取“射影极限”，我们得到：

\[
\varprojlim_{G^{\prime}}R f_{!}G^{\prime}\simeq\varprojlim_{G^{\prime},V}R f_{!}G_{V}^{\prime}\simeq\varprojlim_{V}R f_{!}G_{V}\enspace.
\]

(c) 如果 $G$ 满足(ii)中的条件，则取 $V$ 满足(6.1.28)和(6.1.29)， $Rf_i G_V \to Rf_i G$ 是 $\mathbf{D}^b(X; p_X)$ 中的同构。因此我们得到 $f_i^\mu G \simeq Rf_i G$ 。

(d) 关于 $ f_{*}^{\mu}G $ 的陈述同样可以通过将步骤(a)中的条件(6.1.29)替换为以下条件来证明：

\[
{}^{t}f^{\prime-1}(\mathbf{S S}(R\varGamma_{V}(G)))\cap f_{\pi}^{-1}(p_{X})\subset\{p\}\enspace.
\]

细节留给读者。

(e) 为了证明(i)，我们可以通过步骤(a)假设 $G$ 满足(ii)中的条件。因此，结果来自步骤 (c)（以及  $f_{*}^{\mu}(G)$  的相应结果）。 □

最后让我们陈述以下结果，其证明直接来自命题 5.4.1。

\statementlabel{命题 6.1.11}令  $Y$  和  $X$  为两个流形，并令  $p_Y \in T^*Y$  、  $p_X \in T^*X$  、  $p = (p_X, p_Y) \in T^*(X \times Y)$  。

令  $F_1$  和  $F_2$  （分别为  $G_1$  和  $G_2$  ）属于  $\mathbf{D}^b(X)$  （分别为  $\mathbf{D}^b(Y)$  ）。然后有一个规范的同态：

\[
\begin{aligned}{\operatorname{H o m}_{\mathsf{D}^{b}(X;p_{X})}(F_{1},F_{2})}&{{}\times\operatorname{H o m}_{\mathsf{D}^{b}(Y;p_{Y})}(G_{1},G_{2})}\\ {\to\operatorname{H o m}_{\mathsf{D}^{b}(X\times Y;p)}}&{{}\left(\begin{matrix}{F_{1}\boxtimes G_{1},F_{2}\boxtimes G_{2}}\\ \end{matrix}\right).}\\ \end{aligned}
\]

\subsection{余切丛中的法锥}
在下一节中，我们将研究微支撑在各种操作下的行为。为了表述结果，我们需要首先引入一些关于余切丛中二次曲线子集的新操作。这些结构利用了 IV §1 中引入的“正圆锥”概念和  $ T^*X $  的辛结构（参见附录）。设 X 为流形。我们通过  $ -H $  识别  $ T^*T^*X $  和  $ TT^*X $  ，其中 H 是哈密顿同构。如果  $ (x) = (x_1, \ldots, x_n) $  是 X 上的局部坐标系， $ (x; \xi) $  是  $ T^*X $  上的关联坐标（因此规范的 1-形式  $ \alpha_X $  由  $ \langle \xi, dx \rangle = \sum_j \xi_j dx_j $  给出），我们有：

\[
-H(\langle\lambda,\mathrm{d}x\rangle+\langle\mu,\mathrm{d}\xi\rangle)=\left\langle\lambda,\frac{\partial}{\partial\xi}\right\rangle-\left\langle\mu,\frac{\partial}{\partial x}\right\rangle,
\]

如果  $A$  是  $T^{*}X$  的光滑圆锥拉格朗日子流形，则  $-H$  会产生同构  $T^{*}A \simeq T_{A}T^{*}X$  。特别是如果  $M$  是  $X$  的子流形，我们有同构：

\[
T^{*}T_{M}X\simeq T^{*}T_{M}^{*}X\simeq T_{T_{M}^{*}X}T^{*}X\enspace.
\]

这里第一个同构是通过命题 5.5.1 和  $ E = T_M X $  获得的。令 $ (x', x'') $  为 $ X $  上的局部坐标系，使得 $ M = \{(x', x''); x' = 0\} $  并令 $ (x', x''; \xi', \xi'') $  表示 $ T^*X $  上的关联坐标。通过这些坐标，我们可以识别  $ T_M X $  和  $ X $  ，以及  $ T^*(T_M X) $  和  $ T^*X $  。我们还将  $ T_{M \times X}(T^*X) $  与  $ T^*X $  识别。那么(6.2.2)中的同构描述为：

\[
\begin{array}{c c c c c}{T^{*}T_{M}X}&{\widetilde{\quad\sim\quad}}&{T^{*}T_{M}^{*}X}&{\widetilde{\quad\sim\quad}}&{T_{T_{M}^{*}X}T^{*}X}\\ {}&{}&{}&{}&{}\\ {\uplus}&{}&{\uplus}&{}&{\uplus}\\ \end{array}
\]

\[
(x^{\prime},x^{\prime\prime};\xi^{\prime},\xi^{\prime\prime})\longleftrightarrow(\xi^{\prime},x^{\prime\prime};-x^{\prime},\xi^{\prime\prime})\longleftrightarrow(x^{\prime},x^{\prime\prime};\xi^{\prime},\xi^{\prime\prime}).
\]

令 $p$  表示投影 $T_M^*X \to M$  ， $\dot{p}$  表示其对 $\dot{T}_M^*X$  的限制。我们有映射（参见（5.5.9））：

\[
\begin{aligned}{T^{*}T_{M}^{*}X}&{{}\xleftarrow{\quad t}p^{\prime}\quad T_{M}^{*}X}&{\times}&{{}\left.T^{*}M\xrightarrow[p_{*}]{\quad}\right.}\\ {\quad\cup}&{{}\quad}&{\cup}&{{}}\\ {T^{*}\dot{T}_{M}^{*}X}&{{}\xleftarrow{\quad t}\dot{p}^{\prime}\quad\dot{T}_{M}^{*}X}&{\times}&{{}\left.T^{*}M\xrightarrow[\dot{p}_{*}]{\quad}\right.}\\ \end{aligned}\left.\begin{aligned}{T^{*}M}\\ {\quad\bigg\|\quad}\\ \end{aligned}\right.
\]

通过上面的坐标，  $ T_M^*X \times_M T^*M $  与  $ T^*T_M^*X $  的子集  $ \{x' = 0\} $  同构，而  $ p_\pi $  由  $ (0, x''; \xi', \xi'') \mapsto (x''; \xi'' $  给出。

回想一下（参见（5.5.10）），  $T^*M$  通过  $T_M^*X$  的零部分和映射  $p'$  嵌入到  $T_{T_M^*X}T^*X\simeq T^*T_M^*X$  中。使用与上面相同的坐标，该嵌入由  $(x'';\xi'')\mapsto(0,x'';0,\xi'').$  描述

\statementlabel{引理 6.2.1}令  $ (x', x'') $  为 X 上与  $ M = \{(x', x''); x' = 0\} $  的局部坐标系，并令  $ (x', x''; \xi', \xi'') $  为  $ T^*X $  上的关联坐标。令 A 为  $ T^*X $  的圆锥子集。然后

（一） $ p_{\pi}^{t}p^{\prime-1}(C_{T_{M}^{*}X}(A)) = T^{*}M \cap C_{T_{M}^{*}X}(A) $ 。

(ii)  $ (x''; \xi'')\in\tilde{T}^*M\cap C_{T_M^*X}(A)\Leftrightarrow $  A 中存在序列  $ \{(x'_n, x_n''; \xi'_n, \xi''_n)\} $  使得：

\[
\left\{\begin{aligned}{}&{{}(x_{n}^{\prime\prime};\xi_{n}^{\prime\prime})\xrightarrow[n]{}(x_{\circ}^{\prime\prime};\xi_{\circ}^{\prime\prime})~,}\\ {}&{{}|x_{n}^{\prime}|\xrightarrow[n]{}0~,}\\ {}&{{}|x_{n}^{\prime}||\xi_{n}^{\prime}|\xrightarrow[n]{}0~.}\\ \end{aligned}\right.
\]

(iii)  $ (x''_o; \xi''_o) \in \dot{p}_{\pi^t} \dot{p}^{' - 1}(C_{\dot{T}_M^*X}(A)) \Leftrightarrow $  A 中存在序列  $ \{(x'_n, x''_n; \xi'_n, \xi''_n)\} $  满足 (6.2.4) 并且：

\[
|\xi_{n}^{\prime}|\xrightarrow[n]{}+\infty.
\]

回想一下， $ C_{T_M^*X}(A) $  表示沿着 $ T_M^*X $  到 $ A $  的法向锥体，参见。四§1。

\statementlabel{证明}(i) 集合 $ C_{T_{M}^{*}X}(A) $  是双圆锥的。因此 (i) 是 (5.5.11) 的一个特例。

(ii), (iii) 让 $ (x', x''; \xi', \xi'') $  表示 $ T_{T_M^*X}(T^*X) $  上的坐标。子流形  $ T_M^*X \times_M T^*X $  由方程  $ \{x' = 0\} $  定义，映射  $ p_\pi $  由  $ (0, x''; \xi', \xi'') \mapsto (x''; \xi'') $  给出。嵌入  $ T^*M \hookrightarrow T_{T_M^*X}(T^*X) $  由  $ (x''; \xi'') \mapsto (0, x''; 0, \xi'') $  给出。

(a) 令 $ \left\{(x_{n}^{\prime}, x_{n}^{\prime\prime}; \xi_{n}^{\prime}, \xi_{n}^{\prime\prime})\right\} $  为 $ A $  中满足(6.2.4) 的序列。

首先假设序列 $\{\xi_n^\prime\}$  是有界的。通过提取子序列，我们可以假设  $\xi_n^\prime \to \xi_\circ$  对应于某些  $\xi_\circ^\prime$  。令 $\{t_n\}$  为 $\mathbb{R}^+$  中的序列，使得 $t_n \to 0$  、 $t_n^{-1} x_n^\prime \to 0$  。我们有：

\[
\left\{\begin{aligned}{}&{{}(x_{n}^{\prime},x_{n}^{\prime\prime};t_{n}\xi_{n},t_{n}\xi_{n}^{\prime\prime})\longrightarrow(0,x_{\circ}^{\prime\prime};0,0),}\\ {}&{{}t_{n}^{-1}(x_{n}^{\prime};t_{n}\xi_{n}^{\prime\prime})\mathop{\longrightarrow}_{n}(0;\xi_{\circ}^{\prime\prime}).}\\ \end{aligned}\right.
\]

因此  $ (x_{0}^{\prime\prime};\xi_{0}^{\prime\prime}) $  属于  $ T^{*}M\cap C_{T_{M}^{*}X}(A) $  。

现在假设 $\{\xi_n^\prime\}$  是无界的。通过提取子序列，我们可以假设  $|\xi_n^\prime| \to \infty$  和  $\xi_n'' / |\xi_n'| \to \xi_0^\prime \neq 0$  。设置 $t_n = |\xi_n'|^{-1}$ 。我们有：

\[
\left\{\begin{aligned}{}&{{}(x_{n}^{\prime},x_{n}^{\prime\prime};t_{n}\xi_{n}^{\prime},t_{n}\xi_{n}^{\prime\prime})\xrightarrow[n]{}(0,x_{\circ}^{\prime\prime};\xi_{\circ}^{\prime},0),}\\ {}&{{}t_{n}^{-1}(x_{n}^{\prime};t_{n}\xi_{n}^{\prime\prime})\xrightarrow[n]{}(0;\xi_{\circ}^{\prime\prime}).}\\ \end{aligned}\right.
\]

因此 $(0,x_{0}^{\prime\prime};\xi_{0}^{\prime},\xi_{0}^{\prime\prime})\in C_{\dot{T}_{M}^{*}X}(A)$  和 $(x_{0}^{\prime\prime};\xi_{0}^{\prime\prime})\in\dot{p}_{\pi}^{t}p^{\prime-1}C_{\dot{T}_{M}^{*}X}(A)$  。

(b) 让 $ (x_{0}^{\prime\prime}; \xi_{0}^{\prime\prime}) \in p_{\pi}^{t} p^{t-1}(C_{T_{M}^{*}X}(A)) $  。

存在  $\xi_o$  使得  $(0, x_o^w; \xi_o', \xi_o'') \in C_{T_M^*X}(A)$  。因此， $A$  中存在序列  $\{(x_n^w, x_n^w; \xi_n', \xi_n'')\}$  ，  $\mathbb{R}^+$  中存在  $\{t_n\}$  ，这样：

\[
\left\{\begin{aligned}{}&{{}(x_{n}^{\prime},x_{n}^{\prime\prime};\xi_{n}^{\prime},\xi_{n}^{\prime\prime})\longrightarrow(0,x_{\circ}^{\prime\prime};\xi_{\circ}^{\prime},0),}\\ {}&{{}t_{n}(x_{n}^{\prime};\xi_{n}^{\prime\prime})\longrightarrow\left(0;\xi_{\circ}^{\prime\prime}\right).}\\ \end{aligned}\right.
\]

序列 $\{(x_n', x_n'; t_n\xi_n', t_n\xi_n')\}$  满足(6.2.4)。如果  $(x_o', \xi_o') \in \dot{p}_x^t \dot{p}'^{-1}C_{T_M^*X}(A)$  ，那么我们可以假设  $\xi_o' \neq 0$  。如果  $t_n \to \infty$  ，则满足  $|t_n\xi_n'| \to \infty$  和 (6.2.5)。如果  $\{t_n\}$  有界，则  $\xi_o'' = 0$  。在本例中，采用序列  $\{s_n\}$  和  $s_n | \xi''_n | \to 0 s_n \to \infty $  、 $ s_n |x'_n| \to 0 $  、 $ A $  中的序列 $ \{(x'_n, x''_n; s_n \xi'_n, s_n \xi''_n)\} $  满足(6.2.4) 和(6.2.5)。   $ \square $ 

\statementlabel{备注 6.2.2}我们有：

\[
{}^{t}\dot{p}^{{}^{\prime}-1}(C_{\dot{T}_{M}^{*}X}(A))={}({}^{t}\dot{p}^{{}^{\prime}-1}(C_{T_{M}^{*}X}(A)))\cap\dot{T}_{M}^{*}X\mathop{\times}\limits_{{M}}T^{*}M\enspace.
\]

因此，用  $ \dot{p}^{\prime-1}(C_{T_{M}^{*}X}(A)) $  表示该集合不会产生混淆。

现在让  $Y$  和  $X$  为两个流形，让  $f$  为从  $Y$  到  $X$  的态射。我们将 $Y$  与 $X \times Y$  中的 $f$  的图进行识别。我们用  $q$  表示投影  $T_{Y}^{*}(X \times Y) \to Y$  ，用  $\dot{q}$  表示对  $\dot{T}_{Y}^{*}(X \times Y)$  的限制。请注意， $T_{Y}^{*}(X \times Y)$  通过第一个投影与  $Y \times_{X} T^{*}X$  标识， $q$  与  $f_{\pi}$  标识。

\statementlabel{定义 6.2.3}令  $A$  （分别为  $B$  ）为  $T^{*}X$  （分别为  $T^{*}Y$  ）的圆锥曲线子集。我们设定：

\[
C_{\mu}(A,B)=C_{T_{Y}^{*}(X\times Y)}(A\times B^{a}),
\]

（这是  $ T_{T_Y^*(X \times Y)} T^*(X \times Y) \simeq T^*(Y \times_X T^*X) $  的闭双圆锥子集，

\[
\begin{aligned}f^{\#}(A,B)&=q_{\pi}^{t}q^{\prime-1}(C_{\mu}(A,B))\ ,\\&=T^{*}Y\cap C_{\mu}(A,B)\ ,\end{aligned}
\]

\[
f_{\infty}^{\#}(A,B)=\dot{q}_{\pi}{}^{t}\dot{q}^{\prime-1}(C_{\mu}(A,B)).
\]

(iv) 我们设置 $ f^{\#}(A) = f^{\#}(A, T_{Y}^{*}Y) $  和 $ f_{\infty}^{\#}(A) = f_{\infty}^{\#}(A, T_{Y}^{*}Y) $  。

(v) 如果  $Y = X$  和  $f$  是恒等式，我们设置  $A \widehat{+} B = f^{*}(A, B^{a}) \subset T^{*}X$  和  $A \widehat{+}_{\infty} B = f_{\infty}^{*}(A, B^{a}) \subset T^{*}X$  。

请注意，当 Y = X 且 f 为恒等式时，(i) 中  $ C_{\mu}(A, B) $  的定义与  $ C(A, B) $  一致。

另请注意  $ C_{\mu}(A,B), f^{\#}(A,B), f_{\infty}^{\#}(A,B) $  等是闭集。

\statementlabel{命题 6.2.4}(i) 假设 f 是封闭嵌入。然后：

\[
\begin{aligned}{}&{{}f^{\#}(A)=T^{*}Y\cap C_{T_{\gamma}^{*}X}(A),}\\ {}&{{}f_{\infty}^{\#}(A)=\dot{p}_{\pi}{}^{t}\dot{p}^{\prime-1}(C_{\dot{T}_{\gamma}^{*}X}(A)),}\\ \end{aligned}
\]

其中 p 表示投影  $ T_{Y}^{*}X \to Y $  。

(ii) 让 $ \pi $  表示投影 $ T^*X \to X $  。那么如果  $ A $  和  $ B $  是  $ T^*X $  的圆锥曲线子集：

\[
\hat{A}+B=T^{*}X\cap C(A,B^{a}),
\]

\[
\stackrel{\leftrightarrow}{A}+\stackrel{\leftrightarrow}{\beta}B=\dot{\pi}_{\pi}^{t}\dot{\pi}^{\prime-1}(C(A,B^{a}))
\]

(iii) 设 $(x)$ （分别为 $(y)$ ）为 $X$ （分别为 $Y$ ）上的局部坐标系，并令 $(x;\xi)$ （分别为 $(y;\eta)$ ）为 $T^*X$ （分别为 $T^*Y$ ）上的相关坐标。然后：

(a)  $(y_{\circ};\eta_{\circ})\in f^{\#}(A,B)\Leftrightarrow$  在  $A\times B$  中存在一个序列  $\{(x_{n};\xi_{n}),(y_{n};\eta_{n})\}$  使得：

\[
\left\{\begin{aligned}{}&{{}y_{n}\xrightarrow[n]{}y_{\circ},x_{n}\xrightarrow[n]{}f(y_{\circ})~,}\\ {}&{{}({}^{t}f^{\prime}(y_{n})\cdot\xi_{n}-\eta_{n})\xrightarrow[n]{}{\eta}_{\circ},|x_{n}-f(y_{n})||\xi_{n}|\xrightarrow[n]{}0~,}\\ \end{aligned}\right.
\]

(b)  $(y_{\circ};\eta_{\circ})\in f_{\infty}^{\#}(A,B)\Leftrightarrow$  在  $A\times B$  中存在序列  $\{(x_{n};\xi_{n}),(y_{n};\eta_{n})\}$  满足 (6.2.6) 并且还：

\[
|\xi_{n}|\xrightarrow[n]{}+\infty.
\]

\statementlabel{证明}(i) 和 (ii) 紧随定义而来。

(iii) 在  $ T^{*}(X \times Y) $  上的坐标发生变化后，由引理 6.2.1 得出：

\[
(x,y;\xi,\eta)\mapsto(x-f(y),y;\xi,\eta+{^{\prime \prime}f^{\prime}(y)\cdot\xi})\;.\quad\Box
\]

\statementlabel{备注 6.2.5}令 A 为  $ T^{*}X $  的闭圆锥曲线子集。然后：

\[
{\mathfrak{f}}^{\#}(A)={\mathfrak{f}}^{\prime}f_{\pi}^{-1}(A)\cup{\mathfrak{f}}_{\infty}^{\#}(A)\enspace.
\]

类似地，令  $A$  和  $B$  为  $T^{*}X$  的两个闭圆锥曲线子集。然后：

\[
\hat{A}+\hat{B}=(A+B)\cup(\hat{A}+\hat{B})\ .
\]

\statementlabel{备注 6.2.6}假设  $A$  是  $T^*X$  的闭圆锥曲线子集。那么 $f_\infty^\#(A) = \varnothing$  当且仅当 $f$  对于 $A$  来说是非特征的（参见定义5.4.12）。

类似地  $ A +_{\infty} B = \emptyset $  ，对于  $ T^* X $  的两个闭圆锥子集，当且仅当  $ A \cap B^a \subset T_X^* X $  。

鉴于前面的评论，很自然地设置：

\statementlabel{定义 6.2.7}(i) 令 $A$  为 $T^*X$  的闭圆锥子集， $V$  为 $T^*Y$  的子集。如果  $f_\infty^*(A) \cap V = \emptyset$  ，我们说  $f$  对于  $V$  上的  $A$  来说是非特征的。

(ii) 如果 $F \in \mathrm{Ob}(\mathbb{D}^{b}(X))$  ，我们说 $f$  对于 $V$  上的 $F$  是非特征的，如果 $f$  对于 $V$  上的 $\mathrm{SS}(F)$  来说是非特征的。

如果 f 是封闭嵌入，则也可以说“Y 是非特征的”，而不是“f 是非特征的”。

\statementlabel{备注 6.2.8}命题 6.2.4 的以下特殊情况将特别有用。

(i) 令  $ (x', x'') $  为  $ X $  、  $ M = \{x' = 0\} $  上的局部坐标系，并令  $ j $  表示嵌入  $ M \to X $  。让  $ (x', x''; \xi', \xi'') $  表示  $ T^*X $  上的关联坐标。那么，如果  $ A $  是  $ T^*X $  的闭圆锥子集，我们有：  $ (x_0''; \xi_0'') \in j^\#(A) \Leftrightarrow $  在  $ A $  中存在序列  $ \{(x_n', x_n''; \xi_n', \xi_n'')\} $  ，使得  $ x_n' \to 0 $  、  $ x_n'' \to x_0'' $  、  $ \xi_n'' \to \xi_0'' $  和  $ |x_n'| | \xi_n' | \to 0 $  。

(ii) 令 $(x)$  为 $X$  上的局部坐标系， $(x;\xi)$  为 $T^*X$  上的相关坐标。令  $A$  和  $B$  为  $T^*X$  的两个闭圆锥曲线子集。那么：  $(x_\circ;\xi_\circ)\in A\oplus B\Leftrightarrow$  在  $A$  中存在序列  $\{(x_n;\xi_n)\}$  ，在  $B$  中存在序列  $\{(y_n;\eta_n)\}$  ，使得  $x_n\to x_\circ$  、  $y_n\to x_\circ$  、  $\xi_n+\eta_n\to\xi_\circ$  和  $|x_n-y_n||\xi_n|\to0$  。

另请注意  $ (x_{\circ}, \xi_{\circ}) \in A +_{\infty} B $  当且仅当存在满足上述条件的序列  $ \{(x_{n}; \xi_{n})\} $  和  $ \{(y_{n}; \eta_{n})\} $  以及  $ |\xi_{n}| \to \infty $  。

\subsection{直像}
在 V §4 中，我们研究了微支撑在各种操作下的行为，例如真直像或非特征逆像。在这里我们扩展这些结果。

我们将利用 IV §1 中介绍的普通圆锥的概念，以及上一节中介绍的操作  $f^{\#}$  、  $+,$  等。

\statementlabel{定理 6.3.1}令  $\Omega$  为  $X$  和  $j$  嵌入  $\Omega \subsetneq X$  的开放子集。让 $F \in \mathrm{Ob}(\mathbf{D}^{b}(\Omega))$  。然后：

\[
\mathrm{S S}(R j_{*}F)\subset\mathrm{S S}(F)\hat{+}N^{*}(\varOmega),
\]

\[
\mathrm{SS}(R j,F)\subset\mathrm{SS}(F)\hat{+}N^{*}(\Omega)^{a}.
\]

\statementlabel{证明}我们将通过将  $\Omega$  移动到通用位置来证明这一点，以便  $F$  对于  $N^{*}(\Omega)$  来说是非特征的。

(i) 我们可以假设 $X$  是一个向量空间。让 $(x_o; \xi_o) \notin \mathrm{SS}(F) + N^*(\Omega)$  。为了显示 $(x_o; \xi_o) \notin \mathrm{SS}(Rj_*F)$ ，我们可以假设 $x_o \in \mathrm{supp}(F)$ ， $\xi_o \neq 0$ ， $x_o \in \partial \Omega$ ，因此 $(x_o; \xi_o) \notin \mathrm{SS}(F)$ 。由于  $\xi_o$  不属于  $N_x^*(\Omega)$  ，我们可以假设  $\Omega$  是  $\gamma$  -开闭凸真锥  $\gamma$  ，使得  $N_x^*(\Omega) \subset \mathrm{Int} \gamma^\circ \cup \{0\}$  、  $\mathrm{Int} \gamma \neq \emptyset$  和  $\xi_o \notin \gamma^\circ$  （参见命题 5.3.7）。

让我们用  $ v \in \mathrm{Int}\gamma $  和  $ \langle\xi_{\circ}, v\rangle < 0 $  。对于  $ s > 0 $  、  $ t > 0 $  ，设置：

\[
H_{s}=\left\{x;\left\langle x-x_{o},\xi_{o}\right\rangle>-s\right\}
\]

\[
\Omega_{t,s}=\left\{x;x-t(\langle x-x_{\circ},\xi_{\circ}\rangle+s)v\in\Omega\right\}.
\]

然后：

\[
\bar{\Omega}_{t,s}\cap H_{s}\subset\varOmega,
\]

\[
\begin{array}{r}{\Omega_{t,s}\cap H_{s}\subset\Omega_{t^{\prime},s}\qquad0<t^{\prime}\leqslant t,}\end{array}
\]

\[
\left(\bigcup_{t}\Omega_{t,s}\right)\cap H_{s}=\Omega\cap H_{s}.
\]

因为 $ t \langle v, \xi_0 \rangle < 1 $ ，我们通过以下方式定义转换 $ (x; \xi) \leftrightarrow (y; \eta) $ ：

\[
\left\{\begin{aligned}{y}&{{}=x-t(\langle x-x_{\circ},\xi_{\circ}\rangle+s)v}\\ {\xi}&{{}=\eta-t\langle\eta,v\rangle\xi_{\circ}~.}\\ \end{aligned}\right.
\]

然后我们有：

\[
(x;\xi)\in N^{*}(\Omega_{t,s})\Leftrightarrow(y;\eta)\in N^{*}(\Omega).
\]

现在我们将证明 $(x_o; \xi_o)$ 和 $\varepsilon > 0$ 存在一个开邻域 $U \times W$ ，这样设置 $U_s = U \cap H_s$ ，我们就有 $0 < t < \varepsilon$ ， $0 < s < \varepsilon$ ：

\[
\left\{\begin{aligned}{}&{{}(a)\quad\mathsf{S S}(F)\cap N^{*}(\varOmega_{t,s})^{a}\cap\pi^{-1}(U_{s})\subset T_{X}^{*}X}\\ {}&{{}(b)\quad(\mathsf{S S}(F)+N^{*}(\varOmega_{t,s}))\cap(U_{s}\times W)=\varnothing}\\ \end{aligned}\right.,
\]

事实上假设(a)或(b)是错误的。我们发现序列 $\{t_n\}, \{s_n\}, \{x_n\}, \{\zeta_n\}, \{\zeta_n\}$ 使得：

\[
\left\{\begin{aligned}{}&{{}t_{n}\xrightarrow[n]{}0~,\quad s_{n}\xrightarrow[n]{}0~,\quad~t_{n}>0~,\quad s_{n}>0~,\quad}\\ {}&{{}x_{n}\xrightarrow[n]{}x_{\circ}~,\quad\qquad\xi_{n}\in N_{x_{n}}^{*}(\varOmega_{t_{n},s_{n}})\backslash\{0\}~,\quad}\\ {}&{{}(x_{n};\xi_{n})\in\operatorname{S S}(F)~,\quad}\\ {}&{{}\xi_{n}+\zeta_{n}=c\tilde{\xi}_{n}~,\quad\tilde{\xi}_{n}\to\xi_{\circ}~,\quad\qquad\mathrm{w h e r e}}\\ {}&{{}c=0\quad\mathrm{o r}\quad c=1\qquad(c=0\mathrm{f o r}(a),c=1\mathrm{f o r}(b))~.}\\ \end{aligned}\right.
\]

我们通过以下方式定义 $ (y_n; \eta_n) \in N^*(\Omega) $ ：

\[
\left\{\begin{aligned}{y_{n}}&{{}=x_{n}-t_{n}(\langle x_{n}-x_{\circ},\xi_{\circ}\rangle+s_{n})v}&{,}\\ {\xi_{n}}&{{}=\eta_{n}-t_{n}\langle\eta_{n},v\rangle\xi_{\circ}}&{,}\\ \end{aligned}\right.
\]

我们设置：

\[
\rho_{n}=\zeta_{n}+\eta_{n}=c\tilde{\xi}_{n}+t_{n}\langle\eta_{n},v\rangle\xi_{0}~.
\]

然后 $ \eta_n \in \gamma^\circ $ ，我们有：

\[
\left\{\begin{aligned}{}&{{}\langle\xi_{o},v\rangle<0~,\quad\langle\tilde{\xi}_{n},v\rangle<0~,\quad\langle\eta_{n},v\rangle\geqslant0~,\quad\langle\rho_{n},v\rangle\leqslant0~,\quad}\\ {}&{{}|\eta_{n}|\leqslant c^{\prime}\langle\eta_{n},v\rangle\qquad\quad\text{for some}c^{\prime}>0~,\quad}\\ {}&{{}c^{\prime \prime}|\rho_{n}|\geqslant-\langle\rho_{n},v\rangle\qquad\quad\text{for some}c^{\prime \prime}>0~.\quad}\\ \end{aligned}\right.
\]

最后一个断言是根据  $ \rho_{n} $  的方向收敛于  $ \xi_{0} $  的方向这一事实得出的。因此我们有：

\[
\begin{aligned}{c^{\prime\prime}|\rho_{n}|}&{{}\geq-c\langle\tilde{\xi}_{n},v\rangle-t_{n}\langle\eta_{n},v\rangle\cdot\langle\xi_{\circ},v\rangle}\\ {}&{{}\geq t_{n}\langle\eta_{n},v\rangle|\langle\xi_{\circ},v\rangle|}\\ {}&{{}\geq c^{\prime\prime}t_{n}|\eta_{n}|\qquad\mathrm{f o r~s o m e~}c^{\prime\prime\prime}>0.}\\ \end{aligned}
\]

请注意  $ \rho_{n} \neq 0 $  因为  $ \eta_{n} \neq 0 $  。

因此  $|\eta_n|(|\rho_n|^{-1})t_n$  是有界的并且  $|\eta_n|(|\rho_n|^{-1})|x_n-y_n|\to0$  。由于  $\rho_n/|\rho_n|$  收敛于  $\xi_\circ/|\xi_\circ|$  和  $(x_n;\zeta_n/|\rho_n|)\in\mathrm{SS}(F)$  、  $(y_n;\eta_n/|\rho_n|)\in N^*(\Omega)$  ，所以我们有  $(x_\circ;\xi_\circ/|\xi_\circ|)\in\mathrm{SS}(F)+N^*(\Omega)$  。这与假设相矛盾并证明了（6.3.1）。

现在让  $ j_{t,s} $  表示沉浸  $ \Omega_{t,s} \hookrightarrow X $  ，并设置 s 为固定值，  $ 0 < s < \varepsilon $  ：

\[
\begin{array}{r}{F_{t}=R\ensuremath{j_{t,s*}}(F|_{\Omega_{t,s}})~.}\end{array}
\]

根据命题 5.4.8 和非特征条件 (6.3.1)，我们有：

\[
\mathbb{S S}(F_{t})\cap(U_{s}\times{\cal W})=\emptyset\enspace.
\]

因此，根据命题 5.2.1，对于所有  $ t \in]0, \varepsilon[ $  ，存在一个闭真锥  $ y' $  和  $ y' $  -开子集  $ \Omega_0 \subset \Omega_1 $  ，使得  $ y' \subset \{y; \langle y; \xi_0 \rangle < 0\} \cup \{0\} $  和  $ x_0 \in \mathrm{Int}(\Omega_1 \setminus \Omega_0) \subset H_s $  ，以及  $ R\phi_{\gamma^*} R\Gamma_{(\Omega_1 \setminus \Omega_0)}(F_t) = 0 $  。

应用命题 2.7.1 我们得到  $ R\phi_{\gamma^*}R\Gamma_{(\Omega_1\setminus\Omega_0)}(F) = 0 $  ，这意味着  $ (x_o; \xi_o) \notin \mathrm{SS}(F) $  。

(ii) 证明类似。设置

\[
G_{t}=R j_{t,s!}(F|_{\Omega_{t,s}})
\]

我们发现：

\[
\mathrm{SS}(G_{t})\cap(U_{s}\times W)=\varnothing.
\]

因此，存在命题 5.1.1 中的 $H, L, \gamma$ ，使得

\[
R\varGamma(H\cap(x+\gamma);G_{t})\xrightarrow{\sim}R\varGamma(L\cap(x+\gamma);G_{t})
\]

 $x$  在  $x_{\circ}$  和  $0 < t \ll 1$  附近成立。从 $H^k(H \cap (x + \gamma); R j_! F) = \varinjlim_t H^k(H \cap (x + \gamma); G_t)$ 和 $H^k(L \cap (x + \gamma); R j_! F) = \varinjlim_t H^k(L \cap (x + \gamma); G_t)$ 开始，我们有了 $H^k(H \cap (x + \gamma); R j_! F) \simeq H^k(L \cap (x + \gamma); R j_! F)$ 。这意味着  $(x_{\circ}; \xi_{\circ}) \notin \mathrm{SS}(R j_! F)$  。

\statementlabel{命题 6.3.2}令  $M$  为  $X$  、  $U = X \setminus M$  、  $j$  嵌入  $U \subsetneq X$  的闭子流形，并令  $F \in \mathrm{Ob}(\mathbb{D}^b(U))$  。然后：

\[
\mathbb{S S}(R j_{*}F)\cap\pi^{-1}(M)\subset\mathbb{S S}(F)+\widehat{T_{M}^{*}X}\enspace,
\]

\[
\mathrm{SS}(Rj,F)\cap\pi^{-1}(M)\subset\mathrm{SS}(F)+\mathrm{T}_{M}^{*}X,
\]

\[
\mathsf{S S}((R j_{*}F)|_{M})\subset T^{*}M\cap C_{T_{M}^{*}X}(\mathsf{S S}(F))\enspace.
\]

\statementlabel{证明}令 $p: \widetilde{X}_M \to X$  为 $X$  沿 $M$  的法向变形（参见IV §1）。回想一下，在  $X$  上存在局部坐标系  $(x', x'')$  ，在  $\widetilde{X}_M$  上存在  $(x', x''', t)$  ，这样  $M = \{x' = 0\}$  ，  $p(x', x'', t) = (tx', x'')$  和  $\mathbb{R}^+$  通过  $\lambda(x', x'', t) = (\lambda x', x'', \lambda^{-1}t)$  ，  $(\lambda \in \mathbb{R}^+)$  作用于  $\widetilde{X}_M$  。定义：

\[
\tilde{U}=\left\{(x^{\prime},x^{\prime \prime},t);x^{\prime}\neq0\right\}=\tilde{X}_{M}\backslash(M\times\mathbb{R}),
\]

\[
\widetilde{U}_{\pm}=\widetilde{U}\cap\{\pm t>0\},\qquad\widetilde{M}=\widetilde{U}\cap\{t=0\}=\dot{T}_{M}X,
\]

\[
X^{\prime}=\tilde{U}/\mathbb{R}^{+},\qquad X_{\pm}^{\prime}=\tilde{U}_{\pm}/\mathbb{R}^{+},\qquad M^{\prime}=\tilde{M}/\mathbb{R}^{+}.
\]

用  $\gamma$  表示投影  $\tilde{U} \to X'$  并用  $\rho$  表示映射  $X' \to X$  ，以便  $p = \rho \circ \gamma$  。那么 $X'$  是流形， $\rho$  是真值， $\gamma$  是光滑的。此外 $M' = \rho^{-1}(M)$  是 $X'$  的超曲面，与法向球丛 $S_M X = \dot{T}_M X / \mathbb{R}^+$  同构， $X' = X'_L \sqcup M' \sqcup X_+'$  和 $\rho$  诱导同构 $X_+' \simeq U$  。让我们用  $i$  表示这个同构，并用  $k$  表示嵌入  $X_+' \hookrightarrow X':$ 

\[
\begin{array}{c c c c c c}{\widetilde{M}}&{\longrightarrow}&{\widetilde{M}/\mathbb{R}^{+}=M^{\prime}}&{\longrightarrow}&{M}\\ {}&{}&{}\\ {\cap}&{}&{\cap}&{}&{\cap}\\ {}&{}&{\widetilde{U}\longrightarrow_{\gamma}}&{\widetilde{U}/\mathbb{R}^{+}=X^{\prime}}&{\longrightarrow_{\rho}}&{X}\\ {\cup}&{}&{\displaystyle\mathop{\times}_{k}}&{}&{\displaystyle\mathop{\mathop{\mathop{\operatorname{div}}\strut}a}_{j}}\\ {\widetilde{U}_{+}}&{\longrightarrow}&{\widetilde{U}_{+}/\mathbb{R}^{+}=X_{+}^{\prime}}&{\longrightarrow_{\widetilde{i}}}&{U~.}\\ \end{array}
\]

设置  $ F' = i^{-1}F $  ，我们得到：

\[
R\rho_{*}R k_{*}F^{\prime}\simeq R j_{*}F\quad,
\]

\[
R\rho_{*}R k_{!}F^{\prime}\simeq R j_{!}F\quad.
\]

设定：

\[
S=\rho_{\pi}^{t}\rho^{\prime-1}(\mathrm{S S}(F^{\prime})+{\cal T}_{M^{\prime}}^{*}X^{\prime})~.
\]

应用命题 5.4.4 和定理 6.3.1 我们得到：

\[
\pi^{-1}(M)\cap\mathrm{S S}(R j_{*}F)\subset S\enspace,
\]

以及类似的公式，其中  $j_{*}$  替换为  $j_{!}$  。由于 $\gamma$  是光滑且满射的，我们发现：

\[
\mathrm{S S}(F^{\prime})+{\cal T}_{M^{\prime}}^{*}X^{\prime}=\gamma_{\pi}{}^{t}\gamma^{\prime-1}(\mathrm{S S}(\gamma^{-1}F^{\prime}|_{\tilde{U}_{+}}){\hat{\vphantom{\widetilde{U}}}}+T_{\widetilde{M}}^{*}\widetilde{U})\;.
\]

因此：

\[
S=p_{\pi}{}^{t}p^{{}^{\prime}-1}\left((\mathrm{S S}(p^{-1}F|_{\tilde{U}_{*}}){\widehat{\mathrm{~\tiny~+~}}}T_{\tilde{M}}^{*}\tilde{U})_{{\tilde{X}}_{M}}\times\tilde{U}\right).
\]

让  $(x', x''$  、  $t;\xi', \xi''$  、  $\tau)$  表示  $T^*\tilde{X}_M$  上的坐标。然后  $p_\pi$  和  $t'p'$  分别由  $(x', x'', t;\xi', \xi'')\mapsto(tx', x'';\xi', \xi'')$  和  $(x', x'', t;\xi', \xi'')\mapsto(x', x'', t; t\xi', \xi'', \langle x', \xi' \rangle)$  定义。因此 $(0, x_o^*, \xi_o^*, \xi_o'') \in S$  意味着 $x_o' \neq 0$  的存在：

\[
(x_{\circ}^{\prime},x_{\circ}^{\prime\prime},0;0,\xi_{\circ}^{\prime\prime},\langle x_{\circ}^{\prime},\xi_{\circ}^{\prime}\rangle)\in\mathrm{S S}(p^{-1}F|_{\tilde{U}_{+}})\widehat{+}\;T_{\tilde{M}}^{*}\tilde{U}\enspace.
\]

我们在 $ \mathrm{SS}(p^{-1}F|_{\tilde{U}_+}) $  中得到一个序列 $ \{(x_n', x_n", t_n; \xi_n', \xi_n'', \tau_n)\} $ ，这样：

\[
\left\{\begin{aligned}{x_{n}^{\prime}}&{{}\xrightarrow[n]{}x_{\circ}^{\prime},\quad x_{n}^{\prime\prime}\xrightarrow[n]{}x_{\circ}^{\prime\prime},\quad t_{n}\xrightarrow[n]{}0,\quad\xi_{n}^{\prime}\xrightarrow[n]{}0}\\ {\xi_{n}^{\prime\prime}}&{{}\xrightarrow[n]{}{\xi}_{\circ}^{\prime\prime},\quad t_{n}\tau_{n}\xrightarrow[n]{}0}\\ \end{aligned}\right..
\]

自：

\[
(x^{\prime},x^{\prime\prime},t;\xi^{\prime},\xi^{\prime\prime},\tau)\in\mathrm{S S}(p^{-1}F|_{\tilde{U}_{+}})\Leftrightarrow(t x^{\prime},x^{\prime\prime};t^{-1}\xi^{\prime},\xi^{\prime\prime})\in\mathrm{S S}(F),
\]

\[
t>0~,\quad x^{\prime}\neq0~,\quad t\tau=\langle x^{\prime},\xi^{\prime}\rangle~,
\]

我们发现序列 $ \{(t_nx'_n, x''_n; t_n^{-1}\xi'_n, \xi''_n)\} $ 包含在SS(F)中并且满足：

\[
t_{n}x_{n}^{\prime}\mathop{\longrightarrow}_{n}0,\qquad x_{n}^{\prime\prime}\mathop{\longrightarrow}_{n}x_{\circ}^{\prime\prime},\qquad\xi_{n}^{\prime\prime}\mathop{\longrightarrow}_{n}\xi_{\circ}^{\prime\prime},\qquad|t_{n}x_{n}^{\prime}|\cdot|t_{n}^{-1}\xi_{n}^{\prime}|\mathop{\longrightarrow}_{n}0.
\]

这显示了备注 6.2.8 中的前两个包含内容。为了证明最后的包含，让我们考虑一下区分的三角形：

\[
R\mathbf{j}_{!}F\to R\mathbf{j}_{*}F\to(R\mathbf{j}_{*}F)_{\cal M}\xrightarrow[+1]{}.
\]

然后通过三角不等式 $ \mathrm{SS}((Rj_*F)_M)\subset\mathrm{SS}(F)+\mathrm{T}_M^*X $  。然后第三个包含项来自命题 5.4.4 和引理 6.2.1 (i)。   $ \square $ 

在第七章中，我们将不得不考虑映射的直像，这些图像仅在“微局域”意义上才是正确的。更准确地说，请考虑以下情况。

让 $ q_1: X \times Y \to X $  和 $ p: T^*(X \times Y) \to (T^*X) \times Y $  表示自然投影。

\statementlabel{命题 6.3.3}令 $\Omega$  为 $T^*X$  的开子集，并让 $F \in \mathrm{Ob}(\mathbb{D}^b(X \times Y))$  假设：

投影 $ \overline{p(\mathrm{SS}(F))} \cap (\Omega \times Y) \to \Omega $  是正确的。

然后：

\[
\mathrm{S S}(R q_{1*}{F})\cap\Omega\subset q_{1\pi}{}^{t}q_{1}^{\prime^{-1}}(\mathrm{S S}({F})),
\]

\[
\begin{array}{r}{\mathrm{S S}(R q_{1};F)\cap\Omega\subset q_{1\pi}{}^{t}q_{1}^{\prime^{-1}}(\mathrm{S S}(F))\enspace,}\end{array}
\]

 $ Rq_{1!}F \to Rq_{1*}F $  是  $ \mathbf{D}^{b}(X;\Omega) $  的同构

请注意，（6.3.9）意味着对于任何紧凑子集  $ K \subset \Omega $  都存在一个紧凑子集  $ L \subset Y $  ，这样：

\[
\mathsf{S S}(F)\cap(K\times T^{*}Y)\subset K\times\left(L\mathop{\times}_{\gamma}T^{*}Y\right).
\]

\statementlabel{证明}通过选择将  $Y$  封闭嵌入到向量空间中，我们可以假设  $Y = \mathbb{R}^n$  对于某些  $n$  。由于  $\mathbb{R}^n$  与  $\mathbb{R}^n$  的开球同构，我们现在可以假设  $Y$  是  $\mathbb{R}^n$  的开球。设置  $\tilde{Y} = \mathbb{R}^n$  ，并用  $j$  表示开放嵌入  $X \times Y \subsetneq X \times \tilde{Y}$  ，用  $\tilde{q}_1$  表示投影  $X \times \tilde{Y} \to X$  。然后 $R q_1 * F \simeq R \tilde{q}_1 * R j_* F$  和 $R q_1 * F \simeq R \tilde{q}_1 * R j_* F$  。设置 $Z = X \times \partial Y$ 。根据定理 6.3.1，我们有：

\[
\mathsf{S S}(R j_{*}F)\subset\mathsf{S S}(F)\cup(T_{Z}^{*}(X\times\widetilde{Y})+\mathsf{S S}(F))~,
\]

\[
\mathrm{S S}(R j_{!}F)\subset\mathrm{S S}(F)\cup(T_{Z}^{*}(X\times\tilde{Y})\widehat{+}\mathrm{S S}(F)),
\]

\[
\mathsf{S S}((R j_{*}F)_{Z})\subset T_{Z}^{*}(X\times\widetilde{Y})+\mathsf{S S}(F).
\]

通过考虑区分三角形 \textbackslash{}(Rj\_1F \textbackslash{}to Rj\_*F \textbackslash{}to (Rj\_*F)\_Z \textbackslash{}xrightarrow\{+1\} 我们看到命题的所有结论将来自：

\[
\tilde{q}_{1\pi}{}^{t}\tilde{q}_{1}^{\prime-1}(T_{Z}^{*}(X\times\tilde{Y})\widehat{+}\operatorname{S S}(F))\cap\varOmega=\emptyset\enspace.
\]

让我们证明（6.3.11）。令 K 为  $ \Omega $  的紧凑子集。由(6.3.10)我们得到：

\[
\mathsf{S S}(F)\cap(K\times T^{*}Y)\cap\pi^{-1}(Z)=\varnothing,
\]

因此

\[
K\cap\tilde{q}_{1\pi}{^{\sf t}}\tilde{q}_{1}^{\prime}{^{-1}}(T_{Z}^{*}(X\times\tilde{Y})+{\tt S S}(F))=\varnothing\enspace.\enspace\Box
\]

\subsection{微局部化}
我们要对微观定位的微观支撑给予一定的约束。

\statementlabel{定理 6.4.1}令  $ M $  为  $ X $  的闭子流形，并令  $ F \in \mathrm{Ob}(\mathbb{D}^b(X)) $  。然后：

\[
\mathbb{S S}(\nu_{M}(F))\subset C_{T_{M}^{*}X}(\mathbb{S S}(F))~,
\]

\[
\mathrm{S S}(\nu_{M}(F))=\mathrm{S S}(\mu_{M}(F))~.
\]

在这里，我们通过使用命题 5.5.1（参见（6.2.2）和（6.2.3））来识别 $ T_{T_M^*X}T^*X $ ， $ T^*T_M^*X $ 和 $ T^*T_M X $ 。

\statementlabel{证明}(ii) 由定理 5.5.5 得出。

(i) 我们选择 X 上的局部坐标系  $ (x', x'') $  使得  $ M = \{(x', x''); x' = 0\} $  。如 IV §1 中所示，  $ \widetilde{X}_M $  具有坐标  $ (x', x'', t) $  ，我们表示为

 $p: \widetilde{X}_M \to X$  的映射 $(x', x''$  、 $t) \mapsto (tx', x'')$  ，以及 $\Omega\widetilde{X}_M$  的开集 $\{t > 0\}$  。因此，我们可以在  $\widetilde{X}_M$  中识别  $T_M X$  到  $\{t = 0\}$  。我们用 $j$ 表示嵌入 $\Omega \hookrightarrow \widetilde{X}_M$ ，用 $s$ 表示沉浸 $T_M X \hookrightarrow \widetilde{X}$ 并设置 $\widetilde{p} = p \circ j$ （参见图（4.1.5））。我们用  $(x', x'', t; \xi', \xi'', \tau)$  表示  $T^* \widetilde{X}_M$  的坐标。根据定义， $\nu_M(F) = s^{-1} R j_* \widetilde{p}^{-1} F$  。应用命题 5.4.5 和命题 6.3.2，我们得到：

\[
\begin{aligned}{\mathbb{S S}(\tilde{p}^{-1}F)=\left\{(x^{\prime},x^{\prime \prime},t;\xi^{\prime},\xi^{\prime \prime},\tau);t>0,(t x^{\prime},x^{\prime \prime};t^{-1}\xi^{\prime},\xi^{\prime \prime})\in\mathbb{S S}(F),t\tau-\langle x^{\prime},\xi^{\prime}\rangle=0\right\}}\\ \end{aligned}
\]

因此  $11\left(x_{o}^{\prime},x_{o}^{\prime\prime};\xi_{o}^{\prime},\xi_{o}^{\prime\prime}\right)\in SS(\nu_{M}(F))$  那么 SS  $(\tilde{p}^{-1}F)$  和  $t_{n}\to0,\left|t_{n}\right|\left|\tau_{n}\right|\to0,\left(x_{n}^{\prime},x_{n}^{\prime\prime};\xi_{n}^{\prime},\xi_{n}^{\prime\prime}\right)\to\left(x_{o}^{\prime},x_{o}^{\prime\prime};\xi_{o}^{\prime},\xi_{o}^{\prime\prime}\right)$  中存在序列  $\left\{(x_{n}^{\prime},x_{n}^{\prime\prime},t_{n};\xi_{n}^{\prime},\xi_{n}^{\prime\prime},\tau_{n})\right\}$  。因此 $\left(t_{n}x_{n}^{\prime},x_{n}^{\prime\prime};t_{n}^{-1}\xi_{n}^{\prime},\xi_{n}^{\prime\prime}\right)$  属于SS  $(F)$  。由于 SS  $(F)$  是圆锥曲线，因此  $(t_{n}x_{n}^{\prime},x_{n}^{\prime\prime};\xi_{n}^{\prime},t_{n}\xi_{n}^{\prime\prime})$  属于 SS  $(F)$  。那么  $C_{T_{M}^{*}X}(\mathrm{SS}(F))$  的定义立即意味着  $(x_{o}^{\prime},x_{o}^{\prime\prime};\xi_{o}^{\prime},\xi_{o}^{\prime\prime})\in C_{T_{M}^{*}X}(\mathrm{SS}(F))$  。

请注意，使用与定理 6.4.1 证明中相同的符号，我们发现如果  $(x_{\circ}^{\prime}, x_{\circ}^{\prime\prime}; \xi_{\circ}^{\prime}, \xi_{\circ}^{\prime\prime})$  属于  $\mathrm{SS}(\nu_{M}(F))$  则  $\langle x_{\circ}^{\prime}, \xi_{\circ}^{\prime} \rangle = 0$  。 （在证明过程中，这是从  $|t_{n}| | \tau_{n}| \to 0$  和  $t_{n} \tau_{n} - \langle x_{n}^{\prime}, \xi_{n}^{\prime} \rangle = 0$  得出的。）使用符号 (5.5.7)，这意味着  $\mathrm{SS}(\nu_{M}(F)) \subset S_{T_{M}^{*}X}$  ，这相当于说  $\nu_{M}(F)$  是圆锥曲线。

\statementlabel{推论 6.4.2}令 $f: Y \to X$  为流形的态射。让  $F \in \mathrm{Ob}(\mathbb{D}^b(X))$  和  $G \in \mathrm{Ob}(\mathbb{D}^b(Y))$  。然后：

\[
\begin{array}{r}{\mathrm{S S}(\mu\mathrm{h o m}(G\to F))\subset C_{\mu}(\mathrm{S S}(F),\mathrm{S S}(G)).}\end{array}
\]

\[
\begin{array}{r}{\mathrm{S S}(\mu\mathrm{h o m}(F\gets G))\subset C_{\mu}(\mathrm{S S}(F),\mathrm{S S}(G))^{a}~,}\end{array}
\]

其中  $a$  是相对于  $T^{*}(Y \times_{X} T^{*}X)$  在  $Y \times_{X} T^{*}X$  上的向量丛结构的对映图。

(iii) 假设 G 是上同调可构造的。然后是自然态射：

\[
R\mathcal{H}o m(G,A_{Y})\overset{L}{\otimes}f^{-1}F\otimes\omega_{Y/X}\to R\mathcal{H}o m(G,f^{!}F)
\]

是  $ T^*Y\setminus f_\infty^\#(\mathrm{SS}(F),\mathrm{SS}(G)) $  的同构。

\statementlabel{证明}(i) 和 (ii) 紧随定理 6.4.1 得出。

(iii) 考虑区分三角形

\[
R\pi_{!}H\longrightarrow R\pi_{*}H\longrightarrow R\dot{\pi}_{*}H\xrightarrow[+1]{}
\]

其中 $H = \mu h o m(G \to F)$  。然后根据命题5.5.4和定义6.2.3(iii) $\mathrm{SS}(R\dot{\pi}_*H) \cap V = \emptyset$ ，结果由命题4.4.2得出。 □

\statementlabel{推论 6.4.3}让  $F$  和  $G$  属于  $\mathbf{D}^{b}(X)$  。然后：

\[
\mathrm{S S}(\mu\mathrm{h o m}(G,F))\subset C(\mathrm{S S}(F),\mathrm{S S}(G))~.
\]

(ii) 假设 G 是上同调可构造的。

然后是自然态射：

\[
R\mathcal{H}o m(G,A_{X})\overset{L}{\otimes}F\to R\mathcal{H}o m(G,F)
\]

是  $ T^*X\setminus(\mathrm{SS}(F)+\infty\mathrm{SS}(G)^a) $  的同构。

\statementlabel{证明}这是推论 6.4.2 的一个特例。 ⑨

\statementlabel{推论 6.4.4}令  $f: Y \to X$  为流形的态射，并令  $F \in \mathrm{Ob}(\mathbb{D}^b(X))$  。然后：

\[
\mathrm{SS}(f^{-1}F)\subset f^{\#}(\mathrm{SS}(F))~,
\]

\[
\mathbb{S S}(f^{!}F)\subset f^{\#}(\mathbb{S S}(F))~.
\]

(iii) 令 V 为  $ T^{*}Y $  的子集，并假设 f 对于 V 上的 F 来说是非特征的。那么自然态射：

\[
f^{-1}F\otimes\omega_{Y/X}\to f^!F
\]

是  $V$  的同构，而且：

\[
\mathrm{S S}(f^{-1}F)\cap V\subset{}^{t}f^{\prime}(f_{\pi}^{-1}(\mathrm{S S}(F))).
\]

\statementlabel{证明}将  $Y$  与  $X \times Y$  中的  $f$  的图进行识别，并用  $q$  表示投影  $T_{Y}^{*}(X \times Y) \to Y$  。然后：

\[
\left\{\begin{aligned}{f^{-1}F\otimes\omega_{Y/X}}&{{}\simeq R q_{!}\mu_{Y}(F\boxtimes\omega_{Y})~,}\\ {f^{!}F\simeq R q_{*}\mu_{Y}(F\boxtimes\omega_{Y})~,}\\ \end{aligned}\right.
\]

并且有一个带有 $ H = \mu_Y (F \boxtimes \omega_Y) $ 的显着三角形（6.4.1）。然后结果由命题 5.5.4 得出。   $ \Box $ 

\statementlabel{推论 6.4.5}让  $F$  和  $G$  属于  $\mathbf{D}^{+}(X)$  。然后：

\[
\mathsf{S S}\bigg(F\mathop{\otimes}^{L}G\bigg)\subset\mathsf{S S}(F)\mathrel{\widehat{+}}\mathsf{S S}(G),
\]

\[
\mathrm{SS}(R\mathcal{H o m}(G,F))\subset\mathrm{SS}(F)+\mathrm{SS}(G)^{a}.
\]

\statementlabel{证明}应用命题 5.4.1、5.4.2 和推论 6.4.4，如命题 5.4.14 的证明。 ⑨

\statementlabel{备注 6.4.6}在推论 6.4.4 (iii) 的情况下，假设  $V$  在  $T^*Y$  中开放。那么映射 $t f': t f'^{-1}(V) \cap f^{-1}(\mathrm{SS}(F)) \to V$  是正确的。然而这最后一个反

dition 严格弱于条件“f 对于 V 上的 F 来说是非特征的”，如以下示例所示。

将  $ X = \mathbb{R}^2 $  与坐标  $ (y, t) $  、  $ Y = \{(y, t); t = 0\} $  并设  $ F = A_Z $  其中  $ Z = \{(y, t); t = y^2\} $  。让 $ V = \{(y; \eta) \in T^* Y; \eta > 0\} $  。然后 $ F|_{Y} \simeq A_{\{0\}} $  和 $ f_{\pi}^{-1}(\mathrm{SS}(F)) \cap t'f''^{-1}(V) = \emptyset $  。

\statementlabel{备注 6.4.7}微支撑对于 X 上的  $ C^{1} $  变换是不变的，但  $ \hat{+} $  和  $ f^{*}(\cdot) $  不是不变的（参见练习 VI 7）。这意味着定理 6.3.1 及其推论并不是最好的结果。

\subsection{对合性与传播}
在偏微分方程理论中，奇点沿双特征曲线的传播和特征变化的对合性的结果至关重要。在本节中，我们将在层框架中证明此类结果，作为推论 6.4.3 的应用。

令 $U$  为 $T^*X$  的开子集， $\phi$  为 $U$  上的实 $C^2$  函数。我们设定：

\[
V_{0}=\left\{p\in U;\phi(p)=0\right\}\;,\qquad V_{\pm}=\left\{p\in U;\pm\phi(p)\geqslant0\right\}\;.
\]

让 $p \in V_0$  。哈密​​顿向量场 $H_\phi$ 通过 $p$ 的积分曲线称为 $p$ 发出的 $V_0$ （或 $\phi$ ）的双特征曲线，并用 $b_p$ 表示。还将 $p$  发布的正半双特征曲线 $b_p^+$  定义为 $p$  发布的 $H_\phi$  的正半积分曲线。如果 $(x; \xi)$  是齐次辛坐标系，则 $b_p^+$  是曲线 $(x(t), \xi(t))_{t \geq 0}$ ，使得：

\[
\frac{\partial x_{j}}{\partial t}=\frac{\partial\phi}{\partial\xi_{j}}~,\qquad\frac{\partial\xi_{j}}{\partial t}=-\frac{\partial\phi}{\partial x_{j}}~,\qquad(x(0);\xi(0))=p~.
\]

请注意，  $ b_{p}^{+} $  取决于  $ V_{+} $  ，而不是  $ \phi $  。更准确地说，用 $ h\phi $  替换 $ \phi $  和 $ h(p) > 0 $  不会影响p 邻域中的 $ b_{p}^{+} $ 。

类似地定义了  $ \dot{b}_{p}^{-} $  ，负半双特征曲线。如果 $ \mathrm{d}\phi(p) = 0 $ ，我们将其理解为 $ b_{p}^{\pm} = \{p\} $ 。

令 S 为  $ T^{*}X $  的局部闭子集。

\statementlabel{定义 6.5.1}让 $p \in S$  。我们可以说 $S$  在 $p$  上是对合的，如果对于任何 $\theta \in T_p^*T^*X$  使得法锥 $C_p(S, S)$  包含在超平面 $\{v \in T_p T^*X; \langle v, \theta \rangle = 0\}$  中，并且有：  $-H(\theta) \in C_p(S)$  。

如果 $S$  对每个 $p \in S$  都是对合的，我们就说 $S$  是对合的。

当然，如果  $S$  是光滑子流形，则当且仅当  $T_pS$  对于所有  $p\in S$  而言  $T_pT^*X$  对合，所以  $S$  在定义 6.5.1 的意义上是对合的，因为  $C_p(S,S)=C_p(S)=T_pS$  。

\statementlabel{命题 6.5.2}假设  $S$  是对合的并且封闭在  $T^*X$  的开子集  $U$  中。令  $\phi: U \to \mathbb{R}$  为类  $C^2$  的函数，使得  $\phi|_{S} = 0$  。那么  $S$  是  $\phi$  的双特征曲线的并集。

\statementlabel{证明}让 $p \in S$  。我们可以假设 $\mathrm{d}\phi(p) \neq 0$ ，否则 $p$ 发布的 $H_\phi$ 的积分曲线仅限于 $\{p\}$ 。因此  $S$  包含在光滑超曲面  $\{\phi = 0\}$  中，这意味着：

\[
{\cal C}_{p}(S,S)\subset\left\{v\in T_{p}T^{*}X;\langle v,\mathrm{d}\phi(p)\rangle=0\right\}.
\]

因此  $ H_\phi \in C_p(S) $  ，结果将来自下一个引理。 □

\statementlabel{引理 6.5.3}令 $U$ 为流形，$S$ 为 $U$ 的闭子集，$v$ 为 $U$ 上的 $C^1$ 向量场，$p_0\in S$，并令 $b^+$ 为从 $p_0$ 出发的 $v$ 的正半积分曲线（即 $b^+:[0,+\infty[\to U$ 满足 $b^+(0)=p_0$ 且 $\frac{\partial}{\partial t}b^+(t)=v(b^+(t))$）。假设对所有 $p\in S$ 有 $v(p)\in C_p(S)$。则 $b^+$ 包含于 $S$。

\statementlabel{证明}我们可以假设  $U$  是  $\mathbb{R}^N$  、  $p_\circ = 0$  和  $v = \frac{\partial}{\partial t}$  相对于坐标  $(t, x) \in \mathbb{R} \times \mathbb{R}^{N-1}$  的开子集。然后 $b^+(t) = (t, 0)$ 。

假设  $(a,0)\notin S$  对于某些  $a>0$  。存在 $\varepsilon>0$ ，设置 $B=\{x\in\mathbb{R}^{N-1};|x|\leqslant\varepsilon\}$ ，我们就有 $S\cap{a}\times B=\varnothing$ 。令  $\gamma_{t}$  表示  $\{(t,0)\}\cup(\{a\}\times B)$  的凸包，并令  $t_{\circ}=\inf\{t;\gamma_{t}\cap S=\varnothing\}$  。然后 $\gamma_{t_{\circ}}\cap S\neq\varnothing$  、 $0\leqslant t_{\circ}<a$  和 $\gamma_{t}\cap S=\varnothing$  为 $t_{\circ}<t\leqslant a$  。选择 $p\in\gamma_{t_{\circ}}\cap S$ 。然后 $\frac{\partial}{\partial t}\notin C_{p}(S)$ 。这是一个矛盾。   $\square$ 

\statementlabel{定理 6.5.4（对合性定理）}让 $ F \in \mathrm{Ob}(\mathbf{D}^b(X)) $  。那么 $ \mathrm{SS}(F) $  是对合的。

\statementlabel{证明}让 $ S = \mathrm{SS}(F) $ ， $ p \in S $ ， $ \theta \in T_p^* T^* X $ 。假设：

\[
C_{p}(S,S)\subset\left\{\theta=0\right\}\;,
\]

\[
H(\theta)\notin C_{p}(S)~.
\]

我们将得出一个矛盾。到（6.5.4），我们有 $\theta \neq 0$ ，并且存在 $T^*X$ 的闭子集 $Z$ ，使得 $p \in Z$ 、 $S \subset Z$ 和 $\langle H(\theta), \lambda \rangle < 0$ 对于所有 $\lambda \in N_p^*(Z) \setminus \{0\}$ 。事实上，如果选择  $p$  处的局部坐标系，我们会找到一个凸开锥体  $\gamma$  ，其顶点位于  $p$  、  $\gamma$  包含  $H(\theta)$  ，并且  $\gamma \cap S = \varnothing$  位于  $p$  的邻域；然后输入  $Z = T^*X \setminus \gamma$  就足够了。

另一方面，(6.5.3) 意味着（鉴于推论 6.4.3） $ \mathrm{SS}(\mu hom(F, F)) \cap \pi^{-1}(p) $  包含在集合  $ \{H(\theta) = 0\} $  中（这里  $ \pi $  表示投影  $ T^{*}T^{*}X \to T^{*}X $  ）。因此我们得到：

\[
\mathrm{S S}(\mu\mathcal{h o m}(F,F))\cap N_{p}^{*}(\mathrm{Z})\subset\{0\}~.
\]

由于 $\mu\mathrm{hom}(F,F)_p = (R\Gamma_Z\mu\mathrm{hom}(F,F))_p$ ，我们通过推论5.4.9 获得 $\mu\mathrm{hom}(F,F)_p = 0$ 。因此推论 6.1.3 得出的  $p\notin\mathrm{SS}(F)$  是一个矛盾。 □

\statementlabel{备注 6.5.5}在定义 6.5.1 中，如果我们将条件  $ C_p(V,V) \subset \{\theta = 0\} $  替换为较弱的条件  $ C_p(V) \subset \{\theta = 0\} $  ，则定理 6.5.4 将不再成立。参见练习 VI.2。

现在让  $U$  为  $T^*X$  的开子集， $\phi$  为  $U$  上的实数  $C^2$  函数，并定义  $V_0$  、  $V_+$  、  $V_-$  如 (6.5.1) 中所示。

\statementlabel{命题 6.5.6}让  $F$  和  $G$  属于  $\mathbf{D}^b(X)$  。假设  $\mathrm{SS}(F) \cap U \subset V_+$  和  $\mathrm{SS}(G) \cap U \subset V_-$  。令 $j \in \mathbb{Z}$  和 $u$  成为捆 $H^j(\mu \mathrm{hom}(G, F))$  的一部分。那么 $\mathrm{supp}(u)$ 包含在 $V_0$ 中，并且是正半双特征曲线的并集。

\statementlabel{证明}设置 $K = \mu_{\mathrm{hom}}(G, F)$ 。由推论6.4.3我们知道，在通过同构 $-H$ 识别 $TT^*X$ 和 $T^*T^*X$ 之后，上面的 $U$ ， $\mathrm{supp}(K)$ 包含在 $V_+ \cap V_- = V_0$ 中， $\mathrm{SS}(K)$ 包含在集合 $C(V_+, V_-) = \{v \in TT^*X; \langle v, d\phi \rangle \geqslant 0\}$ 中。因此：

\[
\mathbf{S S}(K)\subset\left\{\theta\in T^{*}T^{*}X;\langle\theta,H_{\phi}\rangle\geqslant0\right\}
\]

（自 $ \langle -H(\theta), \mathrm{d}\phi \rangle = \langle \theta, H_\phi \rangle $ 以来）。

然后结果由下一个引理得出。

\statementlabel{引理 6.5.7}令 $Z$  为流形， $v$  为 $Z$  上的 $C^1$  向量场，并令 $K \in \mathrm{Ob}(\mathbb{D}^b(Z))$  。假设 $\mathrm{SS}(K) \subset \{\theta \in T^*Z; \langle \theta, v \rangle \geq 0\}$ 。令 $j \in \mathbb{Z}$  和 $u$  成为 $Z$  上的层 $H^j(K)$  的一部分。那么  $\mathrm{supp}(u)$  是  $v$  的正半积分曲线的并集。

\statementlabel{证明}令 $p \in Z$  和 $b$  为通过 $p$  的积分曲线。我们可以假设  $b \neq \{p\}$  ，在更改坐标后，我们可以假设  $Z = \mathbb{R} \times Y$  、  $b = \mathbb{R} \times \{y_0\}$  、  $v = \frac{\partial}{\partial t}$  ，其中  $t$  是  $\mathbb{R}$  上的坐标。

令 $\gamma$  为 $\mathbb{R}$  的圆锥体 $\{t \leqslant 0\}$  ，并令 $\tilde{\gamma} = \gamma \times Y$  。应用命题 5.2.3 我们得到  $K \simeq \phi_{\tilde{\gamma}}^{-1} R \phi_{\tilde{\gamma}*} K$  。这意味着  $H^j(K)|_b \simeq \phi_{\gamma}^{-1} L$  对于某些  $L \in \mathrm{Ob}(\mathcal{S}\mathfrak{h}(\mathbb{R}_\gamma))$  ，如果  $u|_b$  是  $H^j(K)|_b$  的一部分，则根据命题 3.5.3，我们发现  $\mathrm{supp}(u|_b)$  是区间  $[a, +\infty[$  的并集。   $\square$ 

示例 6.5.8。令 $t$  为 $\mathbb{R}$  上的坐标， $(t; \tau)$  为 $T^*\mathbb{R}$  上的坐标，并让 $F = A_{\{t > 0\}}$  、 $G = A_{\{t=0\}}$  上的坐标。然后  $\mu\text{hom}(G, F) = A_{\{t=0, \tau \leq 0\}}[-1]$  并且如果  $u$  是  $H^1(\mu\text{hom}(G, F))$  的非零部分，则其支持是区间  $\{t = 0, \tau \in] - \infty, 0]\}$  。这是  $(0, 0)$  发布的  $H_t = -\frac{\partial}{\partial \tau}$  的正半双特征曲线，但它不是负半双特征曲线的并集。

\subsection{对合流形邻域中的层}
令 $f: Y \to X$  为流形的态射。

\statementlabel{命题 6.6.1}假设  $f$  是封闭嵌入，并用  $X$  的子流形识别  $Y$  。让 $p \in T_Y^* X$  并让 $F \in \mathrm{Ob}(\mathbf{D}^b(X))$  。

(i) 假设  $ \mathrm{SS}(F) \subset \pi^{-1}(Y) $  在 p 的邻域内。那么存在 $ G \in \mathrm{Ob}(\mathbb{D}^{b}(Y)) $ 使得 $ F \simeq f_{*}G $ 在 $ \mathbb{D}^{b}(X; p) $ 中。

(ii) 假设  $ \mathrm{SS}(F) \subset T_Y^* X $  在 p 的邻域内。那么存在 $ M \in \mathrm{Ob}(\mathbf{D}^b(\mathfrak{Mod}(A))) $ 使得 $ F \simeq M_Y $ 在 $ \mathbf{D}^b(X; p) $ 中。

\statementlabel{证明}(i) 如果 $p \in T_X^*X$  ，则无需证明。假设 $p \in \dot{T}_Y^*X$ 。通过对 $Y$  的余维进行归纳，我们可以假设 $Y$  是一个超曲面。令  $\{\phi = 0\}$  为  $Y$  和  $p = (x_o; \mathrm{d}\phi(x_o))$  的方程。设置  $\Omega^\pm = \{x \in X; \pm \phi(x) > 0\}$  并用  $j_\pm$  表示开放嵌入  $\Omega^\pm \hookrightarrow X$  。应用定理 6.3.1 我们发现  $p \notin \mathrm{SS}(R j_{-*} j^{-1}(F))$  。因此  $R\Gamma_{\{\phi \geq 0\}}(F) \to F$  是  $\mathbf{D}^b(X; p)$  中的同构，我们可以从一开始就假设  $\mathrm{supp}(F)$  包含在  $\{\phi \geq 0\}$  中。再次根据定理 6.3.1，我们发现  $p \notin \mathrm{SS}(R j_+! j_+^{-1}(F))$  。因此  $F \to F_Y$  是  $\mathbf{D}^b(X; p)$  的同构。自  $F_Y \simeq f_* f^{-1}F$  起，结果如下。

(ii) 我们在  $\mathbf{D}^b(X; p)$  中有  $F = f_* G$  ，对于一些  $G \in \mathrm{Ob}(\mathbf{D}^b(Y))$  。应用命题 5.4.4，我们在  $\pi(p)$  的邻域中得到  $\mathrm{SS}(G) \subset T_Y^* Y$  。让 $g$  成为映射 $Y \to \{pt\}$  。根据命题 5.4.5， $G = g^{-1}M$  ，对于某些  $M \in \mathrm{Ob}(\mathbf{D}^b(\mathfrak{Mod}(A)))$  。因此  $F \simeq M_Y$  中的  $\mathbf{D}^b(X; p)$  。   $\square$ 

\statementlabel{命题 6.6.2}假设  $f$  是光滑的，并用  $T^*Y$  的子流形识别  $Y \times_X T^*X$  。让 $p \in Y \times_X T^*X$  并让 $G \in \mathrm{Ob}(\mathbf{D}^b(Y))$  。假设  $\mathrm{SS}(G) \subset Y \times_X T^*X$  在  $p$  的邻域中。那么存在 $F \in \mathrm{Ob}(\mathbf{D}^b(X))$ 使得 $G \simeq f^{-1}F$ 在 $\mathbf{D}^b(Y; p)$ 中。

\statementlabel{证明}我们可以通过对  $\dim Y - \dim X$  进行归纳来论证，并假设  $Y = \mathbb{R}^n$  、  $X = \mathbb{R}^{n-1}$  、  $f$  是  $(x_1, x') \mapsto x', p = (0; \xi_o)$  与  $\xi_o = (0, \xi_o')$  的投影。

如果 $\xi_\circ = 0$ ，则结果已在命题5.4.5中得到证明。因此我们假设  $\xi_\circ \neq 0$  。令  $H_\varepsilon$  为开半空间  $\{x \in \mathbb{R}^n; \langle x, \xi_\circ \rangle > -\varepsilon\}$  ，令  $\gamma$  为  $\mathbb{R}^n$  中的真闭凸锥，使得  $\gamma^\circ a$  为  $\xi_\circ$  的邻域，最后令  $U$  为  $0$  的邻域， $0$  是  $H_\varepsilon$  与  $\gamma$  的交集 $\mathbb{R}^n$ 。我们可以假设：

\[
\mathrm{S S}(G)\cap(U\times\gamma^{\circ a})\subset Y\times T^{*}X\enspace.
\]

我们推论

\[
((\mathbf{S S}(G)\backslash(U\times\gamma^{\circ a}))\widehat{+}N^{*}(H_{\varepsilon})^{a})\cap(U\times\operatorname{I n t}\gamma^{\circ a})=\varnothing,
\]

\[
((\mathrm{S S}(G)\cap(U\times\gamma^{\circ a}))\mathrel{\hat{+}}N^{*}(H_{\varepsilon})^{a})\cap(U\times\gamma^{\circ a})\subset Y\times T^{*}X.
\]

然后：

\[
\mathbf{S S}(G_{H_{t}})\cap(U\times\gamma^{\circ a})\subset Y\times T^{*}X\quad.
\]

我们设定：

\[
G^{\prime}=\phi_{\gamma}^{-1}R\phi_{\gamma\ast}(G_{I I_{\tau}})~.
\]

足以表明：

\[
\mathrm{S S}(G^{\prime})\cap(U\times\gamma^{\circ a})\subset Y\times T^{*}X\quad.
\]

事实上，假设(6.6.4)，我们根据命题5.2.3得到 $ \mathrm{SS}(G') \cap \pi^{-1}(U) \subset Y \times_X T^*X $ ，并且仍然需要应用命题5.4.5来完成证明，因为根据命题5.2.3在 $ \mathbf{D}^b(Y; p) $ 中 $ G' \simeq G $ 。

最后，(6.6.4) 由(6.6.2) 和下一个引理得出。 ⑨

\statementlabel{引理 6.6.3}令  $E$  为实数有限维向量空间， $\gamma$  为  $E$  与  $\gamma \ni 0$  的闭凸真圆锥，并令  $G \in \mathrm{Ob}(\mathbb{D}^b(E))$  、  $G' = \phi_\gamma^{-1} R \phi_{\gamma*}(G)$  。令  $x \in E$  并假设对于  $x$  的紧邻邻域  $K$  ，  $(K + \gamma) \cap \mathrm{supp}(G)$  是紧邻的。让  $\xi \in E^*$  和  $(x + \gamma; \xi) \cap \mathrm{SS}(G) = \emptyset$  。然后 $(x; \xi) \notin \mathrm{SS}(G')$ 。

\statementlabel{证明}用  $ q_1 $  和  $ q_2 $  表示  $ E \times E $  的第一和第二投影，并用 s 表示映射：

\[
s:E\times E\to E~,\qquad s(x,y)=y-x~.
\]

那么根据命题3.5.4，

\[
G^{\prime}\simeq R s_{*}\bigg(q_{1}^{-1}A_{\gamma}\overset{L}{\otimes}q_{2}^{-1}G\bigg).
\]

现在我们有：

\[
\mathrm{SS}(G^{\prime})\subset\{(x;\xi);\mathrm{there}\mathrm{exists}y,(y,x+y;-\xi,\xi)\in\mathrm{SS}(A_{\gamma})\times\mathrm{SS}(G)\}
\]

这就完成了证明。 ⑨

\subsection{微局部化与逆像}
在本节中，我们将借助微支撑完成 IV §3 中获得的结果。

令  $f: Y \to X$  为流形的态射，并令  $M$  （分别为  $N$  ）为  $X$  （分别为  $Y$  ）和  $f(N) \subset M$  的闭子流形。我们将 IV  $\S3$  中的  $t'f', f_\pi, t'f_N$  、 $f_{N\pi}$  映射表示为：

\[
\begin{aligned}{}&{{}T^{*}Y\xleftarrow{t f^{\prime}}Y\times\overset{T}{}_{X}X\xrightarrow{\quad f_{\pi}\quad}T^{*}X}\\ {}&{{}\quad\quad\quad\quad\quad\quad\quad\quad\quad\quad\quad\quad\quad\quad\quad\quad\quad}\\ {}&{{}T_{N}^{*}Y\xleftarrow{t f_{N}^{\prime}}N\times\overset{T_{M}}{}_{M}X\xrightarrow{\quad f_{N\pi}\quad}T_{M}^{*}X}\\ \end{aligned},
\]

\statementlabel{定理 6.7.1}令  $V$  为  $T_N^*Y$  的开子集，并令  $F\in\mathbf{D}^b(X)$  。假设：

(i) f 对于 V 上的 F 来说是非特征的（参见定义 6.2.7），

(ii)  $ f_{N\pi}|_{\tau_{f_{N}^{-1}(V)}}:^{\tau}f_{N}^{\prime-1}(V) \to T_{M}^{*}X $  对于  $ C_{T_{M}^{*}X}(\mathrm{SS}(F)) $  来说是非特征的，

（三） $ t f^{\prime-1}(V) \cap f_{\pi}^{-1}(\mathrm{SS}(F)) \subset Y \times_{X} T_{M}^{*}X $ 。

然后是自然态射（参见命题 4.3.5）：

\[
R^{\iota}f_{N!}^{\prime}(\omega_{N/M}\otimes f_{N\pi}^{-1}\mu_{M}(F))|_{V}\to\mu_{N}(\omega_{Y/X}\otimes f^{-1}F)|_{V}
\]

and

\[
(\mu_{N}(f^{!}F))|_{V}\to(R^{t}f_{N*}^{\prime}f_{N\pi}^{!}\mu_{M}(F))|_{V}
\]

是同构。

\statementlabel{证明}证明的思路与命题 4.2.5 或 4.3.5 相同。首先，请注意，通过 (i) 和备注 6.4.6， $ f_{\pi}^{-1}(\mathrm{SS}(F)) \to T^*Y $  在  $ V $  的邻域上是正确的。因此，鉴于推论 6.4.4，命题 4.3.5 中的垂直箭头是 (i) 和 (ii) 的 $ V $  的同构。因此，只要证明定理中的第一个态射是同构就足够了。考虑一下我们将保留其符号的图（4.1.11）。然后：

\[
\begin{aligned}{(T_{N}f)^{-1}\nu_{M}(F)}&{{}\simeq s_{Y}^{-1}\tilde{f}^{\prime-1}R j_{X*}\tilde{p}_{X}^{-1}F\enspace,}\\ {\nu_{N}(f^{-1}F)}&{{}\simeq s_{Y}^{-1}R j_{Y*}\tilde{f}^{-1}\tilde{p}_{X}^{-1}F\simeq s_{Y}^{-1}(R j_{Y*}\tilde{f}^{\prime}\tilde{p}_{X}^{-1}F\otimes\omega_{\tilde{Y}_{N}/\tilde{X}_{M}}^{\otimes-1})}\\ {}&{{}\simeq s_{Y}^{-1}(\tilde{f}^{\prime!}R j_{X*}\tilde{p}_{X}^{-1}F\otimes\omega_{\tilde{Y}_{N}/\tilde{X}_{M}}^{\otimes-1})\enspace.}\\ \end{aligned}
\]

考虑一个独特的三角形：

\[
\begin{array}{r}{\omega_{\tilde{\Upsilon}_{N}/\tilde{\Upsilon}_{M}}\otimes\tilde{f}^{\prime-1}R j_{X*}\tilde{p}_{X}^{-1}F\underset{\alpha}{\longrightarrow}\tilde{f}^{\prime!}R j_{X*}\tilde{p}_{X}^{-1}F\underset{+1}{\longrightarrow}H\overset{\tiny}{\longrightarrow}.}\end{array}
\]

然后  $H$  由  $T_N Y$  支持，为了证明该定理，足以证明  $V \cap \mathrm{SS}((s_Y^{-1}H)^\wedge) = \emptyset$  。在这里，我们确定  $T^* T_N Y$  和  $T^* T_N^* Y$  ，如命题 5.5.1 中所示。由推论6.4.4，足以证明：

\[
\left\{\begin{matrix}{{f o r~a n y}\boldsymbol{q}_{\circ}\in V,{t h e r e~e x i s t s}\boldsymbol{q}\in T_{N}Y\times T^{*}\tilde{Y}_{N}{s u c h~t h a t}}\\ {\boldsymbol{q}_{\circ}=^{\prime}s_{Y}^{\prime}(\boldsymbol{q})\in V{a n d}\tilde{f}^{\prime}{i s n o n-c h a r a c t e r i s t i c}{f o r}R\boldsymbol{j}_{X*}\tilde{p}_{X}^{-1}F{a t}\boldsymbol{q}}\\ \end{matrix}\right.
\]

我们在 Y 上选择  $ X $  上的局部坐标系  $ (x)=(x',x'') $  、  $ (y)=(y',y'') $  ，使得  $ M=\{(x',x'');x'=0\} $  、  $ N=\{(y',y'');y'=0\} $  。我们用 $ (x',x'',t;\xi',\xi'',\tau) $ 表示

 $ \widetilde{X}_{M} $  上的坐标， $ (y', y'', t; \eta', \eta'', \tau)  \widetilde{Y}_{N} $  上的坐标。我们写：

\[
f(y^{\prime},y^{\prime \prime})=(g(y^{\prime},y^{\prime \prime}),h(y^{\prime},y^{\prime \prime})),
\]

\[
\tilde{f}(y^{\prime},y^{\prime\prime},t)=(\tilde{g}(y^{\prime},y^{\prime\prime},t),\tilde{h}(y^{\prime},y^{\prime\prime},t),t).
\]

因此：

\[
t\tilde{g}(y^{\prime},y^{\prime \prime},t)=g(ty^{\prime},y^{\prime \prime})~,
\]

\[
\tilde{h}(y^{\prime},y,t)=h(ty^{\prime},y^{\prime \prime})~.
\]

请注意：

\[
\begin{aligned}{}&{{}t f^{\prime}(t y^{\prime},y^{\prime \prime})\cdot(\xi^{\prime},t\xi^{\prime \prime})}\\ {}&{{}\quad=(\tilde{g}_{y^{\prime}}(y,t)\cdot\xi^{\prime}+\tilde{h}_{y^{\prime}}(y,t)\cdot\xi^{\prime \prime},\tilde{g}_{y^{\prime \prime}}(y,t)\cdot t\xi^{\prime}+\tilde{h}_{y^{\prime \prime}}(y,t)\cdot t\xi^{\prime \prime})\;.}\\ \end{aligned}
\]

我们设置 $ q = (0, y_o^o, 0; \eta_o', 0, \tau_o) $ ， $ x_o'' = h(0, y_o'') $ 。然后 $ q_o = (y_o''; \eta'_o) \in V $ 。让我们假设  $ \tilde{f}' $  是  $ R j_x \tilde{p}_x^{-1} F $  在  $ q $  的特征。应用命题 6.2.4 我们发现序列：

\[
\{(y_{n}^{\prime},y_{n}^{\prime\prime},t_{n}^{\prime})\}\mathrm{~i n~}\widetilde{Y}_{N},\qquad\{(x_{n}^{\prime},x_{n}^{\prime\prime},t_{n};\xi_{n}^{\prime},\xi_{n}^{\prime\prime},\tau_{n})\}\mathrm{~i n~}\mathrm{S S}(R j_{X*}\widetilde{p}_{X}^{-1}F),
\]

这样：

\[
\left\{\begin{aligned}{(y_{n}^{\prime},y_{n}^{\prime\prime},t_{n}^{\prime})\xrightarrow[n]{}(0,y_{\circ}^{\prime\prime},0),}\\ {(x_{n}^{\prime},x_{n}^{\prime\prime},t_{n})\xrightarrow[n]{}(0,x_{\circ}^{\prime\prime},0),}\\ \end{aligned}\right.
\]

\[
\tau_{n}+\tilde{g}_{t}(y_{n}^{\prime},y_{n}^{\prime\prime},t_{n}^{\prime})\cdot\xi_{n}^{\prime}+\tilde{h}_{t}(y_{n}^{\prime},y_{n}^{\prime\prime},t_{n}^{\prime})\cdot\xi_{n}^{\prime\prime}\underset{n}{\longrightarrow}\tau_{\circ}~,
\]

\[
\left\{\begin{aligned}{}&{{}\tilde{g}_{y^{\prime}}(y_{n}^{\prime},y_{n}^{\prime\prime},t_{n}^{\prime})\cdot\xi_{n}^{\prime}+\tilde{h}_{y^{\prime}}(y_{n}^{\prime},y_{n}^{\prime\prime},t_{n}^{\prime})\cdot\xi_{n}^{\prime\prime}\mathop{\longrightarrow}_{n}\eta_{\circ}^{\prime},}\\ {}&{{}\tilde{g}_{y^{\prime\prime}}(y_{n}^{\prime},y_{n}^{\prime\prime},t_{n}^{\prime})\cdot\xi_{n}^{\prime}+\tilde{h}_{y^{\prime\prime}}(y_{n}^{\prime},y_{n}^{\prime\prime},t_{n}^{\prime})\cdot\xi_{n}^{\prime\prime}\mathop{\longrightarrow}_{n}0,}\\ \end{aligned}\right.
\]

\[
|\tau_{n}|+|\xi_{n}^{\prime}|+|\xi_{n}^{\prime\prime}|\xrightarrow[n]{}\infty,
\]

\[
|({x}_{n}^{\prime},{x}_{n}^{\prime\prime},t_{n})-\tilde{f}({y}_{n}^{\prime},{y}_{n}^{\prime\prime},t_{n}^{\prime})|\cdot(|\tau_{n}|+|\xi_{n}^{\prime}|+|\xi_{n}^{\prime\prime}|)\xrightarrow[n]{}0.
\]

由(6.7.4)和(6.7.6)我们得到：

\[
|\xi_{n}^{\prime}|+|\xi_{n}^{\prime\prime}|\xrightarrow[n]{}\infty.
\]

此外，(6.7.7) 特别意味着：

\[
|t_{n}-t_{n}^{\prime}|\cdot(|\tau_{n}|+|\xi_{n}^{\prime}|+|\xi_{n}^{\prime\prime}|)\xrightarrow[n]{}0\enspace.
\]

首先假设  $ t_n > 0 $  。然后  $ (t_nx'_n, x''_n; \xi'_n, t_n\xi''_n) \in \mathrm{SS}(F) $  ，我们从 (6.7.2), (6.7.5) 和 (6.7.9) 推出：

\[
t f^{\prime}(t_{n}y_{n}^{\prime},y_{n}^{\prime\prime})\cdot(\xi_{n}^{\prime},t_{n}\xi_{n}^{\prime\prime})\xrightarrow[n]{}(\eta_{\circ}^{\prime},0)\quad.
\]

由于  $f$  对于  $p = (0, y''; \eta'_0, 0)$  处的  $F$  来说是非特征的（假设 (i)），这意味着：

\[
\left\{|\xi_{n}^{\prime}|+|t_{n}\xi_{n}^{\prime\prime}|\right\}_{n}\qquad\mathrm{i s~b o u n d e d}.
\]

因此 $ |\xi_{n}| \to +\infty $  。

由于序列  $\{(\xi_n', t_n \xi_n'')\}$  是有界的，我们可以假设在提取子序列后，  $\xi_n' \to \xi_1', t_n \xi_n'' \to \xi_1''.$  然后由 (6.7.10),  $(0, y_o''; \xi_1', \xi_1'') \in t f'^{-1}(V) \times_X f_\pi^{-1}(\mathrm{SS}(F)).$  由假设 (iii)，我们得到  $\xi_1'' = 0$  。因此 $(\xi_n', t_n \xi_n'') \to_n (\xi_1', 0)$ 。我们可以假设  $\xi_n''/|\xi_n'|$  有一个极限  $\xi_2''.$  我们在  $\mathrm{SS}(F)$  中获得一个序列  $\{(t_n x_n', x_n''; \xi_n', t_n \xi_n'')\}$  使得：

\[
(t_{n}x_{n}^{\prime},x_{n}^{\prime\prime})\xrightarrow[n]{}(0,x_{\circ}^{\prime\prime})~,
\]

\[
\left(\xi_{n}^{\prime},t_{n}\xi_{n}^{\prime\prime}\right)\underset{n}{\longrightarrow}\left(\xi_{1}^{\prime},0\right),
\]

\[
\frac{1}{t_{n}|\xi_{n}^{\prime\prime}|}(t_{n}x_{n}^{\prime},t_{n}\xi_{n}^{\prime\prime})\xrightarrow[n]{}(0,\xi_{2}^{\prime\prime})
\]

那么  $(0,x_{\circ}^{\prime};\xi_{1}^{\prime},\xi_{2}^{\prime})\in C_{T_{M}^{*}X}(\mathrm{SS}(F))$  ，假设 (ii) 意味着  $h_{y^{\prime\prime}}(0,y_{\circ}^{\prime\prime})\cdot\xi_{2}^{\prime\prime}=\tilde{h}_{y^{\prime\prime}}(0,y_{\circ}^{\prime\prime},0)\cdot\xi_{2}^{\prime\prime}\neq0$  。到（6.7.5）我们有（回想一下 $|\xi_{n}^{\prime\prime}|\to\infty):$ 

\[
\tilde{g}_{y^{\prime\prime}}(y_{n}^{\prime},y_{n}^{\prime\prime},t_{n})\cdot\frac{\xi_{n}^{\prime}}{|\xi_{n}^{\prime\prime}|}+\tilde{h}_{y^{\prime\prime}}(y_{n}^{\prime},y_{n}^{\prime\prime},t_{n})\cdot\frac{\xi_{n}^{\prime\prime}}{|\xi_{n}^{\prime\prime}|}\xrightarrow[n]{}0.
\]

我们得到了一个矛盾，因为  $ |\xi_{n}'| $  是有界的。

最后假设  $ t_n = 0 $  。根据定理 6.3.1， $ \mathrm{SS}(\tilde{p}_X^{-1}F) $  中存在双序列  $ \{(x_{n,m}^{\prime}, x_{n,m}^{\prime\prime}, t_{n,m}; \xi_{n,m}^{\prime}, \xi_{n,m}^{\prime\prime}, \tau_{n,m})\} $ ，使得：

\[
(x_{n,m}^{\prime},x_{n,m}^{\prime\prime},t_{n,m};\xi_{n,m}^{\prime},\xi_{n,m}^{\prime\prime},\tau_{n,m})\xrightarrow[m]{}(x_{n}^{\prime},x_{n}^{\prime\prime},0;\xi_{n}^{\prime},\xi_{n}^{\prime\prime},\tau_{n})
\]

 $ t_{n,m} $  是正数。

因此我们可以选择满足(6.7.3)-(6.7.7)的子序列，证明就完成了。   $ \Box $ 

\statementlabel{备注 6.7.2}我们可以通过图来分解f，如下：

\textbackslash{}[\textbackslash{}begin\{对齐\}\&Y\textbackslash{}quad\textbackslash{}hookrightarrow\textbackslash{}quad Y\textbackslash{}times X\textbackslash{}xlongequal\{\textbackslash{}quad\}\textbackslash{}quad Y\textbackslash{}times X\textbackslash{}quad\textbackslash{}longrightarrow\textbackslash{}quad X\textbackslash{}\textbackslash{}\&\textbackslash{}quad\textbackslash{}cup\textbackslash{}quad\textbackslash{}quad

那么就可以从以下两个推论推导出定理 6.7.1。

\statementlabel{推论 6.7.3}在定理 6.7.1 的情况下，假设 (i) 并且：

\[
f|_{N}\colon N\to M{\quad~i s~s m o o t h~}.
\]

那么结论成立。

事实上，在这种情况下，  $ f_{N\pi} $  是平滑的，而  $ f^{\prime-1}(T_N^*Y) \subset Y \times_X T_M^*X $  是平滑的。

\statementlabel{推论 6.7.4}在定理 6.7.1 的情况下，假设  $Y = X$  、  $f$  是恒等式（因此：  $N \subset M \subset X$  和  $N \times_M T_M^* X = T_N^* X \cap T_M^* X$  ）。假设(ii)。然后是自然态射：

\[
\mu_{N}(F)|_{V\cap T_{M}^{*}X}\to(f_{N\pi}^{!}\mu_{M}(F))|_{V\cap T_{M}^{*}X}
\]

是同构。

\statementlabel{证明}假设 (ii) 意味着，在  $V$  上， $C_{T_N^*X}(\mathrm{SS}(F)) \cap T_{N\times_M T_N^*X}^{*}(T_M^*X)$  包含在零部分中。特别是，  $C_{N^*x_M T_N^*X}(\mathrm{SS}(F)) \cap T_N^*X)$  也是如此，因为  $T_{N\times_M T_N^*X}^{*}(T_M^*X) \simeq T_{N^*x_M T_N^*X}(T_N^*X)$  是通过切丛和余切丛来识别的。因此，在  $V \subset T_N^*X$  中存在  $V \cap T_N^*X$  的开邻域  $W$  ，使得  $\mathrm{SS}(F) \cap W \subset T_N^*X$  。那么定理 6.7.1 的假设成立，将  $V$  替换为  $W$  。   $\square$ 

\statementlabel{推论 6.7.5}令 $f: Y \to X$  为流形的态射， $V$  为 $T^*Y$  、 $F \in \mathrm{Ob}(\mathbb{D}^b(X))$  和 $G \in \mathrm{Ob}(\mathbb{D}^b(Y))$  的子集。假设  $f$  对于  $V \cap \mathrm{SS}(G)$  上的  $F$  来说是非特征的。那么自然态射（参见命题 4.4.5）：  $\mu \mathrm{hom}(G, f^!F) \to R^t f'_\star \mu \mathrm{hom}(G \to F)$  和  $\mu \mathrm{hom}(f^{-1}F, G) \to R^t f'_\star \mu \mathrm{hom}(F \leftarrow G)$  是  $V$  上的同构。

\statementlabel{证明}在命题 4.4.5 的证明中使用推论 6.7.3。 □

\statementlabel{推论 6.7.6}在推论 6.7.5 的情况下，还假设  $f$  是封闭嵌入。那么自然态射（参见命题 4.4.6）：  $Rf_!f_{\pi}^{-1}\mu h o m(Rf_!G, F) \to \mu h o m(G, f^!F)$  和  $Rf_!f_{\pi}^{-1}\mu h o m(F, Rf_*G) \to \mu h o m(f^{-1}F, G)$  是  $V$  上的同构。

\statementlabel{证明}使用推论 6.7.5 和命题 4.4.5 来证明命题 4.4.6 (i) 和 (ii)。 ⑨

\subsection{习题}
\statementlabel{练习 六.1}设  $X$  为向量空间  $E$  的开子集， $\gamma$  为  $E$  的真闭凸锥，并设  $F \in \mathrm{Ob}(\mathbf{D}^b(X))$  。假设 $\mathrm{SS}(F) \subset X \times \gamma^{\circ a}$ 。

(i) 证明对于任何 $ x \in X $ ，存在 $ x $  和 $ G \in \mathrm{Ob}(\mathbf{D}^b(E_y)) $  的邻域 $ U $ ，使得 $ F|_{U} \simeq (\phi_y^{-1}G)|_{U} $  。

(ii) 证明对于所有 j， $ \mathrm{SS}(H^{j}(F)) \subset X \times \gamma^{\circ a} $ 。

\statementlabel{练习 六.2}让  $ Z = \{(x, y) \in \mathbb{R}^2, x^2 \geqslant y > -x^2\} $  并让  $ F = A_Z $  。

(i) 证明  $ \mathrm{SS}(F)=\{(x,y;\xi,\eta);\ x^{2}\geqslant y\geqslant-x^{2},\ \xi=\eta=0\ \text{or}\ y=-x^{2},  \xi=2x\eta,\eta\leqslant0\text{or}y=x^{2},\xi=-2x\eta,\eta\leqslant0\} $  。

(ii) 让 $ p = (0,0;0,0) $  。证明 $ C_p(\mathbf{S}\mathbf{S}(F)) = \{ y = 0, \xi = 0, \eta \leq 0 \} $ 。

(iii) 让 $ \theta = \mathrm{dy} $  。显示 $ C_p(\mathrm{SS}(F)) \subset \{\theta^{-1}(0)\} $  但 $ -H_\theta \notin C_p(\mathrm{SS}(F)) $  。 （参见备注 6.5.5。）

\statementlabel{练习 六.3}令 $F \in \mathrm{Ob}(\mathbb{D}^b(X))$  并假设 $H^j(F) = 0$  为 $j < 0$  。令 $u \in H^0(X; F) = \Gamma(X; H^0(F))$  、 $x \in X$  和 $\gamma$  为 $T_x X$  中的闭合凸真圆锥，以及 $\gamma \ni 0$  。假设  $(\mathrm{SS}(F) \cap \pi^{-1}(x)) \cap \gamma^{\circ a} \subset \{0\}$  和  $C_x(\mathrm{supp}(u)) \cap \gamma \subset \{0\}$  。证明 $x \notin \mathrm{supp}(u)$ 。 （提示：使用命题 5.2.1。）

\statementlabel{练习 六.4}令  $f: Y \to X$  为流形的态射，并令  $F_1$  和  $F_2$  属于  $\mathbf{D}^b(X)$  。

(i) 假设  $f$  对于  $\mathrm{SS}(F_1) \oplus \mathrm{SS}(F_2)^a$  来说是非特征的。证明同构：

\[
f^{-1}R\mathcal{H}o m(F_{2},F_{1})\simeq R\mathcal{H}o m(f^{-1}F_{2},f^{-1}F_{1})~.
\]

(ii) 假设  $f$  对于  $\mathrm{SS}(F_1) \nearrow \mathrm{SS}(F_2)$  来说是非特征的。证明同构：

\[
f^{-1}F_{1}\stackrel{L}{\otimes}f^!F_{2}\simeq f^!\biggl(F_{1}\stackrel{L}{\otimes}F_{2}\biggr).
\]

\statementlabel{练习 六.5}令  $M$  为  $X$  的闭子流形， $\widetilde{X}$  覆盖  $X\backslash M$  和  $\rho$  映射  $\widetilde{X}\to X$  （即：  $\rho=i\circ p$  其中  $i:X\backslash M\hookrightarrow X$  和  $p$  是投影  $\widetilde{X}\to X\backslash M$  ）。

致 $F \in \mathrm{Ob}(\mathbb{D}^{b}(X))$ 一位同事：

\[
\begin{aligned}{\rho_{M}(F)}&{{}=R\rho_{*}\rho^{-1}F}\\ {}&{{}\simeq R\mathcal{H}{o m}(R\rho_{!}A_{\tilde{X}},F).}\\ \end{aligned}
\]

(i) 证明：

\[
\mathrm{SS}(\rho_{M}(F))\subset\mathrm{SS}(F)\cup(\mathrm{SS}(F)\hat{+}T_{M}^{*}X).
\]

(ii) 令 $f: Y \to X$  为横截于 $M$  的流形态射并设置 $N = f^{-1}(M)$  。假设  $f$  对于  $F$  和  $\mathrm{SS}(F) \widehat{+} T_{M}^{*}X$  来说是非特征的。证明同构：

\[
f^{-1}\rho_{M}(F)\simeq\rho_{N}(f^{-1}F).
\]

这里  $ \rho_N $  类似地通过使用  $ Y \setminus N $  的覆盖  $ \tilde{X} \times_X Y $  来定义。

\statementlabel{练习 六.6}令 $ X = \mathbb{R}^2 $  和 $ (x, y) $  为 $ X $  上的坐标， $ (x, y; \xi, \eta) $  为 $ T^*X $  上的关联坐标。设置 $ \Omega = \{(x, y; \xi, \eta); \eta > 0\} $ ， $ \Omega' = \Omega \setminus \{(x, y; \xi, \eta);\eta > 0, x \leq 0, \xi = 0\} $ 。

(i) 证明对于  $ F \in \mathrm{Ob}(\mathbb{D}^{b}(X)) $  ，  $ \mathrm{SS}(F) \cap \Omega' = \varnothing $  意味着  $ \mathrm{SS}(F) \cap \Omega = \varnothing $  。

(ii) 证明对于任何  $F$  ,  $G \in \mathrm{Ob}(\mathbb{D}^{b}(X))$  ，

\[
\mathrm{H o m}_{\mathbb{D}^{b}(X;\Omega)}(F,G)\to\mathrm{H o m}_{\mathbb{D}^{b}(X;\varOmega^{\prime})}(F,G)
\]

是同构。

(iii) 证明 $ H^0(\Omega'; \mu \text{hom}(A_{\{0\}}, A_{\{0\}})) \simeq A^2 \text{but Hom}_{\mathbb{D}^b(X; \Omega')}(A_{\{0\}}, A_{\{0\}}) \simeq A $  。 （提示：使用对合定理。）

\statementlabel{练习 六.7}令  $\gamma$  为曲线  $\{(x,y)\in\mathbb{R}^2; y=0, x\leq0\}\cup\{(x,y)\in\mathbb{R}^2; y=x^2, x\geq0\}$  ，其中  $1<\lambda<2$  。让 $A=T_y^*\mathbb{R}^2$  。证明

\[
\varLambda+\widehat{\varLambda}=\varLambda\cup T_{\{0\}}^{*}\mathbb{R}^{2}.
\]

\statementlabel{练习 六.8}让 $ y = \{(x, y) \in \mathbb{R}^2; xy = 0, x \geq 0, y \geq 0\} $  。证明如果  $ F \in \mathrm{Ob}(\mathbb{D}^b(X)) $  和  $ \mathrm{SS}(F) \subset \mathrm{SS}(A_y) $  ，则存在  $ M \in \mathrm{Ob}(\mathbb{D}^b(\mathfrak{Mod}(A))) $  ， $ F \simeq M_y $  存在于  $ \mathbb{D}^b(\mathbb{R}^2) $  中。

\statementlabel{练习 六.9}我们保留命题 6.1.8 的符号。证明同构：

\[
\mathrm{H o m}_{\mathsf{D}^{b}(X;p_{X})^{\wedge}}(f_{!}^{\mu}G,F)\simeq\mathrm{H o m}_{\mathsf{D}^{b}(Y;p_{Y})^{\vee}}(G,f_{\mu}^{!}F)\simeq\mu\mathcal{h}o m(G\to F)_{p}\enspace,
\]

\[
\mathrm{H o m}_{\mathsf{D}^{b}(X;p_{X})^{\vee}}(F,f_{*}^{\mu}G)\simeq\mathrm{H o m}_{\mathsf{D}^{b}(Y;p_{Y})^{\wedge}}(f_{\mu}^{-1}F,G)\simeq\mu\mathcal{h}o m(F\leftarrow G)_{p}\enspace.
\]

\subsection{注记}
余切束中的正锥体的研究（即§2的结果）始于 Kashiwara-Schapira [1]。它的动机是这样的评论：偏微分方程的双曲性的经典概念可以自然地用这样的锥体来表示，事实上，定理 6.3.1、6.4.1 和 6.7.1 是微双曲系统结果的层理论版本，参见。命题 11.5.4 和 11.5.8 如下，进一步发展请参阅 D'Agnolo-Schapira [1]。

 $ \hat{+} $  操作首先在 Kashiwara-Schapira [2] 中引入，现在作为研究乘积（例如分布乘积，参见 Lebeau [1]）或研究分层（参见第八章）时的自然工具出现。

对合性定理（定理 6.5.4）是线性微分方程组相应结果的广泛推广，由 Sato-Kawai-Kashiwara [1]（在 Guillemin-Quillen-Sternberg [1] 的基础论文之后）首先使用解析方法证明。有趣的是，

微分方程的对合性定理现在具有三个完全不同的证明：第一个是解析的，如上所述，第二个是纯代数的，归因于 Gabber [1]，最后一个是纯粹的“几何”（和“实数”），使用定理 6.5.4（参见下面的定理 11.3.3）。该定理首先由 Kashiwara-Schapira [2, 3] 获得，其形式不太精确，并采用了不同的证明。

命题 6.5.6 在研究微微分方程解析奇点的传播，特别是衍射问题时非常有用，这是由 Schapira [3] 提出的。本章的所有其他结果均归功于作者，其中许多结果已在 Kashiwara-Schapira [3] 中得到证明。

\section{接触变换与纯层}
\subsection{概要}
在本章中，我们对层进行接触变换。我们首先将 III §6 中引入的核概念扩展到微局部情况，并开发了核的微局部演算。

令 $X$ 和 $Y$ 为两个流形， $\Omega_X$ 和 $\Omega_Y$ 分别是 $T^*X$ 和 $T^*Y$ 的两个开子集。我们证明，在适当的假设下，函子 $\Phi_K$ 和 $\Psi_K$ 分别从 $\mathbf{D}^b(Y; \Omega_Y)$ 到 $\mathbf{D}^b(X; \Omega_X)$ 和从 $\mathbf{D}^b(X; \Omega_X)$ 到 $\mathbf{D}^b(Y; \Omega_Y)$ 被明确定义，并给出了范畴的等价性。此外，这些等价与函子  $\mu\text{hom}$  兼容。接下来，如果  $\chi: \Omega_X \simeq \Omega_Y$  是接触变换，我们表明在收缩  $\Omega_X$  和  $\Omega_Y$  之后总是可以使用这些内核构造等价  $\mathbf{D}^b(X; \Omega_X) \simeq \mathbf{D}^b(Y; \Omega_Y)$  。现在让  $M$  和  $N$  分别为  $X$  和  $Y$  的两个超曲面，并假设接触变换  $\chi$  互换了  $T_M^*X \cap \Omega_X$  和  $T_N^*Y \cap \Omega_Y$  。如果  $\chi$  的图与  $X \times Y$  的超曲面  $S$  的共形丛相关联，并且如果选择束  $A_S$  作为内核  $K$  ，则可以证明  $\mathbf{D}^b(X; p)$  、  $(p \in \Omega_X)$  中的  $\Phi_K(A_N) \simeq A_M[d]$  ，其中  $d$  是一个移位我们使用惯性指数进行计算。

该计算引出了沿光滑拉格朗日流形  $A$  的纯层概念，这是 Hörmander [2]、[4] 的傅里叶分布概念或 Sato-Kawai-Kashiwara [1] 的简单完整系统概念的层理论类似物。在  $A = T_M^*X$  的情况下，对于  $X$  的闭子流形  $M$  ， $p$  处沿  $A$  的纯束  $F$  只不过是  $L_M[d]$  的  $\mathbf{D}^b(X; p)$  中的图像，其中  $L$  是  $A$  模（因此  $L_M$  是 $X$  上的束由  $M$  支撑， $M$  上的常量与茎  $L$  ) 和  $d$  是一个转变。 （在这种情况下， $F$  是纯的，带有shift  $d + \frac{1}{2} \text{codim} M$  ）当投影 $\pi: A \to X$  的等级不再恒定时， $F$  的移位可能会“跳跃”，其计算需要惯性指数的完整机制。我们通过在获取正像或逆像之后计算两个核的合成的移位以及纯束的移位来结束本章。

除了第十章第 §3 节和第 XI 章第 §4 节之外，本章的内容对于理解本书的其余部分并不是必需的。

我们保持惯例4.0。此外，除非另有说明，否则余切束的所有子流形都应该是局部圆锥的。

\subsection{微局部核}
在本节中，我们将对 III §6 的结构进行“微局部化”。

令 $X$  和 $Y$  为两个流形。我们分别用 $q_1$ 和 $q_2$ 表示投影 $X \times Y \to X$ 和 $X \times Y \to Y$ ，用 $p_1$ 和 $p_2$ 分别表示投影 $T^*(X \times Y) \to T^*X$ 和 $T^*(X \times Y) \to T^*Y$ 。我们还设置 $p_j^a = p_j \circ a$ ，其中“ $a$ ”是对映图。如果 $Z$  是第三个流形，我们用 $q_{ij}$  表示从 $X \times Y \times Z$  到第 $(i,j)$  因子的投影。例如  $q_{13}$  是到  $X \times Z$  的投影。类似地定义了  $T^*(X \times Y \times Z) \simeq T^*X \times T^*Y \times T^*Z$  的投影  $p_{ij}$  。在本章中， $T^*(X \times Y) \times_{T^*Y} T^*(Y \times Z)$  表示投影 $p_1: T^*(Y \times Z) \to T^*Y$  和 $p_2: T^*(X \times Y) \to T^*Y$  的纤维乘积。我们通过以下方式将这个集合标识为 $T^*(X \times Y \times Z)$ ：

\[
(((x;\xi),(y;-\eta)),((y;\eta),(z;\zeta))){\leftrightarrow}((x;\xi),(y;\eta),(z;\zeta))\enspace.
\]

令  $\Omega_X$  、  $\Omega_Y$  、  $\Omega_Z$  分别为  $T^*X$  、  $T^*Y$  、  $T^*Z$  的开子集。一组设置  $\Omega_X^a = a(\Omega_X)$  ，其中  $a$  是对映图， $\Omega_Y^a, \Omega_Z^a$  等类似。

\statementlabel{定义 7.1.1}其中一个用  $\mathbf{N}(X, Y; \Omega_X, \Omega_Y)$  表示  $\mathbf{D}^b(X \times Y; \Omega_X \times T^*Y)$  的满子范畴，由满足以下条件的对象  $K$  组成：

（一） $ \mathrm{SS}(K) \cap (\Omega_X \times T^*Y) \subset \Omega_X \times \Omega_Y^a $ ，

(ii)  $ p_1: \mathrm{SS}(K) \cap (\Omega_X \times T^*Y) \to \Omega_X $  是正确的。

如果不存在混淆的风险，我们会写  $\mathbb{N}(\Omega_{X},\Omega_{Y})$  而不是  $\mathbb{N}(X,Y;\Omega_{X},\Omega_{Y})$  。

当然，如果  $\Omega_X' $  和  $\Omega_Y'$  分别是  $T^*X$  和  $T^*Y$  的开子集，并且  $\Omega_X' \subset \Omega_X$  和  $\Omega_Y' \subset \Omega_Y$  ，则  $\mathrm{Ob}(\mathbb{N}(\Omega_X, \Omega_Y))$  包含在  $\mathrm{Ob}(\mathbb{N}(\Omega_X', \Omega_Y))$  中。如果  $Y = \{\mathrm{pt}\}$  ，则：

\[
\aleph(\Omega_{X},\left\{\mathsf{p t}\right\})\simeq\mathbb{D}^{b}(X;\Omega_{X})~.
\]

让 $ K \in \text{Ob}(\mathbb{D}^b(X \times Y)) $ ， $ L \in \text{Ob}(\mathbb{D}^b(Y \times Z)) $ 。回想一下， $ \text{III} $  §6， $ K \circ L $  是  $ \text{Ob}(\mathbb{D}^b(X \times Z)) $  的对象，定义如下：

\[
K\circ L=R q_{13!}\bigg(q_{12}^{-1}K\otimes{}^{L}q_{23}^{-1}L\bigg)\;.
\]

\statementlabel{命题 7.1.2}假设 $ K \in \mathrm{Ob}(\mathbb{N}(\Omega_X, \Omega_Y)) $  。然后：

(i) 自然态射  $ K \circ L \to Rq_{13*} (q_{12}^{-1} K \otimes^L q_{23}^{-1} L) $  是  $ \mathbf{D}^b(X \times Z; \Omega_X \times T^*Z) $  中的同构，

\[
\begin{array}{r c l}{\mathrm{S S}(K\circ L)\;\cap\;(\Omega_{X}\times T^{*}Z)\;\subset\;p_{13}((\mathrm{S S}(K)\;\cap\;(\Omega_{X}\times\Omega_{Y}^{a}))\;\times_{T^{*}Y}(\mathrm{S S}(L)\;\cap\;}\\ {(\Omega_{Y}\times T^{*}Z))),}\\ \end{array}
\]

(iii) 如果  $ L \in \mathrm{Ob}(\mathbb{N}(\Omega_Y, \Omega_Z)) $  ，则  $ K \circ L \in \mathrm{Ob}(\mathbb{N}(\Omega_X, \Omega_Z)) $  。

特别是，  $ (K,L)\mapsto K\circ L $  是从  $ \mathbb{N}(\Omega_X,\Omega_Y)\times\mathbb{N}(\Omega_Y,\Omega_Z) $  到  $ \mathbb{N}(\Omega_X,\Omega_Z) $  的双函子。

\statementlabel{证明}首先注意到：

\[
(p_{12}^{-1}(\mathrm{S S}(K))\stackrel{\frown}{+}_{p_{23}^{-1}}(\mathrm{S S}(L))\cap(\Omega_{X}\times T^{*}Y\times T^{*}Z)=\emptyset\enspace.
\]

事实上，在 SS(K) 中取序列  $\{(x_n, y_n, \xi_n, \eta_n)\}$  ，在 SS(L) 中取序列  $\{(y'_n, z_n; \eta'_n, \zeta_n)\}$  ，这样  $(x_n; \xi_n) \to (x_o, \xi_o) \in \Omega_X$  、  $(z_n; \zeta_n) \to (z_o; \zeta_o)$  、  $y_n \to y_o$  、  $y'_n \to y_o$  、  $\eta_n + \eta'_n \to \eta_o$  。那么序列  $\{\eta_n\}$  是有界的，因为  $K \in \mathrm{Ob}(\mathbb{N}(\Omega_X, \Omega_Y))$  。这意味着（7.1.3）。根据推论 6.4.5（参见备注 6.2.5），我们得到：

\[
\begin{array}{r}{\mathrm{S S}\left(q_{12}^{-1}K\otimes q_{23}^{-1}L\right)\cap\Omega_{X}\times T^{*}Y\times T^{*}Z\subset p_{12}^{-1}(\mathrm{S S}(K))+p_{23}^{-1}(\mathrm{S S}(L))\enspace.}\end{array}
\]

关于  $K$  的假设还意味着，对于  $\Omega_X \times T^*Z$  的任何紧致子集  $A$  ，都存在  $Y$  的紧致子集  $B$  ，使得：

\[
(x,y,z;\xi,\eta,\zeta)\in(p_{12}^{-1}(\mathbb{S}\mathbb{S}(K))+p_{23}^{-1}(\mathbb{S}\mathbb{S}(L)))\cap p_{13}^{-1}(A)\Rightarrow y\in B\enspace.
\]

通过（7.1.4）和（7.1.5），我们可以将命题 6.3.3 应用到  $ \Omega_X \times T^*Z $  上的  $ q_{12}^{-1}K \otimes^L q_{23}^{-1}L $  上的映射  $ q_{13} $  。我们得到 (i) 并且：

\[
\mathrm{S S}(K\circ L)\cap(\Omega_{X}\times T^{*}Z)\subset\{(x,z;\xi,\zeta);\mathrm{t h e r e~e x i s t s}(y;\eta)\in T^{*}Y\mathrm{w i t h}
\]

\[
(x,y;\xi,-\eta)\in\mathrm{SS}(K),(x;\xi)\in\Omega_{X}\mathrm{and}(y,z;\eta,\zeta)\in\mathrm{SS}(L)\}.
\]

这就证明了下面的(ii)和(iii)。 □

回想一下，在 III §6 中，我们通过设置将  $ K \in \mathrm{Ob}(\mathbb{D}^b(X \times Y)) $  函子  $ \Phi_K $  :  $ \mathbf{D}^b(Y) \to \mathbf{D}^b(X) $  和  $ \psi_K: \mathbb{D}^b(X) \to \mathbb{D}^b(Y) $  关联起来：

\[
\left\{\begin{aligned}{\varPhi_{K}(G)}&{{}=R q_{1!}\bigg(\overset{L}{K}\otimes q_{2}^{-1}G\bigg)\simeq K\circ G,}\\ {\varPsi_{K}(F)}&{{}=R q_{2*}R\mathcal{H o m}(K,q_{1}^{!}F).}\\ \end{aligned}\right.
\]

\statementlabel{定义 7.1.3}让 $ K \in \text{Ob}(\mathbb{N}(\Omega_X, \Omega_Y)) $  。通过设置  $ \Phi_K(G) = K \circ G $  来定义函子  $ \Phi_K: \mathbb{D}^b(Y; \Omega_Y) \to \mathbb{D}^b(X; \Omega_X) $  。

鉴于命题 7.1.2 与  $ Z = \{pt\} $  一起应用，该定义是有意义的，并且与（7.1.6）兼容。

我们将用  $ r: X \times Y \to Y \times X $  表示规范映射：

\[
r(x,y)=(y,x).
\]

\statementlabel{命题 7.1.4}假设  $ K \in \text{Ob}(\mathbb{D}^b(X \times Y)) $  满足：  $ r_* K \in \mathbb{N}(\Omega_Y^a, \Omega_X^a) $  。然后：

(i)  $ \Psi_K $  诱导出一个从  $ \mathbf{D}^b(X; \Omega_X) $  到  $ \mathbf{D}^b(Y; \Omega_Y) $  的明确定义的函子。

(ii) 自然态射  $ R q_{2!}R \mathcal{H} \text{om}(K, q_1^! F) \to \Psi_K(F) $  是  $ \mathbf{D}^b(Y; \Omega_Y) $  中的同构。

（三） $ \mathrm{SS}(\Psi_K(F)) \cap \Omega_Y \subset p_2^a(\mathrm{SS}(K) \cap p_1^{-1}(\mathrm{SS}(F) \cap \Omega_X)) $ 。

由于证明与命题7.1.2几乎相同，我们不再重复。

令 $W$  为第四个流形， $\Omega_W$  为 $T^*W$  的开子集。

\statementlabel{命题 7.1.5}$(K, L, M) \mapsto (K \circ L) \circ M$  和 $(K, L, M) \mapsto K \circ (L \circ M)$  给出的从 $\mathbf{N}(\Omega_X, \Omega_Y) \times \mathbf{N}(\Omega_Y, \Omega_Z) \times \mathbf{N}(\Omega_Z, \Omega_W)$  到 $\mathbf{N}(\Omega_X, \Omega_W)$  的两个函子是同构的。

证据是显而易见的。 ⑨

\statementlabel{命题 7.1.6}(i)  $(K, L, H) \mapsto \Phi_{K \circ L}(H)$  和 $(K, L, H) \mapsto \Phi_K(\Phi_L(H))$  给出的从 $\mathbf{N}(\Omega_X, \Omega_Y) \times \mathbf{N}(\Omega_Y, \Omega_Z) \times \mathbf{D}^b(Z; \Omega_Z)$  到 $\mathbf{D}^b(X; \Omega_X)$  的两个函子是同构的。

(ii)  $(L, K, F) \mapsto \Psi_{r^{-1}(K \circ L)}(F)$  和 $(L, K, F) \mapsto \Psi_{r^{-1}L}(\Psi_{r^{-1}K}(F))$  给出的从 $\mathbb{N}(\Omega_Z^a, \Omega_Y^a) \times \mathbb{N}(\Omega_Y^a, \Omega_X^a) \times \mathbb{D}^b(X; \Omega_X)$  到 $\mathbb{D}^b(Z; \Omega_Z)$  的两个函子是同构的。

这里  $r$  表示规范映射  $X \times Y \to Y \times X$  、  $Y \times Z \to Z \times Y$  、  $X \times Z \to Z \times X$  之一。

\statementlabel{证明}(i) 是命题 7.1.5 的一个特例。

(ii) 由命题 3.6.4 得出。 ⑨

让  $ q_{ij} $  表示  $ X \times Z \times Y \times Z $  上定义的  $ (i,j) $  次投影。我们用  $ i_z $  表示从  $ \mathbf{N}(\Omega_X, \Omega_Y) $  到  $ \mathbf{N}(\Omega_X \times \Omega_Z, \Omega_Y \times \Omega_Z) $  的函子，给出如下：

\[
i_{Z}\colon K\mapsto q_{13}^{-1}K\overset{L}{\otimes}q_{24}^{-1}A_{A_{Z}}.
\]

考虑一下图表：

\[
\begin{aligned}&\mathbf{N}(\Omega_{X},\Omega_{Y})\times\mathbf{N}(\Omega_{Y},\Omega_{Z})\xrightarrow{\quad\beta_{1}\quad}\mathbf{N}(\Omega_{X},\Omega_{Z})\\&\downarrow\left\lfloor_{\substack{\alpha_{1}}}\right\rfloor_{\substack{\alpha_{2}}}\left\lfloor_{\substack{\alpha_{2}}}\right\rfloor\\&_{z},\Omega_{Y}\times\Omega_{Z})\times\mathbf{D}^{b}(Y\times Z;\Omega_{Y}\times T^{*}Z)\xrightarrow{\quad\beta_{2}\quad}\mathbf{D}^{b}(X\times Z;\Omega_{X}\times T^{*}Z)\\ \end{aligned}
\]

其中  $ \alpha_{1}(K,L)=(i_{Z}(K),L),\alpha_{2}(H)=H,\beta_{1}(K,L)=K\circ L $  和  $ \beta_{2}(H,G)=\Phi_{H}(G) $  。

\statementlabel{命题 7.1.7}图 (7.1.9) 是可交换的，即  $ \beta_2 \circ \alpha_1 \simeq \alpha_2 \circ \beta_1 $  。

证据是显而易见的。 ⑨

\statementlabel{命题 7.1.8}让 $ K \in \text{Ob}(\mathbb{D}^b(X \times Y)) $  。假设  $ K \in \text{Ob}(\mathbb{N}(\Omega_X, \Omega_Y)) $  和  $ r^{-1}K \in \text{Ob}(\mathbb{N}(\Omega_Y^a, \Omega_X^a)) $  。那么函子 $ \Phi_K: \mathbb{D}^b(Y; \Omega_Y) \to \mathbb{D}^b(X; \Omega_X) $  和 $ \Psi_K: \mathbb{D}^b(X; \Omega_X) \to \mathbb{D}^b(Y; \Omega_Y) $  是伴随函子。

\statementlabel{证明}我们知道（参见练习 I.14）：

\[
\mathrm{H o m}_{\mathsf{D}^{b}(X;\;\varOmega_{X})}(\varPhi_{K}(G),F)=\varinjlim\mathrm{H o m}_{\mathsf{D}^{b}(X)}(\varPhi_{K}(G^{\prime}),F^{\prime})\enspace,
\]

其中归纳极限取态射范畴  $ G' \to G $  和  $ F \to F' $  使得  $ G' \simeq G $  在  $ \mathbf{D}^b(Y; \Omega_Y) $  和  $ F \simeq F' $  在  $ \mathbf{D}^b(X; \Omega_X) $  。同样：

\[
\mathrm{H o m}_{\mathsf{D}^{b}(Y;\mathcal{Q}_{Y})}(G,\mathcal{W}_{K}(F))=\varinjlim\mathrm{H o m}_{\mathsf{D}^{b}(Y)}(G^{\prime},\mathcal{W}_{K}(F^{\prime})),
\]

其中归纳极限取自同一范畴。因此，结果来自命题 3.6.2。 ⑨

\statementlabel{命题 7.1.9}让 $ K \in \text{Ob}(\mathbb{D}^b(X \times Y)) $  。假设  $ K $  是上同调可构造的，并假设  $ K \in \text{Ob}(\mathbb{N}(\Omega_X, \Omega_Y)) $  。设置 $ K^* = r_* R \mathcal{H} \text{om}(K, \omega_{X \times Y/Y}) $ 。然后  $ r^{-1} K^* \in \text{Ob}(\mathbb{N}(\Omega_X^a, \Omega_Y^a)) $  和  $ \Phi_K \simeq \Psi_{K^*} $  作为从  $ \mathbb{D}^b(Y; \Omega_Y) $  到  $ \mathbb{D}^b(X; \Omega_X) $  的函子。

\statementlabel{证明}由于  $ \mathrm{SS}(r^{-1}K^*) = \mathrm{SS}(K)^a $  ，  $ r^{-1}K^* $  属于  $ \mathbb{N}(\Omega_X^a, \Omega_Y^a) $  。让 $ G \in \mathrm{Ob}(\mathbb{D}^b(Y)) $  。由(7.1.3)可知：

\[
\left(\mathrm{SS}(K^{*})\stackrel{\frown}{+}_{\infty}\mathrm{SS}(q_{2}^{!}G)\right)\cap(\Omega_{X}\times T^{*}Y)=\emptyset\enspace.
\]

应用推论 6.4.3 我们得到  $ \mathbf{D}^{b}(X \times Y; \Omega_{X} \times T^{*}Y) $  中的同构：

\[
R\mathcal{H}{o m}(K^{\ast},q_{2}^{!}G)\simeq R\mathcal{H}{o m}(K^{\ast},q_{2}^{!}A_{Y})\overset{L}{\otimes}q_{2}^{-1}G.
\]

因此根据命题7.1.2：

\[
\begin{aligned}{\varPsi_{K^{*}}(G)}&{{}\simeq R q_{1*}\bigg(K\overset{L}{\otimes}q_{2}^{-1}G\bigg)}\\ {}&{{}\simeq R q_{1!}\bigg(K\overset{L}{\otimes}q_{2}^{-1}G\bigg)}\\ \end{aligned}
\]

在 $ \mathbf{D}^{b}(X;\Omega_{X}) $  中。 ⑨

\statementlabel{命题 7.1.10}让  $ K \in \text{Ob}(\mathbb{N}(\Omega_X, \Omega_Y)) $  和  $ L \in \text{Ob}(\mathbb{N}(\Omega_Y, \Omega_X)) $  。假设  $ \mathbb{D}^b(X \times X; \Omega_X \times T^*X) $  中的  $ K \circ L \simeq A_{\Delta_X} $  和  $ \mathbb{D}^b(Y \times Y; \Omega_Y \times T^*Y) $  中的  $ L \circ K \simeq A_{\Delta_Y} $  。那么 $ \Phi_K: \mathbb{D}^b(Y; \Omega_Y) \to \mathbb{D}^b(X; \Omega_X) $  和 $ \Phi_L: \mathbb{D}^b(X; \Omega_X) \to \mathbb{D}^b(Y; \Omega_Y) $  是范畴的等价，彼此相反。

\statementlabel{证明}这直接源自命题 7.1.6。 ⑨

\statementlabel{命题 7.1.11}让 $ K \in \text{Ob}(\mathbb{D}^b(X \times Y)) $ ， $ F \in \text{Ob}(\mathbb{D}^b(X; \Omega_X)) $ ， $ G \in \text{Ob}(\mathbb{D}^b(Y; \Omega_Y)) $ 。

(i) 假设  $ K \in \mathrm{Ob}(\mathbb{N}(\Omega_X, \Omega_Y)) $  。那么  $ \mathbb{D}^b(\Omega_X) $  中有一个自然的同构：

\[
R p_{1*}{\mu}\mathcal{h}o m(K,R\mathcal{H}o m(q_{2}^{-1}G,q_{1}^{!}F))\simeq{\mu}\mathcal{h}o m(\varPhi_{K}(G),F)\enspace.
\]

(ii) 假设  $ r_K K \in \text{Ob}(\mathbb{N}(\Omega_Y^a, \Omega_X^a)) $  。那么  $ \mathbb{D}^b(\Omega_Y) $  中有一个自然的同构：

\[
R p_{2*}^{a}\mu\mathcal{h}{o m}(K,R\mathcal{H}{o m}(q_{2}^{-1}G,q_{1}^{!}F))\simeq\mu\mathcal{h}{o m}(G,\varPsi_{K}(F))\enspace.
\]

\statementlabel{证明}(i) 考虑映射：

\[
\begin{aligned}{X\times X}&{{}\xleftarrow[q]{}\quad X\times X\times Y\quad\xrightarrow[j]{}\quad X\times Y\times X\times Y}\\ {\quad\cup\quad}&{{}}\\ {\quad\varDelta_{X}}&{{}\quad\xleftarrow{\quad}\quad\varDelta_{X}\times Y\quad\xrightarrow{\quad}\quad\varDelta_{X\times Y}\quad,}\\ \end{aligned}
\]

以及相关的映射：

\[
\begin{array}{c c c}{T_{\varDelta_{X}}^{*}(X\times X)\xleftarrow[q_{\pi}]{}}&{Y\times T_{\varDelta_{X}}^{*}(X\times X)}\\ {\bigg|\wr}&{}&{\bigg|\wr}\\ {T^{*}X}&{\xleftarrow{\quad}}&{Y\times T^{*}X}\\ \end{array}
\]

\textbackslash{}[\textbackslash{}begin\{array\}\{c c c\}\{T\_\{\textbackslash{}mathcal\{A\}\_\{X\}\textbackslash{}times Y\}\textasciicircum{}\{*\}(X\textbackslash{}times X\textbackslash{}times Y)\textbackslash{};\textbackslash{}xleftarrow[\{t\_\{j\textasciicircum{}\{\textbackslash{}prime\}\}]:\}\textbackslash{};\}T\_\{\textbackslash{}mathcal\{A\}\_\{X\textbackslash{}times Y\}\}\textasciicircum{}\{*\}(X\textbackslash{}times Y\textbackslash{}times X\textbackslash{}times Y)\}\textbackslash{}\textbackslash{} \{\textbackslash{}bigg|\textbackslash{}wr\}\&\{\textbackslash{}bigg|\}\&\{\textbackslash{}wr\}\textbackslash{}\textbackslash{} \{(T\textasciicircum{}\{*\}X)\textbackslash{}times Y\textbackslash{};\textbackslash{};\textbackslash{}

与往常一样，1 表示在  $X \times Y \times X \times Y$  上定义的  $j$  投影， $a_i$  表示  $(i, i)$  上定义的投影。设定：

\[
\begin{aligned}{H}&{{}=R\mathcal{H}{o m}(q_{34}^{-1}K,R\mathcal{H}{o m}(q_{2}^{-1}G,q_{1}^{!}F))}\\ {}&{{}\simeq R\mathcal{H}{o m}\bigg(q_{34}^{-1}K\otimes\overset{L}{q_{2}^{-1}G},q_{1}^{!}F\bigg)\;.}\\ \end{aligned}
\]

然后， $ \mu h o m(K, R \mathcal{H} o m(q_2^{-1}G, q_1^! F)) \simeq \mu_{\Delta_X \times Y}(H) $ 。

映射  $j$  对于  $\Omega_X \times T^* Y \subset T_{\Delta_X \times Y}^* (X \times X \times Y)$  上的  $H$  来说是非特征的。事实上，考虑  $\mathrm{SS}(H)$  中的序列  $\{\langle x_n, y_n, x'_n, y'_n; \xi_n, \eta_n, \xi'_n, \eta'_n \rangle\}$  ，以及  $(x_n, y_n, x'_n, y'_n) \to (x_o, y_o, x_o, y_o)$  、  $(\xi_n, \xi'_n, \eta_n + \eta'_n) \to (\xi_o, -\xi_o, \eta_o)$  、  $(x_o; \xi_o) \in \Omega_X$  。然后序列  $\{\eta'_n\}$  由  $K$  上的假设界定。

应用推论 6.7.3，我们得到  $ \Omega_X \times T^*Y $  的同构：

\[
\begin{aligned}{R^{\mathfrak{t}}j_{\ast}^{\prime}\mu\mathcal{h o m}(K,R\mathcal{H}\mathcal{o m}(q_{2}^{-1}G,q_{1}^{\mathfrak{t}}F))\simeq}&{{}~\mu_{\varDelta_{X}\times Y}j^{\mathfrak{t}}R\mathcal{H}\mathcal{o m}\bigg(q_{34}^{-1}K\otimes q_{2}^{-1}G,q_{1}^{\mathfrak{t}}F\bigg)}\\ {\simeq}&{{}~\mu_{\varDelta_{X}\times Y}R\mathcal{H}\mathcal{o m}\bigg(q_{23}^{-1}K\otimes q_{3}^{-1}G,q_{1}^{\mathfrak{t}}F\bigg)\;.}\\ \end{aligned}
\]

现在我们通过 $ q_{\pi} $ 获取双方的直像。根据命题 4.3.4，存在自然态射：

\[
\begin{aligned}{{R q}_{\pi*}{R^{i j}}_{*}^{\prime}\mu{\mathcal{h}}o m(K,{R\mathcal{H}o m}(q_{2}^{-1}G,q_{1}^{!}F))}&{{}\leftarrow\mu_{\varDelta_{X}}{R q}_{*}R\mathcal{H}o m\bigg(q_{23}^{-1}{K\otimes^{L}q_{3}^{-1}G},q_{1}^{!}F\bigg)}\\ {}&{{}\simeq\mu_{\varDelta_{X}}{R\mathcal{H}o m}(q_{2}^{-1}\varPhi_{K}(G),q_{1}^{!}F)}\\ {}&{{}=\mu{\mathcal{h}}o m(\varPhi_{K}(G),F),}\\ \end{aligned}
\]

如果 $G$  具有紧支持，则该态射是同构。自 $q_{\pi} \circ^t j' = p_1$ 以来，我们可以总结如下。

存在自然态射：

\[
\mu\mathcal{h}o m(\varPhi_{K}(G),F)\to R p_{1*}\mu\mathcal{h}o m(K,R\mathcal{H}o m(q_{2}^{-1}G,q_{1}^{!}F))
\]

并且只要 G 具有紧支持，该态射就是同构。

为了结束(i)的证明，我们注意到对于 $\Omega_X$ 的每个紧子集 $A$ ，都存在 $Y$ 的紧子集 $B$ ，使得 $\text{supp}(G) \cap B = \varnothing$ 意味着以下条件：

（一） $ \operatorname{supp}(Rp_{1*}\mu_{\mathrm{hom}}(K, R\mathcal{H}_{\mathrm{om}}(q_{2}^{-1}G, q_{1}^{!}F)))\cap A=\emptyset $ ，

（b） $ \text{supp}(\mu \text{hom}(\Phi_K(G), F)) \cap A = \varnothing $ 。

事实上，由于 $K$  属于 $\mathbb{N}(\Omega_X, \Omega_Y)$  ，(a) 是清楚的，并且(b) 由命题7.1.2 (iii) 得出。为了证明(7.1.12)是同构，需要用 $G'$ 替换 $G$ ，使得 $G'|_{B} = G|_{B}$ 和 $\text{supp}(G')$ 是紧的。

(ii) 证明类似。 ⑨

在下一节中，我们将给出充分的条件，以便  $ \varPhi_{K} $  和  $ \varPsi_{K} $  是范畴的等价。

\subsection{层的接触变换}
假设给定一个闭圆锥子集 $ A \subset \Omega_X \times \Omega_Y^a $ ，其中 $ \Omega_X $ 和 $ \Omega_Y $ 分别是 $ T^*X $ 和 $ T^*Y $ 的开子集，如§1中所示。我们假设：

\[
p_{1}|_{\cal A}\colon{\cal A}\to\Omega_{X}\quad\mathrm{a n d}\quad p_{2}^{a}|_{\cal A}\colon{\cal A}\to\Omega_{Y}\qquad\mathrm{a r e~h o m e o m o r p h i s m s}.
\]

让我们用  $\chi$  表示从  $\Omega_Y$  到  $\Omega_X$  的映射  $p_1|_{\mathcal{A}} \circ (p_2^a|_{\mathcal{A}})^{-1}$  。如果  $A$  是平滑且拉格朗日的，并且如果  $p_j$  是微分同胚，则  $\chi$  是接触变换。

\statementlabel{定理 7.2.1}让 $ K \in \mathrm{Ob}(\mathbb{D}^{b}(X \times Y)) $  。假设 (7.2.1) 并且：

(7.2.2) K 是上同调可构造的。

\[
(p_{1}^{-1}(\Omega_{X})\cup p_{2}^{a-1}(\Omega_{Y}))\cap\mathrm{S S}(K)\subset\varLambda,
\]

(7.2.4) 自然态射  $ A_A \to \mu \text{hom}(K, K)|_A $  是  $ \mathbf{D}^b(A) $  中的同构。

那么 $ \Phi_K: \mathbf{D}^b(Y; \Omega_Y) \to \mathbf{D}^b(X; \Omega_X) $  和 $ \Psi_K: \mathbf{D}^b(X; \Omega_X) \to \mathbf{D}^b(Y; \Omega_Y) $  是范畴的等价，彼此相反。

此外，如果  $G_{1}$  和  $G_{2}$  属于  $\mathbf{D}^{b}(Y;\Omega_{Y})$  ，则  $\mathbf{D}^{b}(\Omega_{X})$  中存在自然同构：

\[
\chi_{*}\mu\mathcal{h}o m(G_{2},G_{1})\simeq\mu\mathcal{h}o m(\varPhi_{K}(G_{2}),\varPhi_{K}(G_{1})).
\]

\statementlabel{证明}根据假设(7.2.1)和(7.2.3)， $K$ 属于 $\mathbf{N}(\Omega_X, \Omega_Y)$ ， $r_* K$ 属于 $\mathbf{N}(\Omega_Y^a, \Omega_X^a)$ ， $r$ 表示从 $X \times Y$ 到 $Y \times X$ 的映射(7.1.7)。考虑笛卡尔平方：

\[
\begin{array}{c}X\times Y\quad\underrightarrow{\subset\subset\underleftarrow{\tilde{j}}}\quad X\times Y\times Y\\\downarrow\quad\scriptstyle{q_{2}}\quad\Box\quad\downarrow\quad q_{23}\\Y\quad\underrightarrow{\subset\underleftarrow{\tilde{j}}}\quad Y\times Y\quad,\end{array}
\]

其中  $j$  和  $\tilde{j}$  是对角嵌入。设置 $E = R\mathcal{H}om(q_{12}^{-1}K, q_{13}^{\prime}K)$ 。这是  $\mathbf{D}^{b}(X \times Y \times Y)$  的对象。根据命题7.1.9，设置 $K^{*} = r^{-1}R\mathcal{H}om(K, \omega_{X \times Y/X})$ ，我们在 $\mathbf{N}(\Omega_{Y}, \Omega_{Y})$ 中有 $K^{*}\circ K \simeq Rq_{23*}E$ 。

另一方面，我们有 $ \tilde{j}^i E \simeq R \mathcal{H} o m(\tilde{j}^{-1} q_{12}^{-1} K, \tilde{j}^i q_{13}^1 K) \simeq R \mathcal{H} o m(K, K) $ 。因此我们得到规范态射：

\[
A_{x\times y}\to R\mathcal{H}o m(K,K)\to\tilde{\mathfrak{j}}^{!}E\enspace,
\]

这导致：

\[
A_{Y}\to R q_{2*}A_{X\times Y}\to R q_{2*}\tilde{j}^{!}E\simeq j^{!}R q_{23*}E.
\]

这样我们就得到了态射：

\[
\alpha:A_{_{\varDelta_{Y}}}\to K^{*}\circ K\qquad\mathrm{~i n~}\qquad\mathbb{N}(\varOmega_{Y},\varOmega_{Y})~.
\]

我们将证明 $\alpha$  是同构。令  $Z$  为另一个流形，令  $F \in \mathrm{Ob}(\mathbb{D}^b(X \times Z; \Omega_X \times T^*Z))$  、  $G \in \mathrm{Ob}(\mathbb{D}^b(Y \times Z; \Omega_Y \times T^*Z))$  。令 $i_Z(K)$  表示在(7.1.8) 中构造的 $\mathbb{N}(\Omega_X \times T^*Z; \Omega_Y \times T^*Z)$  的对象。根据命题 7.1.11，我们在  $\mathbb{D}^b(\Omega_X \times T^*Z)$  中有一个自然的同构：

\[
\chi_{*}(\mu{\mathcal{h}}o m(G,\Psi_{i_{Z}(\mathbb{K})}(F))|_{\Omega_{Y}\times T^{\bullet}Z})\simeq\mu{\mathcal{h}}o m(\varPhi_{i_{Z}(\mathbb{K})}(G),F)|_{\Omega_{X}\times T^{\bullet}Z}.
\]

选择  $Z = Y$  、  $G = A_{\Delta_Y}$  和  $F = K$  。我们得到同构：

\[
\chi_{*}\mu{\mathcal{h}}o m(A_{_{\varDelta_{Y}}},\Psi_{i_{Y}(K)}(K))\simeq\mu{\mathcal{h}}o m(K,K)\enspace.
\]

从 $ \Psi_{i_Y(K)}(K) \simeq K^* \circ K $ 命题7.1.9开始，我们得到了 $ \Omega_Y \subset T_{\Delta_Y}^*(Y \times Y) $ 上的同构：

\[
\mu_{\varDelta}(\alpha)\colon\mu_{\varDelta_{\Upsilon}}(A_{\varDelta_{\Upsilon}})\simeq\mu_{\varDelta_{\Upsilon}}(K^{*}\circ K)\enspace.
\]

由于  $ \mathrm{SS}(K^* \circ K) \cap (\Omega_Y \times T^* Y) \subset T_{\mathcal{A}_Y}^*(Y \times Y) $  ，命题 6.6.1 和同构 (7.2.8) 意味着  $ \alpha $  是  $ \mathbf{N}(\Omega_Y, \Omega_Y) $  中的同构。类似地证明  $ \mathbf{N}(\Omega_X, \Omega_X) $  中的同构  $ A_{\Delta_X} \simeq K \circ K^{*\prime} $  其中  $ K^{*\prime} = r^{-1} R \mathcal{H} o m(K, \omega_{X \times Y/Y}) $  。根据命题 7.1.10，它们暗示 $ \Phi_K : \mathbf{D}^b(Y; \Omega_Y) \to \mathbf{D}^b(X; \Omega_X) $  是范畴的等价物。由于  $ \Psi_K $  是  $ \Phi_K $  的伴随函子，因此  $ \Psi_K $  是  $ \Phi_K $  的拟逆。然后 (7.2.5) 从 (7.2.7) 得出  $ Z = \{ \mathrm{pt}\} $  。   $ \square $ 

我们将证明，如果  $\chi$  是  $\Omega_Y$  和  $\Omega_X$  之间的接触变换，则在缩小  $\Omega_Y$  和  $\Omega_X$  后，可以使用定理 7.2.1 构造  $\mathbf{D}^b(Y; \Omega_Y)$  和  $\mathbf{D}^b(X; \Omega_X)$  之间范畴的等价性。

令 $\Omega_{X}$ 和 $\Omega_{Y}$ 分别为 $T^{*}X$ 和 $T^{*}Y$ 的两个开子集， $\chi: \Omega_{Y} \to \Omega_{X}$ 为接触变换。我们设定

\[
\Lambda=\left\{(x,y;\xi,\eta)\in\Omega_{X}\times\Omega_{Y}^{a};(x;\xi)=\chi(y;-\eta)\right\}.
\]

这是一个圆锥拉格朗日流形，封闭在  $\Omega_X \times \Omega_Y^a$  中。让 $p_Y \in \Omega_Y$ ， $p_X = \chi(p_Y) \in \Omega_X$ 。

\statementlabel{推论 7.2.2}存在开邻域  $ \pi(p_X) $  的  $ X' $  、  $ \pi(p_Y) $  的  $ Y' $  、  $ p_X $  的  $ \Omega_X' $  、  $ p_Y $  的  $ \Omega_Y' $  和  $ \Omega_X' \subset T^*X' \cap \Omega_X $  、  $ \Omega_Y' \subset T^*Y' \cap \Omega_Y $  ，并且存在  $ K \in \text{Ob}(\mathbb{D}^b(X') \times Y') $  )，这样：

(a)  $ \chi $  引发接触变换  $ \Omega_{Y}^{\prime} \simeq \Omega_{X}^{\prime} $  ，

(c)  $ \Phi_K: \mathbb{D}^b(Y'; \Omega_Y') \to \mathbb{D}^b(X'; \Omega_X') $  是范畴的等价物，

(d) 对于属于 $ \mathbf{D}^b(Y'; \Omega_Y) $  的 $ G_1 $  和 $ G_2 $ ，我们在 $ \mathbf{D}^b(\Omega_X') $  中有同构(7.2.5)。

\statementlabel{证明}根据推论A.2.7，在收缩 $\Omega_Y$ 和 $\Omega_X$ 之后，我们可以将 $\chi$ 分解为 $\chi_2 \circ \chi_1$ ，其中每个 $\chi_i$  ( $i = 1,2$ )是接触变换，并且通过(7.2.9)与 $\chi_i$ 关联的拉格朗日流形 $A_i$ 是超曲面的共形丛。由命题7.1.2和7.1.6可知，如果 $K_1 \in \text{Ob}(\mathbb{D}^b(X' \times Z))$ 和 $K_2 \in \text{Ob}(\mathbb{D}^b(Z \times Y'))$ 满足推论的条件(b)、(c)、(d)，则 $K_2 \circ K_1$ 也将满足这些条件。因此，我们可以从一开始就假设存在一个超曲面 $S \subset X \times Y$ ，使得(7.2.9)中定义的拉格朗日流形 $A$ 包含在 $T_S^*(X \times Y)$ 中。由于 $A$  是 $\mathbb{R}^+-$  圆锥曲线，并且 $S$  是超曲面，因此分别存在 $p_X, p_Y, \pi(p_X)$  、 $\pi(p_Y)$  的开邻域 $\Omega_X', \Omega_Y', X', Y'$  ，使得 $\left((\Omega_X' \times T^*Y') \cup (T^*X' \times \Omega_Y')\right) \cap T_S^*(X \times Y) \subset A$  。

那么 $ \mathbf{D}^{b}(X' \times Y') $  中的 $ K = A_{S \cap (X' \times Y')} $  满足定理7.2.1 的所有假设。   $ \square $ 

\statementlabel{定义 7.2.3}在定理7.2.1的情况下，我们说 $ \Phi_{\kappa} $ 和 $ \Psi_{\kappa} $ 是 $ \chi $ 之上的扩展接触变换。

我们将证明，恒等之上的所有扩展接触变换都来自  $\mathbf{D}^{b}(\mathfrak{Mod}(A))$  中范畴的等价性。从这个意义上说， $\Phi_{K}$ 本质上是独一无二的。假设  $\chi$  是  $p \in T^*X$  邻域中的恒等式。根据命题 6.6.1，我们发现  $K \simeq M_{\Delta_X}$  位于  $\mathbf{D}^{b}(X \times X; (p, p^a))$  和  $M \in \mathrm{Ob}(\mathbf{D}^{b}(\mathfrak{Mod}(A))$  中。如果  $G \in \mathrm{Ob}(\mathbf{D}^{b}(X))$  ，则：

\[
\varPhi_{K}(G)=M_{X}\mathop{{\otimes}}^{L}G\qquad\mathrm{i n}\qquad\mathbb{D}^{b}(X;p)~.
\]

\statementlabel{命题 7.2.4}假设函子  $ M_X \otimes^L \cdot \text{defines an equivalence of categories in } \mathbf{D}^b(X; p) $  。然后是函子  $ M \otimes^L \cdot \text{defines an equivalence of categories in } \mathbf{D}^b(\mathfrak{Mod}(A)) $  。

\statementlabel{证明}令  $Y$  为  $X$  的子流形，使得  $p \in T_Y^* X$  。设置 $A = T_Y^* X$ 。然后 $M_X \otimes^L \cdot$  归纳出 $\mathbf{D}_A^b(X; p)$  上范畴的等价性，并且最后一个范畴相当于命题6.6.1 的 $\mathbf{D}^b(\mathfrak{Mod}(A))$ 。   $\square$ 

示例 7.2.5。令 $X$  和 $Y$  为 $\mathbb{R}^n$  的两个副本，分别赋予 $(x)$  和 $(y)$  线性坐标系。让  $(x; \xi)$  和  $(y; \eta)$  表示余切丛上的关联坐标，并考虑接触变换  $\chi: \dot{T}^* Y \simeq \dot{T}^* X$  ：

\[
\chi:(y;\eta)\mapsto(x;\xi)=(y+\eta/|\eta|;\eta)
\]

我们设置了  $|\eta| = (\sum_j \eta_j^2)^{1/2}$  。让 $S = \{(x, y) \in X \times Y; \sum_j (x_j - y_j)^2 \geq 1\}$  ，并让 $A = \mathrm{SS}(A_S) \cap \dot{T}^*(X \times Y)$  。然后：

\[
\Lambda=\left\{(x,y;\xi,\eta);\sum_{j}(x_{j}-y_{j})^{2}=1,\xi=-\eta=\lambda(x-y),\lambda>0\right\}
\]

因此：

\[
(x;\xi)=\chi(y;\eta)\Leftrightarrow(x,y;\xi,-\eta)\in A.
\]

让 $ K = A_S $  。那么定理7.2.1的所有条件都满足， $ \Phi_K: \mathbf{D}^b(Y; \dot{T}^*Y) \to \mathbf{D}^b(X; \dot{T}^*X) $ 是范畴的等价。

如果 $G \in \mathrm{Ob}(\mathbb{D}^b(Y))$ ，我们有 $\Phi_K(G) \simeq R q_{1!}(q_2^{-1}G)_{s}$ 。特别是我们在  $\mathbb{D}^b(X; \dot{T}^*X)$  中找到了  $\Phi_K(A_{\{0\}}) \simeq A_{\{\sum_j x_j^2 \geq 1\}}, \Phi_K(A_Y) \simeq 0, \Phi_K(A_{\{y \neq 0\}}) \simeq A_{\{\sum_j x_j^2 \geq 1\}}[-1]$  。

示例 7.2.6。让我们考虑一个与 7.2.5 类似的关于复流形的例子。

令 $X$ 和 $Y$ 是 $\mathbb{C}^n$ 的两个副本，分别赋予 $\mathbb{C}$ -线性坐标 $(z)$ 和 $(w)$ ，以及 $z = x + \sqrt{-1}y$ 、 $w = u + \sqrt{-1}v$ 。让 $(z; \zeta)$  和 $(w; \theta)$ 

表示复余切丛上的相关坐标。如果  $X^{\mathrm{R}}$  和  $Y^{\mathrm{R}}$  表示  $X$  和  $Y$  的实数底层流形，则  $T^*X^{\mathrm{R}}$  和  $T^*Y^{\mathrm{R}}$  上的规范 1-形式由  $2\operatorname{Re}(\sum_j\zeta_j dz_j)$  和  $2\operatorname{Re}(\sum_j\theta_j dw_j)$  给出。让 $\Omega_X = \{(z;\zeta); \zeta^2 \notin \mathbb{R}^+ \cup \{0\}\}$ ，我们设置 $\zeta^2 = \sum_j \zeta_j^2$ ，并类似地让 $\Omega_Y = \{(w;\theta); \theta^2 \notin \mathbb{R}^+ \cup \{0\}\}$ 。那么  $(-\theta^2)^{1/2}$  是  $\Omega_Y$  上的全纯函数，由  $\operatorname{Re}((-\theta^2)^{1/2}) > 0$  明确定义，我们可以考虑从  $\Omega_Y$  到  $\Omega_X$  的全纯接触变换  $\chi$  给出：

\[
\chi:(w;\theta)\mapsto(z;\zeta)=(w+\theta/(-\theta^{2})^{1/2};\theta)\;.
\]

让 $ Z = \{(z, w); (z - w)^2 = -1\} $  。考虑拉格朗日流形：

\[
{\cal A}^{\pm}=\{(z,w;\zeta,\theta);\zeta^{2}\notin\mathbb{R}^{+}\cup\{0\},\zeta=-\theta,z=w\pm\zeta/(-\zeta^{2})^{1/2}\}.
\]

然后 $(z;\zeta) = \chi(w;\theta) \Leftrightarrow (z,w;\zeta,-\theta) \in \Lambda^+$  和 $T_Z^*(X \times Y) \cap (\Omega_X \times \Omega_Y^a) = A^+ \sqcup A^-$  。设置  $Z_+ = \{(z,w);\mathrm{Im}((z-w)^2) = 0$  、  $\mathrm{Re}((z-w)^2) < -1\}$  并让  $K = A_{Z_+}$  。然后：

\[
\begin{aligned}{\SS(K)}&{{}=\left\{(z,w;\zeta,\theta);(z,w)\in\widetilde{Z}_{+},\zeta=-\theta=k(z-w),k\in\mathbb{C},\mathbf{R e}k\geqslant0,\right.}\\ {}&{{}\quad\left.(1+\mathbf{R e}((z-w)^{2}))\cdot\mathbf{R e}k=0\right\}\;.}\\ \end{aligned}
\]

这给出：

\[
\mathrm{S S}(K)\cap(\Omega_{X}\times T^{*}Y)=\mathrm{S S}(K)\cap(T^{*}X\times\Omega_{Y}^{a})=\varLambda^{+}\enspace.
\]

（请注意  $ \mathrm{SS}(A_Z) \cap (\Omega_X \times T^*Y) = \mathrm{SS}(A_Z) \cap (T^*X \times \Omega_Y^a) = \varLambda^+ \sqcup \varLambda^- $  和 K 和  $ A_Z[-1] $  在  $ \mathbb{D}^b(X \times Y; \varLambda_+) $  中是同构的。）

因此  $\Phi_k$  定义了范畴  $\mathbf{D}^b(Y; \Omega_Y) \xrightarrow{\sim} \mathbf{D}^b(X; \Omega_X)$  的等价。

定义实数子流形  $ N = \{w \in Y; \text{Im } w = 0\} $  和  $ M = \{z \in X; (\text{Im } z)^2 = 1\} $  。让我们计算 $ \Phi_K(A_N) $  。

对于  $ z = x + \sqrt{-1y} $  、  $ (x \in \mathbb{R}^n, y \in \mathbb{R}^n) $  ，我们有  $ \Phi_K(A_N)_z = R\Gamma_c(S_z; A_Y) $  ，其中：

\[
S_{z}=\left\{w\in N;(z,w)\in Z_{+}\right\}=\left\{w\in\mathbb{R}^{n};\left\langle x-w,y\right\rangle=0,(x-w)^{2}<y^{2}-1\right\}\;.
\]

因此 $ y^2 \leq 1 $  和 $ S_z $  的 $ S_z = \varnothing $  与 $ y^2 > 1 $  的 $ (n-1) $  维开球同胚。因此我们得到：

\[
\varPhi_{K}(A_{N})\simeq A_{\{y^{2}>1\}}[1-n]\enspace.
\]

\subsection{核的微局部复合}
令 $X$ 、 $Y$ 和 $Z$ 为三个流形，并让 $p_X$ 、 $p_Y$ 和 $p_Z$ 分别为 $T^*X$ 、 $T^*Y$ 和 $T^*Z$ 的一个点。我们设置  $x_\circ = \pi_X(p_X)$  、  $y_0 = \pi_Y(p_Y)$  和  $z_\circ = \pi_Z(p_Z)$  。我们保留与§1中相同的符号  $p_{ij}$  、  $q_{ij}$  、  $p_{ij}^a$  等。

对于  $\mathbf{D}^b(X \times Y; (p_X, p_Y^a))$  的对象  $K_1$  和  $\mathbf{D}^b(Y \times Z; (p_Y, p_Z^a))$  的对象  $K_2$  的一对  $(K_1, K_2)$  ，如果满足以下条件，我们说  $(K_1, K_2)$  是微局部可组合的（在  $(p_X, p_Y, p_Z)$  处）

\[
\left\{\begin{aligned}{}&{{}\mathsf{S S}(K_{1})\times\mathsf{S S}(K_{2})\cap p_{13}^{a-1}((p_{X},p_{Z}^{a}))\subset\left\{((p_{X},p_{Y}^{a}),(p_{Y},p_{Z}^{a})\right\}}\\ {}&{{}\begin{aligned}{}&{{}{t}*\gamma}\\ {}&{{}{o n~a~n e i g h b o r h o o d~o f}((p_{X},p_{Y}^{a}),(p_{Y},p_{Z}^{a}))\quad.}\\ \end{aligned}}\\ \end{aligned}\right.
\]

请注意，如果 $F$ 属于 $\mathbf{D}^{b}(X; p)$ ，则 $\mathrm{SS}(F)$ 在 $p$ 的萌芽是明确定义的。通过这一评论，我们可以看出（7.3.1）是有道理的。

\statementlabel{命题 7.3.1}令  $ (K_1, K_2) \in \mathrm{Ob}(\mathbb{D}^b(X \times Y)) \times \mathrm{Ob}(\mathbb{D}^b(Y \times Z)) $  为  $ (p_X, p_Y, p_Z) $  处的微局部可组合对。

(i) “  $\lim_{x \to 0} K_1' \circ K_2'$  属于  $\mathbf{D}^b(X \times Z; (p_X, p_Z^a))$  ，其中  $K_1' \to K_1$  范围位于  $(p_X, p_Y^a)$  处的同构范畴， $K_2' \to K_2$  范围位于  $(p_Y, p_Z^a)$  处的同构范畴。

此外，对于 $T^{*}X \times T^{*}Y \times T^{*}Z$  中 $(p_{X}, p_{Y}, p_{Z}^{a})$  的任何邻域 $W$ ，我们有：

\[
\left\{\begin{aligned}{}&{{}\SS(\varinjlim K_{1}^{\prime}\circ K_{2}^{\prime})\subset q_{13}^{a}\left(W\cap\left(\SS(K_{1})\times\SS(K_{2})\right)\right)}\\ {}&{{}{o n~a~n e i g h b o r h o o d~o f}\left(p_{X},p_{Z}^{a}\right).}\\ \end{aligned}\right.
\]

(ii) 如果 $ (K_{1}, K_{2}) $ 还满足以下条件：

\[
\left(\mathrm{S S}(K_{1})\times\mathrm{S S}(K_{2})\right)\cap\left\{p_{X}\right\}\times T_{y_{\circ}}^{*}Y\times\left\{p_{Z}^{a}\right\}\subset\left\{(p_{X},p_{Y},p_{Z}^{a})\right\}\mathrm{,}
\]

\[
\left(\mathtt{S S}(K_{1})\mathop{\times}\limits_{{T^{*}\it{Y}}}\mathtt{S S}(K_{2})\right)\cap\left\{(x_{\circ};0)\right\}\times\dot{T}_{y_{\circ}}^{*}Y\times\left\{(z_{\circ};0)\right\}=\emptyset\mathrm{~,~}
\]

那么“  $\lim\limits_{\substack{\leftarrow\\ V}}\left(K_{1}\right)_{X\times V}\circ K_{2}$  属于 $\mathbf{D}^{b}(X\times Z;\left(p_{X},p_{Z}^{a}\right))$  并且同构于“  $\lim\limits_{V}\left(K_{1}^{\prime}\right)\circ K_{2}^{\prime}$  。在这里， $V$ 覆盖了 $y_{\circ}$ 的开放邻里系统。

(iii) 存在态射 $ K'_1 \to K_1 $ 和 $ K'_2 \to K_2 $ ，使得它们分别在 $ (p_X, p_Y^a) $ 和 $ (p_Y, p_Z^a) $ 处同构，并且 $ (K'_1, K'_2) $ 满足(7.3.3)和(7.3.4)。

(iv) 如果 $ p_X \in \dot{T}^*X $ ，则“ $ \lim\limits_{\leftarrow} $ ” $ K'_1 \circ K_2 $ 属于 $ \mathbf{D}^b(X \times Z $ ；   $ (p_X, p_Z^a) $  ) 并且它与“  $ \lim\limits_{\leftarrow} $  ”  $ K'_1 \circ K'_2 $  同构。

\statementlabel{定义 7.3.2}对于一对 $(K_1, K_2) \in \mathrm{Ob}(\mathbb{D}^b(X \times Y; (p_X, p_Y^a)) \times \mathrm{Ob}(\mathbb{D}^b(Y \times Z; (p_Y, p_Z^a)))$  ，前面命题中给出的 $\mathbb{D}^b(X \times Z; (p_X, p_Z^a))$  的原对象“ $\lim K'_1 \circ K'_2$  ”用 $K_1 \circ_\mu K_2$  表示，称为 $K_1$  和 $K_2$  的微局部组合。

为了证明命题 7.3.1，让我们从以下引理开始。

\statementlabel{引理 7.3.3}符号如命题 7.3.1 所示，假设  $ p_X \in \dot{T}^*X $  。令  $ S $  为  $ T^*(Y \times Z) $  的闭圆锥曲线子集，满足：

\[
\begin{array}{r}{\mathrm{S S}(K_{1})\underset{T^{*}Y}\times S\subset\left\{(p_{X},p_{Y},p_{Z}^{a})\right\}{~o n~a~n e i g h b o r h o o d~o f~}(p_{X},p_{Y},p_{Z}^{a})\enspace.}\end{array}
\]

那么 $ \mathbb{D}^b(X \times Y) $ 中存在态射 $ \varphi: K_1' \to K_1 $ ，使得 $ \varphi $ 是 $ (p_X, p_Z^a) $ 和 $ K_1' $ 的同构满足以下条件：

\[
\begin{array}{r}{\mathrm{S S}(K_{1}^{\prime})\underset{T^{*}Y}{\times}\mathrm{\boldmath~S~}\cap(\{p_{X}\}\times T_{y\circ}^{*}Y\times\{p_{Z}^{a}\})\subset\{(p_{X},p_{Y},p_{Z}^{a})\}\mathrm{~,~}}\end{array}
\]

\[
\mathsf{S S}(K_{1}^{\prime})\cap((T_{X}^{\ast}X)_{x_{\circ}}\times\dot{T}_{y_{\circ}}^{\ast}Y)=\emptyset\enspace.
\]

\statementlabel{证明}让我们在  $T_y^* Y$  中采用  $p_Y$  的紧凑邻域  $L$  ，这样  $(\mathrm{SS}(K_1) \times_{T^* Y} S) \cap (\{p_X\} \times L \times \{p_Z^a\}) \subset \{(p_X, p_Y, p_Z^a)\}$  。令 $G$  为 $p_1: T^*(Y \times Z) \to T^* Y$  的 $S \cap (T_y^* Y \times \{p_Z^a\})$  的图像。然后我们有：

\[
\left(\left\{p_{x}\right\}\times(G\cap L)^{a}\right)\cap\mathsf{S S}(K_{1})\subset\left\{\left(p_{x},p_{Y}^{a}\right)\right\}\;.
\]

在 $ T_{(x_0, y_0)}^* (X \times Y) $ 中取一个适当的闭合凸锥 $ \gamma $ ，这样我们有：

\[
\left(\left\{p_{x}\right\}\times T_{Y_{\circ}}^{*}Y\right)\cap\mathrm{I n t}\gamma\subset\left\{p_{x}\right\}\times L^{a}\quad\mathrm{a n d}\quad(p_{x},p_{Y}^{a})\in\mathrm{I n t}\gamma,
\]

\[
\gamma\cap(\{(x_{\circ};0)\}\times\dot{T}_{y_{\circ}}^{*}Y)=\emptyset.
\]

这里  $ (x_{0};0) $  表示  $ T_{x_{0}}^{*}X $  的起源。

然后，取一个开圆锥曲线集 U，使得

\[
\mathrm{Int}_{\gamma}\supset U\ni(p_{x},p_{Y}^{a})\enspace.
\]

现在，应用精炼的微局域截止引理（命题 6.1.4），我们获得态射  $ \varphi: K'_1 \to K_1 $  使得：

 $ \varphi $  是  $ U $  的同构

\[
\mathrm{SS}(K_{1}^{\prime})\subset U\cup(\gamma\backslash(\{p_{X}\}\times G^{a}))~.
\]

那么 $ \varphi $  满足所需的属性。 □

7.3.1 命题证明首先，我们将证明（7.3.1）、（7.3.3）和（7.3.4）意味着：

\[
\left\{\begin{aligned}{}&{{}\bigg(\mathtt{S S}(K_{1})\times\mathtt{S S}(K_{2})\bigg)\cap\big(\{p_{X}\}\times T^{*}Y\times\{p_{Z}^{a}\}\big)\subset\big\{(p_{X},p_{Y},p_{Z}^{a})\big\}}\\ {}&{{}\notag\mathrm{o n~a~n e i g h b o r h o o d~o f}(x_{\circ},y_{\circ},z_{\circ})~.}\\ \end{aligned}\right.
\]

如果 (7.3.13) 为假，则存在序列  $ \{(y_n; \eta_n)\} \subset T^* Y \setminus \{p_Y\} $  使得  $ y_n \to y $  、  $ (p_X, (y_n; -\eta_n)) \in \mathrm{SS}(K_1) $  、  $ ((y_n; \eta_n), p_Z^a) \in \mathrm{SS}(K_2) $  。如果  $ \{\eta_n\} $  有界，则

\[
\{(y_n;\eta_n)\}\text{converges to } p_Y\text{ by } (7.3.3), \text{which contradicts } (7.3.1). \text{If } \{\eta_n\}\text{ is unbounded, we may assume that } \eta_n/|\eta_n|\text{ converges to } \eta \neq 0. \text{ Then } ((x_\circ;0),(y_0;-\eta)) \in \mathrm{SS}(K_1) \text{ and } ((y_0;\eta),(z_\circ;0)) \in \mathrm{SS}(K_2), \text{which contradicts } (7.3.4).
\]

现在我们将分几步证明这个命题。

(a) 首先我们假设  $ p_X \in T^*X $  然后 (iii) 由引理 7.3.3 得出。此外，为了证明(i)、(ii)、(iv)，我们可以使用相同的引理假设 $ (K_1, K_2) $ 满足(7.3.3)和(7.3.4)。

然后，根据 (7.3.4)， $ \mathrm{SS}(q_{12}^{-1}K_1) \cap \mathrm{SS}(q_{23}^{-1}K_2)^a \subset T_{X \times Y \times Z}^*(X \times Y \times Z) $  位于  $ (x_\circ, y_0, z_\circ) $  的邻域上。因此，根据命题 5.4.14， $ G = q_{12}^{-1}K_1 \otimes^L q_{23}^{-1}K_2 $  满足：

\[
\left\{\begin{aligned}{\SS(G)\subset S=\SS(K_{1})\times T_{Z}^{\ast}Z+T_{X}^{\ast}X\times\SS(K_{2})}\\ {\mathrm{o n~a~n e i g h b o r h o o d~o f}(x_{\circ},y_{\circ},z_{\circ})~.}\\ \end{aligned}\right.
\]

由于  $ (t'q'_{13})^{-1}(S) $  与  $ \mathrm{SS}(K_1) \times_{T^*Y} \mathrm{SS}(K_2) $  到  $ T^*X \times Y \times T^*Z $  的图像同构，所以 (7.3.13) 和 (7.3.14) 意味着：

\[
\left\{\begin{aligned}{}&{{}({}^{t}q_{13}^{\prime})^{-1}(\mathbf{S S}(G))\cap q_{13\pi}^{-1}((p_{X},p_{Z}^{a}))\subset\{(p_{X},y_{\circ},p_{Z}^{a})\}}\\ {}&{{}\operatorname{o n a n e i g h b o r h o o d o f}(p_{X},y_{\circ},p_{Z}^{a})~.}\\ \end{aligned}\right.
\]

因此，我们可以应用命题6.1.10得出结论：“ $\lim_{\substack{\leftarrow \\ V}} Rf_*(G_{X \times V \times Y}) = \lim_{\substack{\leftarrow \\ V}} (K_1)_{X \times V} \circ K_2$ 属于 $\mathbf{D}^b(X \times Z; (p_X, p_Z^a))$ 。现在让 $K'_1 \to K_1$ 和 $K'_2 \to K_2$ 分别在 $(p_X, p_Y^a)$ 和 $(p_Y, p_Z^a)$ 处同构。在 $(p_X, p_Y^a)$ 处存在同构 $K'_1 \to K'_1$ ，使得 $(K'_1, K'_2)$ 和 $(K'_1, K_2)$  满足条件 (7.3.3), (7.3.4)，则根据命题 5.4.14， $q_{12}^{-1}K'_1 \otimes^L q_{23}^{-1}K'_2 \to q_{12}^{-1}K'_1 \otimes^L q_{23}^{-1}K_2 \to q_{12}^{-1}K_1 \otimes^L q_{23}^{-1}K_2$  是  $(p_X, y_0, p_Z^a)$  的同构。 因此，根据命题 6.1.10，我们有：

\[
\varprojlim_{V}\left(K_{1}^{\prime\prime}\right)_{X\times V}\circ K_{2}^{\prime}\xrightarrow{\sim}\varprojlim_{V}\left(K_{1}\right)_{X\times V}\circ K_{2}\enspace.
\]

对  $ K_{1} $  取“射影极限”，我们得到：

\[
\varprojlim_{K_{1}^{\prime}}K_{1}^{\prime}\circ K_{2}^{\prime}\xrightarrow{\quad}\varprojlim_{V}\left(K_{1}\right)_{X\times V}\circ K_{2}\enspace.
\]

设置 $ K'_2 = K_2 $ 我们得到：

\[
\varprojlim_{K_{1}^{\prime}}K_{1}^{\prime}\circ K_{2}\xrightarrow{\quad}\varprojlim_{V}\left(K_{1}\right)_{X\times V}\circ K_{2}\mathrm{~,~}
\]

并取(7.3.17)中 $ K'_2 $ 的“射影极限”，我们得到 $ \lim_{K'_1, K'_2} K'_1 \circ K'_2 \approx \lim_{V} (K_1)_{X \times V} \circ K_2 $ 。

（b） $ p_Z \in \dot{T}^*Z $ 。证明与(a)类似。

(c) 假设  $ p_X \in T_X^*X $  、  $ p_Z \in T_Z^*Z $  和  $ p_Y \in \dot{T}^*Y $  。由于  $ (p_X, tp_Y, p_Z^a) \notin \mathrm{SS}(K_1) \times_{T^*Y} \mathrm{SS}(K_2) $  对于  $ t > 0 $  ，我们有  $ (p_X, p_Y^a) \notin \mathrm{SS}(K_1) $  或  $ (p_Y, p_Z^a) \notin \mathrm{SS}(K_2) $  ，结果如下。

（d） $(p_X,p_Y,p_Z)\in T_X^*X\times T_Y^*Y\times T_Z^*Z$。由于 $\operatorname{supp}(K_1)\times_Y\operatorname{supp}(K_2)\cap\{x_0\}\times Y\times\{z_0\}\subset\{(x_0,y_0,z_0)\}$ 在 $(x_0,y_0,z_0)$ 的邻域上成立，$(K_1,K_2)$ 满足条件 (7.3.3) 和 (7.3.4)。

\statementlabel{命题 7.3.4}让 $q'_1$ 和 $q'_2$ 分别表示从 $(X \times Z) \times (Y \times Z)$ 到 $X \times Z$ 和 $Y \times Z$ 的投影，让 $i$ 表示对角线嵌入 $X \times Y \times Z \hookrightarrow (X \times Z) \times (Y \times Z)$ 。令  $(K_1, K_2) \in \mathrm{Ob}(\mathbb{D}^b(X \times Y; (p_X, p_Y^a)) \times \mathrm{Ob}(\mathbb{D}^b(Y \times Z; (p_Y, p_Z^a)))$  为  $(p_X, p_Y, p_Z)$  处的微局部可组合对。那么  $(i_*(K_1 \boxtimes A_Z), K_2) \in \mathrm{Ob}(\mathbb{D}^b(X \times Z \times Y \times Z; (p_X, p_Z^a, p_Y^a, p_Z))) \times \mathrm{Ob}(\mathbb{D}^b(Y \times Z \times \{pt\}; (p_Y, p_Z^a, pt)))$  可以在  $(((p_X, p_Z^a), (p_Y, p_Z^a), pt)$  进行微局部组合，并且：

\[
K_{1}\mathop{\circ}\limits_{\mu}K_{2}\simeq i_{*}(K_{1}\boxtimes A_{Z})\mathop{\circ}\limits_{\mu}K_{2}\enspace.
\]

\statementlabel{证明}很简单。

\statementlabel{命题 7.3.5}令  $X$  、  $Y$  、  $Z$  和  $W$  为四个流形：  $p_X \in T^*X$  、  $p_Y \in T^*Y$  、  $p_Z \in T^*Z$  、  $p_W \in T^*W$  。让 $K_1 \in \text{Ob}(\mathbb{D}^b(X \times Y; (p_X, p_Y^a)))$ ， $K_2 \in \text{Ob}(\mathbb{D}^b(Y \times Z; (p_Y, p_Z^a)))$ ， $K_3 \in \text{Ob}(\mathbb{D}^b(Z \times W; (p_Z, p_W^a)))$ 。假设：

\[
\left\{\begin{aligned}{}&{{}\operatorname{S S}(K_{1})\times\operatorname{S S}(K_{2})\times\operatorname{S S}(K_{3})=\big\{(p_{X},p_{Y}^{a}),(p_{Y},p_{Z}^{a}),(p_{Z},p_{W}^{a})\big\}}\\ {}&{{}{~i n~a~n e i g h b o r h o o d~o f~t h i s~p o i n t~}.}\\ \end{aligned}\right.
\]

那么  $(K_{1},K_{2}),(K_{2},K_{3}),(K_{1}\circ_{\mu}K_{2},K_{3})$  和  $(K_{1},K_{2}\circ_{\mu}K_{3})$  是微局部可组合的并且：

\[
\left(K_{1}\circ_{_{\mu}}K_{2}\right)_{{}_{\mu}}\circ K_{3}\simeq K_{1}\circ_{_{\mu}}\left(K_{2}\circ_{_{\mu}}K_{3}\right)\mathrm{i n}\mathbb{D}^{b}(X\times W;(p_{X},p_{W}^{a})).
\]

\statementlabel{证明}很简单。

\statementlabel{命题 7.3.6}设 $X,Y,Z,X',Y',Z'$ 为六个流形，并设 $p_W\in T^*W$（$W=X,Y,Z,X',Y',Z'$）。

让  $K_1 \in \mathrm{Ob}(\mathbb{D}^b(X \times Y; (p_X, p_Y^a)))$  、  $K_2 \in \mathrm{Ob}(\mathbb{D}^b(Y \times Z; (p_Y, p_Z^a)))$  、  $K_1' \in \mathrm{Ob}(\mathbb{D}^b(X' \times Y'; (p_X, p_Y^a)))$  和  $K_2' \in \mathrm{Ob}(\mathbb{D}^b(Y' \times Z'; (p_Y, p_Z^a)))$  。假设  $(K_1, K_2)$  和  $(K_1', K_2')$  是微局部可组合的。那么  $(K_1 \boxtimes^L K_1', K_2 \boxtimes^L K_2')$  是微局部可组合的，我们有：

\[
\bigg(K_{1}\mathop{\boxtimes}^{L}K_{1}^{\prime}\bigg)_{\mu}\circ\bigg(K_{2}\mathop{\boxtimes}^{L}K_{2}^{\prime}\bigg)\simeq\bigg(K_{1}\mathop{\circ}_{\mu}K_{2}\bigg)\mathop{\boxtimes}\bigg(K_{1}\mathop{\circ}_{\mu}K_{2}^{\prime}\bigg).
\]

\statementlabel{证明}很简单。

为了结束本节，我们定义了一类可与任何内核进行微局部组合的内核。

\statementlabel{定义 7.3.7}其中一个用  $\mathbf{N}(X, Y; p_X, p_Y)$  表示  $\mathbf{D}^b(X \times Y; (p_X, p_Y^a))$  的完整三角子范畴，由对象  $K$  组成，使得

\[
\mathrm{SS}(\boldsymbol{K})\cap(\left\{p_{x}\right\}\times T^{*}Y)\subset\left\{(p_{x},p_{Y}^{a})\right\}\mathrm{on~a~neighborhood~of}\left(p_{x},p_{Y}^{a}\right)
\]

然后人们立即看到：

\statementlabel{命题 7.3.8}对于任何  $ K \in \text{Ob}(\mathbb{N}(X, Y; p_X, p_Y)) $  和  $ L \in \text{Ob}(\mathbb{D}^b(Y \times Z; (p_Y, p_Z^a))) $  ，  $ (K, L) $  在  $ (p_X, p_Y, p_Z) $  处是微局部可组合的，并且内核的微局部组合诱导函子：

\[
\begin{aligned}{}&{{}\circ\cdot\colon\mathbb{N}(X,Y;p_{X},p_{Y})\times\mathbb{D}^{b}(Y\times Z;(p_{Y},p_{Z}^{a}))\to\mathbb{D}^{b}(X\times Z;(p_{X},p_{Z}^{a}))~,}\\ {}&{{}\mu}\\ {}&{{}\quad\cdot\circ\cdot\colon\mathbb{N}(X,Y;p_{X},p_{Y})\times\mathbb{N}(Y,Z;p_{Y},p_{Z})\to\mathbb{N}(X,Z;p_{X},p_{Z})~.}\\ {}&{{}\quad\mu}\\ \end{aligned}
\]

\subsection{与子流形相关的层的积分变换}
让我们从微分几何的基本结果开始。令  $f: Y \to X$  为流形的态射， $N$  （分别为  $M$  ）为  $Y$  （分别为  $X$  ）的闭子流形。我们假设  $f$  是光滑的，并且我们将  $Y \times_X T^*X$  识别为  $T^*Y$  的封闭对合子流形。我们设定：

\[
V=Y\mathop{\times}\limits_{x}\dot{T}^{*}X~,\qquad\varLambda=\dot{T}_{N}^{*}Y~.
\]

让 $ p \in A \cap V $  。我们假设：

交集  $A \cap V$  在 p 处是干净的。

回想一下映射  $f'$  和  $f_\pi$  ，分别从  $Y \times_x T^*X$  到  $T^*Y$  和  $T^*X$  。通过（7.4.2），映射 $f_{\pi}|_{A\cap V}:A\cap V \to T^*X$ 具有恒定的等级（参见练习A.5），并且对于 $Wp$ 的足够小的开邻域， $f_\pi(W\cap A\cap V)$ 是拉格朗日流形。我们假设存在 $X$ 的子流形 $M$ ，使得：

\[
f_{\pi}(W\cap\varLambda\cap V)=T_{M}^{\ast}X\mathrm{~i n~a~n e i g h b o r h o o d~o f~}f_{\pi}(p).
\]

我们设置  $ y_{\circ}=\pi_{Y}(p) $  、  $ x_{\circ}=\pi_{X}(f_{\pi}(p)) $  和  $ \lambda_{\circ}(p)=T_{p}\pi_{Y}^{-1}(y_{\circ}) $  。

\statementlabel{命题 7.4.1}在前面的情况（f smooth，（7.4.2），（7.4.3））中，我们还假设 $N$ 是Y的超曲面。那么 $X$ 上的局部坐标系 $(x) = (x_1, x', x'')$ 和Y上的 $(x, t, u)$ 分别存在于 $x_\circ$ 和 $y_0$ 处，使得：  $p = (0; \mathrm{d}x_1)$  ，  $f(x, t, u) = x$  ，  $M = \{x \in X; x_1 = x' = 0\}$  、  $N = \{(x, t, u) \in Y; x_1 = q(t) + \langle x', u \rangle\}$  ，其中  $q(t)$  是二次形式，而  $x_1 \in \mathbb{R}$  、  $x', u \in \mathbb{R}^r$  、  $t \in \mathbb{R}^n$  、  $x'' \in \mathbb{R}^m$  则为某些  $r, n, m \in \mathbb{N}$  。此外，如果将假设（7.4.2）替换为：

交集  $A \cap V$  是横向的  $p$  ，

那么二次形式 $q(t)$  是非简并的。

\statementlabel{证明}$ p $  处的投影  $ A \cap V \to Y $  的 corank 等于：

\[
\dim(\lambda_{\circ}(p)\cap T_{p}(V\cap\varLambda))=\dim(\lambda_{\circ}(p)\cap T_{p}V\cap T_{p}\varLambda)=1\enspace.
\]

由于该公式在  $V \cap \varLambda$  中的  $p$  邻域上也成立，因此  $\pi_Y(A \cap V)$  是  $Y$  的平滑子流形  $L$  。此外 $L \simeq (A \cap V)/{\mathbb{R}^+}$ 。 （请注意  $L$  是  $f|_{N}: N \to X$  的判别轨迹）考虑该图：

\begin{center}
\includegraphics[width=0.37\textwidth]{流行上的层_Chn-assets/img_in_image_box_369_551_816_745.png}
\end{center}

我们设置 $ f_{L} = f|_{L} $ ， $ \tilde{f}_{L} = f_{\pi}|_{A\cap V} $ 。

我们选择 $X$ 上的坐标 $(x_1, x', x')$ ，使得 $M = \{x_1 = x' = 0\}$ 、 $x_\circ = 0$ 、 $f_\pi(p) = dx_1$ 。令  $(x_1, x', x''; \xi_1, \xi', \xi'')$  为  $\dot{T}^*X$  上的关联坐标，以及  $\xi_1 \neq 0$  。由于  $L \simeq (A \cap V) / \mathbb{R}^+$  和  $\tilde{f}_L$  是平滑的，因此  $u = \xi' / \xi_1 \circ \tilde{f}_L$  和  $x'' \circ f_L$  定义了  $L$  上的坐标。我们写 $x''$ 而不是 $x'' \circ f_L$ 。然后我们通过 $L$ 上的坐标系 $(x", u, t'')$ 来完成 $L$ 上的坐标 $(x", u)$ ，并将其扩展到 $Y$ ； （请注意，如果交集  $A \cap V$  是横向的，则  $\tilde{f}_L$  是同构， $(x", u)$  是  $L$  上的坐标系）。然后我们得到  $Y$  上的坐标  $(x_1, x', x'', t''$  、  $u)$  ，我们完成  $t'$  消失在  $L$  上。我们设置 $t = (t', t'')$ 。然后：

\[
f(x_{1},x^{\prime},x^{\prime\prime},t^{\prime},t^{\prime\prime},u)=(x_{1},x^{\prime},x^{\prime\prime})
\]

\[
L=\left\{x_{1}=x^{\prime}=t^{\prime}=0\right\},
\]

\[
N=\left\{x_{1}=g(x^{\prime},x^{\prime \prime},t^{\prime},t^{\prime \prime},u)\right\}\quad\mathrm{f o r~s o m e~f u n c t i o n~}g.
\]

由于  $ \partial_{x'} g = \xi' / \xi_1 $  在  $ L $  上，我们可以写：

\[
g=\langle x^{\prime},u\rangle+h(x^{\prime},x^{\prime\prime},t^{\prime},t^{\prime\prime},u),
\]

与：

\[
\partial_{x^{\prime}}h=0\quad\mathrm{o n}\quad\{x^{\prime}=t^{\prime}=0\}~.
\]

包含 $L \subset N$ 意味着：

\[
h=0\quad\mathrm{o n}\quad\{x^{\prime}=t^{\prime}=0\}~,
\]

包含  $ T_N^*Y \times_NL \subset V $  意味着：

\[
\partial_{t^{\prime}}h=0\quad\mathrm{o n}\quad\left\{x^{\prime}=t^{\prime}=0\right\}~.
\]

让我们假设一下：

\[
\partial_{t^{\prime}t}^{2}h(0)\quad\mathrm{i s~n o n-d e g e n e r a t e}.
\]

然后根据莫尔斯引理（参见 Hörmander [4，III 附录 C]），我们可以通过用其他坐标替换  $ t' = (t_1, \ldots, t_r) $  来写：

\[
h(0,x^{\prime\prime},t^{\prime},t^{\prime\prime},u)=\sum_{j=1}^{r}\pm t_{j}^{2}\ ,
\]

我们得到：

\[
h(x^{\prime},x^{\prime \prime},t^{\prime},t^{\prime \prime},u)=q(t^{\prime})+\langle x^{\prime},\varphi(x^{\prime},x^{\prime \prime},t^{\prime},t^{\prime \prime},u)\rangle~,
\]

其中  $ q(t') $  是非简并二次形式（相对于  $ t' $  -变量）。通过（7.4.5）， $ \varphi(0,x",0,t",u)=0 $ 。然后将 u 替换为  $ u-\varphi $  ，我们得到结果。

最后我们证明(7.4.8)。让  $(x,t,u;\xi,\tau,v)$  表示  $T^*Y$  上的坐标，用  $\xi=(\xi_1,\xi',\xi''),\tau=(\tau',\tau'').$  我们有：

\[
\begin{aligned}{A}&{{}=\left\{x_{1}=\left\langle x^{\prime},u\right\rangle+h,-\xi^{\prime}/\xi_{1}=u+\partial_{x^{\prime}}h,-\xi^{\prime\prime}/\xi_{1}=\partial_{x^{\prime\prime}}h,\right.}\\ {}&{{}\quad\left.-\tau/\xi_{1}=\partial_{t}h,-v/\xi_{1}=x^{\prime}+\partial_{u}h\right\}\;,}\\ \end{aligned}
\]

\[
V=\left\{\tau=v=0\right\}.
\]

由(7.4.5)–(7.4.7)我们得到：

\[
\begin{aligned}{T_{p}A}&{{}=\left\{x_{1}=0,-\xi^{\prime}=u+\partial_{x^{\prime}x^{\prime}}^{2}h(0)\cdot x^{\prime}+\partial_{x^{\prime}t^{\prime}}^{2}h(0)\cdot t^{\prime},\right.}\\ {}&{{}\quad-\tau^{\prime}=\partial_{t^{\prime}t^{\prime}}^{2}h(0)\cdot t^{\prime}+\partial_{t^{\prime}x^{\prime}}^{2}h(0)\cdot x^{\prime},\xi^{\prime\prime}=\tau^{\prime\prime}=0,-v=x^{\prime}\}\enspace,}\\ \end{aligned}
\]

\[
T_{p}V=\left\{\tau=v=0\right\}.
\]

因此：

\[
\begin{aligned}{T_{p}A\cap T_{p}V}&{{}=\left\{x_{1}=x^{\prime}=\tau=v=\xi^{\prime\prime}=0,\right.}\\ {}&{{}\left.-\xi^{\prime}=u+\partial_{x^{\prime}t^{\prime}}^{2}h(0)\cdot t^{\prime},0=\partial_{t^{\prime}t^{\prime}}^{2}h(0)\cdot t^{\prime}\right\}\;.}\\ \end{aligned}
\]

由于映射  $T_p(T^* Y) \to T_{y_0} Y$  的  $T_p A \cap T_p V$  图像包含在  $T_{y_0} L$  中，因此我们有：  $(x, t, u; \xi, \tau, v) \in T_p A \cap T_p V \Rightarrow t' = 0$  。因此 $\partial_{t', t'}^2 h(0)$  是非退化的。这样就完成了证明。   $\square$ 

现在我们保留命题 7.4.1 给出的坐标系，并设置：

\[
N^{+}=\left\{(x,t,u);x_{1}\geq q(t)+\langle x^{\prime},u\rangle\right\}.
\]

令 $ n_{+} $ （或 $ n_{-} $ ）为 $ q(t) $ 的正（或负）特征值的数量，并通过以下关系定义 $ n_{0} $ ：

\[
n_{+}+n_{0}+n_{-}=\operatorname{d i m}Y/X-\operatorname{c o d i m}M+1.
\]

换句话说， $ n_{0} $  是 q 的完全各向同性空间的空间维数。

特别是，如果 q 是非简并的，则  $ n_{0} = 0 $ 。

\statementlabel{命题 7.4.2}我们保留命题7.4.1的假设，令 $ L \in \mathrm{Ob}(\mathbb{D}^{b}(\mathfrak{Mod}(A))) $  。

(i)  $ \mathbf{D}^{b}(X) $  的原对象范畴中存在态射：

\[
\stackrel{\mathrm{``L i m}}{\leftarrow}_{U}R f_{!}(L_{N^{+}\cap U})\to L_{M}[\delta]
\]

与  $ \delta = 1 - \operatorname{codim} M - (n_0 + n_-) = n_+ $  -  $ \operatorname{dim} Y/X $  。这里 U 的范围涵盖  $ y_0 $  的开放邻域系列。

(ii)  $ f_i^\mu L_{N^+} \simeq \varprojlim_U R f_i(L_{N^+ \cap U}) $  存在于  $ \mathbf{D}^b(X; f_\pi(p)) $  中，并且 (i) 中的态射是  $ \mathbf{D}^b(X; f_\pi(p)) $  中的同构（参见定义 6.1.7 和命题 6.1.10）。

请注意，（i）和（ii）来自我们将要证明的接下来的两个陈述。

(a) 对于  $y_{\circ}$  的开邻域  $U$  ，  $\mathbf{D}^b(X)$  中存在态射  $Rf_!(L_{N+\circ U}) \to L_M[\delta]$  。

(b) 在 $U$  中存在 $y_{\circ}$  的开邻域系统，使得对于该系统中的任何 $U'$ ，态射 $Rf_!(L_{N^+\cap U^-}) \to L_M[\delta]$  是 $\mathbf{D}^b(X; f_{\pi}(p))$  中的同构。

另请注意，“  $ \lim_{u} Rf_1(L_{N^+ \cap U}) $  不存在于  $ \mathbb{D}^b(X) $  中。

\statementlabel{证明}我们将坐标 $ (t) $ 分解为 $ (t) = (s', s'', t'') $ ，这样

\[
q(t)=s^{\prime2}-s^{\prime2}
\]

我们将 f 分解如下：

\[
f\colon Y\xrightarrow[f_{1}]{}Y_{1}\xrightarrow[f_{2}]{}Y_{2}\xrightarrow[f_{3}]{}Y_{3}\xrightarrow[f_{4}]{}X
\]

与 $f_{1}(x,s^{\prime},s^{\prime\prime},t^{\prime},u)=(x,s^{\prime},t^{\prime},u),f_{2}(x,s^{\prime},t^{\prime},u)=(x,t^{\prime},u),f_{3}(x,t^{\prime},u)=(x,u),f_{4}(x,u)=(x).$ 

让  $ \eta > 0, \varepsilon > 0 $  与  $ \eta \ll \varepsilon $  。设定：

\[
U_{0}=\left\{(x,s^{\prime},s^{\prime \prime},t^{\prime},u);|x|<\eta,|t^{\prime}|<\varepsilon,|u|<\varepsilon,s^{\prime}{}^{2}<\varepsilon,s^{\prime \prime}{}^{2}<\varepsilon\right\}~.
\]

(a) 让我们计算  $ Rf_{1!}L_{N^{+}\cap U_{0}}. $  我们有  $ |x|<\eta,|t''|<\varepsilon,|u|<\varepsilon,s'^{2}<\varepsilon $  ：

\[
f_{1}^{-1}(x,s^{\prime},t^{\prime},u)\cap N_{+}\cap U_{0}=\left\{s^{\prime \prime};\langle x^{\prime},u\rangle+s^{\prime 2}-x_{1}\leqslant s^{\prime \prime}{}^{2}<\varepsilon\right\}.
\]

对于  $ \langle x', u \rangle + s'^2 - x_1 \geq \varepsilon $  该集合是空的，对于  $ \langle x', u \rangle + s'^2 - x_1 \leq 0 $  来说它是一个开球，并且它与该集合同胚（与  $ 0 < a < b $  ）：

\[
S_{a}^{b}=\{a\leqslant s^{\prime\prime2}<b\}\qquad\mathrm{f o r}\qquad0<\langle x^{\prime},u\rangle+s^{\prime2}-x_{1}<\varepsilon.
\]

从 $ R\Gamma_c(S_a^b; L_{\mathbb{R}^{n-}})=0 $ 开始，我们得到：

\[
R f_{1!}L_{N^{+}\cap U_{0}}\simeq L_{N_{1}^{+}\cap U_{1}}[-n_{-]},
\]

其中  $ U_{1}=\{|x|<\eta,|t'|<\varepsilon,|u|<\varepsilon,s^{\prime2}<\varepsilon\} $  和  $ N_{1}^{+}=\{x_{1}\geqslant s^{\prime2}+\langle x^{\prime},u\rangle\} $  。

(b) 让我们计算 $ R f_{2!} L_{N_1^+ \cap U_1} $  。我们有  $ |x| < \eta $  、  $ |t'| < \varepsilon $  、  $ |u| < \varepsilon $  ：

\[
f_{2}^{-1}(x,t^{\prime},u)\cap N_{1}^{+}\cap U_{1}=\left\{s^{\prime};s^{\prime2}<\varepsilon,s^{\prime2}\leqslant x_{1}-\langle x^{\prime},u\rangle\right\}
\]

对于  $0 \leq x_1 - \langle x', u \rangle < \varepsilon$  ，该集合是封闭球或  $\{0\}$  。对于  $x_1 - \langle x', u \rangle < 0$  ，该集合是空的。由于  $|x| < \eta$  和  $\eta \ll \varepsilon$  ，我们不需要考虑  $x_1 - \langle x', u \rangle \geq \varepsilon$  的情况。因此：

\[
R f_{2!}L_{N_{1}^{+}\cap U_{1}}\simeq L_{N_{2}^{+}\cap U_{2}}\;,
\]

其中  $ U_{2}=\{|x|<\eta,|t'|<\varepsilon,|u|<\varepsilon\} $  和  $ N_{2}^{+}=\{x_{1}\geqslant\langle x',u\rangle\} $  。

(c) 让我们计算 $ R f_{3!} L_{N_2^+ \cap U_2} $  。我们立即发现：

\[
R f_{3!}L_{N_{2}^{+}\cap U_{2}}=L_{N_{3}^{+}\cap U_{3}}[-n_{0}],
\]

其中  $ U_{3} = \{|x| < \eta, |u| < \varepsilon\} $  和  $ N_{3}^{+} = \{x_{1} \geqslant \langle x', u \rangle\} $  。

(d) 最后，我们计算  $ Rf_{4!}L_{N_1^+ \cap U_3} $  。考虑一下著名的三角形：

\[
L_{\{|u|<\varepsilon,\langle x^{\prime},u\rangle>x_{1}\}}\xrightarrow{\quad}\:L_{\{|u|<\varepsilon\}\quad}\xrightarrow{\quad}L_{\{|u|<\varepsilon,\langle x^{\prime},u\rangle\leqslant x_{1}\}\quad}+1\quad.
\]

对于 $x_1 \geq \varepsilon |x'|$ ，集合 $\{u; |u| < \varepsilon, \langle x', u \rangle > x_1\}$  为空，否则为非空开凸集。

因此，如果  $ x_1 \geq \varepsilon|x'| $ ， $ (Rf_{4}!(L_{N_3^+ \cap U_3}))_x $  与  $ (Rf_{4}!L_{U_3})_x $  同构，否则为零：

\[
\begin{aligned}{{R f}_{4!}L_{N_{3}^{+}\cap U_{3}}}&{{}\simeq({R f}_{4!}L_{\{|u|<\varepsilon\}})_{\{x_{1}\geqslant\varepsilon|x^{\prime}|\}\cap U_{4}}}\\ {}&{{}\simeq L_{N_{4}^{+}\cap U_{4}}[-r],}\\ \end{aligned}
\]

其中  $ U_4 = \{|x| < \eta\} $  、  $ N_4^+ = \{x_1 \geq \varepsilon |x'|\} $  和  $ r $  是  $ u $  变量的空间维度，即  $ r = \text{codim} M - 1 $  。

(e) 我们已获得：

\[
R f_{!}L_{N^{+}\cap U_{0}}\simeq L_{N_{4}^{+}\cap U_{4}}\left[\delta\right],
\]

其中  $ \delta = 1 - \mathrm{codim} M - n_0 - n_- $  、  $ U_0 $  、  $ U_4 $  、  $ N_4^+ $  已在上面定义。为了完成证明，还需要注意自然态射  $ L_{\{x_1 \geq \varepsilon | x' \mid\}} \to L_M $  是  $ \mathbf{D}^b(X; f_\pi(p)) $  中的同构。 （回想一下 $ f_\pi(p) = (0; \mathrm{d}x_1) $ 。）

我们可以使用惯性指数来解释命题 7.4.2 中出现的转变 $ \delta $ 。有关该指数的属性，请参阅附录。

我们引入一些符号。在命题7.4.1的情况下，我们设定：

\[
\left\{\begin{aligned}{}&{{}E_{p}=T_{p}T^{*}Y,\qquad\lambda_{N}(p)=T_{p}T_{N}^{*}Y,}\\ {}&{{}\lambda_{\circ}(p)=T_{p}\pi_{Y}^{-1}\pi_{Y}(p),\qquad\lambda_{1}(p)=T_{p}f_{\pi}^{-1}\pi_{X}^{-1}\pi_{X}(f_{\pi}(p)).}\\ \end{aligned}\right.
\]

\statementlabel{命题 7.4.3}在命题7.4.2的情况下，我们有：

\[
\begin{array}{l}1-\operatorname{c o d i m}M-(n_{0}+n_{-})\\ \quad=\frac{1}{2}[1+\operatorname{d i m}M-\operatorname{d i m}Y-\operatorname{d i m}((T_{p}V)^{\perp}\cap\lambda_{N}(p))-\tau]\quad,\end{array}
\]

其中 $ \tau = \tau_{E_{p}}(\lambda_{\circ}(p), \lambda_{N}(p), \lambda_{1}(p)) $  。

（回想一下  $\tau_{E_{p}}(\cdot,\cdot,\cdot)$  是辛向量空间  $E_{p}$  中三个拉格朗日线性子空间的惯性指数。）

\statementlabel{证明}我们利用命题7.4.1给出的 $Y$ 上的坐标系 $(x,t,u)$ 、 $X$ 上的 $(x)$ ，使得 $f(x,t,u)=x$ 、 $M=\{x_1=x'=0\}$ 、 $N=\{x_1=\langle x',u\rangle+q(t)\}$ 和 $q(t)=\frac{1}{2}\langle At,t\rangle$ 成为对称矩阵 $A$ 。我们用  $(x,t,u;\xi,\tau,v)$  表示  $T^*Y$  上的相关坐标。

然后：

\[
\begin{aligned}&\lambda_{\circ}(p)=\left\{x=t=u=0\right\},\\ &\\&\lambda_{N}(p)=\left\{x_{1}=\xi^{\prime \prime}=0,\xi^{\prime}=-u,\tau=-At,v=-x^{\prime}\right\},\\ &\\&\lambda,(p)=\left\{x=\tau=v=0\right\}.\\ \end{aligned}
\]

考虑  $\rho = \{x = 0\}^\perp$  定义的  $E_p$  的各向同性空间  $\rho$  。我们可以将  $E_p^\rho$  识别为  $(t, u; \tau, v)$  变量的空间。从  $\rho^\perp \supset \lambda_\circ(p) + \lambda_1(p)$  开始，我们通过定理 A.3.2 得到：

\[
\tau=\tau_{E_{p}^{p}}(\{t=u=0\},\{v=0,\tau=-A t\},\{\tau=v=0\}).
\]

考虑各向同性空间  $ \{v = 0\}^{\perp} $  。我们得到：

\[
\tau=\tau(\{t=0\},\{\tau=-A t\},\{\tau=0\})~,
\]

其中惯性指数现在在  $(t,\tau)$  变量的空间中计算。根据命题 A.3.6，我们得到：

\[
\tau=-\operatorname{s g n}(A),
\]

其中  $ \mathrm{sgn}(A)=n_{+}-n_{-} $  是 q 的正特征值数量减去负特征值数量。

另一方面，我们有：

\[
\begin{array}{r}{n_{0}=\dim(T_{p}V^{\perp}\cap\lambda_{N}(p)).}\end{array}
\]

通过写 $ n_{0} + n_{-} = \frac{1}{2}(n_{0} + n_{-} + n_{+} + n_{0} - (n_{+} - n_{-})) $ ，我们发现由(7.4.10)：

\[
\begin{aligned}&1-\operatorname{codim}M-\tfrac{1}{2}[\operatorname{dim}Y/X-\operatorname{codim}M+1+\operatorname{dim}(T_{p} V^{\perp}\cap\lambda_{N}(p))+\tau]\\&=\tfrac{1}{2}[1-\operatorname{codim}M-\operatorname{dim}Y/X-\operatorname{dim}(T_{p} V^{\perp}\cap\lambda_{N}(p))-\tau],\\ \end{aligned}
\]

这是期望的结果。 ⑨

我们将应用前面的结果来组成与超曲面相关的内核。首先我们需要一个基本引理。

令  $ A_1 \subset T^*(X \times Y) $  和  $ A_2 \subset T^*(Y \times Z) $  为两个（圆锥）拉格朗日子流形，令  $ (p_X, p_Y^a) \in A_1 $  、  $ (p_Y, p_Z^a) \in A_2 $  。我们做出假设：

\[
\left\{\begin{aligned}{p_{2}^{a}|_{\varLambda_{1}}}&{{}:\varLambda_{1}\to T^{*}Y\quad\mathrm{a n d}\quad p_{1}|_{\varLambda_{2}}}&{}&{{}:\varLambda_{2}\to T^{*}Y}\\ {{\mathrm{a r e~t r a n s v e r s a l}}}&{{}}&{}\\ \end{aligned}\right.
\]

如§1中所示， $ T^*(X \times Y) \times_{T^*Y} T^*(Y \times Z) $ 表示 $ p_2^a $ 的纤维乘积： $ T^*(X \times Y) \to T^*Y $ 和 $ p_1: T^*(Y \times Z) \to T^*Y $ 。

\statementlabel{引理 7.4.4}假设(7.4.13)。那么  $\varLambda_1 \times_{T^*Y} \varLambda_2$  是  $T^*(X \times Y) \times_{T^*Y} T^*(Y \times Z)$  的平滑子流形。此外，用 $\varLambda_1 \cap U$ 和 $\varLambda_2 \cap V$ 替换 $\varLambda_1$ 和 $\varLambda_2$ ，其中 $U$ 和 $V$ 分别是 $(p_X, p_Y^a)$ 和 $(p_Y, p_Z^a)$ 足够小的开放邻域，有：

(i)  $ p_{13} $  诱导  $ \varLambda_{1} \times_{T^{*}Y} \varLambda_{2} $  与  $ T^{*}(X \times Z) $  的平滑拉格朗日子流形  $ \varLambda $  的同构。

(ii) 让  $\delta$  表示对角线嵌入  $X \times Y \times Z \hookrightarrow X \times Y \times Z$  。那么  $\delta_{\pi}$  是  $A_1 \times A_2$  的横向，而  $A_{12} = {}^{\mathrm{t}}\delta{}^{\prime}\delta_{\pi}^{-1}(A_1 \times A_2)$  是  $T^{*}(X \times Y \times Z)$  的拉格朗日子流形。此外， $t_{q_{13}}^{\prime}$  横向于 $A_{12}$  和 $q_{13\pi}t_{q_{13}^{\prime}-1}(A_{12}) = A$  。

\statementlabel{证明}该假设意味着  $ A_1 \times_{T^*Y} A_2 $  是  $ T^*(X \times Y \times Z) $  的平滑子流形，维度为  $ \dim(X \times Y) + \dim(Y \times Z) - \dim T^*Y = \dim(X \times Z) $  。

(i) 我们设定：

\[
\left\{\begin{aligned}{}&{{}E_{X}=T_{p_{X}}T^{*}X~,\quad E_{Y}=T_{p_{Y}}T^{*}Y~,\quad E_{Z}=T_{p_{Z}}T^{*}Z~,\quad}\\ {}&{{}\lambda_{1}=T_{(p_{X},p_{Y}^{a})}\varLambda_{1}~,\quad\lambda_{2}=T_{(p_{Y},p_{Z}^{a})}\varLambda_{2}~.}\\ \end{aligned}\right.
\]

如果 $E$  是辛向量空间，则 $E^{a}$  表示具有相反斜对称形式的空间 $E$ 。

我们将用“a”表示由 $x \mapsto x$ 给出的从 $E$ 到 $E^a$ 的映射。然后  $E_Y^a \simeq T_{p_Y^a}(T^*Y)$  、  $E_Z^a \simeq T_{p_Z^a}(T^*Z)$  和  $a^*$  ，我们仍然用  $\lambda_1$  和  $\lambda_2$  分别表示  $E_X \oplus E_Y^a$  和  $E_Y \oplus E_Z^a$  的拉格朗日平面，分别对应于  $\lambda_1$  和  $\lambda_2$  同构  $T_{(p_X, p_Y^a)} T^*(X \times Y) \simeq E_X \oplus E_Y^a$  和  $T_{(p_Y, p_Z^a)} T^*(Y \times Z) \simeq E_Y \oplus E_Z^a$  。

那么根据命题 A.1.4 我们有：

\[
\lambda_{1}\circ\lambda_{2}=p_{13}\left(\lambda_{1}\times\lambda_{2}\atop E_{Y}\right),
\]

其中  $ p_{13} $  是  $ (E_{X} \oplus E_{Y}^{a}) \times_{E_{Y}}(E_{Y} \oplus E_{Z}^{a}) $  到  $ E_{X} \oplus E_{Z}^{a} $  的投影。

由于  $\lambda_1 \circ \lambda_2$  是  $E_X \oplus \bar{E}_Z^a$  和  $\dim(\lambda_1 \times_{E_Y} \lambda_2) = \dim \lambda_1 + \dim \lambda_2 - \dim E_Y = \dim(X \times Z) = \dim \lambda_1 \circ \lambda_2$  的拉格朗日子空间，因此我们看到  $p_{13} : \quad \lambda_1 \times_{T \cdot Y} \lambda_2 \to T^*(X \times Z)$  是嵌入， $p_{13}(\lambda_1 \times_{T \cdot Y} \lambda_2)$  是拉格朗日。

(ii) 足以表明

\[
\begin{aligned}\delta_{\pi}:T_{(p_{X},p_{Y}^{\alpha},p_{Y},p_{Z}^{\alpha})}\Bigg(X\times Y\times Z&\underbrace{\times}_{X\times Y\times Y\times Z}T^{*}(X\times Y\times Y\times Z)\Bigg)\\\hookrightarrow T_{(p_{X},p_{Y}^{\alpha},p_{Y},p_{Z}^{\alpha})}(X\times Y\times Y\times Z)\end{aligned}
\]

横向于  $ \lambda_{1} \times \lambda_{2} \simeq T_{(p_{x}, p_{y}^{a}, p_{y}, p_{z}^{a})} (\varLambda_{1} \times \varLambda_{2}) $  且

\[
{}^{t}q_{13}^{\prime}:T_{(y_{\circ},p_{X},p_{Z}^{a})}(Y\times T^{*}(X\times Z))\hookrightarrow T_{(p_{X},y_{\circ},p_{Z}^{a})}T^{*}(X\times Y\times Z)
\]

横向于  $\mathrm{Im}(\delta_{\pi}^{-1}(\lambda_1 \times \lambda_2) \xrightarrow{\iota_0} T_{(p_X, y_0, p_Z^a)} T^*(X \times Y \times Z))$  ，其中  $y_0 = \pi(p_Y) \in T_Y^* Y$  。现在，令  $\rho' \subset \rho \subset T_{(p_Y^a, p_Y)}(T^*(Y \times Y))$  为由下式定义的各向同性子空间

\[
\rho^{\prime\bot}=T_{(p_{Y}^{\alpha},p_{Y})}(Y\times_{Y\times Y}T^{*}(Y\times Y))\mathrm{a n d}\rho=T_{(p_{Y}^{\alpha},p_{Y})}(T_{Y}^{*}(Y\times Y))\enspace.
\]

然后我们有：

\[
T_{(p_{X},p_{Y}^{a},p_{Y},p_{Z}^{a})}\left(X\times Y\times Z_{\substack{\times\times X\times Y\times X\times Z}}\times T^{*}(X\times Y\times Y\times Z)\right)\simeq E_{Y}\oplus\rho^{\prime\perp}\oplus E_{Z}^{a}~,
\]

\[
T_{(p_{x},y_{\circ},p_{z}^{a})}T^{*}(X\times Y\times Z)\simeq(E_{X}\oplus E_{Y}^{a}\oplus E_{Y}\oplus E_{Z}^{a})^{(0\oplus\rho^{\prime}\oplus0)}\;,
\]

and

\[
T_{(y\circ,p_{x},p_{z}^{a})}(Y\times T^{*}(X\times Z))\simeq(E_{X}\oplus\rho\oplus E_{Z}^{a})^{(0\oplus\rho^{\prime}\oplus0)}.
\]

因此问题被简化为

\[
\begin{aligned}(\lambda_{1}\oplus\lambda_{2})+(E_{X}\oplus\rho^{\prime\bot}\oplus E_{Z}^{a})&=E_{X}\oplus E_{Y}^{a}\oplus E_{Y}\oplus E_{Z}^{a},\\ (\lambda_{1}\oplus\lambda_{2})^{(0\oplus\rho^{\prime}\oplus0)}+(E_{X}\oplus\rho\oplus E_{Z}^{a})^{(0\oplus\rho^{\prime}\oplus0)}&=(E_{X}\oplus E_{Y}^{a}\oplus E_{Y}\oplus E_{Z}^{a})^{(0\oplus\rho^{\prime}\oplus0)}\;.\end{aligned}
\]

这立即来自：

\[
\begin{aligned}{}&{{}(\lambda_{1}\oplus\lambda_{2})\cap(0\oplus\rho^{\prime}\oplus0)\subset(\lambda_{1}\oplus\lambda_{2})\cap(0\oplus\rho\oplus0)}\\ {}&{{}\quad\simeq\operatorname{K e r}\biggl(\lambda_{1}\mathop{\times}_{E_{\gamma}}\lambda_{2}\to E_{X}\oplus E_{Z}^{a}\biggr)=0~.\quad\Box}\\ \end{aligned}
\]

\statementlabel{定义 7.4.5}在引理 7.4.4 的情况下，设置  $ \Lambda = \Lambda_1 \circ \Lambda_2 $  。

当然，  $A_{1} \circ A_{2}$  仅在  $(p_{x}, p_{z}^{a})$  的邻域中定义。更准确地说， $A_{1} \circ A_{2}$  是  $(p_{x}, p_{z}^{a})$  集合的萌芽。现在让 $S_{1} \subset X \times Y$  和 $S_{2} \subset Y \times Z$  成为两个超曲面。我们设定：

\[
\varLambda_{1}=T_{S_{1}}^{*}(X\times Y)\cap(\dot{T}^{*}X\times\dot{T}^{*}Y),
\]

\[
{\cal A}_{2}=T_{S_{2}}^{*}(Y\times Z)\cap(\dot{T}^{*}Y\times\dot{T}^{*}Z)~.
\]

让 $ (p_X, p_Y^a) \in \mathcal{A}_1 $ ， $ (p_Y, p_Z^a) \in \mathcal{A}_2 $ 。我们设置 $ x_\circ = \pi(p_X) $ ， $ y_0 = \pi(p_Y) $ ， $ z_\circ = \pi(p_Z) $ 。那么 $ p_X \notin T_X^*X $ 和 $ p_Z \notin T_Z^*Z $ 意味着：

\[
q_{2}|_{S_{1}}:S_{1}\to Y\quad\mathrm{a n d}\quad q_{1}|_{S_{2}}:S_{2}\to Y\mathrm{a r e~s m o o t h}.
\]

我们假设 (7.4.13) 和以下内容：

\[
\left\{\begin{aligned}{}&{{}A_{1}\circ A_{2}=T_{S}^{*}(X\times Z)\operatorname{o n a\ n e i g h b o r h o o d o f}(p_{X},p_{Z}^{a}),}\\ {}&{{}\operatorname*{f o r\quad a s u b m a n i f o l d}S\operatorname{o f}X\times Z.}\\ \end{aligned}\right.
\]

我们保留符号（7.4.14）并且我们还设置：

\[
\left\{\begin{aligned}{\lambda_{\circ}(p_{X})}&{{}=T_{p_{X}}\pi^{-1}(x_{\circ})}&{,\qquad\lambda_{\circ}(p_{Y})}&{{}=T_{p_{Y}}\pi^{-1}(y_{0})}&{,}\\ {\lambda_{\circ}(p_{Z})}&{{}=T_{p_{Z}}\pi^{-1}(z_{\circ})}&{.}\\ \end{aligned}\right.
\]

我们还应缩写为 $\lambda_{\circ X}$ （分别为 $\lambda_{\circ Y}$ ， $\lambda_{\circ Z}$ ）而不是 $\lambda_{\circ}(p_{X})$ （分别为 $\lambda_{\circ}(p_{Y})$ ， $\lambda_{\circ}(p_{Z}$ ））。这些分别是  $E_{X}$  、  $E_{Y}$  、  $E_{Z}$  的拉格朗日子空间。最后我们设置：

\[
\left\{\begin{aligned}\lambda_{1}(p_{Y})&=p_{2}^{a}(\lambda_{1}\cap p_{1}^{-1}(\lambda_{\circ}(p_{X})))\\ &=T_{p_{Y}}(p_{2}^{a}(\varLambda_{1}\cap\pi^{-1}(x_{\circ})\times T^{*}Y))=\lambda_{\circ}(p_{X})\circ\lambda_{1}^{a},\\ \lambda_{2}(p_{Y})&=p_{1}(\lambda_{2}\cap p_{2}^{a-1}(\lambda_{\circ}(p_{Z})))\\ &=T_{p_{Y}}(p_{1}(\varLambda_{2}\cap T^{*}Y\times\pi^{-1}(z_{\circ})))=\lambda_{2}\circ\lambda_{\circ}(p_{Z}).\end{aligned}\right.
\]

这里 $p_2^a$ 是投影 $E_X \oplus E_Y^a \to E_Y^a$ 或 $E_Y \oplus E_Z^a \to E_Z^a$ ，我们使用命题A.1.4中的符号和命题A.3.8之后给出的符号。请注意，  $\lambda_1(p_Y)$  和  $\lambda_2(p_Y)$  是  $E_Y$  的拉格朗日子空间。

\statementlabel{命题 7.4.6}令  $ (p_X, p_Y, p_Z) \in \dot{T}^*X \times \dot{T}^*Y \times \dot{T}^*Z $  并令  $ S_1 $  和  $ S_2 $  为  $ X \times Y $  和  $ Y \times Z $  的两个超曲面，如上所述。

令 $ L' $  和 $ L'' $  属于 $ \mathbf{D}^b(\mathfrak{Mod}(A)) $  并让 $ L = L' \otimes^L L'' $  。假设横向条件（7.4.13）并假设（7.4.17）。那么  $ (L'_{S_1}, L'_{S_2}) $  ，在  $ (p_X, p_Y, p_Z) $  上是微局部可组合的，并且  $ \mathbf{D}^b(X \times Z; (p_X, p_Z^a)) $  中有一个同构：

\[
L_{S_{1}\;\;\;\;\mu}^{\prime}\circ L_{S_{2}}^{\prime\prime}\simeq L_{S}[\delta]~,
\]

与

\[
\left\{\begin{array}{l}{{\delta=1-\frac{1}{2}[\dim Y+\operatorname{c o d i m}S+\tau]~,}}\\ {{\quad}}\\ {{\tau=\tau_{E_{Y}}(\lambda_{\circ}(p_{Y}),\lambda_{2}(p_{Y}),\lambda_{1}(p_{Y}))~.}}\end{array}\right.
\]

\statementlabel{证明}第一个语句 (7.4.21) 紧随 (7.4.20) 之后。根据 (7.4.16)， $ q_{12}^{-1}S_1 $  和  $ q_{23}^{-1}S_2 $  在  $ X \times Y \times Z $  中横向相交。设定：

\[
S_{12}=q_{12}^{-1}S_{1}\cap q_{23}^{-1}S_{2}~,\qquad W=q_{12}^{-1}S_{1}~,\qquad f=q_{13}|_{w}~.
\]

到（7.4.16） $f$  顺利。此外， $S_{12}$  是 $W$  的超曲面， $f'$  是 $T_{S_{12}}^{*}$  W 的横向， $f_{\pi}$  引发同构 $f^{\prime -1}T_{S_{12}}^{*}W\simeq T_{S}^{*}(X\times Z)$  。正在申请

命题7.4.2和7.4.3我们得到 $ \mathbf{D}^{b}(X \times Z; (p_{x}, p_{z}^{a})) $ 中的同构：

\[
\stackrel{\sim}{\leftarrow}_{U}\mathbf{R}f_{!}(L_{S_{12}\cap U})\simeq L_{S}[\delta]
\]

其中 $U$ 的范围穿过 $W$ 中 $(x_o, y_o, z_o)$ 的开邻域族，并且将在后面计算移位 $\delta$ 。令 $U_1$ 、 $U_2$ 、 $U_3$ 分别为 $x_o$ 、 $y_o$ 、 $z_o$ 在 $X$ 、 $Y$ 、 $Z$ 、 $U = (U_1 \times U_2 \times U_3) \cap W, U_{13} = U_1 \times U_3$ 的开邻域。我们有 $R f_!(L_{S_12 \cap U}) \simeq (R q_{13}!(L'_{S_1} \otimes^L L''_{S_2})_{U_2})_{U_{13}}$ 。

通过引理7.3.3，我们得到“ $ \lim_{U_1} (\tilde{L}_{S_1}' \circ \tilde{L}_{S_2}'')_{U_2} $ ”存在并且与 $ \mathbf{D}^b(X \times Z; (p_X, p_Z^a)) $ 中的 $ L_S[\delta] $ 同构。

根据命题 7.3.1 (ii)，我们有：

\[
L_{S_{1}\setminus\setminus\setminus}^{\prime}\circ_{L_{S_{2}}^{\prime\prime}}\simeq L_{S}[\delta]\qquad\mathrm{i n}\qquad\mathbb{D}^{b}(X\times Z;(p_{X},p_{Z}^{a}))~.
\]

最后我们计算偏移 $ \delta $ 。令  $ \tilde{p} = (p_X, p_Y^a, p_Y, p_Z^a) \in T_{S_1 \times S_2}^* (X \times Y \times Y \times Z) $  、  $ p' = \iota i'(\tilde{p}) \in T_{S_{12}}^* (X \times Y \times Z) $  ，其中  $ \iota $  表示对角嵌入  $ X \times Y \times Z \hookrightarrow X \times Y \times Y \times Z $  ，以及  $ \iota p'' = \iota j' (p') $  ，其中  $ j $  表示嵌入  $ W \hookrightarrow X \times Y \times Z $  。套装

\[
E^{\prime\prime}=T_{p^{\prime\prime}}T^{*}W,\qquad\lambda_{\circ}^{\prime\prime}=T_{p^{\prime\prime}}\pi^{-1}\pi(p^{\prime\prime}),
\]

\[
\lambda_{1}^{\prime\prime}=T_{p^{\prime\prime}}f_{\pi}^{-1}(\pi^{-1}(x_{\circ},z_{\circ}))\quad,\quad\lambda_{2}^{\prime\prime}=T_{p^{\prime\prime}}T_{S_{12}}^{*}W\quad.
\]

根据命题 7.4.3，我们有：

\[
\delta=\tfrac{1}{2}[1+\dim S-(\dim(X\times Y\times Z)-1)-\tau^{\prime \prime}]
\]

与 $\tau'' = \tau_{E}''(\lambda_\circ, \lambda_2, \lambda_1)$  。自  $\frac{1}{2}[1 + \dim S - (\dim(X \times Y \times Z) - 1)] = 1 - \frac{1}{2}(\dim Y + \dim S)$  以来，仍然需要证明  $\tau = \tau'$  。

设置 $E' = T_p', T^*(X \times Y \times Z)$ 、 $\lambda_o' = T_p', \pi^{-1}\pi(p'), \lambda_1' = T_p', T^*_{\{x_o\} \times Y \times \{z_o\}} (X \times Y \times Z)$ 、 $\lambda_2' = T_p', T^*_{S_{12}}(X \times Y \times Z)$ 。请注意  $p' \in T^*X \times T_Y^*Y \times T^*Z$  。考虑  $\rho' = (T_p)(W \times_{X \times Y \times Z} T^*(X \times Y \times Z)))^\perp$  给出的  $E'$  的各向同性空间  $\rho'$  。然后是 $(E')^{\rho'} \simeq E''$ 、 $\rho' \subset \lambda_o \cap \lambda_2$  和 $(\lambda_i')^{\rho'} \simeq \lambda_i''$  (i = 0, 1, 2)。因此：

\[
\tau^{\prime\prime}=\tau_{\mathrm{\scriptsize~d e f}}^{\prime}=\tau_{E^{\prime}}(\lambda_{\circ}^{\prime},\lambda_{2}^{\prime},\lambda_{1}^{\prime})\enspace.
\]

现在设置  $\widetilde{E} = T_{\widetilde{p}}T^{*}(X \times Y \times Y \times Z)$  、  $\widetilde{\lambda}_{\circ} = T_{\widetilde{p}}\pi^{-1}\pi(\widetilde{p}),\widetilde{\lambda}_{1} = T_{\widetilde{p}}T_{\{x_{\circ}\}\times Y\times\{z_{\circ}\}}^{*}(X \times Y \times Y \times Z)$  、  $\widetilde{\lambda}_{2} = T_{p}T_{S_{1}\times S_{2}}^{*}(X \times Y \times Y \times Z)$  。考虑  $\widetilde{\rho} = \{0\} \times (T_{(p_{Y}^{\alpha}, p_{Y})}(T_{Y}^{*}(Y \times Y) \cap \pi^{-1}\pi(p_{Y}^{\alpha}, p_{Y}))) \times \{0\}$  给出的  $\widetilde{E}$  的各向同性空间  $\widetilde{\rho}$  。然后 $\widetilde{E}^{\widetilde{\rho}} \simeq E^{\prime}$ 、 $\widetilde{\rho} \subset \widetilde{\lambda}_{\circ} \cap \widetilde{\lambda}_{1}$ 、 $(\widetilde{\lambda}_{i})^{\widetilde{\rho}} \simeq \lambda_{i}^{\prime}$ 、 $(i=0,1,2)$ 。因此：

\[
\tau^{\prime}=\tilde{\tau}_{\mathrm{\footnotesize\footnotesize def}}=\tau_{\tilde{E}}(\tilde{\lambda}_{\circ},\tilde{\lambda}_{2},\tilde{\lambda}_{1})\quad.
\]

现在我们使用符号 (7.4.18), (7.4.19) 并识别  $ \tilde{E} $  和  $ E_X \oplus E_Y^a \oplus E_Y \oplus E_Z^a $  。然后：

\[
\tilde{\tau}=\tau(\lambda_{\circ,X}\oplus\lambda_{\circ,Y}\oplus\lambda_{\circ,Y}\oplus\lambda_{\circ,Z},\lambda_{1}\oplus\lambda_{2},\lambda_{\circ,X}\oplus\varDelta\oplus\lambda_{\circ,Z})
\]

其中  $\varDelta$  是  $E_{Y}^{a}\oplus E_{Y}$  的对角线。

因此：

\[
\tilde{\tau}=\tau_{E_{Y}^{a}\oplus E_{Y}}(\lambda_{\circ,Y}\oplus\lambda_{\circ,Y},p_{2}^{a}(\lambda_{1}\cap p_{1}^{-1}(\lambda_{\circ,X}))\oplus p_{1}(\lambda_{2}\cap p_{2}^{a-1}(\lambda_{\circ,Z})),\varDelta).
\]

应用命题 A.3.9，我们得到：

\[
\begin{aligned}\tilde{\tau}&=\tau_{E_{Y}}(\lambda_{\circ,Y},\lambda_{\circ,Y},\lambda_{2}(p_{Y}),\lambda_{1}(p_{Y}))\\&=\tau_{E_{Y}}(\lambda_{\circ,Y},\lambda_{2}(p_{Y}),\lambda_{1}(p_{Y}))\\&=\tau~.\quad\Box\end{aligned}
\]

\statementlabel{推论 7.4.7}令 $X$  和 $Y$  为两个相同维度的流形 $n$  ， $S$  为 $X \times Y$  的超曲面。设置  $A = T_S^*(X \times Y)$  ，并让  $p = (p_X, p_Y^a) \in A \cap (\dot{T}X \times \dot{T}Y)$  。假设

\[
\left\{\begin{aligned}{}&{{}p_{1}|_{\varLambda}\colon\varLambda\to T^{*}X\quad{a n d}\quad p_{2}^{a}|_{\varLambda}\colon\varLambda\to T^{*}Y\quad{a r e l o c a l}}\\ {}&{{}{i s o m o r p h i s m s~i n~a~n e i g h b o r h o o d~o f~p~}.}\\ \end{aligned}\right.
\]

设置  $K = A_{S}$  、  $L = r_{*}A_{S}[n-1]$  ，其中  $r$  是规范映射  $X \times Y \to Y \times X$  。然后

\[
K\in\operatorname{O b}(\mathbb{N}(X,Y;p_{X},p_{Y})){~a n d~}L\in\operatorname{O b}(\mathbb{N}(Y,X;p_{Y},p_{X})),
\]

\[
\begin{array}{r}{\mathrm{(i i)}K\circ_{\mu}L\simeq A_{\varDelta_{X}}\:i n\:\mathbb{N}(X,X;p_{X},p_{X})\:a n d\:L\circ_{\mu}K=A_{\varDelta_{Y}}\:i n\:\mathbb{N}(Y,Y;p_{Y},p_{Y}).}\end{array}
\]

\statementlabel{证明}设置 $ x_{\circ} = \pi(p_X) $ ， $ y_{\circ} = \pi(p_Y) $ 。设置 $ A' = T_r(S)(Y \times X) $ ， $ A'' = A \circ A' $ 。然后， $ (x, x'; \xi, \xi') \in A'' $  当且仅当存在  $ (y; \eta) \in T^*Y $  使得  $ (x, y; \xi, -\eta) \in A $  和  $ (y, x'; \eta, \xi') \in A' $  。通过 (7.4.23)，这意味着  $ x = x' $  、  $ \xi = -\xi' $  。因此 $ A'' = T_dx^*(X \times X) $  。同样， $ A' \circ A = T_dy^*(Y \times Y) $  。因此命题7.4.6的所有假设都得到满足，我们得到：

\[
K\circ_{L}\simeq A_{\varDelta_{X}}[n-1+\delta]\quad\mathrm{i n}\quad\mathbb{D}^{b}(X\times X;(p_{X},p_{X}^{a}))~,
\]

与  $ \delta = 1 - \frac{1}{2}(2n + \tau) = 1 - n + \frac{1}{2}\tau $  ，

\[
\tau=\tau_{T_{p_{Y}}T^{\star}Y}(\lambda_{\circ}(p_{Y}),\lambda_{2}(p_{Y}),\lambda_{1}(p_{Y}))~,
\]

\[
\lambda_{\circ}(p_{Y})=T_{p_{Y}}\pi^{-1}\pi(p_{Y}),
\]

\[
\lambda_{2}(p_{Y})=p_{1}(T_{(p_{Y},p_{X}^{a})}T_{r(S)}^{*}(Y\times X)\cap p_{2}^{a-1}T_{p_{X}}\pi^{-1}(x_{\circ}))~,
\]

\[
\lambda_{1}(p_{Y})=p_{2}^{a}(T_{(p_{X},p_{Y}^{a})}T_{S}^{*}(X\times Y)\cap p_{1}^{-1}T_{p_{X}}\pi^{-1}(x_{\circ}))~.
\]

自  $ \mathbb{N}(X, X; p_X, p_X) $  中的  $ \lambda_1(p_Y) = \lambda_2(p_Y) $  、  $ \tau = 0 $  和  $ K \circ_\mu L \simeq A_{\Delta_X} $  开始。同样可以证明 $ L \circ_\mu K \simeq A_{\Delta_Y} $ 。   $ \Box $ 

\statementlabel{备注 7.4.8}存在 $ p_{X} $  的开邻域 $ \Omega_{X} $  和 $ p_{Y} $  的 $ \Omega_{Y} $  ，使得定理7.2.1 的假设(7.2.1) 和(7.2.3) 满足 $ K = A_{S} $  。

因此 $ \Phi_K $  和 $ \Psi_K $  给出 $ \mathbf{D}^b(X; \Omega_X) $  和 $ \mathbf{D}^b(Y; \Omega_Y) $  之间彼此相反的范畴的等价性，并且满足(7.2.5)。

通过命题 7.4.6，我们获得了等价性  $ \mathbf{D}^{b}(X; p_{X}) \simeq \mathbf{D}^{b}(Y; p_{Y}) $  的更直接证明，现在我们能够控制接触变换中出现的位移  $ \delta $  （参见示例 7.2.5、7.2.6）。对这些转变的研究是下一节的主题。

\subsection{纯层}
令  $A$  为  $T^*X$  、  $p\in A$  的（二次曲线）拉格朗日子流形，并令  $\varphi$  为  $X$  上的实函数（ $C^\infty$  类）。我们设定：

\[
{\cal A}_{\varphi}=\left\{(x;\mathrm{d}\varphi(x));x\in X\right\}.
\]

请注意，  $ A_{\varphi} $  是  $ T^{*}X $  的拉格朗日流形，但  $ A_{\varphi} $  通常不是圆锥曲线。

\statementlabel{定义 7.5.1}如果 $\varphi(\pi(p)) = 0$  和流形 $\Lambda$  和 $\Lambda_\varphi$  在 $p$  处横向相交，则 $\varphi$  在 $p$  处与 $\Lambda$  相交。

请注意，当  $A$  是  $X$  的子流形  $M$  的保角丛时，这个条件相当于说  $\pi(p)$  是  $\varphi|_{M}$  的非简并临界点（即它的 Hessian 矩阵是非简并的）。

我们设置为 $ p \in \mathcal{A}_{\varphi} \cap \mathcal{A} $ ，

\[
\lambda_{\circ}(p)=T_{p}\pi^{-1}\pi(p),\qquad\lambda_{A}(p)=T_{p}A,\qquad\lambda_{\varphi}(p)=T_{p}A_{\varphi},
\]

\[
\tau_{\varphi}(p)=\tau_{T_{p}T^{\star}X}(\lambda_{\circ}(p),\lambda_{A}(p),\lambda_{\varphi}(p))\enspace.
\]

如果没有混淆的风险，我们会写  $\lambda_{\circ},\tau_{\varphi},\ldots$  ，而不是  $\lambda_{\circ}(p),\tau_{\varphi}(p)$  等。

如果  $\varphi$  在  $p \in T^*X$  处横向于  $A$  ，则  $\mathrm{d}\varphi(p) \neq 0$  ，且集合  $\{\varphi = 0\}$  是光滑流形。在这种情况下我们设置：

\[
\lambda_{\{\varphi=0\}}(p)=T_{p}T_{\{\varphi=0\}}^{*}X.
\]

然后：

\[
\tau_{\varphi}=\tau(\lambda_{\circ},\lambda_{A},\lambda_{\{\varphi=0\}})~.
\]

事实上，用  $\rho(p)$  表示  $p$  处的欧拉向量场生成的  $T_p T^* X$  线； （回想一下，该向量场是由齐次辛坐标系  $(x; \xi)$  中的  $\sum_j \xi_j \frac{\partial}{\partial \xi_j}$  定义的）。那么  $\rho(p)$  是  $\lambda_\circ(p) \cap \lambda_A(p), (\lambda_\varphi(p))^{\rho(p)} = (\lambda_{\{ \varphi=0\}}(p))^{\rho(p)}$  中包含的各向同性空间，并且 (7.5.5) 由定理 A.3.2 得出。

请注意：

\[
\left\{\begin{aligned}{}&{{}\lambda_{\varphi}(p)\quad{a n d}\quad\lambda_{A}(p)\quad{i n t e r s e c t~t r a n s v e r s a l l y~i f~a n~o n l y~i f}}\\ {}&{{}\lambda_{\{\varphi=0\}}\cap\lambda_{A}(p)=\rho(p).}\\ \end{aligned}\right.
\]

\statementlabel{引理 7.5.2}令 $X$  和 $X'$  为两个相同维度的流形， $S \subset X' \times X$  为超曲面， $A_S = \dot{T}_S^*(X' \times X)$  。设  $\tilde{p} = (p', p^a) \in A_S$  ，并假设  $p_1|_{\mathcal{A}_S}$  和  $p_2^a|_{\mathcal{A}_S}$  是  $\tilde{p}$  邻域中的局部同构。用  $\chi: T^*X \to T^*X'$  表示在  $p$  邻域中定义的接触变换，由  $\chi = p_1|_{\mathcal{A}_S} \circ (p_2^a|_{\mathcal{A}_S})^{-1}$  给出。

令  $\varphi$  为  $X$  上的函数，使得  $A_{\varphi} \ni p$  。假设  $\chi(T_{\{\varphi=0\}}^{*})X)=T_{\{\psi=0\}}^{*}X^{\prime}$  ，对于  $X^{\prime}$  上的实数函数  $\psi$  。我们选择  $\psi$  这样  $\mathrm{d}\psi(\pi(p^{\prime}))=p^{\prime}$  和  $\psi(\pi(p^{\prime}))=0$  。让  $F\in\mathrm{Ob}(\mathbb{D}^{b}(X))$  ，并设置  $K=A_{S}$  。

然后，对于  $ T_p(T^*X) $  的任何拉格朗日平面  $ \lambda $  和任何  $ j \in \mathbb{Z} + \tau_\varphi/2 $  ，我们有：

\[
R\varGamma_{\{\varphi\geqslant0\}}(F)_{\pi(p)}[j+\tau_{\varphi}/2]\simeq R\varGamma_{\{\psi\geqslant0\}}(\varPhi_{K}F)_{\pi(p^{\prime})}[j+(\tau_{\psi}/2)+d]
\]

与  $d = \tfrac{1}{2}(n-1) + \tfrac{1}{2}\tau(\lambda_{\circ}(p), \lambda, \lambda_{1}(p))$  和  $\lambda_{1}(p) = \chi(\lambda_{\circ}(p^{\prime})) \subset T_{p}(T^{*}X), \tau_{\varphi} = \tau(\lambda_{\circ}(p), \lambda, \lambda_{\varphi}(p))$  和  $\tau_{\psi} = \tau(\lambda_{\circ}(p^{\prime}), \chi(\lambda), \lambda_{\psi}(p^{\prime})) = \tau(\lambda_{1}(p), \lambda, \lambda_{\varphi}(p))$  。

\statementlabel{证明}首先要注意的是：

\[
R\varGamma_{\{\varphi\geqslant0\}}(F)_{\pi(p)}=\mu\mathcal{h}o m(A_{\{\varphi=0\}},F)_{p}~.
\]

应用定理7.2.1和命题7.4.6，我们发现：

\[
R\varGamma_{\{\varphi\geqslant0\}}(F)_{\pi(p)}=R\varGamma_{\{\psi\geqslant0\}}(\varPhi_{K}(F))_{\pi(p^{\prime})}[-\delta],
\]

其中  $ \delta $  由  $ \mathbf{D}^b(X'; p') $  中的关系  $ \Phi_K(A_{\{\varphi=0\}}) \simeq A_{\{\psi=0\}}[\delta] $  定义，即（命题 7.4.6）：

\[
\delta=1-\tfrac{1}{2}[\dim X+1+\tau(\lambda_{o}(p),\lambda_{\{\varphi=0\}}(p),\lambda_{1}(p))].
\]

让我们计算 $ d = \tfrac{1}{2}\tau_{\varphi} - \delta - \tfrac{1}{2}\tau_{\psi} $  。由余循环条件（定理A.3.2）我们得到：

\[
\begin{aligned}2d&=\tau(\lambda_{\circ},\lambda,\lambda_{\varphi})-(1-n-\tau(\lambda_{\circ},\lambda_{\varphi},\lambda_{1}))-\tau(\lambda_{1},\lambda,\lambda_{\varphi})\\&=(n-1)+\tau(\lambda_{\circ},\lambda,\lambda_{1})~.\quad\Box\end{aligned}
\]

\statementlabel{命题 7.5.3}假设  $\varphi$  在  $p \in T^*X$  处横向于  $\mathcal{A}$  ，并设  $F \in \mathrm{Ob}(\mathbf{D}^b(X))$  。我们假设  $\mathrm{SS}(F) \subset \mathcal{A}$  在  $p$  的邻域中。令  $j$  为一个数字，使得  $j - \frac{1}{2}(n + \dim(\lambda_o \cap \lambda_A)) \in \mathbb{Z}$  。那么 $R\Gamma_{(\varphi \geq 0)}(F)_{\pi(p)}[j + \tau_\varphi(p)/2]$  不依赖于 $\varphi$  。 （参见 (7.5.3)  $\tau_\varphi(p)$  的定义。）

\statementlabel{证明}(a) 首先假设  $A = T_M^*X$  为  $X$  的子流形  $M$  。让我们选择一个坐标系  $x = (x', x'')$  和  $x' = (x_1, \ldots, x_l)$  ，例如  $M = \{x'' = 0\}$  和  $\pi(p) = 0$  。让  $(x; \xi)$  表示  $T^*X$  上的关联坐标，以及  $\xi = (\xi'; \xi'')$  。然后：

\[
\Lambda_{\varphi}=\left\{(x;\xi);\xi_{j}=\partial\varphi/\partial x_{j}\right\}
\]

\[
T_{p}\varLambda_{\varphi}=\left\{(x;\xi);\xi_{j}=\sum_{k}\partial_{x_{j}x_{k}}^{2}\varphi(0)\cdot x_{k}\right\}
\]

\[
T_{p}T_{M}^{*}X=\{(x;\xi);x^{\prime\prime}=\xi^{\prime}=0\}\quad.
\]

当且仅当矩阵 $ (\partial_{x_j x_k} \varphi(0))_{1 \leq j, k \leq l} $  是非简并的时，交集 $ T_p A_\varphi \cap T_p T_M^* X $  是 $ \{0\} $ 。因此，根据莫尔斯引理（参见 Hörmander [4, III, 附录 C]），我们可以假设坐标更改后  $ \varphi|_{M} = \sum_{1 \leq j \leq l} a_j x_j^2 $  、  $ a_j \in \mathbb{R} $  、  $ a_j \neq 0 $  。我们得到：

\[
\begin{array}{c}{{\tau_{\varphi}=\tau(\{x=0\},\{x^{\prime\prime}=\xi^{\prime}=0\},\{\xi=\partial_{x,x}^{2}\varphi(0)\cdot x\})}}\\ {{{}}}\\ {{\tau(\{x^{\prime}=0\},\{\xi^{\prime}=0\},\{\xi^{\prime}=\partial_{x^{\prime},x^{\prime}}^{2}\varphi(0)\cdot x^{\prime}\})~,}}\end{array}
\]

其中最后一行中的索引是在  $ (x', \xi') $  -变量的空间中计算的。由命题 A.3.6 我们得到：

\[
\begin{aligned}{\tau_{\varphi}}&{{}=-\operatorname{s g n}(\partial_{x^{\prime},x^{\prime}}^{2}\varphi(0))}\\ {}&{{}=\#\left\{j;1\leqslant j\leqslant l,a_{j}<0\right\}-\#\left\{j,1\leqslant j\leqslant l,a_{j}>0\right\}.}\\ \end{aligned}
\]

另一方面，根据命题 6.6.1，存在  $ L \in \text{Ob}(\mathbb{D}^b(\mathfrak{Mod}(A))) $  使得  $ F \simeq L_M $  位于  $ \mathbb{D}^b(X; p) $  中。我们得到：

\[
\begin{aligned}{{R}\varGamma_{\{\varphi\geqslant0\}}(F)_{\pi(p)}[j+\tau_{\varphi}/2]}&{{}\simeq{R}\varGamma_{\{\sum_{1\leqslant i\leqslant l}a_{i}x_{i}^{2}\geqslant0\}}(L_{\mathbb{R}^{l}})_{0}[j+\tau_{\varphi}/2]}\\ {}&{{}\simeq L[j+\tau_{\varphi}/2-q],}\\ \end{aligned}
\]

其中 $ q = \# \{ j; 1 \leq j \leq l, a_j < 0 \} $  。 （这是在证明命题 7.4.2 的过程中执行的计算的一个特殊情况。）由于  $ \tau_{\varphi} - 2q = -l $  ，我们得到了这种情况下的结果。

(b) 我们处理一般情况。我们可以假设 $ p \in T^*X $ （否则 $ A = T_M^*X $ ，参见练习A.2）。令 $ \varphi_1 $  和 $ \varphi_2 $  为 $ p $  处 $ A $  的横向函数。根据推论 A.2.7，我们可以找到一个在  $ p $  邻域中定义的接触变换  $ \chi: T^*X \simeq T^*X' $  ，使得：

(i)  $ \chi = p_1|_{\mathcal{A}_S} \circ (p_2^a|_{\mathcal{A}_S})^{-1} $  、 $ \mathcal{A}_S = T_S^*(X' \times X) $  、 $ S $  是 $ X' \times X $  的超曲面， $ \chi(\mathcal{A}) = T_M^* X' $  是 $ X' $  的子流形 $ M' $  。

(ii)  $ \chi(\dot{T}_{\{\varphi_i=0\}}^{*})X) = \dot{T}_{\{\psi_i=0\}}^{*}X' $  对于某些函数  $ \psi_i $  (  $ i = 1,2 $  )。

然后结果来自证明的第一部分和引理 7.5.2，因为在此引理中获得的移位  $d$  不依赖于  $i$  。   $\square$ 

\statementlabel{定义 7.5.4}令  $\mathcal{A}$  为  $T^*X$  、  $p\in\mathcal{A}$  的拉格朗日子流形，并令  $F\in\mathrm{Ob}(\mathbf{D}^b(X))$  。假设  $\mathrm{SS}(F)\subset\mathcal{A}$  在  $p$  的邻域中，并令  $\varphi$  为  $p$  处  $\mathcal{A}$  的横截函数。令 $d$  为满足以下条件的数：

\[
d\equiv\tfrac{1}{2}\dim(\lambda_{\circ}(p)\cap\lambda_{A}(p))\bmod{\mathbb{Z}}\enspace,
\]

让 $ \tau_{\phi}(p) = \tau(\lambda_{\circ}(p), \lambda_{\wedge}(p), \lambda_{\phi}(p)) $ （参见（7.5.3）），并让 $ L \in \mathrm{Ob}(\mathbb{D}^{b}(\mathfrak{Mod}(A))) $ 。如果有

是同构：

\[
R\varGamma_{\{\varphi\geq0\}}(F)_{\pi(p)}[-d+\textstyle{\frac{1}{2}}\dim X+\textstyle{\frac{1}{2}}\tau_{\varphi}(p)]\simeq L\enspace,
\]

那么可以说 F 是 L 类型，并且在 p 处移位 d。

如果  $ H^{j}(L) = 0 $  对于  $ j \neq 0 $  ，则可以说 F 在 p 处是纯的。此外，如果 L 是等级一的自由 A 模，则称 F 在 p 处是简单的。

在上述情况中， $L$  的同构类称为 $F$  的类型，并在 $p$  处移位 $d$  。请注意，根据命题 7.5.3， $F$  的类型不依赖于 $\varphi$ 。

示例 7.5.5。 (i) 如果 $F$  具有类型 $L$  和shift  $d$  ，则 $F$  具有类型 $L[-k]$  和shift  $d+k$  且 $F[k]$  具有类型 $L$  和shift  $d+k$  。

(ii) 令  $M$  为  $X$  的闭子流形。那么  $A_M$  是一个简单的捆，在每个  $p \in T_M^*X$  处都有移位  $\frac{1}{2}\operatorname{codim} M$  。事实上，选择坐标  $x = (x', x'')$  和  $M = \{x'' = 0\}$  、  $x' = (x_1, \ldots, x_l)$  、  $p = (0; \lambda dx_n)$  。以 $\varphi(x) = \lambda x_n + \sum_{j=1}^l x_j^2$ 为例。然后  $\varphi$  在  $p$  处横向于  $T_M^*X$ ，我们有：

\[
\begin{aligned}(H^{k}_{\{\varphi\geq0\}}(A_{M}))_{0}&=A\quad\text{for}k=0~,\\&=0\quad\text{otherwise}~.\end{aligned}
\]

我们在 (7.5.7) 得到 $ d = \tfrac{1}{2}(\dim X - l) $ 。

(iii) 令  $\psi$  为  $X$  、  $Z = \{x; \psi(x) \geq 0\}$  、  $U = \{x; \psi(x) < 0\}$  、  $p = (x_o; \mathrm{d}\psi(x_o))$  上的实函数，以及  $\psi(x_o) = 0$  、  $\mathrm{d}\psi(x_o) \neq 0$  。然后  $A_Z$  （分别为  $A_U$  ）很简单，在  $p$  处使用  $\mathrm{shift} \, \tfrac{1}{2} (\mathrm{resp.}\, -\tfrac{1}{2})$  。事实上 $A_Z \simeq A_{\{\psi=0\}} \simeq A_U [1]$ 在 $\mathbf{D}^b(X; p)$ 。

(iv) 回忆示例 5.3.4：令  $ Z = \{x \in \mathbb{R}^2; x_1 > 0, -x_1^{3/2} \leq x_2 < x_1^{3/2}\} $  和  $ \varLambda = \{(x; \xi); \xi_2 > 0, x_2 = -(2\xi_1/3\xi_2)^3, x_1 = (2\xi_1/3\xi_2)^2\} $  。通过(iii)，捆 $ A_Z $ 在 $ \varLambda \cap \{\xi_1 > 0\} $ 上移位 $ \frac{1}{2} $ 并在 $ \varLambda \cap \{\xi_1 < 0\} $ 上移位 $ -\frac{1}{2} $ 。让我们计算  $ A_Z $  在  $ (0; dx_2) $  处的移位。选择 $ \varphi(x) = x_2 $ 。然后  $ \varphi $  在  $ (0; dx_2) $  和  $ \tau_\varphi = 0 $  处横向于  $ \varLambda $  。由于  $ (H_{\{x_2 \geq 0\}}^j(A_Z))_0 = A $  对应于  $ j = 1 $  ，否则为 0，因此移位为 0。

\statementlabel{命题 7.5.6}假设  $F$  的类型为  $L$  ，并在  $p$  处沿  $\mathcal{A}$  移动  $d$  。考虑在  $p$  邻域中定义的接触变换  $\chi: T^*X \to T^*X'$  ，并假设  $\chi = p_1|_{\mathcal{A}_S} \circ (p_2^a|_{\mathcal{A}_S})^{-1}$  、  $A_S = T_S^*(X' \times X)$  、  $S$  是超曲面。设置 $K = A_S$ 。那么  $\Phi_K(F)$  的类型为  $L$  ，并在  $\chi(p)$  处沿  $\chi(\mathcal{A})$  移动  $d'$  ，其中：

\[
d^{\prime}=d-\tfrac{1}{2}(n-1)-\tfrac{1}{2}\tau(\lambda_{o}(p),\lambda_{A}(p),\lambda_{1}(p))\quad and
\]

\[
\begin{aligned}\lambda_{1}(p)&=p_{2}^{a}(\lambda_{A_{s}}(\chi(p),p^{a})\cap p_{1}^{-1}(\lambda_{\circ}(\chi(p))))\\&=\chi^{-1}(\lambda_{\circ}(\chi(p)))=\lambda_{\circ}(\chi(p))\circ\lambda_{A_{s}}(\chi(p),p^{a})^{a}\end{aligned}
\]

\statementlabel{证明}我们可以在 $ X $ 上找到一个函数 $ \varphi $ ，使得 $ \varphi $ 在 $ p $ 和 $ \chi(T_{\{\varphi=0\}}^{*})X)=T_{\{\psi=0\}}^{*}X' $ 处横向于 $ A $ ，以获得平滑函数 $ \psi $ 。因此，结果来自引理 7.5.2。   $ \square $ 

\statementlabel{推论 7.5.7}令  $U$  为  $T^*X$  、  $L\in\mathrm{Ob}(\mathbb{D}^b(\mathfrak{Mod}(A)))$  、  $F\in\mathrm{Ob}(\mathbb{D}^b(X))$  的开子集并假设  $\mathrm{SS}(F)\cap U\subset A$  。然后， $A\cap U$  中的点集  $p$  使得  $F$  在  $p$  处的类型为  $L$  在  $A\cap U$  中开闭。

\statementlabel{证明}让 $p \in A \cap U$  。我们可能会发现满足命题 7.5.6 假设的接触变换  $\chi$  并且使得  $\chi(A) = T_N^* X^\prime$  ，对于  $X^\prime$  的超曲面  $N$  。然后  $\phi_K(F) \simeq L_N^\prime$  中的  $\mathbf{D}^b(X^\prime; \chi(p))$  ，对于一些  $L^\prime \in \mathrm{Ob}(\mathbf{D}^b(\mathfrak{M}_\mathrm{rob}(A)))$  。这意味着  $\mathbf{D}^b(X^\prime; \Omega)$  中的  $\phi_K(F) \simeq L_N^\prime$  是  $\chi(p)$  的开放邻域  $\Omega$  。根据命题7.5.6，我们可以简化为 $A = T_N^* X$  的情况。那么推论就很明显了。   $\square$ 

\statementlabel{备注 7.5.8}通过最后一个命题的证明，我们看到如果  $F$  在  $p$  处是  $L$  类型，那么在  $p$  处存在一个对象  $G$  simple 使得  $F \simeq L_X \otimes G$  在  $\mathbf{D}^b(X; p)$  中。

偏移的计算并不那么容易，如例 7.5.5 (iv) 所示，但我们使用以下结果将情况带入通用情况。

考虑一个连通的拓扑空间  $S$  和一个连续映射  $p: S \to \mathcal{A}$  。我们假设给定一个连续的拉格朗日平面族  $\mu(s)$  的  $T_{p(s)} T^* X$  、  $s \in S$  ，满足：

\[
\mu(s)\cap\lambda_{o}(p(s))=\mu(s)\cap\lambda_{A}(p(s))=\{0\}~.
\]

\statementlabel{命题 7.5.9}假设函数 $ d(s) - \frac{1}{2}\tau(\lambda_o(p(s)), \lambda_A(p(s)), \mu(s)) $ 在S上是常量。那么在 $ p(s) $ 处移位 $ d(s) $ 的F的类型是常量。

\statementlabel{证明}(a) 首先假设  $\varLambda = T_M^* X$  ，对于  $X$  的子流形  $M$  。那么  $d(s)$  是局部常数， $\tau(\lambda_o(p(s)), \lambda_A(p(s)), \mu(s))$  也是局部常数，因为  $\dim(\lambda_o(p(s)) \cap \lambda_A(p(s)))$  是常数。

(b) 为了处理一般情况，我们选择在  $p(s_o)$  的邻域中定义的接触变换  $\chi$  满足命题 7.5.6 的假设，并且使得  $\chi(A) = T_N^*X'$  对于  $X'$  的子流形  $N$  ，并且  $\chi(\mu(s_o))$  横向于  $\lambda_o(\chi(p(s)))$  。设置  $p'(s) = \chi(p(s))$  ，并简写为  $p$  或  $p'$  而不是  $p(s)$  或  $p'(s)$  。

让我们通过  $d'(s) = d(s) - \tfrac{1}{2}(n-1) - \tfrac{1}{2}\tau(\lambda_o(p), \lambda_A(p), \lambda_1(p))$  来定义  $d'(s)$ ，其中  $\lambda_1(p) = \chi^{-1}(\lambda_o(p'))$  。

那么根据命题 7.5.6， $F$  在  $p(s)$  处移位  $d(s)$  的类型是  $\Phi_K(F)$  在  $p'(s)$  处移位  $d'(s)$  的类型。我们有：

\[
\tau(\lambda_{\circ}(p^{\prime}),\lambda_{\chi(A)}(p^{\prime}),\chi(\mu(s)))=\tau(\lambda_{1}(p),\lambda_{A}(p),\mu(s))~.
\]

那么由(a)足以证明

\[
\begin{aligned}d(s)-\frac{1}{2}\tau(\lambda_{\circ}(p),\lambda_{A}(p),\lambda_{1}(p))-\frac{1}{2}\tau(\lambda_{1}(p),\lambda_{A}(p),\mu(s))&=d(s)-\frac{1}{2}\tau(\lambda_{\circ}(p),\mu(s),\lambda_{1}(p))\\&-\frac{1}{2}\tau(\lambda_{\circ}(p),\lambda_{A}(p),\mu(s))\end{aligned}
\]

是局部常数。由于 $\mu(s)$  横向于 $\lambda_{\circ}(p)$  和 $\lambda_{1}(p)$  ，因此仍然需要检查 $\dim(\lambda_{\circ}(p)\cap\lambda_{1}(p))$  是否是局部常数（定理A.3.2）。我们有：

\[
\begin{aligned}{\lambda_{\circ}(p)\cap\lambda_{1}(p)}&{{}=\lambda_{\circ}(p)\cap p_{2}^{a}(p_{1}^{-1}(\lambda_{\circ}(p^{\prime}))\cap\lambda_{T_{s}^{*}(X^{\prime}\times X)}(p^{\prime},p^{a}))}\\ {}&{{}\simeq\lambda_{T_{s}^{*}(X^{\prime}\times X)}(p^{\prime},p^{a})\cap\lambda_{\circ}(p^{\prime},p^{a})}\\ \end{aligned}
\]

由于 S 是超曲面，因此该空间的维数为 1。这样就完成了证明。 ⑨

让我们给出命题 7.5.9 的一个应用。

令 $ X_1 $  和 $ X_2 $  为两个流形， $ A_i  T^*X_i $  、 $ p_i \in A_i $  (  $ i = 1,2 $  ) 的拉格朗日子流形。我们用  $ q_i $  表示投影  $ X_1 \times X_2 \to X_i $  。

\statementlabel{命题 7.5.10}设  $ F_i \in \text{Ob}(\mathbf{D}^b(X_i)) $  ，并假设  $ F_i $  的类型为  $ L_i $  ，并沿  $ A_i $  、  $ (i = 1, 2) $  在  $ p_i $  处移位  $ d_i $  。

(i)  $ F_1 \boxtimes^L F_2 $  是  $ L_1 \otimes^L L_2 $  类型，并在  $ (p_1, p_2) $  和  $ \mathcal{A}_1 \times \mathcal{A}_2 $  处移位  $ d_1 + d_2 $  。

(ii)  $ R\mathcal{H}om(q_1^{-1}F_1, q_2^{-1}F_2) $  的类型为  $ R\mathcal{H}om(L_1, L_2) $  ，并在  $ (p_1^a, p_2) $  和  $ \Lambda_1^a \times \Lambda_2 $  处移位  $ d_2 - d_1 $  。

\statementlabel{证明}(i) 在  $A_i$  、  $F_i \simeq (L_i)_{M_i}[d(p'_i) - \frac{1}{2} \text{codim} M_i]$  、  $i = 1, 2$  的通用点  $p'_i$  处，对于  $X_i$  的子流形  $M_i$  。因此 $F_1 \boxtimes^L F_2 \simeq (L_1 \otimes L_2)_{M_1 \times M_2}[d(p'_i) + d(p'_2) - \frac{1}{2} \text{codim} M_1 \times M_2]$  。因此， $F_1 \boxtimes^L F_2$  在  $A_1 \times A_2$  的通用点处是  $L_1 \otimes^L L_2$  类型，因此根据推论 7.5.7 在  $A_1 \times A_2$  的每个点处。为了计算  $(p_1, p_2)$  处的移位，让我们选择两个拉格朗日空间族  $\mu_i(p'_i) \subset T_{p'_i} T^* X_i$  、  $(i = 1, 2)$  、  $p'_i \in A_i$  ，它们满足 (7.5.9)。让  $d(p'_1, p'_2)$  表示  $F_1 \boxtimes^L F_2$  在  $(p'_1, p'_2)$  处的移位。根据命题 7.5.9，我们有：

\[
\begin{aligned}{2(d(p_{1}^{\prime},p_{2}^{\prime})-d(p_{1},p_{2}))}&{{}=\tau(\lambda_{\circ}(p_{1}^{\prime},p_{2}^{\prime}),\lambda_{A_{1}\times A_{2}}(p_{1}^{\prime},p_{2}^{\prime}),\mu(p_{1}^{\prime})\oplus\mu(p_{2}^{\prime}))}\\ {}&{{}\quad-\tau(\lambda_{\circ}(p_{1},p_{2}),\lambda_{A_{1}\times A_{2}}(p_{1},p_{2}),\mu(p_{1})\oplus\mu(p_{2}))~,}\\ \end{aligned}
\]

\[
\begin{aligned}{2(d(p_{i}^{\prime})-d(p_{i}))}&{{}=\tau(\lambda_{\circ}(p_{i}^{\prime}),\lambda_{\varLambda_{i}}(p_{i}^{\prime}),\mu(p_{i}^{\prime}))}\\ {}&{{}\quad-\tau(\lambda_{\circ}(p_{i}),\lambda_{\varLambda_{i}}(p_{i}),\mu(p_{i}))~,\qquad i=1,2~.}\\ \end{aligned}
\]

由于  $ d(p_{1}^{\prime}, p_{2}^{\prime}) = d(p_{1}^{\prime}) + d(p_{2}^{\prime}) $  为通用  $ p_{i}^{\prime} $  ，并且

\[
\begin{aligned}\tau(\lambda_{\circ}(p_{1}^{\prime},p_{2}^{\prime}),\lambda_{A_{1}\times A_{2}}(p_{1}^{\prime},p_{2}^{\prime}),\mu(p_{1}^{\prime})\oplus\mu(p_{2}^{\prime}))&=\tau(\lambda_{\circ}(p_{1}^{\prime}),\lambda_{A_{1}}(p_{1}^{\prime}),\mu(p_{1}^{\prime}))\\&\quad+\tau(\lambda_{\circ}(p_{2}^{\prime}),\lambda_{A_{2}}(p_{2}^{\prime}),\mu(p_{2}^{\prime}))\end{aligned}
\]

对于任何 $ (p_{1}^{\prime}, p_{2}^{\prime}) $ ，我们得到 $ d(p_{1}, p_{2}) = d(p_{1}) + d(p_{2}) $ 。

(ii) 证明类似。 □

为了结束本节，我们将把命题 7.4.6 和 7.5.6 扩展到更一般的情况。令  $X$  、  $Y$  、  $Z$  为三个流形，  $A_1 \subset T^*(X \times Y)$  、  $A_2 \subset T^*(Y \times Z)$  为两个拉格朗日流形。让  $\tilde{p}_1 = (p_X, p_Y^a) \in A_1$  和  $\tilde{p}_2 = (p_Y, p_Z^a) \in A_2$  。我们保留符号 (7.4.14) 和 (7.4.19)，即我们为  $W = X$  、  $Y$  、  $Z$  、  $\lambda_1(p_Y) = \lambda_\circ(p_X) \circ \lambda_1^a$  、  $\lambda_2(p_Y) = \lambda_2 \circ \lambda_\circ(p_Z)$  设置  $\lambda_i = T_{\tilde{p}_i} A_i$  、  $i = 1, 2$  、  $\lambda_\circ(p_W) = T_{p_w} \pi^{-1} \pi(p_W)$  。我们引入符号（参见练习 A.9）：

\[
\tau(\lambda_{1}:\lambda_{2})=\tau(\lambda_{\circ}(p_{Y}),\lambda_{2}(p_{Y}),\lambda_{1}(p_{Y}))~.
\]

\statementlabel{定理 7.5.11}令 $ K_1 \in \mathrm{Ob}(\mathbf{D}^b(X \times Y; (p_X, p_Y^a))) $  和 $ K_2 \in \mathrm{Ob}(\mathbf{D}^b(Y \times Z; (p_Y, p_Z^a))) $  以及 $ \mathrm{SS}(K_i) \subset \mathcal{A}_i $  在 $ \tilde{p}_i $  ( $ i = 1, 2 $ ) 的邻域内。假设  $ p_2^a|_{\mathcal{A}_1} : \mathcal{A}_1 \to T^*Y $  和  $ p_1|_{\mathcal{A}_2} : \mathcal{A}_2 \to T^*Y $  是横向的， $ K_i $  的类型为  $ L_i $  ，并在  $ \tilde{p}_i $  、  $ i = 1, 2 $  处沿  $ \mathcal{A}_i $  移动  $ d_i $  。

那么  $(K_1, K_2)$  是微局部可组合的，  $(p_x, p_z^a)$  邻域中的  $\mathrm{SS}(K_1 \circ_\mu K_2) \subset \mathcal{A} = \mathcal{A}_1 \circ \mathcal{A}_2$  和  $K_1 \circ_\mu K_2$  的类型为  $L_1 \otimes^L L_2$  和  $\mathrm{shift} \, d_1 + d_2 - \frac{1}{2}(\mathrm{dim} \, Y + \tau(\lambda_1 : \lambda_2))$  以及  $(p_x, p_z^a)$  处的  $\mathcal{A}$  。

\statementlabel{证明}(a) 请注意，如果  $\varLambda_{1}$  和  $\varLambda_{2}$  是超曲面的共形丛， $\varLambda$  是子流形的共形丛，那么我们恢复命题 7.4.6。

如果  $Z = \{pt\}$  和  $A_{1}$  是与接触变换相关的超曲面的共形丛，那么我们恢复命题 7.5.6。

(b) 令 $\Delta$  表示 $\mathbb{R} \times \mathbb{R}$  的对角线，并令 $q \in \dot{T}^* \mathbb{R}$  。然后根据命题 7.3.6， $(K_1 \boxtimes A_\Delta) \circ_\mu (K_2 \boxtimes A_\Delta) \simeq (K_1 \circ_\mu K_2) \boxtimes A_\Delta$  。

将  $W, p_w$  替换为  $W \times \mathbb{R}, (p_w, q), (W = X, Y, Z)$  并将  $K_1, K_2$  替换为  $K_1 \boxtimes A_\delta, K_2 \boxtimes A_\delta$  ，我们可以从一开始就假设  $p_w \in \dot{T}^*W (W = X, Y, Z)$  。

(c) 可以通过以下过程将定理简化为  $ Z = \{\text{pt}\} $  的情况。使用与命题 7.3.4 中相同的符号， $ K_1 \circ_\mu K_2 \simeq i_*(K_1 \boxtimes A_Z) \circ_\mu K_2 $  。由于  $ i_*(K_1 \boxtimes A_Z) $  具有类型  $ L_1 $  和 shift  $ d_1 + \frac{1}{2}\dim Z $  ，以及  $ \tau(\lambda_1 : \lambda_2) = \tau(\lambda_1' : \lambda_2) $  和  $ \lambda'_1 = \lambda_1 \times T_Z^*(Z \times Z) $  ，因此  $ (K_1, K_2) $  的定理等效于  $ (i_*(K_1 \boxtimes A_Z), K_2) $  的定理。

(d) 令  $W$  为另一个流形，并令  $\mathcal{A}_3$  为  $T^*(Z \times W)$  和  $K_3 \in \mathrm{Ob}(\mathbf{D}^b(Z \times W; (p_Z, p_W^a)))$  的拉格朗日子流形。假设  $K_3$  具有类型  $L_3$ ，其中  $d_3$  沿  $\mathcal{A}_3$  进行移位，并且还假设映射  $\mathcal{A}_1 \circ \mathcal{A}_2 \to T^*Z$  和  $\mathcal{A}_3 \to T^*Z$  是横向的，并且  $\mathcal{A}_2 \to T^*Z$  和  $\mathcal{A}_3 \to T^*Z$  是横向的。因此， $\mathcal{A}_1 \to T^*Y$  和 $\mathcal{A}_2 \circ \mathcal{A}_3 \to T^*Y$  是横向的。如果定理对于  $(K_1, K_2)$  、  $(K_2, K_3)$  和  $(K_1 \circ_\mu K_2, K_3)$  （分别为  $(K_1, K_2 \circ_\mu K_3)$  ）有效，则它对于  $(K_1, K_2 \circ_\mu K_3)$  （分别为  $(K_1 \circ_\mu K_2, K_3)$  ）有效。这直接从同构  $(K_1 \circ_\mu K_2) \circ_\mu K_3 \simeq K_1 \circ_\mu (K_2 \circ_\mu K_3)$  和公式得出（参见练习 A.9）：

\[
\tau(\lambda_{1}:\lambda_{2})+\tau(\lambda_{1}\circ\lambda_{2}:\lambda_{3})=\tau(\lambda_{1}:\lambda_{2}\circ\lambda_{3})+\tau(\lambda_{2}:\lambda_{3})~.
\]

(e) 如果存在 $A_1 \times_{T \cdot Y} A_2$  的子集 $\Omega$ ，使得 $(p_X, p_Y, p_Z)$  属于 $\Omega$  并且如果该定理在 $\Omega$  的每个点都有效，则该定理在 $(p_X, p_Y, p_Z)$  也有效。事实上，这是从命题 7.5.9 得出的。

令 $p: S \to A_1 \times_{T \cdot Y} A_2$  为连续映射， $p(s_\circ) = (p_X, p_Y, p_Y)$  为 $s_\circ \in S$  。我们写 $p(s) = (p_X(s), p_Y(s), p_Z(s))$ ，然后设置：

\[
\lambda_{1}(s)=T_{(p_{X}(s),p_{Y}(s)^{a})}\varLambda_{1}\quad,\qquad\lambda_{2}(s)=T_{(p_{Y}(s),p_{Z}(s)^{a})}\varLambda_{2}\qquad\mathrm{a n d}
\]

\[
E_{W}(s)=T_{p_{w}(s)}T^{*}W\;,\qquad\lambda_{\circ w}(s)=\lambda_{\circ}(p_{w}(s))\qquad\mathrm{f o r}\;W=X,\;Y,Z\enspace.
\]

我们选择拉格朗日平面 $ \mu_w(s) \subset E_w(s) $  连续依赖 $ s $  ，并且 $ \mu_w(s) \cap \lambda_{\circ w}(s) = 0 $  ( $ W = X, Y, Z $  )，以及 $ (\mu_X(s) \oplus \mu_Y(s)^a) \cap \lambda_1(s) = 0 $  ， $ (\mu_Y(s) \oplus \mu_Z(s)^a) \cap \lambda_2(s) = 0 $  ， $ (\mu_X(s) \oplus \mu_Z(s)^a) \cap (\lambda_1(s) \circ \lambda_2(s)) = 0 $  。然后我们选择

函数  $ d_{1}(s) $  、  $ d_{2}(s) $  、  $ d(s) $  和  $ d_{i}(s_{0}) = d_{i} $  、  $ (i = 1, 2) $  使得以下函数局部常值：

\[
\begin{aligned}&d_{1}(s)-\frac{1}{2}\tau_{E_{X}(s)\oplus E_{Y}(s)^{a}}(\lambda_{\circ X}(s)\oplus\lambda_{\circ Y}(s)^{a},\lambda_{1}(s),\mu_{X}(s)\oplus\mu_{Y}(s)^{a}),\\&d_{2}(s)-\frac{1}{2}\tau_{E_{Y}(s)\oplus E_{Z}(s)^{a}}(\lambda_{\circ Y}(s)\oplus\lambda_{\circ Z}(s)^{a},\lambda_{2}(s),\mu_{Y}(s)\oplus\mu_{Z}(s)^{a}),\\&d(s)-\frac{1}{2}\tau_{E_{X}(s)\oplus E_{Z}(s)^{a}}(\lambda_{\circ X}(s)\oplus\lambda_{\circ Z}(s)^{a},\lambda_{1}(s)\circ\lambda_{2}(s),\mu_{X}(s)\oplus\mu_{Y}(s)^{a}).\\ \end{aligned}
\]

然后根据命题 7.5.9， $ K_{1} $  的类型为  $ L_{1} $ ，在  $ (p_{X}(s), p_{Y}(s)^{a}) $  处具有移位  $ d_{1}(s) $ ， $ K_{2} $  的类型为  $ L_{2} $ ，在  $ (p_{Y}(s), p_{Z}(s)^{a}) $  处有移位  $ d_{2}(s) $ 。

假设  $K_1 \circ_\mu K_2$  的类型为  $L_1 \otimes^L L_2$  ，并且在某个点上有 shift  $d = d_1(s_1) + d_2(s_1) - \frac{1}{2}[\dim Y + \tau(\lambda_{\circ Y}(s_1), \lambda_2(s_1) \circ \lambda_{\circ Z}(s_1), \lambda_{\circ X}(s_1) \circ \lambda_1(s_1)^a)](p_X(s_1), p_Z(s_1)^a)$  。如果  $d(s_1) = d$  ，则  $K_1 \circ_\mu K_2$  的类型为  $L_1 \otimes^L L_2$ ，根据命题 7.5.9，在  $(p_X(s), p_Z(s)^a)$  处有移位  $d(s)$ 。因此，为了证明该定理在 $s = s_\circ$ 处有效，只需证明 $d_1(s) + d_2(s) - d(s) - \frac{1}{2}\tau(\lambda_{\circ Y}(s), \lambda_2(s) \circ \lambda_{\circ Z}(s), \lambda_{\circ X}(s) \circ \lambda_1(s)^a)$ 在 $s$ 处是局部常数即可。这相当于证明下面定义的函数 $\tau(s)$ 是局部常数：

\[
\begin{aligned}{\tau(s)}&{{}=\tau_{E_{X}(s)\oplus E_{Y}(s)^{a}}(\lambda_{\circ X}(s)\oplus\lambda_{\circ Y}(s)^{a},\lambda_{1}(s),\mu_{X}(s)\oplus\mu_{Y}(s)^{a})}\\ {}&{{}\quad+\tau_{E_{Y}(s)\oplus E_{Z}(s)^{a}}(\lambda_{\circ Y}(s)\oplus\lambda_{\circ Z}(s)^{a},\lambda_{2}(s),\mu_{Y}(s)\oplus\mu_{Z}(s)^{a})}\\ {}&{{}\quad-\tau_{E_{X}(s)\oplus E_{Z}(s)^{a}}(\lambda_{\circ X}(s)\oplus\lambda_{\circ Z}(s)^{a},\lambda_{1}(s)\circ\lambda_{2}(s),\mu_{X}(s)\oplus\mu_{Z}(s)^{a})}\\ {}&{{}\quad-\tau_{E_{Y}(s)}(\lambda_{\circ Y}(s),\lambda_{2}(s)\circ\lambda_{\circ Z}(s),\lambda_{\circ X}(s)\circ\lambda_{1}(s)^{a}).}\\ \end{aligned}
\]

通过练习 A.10，我们有：

\[
\tau(s)=\tau_{E_{Y}(s)}(\mu_{Y}(s),\mu_{X}(s)\circ\lambda_{1}(s)^{a},\lambda_{2}(s)\circ\mu_{Z}(s))\enspace.
\]

由于  $\mu_{Y}(s)$  、  $\mu_{X}(s)\circ\lambda_{1}(s)^{\alpha}$  和  $\lambda_{2}(s)\circ\mu_{Z}(s)$  彼此垂直，因此  $\tau(s)$  是局部常数。

选择 $S = \bar{\Omega}$  和 $p$  恒等映射，结果如下。

(f) 当 $ A_1 $  是从 $ T^*X $  到 $ T^*Y $  的接触变换图时，我们将证明该定理。通过(c)，我们可以从一开始就假设 $ Z = \{pt\} $ 。然后，我们将接触变换  $ T^*X \to T^*Y $  分解为  $ T^*X \to T^*X' \to T^*Y $  ，使得  $ T^*X \to T^*X' $  （分别为  $ T^*X'  \to  T^*Y $  ）的图与超曲面  $ S \subset X \times X' $  （分别为  $ S' \subset X' \times Y $  ）的共形丛相关联。根据备注 7.5.8，我们可以假设  $ K_1 = A_S \circ A_S $  。该定理对于  $ (A_S, A_S') $  有效，只要  $ A_1 \to X \times Y $  具有恒定等级（a）。因此，它对于 (e) 处的  $ (A_S, A_S') $  都是有效的。应用(d)和命题7.5.6，我们得到(f)。

(g) 现在我们将在一般情况下证明该定理。

通过(c)，我们可以假设 $ Z = \{\text{pt}\} $ 。然后存在 $ T^*X $  上的接触变换 $ \chi_1 $  和 $ T^*Y $  上的 $ \chi_2 $  使得 $ (\chi_1 \times \chi_2^{-1})(\Lambda_1) $  、 $ \chi_2(\Lambda_2) $  和 $ \chi_1(\Lambda_1 \circ \Lambda_2) $  是超曲面的共形丛（命题A.2.6）。令  $ F_i $  和  $ G_i $  为简单对象，其微支撑分别包含在与  $ \chi_i $  和  $ \chi_i^{-1} $  相关的拉格朗日流形中 (i = 1, 2)，并且满足  $ G_1 \circ_\mu F_1 \simeq A_{\lambda x} $  ，

 $G_{2}\circ_{\mu}F_{2}\simeq A_{\Delta_{\gamma}}$ 。然后 $F_{1}\circ_{\mu}K_{1}\circ_{\mu}G_{2}$  和 $F_{2}\circ_{\mu}K_{2}$  的微支撑包含在超曲面的共形束中。因此，通过 (a) 定理对于  $(F_{1}\circ_{\mu}K_{1}\circ_{\mu}G_{2},F_{2}\circ_{\mu}K_{2})$  是有效的。通过 (f)，该定理对于  $(G_{2},F_{2}\circ_{\mu}K_{2})$  和  $(F_{1}\circ_{\mu}K_{1},G_{2})$  有效。因此，通过 (d)，它对于  $(F_{1}\circ_{\mu}K_{1},K_{2})$  是有效的。再次通过 (f)，该定理对于  $(G_{1},F_{1}\circ_{\mu}K_{1})$  和  $(G_{1},F_{1}\circ_{\mu}K_{1}\circ_{\mu}K_{2})$  有效。因此，通过 (f)，它对于  $(G_{1}\circ_{\mu}F_{1}\circ_{\mu}K_{1},K_{2})$  有效，即  $(K_{1},K_{2})$  。   $\square$ 

请注意，在 (g) 的证明中不可能应用命题 7.5.6，因为该结果仅适用于  $Z = \{pt\}$  时。

作为定理 7.5.11 的应用，我们得到：

\statementlabel{推论 7.5.12}令  $f: Y \to X$  为流形  $p \in Y \times_X T^*X$  、  $p_Y = f'(p)$  、  $p_X = f_\pi(p)$  的态射。令  $A_Y$  为  $T^*Y$  的拉格朗日子流形，使得  $f'$  在  $p_Y$  处横向于  $A_Y$  。让 $G \in \mathrm{Ob}(\mathbf{D}^b(Y))$ 并假设 $\mathrm{SS}(G) \subset A_Y$ 在 $p_Y$ 的邻域上，并且 $G$ 具有类型 $L$ ，并在 $p_Y$ 处沿 $A_Y$ 移动 $d$ 。

(i) 对于  $p$  的足够小的开邻域  $W$  ，  $\Lambda_{X}=f_{\pi}(t f^{\prime-1}(\Lambda_{Y})\cap W)$  是  $f_{\pi}$  与  $t f^{\prime-1}(\Lambda_{Y})\cap W$  同构的拉格朗日流形。

（二） $ f_{1}^{\mu}G = \underbrace{\mathrm{lim}_{U}}_{ \mathrm{type} \quad I. \quad \mathrm{with} \quad \mathrm{shift}} \mathrm{Rf}_{*}G_{U} \text{ exists in } \mathbb{D}^{b}(X; p_{X}) \text{(cf. 命题 6.1。10) and } f_{1}^{\mu}G \text{ has} $ 

L 型带移位

\[
d-\tfrac{1}{2}(\dim Y/X+\tau(\lambda_{\circ}(p_{X}),\lambda_{\wedge_{X}}(p_{X}),f_{\pi}^{\intercal}f^{\prime-1}(\lambda_{\circ}(p_{Y}))))~.
\]

这里 $U$ 涵盖了 $\pi_{Y}(p_{Y})$ 的开放邻里系统，我们写了 $f_{\pi}^{\prime}f^{\prime-1}(\lambda_{\circ}(p_{Y}))$ 而不是 $(\mathrm{d}f_{\pi})(\mathrm{d}^{\prime}f^{\prime})^{-1}(\lambda_{\circ}(p_{Y}))$ 。

\statementlabel{证明}(i) 由练习 A.5 得出。

(ii) 将定理 7.5.11 应用于  $ K_1 = A_d $  、  $ K_2 = G $  ，其中  $ \Delta \subset X \times Y $  是  $ f $  的图。   $ \square $ 

\statementlabel{推论 7.5.13}令  $f: Y \to X$  为流形  $p \in Y \times_X T^*X$  、  $p_Y = f'(p)$  、  $p_X = f_\pi(p)$  的态射。令  $\mathcal{A}_X$  为  $T^*X$  的拉格朗日子流形，使得  $f_\pi$  在  $p_X$  处横向于  $\mathcal{A}_X$  。让 $F \in \mathrm{Ob}(\mathbf{D}^b(X))$ 并假设 $p_X$ 和 $F$ 邻域上的 $\mathrm{SS}(F) \subset \mathcal{A}_X$ 具有类型 $L$ ，并在 $p_X$ 处沿 $\mathcal{A}_X$ 移动 $d$ 。

(i) 对于  $p$  的足够小的开邻域  $W$  ，  $A_{Y}=t f^{\prime}(f_{\pi}^{-1}(A_{X})\cap W)$  是拉格朗日流形，通过  $t f^{\prime}$  同构于  $f_{\pi}^{-1}(A_{X})\cap W$  。

(ii)  $ f_{\mu}^{-1}F = \lim_{\overrightarrow{r_{i}}\rightarrow\overrightarrow{r_{n}}}f^{-1}F' $  和  $ f_{\mu}^{!}F = \lim_{\overrightarrow{x}\rightarrow\overrightarrow{x}^{\prime}}f^{!}F' $  存在于  $ \mathbf{D}^{b}(Y; p_{Y}) $  中，其中  $ F^{\prime}\to F $ 

（分别为  $F \to F'$  ）范围涵盖  $p_X$  和  $\omega_{Y/X} \otimes f_\mu^{-1} F \simeq f_\mu^! F$  处的同构  $s$  范畴（参见命题 6.1.9）。

(iii)  $ f_{\mu}^{-1}F $  的类型为 L，移位为 d。

\statementlabel{证明}(i) 由练习 A.5 得出。

(ii) 是命题 6.1.9 的结果。

(iii) 将定理 7.5.11 应用于  $ K_1 = A_d $  和  $ K_2 = F $  ，其中  $ \Delta \subset Y \times X $  是  $ f $  的图。   $ \square $ 

\subsection{习题}
\statementlabel{练习 七.1}令 $V$  为实数（或复数）有限维向量空间， $V^*$  为对偶向量空间。设置 $S_V = (V \setminus \{0\}) / \mathbb{R}^+$ （分别为 $(V \setminus \{0\}) / \mathbb{C}^*$ ）并让 $Z = \{(x, y) \in S_V \times S_{V^*} : \langle x, y \rangle = 0\}$ 。让 $K = A_Z$ ， $\Omega = \dot{T}^*S_V$ ， $\Omega' = \dot{T}^*S_{V^*}$ 。证明定理 7.2.1 的假设对于 $(\Omega, \Omega')$  上的 $K$  是满足的（参见 Brylinski [2]，详细研究了复流形和可构造层框架中的这种接触变换）。

\statementlabel{练习 七.2}令  $ \tau: E \to Z $  为流形  $ Z $  上的向量丛。如命题 5.5.1 中所示，确定  $ T^*E $  和  $ T^*E^* $ 。让  $ F_1 $  和  $ F_2 $  属于  $ \mathbf{D}_R^b(E) $  。证明同构：

\[
\mu\mathrm{h o m}(F_{2},F_{1})\simeq\mu\mathrm{h o m}(F_{2}^{\wedge},F_{1}^{\wedge})\enspace.
\]

（提示：使用练习 III 17 和 IV 6，表明它相当于在  $ \mathbf{D}_{\mathbf{R}_+}^b(\dot{E}) $  或  $ \mathbf{D}^b(\dot{E}/\mathbb{R}^+) $  中证明  $ F_1 $  和  $ F_2 $  的类似公式。然后使用定理 7.2.1。）

\statementlabel{练习 VII.3}令 $(x)$  和 $(y)$  分别表示 $\mathbb{R}^n$  的两个副本 $X$  和 $Y$  上的两个线性坐标系，并让 $(x;\xi)$  和 $(y;\eta)$  为相关坐标。考虑为  $\xi_n \neq 0$  、  $\eta_n \neq 0$  通过  $y_j = \xi_j \xi_n^{-1}$  (  $p < j < n$  )、  $y_n = (\sum_{j=p+1}^n x_j \xi_j) \xi_n^{-1}$  、  $y_k = x_k$  (  $k \leq p$  )、  $\eta_j = -x_j \xi_n$  (  $p < j < n$  )、  $\eta_n = \xi_n$  定义的接触变换  $\chi$  (部分勒让德变换)，  $\eta_k = \xi_k$ （ $k \leq p$ ）。

(i) 证明 $\chi$  与子流形 $S = \{(x,y); x_k - y_k = 0 \; 1 \leqslant k \leqslant p, x_n - y_n + \sum_{j=p+1}^{n-1} x_j y_j = 0\}$  的共形束相关联。

(11) 通过应用定理7.2.1，从微观角度证明， $ \psi_{K}(A_{M}) \simeq A_{N} $  ，其中 $ N = \{y_{n} = 0\} $  、 $ M = \{x_{p+1} = \cdots = x_{n} = 0\} $  和 $ K = A_{S}[d] $  用于计算要计算的移位d。

\statementlabel{练习 七.4}令 $X$  为复流形，并令 $A$  为 $T^*X$  的复连通拉格朗日子流形。设  $F \in \mathrm{Ob}(\mathbf{D}^b(X))$  ，并假设  $\mathrm{SS}(F) \subset A$  在  $A$  的邻域中。证明如果  $F$  在  $p \in A$  处通过 shift  $d$  是纯的，则  $F$  在任意  $p \in A$  处通过 shift  $d$  是纯的。 （提示：使用命题 7.5.9 和练习 A.7。）

\statementlabel{练习 七.5}在定理 7.5.11 的情况下，假设  $X$  、 $Y$  、 $Z$  和  $A_i$  是复流形， $K_i$  的类型为  $L_i$  且带有移位  $0$  (  $i = 1,2$  )。证明  $K_1 \circ_\mu K_2$  具有类型  $L_1 \otimes^L L_2$  和 shift  $-\dim_C Y$  。

\subsection{注记}
辛几何可以追溯到汉密尔顿和雅可比，物理学家习惯于在“相空间”中工作，即在余切丛中（参见优秀的

Guillemin-Sternberg 的介绍[1]）。然而，直到最近，数学工具才出现，这使得超越纯粹几何的观点成为可能，并在余切丛（一种说法是“微局部”）中与附加到流形 X 的经典对象（例如分布或微分算子）一起工作。第一步是引入“伪微分算子”（这里不回顾这个理论，参考Hörmander [4]），然后是Maslov [1]的“规范算子”，最后是Sato [2]的微函数理论和Hörmander [2]的傅里叶积分算子理论，它特别使我们能够对那些经典对象进行接触变换。

我们在本章中提出的理论是傅里叶积分算子理论的层理论版本，或其复杂的等价物，Sato-Kawai-Kashiwara 的简单完整模理论[1]。它首先在 Kashiwara-Schapira 中通过不同的方法引入[3]。

\section{构造层}
\subsection{概要}
我们首先解释单纯复形上的可构造层的概念，然后回顾有关 Hironaka 子解析集的一些基本事实。

接下来我们介绍以下定义，这是惠特尼分层定义的修改。如果  $ X_\alpha $  是亚解析子流形，则实解析流形  $ X $  的分层  $ X = \bigsqcup_{\alpha \in A} X_\alpha $  是  $ \mu $  分层，并且对于每对  $ (\alpha, \beta) $  和  $ X_\beta \cap \overline{X}_\alpha \neq \emptyset $  ，我们有  $ X_\beta \subset \overline{X}_\alpha $  并且还有：

\[
(T_{X_{\beta}}^{*}X+\stackrel{\frown}{T_{X_{\alpha}}^{*}X})\cap\pi^{-1}(X_{\beta})\subset T_{X_{\beta}}^{*}X\enspace.
\]

（回想一下 VI §2 的运算  $ \hat{+} $ 。）证明给定  $ T^*X $  的闭圆锥亚解析各向同性子集  $ A $  ，存在  $ \mu $  -分层  $ X = \bigsqcup_{\alpha} X_\alpha $  使得  $ A \subset \bigsqcup_{\alpha} T_X^*X $  。

接下来我们介绍以下定义。如果存在由子解析子集局部有限覆盖 $X = \bigcup_j X_j$ ，使得对于所有 $k$  和所有 $j$  而言，层 $H^k(F)|_{X_j}$  是局部常值的，那么 $\mathbf{D}^b(X)$  的对象 $\mathbb{R}$  是弱 $\mathbb{R}$  可构造的（简称w-  $\mathbb{R}$  可构造的）。此外，如果复合体  $F_x$  对于所有  $x \in X$  都是完美的，那么有人说  $F$  是  $\mathbb{R}$  可构造的。

利用 $\mu$  -分层的存在性，我们证明 $F$  是 $w$  -  $\mathbb{R}$  -可构造的当且仅当 $SS(F)$  包含在闭圆锥亚解析各向同性集中，或者等效地如果 $SS(F)$  是亚解析且拉格朗日。换句话说， $w$  -  $\mathbb{R}$  - 可构造是一个微局部属性。然后我们可以充分利用前面几章的结果，立即获得 $w$  -  $\mathbb{R}$  和 $\mathbb{R}$  可构造对象的各种函子性质。利用亚解析集三角剖分的存在性，我们还证明了  $\mathbb{R}$  可构造层的导出范畴等价于由具有  $\mathbb{R}$  可构造上同调的对象组成的  $\mathbf{D}^{b}(X)$  的满子范畴。

当X是复流形时，还可以引入w-C和C-可构造对象的概念。定义类似，将“亚分析子流形”替换为“复分析子流形”。一个有用的结果断言 F 是 w-C-可构造的，当且仅当 F 是 w-R-可构造的（在实底层流形上），而且 SS(F) 通过  $ C^x $  对  $ T^*X $  的作用是不变的。以此

根据定理，C 构造层的许多性质都是从 R 构造层的性质推导出来的。

最后我们介绍了著名的“近循环函子”和“消失循环函子”，并将它们与特化函子和微局部化函子进行了比较。

公约8.0。我们保持惯例 4.0。特别是基环  $A$  满足  $\operatorname{gld}(A) < \infty$  。此外，除非另有说明，所有流形和流形态射在无穷大处都是实解析且可数的。

其中一个用  $\dim X$  （分别为  $\dim_{C} X$  ）表示实数（分别为复数）解析子集  $X$  的维数。

\subsection{单纯复形上的构造层}
\statementlabel{定义 8.1.1}单纯复形  $ \mathbf{S} = (S, \varDelta) $  是由  $ S $  的集合  $ S $  和  $ S $  的子集的集合  $ \varDelta $  组成的数据，满足以下公理。

(S.1) 任何  $ \sigma \in \Delta $  都是 S 的有限非空子集。

(S.2) 如果  $ \tau $  是  $ \Delta $  的元素  $ \sigma $  的非空子集，则  $ \tau $  属于  $ \Delta $  。

(S.3) 对于任何  $ p \in S $  ，  $ \{p\} $  属于  $ \Delta $  。

(S.4) 对于任何  $ p \in S $  ，集合  $ \{\sigma \in \Delta; p \in \sigma\} $  是有限的。

 $\Delta$  的元素称为单纯形， $S$  的元素称为顶点。

让  $ \mathbb{R}^S $  表示从  $ S $  到  $ \mathbb{R} $  的映射集。回想一下，元素  $ x \in \mathbb{R}^S $  只是一个族  $ x(p) \in \mathbb{R} $  ，由  $ p \in S $  索引。我们为 $ \mathbb{R}^S $ 配备了产品拓扑。对于 $ \sigma \in \mathcal{A} $ ，我们通过以下方式将 $ \mathbb{R}^S $  的子集 $ |\sigma| $  关联起来：

(8.1.1) $ |\sigma| = \left\{x \in \mathbb{R}^S; x(p) = 0 \text{ for } p \notin \sigma, x(p) > 0 \text{ for } p \in \sigma \text{ and } \sum_p x(p) = 1\right\} $ 

请注意， $ |\sigma| $  彼此不相交。还规定：

\[
|\mathbf{S}|=\bigcup_{\sigma\in\varDelta}|\sigma|~,
\]

\[
U(\sigma)=\bigcup_{\tau\in\varDelta,\tau\supset\sigma}|\tau|\;,
\]

对于  $ x \in |S| $  ：

\[
U(x)=U(\sigma(x)),
\]

其中  $ \sigma(x) $  是唯一的单纯形，使得  $ x \in |\sigma(x)| $  ，即  $ \sigma(x) = \{p; x(p) \neq 0\} $  。

\[
U(\sigma)=\{x\in|\mathbf{S}|;x(p)>0\mathrm{f o r}p\in\sigma\}\enspace,
\]

 $ U(x) = \{ y \in |S|; y(p) > 0 $  对于任何  $ p $  使得  $ x(p) > 0\} $  。

示例 8.1.2。取  $S = \{1,2,3\}$  ，并令  $\mathcal{A}$  为  $S$  的非空子集族。人们可以通过下面的图片可视化集合  $|\sigma|$  、  $U(\sigma)$  和  $U(x)$  ，其中  $\sigma = \{1,2\}$  、  $x = (1,0,0)$  。

\begin{center}
\includegraphics[width=0.19\textwidth]{流行上的层_Chn-assets/img_in_image_box_118_321_349_565.png}
\end{center}

\begin{center}
\includegraphics[width=0.47\textwidth]{流行上的层_Chn-assets/img_in_image_box_473_320_1044_556.png}
\end{center}

\begin{center}图。 8.1.1\end{center}

我们赋予 $ |\mathbf{S}| $ 来自 $ \mathbb{R}^S $ 的归纳拓扑。那么 $ U(\sigma) $  和 $ U(x) $  是 $ |\mathbf{S}| $  的开子集，如(8.1.5)、(8.1.6) 所示。定义：

\[
S(\sigma)=\left\{p\in S;\left\{p\right\}\cup\sigma\in\varDelta\right\}.
\]

那么  $S(\sigma)$  是  $S$  的公理 (S.4) 的有限子集，并且  $U(\sigma)$  包含在  $\mathbb{R}^{S(\sigma)}$  中，因此  $U(\sigma)$  同胚于  $\mathbb{R}^{l}$  的局部闭子集  $l = \#S(\sigma)$  。由于  $|\sigma| = \{x \in U(\sigma); x(p) = 0$  对于  $p \notin \sigma\}$  ，  $|\sigma|$  是  $U(\sigma)$  的闭子集，因此是  $|\mathbf{S}|$  的局部闭子集。人们还可以轻松验证：

\[
\begin{aligned}|\overline{\sigma}|&=\left\{x\in\mathbb{R}^{S};x(p)=0\text{if}p\notin\sigma,x(p)\geqslant0\text{and}\sum_{p}x(p)=1\right\}\\&=\bigcup_{\tau\subset\sigma}|\tau|,\\ \end{aligned}
\]

其中  $ |\sigma| $  表示  $ |\mathbf{S}| $  中  $ |\sigma| $  的闭包。另请注意，对于  $ \Delta $  中的  $ \sigma $  和  $ \tau $  ：

\[
\left\{\begin{aligned}{U(\sigma)\cap U(\tau)}&{{}=U(\sigma\cup\tau)}&{}&{{}\mathrm{i f}\sigma\cup\tau\in\varDelta}\\ {}&{{}}&{}\\ {}&{{}=\varnothing}&{}&{{}\mathrm{i f}\sigma\cup\tau\notin\varDelta~.}\\ \end{aligned}\right.
\]

\[
U(\sigma)\subset U(\tau)\qquad\mathrm{~i f~a n d~o n l y~i f~}\tau\subset\sigma.
\]

公理 (S.4) 与 (8.1.9) 一起意味着族  $ \{U(\sigma)\}_{\sigma \in \mathcal{A}} $  是  $ |S| $  的局部有限开覆盖。因此 $ |S| $  是一个仿紧拓扑空间。

\statementlabel{定义 8.1.3}让 $ F \in \mathrm{Ob}(\mathbb{D}^{b}(|\mathbf{S}|)) $  。

(i) 如果对于所有 $j \in \mathbb{Z}$  和所有 $\sigma \in \Delta$  ，层 $H^j(F)|_{|\sigma|}$  是恒定的，则可以说 $F$  是弱 $\mathbf{S}$  可构造的（简称为 $\mathbf{S}$  可构造的）。

(ii) 如果 $F$  是w-S-可构造的，并且对于每个 $x \in X$  ， $F_x$  是完美复形，则可以说 $F$  是S-可构造的。

（对于完美复形的概念，参见练习 I 30，并回想一下，如果  $A$  是诺特式的，则  $F_x$  是完美的，当且仅当对所有  $j \in \mathbb{Z}$  而言， $H^j(F)_x$  是有限生成的。回想一下，我们假设  $\mathrm{gld}(A) < \infty$  。）

如果是这样，则束  $F$  被称为 w-S-可构造（分别为 S-可构造），并被视为  $\mathbf{D}^{b}(|\mathbf{S}|)$  的对象。其中一个用 $\mathbf{D}_{w-S-c}^{b}(|\mathbf{S}|)$ （或 $\mathbf{D}_{S-c}^{b}(|\mathbf{S}|)$ ）表示 $\mathbf{D}^{b}(|\mathbf{S}|)$ 的完整三角子范畴，由w-S-可构造（或S-可构造）对象组成。

其中一个用w-Cons(S)（或Cons(S)）表示 $ \mathrm{Mod}(A_{|S|}) $ 的满子范畴，由w-S-可构造（或S-可构造）层组成。

令 $ u: F \to G $  为 $ |S| $  、 $ F $  和 $ G $  上的层态射， $ w $  可-S-构造。然后立即证明 Ker  $ u $  、 Im  $ u $  、 Coker  $ u $  是  $ w $  -S-可构造的。此外，如果  $ 0 \to F' \to F \to F'' \to 0 $  是  $ |S| $  上的层的精确序列，并且  $ F' $  和  $ F'' $  是  $ w $  -S-可构造的，那么  $ F $  也是如此。因此 $ w $  -Cons(S) 是一个阿贝尔范畴。

如果基环  $ A $  是诺特环，则  $ \text{Cons}(S) $  的结果相同。

\statementlabel{命题 8.1.4}假设  $F$  是一个 w-S-可构造的层。那么对于任何  $\sigma \in \Delta$  和任何  $x \in |\sigma|$  我们都有同构：

\[
H^{0}(U(\sigma);F)\xrightarrow{\sim}H^{0}(|\sigma|;F)\xrightarrow{\sim}F_{x}\enspace,
\]

\[
H^{j}(U(\sigma);F)=H^{j}(|\sigma|;F)=0\qquad{f o r~j\neq0~}.
\]

\statementlabel{证明}对于 $0 < \varepsilon < 1$ ，设置 $I_\varepsilon = \{t \in \mathbb{R}; \varepsilon \leqslant t \leqslant 1\}$ 并通过以下方式定义映射 $\pi_\varepsilon: I_\varepsilon \times U(\sigma) \to U(\sigma)$ ：

\[
\pi_{\varepsilon}(t,y)(p)=t y(p)+(1-t)x(p),\qquad\mathrm{f o r~}p\in S.
\]

映射  $\pi_{\varepsilon}$  是连续且满射的， $\pi_{\varepsilon}(1,\cdot)$  是恒等式  $\{1\}\times U(\sigma)\simeq U(\sigma)$  ， $\pi_{\varepsilon}(\varepsilon,\cdot)$  是同胚  $\{\varepsilon\}\times U(\sigma)\xrightarrow{\sim}\pi_{\varepsilon}(\{\varepsilon\}\times U(\sigma))$  。

用 $q_2$  表示从 $I_\varepsilon \times U(\sigma)$  到 $U(\sigma)$  的投影。由于  $y \in |\tau|$  意味着  $\pi_\varepsilon(t,y) \in |\tau|$  ，我们得到  $\pi_\varepsilon^{-1}F$  是  $q_2$  纤维上的常值层，因此  $\pi_\varepsilon^{-1}F = q_2^{-1}G$  ，对于  $U(\sigma)$  上的某些束  $G$  。 （以  $G = q_{2*} \pi_\varepsilon^{-1}F$  为例。）应用推论 2.7.7，我们得到对于每个  $t \in I_\varepsilon$  ，态射  $R\Gamma(I_\varepsilon \times U(\sigma); \pi_\varepsilon^{-1}F) \to R\Gamma(\{t\} \times U(\sigma); \pi_\varepsilon^{-1}F)$  是同构。因此：

\[
\begin{aligned}{R\varGamma(U(\sigma);F)}&{{}\simeq R\varGamma(\{\varepsilon\}\times U(\sigma);\pi_{\varepsilon}^{-1}F)}\\ {}&{{}\simeq R\varGamma(\pi_{\varepsilon}(\{\varepsilon\}\times U(\sigma));F)\;.}\\ \end{aligned}
\]

现在我们注意到， $\{\pi_\varepsilon(\{\varepsilon\}\times U(\sigma))\}_{\varepsilon>0}$ 家族在 $X$ 中形成了 $x$ 的邻里系统。取 (8.1.12) 两边的上同调和  $\varepsilon>0$  的归纳极限，我们得到  $H^j(U(\sigma); F)$  对于  $j>0$  为零，对于  $j=0$  与  $F_x$  同构。如果我们将相同的参数应用于单纯复形  $\sigma$  和层  $F|_{\sigma|}$  ，我们会得到相同的结果，将  $U(\sigma)$  替换为  $|\sigma|$  ，从而完成证明。   $\square$ 

\statementlabel{推论 8.1.5}(i) 函子 $ \Gamma(U(\sigma); \cdot) $  精确于范畴 $ w $  -Cons(S)。

(ii) 如果  $F$  是 w-S-可构造层，并且  $\Gamma(U(\sigma); F) = 0$  对于所有  $\sigma \in \Delta$  ，则  $F = 0$  。

假设  $|\mathbf{S}|$  上的一个层  $F$  是 S-无环的，如果它满足所有  $k > 0$  和所有  $\sigma \in \mathcal{A}$  的  $H^k(U(\sigma); F) = 0$  。那么根据命题 8.1.4，w-S-可构造层是 S-非循环的。当然， $|\mathbf{S}|$  上的松弛层也是 S-无环的。

让 $F$  成为 $|\mathbf{S}|$  上的一个层。我们将在功能上将  $F$  与  $w$  -S-可构造束  $\beta(F)$  和态射  $s(F):\beta(F) \to F$  关联起来。为此，首先定义

\[
\alpha(F)=\bigoplus_{\sigma\in\varDelta}\left(F(U(\sigma))\right)_{U(\sigma)}.
\]

如果  $M$  是  $A$  模，则  $\mathrm{Hom}(M_{U(\sigma)}, F) = \mathrm{Hom}(M, \Gamma(U(\sigma); F))$  。因此，存在自然态射  $i(F): \alpha(F) \to F$  ，并且该态射对于  $F$  是函数式的。令 $F'$  表示 $i(F)$  的核，并令 $\beta(F)$  表示由复合 $\alpha(F') \to F' \to \alpha(F)$  获得的态射 $\alpha(F') \to \alpha(F)$  的核。我们得到一个态射  $s(F): \beta(F) \to F$  ，它是关于  $F$  的函子。

\begin{center}
\includegraphics[width=0.54\textwidth]{流行上的层_Chn-assets/img_in_image_box_243_785_892_1223.png}
\end{center}

请注意， $ \alpha(F) $  和 $ \beta(F) $  是可W-S 构造的。

\statementlabel{引理 8.1.6}对于  $ \sigma \in \Delta $  ，自然态射  $ \Gamma(U(\sigma); \beta(F)) \to \Gamma(U(\sigma); F) $  是同构。

\statementlabel{证明}考虑一下图表：

\begin{center}
\includegraphics[width=0.75\textwidth]{流行上的层_Chn-assets/img_in_image_box_113_1510_1012_1690.png}
\end{center}

由于  $\Gamma(U(\sigma); \cdot)$  是  $w$  -Cons(S) 上的精确函子，因此顶行是精确的。由于态射 $\alpha(F)(U(\sigma)) \to F(U(\sigma))$  是满射的，所以底行是精确的。那么 $\alpha(F^{\prime})(U(\sigma)) \to F^{\prime}(U(\sigma))$  的满射性就意味着期望的结果。   $\square$ 

\statementlabel{引理 8.1.7}假设 $F$  是弱S 可构造的。那么 $s(F):\beta(F)\to F$  是同构。

\statementlabel{证明}层 $ \operatorname{K e r} s(F) $  和 $ \operatorname{C o k e r} s(F) $  是w-S-可构造的，并且根据引理8.1.6 和推论8.1.5 (i)，它们在 $ U(\sigma) $  上的截面为零。那么根据推论 8.1.5 (ii)，这些层为零。   $ \Box $ 

根据此引理，  $ \beta: \mathfrak{Mod}(A_{|S|}) \to w\text{-Cons}(S) $  是包含函子  $ w\text{-Cons}(S) \to \mathfrak{Mod}(A_{|S|}) $  的右伴随函子。

\statementlabel{命题 8.1.8}(i) 从 $ \mathscr{S}\mathfrak{h}(|\mathscr{S}|) $  到 $ w $  的函子 $ \beta $  -Cons( $ S $  ) 保持精确。

(ii)  $ \Gamma(U(\sigma); R^{k}\beta(F)) = H^{k}(U(\sigma); F) $  为所有 $ \sigma \in \Delta $  、 $ k \geqslant 0 $  、F 捆 $ |S| $  。

(iii)  $ R^k\beta(F) = 0 $  对于任何  $ k > 0 $  当且仅当  $ F $  是  $ \mathbf{S} $  -非循环。

\statementlabel{证明}(i) 令 $0 \to F' \to F \to F'' \to 0$  为 $|\mathbf{S}|$  上层的精确序列。层  $\beta(F'), \beta(F), \beta(F'')$  是 w-S-可构造的，为了看出序列  $0 \to \beta(F') \to \beta(F) \to \beta(F'')$  是精确的，鉴于推论 8.1.5 (ii)，足以表明对于所有  $\sigma \in \mathcal{A}$  ，序列 $0 \to \varGamma(U(\sigma); \beta(F')) \to \varGamma(U(\sigma); \beta(F)) \to \varGamma(U(\sigma); \beta(F''))$ 是精确的。这由引理 8.1.6 和函子 $\varGamma(U(\sigma);\cdot)$ 的左正合性推出。

(ii) 令 $F^{*}$ 为 $F$ 的内射分解。然后 $R^{k}\beta(F)=H^{k}(\beta(F^{\cdot}))$ 。此外：

\[
\begin{aligned}\varGamma(U(\sigma);R^{k}\beta(F))&=\varGamma(U(\sigma);H^{k}(\beta(F^{*})))\\&=H^{k}(\varGamma(U(\sigma);\beta(F^{*})))\\&=H^{k}(\varGamma(U(\sigma);F^{*}))\\&=H^{k}(U(\sigma);F)\enspace.\\ \end{aligned}
\]

（我们使用了 w-S-可构造层是 S-非循环这一事实，以及引理 8.1.6。）

(iii) 由 (ii) 和推论 8.1.5 (ii) 得出。 □

\statementlabel{命题 8.1.9}令  $F'$  为从 S-无环层下方界定的复数。假设对于所有  $n \in \mathbb{Z}$  ，  $H^n(F')$  都是 w-S-可构造的。那么态射 $\beta(F') \to F'$  是拟同构。

\statementlabel{证明}令  $d'$  表示复数  $F^*,$  的微分（即：  $d^n: F^n \to F^{n+1}$  ），并设置  $Z^n(F^*) = \mathrm{Ker} d^n, B^n(F^*) = \mathrm{Im} d^{n-1}$  。我们得到确切的序列：

\[
0\to Z^{n-1}(F^{\cdot})\to F^{n-1}\to B^{n}(F^{\cdot})\to0,
\]

(8.1.15)

\[
0\to B^{n}(F^{\cdot})\to Z^{n}(F^{\cdot})\to H^{n}(F^{\cdot})\to0~.
\]

假设 $ Z^{n-1}(F') $  是S-无环的。将函子 $ \beta $ 应用于(8.1.14)，我们发现 $ R^k\beta(B^n(F')) = 0 $ 对于 $ k > 0 $ ，这意味着 $ B^n(F') $ 根据命题8.1.8是S-无环的。应用于序列 (8.1.15) 的相同参数给出  $ Z^n(F') $  是 S-无环的。通过归纳，这证明所有 $ Z^n(F') $  和 $ B^n(F') $  都是S-无环的。

函子  $\beta$  保持精确，我们有  $\beta(Z^n(F^c)) = Z^n(\beta(F^c))$  。由于  $R^1\beta(Z^{n-1}(F^c)) = 0$  ，序列 (8.1.14) 产生  $\beta(B^n(F^c)) = B^n(\beta(F^c))$  。类似地，由于  $R^1\beta(B^n(F^c)) = 0$  ，序列 (8.1.15) 产生  $\beta(H^n(F^c)) = H^n(\beta(F^c))$  。根据假设和引理 8.1.7， $\beta(H^n(F^c)) \simeq H^n(F^c)$  。因此我们证明了同构 $H^n(\beta(F^c)) \simeq H^n(F^c)$ 。   $\square$ 

为了在 $ \mathbf{D}^{b}(|\mathbf{S}|) $ 范畴中工作，我们将做出以下假设。

\[
\left\{\begin{matrix}{{T h e r e~e x i s t s~a n~i n t e g e r~}n{~s u c h~t h a t}}\\ {\#{\sigma}\leqslant n+1{~f o r\quad a l l~}{\sigma}\in\varDelta}\\ \end{matrix}\right..
\]

我们称满足（8.1.16）的最小 $n$ 为S的维数。在假设（8.1.16）下，鉴于命题8.1.8和3.2.2，我们有 $R^{k}\beta = 0$ 为 $k > n$ 。然后考虑函子（其中  $\delta$  是自然函子）：

\[
\begin{array}{r}{w\mathrm{-}\mathbf{C o n s}(\mathbf{S})\xleftarrow[\beta]{\delta}\mathfrak{M o d}(A_{|\mathbf{S}|})\enspace.}\end{array}
\]

他们诱导函子

\[
\mathbf{D}^{b}(w\mathbf{-C o n s}(\mathbf{S}))\overset{\delta}{\xleftrightarrow[R\beta]{\phantom{0}}}\mathbf{D}_{w-\mathbf{S}-c}^{b}(|\mathbf{S}|)\enspace.
\]

\statementlabel{定理 8.1.10}假设(8.1.16)。那么(8.1.17)中的函子 $ \delta $ 和 $ R\beta $ 是范畴的等价，彼此相反。

\statementlabel{证明}$ \mathbf{D}^{b}(w\text{-}\mathbf{Cons}(\mathbf{S})) $  中的同构  $ R\beta \circ \delta \simeq id $  源自  $ w\text{-}\mathbf{Cons}(\mathbf{S}) $  中的同构  $ \beta \circ \delta \simeq id $  。

 $ \mathbf{D}_{w-s-c}^b(|\mathbf{S}|) $  中的同构  $ \delta \circ R\beta \simeq id $  由命题 8.1.9 得出。

\statementlabel{定理 8.1.11}假设(8.1.16) 并假设基环 $ A $  是诺特环。那么自然函子 $ \delta : \mathbf{D}^b(\text{Cons}(\mathbf{S})) \to \mathbf{D}_{\mathbf{S}-c}^b(|\mathbf{S}|) $  就是范畴的等价。

\statementlabel{证明}由命题1.7.11（和备注1.7.12）和定理8.1.10，足以证明以下引理。

\statementlabel{引理 8.1.12}令 $u: F \to G$  为w-S-可构造层的外同态，并假设 $G$  是S-可构造的。那么存在一个可构造S的层 $H$ 和一个态射 $v: H \to F$ ，使得 $u \circ v$ 是一个外态。

\statementlabel{证明}对于任何  $\sigma \in \varDelta$  ，  $F(U(\sigma)) \to G(U(\sigma))$  根据推论 8.1.5 是满射的。由于  $G(U(\sigma)) \simeq G_x$  对于  $x \in |\sigma|$  ，根据命题 8.1.4， $G(U(\sigma))$  是有限生成的。因此，存在有限生成的自由模  $H(\sigma)$  和外同态  $H(\sigma) \to G(U(\sigma))$  。这个态射分裂，我们得到态射  $H(\sigma) \to F(U(\sigma)) \to G(U(\sigma))$  。设置 $H = \bigoplus_{\sigma} H(\sigma)_{U(\sigma)}$ 。那么  $H$  是一个  $\mathbf{S}$  可构造的层，态射  $H(\sigma) \to F(U(\sigma))$  定义了态射  $H \to F$  ，复合  $H \to F \to G$  是一个外态射。   $\square$ 

请注意，通过这一论证，我们可以看出，w-Cons(S) 具有足够的射影，并且如果 A 是诺特式的，则范畴 Cons(S) 也具有足够的射影。

\subsection{次解析集}
在本节中，我们回顾 Hironaka 亚解析集的定义及其主要性质，但没有证明。参考 Hironaka [1,2]、Hardt [1,2] 或 Bierstorne-Milman [1]。

设 X 为流形（回想一下，所有流形都假定为实解析流形），并设 Z 为 X 的子集。

\statementlabel{定义 8.2.1}有人说  $Z$  在  $x \in X$  处是亚解析的，如果存在  $x$  的开邻域  $U$  、紧流形  $Y_j^i$  (  $i = 1,2,1 \leq j \leq N$  ) 和态射  $f_j^i: Y_j^i \to X$  ，使得：

\[
Z\cap U=U\cap\bigcup_{j=1}^{N}(f_{j}^{1}(Y_{j}^{1})\backslash f_{j}^{2}(Y_{j}^{2}))~.
\]

如果  $Z$  在每个  $x \in X$  都是亚解析的，则可以说  $Z$  在  $X$  是亚解析的。

亚解析集继承以下属性。

\statementlabel{命题 8.2.2}(i) 假设  $Z$  在  $X$  中是亚解析的。然后  $\bar{Z}$  和 Int(  $Z$  ) 在  $X$  中进行亚解析。此外， $Z$  的连通分量是局部有限的和亚解析的。

(ii) 假设  $Z_1$  和  $Z_2$  在  $X$  中是亚解析的。那么  $Z_1 \cup Z_2$  、  $Z_1 \backslash Z_2$  、  $Z_1 \cap Z_2$  是亚解析的。

(iii) 令 $f: Y \to X$  为流形的态射。如果  $Z \subset X$  在  $X$  中是亚解析的，那么  $f^{-1}(Z)$  在  $Y$  中是亚解析的。如果  $W \subset Y$  在  $Y$  中是亚解析且  $f$  在  $\overline{W}$  上是真映射，则  $f(W)$  在  $X$  中是亚解析。

(iv) 令  $Z$  为  $X$  的封闭子解析子集。那么存在流形  $Y$  和真态射  $f: Y \to X$  使得  $f(Y) = Z$  。

现在我们陈述曲线选择引理。

\statementlabel{命题 8.2.3}令  $Z$  为  $X$  的子分析子集，并令  $x_\circ \in \overline{Z}$  。然后存在解析曲线  $t \mapsto x(t)$  、  $] - 1, 1[ \to X$  ，使得  $x(0) = x_\circ$  和  $x(t) \in Z$  对于  $t \neq 0$  。

实际解析情况中存在去奇异化定理。

\statementlabel{命题 8.2.4}令  $\varphi: X \to \mathbb{R}$  为实解析函数，在  $X$  的每个连通分量上不完全相同为零。设置 $Z = \{x \in X; \varphi(x) = 0, \mathrm{d}\varphi(x) = 0\}$ 。然后存在流形  $f: Y \to X$  的真态射，它导致同构  $Y \setminus f^{-1}(Z) \simeq X \setminus Z$  并且使得在每个  $y_0 \in f^{-1}(Z)$  的邻域中存在一个局部坐标系  $(y_1, \ldots, y_n)$  和  $\varphi \circ f = \pm y_1^{r_1} \ldots y_n^{r_n}$  ，  $r_j$  是非负整数。

令  $Z$  为  $X$  的子分析子集。将  $Z_{\text{reg}}$  定义为点  $x \in Z$  的子集，使得  $X$  中存在  $x$  的开邻域  $U$  ，使得  $U \cap Z$  是  $U$  的闭子流形。一组 $Z_{\text{sing}} = Z \setminus Z_{\text{reg}}$ 。然后集合  $Z_{\text{reg}}$  和  $Z_{\text{sing}}$  在  $X$  中进行亚解析，并且  $\overline{Z_{\text{reg}}}$  包含  $Z$  。

如果  $ x \in Z_{\text{reg}} $  ，则  $ \tilde{Z} $  在 x 处的维度（表示为  $ \dim_x(Z) $  ）已明确定义。一套：

\[
\operatorname{d i m}(Z)=\operatorname*{s u p}_{x\in Z_{r e g}}\operatorname{d i m}_{x}(Z)\enspace.
\]

下面的三角剖分定理将把沿着亚解析分层的可构造层的研究减少到单纯复形上的可构造层的研究。

\statementlabel{命题 8.2.5}令 $ X = \bigsqcup_{\alpha \in A} X_\alpha $  为 $ X $  的子解析子集的局部有限划分。那么存在一个单纯复形  $ \mathbf{S} = (\mathcal{S}, \varDelta) $  和一个同胚  $ i: |\mathbf{S}| \xrightarrow{\sim} X $  使得：

(i) 对于任何  $\sigma \in \Delta$  ，  $i(|\sigma|)$  是  $X$  的亚解析子流形，

(ii) 对于任何  $ \sigma \in \Delta $  都存在  $ \alpha \in A $  使得  $ i(|\sigma|) \subset X_{\alpha} $  。

\subsection{次解析迷向集与 $\mu$-分层}
在本节中，我们将收集稍后需要的有关亚解析各向同性集的所有结果，并且我们将引入  $ \mu $  分层的新概念。

令 S 为流形 X 的局部闭子解析子集。我们设置：

\[
T_{S}^{*}X=\overline{{T_{S_{\mathrm{r e g}}}^{*}}}X\cap\pi^{-1}(S)\enspace.
\]

\statementlabel{命题 8.3.1}集合  $ T_S^* X $  在  $ T^* X $  中是亚解析的。

\statementlabel{证明}我们可以假设 S 关闭。

令  $f: Y \to X$  为流形的真态射，使得  $f(Y) = S$  。设定：

\[
Y_{\circ}=\left\{y\in f^{-1}(S_{\sf r e g});f_{y}^{\prime}:T_{y}Y\to T_{f(y)}S{~i s~s u r j e c t i v e}\right\}\;.
\]

那么  $Y_{\circ}$  是开放的、亚解析的，而  $S_{\circ}=f(Y_{\circ})$  在  $S$  中是开放的、稠密的和亚解析的。集合 $P=T^{*}X\backslash f_{\pi}^{t}f^{\prime-1}(\dot{T}^{*}Y_{\circ})$  是亚解析的，我们有：

\[
{P=\left\{(x;\xi)\in T^{*}X,\mathrm{f o r a n y}y\in Y_{\circ}\cap f^{-1}(x),t f^{\prime}(y)\cdot\xi=0\right\}.}
\]

因此 $ P \cap \pi^{-1}(S_{\circ}) = T_{S_{\circ}}^{*}X $ 并且这个集合是亚解析的。因此 $ T_{S}^{*}X = \overline{T_{S_{\circ}}^{*}X} $  是亚解析的。   $ \square $ 

\statementlabel{命题 8.3.2}(i) 设 M 为闭子流形，S 为 X 的亚解析子集。则法锥  $ C_M(S) $  在  $ T_M X $  中是亚解析的。

(ii) 令 $S_{1}$  和 $S_{2}$  为 $X$  的两个子分析子集。然后法锥 $C(S_{1}, S_{2})$  在 $TX$  中进行亚解析。

\statementlabel{证明}(ii) 由 (i) 得出，并且 (i) 直接由定义 4.1.1 得出。 □

\statementlabel{命题 8.3.3}令  $S_1$  和  $S_2$  为  $\mathbb{R}^n$  的两个子解析子集，并令  $v \in C_x(S_1, S_2)$  。那么存在两条实解析曲线 $x_i(t), t \in]-1,1[,\,i=1,2$ ，使得 $x_i(0)=x, x_i(t) \in S_i\,(t \neq 0, i=1,2)$  和 $x_1(t)-x_2(t)=t^k v+O(t^{k+1})$  对于某个整数 $k>0$  。

\statementlabel{证明}令 $p$  为映射 $\mathbb{R}^n \times \mathbb{R}^n \times \mathbb{R} \to \mathbb{R}^n \times \mathbb{R}^n$  、 $(x, y, s) \mapsto (x, x - sy)$  。然后将 $TX$ 识别为 $\mathbb{R}^n \times \mathbb{R}^n \times \mathbb{R}$ 的集合 $\{s = 0\}$ ，圆锥 $C(S_1, S_2)$ 被识别为 $p^{-1}(S_1 \times S_2) \cap \Omega \cap TX$ ，其中 $\Omega = \{s > 0\}$ 。

让 $v \in C_x(S_1, S_2)$  。根据曲线选择引理，存在一条曲线  $u(t) = (x(t), y(t), s(t))$  、  $-1 < t < 1$  ，使得  $x(0) = x$  、  $y(0) = v$  、  $s(0) = 0$  和  $x(t) \in S_1$  、  $x(t) - s(t)y(t) \in S_2$  、  $s(t) > 0$  ，对于  $t \neq 0$  。设置 $x_1(t) = x(t)$ ， $x_2(t) = x(t) - s(t)y(t)$ ，我们得到 $x_1(t) - x_2(t) = s(t)y(t) = t^k v + O(t^{k+1})$ ，对于一些 $k > 0$ 。   $\square$ 

\statementlabel{命题 8.3.4}令 S 为 X 的子解析子集，并令  $\theta$  为 X 上的（实解析）1-形式。则以下条件是等效的：

（一） $ \theta|_{S_{reg}} = 0 $ 

(ii) 对于任何 $ x \in X $  、 $ \theta|_{C_x(S)} = 0 $  。

条件 (ii) 意味着  $\theta$  给出的  $T_x X$  上的线性函数在  $C_x(S)$  上消失。

\statementlabel{证明}(ii)  $\Rightarrow$  (i) 是显而易见的。

假设 (i) 并设  $ x_o \in \overline{S} $  、  $ v \in C_{x_o}(S) $  、  $ v \neq 0 $  。根据前面的结果，存在一条曲线 $ x(t) $ ，-1 < t < 1，使得 $ x(t) \in S_{\text{reg}} $  对于 $ t \neq 0 $  和 $ x(t) = x_o + t^k v + O(t^{k+1}) $  。选择局部坐标系并写入  $ \theta = \sum_j a_j(x) dx_j $  ,  $ v = (v_1, \ldots, v_n) $  。然后  $ \sum_j a_j(x(t)) (\partial x_j / \partial t) = 0 $  为  $ t \neq 0 $  ，因此  $ k \sum_j a_j(x_o) v_j t^{k-1} + O(t^k) = 0 $  ，我们得到  $ \langle \theta(x_o), v \rangle = 0 $  。   $ \square $ 

\statementlabel{定义 8.3.5}如果命题 8.3.4 的等价条件满足，我们说  $ \theta $  在 S 上消失，我们写成  $ \theta|_{S} = 0 $  。

\statementlabel{推论 8.3.6}(i) 令 $ h: X' \to X $  为流形的态射，并令 $ S \subset X $  、 $ S' \subset X' $  为 $ h(S') \subset S $  的亚解析子集。令  $ \theta $  为  $ X $  上的 1-形式。那么  $ \theta|_{S} = 0 $  意味着  $ h^* \theta|_{S'} = 0 $  。

(ii) 令 $ (S_j)_{j \in J} $  为 $ X $  的亚解析子集的局部有限族，并令 $ \theta $  为 $ X $  上的1-形式。如果  $ \theta|_{S_j} = 0 $  代表所有  $ j $  ，则  $ \theta|_{\bigcup_j S_j} = 0 $  。

\statementlabel{证明}(i) 对于所有  $ x' \in X' $  ，有：

\[
h^{\prime}(C_{x^{\prime}}(S^{\prime}))\subset C_{h(x^{\prime})}(S)\enspace.
\]

(ii) 这是从

\[
C_{x}\left(\bigcup_{j}S_{j}\right)=\bigcup_{j}C_{x}(S_{j})~.\quad\Box
\]

\statementlabel{命题 8.3.7}令 $h: X' \to X$  为流形的态射，并令 $S \subset X$  和 $S' \subset X'$  为亚解析子集。假设 $h(S') = S$ 。令  $\theta$  为  $X$  上的 1-形式。那么下面的条件是等价的。

（一） $ \theta|_{S} = 0 $ ，

（二） $ h^*\theta|_{S'} = 0 $ 。

\statementlabel{证明}(i)  $\Rightarrow$  (ii) 由推论 8.3.6 (i) 得出。

(ii)  $\Rightarrow$  (i) 我们可以假设 $S$  是非奇异的。令  $g: Y \to X'$  为真态射，使得  $g(Y) = S$  。然后 $g^*h^*\theta = 0$ 。

设置  $f = g \circ h: Y \to X$  并定义  $Y_{\circ} = \{ y \in Y; f_{y}': T_{y} Y \to T_{f(y)} S \text{ is surjective} \}$  。由于  $f(Y_{\circ})$  在  $S$  中是密集的，因此证明已完成。 （请注意，我们假设  $Y$  在无穷大处可数。） $\square$ 

我们需要一个关于向量丛中圆锥亚解析子集的结果。

\statementlabel{命题 8.3.8}令 $ \tau: E \to X $  为向量丛。

(i) 令 A 为  $ \dot{E} $  的圆锥子集，在  $ \dot{E} $  中进行亚解析。那么A在E中是亚解析的。

(ii) 如果 $\varLambda$  是 $E$  的子解析子集，使得 $\varLambda \to X$  是真映射，则 $\mathbb{R}^{+} \cdot \varLambda$  是子解析子集。

(iii) 如果  $\varLambda$  是  $E$  的圆锥亚解析子集，则  $\tau(\varLambda)$  是  $X$  中的亚解析。

(iv) 令 $ \tau': E' \to X $  为另一个向量丛， $ f: E \to E' $  为向量丛的态射。如果  $ \varLambda $  是  $ E $  的圆锥亚解析子集，则  $ f(\varLambda) $  是  $ E' $  中的亚解析。

\statementlabel{证明}令 $ \gamma: \dot{E} \to \dot{E}/\mathbb{R}^+ $  为投影， $ i: X \to E $  为零部分。让我们采用 $ \dot{E} $  的子流形 $ S $ ，使得 $ S \to \dot{E}/\mathbb{R}^+ $  是同构。

(i) 令 $ \mu: \mathbb{R}^{+} \times E \to E $  为乘法映射。然后 $ A' = \mu([0,1] \times (S \cap A)) $ 

是  $E$  中的亚解析，并且 (i) 从  $A' \setminus i(X) = A$  在  $i(X)$  邻域上的事实得出。

(ii)  $ \mu:[0,1] \times \bar{A} \to E $  是真映射，因为  $ \bar{A} $  相对于  $ X $  是正确的。然后 (ii) 从  $ \mathbb{R}^+\cdot \mathcal{A} = \mu([0,1] \times \mathcal{A}) \cup p(\mu^{-1}(\mathcal{A}) \cap [0,1] \times E) $  得出，其中  $ p $  是第二个投影  $ \mathbb{R} \times E \to E $ 

(iii) 是 (iv) 与  $E^{\prime}=X$  的特例。

(iv) 令 $ A \subset E $  为 $ E $  零部分的亚解析开邻域，使得 $ \overline{A} \to X $  是正确的。那么 $ f(A \cap \varLambda) $  是亚解析的。然后 (iv) 由 (ii) 得出，因为  $ f(\varLambda) = \mathbb{R}^+\cdot f(A \cap \varLambda) $  。   $ \square $ 

现在我们将研究  $ T^*X $  的亚解析各向同性子集。回想一下，  $ \alpha_X $  （或简称  $ \alpha $  ）表示  $ T^*X $  上的规范 1-形式。

\statementlabel{定义 8.3.9}令  $A$  为  $T^*X$  的圆锥亚解析子集。

(i) 如果  $\alpha|_{\Lambda} = 0$  ，我们说  $\Lambda$  是各向同性的。

(ii) 如果 $\mathcal{A}$  既是各向同性的又是对合的，我们就说 $\mathcal{A}$  是拉格朗日量（参见定义6.5.1）。

\statementlabel{命题 8.3.10}(i) 令  $A$  为  $T^*X$  的闭圆锥二次解析子集。以下条件是等效的。

(a) A 是各向同性的。

(b) 存在 $X$  的亚解析子集的局部有限族 $\{X_j\}$ ，使得 $A \subset \bigcup_j T_{X_i}^* X$  。

(c) 与 (b) 相同，具有有限的亚解析子流形族  $ X_{j} \subset \pi(A) $  。

(ii) 此外，如果  $A$  是各向同性且  $Y$  是  $X$  的亚解析子集，则  $Y$  中存在亚解析子流形  $Y_{\circ} \subset Y$  开稠密，使得  $A \cap \pi^{-1}(Y_{\circ}) \subset T_{Y_{\circ}}^{*}X$  。

\statementlabel{证明}(i) (b)  $\Rightarrow$  (a) 根据推论 8.3.6 足以证明  $\alpha|_{T_{X_j}^{*}X}=0$  ，它是从  $X_{j,\mathrm{reg}}$  自  $T_{X_j}^{*}X=\overline{T_{X_{j,\mathrm{reg}}}^{*}}X$  以来的相应结果得出的。

(a)  $\Rightarrow$  (c) 令 $d$  为投影 $\pi: \mathcal{A}_{\text{reg}} \to X$  的秩的最大值。我们将通过 $d$  进行归纳论证。让 $\mathcal{A}_o = \{p \in \mathcal{A}_{\text{reg}}; \text{rank of } \pi \text{ at } p \text{ is } d\}$  。然后  $\mathcal{A}_o$  在  $\mathcal{A}$  中是开放和亚解析的。设置  $X_o = (\pi(\mathcal{A}))_{\text{reg}}$  和  $\mathcal{A}_o' = \mathcal{A}_o \cap \pi^{-1}(X_o)$  。那么  $\mathcal{A}_o'$  在  $\mathcal{A}$  中是开放和亚解析的，在  $\mathcal{A}_o$  中是稠密的，并且投影  $\pi: \mathcal{A}_o' \to X_o$  的微分在每个  $p \in \mathcal{A}_o$  处是满射的。让我们证明：

\[
\varLambda^{\prime}_{\circ}\subset T_{X_{\circ}}^{*}X~.
\]

选择局部坐标系 $(x_1, \ldots, x_n) = (x', x'')$  和 $x' = (x_1, \ldots, x_p)$ ，使得 $X_\circ = \{x'' = 0\}$  并让 $\xi = (\xi', \xi'')$  表示对偶坐标。然后  $\alpha|_{\pi^{-1}(X_\circ)} = \xi' dx'$  和  $dx_1, \ldots, dx_p$  线性独立于  $X_\circ$  和  $\pi$  平滑，我们在  $A_\circ$  上得到  $\xi' = 0$  ，这证明了 (8.3.3)。现在设置 $A' = \overline{A \setminus T_x^* X}$ 。从 $A_\circ \cap A' = \varnothing$ 到(8.3.3)，归纳继续进行。

(c)  $ \Rightarrow $  (b) 显而易见。

(ii) 分别用 $\mathcal{A} \cap \pi^{-1}(\overline{Y})$ 和 $\overline{Y}$ 替换 $\mathcal{A}$ 和 $Y$ ，我们可以从一开始就假设 $\mathcal{A} \subset \pi^{-1}(Y)$ 。然后选择 $\{X_j\}$ 如(i)(c)中那样，将 $Y_\circ$ 作为 $Y_{\text{reg}}$ 中打开的 $Y_{\text{reg}} \cap X_j$ 的连通分量的并集就足够了。

\statementlabel{命题 8.3.11}令 $f: Y \to X$  为流形的态射。

(i) 令 $ A \subset T^*Y $  为圆锥亚解析各向同性集，使得 $ f_\pi: \overline{f^{'}-1}(\overline{A}) \to T^*X $  是正确的。那么  $ f_\pi^{'f^{'}-1}(A) $  是  $ T^*X $  的圆锥亚解析各向同性子集。

(ii) 令 $ A \subset T^*X $  为圆锥亚解析各向同性集。那么  $ {}^{t}f'f_{\pi}^{-1}(A) $  是  $ T^*Y $  的圆锥亚解析各向同性子集。

\statementlabel{证明}(i) 首先注意  $ f'*\alpha_Y = f_\pi^*\alpha_X $  。接下来应用命题 8.3.7。我们得到的含义是：

\[
\alpha_{Y}|_{\cal A}=0\Rightarrow(^{t}f^{\prime})^{*}\alpha_{Y}|_{^{t}f^{\prime-1}({\cal A})}=0\Leftrightarrow f_{\pi}^{*}\alpha_{X}|_{^{t}f^{\prime-1}({\cal A})}=0\Rightarrow\alpha_{X}|_{f_{\pi}^{t}f^{\prime-1}({\cal A})}=0\enspace.
\]

(ii) 除了  $ f'f_{\pi}^{-1}(A) $  的亚解析性之外，证明是类似的，这是从命题 8.3.8 得出的。   $ \square $ 

以下结果将在证明有限性定理时发挥重要作用。

\statementlabel{命题 8.3.12（微局域 Bertini-Sard 定理）}令 $\varphi: X \to \mathbb{R}$  为实解析函数，并令 $A \subset T^*X$  为闭圆锥亚解析各向同性子集。让 $S = \{t \in \mathbb{R}; t = \varphi(x), \mathrm{d}\varphi(x) \in A$  进行一些 $x \in X\}$  。假设  $\varphi$  对于  $\pi(A)$  是正确的。那么 $S$  是离散的。

\statementlabel{证明}我们有 $S = \{t \in \mathbb{R}; (t; \mathrm{d}t) \in \varphi_{\pi}^{t} \varphi^{\prime-1}(\varLambda)\}$ 。由于 $\varphi_{\pi}^{t} \varphi^{\prime-1}(\varLambda)$  是封闭的亚解析各向同性子集，因此 $S$  是离散的。   $\square$ 

\statementlabel{命题 8.3.13}令  $\mathcal{A}$  和  $\mathcal{A}_{\circ}$  为  $T^{*}X$  和  $\mathcal{A}\subset\mathcal{A}_{\circ}$  的两个局部闭圆锥曲线子集。假设  $\mathcal{A}_{\circ}$  是亚解析且各向同性的，而  $\mathcal{A}$  是对合的且在  $\mathcal{A}_{\circ}$  中是封闭的（但我们不假设  $\mathcal{A}$  是亚解析的）。那么 $\mathcal{A}$  是亚解析和拉格朗日。

\statementlabel{证明}我们需要两个引理。

\statementlabel{引理 8.3.14}令 $\mathcal{A}_{\circ}$  为 $X$  的各向同性子流形， $\mathcal{A}$  为对合子集，封闭在 $\mathcal{A}_{\circ}$  中。然后  $\mathcal{A}$  在  $\mathcal{A}_{\circ}$  和  $\dim \mathcal{A} = \dim X$  中打开。

引理 8.3.14 的证明。根据命题 A.2.9，@@​​MATH0@@ 局部包含在拉格朗日流形中。因此我们可以假设 $A_{\circ}$ 拉格朗日。对于在  $A_{\circ}$  上消失的任何实函数  $\phi$  ，  $A$  是由命题 6.5.2 得到的  $H_{\varphi}$  积分曲线的并集，而在  $x$  点发出的  $H_{\varphi}$  积分曲线的并集是  $x$  的邻域。

\statementlabel{引理 8.3.15}令 $S$  为 $T^*X$  的圆锥局部闭亚解析各向同性子集，并令 $V$  为局部闭卷合子集。假设  $V \subset S$  和  $\dim S < \dim X$  。然后 $V = \emptyset$ 。

引理 8.3.15 的证明。通过对 dim S 进行归纳论证，足以证明  $ V \cap S_{\text{reg}} = \emptyset $  。因此我们可以假设 S 是平滑的。让 $ x \in V $  。通过引理 8.3.14， $ \dim_x V = \dim X $  。这与 $ \dim S < \dim X $ 相矛盾。   $ \square $ 

命题 8.3.13 的证明结束。设置 $ A_{\circ}^{\prime} = (A_{\circ})_{\text{reg}} $ ， $ A^{\prime} = A \cap A_{\circ}^{\prime} $ 。根据引理 8.3.14，我们发现  $ A^{\prime} $  在  $ A_{\circ}^{\prime} $  中既开又闭。   $ \text{Since } A_{\circ}^{\prime} $  是亚解析的， $ A^{\prime} $  也是亚解析的。那么就足以证明  $ A = A^{\prime} \cap A_{\circ} $  。设置 $ V = A \setminus \overline{A}^{\prime} $  。这是  $ A_{\circ} \setminus (A_{\circ})_{\text{reg}} $  中包含的  $ T^*X $  的对合子集。那么根据引理 8.3.15  $ V $  为空。   $ \square $ 

我们将研究沿着拉格朗日子流形的各向同性子集的法锥。首先，我们需要一个引理。

\statementlabel{引理 8.3.16}设  $ X = \mathbb{R}^n \times \mathbb{R} $  ，坐标为  $ (x, t) $  ，设  $ Y $  为超曲面  $ \{t = 0\} $  ，并设  $ Z $  为  $ X $  的亚解析子集，使得  $ Z \subset \overline{Z} \setminus Y $  。令  $ \alpha $  为  $ X $  上的 1-形式，并令  $ \theta = t\alpha + b \, dt $  ，其中  $ b $  是  $ X $  上的实解析函数。那么  $ \theta|_{z} = 0 $  意味着  $ \alpha|_{z \cap Y} = 0 $  和  $ b|_{z \cap Y} = 0 $  。

\statementlabel{证明}人们可以假设 Z 是闭合的。

令  $f: X' \to X$  为流形的真态射，使得  $f(X') = Z$  ，并设置  $Y' = f^{-1}(Y)$  。令  $X''$  为  $X'$  上  $t \circ f = 0$  的连通分量的并集。自 $f(X' \setminus X'') \supset Z \setminus Y$  、 $f(X' \setminus X'') = Z$  开始。因此，我们可以通过用  $X' \setminus X'$  替换  $X'$  来假设  $Y'$  并不密集。通过将命题8.2.4应用于函数 $t' = t \circ f$ ，我们可以从一开始就假设在 $Y$ 的每个点的邻域中，存在一个局部坐标系 $(t_1, \ldots, t_N)$ ，使得 $t' = \pm \prod_{j=1}^l t_j^{a_j}$ ，其中 $1 \leq l \leq N$ 和 $a_j$ 是正整数。设置 $\alpha' = f^* \alpha$ ， $b' = f^* b$ 。然后 $t' \alpha' + b' \, dt' = 0$ 。自 $\mathrm{dt}'/t' = \sum_j(a_j/t_j) \, \mathrm{dt}_j$ 以来，我们得到：

\[
\alpha^{\prime}=-b^{\prime}\left(\sum_{j}\left(a_{j}/t_{j}\right)\mathrm{d}t_{j}\right).
\]

通过这个公式，我们看到  $ t_1 \ldots t_l $  除  $ b' $  ，因此  $ b'|_{Y'} = 0 $  和  $ b|_{Y \cap Z} = 0 $  。

我们也可以将(8.3.4)写为：

\[
\alpha^{\prime}=(-b^{\prime}/t_{1}\ldots t_{l})\left(\sum_{j}t_{1}\ldots t_{l}(a_{j}/t_{j})\operatorname{d}t_{j}\right).
\]

因此  $ \alpha'|_{Y'_{reg}} = 0 $  ，这意味着命题 8.3.7 中的  $ \alpha'|_{Y'} = 0 $  和  $ \alpha|_{Y \cap Z} = 0 $ 。

\statementlabel{定理 8.3.17}令  $ A \subset T^*X $  为（平滑、圆锥）拉格朗日子流形，并令  $ S $  为  $ T^*X $  的圆锥亚解析各向同性子集。我们识别 $ T_A T^*X $ 和 $ T^*A $ 

通过 -H，其中 H 是哈密顿同构： $ TT^*X \simeq T^*T^*X $  。然后：

(i) 法锥  $ C_{A}(S) $  是  $ T^{*}A $  的圆锥亚解析各向同性子集，

(ii)  $ C_{A}(S) $  包含在超平面  $ \{\alpha_{X}=0\} $  中。

（回想一下， $ \mathcal{A} $  上的  $ \alpha_X $  为 0， $ \alpha_X $  定义了  $ T_\mathcal{A}T^*X $  上的线性函数。因此，超平面  $ \{\alpha_X=0\} $  在  $ T_\mathcal{A}T^*X\simeq T^*\mathcal{A} $  中得到了明确的定义。）

\statementlabel{证明}我们可以假设 $A$  是子流形的保形束（命题 A.2.9 和练习 A.2）。为了简单起见，我们假设 $A = \{(x; \xi) \in T^* \mathbb{R}^n; x = 0\}$ （在一般情况下，证明是类似的）。考虑映射 $p: T^*X \times \mathbb{R} \to T^*X$ ， $(x; \xi), t) \mapsto (tx; \xi)$ 。然后 $T_A T^*X \simeq \{t = 0\}$ 、 $C_A(S) \simeq \overline{p^{-1}(S) \cap \{t > 0\}} \cap \{t = 0\}$  和 $p^*(\alpha_X) = \sum_j \xi_j \mathrm{d}(tx_j) = t(\sum_j \xi_j \mathrm{d}x_j) + \langle x, \xi \rangle \mathrm{d}t$  。由于 $p^*(\alpha_X)|_{p^{-1}(S)} = 0$  由命题 8.3.7 得出，我们由引理 8.3.16 得到：

\[
\left(\sum_{j}\xi_{j}\mathrm{d}x_{j}\right)\bigg|_{C_{A}(\mathbb{S})}=0\qquad\mathrm{a n d}\qquad\langle x,\xi\rangle|_{C_{A}(\mathbb{S})}=0.
\]

因此  $ (\sum_j x_j \, \mathrm{d}\xi_j)|_{C_A(S)} = 0 $  ，证明已经完成，因为  $ -\sum_j x_j \, \mathrm{d}\xi_j $  是  $ T^* A $  上的规范 1-形式。   $ \square $ 

\statementlabel{推论 8.3.18}(i) 令 $\Lambda_1$  和 $\Lambda_2$  为 $T^*X$  的两个圆锥亚解析各向同性子集。那么  $\Lambda_1 + \Lambda_2$  是  $T^*X$  的圆锥亚解析各向同性子集。

(ii) 令 $f: Y \to X$  为流形的态射，并令 $A \subset T^*X$  为圆锥亚解析各向同性子集。那么 $f^\#(A)$  是 $T^*Y$  的圆锥亚解析各向同性子集（参见定义6.2.3）。

现在我们将通过亚解析流形来研究 X 的分层。让我们首先回顾一些基本定义。

让  $(X_j)_{j \in J}$  成为  $X$  的子集族。有人说这个家族是一个覆盖如果 $X = \bigcup_{j \in J} X_j$ 。此外，如果  $X_j$  彼此不相交，则有人说该覆盖物是一个分区，而有人则写为  $X = \bigcup_{j \in J} X_j$  。有人说覆盖  $X = \bigcup_{i \in I} X_i$  比覆盖  $X = \bigcup_{j \in J} X_j$  更好，如果对于每个  $i \in I$  都存在  $j \in J$  和  $X_i \subset X_j$  。

\statementlabel{定义 8.3.19}(i) 对于 X 的闭合子解析子集 Y，分区  $ Y = \bigsqcup_{\alpha \in A} X_\alpha $  被称为 Y 的子解析分层，如果它是局部有限的，则  $ X_\alpha $  是子解析子流形，并且对于所有对  $ (\alpha, \beta) \in A \times A $  使得  $ \overline{X}_\alpha \cap X_\beta \neq \emptyset $  其中一个具有  $ X_\beta \subset \overline{X}_\alpha $  。每个 $ X_\alpha $  称为一个层。

(ii) 对于 X 的两个子流形 N 和 M，我们说  $ (M, N) $  满足  $ \mu $  条件，如果：

\[
(T_{M}^{*}X\stackrel{\frown}{+}T_{N}^{*}X)\cap\pi^{-1}(N)\subset T_{N}^{*}X\enspace.
\]

(iii) 分区 $X = \bigsqcup_{\alpha \in A} X_\alpha$  是一个 $\mu$  分层，如果它是亚解析分层并且此外对于所有对 $(\alpha, \beta) \in A \times A$  使得 $X_\beta \subset \overline{X}_\alpha \setminus X_\alpha$  、 $(X_\alpha, X_\beta)$  满足 $\mu$  条件。

请注意，如果  $X = \bigsqcup_{\alpha} X_{\alpha}$  是  $\mu$  分层，则集合  $\mathcal{A} = \bigsqcup_{\alpha} T_{X_{\alpha}}^{*} X$  是  $T^{*}X$  的闭圆锥亚解析各向同性子集。事实上，如果  $X_{\beta} \subset \overline{X}_{\alpha}$  ，则  $T_{X_{\alpha}}^{*} X \cap \pi^{-1}(X_{\beta})$  包含在  $T_{X_{\beta}}^{*} X$  中。

\statementlabel{定理 8.3.20}令  $X = \bigcup_{j \in J} X_j$  为子解析子集对  $X$  的局部有限覆盖。然后存在比覆盖 $X = \bigcup_{j \in J} X_j$ 更精细的 $\mu$ 分层 $X = \bigsqcup_{\alpha \in A} X_\alpha$ 。

为了证明这个结果，我们需要一个引理。

\statementlabel{引理 8.3.21}(i) 令 Y 为  $X$  的闭合子解析子集， $Y = \bigcup_{j \in J} X_j$  为子解析子集的局部有限覆盖。然后存在  $Y$  的更精细的子分析分层。

(ii) 令 $N$  和 $M$  为 $X$  的两个亚解析子流形。然后 $N$  中的点 $x$  的集合 $\Omega$  使得 $(M,N)$  满足 $x$  邻域上的 $\mu$  -条件在 $N$  中是稠密的，在 $X$  中是亚解析的。

(iii) 令 $ X' $  为 $ X $  、 $ X' = \bigsqcup_{\alpha \in A} X_\alpha $  的子解析开子集和 $ X' $  的子解析分层，其中 $ X_\alpha $  子解析位于 $ X $  中。那么存在  $ X' $  的最大开子集  $ \Omega $  使得  $ \Omega = \bigsqcup_{\alpha \in A} (X_\alpha \cap \Omega) $  是  $ \Omega $  的  $ \mu $  分层。此外， $ \Omega $  在 $ X $  中是亚解析的。

引理 8.3.21 的证明。 (i) 令  $A$  为  $J$  的有限子集的集合，并设置  $\alpha \in A$  、  $Z_\alpha = (\bigcap_{j \in \alpha} X_j) \setminus (\bigcup_{j \notin \alpha} X_j)$  。那么 $Y = \bigsqcup_{\alpha} Z_\alpha$  是一个局部有限的子解析分区。因此，我们可以从一开始就假设 $Y = \bigsqcup_{j} X_j$  是局部有限的子解析分区。我们将通过对 dim  $Y$  的归纳进行论证。令  $X_j'$  为  $(X_j)_{\text{reg}} \cap Y_{\text{reg}}$  的连通分量的并集，在  $Y_{\text{reg}}$  中打开。设置 $Y' = Y \setminus (\bigcup_{j} X_j)$ 。由于 $Y \setminus Y' = \bigsqcup_{j} X_j'$  是一个不相交的并集，因此它是一个亚分析分层。令 $A$  为 $J$  的有限子集的集合。然后根据归纳假设，存在比以下更精细的亚分析分层 $Y' = \bigsqcup_{j \in J'} Y_j$ 

\[
Y^{\prime}=\bigcup_{j\in J,\alpha\in A}\left[Y^{\prime}\cap X_{j}\cap\left(\bigcap_{k\in\alpha}(\overline{{X_{k}\setminus Y^{\prime}}})\right)\middle\backslash\left(\bigcup_{k\in J\setminus\alpha}\overline{{X_{k}\setminus Y^{\prime}}}\right)\right]~.
\]

那么 $ Y = (\bigsqcup_{j \in J}(X_j \setminus Y')) \sqcup (\bigsqcup_{j \in J'} Y_j) $ 就是一个分层。

为了看到这一点，只要证明对于任何  $j$  和任何  $k \in J^{\prime}$  ，  $\overline{X_j} \setminus \overline{Y^{\prime}} \cap Y_k \neq \emptyset$  意味着  $\overline{X_j} \setminus Y^{\prime} \supset Y_k$  就足够了。然后是  $\alpha \in A$  这样  $Y_k \subset \bigcap_{i \in \alpha} (\overline{X_i} \setminus \overline{Y^{\prime}}) \setminus (\bigcup_{i \in J \setminus \alpha} \overline{X_i} \setminus \overline{Y^{\prime}})$  。因此  $\overline{X_j} \setminus Y^{\prime} \cap Y_k \neq \emptyset$  意味着  $j \in \alpha$  。这证明了想要的结果。

(ii)  $\Omega$  是  $N\backslash\pi(Z)$  ，其中  $Z=\overline{(T_N^*X+\overline{T_M^*X})\backslash T_N^*X}$  。因此它是亚分析的。由于 $T_N^*X\widehat{+}T_M^*X$  是各向同性的，因此存在 $N$  的开稠子集 $U$ ，使得 $(T_N^*X+\overline{T_M^*X})\cap\pi^{-1}(U)\subset T_N^*X$  （命题8.3.10 (ii)）。然后 $\pi(Z)\cap U=\varnothing$ 。

(iii) 集合 $\Omega$  是 $\bigcup_{(\alpha,\beta)} X_\beta \cap \pi((T_{X_\alpha}^* X + T_{X_\beta}^* X) \setminus T_{X_\beta}^* X)$  的补集，其中 $(\alpha,\beta) \in A \times A, X_\beta \subset \overline{X}_\alpha \setminus X_\alpha$  。

定理的证明 8.3.20。我们可以假设 $ X = \bigsqcup_j X_j $  是一个亚分析分层。通过对暗淡 Y 进行归纳，足以证明如果 Y 是闭子解析

 $X$  的子集，使得  $X \setminus Y = \bigsqcup_{j \in J}(X_j \setminus Y)$  是  $\mu$  分层，则存在  $Y$  的无处密集子分析子集  $Y'$  和更精细的子分析分层  $X = \bigsqcup_i X_i$  ，使得  $X \setminus Y' = \bigsqcup_i (X_i' \setminus Y')$  是  $\mu$  分层。为了证明这一点，让我们通过比分区  $Y = \bigsqcup_j (X_j \cap Y)$  更精细的亚解析流形对  $Y$  进行分层  $Y = \bigsqcup_k Y_k$  。那么  $X = (\bigsqcup_j (X_j \setminus Y)) \sqcup (\bigsqcup_k Y_k)$  是一个亚分析分层，我们可以从一开始就假设  $J = J' \sqcup J''$  、  $X \setminus Y = \bigsqcup_{j \in J'} X_j$  和  $Y = \bigsqcup_{j \in J''} X_j$  。令 $\Omega$  为 $X$  的最大开子集，使得 $\Omega = \bigsqcup_j (X_j \cap \Omega)$  是 $\mu$  分层（引理8.3.21 (ii)）。那么  $\Omega$  是亚解析的并且包含  $X \setminus Y$  。因此，足以证明  $\Omega \cap Y$  在  $Y$  中是密集的。但这是直接从引理 8.3.21 (ii) 得出的。

\statementlabel{推论 8.3.22}令  $A$  为  $T^*X$  的闭圆锥亚解析各向同性子集。然后存在  $\mu$  -分层  $X = \bigsqcup_{\alpha} X_\alpha$  使得  $A \subset \bigsqcup_{\alpha} T^*X$  。

\statementlabel{证明}根据命题 8.3.10，存在由  $ A \subset \bigcup_j T_{X_j}^* X $  的子解析子集覆盖  $ X = \bigcup_j X_j $  的局部有限。应用定理 8.3.20，我们发现  $ \mu $  分层  $ X = \bigsqcup_{\alpha} X_\alpha $  比覆盖层更精细。那么  $ \bigcup_j T_{X_j}^* X $  包含在  $ \bigsqcup_{\alpha} T_{X_\alpha}^* X $  中，证明如下。   $ \square $ 

\statementlabel{命题 8.3.23}令  $M$  为  $X$  的子流形， $\Lambda$  为  $T^*X$  的二次曲线亚解析各向同性子集。如果  $T_M^*X \cap \Lambda = \emptyset$  和  $(\Lambda + T_M^*X) \cap \pi^{-1}(M) \subset T_M^*X$  ，则  $\bar{\Lambda} \cap T_M^*X$  在  $T_M^*X \to M$  的每条纤维上都不密集。

\statementlabel{证明}我们将证明，对于  $ x_{\circ} \in M $  ，  $ \pi^{-1}(x_{\circ}) \cap \bar{\Lambda} $  在  $ T_M^*X $  中并不密集。让我们在  $ X $  上采用局部坐标系  $ x = (x', x'') $ ，使得  $ M = \{x' = 0\} $  和  $ x_{\circ} $  为原点。我们将在  $ T^*T_M^*X $  和  $ T_{T_M^*X} T^*X $  上使用关联的坐标系，如 (6.2.3) 中所示。那么假设 $ (\Lambda + T_M^*X) \cap \pi^{-1}(M) \subset T_M^*X $ 意味着：

\[
\left\{x^{\prime}=0\right\}\cap C_{T_{M}^{*}X}(\varLambda)\subset\left\{\xi^{\prime\prime}=0\right\}.
\]

另一方面，根据定理 8.3.17， $C_{T_M^*X}(\varLambda)$  是各向同性的，因此根据命题 8.3.10 (ii)，存在  $\pi^{-1}(x_o) \cap T_M^*X$  的开稠子集  $U$  ，使得  $C_{T_M^*X}(\varLambda) \times T_M^*XU \subset T_{\pi^{-1}(x_o) \cap T_M^*X}^*(T_M^*X)$  。由于  $T_{\pi^{-1}(x_o) \cap T_M^*X}^*(T_M^*X)$  由  $\{x'' = 0, x' = 0\}$  给出，因此  $C_{T_M^*X}(\varLambda) \times T_M^*XU$  包含在  $\{x' = \xi'' = 0\}$  中，即  $T_{T_M^*X}^*X^*X$  的零部分中。这意味着  $\bar{\varLambda} \cap U$  包含在  $U$  邻域上的  $T_M^*X$  中，并且  $T_M^*X \cap \varLambda$  为空，我们得到  $\bar{\varLambda} \cap U = \varnothing$  。

\statementlabel{推论 8.3.24}让 $ X = \left| \right|_{\alpha} X_{\alpha} $  成为 $ \mu $  分层。然后：

\[
\pi\left(T_{X_{\alpha}}^{*}X\setminus\bigcap_{\beta\neq\alpha}\overline{T_{X_{\beta}}^{*}X}\right)=X_{\alpha}.
\]

让我们引入另一个概念，类似于分层，但有时更方便。

\statementlabel{定义 8.3.25}(i) 拓扑空间 $X$  上的 $A$ （递减）过滤是 $X$  的闭子集 $\{X_j\}$  的递减序列 $\{X_j\}$ ，使得 $X_j = X$  对于 $j \ll 0$  和 $X_j = \varnothing$  对于 $j \gg 0$  。

(ii) 实解析流形 $X$  上的亚解析过滤是过滤 $\{X_j\}$ ，使得 $X_j$  是 $X$  中的亚解析， $X_i \setminus X_{i+1}$  是实分析子流形。

(iii) 如果  $(X_j\backslash X_{j+1},X_k\backslash X_{k+1})$  满足任何  $j$  和  $k$  以及  $j>k$  的  $\mu$  -条件，则亚分析过滤称为  $\mu$  过滤。

同样，可以考虑增加过滤。

\statementlabel{命题 8.3.26}令  $X = \bigcup_{\alpha} X_{\alpha}$  为  $X$  闭子解析子集的局部有限覆盖。然后存在  $X$  的亚解析  $\mu$  过滤  $\{X_j\}$  ，使得对于任何  $j$  ，  $X_j \setminus X_{j+1}$  的任何连通分量都包含在某些  $X_\alpha$  中。

\statementlabel{证明}我们可以假设 $ X = \bigsqcup_{\alpha} X_{\alpha} $  是 $ \mu $  分层。设定：

\[
X_{j}=\bigcup_{\dim X_{\alpha}\leq-j}X_{\alpha}.
\]

那么 $X_j\backslash X_{j+1}$  是 $(-j)$  维流形的不相交并，并且可以立即检查过滤 $\{X_j\}$  是否满足所需的属性。 □

为了结束本节，让我们证明莫尔斯函数关于各向同性子集的存在性，我们将在第十章中使用这个结果。

\statementlabel{命题 8.3.27}令  $X$  为  $\mathbb{R}^N$  的  $n$  维闭子流形，并令  $A$  为  $T^*X$  的闭圆锥亚解析各向同性子集。令  $A_\circ$  为  $A$  中包含的亚解析  $n$  维子流形，并假设  $\dim(A\setminus A_\circ) < n$  。然后存在 $x_\circ \in \mathbb{R}^N$ ，这样，为 $x \in X$ 和 $A_\phi = \{(x; d\phi(x)); x \in X\} \subset T^*X$ 设置 $\phi(x) = |x - x_\circ|^2$ ，我们有：

\[
A_{\phi}\cap A\subset A_{\circ}\qquad\mathrm{a n d~t h e~i n t e r s e c t i o n~i s~t r a n s v e r s a l}.
\]

\statementlabel{证明}让  $q$  表示投影  $X \times_{\mathbb{R}^N} T^*\mathbb{R}^N \to T^*X$  。将  $X, A, A_\circ$  替换为  $\mathbb{R}^N, q^{-1}A, q^{-1}A_\circ$  ，我们可以从一开始就假设  $X = \mathbb{R}^n$  。我们用 $(x; \xi)$  表示 $T^*\mathbb{R}^n$  上的齐次辛坐标。设置 $g(x, \xi) = x - \xi/2$ ，我们有 $A_\phi = g^{-1}(x_\circ)$ 。现在让 $f$  成为复合映射：

\[
f\colon\varLambda\hookrightarrow T^{*}\mathbb{R}^{n}\xrightarrow[g]{}\mathbb{R}^{n}\enspace.
\]

由于  $ \dim(\varLambda \setminus \mathcal{A}_o) \leqslant n - 1 $  ， $ f(\varLambda \setminus \mathcal{A}_o) $  的测量值为零。集合 $ G = \{p \in \mathcal{A}_o; T_p f $  不是满射\textbackslash{}\}。根据 Sard 定理（参见 Guillemin-Pollack [1]）， $ f(G) $  的测度为零。因此 $ f(\varLambda \setminus \mathcal{A}_o) \cup f(G) \neq \mathbb{R}^n $  和任何 $ x_o \in \mathbb{R}^n \setminus (f(\varLambda \setminus \mathcal{A}_o)) \cup f(G)) $  都具有所需的属性。   $ \square $ 

关于分层空间上的莫尔斯函数，参见。 Lazzeri [1] Pignogni [1] 和 Goresky-MacPherson [2]。

\subsection{$\mathbb{R}$-构造层}
在本节中，我们将研究沿亚解析分层局部常值的层。

\statementlabel{命题 8.4.1}令  $ X = \bigsqcup_{\alpha \in A} X_\alpha $  为  $ \mu $  分层并令  $ F \in \mathrm{Ob}(\mathbf{D}^b(X)) $  。以下条件是等效的：

(i) 对于所有  $ j \in \mathbb{Z} $  、所有  $ \alpha \in A $  ，层  $ H^j(F)|_{X_x} $  局部常值，

（二） $ \mathrm{SS}(F) \subset \bigsqcup_{\alpha \in A} T_{X_x}^* X $ 。

\statementlabel{证明}(ii) $\Rightarrow$  (i)。根据推论 6.4.5，我们知道  $\mathrm{SS}(F_{X_s})$  包含在  $X_\alpha$  邻域上的  $\mathrm{SS}(F) \widehat{+} T_{X_s}^* X$  中。因此 $\mathrm{SS}(F_{X_s}) \cap \pi^{-1}(X_\alpha)$ 包含在 $T_{X_s}^* X$ 中，我们通过命题5.4.4得到 $\mathrm{SS}(F|_{X_s}) \subset T_{X_s}^* X_\alpha$ 。这意味着根据命题 5.4.5， $H^j(F)|_{X_s} = H^j(F|_{X_s})$  是局部常数。

（一） $\Rightarrow$ （二）。问题是  $X$  上的局部问题，我们可以假设分层是有限的。令  $B$  为  $A$  的子集，使得  $Y = \bigsqcup_{\alpha \in B} X_\alpha$  在  $X$  中闭合，并假设我们已经知道  $\mathrm{SS}(F) \cap \pi^{-1}(X \setminus Y)$  包含在  $\bigsqcup_{\alpha \in A} T_{X_\alpha}^* X$  中。令  $\alpha_0 \in B$  使得  $X_{\alpha_0}$  在  $Y$  中打开并设置  $Y' = Y \setminus X_{\alpha_0}$  。通过归纳论证，只需证明  $\mathrm{SS}(F) \cap \pi^{-1}(X \setminus Y')$  包含在  $\bigsqcup_{\alpha \in A} T_{X_\alpha}^* X$  中即可。因此我们可以假设  $Y = X_{\alpha_0}$  。让  $j$  表示开放嵌入  $X \setminus X_{\alpha_0} \hookrightarrow X$  。我们得到一个显着的三角形：

\[
R j_{!}j^{-1}F\longrightarrow F\longrightarrow F_{X_{\alpha_{0}}}\xrightarrow[+1]{}.
\]

根据假设和命题 5.4.4， $ \mathrm{SS}(F_{X_s}) \subset T_{X_s}^* X $  。根据命题 6.3.2，我们有：

\[
\begin{aligned}{\mathbb{S S}(R j_{!}j^{-1}F)\cap\pi^{-1}(X_{\alpha_{0}})}&{{}\subset\left(\bigsqcup_{\alpha\in A}T_{X_{z}}^{*}X\right)\widehat{+}T_{X_{z_{0}}}^{*}X}\\ {}&{{}=\bigcup_{\alpha\in A}\left(T_{X_{z}}^{*}X\widehat{+}T_{X_{z_{0}}}^{*}X\right)}\\ {}&{{}\subset\bigsqcup_{\alpha\in A}T_{X_{z}}^{*}X.}\\ \end{aligned}
\]

然后结果如下。 ⑨

\statementlabel{定理 8.4.2}让 $ F \in \mathrm{Ob}(\mathbb{D}^{b}(X)) $  。以下条件是等效的。

(i) 存在由子解析子集覆盖 $ X = \bigcup_{i \in I} X_i $  的局部有限，使得对于所有 $ j \in \mathbb{Z} $  、所有 $ i \in I $  ，层 $ H^j(F)|_{X_i} $  是局部常值的。

(ii) SS(F) 包含在闭圆锥亚解析各向同性子集中。

(iii) SS(F) 是闭圆锥亚解析拉格朗日子集。

\statementlabel{证明}(i)  $\Rightarrow$  (ii) 根据命题 8.3.20、8.4.1 和 8.3.10。

(ii)  $ \Rightarrow $  (i) 由推论 8.3.22 和命题 8.4.1 得出。

(iii)  $ \Rightarrow $  (ii) 是显而易见的。

(ii)  $\Rightarrow$  (iii) 通过对合性定理（定理 6.5.4）和命题 8.3.13。

\statementlabel{定义 8.4.3}让  $F$  成为  $\mathbf{D}^{b}(X)$  的对象。

(i) 如果 F 满足定理 8.4.2 的等价条件，则称 F 是弱  $ \mathbb{R} $  可构造的（简称为 w-  $ \mathbb{R} $  可构造的）。

(ii) 如果 $F$  是可构造的，并且对于每个 $x \in X$  、 $F_x$  是完美复形，则可以说 $F$  是可构造的 $\mathbb{R}$  。

注意，当基环是诺特环时，条件(ii)相当于说所有上同调群 $ H^{j}(F_{x}) $ 都是有限生成的。

其中一个用 $ \mathbf{D}_{w-\mathbf{R}-c}^{b}(X) $ （或 $ \mathbf{D}_{\mathbf{R}-c}^{b}(X) $ ）表示 $ \mathbf{D}^{b}(X) $ 的完整三角子范畴，由w- $ \mathbf{R} $ 可构造（或 $ \mathbf{R} $  -可构造）对象组成。

 $X$  上的束  $F$  被称为  $w\mathbb{R}$  可构造（分别为  $\mathbb{R}$  -可构造），如果是这样，则被视为  $\mathbf{D}^{b}(X)$  的对象。我们用 $w\mathbb{R}\text{-}\text{Cons}(X)$ （或 $\mathbb{R}\text{-}\text{Cons}(X)$ ）表示 $\mathfrak{Mod}(A_{X})$ 的满子范畴，由 $w\mathbb{R}$  -可构造（或 $\mathbb{R}$  -可构造）层组成。

令 $ u: F \to G $  为 $ X $  、 $ F $  和 $ G $  上的层态射， $ w $  -R 可构造。可以立即证明 Ker  $ u $  、 Im  $ u $  和 Coker  $ u $  是  $ w $  -R 可构造的。此外，如果  $ 0 \to F' \to F \to F'' \to 0 $  是  $ X $  上的层的精确序列，并且  $ F' $  和  $ F'' $  是  $ w $  -R 可构造的，那么  $ F $  也是可构造的。因此  $ w $  -R-Cons(  $ X $  ) 是阿贝尔范畴。

如果基环  $A$  是诺特环，则将  $w\text{-}\mathbb{R}\text{-}\mathrm{Cons}(X)$  替换为  $\mathbb{R}\text{-}\mathrm{Cons}(X)$  时，结果相同。

示例 8.4.4。令 Z 为 X 的局部闭子解析子集。则层 $ A_{Z} $  是 $ \mathbb{R} $  可构造的。

\statementlabel{定理 8.4.5}(i) 自然函子 $ \mathbf{D}^{b}(w\text{-}\mathbb{R}\text{-}\mathbb{C}\text{ons}(X)) \to \mathbf{D}_{w\text{-}\mathbb{R}-c}^{b}(X) $  是范畴的等价。

(ii) 假设基环 $A$  是诺特环。那么自然函子 $\mathbf{D}^{b}(\mathbb{R}_-\text{Cons}(X)) \to \mathbf{D}_{\mathbf{R}-c}^{b}(X)$  就是范畴的等价。

\statementlabel{证明}(i) 我们必须证明以下断言 (a) 和 (b)：

(a) 对于任何  $ F \in \mathrm{Ob}(\mathbb{D}_{w-\mathbb{R}-c}^b(X)) $  ，都存在  $ G \in \mathrm{Ob}(\mathbb{D}^b(w-\mathbb{R}-\mathrm{Cons}(X))) $  ，它与  $ \mathbb{D}^b(X) $  中的  $ F $  同构，

(b) 对于  $ \mathbf{D}^b(w\text{-}\mathbb{R}\text{-}\text{Cons}(X)) $  中的任何对象  $ F $  和  $ G $  ，我们有同构：

\[
\mathrm{H o m}_{{\tt D}^{b}(w-{\tt R}-\mathfrak{C o n s}(X))}(F,G)\xrightarrow{\sim}\mathrm{H o m}_{{\tt D}^{b}(X)}(F,G)\enspace.
\]

让我们证明(a)。

我们选择子解析集  $ X = \bigcup_{j \in J} X_j $  的局部有限覆盖，使得  $ H^k(F)|_{X_j} $  对于所有  $ k $  和所有  $ j $  都是局部常数。应用命题 8.2.5，我们发现一个单纯复形  $ \mathbf{S} = (S, \mathcal{A}) $  和一个同胚  $ i: |S| \approx X $  使得：

\[
i(|\sigma|)\mathrm{i s~s u b a n a l y t i c~f o r~e v e r y~}\sigma\in\varDelta,
\]

\[
\mathrm{f o r~a n y}~\sigma\in\varDelta,\mathrm{t h e r e~e x i s t s}~j\in J~\mathrm{w i t h}~i(|\sigma|)\subset X_{j}
\]

因此  $ i^{-1}(F) $  是  $ \mathbf{D}_{w-S-c}^{b}(|\mathbf{S}|) $  的对象，并且根据定理 8.1.10 存在  $ G \in \mathrm{Ob}(\mathbf{D}^{b}(w-\mathrm{Cons}(\mathbf{S}))) $  同构于  $ i^{-1}(F) $  。那么  $ i_*G $  是  $ \mathbf{D}^{b}(w-\mathbb{R}-\mathrm{Cons}(X)) $  与  $ F $  同构的对象。

让我们证明(b)。

令 $F'$  和 $G'$  为 $w$  -  $\mathbb{R}$  可构造层的两个有界复合体。存在一个单纯复形  $\mathbf{S} = (S, \mathcal{A})$  和一个同胚  $i: |S| \simeq X$  满足 (8.4.2) 并且还：

\[
\begin{cases}{i^{-1}(F^{\cdot}){~a n d~}i^{-1}(G^{\cdot}){~a r e~c o m p l e x e s~o f~}w\text{-}\mathbb{S}\text{-c o n s t r u c t i b l e}}\\ {s h e a v e s\enspace.}\\ \end{cases}
\]

考虑一下图表：

\[
\begin{array}{c c c}{\operatorname{H o m}_{\mathsf{D}^{b}(w-\mathfrak{C o n s}(\mathbb{S}))}(i^{-1}F^{*},i^{-1}G^{*})}&{\xrightarrow{\quad}}&{\operatorname{H o m}_{\mathsf{D}^{b}(w-\mathsf{R}-\mathfrak{C o n s}(X))}(F^{*},G^{*})}\\ {\bigg\downarrow^{u}}&{\bigg\downarrow^{w}}\\ {\operatorname{H o m}_{\mathsf{D}^{b}(|\mathbb{S}|)}(i^{-1}F^{*},i^{-1}G^{*})\quad}&{\xrightarrow{\quad\sim\quad}}&{\operatorname{H o m}_{\mathsf{D}^{b}(X)}(F^{*},G^{*})}\\ \end{array}
\]

根据定理 8.1.10， $u$  是同构。因此 $w$  是满射的。让我们证明 $w$  是单射的。让  $\varphi \in \text{Hom}_{\mathbb{D}^{b}(w - \mathbb{R} - \mathbb{C}_{\text{ons}(X)})} (F^*, G^*)$  与  $w(\varphi) = 0$  。一个通过准同构  $G' \cong G''$  和态射  $\varphi': F' \to G'$  来表示  $\varphi$  。将  $\varphi$  替换为  $\varphi'$  我们可以从一开始就假设  $\varphi$  是由从  $F'$  到  $G'$  的态射给出的。

那么  $ \varphi = v \circ i^{-1}(\varphi) $  和  $ w(\varphi) = 0 $  意味着  $ i^{-1}(\varphi) = 0 $  ，因此  $ \varphi = 0 $  。

(ii) 借助定理 8.1.11 的证明是类似的。 □

根据定理 8.4.2，我们看到 w-R-可构造是一个微局部属性。因此我们可以充分利用第五章和第六章的结果来研究这些对象的功能特性。

\statementlabel{命题 8.4.6}(i) 令 $f: Y \to X$  为流形的态射。

（一） $ \text{令 } F \in \text{Ob}(\mathbb{D}_{w-\mathbb{R}-c}^b(X)) $ 。那么  $ f^{-1} F $  和  $ f^! F $  属于  $ \mathbb{D}_{w-\mathbb{R}-c}^b(Y) $  。

(i)ₙ 令  $ G \in \text{Ob}(\mathbb{D}_{w-\mathbb{R}-c}^b(Y)) $  并假设  $ f $  对  $ \text{supp}(G) $  是正确的。然后 $ Rf_s G \in \text{Ob}(\mathbb{D}_{w-\mathbb{R}-c}^b(X)) $ 。

(ii) 让 $F$  和 $G$  属于 $\mathbf{D}_{w-\mathbf{R}-c}^{b}(X)$  。那么  $G\otimes^L F$  和  $R\mathcal{H}om(G,F)$  属于  $\mathbf{D}_{w-\mathbf{R}-c}^{b}(X)$  ，  $\mu\mathcal{H}om(G,F)$  属于  $\mathbf{D}_{w-\mathbf{R}-c}^{b}(T^*X)$  。

(iii) 令 $ E \to X $  为向量丛并令 $ F \in \text{Ob}(\mathbb{D}_{\mathbb{R}^+}^b(E)) $  。假设  $ F $  是  $ w $  -  $ \mathbb{R} $  -可构造的。那么  $ F^\wedge $  ，它的 Fourier-Sato 变换，是  $ w $  -  $ \mathbb{R} $  -可构造的。

\statementlabel{证明}使用定理 8.3.17 和命题 8.3.11， (i)  $ _{a} $  由推论 6.4.4 得出； (i)  $ _{b} $  由命题 5.4.4 得出，(ii) 由推论 6.4.5 和推论 6.4.3 得出，(iii) 由定理 5.5.5 得出。 ⑨

为了研究 R 可构造层，我们需要一个引理。首先我们介绍一下：

\[
B_{\varepsilon}=\left\{x\in\mathbb{R}^{n};\left|x\right|<\varepsilon\right\},
\]

\[
B_{r},=\left\{x\in\mathbb{R}^{n},\varepsilon<\vert x\vert<\varepsilon\right\},
\]

\[
S_{\varepsilon}=\{x\in\mathbb{R}^{n};|x|=\varepsilon\}\ ,
\]

我们分别用  $ \bar{B}_{\varepsilon} $  和  $ \bar{B}_{\varepsilon^{\prime},\varepsilon} $  表示  $ B_{\varepsilon} $  和  $ B_{\varepsilon^{\prime},\varepsilon} $  的闭包。

\statementlabel{引理 8.4.7}令 $ F \in \text{Ob}(\mathbb{D}_{w-\mathbb{R}-c}^b(X)) $  和 $ \varphi: X \to \mathbb{R}^n $  为实解析函数。假设 $ \varphi|_{\text{supp}(F)} $  是正确的。那么我们就有了自然同构：

\[
R\varGamma(\varphi^{-1}(\bar{B}_{\varepsilon});F)\xrightarrow{\sim}R\varGamma(\varphi^{-1}(B_{\varepsilon});F)\xrightarrow{\sim}R\varGamma(\varphi^{-1}(0);F)~,\qquad{f o r~}0<\varepsilon\ll1\quad.
\]

\[
R\Gamma_{\varphi^{-1}(0)}(X;F)\xrightarrow{\simeq}R\Gamma_{\varphi^{-1}(\bar{B}_{\varepsilon})}(X;F)\xrightarrow{\simeq}R\Gamma_{c}(\varphi^{-1}(B_{\varepsilon});F)~,\quad{~f o r~}0<\varepsilon^{\prime}<\varepsilon\ll1\quad.
\]

\[
R\varGamma(\varphi^{-1}(B_{0,\varepsilon});F)\xrightarrow{\sim}R\varGamma(\varphi^{-1}(B_{\varepsilon^{\prime\prime},\varepsilon^{\prime}});F)\xrightarrow{\sim}R\varGamma(S_{\varepsilon^{\prime\prime}};F)\enspace,
\]

\[
for0\leq\varepsilon<\varepsilon^{\prime \prime \prime}<\varepsilon^{\prime}\leq\varepsilon\ll1.
\]

\statementlabel{证明}通过将  $\varphi$  替换为  $|\varphi|^2$  ，我们可以假设  $n=1$  。然后应用微局域 Bertini-Sard 定理（命题 8.3.12）和  $A=\mathrm{SS}(F)$  以及微局域 Morse 引理（推论 5.4.19）。   $\square$ 

请注意，当  $X$  在  $\mathbb{R}^n$  中打开，以及  $\varphi(x) = x - x_\circ$  对于某些  $x_\circ \in X$  时，此引理尤其适用。

\statementlabel{命题 8.4.8}令  $f: Y \to X$  为流形的态射，并令  $G \in \mathrm{Ob}(\mathbf{D}_{\mathbf{R}-c}^b(Y))$  。假设  $f$  对于  $\mathrm{supp}(G)$  是正确的。然后 $R f_* G \in \mathrm{Ob}(\mathbf{D}_{\mathbf{R}-c}^b(X))$ 。

\statementlabel{证明}根据命题8.4.6，我们只需证明对于每个 $x \in X$ ， $R\Gamma(f^{-1}(x); G|_{f^{-1}(x)})$  都是完美复形。通过用图表分解 $f$ ，我们可以假设 $f$ 是平滑的，因此我们可以假设 $X = \{\text{pt}\}$ 。 （我们使用明显的注释，如果  $S \subset Y$  是子流形，则  $G|_{S}$  属于  $\mathbf{D}_{\mathbf{R}-c}^{b}(S)$  。）通过封闭嵌入  $Y \hookrightarrow \mathbb{R}^n$  ，我们可以假设  $Y = \mathbb{R}^n$  。

最后，通过归纳论证，我们可以假设  $n = 1$  。换句话说，还有待证明如果 $Y = \mathbb{R}$ 和 $G$ 有紧凑的支持， $R\Gamma(Y; G)$ 是完美的。由于  $SS(G)$  是亚解析且各向同性的，因此  $\mathbb{R}$  中存在有限序列  $t_1 < \cdots < t_N$  ，使得  $G$  由  $[t_1, t_N]$  支持，并且  $G|_{t_{i-1}, t_i}$  对于所有  $i \leq N$  都是常数。通过考虑不同的三角形：

\[
G_{]-\infty,t_{i}[}\longrightarrow G\longrightarrow G_{[t_{i},+\infty[}\xrightarrow[+1]{}\qquad{\mathrm{a n d}}\qquad G_{]t_{i},+\infty[}\longrightarrow G_{[t_{i},+\infty[}\longrightarrow G_{\{t_{i}\}}\xrightarrow[+1]{},
\]

并且通过归纳论证，足以证明 $ R\Gamma(Y; G)_{t_{i-1}, t_i} $ 是完美的。由于对于任何  $ t \in [t_{i-1}, t_i] $  该复形与  $ R\Gamma_{t} $  (Y; G) 同构，因此结果来自区分的三角形

\[
R\varGamma_{\{t\}}(Y;G)\longrightarrow R\varGamma(\rbrack t-\varepsilon,t+\varepsilon[;G)\longrightarrow R\varGamma(\rbrack t-\varepsilon,t[\cup]\underset{}{t,t+\varepsilon}{;G})\underset{++}{\longrightarrow}\ .
\]

\statementlabel{命题 8.4.9}让 $ F \in \mathrm{Ob}(\mathbb{D}_{w-\mathbb{R}-c}^b(X)) $  。考虑以下条件：

(i) F 是  $ \mathbb{R} $  可构造的。

(ii) F 是上同调可构造的。

(iii) 对于任何  $ x \in X $  ， $ R\Gamma_{\{x\}}(X; F) $  都是完美的

(iv) DF =  $ R\mathcal{H}om(F, \omega_X) $  是  $ \mathbb{R} $  可构造的。

然后 (i)  $\Leftrightarrow$  (ii)  $\Rightarrow$  (iii)  $\Rightarrow$  (iv)。如果基环  $A$  是一个域或  $\mathbb{Z}$  ，则这四个条件是等效的。

\statementlabel{证明}(i)  $\Rightarrow$  (ii) 根据引理 8.4.7，仍需证明 (i)  $\Rightarrow$  (iii)。利用引理 8.4.7 的符号，在  $x_o \in X$  邻域中选择坐标系后，我们得到了  $0 < \varepsilon < 1$  的可区分三角形：

\[
R\varGamma_{\{x_{\circ}\}}(X;F)\longrightarrow R\varGamma(B_{\varepsilon};F)\longrightarrow R\varGamma(B_{\varepsilon}\backslash\{x_{\circ}\};F)\xrightarrow[+1]{}
\]

和同构  $ R\Gamma(B_\varepsilon; F) \simeq F_{x_0}, R\Gamma(B_\varepsilon \setminus \{x_0\}); F) \simeq R\Gamma(S_\varepsilon; F) (0 < \varepsilon' < \varepsilon) $  。根据命题 8.4.8，复数  $ R\Gamma(S_\varepsilon; F) $  是完美的，结果如下。

(ii)  $\Rightarrow$  (i) 和 (ii)  $\Rightarrow$  (iii) 是显而易见的。

(iii)  $ \Rightarrow $  (iv) 通过与命题 3.4.3 相同的证明，我们看到

\[
\begin{array}{r}{(\mathrm{D}F)_{x}\simeq R\mathrm{H o m}(R\varGamma_{\{x\}}(X;F),A)\enspace.}\end{array}
\]

复数 $ (DF)_{x} $  是完美复数的对偶，它本身就是完美的。

(iv)  $\Rightarrow$  (i) 对于 $x \in X$  、 $R\Gamma_{\{x\}}(X; DF) = RHom(F_x, A)$  。由于这个复数是完美的， $F_x$  也是完美的（参见练习 I.31 的  $A = \mathbb{Z}$  ）。   $\square$ 

\statementlabel{命题 8.4.10}(i) 令  $f: Y \to X$  为流形的态射，并令  $F \in \mathrm{Ob}(\mathbf{D}_{\mathbf{R}-c}^b(X))$  。那么  $f^{-1}F$  和  $f^! F$  属于  $\mathbf{D}_{\mathbf{R}-c}^b(Y)$  。

(ii) 让 $F$  和 $G$  属于 $\mathbf{D}_{\mathbb{R}-c}^{b}(X)$  。那么  $G\otimes^L F$  和  $R\mathcal{H}om(G,F)$  属于  $\mathbf{D}_{\mathbb{R}-c}^{b}(X)$  。

\[
f^{-1}F,\;f^{!}F
\]

(i) 自从 $ (f^{-1}F)_{y}=F_{f(y)} $ 以来，这个综合体是完美的。 （我们在证明命题8.4.8的过程中已经说过了这一点。）

我们有 $f^!D_X D_X F = f^!R\mathcal{H}om(D_X F, \omega_X) \simeq R\mathcal{H}om(f^{-1}D_X F, \omega_Y) \simeq D_Y(f^{-1}D_X F)$ 。因此 $f^!F$  是 $\mathbb{R}$  -可由命题8.4.9 构造。

(ii) 显然 $ G \otimes^L F $  是 $ \mathbb{R} $  可构造的。

由于  $ R\mathcal{H}om(G,F) = D(DF \otimes^L G) $  （命题 3.4.6），该对象是  $ \mathbb{R} $  -可由命题 8.4.9 构造。   $ \square $ 

\statementlabel{推论 8.4.11}让 $ F \in \mathrm{Ob}(\mathbb{D}_{R-c}^{b}(X)) $  。

(i) 令 K 为 X 的紧亚解析子集。则  $ RG(X; F) $  和  $ RG(K; F) $  是完美复形。

(ii) 令  $\Omega$  为  $X$  的相对紧凑的开放子解析子集。那么 $R\Gamma(\Omega; F)$  和 $R\Gamma_{c}(\Omega; F)$  就是完美的复合体。

\statementlabel{证明}应用命题 8.4.8 和 8.4.10。 （回想练习 II.19 的同构。） ☐

\statementlabel{命题 8.4.12}(i) 设  $ \tau: E \to Z $  为向量束，并设  $ F \in \text{Ob}(\mathbf{D}_{\mathbf{R}^+}^b(E)) $  。假设  $ F $  是  $ \mathbb{R} $  可构造的。那么  $ F^\wedge $  ，它的 Fourier-Sato 变换，是  $ \mathbb{R} $  可构造的。

(ii) 让 $ F $  和 $ G $  属于 $ \mathbf{D}_{\mathbf{R}-c}^b(X) $  。那么 $ \mu h o m(G, F) $  属于 $ \mathbf{D}_{\mathbf{R}-c}^b(T^*X) $  。

\statementlabel{证明}(i) 我们保留 III §7 的注释。然后 $ F^\wedge = Rp_{2*}R\Gamma_P(p_1^{-1}F) $ 。用 $ i $  表示零嵌入 $ E^* \hookrightarrow E \times_Z E^* $  。由于对象  $ R\Gamma_P(p_1^{-1}F) $  在向量丛  $ E \times_Z E^* \to E^* $  上是圆锥曲线，因此根据命题 3.7.5，我们也有  $ F^\wedge \simeq i^{-1}R\Gamma_P(p_1^{-1}F) $ 。因此 $ F^\wedge $  是 $ \mathbb{R} $  -可由命题8.4.10 构造。

(ii) 根据命题 8.4.10 和前面的结果，还需要证明如果  $M$  是  $X$  的子流形，则  $v_M(F)$  是  $\mathbb{R}$  可构造的。

通过构造特化 $ \nu_M(F) $ ，就足以证明 $ R j_* j^{-1}F $ 是 $ \mathbb{R} $ 可构造的，当 $ j $ 是 $ \Omega \hookrightarrow X $ 和 $ \Omega = \{x \in X; \varphi(x) > 0\} $ 的开放嵌入 $ \Omega \hookrightarrow X $ 和 $ \Omega = \{x \in X; \varphi(x) > 0\} $ 用于 $ X $ 上的实函数 $ \varphi $ 和 $ d\varphi \neq 0 $ 时。这是从  $ R j_* j^{-1}F \simeq R \mathcal{H} o m(A_\Omega, F) $  和命题 8.4.10 (ii) 得出的。   $ \square $ 

为了结束本节，我们将研究对偶函子和函子  $ \mu h o m $  之间的一些关系。

\statementlabel{命题 8.4.13}令  $M$  为  $X$  的子流形，并令  $F \in \mathrm{Ob}(\mathbf{D}_{R-c}^{b}(X))$  。那么我们就有了自然同构：

（一） $ \nu_M(\mathbf{D}_X F) \simeq \mathbf{D}_{T_M X}(\nu_M(F)) $ ，

（二） $ \mu_M(D_X F) \simeq D_{T_M^*X}(\mu_M(F))^a \otimes \omega_{M/X} $ 。

（回想一下，“a”表示  $ T_{M}^{*}X $  上对映图的直像。）

\statementlabel{证明}我们使用 IV §2 的符号，并应用引理 4.2.1。然后：

\[
\begin{aligned}{\nu_{M}(\mathbf{D}_{X}F)}&{{}=s^{-1}R j_{*}\tilde{p}^{-1}R\mathcal{H}o m(F,\omega_{X})}\\ {}&{{}\simeq s^{-1}R j_{*}R\mathcal{H}o m(\tilde{p}^{!}F,\tilde{p}^{!}\omega_{X})}\\ {}&{{}\simeq s^{-1}R\mathcal{H}o m(j_{!}\tilde{p}^{!}F,\omega_{\tilde{X}_{M}})}\\ {}&{{}\simeq R\mathcal{H}o m(s^{!}j_{!}\tilde{p}^{!}F,\omega_{T_{M}X})}\\ {}&{{}\simeq\mathbf{D}_{T_{M}X}(\nu_{M}(F)).}\\ \end{aligned}
\]

在证明过程中，我们使用  $ G = R\mathcal{H}om(j_! \tilde{p}^! F, \omega_{\tilde{X}_M}) $  是 \textbackslash{}textit\{R-constructible\} 的事实，这是从命题 8.4.12 的证明得出的，以及公式  $ s^{-1}DG \simeq Ds^! G $  ，（参见练习 VIII.3）。

(ii) 我们将应用(i)、公式(3.7.9)和命题3.7.12。然后：

\[
\begin{aligned}{\mu_{M}(\mathbf{D}F)}&{{}=(\nu_{M}(\mathbf{D}F))^{\vee a}\otimes\omega_{T_{M}^{*}X/M}^{\otimes-1}}\\ {}&{{}\simeq(\mathbf{D}\nu_{M}(F))^{\vee a}\otimes\omega_{M/X}}\\ {}&{{}\simeq\mathbf{D}\mu_{M}(F)^{a}\otimes\omega_{M/X}.}\\ \end{aligned}
\]

\statementlabel{命题 8.4.14}假设  $F$  和  $G$  属于  $\mathbf{D}_{\mathbf{R}-c}^{b}(X)$  。那么我们就有了自然同构：

（一） $ \mu h o m(F,G) \simeq \mu h o m(\mathbf{D}_{X}G,\mathbf{D}_{X}F)^{a} $ ，

（二） $ \mathbf{D}_{T^{*}X}(\mu h o m(F,G)) \simeq \mu h o m(G,F) \otimes \omega_{X} $ 。

\statementlabel{证明}(i) 事实上，这个公式在 F 和 G 是上同调可构造的较弱假设下是正确的（参见练习 IV.4），并且可以直接从命题 3.4.6 得出。

(ii) 根据命题 8.4.13，我们有：

\[
\begin{aligned}{\operatorname{D}_{T^{*}X}(\mu h o m(F,G))}&{{}\simeq\operatorname{D}_{T^{*}X}\mu_{\varDelta}(G\boxtimes D F)}\\ {}&{{}\simeq\mu_{\varDelta}(\operatorname{D}_{X\times X}(G\boxtimes D F))^{a}\otimes\omega_{\varDelta/X\times X}^{\otimes-1}}\\ {}&{{}\simeq\mu_{\varDelta}(\operatorname{D}G\boxtimes F)^{a}\otimes\omega_{X}}\\ {}&{{}\simeq\mu_{\varDelta}(F\boxtimes\operatorname{D}G)\otimes\omega_{X}}\\ {}&{{}\simeq\mu h o m(G,F)\otimes\omega_{X}~.\quad\Box}\\ \end{aligned}
\]

\subsection{$\mathbb{C}$-构造层}
在本节中，我们将考虑复解析流形。由于复杂的几何超出了本书的范围，因此我们将相当简短，并且我们不会给出所有详细的证明。

如果  $X$  是一个复流形，我们用  $X^R$  表示实底层流形，并且我们参考  $XI\S1$  来研究  $X$  和  $X^R$  之间的关系， $(T^*X)^R$  和  $T^*(X^R)$  之间的同构等。如果没有混淆的风险，我们应该写  $X$  而不是  $X^R$  。

令  $S$  为  $X$  的局部闭子集。如果 $\overline{S}$  和 $\overline{S} \setminus S$  是复解析子集，我们就说 $S$  是 $\mathbb{C}$  解析的。 （关于解析几何的主要概念，我们参考 Cartan [2]。）特别是，如果  $S$  是  $X$  中的  $\mathbb{C}$  解析，则  $S$  是  $X^R$  中的亚解析。

让  $\mathcal{A}$  成为  $T^*X$  的子集。如果  $\mathcal{A}$  在  $T^*(X^\mathbb{R})$  中具有相应的属性，我们就说  $\mathcal{A}$  是  $\mathbb{R}^+-$  圆锥曲线。如果  $\mathbb{C}^\times$  的作用使  $\mathcal{A}$  局部不变，那么我们说  $\mathcal{A}$  是局部  $\mathbb{C}^*-$  圆锥曲线，也就是说，对于任何  $\mathbb{C}^*-$  轨道  $S$  ， $\mathcal{A}\cap S$  在  $S$  中开。如果 $\mathcal{A}$  是 $\mathbb{C}^*-$  轨道的并集，我们就说 $\mathcal{A}$  是 $\mathbb{C}^*-$  圆锥曲线。

如果  $\mathcal{A}$  是复数解析子流形，则  $\mathcal{A}$  是对合子空间（分别是拉格朗日、各向同性），如果  $T_p\mathcal{A}$  是  $T^*X$  在每个  $p\in\mathcal{A}$  处的复对合子空间（分别是拉格朗日、各向同性）。假设  $\mathcal{A}$  是  $T^*X$  的封闭  $\mathbb{C}$  解析子集。那么 $\mathcal{A}$  是 $\mathbb{R}^+$  -圆锥曲线当且仅当 $\mathcal{A}$  是 $\mathbb{C}^*$  -圆锥曲线，在这种情况下我们简单地说 $\mathcal{A}$  是圆锥曲线。

另请注意，如果 $A$  是复解析且 $A_{\mathrm{reg}}$  是拉格朗日，则 $A$  是拉格朗日（参见练习VIII 8）。

\statementlabel{命题 8.5.1}令 S 和 Y 为 X 的两个 C 解析子集。则正锥  $ C(Y, S) $  （在  $ TX^{R} $  中定义，使用 X 的真实结构）在 TX 中是 C 解析的。

对于证明，我们参考 Whitney [2]。

请特别注意，  $C(Y,S)$  对于  $\mathbb{C}^{\times}$  的操作是不变的。

\statementlabel{命题 8.5.2}令 $\Omega$  为 $T^*X$  的开子集，并让 $A$  为 $\Omega$  的局部 $\mathbb{C}^\times$  圆锥对合闭子集。假设  $A$  包含在  $\Omega$  的闭合  $\mathbb{R}^+$  圆锥亚解析各向同性子集中。那么 $A$ 就是一个复杂的解析集。

\statementlabel{证明}根据命题 8.3.13， $A$  在  $\Omega$  中是亚解析和拉格朗日的。对于任何  $p \in A_{\text{reg}}$  ，  $T_p A$  是  $T_p T^* X$  的实拉格朗日平面。由于  $A$  是局部  $\mathbb{C}^\times$  -圆锥曲线，因此  $T_p A = (T_p A)^\perp$  包含  $\mathbb{C} \cdot H(\alpha)$  ，其中  $\alpha$  是  $T^* X$  上的复数 1-形式。因此  $\alpha|_{A_{\text{reg}}} = 0$  ，意味着  $\mathrm{d}\alpha|_{A_{\text{reg}}} = 0$  和  $\mathrm{Re}(\mathrm{d}\alpha)|_{T_p A + \sqrt{-1} T_p A} = 0$  ，读作  $\sqrt{-1} T_p A \subset (T_p A)^\perp = T_p A$  。这表明  $T_p A$  是任何  $p \in A_{\text{reg}}$  的复向量子空间。因此 $A_{\text{reg}}$  是一个复流形。设置 $n = \dim_C X$ 。然后 $\dim_R(A \setminus A_{\text{reg}}) \leqslant 2n - 1$ 。令  $S$  为  $(A \setminus A_{\text{reg}})_{\text{reg}}$  的  $(2n - 1)$  维连通分量的并集，接下来，令  $S'$  为  $S$  的最大开子集，其中  $(A_{\text{reg}}, S)$  满足  $\mu$  条件。由于  $S'$  在  $S$  、  $\dim((\mathcal{A} \setminus A_{\text{reg}})\setminus S') \leqslant 2n - 2$  中是密集的。

对于任何  $p \in S^{\prime}$  ，选择  $\mathcal{A}_{\text{reg}}$  中收敛到  $p$  的序列  $\{p_n\}$  ，使得  $T_{p_n} \mathcal{A}_{\text{reg}}$  收敛到  $\tau \subset T_p T^* X$  ，以及  $\tau \supset T_p S^{\prime}$  （参见练习 VIII.12）。由于  $T_{p_n} \mathcal{A}_{\text{reg}}$  是复向量空间，因此  $\tau$  也是复向量空间。因此我们得到：

\[
\mathrm{d i m}_{\mathbf{c}}(T_{p}S^{\prime}+\sqrt{-1}T_{p}S^{\prime})=n~,\qquad\mathrm{f o r~a n y~}p\in S^{\prime}~.
\]

让  $ S'_C $  成为  $ S' $  的复化。嵌入  $ S' \hookrightarrow T^*X $  扩展到映射  $ S'_C \to T^*X $  ，并且如果必要的话缩小  $ S'_C $  ，该映射具有恒定的等级 (8.5.1)。因此，它的图像是  $ S' $  邻域上的复杂子流形。让我们用 Z 表示它。同样的论证表明，任何包含  $ S' $  的复子流形都包含 Z，位于  $ S' $  的邻域上。请注意，Z 是各向同性的，因此是拉格朗日函数。现在我们将进行如下操作。我们将证明：

(a)  $ S' \cap \overline{A_{\mathrm{reg}}}\backslash Z $  在  $ S' $  中并不密集，

（b） $ S = \varnothing $ ，

(c) A 是复解析的。

(a) 让我们用矛盾来论证。否则，  $S'\cap\mathcal{A}_{\text{reg}}\setminus Z$  包含  $S'$  的非空开子集。因此我们可以假设  $\mathcal{A}_{\text{reg}}\setminus Z\to S$  。由于  $\mathcal{A}_{\text{reg}}\setminus Z$  是亚解析的，因此存在适当的实解析图  $f:W\to T^*X$  使得  $f(W)=\mathcal{A}_{\text{reg}}\setminus Z$  。我们可以假设  $f^{-1}(\mathcal{A}_{\text{reg}}\setminus Z)$  在  $W$  中是开放且密集的。现在让我们对  $W$  进行复化  $W_C$  并将  $f$  扩展为全纯映射  $\tilde{f}\colon W_C\to T^*X$  。对于  $f^{-1}(\mathcal{A}_{\text{reg}}\setminus Z)$  的开稠子集中的  $w$  ，  $\text{Im}(T_w W_C\to T_{f(w)}T^*X)$  与  $T_{f(w)}(\mathcal{A}_{\text{reg}})$  一致，我们得到：

\[
\mathrm{the~rank~of~}\tilde{f}\mathrm{~is~equal~or~less~than~n~}.
\]

由于  $f(W) \supset S$  ，存在  $w \in W$  使得  $T_w W \to T_{f(w)} T^* X$  的图像包含  $T_{f(w)} S$  。因此，在这样的点  $w$  ，  $\tilde{f}$  的秩是  $n$  ，并且  $\tilde{f}$  在  $w$  的邻域中具有恒定的秩。对于 $w$  的足够小的邻域 $U$  来说， $\tilde{f}(U)$  是一个 $n$  维复数子流形，包含 $f(w)$  邻域上的 $S$  。因此  $\tilde{f}(U) = Z$  位于  $f(w)$  和  $\mathcal{Q} \neq \tilde{f}(U \cap f^{-1}(A_{\text{reg}} \setminus Z)) \subset Z$  的邻域上，这是一个矛盾。

(b) 假设  $S \neq \emptyset$  。然后存在 $p \in S' \backslash (A_{\text{reg}} \setminus Z)$ 。在  $p$  、  $A_{\text{reg}} \subset Z$  和  $A = A_{\text{reg}}$  、  $A \subset Z$  的邻域上。根据引理 8.3.14， $A$  在  $Z$  中打开。因此  $A = Z$  位于  $p$  的邻域，这与  $p \in S' \subset A \setminus A_{\text{reg}}$  相矛盾。

（c）由（b）， $ \dim_{\mathbb{R}}(\varLambda \setminus \varLambda_{\mathrm{reg}}) \leq 2n - 2 $ 。将 $ \overline{\mathrm{Remmert-Stein}} $  [1]的定理应用于复解析集的扩展，我们得到 $ \varLambda_{\mathrm{reg}} $ 是复解析集。   $ \square $ 

作为最后结果的应用，考虑 X 的闭合 C 解析子集 S。集合  $ T_s^*X $  由 (8.3.1) 定义。在 X 上局部，S 被定义为  $ \{x \in X; f_j(x) = 0, j = 1, \ldots, p\} $ ，其中  $ f_j $  是全纯函数。因此

\[
T_{S}^{*}X=\overline{\left\{\left(x;\sum_{j=1}^{p}\lambda_{j}\mathbf{d}f_{j}(x)\right);\lambda_{j}\in\mathbb{C},f_{j}(x)=0\forall j\right\}}
\]

根据命题 8.3.1 和 8.5.2，我们看到  $ T_S^*X $  是  $ T^*X $  的闭圆锥 C 解析和 C 拉格朗日子集。

现在我们可以将§3 的大部分结果扩展到复杂的情况。

\statementlabel{命题 8.5.3}令  $A$  为  $T^*X$  的闭圆锥曲线  $\mathbb{C}$  解析子集。那么存在一个有限族 $\{X_j\}$  的封闭 $\mathbb{C}$  的解析子集 $X$  使得 $A \subset \bigcup_j T_{X_j}^*X$  。此外，对于  $X$  的任何  $\mathbb{C}$  解析子集  $Y$  ，都存在  $\mathbb{C}$  解析流形  $Y_\circ \subset Y$  ，在  $Y$  中开且稠，使得  $A \cap \pi^{-1}(Y_\circ) \subset T_Y^*X$  。

证明与命题 8.3.10 相同。

\statementlabel{命题 8.5.4}令 $X = \bigcup_{j \in J} X_j$  为 $X$  和 $\mathbb{C}$  解析子集的局部有限覆盖。然后存在比覆盖层更精细的 $\mu$  -分层 $X = \bigsqcup_{\alpha \in A} X_\alpha$  ，因此对于所有 $\alpha \in A$  ， $X_\alpha$  都是复流形。

证明类似于定理8.3.20，使用命题8.5.3 和8.5.1。

现在我们可以陈述本节的主要结果。

\statementlabel{定理 8.5.5}让 $ F \in \mathrm{Ob}(\mathbb{D}^{b}(X)) $  。以下条件是等效的。

(i) 存在由 C 解析子集覆盖 $ X = \bigcup_{j \in J} X_j $  的局部有限，使得对于所有 $ j \in J $  、所有 $ k \in \mathbb{Z} $  ，层 $ H^k(F)|_{X_i} $  是局部常值的。

(ii) SS(F) 包含在封闭的  $\mathbb{C}^{\times}$  -圆锥亚解析  $\mathbb{R}$  -各向同性子集  $\Lambda$  中。

(iii) SS(F) 是闭圆锥 C 解析拉格朗日子集。

(iv)  $ F \in \text{Ob}(\mathbf{D}_{w-\mathbb{R}-c}^b(X)) $  和 $ \text{SS}(F) $  是 $ \mathbb{C}^\times $  -圆锥曲线。

\statementlabel{证明}(i)  $\Rightarrow$  (ii) 应用命题 8.5.4 和 8.4.1。

(ii)  $ \Rightarrow $  (iii) 应用对合性定理 6.5.4 和命题 8.5.2。

(iii) ⇔ (iv) 根据命题 8.5.2。

(iii)  $ \Rightarrow $  (i) 根据命题 8.5.3、8.5.4 和 8.4.1。 ⑨

\statementlabel{定义 8.5.6}让 $ F \in \mathrm{Ob}(\mathbb{D}^{b}(X)) $  。

(i) 如果 F 满足定理 8.5.5 的等价条件，则称 F 是弱 C 可构造的（简称 w-C 可构造的）。

(ii) 如果 $F$  是 $w$  -  $\mathbb{C}$  可构造的，而且对于每个 $x \in X$  ， $F_x$  是完美复形，有人说 $F$  是 $\mathbb{C}$  可构造的。

其中一个用 $ \mathbf{D}_{w-C-c}^{b}(X) $ （或 $ \mathbf{D}_{C-c}^{b}(X) $ ）表示 $ \mathbf{D}^{b}(X) $ 的完整三角子范畴，由 $ w $  -C可构造（或C可构造）对象组成。

 $X$  上的束  $F$  被称为  $w$  -  $\mathbb{C}$  -可构造的（分别为  $\mathbb{C}$  -可构造的），如果是这样，则被视为  $\mathbf{D}^{b}(X)$  的对象。其中一个用 $w$  -  $\mathbb{C}$  -  $\mathsf{Cons}(X)$ （分别为 $\mathbb{C}$  -  $\mathsf{Cons}(X)$ ）表示此类层的范畴。请注意，自然态射  $\mathbf{D}^{b}(w$  -  $\mathbb{C}$  -  $\mathsf{Cons}(X)) \to \mathbf{D}_{w-C-c}^{b}(X)$  通常不是等价的。 （由于三角剖分定理在复杂情况下不成立，因此证明失败。）

\statementlabel{命题 8.5.7}(i) 令 $f: Y \to X$  为复流形的态射。

(a) 让 $ F \in \mathrm{Ob}(\mathbb{D}_{C-c}^{b}(X)) $  。那么  $ f^{-1}F $  和  $ f^!F $  属于  $ \mathbb{D}_{C-c}^{b}(Y) $  。

(b) 设  $ G \in \text{Ob}(\mathbb{D}_{\mathbb{C}-c}^b(Y)) $  并假设  $ f $  对于  $ \text{supp}(G) $  是正确的。然后 $ R f_* G \in \text{Ob}(\mathbb{D}_{\mathbb{C}-c}^b(X)) $ 。

(ii) 让 $F$  和 $G$  属于 $\mathbf{D}_{\mathbf{C}-c}^{b}(X)$  。那么  $G\otimes^{L}F$  和  $R\mathcal{H}om(G,F)$  属于  $\mathbf{D}_{\mathbf{C}-c}^{b}(X)$  ，  $\mu\mathcal{H}om(G,F)$  属于  $\mathbf{D}_{\mathbf{C}-c}^{b}(T^{*}X)$  。

(iii) 设  $ E \to X $  为复向量丛，并设  $ F \in \text{Ob}(\mathbf{D}_{\mathbf{R}^+}^b(E)) $  。假设  $ F $  是  $ \mathbb{C} $  可构造的。那么  $ F^\wedge $  ，它的 Fourier-Sato 变换，是  $ \mathbb{C} $  可构造的。而且(i)、(ii)、(iii)仍然保持将 $ \mathbf{D}_{\mathbf{C}-c}^b(\cdot) $ 处处替换为 $ \mathbf{D}_{\mathbf{w}-\mathbf{C}-c}^b(\cdot) $ 。

\statementlabel{证明}通过命题 8.4.8、8.4.10、8.4.12 足以证明弱 C 可构造对象的结果。

(i) 根据推论 6.4.4，SS(  $ f^{-1}F $  ) 包含在  $ f^{\#}(\text{SS}(F)) $  中，并且最后一个集合是  $ \mathbb{R}^{+} $  圆锥曲线，且各向同性。此外，由于 SS(F) 是  $ \mathbb{C}^{\times} $  -圆锥曲线，因此该集合是  $ \mathbb{C}^{\times} $  -圆锥曲线，而  $ f^{-1}F $  是 w-  $ \mathbb{C} $  -可由定理 8.5.5 构造。同样的证明也适用于  $ f^1F $  。 (ii) 和 (iii) 也同样得到证明。   $ \square $ 

还可以处理非真直像。

令 $f: Y \to X$  为复流形的态射，让 $\varphi: Y \to \mathbb{R}$  为实解析函数。设置如(5.4.13)：

\[
Y_{t}=\left\{y\in Y;\varphi(y)<t\right\}\;,\qquad\widetilde Y_{t}=\left\{y\in Y;\varphi(y)\leqslant t\right\}\;.
\]

用  $ j_t $  （分别为  $ \tilde{j}_t $  ）表示嵌入  $ Y_t \subsetneq Y $  （分别为  $ \tilde{Y}_t \subsetneq Y $  ）并设置  $ f_t = f \circ j_t $  、  $ \tilde{f}_t = f \circ \tilde{j}_t $  。让 $ t_\circ \in \mathbb{R} $  。

\statementlabel{命题 8.5.8}让  $ G \in \text{Ob}(\mathbf{D}_{C-c}^{b}(Y)) $  并假设：

(i) support(G) ∩ Y\_\{t\} 对于所有 t 在 X 上是真映射，

(ii) 对于所有  $ y \in Y \setminus Y_t $  ，我们有：

\[
\mathrm{d}\varphi(y)\notin\mathrm{S S}(G)+{}^{t}f^{\prime}\left(Y_{\substack{\times\\ X}}T^{*}X\right).
\]

那么命题5.4.17的结论(a)、(a')、(b)、(c)、(c')、(d)成立，并且 $ Rf_G $ 和 $ Rf_1G $ 属于 $ \mathbb{D}_{C-c}^b(X) $  。

\statementlabel{证明}由于 SS(G) 是  $\mathbb{C}^{\times}$  圆锥曲线，因此我们也有  $y\in Y\setminus Y_t$  ：

\[
-\mathrm{d}\varphi(y)\notin\mathrm{SS}(G)+t f^{\prime}\left(Y\times T^{*}X\right).
\]

因此我们可以应用命题 5.4.17 (i) 和 (ii)。此外，由于  $ R f_* G \simeq R f_*(G_{\tilde{Y}_t}) $  和  $ R f_! G \simeq R f_*(R\Gamma_{\tilde{Y}_t}(G)) $  ，根据命题 8.4.8，这些对象属于  $ \mathbf{D}_{\mathbf{R}-c}^b(X) $  。由于它们的微支撑包含在集合  $ f_{\pi}^{t} f^{t-1}(\text{SS}(G) \cap \pi^{-1}(\tilde{Y}_{t,c})) $  中，并且该集合是亚解析的、各向同性的和  $ \mathbb{C}^\times $  -圆锥的，因此仍然需要应用定理 8.5.5。   $ \square $ 

当  $ \dim X = 1 $  时，命题 8.5.8 的假设总是“局部”满足。更准确地说，有：

\statementlabel{命题 8.5.9}令  $f: Y \to X$  为复流形的态射，与  $\dim_{\mathbb{C}} X = 1$  ，并令  $G \in \mathrm{Ob}(\mathbb{D}_{\mathbb{C}-c}^b(Y))$  。令  $x_o \in X$  和  $K$  为  $f^{-1}(x_o)$  的紧凑子集。存在 $x_o$ （分别 $K$ ）和 $V \subset f^{-1}(U)$ 的开邻域 $U$ （分别 $V$ ），这样，用 $f_V$ 表示由 $f$ 引起的态射 $f_V: V \to U$ ，有：

(i)  $ Rf_V!G $  和  $ Rf_V*G $  属于  $ \mathbf{D}_{\mathbb{C}-c}^b(U) $  ，

（二） $ \mathrm{SS}(Rf_V!G) \cup \mathrm{SS}(Rf_V*G) \subset f_{V\pi}^{t} f_V^{\prime-1}(\mathrm{SS}(G)) $ 。

此外，当 Y 是  $\mathbb{C}^n$  和  $K = \{0\}$  的开子集时，我们可以取以 0 为中心、半径为  $\varepsilon \ll 1$  的开球作为  $V$  。

\statementlabel{证明}通过将  $f$  分解为  $Y \hookrightarrow Y \times X \to X$  ，我们可以从一开始就假设  $Y = Z \times X$  对于复流形  $Z$  和  $f$  是第二个投影。设置  $Z_{\circ} = Z \times \{x_{\circ}\} (= f^{-1}(x_{\circ}))$  ，并用  $j$  表示嵌入  $Z_{\circ} \hookrightarrow Y$  。设置 $A = \mathrm{SS}(G)$ 并引入：

\[
A_\circ=j^{\#}A\qquad(\mathrm{cf.~VI~}\S2)\enspace.
\]

\statementlabel{引理 8.5.10}其中一个有：

\[
\mathrm{i}j^{\prime}j_{\pi}^{-1}\left(\varLambda+{}^{t}f^{\prime}\left(\boldsymbol{Y}_{\times}\times\boldsymbol{T}^{*}\boldsymbol{X}\right)\right)=\varLambda_{\circ}\ .
\]

引理 8.5.10 的证明。我们在 $ Z \times X $ 上选择局部坐标系 $ (z,x) $ ，并用 $ (z,x;\zeta,\xi) $ 表示 $ T^{*}(Z \times X) $ 上的\textbackslash{}textit\{关联\}坐标。那么 $ (z,;\zeta,0) $  属于 $ {}^{t}jj_{\pi}^{-1}(A + {}^{t}f'(Y \times_{X} T^{*}X)) $  当且仅当 $ A $  中存在序列 $ \{(z_{n}, x_{n}; \zeta_{n}, \xi_{n})\} $  且：

\[
\left(z_{n},x_{n},\zeta_{n}\right)\xrightarrow[n]{}\left(z_{\circ},x_{\circ},\zeta_{\circ}\right)\;.
\]

因此，只要证明对于这样的序列，也有

\[
|\xi_{n}||x_{n}-x_{\circ}|\xrightarrow[n]{}0.
\]

假设(8.5.5)不满足。通过曲线选择引理（在全纯情况下），我们找到一个全纯映射

\[
\theta:\left\{t\in\mathbb{C};0<|t|<1\right\}\to A
\]

这样，设置  $ \theta(t)=(z(t),x(t);\zeta(t),\xi(t)) $  ，当  $ t\to0 $  时：

\[
\left\{\begin{aligned}{}&{{}(z(t);\zeta(t))\to(z_{\circ};\zeta_{\circ})}\\ {}&{{}x(t)-x_{\circ}\sim t^{s}}\\ {}&{{}\xi(t)\sim t^{-r},\quad1\leqslant s\leqslant r.}\\ \end{aligned}\right.
\]

由于 $A$  是各向同性的，我们还有：

\[
\left\langle\zeta(t),\frac{\mathrm{d}z(t)}{\mathrm{d}t}\right\rangle+\xi(t)\frac{\mathrm{d}x(t)}{\mathrm{d}t}=0.
\]

因此 $ \xi(t)\frac{dx(t)}{dt} $ 是有界的，这是一个矛盾。 □

命题 8.5.9 的证明结束。令 $ \varphi: Z_o \to \mathbb{R} $  为实解析函数，并假设 $ \varphi $  为正且正确。设置  $ s_o = \sup_K |\varphi| $  ，并使用命题

8.3.12，选择实数 $ s_{0} < s_{1} < s_{2} $ ，使得：

\[
\mathrm{d}\varphi(z)\notin\varLambda_{o}\qquad\mathrm{f o r}\qquad s_{1}\leqslant\varphi(z)\leqslant s_{2}\enspace.
\]

设置 $ Z_{\circ}^{s} = \{z \in Z_{\circ}; \varphi(z) < s\} $ 。我们声称，对于  $ X $  中  $ x_{\circ} $  的足够小的开邻域  $ U $ ，命题 8.5.8 的假设在  $ V = Z_{0}^{s_{2}} \times U $  上得到满足， $ \varphi $  被视为  $ Z \times X $  上的函数（即：设置  $ \varphi(z, x) = \varphi(z) $  ），并且  $ t_{\circ} = s_{1} $  。事实上，否则我们会在 $ Z \times X $ 中找到一个序列 $ \{(z_{n}, x_{n})\} $ ，这样

\[
\left\{\begin{aligned}{}&{{}(z_{n},x_{n};\mathbf{d}\varphi(z_{n}),0)\in\varLambda+Z\times T^{*}X}\\ {}&{{}x_{n}\xrightarrow[n]{}\mathbf{x}_{\circ},\quad s_{1}\leqslant\varphi(z_{n})\leqslant s_{2}\enspace,}\\ \end{aligned}\right.
\]

我们可以假设  $ z_n \to z_\circ $  。然后根据引理 8.5.10， $ (z_\circ; \mathrm{d}\varphi(z_\circ)) \in \mathcal{A}_\circ $  这是一个矛盾。   $ \square $ 

\statementlabel{备注 8.5.11}如果提出  $ \dim X = 1 $  假设，命题 8.5.9 就不成立。当  $ \dim X > 1 $  时，我们参考 Henry-Merle-Sabbah [1] 对这个问题进行几何研究。

\subsection{附近循环函子与消失循环函子}
在第四章中，我们分别定义了特化和微局部化的函子 $v_{Y}$ 和 $\mu_{Y}$ 。当 $X$  和 $Y$  是复流形而 $Y$  是超曲面时，可以用不同的方式表示它们。

令 $X$  为复流形， $f: X \to \mathbb{C}$  为全纯映射。设置  $Y = f^{-1}(0)$  并用  $i$  表示嵌入  $Y \subsetneq X$  。

令  $\tilde{\mathbb{C}}^*$  为  $\mathbb{C}^* = \mathbb{C} \setminus \{0\}$  的通用覆盖， $p: \tilde{\mathbb{C}}^* \to \mathbb{C}$  为投影； （例如，采用  $\tilde{\mathbb{C}}^* = \mathbb{C}$  和  $p(z) = \exp(2\pi\sqrt{-1}z)$  ）。设置 $\tilde{X}^* = X \times_C \tilde{\mathbb{C}}^*$ 并用 $\tilde{p}$ 表示与 $p$ 关联的投影 $\tilde{X}^* \to X$ （即 $\tilde{p} = \mathrm{id} \times_C p$ ）：

\begin{center}
\includegraphics[width=0.28\textwidth]{流行上的层_Chn-assets/img_in_image_box_405_1316_747_1512.png}
\end{center}

请注意， $ \tilde{p}_{1} $  是一个精确函子。

\statementlabel{定义 8.6.1}对于  $ F \in \text{Ob}(\mathbb{D}^b(X)) $  ，一组：

\[
\psi_{f}(F)=i^{-1}R\tilde{p}_{*}\tilde{p}^{-1}(F)\quad,
\]

并调用  $\psi_{f}$  附近循环函子。

请注意， $ \psi_f(F) $  属于 $ \mathbf{D}^b(Y) $  并且仅依赖于 $ F|_{X\setminus Y} $  。从  $ \tilde{p}^{-1} \simeq \tilde{p}^t $  开始，我们得到了 Poincaré-Verdier 对偶性：

\[
\begin{aligned}{\tilde{p}_{*}\tilde{p}^{-1}F}&{{}\simeq R\mathcal{H}{o m}(\tilde{p}_{!}A_{\tilde{X}^{*}},F)}\\ {}&{{}\simeq R\mathcal{H}{o m}(f^{-1}p_{!}A_{\tilde{C}^{*}},F)~,}\\ \end{aligned}
\]

因此：

\[
\psi_{f}(F)\simeq i^{-1}R\mathcal{H}o m(f^{-1}p_{!}A_{\tilde{c}^{*}},F).
\]

现在  $1\in\mathbb{Z}$  对  $\tilde{\mathbb{C}}^{*}$  的作用引发了  $p_1A_{\tilde{C}^*}$  的自同构  $T$  ，因此引发了  $\psi_f(F)$  的自同构。这种自同构称为  $\psi_f(F)$  的单调，并用  $M$  表示。接下来，考虑复杂的情况：

\[
K:0\longrightarrow p_{!}A_{\widetilde{\mathbf{c}}\bullet}\mathop{\longrightarrow}_{\mathrm{t r}}A_{\mathbf{c}}\longrightarrow0
\]

其中  $ A_C $  代表零度，微分  $ t_r $  是迹态射  $ p_! A_{\tilde{C}} \cdot \simeq p_! p^! A_C \to A_C $  。

\statementlabel{定义 8.6.2}对于  $ F \in \mathrm{Ob}(\mathbb{D}^{b}(X)) $  ，一组：

\[
\phi_{f}(F)=i^{-1}R\mathcal{H}o m(f^{-1}K,F)\enspace,
\]

并调用  $\phi_{f}$  消失循环函子。

 $T$  对  $p_1A_{\tilde{C}}$  的操作。  $A_C$  上的恒等性导致  $K$  的自同构，因此  $\phi_f(F)$  的自同构。这种自同构称为  $\phi_f(F)$  的单调，并且仍用  $M$  表示。

考虑复合体的确切序列（由行表示）：

\begin{center}
\includegraphics[width=0.55\textwidth]{流行上的层_Chn-assets/img_in_image_box_248_1133_908_1326.png}
\end{center}

以及：

\begin{center}
\includegraphics[width=0.57\textwidth]{流行上的层_Chn-assets/img_in_image_box_241_1395_923_1585.png}
\end{center}

请注意，(8.6.5) 顶行的准确性源自以下事实：对于  $ a \in \mathbb{C}^* $  ，

\[
\big(p_{!}A_{\tilde{\mathbf{c}}^{*}}\big)_{a}\simeq A^{(\mathbf{Z})},
\]

（其中 $A^{(\mathbb{Z})}$ 表示 $A$ 中的序列 $\{a_n\}_{n\in\mathbb{Z}}$ 的集合，使得 $a_n=0$ 除了有限多个 $n$ 之外）， $T$ 由 $\{a_n\}_{n}\to\{a_{n+1}\}_n$ 给出，而 $\mathrm{tr}:(p_!A_{\tilde{C}^*})_a\to(A_{\mathbb{C}^*})_a$ 由 $\{a_n\}_{n}\to\sum_na_n$ 给出。

因此我们在  $ \mathbf{D}^{b}(A_{\mathrm{c}}) $  中获得了两个不同的三角形：

\[
\left\{\begin{aligned}& A_{\mathbf{c}}\longrightarrow K\longrightarrow p_{!}A_{\tilde{\mathbf{c}}^{*}}[1]\xrightarrow[+1]{}\\ & p_{!}A_{\tilde{\mathbf{c}}^{*}}[+1]\longrightarrow K\longrightarrow A_{\{0\}}\xrightarrow[+1]{} \end{aligned}\right.,
\]

（注意  $ A_{\{0\}} $  与  $ \mathbf{D}^b(A_C) $  中的复数  $ A_{C^*} \to A_C $  同构），这导致了  $ \mathbf{D}^b(Y) $  中的不同三角形：

\[
\left\{\begin{aligned}&\psi_{f}(F)[-1]\stackrel{\mathrm{c a n}}{\longrightarrow}\phi_{f}(F)\longrightarrow i^{-1}F\stackrel{\longrightarrow}{+1},\\ &i^{!}F\longrightarrow\phi_{f}(F)\stackrel{\mathrm{v a r}}{\longrightarrow}\psi_{f}(F)[-1]\stackrel{\longrightarrow}{+1}.\end{aligned}\right.
\]

请注意，通过这种构造：

\[
\left\{\begin{aligned}{}&{{}\mathbf{c a n}\circ\mathbf{v a r}=1-\mathbf{M}\qquad\mathrm{i n E n d}(\phi_{f}(F))~,}\\ {}&{{}\mathbf{v a r}\circ\mathbf{c a n}=1-\mathbf{M}\qquad\mathrm{i n E n d}(\psi_{f}(F))~.}\\ \end{aligned}\right.
\]

让我们研究当 $Y$  为非奇异时 $\phi_{f}, \psi_{f}, v_{Y}$  和 $\mu_{Y}$  之间的关系。在本例中， $df$  定义了一个函数

\[
\tilde{f}\colon T_{Y}X\to\mathbb{C}\enspace.
\]

我们用  $s$  表示由  $\tilde{f}^{-1}(1)$  给出的  $T_Y X \to Y$  部分，用  $s'$  表示由  $df$  给出的  $T_Y^* X \to Y$  部分。

\statementlabel{命题 8.6.3}假设  $Y$  是非奇异的，并设  $F \in \mathrm{Ob}(\mathbb{D}_{w-R-c}^b(X))$  。在  $p^{-1}(1) \in \tilde{\mathbb{C}^*}$  中选择一个点，并将  $Y$  与  $T_Y X$  的零部分识别。那么我们就有了同构：

\[
\left\{\begin{aligned}{\psi_{f}(F)}&{{}\simeq\psi_{\tilde{f}}(\nu_{Y}(F))}\\ {\phi_{f}(F)}&{{}\simeq\phi_{\tilde{f}}(\nu_{Y}(F))}\\ \end{aligned}\right..
\]

此外，如果  $F \in \mathrm{Ob}(\mathbf{D}_{w-C-c}^{b}(X))$  ，我们有同构：

\[
\left\{\begin{aligned}{\psi_{f}(F)}&{{}\simeq s^{-1}\nu_{Y}(F),}\\ {\phi_{f}(F)}&{{}\simeq s^{{^{\prime}}-1}\mu_{Y}(F).}\\ \end{aligned}\right.
\]

请注意，如果  $F$  是（弱）C 构造的，那么  $\psi_f(F)$  和  $\phi_f(F)$  也是如此，并且

\[
\operatorname{s u p p}(\phi_{f}(F))\subset\{x\in Y;\mathbf{d}f(x)\in\mathrm{S S}(F)\}.
\]

\statementlabel{证明}(a) 首先假设  $ X = \mathbb{C} \times Y $  且 f 是第一个投影。考虑一下图表：

\begin{center}
\includegraphics[width=0.53\textwidth]{流行上的层_Chn-assets/img_in_image_box_247_249_890_571.png}
\end{center}

这里  $j'$  是包含映射， $k'$  是  $k'(z,y) = (0,z,y)$  定义的包含映射， $h$  是映射  $(t,z,y) \mapsto (tz,y)$  。映射 $\pi', \pi, \pi''$  是由 $p: \tilde{\mathbb{C}}^* \to \mathbb{C}$  诱导的。

根据 $ v_{Y} $  的定义，我们有：

\[
\nu_{Y}(F)|_{\dot{T}_{Y}X}\simeq k^{-1}R j_{*}h^{\prime-1}F\quad.
\]

 $ p^{-1}(1) \in \tilde{\mathbb{C}}^* $  的选择使我们能够识别  $ \tilde{\mathbb{C}}^* $  和  $ \mathbb{C} $  ，以及  $ p: \tilde{\mathbb{C}}^* \to \mathbb{C} $  和  $ p(z) = \exp(2\pi\sqrt{-1}z) $  。

我们得到一个笛卡尔图：

\begin{center}
\includegraphics[width=0.43\textwidth]{流行上的层_Chn-assets/img_in_image_box_307_963_830_1151.png}
\end{center}

其中 $ h''(t,z,y)=(z+(1/2\pi\sqrt{-1})\log t,y) $  。

由于  $ h' $  和  $ h'' $  是平滑的，因此我们有  $ F \in \mathrm{Ob}(\mathbb{D}^{b}(X)) $  ：

\[
\begin{aligned}{h^{\prime-1}R\tilde{p}_{*}\tilde{p}^{-1}F}&{{}\simeq\pi_{*}^{\prime}h^{\prime\prime-1}\tilde{p}^{-1}F}\\ {}&{{}\simeq\pi_{*}^{\prime}\pi^{\prime-1}h^{\prime-1}F\enspace.}\\ \end{aligned}
\]

回到图（8.6.13），我们得到：

\[
\begin{aligned}{k^{-1}R j_{*}{h^{\prime}}^{-1}R\tilde{p}_{*}\tilde{p}^{-1}F}&{{}\simeq k^{-1}R j_{*}\pi_{*}^{\prime}\pi^{\prime -1}{h^{\prime}}^{-1}F}\\ {}&{{}\simeq k^{-1}\pi_{*}R j_{*}^{\prime}\pi^{\prime -1}{h^{\prime}}^{-1}F}\\ {}&{{}\simeq k^{-1}\pi_{*}\pi^{-1}R j_{*}h^{\prime -1}F.}\\ \end{aligned}
\]

设置 $ G = Rj_* h'^{-1} F $ 并回想一下我们假设 $ F $ 弱 $ \mathbb{R} $ 可构造。

我们要证明

\[
\left\{\begin{aligned}{}&{{}k^{-1}R\pi_{*}\pi^{-1}G\to R\pi_{*}^{\prime\prime}k^{\prime-1}\pi^{-1}G\simeq R\pi_{*}^{\prime\prime}\pi^{\prime\prime-1}k^{-1}G}\\ {}&{{}{i s~a n~i s o m o r p h i s m}}\\ \end{aligned}\right..
\]

事实上，任何 $ x \in \mathbb{C}^* \times Y $ 都有一个开放简单连通邻域 $ U $ 的基本系统，使得 $ G_x \simeq R\Gamma(U; G) $ （因为 $ G $ 是弱 $ \mathbb{R} $ 可构造的）。自 $ \pi^{-1}(U) \simeq U \times \mathbb{Z} $ 以来，我们有：

\[
\begin{aligned}{{R}\varGamma(U;\pi_{\ast}\pi^{-1}G)}&{{}\simeq{R}\varGamma(\pi^{-1}(U);\pi^{-1}G)}\\ {}&{{}\simeq{R}\varGamma(U;G)^{\mathbb{Z}}.}\\ \end{aligned}
\]

取电感极限，我们得到：

\[
(\pi_{*}\pi^{-1}G)_{x}\simeq(G_{x})^{\mathbb{Z}}.
\]

同理可证明：

\[
(\pi_{*}^{\prime\prime}\pi^{\prime\prime^{-1}}k^{-1}G)_{x}\simeq((k^{-1}G)_{x})^{2}~.
\]

这给出了(8.6.15)。

结合（8.6.15）和（8.6.14），我们得到：

\[
\begin{aligned}{\nu_{Y}(\tilde{p}_{*}\tilde{p}^{-1}F)|_{\dot{T}_{Y}X}}&{{}\simeq\pi_{*}^{{^{\prime \prime}}}\pi^{{^{\prime \prime}}-1}k^{-1}R j_{*}{h^{{^{\prime}}-1}F}}\\ {}&{{}\simeq\pi_{*}^{{^{\prime \prime}}}\pi^{{^{\prime \prime}}-1}(\nu_{Y}(F)|_{\dot{T}_{Y}X}).}\\ \end{aligned}
\]

应用 $ R\dot{\tau}_* $ ，我们在左侧得到 $ \psi_f(F) $ ，在右侧得到 $ \psi_f(v_Y(F)) $ （参见定理4.2.3）。

这证明了(8.6.10)中的第一个同构。为了得到第二个，考虑可区分三角形的交换图就足够了：

\[
\begin{array}{c c c c c}{{\psi_{f}(F)[-1]}}&{{\xrightarrow{\mathrm{c a n}}}}&{{\phi_{f}(F)}}&{{\xrightarrow{\quad}}}&{{i^{-1}F\quad}}\\ {{\bigg\downarrow\bigg\downarrow}}&{{}}&{{\bigg\downarrow}}&{{}}&{{\bigg\downarrow\bigg\downarrow}}\\ {{\psi_{f}(\nu_{Y}(F))[-1]}}&{{\xrightarrow{\mathrm{c a n}}}}&{{\phi_{f}(\nu_{Y}(F))}}&{{\xrightarrow{\quad}}}&{{i^{-1}(\nu_{Y}(F))\quad}}\\ \end{array}.
\]

现在假设 $F$ 是弱 $\mathbb{C}$ 可构造的，我们将证明（8.6.11）。在这种情况下， $\tilde{p}^{-1}\nu_{Y}(F)$  在 $\tau\circ\tilde{p}$  的纤维上是恒定的： $\tilde{\mathbb{C}}^{*}\times_{\mathbb{C}}T_{Y}X\to Y$  。令 $\tilde{s}$  为映射 $\tilde{\mathbb{C}}^{*}\times_{\mathbb{C}}T_{Y}X\to Y$  的一部分，使得 $\tilde{p}\circ\tilde{s}=s$  。应用推论 2.7.7，我们得到：

\[
\begin{aligned}{\psi_{\tilde{f}}(\nu_{Y}(F))}&{{}\simeq R\dot{\tau}_{*}(\tilde{p}_{*}\tilde{p}^{-1}\nu_{Y}(F)|_{\dot{T}_{Y}X})}\\ {}&{{}\simeq\tilde{s}^{-1}\tilde{p}^{-1}\nu_{Y}(F)}\\ {}&{{}\simeq s^{-1}\nu_{Y}(F).}\\ \end{aligned}
\]

将  $ T_Y X $  与  $ \mathbb{C} \times Y $  和  $ T_Y^* X $  与  $ \mathbb{C} \times Y $  识别，我们有  $ s^{-1} \mu_Y(F) \simeq R\pi_* R \mathcal{H} o m(A_{U \times Y}, \mu_Y(F)) $  ，因为  $ \mu_Y(F) $  是  $ \mathbb{C}^\times $  -圆锥曲线。这里 $ U = \{z \in \mathbb{C}; \text{Re} z > 0\} $ 。

由定理3.7.9，我们得到：

\[
\begin{aligned}{s^{\prime-1}\mu_{Y}(F)}&{{}\simeq R\tau_{*}R\mathcal{H}{o m}(A_{U\times Y}^{\vee},\nu_{Y}(F))}\\ {}&{{}\simeq R\tau_{*}R\mathcal{H}{o m}(A_{Z\times Y},\nu_{Y}(F)),}\\ \end{aligned}
\]

其中 $ Z = \{z \in \mathbb{C}; \text{Im } z = 0, \text{Re } z \geq 0\} $  。

然后我们通过使用包含  $\mathbb{C} \hookrightarrow Z \hookrightarrow \tilde{\mathbb{C}}^{*}$  和  $\mathbb{C}$  来得到一个交换图：

\begin{center}
\includegraphics[width=0.19\textwidth]{流行上的层_Chn-assets/img_in_image_box_478_478_716_651.png}
\end{center}

由于复数  $ A_{C\backslash Z} \to A_C $  与  $ A_Z $  同构，因此足以证明：

\[
(8.6.16)\quad R\tau_{*}R\mathcal{H}o m(p_{!}A_{\widetilde{C}^{*}}\boxtimes A_{Y},\nu_{Y}(F))\xrightarrow{\sim}R\tau_{*}R\mathcal{H}o m(A_{C\setminus Z}\boxtimes A_{Y},\nu_{Y}(F))~.
\]

由于  $ \nu_Y(F) $  在  $ (\mathbb{C} \setminus Z) \times Y \to Y $  的纤维上是局部常数，所以我们有  $ s^{-1} \nu_Y(F) \simeq R \tau_* R \mathcal{H} o m(A_{\mathbb{C} \setminus Z} \boxtimes A_Y, \nu_Y(F)) $  。这显示了 (8.6.16) 并且我们得到  $ s'^{-1} \mu_Y(F) \simeq \phi_f(F) $  。

一般情况下，通过图表 $f: X \subsetneqq_{\mathfrak{g}} \mathbb{C} \times X \to \mathbb{C}$ 将 $X$ 嵌入到 $\mathbb{C} \times X$ 中，并注意：

\[
\psi_{f}(F)\simeq\psi_{t}(g_{*}F)\mathrm{~,~}
\]

\[
\phi_{f}(F)\simeq\phi_{t}(g_{*}F)\mathrm{~,~}
\]

（参见练习 VIII.15）。 □

我们可以从消失循环函子中恢复 SS(F)。

\statementlabel{命题 8.6.4}让  $ F \in \text{Ob}(\mathbb{D}_{w-C-c}^b(X)) $  并让  $ p \in T^*X $  。那么下面两个条件是等价的。

（一） $ p \notin SS(F) $ 。

(ii) 存在  $p$  的开邻域  $U$ ，使得对于在  $x$  与  $f(x) = 0$  和  $\mathrm{d}f(x) \in U$  的邻域中定义的任何  $x \in X$  和任何全纯函数  $f$  ，都有  $\phi_{f}(F)_{x} = 0$  。

\statementlabel{证明}(i) 由命题 8.6.3 和推论 5.4.10 暗示 (ii)。相反，假设 (ii) 并让我们证明  $ U \cap \text{SS}(F) = \emptyset $  。根据定理 8.5.5， $ \varLambda = \text{SS}(F) $  是  $ T^*X $  的复解析拉格朗日子集。在 $ \varLambda \cap U $  的通用点 $ p' $  上，对于某些 $ L \in \text{Ob}(\mathbf{D}^b(\mathfrak{Mod}(A)) $  和 $ X $  的复子流形 $ Y $ ， $ F $  在 $ \mathbf{D}^b(X; p') $  到 $ L_Y $  中是同构的（命题6.6.1）。我们将展示 $ L = 0 $ 。如果 $ p' \in T_X^*X $ ，则 $ Y = X $ ，我们有 $ L = 0 $ （采用 $ f = 0 $ ）。否则选择局部坐标

 $X$  上的系统  $(z_1,\ldots,z_n)$  使得  $p' = (0; \mathrm{d}z_1)$  和  $Y = \{z_1 = \cdots = z_l = 0\}$  。让 $f(z) = z_1 + \sum_{j=l+1}^n z_j^2$  。然后是  $\phi_f(F)_0 \simeq \mu_{\{f=0\}}(F)_{p'} \simeq \mu_{\{f=0\}}(L_Y)_{p'} \simeq \phi_f(L_Y)_0 \simeq \phi_{f|Y}(L_Y)_0 \simeq L[l-n]$  ，（参见练习 VIII.14 和 VIII.15）。这样就完成了证明。   $\square$ 

\subsection{习题}
\statementlabel{练习 VIII.1}令 $ \mathbf{S} = (\mathcal{S}, \mathcal{D}) $  为单纯复形。通过设置 $ \mathrm{Ob}(\Delta) = \{\sigma; \sigma \in \mathcal{D}\} $  和 $ \mathrm{Hom}(\sigma, \tau) = \mathcal{Q} $ （如果 $ \sigma \not\subset \tau $ ）、 $ \mathrm{Hom}(\sigma, \tau) = \{\mathrm{pt}\} $ （如果 $ \sigma \subset \tau $ ）来定义范畴 $ \Delta $ 。令 $ \Delta(A) $  表示从 $ \Delta $  到 $ \mathfrak{Mod}(A) $  的协变函子范畴，并令 $ \delta $  为 $ F \mapsto \{\sigma \mapsto F(U(\sigma))\} $  给出的从 $ w\text{-}\text{Cons}(\mathbf{S}) $  到 $ \Delta(A) $  的函子。证明  $ \delta $  是范畴的等价。

\statementlabel{练习 VIII.2}让  $ X = \mathbb{R}^3 $  赋予坐标  $ (t, x, y) $  并让  $ Z = \{y^2 - t^2x^2 - x^3 = 0\} $  、  $ Z_2 = \{x = y = 0\} $  、  $ Z_1 = Z \backslash Z_2 $  、  $ Z_0 = X \backslash Z $  。证明 $ X = Z_0 \sqcup Z_1 \sqcup Z_2 $  是亚解析分层，但不是 $ \mu $  分层。

\statementlabel{练习 VIII.3}令  $f: Y \to X$  为流形  $F \in \mathrm{Ob}(\mathbf{D}^{b}(X))$  、  $G \in \mathrm{Ob}(\mathbf{D}^{b}(Y))$  的态射。

(i) 证明同构：

\[
f^{!}({\bf D}_{X}F)\simeq{\bf D}_{Y}(f^{-1}F)\enspace,
\]

\[
{R f}_{*}(\mathbf{D}_{Y}G)\simeq\mathbf{D}_{X}({R f}_{!}G)\enspace.
\]

(ii) 假设 $F$  是 $\mathbb{R}$  可构造的。证明同构：

\[
f^{-1}(\mathbf{D}_{X}F)\simeq\mathbf{D}_{Y}(f^{!}F)\enspace.
\]

(iii) 假设 $G$  和 $Rf_1(D_Y G)$  是 $\mathbb{R}$  可构造的。证明同构：

\[
R\vec{f}_{!}({\bf D}_{Y}G)\simeq{\bf D}_{X}(R f_{*}G)\mathrm{~,~}
\]

并证明  $ Rf_*G $  是  $ \mathbb{R} $  可构造的。

\statementlabel{练习 VIII.4}在命题 5.4.17 的情况下，证明如果  $G$  是  $\mathbb{R}$  可构造的，则  $R f_{*}G$  （或  $R f_{!}G$  ）是  $\mathbb{R}$  可构造的。

\statementlabel{练习 八.5}令  $ X = \bigsqcup_{\alpha} X_{\alpha} $  为  $ \mu $  分层并设置  $ \varLambda = \bigsqcup_{\alpha} T_{X_{\alpha}}^{*}X $  。让 $ F \in \mathrm{Ob}(\mathbb{D}^{b}(X)) $  。证明  $ \mathrm{SS}(F) \subset \varLambda $  当且仅当  $ \mathrm{SS}(H^{j}(F)) \subset \varLambda $  对于所有  $ j $  。

\statementlabel{练习 八.6}让  $\Omega$  成为  $T^*X$  的子集。定义了  $\mathbf{D}^b(X; \Omega$  的满子范畴  $\mathbf{D}_{w-R-c}^b(X; \Omega)$  （分别为  $\mathbf{D}_{R-c}^b(X; \Omega)$  ），由满足以下条件的对象  $F$  组成：在任何点  $p \in \Omega$  都存在  $F' \in \mathrm{Ob}(\mathbf{D}_{w-R-c}^b(X))$  （分别为  $F' \in \mathrm{Ob}(\mathbf{D}_{R-c}^b(X))$  ），使得  $F$  和  $F'$  在  $\mathbf{D}^b(X; p)$  中是同构的。

(i) 证明  $ F \in \text{Ob}(\mathbf{D}^b(X; \Omega)) $  属于  $ \mathbf{D}_{w-R-c}^b(X; \Omega) $  当且仅当存在  $ \Omega $  的开邻域  $ U $  使得  $ \text{SS}(F) \cap U $  包含在  $ U $  的闭亚解析各向同性子集中。

(ii) 在命题 6.3.3 的情况下，证明如果  $F$  是  $\mathbb{R}$  可构造的，则  $Rq_1, F$  和  $Rq_1, F$  属于  $\mathbf{D}_{\mathbf{R}-c}^b(X; \Omega)$  。

\statementlabel{练习 八.7}在定理 7.2.1 的情况下，假设  $K$  是  $\mathbb{R}$  可构造的。证明  $\Phi_K$  引发等价  $\mathbf{D}_{\mathbb{R}-c}^b(Y; \Omega_Y) \simeq \mathbf{D}_{\mathbb{R}-c}^b(X; \Omega_X)$  。 （提示：使用练习 VIII.6。）

\statementlabel{练习 八.8}令 $X$  为复流形， $A$  为 $T^*X$  的闭圆锥曲线 $\mathbb{C}$  解析子集。假设 $A_{\text{reg}}$  是对合的。证明 $A$  在 $T^*X^R$  中是对合的（在定义6.5.1 的意义上）。 （提示：使用 Kashiwara-Monteiro-Fernandes [1] 的结果。）

\statementlabel{练习 八.9}令 $ \mathbf{S} = (S, \mathcal{A}) $  为有限维的单纯复形（即满足（8.1.16）），并让 $ X = |S| $  表示相关的拓扑空间。设定：

\[
X_{k}=\left\{x\in X;x\in|\sigma|\mathrm{f o r~s o m e}\sigma\in\varDelta\mathrm{w i t h}\#{\sigma}\leqslant-k+1\right\}
\]

(i) 证明 $ H_{X_k \setminus X_{k+1}}^j(\omega_X) = 0 $  对于 $ j \neq k $  。

(ii) 使用练习 II.21，证明  $\omega_X$  与  $\mathbf{D}^b(X)$  中的复数  $\{H_{X_k \setminus X_{k+1}}^k (\omega_X)\}$  同构，并描述该复数。

\statementlabel{练习 八.10}令 $X$  为复流形， $S$  为闭合 $\mathbb{C}$  解析子集， $U = X \setminus S$  和 $j$  为开放嵌入 $U \subsetneq X$  。让  $F \in \mathrm{Ob}(\mathbf{D}_{\mathbb{C}-c}^{b}(U))$  与  $\mathrm{SS}(F) \subset T_U^* U$  。证明  $R j_! F$  和  $R j_* F$  属于  $\mathbf{D}_{\mathbb{C}-c}^{b}(X)$  。

\statementlabel{练习 八.11}令  $X = \bigsqcup_{\alpha} X_{\alpha}$  为实解析流形  $X$  的亚解析分层。证明如果 $A = \bigsqcup_{\alpha} T_{X_{\alpha}}^{*}$  是闭集，则 $A$  是对合的（在定义6.5.1 的意义上）。

\statementlabel{练习 八.12}令  $N$  和  $M$  为  $\mathbb{R}^n$  的亚解析子流形，以及  $N \subset \overline{M} \setminus M$  。考虑惠特尼提出的以下条件 (a) 和 (b)。

(a) 对于  $M$  中收敛到  $x \in N$  的任何序列  $\{x_n\}$  ，如果  $\{T_{x_n}M\}$  收敛到平面  $\tau$  ，则  $\tau \supset T_x N$  。

(b) 对于  $M$  中的任何序列  $\{x_n\}$  和  $N$  中的任何序列  $\{y_n\}$  ，两者都收敛到  $x \in N$  ，如果线序列  $\{\mathbb{R}(x_n, y_n)\}$  收敛到  $l$  且序列  $\{T_{x_n}M\}$  收敛到  $\tau$  ，则  $\tau \supset l$  。

证明 (b) 蕴含 (a) 并且  $ \mu $  -条件蕴含 (b)。 （提示：使用曲线选择引理。）

\statementlabel{练习 八.13}令  $f$  为复流形  $X$  上的全纯函数，并令  $F \in \mathrm{Ob}(\mathbf{D}_{w-C-c}^{b}(X))$  。证明存在规范同构  $R\Gamma_{\{\mathrm{Re}f \geq 0\}}(F)|_{f^{-1}(0)} \stackrel{\sim}{\to} \phi_f(F)$  。

\statementlabel{练习 八.14}让 $ X = \mathbb{C}^n $ ， $ f(x) = \sum_j x_j^2 $ ，并让 $ M \in \mathrm{Ob}(\mathbb{D}^b((\mathfrak{Mod}(A)))) $ 。证明同构 $ \phi_f(M_X) \simeq M_{\{0\}}[-n] $ 。

\statementlabel{练习 八.15}令  $f: Y \to X$  为复流形的态射， $t \cdot X \to \mathbb{C}$  为  $X$  上的全纯函数。令 $G \in \mathbf{D}_{w-C-c}^{b}(Y)$  ，并假设 $f$  在supp(G) 上是正确的。用 $f_o$ 表示 $f$ 对 $\{t \circ f = 0\}$ 的限制，证明：

\[
\psi_{t}(R f_{*}G)\simeq R f_{\circ*}\psi_{t\circ f}(G)\enspace,
\]

\[
\phi_{t}(R f_{*}G)\simeq R f_{\circ\ast}\phi_{t\circ f}(G)\enspace.
\]

\subsection{注记}
亚解析集理论由 Gabrielov [1] 和 Hironaka [1,2] 发起，起源于 Lojasiewicz [1,2] 关于半解析集的工作。主要结果归功于 Hironaka（上文引用），但还有许多其他重要贡献，特别是 Hardt [1,2]、Tamm [1] 和 Teissier [1]。这个理论在技术上长期以来一直存在困难，现已被大大简化（例如参见 Denkowska-Lojasiewicz-Stasica [1]），现在 Bierstone-Milman [1] 发表了一篇简短的论文，其中包含主要定理和完整的证明。

分层空间的概念第一次出现似乎是在1955年的Cartan-Chevalley[1]中，而Whitney[1,2]首先引入了层的正则性条件（著名的“惠特尼条件（a）和（b）”），并证明对于给定的分层，存在满足这些条件的更精细的分层，（结果类似于定理8.3.20）。真实分层有多种规律性条件（参见 Trotman [1] 的综述），特别是 Verdier [3] 于 1976 年引入的 w 分层概念，它是 Kuo [1] 提出的另一种条件的变体。 Verdier 在他的论文（loc. cit.）中证明，首先，w 条件比 Whitney 条件更强，其次，总是存在 w 分层。最后一个结果实际上等价于定理 8.3.20，因为 Trotman [2] 最近证明了 w 分层和  $ \mu $  分层的概念是等价的。我们还要提一下 Delort [1] 的一篇论文，他给出了 + 操作的概括。

在复杂的相对情况下，Thom 的  $A_{f}$  条件似乎是研究直像的有效工具（参见 Henry-Merle-Sabbah [1]），而 Hironaka [3] 证明，如果目标空间具有一维，则“局部”满足该条件：这本质上等价于命题 8.5.9。

分层的微观局部研究（即通过使用相关的拉格朗日子集来处理分层的想法，现在称为“共正规几何”）始于柏原市 [3, 5]，然后在柏原-夏皮拉 [1, 3] 得到发展。

Grothendieck（在 [SGA 4]，Artin 的 exposé IX 中）首先引入了可构造层的概念并研究了它们的函数特性（在 étale 上同调的背景下）。此理论在之后重新引起人们的兴趣

Kashiwara [3] 与完整 D 模建立了联系（参见下文第 XI 章第 3 节），并由 Goresky-MacPherson [1] 发现了相交上同调，并由 Gabber 和 Beilinson-Bernstein-Deligne [1] 发现了反常层（参见 Brylinski [1] 进行了评论）。

邻近循环和消失循环函子的理论（参见§6）源自 Grothendieck（参见 [SGA 7] exposé I）和 Deligne [2]。它起源于皮卡德和莱夫谢茨对消失循环和单一性的研究。

如命题 8.6.4 所示，复流形上 C 构造层的微支撑可以使用消失循环函子来获得。正如 Brylinski（前述）所指出的，这一事实已经隐含在 Kashiwara [5] 中。

§1 的结果以及定理 8.4.5 摘自 Kashiwara [6]。 §3、4、5 的大部分结果首先是通过 Kashiwara-Schapira [3] 中的类似方法获得的，但当然，可构造层上的操作早已为人所知（参见例如 Verdier [4]）。然而，定义8.3.19和命题8.4.14以及命题8.4.1的证明都是新的。 （请注意，在此证明中我们没有使用 Thom-Mather 的同位素定理。）

\section{特征圈}
\subsection{概要}
在实解析流形 X 上，我们首先借助对偶复形  $ \omega_{x} $  引入亚解析链和亚解析循环的概念。然后我们定义两个循环的交集。

如果 $F$ 是 $\mathbf{D}^{b}(X)$ 的 $\mathbb{R}$-可构造对象，我们构造 $F$ 的特征圈 $\mathrm{CC}(F)$，它是 $H_{\mathrm{SS}(F)}^{0}(T^{*}X;\pi^{-1}\omega_X)$ 中 $\mathrm{id}_F\in\mathrm{Hom}(F,F)$ 的自然像。例如，若 $M$ 是 $X$ 的闭子流形，则 $\mathrm{CC}(V_M)=m[T_M^*X]$。

接下来我们研究“Lefschetz不动点公式”的一个版本。如果  $f$  是  $X$  的自同态，并且如果给定态射  $\varphi \in \mathrm{Hom}(f^{-1}F, F)$  ，则定义一个特征类  $C(\varphi) \in H^0(X; \omega_X)$  ，其度数计算  $\Gamma(X; \varphi) \in \mathrm{Hom}(R\Gamma(X; F), R\Gamma(X, F))$  的迹  $\mathrm{tr}(\varphi)$  （假设  $F$  具有紧支持）。当 $f$  只有有限多个不动点并且横向于恒等式时，我们展示如何从局部公式推导 $\mathrm{tr}(\varphi)$ 。

最后我们研究  $ \mathbf{D}_{\mathbf{R}-c}^{b}(X) $  的 Grothendieck 群。我们证明该群同构于 $ T^*X $ 上的拉格朗日循环群（同构由 $ F \mapsto \operatorname{CC}(F) $ 定义）以及 $ \mathbb{R} $ 上的可构造函数群 $ X $ （同构由 $ F \mapsto \chi(F)(x) = \chi(F_x) $ 定义）。这产生了  $ X $  上可构造函数代数的新微积分。

我们保留惯例 8.0。此外，在第 9.1 节中，从第 9.4 节到本章结束，基环 A 将是特征为零的（交换）域，我们将用 k 表示它。

\statementlabel{备注 9.0}图的交换性在第 1 节中被仔细检查，并在本章的其余部分中作为此类证明的原型，我们将在其中

有时会冒昧地向读者留下类似的验证。此外，如果要遵循所有计算中的符号，尤其是本章中的符号，会使本书变得太重。换句话说，我们并不总是区分交换图和反交换图。这并不影响对内容的理解。通常对于应用来说，通过查看简单的示例可以轻松获得公式中的正确符号。

\subsection{指数公式}
在本节中，我们假设基环  $ A $  是特征零的交换域，并用  $ k $  表示。如果  $ V \in \text{Ob}(\mathbf{D}^b(\mathfrak{Mod}^f(k))) $  ，其索引  $ \gamma(V) $  定义为：

\[
\chi(V)=\sum_{j}(-1)^{j}\dim H^{j}(V)\ .
\]

现在让  $X$  成为（实解析）流形，并让  $F \in \mathbf{D}_{R-c}^{b}(X)$  。让 $x \in X$  。

\statementlabel{定义 9.1.1}一套：

\[
\chi(F)(x)=\chi(F_{x}),
\]

\[
\chi_{c}(F)(x)=\chi(R\Gamma_{\{x\}}(X;F))=\chi(\mathbf{D}(F))(x)\enspace.
\]

如果  $ R\Gamma(X;F) $  （分别为  $ R\Gamma_c(X;F) $  ）属于  $ \mathbf{D}^b(\mathfrak{Mod}^f(k)) $  ，则一组：

\[
\chi(X;F)=\chi(R\Gamma(X;F))
\]

\[
({\tt r e s p.}\;\chi_{c}(X;F)=\chi(R\varGamma_{c}(X;F)))\enspace.
\]

整数 $ \chi(X; F) $ 称为F的欧拉-庞加莱指数，函数 $ \chi(F)(x) $ 称为F的局部欧拉-庞加莱指数。

在本节中，我们将介绍第一个计算这些指数的工具。

恒等式  $ \mathrm{Hom}(F, F) = \mathrm{Hom}(k_X, R_{\mathcal{Hom}}(F, F)) $  定义了态射：

\[
k_{x}\to R\mathcal{H}o m(F,F)~.
\]

对偶函子  $D_x$  被定义为  $R\mathcal{H}om(\cdot, \omega_x)$  ，通过收缩我们得到态射：

\[
\mathrm{t r}_{X}\colon F\otimes\mathrm{D}_{X}F\to\omega_{X}\enspace.
\]

我们称 $\operatorname{tr}_X$ 为“迹态射”。让 $\delta_X: X \to X \times X$  表示对角线嵌入。 （如果没有混淆的风险，我们写  $\delta$  而不是  $\delta_X$  ，类似地， D 而不是  $D_X$  等）然后有一个独特的态射：

\[
\delta^{!}\to\delta^{-1}
\]

使得下图可以通勤：

\[
\begin{array}{c}\delta^{-1}\delta_{!}\delta^{!}\quad\longrightarrow\quad\delta^{-1}\quad\stackrel{\sim}{\longrightarrow}\quad\delta^{!}\delta_{!}\delta^{-1}\\\downarrow\quad\downarrow\quad\downarrow\quad\downarrow\\\delta^{-1}\delta_{*}\delta^{!}\quad\stackrel{\sim}{\longrightarrow}\quad\delta^{!}\quad\longrightarrow\quad\delta^{!}\delta_{*}\delta^{-1}.\end{array}
\]

因此我们得到态射链：

\[
k_{X}\longrightarrow R\mathcal{H}o m(F,F)\simeq\delta^{!}(F\boxtimes\mathbf{D}F)\longrightarrow\delta^{-1}(F\boxtimes\mathbf{D}F)\simeq F\otimes\mathbf{D}F\xrightarrow[\mathtt{t r}_{X}]{\quad}\omega_{X}\quad.
\]

\statementlabel{定义 9.1.2}$1 \in \Gamma(X; k_X)$  在  $H_{\mathrm{supp}(F)}^0(X; \omega_X)$  中通过上述态射链的图像称为  $F$  的特征类，记为  $\mathrm{C}(F)$  。对于闭子集 $S \supset \mathrm{supp}(F)$  ，我们还用相同的符号 $\mathrm{C}(F)$  来表示 $F$  的特征类的像射 $H_{\mathrm{supp}(F)}^0(X; \omega_X) \to H_S^0(X; \omega_X)$  。

当 $X$  是一个点时，则 $F \in \mathrm{Ob}(\mathbb{D}^b(\mathfrak{Mod}^f(k)))$  和 $\mathrm{C}(F) = \chi(F)$ （参见练习I.32）。

我们将研究特征类的直像。

令 $f: Y \to X$  为流形的态射，令 $G \in \mathrm{Ob}(\mathbf{D}_{\mathbf{R}-c}^{b}(Y))$  并假设 $f$  在supp(G) 上是正确的。考虑下图：

\[
\begin{array}{ccccc}R f_{*}^{k_{Y}}\to&R f_{*}^{}R\mathcal{H}o m(G,G)&\simeq&R f_{!}^{}\delta_{Y}^{!}(G\boxtimes\mathbb{D}_{Y}^{}G)\quad\to&R f_{!}^{}(G\otimes\mathbb{D}_{Y}^{}G)\quad\to R f_{!}^{}\omega_{Y}\\\downarrow\quad&1&&\downarrow\quad2\quad\downarrow\quad3\quad\downarrow\quad4\quad\downarrow\\k_{X}&\to R\mathcal{H}o m(R f_{!}^{}G,R f_{!}^{}G)&\simeq\delta_{X}^{!}(R f_{!}^{}G\boxtimes\mathbb{D}_{X}^{}R f_{!}^{}G)\to(R f_{!}^{}G)\otimes(\mathbb{D}_{X}^{}R f_{!}^{}G)\to&\omega_{X}\end{array}
\]

\statementlabel{命题 9.1.3}图(9.1.6)是可交换的。

\statementlabel{证明}为了简单起见，我们将省略为派生函子写“R”。

(i) (9.1.6) 中的平方 1 显然是可交换的。

(ii) 考虑映射的交换图：

\textbackslash{}[Y\textbackslash{}times Y\textbackslash{}xrightarrow\{\textbackslash{}quad f\_\{1\}\textbackslash{}quad\}X\textbackslash{}times Y\textbackslash{}xrightarrow\{\textbackslash{}quad f\_\{2\}\textbackslash{}quad\}X\textbackslash{}times X\textbackslash{}quad\textbackslash{}scriptstyle\{\textbackslash{}bigg\textbackslash{}downarrow\textbackslash{}quad\}\textbackslash{}delta\_\{Y\}\textbackslash{}quad\textbackslash{}scriptstyle\{\textbackslash{}bigg\textbackslash{}downarrow\textbackslash{}quad\}\textbackslash{}delta\_\{Y\}\textbackslash{}quad\textbackslash{}scriptstyle\{\textbackslash{}bigg\textbackslash{}downarrow\textbackslash{}quad\}\textbackslash{}delta\_\{X\}\textbackslash{}quad\textbackslash{}scriptstyle\{\textbackslash{}bigg\textbackslash{}downarrow\textbackslash{}quad\}\textbackslash{}delta\_\{X\}\textbackslash{}quad\textbackslash{}scriptstyle\{\textbackslash{}bigg\textbackslash{}downarrow\textbackslash{}quad\}\textbackslash{}delta\_\{X\}\textbackslash{}

其中平方是笛卡尔平方的。

然后(9.1.6)中的方格2和方格3分解如下图(9.1.8)。

\textbackslash{}[\textbackslash{}begin\{array\}\{c\}f\_\{*\}\textbackslash{}mathcal\{H\}o m(G,G)\textbackslash{}quad\textbackslash{}widetilde\{\textbackslash{}sim\}\textbackslash{}quad f\_\{*\}\textbackslash{}delta\_\{Y\}\textasciicircum{}\{1\}\textbackslash{}mathcal\{H\}o m(A\_\{Y\}\textbackslash{}boxtimes G,G\textbackslash{}boxtimes\textbackslash{}omega\_\{Y\})\textbackslash{}quad\textbackslash{}widetilde\{\textbackslash{}longleftarrow\}\textbackslash{}quad f\_\{*\}\textbackslash{}delta\_\{Y\}\textasciicircum{}\{1\}(G\textbackslash{}boxtimes\textbackslash{}mathbb\{D\}\_\{Y\}G)\textbackslash{}quad\textbackslash{}longrightarrow\textbackslash{}quad f\_\{*\}\textbackslash{}delta\_\{Y\}\textasciicircum{}\{-1\}(G\textbackslash{}boxtimes\textbackslash{}mathbb\{D\}\_\{Y\}G)\textbackslash{}quad\textbackslash{}widetilde\{\textbackslash{}longleftarrow\}\textbackslash{}quad f\_\{!\}(G\textbackslash{}boxtimes\textbackslash{}mathbb\{D\}\_\{Y\}G)\textbackslash{}\textbackslash{}\textbackslash{}向下箭头\textbackslash{}

(9.1.8) 中所有图的交换性都很容易检查，除了方块 (a) 和 (b) 之外，它们分别来自下面的引理 9.1.4 和 9.1.5。

\statementlabel{引理 9.1.4}令  $ Z \xrightarrow{g} Y \xrightarrow{t} X $  为具有有限  $ c $  软维数的局部紧空间态射链，并设置  $ h = f \circ g $  。假设  $ g $  和  $ h $  是封闭嵌入。然后如果  $ G \in \text{Ob}(\mathbf{D}^b(Y)) $  ，则下图可通勤。

\[
\begin{array}{ccc}g^1G&\xrightarrow{\quad}&g^{-1}G\\\downarrow&&\downarrow\\g^!f^!f_!G&&g^{-1}f^{-1}f_*G\\\downarrow&&\downarrow\\h^!f_!G&\xrightarrow{\quad}&h^{-1}f_!G\end{array}
\]

引理 9.1.4 的证明。由于  $ h_1 = h_* : \mathbb{D}^b(Z) \to \mathbb{D}^b(X) $  是忠实的，因此通过应用函子  $ h_1(\simeq h_*) $  足以证明从 (9.1.9) 推导出的图的交换性。

然后将生成的图嵌入到图（9.1.10）中：

\begin{center}
\includegraphics[width=0.86\textwidth]{流行上的层_Chn-assets/img_in_image_box_57_930_1082_1282.png}
\end{center}

(9.1.10) 的交换律源自组合  $ f_1G \to f_1f^1f_1G \to f_1G $  和  $ f_sG \to f_s f^{-1}f_sG \to f_sG $  是恒等式这一事实。   $ \square $ 

\statementlabel{引理 9.1.5}考虑具有有限 c-软维数的局部紧空间的笛卡尔平方 (9.1.11)，其中映射 g 是闭合嵌入：

\[
\begin{array}{c}Y^{\prime}\xrightarrow{\quad g^{\prime}\quad}Y\\f^{\prime}\bigg|\quad\Box\quad\bigg|\neq f\\X^{\prime}\xrightarrow{\quad g\quad}X\quad.\end{array}
\]

然后对于  $ G \in \text{Ob}(\mathbb{D}^b(Y)) $  ，下面的图 (9.1.12) 是可交换的。

\[
\begin{array}{c}f_{*}g^{\prime!}G\quad\longrightarrow\quad f_{*}g^{\prime-1}G\\\downarrow\quad\uparrow\quad\uparrow\\g^{\prime}f_{*}G\quad\longrightarrow\quad g^{-1}f_{*}G\quad.\end{array}
\]

引理 9.1.5 的证明。证明应用忠实函子  $ g_! \simeq g_* $  得到的图的交换性就足够了。将所得图嵌入到以下交换图中，从而完成证明。

\begin{center}
\includegraphics[width=0.66\textwidth]{流行上的层_Chn-assets/img_in_image_box_186_563_973_1022.png}
\end{center}

命题 9.1.3 的证明结束。剩下的就是证明(9.1.6)中平方4的交换律。这是从下一个引理得出的。   $ \Box $ 

\statementlabel{引理 9.1.6}令  $f: Y \to X$  为具有有限  $c$  软维数的局部紧空间的态射，并令  $F \in \mathrm{Ob}(\mathbf{D}^{b}(X))$  、  $G \in \mathrm{Ob}(\mathbf{D}^{b}(Y))$  。那么如下图通勤。

\begin{center}
\includegraphics[width=0.78\textwidth]{流行上的层_Chn-assets/img_in_image_box_109_1273_1043_1446.png}
\end{center}

\statementlabel{证明}很简单。

根据命题 9.1.3 我们得到：

\statementlabel{定理 9.1.7}令  $f: Y \to X$  为流形的态射，令  $G \in \text{Ob}(\mathbb{D}_{\mathbb{R}-c}^b(Y))$  ，并假设  $f$  在  $\text{supp}(G)$  上是正确的。那么 $\text{C}(R f_* G)$  就是 $\text{morphism} H_{\text{supp}(G)}^0(Y; \omega_Y) \to H_f(\text{supp}(G))(X; \omega_X)$  的 $\text{C}(G)$  的图像。

特别是如果  $F \in \mathbf{D}_{\mathbf{R}-c}^{b}(X)$  具有紧凑支持，我们发现（通过考虑态射  $a_X: X \to \{\mathrm{pt}\}$  ）：

\[
\chi(X;F)=\int_{X}\mathrm{C}(F),
\]

其中 $\int_X$ 是 III §3 中定义的映射 $H_c^0(X,\omega_X)\to k$。在后续部分中，我们将从两个方向概括这些结构：一方面，我们将给出  $\mathrm{C}(F)$  的微本地版本。 另一方面，我们将把指标公式推广为迹公式。为此，我们需要一些准备。

\subsection{次解析链与次解析圈}
奇异同源性是通过奇异链定义的。在本书中，我们仅限于亚分析案例。设 X 为流形。 （回想一下，流形的所有流形和态射都是实解析的。）对于每个整数  $p$  ，令  $CS_p'(X)$  为由符号 [S] 生成的  $A$  模，其中  $S$  的范围贯穿  $X$  的子解析  $p$  维定向子流形族，其关系为 (9.2.1)， (9.2.2)、(9.2.3)如下：

\[
[S_{1}\sqcup S_{2}]=[S_{1}]+[S_{2}]~,\qquad{(S_{1}{\mathrm{~a n d~}}S_{2}{\mathrm{~a r e~d i s j o i n t}})~.}
\]

\[
\begin{aligned}[S]&=[S^{\prime}],\quad \text{if }S^{\prime}\text{ is an open dense subanalytic subset of }S,\\ &\text{endowed with the induced orientation.}\end{aligned}
\]

\[
\begin{aligned}[S^{a}]&=-[S],\quad \text{where }S^{a}\text{ is }S\text{ endowed with}\\ &\text{the opposite orientation.}\end{aligned}
\]

我们用 $ \mathcal{CS}_p^X $ （或 $ \mathcal{CS}_p $ ，如果没有混淆的风险）表示与预层 $ U \mapsto CS'_p(U) $ 相关联的 $ X $ 上的层。如果  $ F $  是  $ X $  上的一个层，我们还设置：

\[
\mathcal{C S}_{p}(F)=\mathcal{C S}_{p}\otimes F~.
\]

\statementlabel{定义 9.2.1}束  $ \mathcal{C}\mathcal{S}_p(F) $  称为亚解析 p 链束，其值在  $ F $  中。

为了描述  $ \mathcal{CS}_p(F) $  ，我们需要对偶复数  $ \omega_X $  的一些补充结果。

对于 X 的子解析子集 S，用  $j_{S}$  表示嵌入  $S \hookrightarrow X$  。 如果 $S$ 的维数 $\leq p$，则 $S_{\text{reg},p}$ 表示 $S_{\text{reg}}$ 的 $p$ 维连通分支的并。如果 S 是局部闭的，我们设置：

\[
\partial{\cal S}=\bar{\cal S}\backslash{\cal S}~.
\]

这是一个闭集。

\statementlabel{命题 9.2.2}令 S be a closed subanalytic subset of dimension   $ \leq p $  .

(i) 对于 S 上的任何层 F，有：

\[
H_{c}^{j}(S;\,F)=H^{j}(S,F)=0\qquad{f o r\quad}\quad j>p.
\]

ii)  $ H^{j}(\omega_{S}) = 0 $  对于 j < -p。

（三） $ H^{-p}(S;\omega_S) \simeq H^0(S;H^{-p}(\omega_S)) \simeq \mathrm{Hom}(H_c^p(S;A_S), A) $ 。

\statementlabel{证明}(i) 选择过滤  $S = S_p \supset \cdots \supset S_0$  ，其中  $S_k$  是  $S$  的闭合子分析子集， $S_k \setminus S_{k-1}$  是平滑的，维度为  $k$  。然后：

\[
H^{j}(S;F_{S_{k}\setminus S_{k-1}})=\varinjlim_{U}H^{j}(S_{k}\setminus S_{k-1};F_{U})
\]

其中  $U$  的范围贯穿  $S_k\backslash S_{k-1}$  的开放子集族，使得  $\overline{U}\cap S_{k-1}=\varnothing$  。 By 命题 3.2。2 we find   $H^j(S;F_{S_k\backslash S_{k-1}})=0$   for   $j>k$  .然后我们通过 $p$  进行归纳论证并考虑长的精确序列：

\[
\cdots\to H^{j}(S;F_{S\setminus S_{p-1}})\to H^{j}(S;F)\to H^{j}(S;F_{S_{p-1}})\to\cdots.
\]

从  $H^{j}(S;F_{S_{p^{-1}}})=H^{j}(S_{p^{-1}};F_{S_{p^{-1}}})$  开始， $j>p$  的所有项都消失，并且我们获得  $j>p$  的  $H^{j}(S;F)=0$  。  $j>p$  的  $H_{c}^{j}(S;F)$  的消失源自等式  $H_{c}^{j}(S;F)=\lim_{\overrightarrow{U}}H^{j}(S;F_{U})$  ，其中  $U$  的范围涵盖  $S$  的相对紧凑的开放子集族。

(ii) 和 (iii) 通过 (3.1.8) 我们有：

\[
R\varGamma(S;\omega_{S})\simeq R\mathrm{H o m}(R\varGamma_{c}(S;A_{S}),A)\enspace.
\]

由于 $ R\Gamma_c(S; A_S) $  集中在度数 $ \leqslant p $  上， $ R\Gamma(S;\omega_S) $  集中在度数 $ \geqslant -p $  上，并且 $ H^{-p}(R\Gamma(S;\omega_S)) \simeq H^{-p}(RH_{\mathrm{Om}}(R\Gamma_c(S; A_S); A))  \simeq $  Hom(  $ H_c^p(S; A_S) $  , A) 上。   $ \square $ 

\statementlabel{命题 9.2.3}令 令 $S$ 为 $X$ 中维数 $\leq p$ 的闭次解析子集，并令 $S_0$ 为 $S$ 的局部闭次解析子集。则：

（一） $ j_{S.*}H^{-p}(\omega_{S}.) \simeq H_{S.}^{0}(j_{S*}H^{-p}(\omega_{S})), $ 

(ii) 如果  $ S_{\circ} $  是闭集且  $ S\setminus S_{\circ} $  具有维度  $ <p $  ，则  $ j_{S\circ*}H^{-p}(\omega_{S\circ})\simeq j_{S*}H^{-p}(\omega_{S}) $  ，

(iii) 如果  $ S_{0} $  在  $ S_{reg,p} $  中是开稠密的，则存在精确序列：

\[
0\to j_{S*}{\cal H}^{-p}(\omega_{S})\to j_{S,\bullet}\mathcal{o n}_{S_{\bullet}}\to j_{S\setminus S,\bullet}{\cal H}^{-p+1}(\omega_{S\setminus S_{\bullet}})~.
\]

\statementlabel{证明}(i) 我们有 $ R j_{S.*} \omega_{S} \simeq R \Gamma_{S} (\omega_{S}) $ 。由于 $ \omega_{S} $  和 $ \omega_{S} $  集中在度 $ \geq -p $  上，我们得到(i)。

(ii) 从  $ \omega_{S\mid S\setminus S} $  集中在 > - p 度的事实得出。

(iii) 考虑区分的三角形：

\[
R\varGamma_{S\setminus S\circ}(\omega_{S})\longrightarrow\omega_{S}\longrightarrow R j_{S\circ*}j_{S\circ}^{-1}\omega_{S}\mathop{\longrightarrow}_{+1}.
\]

它产生了长序列（在注意到  $ \omega_{S\backslash S_s} \simeq R\Gamma_{S\backslash S_s}(\omega_S) $  和  $ \omega_S|_{\mathcal{S}_s} \simeq or_{\mathcal{S}_s}[-p] $  之后）：

\[
H^{-p}(\omega_{S\setminus S_{\circ}})\to H^{-p}(\omega_{S})\to j_{S,\bullet}o\imath_{S_{\circ}}\to H^{-p+1}(\omega_{S\setminus S_{\circ}}).
\]

自 $ S \setminus S $ 以来。具有维度  $ < p $  ，第一项消失。 ⑨

现在我们将回到层的研究 $ \mathcal{CS}_p $ 。对于非负整数  $ p $  ，让  $ LCLS_p(X) $  表示维度  $ \leq p $  的  $ X $  的局部闭子解析子集的集合。 We introduce a partial order   $ < $   on   $ LCLS_p(X) $   as follows.对于 $ LCLS_p(X) $  中的 $ S_1 $  和 $ S_2 $  ， $ S_1 < S_2 $  当且仅当存在 $ S_1 \cap S_2 $  的子分析子集 $ S $  使得：

\[
\begin{cases}{S{~i s~o p e n~i n~}S_{1}{~a n d~c l o s e d~i n~}S_{2}{~a n d~}S_{1}\setminus S}\\ {{h a s~d i m e n s i o n~}<p\enspace.}\\ \end{cases}
\]

请注意，  $LCLS_p(X)$  是有向有序集，即对于  $LCLS_p(X)$  中的  $S_1$  和  $S_2$  ，始终存在  $S$  与  $S_1 < S$  和  $S_2 < S$  。事实上，选择  $S = (S_1 \setminus \partial S_2) \cup (S_2 \setminus \partial S_1)$  。请注意， $S < S_{\text{reg},p}$  对于任何 $S \in LCLS_p(X)$  。如果  $S_1 < S_2$  ，我们有一个规范态射：

\[
j_{S_{1}*}H^{-p}(\omega_{S_{1}})\to j_{S_{2}*}H^{-p}(\omega_{S_{2}})~.
\]

该态射由以下组合获得：

\[
j_{S_{1}*}H^{-p}(\omega_{S_{1}})\to j_{S*}H^{-p}(\omega_{S})\to j_{S_{2}*}H^{-p}(\omega_{S_{2}})~,
\]

其中  $S$  满足 (9.2.5)。这不取决于  $S$  的选择。事实上，如果  $S'$  是另一个满足 (9.2.5) 的集合，则  $S'' = S \cap S'$  仍然满足 (9.2.5) 并且态射  $j_{S*}H^{-p}(\omega_S) \to j_{S''*}H^{-p}(\omega_{S''})$  是明确定义的，并且是同构，因为  $S''$  在  $S$  中是闭和开的，而  $S\backslash S''$  的维度 <  $p$  。

\statementlabel{引理 9.2.4}我们有同构：

\[
\mathcal{C S}_{p}^{X}\simeq\varinjlim_{S\in{L C L S}_{p}(X)}j_{S*}{\cal H}^{-p}(\omega_{S})\simeq\varinjlim_{S^{\prime}*}{}^{{o r}_{S^{\prime}}},
\]

其中最后一个归纳极限中的  $S'$  范围贯穿  $p$  维亚解析子流形族  $X$  。

\statementlabel{证明}最后一个同构是显而易见的，因为对于  $ S \in LCLS_p(X) $  来说，有  $ S < S_{\text{reg}, p} $  。

令  $U$  为  $X$  的开子集，并让  $V$  为  $X$  与  $\overline{V} \subset U$  的开子分析子集。我们定义：

\[
\alpha:C S_{p}^{\prime}(U)\to\varinjlim_{S}\varGamma(V;j_{S*}{^{\circ}}\varkappa_{S})\to\varGamma\bigg(V;\varinjlim_{S}j_{S*}{^{\circ}}\varkappa_{S}\bigg)
\]

如下。

如果  $S$  是有向  $p$  维亚解析子流形，则  $S$  的方向定义  $\sigma r_{S \cap V}$  上  $V$  上的  $s$  部分。然后  $\alpha$  由  $[S] \mapsto s \in \Gamma(V; j_{S*} \circ r_{S})$  给出。取  $U$  和  $V$  的归纳极限，我们得到态射  $\mathcal{C} \mathcal{S}_{p}^{X} \to \varinjlim j_{S*} \circ r_{S}$  。这显然是同构。   $\square$ 

让 $CLS_p(X)$  表示维度 $\leq p$  的 $X$  的闭合子解析子集的集合。我们通过包含关系对 $CLS_p(X)$ 进行排序。如果  $S_1$  和  $S_2$  属于  $CLS_p(X)$  和  $S_1 \subset S_2$  ，则根据命题 9.2.3，我们有单射态射。

\[
j_{S_{1}*}{\cal H}^{-p}(\omega_{S_{1}})\to j_{S_{2}*}H^{-p}(\omega_{S_{2}})~.
\]

\statementlabel{定义 9.2.5}一套：

\[
\mathcal{D S}_{p}^{X}=\varinjlim_{S\in{C L S}_{p}(X)}j_{S*}{H}^{-p}(\omega_{S}),
\]

 $ \mathcal{L}_{\mathcal{S}_{p}^{X}} $  称为  $ X $  上的一束亚解析 p 循环（在基环  $ A $  上）。

如果不存在混淆的风险，我们会写  $ \mathcal{X}\mathcal{S}_{p} $  而不是  $ \mathcal{X}\mathcal{S}_{p}^{X} $  。

对于  $X$  上的  $A_X$  -模  $F$  ，设置  $\mathscr{L}_{\mathscr{S}_p}(F) = \mathscr{L}_{\mathscr{S}_p} \otimes F$  ，并将其称为具有  $F$  中的值的子分析  $p$  循环层。由于  $CLS_p(X)$  包含在  $LCLS_p(X)$  中，因此我们有一个精确的序列：

\[
0\to\mathcal{L S}_{p}\to\mathcal{C S}_{p}.
\]

让 $ S \in LCLS_p(X) $  。我们有一个显着的三角形：

\[
\omega_{\partial S}\longrightarrow\omega_{\overline{{S}}}\longrightarrow R j_{S*}\omega_{S}\mathop{\longrightarrow}_{+1},
\]

这产生了长的精确序列：

\[
0\to H^{-p}(\omega_{\overline{{S}}})\to j_{S*}H^{-p}(\omega_{S})\to H^{1-p}(\omega_{\partial S})\enspace.
\]

通过归纳极限，我们得到精确序列：

\[
0\to\mathcal{X S}_{p}\to\mathcal{C S}_{p}\to\mathcal{X S}_{p-1}\enspace,
\]

第一个箭头与（9.2.8）中的态射一致。然后我们定义边界算子。

\[
\partial_{p}\colon\mathcal{C S}_{p}\to\mathcal{C S}_{p-1}
\]

作为复合材料  $\mathcal{C}\mathcal{S}_p \to \mathcal{L}\mathcal{S}_{p-1} \to \mathcal{C}\mathcal{S}_{p-1}$  。如果不存在混淆的风险，我们会写  $\partial$  而不是  $\partial_p$  。

\statementlabel{命题 9.2.6}(i) C.S. 是层复合体（参见符号 1.3.9）。

（二） $ \mathcal{D}S_p \simeq \mathrm{Ker} \partial_p $ 。

(iii) 对于 X 上的任何层 F， $ \mathcal{C}\mathcal{S}_{p}(F) $  是软层。

(iv) 存在一个规范态射  $ \omega_X \to \mathcal{C}\mathcal{S} $  。在 $ \mathbf{D}^b(X) $ 。

(v) 如果  $ \Omega $  是  $ X $  的开子集，则  $ \mathcal{C}\mathcal{S}_p^X|_{\Omega} \simeq \mathcal{C}\mathcal{S}_p^\Omega $  和  $ \mathcal{Z}\mathcal{S}_p^X|_{\Omega} \simeq \mathcal{Z}\mathcal{S}_p^\Omega $  。

\statementlabel{证明}(i) 和 (ii) 是明确的。

(iii) 对于 $ X $  的开放子解析子集 $ W $ ，让 $ \rho_W $  表示在 $ j_{S*}H^{-p}(\omega_S) $  上通过态射定义的束 $ \mathcal{C}\mathcal{S}_p $  的自同态：

\[
j_{S*}{\cal H}^{-p}(\omega_{S})\to j_{S\cap\cal W*}{\cal H}^{-p}(\omega_{S\cap\cal W})\;.
\]

然后是 $X \setminus \overline{W}$  上的 $\rho_{\mathbf{w}}|_{\mathbf{w}} = \mathrm{id}$  、 $\rho_{\mathbf{w}}^2 = \rho_{\mathbf{w}}$  和 $\rho_{\mathbf{w}} = 0$  。现在，如果  $Z$  是  $X$  和  $s \in \Gamma(Z; \mathcal{C}\mathcal{S}_p(F))$  的闭子集，则存在  $Z$  的开邻域  $U$  和  $\tilde{s} \in \Gamma(U; \mathcal{C}\mathcal{S}_p(F))$  节，其对  $Z$  的限制是  $s$  。令  $W$  为  $X$  与  $Z \subset W \subset \overline{W} \subset U$  的子解析开子集。那么  $\rho_{\mathbf{w}}(\tilde{s}) \in \Gamma_{\overline{W}}(X; \mathcal{C}\mathcal{S}_p(F))$  是  $s$  的全局扩展。

(v) 由类似的论证证明。

(iv) 从  $ H^{-n}(\mathcal{CL}) = \mathrm{Ker} \partial_n = \mathcal{L}\mathcal{L}_n = H^{-n}(\omega_X) $  开始，我们获得态射  $ H^{-n}(\omega_X)[n] \to \mathcal{CL} $  。 □

\statementlabel{命题 9.2.7}令 L 为 X 上有限秩的局部自由束。

(i) 对于任何 $\alpha \in \Gamma(X; \mathcal{C}\mathcal{S}_p(L))$  ， $\text{supp}(\alpha)$  是纯维度 $p$  的闭合子解析子集。

(ii) 让 $ S \in CLS_p(X) $  。然后：

\[
\Gamma_{S}(X;\mathcal{X S}_{p}(L))\simeq H^{-p}(S;\omega_{S}\otimes L)\enspace.
\]

(iii) 让 $ S \in LCLS_p(X) $  。然后：

\[
\begin{aligned}{\varGamma(X;j_{S*}\mathbb{H}^{-p}(\omega_{S}\otimes L))}&{{}\simeq H^{-p}(S;\omega_{S}\otimes L)}\\ {}&{{}\simeq\left\{\alpha\in\varGamma_{\overline{{S}}}(X;\mathcal{C S}_{p}(L));{s u p p}(\partial\alpha)\subset\partial S\right\}.}\\ \end{aligned}
\]

\statementlabel{证明}(i) 局部而言，对于某些 p 维子解析子流形 S， $ \alpha $  部分属于  $ j_{S*}\circ\iota_{S} $ 。

(ii) 让  $ S' \in CLS_p(X) $  与  $ S \subset S' $  。然后 $ \Gamma_S(j_{S'^*}H^{-p}(\omega_S))\simeq j_{S^*}H^{-p}(\omega_S) $ 。这意味着(ii)。

(iii) 如果  $\alpha \in H^{-p}(S; \omega_S \otimes L)$  ，则明确  $\operatorname{supp}(\alpha) \subset \overline{S}$  和  $\operatorname{supp}(\partial \alpha) \subset \partial S$  。我们必须证明相反的情况。问题是局部的，我们可以假设  $L = A_X$  和  $\alpha \in \Gamma(S'; H^{-p}(\omega_S,)))$  对于某些  $S' \in LCLS_p(X)$  。然后 $\Gamma_{\overline{S}}(j_S, \star H^{-p}(\omega_S,))\simeq j_{\overline{S} \cap S' \star} H^{-p}(\omega_{\overline{S} \cap S'})$ 。因此我们可以假设  $S' \subset \overline{S}$  。将  $S'$  替换为  $(S \cap S')_{\operatorname{reg}, p} \cup (S \setminus \overline{S}')_{\operatorname{reg}, p}$  ，我们甚至可以假设  $S'$  在  $S$  和  $\operatorname{dim}(S \setminus S') < p$  中打开。  $S : \omega_S \setminus S' \to \omega_S \to \omega_S' \xrightarrow{+1}$  上的显着三角形产生精确的序列：

\[
0\to H^{-p}(\omega_{S})\to j_{S^{\prime}*}H^{-p}(\omega_{S^{\prime}})|_{S}\to j_{S\setminus S^{\prime}*}H^{-p+1}(\omega_{S\setminus S^{\prime}})|_{S}.
\]

由于  $ \partial\alpha|_{S} = 0 $  ，  $ \alpha $  属于  $ H^{-p}(\omega_S) $  。 □

\statementlabel{推论 9.2.8}令 L 为 X 上有限秩的局部自由束。则：

\[
\varGamma(X;\mathcal{C S}_{p}(L))\simeq\varinjlim_{S\in{L C L S}_{p}(X)}\varGamma(S;H^{-p}(\omega_{S})\otimes L)\enspace,
\]

\[
\begin{array}{r}{\varGamma(X;\mathcal{X S}_{p}(L))\simeq\varliminf_{S\in\overrightarrow{C L S}_{p}(X)}\varGamma(S;H^{-p}(\omega_{S})\otimes L)\enspace.}\end{array}
\]

为了证明下面的定理9.2.10，让我们研究一下亚解析链上的一些函子运算。

令 $f: Y \to X$  为流形的态射，并让 $\mathcal{GL}^Y$  和 $\mathcal{GL}^X$  分别表示 $Y$  和 $X$  上的亚解析链的复形。我们有：

\[
f_{!}^{\mathbf{\mathcal{C}}}\mathcal{S}_{p}^{\mathbf{Y}}=\operatorname*{l i m}_{\overrightarrow{s}}j_{S*}{H}^{-p}(\omega_{S}),
\]

其中  $S$  属于  $LCLS_p(Y)$  ，而  $f$  则适用于  $\overline{S}$  。对于这样的  $S$  ，集合  $S' = f(\overline{S})\backslash f(\partial S)$  属于  $LCLS_p(X)$  并且映射  $S \cap f^{-1}(S') \to S'$  是正确的。因此我们得到：

\[
f_{!}j_{S*}{\cal H}^{-p}(\omega_{S})\to j_{S^{\prime}*}{\cal H}^{-p}(\omega_{S^{\prime}})\to\mathcal{C S}_{p}^{X},
\]

其中第一个箭头是通过采用态射  $Rf_1\omega_{\mathfrak{S}}\to\omega_{f(\mathfrak{S})}$  的  $p$  次上同调获得的。通过归纳极限，我们得到态射 $f(\mathcal{C}\mathcal{S}_p^Y\to\mathcal{C}\mathcal{S}_p^X)$ 。由于该态射与微分交换，我们得到复形的态射：

\[
f_{!}\mathcal{C S}^{Y}\to\mathcal{C S}^{X}~.
\]

如果 F 是 X 上的一个层，则定义：

\[
f_{!}^{\phantom{\dag}}\mathcal{C S}_{.}^{Y}(f^{-1}F)\to\mathcal{C S}_{.}^{X}(F)\enspace.
\]

如果 $\alpha \in \Gamma(Y; \mathcal{C} \mathcal{S}_p(f^{-1}F))$  和 $f$  在supp(\textbackslash{}alpha) 上是真映射，我们用 $f_*(\alpha)$  表示 $\alpha$  的态射(9.2.13) 的像，并将其称为 $\alpha$  的直接像。

两个循环 $ \alpha \in \Gamma(X; \mathcal{C}\mathcal{S}_p^X(F)) $  和 $ \beta \in \Gamma(Y; \mathcal{C}\mathcal{S}_q^Y(G)) $  的外积 $ \alpha \boxtimes \beta $  由态射定义：

\[
j_{S_{1}*}H^{-p}(\omega_{S_{1}})\boxtimes j_{S_{2}*}H^{-q}(\omega_{S_{2}})\to j_{S_{1}\times S_{2}*}H^{-p-q}(\omega_{S_{1}\times S_{2}})
\]

其中  $S_{1} \in LCLS_{p}(X)$  和  $S_{2} \in LCLS_{q}(Y)$  。由于这些态射与微分交换，我们得到态射：

\[
\mathcal{C S}_{.}^{X}(F)\boxtimes\mathcal{C S}_{.}^{Y}(G)\to\mathcal{C S}_{.}^{X\times Y}(F\boxtimes G)\enspace.
\]

人们还可以定义亚解析循环的同伦概念。

\statementlabel{定义 9.2.9}让  $\gamma_0$  和  $\gamma_1$  属于  $\Gamma(X;\mathcal{L}_{\mathcal{S}_p^X})$  。有人说 $\gamma_0$  和 $\gamma_1$  是同伦的，如果存在 $\tau\in\Gamma_{X\times[0,1]}(X\times\mathbb{R};\mathcal{L}_{\mathcal{S}_{p+1}^{X\times\mathbb{R}}})$  使得：

\[
\partial\tau=i_{0*}\gamma_{0}-i_{1*}\gamma_{1},
\]

其中  $ i_t $  是包含  $ X \subsetneq X \times \mathbb{R} $  、  $ x \mapsto (x, t) $  。

令 $q$  表示 $X \times \mathbb{R} \to X$  的投影，让 $S$  表示 $\mathrm{supp}(\tau)$  与 $q$  的投影。将  $q_*$  应用于 (9.2.15) 我们得到  $\gamma_0 - \gamma_1 = \partial q_* \tau$  。因此  $\gamma_0$  和  $\gamma_1$  在  $H_p(\varGamma_S(X; \mathcal{C}\mathcal{S}_x))$  中将具有相同的图像。现在我们可以证明：

\statementlabel{定理 9.2.10}规范态射（参见命题 9.2.6），  $\omega_X \to \mathcal{C}^{\mathcal{L}_X}$  是  $\mathbf{D}^b(X)$  中的同构。

\statementlabel{证明}让 $n = \dim X$  。由于  $H^{-n}(\mathcal{C}\mathcal{S}_x^X) = \mathcal{D}\mathcal{S}_n = H^{-n}(\omega_X)$  ，足以证明  $H^{-p}(\mathcal{C}\mathcal{S}_x^X) = 0$  对于  $0 \leq p < n$  。让  $x \in X$  并让  $\alpha \in (\mathcal{D}\mathcal{S}_p)_x$  和  $0 \leq p < n$  。然后，如有必要，缩小 $X$ ，存在维度 $\leq p$ 的闭合子解析子集 $S$ ，使得 $\alpha$ 是 $H^{-p}(S; \omega_S)$ 的一部分（我们仍然用 $\alpha$ 表示）的图像。我们可以找到一个 $(n-1)$ 维流形 $Y$ ，使得 $X \simeq \mathbb{R} \times Y$ 位于 $x$ 的邻域， $f$ 表示投影 $X \to Y$ ， $f$ 在 $S$ 上是正确的。 （这一点留作练习。）设  $\beta \in H_{[0, \infty]}^0(\mathbb{R}; \mathcal{C}\mathcal{S}_1^{\mathbb{R}})$  为边界  $\partial\beta$  为  $[ \{0\} ]$  的规范元素，并设置  $\gamma = \beta \boxtimes \alpha$  。然后  $\gamma \in \Gamma(\mathbb{R} \times X; \mathcal{C}\mathcal{S}_{p+1})$  和  $\partial\gamma = \partial\beta \boxtimes \alpha = i_* \alpha$  ，其中  $i$  是映射  $X \hookrightarrow \mathbb{R} \times X$  ，  $x \mapsto (0, x)$  。让 $\varphi$  表示映射 $\mathbb{R} \times X \to X$  、 $(t, (s, y)) \mapsto (t + s, y)$  。那么 $\varphi : \mathbb{R} \times S \to X$  是真映射，因此 $\varphi|_{\operatorname{supp}(\gamma)}$  是正确的。因此  $\varphi_* \gamma \in \Gamma(X; \mathcal{C}\mathcal{S}_{p+1})$  和  $\partial\varphi_* \gamma = \varphi_* \partial\gamma = \varphi_* i_* \alpha = \alpha$  。这证明了  $H^{-p}(\mathcal{C}\mathcal{S}_x^X) = 0$  对于  $0 \leq p < n$  。   $\square$ 

\statementlabel{推论 9.2.11}(i)  $\mathcal{CS}_p$  和  $\mathcal{LS}_p$  的茎是扁平  $A$  模，并且对于  $X$  上的任何束  $F$  ，  $\mathcal{LS}_p(F)$  是态射  $\mathcal{CS}_p(F) \to \mathcal{CS}_{p-1}(F)$  的内核。

(ii)  $ \mathbf{D}^b(X) $  中存在自然同构  $ \omega_X \otimes F \simeq \mathcal{C}\mathcal{S}_(F) $  。

\statementlabel{证明}(i)  $\mathcal{C}\mathcal{S}_p$  的茎是自由 $A$  模的归纳极限，因此它是平坦的。考虑  $p \leq \dim X$  的确切序列  $0 \to \mathcal{L}\mathcal{S}_p \to \mathcal{C}\mathcal{S}_p \to \mathcal{L}\mathcal{S}_{p-1} \to 0$  。然后，如果  $\mathcal{L}\mathcal{S}_{p-1}$  是平坦的，则  $\mathcal{L}\mathcal{S}_p$  是平坦的，我们得到精确的序列  $0 \to \mathcal{L}\mathcal{S}_p(F) \to \mathcal{C}\mathcal{S}_p(F) \to \mathcal{L}\mathcal{S}_{p-1}(F) \to 0$  。因此，我们通过  $p$  归纳得到 (i) 。

(ii) 由定理 9.2.10 由 (i) 得出。 □

最后，让我们定义两个循环的交集。令  $S_{j}$  为  $n$  维流形  $X$  的维  $p_{j}$  (  $j \to 1, 2$  ) 的闭合子解析子集。令 $L_{1}$  和 $L_{2}$  为两个有限阶的局部自由层。考虑态射链：

\[
\begin{aligned}{H_{\mathcal{S}_{1}}^{-p_{1}}(X;\omega_{X}\otimes L_{1})}&{{}\otimes H_{\mathcal{S}_{2}}^{-p_{2}}(X;\omega_{X}\otimes L_{2})}\\ {}&{{}\simeq H_{\mathcal{S}_{1}}^{\bar{n}-p_{1}}(X;{\sigma}t_{X}\otimes L_{1})\otimes H_{\mathcal{S}_{2}}^{\bar{n}-p_{2}}(X;{\sigma}t_{X}\otimes L_{2})}\\ {}&{{}\to H_{\mathcal{S}_{1}\cap\mathcal{S}_{2}}^{2\bar{n}-p_{1}-p_{2}}(X;{\sigma}t_{X}\otimes L_{1}\otimes{\sigma}t_{X}\otimes L_{2})}\\ {}&{{}\simeq H_{\mathcal{S}_{1}\cap\mathcal{S}_{2}}^{\bar{n}-p_{1}-p_{2}}(X;\omega_{X}\otimes L_{1}\otimes L_{2}\otimes{\sigma}t_{X}).}\\ \end{aligned}
\]

\statementlabel{定义 9.2.12}让 $ C_j \in H_{S_j}^{p_j}(X; \omega_X \otimes L_j) \quad (j = 1, 2) $  。  $ C_1 \otimes C_2 $  中态射 (9.2.16) 的  $ H_{S_1 \cap S_2}^{n-p_1-p_2}(X; \omega_X \otimes L_1 \otimes L_2 \otimes \sigma t_X) $  中的图像称为  $ C_1 $  和  $ C_2 $  的交集，并表示为  $ C_1 \cap C_2 $  。如果 $ S_1 \cap S_2 $  是紧凑的，

 $ n = p_1 + p_2 $  和 $ L_1 \otimes L_2 \simeq \sigma r_X $  ，标量 $ \int_X C_1 \cap C_2 $  称为 $ C_1 $  和 $ C_2 $  的交集数，并表示为 $ \#(C_1 \cap C_2) $  。

回想一下  $ \int_{X} $  是态射  $ H_{c}^{0}(X;\omega_{X})\to A $  。

\statementlabel{备注 9.2.13}让 $ C_j \in H_{\mathcal{S}_i}^{p_j}(X; \omega_X) $ ， $ j = 1, 2 $ 。然后：

\[
C_{1}\cap C_{2}=(-1)^{(n-p_{1})(n-p_{2})}C_{2}\cap C_{1}
\]

在 $ H_{\mathcal{S}_1 \cap \mathcal{S}_2}^{n-p_1-p_2}(X; \omega_X \otimes \sigma t_X) $  中。这是因为如果  $ K_1 $  和  $ K_2 $  是两个复数，则该图：

\[
\begin{array}{c}{H^{p_{1}}(K_{1})\otimes H^{p_{2}}(K_{2})\xrightarrow{\phantom{H^{p_{1}}(K_{1})\otimes H^{p_{2}}(K_{2})}}\;H^{p_{1}+p_{2}}(K_{1}\otimes K_{2})}\\ {\bigg\downarrow\quad\bigg\downarrow}\\ {H^{p_{2}}(K_{2})\otimes H^{p_{1}}(K_{1})\xrightarrow{\phantom{H^{p_{2}}(K_{2})\otimes H^{p_{1}}(K_{1})}}\;H^{p_{1}+p_{2}}(K_{2}\otimes K_{1})}\\ \end{array}
\]

根据  $ p_{1}p_{2} $  是偶数或奇数，进行通勤或反通勤（参见备注 1.10.16）。

如果  $S_{1} \cap S_{2}$  具有维度  $\leqslant p_{1} + p_{2} - n$  ，则态射 (9.2.16) 定义（参见命题 9.2.6）：

(9.2.18)

\[
\begin{array}{r}{\varGamma_{\mathsf{S}_{1}}(X;\mathcal{D S}_{p_{1}}(L_{1}))\otimes\varGamma_{\mathsf{S}_{2}}(X;\mathcal{D S}_{p_{2}}(L_{2}))\to\varGamma_{\mathsf{S}_{1}\cap\mathsf{S}_{2}}(X;\mathcal{D S}_{p_{1}+p_{2}-n}(L_{1}\otimes L_{2}\otimes\mathfrak{o r}_{X}))~.}\end{array}
\]

我们仍然将这种态射称为交交态射，并且仍然用  $ \cap $  来表示它。

\statementlabel{备注 9.2.14}将 (9.2.12) 应用于  $\mathbf{a}_X: X \to \{\mathrm{pt}\}$  ，我们得到  $\mathbf{a}_X! \mathcal{C}\mathcal{S}^X \to A$  ，然后这给出了  $\mathbf{D}^b(X)$  中的态射  $\mathcal{C}\mathcal{S}^X \to \omega_X$  。这与定理 9.2.10 中构造的同构（直到符号）一致。

\subsection{拉格朗日圈}
在  $ T^{*}X $  中，亚解析各向同性子集支持的循环特别令人感兴趣。让 $ n = \dim X $  。

\statementlabel{定义 9.3.1}我们设定：

\[
\mathcal{L}_{X}=\varinjlim_{A}H_{A}^{0}(\pi^{-1}\omega_{X})
\]

其中  $A$  范围涵盖  $T^*X$  的所有闭圆锥亚解析各向同性子集族。我们将 $\mathcal{L}_X$  称为 $T^*X$  上的拉格朗日循环层。

从  $ \pi^{-1}\omega_X \simeq \omega_{T^*X} \otimes \sigma_{T^*X/X}[-n] $  开始，拉格朗日循环是一个 n 循环，其值在  $ \sigma_{T^*X/X} $  中。作为§2结果的一个特例，我们得到：

(a) 令 $A$  为二次曲线局部闭亚解析各向同性子集。然后：

\[
\begin{aligned}{H_{\varLambda}^{0}(\mathcal{L}_{X})}&{{}\simeq H_{\varLambda}^{0}(\pi^{-1}\omega_{X})}\\ {}&{{}\simeq j_{\varLambda*}H^{-n}(\omega_{\varLambda}\otimes\mathcal{o r}_{T^{*}X/X})\simeq H_{\varLambda}^{0}(\mathcal{Z S}_{n}(\mathcal{o r}_{T^{*}X/X}))\enspace.}\\ \end{aligned}
\]

特别是如果  $A$  维度  $n$  是平滑的：

\[
H_{\varLambda}^{0}({\mathcal{L}}_{\boldsymbol{X}})|_{\varLambda}\simeq o r_{\varLambda}\otimes o r_{{\mathbb{T}}^{*}{\boldsymbol{X}}/{\boldsymbol{X}}}\;.
\]

(b) 令  $s$  为  $\mathcal{L}_X$  在  $T^*X$  的开子集  $U$  上的一部分。那么  $\text{supp}(s)$  是纯维度  $n$  的闭圆锥亚解析各向同性子集。我们不知道  $\text{supp}(s)$  是否是定义 6.5.1 意义上的对合性。

让我们研究拉格朗日循环的一些功能特性。 令 $Y$ 为另一流形，并令 $m=\dim Y$。令 $\varLambda_X$（分别地，$\varLambda_Y$）为 $T^*X$（分别地，$T^*Y$）中的闭锥次解析迷向子集。通过考虑态射链来定义两个拉格朗日循环的外积：

\[
H_{A_{X}}^{0}(\pi^{-1}\omega_{X})\boxtimes H_{A_{Y}}^{0}(\pi^{-1}\omega_{Y})\to H_{A_{X}\times A_{Y}}^{0}(\pi^{-1}\omega_{X}\boxtimes\pi^{-1}\omega_{Y})\simeq H_{A_{X}\times A_{Y}}^{0}(\pi^{-1}\omega_{X\times Y})\enspace.
\]

对  $A_{x}$  和  $A_{Y}$  求归纳极限，我们得到：

\[
\mathcal{L}_{X}\boxtimes\mathcal{L}_{Y}\to\mathcal{L}_{X\times Y}
\]

现在让 $f: Y \to X$  成为流形的态射。回想一下映射 $f': Y \times_X T^*X \to T^*Y$  和 $f_\pi: Y \times_X T^*X \to T^*X$ （参见（4.3.2））。为了避免混淆，我们有时用 $\pi_Y$ （分别为 $\pi_X$ ，分别为 $\pi$ ）来表示投影 $T^*Y \to Y$ （分别为 $T^*X \to X$ ，分别为 $Y \times_X T^*X \to Y$ ）。

如前所述，令 $A_{X}$  和 $A_{Y}$  分别为 $T^{*}X$  和 $T^{*}Y$  的闭圆锥亚解析各向同性子集。

\statementlabel{命题 9.3.2}(i) 假设  $f_{\pi}$  对于  $t f^{\prime-1}(\varLambda_{\Upsilon})$  是正确的。那么就有一个自然态射：

\[
H_{A_{Y}}^{0}(T^{*}Y;\pi_{Y}^{-1}\omega_{Y})\to H_{f_{*}^{t f^{\prime}-1}(A_{Y})}^{0}(T^{*}X;\pi_{X}^{-1}\omega_{X})\enspace.
\]

(ii) 假设  $ f' $  对于  $ f_{\pi}^{-1}(\Lambda_{X}) $  是正确的。那么就有一个自然态射：

\[
H_{A x}^{0}(T^{*}X;\pi_{X}^{-1}\omega_{X})\to H_{{\mathfrak{t}}f^{\prime}f_{x}^{-1}(A x)}^{0}(T^{*}Y;\pi_{Y}^{-1}\omega_{Y})\enspace.
\]

首先请注意，对于假设 (i)， $ f_{\pi}^{t}f^{\prime-1}(\varLambda_{Y}) $  是  $ T^{*}X $  的闭圆锥亚解析子集，并且根据命题 8.3.11，该集合是各向同性的。类似地，根据假设 (ii)， $ t^{t}f_{\pi}^{-1}(\varLambda_{X}) $  是  $ T^{*}Y $  的闭圆锥亚解析各向同性子集。

\statementlabel{证明}(i) 令  $ \pi: Y \times_X T^* X \to Y $  为投影，考虑态射：

\[
\begin{aligned}{H_{\varLambda_{\gamma}}^{0}(T^{*}Y;\pi_{Y}^{-1}\omega_{Y})}&{{}\to H_{t f^{\prime-1}(\varLambda_{Y})}^{0}\left(Y\mathop{\times}_{X}T^{*}X;\pi^{-1}\omega_{Y}\right)}\\ {}&{{}\to H_{f_{\pi}^{t}f^{\prime-1}(\varLambda_{Y})}^{0}(T^{*}X;R f_{\pi!}\pi^{-1}\omega_{Y})}\\ \end{aligned}
\]

并使用态射：

\[
R f_{\pi!}\pi^{-1}\omega_{Y}\simeq\pi_{X}^{-1}R f_{!}\omega_{Y}\to\pi_{X}^{-1}\omega_{X}\enspace.
\]

(ii) 考虑态射：

\[
\begin{aligned}{H_{A x}^{0}(T^{*}X;\pi_{X}^{-1}\omega_{X})}&{{}\to H_{f_{\pi}^{-1}(A x)}^{0}\left(Y\mathop{\times}_{X}T^{*}X;\pi^{-1}f^{-1}\omega_{X}\right)}\\ {}&{{}\to H_{t f^{\prime}f_{x}^{-1}(A x)}^{0}(T^{*}Y;R^{t}f_{!}^{\prime}\pi^{-1}f^{-1}\omega_{X})\enspace.}\\ \end{aligned}
\]

总而言之，仍然需要构造态射：

\[
R^{\mathfrak{t}}f_{!}^{\prime}\pi^{-1}f^{-1}\omega_{X}\to\pi_{Y}^{-1}\omega_{Y}\enspace.
\]

我们按如下方式进行。考虑一下映射的图表：

\[
\begin{array}{ccc}Y\times Y&\xrightarrow{f_{1}}&Y\times X\\\left\langle\delta_{Y}\right.&&\left\langle\delta\right.\\\left.\varDelta_{Y}\right.&\xrightarrow{\sim}&\varDelta_{f}\end{array}
\]

并注意  $ t'f'_1: T_{\Delta_f}^*(Y \times X) \to T_{\Delta_Y}^*(Y \times Y) $  与  $ t'f' $  同构。将命题 4.3.5 应用于层  $ \delta_* f^{-1} \omega_X $  ，我们得到：

\[
\mathfrak{t}_{!}^{\prime}\mu_{\varDelta_{f}}(\delta_{\ast}f^{-1}\omega_{X})\to\mu_{\varDelta_{Y}}(\omega_{Y/X}\otimes f_{1}^{-1}\delta_{\ast}f^{-1}\omega_{X})\to\mu_{\varDelta_{Y}}(\delta_{Y\ast}\omega_{Y})
\]

(9.3.5)如下。 ⑨

注意态射：

\[
\pi^{-1}f^{-1}\omega_{X}\to{}^{t}f^{\prime!}\pi_{Y}^{-1}\omega_{Y}
\]

从(9.3.5)推导出来的是同构。

\statementlabel{定义 9.3.3}在命题9.3.2的情况下，1用 $ f_{*} $ 表示(9.3.3)中的态射，用 $ f^{*} $ 表示(9.3.4)中的态射。

示例 9.3.4。 (i) 如果  $X$  是 0 维流形  $\{pt\}$  ，则  $\mathcal{L}_X = A$  。 1 用 $[pt]$  表示对应于 $1 \in A$  的循环。

(ii) 设  $ a_x $  为态射  $ X \to \{\text{pt}\} $  。一套：

\[
[T_{x}^{*}X]=a_{x}^{*}[{\mathrm{p t}}]~.
\]

(iii) 令 $ f: Y \hookrightarrow X $  为封闭嵌入。一套：

\[
[T_{Y}^{\ast}X]=f_{\ast}[T_{Y}^{\ast}Y]\enspace.
\]

(iv) 令  $f: Y \to X$  为流形的态射，并假设  $f$  横向于  $X$  的闭子流形  $Z$  。然后：

\[
f^{*}[T_{Z}^{*}X]=[T_{f^{-1}(Z)}^{*}Y]\enspace.
\]

让我们定义拉格朗日循环的交集概念。

令  $ \sigma: X \to T^*X $  为  $ \pi $  的连续部分（即：  $ \pi \circ \sigma = \mathrm{id}_X $  ）。我们有一个同构：

\[
\begin{aligned}{H_{\sigma(X)}^{0}(T^{*}X;\pi^{!}A_{X})}&{{}\simeq H^{0}(X;\sigma^{!}\pi^{!}A_{X})}\\ {}&{{}=H^{0}(X;A_{X})\enspace.}\\ \end{aligned}
\]

请注意  $ \pi^!A_X \simeq \omega_{T^*X} \otimes \pi^{-1} \circ \tau_X[-n] $  ，其中  $ n = \dim X $  。如果选择 $ T^*X $  方向，则 $ \pi^!A_X \simeq \pi^{-1}\omega_X $  。

\statementlabel{定义 9.3.5}(i)  $ 1 \in H^0(X; A_X) $  的同构(9.3.9)在 $ H_{\sigma(X)}^0(T^*X; \pi^!A_X) $  中的图像用[ $ \sigma $  ]表示。

(ii) 1 用 $ [\sigma_{0}] $  表示与零部分相关的循环。

令  $S_{1}$  和  $S_{2}$  为  $T^{*}X$  的两个闭子集。考虑态射链：

\[
\begin{aligned}{H_{S_{1}}^{0}(T^{*}X;\pi^{!}A_{X})}&{{}\otimes H_{S_{2}}^{0}(T^{*}X;\pi^{-1}\omega_{X})}\\ {}&{{}\to H_{S_{1}\cap S_{2}}^{0}(T^{*}X;\pi^{!}A_{X}\otimes\pi^{-1}\omega_{X})}\\ {}&{{}\to H_{S_{1}\cap S_{2}}^{0}(T^{*}X;\pi^{!}\omega_{X})}\\ {}&{{}\simeq H_{S_{1}\cap S_{2}}^{0}(T^{*}X;\omega_{T^{*}X})\enspace.}\\ \end{aligned}
\]

如定义 9.2.12 中所示，如果  $ \gamma \in H_S^0(T^*X; \pi^!A_X) $  和  $ \delta \in H_S^0(T^*X; \pi^{-1}\omega_X) $  ，我们用  $ \gamma \cap \delta $  表示  $ \gamma \otimes \delta $  通过 (9.3.10) 中的态射链的图像，我们将其称为  $ \gamma $  和  $ \delta $  的交集。如果  $ S_1 \cap S_2 $  是紧致的，则 1 用  $ \#(\gamma \cap \delta) $  表示标量  $ \int_{T \cdot X} \gamma \cap \delta $  。如果  $ p $  是  $ S_1 \cap S_2 $  的孤立点，则  $ \#(\gamma \cap \delta)_p $  表示标量  $ \int_U \gamma \cap \delta $  ，其中  $ U $  是  $ d $  的足够小的邻域。

令 $S$  为 $T^*X$  的闭子集，使得 $\pi$  在 $S$  上是正确的。态射 $R\pi_1\omega_{T^*X} \to \omega_X$  定义态射：

\[
\alpha:H_{S}^{0}(T^{\ast}X;\omega_{T^{\ast}X})\to H_{\pi(S)}^{0}(X;\omega_{X})\enspace.
\]

另一方面，令 $A$  为闭圆锥曲线集。态射 $R\pi_* \pi^{-1}\omega_X \xrightarrow{\sim} \omega_X$  定义态射：

\[
\beta:H_{\varLambda}^{0}(T^{\ast}X;\pi^{-1}\omega_{X})\to H_{\pi(A)}^{0}(X;\omega_{X})\enspace.
\]

\statementlabel{命题 9.3.6}令 $ \sigma $ 为 $ \pi $ 的连续部分，并令 $ \lambda $ 为A支持的拉格朗日循环。则：

\[
\alpha([\sigma]\cap\lambda)=\beta(\lambda).
\]

\statementlabel{证明}考虑下面的交换图（9.3.13）就足够了，其中 $R$ 被省略。这里，左下方格的交换律由 $\sigma^1\pi^1 \simeq$  id 得出，右方方格的交换律由 $\pi_* \pi^{-1} \to$  id 和 $\pi_* \pi^{-1} \to \pi_* \sigma_* \sigma^{-1} \pi^{-1} \simeq$  id 相等这一事实得出。   $\square$ 

\[
\begin{array}{ccc}\Gamma(X;A_{X})\otimes\varGamma_{\varLambda}(T^{*}X;\pi^{-1}\omega_{X})&\longrightarrow&\varGamma(X;A_{X})\otimes\varGamma_{\pi(A)}(X;\pi_{*}\pi^{-1}\omega_{X})\\\downarrow&&\downarrow\\\varGamma_{\sigma(X)}(T^{*}X;\pi^{1}A_{X})\otimes\varGamma_{\varLambda}(T^{*}X;\pi^{-1}\omega_{X})&\longrightarrow&\varGamma(X;A_{X})\otimes\varGamma_{\sigma^{-1}(A)}(X,\sigma^{-1}\pi^{-1}\omega_{X})\\\downarrow&&\downarrow\\\varGamma_{\sigma(X)\cap\varLambda}(T^{*}X;\omega_{T^{*}X})&&\downarrow\\\alpha&&\varGamma_{\pi(A)}(X;\omega_{X})\\\varGamma_{\sigma^{-1}(A)}(X;\omega_{X})&\longrightarrow&\varGamma_{\pi(A)}(X;\omega_{X})\end{array}
\]

\subsection{特征圈}
从现在开始，直到本章结束，我们假设基环 $ A $ 是一个特征为零的域，并用 $ k $ 表示它。让 $ F \in \text{Ob}(\mathbf{D}_{\mathbf{R}-c}^b(X)) $  。考虑态射链：

\[
\begin{aligned}{{R\mathcal{H}o m}(F,F)}&{{}\simeq{R\pi_{*}\mu\mathcal{h}o m}(F,F)}\\ {}&{{}\simeq{R\pi_{*}}R\varGamma_{\mathbb{S}\mathbb{S}(F)}\mu\mathcal{h}o m(F,F)}\\ {}&{{}\simeq{R\pi_{*}}R\varGamma_{\mathbb{S}\mathbb{S}(F)}\mu_{\varDelta}(F\boxtimes\mathbf{D}F)}\\ {}&{{}\to{R\pi_{*}}R\varGamma_{\mathbb{S}\mathbb{S}(F)}\mu_{\varDelta}(\delta_{*}(F\otimes\mathbf{D}F))}\\ {}&{{}\to{R\pi_{*}}R\varGamma_{\mathbb{S}\mathbb{S}(F)}\mu_{\varDelta}(\delta_{*}\omega_{\chi})}\\ {}&{{}\simeq{R\pi_{*}}R\varGamma_{\mathbb{S}\mathbb{S}(F)}(\pi^{-1}\omega_{\chi})\enspace.}\\ \end{aligned}
\]

\statementlabel{定义 9.4.1}$\mathrm{id}_{F} \in \mathrm{Hom}(F, F)$  在 $H_{\mathrm{SS}(F)}^0(T^*X; \pi^{-1}\omega_X) \subset H^0(T^*X; \mathcal{L}_X)$  中的图像称为 $F$  的特征环，记为 $\mathrm{CC}(F)$  。

让我们比较一下 R 构造层上的操作和拉格朗日循环上的操作。

令  $Y$  为另一个流形，并令  $G \in \mathrm{Ob}(D_{\mathbb{R}-c}^{b}(Y))$  。使用命题 4.4.8，可以轻松检查公式：

\[
\mathrm{C C}(F\boxtimes G)=\mathrm{C C}(F)\boxtimes\mathrm{C C}(G)\enspace.
\]

现在让 $f: Y \to X$  成为流形的态射。

\statementlabel{命题 9.4.2}假设 $f$  在supp(G) 上是正确的。然后：

\[
\mathrm{C C}(R f_{\ast}G)=f_{\ast}\mathrm{C C}(G).
\]

\statementlabel{证明}用  $ \delta_X(\text{resp.} \ \delta_Y) $  表示对角线嵌入  $ X \hookrightarrow X \times X $  （分别为  $ Y \hookrightarrow Y \times Y $  ），用  $ \delta $  表示图嵌入  $ Y \hookrightarrow X \times Y $  ，并考虑该图（参见 (4.4.3)）：

\[
\begin{array}{ccc}Y\times Y&\xrightarrow[f_{1}]{}&X\times Y\xrightarrow[f_{2}]{}X\times X\\\downarrow\delta_{Y}&&\downarrow\delta\\\quad\quad\quad\quad\quad\quad\quad\quad\quad&\quad\quad\quad\quad&\quad\quad\quad\quad\\\quad\quad\quad\quad\quad\quad\quad\quad\quad&&\quad\quad\quad\quad\\\varDelta_{Y}&\xrightarrow{\sim}&\varDelta_{f}\quad\xrightarrow[f]{\quad}\quad\varDelta_{X}\quad.\end{array}
\]

用 $\pi_{Y}$ （分别 $\pi$ ， $\pi_{X}$ ）表示投影 $T^{*}Y \to Y$ （分别 $Y \times_{X} T^{*}X \to Y$ ， $T^{*}X \to X$ ），并设置为简称 $S = \mathrm{SS}(G)$ ， $S' = {}^{t}f^{' - 1}(S)$ ， $S'' = f_{\pi}(S')$ 。存在自然态射（参见 IV §3）：

\[
{}^{t}f_{1}^{\prime}\circ\mu_{\varDelta_{Y}}\to\mu_{\varDelta_{f}}\circ R f_{1!}~,\qquad f_{2\pi!}\circ\mu_{\varDelta_{f}}\to\mu_{\varDelta_{Y}}\circ R f_{2!}
\]

这产生了下面的交换图（9.4.4），其中我们省略了简写 R 来表示导出函子。

因此  $\varGamma_{S'}(T^*X; \pi_X^{-1}\omega_X)$  通过  $\varGamma_{S'}(Y \times_X T^*X; \pi^{-1}\omega_Y)$  中的  $\mathrm{CC}(G) \in \varGamma_S(T^*Y; \pi_Y^{-1}\omega_Y)$  的图像与  $\mathrm{CC}(Rf_!G)$  一致。   $\square$ 

\statementlabel{命题 9.4.3}假设  $f'$  对于  $f_{\pi}^{-1}(\mathrm{SS}(F))$  是正确的（即：  $f$  对于  $F$  来说是非特征的）。然后：

\[
CC(f^{-1}F)=f^{\ast}C C(F).
\]

\statementlabel{证明}我们设置 $ S = SS(F) $ ， $ S' = f_{\pi}^{-1}(S) $ ， $ S'' = 'f'(S') $ 。

考虑一下图表：

\[
\begin{array}{c c c c c}{Y\times Y\xrightarrow{f_{1}}}&{Y\times X}&{\xrightarrow{f_{2}}}&{X\times X}&{}\\ {\displaystyle\bigoplus_{\delta_{Y}}}&{\displaystyle\bigoplus_{\delta}}&{\displaystyle\bigoplus_{\delta_{X}}}&{\displaystyle\bigoplus_{\delta_{X}}}&{}\\ {\varDelta_{Y}}&{\xrightarrow{\sim}}&{\varDelta_{f}}&{\xrightarrow[f]{}}&{\varDelta_{X}}\\ \end{array}.
\]

并如上所述，分别用  $ \pi_{Y}, \pi, \pi_{X} $  表示投影  $ T^{*}Y \to Y, Y \times_{X} T^{*}X \to Y $  和  $ T^{*}X \to X $  。存在自然态射：

\[
f_{2\pi}^{-1}\circ\mu_{\varDelta_{X}}\to\mu_{\varDelta_{f}}\circ f_{2}^{-1}~,\qquad{}^{t}f_{1!}^{\prime}\circ\mu_{\varDelta_{f}}\to\mu_{\varDelta_{Y}}\circ\left(\omega_{Y/X}\otimes f_{1}^{-1}\right)~.
\]

\begin{center}
\includegraphics[width=0.54\textwidth]{流行上的层_Chn-assets/img_in_image_box_284_253_929_1643.png}
\end{center}

由于  $f'$  在  $S'$  上是真映射，所以我们有下面的交换图 (9.4.6)。因此 $\mathrm{CC}(f^{-1}F)$ 是 $\mathrm{CC}(F) \in H_S^0(T^*X; \pi_X^{-1} \omega_X)$ 通过 $H_S^0(Y \times_X T^*X; f_{\pi}^{-1} \pi_X^{-1} \omega_X)$ 在 $H_S^0(T^*Y; \pi_Y^{-1} \omega_Y)$ 中的图像。   $\square$ 

让“a”表示 $ T^*X $  上的对映图。由于“a”将亚解析各向同性子集发送到亚解析各向同性子集，因此它定义了一个态射：

\[
a_{*}:a^{-1}\mathcal{L}_{X}\to\mathcal{L}_{X}\enspace.
\]

如果  $\lambda$  是拉格朗日循环，我们用  $\lambda^{a}$  表示其镜像  $a_{*}$  。请注意，对于  $X$  的子流形  $Y$  ，

\[
[T_{Y}^{\ast}X]^{a}=(-1)^{\dim Y}[T_{Y}^{\ast}X]~.
\]

事实上，在识别 $a^!$ 和 $a^{-1}$ 之后， $RG_{T\notin X}(\pi^{-1}\omega_X) \simeq \omega_{T\notin X/T^*X} \otimes \pi^{-1}\omega_X$ 、 $a^{-1}\omega_{T\notin X/T^*X} \to \omega_{T\notin X/T^*X}$ 和 $a^!\omega_{T_Y^*X/T^*X} \to \omega_{T_Y^*X/T^*X}$ 等于 $(-1)^{\dim Y}$ 。

\statementlabel{命题 9.4.4}让 $ F \in \mathrm{Ob}(\mathbb{D}_{R-c}^b(X)) $  。然后：

\[
\mathrm{C C}(\mathbf{D}_{X}F)=(\mathrm{C C}(F))^{a}.
\]

\statementlabel{证明}其中一个有  $ \mu_{\Delta}(\mathbf{DF} \boxtimes \mathbf{DDF}) \simeq a^{-1} \mu_{\Delta}(F \boxtimes \mathbf{DF}) $  。 □

让我们注意， $\mathrm{CC}(F)$  是一个微局部概念，在以下意义上。如果  $\Omega$  是  $T^*X$  的开子集，并且  $F$  和  $F'$  在  $\mathbf{D}^b(X;\Omega)$  中是同构的，那么我们有  $\mathrm{CC}(F)|_{\Omega}=\mathrm{CC}(F')|_{\Omega}$  。这是从  $\mu\mathrm{hom}(F,F)|_{\Omega}\simeq\mu\mathrm{hom}(F',F')|_{\Omega}$  得出的。利用这个注释，我们可以更明确地计算 $\mathrm{CC}(F)$ 。如果  $X$  是一个点，我们有  $\mathrm{CC}(F)=\chi(F)=\sum_{j}(-1)^{j}\dim H^j(F)$  。如果  $Y$  是一个闭子流形，并且如果  $H^j(F)$  是所有  $j$  的局部常数秩  $m_j$  束，则  $\mathrm{CC}(F)=\sum_{j}(-1)^{j}m_j[T^*X]$  。事实上，让  $i:Y\hookrightarrow X$  表示  $Y$  的嵌入，并让  $a_Y:Y\to\{\mathrm{pt}\}$  。然后 $F=Ri_*a_Y^{-1}(V)$  对于一些 $V\in\mathrm{Ob}(\mathbf{D}^b(\mathfrak{Mod}^f(k)))$  ，其结果遵循命题9.4.2 和9.4.3。特别是，我们有  $X$  的闭合子流形  $Y$  ：

\[
\mathbf{C C}(k_{Y})=[\mathbf{\nabla}T_{Y}^{*}\mathbf{X}].
\]

为了处理一般情况，我们引入层：

\[
\mathcal{C L}_{X}=\varinjlim_{\lambda}H_{\lambda}^{0}(\pi^{-1}\omega_{X}),
\]

其中  $A$  范围涵盖  $T^*X$  的局部闭亚解析圆锥各向同性子集族。我们称 $\mathcal{CL}_X$ 为拉格朗日链束。我们有以下内容：

\[
\mathcal{L}_{X}\subset\mathcal{C L}_{X}\subset\mathcal{C S}_{n}^{T^{*}X}(o\tau_{T^{*}X/X})~.
\]

\[
\begin{aligned}\operatorname{Hom}(F,F)&\xrightarrow{\quad\uparrow\quad}\quad\operatorname{Hom}(f^{-1}F,f^{-1}F)\\\mathcal{F}_{\delta}(X\times X;F\boxtimes\operatorname{DF})&\xrightarrow{\quad\downarrow\quad}\quad\mathcal{F}_{\delta^{\prime}}(Y\times X;f^{-1}F\boxtimes\operatorname{DF})\xrightarrow{\quad\downarrow\quad}\quad\mathcal{F}_{\delta^{\prime}}(Y\times Y;f^{-1}F\boxtimes\operatorname{D}(f^{-1}F))\\\downarrow&\quad\downarrow\quad\downarrow\\\mathcal{F}_{\delta}(T^{*}X;\mu_{\delta,x}(F\boxtimes\operatorname{DF}))&\xrightarrow{\quad\downarrow\quad}\quad\mathcal{F}_{\delta^{\prime}}(Y\times T^{*}X;\mu_{\delta,f}(f^{-1}F\boxtimes\operatorname{DF}))&\xrightarrow{\quad\downarrow\quad}\quad\mathcal{F}_{\delta^{\prime\prime}}(T^{*}Y;\mu_{\delta,f}(f^{-1}F\boxtimes\operatorname{DF}))&\xrightarrow{\quad\downarrow\quad}\quad\mathcal{F}_{\delta^{\prime\prime}}(T^{*}Y;\mu_{\delta,f}(f^{-1}F\boxtimes f^{1}\mathbf{DF}))\\\downarrow&\quad\downarrow\quad\downarrow\\\mathcal{F}_{\delta}(T^{*}X;\mu_{\delta,x}(\delta_{X*}(F\otimes\mathbf{DF})))&\xrightarrow{\quad\downarrow\quad}\quad\mathcal{F}_{\delta^{\prime}}(Y\times T^{*}X;\mu_{\delta,f}(\delta_{X*}(f^{-1}F\otimes f^{-1}\mathbf{DF})))&\xrightarrow{\quad\downarrow\quad}\quad\mathcal{F}_{\delta^{\prime\prime}}(T^{*}Y;\mu_{\delta,f}(\delta_{Y*}(f^{-1}F\otimes f^{1}\mathbf{DF})))\\\downarrow&\quad\downarrow\quad\downarrow\\\mathcal{F}_{\delta}(T^{*}X;\mu_{\delta,x}(\delta_{X*}\omega_{X}))&\xrightarrow{\quad\downarrow\quad}\quad\mathcal{F}_{\delta^{\prime}}(Y\times T^{*}X;\mu_{\delta,f}(\delta_{X*}(f^{-1}\omega_{X}))&\xrightarrow{\quad\downarrow\quad}\quad\mathcal{F}_{\delta^{\prime\prime}}(T^{*}Y;\mu_{\delta,f}(\delta_{Y*}\omega_{Y}))\\\downarrow&\quad\downarrow\quad\downarrow\\\mathcal{F}_{\delta}(T^{*}X;\pi_{X}^{-1}\omega_{X})&\xrightarrow{\quad\downarrow\quad}\quad\mathcal{F}_{\delta^{\prime}}(Y\times T^{*}X;\pi^{-1}f^{-1}\omega_{X})&\xrightarrow{\quad\downarrow\quad}\quad\mathcal{F}_{\delta^{\prime\prime}}(T^{*}Y;\pi_{Y}^{-1}\omega_{Y})\\\text{Diagram(9.4.6)}&\end{aligned}
\]

现在让  $ F \in \text{Ob}(\mathbb{D}_{\mathbb{R}-c}^b(X)) $  并让  $ \varLambda = \text{SS}(F) $  。存在一个开密集子集 $ \varLambda_o \subset \varLambda $ ，使得：

\[
\begin{aligned}&\varLambda_{\circ}\text{ is a subanalytic submanifold, and denoting by}\\ &\{\varLambda_{\alpha}\}_{\alpha}\text{ the connected components of }\varLambda_{\circ},\text{ for each }\alpha\\ &\text{there is a submanifold }X_{\alpha}\text{ of }X\text{ such that}\\ &\varLambda_{\alpha}\subset T_{X_{\alpha}}^{*}X.\end{aligned}
\]

对于这样的 $\mathcal{A}_\alpha$ ，我们用 $[A_\alpha]$ 表示由 $\overline{\mathcal{A}}_\alpha$ 支持的链，即 $\mathcal{A}_\alpha$ 上的 $[T_{X_\alpha}^*X]$ 。如果  $U$  是一个开子集，使得  $\mathcal{A}_\alpha = T_{X_\alpha}^*X \cap U$  ，我们有时也会写  $[T_{X_\alpha}^*X] \cap U$  而不是  $[A_\alpha]$  。让 $p \in \mathcal{A}_\alpha$  。根据命题 6.6.1(ii)，存在  $V \in \mathrm{Ob}(\mathbf{D}^b \mathfrak{M}_\alpha \mathrm{Mod}^f(k))$  使得  $F \simeq V_{X_\alpha}$  位于  $\mathbf{D}^b(X; p)$  中。因此  $\mathrm{CC}(F) = m_\alpha[\mathcal{A}_\alpha]$  位于  $p$  的邻域，以及  $m_\alpha = \chi(V)$  。由于  $\mathcal{A}_\alpha$  是相连的，所以我们在  $\mathcal{A}_\alpha$  的邻域中找到  $\mathrm{CC}(F) = m_\alpha[\mathcal{A}_\alpha]$  ，并且如果  $\Omega$  是  $T^*X$  的开稠子集，这样  $\Omega \cap \mathcal{A} = \mathcal{A}_\circ$  ，我们就得到  $\mathrm{CC}(F)|_{\Omega} = \sum_{\alpha} m_\alpha[\mathcal{A}_\alpha]|_{\Omega}$  。因此：

\[
\mathrm{C C}(F)=\sum_{\alpha}m_{\alpha}[\varLambda_{\alpha}]\quad\mathrm{i n}\quad\mathcal{C L}_{X}\enspace.
\]

\statementlabel{命题 9.4.5}(i) 让 $ F \in \mathrm{Ob}(\mathbf{D}_{\mathbf{R}-c}^{b}(X)) $  。那么  $ \mathrm{CC}(F) $  是环  $ \mathbb{Z} $  上的拉格朗日循环。

(ii) 令 $ F' \longrightarrow F \longrightarrow F'' \xrightarrow[+1]{} $  为 $ \mathbf{D}_{\mathbf{R}-c}^{b}(X) $  中的一个显着三角形。然后：

\[
\mathrm{C C}(F)=\mathrm{C C}(F^{\prime})+\mathrm{C C}(F^{\prime \prime})~.
\]

特别是：

\[
\mathbf{C C}(F[k])=(-1)^{k}\mathbf{C C}(F).
\]

\statementlabel{证明}(i) 由(9.4.10)得出。

(ii) 令 $A$  为包含 $\mathrm{SS}(F') \cup \mathrm{SS}(F'')$  的 $T^*X$  的闭圆锥亚解析各向同性子集，并令 $A_o \subset A$  为满足(9.4.9) 的开稠密子集。然后  $\mathrm{CC}(F) = \sum_{\alpha} m_\alpha [A_\alpha]$  并与  $F'$  和  $F''$  类似，分别用  $m_\alpha'$  和  $m_\alpha''$  替换  $m_\alpha$  。让 $p \in T_X^*X$  。根据命题 6.6.1 (ii)， $\mathbf{D}^b(\mathfrak{M}_\alpha\mathfrak{M}^b(f)(k))$  中存在一个显着的三角形  $V' \longrightarrow V \longrightarrow V'' \xrightarrow{+1}$ ，使得  $F' \longrightarrow F \longrightarrow F'' \xrightarrow{+1}$  与  $\mathbf{D}^b(X; p)$  中的  $V_{X_s}' \longrightarrow V_{X_s} \longrightarrow V_{X_s}'' \xrightarrow{+1}$  同构。因此 $m_\alpha = m_\alpha' + m_\alpha'$  。

示例 9.4.6。用 $t$  表示 $X = \mathbb{R}$  上的坐标并考虑这些集合：

\[
Z_{\pm}=\left\{\pm t\geqslant0\right\},\qquad U_{\pm}=\left\{\pm t>0\right\}.
\]

让  $(t,\tau)$  表示  $T^*X$  上的关联坐标。考虑拉格朗日链：

\[
\alpha_{\pm}=\left[T_{X}^{\ast}X\right]\cap\left\{\pm t>0\right\}\;,\qquad\beta_{\pm}=\left[T_{\{0\}}^{\ast}X\right]\cap\left\{\pm\tau>0\right\}\;.
\]

（在这里和后续中，如果  $\varphi$  是  $X$  上的实函数，我们通常写成  $\{\varphi > 0\}$  来表示  $T^*X$  的子集  $\pi^{-1}\{x; \varphi(x) > 0\}$  ，简称。）然后：

\[
H_{\{t\tau=0\}}^{0}(T^{*}X;\mathcal{L}_{X})=\mathbb{Z}[T_{X}^{*}X]\oplus\mathbb{Z}[T_{\{0\}}^{*}X]\oplus\mathbb{Z}(\alpha_{+}+\beta_{+})~.
\]

（我们考虑 Z 中系数的亚解析链。）现在我们有：

（一） $ CC(k_X) = [T_X^* X] = \alpha_+ + \alpha_- $ 

（二） $ CC(k_{\{0\})} = [T_{\{0\}}^{*}X] = \beta_{+} + \beta_{-} $ ，

（三） $ CC(k_{z_{+}}) = \alpha_{\pm} + \beta_{\pm} $ ，

（四） $ CC(k_{U_{\pm}}) = \alpha_{\pm} - \beta_{\pm} $ 

事实上前两个公式已经被证明（参见（9.4.8））。让我们计算 $\mathrm{CC}(k_{U_+})$  。该循环由集合  $\{t \geq 0, \tau = 0\}\cup \{t = 0, \tau \leq 0\}$  支持。

关于 $\{t > 0, \tau = 0\}$ ， $k_{U_+} \simeq k_X$ 。关于 $\{t = 0, \tau < 0\}$ 、 $k_{U_+} \simeq k_{\{0\}}[-1]$ 。这给出了  $\mathrm{CC}(k_{U_+}) = \alpha_+ - \beta_-$  。其他公式也同样证明。

示例 9.4.7。令 $ (t,x)=(t,x_1,\ldots,x_n) $  为 $ X = \mathbb{R}^{1+n} $  上的坐标。考虑集合：

\[
Z_{\pm}=\left\{\pm t\geqslant\left|x\right|\right\},\qquad Z_{0}=\left\{\left|t\right|\leqslant\left|x\right|\right\},
\]

\[
U_{\varepsilon}=\operatorname{I n t}Z_{\varepsilon}~,\qquad(\varepsilon=+,-,0)~,~\quad S_{\pm}=\left\{\pm t=|x|>0\right\}~,
\]

其中 $ |x|^2 = \sum_j x_j^2 $  。

让  $(t, x; \tau, \xi)$  表示  $T^*X$  上的关联坐标。定义拉格朗日链：

\[
\sigma_{\pm}=[T_{X}^{*}X]\cap U_{\pm}~,
\]

\[
\tau_{\varepsilon_{1}\varepsilon_{2}}=[T_{S_{\varepsilon_{1}}}^{*}X]\cap\{\varepsilon_{2}\tau>0\}~,qquad\varepsilon_{1}=\pm1~,\quad\varepsilon_{2}=\pm1~,
\]

\[
\gamma_{\pm}=[T_{\{0\}}^{*}X]\cap\{\pm\tau>|\xi|\}\ ,
\]

\[
\gamma_{0}=[T_{\{0\}}^{*}X]\cap\{|\tau|<|\xi|\}\enspace.
\]

然后我们有：

（一） $ CC(k_x) = \sigma_+ + \sigma_0 + \sigma_- $ 

（二） $ CC(k_{\{0\}}) = \gamma_{+} + \gamma_{0} + \gamma_{-} $ 

（三） $ CC(k_{z_{\pm}}) = \sigma_{\pm} + \tau_{\pm\pm} + \gamma_{\pm} $ 

（四） $ CC(k_{U_{+}}) = \sigma_{\pm} - \tau_{\pm\mp} + (-1)^{n+1}\gamma_{\mp} $ 

（五） $ CC(k_{Z_0}) = \sigma_0 + \tau_{+-} + \tau_{-+} + (-1)^n(\gamma_+ + \gamma_-) $ 

（六） $ CC(k_{U_0}) = \sigma_0 - (\tau_{++} + \tau_{--}) + \gamma_0 $ 

（七） $ CC(k_{S_{+}}) = \tau_{\pm+} + \tau_{\pm-} - \gamma_{0} + ((-1)^{n} - 1)\gamma_{\mp} $ 

让我们检查一下例子（iii）。我们有：

\[
\mathrm{S S}(k_{Z_{\pm}})\subset\mathrm{s u p p}(\sigma_{\pm})\cup\mathrm{s u p p}(\tau_{\pm\pm})\cup\mathrm{s u p p}(\gamma_{\pm})~.
\]

此外， $ k_{Z_{\pm}} $  在  $ \mathrm{supp}(\sigma_{\pm}) $ （分别是  $ \mathrm{supp}(\tau_{\pm\pm}) $  和  $ \mathrm{supp}(\gamma_{\pm}) $  ）的通用点上与  $ k_{X} $  （分别是  $ k_{S_{\pm}} $  和  $ k_{\{0\}} $  ）是微局部同构的。

其他例子也可以类似处理。

\subsection{微局部指数公式}
令  $X$  为流形，并令  $F \in \mathrm{Ob}(\mathbf{D}_{\mathbf{R}-c}^b(X))$  。我们将通过交集公式计算 $\chi(X; F)$ （和 $\chi(F)(x)$ 等）。回想一下公式 (9.3.11) 和 (9.3.12) 的态射  $\alpha$  和  $\beta$ 。

\[
\alpha:H_{S}^{0}(T^{\ast}X;\omega_{T^{\ast}X})\to H_{\pi(S)}^{0}(X;\omega_{X})
\]

（S 是闭集， $ \pi $  在 S 上是正确的）。

\[
\beta:H_{\mathcal{A}}^{0}(T^{\ast}X;\pi^{-1}\omega_{X})\to H_{\mathcal{A}\cap X}^{0}(X;\omega_{X})
\]

（A 是封闭且圆锥曲线）。

\statementlabel{命题 9.5.1}令  $ \sigma $  为  $ \pi $  的连续部分。我们在 $ H_{\mathrm{supp}(F)}^0(X; \omega_X) $ 中有：

\[
\begin{aligned}{\mathbf{C}(F)}&{{}=\alpha([\sigma]\cap\mathbf{C C}(F))}\\ {}&{{}}\\ {}&{{}=\beta(\mathbf{C C}(F))~.}\\ \end{aligned}
\]

\statementlabel{证明}根据命题 9.3.6 足以证明  $ \mathrm{C}(F) = \beta(\mathrm{CC}(F)) $  。这是从下面的交换图 (9.5.1) 得出的，我们在其中设置  $ \varLambda = \mathrm{SS}(F) $  。

\begin{center}
\includegraphics[width=0.62\textwidth]{流行上的层_Chn-assets/img_in_image_box_183_942_927_1452.png}
\end{center}

\begin{center}图表 (9.5.1)\end{center}

\statementlabel{推论 9.5.2}假设 F 有紧支撑。然后：

\[
\chi(X;F)=\#\left(\left[\sigma\right]\cap\mathrm{C C}(F)\right).
\]

\statementlabel{证明}应用命题9.5.1和公式(9.1.3)。 ⑨

请注意，该推论也可以通过使用映射  $ a: X \to \{\text{pt}\} $  从 (9.1.14) 推导出来。

我们将对最后一个结果进行概括。令  $I$  为  $\mathbb{R}$  的开区间，让  $\varphi: X \to I$  为实解析函数，让  $\sigma_\varphi$  表示部分  $x \mapsto \mathrm{d}\varphi(x)$  。让 $A_\varphi = \sigma_\varphi(X) = \{(x; \mathrm{d}\varphi(x)); x \in X\}$  。因此， $\varphi$  与一个循环相关联（参见定义 9.3.5）：

\[
[\sigma_{\varphi}]\in H_{\Lambda_{\varphi}}^{0}(T^{*}X;\pi^{\prime}k_{X})\enspace.
\]

\statementlabel{定理 9.5.3}让 $ F \in \mathrm{Ob}(\mathbf{D}_{R-c}^b(X)) $  。假设：

\[
\{x\in\mathbf{s u p p}(F);\varphi(x)\leqslant t\}\qquad{i s~c o m p a c t~f o r~a l l~t\in I~},
\]

\[
\mathrm{SS}(F)\cap A_{\varphi}\qquad\mathrm{is~compact}.
\]

那么对于所有  $ j \in \mathbb{Z} $  ，空间  $ H^j(X; F) $  都是有限维的并且：

\[
\chi(X;F)=\#([\sigma_{\phi}]\cap\mathrm{C C}(F))~.
\]

\statementlabel{证明}让  $\Omega_t = \{x \in X; \varphi(x) < t\}$  ，让  $j_t$  表示开放嵌入  $\Omega_t \hookrightarrow X$  并设置  $F_t = R j_{t*} j_t^{-1} F$  。应用推论 5.4.19，我们得到  $R\Gamma(X; F) \xrightarrow{\sim} R\Gamma(X; F_t)$  对于  $t \gg 0$  。因此，根据推论 9.5.2，空间  $H^j(X; F)$  是有限维的，而  $\chi(X; F) = \chi(X; F_t) = \#([\sigma_\varphi] \cap \mathrm{CC}(F_t))$  是有限维的。由于  $\mathrm{SS}(F) \cap \pi^{-1}(\Omega_t) = \mathrm{SS}(F_t) \cap \pi^{-1}(\Omega_t)$  和  $\mathrm{SS}(F) \cap \Lambda_\varphi \subset \pi^{-1}(\Omega_t)$  为  $t \gg 0$  ，根据假设 (9.5.5)，仍需证明：

\[
\mathrm{S S}(F_{t})\cap\varLambda_{\varphi}\subset\pi^{-1}(\varOmega_{t})\qquad\mathrm{f o r}\qquad t\gg0~.
\]

让 $ K = \pi(\mathrm{SS}(F) \cap \mathcal{A}_\varphi) $  。那么 $ x \in \mathrm{supp}(F) \setminus K $  意味着 $ \mathrm{d}\varphi(x) \neq 0 $ ，我们得到命题9.5.8：

\[
\mathbf{S S}(F_{t})\subset\mathbf{S S}(F)+\mathbb{R}_{\leqslant0}A_{\varphi}~,\qquad\mathrm{f o r}\qquad t\gg0~.
\]

该包含与（9.5.5）一起意味着（9.5.6）。 ⑨

\statementlabel{推论 9.5.4}让 $ F \in \mathrm{Ob}(\mathbb{D}_{R-c}^b(X)) $  。假设 (9.5.4) 并且：

\[
\mathrm{SS}(F)^{a}\cap\varLambda_{\varphi}\qquad\mathrm{is~compact}.
\]

那么对于所有  $j$  ，空间  $H_c^j(X; F)$  都是有限维的并且：

\[
\begin{aligned}{\chi_{c}(X;F)}&{{}=\#([\sigma_{\varphi}]\cap\operatorname{C C}(F)^{a})}\\ {}&{{}=\#([\sigma_{-\varphi}]\cap\operatorname{C C}(F))~.}\\ \end{aligned}
\]

\statementlabel{证明}我们有 $ \operatorname{Hom}(R\Gamma_c(X;F),k) \simeq \operatorname{Hom}(F,\omega_X) \simeq R\Gamma(X;\mathrm{D}_X F) $ 。因此  $ \chi_c(X;F) = \chi(X;\mathrm{D}_X F) $  ，结果来自命题 9.4.4 和定理 9.5.3。   $ \square $ 

\statementlabel{备注 9.5.5}令 $\varphi: X \to \mathbb{R}$  为实解析函数，并假设 $\varphi(x_\circ) = 0$  、 $\mathrm{d}\varphi(x_\circ) = 0$  和 $\varphi$  在 $x_\circ$  处的Hessian 矩阵是正定的。 （这意味着在  $x_\circ$  邻域的局部图中，可以写成  $\varphi(x) = \sum_{j=1}^{n} x_j^2$  ）。）应用微局域 Bertini-Sard 定理（命题 8.3.12），可以得到  $x_\circ$  是  $\mathrm{SS}(F) \cap \mathcal{A}_\varphi$  的孤立点。那么根据定理9.5.3我们发现：

\[
\chi(F)(x_{\circ})\;=\;\#\left(\left[\sigma_{\varphi}\right]\cap\mathrm{C C}(F)\right)_{x_{\circ}}\;.
\]

类似地，根据推论 9.5.4：

\[
\chi_{c}(F)(x_{\circ})\;=\;\#\big(\left[\sigma_{-\varphi}\right]\cap\mathrm{C C}(F)\big)_{\mathbf{x}_{\circ}}\;.
\]

在这里和续集中，  $ \#(c)_x $  ，对于  $ c \in H_S^0(X; \omega_X) $  和闭子集  $ S $  的孤立点  $ x $  ，意味着  $ c $  通过映射的图像：  $ H_S^0(X; \omega) \to H_{\{x\}}^0(X; \omega) \to k $  。

事实上，我们现在将证明一个更普遍的结果。

\statementlabel{定理 9.5.6}令 $ x_o \in X $ 、 $ F \in \mathrm{Ob}(\mathbb{D}_{R-c}^b(X)) $  和 $ \varphi $  成为 $ X $  上的实分析函数。假设：

\[
{\cal A}_{\varphi}\cap\mathrm{S S}(F)\subset\left\{\mathrm{d}\varphi(x_{\circ})\right\}\;.
\]

然后我们有：

\[
\chi(R\varGamma_{\{\varphi\geqslant\varphi(x_{\circ})\}}(F))(x_{\circ})=\mathrm{\normalfont\#}([\sigma_{\varphi}]\cap\mathrm{C C}(F))\enspace.
\]

\statementlabel{证明}我们可以假设  $ X = \mathbb{R}^n $  、  $ x_\circ = 0 $  、  $ \varphi(x_\circ) = 0 $  。设置 $ A = \mathrm{SS}(F) $ ， $ \psi(x) = \sum_{j=1}^n x_j^2 $ 。如果  $ \theta $  是  $ \mathbb{R} $  上的实函数，那么在不存在混淆风险的情况下，让我们写成  $ \{\theta > 0\} $  （例如）来表示集合  $ \pi^{-1}(\{x \in X; \theta(x) > 0\}) $  （正如我们在示例 9.4.6、9.4.7 中所做的那样）。根据微局域 Bertini-Sard 定理（命题 8.3.12），存在  $ \delta_\circ > 0 $  使得：

\[
A\cap A_{\psi}\cap\{0<|x|\leqslant\delta_{\circ}\}=\varnothing,
\]

\[
(\hat{A}+\mathop{T_{\{\phi=0\}}^{*}}X)\cap\varLambda_{\psi}\cap\{0<|x|\leqslant\delta_{\circ}\}=\emptyset.
\]

（回想一下推论 8.3.18， $ A + T_{\{\varphi=0\}}^{*} $  是各向同性的。）此外，存在  $ \delta_1 > 0 $  使得：

\[
k\geqslant0,\qquad0<|x|\leqslant\delta_{1},\qquad\varphi(x)>0\Rightarrow(x;k\mathsf{d}\psi(x)+\mathsf{d}\varphi(x))\notin A.
\]

事实上，如果不满足(9.5.13)，则存在一条实解析曲线 $ t \mapsto (x(t), a(t), b(t)) \in X \times \mathbb{R} \times \mathbb{R} $ ，使得 $ a(t) \geq 0 $ 、 $ b(t) > 0 $ 、 $ x(0) = 0 $ 和 $ \varphi(x(t)) > 0 $ 对于 $ t \neq 0 $ 和 $ (x(t); a(t)\,\mathrm{d}\psi(x(t)) + b(t)\,\mathrm{d}\varphi(x(t))) \in \varLambda $ 。由于 $ \varLambda $  是各向同性的，这意味着：

\[
a(t)\frac{\mathrm{d}}{\mathrm{d}t}\psi(x(t))+b(t)\frac{\mathrm{d}}{\mathrm{d}t}\varphi(x(t))\equiv0~.
\]

由于  $ \frac{d}{dt}\psi \geq 0 $  和  $ \frac{d}{dt}\varphi > 0 $  对于  $ 0 < t \ll 1 $  ，这是一个矛盾。这就证明了(9.5.13)。

设置 $ B(\delta) = \{x \in X; |x| < \delta\} $ 。由于 $ RF_{\{\varphi \geq 0\}}(F) $ 是 $ \mathbb{R} $ 可构造的，因此存在 $ \delta > 0 $ 和 $ \delta < \inf(\delta_0, \delta_1) $ ，使得：

\[
R\varGamma(B(\delta);R\varGamma_{\{\varphi\geqslant0\}}(F))\xrightarrow{\sim}(R\varGamma_{\{\varphi\geqslant0\}}(F))_{0}\enspace.
\]

我们修复这样一个  $\delta$  ，我们用  $j$  表示开放嵌入  $B(\delta) \hookrightarrow X$  并设置  $F' = R j_* j^{-1} F$  ，

\[
A^{\prime}=A\cup(A+T_{\partial B(\delta)}^{*}X)\enspace.
\]

请注意：

\[
\mathsf{S S}(F^{\prime})\subset\mathsf{S S}(F)+\mathbb{R}_{\leqslant0}A_{\downarrow}
\]

特别是 $ \mathrm{SS}(F') \subset A' $ 。由于  $ A' $  是各向同性的，根据微局域 Bertini-Sard 定理，存在  $ \varepsilon > 0 $  使得：

\[
A^{\prime}\cap A_{\varphi}\cap\left\{0<|\varphi(x)|\leq\varepsilon\right\}=\varnothing.
\]

应用微局部莫尔斯引理（推论 5.4.19），我们得到：

\[
R\varGamma_{\{\varphi\geq0\}}(X;F^{\prime})\simeq R\varGamma_{c}(X\cap\{\varphi>-\varepsilon\};F^{\prime})\enspace.
\]

截至（9.5.15）， $ \mathrm{SS}(F') \cap \mathcal{A}_{\varphi} \cap \{0 > \varphi > -\varepsilon\} = \emptyset $ 。由此我们得到推论 9.5.4：

\[
\chi_{c}(X\cap\{\varphi>-\varepsilon\};F^{\prime})=\#([\sigma_{\varphi}]\cap\mathrm{C C}(F^{\prime}|_{\{\varphi>-\varepsilon\}}))\enspace.
\]

（(9.5.17) 的右边表示  $ T^*(X \cap \{\varphi > -\varepsilon\} $  中  $ [\sigma_\varphi] $  和  $ \mathrm{CC}(F') $  的交集数）， 这是有意义的，因为 $A_\varphi\cap \mathrm{SS}(F')\cap\{\varphi>-\varepsilon\}$ 是紧的。）

由此我们得到：

\[
\chi(R\varGamma_{\{\varphi\geqslant0\}}(F))(x_{\circ})=~\#([\sigma_{\varphi}]\cap\mathrm{C C}(F^{\prime}|_{\{\varphi>-\varepsilon\}}))~.
\]

仍有待证明：

\[
A_{\varphi}\cap\mathrm{S S}(F^{\prime})\cap\left\{\varphi>-\varepsilon\right\}\cap\left\{|x|=\delta\right\}=\varnothing.
\]

该关系在集合 $\{-\varepsilon < \varphi < 0\}$  上由(9.5.15) 满足，并且在集合 $\{\varphi > 0\}$  上由(9.5.13) 满足。假设对  $\{\varphi = 0\}$  不满意。我们找到  $p \in \mathrm{SS}(F)$  和  $c \geqslant 0$  这样  $p - c \, \mathrm{d}\psi = \mathrm{d}\varphi$  ,  $|\pi(p)| = \delta$  。如果 $c = 0$  这与假设相矛盾。如果 $c > 0$ ，这意味着 $(\mathcal{A} \backslash \mathcal{A}_\varphi) \cap \{\varphi = 0\} \cap \mathcal{A}_\psi \neq \emptyset$ ，这与(9.5.12)相矛盾。至此证明完成。   $\square$ 

示例 9.5.7。令 Y 为维度  $ l $  的闭子流形，令  $ x_0 \in Y $  和  $ \varphi $  为实解析函数，使得  $ \varphi(x_o) = 0 $  以及流形  $ T_Y^* X $  和  $ \Lambda_\varphi $ 

横向相交于  $p = \mathrm{d}\varphi(x_{\circ})$  。根据莫尔斯引理， $Y$  上存在局部坐标系  $(x_1, \ldots, x_l)$  ，使得  $\varphi|_{Y} = \sum_{j=1}^l a_j x_j^2$  和  $a_j \neq 0 \forall j$  。让 $q = \#\{j; a_j < 0\}$  。从 $R\Gamma_{\{\varphi \geq 0\}}(k_Y)_{x_0} \simeq k[-q]$ （参见VII §5）开始，我们通过定理9.5.6得到：

\[
\#([\sigma_{\varphi}]\cap[T_{Y}^{*}X])=(-1)^{q}.
\]

更一般地，令 $ F \in \mathrm{Ob}(\mathbb{D}_{\mathbb{R}-c}^b(X)) $  和 $ \phi: X \to \mathbb{R} $  为实分析函数。设置  $ A = \mathrm{SS}(F) $  ，假设 (9.5.4) 并且：

(9.5.19)  $ A_{\phi} \cap A $  是有限的，包含在  $ A_{\text{reg}} $  中，并且交集是横向的。

那么很自然地将 $\phi$ 称为“关于 $A$ 的摩尔斯函数”。事实上，如果 $F = k_X$ ，人们就恢复了莫尔斯函数的经典概念。让  $A_\phi \cap A = \{p_1, \ldots, p_N\}$  并假设：

(9.5.20) F 是纯的，且具有移位  $ d_{i} $  和重数  $ m_{i}  p_{i}, i = 1, \ldots, N $ 

根据定理 9.5.6 和纯层的定义我们推导出：

\[
\#\left(\left[\sigma_{\phi}\right]\cap\mathrm{C C}(F)\right)_{p_{i}}=(-1)^{d_{i}^{\prime}}\cdot m_{i}
\]

与：

\[
d_{i}^{\prime}=d_{i}-\tfrac{1}{2}\dim X-\tfrac{1}{2}\tau(\lambda_{\circ}(p_{i}),\lambda_{\wedge}(p_{i}),\lambda_{\phi}(p_{i}))~.
\]

利用定理 9.5.3，我们得到：

\[
\chi(X;F)=\sum_{i=1}^{N}(-1)^{d_{i}^{\prime}}\cdot m_{i}\ .
\]

请注意，最后一个结果也可以使用命题 5.4.20 获得，甚至可以细化（用“莫尔斯不等式”代替这个等式）。

\statementlabel{备注 9.5.8}由于  $ \mathcal{L}_X $  是  $ \mathcal{D}\mathcal{S}_n \otimes or_{T \cdot X/X} $  的子层，因此拉格朗日循环可以被视为  $ or_{T \cdot X/X} $  值的子解析循环。

我们将用这种语言描述拉格朗日循环。

对于  $n$  维流形  $M$  和  $X$  上的无处消失的实数  $n$  -形式  $v$  ，我们用  $\mathrm{sgn}(v)$  表示  $v$  给出的方向。因此  $\mathrm{sgn}(v)$  被视为  $\sigma_{tX}$  的一部分。现在，令 $(x_1, \ldots, x_n)$  为 $X$  的坐标系，并让 $(x; \xi)$  为 $T^*X$  的关联坐标系。我们通过  $\mathrm{sgn}(\mathrm{d}\xi) \leftrightarrow \mathrm{sgn} \, \mathrm{dx}$  识别  $\sigma_{tT^*X/X}$  和  $\sigma_{tX}$  ，其中  $\mathrm{d}\xi = \mathrm{d}\xi_1 \wedge \cdots \wedge \mathrm{d}\xi_n$  和  $\mathrm{dx} = \mathrm{dx}_1 \wedge \cdots \wedge \mathrm{dx}_n$  。然后，通过适当的识别，我们有：

\[
\left\{\begin{aligned}{}&{{}[T_{\{x_{1}=\cdots=x_{l}=0\}}^{*}(X)]\qquad{\mathrm{i s~t h e~c y c l e}}}\\ {}&{{}\left\{(x;\xi);x_{1}=\cdots=x_{l}=\xi_{l+1}=\cdots=\xi_{n}=0\right\}\qquad{\mathrm{w i t h~t h e~o r i e n t a t i o n}}}\\ {}&{{}(-1)^{l}\operatorname{s g n}(\mathbf{d}\xi_{1}\wedge\cdots\wedge\mathbf{d}\xi_{l}\wedge\mathbf{d}x_{l+1}\wedge\cdots\wedge\mathbf{d}x_{n})\otimes\operatorname{s g n}(\mathbf{d}\xi)\enspace.}\\ \end{aligned}\right.
\]

对于 $X$  上的函数 $\varphi$ ， $[\sigma_\varphi]$  是循环 $A_\varphi$ ，其方向由 $X$  通过投影 $A_\varphi \simeq X$  诱导。

我们将  $T^*X$  中两个  $n$  循环的交集定义如下。令  $Z_1$  和  $Z_2$  为  $n$  的  $T^*X$  横向相交于  $p\in T^*X$  的维子流形。令  $\theta_i (i=1,2)$  为  $n$  -形式，使得  $\theta_1|_{Z_2}=0$  和  $\mathrm{sgn}(\theta_1\wedge\theta_2)=\mathrm{sgn}(\mathrm{d}\xi\wedge\mathrm{d}x)=(-1)^{n(n-1)/2}\mathrm{sgn}((\mathrm{d}\alpha_X)^n)$  。如果  $\alpha_i$  是  $n$  -周期  $[Z_i]$  且方向为  $\mathrm{sgn}(\theta_i|_{Z_i})$  ，则  $\#(\alpha_1\cap\alpha_2)=1$  。根据这一约定，推论 9.5.2、定理 9.5.3、推论 9.5.4 以及 (9.5.8) 和 (9.5.9) 仍然成立。

\subsection{莱弗谢茨不动点公式}
考虑流形的两个态射：

\[
Y\underset{g}{\overset{f}{\rightleftarrows}}X
\]

并设置  $ h = (f, g) $  :  $ Y \to X \times X $  。让  $ F \in \mathrm{Ob}(\mathbb{D}_{R-c}^{b}(X)) $  并假设给出：

\[
\varphi\in\mathrm{Hom}(f^{-1}F,g^!F).
\]

我们设定：

\[
S=\operatorname{s u p p}(F),\qquad T=f^{-1}(S)\cap g^{-1}(S),
\]

\[
Z=T\cap h^{-1}(\varDelta_{X})=\left\{y\in Y;f(y)=g(y)\in S\right\}\;,
\]

我们假设：

T 是紧凑的。

请注意， $\varphi$  的支持包含在 $T$  中。自然态射：  $\mathrm{id} \to R f_* f^{-1}$  和  $R g_! g^1 \to \mathrm{id}$  ，定义态射：

\[
R\varGamma(X;F)\to R\varGamma_{f^{-1}(\mathbb{S})}(Y;f^{-1}F)~,\qquad R\varGamma_{T}(Y;g^{!}F)\to R\varGamma(X,F)~.
\]

与  $\varphi$  结合，我们得到一个态射，我们仍然用  $\varphi$  来简称：

\[
\varphi:R\varGamma(X;F)\to R\varGamma(X;F)\enspace.
\]

本节的目的是计算  $ \mathrm{tr}(\varphi) $  ，即  $ \varphi $  的迹，其中：

\[
\mathrm{t r}(\varphi)=\sum_{j}(-1)^{j}\mathrm{t r}(H^{j}(\varphi))\ ,
\]

 $ H^j(\varphi) $  表示从 (9.6.4) 推导出的  $ H^j(X; F) $  的自同态。由于 $ H^j(\varphi) $  分解为 $ H^j(X; F) \to H^j_{g(T)}(X; F) \to H^j(X; F) $  和 $ H^j_{g(T)}(X; F) $  是有限的-

维度， $ \operatorname{tr}(H^j(\varphi)) $  有一个含义（参见练习 I.33）。让  $ \delta $  表示对角线嵌入  $ X \hookrightarrow X \times X $  ，并考虑态射链（这里我们不需要假设 (9.6.3)）：

\[
\begin{aligned}{R\mathcal{H}o m(F,F)\xleftarrow{\sim}}&{{}\delta^{!}(F\boxtimes\mathbf{D}F)}\\ {\longrightarrow}&{{}\delta^{!}R h_{*}h^{-1}(F\boxtimes\mathbf{D}F)}\\ {\xrightarrow{\sim}}&{{}R\varGamma_{\varDelta}R h_{*}R\varGamma_{T}(f^{-1}F\otimes\mathbf{D}g^{!}F)}\\ {\xrightarrow{\sim}}&{{}R h_{*}R\varGamma_{Z}(f^{-1}F\otimes\mathbf{D}g^{!}F)}\\ {\xrightarrow[\operatorname{t r}]{\quad}}&{{}R h_{*}R\varGamma_{Z}(g^{!}F\otimes\mathbf{D}g^{!}F)}\\ {\xrightarrow{\quad}}&{{}R h_{*}R\varGamma_{Z}(\omega_{Y})\enspace.}\\ \end{aligned}
\]

这导致：

\[
\mathrm{H o m}(F,F)\to H_{Z}^{0}(Y;\omega_{Y})~.
\]

\statementlabel{定义 9.6.1}$H_{Z}^{0}(Y;\omega_{Y})$  中 $\mathrm{id}_{F}$  的图像称为 $\varphi$  的特征类，记为 $C(\varphi)$  。

\statementlabel{命题 9.6.2（Lefschetz 不动点定理）}令 $f, g, \varphi$  如上，并假设(9.6.3) 和 $\mathrm{supp}(F)$  是紧致的。然后：

\[
\mathrm{tr}(\boldsymbol{\varphi})=\int_{\mathbf{Y}}\mathrm{C}(\boldsymbol{\varphi})~.
\]

\statementlabel{证明}考虑下面的交换图（9.6.6）。由于 $\varphi$  分解为 $R\Gamma(X; F) \xrightarrow{\alpha} R\Gamma_T(Y; g^1F) \xrightarrow{\beta} R\Gamma(X; F)$  ，所以 $\varphi$  的迹等于 $\beta \circ \alpha$  的迹，并且它是 $id_F$  在 $k$  中通过右侧列的态射链的图像。这样就完成了证明。   $\square$ 

请注意，上面的命题在  $ \operatorname{supp}(F) $  是紧的条件下也成立（参见练习 IX.9）。

现在让  $y \in Z$  并假设  $y$  是  $Z$  的一个孤立点。设置 $Z' = Z \setminus \{y\}$ 。然后：

\[
H_{Z}^{0}(Y;\omega_{Y})\simeq H_{\{y\}}^{0}(Y;\omega_{Y})\oplus H_{Z^{\prime}}^{0}(Y;\omega_{Y})\enspace.
\]

\statementlabel{定义 9.6.3}在上述情况中， $C_{y}(\varphi)$  表示 $C(\varphi)$  的态射的像：

\[
H_{Z}^{0}(Y;\omega_{Y})\to H_{\{y\}}^{0}(Y;\omega_{Y})\simeq k\enspace.
\]

另外，如果 $ y \notin Z $ ，我们设置 $ C_y(\varphi) = 0 $ 。

当然，如果Z是有限的，那么：

\[
\mathrm{t r}(\varphi)=\sum_{y\in\mathbb{Z}}\mathrm{C}_{y}(\varphi)\enspace.
\]

\begin{center}
\includegraphics[width=0.70\textwidth]{流行上的层_Chn-assets/img_in_image_box_149_162_990_1226.png}
\end{center}

注释 9.6.4。从现在开始，我们将假设 $Y = X$ ， $g = \mathrm{id}_{X}$ 。因此  $Z$  是  $\mathrm{supp}(F)$  和  $f$  不动点集的交集。

如果  $x$  是  $f$  的不动点，我们用  $\varphi_{x}$  表示作为复合态射获得的  $F_{x}$  的自同态：

\[
F_{x}\simeq(f^{-1}F)_{x}\mathop{\longrightarrow}_{\varphi}F_{x}.
\]

当  $x$  是  $f$  的不动点且  $x$  是  $f^{-1}(x) \cap T = f^{-1}(x) \cap \mathrm{supp}(F) \cap f^{-1}(\mathrm{supp}(F))$  的孤立点时，我们用  $\varGamma_{\{x\}}(\varphi)$  表示作为复合态射获得的  $R\varGamma_{\{x\}}(X; F)$  的自同态：

\[
R\varGamma_{\{x\}}(X;F)\longrightarrow R\varGamma_{\{f^{-1}(x)\}}(X;f^{-1}F)\underset{\varphi}{\longrightarrow}R\varGamma_{\{f^{-1}(x)\}\cap T}(X;F)\longrightarrow R\varGamma_{\{x\}}(X;F)\enspace.
\]

应注意  $C_{x}(\varphi)$  、  $\mathrm{tr}(\varphi_{x})$  和  $\mathrm{tr}(\varGamma_{\{x\}}(\varphi))$  一般而言是不同的。

示例 9.6.5。令  $f \to \mathbb{R} \to \mathbb{R}$  为仅具有两个不动点  $t_0$  和  $t_1$  、  $t_0 < t_1$  的解析微分同胚，以便  $f$  引发区间  $[t_0, t_1]$  的同胚。让  $F = k_{[t_0, t_1]}$  和  $\varphi$  表示规范同构  $f^{-1}k_{[t_0, t_1]}\approx k_{[t_0, t_1]}$  。然后 $\mathrm{tr}(\varphi) = \chi(\mathbb{R}; F) = 1 = C_{t_0}(\varphi) + C_{t_1}(\varphi)$ 。另一方面， $\varphi_{t_i}$  充当 $F_{t_i}$  ( $i = 0, 1$ ) 上的身份。因此  $\mathrm{tr}(\varphi_{t_i}) = 1$  和  $C_{t_i}(\varphi) \neq \mathrm{tr}(\varphi_{t_i})$  代表  $i = 0$  或  $i = 1$  。  $\Gamma_{\{t_i\}}(\varphi)$  也有类似的论点。正如我们稍后将看到的， $C_{t_i}(\varphi) = 0$  或  $1$  根据  $f(t) - t$  在  $t_i$  处增加或减少。

示例 9.6.6。令  $X$  为  $n$  维紧流形，令  $F = k_X$  并令  $\varphi$  为规范态射  $f^{-1}k_X \cong k_X$  。考虑态射链：

\[
\begin{aligned}{k_{X}}&{{}\simeq R\mathcal{H o m}(k_{X},k_{X})}\\ {}&{{}\simeq\delta^{!}(k_{X}\boxtimes\omega_{X})}\\ {}&{{}\to\delta^{!}h_{*}\omega_{X}.}\\ \end{aligned}
\]

让  $ \Gamma_f = \{(f(x), x); x \in X\} $  和  $ j: \Gamma_f \hookrightarrow X \times X $  嵌入  $ \Gamma_f $  。那么通过它的构造， $ \mathrm{C}(\varphi) $  是 $ 1 \in \Gamma(X; k_X) $  的态射的图像：

\[
\begin{aligned}{\varGamma(X;k_{X})}&{{}\stackrel{\alpha}{\longrightarrow}H_{\varDelta}^{0}(X\times X;k_{X}\boxtimes\omega_{X})}\\ {}&{{}\stackrel{\beta}{\longrightarrow}H_{j^{-1}(\varDelta)}^{0}(\varGamma_{f};j^{-1}(k_{X}\boxtimes\omega_{X}))}\\ {}&{{}\stackrel{\sim}{\longrightarrow}H_{j^{-1}(\varDelta)}^{0}(\varGamma_{f};\omega_{\varGamma_{f}})}\\ {}&{{}\stackrel{\sim}{\longrightarrow}H_{j^{-1}(\varDelta)}^{0}(X\times X;\omega_{X\times X}).}\\ \end{aligned}
\]

令 $ [\Gamma_f] \in H^0_{\Gamma_f}(X \times X; \mathcal{L} \mathcal{S}_n(k_X \boxtimes \sigma \iota_X)) $  和 $ [\varDelta] \in H^0(X \times X; \mathcal{L} \mathcal{S}_n(k_X \boxtimes \sigma \iota_X)) $  分别为与 $ \Gamma_f $  和 $ \varDelta $  相关的循环。然后，识别 $ (o \iota_X \boxtimes k_X)|_{\varDelta} $ 和 $ (k_X \boxtimes o \iota_X)|_{\varDelta} $ ，我们有 $ \alpha(1) = [\varDelta] $ 。然后我们得到：

\[
\beta([\varDelta])=[{\varGamma}_{f}]\cap[\varDelta],
\]

这给出：

\[
{\bf C}(\varphi)=[{\it\Gamma}_{f}]\cap[{\it\Delta}]\enspace.
\]

由此我们得到著名的Lefschetz不动点公式：

\[
\sum_{j}(-1)^{j}\mathrm{t r}(H^{j}(\varphi))=\#([\varGamma_{f}]\cap[\varDelta])~,
\]

（参见备注 9.0）。

换句话说， $\phi$  的迹被计算为 $f$  的图形与对角线的交点数。

\statementlabel{备注 9.6.7}令  $x$  为  $Z$  的孤立点。那么  $C_x(\varphi)$  是以下意义上的局部不变量。令  $F$  和  $F'$  属于  $\mathbf{D}_{\mathbf{R}-c}^b(X)$  、  $\varphi: f^{-1}F \to F$  、  $\varphi': f^{-1}F' \to F'$  并假设存在  $x$  的邻域  $W$  和同构  $\psi: F|_{w} \to F'|_{w}$ 

如下图所示：

\[
\begin{array}{c}f^{-1}F|_{\mathcal{W}\cap f^{-1}(\mathcal{W})}\xrightarrow{\quad\varphi\quad}F|_{\mathcal{W}\cap f^{-1}(\mathcal{W})}\\f^{-1}(\psi)\bigg\downarrow\quad\bigg\downarrow\quad\bigg\downarrow\quad\bigg\downarrow\quad\bigg\downarrow\quad\bigg\downarrow\\f^{-1}F^{\prime}|_{\mathcal{W}\cap f^{-1}(\mathcal{W})}\xrightarrow{\quad\varphi^{\prime}\quad}F^{\prime}|_{\mathcal{W}\cap f^{-1}(\mathcal{W})}\quad.\end{array}
\]

然后 $ C_{x}(\varphi) = C_{x}(\varphi') $ 。证据是立竿见影的。

我们现在将解释在某些情况下如何计算  $C(\varphi)$  。我们首先证明  $C(\varphi)$  的同伦不变性。

设  $ I = [0,1] $  、  $ f: X \times I \to X $  流形态射的限制  $ X \times \mathbb{R} \to X $  、  $ p: X \times I \to X $  投影。令 $ F \in \mathrm{Ob}(\mathbb{D}_{R-c}^b(X)) $  和 $ \varphi: f^{-1}F \to p^{-1}F $  为态射。对于  $ t \in I $  ，令  $ i_t: X \hookrightarrow X \times I $  为注入  $ x \mapsto (x,t) $  。设置 $ f_t = f \circ i_t $ 。   $ \varphi_t = i_t^{-1}(\varphi): f_t^{-1}F \to F $  ，最后设置：

\[
\tilde{Z}=\left\{(x,t)\in\mathrm{supp}(F)\times I;f(x,t)=x\right\}\;.
\]

\statementlabel{命题 9.6.8}在前面的情况下，假设  $ \tilde{Z} $  是紧凑的。那么 $ H^0_{p(\tilde{Z})}(X;\omega_X) $  中 $ C(\varphi_t) $  的图像不依赖于 $ t $  。

\statementlabel{证明}考虑交换图：

\begin{center}
\includegraphics[width=0.65\textwidth]{流行上的层_Chn-assets/img_in_image_box_140_910_919_1520.png}
\end{center}

其中  $ i_t^\# $  是 id  $ \to i_{t*}\circ i_t^{-1} $  引发的态射。那么 $ C(\varphi_t) $  就是 $ i_t^\# $  的 $ c\in H_2^0(X\times I;p^{-1}\omega_X) $  的图像。由于态射

\[
H_{p(\tilde{Z})}^{0}(X;\omega_{X})\xrightarrow{p^{\#}}H_{p^{-1}p(\tilde{Z})}^{0}(X\times I;p^{-1}\omega_{X})\xrightarrow{i_{t}^{\#}}H_{p(\tilde{Z})}^{0}(X;\omega_{X})
\]

是同构， $ i_{t}^{\#} \circ p^{\#} = \mathrm{id}, i_{t}^{\#} $  不依赖于 t。 □

现在设 x 为 Z 的孤立点，设 V 为 x 的局部闭亚解析邻域，满足：

\[
V\cap f^{-1}(V)\text{is closed in}V\text{and open in}f^{-1}(V).
\]

让 $ \varGamma_{V}(\varphi) $  表示从 $ f^{-1}R\varGamma_{V}(F) $  到 $ R\varGamma_{V}(F) $  作为复合获得的态射：

\[
\begin{aligned}{f^{-1}R\varGamma_{V}(F)}&{{}\longrightarrow R\varGamma_{f^{-1}(V)}(f^{-1}F)\mathop{\longrightarrow}_{\varphi}R\varGamma_{f^{-1}(V)}(F)}\\ {}&{{}\longrightarrow R\varGamma_{f^{-1}(V)\cap V}(F)\longrightarrow R\varGamma_{V}(F)\enspace.}\\ \end{aligned}
\]

\statementlabel{命题 9.6.9}假设(9.6.8)。然后 $C_x(\varphi) = C_x(\Gamma_V(\varphi))$ 。特别是如果  $V$  相对紧凑且  $Z \cap \bar{V} = \{x\}$  、  $C_x(\varphi) = \mathrm{tr}(\Gamma_V(\varphi))$  。

\statementlabel{证明}这是根据局部不变性（备注 9.6.7）和 Lefschetz 定理（命题 9.6.2）得出的。   $ \Box $ 

我们也有以下类似的结果。

设 V 为 x 的局部闭亚解析邻域，满足：

\[
V\cap f^{-1}(V)\text{is open in}V\text{and closed in}f^{-1}(V).
\]

令  $ \varphi_{V} $  为从  $ f^{-1}F_{V} $  到  $ F_{V} $  的复合态射：

\[
f^{-1}F_{V}\longrightarrow(f^{-1}F)_{f^{-1}(V)}\xrightarrow[\varphi]{}F_{f^{-1}(V)}\longrightarrow F_{V\cap f^{-1}(V)}\longrightarrow F_{V}\enspace.
\]

\statementlabel{命题 9.6.10}假设(9.6.9)。然后 $C_x(\varphi) = C_x(\varphi_V)$ 。特别是，如果  $V$  相对紧凑，而  $Z \cap \bar{V} = \{x\}$  则相对紧凑，则  $C_x(\varphi) = \mathrm{tr}(\varphi_V)$  。

接下来，我们计算  $\varphi$  在  $f$  的固定点  $x$  上的特化迹。我们将简写为  $v_{x}F$  而不是  $v_{\{x\}}(F)$  。我们用  $v_{x}\varphi$  表示复合态射（参见命题 4.2.5）：

\[
\nu_{x}\varphi:(T_{x}f)^{-1}\nu_{x}F\longrightarrow\nu_{x}f^{-1}F\xrightarrow[\varphi]{}\nu_{x}F.
\]

假设：

 $ \Gamma_f $  和 $ \Delta $  横向相交于 $ x \in X $  。

根据这个假设，0 是映射  $ T_{x}f: T_{x}X \to T_{x}X $  的唯一固定点。

\statementlabel{命题 9.6.11}假设(9.6.11)。然后 $ C_x(\varphi) = C_0(\nu_x\varphi) $ 。

\statementlabel{证明}考虑下面的交换图（9.6.12）就足够了，其中我们使用命题8.4.13的同构 $ \mathrm{D}\nu_x F \simeq \nu_x DF $ 。   $ \square $ 

\begin{center}
\includegraphics[width=0.86\textwidth]{流行上的层_Chn-assets/img_in_image_box_67_166_1102_790.png}
\end{center}

图(9.6.12)

根据这个命题，为了在满足(9.6.11)时计算 $ C_x(\varphi) $ ，只需考虑以下情况就足够了。

令 $V$  为实数有限维向量空间， $u: V \to V$  为线性自同态， $F$  为 $\mathbf{D}_{\mathbf{R}-c}^b(V)$  的圆锥对象，并令 $\varphi \in \mathrm{Hom}(u^{-1}F, F)$  。我们假设：

1 不是 u 的特征值

令  $V^C$  表示复数空间  $V\otimes_R\mathbb{C}$  ，对于  $\lambda\in\mathbb{C}$  ，用  $V_\lambda^C$  表示  $V^C$  的广义特征空间和广义特征值  $\lambda$  ，即  $x\in V_\lambda^C$  当且仅当存在  $m>0$  和  $(u-\lambda)^m(x)=0$  。如果满足以下条件，我们可以说  $V$  的线性子空间  $V_s$  是一个“收缩空间”：

\[
\begin{cases}{u(V_{s})\subset V_{s}\;,}\\ {u|_{V_{s}}\quad{h a s~n o~e i g e n v a l u e s~}\lambda\in\mathbb{R}{~w i t h~}\lambda>1\;,}\\ {u|_{V/V_{s}}\quad{h a s~n o~e i g e n v a l u e s~}\lambda\in\mathbb{R}{~w i t h~}0\leqslant\lambda<1\;.}\\ \end{cases}
\]

这相当于：

\[
\left\{\begin{aligned}{}&{{}u(V_{s})\subset V_{s},}\\ {}&{{}\bigoplus_{0\leqslant\lambda<1}V_{\lambda}^{\mathbb{C}}\subset V_{s}\otimes_{\mathbb{R}}\mathbb{C}\subset\bigoplus_{\lambda\notin]1,+\infty[}V_{\lambda}^{\mathbb{C}}.}\\ \end{aligned}\right.
\]

此条件意味着  $ u^{-1}(V_{s}) = V_{s} $  。

类似地，如果满足以下条件，我们可以说线性子空间  $V_{e}$  是一个“扩展空间”：

\[
\begin{cases}{u(V_{e})\subset V_{e},}\\ {u|_{V_{e}}}&{{h a s~n o~e i g e n v a l u e s~}\lambda\in\mathbb{R}{~w i t h~}0\leqslant\lambda<1}\\ {u|_{V/V_{e}}}&{{h a s~n o~e i g e n v a l u e s~}\lambda\in\mathbb{R}{~w i t h~}\lambda>1.}\\ \end{cases},
\]

这相当于：

\[
\left\{\begin{aligned}{}&{{}u(V_{e})\subset V_{e},}\\ {}&{{}\bigoplus_{1<\lambda}V_{\lambda}^{\mathbb{C}}\subset V_{e}\otimes_{\mathbb{R}}\mathbb{C}\subset\bigoplus_{\lambda\notin[0,1]}V_{\lambda}^{\mathbb{C}}.}\\ \end{aligned}\right.
\]

我们将集中研究扩大空间。

让 $V_e$ 成为一个扩展空间。那么 $u|_{V_e}: V_e \to V_e$  是同构。令  $\Gamma_c(\varphi_{V_e})$  为作为组合获得的自同态：

\[
R\Gamma_{c}(V_{e};F|_{V_{e}})\to R\Gamma_{c}(V_{e};u^{-1}F|_{V_{e}})\xrightarrow{\varphi}R\Gamma_{c}(V_{e};F|_{V_{e}})
\]

\statementlabel{命题 9.6.12}令 $u: V \to V$  为没有非平凡不动点的线性映射， $V_e$  为 $V(cf.(9.6.13)).$  的扩展子空间

令  $F$  为  $\mathbf{D}_{\mathbf{R}-c}^{b}(V)$  的圆锥对象，并令  $\varphi \in \mathrm{Hom}(u^{-1}F, F)$  。然后 $C_{0}(\varphi) = \mathrm{tr}(\Gamma_{c}(\varphi_{V_{e}}))$ 。

\statementlabel{证明}(a) 首先假设  $u$  与  $|\lambda| = 1$  没有特征值  $\lambda \in \mathbb{C}$  。考虑分解 $V = V_+ \oplus V_-$ ：

\[
V_{+}=\left(\bigoplus_{|\lambda|>1}V_{\lambda}^{\mathbf{c}}\right)\cap V\;,\qquad V_{-}=\left(\bigoplus_{|\lambda|<1}V_{\lambda}^{\mathbf{c}}\right)\cap V
\]

并在  $V$  上选择一个度量，使得存在常量  $c_1$  和  $c_2$  以及  $0 < c_1 < 1 < c_2$  和  $|u(x)| \geq c_2|x|$  在  $V_+$  上， $|u(x)| \leq c_1|x|$  在  $V_-$  上。

设置 $ Z_{ab} = \{x \in V_+; |x| < a\} \times \{x \in V_-; |x| \leq b\} $ 。然后  $ u^{-1}(Z_{ab}) \cap Z_{ab} $  在  $ Z_{ab} $  中打开并在  $ u^{-1}(Z_{ab}) $  中关闭。应用命题 9.6.9，我们发现：

\[
\mathrm{C}_{0}(\varphi)=\operatorname{t r}(\varGamma(\varphi_{Z_{a b}}))~.
\]

这里  $ \varGamma(\varphi_{Z_{a,b}}) $  是  $ \varphi_{Z_{ab}} $  引起的  $ RF(X; F_{Z_{ab}}) \simeq RF_{c}(Z_{ab}; F|_{Z_{ab}}) $  的自同态。

\statementlabel{引理 9.6.13}对于任何  $b>0$  ，都存在  $a_0>0$  ，使得  $R\Gamma_c(Z_{a,b};F)\to R\Gamma_c(Z_{\infty,b};F)$  是  $a\geq a_0$  的同构。

\statementlabel{证明}设置  $ Y = \{(t, x_{+}) \in \mathbb{R} \times V_{+}; t^2 + |x_{+}|^2 = 1\} $  并通过  $ x_{+} \mapsto \left( \frac{1}{\sqrt{1 + |x_{+}|^2}}, \frac{x_{+}}{\sqrt{1 + |x_{+}|^2}} \right) $  将  $ V_{+} $  嵌入到 Y 中。

那么  $V_{+}$  是  $t > 0$  定义的  $Y$  的开子集。令 $j: V_{+} \times V_{-} \to Y \times V_{-}$  为开放嵌入。那么  $F' = Rj_! F_{\{|x_{-}| \leq b\}}$  是  $w$  -  $\mathbb{R}$  - 可构造的。现在，我们可以使用  $\varphi = t$  应用引理 8.4.7。事实上，该断言相当于说  $R\Gamma(\{|t| \leq \varepsilon\}; F') = 0$  代表  $0 < \varepsilon \ll 1$  。   $\square$ 

命题 9.6.12 的证明结束。令  $q: V \to V_{-}$  为投影。然后，由于  $Rq_{1}F$  是圆锥曲线，因此对于任何  $b > 0$  ：

\[
\begin{aligned}{H_{c}^{k}(V_{+};F)}&{{}\simeq H^{k}(R q_{!}F)_{0}\simeq H^{k}(\{x_{-}\in V_{-};|x_{-}|\leqslant b\};R q_{!}F)}\\ {}&{{}\simeq\varinjlim_{a\to\infty}H_{c}^{k}(Z_{a,b};F)\enspace.}\\ \end{aligned}
\]

然后根据前面的引理，这与  $ a \gg 1 $  的  $ H_c^k(Z_{a,b}; F) $  同构。因此我们得到了一条同构链 $ R\Gamma_c(V_+; F) \simeq R\Gamma_c(Z_{\infty,b}; F) \simeq R\Gamma_c(Z_{a,b}; F) $ 。

这证明了  $ V_e = V_+ $  且不存在特征值  $ \lambda $  和  $ |\lambda| = 1 $  的特殊情况下的结果。

(b) 将  $ u $  替换为  $ tu $  、  $ t \in \mathbb{R} $  、  $ |1 - t| \ll 1 $  ，不会影响命题 9.6.8 的 $ C_0(\varphi) $ （参见下面的备注 9.6.15），也不会影响命题 2.7.5 的 $ \operatorname{tr}(\Gamma_c(\varphi_{V_s})) $ 。因此，我们可以假设所有特征值  $ \lambda $  满足  $ |\lambda| \neq 1 $  。

(c) 设 $ W = \bigoplus_{\lambda > 1} (V_{\lambda}^C \cap V) $  和 $ V_e $  为扩展空间。由(a)足以证明：

\[
\mathrm{t r}(\varGamma_{c}(\varphi_{W}))=\mathrm{t r}(\varGamma_{c}(\varphi_{V_{e}}))~,
\]

因为 (9.6.20) 将特别暗示  $ \mathrm{tr}(\Gamma_c(\varphi_w)) = \mathrm{tr}(\Gamma_c(\varphi_{V_+})) $  。为了证明（9.6.20），让我们用 $ V' $ 表示空间 $ V_e/W $ ，用 $ u' $ 表示由 $ u $ 导出的映射 $ V' \to V' $ ，并用 $ p: V_e \to V_e/W $ 表示投影。将 $ V $ 替换为 $ V' $ ， $ u $ 替换为 $ u' $ ， $ F $ 替换为 $ Rp_t F $ ，就足以证明：

\[
\mathrm{t r}(\varGamma_{c}(\varphi_{V}))=\mathrm{t r}(\varphi_{0})
\]

假设不存在特征值  $ \lambda \in [0, \infty[ $  。

由于 $ R\Gamma_c(V;F) \simeq R\Gamma_{\{0\}}(V;F) $  和 $ R\Gamma(V;F) \simeq F_0 $  ，(9.6.21) 相当于说 $ \varphi $  作用于 $ R\Gamma(V\setminus{0};F) $  的迹为零。由于  $ F $  是圆锥曲线，因此该迹线也是  $ \gamma_* \varphi $  作用于  $ R\gamma_*F $  的迹线，其中  $ \gamma $  是映射  $ V\setminus{0} \to S_V = (V\setminus{0}) / \mathbb{R}^+ $  。根据假设， $ \gamma_* u $ （由 $ u $  导出的 $ S_V $  上的映射）没有不动点。因此，命题 9.6.2 提出了 $ \operatorname{tr}(\gamma_* \varphi) = 0 $ 。   $ \square $ 

同样：

\statementlabel{命题 9.6.14}令 $u: V \to V$  为没有非平凡不动点的线性映射，并令 $V_s$  为收缩空间（参见（9.6.14））。令  $F$  为  $\mathbf{D}_{\mathbf{R}-c}^b(V)$  的圆锥对象，并令  $\varphi \in \mathrm{Hom}(u^{-1}F, F)$  。然后 $C_0(\varphi) = \mathrm{tr}(\Gamma_{V_s}(\varphi))$ 。

由于证明与命题9.6.13相似，我们不再重复。请注意， $ R\Gamma_{V}(V;F) $  和  $ R\Gamma_{c}(V_{e};F) $  不一定是同构的。

\statementlabel{备注 9.6.15}令  $V$  为实数有限维向量空间，并像往常一样用  $p$  表示投影  $V \times \mathbb{R} \to V$  并用  $i_t$  表示嵌入  $V \hookrightarrow V \times \mathbb{R}$  、  $x \mapsto (x, t)$  。令  $u: V \times \mathbb{R} \to V$  为连续映射，相对于  $x \in V$  呈线性，

令  $F$  为  $\mathbf{D}_{\mathbf{R}-c}^b(V)$  的圆锥对象，让  $\varphi$  为从  $u^{-1}F$  到  $p^{-1}F$  的态射。设置  $u_t = i_t \circ u$  并让  $\varphi_t = i_t^{-1} \circ \varphi : u_t^{-1}F \to F$  。那么，如果  $u_t$  没有 1 作为特征值，则  $C_0(\varphi_t)$  不依赖于  $t$  。这是从命题 9.6.8 得出的。

\statementlabel{推论 9.6.16}假设  $X$  是复流形， $F \in \mathrm{Ob}(\mathbb{D}_{\mathbb{C}-c}^b(X))$  。然后假设 (9.6.11)，我们有  $\mathrm{C}_x(\varphi) = \mathrm{tr}(\varphi_x)$  。此外，如果  $0$  不是  $f'(x)$  的特征值（即：  $f$  是  $x$  的局部同构），那么我们还有  $\mathrm{C}_x(\varphi) = \mathrm{tr}(\Gamma_{\{x\}}(\varphi))$  。

\statementlabel{证明}将  $F$  替换为  $\nu_x F$  ，我们可以从一开始就假设  $X$  是复向量空间， $f$  是线性的， $F$  是  $\mathbb{C}^\times$  圆锥曲线，因为  $F$  在  $X$  中的  $\mathbb{C}^\times$  轨道上是局部常数。将  $f$  替换为  $\lambda f$  和  $\lambda \in \mathbb{C}$  、  $|\lambda - 1| \ll 1$  ，我们可以假设  $f$  的任何非零特征值都不是真实的。然后我们采取 $V_e = \{0\}$ 。如果  $0$  不是特征值，我们可以采用  $V_e = T_x X$  然后我们使用  $R\Gamma_c(T_x X; \nu_x F) \simeq R\Gamma_{\{x\}}(X; F)$  这一事实。   $\square$ 

示例 9.6.17。令 $X$  为 $n$  维流形， $F = k_X$  并令 $\varphi: f^{-1}F \to F$  为规范映射。假设  $\Gamma_f$  和  $\Delta$  横向相交于  $x_o \in X$  。然后：

\[
\mathrm{C_{x}}_{\circ}(\varphi)=\operatorname{s g n}(\operatorname*{d e t}(1-f^{\prime}(x_{\circ})))~.
\]

事实上，根据命题 9.6.11，我们可以假设  $X$  是向量空间， $f$  是线性的。采取 $V_e = \bigoplus_{\lambda > 1}(V_\lambda^c \cap V)$ 并应用命题9.6.12。然后 $\det_{V_e}(f|_{V_e}) > 0$ 。根据命题 3.2.3 (iv)， $\varphi_{V_e}$  充当  $R\Gamma_c(V_e; k_X) \simeq k[-\dim V_e]$  上的身份。因此 $C_0(\varphi) = \mathrm{tr}(\Gamma_c(\varphi_{V_e})) = (-1)^{\dim V_e} = \mathrm{sgn}(\det(1 - f'(x_e)))$  。 （最后一个等式来自：  $\det_V(1 - f'(x_e)) = \det_{V_e}(1 - f'(x_e)) \cdot \det_{V/V_e}(1 - f'(x_e))$  和  $\det_{V/V_e}(1 - f'(x_e)) > 0$  、  $\det_{V_e}(f'(x_e) - 1) > 0$  。）

示例 9.6.18。令  $f, F, \varphi$  如例 9.6.5 所示，并假设  $f'(t_0) \neq 1$  。然后：

\[
C_{t_{0}}(\varphi)=\left\{\begin{aligned}{}&{{}0\quad{i f}f^{\prime}(t_{0})>1,}\\ {}&{{}1\quad{i f}0\leqslant f^{\prime}(t_{0})<1.}\\ \end{aligned}\right.
\]

事实上， $C_{t_0}(\varphi) = \mathrm{tr}(\Gamma_c(\varphi_{V_e}))$ ，如果 $f'(t_0) > 1$ ，我们采用 $V_e = T_x X$ ，如果 $0 \leq f'(t_0) < 1$ ，我们采用 $V_e = \{0\}$ 。从  $v_{t_0} F \simeq k_{\{t \geq 0\}}$  开始，我们有了  $R\Gamma_c(T_{\{t_0\}} X; v_{t_0} F) = 0$  和  $R\Gamma_c(\{0\}; v_{t_0} F) \simeq k$  。然后是(9.6.23)。

\subsection{构造函数与拉格朗日圈}
 $X$  上的可构造函数  $\varphi$  是一个  $\mathbb{Z}$  值函数，使得：

对于每个  $ m \in \mathbb{Z} $  ，  $ \varphi^{-1}(m) $  是亚解析的

并且族 $ \{\varphi^{-1}(m)\}_{m \in \mathbb{Z}} $  是局部有限的。

根据第八章的结果，当且仅当存在  $\mu$  分层  $X = \bigsqcup_{\alpha} X_{\alpha}$  使得  $\varphi|_{X_{\alpha}}$  对于每个  $\alpha$  都是常数时， $\mathbb{Z}$  值函数  $\varphi$  才是可构造的。令 $1_{A}$  表示 $X$  的子集 $A$  的特征函数。根据三角剖分定理，当且仅当存在局部有限覆盖  $X = \bigcup_{\alpha} X_{\alpha}$  时，函数  $\varphi$  才是可构造的，其中  $X_{\alpha}$  是紧致的、亚解析的和可收缩的，并且存在整数  $m_{\alpha}$  和  $\varphi = \sum_{\alpha} m_{\alpha}1_{X_{\alpha}}$  。

 $X$ 上的可构造函数集自然地被赋予了代数的结构，我们用 $\mathrm{CF}(X)$ 来表示这个代数。预层  $U \mapsto \mathrm{CF}(U)$  、  $(U \text{ open in } X)$  显然是一个捆（组）。我们用  $\mathcal{C} \mathcal{F}_X$  表示 if ，并将其称为  $X$  上的  $(\mathbb{Z}-\text{valued})$  可构造函数的束。请注意， $\mathcal{C} \mathcal{F}_X$  是软的。

现在让  $F \in \mathrm{Ob}(\mathbf{D}_{\mathbf{R}-c}^{b}(X))$  ，（基环又是特征零的场  $k$  ）。它的局部欧拉-庞加莱指数 $\chi(F)(x) = \sum_{j} (-1)^{j} \dim H^{j}(F)_x$ 显然是一个可构造的函数。此外：

\[
\left\{\begin{aligned}\chi(F\oplus G)&=\chi(F)+\chi(G)\\ \chi(F\otimes G)&=\chi(F)\cdot\chi(G)\end{aligned}\right.
\]

如果 $ F' \to F \to F'' \xrightarrow[+1]{} $  是一个独特的三角形，那么：

\[
\chi(F)=\chi(F^{\prime})+\chi(F^{\prime \prime})
\]

（参见练习 I.34）。

我们将用 $ \mathbf{K}_{\mathbf{R}-c}(X) $  表示 $ \mathbf{D}_{\mathbf{R}-c}^{b}(X) $  的格洛腾迪克群（参见练习I.27）。如果存在显着的三角形  $ F' \to F \to F'' \xrightarrow{+1} $  ，则通过关系  $ F = F' + F'' $  获得该群作为由  $ \mathrm{Ob}(\mathbf{D}_{\mathbf{R}-c}^{b}(X)) $  生成的自由阿贝尔群的商。根据定理 8.1.11， $ \mathbf{K}_{\mathbf{R}-c}(X) $  也是阿贝尔范畴  $ \mathbb{R}\text{-Cons}(X) $  的格洛腾迪克群。根据 (9.7.3)，局部欧拉-庞加莱指数  $ \chi $  引发群同态，我们仍然用  $ \chi $  表示：

\[
\chi:\mathbf{K}_{\mathbf{R}-c}(X)\to\mathbf{C F}(X)
\]

\statementlabel{定理 9.7.1}(9.7.4) 中的态射 $ \chi $  是同构。

\statementlabel{证明}(a) 让 $\varphi \in \mathrm{CF}(X)$  。选择子分析分层  $X = \bigsqcup_{\alpha \in A} X_\alpha$  使得  $\varphi = \sum_{\alpha \in A} m_\alpha 1_{X_\alpha}$  、  $m_\alpha \in \mathbb{Z}$  。让  $\varepsilon_\alpha = \mathrm{sgn}(m_\alpha)$  为  $m_\alpha \neq 0$  并定义：

\[
F=\bigoplus_{\alpha\in A^{^{\prime}}}k_{X_{\alpha}}^{|m_{\alpha}|}\left[\frac{1-\varepsilon_{\alpha}}{2}\right]
\]

其中 $ A' = \{\alpha \in A; m_\alpha \neq 0\} $  。

显然  $F$  属于  $\mathbf{D}_{\mathbf{R}-c}^{b}(X)$  和  $\chi(F)=\varphi$  。这证明了 $\chi$  的满射性。 （请注意，  $F$  满足：  $\operatorname{supp}(F)=\operatorname{supp}(\varphi)$  。我们将在后面使用此注释。）

(b) 让我们证明 $\chi$  是单射的。  $\mathbf{K}_{\mathbf{R}-c}(X)$  的元素 $u$  由有限和 $u = \sum_j a_j [F_j]$  表示， $a_j \in \mathbb{Z}$  和 $[F_j]$  表示 $F_j \in \mathrm{Ob}(\mathbf{D}_{\mathbf{R}-c}^b(X))$  在 $\mathbf{K}_{\mathbf{R}-c}(X)$  中的图像。设置 $F \in \mathrm{Ob}(\mathbf{D}_{\mathbf{R}-c}^b(X))$ 、 $F^k = F \oplus \cdots \oplus F$  ( $k$ 

次）对于  $ k \in \mathbb{N} $  和  $ F^k = F^{-k}[1] $  对于  $ k < 0 $  ，我们还有：  $ u = [\bigoplus_j F^{a_j}] $  。因此，任何  $ u \in \mathbf{K}_{\mathbf{R}-c}(X) $  都由  $ \mathbf{D}_{\mathbf{R}-c}^b(X) $  的单个对象  $ F $  表示。

令  $X = \bigsqcup_{\alpha} Z_\alpha$  为亚解析分层，使得  $H^j(F)|_{Z_\alpha}$  是所有  $j$  、所有  $\alpha$  的常量。让 $X_k$  表示 $k$  共维层的并集。从显着的三角形 $F_{X_0} \longrightarrow F \longrightarrow F_{X \setminus X_0} \xrightarrow[+1]{}$ ，我们推导出：

\[
[{\cal F}]=[{\cal F}_{X_{0}}]+[{\cal F}_{X\setminus X_{0}}]\quad.
\]

继续这个过程，我们得到  $ [F] = \sum_k [F_{X_k}] $  ，因此（参见练习 I.27）：

\[
[{\cal F}]=\sum_{j,k}(-1)^{j}[H^{j}(F)_{X_{k}}]~.
\]

因此  $ \chi(F) = 0 $  意味着对于任何  $ \alpha $  ：

\[
\mathrm{d i m}_{\mathrm{\scriptsize~j~e v e n}}\oplus H^{j}(F)_{\mathrm{Z_{a}}}=\mathrm{d i m}_{\mathrm{\scriptsize~j~o d d}}\oplus H^{j}(F)_{\mathrm{Z_{a}}}\;.
\]

因此：

\[
\bigoplus_{j\mathsf{e v e n}}H^{j}(F)_{Z_{\alpha}}\simeq\bigoplus_{j\mathsf{o d d}}H^{j}(F)_{Z_{\alpha}}
\]

and

\[
\bigoplus_{j\mathsf{e v e n}}H^{j}(F)_{X_{k}}\simeq\bigoplus_{j\mathsf{o d d}}H^{j}(F)_{X_{k}}
\]

（因为  $X_{k}$  是  $Z_{\alpha}$  的不相交并集）。这显示 $[F]=0$ 。 □

现在我们将研究可构造函数的一些自然运算。令 X 和 Y 为两个流形。

(a) 对“外部产品”的定义：

\[
\mathcal{C F}_{X}\boxtimes\mathcal{C F}_{Y}\to\mathcal{C F}_{X\times Y}
\]

通过设置  $ (\varphi \otimes \psi)(x, y) = \varphi(x) \cdot \psi(y) $  。显然，如果  $ \varphi = \chi(F) $  、  $ \psi = \chi(G) $  ，则  $ \varphi \boxtimes \psi = \chi(F \boxtimes G) $  。

(b) 令 $f: Y \to X$  为流形的态射。定义“逆像”：

\[
f^{*}\colon f^{-1}\mathcal{C F}_{X}\to\mathcal{C F}_{Y}
\]

通过设置  $ (f^*\varphi)(y) = \varphi(f(y)) $  。如果  $ \varphi = \chi(F) $  ，那么显然是  $ f^*\varphi = \chi(f^{-1}F) $  。

(c) 让 $ \varphi \in \Gamma_c(X; \mathcal{C} \mathcal{F}_X) $  。定义  $ \varphi $  的“积分”，表示为  $ \int_X \varphi $  ，如下所示。选择具有紧凑支持的  $ F \in \mathrm{Ob}(\mathbf{D}_{\mathbf{R}-c}^b(X)) $  ，例如  $ \chi(F) = \varphi $  。一套：

\[
\int\limits_{X}\varphi=\chi(X;F).
\]

这个数字仅取决于 $\varphi$ ，通过函子 $\chi(X;\cdot)$ 关于可区分三角形的可加性（和定理9.7.1）。

注意， $\int_{X}\varphi$ 也可以如下计算。选择相对紧凑的局部封闭子解析子集的有限族  $\{X_{\alpha}\}$  ，例如  $\varphi=\sum_{\alpha}m_{\alpha}\mathbf{1}_{X_{\alpha}}$  、  $m_{\alpha}\in\mathbb{Z}$  。然后，我们有：

\[
\int\limits_{X}\varphi=\sum_{\alpha}m_{\alpha}\quad\chi(X;k_{X_{\alpha}})\quad.
\]

特别是，如果选择  $ X_\alpha $  的紧凑型和可收缩性，则  $ \int_X \varphi = \sum_\alpha m_\alpha $  。

(d) 令 $f: Y \to X$  为流形的态射。人们定义“直像”：

\[
f_{*}:f_{!}\mathcal{C F}_{Y}\to\mathcal{C F}_{X}
\]

通过设置：

\[
(f_{*}\psi)(x)=\int_{\dot{Y}}\psi\cdot\mathbf{1}_{f^{-1}(x)}.
\]

为了检查  $ f_* \psi $  是否可构造，应用定理 9.7.1 并用  $ \chi(G) $  和  $ G \in \mathrm{Ob}(\mathcal{D}_{\mathbb{R}-c}^b(Y)) $  表示  $ \psi $  ，这样  $ f $  在  $ \mathrm{supp}(G) $  上是正确的。然后

\[
f_{*}\chi(G)=\chi(Rf_{!}G)
\]

自 $ \chi(Rf_!G)(x) = \chi(Y; G \otimes k_{f^{-1}(x)}) $ 以来。

(e) 最后，如果 $\varphi \in \mathcal{C} \mathcal{F}_x$ ，我们将介绍它的“双重功能”，表示为 $D_x \varphi$ （或简称 $D \varphi$ ）。我们将  $\varphi$  表示为  $\varphi = \chi(F)$  并设置  $D_x \varphi = \chi(D_x F)$  。由于 $D_x$  是三角函子，而 $\chi$  是加法函数，因此这个定义是有意义的。请注意  $(D_x \varphi)(x) = \chi_c(F)(x)$  。因此 $(D_x \varphi)(x)$ 也可以计算如下。选择以  $x$  为中心的局部坐标系，并用  $B(x, \varepsilon)$  表示以  $x$  为中心、以  $\varepsilon$  为半径的开球。对于足够小的  $\varepsilon > 0$ ，我们通过引理 8.4.7 得到：

\[
\begin{aligned}{R\varGamma_{\{x\}}(X;F)}&{{}\simeq R\varGamma_{c}(B(x;\varepsilon);F)}\\ {}&{{}\simeq R\varGamma(X;F\otimes k_{\mathcal{B}(x,\varepsilon)})~.}\\ \end{aligned}
\]

因此：

\[
(\mathrm{D}_{X}\varphi)(x)=\int\limits_{X}\mathbf{1}_{B(x,\varepsilon)}\varphi
\]

对于 $0 < \varepsilon \ll 1$ 。请注意，  $(D_x\varphi)(x)$  仅取决于  $\varphi$  在  $x$  的萌芽。因此  $D_x$  是  $\mathcal{C}F_x$  的层自同态。

\[
\mathrm{D}_{X}\colon\mathcal{C F}_{X}\to\mathcal{C F}_{X}~.
\]

\statementlabel{命题 9.7.2}(i) 让 $ \varphi \in \mathrm{CF}(X) $  。然后 $ D_X \circ D_X \varphi = \varphi $ 。

(ii) 令 $f: Y \to X$  为流形的态射并令 $\psi \in \Gamma(X; f_! \mathcal{C} \mathcal{F}_Y)$  。然后 $f_* \mathrm{D}_Y \psi = \mathrm{D}_X f_*\psi$ 。

\statementlabel{证明}这是根据层的相应属性和定理 9.7.1 得出的。 ⑨

示例 9.7.3。令  $Z$  为  $X$  的局部闭子分析子集，并假设对于任何  $x \in \overline{Z}$  ，在  $X$  中存在  $x$  的开邻域  $U$  的基本系统，使得  $U \cap Z$  同态于  $\mathbb{R}^d$  。然后就可以轻松证明  $D_x k_Z \simeq k_{\overline{Z}}[d]$  。因此：

\[
\mathrm{D}_{X}\mathbf{1}_{Z}=(-1)^{d}\mathbf{1}_{\widetilde{Z}}\ .
\]

此外，假设  $\bar{Z}$  是紧凑的并且  $\int_{X} \mathbf{1}_{\bar{Z}} = 1$  。设置 $\partial Z = \bar{Z} \setminus Z$ 。从  $\int_{X} \mathbf{1}_{\partial Z} = \int_{X} \mathbf{1}_{\bar{Z}} - \int_{X} \mathbf{1}_{Z}$  和  $\int_{X} \mathbf{1}_{Z} = \int_{X} D_{X} \mathbf{1}_{Z}$  开始，我们得到：

\[
\chi(\partial Z;k_{\partial Z})=\int\limits_{X}\mathbb{1}_{\partial Z}=1-(-1)^{d}.
\]

这是欧拉公式的推广。

有了这些操作，我们现在可以为可构造函数构造与我们为层定义的操作类似的操作。

令 $f: Y \to X$  为流形的态射。我们定义：

\[
f^{!}\colon f^{-1}\mathcal{C F}_{X}\to\mathcal{C F}_{Y}
\]

by:

\[
f^{!}=\mathbf{D}_{Y}\circ f^{*}\circ\mathbf{D}_{X}~.
\]

如果 $ \varphi = \chi(F) $ ， $ F \in \mathrm{Ob}(\mathbb{D}_{R-c}^b(X)) $ ，我们得到 $ f^! \chi(F) = \chi(f^! F) $ 。同样，我们定义：

\[
h o m:\mathcal{C F}_{X}\times\mathcal{C F}_{X}\to\mathcal{C F}_{X}
\]

通过设置：

\[
h o m(\psi,\varphi)=\mathbf{D}_{X}(\psi\cdot(\mathbf{D}_{X}\varphi))~.
\]

根据命题 3.4.6，我们得到  $F$  和  $G$  中的  $\mathbf{D}_{\mathbf{R}-c}^{b}(X)$  ，即

\[
h o m(\chi(G),\chi(F))=\chi(R\mathcal{H}o m(G,F))~.
\]

为了定义可构造函数的微局域化，我们需要研究向量丛上的齐次可构造函数。

令 $ \tau: E \to Z $  为向量丛。其中一个用  $ \mathcal{CF}_{E/\mathbb{R}^+} $  表示  $ \mathcal{CF}_E $  的子层，由在  $ \mathbb{R}^+ $  的轨道上恒定的函数组成。一套：

\[
\begin{array}{r}{\mathrm{C F}_{\mathbb{R}^{+}}(E)=\varGamma(E;\mathcal{C F}_{E/\mathbb{R}^{+}})}\end{array}
\]

令 $ \mathbf{D}_{\mathbf{R}-c,\mathbf{R}+}^{b}(E) $  表示由圆锥对象组成的 $ \mathbf{D}_{\mathbf{R}-c}^{b}(E) $  的完整三角子类，并令 $ \mathbf{K}_{\mathbf{R}-c,\mathbf{R}+}(E) $  表示其格洛腾迪克群。

\statementlabel{引理 9.7.4}局部欧拉庞加莱指数  $ \chi $  引发同构：

\[
\chi:\mathbf{K}_{\mathbf{R}-c,\mathbf{R}^{+}}(E)\xrightarrow{\sim}\mathbf{C F}_{\mathbf{R}^{+}}(E)
\]

\statementlabel{证明}让 $i:Z \hookrightarrow E$ 表示零嵌入，让 $j:\dot{E} \hookrightarrow E$ 和 $\gamma:\dot{E} \to \dot{E}/\mathbb{R}^{+}$ 表示自然映射。简称 $S_{E} = \dot{E}/\mathbb{R}^{+}$ 。那么我们有一个交换图：

\[
\begin{array}{c}0\longrightarrow\mathrm{K}_{\mathrm{R}-c}(Z)\xrightarrow{\alpha}\mathrm{K}_{\mathrm{R}_{c,\mathrm{R}+}}(E)\xrightarrow{\beta}\mathrm{K}_{\mathrm{R}-c}(S_{E})\longrightarrow\longrightarrow0\\\downarrow\quad\quad\quad\quad\downarrow\quad\quad\quad\quad\downarrow\quad\quad\quad\quad\downarrow\quad\quad\quad\downarrow\\0\longrightarrow\mathrm{CF}(Z)\xrightarrow{\alpha^{\prime}}\mathrm{CF}_{\mathrm{R}+}(E)\xrightarrow{\beta^{\prime}}\mathrm{CF}(S_{E})\longrightarrow\longrightarrow0\end{array}
\]

其中  $\alpha$  由函子  $Ri_* : \mathbf{D}_{\mathbf{R}-c}^b(Z) \to \mathbf{D}_{\mathbf{R}-c,\mathbf{R}+}^b(E)$  定义， $\beta$  由函子  $R\gamma_* \circ j^{-1} : \mathbf{D}_{\mathbf{R}-c,\mathbf{R}+}^b(E) \to \mathbf{D}_{\mathbf{R}-c}^b(S_E)$  定义， $\alpha'$  是“零扩展”态射， $\beta'$  是自然态射。那么(9.7.17)中的第一个和第三个垂直箭头是定理9.7.1的同构，并且底行显然是一个精确序列。因此足以表明顶行是一个精确的序列。令  $\beta'' : \mathbf{K}_{\mathbf{R}-c}(S_E) \to \mathbf{K}_{\mathbf{R}-c,\mathbf{R}+}(E)$  为源自  $F \mapsto R j_! \gamma^{-1}F$  的映射（参见命题 8.3.8(i)）， $\alpha'' : \mathbf{K}_{\mathbf{R}-c,\mathbf{R}+}(E) \to \mathbf{K}_{\mathbf{R}-c}(Z)$  为源自  $F \mapsto i^{-1}F$  的映射。然后，对于 $F \in \mathbf{D}_{\mathbf{R}-c,\mathbf{R}+}^b(E)$ ，我们有一个显着的三角形 $R j_! \gamma^{-1}R\gamma_* F \longrightarrow F \longrightarrow i_* i^{-1}F \xrightarrow[+1]{}$ ，因此我们得到 $1 = \beta'' \circ \beta + \alpha \circ \alpha''$ ， $\alpha'' \circ \beta'' = 0$ 。这显示了第一条水平线的准确性。

然后我们定义态射：

\[
\tau_{!}\qquad\mathrm{a n d}\qquad\tau_{*}\colon\tau_{*}\mathcal{C F}_{E/\mathbb{R}^{+}}\to\mathcal{C F}_{Z}
\]

通过设置：

\[
\tau_{!}\varphi=i^{!}\varphi\qquad\mathrm{a n d}\qquad\tau_{*}\varphi=i^{*}\varphi
\]

其中  $i$  是零嵌入  $Z \hookrightarrow E$  。

我们可以用另一种方式来描述 $\tau_!$ 和 $\tau_*$ 。在向量丛  $E$  上选择一个度量，并让  $B(\varepsilon)$  为  $E$  的开集，这样对于每个  $z \in Z$  ，  $B(\varepsilon) \cap \tau^{-1}(z)$  是半径为  $\varepsilon$  且中心为 0 的开球（诱导范数）。让 $\varphi \in \mathrm{CF}_{\mathbb{R}^+}(E)$  。由命题 3.7.5 可知：

\[
\tau_{!}\varphi=\tau_{*}(\varphi\cdot\mathbf{1}_{B(\varepsilon)})\qquad\mathrm{a n d}\qquad\tau_{*}\varphi=\tau_{*}(\varphi\cdot\mathbf{1}_{\overline{{B(\varepsilon)}}})\ .
\]

我们现在可以定义可构造函数的傅立叶变换。我们遵循 III §7 的符号，用  $\pi: E^* \to Z$  表示对偶向量丛，用  $p_1$  和  $p_2$  分别表示从  $E \times_Z E^*$  到  $E$  和  $E^*$  的投影。我们设置  $P' = \{(x, y) \in E \times_Z E^*; \langle x, y \rangle \leqslant 0\}$  并用  $i_E$  表示零嵌入  $E^* \simeq Z \times_Z E^* \hookrightarrow E \times_Z E^*$  。

\statementlabel{定义 9.7.5}让 $ \varphi \in \mathrm{CF}_{\mathbb{R}^{+}}(E) $  。

一套：

\[
\varphi^{\wedge}=p_{2}!((p_{1}^{*}\varphi)\cdot\mathbf{1}_{P^{\prime}})\quad,
\]

\[
\varphi^{\vee}=(-1)^{n}\varphi^{\wedge a},
\]

其中 $n$ 是 $E$ 的纤维尺寸， $\varphi^{\wedge}=a^{*}\varphi^{\wedge}$ 是 $\varphi^{\wedge}$ 的对映图的图像。

我们将  $ \varphi^{\wedge} $  （分别为  $ \varphi^{\vee} $  ）称为  $ \varphi $  的傅里叶变换（分别为傅里叶逆变换）。

请注意，还有：

\[
\begin{array}{r}{\varphi^{\wedge}=i_{E}^{!}((p_{1}^{*}\varphi)\cdot\mathbf{1}_{P^{\prime}}).}\end{array}
\]

通过 $\wedge$ 的构造，如果 $ F \in \text{Ob}(\mathbf{D}_{R-c}^b(E)) $ 和 $ F $ 是圆锥曲线，则：

\[
\chi(F)^{\wedge}=\chi(F^{\wedge})~.
\]

因此，根据引理 9.7.4 和定理 3.7.9，我们得到：

\statementlabel{命题 9.7.6}让 $ \varphi \in \mathrm{CF}_{\mathbb{R}^+}(E) $  。那么 $ \varphi^\wedge $  和 $ \varphi^\vee $  属于 $ \mathrm{CF}_{\mathbb{R}^+}(E^*) $  而且：

\[
\varphi^{\wedge\vee}=\varphi^{\vee\wedge}=\varphi~.
\]

“翻译”命题 3.7.13、3.7.14 和 3.7.15 以获得可构造函数的类似结果是一个简单的练习（留给读者）。

我们只需要注意，傅里叶变换会导致基础变化。特别是，可以按纤维计算  $ \varphi^\wedge $ ，即，如果  $ \varphi \in \mathrm{CF}_{\mathbf{R}^+}(E) $  和  $ z \in Z $  ，则：

\[
\left(\varphi^{\wedge}\right)|_{\pi^{-1}(z)}=(\varphi|_{\tau^{-1}(z)})^{\wedge}.
\]

示例 9.7.7。 (i) 为简单起见，假设  $E$  是向量空间。让 $y \in E^*$  。那么根据它的定义：

\[
\varphi(y)=\int\limits_{E}(\varphi(x)\cdot\mathbf{1}_{|x|<\varepsilon}\cdot\mathbf{1}_{\langle x,y\rangle\leqslant0})~.
\]

(ii) 令 $U$  为 $E$  的开凸锥，并令 $\gamma$  为包含 $Z$  的 $E$  的闭凸真锥。然后：

\[
({\bf1}_{U})^{\wedge}=(-1)^{n}{\bf1}_{U^{\circ a}},
\]

\[
({\bf1}_{\gamma})^{\wedge}={\bf1}_{\mathtt{I n t}\;\gamma^{\circ}}\;.
\]

这是从引理 3.7.10 得出的（或者可以直接计算）。

最后我们将定义可构造函数的专业化和微局部化。

令 $M$  为 $X$  的闭子流形， $i$  嵌入 $M \hookrightarrow X$  并回想一下第四章的映射 $p, j, s$  和图(4.1.5)。

\statementlabel{定义 9.7.8}(i) 让 $ \varphi \in \mathrm{CF}(X) $  。

一套：

\[
\nu_{M}(\varphi)=-s^{!}((p^{*}\varphi)\cdot\mathbf{1}_{\Omega}),
\]

\[
\mu_{M}(\varphi)=\nu_{M}(\varphi)^{\wedge}
\]

(ii) 令 $ \varphi $  和 $ \psi $  属于CF(X)。

一套：

\[
\mu\mathrm{h o m}(\psi,\varphi)=\mu_{\varDelta}(\varphi\boxtimes\mathbf{D}_{X}\psi).
\]

因此我们定义了态射：

\[
\begin{array}{r}{\nu_{M}\colon i^{-1}\mathcal{C F}_{X}\to\tau_{*}\mathcal{C F}_{T_{M}X/\mathbb{R}^{+}}\mathrm{~,~}}\end{array}
\]

\[
\mu_{M}\colon i^{-1}\mathcal{C F}_{X}\to\pi_{*}\mathcal{C F}_{T_{M}^{*}X/\mathbb{R}^{+}}\enspace,
\]

\[
\mu\mathrm{h o m}\colon\mathcal{C F}_{X}\times\mathcal{C F}_{X}\to\pi_{*}\mathcal{C F}_{T^{*}X/\mathbb{R}^{+}}.
\]

当然，这些态射“通勤”到 $\chi$ 。例如

\[
\chi(\mu\mathcal{h}o m(G,F))=\mu\mathcal{h}o m(\chi(G),\chi(F))
\]

对于 $\mathbf{D}_{R-c}^{b}(X)$  中的 $F$  和 $G$  。

示例 9.7.9。 (i) 设  $ (x', x'') $  为  $ X $  上的局部坐标系，使得  $ M = \{x' = 0\} $  。令 $ t $  为 $ \mathbb{R} $  上的坐标，让 $ j $  表示嵌入 $ X \times \mathbb{R}^+ \hookrightarrow X \times \mathbb{R} $  。然后：

\[
(\nu_{M}(\varphi))(x^{\prime},x^{\prime\prime})=-\mathbf{D}_{X}((\mathbf{D}_{X\times\mathbb{R}}(\varphi(t x^{\prime},x^{\prime\prime})\cdot\mathbb{1}_{\{t>0\}}))|_{t=0})~.
\]

(ii) 令 $(x,y)$  表示 $\mathbb{R}^2$  上的坐标，并设置 $M=\{0\}$  、 $A=\{y=0,x\geqslant0\}\cup{y=x^2,x\geqslant0}$  。然后 $C_M(A)=\{y=0,x\geqslant0\}$  但 $v_M(\mathbb{1}_A)=2\mathbb{1}_{C_M}(A)-\mathbb{1}_{\{0\}}$  。

因此，一般来说， $ \nu_{M}(1_{A}) \neq 1_{C_{M}(A)} $  。 （但是  $ \nu_{M}(1_{A}) $  的支持始终包含在  $ C_{M}(A) $  中。）

(iii) 令 $ \tau: E \to Z $  为向量丛并令 $ \varphi \in \mathrm{CF}_{\mathbb{R}^+}(E) $  。然后 $ \nu_Z(\varphi) = \varphi $ 。

现在让我们研究  $ T^{*}X $  上的拉格朗日循环群。

让 $ \mathcal{L}_X $  表示环 $ \mathbb{Z} $  上的拉格朗日循环 $ T^*X $  上的层。根据命题 9.4.5，特征环  $ \operatorname{CC}(\cdot) $  定义了同态（仍记为 CC）：

\[
{\mathrm{C C}}\colon\mathbf{K}_{\mathbf{R}-c}(X)\to H^{0}(T^{*}X;\mathcal{L}_{X})~.
\]

\statementlabel{定理 9.7.10}(9.7.25)中的同态CC是同构。

\statementlabel{证明}(a) 令 $\lambda$  为拉格朗日循环。我们将用  $\mathrm{CC}(F) = \lambda$  构造  $F \in \mathrm{Ob}(\mathbb{D}_{R-c}^{b}(X))$  。让 $A = \mathrm{supp}(\lambda)$  。这是  $T^*X$  的圆锥亚解析各向同性子集。

我们将通过 $\dim \pi(A)$  进行归纳论证。存在  $X$  的亚解析子流形  $Y$  ，使得  $Y$  是  $\pi(A)$  、  $\dim(\pi(A)\setminus Y) < \dim \pi(A)$  和  $A \cap \pi^{-1}(Y) \subset T_Y^* X$  的开子集。因此  $\Gamma_{T_Y^*X}(\mathcal{L}_X)$  是  $\pi^{-1}(Y)$  上排名第一的局部常数。因此存在  $F \in \mathrm{Ob}(\mathbf{D}_{\mathbf{R}-c}^b(X))$  使得  $\mathrm{supp}(F) \subset \overline{Y}, H^j(F)|_Y$  是局部常数，并且  $\mathrm{CC}(F) = \lambda$  在  $\pi^{-1}(Y)$  上。设置 $\mu = \lambda - \mathrm{CC}(F)$ ，我们有 $\dim \pi(\mathrm{supp}(\mu)) < \dim \pi(A)$ ，并且用 $\mu$ 替换 $\lambda$ ，归纳继续。

(b) 设  $ F \in \mathrm{Ob}(\mathbf{D}_{\mathbb{R}-c}^b(X)) $  ，并假设  $ \mathrm{CC}(F) = 0 $  。我们将在 $ K_{\mathbb{R}-c}(X) $ 中证明 $ F = 0 $ 。通过定理 9.7.1，足以证明  $ \chi(F)(x) = 0 $  对于所有  $ x $  。但这是从备注 9.5.5 得出的。   $ \square $ 

人们定义了“欧拉态射”：

\[
Eu:\pi_{*}\mathcal{L}_{X}\to\mathcal{C F}_{X}
\]

如下。令 $ x \in X $ 并让 $ \phi $ 满足 $ \phi(x) = 0 $ 、 $ \mathrm{d}\phi(x) = 0 $ 且 $ \phi $ 在 $ x $ 的Hessian矩阵是正定的。让  $ \lambda $  成为  $ \pi_*\mathcal{L}_X $  的一部分。一套

\[
E u(\lambda)(x)=\#([\sigma_{\phi}]\cap\lambda)_{x}.
\]

由定理9.7.10和备注9.5.5， $ Eu $ 是明确定义的，并且我们得到：

\statementlabel{定理 9.7.11}示意图：

\begin{center}
\includegraphics[width=0.37\textwidth]{流行上的层_Chn-assets/img_in_image_box_352_934_800_1122.png}
\end{center}

是可交换的，并且箭头是同构的。

\subsection{习题}
\statementlabel{练习 IX.1}令  $F$  为流形  $X$  上的  $\mathbb{R}$  可构造层。证明  $\mathcal{Hom}(F,\mathcal{CS}^X)$  与  $\mathbf{D}^b(X)$  中的  $R\mathcal{Hom}(F,\omega_X)$  同构。 （提示：使用练习 IX.2。）

\statementlabel{练习 九.2}令  $S$  为流形  $X$  的封闭子解析子集。  $S$  上的连续实值函数  $\varphi$  如果其图在  $X \times \mathbb{R}$  中是亚解析的，则称为亚解析函数。

(i) 证明 $ S $  上的亚解析函数的束 $ \mathcal{S}_S $  是软环层，并且态射 $ \mathcal{S}_X \to \mathcal{S}_S $  是满射的。

(ii) 如果环化空间  $ (X; \mathcal{S}) $  与赋予环层的实解析流形的亚解析子集局部同构，则称其为亚解析空间

的亚解析函数。定义  $X$  、 $\mathbb{R}$  可构造层中的子解析子集、子解析循环  $\mathcal{C}^{\mathcal{X}}$  的层复合体，并证明  $\mathcal{C}^{\mathcal{X}}$  与  $\mathbf{D}^{b}(X)$  中的  $\omega_{X}$  同构。还证明与练习 IX.1 类似的陈述。 （提示：使用练习 IX.3 中 Hironaka 的结果。）

\statementlabel{练习 IX.3}设 X 为 n 维流形，假设基环 A 为 C。证明存在复形态射：

\[
\mathcal{C S}^{x}\overset{\psi}{\longrightarrow}\mathcal{D}\mathcal{C}_{X}\otimes{c r}_{X}[n]~,
\]

这是一个拟同构。这里， $ \mathcal{D}\ell_x $  是具有分布系数的 de Rham 复形。 （提示：应使用由 Hironaka 提供的以下结果——作者感谢他允许他们将其插入此处——：

“让  $S$  是维度  $p$  的  $X$  的局部闭实解析子流形，并在  $X$  中进行子解析。然后存在一个  $p$  维实解析流形  $Y$  、实解析流形  $f: Y \to X$  的态射和开放嵌入  $i: S \hookrightarrow Y$  使得 $i(S)$  在  $Y$  中是亚解析的，  $f$  在  $i(\overline{S})$  上是真映射，  $f \circ i$  与包含  $S \hookrightarrow X$  一致（甚至可以假设  $i(S)$  是由  $Y$  中的不等式  $\{y_1 > 0, \ldots, y_k > 0\}$  局部定义的，其中  $(y_1, \ldots, y_p)$  是局部的 $Y)$  上的坐标系”。

然后在(9.E.1)中定义 $\psi$ 如下。对于  $p$  维导向的子分析子流形  $S$  和  $\theta \in \Gamma_c(X; \mathcal{C}_X^\infty,(p))$  集：  $\int \psi([S])\theta = \int_S \theta$  。

\statementlabel{练习 九.4}令 $f: Y \to X$  为流形态射，让 $\varphi: Y \to \mathbb{R}$  为实解析函数。对于  $\lambda \in \mathbb{R}$  通过  $p \mapsto^t f'(p) + \lambda \mathrm{d}\varphi(\pi(p))$  定义映射  $\rho_\lambda: Y \times_X T^*X \to T^*Y$  。如果  $S$  是  $T^*Y$  的闭子集，则定义：

\[
\alpha_{\lambda}\colon H_{S}^{0}(T^{*}Y;\pi_{Y}^{-1}\omega_{Y})\to H_{\rho_{\lambda}^{-1}(S)}^{0}\left(\begin{matrix}{Y\times T^{*}X;\pi^{-1}\omega_{Y}}\\ {x}\\ \end{matrix}\right)
\]

通过使用态射  $ \rho_{\lambda}^{-1}\pi_{Y}^{-1}\omega_{Y} \simeq \pi^{-1}\omega_{Y} $  。 If T is a closed subset of   $ Y \times_X T^*X $   such that   $ f_\pi $   is proper on T, define:

\[
\beta:H_{T}^{0}\left(Y\mathop{\times}_{\begin{array}{c}{x}\\ \end{array}}T^{*}X;\pi^{-1}\omega_{Y}\right)\to H_{f_{x}(T)}^{0}(T^{*}X;\pi_{X}^{-1}\omega_{X})
\]

通过使用态射：  $ Rf_{\pi!}\pi^{-1}\omega_{X}\simeq\pi_{X}^{-1}Rf_{!}\omega_{Y}\to\pi_{X}^{-1}\omega_{X} $  。

现在让  $G \in \mathrm{Ob}(\mathbf{D}_{\mathbf{R}-c}^{b}(Y))$  ，并假设  $f$  对于所有  $t$  ，在  $\mathrm{supp}(G) \cap \overline{Y}_{t}$  上是真映射，其中  $Y_{t} = \{\varphi < t\}$  。让  $t_{\circ} \in \mathbb{R}$  ，并设置一些  $\lambda > 0$  固定：

\[
T_{\pm}=\bigcup_{0\leqslant\pm\mu\leqslant\lambda}\rho_{\mu}^{-1}(\mathbf{S S}(G)\cap\pi^{-1}(\widetilde{Y}_{t_{\circ}}))~.
\]

(i) 假设对于  $ y \in Y \setminus Y_{t_0} $  ，

\[
\mathsf{d}\varphi(y)\notin\mathsf{S S}(G)+{}^{t}f^{\prime}\left(Y\mathop{\times}_{\substack{x}}T^{*}X\right).
\]

证明  $ \mathcal{R}f_* G \in \mathrm{Ob}(\mathbf{D}_{\mathbf{R}-c}^b(X)) $  、  $ \rho_\lambda^{-1}(\mathrm{SS}(G)) $  包含在  $ T_+ $  中， $ CC(Rf_* G) = \beta \circ \alpha_\lambda(CC(G)) $  包含在  $ H_{f_*(T_+)}^0(T^* X; \pi_X^{-1} \omega_X) $  中。

(ii) 证明类似的结果，将 $ \mathrm{d}\varphi(y) $ 替换为 $ -\mathrm{d}\varphi(y) $ ， $ Rf_G $ 替换为 $ Rf_G $ ， $ \alpha_\lambda $ 替换为 $ \alpha_{-\lambda} $ ， $ T_+ $ 替换为 $ T_{-} $ 。

\statementlabel{练习 九.5}令 $ \tau: E \to Z $  为向量丛，令 $ \theta $  为 $ E $  上的欧拉向量场并设置 $ P_\pm = \{\pm\theta \geqslant 0\} \subset T^*E $  。令  $ A $  为  $ P_+ \cap P_- $  的闭圆锥子集，并用  $ i $  表示嵌入  $ T^*Z \hookrightarrow P_\pm \subset T^*E $  。

(i) 定义态射：

\[
\varphi_{\pm}:R\varGamma_{A}(\pi_{E}^{-1}\omega_{E})\to i_{*}\pi_{Z}^{-1}\omega_{Z}\enspace.
\]

(Hint: consider the morphisms:   $ \mathrm{R}\Gamma_{A}(\pi_{E}^{-1}\omega_{E}) \simeq \mathrm{R}\Gamma_{A}\mathrm{R}\Gamma_{P_{\pm}}(\pi_{E}^{-1}\omega_{E}) \to \mathrm{R}\Gamma_{A}(A_{P_{\pm}}\otimes\mathrm{R}\Gamma_{P_{\pm}}(\pi_{E}^{-1}\omega_{E})) $   and use   $ A_{P_{\pm}} \otimes \mathrm{R}\Gamma_{P_{\pm}}(\pi_{E}^{-1}\omega_{E}) \simeq i_{*}\pi_{Z}^{-1}\omega_{Z}. $ 

(ii) 定义： $ \psi_{\pm}: H_{\lambda}^{0}(T^{*}E;\pi_{E}^{-1}\omega_{E}) \to H_{i^{-1}(\lambda)}^{0}(T^{*}Z;\pi_{Z}^{-1}\omega_{Z}) $  。

(iii) 令 $F$  为 $\mathbf{D}_{\mathbf{R}-c}^{b}(E)$  的圆锥对象。证明：

\[
\mathrm{C C}(R\tau_{*}F)=\psi_{-}(\mathrm{C C}(F))~,
\]

\[
\mathrm{C C}(R\tau_{!}F)=\psi_{+}(\mathrm{C C}(F))~.
\]

（提示：使用练习 IX 4。）

\statementlabel{练习 IX.6}令 X 为奇数维紧流形。证明 $ \chi(X; k_{X}) = 0 $ 。

\statementlabel{练习 九.7}令  $F$  为  $\mathbf{D}^{b}(E)$  的二次曲线  $\mathbb{R}$  可构造对象，其中  $E$  是向量丛。使用练习VII.2，证明 $\mathrm{CC}(F) = \mathrm{CC}(F^{\wedge})$ ，在如命题5.5.1中识别 $T^{*}E$ 和 $T^{*}E^{*}$ 之后，以及在识别 $\pi_{E}^{-1} \omega_{E}$ 和 $\pi_{E}^{-1} \omega_{E^*}$ 之后。

\statementlabel{练习 九.8}（霍普夫指数定理）。 令 $X$ 为紧流形，$v$ 为 $X$ 上仅有有限多个零点 $x_1,\ldots,x_N$ 的实解析向量场。假设每个  $x_k$  (  $k = 1, \ldots, N$  ) 是非简并的（即：在  $x_k$  、  $v = \sum_{ij} a_{ij}^k x_i \frac{\partial}{\partial x_j} + O(|x|^2$  邻域的局部图中，矩阵  $A_k = (a_{ij}^k)$  是非简并的）。设置 $\varepsilon_k = \text{sgn}\det(A_k)$ 。然后证明：

\[
\chi(X;\mathbb{Q}_{X})=\sum_{k=1}^{N}\varepsilon_{k}\enspace.
\]

（提示：考虑v给定的 $ TX \to X $ 的部分s，并计算s与零部分的交集数。）

\statementlabel{练习 九.9}去掉命题 9.6.2 中 $\operatorname{supp}(F)$ 紧的假设，并证明同样的结论。

\statementlabel{练习 九.10}令 $E$  为有限维向量空间，并令 $\mathcal{A} = \Gamma_c(E; \mathcal{C} \mathcal{F}_E)$  为 $E$  上紧支持的可构造函数组。

(i) 通过公式 $\varphi * \psi = s_*(\varphi \boxtimes \psi)$  定义 $\mathcal{A}$  上的卷积运算，其中 $s$  是映射 $E \times E \to E$  、 $(x, y) \mapsto x + y$  。

(ii) 证明 $ (\mathcal{A},*) $  是一个交换代数，单位为 $ 1_{\{0\}} $  。

(iii) 令 Z 为 E 的凸紧致子解析子集。证明：

\[
\mathbf{1}_{Z}{*}D(\mathbf{1}_{-Z})=\mathbf{1}_{\{0\}}.
\]

(iv) 证明 $ \mathrm{D}(\varphi \star \psi) = \mathrm{D}\varphi \star \mathrm{D}\psi $  和 $ \int_{E}(\varphi \star \psi) = (\int_{E}\varphi)(\int_{E}\psi) $  。

\statementlabel{练习 九.11}令 $k$ 为域，令 $V$ 为实向量空间，$V^*$ 为其对偶空间，$q$ 为投影 $T^*V\simeq V\times V^*\to V^*$。令 $F\in\mathrm{Ob}(\mathbf{D}_{\mathbf{R}-c}^b(k_V))$ 具有紧支撑，并假设： (i)   $F \simeq k_{\{x_0\}}^m$   in   $\mathbf{D}^b(V; p)$  , where   $p = (x_0; \xi_0) \in T^*V$  , and (ii)   $q^{-1}(\xi_0) \cap \mathrm{SS}(F) = \{p\}$  .

证明 $ q(\mathrm{SS}(F)) = V^* $ 。

（提示：证明 $ \operatorname{supp}(q_* \operatorname{CC}(F)) = V^* $ 。）

\statementlabel{练习 九.12}令 X 为复流形。

(i) 让 $ F \in \mathrm{Ob}(\mathbb{D}_{\mathbb{C}-c}^b(X)) $  。 Prove that   $ \chi(F_x) = \chi(RF_{\{x\}}(X; F)) $   for any   $ x \in X $   (cf. Sullivan [1]).

(ii) 令  $ \phi $  为  $ X $  上的 C 可构造函数，即存在复杂分层  $ X = \bigsqcup_{\alpha} X_{\alpha} $  使得  $ \phi|_{X_{\alpha}} $  为常数。证明 $ D_{X}(\phi) = \phi $ 。

（提示：使用特化函子，简化为  $X$  是复向量空间且  $F$  在  $\mathbb{C}^\times$  轨道上局部常值的情况。然后使用  $\mathbf{S}^1$  的欧拉-庞加莱指数为零这一事实。）

\subsection{注记}
“圈”的概念是代数几何的基础，我们参考富尔顿[1]的书进行回顾。在真实的解析情况下，§2的构造是同调理论专家所熟知的结果的上同调版本，特别是在Borel-Moore [1]、Borel-Haefliger [1]、Herrera [1]、Bloom-Herrera [1]、Poly [1]和Verdier [4]的工作之后。

格洛腾迪克 (Grothendieck) 猜想出了复杂簇上可构造层的黎曼-罗赫公式。这是由 MacPherson [1] 获得的，他引入了所谓的“局部欧拉阻碍”（也参见 M-H. Schwartz [1] 的开创性工作），并证明了循环群与可构造函数之间的同构与直像交换。然后Sabbah [1]和Ginsburg [1,2]构建了拉格朗日循环上的函子运算，并在复杂情况下得到了与§3类似的结果。然而，我们在§3中的方法本质上是不同的。

Kashiwara [2,5] 首次证明了完整 D 模的索引定理，他引入了一个不变量，该不变量被证明等价于

麦克弗森（参见 Brylinski-Dubson-Kashiwara [1]）。然后复杂情况下的 Dubson [1] 和真实情况下的 Kashiwara [7] 表明，指数公式可以解释为交集公式。然而，§4 的特征循环的构造是最近的（前一个使用公式（9.4.10）），并且是由于 Kashiwara [8]。

我们不会在这里回顾 Lefschetz 公式，而只是回顾一下，Verdier 首先发现 Grothendieck 的对偶形式主义可以应用于在 étale 上下文中导出可构造层之间对应关系的一般迹公式（参见 [SGA 5]，exposé III）。除此之外，§6的结果本质上是新的，并且是由Kashiwara [8]在88年宣布的。然而，应该提到的是，有边界流形的Lefschetz公式首先由Brenner-Shubin [1]证明，并且扩展或收缩空间的概念已经出现在他们的论文中。

最后请注意，特征类的构造最近已被 Schapira-Schneiders [1]（也参见 Schapira [4]）扩展到包括 D 模在内的更广泛的层类。

\section{反常层}
\subsection{概要}
尽管反常层的历史很短，但它们在数学的各个分支中发挥着重要作用，例如代数几何或群表示。这一理论现在已被很好地理解，但在文献中很难找到对其进行系统处理的分析案例。

在本章中，我们在实解析流形  $X$  上定义反常层，并证明它们形成  $\mathsf{D}_{\mathsf{R}-c}^{b}(X)$  的阿贝尔子范畴。然后我们研究 $X$  是复流形的情况。我们证明反常性是一个微局部属性，也就是说， $\mathsf{D}_{\mathsf{C}-c}^{b}(X)$  的对象是反常的当且仅当它在其微支撑的通用点上是纯粹的移位 $\to\dim_{C}X$  时。由于莫尔斯理论很好地适应了微局域表征，因此我们可以证明斯坦因流形上反常层的消失定理。然后我们证明反常性可以通过各种操作，特别是通过专业化、微局部化和傅里叶佐藤变换来保留。为了实现这个程序，我们需要同调代数的一些补充。在第 1 节中，我们解释了  $t$  结构的概念，它允许在三角范畴中构造阿贝尔子范畴。

参考 Beilinson-Bernstein-Deligne [1] 和 Goresky-MacPherson [2,3]（也参见 Borel 等人 [1] 和 Gelfand-Manin [1]）。

我们保留 8.0 的惯例。

\subsection{$t$-结构}
令 D 为三角范畴（参见 I §5）。

\statementlabel{定义 10.1.1}让  $ \mathbf{D}^{\leq0} $  和  $ \mathbf{D}^{\geq0} $  成为  $ \mathbf{D} $  的满子范畴。如果满足以下条件，我们说 $ (\mathbf{D}^{\leq0}, \mathbf{D}^{\geq0}) $ 是 $ \mathbf{D} $ 上的t结构。在这里， $ \mathbf{D}^{\leq n} = \mathbf{D}^{\leq0}[-n] $  和 $ \mathbf{D}^{\geq n} = \mathbf{D}^{\geq0}[-n] $  。

(i)  $ \mathbf{D}^{≤-1} \subset \mathbf{D}^{≤0} $  和 $ \mathbf{D}^{\geq1} \subset \mathbf{D}^{\geq0} $  。

(ii)  $ \mathrm{Hom}_{\mathbf{D}}(X, Y) = 0 $  用于  $ X \in \mathrm{Ob}(\mathbf{D}^{\leq 0}) $  和  $ Y \in \mathrm{Ob}(\mathbf{D}^{\geq 1}) $  。

(iii) 对于任何 $ X \in \text{Ob}(\mathbb{D}) $  ， $ \mathbb{D} $  中存在一个与 $ X_0 \in \text{Ob}(\mathbb{D}^{\leq 0}) $  和 $ X_1 \in \text{Ob}(\mathbb{D}^{\geq 1}) $  不同的三角形 $ X_0 \longrightarrow X \longrightarrow X_1 \xrightarrow{+1} X $  。

完整的子范畴 $ \mathcal{C} = \mathbf{D}^{\leq 0} \cap \mathbf{D}^{\geq 0} $  称为 $ t $  结构的核心。

\statementlabel{备注 10.1.2}(i)  $t$  结构的概念是自对偶的，即如果  $(\mathbf{D}^{\leq 0}, \mathbf{D}^{\geq 0})$  是  $\mathbf{D}$  上的  $t$  结构，则  $(( \mathbf{D}^{\geq 0})^{\circ}, ( \mathbf{D}^{\leq 0})^{\circ})$  是相反三角范畴  $\mathbf{D}^{\circ}$  上的  $t$  结构。

(ii) 如果  $ (\mathbf{D}^{\leq 0}, \mathbf{D}^{\geq 0}) $  是  $ \mathbf{D} $  上的  $ t $  结构，则  $ (\mathbf{D}^{\leq n}, \mathbf{D}^{\geq n}) $  也是  $ \mathbf{D} $  上的  $ t $  结构。

让我们举一些例子。

Examples 10.1.3. (i) 令 $\mathcal{C}$  为阿贝尔范畴，并令 $\mathbf{D} = \mathbf{D}(\mathcal{C})$  表示其导出范畴。用 $\mathbf{D}^{\leq 0}(\mathcal{C})$ （或 $\mathbf{D}^{\geq 0}(\mathcal{C})$ ）表示 $\mathbf{D}$ 的满子范畴，由满足 $j > 0$ （或 $j < 0$ ）的 $H^j(X) = 0$ 的对象 $X$ 组成。那么  $(\mathbf{D}^{\leq 0}(\mathcal{C}), \mathbf{D}^{\geq 0}(\mathcal{C}))$  是  $\mathbf{D}(\mathcal{C})$  上的  $t$  结构。事实上，定义 10.1.1 中的公理 (iii) 通过考虑区分三角形  $\tau^{\leq 0} X \longrightarrow X \longrightarrow \tau^{\geq 1} X \longrightarrow \cdots$  得到了验证。在这种情况下， $\mathbf{D}(\mathcal{C})$ 的核心相当于 $\mathcal{C}$ 。

(ii) 令 $A$  为与 $\mathrm{gld}(A) < \infty$  的诺特交换环，并令 $X$  为实解析流形。然后  $\mathbf{D}_{\mathbf{R}-c}^{b}(X) = \mathbf{D}^{b} \leq 0(A_X) \cap \mathbf{D}_{\mathbf{R}-c}^{b}(X)$  和  $\mathbf{D}_{\mathbf{R}-c}^{b}(X) = \mathbf{D}^{b} \geq 0(A_X) \cap \mathbf{D}_{\mathbf{R}-c}^{b}(X)$  在  $\mathbf{D}_{\mathbf{R}-c}^{b}(X)$  上给出  $t$  结构。令  $\mathbf{D}_X$  为对偶函子（定义 3.1.16）。然后 $\mathrm{D}_X$ 的 $\mathbf{D}_{\mathbf{R}-c}^{b}(X)$ 和 $\mathbf{D}_{\mathbf{R}-c}^{b}(X)$ 的图像在 $\mathbf{D}_{\mathbf{R}-c}^{b}(X)$ 上形成 $t$ 结构。

在 X§1 的其余部分中，  $ (\mathbf{D}^{\leq 0}, \mathbf{D}^{\geq 0}) $  是三角范畴  $ \mathbf{D} $  上的  $ t $  结构。

\statementlabel{命题 10.1.4}(i) 包含  $ \mathbf{D}^{\leq n} \to \mathbf{D} $  （分别为  $ \mathbf{D}^{\geq n} \to \mathbf{D} $  ）具有右伴随函子  $ \tau^{\leq n} : \mathbf{D} \to \mathbf{D}^{\leq n} $  （分别为左伴随函子  $ \tau^{\geq n} : \mathbf{D} \to \mathbf{D}^{\geq n} $  ），即存在态射  $ \tau^{\leq n} \to \mathrm{id}_{\mathbf{D}} $  （分别为  $ \mathrm{id}_{\mathbf{D}} \to \tau^{\geq n} $  ），使得：

\[
\mathrm{H o m}_{\mathsf{D}^{\leqslant n}}(X,\tau^{\leqslant n}Y)\to\mathrm{H o m}_{\mathsf{D}}(X,Y)
\]

是任何  $ X \in \text{Ob}(\mathbf{D}^{\leq n}) $  和  $ Y \in \text{Ob}(\mathbf{D}) $  的同构（分别是  $ \text{Hom}_{\mathbf{D}^{\geq n}}(\tau^{\geq n}X, Y) \to \text{Hom}_{\mathbf{D}}(X, Y) $  是任何  $ X \in \text{Ob}(\mathbf{D}) $  和  $ Y \in \text{Ob}(\mathbf{D}^{\geq n}) $  的同构）。

(ii) 存在唯一的态射 $\mathrm{d}:\tau^{\geq n+1}(X)\to\tau^{\leq n}(X)[1]$ ，使得 $\tau^{\leq n}(X)\to X\to\tau^{\geq n+1}(X)\to\tau^{\leq n}(X)[1]$  是一个独特的三角形。此外  $d$  是函子  $\tau^{\geq n+1}\to[1]\circ\tau^{\leq n}$  的态射。

函子 $\tau^{\leq n}$  和 $\tau^{\geq n}$  被称为相对于 $t$  结构的截断函子。

\statementlabel{证明}我们可以假设  $n = 0$  。根据定义 10.1.1 (iii)，对于  $X \in \mathrm{Ob}(\mathbf{D})$  ，存在一个与  $X_0 \in \mathrm{Ob}(\mathbf{D}^{\leq 0})$  和  $X_1 \in \mathrm{Ob}(\mathbf{D}^{\geq 1})$  不同的三角形  $X_0 \longrightarrow X \longrightarrow X_1 \xrightarrow[+1]{}$  。

我们将证明  $ \text{Hom}_{\mathbf{D}}(Y, X_0) \to \text{Hom}_{\mathbf{D}}(Y, X) $  是任何  $ Y \in \text{Ob}(\mathbf{D}^{\leq 0}) $  的同构（分别是  $ \text{Hom}_{\mathbf{D}}(X_1, Z) \to \text{Hom}_{\mathbf{D}}(X, Z) $  是任何  $ Z \in \text{Ob}(\mathbf{D}^{\geq 1}) $  的同构）。这是从长的精确序列得出的：

\[
\operatorname{Hom}(Y,X_{1}[-1])\to\operatorname{Hom}(Y,X_{0})\to\operatorname{Hom}(Y,X)\to\operatorname{Hom}(Y,X_{1})
\]

（分别）

\[
\operatorname{H o m}(X_{0}[1],Z)\to\operatorname{H o m}(X_{1},Z)\to\operatorname{H o m}(X,Z)\to\operatorname{H o m}(X_{0},Z))\enspace.
\]

这给出了 (i)  $ X_0 \simeq \tau^{\leq 0} X $  和  $ X_1 \simeq \tau^{\geq 1} X $ （参见练习 I.2）。 (ii) 的第一部分源自定义 10.1.1 (iii) 和下一个引理。

\statementlabel{引理 10.1.5}令  $\mathbf{D}$  为三角范畴，并令  $X\xrightarrow{f}Y\xrightarrow{g}Z\xrightarrow{h_{i}}X[1]$  (  $i=1,2$  ) 为两个不同的三角形。如果  $\operatorname{Hom}_{\mathbf{D}}(X[1],Z)=0$  ，则  $h_{1}=h_{2}$  。

\statementlabel{证明}通过 (TR4)（参见命题 1.4.4），我们有三角形的态射：

\[
\begin{array}{c}X\xrightarrow{f}Y\xrightarrow{g}Z\xrightarrow{h_{1}}X[1]\\\downarrow\quad\downarrow\quad\downarrow\quad\phi\quad\downarrow\quad\downarrow\\\left.\phantom{\frac{f}{g}}\right.X\xrightarrow{f}Y\xrightarrow{g}Z\xrightarrow{h_{2}}X[1]\end{array}
\]

因此  $g = \phi \circ g$  和  $h_1 = h_2 \circ \phi$  。由于  $(\mathrm{id}_z - \phi) \circ g = 0$  ，根据命题 1.5.3 存在  $\psi: X[1] \to Z$  使得  $\mathrm{id}_z - \phi = \psi \circ h_1$  。根据假设，  $\psi$  必须为零，因此  $\phi = \mathrm{id}_z$  、  $h_1 = h_2$  。   $\square$ 

最后让我们证明 d 是函子的态射。 对于 $f:X\to Y$，由 (TR4) 存在一个可区分三角形的态射：

\[
\begin{array}{c}\tau^{\leq n}X\quad\longrightarrow\quad X\quad\longrightarrow\quad\tau^{\geq n+1}X\quad\xrightarrow{\mathrm{d}(X)}\quad\tau^{\leq n}X[1]\\\left\downarrow_{\tau^{\leq n}f}\quad\left\downarrow f\quad\right.\quad\left.\quad\right|\phi\quad\right.\quad\left\downarrow_{\tau^{\leq n}f[1]}\quad\left.\quad\right|\phi\quad\right.\\\tau^{\leq n}Y\quad\longrightarrow\quad Y\quad\longrightarrow\quad\tau^{\geq n+1}Y\quad\xrightarrow{\mathrm{d}(Y)}\quad\tau^{\leq n}Y[1]\quad.\end{array}
\]

中间正方形的交换律意味着  $ \phi = \tau^{\geq n+1}f $  。 □

请注意，我们有：

\[
\left\{\begin{aligned}{\tau^{\leqslant n}(X[m])}&{{}\simeq\tau^{\leqslant n+m}(X)[m],}\\ {\tau^{\geqslant n}(X[m])}&{{}\simeq\tau^{\geqslant n+m}(X)[m].}\\ \end{aligned}\right.
\]

我们还应分别编写  $\tau^{>n}$  和  $\tau^{<n}$  来代替  $\tau^{\geq n+1}$  和  $\tau^{\leq n-1}$  ， $\mathbf{D}^{<n}$  和  $\mathbf{D}^{>n}$  也类似。

\statementlabel{命题 10.1.6}(i) 如果  $ X \in \text{Ob}(\mathbb{D}^{\leq n}) $  （分别为  $ X \in \text{Ob}(\mathbb{D}^{\geq n}) $  ），则态射  $ \tau^{\leq n} X \to X $  （分别为  $ X \to \tau^{\geq n} X $  ）是同构。

(ii) 让 $ X \in \mathrm{Ob}(\mathbb{D}) $  。然后 $ X \in \mathrm{Ob}(\mathbb{D}^{\leq n}) $ （或 $ \mathrm{Ob}(\mathbb{D}^{\geq n}) $ ）当且仅当 $ \tau^{>n}X = 0 $ （或 $ \tau^{<n}X = 0 $ ）。

\statementlabel{证明}(i) 源自命题 10.1.4。

(ii) 由著名的三角形  $ \tau^{\leq n}X \longrightarrow X \longrightarrow \tau^{>n}X \xrightarrow[+1]{} \square $  得出

\statementlabel{命题 10.1.7}让 $ X' \longrightarrow X \longrightarrow X'' \xrightarrow{+1} be a distinguished triangle in D $  。 If   $ X' $   and   $ X'' $   belong to   $ \mathbf{D} \geq 0 $   (resp.   $ \mathbf{D} \leq 0 $  ) then so does X.

\statementlabel{证明}It is enough to prove the statement for   $\mathbf{D}^{\geq 0}$  .假设  $X'$  和  $X''$  属于  $\mathbf{D}^{\geq 0}$  。 Then   $\mathrm{Hom}(\tau^{<0}X, X') = \mathrm{Hom}(\tau^{<0}X, X'') = 0$  , by 定义 10.1。1 (i) and (ii).因此 $\mathrm{Hom}(\tau^{<0}X, X) = 0$  。这给出了  $\tau^{<0}X = 0$  。然后应用命题 10.1.6 (ii)。   $\square$ 

\statementlabel{命题 10.1.8}设a和b是两个整数。

(i) 如果  $b \geqslant a$  ，则  $\tau^{\geqslant b} \circ \tau^{\geqslant a} \simeq \tau^{\geqslant a} \circ \tau^{\geqslant b} \simeq \tau^{\geqslant b}$  和  $\tau^{\leqslant b} \circ \tau^{\leqslant a} \simeq \tau^{\leqslant a} \circ \tau^{\leqslant b} \simeq \tau^{\leqslant a}$  。

(ii) 如果  $a > b$  ，则  $\tau^{\leqslant b} \circ \tau^{\geqslant a} = \tau^{\geqslant a} \circ \tau^{\leqslant b} = 0$  。

（三） $ \tau^{\geq a} \circ \tau^{\leq b} \simeq \tau^{\leq b} \circ \tau^{\geq a} $ 。 更准确地说，对 $X\in\text{Ob}(\mathbb{D})$，存在唯一的态射 $\phi:\tau^{\geq a}\circ\tau^{\leq b}X\to\tau^{\leq b}\circ\tau^{\geq a}X$，使下图交换：

\begin{center}
\includegraphics[width=0.38\textwidth]{流行上的层_Chn-assets/img_in_image_box_382_539_836_714.png}
\end{center}

通勤，而且  $ \phi $  是同构。

\statementlabel{证明}(i)   $\tau^{\geq a} \circ \tau^{\geq b} \simeq \tau^{\geq b}$   is a consequence of 命题 10.1。6 (i). For any   $X \in \mathrm{Ob}(\mathbf{D})$  ,   $Y \in \mathrm{Ob}(\mathbf{D}^{\geq b})$  , we have   $\mathrm{Hom}_{\mathbf{D}^{\geq b}}(\tau^{\geq b} \tau^{\geq a} X, Y) \simeq \mathrm{Hom}_{\mathbf{D}}(\tau^{\geq a} X, Y) \simeq \mathrm{Hom}_{\mathbf{D}}(X, Y) \simeq \mathrm{Hom}_{\mathbf{D}^{\geq b}}(\tau^{\geq b} X, Y)$   and hence   $\tau^{\geq b} \tau^{\geq a} X \simeq \tau^{\geq b} X$  .

对偶陈述给出了其他同构。

(ii) 紧随命题 10.1.6 (ii)。

(iii) 通过 (ii)，我们可以假设  $ b \geq a $  。通过 (i) 存在可区分的三角形：

\[
\left\{\begin{aligned}&\tau^{\leqslant b}\tau^{\geqslant a}X\longrightarrow\tau^{\geqslant a}X\longrightarrow\tau^{>b}X\xrightarrow[+1]{}\\ &\tau^{<a}X\longrightarrow\tau^{\leqslant b}X\longrightarrow\tau^{\geqslant a}\tau^{\leqslant b}X\xrightarrow[+1]{}\\ \end{aligned}\right.,
\]

因此根据前面的命题 $\tau^{\leq b}\tau^{\geq a}X$ 和 $\tau^{\geq a}\tau^{\leq b}X$ 属于 $\mathbf{D}^{\geq a}\cap\mathbf{D}^{\leq b}$ 。这给出了  $\phi$  的存在性和唯一性。现在我们将证明 $\phi$  是同构。将八面体公理 TR5 应用于态射  $\tau^{<a}X\to\tau^{\leq b}X\to X$  ，（参见 (10.1.4)）。

\begin{center}
\includegraphics[width=0.30\textwidth]{流行上的层_Chn-assets/img_in_image_box_396_1306_763_1747.png}
\end{center}

我们得到了显着的三角形：

\[
\tau^{\geq a}\tau^{\leq b}X\xrightarrow{f}\tau^{\geq a}X\xrightarrow{g}\tau^{>b}X\xrightarrow[+1]{}
\]

（请注意  $f$  和  $g$  必须是唯一的。）因此  $\tau^{\geq a}\tau^{\leq b}X\simeq\tau^{\leq b}\tau^{\geq a}X$  。 □

让 $ \mathcal{C} $ 成为D的核心。

\statementlabel{定义 10.1.9}我们通过以下方式定义函子 $ H^0: \mathbf{D} \to \mathcal{C} $ ：

\[
H^{0}(X)=\tau^{\geq0}\tau^{\leq0}X\simeq\tau^{\leq0}\tau^{\geq0}X\quad.
\]

我们还设置：

\[
H^{n}(X)=H^{0}(X[n])\simeq(\tau^{\geq n}\tau^{\leq n}X)[n]\enspace.
\]

\statementlabel{命题 10.1.10}让  $ X \in \text{Ob}(\mathbb{D}) $  并假设  $ X $  对于某些  $ a $  属于  $ \mathbb{D}^{\geq a} $  （分别为  $ \mathbb{D}^{\leq a} $  ）。那么  $ X $  属于  $ \mathbb{D}^{\geq 0} $  （分别为  $ \mathbb{D}^{\leq 0} $  ）当且仅当  $ H^n(X) = 0 $  为  $ n < 0 $  （分别为  $ n > 0 $  ）。

\statementlabel{证明}足以证明如果  $ X \simeq \tau^{\geq a} X $  和  $ H^a(X) = 0 $  ，则  $ X \simeq \tau^{>a} X $  。这直接从著名的三角形得出：

\[
H^{a}(X)[-a]\longrightarrow\tau^{\geq a}X\to\tau^{>a}X\xrightarrow[+1]{}\,.\quad\Box
\]

\statementlabel{命题 10.1.11}(i) 心 $ \mathcal{C} = \mathbf{D}^{\leq 0} \cap \mathbf{D}^{\geq 0} $  是一个阿贝尔范畴。

(ii) 如果  $ X' \longrightarrow X \longrightarrow X'' \xrightarrow{+1} $  是 D 中的一个显着三角形，并且如果  $ X' $  和  $ X'' $  属于  $ \mathcal{C} $  ，则 X 属于  $ \mathcal{C} $  。

(iii) 如果  $0 \to X \to Y \to Z \to 0$  是  $\mathcal{C}$  中的精确序列，则存在唯一的  $h: Z \to X[1]$  使得  $X \to Y \to Z \to X[1]$  是  $\mathbf{D}$  中的一个显着三角形。

\statementlabel{证明}(ii) 源自命题 10.1.7。

(i) 通过考虑显着的三角形  $ X \longrightarrow X \oplus Y \longrightarrow Y \xrightarrow{+1} $  ，我们得到 (ii)  $ \mathcal{C} $  是一个加法范畴。让  $ f: X \to Y $  成为  $ \mathcal{C} $  中的态射，并让我们将  $ f $  嵌入到一个独特的三角形  $ X \xrightarrow{f} Y \longrightarrow Z \xrightarrow{+1} $  中。那么，根据命题 10.1.7， $ Z $  属于  $ \mathbf{D}^{\leq 0} \cap \mathbf{D}^{\geq -1} $  。我们将证明：

\[
\left\{\begin{aligned}{H^{0}(Z)}&{{}\simeq\tau^{\geq0}Z\simeq\operatorname{C o k e r}f\qquad\mathrm{a n d}}\\ {H^{0}(Z[-1])}&{{}\simeq\tau^{\leq0}(Z[-1])\simeq\operatorname{K e r}f\;.}\\ \end{aligned}\right.
\]

为此，采用  $ W \in \mathrm{Ob}(\mathcal{C}) $  并考虑长的精确序列：

\[
\mathrm{H o m}(X[1],W)\to\mathrm{H o m}(Z,W)\to\mathrm{P}\curvearrowright\mathrm{m}(Y,W)\to\mathrm{H o m}(X,W)\qquad\mathrm{a n d}
\]

\[
\operatorname{H o m}(W,Y[-1])\to\operatorname{H o m}(W,Z[-1])\quad\triangleright\operatorname{H o m}(W,X)\to\operatorname{H o m}(W,Y)\enspace.
\]

由于  $\mathrm{Hom}(X[1], W) = \mathrm{Hom}(W, Y[-1]) = 0$  和  $\mathrm{Hom}(Z, W) \simeq \mathrm{Hom}(\tau^{\geq 0} Z, W)$  、  $\mathrm{Hom}(W, Z[-1]) \simeq \mathrm{Hom}(W, \tau^{\leq 0}(Z[-1]))$  ，（因为  $W \in \mathrm{Ob}(\mathcal{C}))$  ，我们得到 (10.1.7)。

现在让我们证明规范态射 $ \operatorname{Coim}f \to \operatorname{Im}f $  是同构。

我们将  $ Y \to \tau^{\geq 0} Z $  嵌入到一个独特的三角形  $ I \longrightarrow Y \longrightarrow \tau^{\geq 0} Z \xrightarrow{+1} $  中。那么根据命题10.1.7， $ I $ 属于 $ \mathbf{D}^{\geq 0} $ 。

我们将八面体公理应用于 $ Y \to Z \to \tau^{\geq 0}Z $ （参见（10.1.8））。

\begin{center}
\includegraphics[width=0.27\textwidth]{流行上的层_Chn-assets/img_in_image_box_412_341_736_784.png}
\end{center}

这样我们就得到了区分三角形 $ \tau^{\leq0}(Z[-1]) \longrightarrow X \longrightarrow I \xrightarrow{+1} $ 。因此  $ I $  属于  $ \mathbf{D}^{\leq0} $  ，因此  $ I \in \mathrm{Ob}(\mathcal{C}) $  。从  $ \tau^{\leq0}(Z[-1]) \simeq \mathrm{Ker} f $  开始，我们通过将 (10.1.7) 应用于区分三角形  $ \mathrm{Ker} f \longrightarrow X \longrightarrow I \xrightarrow{+1} $  发现  $ I \simeq \mathrm{Coim} f $  。类似地，将（10.1.7）应用于区分三角形 $ I \longrightarrow Y \longrightarrow \mathrm{Coker} f \xrightarrow{+1} $ ，我们得到 $ I \simeq \mathrm{Im} f $ 。这就证明了(ii)。

(iii) 由引理 10.1.5 和 (10.1.7) 得出。 ⑨

\statementlabel{命题 10.1.12}函子 $ H^0: \mathbf{D} \to \mathcal{C} $ （参见定义10.1.9）是一个上同调函子。

\statementlabel{证明}令  $X \longrightarrow Y \longrightarrow Z \xrightarrow[+1]{}$  为  $\mathbf{D}$  中的一个显着三角形。我们必须证明 $H^{0}(X) \to H^{0}(Y) \to H^{0}(Z)$  是准确的。我们将证明分解为三个步骤。

(a) 假设  $X$  、  $Y$  、  $Z$  属于  $\mathbf{D}^{\geq 0}$  ，并让我们证明  $0 \to H^0(X) \to H^0(Y) \to H^0(Z)$  是准确的。

让 $ W \in \text{Ob}(\mathcal{C}) $  。然后 $ \text{Hom}_{\mathcal{C}}(W, H^0(X)) \simeq \text{Hom}_{\mathcal{D}}(W, \tau^{\geq 0}X) \simeq \text{Hom}_{\mathcal{D}}(W, X) $ 。因此， $ \text{Hom}_{\mathcal{D}}(W, Z[-1]) = 0 $  和长精确序列

\[
\operatorname{H o m}_{\mathsf{D}}(W,Z[-1])\to\operatorname{H o m}_{\mathsf{D}}(W,X)\to\operatorname{H o m}_{\mathsf{D}}(W,Y)\to\operatorname{H o m}_{\mathsf{D}}(W,Z)
\]

给出想要的结果。

(b) 这里我们只假设 $Z \in \mathrm{Ob}(\mathbf{D}^{≥0})$ ，并且我们将证明 $0 \to H^0(X) \to H^0(Y) \to H^0(Z)$ 是准确的。对于任何  $W \in \mathrm{Ob}(\mathbf{D}^{<0})$  、  $\mathrm{Hom}_{\mathbf{D}}(W,Z) = \mathrm{Hom}_{\mathbf{D}}(W,Z[-1]) = 0$  ，都会给出  $\mathrm{Hom}_{\mathbf{D}}(W,X) \simeq \mathrm{Hom}_{\mathbf{D}}(W,Y)$  。因此 $\tau^{<0}X \to \tau^{<0}Y$  是同构。应用八面体公理（参见（10.1.9）），我们得到一个显着的三角形 $\tau^{\geq0}X \longrightarrow \tau^{\geq0}Y \longrightarrow Z \xrightarrow{+1}$ ，并且应用（a）就足够了。

\begin{center}
\includegraphics[width=0.25\textwidth]{流行上的层_Chn-assets/img_in_image_box_420_165_721_611.png}
\end{center}

(c) (b) 的双重陈述给出，如果  $ X $  属于  $ \mathbf{D}^{\leq 0} $  ，则  $ H^0(X) \to H^0(Y) \to H^0(Z) \to 0 $  是精确的。

(d) 我们处理一般情况。应用于  $ \tau^{\leq0}X \to X \to Y $  的八面体公理（参见 (10.1.10)）给出了不同的三角形：  $ \tau^{\leq0}X \longrightarrow Y \longrightarrow W \xrightarrow[+1]{} $  和  $ \tau^{>0}X \longrightarrow W \longrightarrow Z \xrightarrow[+1]{} $  。

\begin{center}
\includegraphics[width=0.25\textwidth]{流行上的层_Chn-assets/img_in_image_box_417_839_717_1283.png}
\end{center}

将 (c) 应用于第一个三角形，我们发现  $ H^0(X) \to H^0(Y) \to H^0(W) $  是精确的。

将 (b) 应用于三角形 $ W \longrightarrow Z \longrightarrow \tau^{>0}X[1]_{\substack{+1 \\ +1}} $ ，我们发现 $ 0 \to H^0(W) \to H^0(Z) $  是精确的。因此 $ H^0(X) \to H^0(Y) \to H^0(Z) $  是准确的。   $ \square $ 

\statementlabel{定义 10.1.13}令  $\mathbf{D}_i$  (  $i = 1,2$  ) 为两个赋予  $t$  结构 (  $\mathbf{D}_i^{\leq 0}$  ,  $\mathbf{D}_i^{\geq 0}$  ) 的三角范畴，并令  $\mathcal{C}_i$  为  $\mathbf{D}_i$  的核心，  $\varepsilon_i: \mathcal{C}_i \to \mathbf{D}_i$  为包含函子。让 $F: \mathbf{D}_1 \to \mathbf{D}_2$  成为三角范畴的函子。

(i) 如果 $F(\mathbf{D}_1^{\geq 0}) \subset \mathbf{D}_2^{\geq 0}$ （或 $F(\mathbf{D}_1^{\leq 0}) \subset \mathbf{D}_2^{\leq 0}$ ），有人说 $F$  是左（或右） $t$  -精确，如果 $t$  既是右又是左 $t$  -精确，则有人说 $F$  是 $t$  -精确。

(ii) 一套：

\[
{}^{p}\mathbf{F}=\mathbf{H}^{0}\circ\mathbf{F}\circ\varepsilon_{1}:\mathcal{C}_{1}\to\mathcal{C}_{2}\enspace.
\]

\statementlabel{命题 10.1.14}令 $\mathbf{D}_{1}, \mathbf{D}_{2}$  和 $F$  如定义10.1.13 中所示，并假设 $F$  是左（右） $t$  -精确的。然后：

(i) 对于 $ X \in \text{Ob}(\mathbb{D}^{\geq 0}) $ （分别为 $ \mathbb{D}^{\leq 0} $ ）， $ H^0(F(X)) \simeq^p F(H^0(X)) $ ，

(ii)  $ ^pF: \mathcal{C}_1 \to \mathcal{C}_2 $  为左（或右）精确。

\statementlabel{证明}证明左 t 精确函子的结果就足够了。

(i) 设  $ X \in \mathrm{Ob}(\mathbf{D}^{>0}) $  ，并将  $ F $  应用于区分的三角形  $ H^0(X) \longrightarrow X \longrightarrow \tau^{>0}X \xrightarrow{+1} $  。我们得到了杰出的三角形 $ F(H^0(X)) \longrightarrow F(X) \longrightarrow F(\tau^{>0}X) \xrightarrow{+1} $ ， $ F(\tau^{>0}X) $ 属于 $ \mathbf{D}_2^{>0} $ 。然后将上同调函子  $ H^0(\cdot) $  应用于该三角形即可得出结果。

(ii) 令  $0 \to X \to Y \to Z \to 0$  为  $\mathcal{C}_1$  中的精确序列。根据命题 10.1.11，它产生了一个杰出的三角形 $X \longrightarrow Y \longrightarrow Z \xrightarrow{+1} \mathrm{in} \mathbf{D}_1$ 。应用 $F$ ，我们得到一个显着的三角形 $F(X) \longrightarrow F(Y) \longrightarrow F(Z) \xrightarrow{+1}$ 。由于  $F(X)$  、  $F(Y)$  、  $F(Z)$  属于  $\mathbf{D}_2^{\geq 0}$  ，我们得到确切的序列：  $0 \to H^0(F(X)) \to H^0(F(Y)) \to H^0(F(Z))$  。然后剩下的是应用（i）。

\statementlabel{备注 10.1.15}如果 $F: \mathbf{D}_1 \to \mathbf{D}_2$  是 $t$  -精确的，则 $F$  将 $\mathcal{C}_1$  发送到 $\mathcal{C}_2$  并且 $F|_{\mathcal{C}_1}$  是精确的。此外， $F|_{\mathcal{C}_1} \simeq^p F$  和 $F(H^n(X)) \simeq H^n(F(X))$  对于任何 $X \in \mathrm{Ob}(\mathbf{D}_1)$  。

\statementlabel{备注 10.1.16}(i) 令 D 为具有示例 10.1.3 (i) 中给出的 $ t $  结构的三角范畴。那么截断函子 $ \tau^{\leq n} $ 和 $ \tau^{\geq n} $ 与备注1.7.6中给出的一致。

(ii) 令 $F: \mathcal{C}_1 \to \mathcal{C}_2$  为阿贝尔范畴的加法函子，并假设 $F$  是左精确的且 $\mathcal{C}_1$  具有足够的单射。然后 $RF: \mathbf{D}^+(\mathcal{C}_1) \to \mathbf{D}^+(\mathcal{C}_2)$  留下 $t$  -精确，其中 $\mathbf{D}^+(\mathcal{C}_i)$  上的 $t$  -结构是自然结构，由例10.1.3 (i)给出的 $\mathbf{D}(\mathcal{C}_i)$  上的 $t$  -结构诱导。

\statementlabel{命题 10.1.17}令 $\tilde{\mathbf{D}}_i$  为三角范畴， $\mathbf{D}_i$  为完整三角化子范畴， $(\mathbf{D}_i^{\leq0}, \mathbf{D}_i^{\geq0})$  为 $\mathbf{D}_i$  (  $i=1,2$  ) 上的 $t$  结构。令  $f:\tilde{\mathbf{D}}_1 \to \tilde{\mathbf{D}}_2$  和  $g: \tilde{\mathbf{D}}_2 \to \tilde{\mathbf{D}}_1$  为三角范畴的函子， $f$  为  $g$  的左伴随。假设  $f(\mathbf{D}_1) \subset \mathbf{D}_2$  （分别为  $g(\mathbf{D}_2) \subset \mathbf{D}_1$  ）和  $f|_{\mathbf{D}_1}$  是右  $t$  -精确（分别为  $g|_{\mathbf{D}_2}$  是左  $t$  -精确）。那么对于任何  $Y \in \mathrm{Ob}(\mathbf{D}_2^{\geq0})$  使得  $g(Y) \in \mathrm{Ob}(\mathbf{D}_1)$  ，（分别  $X \in \mathrm{Ob}(\mathbf{D}_1^{\leq0})$  使得  $f(X) \in \mathrm{Ob}(\mathbf{D}_2)$  ），一个有  $g(Y) \in \mathrm{Ob}(\mathbf{D}_1^{\geq0})$  ，（分别  $f(X) \in \mathrm{Ob}(\mathbf{D}_2^{\leq0}$  ））。

\statementlabel{证明}对于任何 $ X \in \text{Ob}(\mathbf{D}_1^{<0}) $ （或 $ Y \in \text{Ob}(\mathbf{D}_2^{>0}) $ ），有：

\[
\mathrm{H o m}_{{\tt D}_{1}^{>0}}(X,\tau^{<0}g(Y))\simeq\mathrm{H o m}_{{\tilde{\bf d}}_{1}}(X,g(Y))\simeq\mathrm{H o m}_{{\tilde{\bf d}}_{2}}(f(X),Y)=0\enspace,
\]

（分别 $\mathrm{Hom}_{\mathfrak{D}_2^{>0}}(\tau^{>0}f(X), Y) \simeq \mathrm{Hom}_{\mathfrak{D}_2}(f(X), Y) \simeq \mathrm{Hom}_{\mathfrak{D}_1}(X, g(Y)) = 0$ ）。因此  $\tau^{<0}g(Y) = 0$  ，（分别为  $\tau^{>0}f(X) = 0$  ）。   $\square$ 

\statementlabel{推论 10.1.18}令  $\mathbf{D}_i$  为具有  $t$  结构 (i = 1, 2) 的三角范畴，并令  $f: \mathbf{D}_1 \to \mathbf{D}_2$  和  $g: \mathbf{D}_2 \to \mathbf{D}_1$  为三角范畴的函子， $f$  为  $g$  的左伴随。那么 $f$  是右 $t$  -精确当且仅当 $g$  是左 $t$  -精确。

\subsection{实流形上的反常层}
让 $p$  成为从 $\mathbb{Z}$  到 $\mathbb{Z}$  的映射。对偶映射 $p^*$  定义为：

\[
p^{*}(n)=-p(n)-n
\]

如果 p 和 $p^{*}$  都在减少，我们称 $p$  为反常现象。因此  $p$  是一个反常的 iff：

\[
p(n)-p(n+1)=0\mathrm{o r}1\qquad\mathrm{f o r~a n y~}n~,
\]

或等价：

\[
0\leqslant p(n)-p(m)\leqslant m-n\qquad\text{for any}n,m\text{with}n\leqslant m
\]

我们通过以下方式定义 $p[k]$ ：

\[
p[k](n)=p(k+n).
\]

然后：

\[
p^{*}[k]^{*}(n)=p[k](n)+k.
\]

现在让 $X$  成为一个（实解析）流形。回想一下，  $\mathbf{D}_{w-R-c}^{b}(X)$  表示  $\mathbf{D}^{b}(X)$  的弱  $\mathbb{R}$  可构造对象的满子范畴。对于 $X$  的子集 $S$  ，我们将用 $i_{S}$  表示包含图 $S \hookrightarrow X$  。如果  $S$  是亚解析的，则其维度  $\dim S$  是明确定义的（参见 VIII §2）。如果 $S = \varnothing$ ，我们设置 $\dim S = -\infty$ 。

\statementlabel{定义 10.2.1}设 p 为反常现象。

(i)  $ ^p\mathbf{D}_{w-\mathbf{R}-c}^{\leq0}(X) $  是  $ \mathbf{D}_{w-\mathbf{R}-c}^b(X) $  的满子范畴，由满足以下条件的对象  $ F $  组成：

\[
0.2.6)\quad\operatorname{d i m}(\operatorname{s u p p}(H^{j}(F)))<k\qquad{f o r~a n y~j~a n d~k~w i t h~j>p(k)~},
\]

（或等效于  $ \dim(\operatorname{supp}\tau^{>p(k)}(F))<k $  ）。

(ii)  $ ^p\mathbf{D}_{w-\mathbf{R}-c}^{\geq0}(X) $  是  $ \mathbf{D}_{w-\mathbf{R}-c}^b(X) $  的满子范畴，由满足以下条件的对象  $ F $  组成：

\[
\begin{cases}H^{j}(i_{S}^{!}(F))=0\quad for any locally closed\\ subanalytic subset S and any j with\\ j<p(\dim S)\ .\end{cases}
\]

我们设定：

\[
{}^{p}\mathbf{D}_{\mathbf{w}-\mathbf{R}-c}^{\geq n}(X)={}^{p}\mathbf{D}_{\mathbf{w}-\mathbf{R}-c}^{\geq0}(X)\left[-n\right]\;,
\]

\[
{}^{p}\mathbf{D}_{w-\mathbf{R}-c}^{≤n}(X)={}^{p}\mathbf{D}_{w-\mathbf{R}-c}^{≤0}(X)[-n],
\]

\[
{}^{p}\mathbf{D}_{w-\mathbf{R}-c}^{0}(X)={}^{p}\mathbf{D}_{w-\mathbf{R}-c}^{\leqslant0}(X)\cap{}^{p}\mathbf{D}_{w-\mathbf{R}-c}^{\geqslant0}(X)\enspace.
\]

我们还设置：

\[
{}^{p}\mathbf{D}_{\mathbf{R}-c}^{\geq n}(X)={}^{p}\mathbf{D}_{\mathbf{w}-\mathbf{R}-c}^{\geq n}(X)\cap\mathbf{D}_{\mathbf{R}-c}^{b}(X)
\]

对于  $ ^p\mathbf{D}_{\mathbf{R}-c}^{\leq n}(X) $  、  $ ^p\mathbf{D}_{\mathbf{R}-c}^{0}(X) $  也类似。

\statementlabel{备注 10.2.2}(i)  $ ^p\mathbf{D}_{\mathbf{w}-\mathbf{R}-c}^{\geq 0}(X) $  和  $ ^p\mathbf{D}_{\mathbf{w}-\mathbf{R}-c}^{\leq 0}(X) $  的定义仅取决于  $ 0 \leq n \leq \dim X $  的  $ p(n) $  的值。

(ii) 如果  $p = 0$  ，则  $^p\mathbf{D}_{w-\mathbf{R}-c}^{≤0}(X) = \mathbf{D}_{w-\mathbf{R}-c}^{≤0}(X) = \{F \in \mathbf{D}_{w-\mathbf{R}-c}^{b}(X);\, H^j(F) = 0$  对于任何  $j > 0\}$  和  $^p\mathbf{D}_{w-\mathbf{R}-c}^{≥0}(X) = \mathbf{D}_{w-\mathbf{R}-c}^{≥0}(X) = \{F \in \mathbf{D}_{w-\mathbf{R}-c}^{b}(X);\, H^j(F) = 0$  对于任何  $j < 0\}$  。

(iii) 如果 $ p' $  是另一个反常满足， $ p(n) \leq p'(n) $  对于任何n，那么

\[
{}^{p}\mathbf{D}_{\mathbf{w}-\mathbf{R}-c}^{\leqslant0}(X)\subset{}^{p^{\prime}}\mathbf{D}_{\mathbf{w}-\mathbf{R}-c}^{\leqslant0}(X)\qquad\mathrm{a n d}\qquad{}^{p^{\prime}}\mathbf{D}_{\mathbf{w}-\mathbf{R}-c}^{\geqslant0}(X)\subset{}^{p}\mathbf{D}_{\mathbf{w}-\mathbf{R}-c}^{\geqslant0}(X)\enspace.
\]

特别是，如果  $a \leq p(n) \leq b$  对于任何  $n$  ，则：

\[
\left\{\begin{aligned}{\mathbb{D}_{w-\mathbb{R}-c}^{\leq a}(X)}&{{}\subset{}^{p}\mathbb{D}_{w-\mathbb{R}-c}^{\leq0}(X)\subset\mathbb{D}_{w-\mathbb{R}-c}^{\leq b}(X)\qquad\mathrm{a n d}}\\ {\mathbb{D}_{w-\mathbb{R}-c}^{\geq b}(X)}&{{}\subset{}^{p}\mathbb{D}_{w-\mathbb{R}-c}^{\geq0}(X)\subset\mathbb{D}_{w-\mathbb{R}-c}^{\geq a}(X).}\\ \end{aligned}\right.
\]

(iv) 假设 $F \in \mathrm{Ob}(\mathbb{D}^b(X))$  具有局部常值的上同调对象，并且这些对象是平坦的 $A$  模。然后 $F \in \mathrm{Ob}({}^p\mathbb{D}_{w-\mathbb{R}-c}^{\geq0}(X))$  如果 $H^j(F) = 0$  for  $j < p$  (暗淡 $X$ )。事实上，对于 $S$  、 $i_S^1 F \simeq i_S^1 \omega_X \otimes i_S^{-1} F \otimes \omega_S [-\dim X] \simeq \omega_S \otimes i_S^{-1} F \otimes \omega_X [-\dim X]$  和 $H^j(\omega_S) = 0$  对于 $j < -\dim S$  的任意局部封闭子解析集，根据命题9.2.2。因此  $H^j(i_S^1 F) = 0$  代表  $j < p$  (dim  $S) \leq p$  (dim  $X) + \dim X - \dim S$  .

\statementlabel{备注 10.2.3}这个定义也可以推广到 X 是亚解析空间的情况（参见练习 IX.2），并且这里陈述的大多数结果在这种推广下仍然正确。

稍后我们将证明 $ (^p\mathbf{D}_{w-\mathbf{R}-c}^{\leq0}(X),^p\mathbf{D}_{w-\mathbf{R}-c}^{\geq0}(X)) $ 形成一个t结构。为此，我们需要一些初步结果。

\statementlabel{命题 10.2.4}令 $ F \in \mathrm{Ob}(\mathbb{D}_{w-\mathbb{R}-c}^b(X)) $  和 $ X = \bigsqcup_a X_a $  为由等维层组成的亚解析分层。

(i) 假设 $ i_{X_a}^{-1} F $  对于所有 $ \alpha $  具有局部常值的上同调。那么 $ F $  属于 $ ^p \mathbf{D}_{w-R-c}^0(X) $  当且仅当 $ H^j(i_{X_a}^{-1} F) = 0 $  对于任何 $ \alpha $  和 $ j $  以及 $ j > p $  (  $ \dim X_a $  )。

(ii) 假设 $ i_{X_a}^1 $  F 对于所有 $ \alpha $  具有局部常值的上同调。那么 F 属于  $ ^p\mathbf{D}_{w-R-c}^{≥0}(X) $  当且仅当  $ H^j(i_{X_a}^1F) = 0 $  对于任何  $ \alpha $  并且 j 与  $ j < p $  (  $ dim X_\alpha $  )。

\statementlabel{证明}(i) 得出：

\[
\operatorname{d i m}(\operatorname{s u p p}(H^{j}(F)))=\operatorname*{s u p}\left\{\operatorname{d i m}X_{\alpha};H^{j}(F)|_{X_{\alpha}}\neq0\right\}\;.
\]

(ii) 根据定义，如果  $F$  属于  $^p\mathbf{D}_{w-\mathbf{R}-c}^{\geq0}(X)$  ，则  $H^j(i_{X_a}^1(F)) = 0$  属于  $j < p$  （暗淡  $X_\alpha$  ）。让我们证明相反的情况。

令 $S$  为子解析闭子集。根据假设， $i_{\chi_a}^1(F)$  集中到  $\geq p(\dim X_\alpha)$  度，并且具有局部常值的上同调层。因此，根据备注 10.2.2 (iv)， $i_{\chi_a}^1(F)$  属于  $\vec{p} \mathbf{D}_{\mathbf{w}-\mathbf{R}-c}^0(X_\alpha)$  ，其中：

\[
\tilde{p}(n)=-n+p(\dim X_{\alpha})+\dim X_{\alpha}.
\]

根据定义，这意味着  $ i_{S \cap X_\alpha}^1(F) = i_{S \cap X_\alpha}^1 i_{X_\alpha}^1 F $  集中到  $ \geq \tilde{p}(\dim(S \cap X_\alpha)) = -\dim(S \cap X_\alpha) + p(\dim X_\alpha) + \dim X_\alpha \geq p(\dim(S \cap X_\alpha)) \geq p(\dim S) $  程度。

设置  $ X_k = \bigsqcup_{\dim X_a \leq k} X_a $  ，我们有：

\[
H_{X_{k}\cap S}^{j}(F)|_{(X_{k}\cap S\setminus X_{k-1}\cap S)}\simeq\bigoplus_{\substack{\mathbf{d i m}X_{\alpha}=k}}H_{X_{\alpha}\cap S}^{j}(F)
\]

 $j < p$  的值为 0（暗淡  $S$  ）。因此，我们可以应用练习 II.25 的结果来得出  $H_{S}^{j}(F) = 0$  对于  $j < p$  的结论（暗淡  $S$  ）。   $\square$ 

乍一看，条件（10.2.6）和（10.2.7）并不对称，但我们将证明它们实际上是对称的。

为此，我们引入一个定义。我们设定：

\[
\mathrm{c o s u p p}^{j}(F)=\left\{x;H^{j}(i_{x}^{!}F)\neq0\right\}\;.
\]

请注意，如果基环  $A$  是一个字段，则  $i_x^{-1}H^j(\mathbf{D}_X F) \simeq H^{-j}(i_x^! F)^*$  ，这意味着：

\[
\overline{\cos u p p^{j}(F)}=\sup(H^{-j}(D_{X}F)).
\]

\statementlabel{命题 10.2.5}让 $ F \in \text{Ob}(\mathbb{D}_{w-\mathbb{R}-c}^b(X)) $  。那么 $ \text{cosupp}^j(F) $  是 $ X $  的子解析子集， $ F $  属于 $ ^p\mathbb{D}_{w-\mathbb{R}-c}^{\geq 0}(X) $  当且仅当：

\[
\operatorname{d i m}(\operatorname{c o s u p p}^{j}(F))<k{~f o r~a n y~}j{~a n d~}k{~w i t h~}j<p(k)+k\enspace.
\]

\statementlabel{证明}让我们进行分层 $ X = \bigsqcup_{\alpha} X_{\alpha} $ ，使得 $ i_{X_{\alpha}}^{\dagger} F $  具有局部常值的上同调。

然后我们有：

\[
i_{x}^{\downarrow}F\simeq i_{x}^{\downarrow}i_{X_{x}}^{\downarrow}F\simeq(o\mathfrak{r}_{X_{x}}\otimes i_{X_{x}}^{\downarrow}F)_{x}\;[-{\operatorname{d i m}}X_{x}]\enspace.
\]

因此：

\[
\mathrm{c o s u p p}^{j}(F)=\bigsqcup_{\alpha}\mathrm{s u p p}(H^{j-\dim X_{\alpha}}(i_{X_{\alpha}}^{!}F))\enspace.
\]

这立即表明  $ \operatorname{cosup}^{j}(F) $  是亚解析的，并且命题 10.2.4 (ii) 的 (10.2.7) 和 (10.2.10) 等价。   $ \Box $ 

请注意，条件  $ j < p(k) + k $  等效于 -j > p\textasciicircum{}\{*\}(k)。

\statementlabel{推论 10.2.6}令  $F \in \mathrm{Ob}(\mathbf{D}_{\mathrm{w}-\mathrm{R}-c}^b(X))$  并让  $X = \bigsqcup_{\alpha} X_\alpha$  成为  $\mu$  分层，使得  $\mathrm{SS}(F) \subset \bigsqcup_{\alpha} T_{X_\alpha}^* X$  。令 $Y$  为 $d$  横向子流形所有层 $X_\alpha$  。假设  $F$  属于  $^p\mathbf{D}_{\mathrm{w}-\mathrm{R}-c}^b(X)$  （分别为  $^p\mathbf{D}_{\mathrm{w}-\mathrm{R}-c}^{b\leq0}(X)$  ）。那么  $i_Y^{-1}F$  （分别为  $i_Y^1F$  ）属于  $^p[d]\mathbf{D}_{\mathrm{w}-\mathrm{R}-c}^{b\geq0}(Y)$  （分别为  $^p[d]\mathbf{D}_{\mathrm{w}-\mathrm{R}-c}^{b\leq0}(Y)$  ）。

\statementlabel{命题 10.2.7}让  $ F \in \mathrm{Ob}(^p\mathbf{D}_{\mathrm{w}-\mathrm{R}-\mathrm{c}}^{\leq 0}(X)) $  和  $ G \in \mathrm{Ob}(^p\mathbf{D}_{\mathrm{w}-\mathrm{R}-\mathrm{c}}^{\geq 0}(X)) $  。然后：

\[
H^{j}(R{\mathcal{H}}o m(F,G))=0\qquad{f o r~j<0~},
\]

\[
U\mapsto\operatorname{H o m}_{\mathsf{D}(U)}(F|_{U},G|_{U})\qquad{~i s~a~s h e a f~}.
\]

\statementlabel{证明}设置  $ S = \bigcup_{j<0} \operatorname{supp}(H^j(R \mathcal{H} o m(F,G))) $  ，并假设  $ S \neq \emptyset $  。我们将得出一个矛盾。

对于 j < 0，我们有：

\[
\begin{aligned}{i_{\mathsf{S}}^{!}H^{j}(R\mathcal{H}o m(F,G))}&{{}\simeq H^{j}(i_{\mathsf{S}}^{!}R\mathcal{H}o m(F,G))}\\ {}&{{}\simeq H^{j}(R\mathcal{H}o m(i_{\mathsf{S}}^{-1}F,i_{\mathsf{S}}^{!}G))\;.}\\ \end{aligned}
\]

让 $ k = \dim S $  。然后  $ i_S^1 G $  被集中到度  $ \geq p(k) $  。另一方面， $ F $  的假设意味着  $ S' = \bigcup_{\text{def}} j > p(k) $  \textbackslash{}textit\{supp\}(  $ H^j(i_S^{-1}F) $  ) 具有维度  $ < k $  。因此 $ S \setminus S' \neq \emptyset $ 。由于 $ i_S^{-1} F|_{S \setminus S'} $  集中到度数 $ \leq p(k) $  ， $ R \mathcal{H} o m(i_S^{-1}F, i_S^1 G)|_{S \setminus S'} $  集中到度数 $ \geq 0 $  。这是一个矛盾。

自  $ \mathrm{Hom}_{D(U)}(F|_{U}, G|_{U}) \simeq H^0(U; R\mathcal{H}om(F, G)) \simeq \Gamma(U; H^0(R\mathcal{H}om(F, G))) $  以来，属性 (10.2.13) 源自 (10.2.12)。   $ \square $ 

我们现在可以证明本节的主要结果。

\statementlabel{定理 10.2.8}对于任何反常的  $p$  ，  $(^p\mathbb{D}_{\mathrm{w}-\mathbb{R}-c}^{\leq0}(X),^p\mathbb{D}_{\mathrm{w}-\mathbb{R}-c}^{\geq0}(X))$  都是  $\mathbb{D}_{\mathrm{w}-\mathbb{R}-c}^b(X)$  上的  $t$  结构。此外，如果  $A$  是诺特式的，则  $(^p\mathbb{D}_{\mathbb{R}-c}^{\leq0}(X),^p\mathbb{D}_{\mathbb{R}-c}^{\geq0}(X))$  是  $\mathbb{D}_{\mathbb{R}-c}^b(X)$  上的  $t$  结构。

\statementlabel{证明}首先我们处理  $ \mathbf{D}_{w-\mathbf{R}-c}^{b}(X) $  的情况。我们必须证明：

(a) 对于  $ F \in \mathrm{Ob}(^p\mathbf{D}_{\mathrm{w}-\mathrm{R}-\mathrm{c}}^{≤0}(X)) $  和  $ G \in \mathrm{Ob}(^p\mathbf{D}_{\mathrm{w}-\mathrm{R}-\mathrm{c}}^{≥1}(X)) $  、 $ \mathrm{Hom}_{\mathbf{D}^b(X)}(F,G) = 0 $  ，

(b)  $ {}^{p}\mathbf{D}_{w-\mathbf{R}-c}^{\leq0}(X) \subset {}^{p}\mathbf{D}_{w-\mathbf{R}-c}^{\leq1}(X) $  和 $ {}^{p}\mathbf{D}_{w-\mathbf{R}-c}^{\geq0}(X) \supset {}^{p}\mathbf{D}_{w-\mathbf{R}-c}^{\geq1}(X) $  ，

(c) 对于任何属于  $\mathbf{D}_{w-\mathbb{R}-c}^{b}(X)$  的  $F$  ，存在一个与  $F'\in\mathrm{Ob}({}^{p}\mathbf{D}_{w-\mathbb{R}-c}^{\leq0}(X))$  和  $F''\in\mathrm{Ob}({}^{p}\mathbf{D}_{w-\mathbb{R}-c}^{\geq1}(X))$  不同的三角形  $F'\longrightarrow F\longrightarrow F''\xrightarrow[+1]{}$  。

由于（a）由前面的命题得出，并且（b）是显而易见的，所以仍然需要证明（c）。

让我们选择  $ X $  的递增  $ \mu $  过滤  $ \{X_k\}_{k} $ （参见定义 8.3.25），使得：

\[
\left\{\begin{array}{l}X_{k}\backslash X_{k-1}{~i s~a~k-d i m e n s i o n a l~m a n i f o l d~f o r~a n y~}k\\ {a n d~}F|_{X_{k}\backslash X_{k-1}}{~h a s~l o c a l l y~c o n s t a n t~c o h o m o l o g y~}.{}\end{array}\right.
\]

考虑条件：

\[
\begin{cases}\text{存在 }\mathbf{D}^{b}_{w-\mathrm{R}-c}(X\setminus X_k)\text{ 中的一个可区分三角形；}\\F^{\prime}\longrightarrow i_{X\setminus X_{k}}^{-1}F\longrightarrow F^{\prime \prime}\xrightarrow[+1]{}\text{其中 }F^{\prime}\in\mathrm{Ob}({}^{p}\mathbf{D}^{0}_{w-\mathrm{R}-c}(X\setminus X_k))\\\text{且 }F^{\prime\prime}\in\mathrm{Ob}({}^{p}\mathbf{D}^{0}_{w-\mathrm{R}-c}(X\setminus X_k)),\ F^{\prime}|_{X_j\setminus X_{j-1}}\text{ 与}\\F^{\prime\prime}|_{X_j\setminus X_{j-1}}\text{ 对任意 }j\geq k\text{ 具有局部常值上同调。}\end{cases}
\]

由于  $ (10.2.15)_k $  满足  $ k \gg 0 $  ，因此在假设  $ (10.2.15)_k $  时证明  $ (10.2.15)_{k-1} $  就足够了。让 $ F' \longrightarrow i_{X \setminus X_k}^{-1} F \longrightarrow F'' \xrightarrow{+1} be $  成为一个与 $ (10.2.15)_k $  中相同的三角形。令  $ j: X \setminus X_k \subsetneq X \setminus X_{k-1} $  为开放嵌入， $ i: X_k \setminus  X_{k-1} \hookrightarrow X \setminus X_{k-1} $  封闭嵌入。态射  $ F' \to i_{X \setminus X_k}^{-1} F $  给出  $ j_! F' \to i_{X \setminus X_{k-1}}^{-1} F $  。我们将这个态射嵌入到一个独特的三角形中：

\[
j_{!}F^{\prime}\longrightarrow i_{X\setminus X_{k-1}}^{-1}F\to G\xrightarrow[+1]{}~.
\]

现在我们有一个态射链 $ \tau^{\leq p(k)}i_*i^!G\to i_*i^!G\to G $ 。我们将构图嵌入到一个独特的三角形中：

\[
\tau^{\leqslant p(k)}i_{*}i^{!}G\longrightarrow G\longrightarrow\widetilde{F}^{\prime\prime}\xrightarrow[{}+1{}]{}
\]

最后我们将组合 $ i_{X \setminus X_{k-1}}^{-1} F \to G \to \widetilde{F} $ 嵌入到一个独特的三角形中：

\[
\widetilde{F}^{\prime}\longrightarrow i_{X\setminus X_{k-1}}^{-1}F\longrightarrow\widetilde{F}^{\prime\prime}\xrightarrow[{}+1]{}.
\]

我们将展示：

\[
\widetilde{F}^{\prime\prime}\in\mathrm{O b}(^{p}\mathbb{D}_{\mathbf{w}-\mathbb{R}-c}^{\geq1}(X\backslash X_{k-1}))~,
\]

\[
\widetilde{F}^{\prime}\in\mathrm{O b}(^{p}\mathbf{D}_{\mathbf{w}-\mathbf{R}-c}^{\leq0}(X\backslash X_{k-1}))~.
\]

通过构造， $ j^{-1}\tilde{F}'' \simeq F'' $  和 $ j^{-1}\tilde{F}' \simeq F' $  。因此，根据命题 10.2.4，足以表明：

\[
i^{\prime}\tilde{F}\mathrm{~i s~c o n c e n t r a t e d~t o~d e g r e e s}\geq p(k)+1,
\]

\[
i^{-1}\tilde{F}^{\prime}\text{is concentrated to degrees}\leq p(k).
\]

通过将函子  $i^{\prime}$  应用于区分三角形 (10.2.17)，我们得到区分三角形：

\[
\tau^{\lessdot p(k)}i^{!}G\longrightarrow i^{!}G\longrightarrow\to i^{!}\widetilde{F}^{\prime\prime}\xrightarrow[{}+1{}]{}.
\]

因此  $ i^! \tilde{F}'' \simeq \tau^{\geq p(k)+1} i^! G $  ，显示 (10.2.21)。

现在让我们将八面体公理 (I §4) 应用于区分的三角形 (10.2.16)、(10.2.17)、(10.2.18)。 （参见（10.2.24））：

\begin{center}
\includegraphics[width=0.34\textwidth]{流行上的层_Chn-assets/img_in_image_box_360_1271_769_1723.png}
\end{center}

我们得到一个显着的三角形：

(10.2.25)

\[
j_{!}F^{\prime}\longrightarrow\widetilde{F}^{\prime}\longrightarrow\tau^{\leqslant p(k)}i_{*}i^{!}G\xrightarrow[+1]{}\;.
\]

因此 $ i^{-1}\tilde{F}' \stackrel{...}{\to} \tau^{\leq p(k)} i^! G $  。这显示 (10.2.22)。

由于可以轻松检查 $ \tilde{F}' $  和 $ \tilde{F}'' $  对于任何 $ j \geq k $  在 $ X_j \setminus X_{j-1} $  上具有局部常值上同调，因此归纳继续进行。如果  $ A $  是诺特式的，则同样的证明也适用于  $ \mathbf{D}_{\mathbf{R}-c}^b(X) $  。   $ \square $ 

在  $X)$  中打开的函子  $U \mapsto^p \mathbf{D}_{w-R-c}^0(U)$  、  $(U$  的行为就像一个层。更准确地说，我们有：

\statementlabel{命题 10.2.9}$ ^p\mathbf{D}_{w-\mathbf{R}-c}^0(X) $  是一个堆栈，即给定一个开覆盖  $ X = \bigcup_{i \in I} U_i $  、  $ F_i \in \mathrm{Ob}(^p\mathbf{D}_{w-\mathbf{R}-c}^0(U_i)) $  和同构：  $ f_{ij} : F_j|_{U_{ij}} \to F_i|_{U_{ij}} $  满足余循环条件  $ (f_{ij}|_{U_{ijk}}) \circ (f_{jk}|_{U_{ijk}}) = f_{ik}|_{U_{ijk}} $  对于任何  $ i $  、  $ j $  、  $ k $  ，存在  $ F \in \mathrm{Ob}(^p\mathbf{D}_{w-\mathbf{R}-c}^0(X)) $  和  $ f_i : F|_{U_i} \to F_i $  ，这样 $ f_{ij} \circ f_j|_{U_{ij}} = f_i|_{U_{ii}} $ 。

此外， $(F,\{f_{i}\}_{i})$  系列在同构方面是唯一的。 （这里我们设置 $U_{ij}=U_{i}\cap U_{j}$ 和 $U_{ijk}=U_{i}\cap U_{j}\cap U_{k}$ 。）

\statementlabel{证明}(a) 唯一性源自命题 10.2.7。

(b) 我们首先要证明当 $I$  有限时 $(F, \{f_i\}_{i})$  的存在性。对 $\#I$ 进行归纳论证，只要证明 $\#I = 2$ 时的结果，说 $I = \{1,2\}$ 即可。在这种情况下，通过一个可区分的三角形来定义  $F$  就足够了：

\[
i_{U_{12}!}(F_{1}|_{U_{12}})\longrightarrow i_{U_{1}}!F_{1}\oplus i_{U_{2}}!F_{2}\longrightarrow F\xrightarrow[+1]{}
\]

其中态射  $ i_{U_{12}}!(F_1|_{U_{12}}) \to i_{U_2}! F_2 $  由  $ f_{21} $  给出。

(c) 现在我们处理一般情况。通过(a)和(b)，我们可以假设 $ I = \mathbb{N} $ 和 $ U_n \subset U_{n+1} $ 对于所有 $ n $ 。用单射复形  $ I_n \in \mathrm{Ob}(\mathbf{C}^+\left(\mathfrak{Mod}(A_n)\right)) $  表示  $ F_n \in \mathrm{Ob}(^p \mathbf{D}_{w-R-c}^0(U_n)) $  。然后同构  $ F_n \rightleftarrows F_{n+1}|_{U_n} $  给出同态  $ I_n \to I_{n+1}|_{U_n} $  ，因此同态  $ \varphi_n: i_{U_n} | I_n \to i_{U_{n+1}} | I_{n+1} $  。现在让 $ F = \lim_{n} I_n $  。足以证明自然态射 $ I_n \to i_{U_n}^{-1} F $ 是拟同构。但这是从以下事实得出的：对于  $ x \in U_n $  ，  $ (I_n)_x $  与  $ m \geq n $  的  $ (I_m)_x $  准同构。   $ \square $ 

我们将用 $p_{\tau} \leq n$  和 $p_{\tau} \geq n$  表示与 $t$  结构相关的截断函子（ $p \mathbf{D}_{w-\mathbf{R}-c}^{ \leq 0}(X)$  、 $p \mathbf{D}_{w-\mathbf{R}-c}^{ \geq 0}(X)$  ）。因此：

\[
\left\{\begin{aligned}{^{p}\tau^{\leqslant n}}&{{}:\mathbf{D}_{w-\mathbf{R}-c}^{b}(X)\to{^{p}\mathbf{D}_{w-\mathbf{R}-c}^{\leqslant n}}(X)\enspace,}\\ {^{p}\tau^{\geqslant n}}&{{}:\mathbf{D}_{w-\mathbf{R}-c}^{b}(X)\to{^{p}\mathbf{D}_{w-\mathbf{R}-c}^{\geqslant n}}(X)\enspace.}\\ \end{aligned}\right.
\]

我们将用 $^pH^n$  来表示关于相同 $t$  结构的 $n$  上同调。因此：

\[
{}^{p}H^{n}:\mathbf{D}_{w-\mathbb{R}-c}^{b}(X)\to{}^{p}\mathbf{D}_{w-\mathbb{R}-c}^{0}(X)\enspace.
\]

当 $ A $  是诺特矩阵时，我们对 $ ^p\mathbf{D}_{R-c}^0(X) $  保留相同的符号。

为了结束本节，让我们研究与反常行为相关的 t 结构的功能特性。

\statementlabel{命题 10.2.10}令 Y 为 X 的局部闭子解析子集。

(i)  $ i_{Y!}i_{Y}^{-1} $  将  $ ^p\mathbf{D}_{w-\mathbf{R}-c}^{\leq0}(X) $  发送给自身（即这个函子是正确的）。

(ii)  $ Ri_{Y*}i_Y^1 $  将  $ ^p\mathbf{D}_{w-R-c}^{≥0}(X) $  发送给自身（即该函子保持 t 精确）。

\statementlabel{证明}(i) 由于  $ i_{Y!} $  和  $ i_{Y}^{-1} $  是精确函子，因此它们可转换为函子  $ H^{j}(\cdot) $  。因此，结果来自  $ \operatorname{supp}(i_{Y!}i_{Y}^{-1}(F)) \subset \operatorname{supp}(F) $  。

(ii) 其中一个有  $ (i_S^1 i_Y* i_Y^1 F)_{Y \cap S} \simeq i_{Y \cap S}^1 F $  。因此，结果来自定义。   $ \square $ 

\statementlabel{命题 10.2.11}令 $f: Y \to X$  为流形态射，并令 $d$  为整数。假设  $\dim f^{-1}(x) \leq d$  对于任何  $x \in X$  。然后：

(i)  $ f^{-1} $  将 $ ^p\mathbf{D}_{w-\mathbb{R}-c}^{≤0}(X) $  发送至 $ ^{p[-d]}\mathbf{D}_{w-\mathbb{R}-c}^{≤0}(Y) $  。

(ii)  $ f^! $  将 $ ^p\mathbf{D}_{w-\mathbb{R}-c}^{\geq 0}(X) $  发送至 $ ^{p[-d]}\mathbf{D}_{w-\mathbb{R}-c}^{\geq -d}(Y) $  。

\statementlabel{证明}(i) 得出：

\[
\begin{aligned}{\operatorname{d i m}(\operatorname{s u p p}(H^{j}(f^{-1}F)))}&{{}=\operatorname{d i m}(f^{-1}(\operatorname{s u p p}(H^{j}(F))))}\\ {}&{{}\leqslant\operatorname{d i m}(\operatorname{s u p p}(H^{j}(F)))+d.}\\ \end{aligned}
\]

(ii) 由命题 10.2.5 得出以下关系：

\[
\begin{aligned}{\operatorname{c o s u p p}^{j}(f^{!}F)}&{{}=\{y\in Y;H^{j}(i_{y}^{!}f^{!}F)\neq0\}}\\ {}&{{}=\{y\in Y;H^{j}(i_{f(y)}^{!}F)\neq0\}}\\ {}&{{}=f^{-1}(\operatorname{c o s u p p}^{j}(F))~,}\\ \end{aligned}
\]

这意味着

\[
\operatorname{d i m}(\operatorname{c o s u p p}^{j}(f^{!}F))\leqslant\operatorname{d i m}(\operatorname{c o s u p p}_{!}^{j}(F))+d\enspace.\quad\Box
\]

\statementlabel{命题 10.2.12}令 $f: Y \to X$  为流形态射，并令 $d$  为整数。假设  $\dim f^{-1}(x) \leq d$  对于任何  $x \in X$  。

(i) 设  $ G \in \text{Ob}(^p \mathbb{D}_{\mathbf{w} - \mathbb{R} - c}^0(Y)) $  并假设  $ Rf_1 G \in \text{Ob}(\mathbb{D}_{\mathbf{w} - \mathbb{R} - c}^b(X)) $  。然后 $ Rf_1 G \in \text{Ob}(^p [d] \mathbb{D}_{\mathbf{w} - \mathbb{R} - c}^d(X)) $ 。

(ii) 设  $ G \in \text{Ob}(^p \mathbb{D}_{\mathrm{w}-\mathbb{R}-c}^{ \geq 0}(Y)) $  并假设  $ R f_* G \in \text{Ob}(\mathbb{D}_{\mathrm{w}-\mathbb{R}-c}^{b}(X)) $  。然后 $ R f_* G \in \text{Ob}(^{p[d]} \mathbb{D}_{\mathrm{w}-\mathbb{R}-c}^{ \geq 0}(X)) $ 。

\statementlabel{证明}这直接源于前面的命题，命题 10.1.17 以及  $ f^{-1} $  和  $ Rf_* $  以及  $ Rf_1 $  和  $ f^1 $  是伴随函子的事实。   $ \square $ 

\statementlabel{命题 10.2.13}假设基环 $A$  是一个字段。然后对偶函子 $\mathrm{D}_X$ （参见定义3.1.16）分别将 $^p\mathbf{D}_{\mathbf{R}-c}^0(X)$ 和 $^p\mathbf{D}_{\mathbf{R}-c}^{\geq 0}(X)$ 发送到 ${}^p\mathbf{D}_{\mathbf{R}-c}^{\geq 0}(X)$ 和 ${}^p\mathbf{D}_{\mathbf{R}-c}^{\leq 0}(X)$ 。

\statementlabel{证明}这是从  $ i_x^{-1}(D_X F) \simeq \text{Hom}(i_x^1 F, A) $  和  $ i_x^1(D_X F) \simeq \text{Hom}(i_x^{-1} F, A) $  以及命题 10.2.5 得出的。   $ \square $ 

最后我们注意到以下有用的结果。

\statementlabel{命题 10.2.14}令 Y 为 X 的闭子流形，并令 F 属于  $ \mathbf{D}_{w-R-c}^{b}(X) $  。

(i)  $ F \in \text{Ob}(^p\mathbb{D}_{\mathbf{w}-\mathbb{R}-c}^{\leq0}(X)) $  当且仅当  $ F|_{X \setminus Y} \in \text{Ob}(^p\mathbb{D}_{\mathbf{w}-\mathbb{R}-c}^{\leq0}(X \setminus Y)) $  和  $ i_Y^{-1} F \in \text{Ob}(^p\mathbb{D}_{\mathbf{w}-\mathbb{R}-c}^{\leq0}(Y)) $  。

(ii)  $ F \in \mathrm{Ob}(^p\mathbb{D}_{w-\mathbb{R}-c}^{\geq0}(X)) $  当且仅当  $ F|_{X\setminus Y} \in \mathrm{Ob}(^p\mathbb{D}_{w-\mathbb{R}-c}^{\geq0}(X\setminus Y)) $  和  $ i_Y^! F \in \mathrm{Ob}(^p\mathbb{D}_{w-\mathbb{R}-c}^{\geq0}(Y)) $  。

\statementlabel{证明}(i) 源自支持 $ H^j(F) = \overline{\operatorname{supp} H^j(F|_{X\setminus Y})} \cup \operatorname{supp} H^j(i_Y^{-1}F) $ 。

(ii) 从  $ \cosupp^j(F) = \cosupp^j(F|_{X\setminus Y}) \cup \cosupp^j(i_Y^! F) $  得出。 □

\subsection{复流形上的反常层}
令 X 为复流形。我们设定：

\[
{}^{p}\pmb{\mathrm{D}}_{w-\mathrm{C}-c}^{≤0}(X)={}^{p}\pmb{\mathrm{D}}_{w-\mathrm{R}-c}^{≤0}(X)\cap\pmb{\mathrm{D}}_{w-\mathrm{C}-c}(X)\ ,
\]

我们同样定义  $^p\mathbf{D}_{w-c-c}^{≥0}(X)$  、  $^p\mathbf{D}_{w-c-c}^{0}(X)$  、  $^p\mathbf{D}_{c-c}^{≤0}(X)$  、  $^p\mathbf{D}_{c-c}^{≥0}(X)$  、  $^p\mathbf{D}_{c-c}^{0}(X)$  等。

在命题 10.2.4 中，我们可以选择  $ X_\alpha $  的复流形，因此  $ ^p\mathbf{D}_{w-c-c}^{\leq 0}(X) $  等仅取决于偶数处的  $ p $  的值。在续集中，我们选择 $ p $ ，这样 $ p = p^* $ 就在 $ 2\mathbb{Z} $ 上，即：

\[
p(n)=-n/2\qquad\mathrm{f o r~e v e n~}n.
\]

我们称 $p$ 为中等偏态。

公约 10.3。在本节中，除非另有说明，所有流形和流形态射都是复解析的。对于复数解析集 S，dim S 表示 S 的复数维度，我们将实数维度写为 $\dim_{\mathbf{R}}$ 。

因此，我们有  $\mathbf{D}_{w-C-c}^{b}(X)$  的对象  $F$  ：

\[
\begin{cases}{F\in\mathbf{O b}({}^{p}\mathbf{D}_{w-\mathbb{C}-c}^{\leq0}(X)){~i f~a n d~o n l y~i f~}\operatorname{d i m}\operatorname{s u p p}(H^{j}(F))\leqslant-j}\\ {{f o r\quad a l l~j~}.}\\ \end{cases}
\]

\[
\left\{\begin{matrix}{F\in\operatorname{O b}(^{p}\mathbf{D}_{\mathbf{w}-\mathbf{C}-\mathbf{c}}^{\geq0}(X)){~i f~a n d~o n l y~i f~f o r~a n y~l o c a l l y~c l o s e d}}\\ {{c o m p l e x~a n a l y t i c~s u b s e t~}S{~o f~}X,H_{S}^{i}(F)|_{S}=0{~f o r~}j<-\operatorname{d i m}S.}\\ \end{matrix}\right.
\]

（最后一个条件相当于  $ \dim \cosupp^j(F) \leq j $  。）

选择一个复杂的  $ \mu $  -分层  $ X = \bigsqcup_{\alpha} X_{\alpha} $  使得  $ \mathrm{SS}(F) \subset \bigsqcup_{\alpha} T_{X_{\alpha}}^{*} X $  。然后：

\[
\begin{aligned}F\in\operatorname{Ob}({}^{p}\mathbf{D}_{w-\mathbb{C}-c}^{\leq0}(X))\quad\Longleftrightarrow\quad F_x\in\operatorname{Ob}(\mathbf{D}^{\leq-\dim X_{\alpha}}(\mathfrak{Mod}(A)))\\\text{for any }x\in X_\alpha,\text{ any }\alpha.\end{aligned}
\]

\[
\begin{aligned}F\in\operatorname{Ob}({}^{p}\mathbf{D}_{w-\mathbb{C}-c}^{\geq0}(X))\quad\Longleftrightarrow\quad (R\varGamma_{\{x\}}(F))_x\in\operatorname{Ob}(\mathbf{D}^{\geq\dim X_\alpha}(\mathfrak{Mod}(A)))\\\text{for any }x\in X_\alpha,\text{ any }\alpha.\end{aligned}
\]

\statementlabel{定义 10.3.1}$ ^p\mathbf{D}_{C-c}^0(X) $  的对象称为反常层。

请注意，一般来说，反常捆不是捆，但鉴于命题 10.2.9 中所述的局部属性，人们仍然称它们为“捆”。

示例 10.3.2。设Y是X的闭子流形。那么 $ A_{Y}[\dim Y] $ 是反常的。

我们可以将命题10.2.11和10.2.12翻译如下。

\statementlabel{命题 10.3.3}令  $d$  为整数， $f: Y \to X$  为复流形的态射，使得  $\dim f^{-1}(x) \leq d$  对于任何  $x$  。然后：

(i)  $ f^{-1} $  将 $ ^p\mathbf{D}_{w-C-c}^{≤0}(X) $  发送至 $ ^p\mathbf{D}_{w-C-c}^{≤d}(Y) $  。

(ii)  $ f^! $  将 $ ^p\mathbf{D}_{w-C-c}^{\geq0}(X) $  发送至 $ ^p\mathbf{D}_{w-C-c}^{\geq-d}(Y) $  。

(iii) 对于 $ G \in \text{Ob}(^p\mathbb{D}_{w-\mathbb{C}-c}^{≤0}(Y)) $  ，如果 $ Rf_i G $  属于 $ \mathbb{D}_{w-\mathbb{C}-c}^b(X) $  ，则 $ Rf_i G $  属于 $ ^p\mathbb{D}_{w-\mathbb{C}-c}^{≤d}(X) $  。

(iv) 对于 $G \in \mathrm{Ob}(^p\mathbb{D}_{w-C-c}^{≥0}(Y))$  ，如果 $Rf_*G$  属于 $\mathbb{D}_{w-C-c}^b(X)$  ，则 $Rf_*G$  属于 $^p\mathbb{D}_{w-C-c}^{≥-d}(X)$  。

此外，定理 10.2.8 的结果在复杂情况下仍然有效，通过在该定理的证明中选择  $X$  的过滤  $\{X_j\}$  ，使得所有  $X_j$  都是复解析的。

\statementlabel{定理 10.3.4}(  $^{p}\mathbf{D}_{\mathbf{w}-\mathbf{C}-c}^{≤0}(X)$  ,  $^{p}\mathbf{D}_{\mathbf{w}-\mathbf{C}-c}^{≥0}(X)$  ) 是  $\mathbf{D}_{\mathbf{w}-\mathbf{C}-c}^{b}(X)$  上的  $t$  结构。此外，如果基环  $A$  是诺特环，则 (  $^{p}\mathbf{D}_{\mathbf{C}-c}^{≤0}(X)$  ,  $^{p}\mathbf{D}_{\mathbf{C}-c}^{≥0}(X)$  ) 是  $\mathbf{D}_{\mathbf{C}-c}^{b}(X)$  上的  $t$  结构。

以下是命题 10.2.13 的结果。

\statementlabel{命题 10.3.5}如果基环  $ A $  是一个域，则对偶函子  $ \mathrm{D}_X $  会互换  $ ^p\mathbf{D}_{\mathbb{C}-c}^{\geq 0}(X) $  和  $ ^p\mathbf{D}_{\mathbb{C}-c}^{\leq 0}(X) $  。

下面的命题也很容易证明。

\statementlabel{命题 10.3.6}令 X 和 Y 为复流形。

\[
\begin{array}{r l r l r l r l}{\mathrm{(i)}}&{{}{I f}}&{F\in\mathrm{O b}(^{p}\mathbb{D}_{\mathsf{w}-\mathbb{C}-\mathsf{c}}^{\leqslant0}(X))}&{{a n d}}&{G\in\mathrm{O b}(^{p}\mathbb{D}_{\mathsf{w}-\mathbb{C}-\mathsf{c}}^{\leqslant0}(Y)),}&{{t h e n}}&{F\boxtimes^{L}G\in}\\ {}&{{}}&{\mathrm{O b}(^{p}\mathbb{D}_{\mathsf{w}-\mathbb{C}-\mathsf{c}}^{\leqslant0}(X\times Y)).}\\ \end{array}
\]

(ii) 如果基环是一个域，并且如果  $ F \in \text{Ob}(^p\text{D}_{\mathbb{C}-\mathbb{c}}^{\geq 0}(X)) $  和  $ G \in \text{Ob}(^p\text{D}_{\mathbb{C}-\mathbb{c}}^{\geq 0}(Y)) $  则  $ F \boxtimes G \in \text{Ob}(^p\text{D}_{\mathbb{C}-\mathbb{c}}^{\geq 0}(X \times Y)) $  。

(iii) 如果  $F \in \mathrm{Ob}(^p\mathrm{D}_{w-C-c}^{≤0}(X))$  和  $G \in \mathrm{Ob}(^p\mathrm{D}_{w-C-c}^{≥0}(Y))$  ，则  $R \mathcal{H}om(q_1^{-1}F, q_2^!G) \in \mathrm{Ob}(^p\mathrm{D}_{w-C-c}^{≥0}(X \times Y))$  。

（这里 $q_1$ 和 $q_2$ 分别是从 $X \times Y$ 到 $X$ 和 $Y$ 的投影。）

\statementlabel{证明}(i) 我们选择复杂的分析 $ \mu $  -分层 $ X = \bigsqcup_{\alpha} X_{\alpha} $  和 $ Y = \bigsqcup_{\beta} Y_{\beta} $  使得 $ \mathrm{SS}(F) \subset \bigsqcup_{\alpha} T_{X_{\alpha}}^{*} X $  和 $ \mathrm{SS}(G) \subset \bigsqcup_{\beta} T_{Y_{\beta}}^{*} Y $  。那么结果来自命题 10.2.4，因为：

\[
i_{X_{\alpha}\times Y_{\beta}}^{-1}\biggl(F\boxtimes G\biggr)\simeq i_{X_{\alpha}}^{-1}F\boxtimes i_{Y_{\beta}}^{-1}G\enspace.
\]

(ii) 由 (i) 得出，并且事实是  $ F \boxtimes G \simeq D_{X \times Y}(D_X F \boxtimes D_Y G) $  。

(iii) 一个人有：

\[
\begin{array}{r}{i_{X_{\alpha}\times Y_{\beta}}^{!}R\mathcal{H}o m(q_{1}^{-1}F,q_{2}^{!}G)\simeq R\mathcal{H}o m(i_{X_{\alpha}}^{-1}F\boxtimes A_{Y_{\beta}},\omega_{X_{\alpha}}\boxtimes i_{Y_{\beta}}^{!}G)\enspace.}\end{array}
\]

根据假设， $ i_{X_{\alpha}}^{-1}F $  集中到 $ \leqslant -\dim X_{\alpha} $  度， $ \omega_{X_{\alpha}}\boxtimes i_{Y_{\beta}}^{!}G $  集中到 $ \geqslant -2\dim X_{\alpha}-\dim Y_{\beta} $  度。因此  $ i_{X_{\alpha}\times Y_{\beta}}^{!}R\mathcal{H}om(q_{1}^{-1}F,q_{2}^{!}G) $  集中到  $ \geqslant -\dim(X_{\alpha}\times Y_{\beta}) $  度。   $ \square $ 

现在我们要介绍一种微局部反常现象。

\statementlabel{定义 10.3.7}$ \mathbf{D}_{\mathbf{w}-\mathbf{C}-c}^{b}(X) $ （或  $ {}^{u}\mathbf{D}_{\mathbf{w}-\mathbf{C}-c}^{b\geq n}(X) $  ）是  $ \mathbf{D}_{\mathbf{w}-\mathbf{R}-c}^{b}(X) $  的完整子类，由对象 F 组成，使得 F 的类型 L 在 SS(F) 的任何非奇异点处移位 0 满足  $ H^{j}(L)=0 $  对于 j>n-dim X（或 j<n-dim X）。

根据命题 7.5.9 和练习 A.7（另见练习 VII.4）， $ F \in \mathrm{Ob}(\mathbb{D}_{w-C-c}^b(X)) $  的以下三个条件是等效的。

\[
{ \bf F } {  b e l o n g s }  t o  { { } ^ { \mu } } { \bf D } _ { w - \mathbb { C } - \mathbb { c } } ^ { \leqslant 0 } ( X ) \left( {  r e s p . } {  } { { } ^ { \mu } } { \bf D } _ { w - \mathbb { C } - \mathbb { c } } ^ { \geqslant 0 } ( X ) \right) .
\]

\[
\begin{cases}{{F o r~e v e r y~n o n-s i n g u l a r~p o i n t~}p{~o f~}\mathbf{S S}(F){~s u c h~t h a t}}\\ {\pi:\mathbf{S S}(F)\to X{~h a s~c o n s t a n t~r a n k~o n~a~n e i g h b o r h o o d}}\\ {{o f~}p{,t h e r e~e x i s t s~a~s u b m a n i f o l d~}Y{~a n d}}\\ {L\in\operatorname{O b}(\mathbf{D}^{b}(\mathfrak{R o b}(A))){~s u c h~t h a t~}F\simeq L_{Y}[\operatorname{d i m}Y]{~i n~}\mathbf{D}^{b}(X;p)}\\ {{a n d~}H^{j}(L)=0{~f o r~}j>0{({r e s p.~}j<0)~}.}\\ \end{cases}
\]

(10.3.7) 的断言对于 SS(F) 的任何不可约分量上的某个点  $p$  都是正确的。

本节的目的之一是证明  $ \mu \mathbf{D}_{w-C-c}^{\leq 0}(X) = {}^{p}\mathbf{D}_{w-C-c}^{\leq 0}(X) $  和  $ \mu \mathbf{D}_{w-C-c}^{\geq 0}(X) = {}^{p}\mathbf{D}_{w-C-c}^{\geq 0}(X) $  。

首先，我们需要斯坦因流形上反常层的消失定理。我们不会在这里回顾 Stein 流形的理论（例如，参见 Hörmander [1]），但我们记得对于某些  $ N $ ，这样的流形可以全纯嵌入到  $ \mathbb{C}^{N} $  中作为闭合子流形。

\statementlabel{定理 10.3.8}令 X 为斯坦因流形。

\[
\begin{array}{r}{\mathrm{(i)}\mathrm{~F o r~a n y~}F\in\mathrm{O b}(\mathbf{\mu}\mathbf{D}_{w-c-c}^{\leq0}(X)),H^{j}(X;F)=0\mathrm{~f o r~}j>0.}\end{array}
\]

\[
\mathrm{(i i)}~{F o r~a n y}~F\in\operatorname{O b}(\mathbf{\mu}\mathbf{D}_{w-\mathbb{C}-c}^{\geq0}(X)),~H_{c}^{j}(X;F)=0~{f o r~}j<0.
\]

依赖于莫尔斯理论的证明类似于命题5.4.20。

\statementlabel{证明}我们将  $X$  作为闭合子流形嵌入到  $\mathbb{C}^N$  中。令  $F \in \mathrm{Ob}(\mathbb{D}_{w-\mathbb{C}-c}^b(X))$  、  $A = \mathrm{SS}(F)$  和  $A_\circ = \{p \in A_{\mathrm{reg}}; \pi: A_{\mathrm{reg}} \to X$  在  $p\}$  的邻域中具有恒定的等级。那么  $A_\circ$  是一个亚解析子流形， $\dim_R(A \setminus A_\circ) < \dim_R X$  。应用命题8.3.27，我们发现 $z_\circ \in \mathbb{C}^N$ ，这样，为 $x \in X$ 和 $A_\varphi = \{(x; \mathrm{d}\varphi(x)); x \in X\} \subset T^*X$ 设置 $\varphi(x) = |x - z_\circ|^2$ ，我们有：

\[
A_{\varphi}\cap A\subset A_{\circ}{~a n d~t h e~i n t e r s e c t i o n~i s~t r a n s v e r s a l~}.
\]

取一个点  $p \in A_{\varphi} \cap A$  并设置  $x_{\circ} = \pi(p)$  。由于  $p \in A_{\circ}$  、  $\mathbf{D}^{b}(X; p)$  和  $\mathbf{D}^{b}(X; p^{a})$  中的  $F \simeq L_{Y}[\dim Y]$  ，对于  $X$  的复杂子流形  $Y$  和某些  $L \in \mathrm{Ob}(\mathbf{D}^{b}(\mathfrak{M}_{\mathfrak{D}\mathfrak{D}}(A)))$  。由于  $A_{\circ} = T_{Y}^{*}X$  在  $p$  的邻域上，并且  $A$  和  $A_{\varphi}$  的交集在  $p$  处是横向的， $\varphi|_{Y}$  在  $x_{\circ}$  处的 Hessian 矩阵是非简并的（参见 VII §5）。另一方面，  $\partial\overline{\partial}\varphi$  在  $T_{x_{\circ}} \mathbb{C}^{N}$  上是正定的，因此在  $T_{x_{\circ}}Y$  上也是正定的。这意味着：

\[
\left\{\begin{matrix}{\begin{matrix}{t h e}\\ {n u m b e r}\\ {o f p o s i t i v e}\\ {e i g e n v a l u e s}\\ {o f H e s s}\\ \end{matrix}(\phi|_{\mathcal{Y}})}\\ {\begin{matrix}{t h e}\\ {s a t}\\ {l e a s t}\\ {d i m}\\ {Y}\\ \end{matrix}.}\\ \end{matrix}\right.
\]

事实上，(10.3.10) 是以下线性代数基本结果的结果：

令  $V$  为  $n$  维复向量空间，  $B$  为实基础空间  $V^R$  上的实对称双线性形式，并将  $B$  扩展到  $\widetilde{B}V^R \otimes_R \mathbb{C} \simeq V \oplus \overline{V}$  上的  $\mathbb{C}$  双线性形式。然后，如果埃尔米特形式  $\widetilde{B}(x, \overline{x})$  在  $V$  上是正定的，我们得到  $B$  至少有  $n$  正特征值。

因此，通过在 Y 上选择一个真实的局部坐标系 $ (x_{1}, \ldots, x_{2d}) $ ，我们可以写：

\[
\varphi(x)=x_{1}^{2}+\cdots+x_{l}^{2}-\cdots-x_{2d}^{2}+\varphi(x_{\circ})
\]

与  $ d = \dim Y $  和  $ l \geq d $  。

然后我们得到：

\[
\begin{aligned}{(H_{\{x;\;\varphi(x)\geqslant\varphi(x_{\circ})\}}^{j}(F))_{x_{\circ}}}&{{}\simeq(H_{\{x_{1}^{2}+\cdots+x_{1}^{2}-\cdots-x_{2d}^{2}\geqslant0\}}^{j+d}(L_{\mathbb{R}^{2d}}))_{0}}\\ {}&{{}\simeq H^{j+d-(2d-l)}(L)}\\ {}&{{}=H^{j+l-d}(L),}\\ \end{aligned}
\]

类似地：

\[
(H_{\left\lbrace{\mathbf{x};\mathbf{\varphi  }\mathbf{\varphi  }(x)\leqslant\mathbf{\varphi  }(x\circ))}\right\rbrace}^{j}(F))_{\mathbf{x}\circ}\simeq H^{j+d-l}(L)\enspace.
\]

从 $ l \geqslant d $ 开始，我们得到：

\[
\begin{cases}{{I f~}F\in\operatorname{O b}(\mathbf{\mu}\mathbf{D}_{w-\mathbb{C}-c}^{{{\leq}0}}(X))}&{{t h e n}}\\ {({H}_{\{x;\mathbf{\varphi}((x)\geqslant\mathbf{\varphi}(x_{\circ})\}}^{j}(F))_{x_{\circ}}=0}&{{f o r~j>0~},}\\ \end{cases}
\]

\[
\begin{cases}{{I f~}F\in\mathbf{O b}\big(\mathbf{\Sigma}^{\mu}\mathbf{D}_{\mathbf{w}-\mathbf{C}-c}^{\geq0}(X)\big),}&{{t h e n}}\\ {(H_{\{x;\mathbf{\varphi}(x)\leq\mathbf{\varphi}(x_{\circ})\}}^{j}(F))_{x_{\circ}}=0}&{{f o r~j<0~}.}\\ \end{cases}
\]

现在  $ \varphi(\pi(\Lambda_\varphi \cap \Lambda)) $  是微局域 Bertini-Sard 定理（命题 8.3.12）的离散集，包含在  $ \mathbb{R}_{\geq 0} $  中。我们设置  $ t_0 = -\infty $  和  $ \varphi(\pi(\Lambda_\varphi \cap \Lambda)) = \{t_1, t_2, \ldots\} $  。

根据命题 5.4.4， $ \mathrm{SS}(R\varphi_*F) \subset \bigcup_i T_{\{t_i\}}^{*} \mathbb{R} $ 。因此  $ R\varphi_* F $  在任何开区间  $ [t_{i-1}, t_i[ (i \geq 1) $  上都是恒定的。设置：

\[
V_{i}=R\varGamma([]-\infty,t_{i+1}[:,R\varphi_{*}F)\enspace,
\]

\[
V_{i}^{c}=R\varGamma_{c}(]-\infty,t_{i+1}[:,R\varphi_{\ast}F),
\]

我们得到 $ i \geqslant 1 $  的显着三角形：

\[
(R\varGamma_{[t_{i},+\infty[}(R\varphi_{*}F))_{t_{i}}\longrightarrow V_{i}\longrightarrow V_{i-1}\xrightarrow[+1]{},
\]

\[
V_{i-1}^{c}\longrightarrow V_{i}^{c}\longrightarrow(R\varGamma_{]-\infty,t_{i}]}(R\varphi_{\ast}F))_{t_{i}}\xrightarrow[+1]{}~.
\]

现在假设  $ F \in \mathrm{Ob}(\mu \mathbf{D}_{w-c-c}^{≤0}(X)) $  。

由于  $ (R\Gamma_{\{\phi(x)\geq t_i\}}(F))|_{\phi^{-1}(t_i)} $  由  $ \varphi^{-1}(t_i)\cap\pi(\mathrm{SS}(F)\cap\varLambda_\varphi) $  支持，我们有：

\[
H^{j}(R\varGamma_{\{\mathfrak{t}_{i},+\infty\}}(R\varphi_{*}F))_{\mathfrak{t}_{i}}\simeq\bigoplus_{p}H^{j}(R\varGamma_{\{\varphi(x)\geq\mathfrak{t}_{i}\}}(F))_{p}
\]

其中  $p$  的范围是有限集  $\varphi^{-1}(t_i) \cap \pi(\mathrm{SS}(F) \cap \mathcal{A}_\varphi)$  。对于  $j > 0$ ，这些组在 (10.3.14) 处消失，并且区分三角形 (10.3.16) 得出：

\[
H^{j}(V_{i})\simeq H^{j}(V_{i-1})\qquad\mathrm{f o r}\qquad j>0\qquad\mathrm{a n d}
\]

\[
H^{0}(V_{i})\to H^{0}(V_{i-1})\qquad\mathrm{i s~s u r j e c t i v e}.
\]

因此，族 $ \{H^{j-1}(V_i)\}_i $ 满足 $ j > 0 $ 的Mittag-Leffler条件，我们通过命题2.7.1得到：

\[
\begin{aligned}{H^{j}(X;F)}&{{}\simeq H^{j}(\mathbb{R};R\varphi_{\ast}F)}\\ {}&{{}\simeq\varprojlim_{i}H^{j}(V_{i})}\\ {}&{{}\simeq H^{j}(V_{0})=0\qquad\mathrm{f o r}\qquad j>0.}\\ \end{aligned}
\]

类似地，如果假设  $F$  属于  $\mu\mathbf{D}_{\mathbf{w}-\mathbf{C}-\mathbf{c}}^{\geq 0}(X)$  ，则为  $j<0$  得到  $H^{j}(R\Gamma_{\{t\leq t_{i}\}}(F))_{t_{i}}=0$  ，因此：

\[
H^{j}(V_{i^{-1}}^{c})\simeq H^{j}(V_{i}^{c})\qquad\mathrm{f o r}\qquad j<0.
\]

这给出：

\[
\begin{aligned}{H_{c}^{j}(X;F)}&{{}\simeq H_{c}^{j}(\mathbb{R};R\varphi_{*}F)}\\ {}&{{}\simeq\varinjlim_{i}H^{j}(V_{i}^{c})}\\ {}&{{}\simeq H^{j}(V_{0}^{c})=0\qquad\mathrm{f o r~\quad~j<0~.\quad}\Box}\\ \end{aligned}
\]

\statementlabel{引理 10.3.9}令  $F \in \mathrm{Ob}(\mathbf{D}_{w-\mathbb{C}-c}^{b}(X))$  和  $X = \bigsqcup_{\alpha} X_\alpha$  成为分层，使得  $\mathrm{SS}(F) \subset \bigsqcup_{\alpha} T_{X_\alpha}^* X$  。令 $Y$  为 $X$  的子流形，它与所有 $X_\alpha$  横向相交。

(i) 假设 $F$  属于 $\mu\mathbf{D}_{w-C-c}^{≤0}(X)$  。那么  $i_{Y}^{-1}F$  属于  $\mu\mathbf{D}_{w-C-c}^{≤-\mathrm{codim}Y}(Y)$  ，  $i_{Y}^{+}F$  属于  $\mu\mathbf{D}_{w-C-c}^{≤\mathrm{codim}Y}(Y)$  。

(ii) 假设 $F$  属于 $\mu\mathbf{D}_{\mathbf{w}-\mathbb{C}-\mathbb{C}}^{\geq 0}(X)$  。那么  $i_{Y}^{-1}F$  属于  $\mu\mathbf{D}_{\mathbf{w}-\mathbb{C}-\mathbb{C}}^{\geq -\operatorname{codim}Y}(Y)$  ，  $i_{Y}^{\dagger}F$  属于  $\mu\mathbf{D}_{\mathbf{w}-\mathbb{C}-\mathbb{C}}^{\geq \operatorname{codim}Y}(Y)$  。

\statementlabel{证明}根据横向条件， $ T_Y^* X \cap (\bigsqcup_\alpha T_{X_x}^* X) \subset T_X^* X $  。因此  $ i_Y $  对于  $ F $  来说是非特征的，我们可以应用命题 5.4.13 来获得：

\[
i_{Y}^{-1}F\otimes\omega_{Y/X}\simeq i_{Y}^{!}F\mathrm{~,~}
\]

\[
\mathrm{SS}(i_{Y}^{-1}F)\subset\bigsqcup_{\alpha}T_{Y\cap X_{\alpha}}^{*}Y.
\]

让 $y \in Y \cap X_{\alpha}$  。根据推论 8.3.24，存在  $p \in Y \times_X T_{X^*} X \setminus \bigcup_{\beta \neq \alpha} T_{X_\beta}^* X$  和  $\pi(p) = y$  。然后  $F \simeq L_{X_\alpha}$  中的  $\mathbf{D}^b(X; p)$  中的一些  $L \in \mathbf{D}^b(\mathfrak{M}_\partial \partial(A))$  ，以及  $i_Y^{-1} F \simeq L_{X_\alpha \cap Y}$  中的  $\mathbf{D}^b(Y; t_{i_Y'}(p))$  ，根据命题 6.1.9。

然后 (i) 和 (ii) 等价于 (10.3.6)、(10.3.7) 和 (10.3.8)。 ⑨

接下来我们研究反常层的特化。

\statementlabel{命题 10.3.10}令 Y 为 X 的平滑超曲面，并令  $F \in \mathrm{Ob}(\mathbf{D}_{w-C-c}^{b}(X))$  。假设  $F|_{X \setminus Y} \in \mathrm{Ob}(\mu \mathbf{D}_{w-C-c}^{s_0}(X \setminus Y))$  （分别为  $\mu \mathbf{D}_{w-C-c}^{s_0}(X \setminus Y)$  ）。那么  $\nu_Y(F)|_{\dot{T}_Y X}$  属于  $^p \mathbf{D}_{w-C-c}^{s_0}(\dot{T}_Y X)$  （分别为  $^p \mathbf{D}_{w-C-c}^{s_0}(\dot{T}_Y X)$  ）。

\statementlabel{证明}我们分三步分解证明。

(a) 假设  $X$  是  $\mathbb{C}$  和  $Y = \{0\}$  的开子集。如果 $F \in \mathrm{Ob}(\mathbf{D}_{w-C-c}^{b}(X))$  和 $F|_{X \setminus Y}$  集中到 $\leqslant -1$  度（或 $\geqslant -1$  ），那么 $\nu_Y(F)$  显然具有相同的属性。

(b) 让我们证明，在该命题的假设下，对于任何  $ p \in \dot{T}_Y X $  ：

\[
\left\{\begin{aligned}{}&{{}i_{p}^{-1}\nu_{Y}(F)\in\operatorname{O b}(\mathbf{D}^{\leq-1}(\mathfrak{M o b}(A))),}\\ {}&{{}({r e s p.}i_{p}^{!}\nu_{Y}(F)\in\operatorname{O b}(\mathbf{D}^{\geq1}(\mathfrak{M o b}(A))))}\\ \end{aligned}\right..
\]

让我们在  $ X $  上选择一个局部全纯坐标系  $ (z_1, \ldots, z_n) $  ，使得  $ Y = \{z_1 = 0\} $  和  $ p = \left(0; \frac{\partial}{\partial z_1}\right) $  。令  $ f: X \to \mathbb{C} $  为第一个坐标  $ z_1 $  给出的态射，用  $ B(\varepsilon) $  表示  $ X $  中以 0 为中心、半径为  $ \varepsilon > 0 $  的开球， $ j_\varepsilon $  表示嵌入  $ B(\varepsilon) \hookrightarrow X $  ， $ f_\varepsilon $  表示映射  $ f \circ j_\varepsilon $  。最后用 $ T_Y $  表示映射 $ T_Y X \to T_{\{0\}} \mathbb{C} $  。应用命题 4.2.4 我们得到同构：

\[
\left\{\begin{aligned}{}&{{}(T_{Y}f)_{*}\nu_{Y}(R j_{\varepsilon*}j_{\varepsilon}^{-1}F)\simeq\nu_{\{0\}}(R f_{*}R j_{\varepsilon*}j_{\varepsilon}^{-1}F),}\\ {}&{{}(T_{Y}f)_{!}\nu_{Y}(R j_{\varepsilon!}j_{\varepsilon}^{-1}F)\simeq\nu_{\{0\}}(R f_{!}R j_{\varepsilon!}j_{\varepsilon}^{-1}F).}\\ \end{aligned}\right.
\]

另一方面，根据引理 8.4.7 我们有：

\[
\left\{\begin{aligned}{}&{{}i_{p}^{-1}\nu_{Y}(F)\simeq\underrightarrow{\operatorname*{l i m}}_{U}R\varGamma(U;\nu_{Y}(F)),}\\ {}&{{}i_{p}^{\downarrow}\nu_{Y}(F)\simeq\underrightarrow{\operatorname*{l i m}}_{U}R\varGamma_{c}(U;\nu_{Y}(F)),}\\ \end{aligned}\right.
\]

其中  $U$  范围涵盖  $T_Y X$  中  $p$  的开放邻域系列。鉴于 (10.3.19)，这表明设置  $v = \left(0; \frac{\partial}{\partial z_x}\right) \in T_{\{0\}} \mathbb{C}$  ：

\[
\left\{\begin{aligned}{}&{{}i_{p}^{-1}\nu_{Y}(F)\simeq\underrightarrow{\operatorname*{l i m}}_{\varepsilon}\nu_{\{0\}}(R f_{\varepsilon*}j_{\varepsilon}^{-1}F)_{v},}\\ {}&{{}i_{p}^{!}\nu_{Y}(F)\simeq\underrightarrow{\operatorname*{l i m}}_{\varepsilon^{\prime}}\mathit{R}\varGamma_{\{v\}}\nu_{\{0\}}(R f_{\varepsilon!}j_{\varepsilon}^{-1}F).}\\ \end{aligned}\right.
\]

根据命题8.5.9，我们知道对于 $U$ ， $0 \in \mathbb{C}$ 、 $R f_{\varepsilon*} j_{\varepsilon}^{-1} F|_{U}$  和 $R f_{\varepsilon!} j_{\varepsilon}^{-1} F|_{U}$  的足够小的开邻域属于 $\mathbf{D}_{w-C-c}^{b}(U)$ 。

我们将在该命题的假设下证明：

\[
\begin{cases}{R f_{\varepsilon\ast}j_{\varepsilon}^{-1}F\in\operatorname{O b}(^{p}\mathbb{D}_{w^{-}\mathbb{C}^{-}c}^{\leq0}(U\backslash\{0\}))~,}\\ {(r e s p.R f_{\varepsilon!}j_{\varepsilon}^{-1}F\in\operatorname{O b}(^{p}\mathbb{D}_{w^{-}\mathbb{C}^{-}c}^{\geq0}(U\backslash\{0\}))~.}\\ \end{cases}
\]

我们选择一个复杂的  $ \mu $  -分层  $ X = \bigsqcup_{\alpha} X_{\alpha} $  使得  $ \mathrm{SS}(F) \subset \bigsqcup_{\alpha} T_{X_{\alpha}}^{*} X $  。然后根据微局域 Bertini-Sard 定理（命题 8.3.12），

\[
\bigsqcup_{\alpha}\left(B(\varepsilon)\times_{X}T_{X_{\alpha}}^{*}X\right)\cap T_{\{f^{-1}(t)\}}^{*}X\subset T_{X}^{*}X\quad\mathrm{f o r}\quad0<|t|\ll\varepsilon.
\]

应用引理 10.3.9，我们发现如果  $F$  属于  $\mathcal{U}_{\mathbf{w}-\mathbf{C}-\mathbf{c}}^{\leq 0}(X)$  （分别为  $\mathcal{U}_{\mathbf{w}-\mathbf{C}-\mathbf{c}}^{\geq 0}(X)$  ），则  $i_{f_{\varepsilon}^{-1}(t)}^1 j_{\varepsilon}^{-1} F$  （分别为  $i_{f_{\varepsilon}^{-1}(t)}^{-1} j_{\varepsilon}^{-1} F$  ）属于  $\mathcal{U}_{\mathbf{w}-\mathbf{C}-\mathbf{c}}^{\leq 1}(f_{\varepsilon}^{-1}(t))$  （分别为  $\mathcal{U}_{\mathbf{w}-\mathbf{C}-\mathbf{c}}^{\geq -1}(f_{\varepsilon}^{-1}(t))$  ）。由于 $f_{\varepsilon}^{-1}(t)$  是斯坦因流形，我们可以应用定理10.3.8。我们得到：

\[
R\varGamma(f_{\varepsilon}^{-1}(t);i_{f_{\varepsilon-1}(t)}^{!}j_{\varepsilon}^{-1}F)\in\operatorname{O b}(\mathbf{D}^{\leq1}(\mathfrak{M o b}(A)))
\]

（分别 $ R\Gamma_c(f_\varepsilon^{-1}(t); i_{f_\varepsilon^{-1}(t)}^{-1}j_\varepsilon^{-1}F) \in \mathrm{Ob}(\mathbb{D}^{\geq -1}(\mathfrak{M}_\mathrm{od}(A))) $ ）。

自从

\[
(R f_{\varepsilon*}j_{\varepsilon}^{-1}F)_{t}\simeq(R\varGamma_{\{t\}}R f_{\varepsilon*}j_{\varepsilon}^{-1}F)_{\varepsilon}[2]\simeq R\varGamma(f_{\varepsilon}^{-1}(t);i_{f_{\varepsilon}^{-1}(t)}^{!}j_{\varepsilon}^{-1}F)[2],
\]

and

\[
(R f_{\varepsilon!}j_{\varepsilon}^{-1}F)_{t}\simeq R\varGamma_{c}(f_{\varepsilon}^{-1}(t);i_{f_{\varepsilon}^{-1}(t)}^{-1}F)\enspace,
\]

我们得到（10.3.22）。这与（10.3.21）和步骤（a）一起意味着（10.3.18）。

(c) 我们将完成命题的证明。

让我们选择一个 $ \mu $  -分层 $ X = \bigsqcup_{\alpha \in A} X_\alpha $ ，使得存在 $ A_\circ \subset A $ ：

\[
\left\{\begin{aligned}{}&{{}\SS(F)\subset\bigsqcup_{\alpha\in A}T_{X_{\alpha}}^{*}X\enspace,}\\ {}&{{}Y=\bigsqcup_{\alpha\in A_{0}}X_{\alpha}\enspace,}\\ {}&{{}\SS(\nu_{Y}(F)|_{\dot{T}_{\mathcal{Y}}X})\subset\bigsqcup_{\alpha\in A_{0}}(T_{X_{\alpha}}^{*}Y\times\dot{T}_{Y}X)\enspace.}\\ \end{aligned}\right.
\]

令  $ x \in X_\alpha $  并取  $ X $  的子流形  $ Z $  ，该子流形与  $ X_\alpha $  横向相交于  $ x $  ，使得  $ \dim Z + \dim X_\alpha = \dim X $  。将逆傅立叶-佐藤变换应用于定理 6.7.1 中获得的同构（或者更好，通过查看该定理的证明），可以得到同构：

\[
\left\{\begin{aligned}{}&{{}\nu_{\mathbf{Y}}(F)|_{\mathbf{Z}\times_{\mathbf{Y}}T_{\mathbf{Y}}\mathbf{X}}\simeq\nu_{\mathbf{Z}\cap\mathbf{Y}}(i_{\mathbf{Z}}^{-1}F),}\\ {}&{{}R\varGamma_{\mathbf{Z}\times_{\mathbf{Y}}T_{\mathbf{Y}}\mathbf{X}}\nu_{\mathbf{Y}}(F)\simeq\nu_{\mathbf{Z}\cap\mathbf{Y}}(i_{\mathbf{Z}}^{!}F).}\\ \end{aligned}\right.
\]

让 $ p \in \{x\} \times_Y \dot{T}_Y X \simeq \{x\} \times_{Z \cap Y} \dot{T}_{Z \cap Y} Z $  。

然后：

\[
\left\{\begin{aligned}{}&{{}\nu_{Y}(F)_{p}\simeq\nu_{Z\cap Y}(i_{Z}^{-1}F)_{p},}\\ {}&{{}R\varGamma_{\{p\}}\nu_{Y}(F)\simeq R\varGamma_{\{p\}}\nu_{Z\cap Y}(i_{Z}^{!}F).}\\ \end{aligned}\right.
\]

应用引理 10.3.9，我们得到  $ i_Z^{-1}(F)|_{Z \setminus Y} $  （分别为  $ i_Z^!, (F)|_{Z \setminus Y} $  ）属于  $ \mu \mathbf{D}_{w-C-c}^{\leq -\operatorname{codim} Z}(Z \setminus Y) $  （分别为  $ \mu \mathbf{D}_{w-C-c}^{\geq \operatorname{codim} Z}(Z \setminus Y) $  ）。应用步骤（b）我们得到：

\[
\nu_{Z\cap Y}(i_{Z}^{-1}F)_{p}\in\operatorname{O b}(\mathbb{D}^{-1-\operatorname{d i m}X_{\mathfrak{s}}}(\mathfrak{M o b}(A)))
\]

（分别 $ R\Gamma_{\{p\}}\nu_{Z\cap Y}(i_{Z}^{-1}F)_{p}\in\mathrm{Ob}(\mathbb{D}^{\geq1+\dim X_{\alpha}}(\mathfrak{M o b}(A))) $ ）。

然后得出结论（10.3.25）和（10.3.4）（分别为（10.3.25）和（10.3.5））。 □

\statementlabel{推论 10.3.11}令 Y 为由方程  $ \{f = 0\} $  定义的 X 的平滑超曲面，并令  $ F \in \mathrm{Ob}(\mathbb{D}_{w-C-c}^b(X)) $  。假设 $ F|_{X \setminus Y} $ 属于 $ \mu \mathbb{D}_{w-C-c}^b(X \setminus Y) $ 

（分别 $ \mu\mathbf{D}_{w-\mathbb{C}-c}^{\geq0}(X\setminus Y) $ ）。然后：

\[
\psi_{f}(F)[-1]{~b e l o n g s~t o~}^{p}\mathbb{D}_{\mathbf{w}-\mathbb{C}-\mathbb{c}}^{\leq0}(Y)({r e s p.~}^{p}\mathbb{D}_{\mathbf{w}-\mathbb{C}-\mathbb{c}}^{\geq0}(Y)).
\]

(ii) 如果 此外，$i_Y^{-1}(F)$（分别地，$i_Y^\dagger(F)$）  属于 $^p\mathbf{D}_{w-C-c}^{\leq0}(Y)$  （分别地，$^p\mathbf{D}_{w-C-c}^{\geq0}(Y)$），则 $\phi_f(F)$  属于 $^p\mathbf{D}_{w-C-c}^{\leq0}(Y)$  （分别地，$^p\mathbf{D}_{w-C-c}^{\geq0}(Y)$）。

\statementlabel{证明}(i) 由命题 8.6.3 和 10.3.10 得出。

(ii) 由 (i) 和区分三角形 (8.6.7) 得出。 □

我们现在可以证明本节的主要结果之一。

\statementlabel{定理 10.3.12}其中一个有  $ \mu\mathbf{D}_{\mathbf{w}-\mathbb{C}-\mathbb{C}}^{0}(X) = \mu\mathbf{D}_{\mathbf{w}-\mathbb{C}-\mathbb{C}}^{0}(X) $  和  $ \mu\mathbf{D}_{\mathbf{w}-\mathbb{C}-\mathbb{C}}^{0}(X) = \mu\mathbf{D}_{\mathbf{w}-\mathbb{C}-\mathbb{C}}^{0}(X) $  。

\statementlabel{证明}(a) 令  $t: X \times \mathbb{C} \to \mathbb{C}$  为投影。如果  $F$  属于  $\mu \mathbf{D}_{w-C-c}^{≤0}(X)$  （分别为  $\mu \mathbf{D}_{w-C-c}^{≥0}(X)$  ），则  $F \boxtimes A_c[1]$  属于  $\mu \mathbf{D}_{w-C-c}^{≤0}(X \times \mathbb{C})$  （分别为  $\mu \mathbf{D}_{w-C-c}^{≥0}(X \times \mathbb{C}$  ）。那么根据命题 10.3.10， $\psi_t(F \boxtimes A_c[1])[-1] \simeq F$  属于  $^p \mathbf{D}_{w-C-c}^{≤0}(X)$  （分别为  $^p \mathbf{D}_{w-C-c}^{≥0}(X)$  ）。

(b) 相反，假设  $F$  属于  $^p\mathbf{D}_{\mathbf{w}-\mathbf{C}-\mathbf{c}}^{\leq 0}(X)$  （分别为  $^p\mathbf{D}_{\mathbf{w}-\mathbf{C}-\mathbf{c}}^{\geq 0}(X)$  ）。

让我们采用  $\mu$  分层  $X = \bigsqcup_{\alpha} X_{\alpha}$  使得  $\mathrm{SS}(F) \subset \bigsqcup_{\alpha} T_{X_{\alpha}}^{*} X$  。通过归纳论证足以表明：

\[
\begin{cases}\text{若 }X_\alpha\text{ 是闭层且 }F|_{X\setminus X_\alpha}\in\mu\mathbf{D}_{w-C-c}^{\leq0}(X\setminus X_\alpha),\\\text{则 }F\text{ 在 }T_{X_\alpha}^*X\text{ 的一般点处、移位 }0\text{ 的型 }L\text{ 满足}\\H^j(L)=0\quad\text{对 }j>-\dim X\text{（分别地 }j<-\dim X\text{）成立。}\end{cases}
\]

我们选择一个子流形 $Z$ ，它与 $X_\alpha$ 横向相交于 $x$ 。由于  $\mu$  -perversities 和  $p$  -perversities 被  $i_Z^{-1}$  保留，具有相同的转变（推论 10.2.6 和引理 10.3.9），我们可以假设，通过用  $Z$  替换  $X$  和用  $i_Z^{-1}F$  替换  $F$  ，  $X_\alpha = \{x\}$  。然后采用全纯函数  $f$  使得  $f(x) = 0$  、  $\mathrm{d}f(x) = p \in \dot{T}_{\{x\}}^{*} X \setminus \bigcup_{\beta \neq \alpha} T_{X_\beta}^{*} X$  。由于  $\mathrm{supp}(\phi_f(F))$  包含在  $\pi(\mathrm{SS}(F) \cap \mathcal{A}_f)$  中，因此它集中到  $\{x\}$  。此外 $\phi_f(F)_x \simeq \phi_f(L_{\{x\}}) = L$ 。因此，我们获得了所需的结果，因为根据推论 10.3.11， $\phi_f(F)$  属于  $^p\mathbf{D}_{w-C-c}^{≤0}(f^{-1}(0))$ （分别为  $^p\mathbf{D}_{w-C-c}^{≥0}(f^{-1}(0))$  ）。   $\square$ 

\statementlabel{推论 10.3.13}从 $\mathbf{D}_{\mathbf{w}-\mathbb{C}-\mathbb{C}}^b(X)$  到 $\mathbf{D}_{\mathbf{w}-\mathbb{C}-\mathbb{C}}^b(f^{-1}(0))$  的函子 $\psi_f[-1]$  和 $\phi_f$  是 $t$  - 相对于中间反常给出的 $t$  结构而言是精确的。

\statementlabel{证明}当  $ df \neq 0 $  时，这由推论 10.3.11 和定理 10.3.12 得出。否则，将练习 VIII.15 应用于 f 的图嵌入。 ⑨

回想一下，设置  $^pH^0 = ^p\tau^{\leq 0} \circ ^p\tau^{\geq 0}$  和  $^pH^n = ^pH^0 \circ [n]$  。

\statementlabel{推论 10.3.14}让 $ F \in \mathrm{Ob}(\mathbb{D}_{w-c-c}^{b}(X)) $  。然后：

\[
\mathrm{S S}({}^{p}H^{j}(F))\subset\mathrm{S S}(F)\enspace.
\]

\statementlabel{证明}这是从  $^p H^j(\phi_f(F)) \simeq \phi_f^p H^j(F)$  （推论 10.3.13）和命题 8.6.4 得出的。 □

我们现在将研究反常层的函子运算。

\statementlabel{命题 10.3.15}（ $ f^{-1} $  的左侧精确性和 $ f' $  的右侧精确性）。令 $ f: Y \to X $  为流形的态射。

(i)  $ f^{-1} $  将 $ ^p\mathbf{D}_{w-C-c}^{≥0}(X) $  发送至 $ ^p\mathbf{D}_{w-C-c}^{≥\dim Y/X}(Y) $  。

(ii)  $ f^! $  将 $ ^p\mathbf{D}_{w-C-c}^0(X) $  发送至 $ ^p\mathbf{D}_{w-C-c}^{≤-\dim Y/X}(Y) $  。

\statementlabel{证明}通过对f进行图分解，我们可以分别考虑f是平滑的情况和f是闭嵌入的情况。

如果  $f$  是平滑的，则结果来自命题 10.3.3，因为在这种情况下是  $f^! \simeq f^{-1} \otimes o \operatorname{tr}_{Y/X}[2 \dim Y/X]$  。

如果  $f$  是一个封闭嵌入，我们通过归纳法进行论证，并将证明简化为  $Y$  是由方程  $\{g=0\}$  给出的超曲面的情况。然后从区分的三角形 (8.6.7) 和推论 10.3.13 得出结果。   $\square$ 

\statementlabel{推论 10.3.16}设  $F \in \mathrm{Ob}(\mathbf{D}_{\mathrm{w}-\mathrm{C}-\mathrm{c}}^{b}(X))$  ，并假设  $f: Y \to X$  对于  $F$  来说是非特征的。然后：

(i)  $f$  对于  $^pH^j(F)$  、  $^p\tau^{\geq j}F$  和  $^p\tau^{\leq j}F$  来说是非特征的。

(ii) 如果  $ F \in \mathrm{Ob}(^p\mathbb{D}_{w-\mathbb{C}-c}^0(X)) $  ，则  $ f^{-1}F \in \mathrm{Ob}(^p\mathbb{D}_{w-\mathbb{C}-c}^{\dim Y/X}(Y)) $  和  $ f^!F \in \mathrm{Ob}(^p\mathbb{D}_{w-\mathbb{C}-c}^{-\dim Y/X}(Y)) $  。

(三)

\[
{}^{p}H^{j}(f^{-1}(F))\simeq f^{-1}({}^{p}H^{j-\dim Y/X}(F))[\dim Y/X]~,
\]

\[
{}^{p}H^{j}(f^{!}(F))\simeq f^{!}({}^{p}H^{j+\dim Y/X}(F))[-\dim Y/X]\enspace.
\]

\statementlabel{证明}(i) 由推论 10.3.14 得出

(ii) 自  $ f^1(F) \simeq f^{-1}(F) \otimes o t_{Y/X} [2 \dim Y/X] $  以来，遵循命题 10.3.15。

(iii) 很容易从(i) 和(ii) 得出。 □

\statementlabel{命题 10.3.17}令 $f: Y \to X$  为流形的态射。假设 $X$  的任何点都有一个开邻域 $U$ ，使得 $f^{-1}(U)$  是斯坦因流形。

(i) 设  $ G \in \text{Ob}(^p\mathbb{D}_{\mathbf{w}-\mathbb{C}-c}^0(Y)) $  ，并假设  $ Rf_* G \in \text{Ob}(\mathbb{D}_{\mathbf{w}-\mathbb{C}-c}^b(X)) $  。然后 $ Rf_* G \in \text{Ob}(^p\mathbb{D}_{\mathbf{w}-\mathbb{C}-c}^0(X)) $ 。

(ii) 设  $ G \in \text{Ob}(^p\mathbb{D}_{\mathbf{w}-\mathbb{C}-\mathbb{C}}^{\geq 0}(Y)) $  ，并假设  $ Rf_1G \in \text{Ob}(\mathbb{D}_{\mathbf{w}-\mathbb{C}-\mathbb{C}}^b(X)) $  。然后 $ Rf_1G \in \text{Ob}(^p\mathbb{D}_{\mathbf{w}-\mathbb{C}-\mathbb{C}}^{\geq 0}(X)) $ 。

\statementlabel{证明}由于闭嵌入的结果是清楚的，因此通过图分解 f 来考虑 f 是平滑的情况就足够了。

令 $G$ 如(i)（或(ii)）中那样，并让我们选择 $\mu$  -分层 $X = \bigsqcup_{\alpha} X_{\alpha}$ ，使得 $SS(R f_{\star} G)$ （或 $SS(R f_1 G)$ ）包含在 $\bigsqcup_{\alpha} T_{X_{\alpha}}^* X$ 中。对于  $x \in X_{\alpha}$  选择  $X$  的子流形  $Z$  和  $\dim Z + \dim X_{\alpha} = \dim X$ ，它与  $X_{\alpha}$  横向相交于  $x$  。则由(10.3.4)和(10.3.5)表示即可：

\[
\left\{\begin{aligned}i_x^{-1}i_Z^!R f_*G&\in\operatorname{Ob}(\mathbf{D}^{\leq\dim X_x}(\mathfrak{Mod}(A))),\\ \text{resp. }i_x^!i_Z^{-1}R f_!G&\in\operatorname{Ob}(\mathbf{D}^{\geq-\dim X_x}(\mathfrak{Mod}(A))).\end{aligned}\right.
\]

令  $ \tilde{f}: f^{-1}(Z) \to Z $  为  $ f $  的限制。然后 $ i_Z^! R f_* G \simeq R \tilde{f}_* i_{f^{-1}(Z)}^! G $  和 $ i_Z^{-1} R f_! G \simeq R \tilde{f}_! i_{f^{-1}(Z)}^{-1} G $  。因此，我们找到  $ x $  的足够小的 Stein 开邻域  $ U $  ：

\[
\left\{\begin{aligned}{i_{\mathbf{x}}^{-1}i_{\mathbf{Z}}^{!}R f_{\ast}G}&{{}\simeq R\varGamma(U;i_{\mathbf{Z}}^{!}R f_{\ast}G)\simeq R\varGamma(\tilde{f}^{-1}(U);i_{f^{-1}(\mathbf{Z})}^{!}G)}\\ {(\operatorname{r e s p.}~i_{\mathbf{x}}^{!}i_{\mathbf{Z}}^{-1}R f_{!}G}&{{}\simeq R\varGamma_{c}(U;i_{\mathbf{Z}}^{-1}R f_{!}G)\simeq R\varGamma_{c}(\tilde{f}^{-1}(U);i_{f^{-1}(\mathbf{Z})}^{-1}G))\enspace.}\\ \end{aligned}\right.
\]

由于根据命题 10.3.15， $ i_{f^{-1}(Z)}^G $  属于  $ ^p\mathbf{D}_{w-C-c}^\leq\dim X_w(f^{-1}(Z)) $ （相应地， $ i_{f^{-1}(Z)}^{-1} $  G 属于  $ ^p\mathbf{D}_{w-C-c}^\geq-\dim X_w $  (  $ f^{-1}(Z) $  )），因此 (10.3.28) 由 (10.3.29) 和定理 10.3.8 得出。   $ \square $ 

\statementlabel{命题 10.3.18}令  $E$  为  $X$  上的（复）向量丛，纤维维度为  $n$  。然后，傅立叶-佐藤变换将  $^p\mathbf{D}_{w-C-c}^{≤0}(E)\cap\mathbf{D}_{\mathbf{R}+}^{b}(E)$  和  $^p\mathbf{D}_{w-C-c}^{≥0}(E)\cap\mathbf{D}_{\mathbf{R}+}^{b}(E)$  分别发送到  $^p\mathbf{D}_{w-C-c}^{≤n}(E^{*})\cap\mathbf{D}_{\mathbf{R}+}^{b}(E^{*})$  和  $^p\mathbf{D}_{w-C-c}^{≥n}(E^{*})\cap\mathbf{D}_{\mathbf{R}+}^{b}(E^{*})$  。

\statementlabel{证明}让 $F$  属于 $\mathbf{D}_{w-C-c}^{b}(E) \cap \mathbf{D}_{\mathbf{R}+}^{b}(E)$  。考虑一下图表：

\begin{center}
\includegraphics[width=0.46\textwidth]{流行上的层_Chn-assets/img_in_image_box_296_873_854_1072.png}
\end{center}

其中  $p_1$  和  $p_2$  （分别为 pr 和 t）表示在  $E \times_X E^*$  （分别为  $E^* \times \mathbb{C}$  ）上定义的第一个和第二个投影， $\gamma$  是映射  $(x, y) \mapsto (y, \langle x, y \rangle)$  。那么根据定义 3.7.8 有：

\[
\left\{\begin{aligned}{F^{\wedge}}&{{}\simeq R p_{2*}{R}\varGamma_{\gamma^{-1}(E^{*}\times(\mathbb{R e})^{-1}(\mathbb{R}_{\geq0}))}(p_{1}^{-1}F)}\\ {}&{{}\simeq R p_{2}!({(p_{1}^{-1}F)_{\gamma^{-1}(E^{*}\times(\mathbb{R e})^{-1}(\mathbb{R}_{<0}))}}).}\\ \end{aligned}\right.
\]

使用(10.3.30)中的第一个同构，我们得到：

\[
\begin{aligned}{F^{\wedge}}&{{}\simeq R\mathbf{p r}_{*}{R\gamma}_{*}{R\Gamma_{\gamma^{-1}(E^{*}\times(\mathtt{R e})^{-1}(\mathtt{R}_{\geq0})})}(p_{1}^{-1}F),}\\ {}&{{}\simeq R\mathbf{p r}_{*}{R\Gamma_{E^{*}\times(\mathtt{R e})^{-1}(\mathtt{R}_{\geq0})}}(R\gamma_{*}p_{1}^{-1}F).}\\ \end{aligned}
\]

由于  $ R\gamma_{*}p_{1}^{-1}F $  相对于  $ \mathbb{C}^{\times} $  对  $ \mathbb{C} $  的作用是  $ \mathbb{C}^{\times} $  圆锥曲线，因此可以轻松推断出：

\[
{\cal F}^{\wedge}\simeq\phi_{t}R\gamma_{*}p_{1}^{-1}F\quad.
\]

类似地，使用（10.3.30）中的第二个同构，我们得到：

\[
\begin{aligned}{F^{\wedge}}&{{}\simeq R\mathbf{p r}_{!}((R\gamma_{!}p_{1}^{-1}F)_{E^{*}\times(\mathbb{R e})^{-1}(\mathbb{R}_{<0})})}\\ {}&{{}\simeq\phi_{t}(R\gamma_{!}p_{1}^{-1}F)~.}\\ \end{aligned}
\]

如果根据命题 10.3.3， $F$  属于  $^p\mathbf{D}_{w-C-c}^{≤0}(E)\cap\mathbf{D}_{R^+}^{b}(E)$  （分别为  $^p\mathbf{D}_{w-C-c}^{≥0}(E)\cap\mathbf{D}_{R^+}^{b}(E))$  ，则  $p_1^{-1}F$  属于  $^p\mathbf{D}_{w-C-c}^{≤n}(E\times_X E^*)$  （分别为  $^p\mathbf{D}_{w-C-c}^{≥n}(E\times_X E^*)$  ）。那么我们可以将命题 10.3.17 应用于映射  $\gamma$  并通过  $t$  的精确性得出结论 $\phi_t$ （推论 10.3.13）。

\statementlabel{命题 10.3.19}令  $Y$  为  $X$  的闭子流形，余维为  $n$  。那么 $\nu_Y: \mathbf{D}_{w-c-c}^b(X) \to \mathbf{D}_{w-c-c}^b(T_Y X)$  和 $\mu_Y[\operatorname{codim} Y]: \mathbf{D}_{w-c-c}^b(X) \to \mathbf{D}_{w-c-c}^b(T_Y^*X)$  是 $t$  -精确的。

\statementlabel{证明}我们首先描述 IV §1 中引入的法线变形的复杂版本。

我们可以假设  $X$  在  $\mathbb{C}^n \times \mathbb{C}^m$  和  $Y = \{z' = 0\}$  中打开，其中  $z = (z', z')$  、  $z' \in \mathbb{C}^n$  、  $z'' \in \mathbb{C}^m$  。令  $f$  为  $(t, z', z'') \mapsto (tz', z'')$  给出的  $\mathbb{C} \times \mathbb{C}^n \times \mathbb{C}^m$  的映射，并令  $\tilde{X} = f^{-1}(X)$  。我们仍然用  $f$  表示  $f$  到  $\tilde{X}$  的限制。

让 $ Z = t^{-1}(0) \subset \tilde{X} $  。映射  $ T_f $  给出态射  $ T_z f: T_Z \tilde{X} \to T_Y X $  。将 $ T_Y X $  与 $ Y \times \mathbb{C}^n $  以及 $ T_Z \tilde{X} $  与 $ Z \times \mathbb{C} $  进行识别， $ T_Z f $  由 $ (z', z'',\tau) \mapsto (z'', \tau z') $  给出。令  $ s: Z \to T_Z \tilde{X} $  为  $ (z', z'') \mapsto (z', z'',\,1) $  给出的部分。然后  $ T_Z f $  标识  $ s(Z) $  和  $ T_Y X $  。

现在将命题 4.2.5 应用于平滑映射  $T_{zf}|_{\dot{T}_{z}\tilde{x}}$  ，我们得到同构：

\[
\begin{aligned}{\nu_{Y}(F)}&{{}\simeq(T_{Z}f\circ s)^{-1}\nu_{Y}(F)}\\ {}&{{}\simeq s^{-1}(T_{Z}f)^{-1}\nu_{Y}(F)}\\ {}&{{}\simeq s^{-1}\nu_{Z}(f^{-1}F)}\\ {}&{{}\simeq\psi_{t}(f^{-1}F).}\\ \end{aligned}
\]

由于 $ F \mapsto (f^{-1}F[1])|_{t \neq 0} $ 和 $ \psi_t[-1] $ 保留了与中间反常相关的 $ t $ 结构，因此关于 $ \nu_Y $ 的结果如下。然后我们应用命题10.3.18来完成证明。   $ \square $ 

\statementlabel{推论 10.3.20}(i) 设  $F \in \mathrm{Ob}(^p\mathbb{D}_{\mathrm{w}-\mathrm{C}-\mathrm{c}}^{\leq 0}(X))$  和  $G \in \mathrm{Ob}(^p\mathbb{D}_{\mathrm{w}-\mathrm{C}-\mathrm{c}}^{\geq 0}(X))$  。那么 $\mu_{\mathrm{hom}}(F, G)[\dim X]$  属于 $^p\mathbb{D}_{\mathrm{w}-\mathrm{C}-\mathrm{c}}^{\geq 0}(T^*X)$  。

(ii) 假设基环 $A$  是一个域。那么如果  $F$  和  $G$  属于  $^p\mathbf{D}_{\mathbb{C}-c}^0(X)$  ，  $\mu\text{hom}(F,G)[\dim X]$  属于  $^p\mathbf{D}_{\mathbb{C}-c}^0(T^*X)$  。

\statementlabel{证明}(i) 由命题 10.3.6 得出，并将前面的命题应用于  $ X \times X $  的对角线。

(ii) 由 (i) 和命题 8.4.14 得出，因为

\[
\mathbf{D}_{T^{*}X}(\mu\mathcal{h}o m(F,G)[\dim X])\simeq\mu\mathcal{h}o m(G,F)\otimes\mathfrak{o r}_{X}[\dim X]~.\quad\Box
\]

\subsection{习题}
\statementlabel{练习 X.1}令  $\mathbf{D}$  为具有  $t$  结构的三角范畴。对于 $X \in \mathrm{Ob}(\mathbf{D})$  ，令 $\mathrm{d}^n(X):\tau^{>n}X \to (\tau^{\leq n}X)[1]$  为命题10.1.4 中构造的态射d。

(i) 证明，对于 $ b \geq a $ 

\textbackslash{}[\textbackslash{}begin\{array\}\{c\}\textbackslash{}tau\textasciicircum{}\{>a\}X\textbackslash{}quad\textbackslash{}xrightarrow\{\textbackslash{}mathrm\{d\}\textasciicircum{}\{a\}(X)\}(\textbackslash{}tau\textasciicircum{}\{\textbackslash{}leq a\}X)[1]\textbackslash{}\textbackslash{}\textbackslash{}downarrow\textbackslash{}quad\textbackslash{}downarrow\textbackslash{}quad\textbackslash{}downarrow\textbackslash{}quad\textbackslash{}downarrow\textbackslash{}quad\textbackslash{}downarrow\textbackslash{}

\[
\tau^{>b}X\xrightarrow{\quad}\overrightarrow{\mathrm{d}^{b}(X)}\quad(\tau^{\lessgtr b}X)\left[1\right]
\]

通勤。

(ii) 证明：

\[
\begin{array}{c c}{(\tau^{>n}X)[1]}&{\xrightarrow{\mathrm{d}^{n}(X)[1]}(\tau^{\leqslant n}X)[1]}\\ {\bigg\downarrow}&{\bigg\downarrow}\\ {\tau^{>n-1}(X[1])}&{\xrightarrow{\mathrm{d}^{n-1}(X[1])}\tau^{\leqslant n-1}(X[1])}\\ \end{array}
\]

反通勤。

\statementlabel{练习 X.2}令  $(X, \mathcal{O}_X)$  为诺特分离方案，使得  $\mathcal{O}_{X,x}$  是任何  $x$  的常规环，即  $\mathrm{gld}(\mathcal{O}_{X,x}) < \infty$  。我们假设  $X$  具有等维  $d$  。令  $\mathbf{D}_{\mathrm{coh}}^b(\mathcal{O}_X)$  为  $\mathbf{D}^b(\mathcal{O}_X)$  的满子范畴，由具有相干上同调的对象组成。定义从  $\mathbf{D}_{\mathrm{coh}}^b(\mathcal{O}_X)$  到  $\mathbf{D}_{\mathrm{coh}}^b(\mathcal{O}_X)$  的对偶函子  $\mathrm{D}(F) = \mathcal{R}\mathcal{H}\mathcal{o}m_{\mathcal{C}_x}(F, \mathcal{O}_X)[d]$  。

(i) (a) 证明  $ \mathbf{D}^{b}(\mathrm{Coh}(\mathcal{O}_{X})) \simeq \mathbf{D}_{\mathrm{coh}}^{b}(\mathcal{O}_{X}) $  ，其中  $ \mathrm{Coh}(\mathcal{O}_{X}) $  是相干  $ \mathcal{O}_{X} $  模的阿贝尔范畴； （这适用于所有诺特方案  $ X $  ）。

(b) 证明 $ D \circ D \simeq id $  。

(ii) 让我们定义  $ \mathbf{D}_{\mathrm{coh}}^{b}(\mathcal{O}_{X}) $  的满子范畴

\[
{}^{p}\mathbf{D}_{\mathrm{c o h}}^{{\leq}0}(\mathcal{O}_{X})=\{F\in\mathbf{D}_{\mathrm{c o h}}^{b}(\mathcal{O}_{X});\mathrm{d i m}\mathrm{s u p p}(\tau^{>p(k)}F)<k\mathrm{f o r~a n y}k\}
\]

 $ ^{p}\mathbf{D}_{\mathrm{coh}}^{\geq0}(\mathcal{O}_{X}) =  \{F \in \mathbf{D}_{\mathrm{coh}}^{b}(\mathcal{O}_X); H_Z^j(F) = 0 \text{ for any closed subset } Z \text{ and } j < p(\dim Z)\} $ 。

（这里 p 表示反常。）

(a) 证明 D 将  $ ^p\mathbf{D}_{\mathrm{coh}}^{\leq 0}(\mathcal{O}_X) $  和  $ ^p\mathbf{D}_{\mathrm{coh}}^{\geq 0}(\mathcal{O}_X) $  分别发送到  $ ^p\mathbf{D}_{\mathrm{coh}}^{\geq 0}(\mathcal{O}_X) $  和  $ ^p\mathbf{D}_{\mathrm{coh}}^{\leq 0}(\mathcal{O}_X) $  中。

(b) 证明 $ ^{p}\mathbf{D}_{\mathrm{coh}}^{\leq0}(\mathcal{O}_{X}),^{p}\mathbf{D}_{\mathrm{coh}}^{\geq0}(\mathcal{O}_{X}) $  是一个t 结构。

(c) 证明 $ U \mapsto^p \mathbf{D}_{\text{coh}}^0(\mathcal{O}_U) $  是一个堆栈（参见命题10.2.9）。

（提示：使用 A. Grothendieck [5] 的结果：对于任何相干理想 I 和

任何具有准相干上同调层的 $F \in \mathrm{Ob}(\mathbf{D}^b(\mathcal{O}_X))$ ，设置 $Z = \mathrm{supp}(\mathcal{O}_X/I)$ ，我们有 $H_Z^n(F) \lesssim \lim_{\rightarrow} \mathcal{E} x \ell^n_{\mathcal{O}_X}(\mathcal{O}_X/I^k, F).$ 

\statementlabel{练习 X.3}设 A 为离散估值环。

证明  $\mathbf{D}^{b}(\mathfrak{Mod}^{f}(A))$  上的任何  $t$  -结构要么是自然结构，要么是其对偶向上移位（换句话说，设置  $X = \text{Spec}(A)$  ，对于某些反常  $p$  来说，它是  $(^p\mathbf{D}_{\text{coh}}^{≤0}(\mathcal{O}_X),^p\mathbf{D}_{\text{coh}}^{≥0}(\mathcal{O}_X))$  ，（参见练习 X.2）。（提示：使用练习 I.18 并用于  $\mathbf{D}^{b}(\mathfrak{Mod}^{f}(A))$  上的  $t$  -结构  $(\mathbf{D}^{\leq0},\mathbf{D}^{\geq0})$  ，首先证明 $A\in\text{Ob}(\mathbf{D}^{\geq0})$ ， $A\notin\text{Ob}(\mathbf{D}^{\geq1})$ 暗示 $A\in\text{Ob}(\mathbf{D}^{\leq0})$ 。）

\statementlabel{练习 X.4}令  $X$  为实解析流形，并令  $F \in \mathrm{Ob}(\mathbb{D}_{w-R-c}^{b}(X))$  。设置 $d = \mathrm{dim}\,\mathrm{supp}(F)$ 。假设 $F \in \mathrm{Ob}({}^{p}\mathbb{D}_{w-R-c}^{≤0}(X))$ （或 $^{p}\mathbb{D}_{w-R-c}^{≥0}(X)$ ）证明 $H^{j}(X; F) = 0$ 对于 $j > p(d) + d$ （或 $H_{c}^{j}(X; F) = 0$ 对于 $j < p(d)$ ）。

\statementlabel{练习 X.5}令 $X$  为实解析流形， $U$  为子解析开子集， $j$  为嵌入 $U \hookrightarrow X$ ，并令 $p$  为反常。对于  $F \in \mathrm{Ob}({}^p\mathbf{D}_{w-\mathrm{R}-c}^0(X))$  ，定义了  $p_j, j^{-1}F = ^{p}H^0(j, j^{-1}F)$  和  $p_j, j^{-1}F = ^{p}H^0(Rj, j^{-1}F)$  ，（参见 X§1）。证明以下两个条件等价：

(i)  $F$  与  $p \mathbf{D}_w - \mathbf{R} - c(X)$  中的  $p j_! j^{-1} F \to p j_* j^{-1} F$  的图像同构，

(ii)  $F$  既没有  $^p\mathbf{D}_{w-R-c}^0(X)$  中的  $X \setminus U$  支持的非零子对象，也没有非零商对象。

\statementlabel{练习 X.6}令 $X$  为复流形。让 $F \in \mathrm{Ob}(\mathbf{D}_{w-C-c}^{b}(X))$  。证明 $\mathrm{SS}(F) = \bigcup_{j} \mathrm{SS}(^{p} H^{j}(F))$ 。

\statementlabel{练习 X.7}令  $X$  为复流形，让  $0 \to F' \to F \to F'' \to 0$  为  $p\mathbf{D}_{w-c-c}^0(X)$  中的精确序列。证明 $\mathrm{SS}(F) = \mathrm{SS}(F') \cup \mathrm{SS}(F'')$ 。

\statementlabel{练习 X.8}我们假设基环A是一个场。设 X 为 n 维复流形，F 为 X 上的反常层。

(i) 令  $Z$  为  $X$  的闭子集。证明  $H^{-n}(F)|_{Z}$  满足“解析连续原理”，即该束的任意部分对  $Z$  的开子集  $U$  的支持，在  $U$  中既开又闭。

(ii) 现在假设  $Z$  是  $X$  的闭子解析子集，并让  $x \in Z$  是  $Z$  的非孤立点。证明 $(H_{Z}^{n}(F))_{x}=0$ 。 （提示：使用对偶论证。）

\statementlabel{练习 X.9}设 Y 为复流形 X 的复超曲面。证明  $ M_{Y}[\dim X-1] $  是任意 A 模 M 的反常层。

\statementlabel{练习 X.10}设  $F \in \mathrm{Ob}(\mathbb{D}_{w-R-c}^b(X))$  ，设  $Z$  为  $X$  的局部封闭子解析子集，并设  $j \in \mathbb{Z}$  。证明  $H_Z^j(F) = 0$  对于所有  $j < r$  当且仅当  $\dim(Z \cap \cos \mathrm{supp}^j(F)) < j - r$  对于任何  $j$  。

\subsection{注记}
在反常层理论的起源中，人们一方面发现了 Goresky-MacPherson [1] 的“交交上同调”，这是将庞加莱对偶性推广到奇异空间的成功尝试，另一方面发现了完整  $\mathcal{D}$  模理论（这两种理论密切相关，参见第十一章）。事实上，早在 1975 年，Kashiwara [3] 就表明，如果  $\mathcal{M}$  是复流形  $X$  上的完整  $\mathcal{D}$  -模，那么复数  $R\mathcal{H}_{\mathcal{O}m_{\mathcal{D}}}(\mathcal{M}, \mathcal{O}_X)$  [  $\dim_C X$  ] 是反常的，他给出了“黎曼-希尔伯特问题”的以下表述（参见 Ramis [1, p.287]）：完整  $\mathcal{D}$  模的阿贝尔范畴的阿贝尔子范畴“holreg”，使得函子  $R\mathcal{H}_{\mathcal{O}m_{\mathcal{D}}}(\ast, \mathcal{O}_X)$  诱导从  $D^b(\text{holreg})$  到  $D_{\mathcal{C}-c}^b(X)$  的范畴的等价性。请注意，这个问题在 80 年代得到了解决（参见 Kashiwara [6] 和 Mebkhout [1]）。这样的对应关系立即意味着范畴 $D_{\mathcal{C}-c}^b(X)$ 包含一个阿贝尔子范畴，该子范畴对应于 $D^b(\text{holreg})$ 的阿贝尔子范畴holreg：这是反常层的范畴。

反常层理论源自 Gabber 和 Beilinson-Bernstein-Deligne [1]，本章的大部分结果都是众所周知的。 §1 中有关  $t$  结构的内容摘自 Beilinson-Bernstein-Deligne（前述）以及定理 10.2.8。断言消失循环函子保留了反常性的关键结果来自 Goresky-MacPherson [3]（在 étale 案例中的 Beilinson-Bernstein-Deligne（前述）之后）。当基环为  $\mathbb{C}$  时，在 Kashiwara-Schapira [3] 中进行了反常的微局部解释。最后，我们还要提一下 MacPherson-Vilonen [1] 中反常层的有趣构造。

\section{对层模与微分算子模的应用}
\subsection{概要}
根据其定义，复流形 $X$ 被赋予了全纯函数的环层 $\mathcal{O}_X$ 。  $\mathcal{O}_X$  的结构和  $\mathcal{O}_X$  模的理论现在已经很好理解了，我们的目的并不是从一开始就解释这个理论。我们将满足于有关  $\mathcal{O}_X$  的代数结构、其松弛维数以及  $\mathcal{O}_X$  的运算的一些基本事实。参考Banica-Stanasila [1]、Cartan [2]、Hörmander [1]、Serre [1]。

接下来，我们介绍X上的有限阶全纯微分算子的环 $ \mathcal{D}_x $ 。这里 $ \mathcal{D}_x $ 模的理论现在已经很好理解了，我们将相当简短地讨论这个主题，记住要理解主要概念，特别是特征变化的概念，并解释 $ \mathcal{D}_x $ 模上的运算。我们还应回顾经典的柯西-科瓦勒夫斯基定理及其对  $ \mathcal{D}_x $  模的扩展，并且我们将推导出公式：

\[
\mathrm{S S}(R\mathcal{H o m}_{\mathcal{D}_{x}}(\mathcal{M},\mathcal{O}_{X}))\subset\mathrm{c h a r}(\mathcal{M})\enspace,
\]

其中  $ \operatorname{char}(\mathcal{M}) $  是  $ \mathcal{D}_x $  -模  $ \mathcal{M} $  的特征变体。 （事实上​​，这种包含是一种等式，参见§4。）作为（11.0.1）的应用，借助VIII§5的结果，当 $ \mathcal{M} $ 是完整的时，我们立即获得了复数 $ R\mathcal{H}om_{\mathcal{Q}_x}(\mathcal{M},\mathcal{O}_x) $ 的可构造性，并且也很容易证明这个复数是反常的。

对于  $\mathcal{D}_X$  模理论的更详细阐述，我们参考 Björk [1]、Kashiwara [5] 和 Schapira [2]。

然后我们“微观地”研究层 $ \mathcal{O}_X $ 。在介绍了微局域算子的环  $ \delta_X^R $  之后，我们勾画了一个重要定理的证明，该定理断言可以在  $ \mathcal{O}_X $  上局部“量化”全纯接触变换。我们通过介绍实解析流形  $ M $  上 Sato 微函数的束  $ \mathcal{C}_M $  来结束本章。使用（11.0.1）以及第五章和第六章的结果，可以很容易地恢复线性偏微分方程理论的许多经典结果，特别是有关椭圆方程或解析波前集、微双曲系统和奇点传播的结果。

应该清楚的是，本章的目的不是对解析（微）微分方程理论进行完整或系统的处理，而是

而是向读者介绍它，特别是让读者在层微支撑理论的指导下更好地理解 Sato-Kawai-Kashiwara [1] 的基本论文。

在本章中，除非另有说明，所有层都是 C 向量空间的层。

\subsection{结构层}
令  $X$  为复维度  $n$  的复流形，并令  $\mathcal{O}_X$  为  $X$  上的全纯函数环层。首先，我们将描述  $X$  和底层实解析流形  $X^R$  之间的关系。

其中  $\bar{X}$  表示  $X$  的复共轭。这是一个复流形，例如  $\bar{X}^{\mathbf{R}} = X^{\mathbf{R}}$  ，但  $\bar{X}$  上的全纯函数是  $X$  上的反全纯函数。因此，我们有环的同构：  $\mathcal{O}_{X} \overset{\sim}{\to} \mathcal{O}_{\bar{X}}$  、  $u \mapsto \bar{u}$  和  $\bar{au} = \bar{au}$  对于  $a \in \mathbb{C}$  和  $u \in \mathcal{O}_{X}$  。

用  $ X \times \bar{X} $  中的对角线识别  $ X^{\mathbb{R}} $  。那么  $ X \times \bar{X} $  是  $ X^{\mathbb{R}} $  的复化。如果  $ \mathcal{A}_{X^{\mathbb{R}}} $  表示  $ X^{\mathbb{R}} $  上的一组（复值）实解析函数，则  $ \mathcal{A}_{X^{\mathbb{R}}} = \mathcal{O}_{X \times \bar{X}}|_{X^{\mathbb{R}}} $  。

而且

\[
T X^{\mathbb{R}}\mathop{{\otimes}}_{\mathbb{R}}\mathbb{C}\simeq T X\mathop{{\times}}_{X}T\bar{X}\enspace.
\]

合成图

\[
T X^{\mathbb{R}}\to T X_{\underbrace{\mathbb{R}}}^{\mathbb{R}}\otimes\mathbb{C}\simeq\mathop{T X}\limits_{\substack{\times X}}\bar{T X}\to T X
\]

定义实解析向量束在  $ X^{R} $  上的同构：

\[
T X^{\mathbb{R}}\simeq(T X)^{\mathbb{R}}
\]

并通过对偶性同构：

\[
T^{*}X^{\mathbb{R}}\xrightarrow{\sim}(T^{*}X)^{\mathbb{R}}.
\]

对于  $X^R$  上的实  $C^1$  -函数  $\phi$  ，最后一个同构将向量  $\partial\phi(x) \in T_x^*X$  与向量  $\mathrm{d}\phi(x) \in T_x^*X^R$  相关联，其中  $\mathrm{d}\phi$  是实微分， $\partial\phi$  是其全纯分量。

令 $ \alpha_X $ （分别为 $ \alpha_{X^R} $ ）为 $ T^*X $ （分别为 $ T^*X^R $ ）上的复数（分别为实数）规范1-形式。然后 $ \alpha_{\bar{X}} = \overline{\alpha_X} $  和 $ \alpha_{X \times \bar{X}} = \alpha_X + \alpha_{\bar{X}} $  。因此：

\[
\alpha_{X^{\mathtt{R}}}=2\mathrm{R e}\alpha_{X}~.
\]

令 $ z = (z_1, \ldots, z_n) $  为 $ X $  上的全纯坐标系， $ (z; \zeta) $  为 $ T^*X $  上的相关坐标，以便 $ \alpha_X = \sum_j \zeta_j \, dz_j $  。如果  $ z = x + \sqrt{-1}y $  、 $ \zeta = \xi + \sqrt{-1}\eta $  ，则：

\[
\alpha_{x^{n}}=2\sum_{j}\left(\xi_{j}\mathrm{d}x_{j}-\eta_{j}\mathrm{d}y_{j}\right).
\]

对于 $ T^*X^R $  上的实值函数 $ h $ ，它导出哈密顿场 $ H_h^R $  。

从现在开始，如果没有混淆的风险，我们将识别 $ (T*X)^{\mathbb{R}} $ 和 $ T*X^{\mathbb{R}} $ ，甚至有时写 $ T*X $ 或 $ X $ 而不是 $ T*X^{\mathbb{R}} $ 或 $ X^{\mathbb{R}} $ 。

接下来，让我们回顾一下与一致性相关的一些基本概念。

令 $(X, \mathcal{A}_X)$  为环化空间。除非另有说明， $\mathcal{A}_X$  模表示左 $\mathcal{A}_X$  模。

如果  $\mathcal{A}_x$  -模  $\mathcal{M}$  对于某些  $N \in \mathbb{N}$  与  $\mathcal{A}_x^N$  同构，则它是有限自由的。如果每个 $x \in X$  都有一个开邻域 $U$ ，使得 $\mathcal{M}|_U$  是 $\mathcal{A}_x|_{U}$  -有限自由，则它是局部有限自由的。

 $\mathcal{A}_{X}$  -模  $M$  的 s-表示是  $\mathcal{A}_{X}$  -模的精确序列：

\[
{\mathcal{M}}_{s}\to\cdots\to{\mathcal{M}}_{0}\to{\mathcal{M}}\to0~.
\]

如果  $ \mathcal{M}_{j} $  具有相应的属性，则可以说此表示是有限自由或局部有限自由等。

长度  $m$  的分辨率是  $\infty$  的  $\infty$  呈示，其中  $\mathcal{M}_{j}=0$  表示  $j>m$  。

如果  $\mathcal{A}_X$  -模  $\mathcal{M}$  局部承认有限类型（或有限表示），则它局部承认有限自由 0 表示（或 1 表示），即，如果局部在  $X$  上，则存在精确序列  $\mathcal{A}_X^N \to \mathcal{M} \to 0$ （或  $\mathcal{A}_X^{N_1} \to \mathcal{A}_X^{N_0} \to \mathcal{M} \to 0$  ）。

\statementlabel{定义 11.1.1}(i) 如果 $\mathcal{M}$  是局部有限类型的，并且如果对于任何开集 $U$ ，任何局部有限类型的子 $\mathcal{A}_{U}$  模是局部有限表示的，则可以说 $\mathcal{A}_{x}$  模 $\mathcal{M}$  是一致的。

(ii) 如果 $\mathcal{A}_{X}$  作为左 $\mathcal{A}_{X}$  模是一致的，则可以说 $\mathcal{A}_{X}$  是一致的。

(iii) 如果  $\mathcal{A}_X$  是相干的，则可以说  $\mathcal{A}_X$  是诺特式的，对于每个  $x \in X$  ，茎  $\mathcal{A}_{X,x}$  是诺特式的，最后，对于  $X$  的任何开放子集  $U$  ，相干  $\mathcal{A}_X|_{U}$  模的任何增加的相干子模族都是局部平稳的。

让 $\mathrm{Mod}_{\mathrm{coh}}(\mathcal{A}_X)$  表示相干环 $\mathcal{A}_X$  上的左相干模的范畴。那么 $\mathrm{Mod}_{\mathrm{coh}}(\mathcal{A}_X)$  是一个阿贝尔范畴。

让我们回到 X 是复维度 n 的复流形的情况。

\statementlabel{定理 11.1.2}(i) 该环层 $ \mathcal{O}_{X} $  是诺特环的（特别是连贯的）。

(ii) 让 $\mathcal{M}$  成为一个连贯的 $\mathcal{O}_{X}$  模。然后在  $X$  上本地， $\mathcal{M}$  承认长度为  $n$  的有限自由解析。

 $ \mathcal{O}_{X} $  的一致性是 Oka [1] 的结果。性质 (ii) 被称为希尔伯特 syzygy 定理。

\statementlabel{定理 11.1.3}令  $Z$  为  $X$  的局部闭子集，并令  $x \in X$  、  $x \notin \mathrm{Int} Z$  。然后：

\[
H_{Z}^{j}(\mathcal{O}_{X})_{x}=0\qquad{f o r\quad}\quad j\notin[1,n]\enspace.
\]

当然， $ H_Z^0(\mathcal{O}_X)_x $ 的消失相当于众所周知的解析连续原理。  $ j > n $  的  $ H_Z^j(\mathcal{O}_X)_x $  消失是由于 Malgrange [1]（参见 II §9.11）。

最后我们将讨论  $\mathcal{O}_{x}$  上的操作。

让 $ (Y,\mathcal{O}_{Y}) $  成为另一个复流形。  $ X\times Y\colon $  上存在层的自然态射

\[
\mathcal{O}_{X}\boxtimes\mathcal{O}_{Y}\to\mathcal{O}_{X\times Y}\enspace.
\]

现在假设给出一个全纯映射  $f: Y \to X$  。  $Y$  上有层的自然态射：

\[
f^{-1}\mathcal{O}_{X}\to\mathcal{O}_{Y}~.
\]

这个态射只不过是  $f$  的组合，它将  $X$  上的全纯函数  $\phi$  发送到  $Y$  上的全纯函数  $\phi \circ f$  。直接的形象并不那么容易描述。

让  $ \mathcal{O}_X^{(p)} $  表示  $ X $  上的全纯 p 型束，并设置：

\[
\Omega_{x}=\mathcal{O}_{X}^{(n)}\otimes\mathcal{o r}_{X}
\]

其中  $ or_{X} $  表示  $ X^{R} $  和  $ n = \dim_{C} X $  上的方向束。

当然，复流形 X 是可定向的，并且大多数时候，人们可能会忘记  $ o t_{X} $  。

束  $\Omega_{X}$  是一个可逆的  $\mathcal{O}_{X}$  模，即它在  $\mathcal{O}_{X}$  上局部没有排名第一。如果  $\mathcal{L}$  是可逆的  $\mathcal{O}_{X}$  -模，则设置：

\[
\mathcal{L}^{\otimes-1}=\mathcal{H o m}_{\mathcal{O}_{X}}(\mathcal{L},\mathcal{O}_{X})\enspace.
\]

当然 $ \mathcal{L}^{\otimes-1}\otimes_{\mathcal{O}_x}\mathcal{L}\simeq\mathcal{O}_x $ 。

体积元素是  $\Omega_{x}$  和  $\mathcal{O}_{x}$  上的生成器。体积元素的存在仅在本地得到保证。

\statementlabel{定理 11.1.4}$ \mathbf{D}^{b}(X) $  中存在自然态射：

\[
R f_{!}\Omega_{Y}[\dim_{\mathbb{C}}Y]\to\Omega_{X}[\dim_{\mathbb{C}}X].
\]

此外，这种态射对于映射的组成是有用的。

回想一下（参见 II.§9），该态射可以通过如下方式获得。

令 $\mathcal{D}\ell_{Y}^{(p,q)}$  为关于 $Y$  和 $\bar{Y}$  的二度 $(p,q)$  的 $Y^{\mathbf{R}}$  上的分布表。设置 $m = \dim_{\mathbb{C}}Y$ 、 $n = \dim_{\mathbb{C}}X$ 、 $l = m - n$ 。存在“整合”态射

\[
f_{!}\mathcal{D}\ell_{{\mathbb{Y}}^{{\mathfrak{R}}}}^{(p,q)}\otimes o\mathfrak{r}_{Y}\to\mathcal{D}\ell_{{\mathbb{X}}^{{\mathfrak{R}}}}^{(p-l,q-l)}\otimes o\mathfrak{r}_{X}
\]

通过用  $\mathcal{D}\ell_Y^{(m,\cdot)}$  和  $\mathcal{D}\ell_X^{(n,\cdot)}$  中的系数分别替换  $\Omega_Y$  和  $\Omega_X$  的 Dolbeault 分辨率，可以得到所需的态射。

作为该定理的应用，让我们构造与（复）余维 d 的 X 的子流形 Z 相关的基本类。

令  $\mathcal{I}$  为  $X$  中  $Z$  的定义理想，即由  $Z$  上消失的部分生成的  $\mathcal{O}_{X}$  的理想束。存在自然态射：

\[
\mathcal{I}/\mathcal{I}^{2}\to\mathcal{O}_{X}^{(1)}\otimes_{\mathcal{C}_{X}}\mathcal{O}_{Z}
\]

由 $f \mapsto df$  给出。从这个态射可以得到：

\[
A^{\mathsf{d}}(\mathcal{I}/\mathcal{I}^{2})\to\mathcal{O}_{X}^{(\mathsf{d})}\otimes_{\mathcal{O}_{X}}\mathcal{O}_{Z}.
\]

但是 $ \bigwedge^d(\mathcal{F}/\mathcal{F}^2) \simeq \mathcal{O}_Z^{(n-d)\otimes-1} \otimes_{\mathcal{O}_X} \mathcal{O}_X^{(n)} $ 。因此我们得到态射：

\[
\mathcal{O}_{Z}\to\mathcal{O}_{Z}^{(n-d)}\otimes_{\mathcal{O}_{X}}\mathcal{O}_{X}^{(d)}\otimes\mathcal{O}_{X}^{(n)\otimes-1}\enspace.
\]

另一方面，定理 11.1.4 定义了态射：

\[
\Omega_{z}\to H_{z}^{d}(\Omega_{x})\;.
\]

将其张量积与  $ \mathcal{O}_X^{(d)} \otimes \mathcal{O}_X^{(n) \otimes^{-1}} $  结合，并与 (11.1.12) 结合，我们得到态射：

\[
\mathcal{O}_{Z}\to H_{Z}^{d}(\mathcal{O}_{X}^{(d)})\otimes\mathfrak{o r}_{Z/X}\;.
\]

\statementlabel{定义 11.1.5}$O_{Z}$  的第 1 部分通过态射 (11.1.13) 的图像称为 X 中 Z 的基本类，并表示为  $\delta_{Z}$  。

示例 11.1.6。假设  $Z$  由  $f_1 = \cdots = f_d = 0$  和  $Z$  上的  $\mathrm{d}f_1 \wedge \cdots \wedge \mathrm{d}f_d \neq 0$  方程定义。那么  $\delta_Z$  就是  $\frac{1}{(2\pi i)^d}\left(\frac{\mathrm{d}f_1}{f_1} \wedge \cdots \wedge \frac{\mathrm{d}f_d}{f_d}\right)$  的类。

这里 $1/f_1 \ldots f_d$ 是通过考虑 $X \setminus Z$ 的开覆盖 $\{f_j \neq 0\}_{j=1,\ldots,d}$ 而获得的 $H^{d-1}(X \setminus Z; \mathcal{O}_X)$ 元素的 $H_Z^d(X; \mathcal{O}_X)$ 中的图像（参见练习II.26）。

\subsection{微分算子模}
首先，让我们回顾一下有关滤波环和模的一些基本事实（参见 Schapira [2]）。

所有环都是酉的（即单位为 1），除非另有说明，“模”表示“左酉模”（即 1 的作用为恒等式的左模）。

由  $ \mathbb{Z} $  索引的过滤环  $ A $  是一个赋予一系列子群  $ \{A_k\}_{k \in \mathbb{Z}} $  的环（仍用  $ A $  表示），使得：

\[
\left\{\begin{matrix}{A_{k}\subset A_{k+1}{~f o r~a n y~}k,1\in A_{0},A_{k}\cdot A_{l}\subset A_{k+l}{~f o r~a n y~}k,l{~a n d}}\\ {A=\bigcup_{k}A_{k}.}\\ \end{matrix}\right.
\]

 $A$  上的过滤模 $M$  是 $A$  模 $M$ ，赋予一系列子组 $\{M_k\}_{k\in\mathbb{Z}}$ ，使得：

\[
M_{k}\subset M_{k+1}{~f o r~a n y~}k,A_{l}\cdot M_{k}\subset M_{k+l}{~f o r~a n y~}k,l{~a n d~}M=\bigcup_{k}M_{k}.
\]

如果存在  $r\in\mathbb{N}$  使得  $M_{k-r}\subset M_k'\subset M_{k+r}$  对于所有  $k\in\mathbb{Z}$  ，则  $M$  上的两个过滤  $\{M_k\}_{k}$  和  $\{M'_k\}_{k}$  被认为是等效的。特别是  $M$  上的过滤等效于其  $r$  移位过滤，表示  $M_{[r]}$  ，定义为：

\[
(M_{[\mathfrak{r}]})_{k}=M_{k+r}\qquad{f o r\quad}\quad k\in\mathbb{Z}\enspace.
\]

过滤  $A$ -模  $\psi: M \to N$  的态射是  $A$ -模 的态射，使得  $\psi(M_k) \subset N_k$  对于所有  $k$  。

考虑  $A$  -模  $0 \to L \to M \xrightarrow[\psi]{\phantom{z}}\;N \to 0$  的精确序列，并假设  $M$  具有过滤功能。  $L$  上的诱导过滤是通过设置  $L_k = L \cap M_k$  来定义的。  $N$  上的图像过滤是通过设置 $N_k = \psi(M_k)$  来定义的。

如果过滤后的 A 模  $ 0 \to L \to M \to N \to 0 $  的序列是 A 模的精确序列，并且  $ L $  （分别为  $ N $  ）上的过滤是诱导（分别为图像）过滤，则该序列被称为严格精确。还可以通过设置定义两个滤波模  $ M $  和  $ M' $  的滤波直和  $ M \oplus M' $ ：

\[
(M\oplus M^{\prime})_{k}=M_{k}\oplus M_{k}^{\prime}~.
\]

因此，过滤后的 A 模的范畴是可加的。人们应该注意这一范畴不是阿贝尔的：例如过滤态射 $ A \to A_{[1]} $  是单态和泛同态，但通常不是过滤A 模的同构。

根据定义，有限自由过滤  $ A $  模是过滤  $ A $  模，同构于  $ A_{[r]} $  类型模的有限直和。

过滤 A 模 M 的有限自由 s 表示是过滤 A 模的严格精确序列：

\[
L_{s}\to\cdots\to L_{0}\to M\to0
\]

其中  $ L_{j} $  是有限自由的。

如果过滤模  $ M $  允许有限自由 0 表示，则称其为有限类型。在这种情况下，有人说 $ M $  上的过滤效果很好。

如果过滤 $A$  模范畴中有限类型模的任何子对象是有限类型，则过滤环 $A$  被称为诺特环（或“过滤诺特环”）。这相当于说，对于任何被赋予良好过滤的 $A$  -模 $M$  和被赋予过滤的 $M$  的任何子模 $N$ ，使得 $N \to M$  是经过过滤的态射，这种过滤是好的。如果滤波环  $A$  是诺特环，而且对于任何  $\alpha \in A_{-1}$  ，  $1 - \alpha$  是可逆的，那么就说  $A$  是扎里斯克环。

滤波诺特环的定义与 Schapira [2] 第二章的定义 1.1.2 略有不同。事实上，正如 Van Oystaeyen 所指出的，Schapira（上文引用的）的定义太弱，不足以暗示  $ \operatorname{gr}(A) $  是诺特主义，特别是该作者的命题 1.1.7 部分是错误的。 （这个错误对本书的其余部分没有影响。要纠正它，用我们上面刚刚给出的定义替换滤波诺特环的弱定义就足够了。）

分级环  $ \operatorname{gr}(A) $  的定义如下：

\[
\mathbf{gr}(A)=\bigoplus_{k}A_{k}/A_{k-1}\ .
\]

类似地，如果  $M$  是经过过滤的  $A$  -模，则通过以下方式定义（分级） $\mathrm{gr}(A)$  -模  $\mathrm{gr}(M)$ ：

\[
\operatorname{g r}(M)=\bigoplus_{k}M_{k}/M_{k-1}.
\]

元素  $ u \in M $  的阶数是  $ k \in \mathbb{Z} \cup \left\{ -\infty \right\} $  中最小的  $ u \in M_k $  ，其中按照惯例为  $ M_{-\infty} = 0 $  。它由  $ \mathrm{ord}(u) $  表示。

1 用 $ \sigma_k $  表示投影 $ M_k \to \mathrm{gr}(M) $  。如果  $ \mathrm{ord}(u) = k $  ，则写  $ \sigma(u) $  而不是  $ \sigma_k(u) $  。

\statementlabel{命题 11.2.1}假设  $ \operatorname{gr}(A) $  是诺特环（作为分级环）， $ A_{k}=0 $  是 k<0。然后A被Noetherian过滤。

证明是一个简单的练习。

现在假设 $ \mathrm{gr}(A) $  是可交换的。通过设置  $ \bar{a} \in A_k/A_{k-1} $  、  $ \bar{b} \in A_l/A_{l-1} $  来定义  $ \mathrm{gr}(A) $  上的泊松括号：

\[
\{\bar{a},\bar{b}\}=\sigma_{k+l-1}([a,b])
\]

其中  $a$  （分别为  $b$  ）是  $A_k$  （分别为  $A_l$  ）和  $\sigma_k(a) = \bar{a}$  （分别为  $\sigma_l(b) = \bar{b}$  ）、  $[a,b] = ab - ba$  的任意元素。

令 $M$  为有限类型的 $A$  模。人们可以通过选择一组生成器 $(v_j)_{j=1}^N$ 并设置 $M_k = \sum_j A_k v_j$ 来赋予它良好的过滤。此外，有限类型模上的两个良好过滤显然是等效的。这样就很容易证明 $\mathrm{gr}(M)$ 的歼灭子根式是 $\mathrm{gr}(A)$ 的分级理想，它只依赖于 $M$ ，而不依赖于良好的过滤。其中之一用  $\mathrm{Icar}(M)$  表示。因此：

\[
\begin{cases}{\underset{}{\bar{a}\in\operatorname{I c a r}(M),}{\bar{a}}{~h o m o g e n e o u s~o f~o r d e r~k}\Leftrightarrow\exists l\in\mathbb{N},\forall a\in A_{k l}}\\ {{w i t h~}\sigma(a)=\underset{}{\bar{a}^{k},a}{M}_{n}\subset M_{n+k l-1}\forall n\enspace.}\\ \end{cases}
\]

假设 A 是经过过滤的诺特矩阵，并考虑一个精确序列：

\[
0\to L\to M\to N\to0
\]

有限类型的 A 模。人们可以赋予 $M$ 良好的过滤， $L$ （分别 $N$ ）具有诱导（分别图像）过滤。那么(11.2.8)就变成了过滤模的精确序列，并且该序列：

\[
0\to\mathbf{g r}(L)\to\mathbf{g r}(M)\to\mathbf{g r}(N)\to0
\]

本身是准确的。这意味着：

\[
\mathrm{Icar}(M)=\mathrm{Icar}(L)\cap\mathrm{Icar}(N)\enspace.
\]

此外，由于 Gabber [1]，我们得到了重要的结果。

\statementlabel{定理 11.2.2}假设  $ \operatorname{gr}(A) $  是诺特交换 Q 代数。那么对于任何有限类型的 A 模 M， $ \operatorname{Icar}(M) $  是对合的，即：

\[
\left\{\operatorname{Icar}(M),\operatorname{Icar}(M)\right\}\subset\operatorname{Icar}(M).
\]

现在让 $X$  成为一个拓扑空间。  $X$  上的滤过环层  $\mathcal{A}$  是环层  $\mathcal{A}$  赋予了一系列子群  $\{\mathcal{A}_k\}_{k\in\mathbb{Z}}$  ，使得  $\mathcal{A}_k \subset \mathcal{A}_{k+1}$  , 1 是  $\mathcal{A}_0$  的一部分， $\mathcal{A}_k \cdot \mathcal{A}_l \subset \mathcal{A}_{k+l}$  和  $\mathcal{A}$  是  $\mathcal{A}_k$  的并集。类似地定义了过滤 $\mathcal{A}$ 模、有限自由过滤 $\mathcal{A}$ 模等的概念。

如果在  $X$  上局部存在过滤后的精确序列  $\mathcal{L} \to \mathcal{M} \to 0$  ，其中  $\mathcal{L}$  是有限自由的，则  $\mathcal{A}$  -模  $\mathcal{M}$  上的过滤被认为是好的。

\statementlabel{命题 11.2.3}让 $\mathcal{A}$  成为 $X$  上的一束过滤环。假设  $\mathrm{gr}(\mathcal{A})$  是相干的（作为分级环）（分别为诺特环），并且对于每个  $x\in X$  ，滤波环  $\mathcal{A}_{X,x}$  是 Zariskian。然后：

(i)  $ \mathcal{A} $  作为环层是连贯的（分别是诺特式的）。

(ii) 令 $\mathcal{M}$  成为具有良好过滤功能的 $\mathcal{A}$  模。那么 $\mathcal{M}$  作为 $\mathcal{A}$  模是一致的当且仅当 $\mathrm{gr}(\mathcal{M})$  是一致的（作为分级模）。

(iii) 假设 $\mathcal{M}$  是连贯的并且具有良好的过滤能力。让  $\mathcal{N}$  成为  $\mathcal{M}$  的连贯子模。那么 $\mathcal{N}$ 上的诱导过滤就好了。

对于证明，参见。比约克 [1] 或夏皮拉 [2]。

现在假设  $X$  是复维度  $n$  的复流形。我们用 $\mathcal{D}_x$  表示有限阶全纯微分算子的环层。这是由 $\mathcal{O}_X$  和向量场生成的 $\mathcal{Hom}_C(\mathcal{O}_X, \mathcal{O}_X)$  的子代数。

 $ \mathcal{D}_x $  的过滤通过设置递归定义：

\[
\left\{\begin{aligned}{}&{{}\mathcal{D}_{X}(m)=0\quad\mathrm{f o r}\quad m<0~,}\\ {}&{{}\mathcal{D}_{X}(m)=\left\{P\in\mathcal{D}_{X};[P,\mathcal{O}_{X}]\subset\mathcal{D}_{X}(m-1)\right\}~.}\\ \end{aligned}\right.
\]

特别是 $ \mathscr{D}_X(0) = \mathscr{O}_X $  。

分级环  $ \mathrm{gr}(\mathcal{D}_X) $  自然同构于  $ \mathcal{O}_{[T^*X]} $  ，其中  $ \mathcal{O}_{[T^*X]} $  表示  $ \pi_*(\mathcal{O}_{T^*X}) $  的子环，由向量束  $ \pi: T^*X \to X $  的纤维上的多项式部分组成。

令 $(x) = (x_1, \ldots, x_n)$  为 $X$  上的局部全纯坐标系， $(x; \xi)$  为 $T^*X$  上的相关坐标。阶  $m$  的微分算子  $P$  被写为：

\[
P(x;D_{x})=\sum_{|\alpha|\leqslant m}a_{\alpha}(x)D_{x}^{\alpha}~,
\]

其中  $ a_{\alpha} $  是全纯函数  $ \alpha = (\alpha_1, \ldots, \alpha_n) $  、  $ |\alpha| = \alpha_1 + \cdots + \alpha_n $  、  $ D_x^{\alpha} = D_1^{\alpha_1} \ldots D_n^{\alpha_n} $  和  $ D_j = \frac{\partial}{\partial x_j} $  （也表示为  $ D_{x_j} $  ）。

主符号  $ \sigma_{m}(P) $  是  $ T^{*}X $  上的函数：

\[
\sigma_{m}(P)(x;\xi)=\sum_{|\alpha|=m}a_{\alpha}(x)\xi^{\alpha}~.
\]

主符号  $ \sigma_{m}(P) $  本质上是在  $ T^{*}X $  上定义的。还可以考虑P的总符号，

\[
P(x,\xi)=\sum_{|\alpha|\leqslant m}a_{\alpha}(x)\xi^{\alpha}
\]

但这个函数取决于局部坐标  $ (x_{1}, \ldots, x_{n}) $  。两个微分算子的组合  $ P \circ Q $  由莱布尼茨公式给出（在选择局部坐标的情况下）：

\[
(P\circ Q)(x,\xi)=\sum_{\alpha\in\mathbb{N}^{n}}\frac{1}{\alpha!}\frac{\partial^{\alpha}}{\partial\xi^{\alpha}}P(x;\xi)\cdot\frac{\partial^{\alpha}}{\partial x^{\alpha}}Q(x;\xi)\quad,
\]

其中  $ \alpha! = \alpha_1! \ldots \alpha_n! $  代表  $ \alpha = (\alpha_1, \ldots, \alpha_n) \in \mathbb{N}^n $  。

应用命题 11.2.3 和 $ \mathcal{O}_{[T^{\star}x]} $  上的众所周知的结果，我们得到：

\statementlabel{命题 11.2.4}(i) 环 $ \mathcal{D}_X $  左右相干且诺特环。

(ii) 令 $\mathcal{M}$  成为具有良好过滤功能的连贯 $\mathcal{D}_X$  模。那么 $\mathrm{gr}(\mathcal{M})$ 是连贯的。此外，如果 $\mathcal{N}$ 是 $\mathcal{M}$ 的相干子模， $\mathcal{N}$ 上的诱导过滤是好的。

让 $\mathcal{M}$  成为一个连贯的 $\mathcal{D}_X$  模。我们可以局部赋予 $\mathcal{M}$ 良好的过滤并定义 $\mathrm{gr}(\mathcal{D}_X)$ 的分级理想 $\mathrm{Icar}(\mathcal{M})$ 。这是一个连贯的分级理想，仅取决于 $\mathcal{M}$ ，因此是全局定义的。它在  $T^*X$  中的零变体，即该理想的所有部分的公共零点的集合，称为  $\mathcal{M}$  的特征变体，并表示为 char  $(\mathcal{M})$  。通过（11.2.10），我们发现如果 $0 \to \mathcal{L} \to \mathcal{M} \to \mathcal{N} \to 0$ 是连贯的 $\mathcal{D}_X$ 模的精确序列，那么：

\[
\operatorname{c h a r}(\mathcal{M})=\operatorname{c h a r}(\mathcal{L})\cup\operatorname{c h a r}(\mathcal{N})~.
\]

假设 M 具有良好的过滤性并设置：

\[
\widetilde{{\bf g r}}({\mathcal{M}})={\mathcal{O}}_{T^{\bullet}X}\mathop{{\otimes}_{\pi^{-1}}{\mathcal{O}}_{[T^{\bullet}X]}}\pi^{-1}{\bf g r}({\mathcal{M}})\enspace.
\]

函子  $ \mathcal{O}_{T^{*}X} \otimes_{\pi^{-1}\mathcal{O}_{[T^{*}X]}} \pi^{-1}(\cdot) $  与  $ \mathcal{O}_{[T^{*}X]} $  模的范畴完全相同，其中一个具有：

\[
\mathrm{char}(\mathcal{M})=\mathrm{supp}(\widetilde{\mathrm{gr}}(\mathcal{M}))
\]

根据定理 11.2.2 和命题 11.2.3， $ \operatorname{char}(\mathcal{M}) $  是  $ T^*X $  的闭合解析子集，对  $ \mathbb{C}^\times $  的作用具有对合和二次曲线（参见练习 VIII.8）

特别是 $ \dim_{\mathbb{C}}(\text{char}(\mathcal{M})) \geq n $ 。当  $ \dim_{\mathbb{C}}(\text{char}(\mathcal{M})) \leq n $  时，有人说  $ \mathcal{M} $  是完整的。

示例 11.2.5。假设  $\mathcal{M}$  有一个生成器  $u$  ，其定义关系为  $P_j u = 0, j = 1, \ldots, N$  。让 $\mathcal{I}$  表示由 $P_j$  生成的 $\mathcal{D}_x$  的左理想。然后 $\mathcal{M} \simeq \mathcal{D}_v / \mathcal{I}$  ，并且：

\[
\mathrm{char}(\mathcal{M})=\left\{(x;\xi);\sigma(P)(x;\xi)=0\mathrm{for any}P\in\mathcal{P}\right\}
\]

当然，对于 $ j=1,\ldots,N\} $ ，集合 $ \operatorname{char}(\mathcal{M}) $ 可能严格小于集合 $ \{(x;\xi);\sigma(P_j)(x;\xi)=0 $ 。事实上，虽然  $ P_j $  生成  $ \mathcal{I} $  ，但  $ \sigma(P_j) $  可能不会生成  $ \operatorname{gr}(\mathcal{I}) $  。然而，理想的 $ \operatorname{gr}(\mathcal{I}) $  是局部有限生成的。

作为  $ \operatorname{char}(\mathcal{M}) $  的对合性的推论，我们得到：

\statementlabel{命题 11.2.6}让 $\mathcal{M}$  成为一个连贯的 $\mathcal{D}_X$  模。然后在  $X$  上本地， $\mathcal{M}$  承认长度为  $n$  的有限自由解析。

换句话说，在X本地，存在一个复形：

\[
0\longrightarrow\mathcal{D}_{X}^{N_{n}}\longrightarrow\cdots\longrightarrow\mathcal{D}_{X}^{N_{1}}\xrightarrow[P_{0}]{}\mathcal{D}_{X}^{N_{0}}\longrightarrow0\enspace,
\]

除了 0 度和  $ \mathcal{M} \simeq \mathcal{D}_{X}^{N_{0}} / \mathcal{D}_{X}^{N_{1}} P_{0} $  之外，这是精确的。

示例 11.2.7。设 Z 为 X 中复余维 d 的闭子流形。

同态 $ \mathcal{E}x\ell_{\mathcal{O}_X}^d(\mathcal{O}_Z,\mathcal{O}_X) \simeq H_Z^d(R\mathcal{H}om_{\mathcal{O}_X}(\mathcal{O}_Z,\mathcal{O}_X)) \to H_Z^d(\mathcal{O}_X) $  是单射的，并且用 $ \mathcal{B}_{Z|X} $  或 $ H_{[Z]}^d(\mathcal{O}_X) $  表示由其图像生成的 $ \mathcal{D}_X $  模。如果  $ (x_1, \ldots, x_n) $  是  $ X $  上的局部坐标系，使得  $ Z = \{x_1 = \cdots = x_d = 0\} $  ，则  $ \mathcal{B}_{Z|X} $  由  $ \frac{1}{x_1 \ldots x_d} $  的类  $ u $  生成，并且该生成器满足以下关系：

\[
x_{1}u=\cdots=x_{d}u=D_{d+1}u=\cdots=D_{n}u=0.
\]

换句话说：

\[
\mathcal{B}_{Z|X}\simeq\mathcal{D}_{X}/(\mathcal{D}_{X}x_{1}+\cdots+\mathcal{D}_{X}x_{d}+\mathcal{D}_{X}D_{d+1}+\cdots+\mathcal{D}_{X}D_{n})\enspace.
\]

请注意，  $\mathcal{B}_{Z|X}$  是连贯的，而  $\operatorname{char}(\mathcal{B}_{Z|X}) = T_Z^* X$  是连贯的。通过考虑  $\mathcal{D}_X$  中的规则序列  $(x_1, \ldots, x_d, D_{d+1}, \ldots, D_n)$  和相关的 Koszul 复合体（例如 Schapira [2 附录 B.4]），我们可以自由解析长度  $n$  的  $\mathcal{B}_{Z|X}$  。

令  $ \text{Mod}(\mathcal{D}_X) $  （分别为  $ \text{Mod}(\mathcal{D}_X^p) $  ）表示左（分别为右） $ \mathcal{D}_X $  模的阿贝尔范畴。

范畴 $ \operatorname{Mod}(\mathcal{D}_X) $  和 $ \operatorname{Mod}(\mathcal{D}_X^{\mathrm{op}}) $  是等效的。事实上，存在一种自然的同构：

\[
\mathcal{O}_{X}^{(n)}\simeq\mathcal{E}x\mathcal{t}_{\mathcal{D}_{X}}^{n}(\mathcal{O}_{X},\mathcal{D}_{X}),
\]

这使得  $ \mathcal{O}_X^{(n)} $  成为一个正确的连贯  $ \mathcal{D}_X $  模。如果 $ \mathcal{M} $  是左 $ \mathcal{D}_X $  模，则赋予 $ \mathcal{O}_X^{(n)} \otimes_{\mathcal{O}_X} \mathcal{M} $  右 $ \mathcal{D}_X $  模的结构，如下所示。

对于  $ v \in \mathcal{O}_X^{(n)} $  、  $ u \in \mathcal{M} $  和向量场  $ \theta $  ，一组：

\[
(v\otimes u)\theta=(v\theta)\otimes u-v\otimes\theta u,
\]

然后可以将此操作扩展到  $\mathcal{D}_X$  对  $\mathcal{O}_X^{(n)} \otimes_{\mathcal{O}_X} M$  的操作。类似地，为右  $\mathcal{D}_X$  -模  $\mathcal{N}$  定义  $\mathcal{D}_X$  对  $\mathcal{N} \otimes_{\mathcal{O}_X} \mathcal{O}_X^{(n) \otimes^{-1}}$  的左操作。

函子 $ \mathcal{M} \mapsto \mathcal{O}_X^{(n)} \otimes_{\mathcal{O}_X} \mathcal{M} $  和 $ \mathcal{N} \mapsto \mathcal{N} \otimes_{\mathcal{O}_X} \mathcal{O}_X^{(n) \otimes^{-1}} $  提供了所需的等价性。

用 $ \mathbf{D}_{\mathrm{coh}}^{b}(\mathcal{D}_{X}) $ （或 $ \mathbf{D}_{\mathrm{coh}}^{b}(\mathcal{D}_{X}^{op}) $ ）表示导出范畴 $ \mathbf{D}^{b}(\mathfrak{Mod}(\mathcal{D}_{X})) $ （或 $ \mathbf{D}^{b}(\mathfrak{Mod}(\mathcal{D}_{X}^{op})) $ ）的满子范畴，由具有相干上同调的对象组成。

范畴  $\mathfrak{Mod}(\mathcal{D}_X)$  具有足够的单射（根据命题 2.4.3），并且对于每个  $x \in X$  、  $\operatorname{gld}(\mathcal{D}_{X,x}) \leq n$  。 （这很容易从命题 11.2.6 中推导出来。）因此，函子  $R\mathcal{Hom}_{\mathcal{D}_X}(\cdot, \cdot)$  和  $\otimes_{\mathcal{D}_X}^L \cdot$  在  $\mathbf{D}^b(\mathcal{D}_X)^\circ \times \mathbf{D}^b(\mathcal{D}_X)$  和  $\mathbf{D}^b(\mathcal{D}_X^\circ) \times \mathbf{D}^b(\mathcal{D}_X)$  上得到了明确定义，值分别在  $\mathbf{D}^+(C_X)$  和  $\mathbf{D}^b(C_X)$  中。

为了结束本节，我们将简要回顾一下  $\mathcal{D}_{X}$  模的主要操作。

 $ \mathbf{D}_{\mathrm{coh}}^{b}(\mathcal{D}_{X}) $  上有一个有趣的对合（因此也在  $ \mathbf{D}_{\mathrm{coh}}^{b}(\mathcal{D}_{X}^{op}) $  上），它类似于层理论中的对偶函子（参见 III §4）。一套：

\[
\mathcal{K}_{X}=\underset{\mathcal{O}_{X}}{\mathcal{D}_{X}}\otimes\Omega_{X}^{\otimes-1}[\operatorname{d i m}_{\mathbb{C}}X]~.
\]

请注意， $ \mathcal{H}_X $  自然地具有左 $ \mathcal{D}_X $  模的两个结构。如果  $ \mathcal{M} \in \mathbf{D}_{\text{coh}}^b(\mathcal{D}_X) $  定义：

\[
\begin{array}{r}{\mathcal{M}^{*}=\mathrm{R H o m}_{\mathcal{D}_{X}}(\mathcal{M},\mathcal{H}_{X})~.}\end{array}
\]

这显然是  $ \mathbf{D}^b(\mathfrak{Mod}(\mathcal{D}_X)) $  的一个对象，这要归功于  $ \mathcal{K}_X $  的双模结构（参见练习 II.23），并且可以立即检查  $ \mathcal{M}^* \in \mathrm{Ob}(\mathbf{D}_{\mathrm{coh}}^b(\mathcal{D}_X)) $  。此外：

\[
{\mathcal{M}}^{**}\simeq{\mathcal{M}}.
\]

现在，让 Y 成为另一个流形。

如果 $\mathcal{M}$ （分别为 $\mathcal{N}$ ）是左 $\mathcal{D}_{X}$ 模（分别为 $\mathcal{D}_{Y}$ 模），则通过设置定义左 $\mathcal{D}_{X \times Y}$ 模：

\[
\mathcal{M}_{\mathbf{\bar{~}u}}\mathcal{N}=\mathcal{D}_{X\times Y}\mathop{{\bigotimes}}_{\mathcal{D}_{X}\boxtimes\mathcal{D}_{Y}}(\mathcal{M}\mathbf{\bar{~}u}\mathcal{N})\;.
\]

很容易验证，如果 $\mathcal{M}$ 和 $\mathcal{N}$ 是一致的，那么 $\mathcal{M} \boxtimes \mathcal{N}$ 也是一致的，

\[
\operatorname{char}(\mathcal{M}\boxtimes\mathcal{N})=\operatorname{char}(\mathcal{M})\times\operatorname{char}(\mathcal{N})
\]

最后让 $f: Y \to X$ 成为全纯映射。

赋予束 $ \mathcal{O}_{Y} \otimes_{f^{-1}\mathcal{O}_{X}} f^{-1}\mathcal{D}_{X} $ 以右 $ f^{-1}\mathcal{D}_{X} $ 模的自然结构。还赋予它左 $ \mathcal{D}_{Y} $ 模的结构，如下所示。

令  $\Theta_X$  （分别为  $\Theta_Y$  ）表示  $X$  （分别为  $Y$  ）上的全纯向量场束。如果  $v$  是  $\Theta_Y$  的一个部分，则应用于  $v$  的微分  $f': TY \to Y \times_X TX$  定义了  $\mathcal{O}_Y \otimes_{f^{-1}\mathcal{O}_X} f^{-1}\Theta_X$  的一个部分  $f'(v)$ ，该部分在  $Y$  上本地可以写为有限和  $\sum_j a_j \otimes w_j$  。考虑  $\mathcal{O}_Y \otimes_{f^{-1}\mathcal{O}_X} f^{-1}\mathcal{D}_X$  类型  $a \otimes u$  的部分。一套：

\[
v(a\otimes u)=v(a)\otimes u+\sum_{j}a a_{j}\otimes w_{j}\circ u~.
\]

如果在  $ X $  上选择局部坐标系  $ (x_1, \ldots, x_n) $  ，在  $ Y $  上选择 $ (y_1, \ldots, y_m) $  以便  $ f = (f_1, \ldots, f_n) $  ，则：

\[
D_{y_{k}}(a\otimes u)=\frac{\partial a}{\partial y_{k}}\otimes u+\sum_{j=1}^{n}a\frac{\partial f_{j}}{\partial x_{k}}\otimes D_{x_{j}}\circ u~.
\]

 $\boldsymbol{\theta}_{Y}$  对  $\mathcal{O}_{Y} \otimes_{f^{-1}\mathcal{O}_{X}} f^{-1}\mathcal{D}_{X}$  的左操作自然地扩展到  $\mathcal{D}_{Y}$  的左操作。

\statementlabel{定义 11.2.8}赋予 $(\mathcal{D}_{Y},f^{-1}\mathcal{D}_{X})$  -bimodule 结构的层 $\mathcal{O}_{Y}\otimes_{f^{-1}\mathcal{O}_{X}}f^{-1}\mathcal{D}_{X}$  由 $\mathcal{D}_{Y\to X}$  表示。其规范部分  $1\otimes1$  由  $1_{Y\to X}$  表示。

双模  $ \mathcal{D}_{Y\to X} $  定义了两个函子：

\[
\begin{aligned}{}&{{}\mathcal{D}_{Y\to X}\mathop{\otimes}_{f^{-1}\mathcal{D}_{X}}^{L}f^{-1}(\cdot):\mathsf{D}^{b}(\mathcal{D}_{X})\to\mathsf{D}^{b}(\mathcal{D}_{Y})}\\ {}&{{}\cdot\mathop{\otimes}_{\mathcal{D}_{Y}}^{L}\mathcal{D}_{Y\to X}\colon\mathsf{D}^{b}(\mathcal{D}_{Y}^{{o p}})\to\mathsf{D}^{b}(f^{-1}\mathcal{D}_{X}^{{o p}}).}\\ \end{aligned}
\]

如果  $g: Z \to Y$  是另一个全纯映射，则存在规范同构：

\[
\mathcal{D}_{\mathbf{Z}\to\mathbf{Y}}\overset{L}{\otimes}_{\mathbf{g}^{-1}\mathcal{D}_{\mathbf{Y}}}g^{-1}\mathcal{D}_{\mathbf{Y}\to\mathbf{X}}\simeq\mathcal{D}_{\mathbf{Z}\to\mathbf{X}}.
\]

示例 11.2.9。 (i) 设  $ x = (x', x'') $  为  $ X $  上的坐标系，并设  $ Y = \{x \in X; x' = 0\} $  。然后作为  $ \mathcal{D}_X^{op} $  模，  $ \mathcal{D}_{Y \to X} \simeq \mathcal{D}_X/(x') \cdot \mathcal{D}_X $  其中  $ (x') \cdot \mathcal{D}_X $  是由坐标  $ (x') $  生成的右理想。该层的一部分可以写为 $ P(x", D_{x'}', D_{x''}) $ （即： $ X $ 上的微分运算符不依赖于 $ x' $ ）。

(II) 令 $ x = (x', x'') $  为Y 上的坐标系，并令 $ f: Y \to X $  为投影 $ (x', x'') \mapsto (x'' $  。然后  $ \mathcal{D}_{Y \to X} \simeq \mathcal{D}_Y / (D_x') \cdot \mathcal{D}_Y $  和该层的一部分可以写为  $ P(x, x'', D_x''), $  （即： Y 上的微分运算符不依赖于  $ D_x' $  ）。

\statementlabel{定义 11.2.10}(i) 令 $\mathcal{M}$  为左 $\mathcal{D}_X$  模（或更一般地， $\mathcal{M} \in \mathrm{Ob}(\mathcal{D}^b(\mathcal{D}_X))$  ）。通过  $f$  定义  $\mathcal{M}$  的逆像，表示为  $\underline{f}^{-1}\mathcal{M}$  ，通过：

\[
\underline{{f}}^{-1}\mathcal{M}=\mathcal{D}_{Y\to X}\bigotimes_{f^{-1}\mathcal{D}_{X}}^{L}f^{-1}\mathcal{M}\in\mathsf{O b}(\mathbb{D}^{b}(\mathcal{D}_{Y}))~.
\]

(ii) 令 $\mathcal{N}$  为右 $\mathcal{D}_{\mathcal{Y}}$  模（或更一般地， $\mathcal{N} \in \mathrm{Ob}(\mathbb{D}^{b}(\mathcal{D}_{\mathcal{Y}}^{op}))$  ）。通过  $f$  定义  $\mathcal{N}$  的直像（或正确的直像），通过以下方式表示  $\underline{f}_{\star} \mathcal{N}$  （或  $\underline{f}_{\star} \mathcal{N}$  ）：

\[
\underline{{\mathcal{f}}}_{\ast}\mathcal{N}=R\mathcal{f}_{\ast}\left(\begin{matrix}{\mathcal{N}_{\bigotimes_{\mathcal{D}_{\mathtt{Y}}}^{L}\mathcal{D}_{\mathtt{Y}\to X}}^{L}}\\ \end{matrix}\right)\in\mathsf{O b}(\mathbb{D}^{b}(\mathcal{D}_{X}^{o p}))
\]

（分别）

\[
\underline{{f}}_{\mathrm{!}}\mathcal{N}=R f_{\mathrm{!}}\left(\mathcal{N}\mathop{\bigotimes}_{\mathcal{D}_{Y}}^{L}\mathcal{D}_{Y\to X}\right)\in\operatorname{O b}(\mathbb{D}^{b}(\mathcal{D}_{X}^{o p})))~.
\]

让我们在没有证明的情况下回顾一下 Kashiwara [5] 提出的关于逆像的主要结果。 （关于直像，参见 Kashiwara [4]、Houzel-Schapira [1]、Schneiders [1] 和 Boutet de Monvel-Malgrange [1]。）

像往常一样，我们将分别用 $ f_{\pi} $ 和 $ f' $ 表示从 $ Y \times_X T^*X $ 到 $ T^*X $ 和 $ T^*Y $ 的自然映射。

\statementlabel{定义 11.2.11}令 M 为左相干  $ \mathcal{D}_X $  模。如果满足以下条件，则称 f 对于 M 而言是非特征值：

\[
T_{Y}^{\ast}X\cap f_{\pi}^{-1}(\mathbf{c h a r}(\mathcal{M}))\subset Y\times_{X}T_{X}^{\ast}X\enspace.
\]

\statementlabel{命题 11.2.12}假设  $f$  对于连贯的  $\mathcal{D}_X$  -模  $\mathcal{M}$  来说是非特征的。然后：

(i)  $ H^{j}(f^{-1}\mathcal{M}) = 0 $  为  $ j \neq 0 $  ，

(ii)  $ H^{0}(\overline{f}^{-1}\mathcal{M}) $  （简单地表示为  $ f^{-1}\mathcal{M} $  ）是一个连贯的  $ \mathcal{D}_{Y} $  模，

（三） $ \mathrm{char}(\underline{f}^{-1}\mathcal{M}) = \mathrm{t}f'f_{\pi}^{-1}(\mathrm{char}(\mathcal{M})). $ 

\subsection{微分算子模的全纯解}
让我们首先回顾一下经典的柯西-科瓦勒夫斯基定理（无需证明），其由 Leray 提出的改进版本 [3]。 （参见例如 Schapira [2]）。

令 $ x = (x_1, \ldots, x_n) $  为 $ \mathbb{C}^n $  上的全纯坐标， $ (x; \xi) $  为 $ T^*\mathbb{C}^n $  上的关联坐标。我们赋予 $ \mathbb{C}^n $ 其通常的埃尔米特结构( $ \langle x, x' \rangle = \sum_j x_j \overline{x}_j' $ )并用 $ B(x_0, \rho) $ 表示以 $ x_0 $ 为中心、以 $ \rho $ 为半径的开球。我们还使用符号  $ x = (x_1, x',) $  ，并用  $ B'(x_0, \rho) $  表示  $ B(x_0, \rho) $  与超平面  $ \{x \in \mathbb{C}^n; x_1 = x_{o,1}\} $  的交集，其中  $ x_0 = (x_{o,1}, \ldots, x_{o,n}) $  。

令  $X$  为  $\mathbb{C}^n$  的开子集，并令  $P$  为  $X$  上定义的阶  $m$  的全纯微分算子。那么 $P$ 写成：

\[
P=\sum_{|\alpha|\leqslant m}a_{\alpha}(x)D_{x}^{\alpha}.
\]

有人提出假设：

\[
a_{(m,0,\ldots,0)}\equiv1.
\]

让  $ x_{0} \in X $  ，并设置：

\[
Y_{x_{0}}=\left\{x\in X;x_{1}=x_{0,1}\right\}
\]

对于  $f$  的  $\mathcal{O}_X$  的一部分，定义  $f$  上  $Y_{x_o}$  （ $(\mathcal{O}_{Y_{x_o}}^{m})$  的一部分）的  $m$  第一迹线，方法如下：

\[
\gamma_{x_{0}}(f)=(f|_{\Upsilon_{x_{0}}},D_{x_{1}}f|_{\Upsilon_{x_{0}}},\ldots,D_{x_{1}}^{m-1}f|_{\Upsilon_{x_{0}}}).
\]

考虑柯西问题：

\[
\left\{\begin{aligned}{P f}&{{}=g}\\ {\gamma_{\mathbf{x}\circ}(f)}&{{}=(h)}\\ \end{aligned}\right..
\]

\statementlabel{定理 11.3.1}让 $x \in X$  。存在  $r > 0$  、  $\rho_o > 0$  、  $\delta > 0$  ，使得对于任何  $\rho > 0$  和  $\rho \leq \rho_o$  、任何  $x_o \in X$  和  $|x - x_o| \leq r$  、任何  $g \in \mathcal{O}_X(B(x_o, \rho))$  、任何  $(h) \in \mathcal{O}_{Y_{x_0}}^m(B'(x_o, \rho))$  ，柯西问题 (11.3.4) 都有唯一解  $f \in \mathcal{O}_X(B(x_o, \delta\rho))$  。

我们首先根据 Zerner [1] 推导出一个有用的扩展结果（参见 Hörmander [3 定理 11.4.7]）。

\statementlabel{命题 11.3.2}令  $\phi$  成为  $X$  上的实数  $C^1$  -函数，使得  $X$  上的  $\sigma(P)(x; \partial \phi(x)) \neq 0$  。令 $\Omega = \{x \in X; \phi(x) < 0\}$  和 $f \in \mathcal{O}_X(\Omega)$  使得 $P f$  在 $x_\circ \in \partial \Omega$  的邻域上全纯扩展。然后  $f$  在  $x_\circ$  的邻域中全纯延伸。

\statementlabel{证明}如果 $\mathrm{d}\phi(x_{\circ})=0$  ， $P$  在 $x_{\circ}$  的邻域内可逆，结论如下。假设  $\mathrm{d}\phi(x_{\circ})\neq0$  并选择  $x_{\circ}$  附近的局部坐标系  $(x_{1},\ldots,x_{n})$  ，例如  $x_{\circ}=0$  、  $\phi(x)=\mathrm{Re}x_{1}-\psi(\mathrm{Im}x_{1},x^{\prime})$  和  $\mathrm{d}\psi(0)=0$  。我们可以假设 $P$ 满足(10.3.1)和(10.3.2)。

设置  $ x_{\varepsilon}=(-\varepsilon,0,\ldots,0) $  ，并保留定理 11.3.1 的符号。考虑柯西问题：

\[
\left\{\begin{aligned}{}&{{}P f_{\varepsilon}=P f}\\ {}&{{}\gamma_{\mathbf{x}_{t}}(f_{\varepsilon})=\gamma_{\mathbf{x}_{t}}(f)}\\ \end{aligned}\right..
\]

那么  $f_{\varepsilon}=f$  和  $f_{\varepsilon}$  在球  $B(x_{\varepsilon},\delta R)$  中是全纯的，其中  $\delta$  可以独立于  $\varepsilon$  选择，并且  $R$  满足。

\[
-\varepsilon<\psi(0,x^{\prime})\qquad\mathrm{f o r}\qquad|x^{\prime}|<R~.
\]

由于  $ \mathrm{d}\psi(0)=0 $  、  $ \psi(0,x')=o(|x'|) $  和  $ B(x_{\varepsilon},\delta R) $  对于足够小的  $ \varepsilon $  来说是 0 的邻域。   $ \square $ 

\statementlabel{定理 11.3.3}设 X 是一个复流形，M 是一个连贯的  $ \mathcal{D}_{X} $  模。然后：

\[
\mathrm{S S}(R\mathcal{H o m}_{\mathcal{D}_{X}}(\mathcal{M},\mathcal{O}_{X}))=\operatorname{c h a r}(\mathcal{M})\enspace.
\]

\statementlabel{证明}我们只在这里证明包含 $ \cdot \subset \cdot $ 。其他包含内容见备注11.4.5。

(a) 首先假设  $ \mathcal{M} = \mathcal{D}_X / \mathcal{D}_X P $  。让  $ p \in T^*X $  和  $ \sigma(P)(p) \neq 0 $  。如果 $ p \in T_X^*X $  ，P 是p 邻域上的可逆函数，而 $ R \mathcal{H}om_{\mathcal{G}_X}(\mathcal{M}, \mathcal{O}_X) = 0 $  是该邻域上的可逆函数。假设  $ p \notin T_X^*X $  并选择一个局部坐标系，例如  $ p = (x_o; \xi_o) $  和  $ \xi_o = (1, 0, \ldots, 0) $  。

我们将应用命题 5.1.1。设定：

\[
\begin{array}{r}{H_{\varepsilon}=\left\{x\in\mathbb{C}^{n};\mathbb{R e}\langle x-x_{\circ},\xi_{\circ}\rangle\geqslant-\varepsilon\right\},}\end{array}
\]

\[
L_{\varepsilon}=\left\{x\in\mathbb{C}^{n};\mathbf{Re}\langle x-x_{\circ},\xi_{\circ}\rangle=-\varepsilon\right\}
\]

\[
\gamma_{\delta}=\left\{x\in\mathbb{C}^{n};\operatorname{I m}x_{1}=0,\mathbf{R e}x_{1}\geqslant\delta|x^{\prime}|\right\}\;.
\]

我们选择  $ 0 < R \ll 1 $  、  $ \delta \gg 0 $  这样：

\[
\sigma(P)(x;\xi)\neq0\qquad\mathrm{f o r}\qquad|x-x_{\circ}|\leqslant R~,\qquad\xi\in\gamma_{\delta}^{\circ}\backslash\{0\}~.
\]

然后我们选择 $ 0 < \varepsilon \ll 1 $ ，这样 $ (x + \gamma_\delta) \cap L_\varepsilon \subset B(x_o, R) $  就成为 $ |x - x_o| \leq \frac{R}{2} $  。

对象 $ R\mathcal{H}om_{\mathcal{D}_X}(\mathcal{M}, \mathcal{O}_X) $ 可以用复数表示：

\[
0\longrightarrow\mathcal{O}_{X}\mathop{\longrightarrow}_{P}\mathcal{O}_{X}\longrightarrow0~.
\]

此外，如果  $K$  是凸且紧的，则众所周知的结果断言  $H^i(K;\mathcal{O}_X)=0$  对于  $i>0$  ，即  $R\Gamma(K;\mathcal{O}_X)\simeq\mathcal{O}_X(K)$  （参见 Hörmander [1]）。为了证明 $p\notin\mathrm{SS}(R\mathcal{H}o m_{\mathcal{D}_X}(\mathcal{M},\mathcal{O}_X))$ ，证明这两个复合体就足够了：

\[
0\longrightarrow\mathcal{O}_{X}((x+\gamma_{\delta})\cap H_{\varepsilon})\xrightarrow[p]{}\mathcal{O}_{X}((x+\gamma_{\delta})\cap H_{\varepsilon})\longrightarrow0
\]

and

\[
0\longrightarrow\mathcal{O}_{X}((x+\gamma_{\delta})\cap L_{\varepsilon})\xrightarrow[p]{}\mathcal{O}_{X}((x+\gamma_{\delta})\cap L_{\varepsilon})\longrightarrow0
\]

与  $ |x - x_{\circ}| \leq \frac{R}{2} $  准同构。

由于  $P$  在空间  $\mathcal{O}_X((x + \gamma_\delta) \cap L_\varepsilon)$  上是满射，根据定理 11.3.1，还需要证明：

(11.3.6)

\[
\begin{array}{l}{1.3.6)}\\ {\left\{\begin{aligned}{}&{{}f\in\mathcal{O}_{X}((x+\gamma_{\delta})\cap L_{\varepsilon})~,}\\ {}&{{}{P f}\in\operatorname{I m}(\mathcal{O}_{X}((x+\gamma_{\delta})\cap H_{\varepsilon})\to\mathcal{O}_{X}((x+\gamma_{\delta})\cap L_{\varepsilon}))\Rightarrow f\in\mathcal{O}_{X}((x+\gamma_{\delta})\cap H_{\varepsilon})~.}\\ \end{aligned}\right.}\\ \end{array}
\]

但是，考虑到（11.3.5），通过与命题 5.1.1 的证明类似的论证，可以很容易地从命题 11.3.2 得出（11.3.6）。

(b) 接下来，我们处理一般情况。让 $p \in T^*X$ ， $p \notin \text{char}(\mathcal{M})$ 。我们选择  $\pi(p)$  邻域内  $\mathcal{M}$  的生成器  $(u_1, \ldots, u_N)$  系统。对于每个  $j$  ，都存在一个运算符  $P_j$  ，使得  $P_j u_j = 0$  和  $\sigma(P_j)(p) \neq 0$  。设置  $\mathcal{L} = \bigoplus_{j=1}^N \mathcal{D}_X / \mathcal{D}_X P_j$  ，并通过设置  $\psi(1 \bmod \mathcal{D}_X P_j) = u_j$  定义  $\mathcal{D}_X$  -线性态射  $\psi: \mathcal{L} \to \mathcal{M}$  。让 $\mathcal{K} = \text{Ker} \psi$  。我们得到了连贯的  $\mathcal{D}_X$  模的精确序列：

\[
0\to\mathcal{H}\to\mathcal{L}\to\mathcal{M}\to0.
\]

令 $U$  成为 $p$  的邻域，使得 $U \cap \mathrm{char}(\mathcal{M}) = \emptyset$  、 $U \cap \mathrm{char}(\mathcal{L}) = \emptyset$  、（因此 $U \cap \mathrm{char}(\mathcal{H}) = \emptyset$  ）。令  $\phi$  成为  $X$  和  $x_\circ \in X$  上的实  $C^1$  函数，使得  $\phi(x_\circ) = 0$  、  $\mathrm{d}\phi(x_\circ) \in U$  。

简称为  $ S_{\phi}(\cdot) $  函子  $ (R\Gamma_{\{\phi\geq0\}}R\mathcal{H}om_{\mathcal{D}_X}(\cdot,\mathcal{O}_{\tilde{X}}))_{x\circ} $  。通过证明的第一部分， $ S_{\phi}(\mathcal{L})=0 $  。将  $ S_{\phi}(\cdot) $  应用于精确序列 (11.3.7) 并取上同调，我们得到长精确序列：

\[
0\to H^{0}(S_{\phi}(\mathcal{M}))\to H^{0}(S_{\phi}(\mathcal{L}))\to H^{0}(S_{\phi}(\mathcal{H}))\to H^{1}(S_{\phi}(\mathcal{M}))\to\cdots
\]

然后 $ H^0(S_\phi(\mathcal{M})) = 0 $ 。由于  $ \mathcal{H} $  满足与  $ \mathcal{M} $  、 $ H^0(S_\phi(\mathcal{H})) = 0 $  相同的假设。通过归纳论证，我们得到所有  $ j $  的  $ H^j(S_\phi(\mathcal{M})) = 0 $  ，这就完成了证明。   $ \square $ 

\statementlabel{备注 11.3.4}上面证明的包含在应用中将是最有用的。

现在我们将研究 $\mathcal{D}$ 模的全纯解的一些操作。令 $\mathcal{M}$ （分别为 $\mathcal{N}$ ）为左 $\mathcal{D}_X$ 模（分别为 $\mathcal{D}_Y$ 模）。存在一个规范态射：

\[
(11.3.8)\quad R\mathcal{H}o m_{\mathfrak{D}_{X}}(\mathcal{M},\mathcal{O}_{X})\boxtimes R\mathcal{H}o m_{\mathfrak{D}_{Y}}(\mathcal{N},\mathcal{O}_{Y})\to R\mathcal{H}o m_{\mathfrak{D}_{X\times Y}}(\mathcal{M}\boxtimes\mathcal{N},\mathcal{O}_{X\times Y})~.
\]

令 $f: Y \to X$  为全纯映射。如果  $\mathcal{M}$  和  $\mathcal{L}$  是两个左  $\mathcal{D}_X$  -模，则  $\mathbf{D}^b(\mathbb{C}_Y)$  中存在自然态射：

\[
\begin{array}{r}{f^{-1}\mathrm{R}\mathcal{H}\mathrm{o m}_{\mathfrak{D}_{X}}(\mathcal{M},\mathcal{L})\to\mathrm{R}\mathcal{H}\mathrm{o m}_{\mathfrak{D}_{Y}}(\underline{{f}}^{-1}\mathcal{M},\underline{{f}}^{-1}\mathcal{L})\enspace.}\end{array}
\]

该态射的获得如下（参见练习 II.23）。首先用单射 $\mathcal{D}_X$ 模的复杂 $\mathcal{L}^{\prime}$ 替换 $\mathcal{L}$ 。接下来通过  $(\mathcal{D}_Y, f^{-1}\mathcal{D}_X)$  -bimodules 选择  $\mathcal{D}_{Y\to X}$  的有界分辨率  $\mathcal{P}^{\prime}$  ，平坦于  $f^{-1}\mathcal{D}_X$  。最后选择  $\mathcal{P}^{\prime}\otimes_{f^{-1}\mathcal{D}_X}\mathcal{L}^{\prime}$  的  $\mathcal{D}_Y$  内射分辨率  $\mathcal{I}^{\prime}$  。然后我们就有了态射：

\[
\begin{aligned}{\mathfrak{f}^{-1}\mathsf{R H o m}_{\mathcal{D}_{X}}(\mathcal{M},\mathcal{L})}&{{}\xrightarrow{\sim}\mathfrak{f}^{-1}\mathcal{H o m}_{\mathcal{D}_{X}}(\mathcal{M},\mathcal{L}^{\prime})}\\ {}&{{}\to\mathcal{H o m}_{f^{-1}\mathcal{D}_{X}}(\mathfrak{f}^{-1}\mathcal{M},\mathfrak{f}^{-1}\mathcal{L}^{\prime})}\\ {}&{{}\to\mathcal{H o m}_{\mathcal{D}_{Y}}\bigg(\mathcal{P}^{\prime}\mathop{\otimes}_{\mathfrak{f}^{-1}\mathcal{D}_{X}}\mathfrak{f}^{-1}\mathcal{M},\mathcal{P}^{\prime}\mathop{\otimes}_{\mathfrak{f}^{-1}\mathcal{D}_{X}}\mathfrak{f}^{-1}\mathcal{L}^{\prime}\bigg)}\\ {}&{{}\to\mathcal{H o m}_{\mathcal{D}_{Y}}\bigg(\mathcal{P}^{\prime}\mathop{\otimes}_{\mathfrak{f}^{-1}\mathcal{D}_{X}}\mathfrak{f}^{-1}\mathcal{M},\mathcal{I}^{\prime}\bigg)}\\ {}&{{}\simeq\mathsf{R H o m}_{\mathcal{D}_{Y}}(\underline{{\mathfrak{f}}}^{-1}\mathcal{M},\underline{{\mathfrak{f}}}^{-1}\mathcal{L})\enspace.}\\ \end{aligned}
\]

可以将柯西-科瓦勒夫斯基定理扩展到系统。

\statementlabel{定理 11.3.5}假设 f 对于 M 来说是非特征的。那么自然态射

\[
\begin{array}{r}{f^{-1}\boldsymbol{\mathrm{R}}\mathcal{H o m}_{\mathcal{D}_{X}}(\mathcal{M},\mathcal{O}_{X})\to\boldsymbol{\mathrm{R}}\mathcal{H o m}_{\mathcal{D}_{Y}}(\underline{{f}}^{-1}\mathcal{M},\mathcal{O}_{Y})}\end{array}
\]

是同构。

\statementlabel{证明}考虑一下映射：

\[
Y\underset{g}{\hookrightarrow}Y\times\underset{p}{\xrightarrow{X}}X
\]

其中  $g$  是图形映射  $(y \mapsto (y, f(y)))$  和  $p$  投影。鉴于（11.2.24），分别证明  $p$  和  $g$  的结果就足够了。对于 $p$ ，我们必须证明：

\[
p^{-1}\mathtt{R}\mathcal{H}o m_{\mathcal{D}_{X}}(\mathcal{M},\mathcal{O}_{X})\simeq\mathtt{R}\mathcal{H}o m_{\mathcal{D}_{Y\times X}}(\mathcal{O}_{Y}\boxtimes\mathcal{M},\mathcal{O}_{Y\times X})\enspace.
\]

这可以很容易地简化为 $ \mathcal{M} = \mathcal{D}_X $  的情况，那么这只不过是庞加莱引理。

因此，人们可能从一开始就假设 $f$  是一种沉浸。通过对 codim  $Y$  的归纳论证，我们甚至可以假设  $Y$  是一个超曲面。然后与定理 11.3.3 相同的证明简化为  $\mathcal{M} = \mathcal{D}_X / \mathcal{D}_X P$  的情况。

我们选择局部坐标系 $ (x_1, \ldots, x_n) $ 和 $ Y = \{x_1 = 0\} $ ，并且我们可以假设 $ P $ 满足(11.3.1)，(11.3.2)。让我们计算诱导系统

 $\mathcal{D}_{Y\to X}\otimes_{\mathcal{D}_X}(\mathcal{D}_X/\mathcal{D}_XP) \simeq \mathcal{D}_X/(x_1\mathcal{D}_X + \mathcal{D}_XP)$ 。 Weierstrass 除法定理（对于微分运算符）断言对于任何  $S \in \mathcal{D}_{X,x}$  都存在唯一的对  $(Q, R) \in \mathcal{D}_{X,x}^2$ ，使得

\[
\left\{\begin{aligned}{S}&{{}=Q\circ P+R,}\\ {R}&{{}=\sum_{j=0}^{m-1}R_{j}(x,D_{x^{\prime}})D_{1}^{j}.}\\ \end{aligned}\right.
\]

因此S可以唯一地写为：

\[
S=Q\circ P+x_{1}T+\sum_{j=0}^{m-1}\widetilde{R}_{j}(x^{\prime},D_{x^{\prime}})D_{x_{1}}^{j}.
\]

由此可见，  $ \mathcal{D}_{Y\to X} \otimes_{\mathcal{D}_X} (\mathcal{D}_X / \mathcal{D}_X P) $  是  $ \mathcal{D}_Y $  之上等级为  $ m $  的自由模，由  $ 1_{Y\to X} \otimes u, \ldots, 1_{Y\to X} \otimes D_{x_i}^{m-1} u $  生成，其中  $ u $  是  $ \mathcal{D}_X / \mathcal{D}_X P $  的生成器  $ 1 \bmod \mathcal{D}_X P $  。

那么定理11.3.4翻译为：

\[
\left\{\begin{aligned}{\operatorname{K e r}(\mathcal{O}_{X}\xrightarrow[P]{}\mathcal{O}_{X})|_{Y}\xrightarrow[\gamma_{Y}]{\sim}\mathcal{O}_{Y}^{m},}\\ {\operatorname{C o k e r}(\mathcal{O}_{X}\xrightarrow[P]{}\mathcal{O}_{X})|_{Y}=0,}\\ \end{aligned}\right.
\]

这是定理 10.3.1 的一个特例。 ⑨

\statementlabel{备注 11.3.6}Schneiders [1] 对于直像也存在类似的结果。

我们将通过证明完整  $ \mathcal{D}_x $  模的全纯解复形是反常的来结束本节。

\statementlabel{定理 11.3.7}假设  $\mathcal{M}$  是一个完整的  $\mathcal{D}_X$  模。那么  $R\mathcal{H}om_{\mathcal{D}_X}(\mathcal{M},\mathcal{O}_X)$  属于  $\mathsf{D}_{\mathbb{C}-c}^b(X)$  （即：是  $\mathbb{C}$  -可构造的），而且  $R\mathcal{H}om_{\mathcal{D}_X}(\mathcal{M},\mathcal{O}_X)[\dim_{\mathbb{C}}X]$  是反常的。

\statementlabel{证明}设置 $ F = \mathcal{R}\mathcal{H}om_{\mathcal{D}_X}(\mathcal{M}, \mathcal{O}_X)[\dim_C X] $ 。

(a) 首先我们证明 $F$  是 $\mathbb{C}$  可构造的。

根据定理 11.3.3 和定理 8.5.5， $F$  是弱  $\mathbb{C}$  可构造的。因此，仍然需要证明对于任何  $x_o \in X, j \in \mathbb{Z}$  ，向量空间  $H^j(F)_{x_o}$  是有限维的。

选择  $x_\circ$  邻域内的局部图，并用  $B(x_\circ,\varepsilon)$  表示以  $x_\circ$  为中心、以  $\varepsilon$  为半径的开球。设置 $F_\varepsilon = R\Gamma(B(x_\circ,\varepsilon); F)$ 。根据引理 8.4.7，自然态射  $F_\varepsilon \to F_\varepsilon'$  是  $0 < \varepsilon' \leqslant \varepsilon \ll 1$  的准同构。为了计算  $F_\varepsilon$ ，我们选择（使用命题 11.2.6） $\mathcal{M}$  的有限自由表示：

\[
0\longrightarrow\mathcal{D}_{X}^{N_{n}}\longrightarrow\cdots\longrightarrow\mathcal{D}_{X}^{N_{0}}\longrightarrow\mathcal{M}\longrightarrow0~.
\]

那么  $F_{\varepsilon}$  由复数表示：

\[
0\longrightarrow\mathcal{O}_{X}^{N_{0}}(B(x_{\circ},\varepsilon))\xrightarrow[P_{0}]{}\cdots\longrightarrow\mathcal{O}_{X}^{N_{n}}(B(x_{\circ},\varepsilon))\longrightarrow0
\]

（从  $H^j(B(x_\circ,\varepsilon);\mathcal{O}_X)=0$  到  $j\neq0$  ）。空间  $\mathcal{O}_X(B(x_\circ,\varepsilon))$  自然地赋予了 Fréchet 空间的结构，并且如果  $0<\varepsilon'<\varepsilon$  则限制态射  $\mathcal{O}_X(B(x_\circ,\varepsilon))\to\mathcal{O}_X(B(x_\circ,\varepsilon'))$  是紧的。因此，结果来自下一个引理，其证明是施瓦茨定理的直接应用[1]。

\statementlabel{引理 11.3.8}令 $F^-$  和 $G^-$  为Fréchet 空间和线性连续映射的两个有界复形， $u:F^-\to G^-$  为连续线性态射。假设对于每个  $j$  ，线性映射  $u^j:F^j\to G^j$  是紧凑的，并假设  $u$  是准同构。那么对于每个  $j\in\mathbb{Z}$  ，空间  $H^j(F^-)$  都是有限维的。

(b) 最后我们将证明 $F$  是反常的。让 $n = \dim_{\mathbb{C}} X$  。令  $S$  为复余维  $d$  的复子流形。然后  $H_{\mathcal{S}}^{j}(F) = 0$  对应  $j < d - n$  ，因为  $H_{\mathcal{S}}^{j}(\mathcal{O}_X) = 0$  对应  $j < d$  （参见 (2.9.14)）。因此  $F$  属于  $^p\mathbf{D}_{\mathbb{C}-c}^{\geq 0}(\mathbb{C}_X)$  。

现在固定 $ j \in \mathbb{N} $ ，并设置 $ Y = \operatorname{supp}(H^j(F)) $ 。

由于 $F$  是 $\mathbb{C}$  可构造的，所以 $Y$  是一个封闭的复解析集。我们将证明 $\dim_{\mathbb{C}}Y \leq -j$ 。

由于在 Y 的泛点上证明这种不等式就足够了，我们可以根据命题 8.2.10 假设  $ \operatorname{char}(\mathcal{M}) \cap \pi^{-1}(Y) \subset T_Y^* X $  。设 x \textbackslash{}in Y 并选择在 x 处与 Y 相交的子流形 Z，其中 \textbackslash{}operatorname\{dim\}\_C Y + \textbackslash{}operatorname\{dim\}\_C Z = n。嵌入 f: Z \textbackslash{}hookrightarrow X 对于 \textbackslash{}mathcal\{M\} 来说是非特征的。因此，根据定理 11.3.5，F|\_\{Z\} \textbackslash{}simeq R\textbackslash{}mathcal\{H\}om\_\{\textbackslash{}mathcal\{D\}\_2\}(f\textasciicircum{}\{-1\}\textbackslash{}mathcal\{M\}, \textbackslash{}mathcal\{O\}\_Z)[n]。

然后  $ \underline{f}^{-1}\mathcal{M} $  根据命题 11.2.6 承认长度  $ \dim_{\mathbb{C}} Z $  的有限自由解析，因此  $ H^k(F|_{z})_x = H^{k+n}(R\mathcal{H}om_{\mathcal{D}_Z}(\underline{f}^{-1}\mathcal{M},\mathcal{O}_Z)) = 0 $  对于  $ k + n > \dim_{\mathbb{C}} Z $  。从  $ H^j(F|_{z})_x \neq 0 $  开始，我们获得了  $ j \leq \dim_{\mathbb{C}} Z - n = -\dim_{\mathbb{C}} Y $  ，这是期望的结果。   $ \square $ 

\subsection{结构层的微局部研究}
令 X 为复维 n 的复流形，S 为复余维 d 的复子流形。

\statementlabel{命题 11.4.1}复数  $ \mu_S(\mathcal{O}_X) $  集中在度  $ d $  上。

\statementlabel{证明}$ j < d $  的  $ H^j(\mu_S(\mathcal{O}_X)) $  消失源自断言 (2.9.14)。  $ j > d $  的消失是一个简单的练习（参见 Sato-Kawai-Kashiwara [1]）。   $ \square $ 

令 Y 为另一个复流形。我们引入记号：

\[
\begin{array}{r}{\Omega_{X\times Y/Y}=\mathcal{O}_{X\times Y}\otimes_{q_{1}^{-1}\mathcal{O}_{X}}q_{1}^{-1}\Omega_{X}.}\end{array}
\]

\statementlabel{定义 11.4.2}(i) 一套：

\[
\mathcal{C}_{S|X}^{\bf R}=H^{d}(\mu_{S}(\mathcal{O}_{X}))\enspace.
\]

(ii) 令 $f: Y \to X$  为复流形的态射。一套：

\[
\mathcal{\tilde{E}}_{Y\to X}^{\mathtt{R}}=H^{n}(\mu_{\varDelta_{f}}(\Omega_{X\times Y/Y}))~,
\]

\[
\mathcal{E}_{X\leftarrow Y}^{\bf R}=H^{n}(\mu_{\varDelta_{f}}(\Omega_{X\times Y/X})^{a})\quad,
\]

其中“a”是对映图。当 $ f = \mathrm{id}_{X} $ 时，写 $ \mathcal{E}_{X}^{\mathsf{R}} $ 而不是 $ \mathcal{E}_{X\to X}^{\mathsf{R}} $ 。

 $ X \times Y $  中的  $ \mathcal{A}_f $  基本类给出了  $ H_{\mathcal{A}_f}^n(\Omega_{X \times Y/Y}) $  的一个部分，并定义了  $ \mathcal{E}_{Y \to X}^R $  的全局部分。用 $ 1_{Y \to X} $ 表示（事实上这个定义与定义11.2.8一致）。当 $ f = \mathrm{id}_X $ 时，写 $ 1_X $ 而不是 $ 1_{X \to X} $ 。层 $ \mathcal{E}_{Y \to X}^R $  和 $ \mathcal{E}_{X \to Y}^R $  由 $ T_{\mathcal{A}_f}^*(X \times Y) \simeq Y \times_X T^*X $  支持。

我们将研究层的组成 $ \mathcal{E}_{Y\rightarrow X}^{\mathbf{R}} $ 。

如果 $X$ 、 $Y$ 、 $Z$ 是三个流形，我们用 $p_{ij}$ 表示从 $T^{*}X \times T^{*}Y \times T^{*}Z$ 到 $(i,j)$ 分量的投影，用 $p_{ij}^{a}$ 表示 $p_{ij}$ 的组合和 $j$ 分量上的对映图。我们类似地为四个流形  $W, X, Y, Z$  定义了在乘积  $T^{*}W \times T^{*}X \times T^{*}Y \times T^{*}Z$  上定义的投影  $p_{ij}^{a}$  。

\statementlabel{引理 11.4.3}(i) 对于  $ K_1 \in \text{Ob}(\mathbf{D}^b(X \times Y)) $  和  $ K_2 \in \text{Ob}(\mathbf{D}^b(Y \times Z)) $  ，存在规范态射：

\[
\begin{aligned}{R p_{13}^{a}}&{{}(p_{12}^{a-1}\mu\mathcal{h o m}(K_{1},\varOmega_{X\times Y/Y})\otimes p_{23}^{a-1}\mu\mathcal{h o m}(K_{2},\varOmega_{Y\times Z/Z}))}\\ {}&{{}\to\mu\mathcal{h o m}(K_{1}\circ K_{2},\varOmega_{X\times Z/Z})[-\dim_{\mathbb{C}}Y].}\\ \end{aligned}
\]

(ii) 对于  $K_1 \in \mathrm{Ob}(\mathbb{D}^b(W \times X)), K_2 \in \mathrm{Ob}(\mathbb{D}^b(X \times Y))$  和  $K_3 \in \mathrm{Ob}(\mathbb{D}^b(Y \times Z))$  ，下图可交换：

\begin{center}
\includegraphics[width=0.80\textwidth]{流行上的层_Chn-assets/img_in_image_box_92_1094_1047_1584.png}
\end{center}

这里  $F_1 = Rp_{13}^a!(p_{12}^{a-1}\mu\mathrm{hom}(K_1, \Omega_{W \times X/Y})) \otimes p_{23}^{a-1}\mu\mathrm{hom}(K_2, \Omega_{X \times Y/Y})) \in \mathrm{Ob}(\mathbb{D}^b(T^*X \times T^*Y)))$  和  $F_2 = Rp_{13}^a!(p_{12}^{a-1}\mu\mathrm{hom}(K_2, \Omega_{X \times Y/Y})) \otimes p_{23}^{a-1}\mu\mathrm{hom}(K_3, \Omega_{Y \times Z/Z})) \in \mathrm{Ob}(\mathbb{D}^b(T^*Y \times T^*Z)))$  ，态射  $\alpha$  和  $\beta$  是诱导的

由 (i) 中给出的态射：

\[
F_{1}\to\mu{\mathcal{h}}o m(K_{1}\circ K_{2},\Omega_{W\times Y/Y})[-{\operatorname{d i m}}X]\qquad{~a n d~}
\]

\[
F_{2}\to\mu{\dot{o}}m(K_{2}\circ K_{3},\Omega_{X\times Z/Z})[-{\tt-d i m}Y],\quad{\mathrm{r e s p e c t i v e l y.}}
\]

\statementlabel{证明}我们可以应用命题4.4.11来获得态射：

\[
\begin{aligned}{}&{{}{R p}_{13,!}^{a}(p_{12}^{a-1}\mu{\mathcal{h}}o m(K_{1},\varOmega_{X\times Y/Y})\otimes p_{23}^{a-1}\mu{\mathcal{h}}o m(K_{2},\varOmega_{Y\times Z/Z}))}\\ {}&{{}\quad\to\mu{\mathcal{h}}o m(K_{1}\circ K_{2},\varOmega_{X\times Y/Y}\circ\varOmega_{Y\times Z/Z})~.}\\ \end{aligned}
\]

与态射结合（参见定理 11.1.4）：

\[
\begin{aligned}{\Omega_{X\times Y/Y}\circ\Omega_{Y\times Z/Z}}&{{}\simeq R p_{13!}(q_{12}^{-1}\Omega_{X\times Y/Y}\otimes q_{23}^{-1}\Omega_{Y\times Z/Z})}\\ {}&{{}\to R q_{13!}(\Omega_{X\times Y\times Z/Z})}\\ {}&{{}\to\Omega_{X\times Z/Z}[-\mathsf{d i m}_{\mathbb{C}}Y],}\\ \end{aligned}
\]

我们得到了结果。

(ii) 证明留给读者。 □

令 $g: Z \to Y$  和 $f: Y \to X$  为复流形的态射。我们保留与命题 4.4.9 相同的符号（特别是参见图 4.4.13，其中  $f'$  是映射  $Z \times_X T^*X \to Z \times_Y T^*Y$  和  $g_\pi$  是映射  $Z \times_X T^*X \to Y \times_X T^*X$  ）。

\statementlabel{命题 11.4.4}(i)  $ Z \times_X T^*X $  上存在自然态射：

\[
\iota_{f^{\prime-1}}\mathcal{E}_{Z\to Y}^{\tt R}\otimes g_{\pi}^{-1}\mathcal{E}_{Y\to X}^{\tt R}\to\mathcal{E}_{Z\to X}^{\tt R}\enspace,
\]

\[
\begin{array}{r}{\mathcal{G}_{\pi}^{-1}\mathcal{E}_{X\leftarrow Y}^{\tt R}\otimes\mathfrak{t f}^{\prime-1}\mathcal{E}_{Y\leftarrow Z}^{\tt R}\to\mathcal{E}_{X\leftarrow Z}^{\tt R}\enspace,}\end{array}
\]

 $1_{Z\to Y}\otimes1_{Y\to X}$  通过第一个态射的图像是  $1_{Z\to X}$  。此外，如果 $h:W\to Z$  是复流形的另一种态射，则 $p_3^{-1}\mathcal{E}_W\to_Z\otimes p_2^{-1}\mathcal{E}_Z\to_Y\otimes p_1^{-1}\mathcal{E}_Y\to_X$  上的态射的复合是结合律，并且与箭头相反的情况类似。这里 $p_1$ 、 $p_2$ 和 $p_3$ 分别是从 $W\times_X T^*X$ 到 $Y\times_X T^*X$ 、 $Z\times_Y T^*Y$ 和 $W\times_Z T^*Z$ 的映射。

(ii) 束 $ \mathcal{E}_X^\mathbb{R} $  天然地具有具有单位 $ 1_X $  的环结构，并且 $ \mathcal{E}_{Y\to X}^\mathbb{R} $ （分别 $ \mathcal{E}_{X\to Y}^\mathbb{R} $  ）是 $ (f^{\prime-1}\mathcal{E}_Y^\mathbb{R}, f_\pi^{-1}\mathcal{E}_X^\mathbb{R}) $  -双模（分别是 $ (f_\pi^{-1}\mathcal{E}_X^\mathbb{R}, f^{\prime-1}\mathcal{E}_Y^\mathbb{R}) $  -双模）。

(iii) 令 $F \in \mathrm{Ob}(\mathbb{D}^b(\mathbb{C}_X))$  并令 $j \in \mathbb{Z}$  。那么 $H^j(\mu_{\mathrm{hom}}(F, \mathcal{O}_X))$ （或 $H^j(\mu_{\mathrm{hom}}(F, \Omega_X))$ ）自然地被赋予左（或右） $\mathcal{E}_X^\mathrm{R}$  -模的结构。

\statementlabel{证明}使用引理 11.4.3 及其符号，我们得到态射：

\[
\begin{aligned}{R p_{13!}^{a}}&{{}(p_{12}^{a-1}\mu\mathcal{h o m}(\mathbb{C}_{\varDelta_{f}},\varOmega_{X\times Y/Y})\times p_{23}^{a-1}\mu\mathcal{h o m}(\mathbb{C}_{\varDelta_{g}},\varOmega_{Y\times Z/Z}))}\\ {}&{{}\to\mu\mathcal{h o m}(\mathbb{C}_{\varDelta_{f}}\circ\mathbb{C}_{\varDelta_{g}},\varOmega_{X\times Z/Z})[-\operatorname{d i m}_{\mathbb{C}}Y].}\\ \end{aligned}
\]

现在 $\mathbb{C}_{\mathcal{A}_{f}}\circ\mathbb{C}_{\mathcal{A}_{g}}$  与 $\mathbb{C}_{\mathcal{A}_{f\circ g}}$  同构，并且它仍然取两边的上同调。组合的结合性留给读者， $1_{Z\to Y}\otimes1_{Y\to X}$  给出 $1_{Z\to X}$  的事实是基本类的经典微积分。第二个态射的定义类似。

(ii) 由 (i) 得出。

(iii) 证明与(i)类似。 □

\statementlabel{备注 11.4.5}$ T^*X $  上的环层  $ \mathcal{E}_X^{\mathbb{R}} $  称为  $ X $  上的微局部算子环。它对 $ T^*X $  的零节 $ X $  的限制表示为 $ \mathcal{D}_X^\infty $ ，称为无限阶微分算子环。最后一个环自然包含第 2 节的（有限阶）微分算子的环  $ \mathcal{D}_X $ 。事实上，环  $ \mathcal{E}_X^{\mathbb{R}} $  包含一个子环，表示为  $ \mathcal{E}_X $  并称为（有限阶）微微分算子环，其对  $ X $  的限制是  $ \mathcal{D}_X $  。我们无意在本书中介绍环 $ \mathcal{E}_X $ ，我们参考了 Sato-Kawai-Kashiwara [1]、Björk [1] 或 Schapira [2] 对其进行研究。我们只提一下，环  $ \mathcal{E}_X $  比  $ \mathcal{E}_X^{\mathbb{R}} $  更容易操作。特别是它是连贯的和诺特式的。

让 $\mathcal{M}$  成为一个连贯的 $\mathcal{D}_X$  模。使用捆  $\mathcal{E}_X$  （以及  $\mathcal{E}_X^\mathbb{R}$  忠实地平坦于  $\mathcal{E}_X$  的事实）可以显示：

\[
\operatorname{char}(\mathcal{M})=\operatorname{supp}\left(\mathcal{E}_{X}^{\mathbb{R}}\bigotimes_{\pi^{-1}\mathcal{D}_{X}}\pi^{-1}\mathcal{M}\right).
\]

从 $ \mathrm{char}(\mathcal{M}) = \mathrm{char}(\mathcal{M}^*) $ （参见（10.2.19））开始，我们得到：

\[
\begin{aligned}{\operatorname{c h a r}(\mathcal{M})}&{{}=\operatorname{s u p p}\biggl(\mathcal{E}_{X\atop\pi^{-1}\mathcal{D}_{X}}^{\mathtt{R}}\otimes\pi^{-1}\mathcal{M}^{*}\biggr)}\\ {}&{{}=\operatorname{s u p p}(R\mathcal{H}o m_{\pi^{-1}\mathcal{D}_{X}}(\pi^{-1}\mathcal{M},\mathcal{E}_{X}^{\mathtt{R}}))}\\ {}&{{}=\operatorname{s u p p}(\mu_{\varDelta}R\mathcal{H}o m_{\mathcal{D}_{X}}(\mathcal{M},\mathcal{O}_{X\times X})).}\\ \end{aligned}
\]

然后就可以轻松完成定理 11.3.3 的证明（参见 Kashiwara-Schapira，[3，定理 10.1.1]）。

现在我们应用定理7.5.11。令  $X, Y, Z$  为三个复流形，  $A_1 \subset T^*(X \times Y)$  和  $A_2 \subset T^*(Y \times Z)$  两个圆锥拉格朗日复流形，  $\tilde{p}_1 = (p_X, p_Y^a) \in A_1$  ，  $\tilde{p}_2 = (p_Y, p_Z^a) \in A_2$  。我们做出假设：

 $ p_2^a|_{\mathcal{A}_1}: \mathcal{A}_1 \to T^*Y $  和 $ p_1|_{\mathcal{A}_2}: \mathcal{A}_2 \to T^*Y $  是横向的。

\statementlabel{命题 11.4.6}令 $ K_1 \in \text{Ob}(\mathbb{D}^b(X \times Y; (p_X, p_Y))) $  和 $ K_2 \in \text{Ob}(\mathbb{D}^b(Y \times Z; (p_Y, p_Z^a))) $  与 $ \text{SS}(K_i) \subset A_i $  位于 $ \tilde{p}_i $  (  $ i = 1, 2 $  ) 的邻域内。假设 (11.4.4) 并且：  $ K_i $  很简单，只需沿着  $ A_i $  (  $ i = 1, 2 $  ) 移动零即可。然后：

(a)  $ K_{1} \circ_{\mu} K_{2} [- \dim_{\mathbb{C}} Y] $  很简单，只需沿着  $ \Lambda_{1} \circ \Lambda_{2} $  移动零，

(b)  $ (p_1^{-1}\mathcal{E}_X^R, p_3^{-1}\mathcal{E}_Z^R) $  -bimodule 在  $ (p_X, p_Z^a) $  的邻域中存在自然态射：

\[
\begin{aligned}{p_{13*}^{a}}&{{}\Biggl(p_{12}^{a-1}H^{0}(\mu\mathcal{h}o m(K_{1},\Omega_{X\times Y/Y}))\mathop{\otimes}_{p_{\overline{{2}}^{-1}}\delta_{\Upsilon}^{\mathbb{R}}}p_{23}^{a-1}H^{0}(\mu\mathcal{h}o m(K_{2},\Omega_{Y\times Z/Z}))\Biggr)}\\ {}&{{}\to H^{0}\Biggl(\mu\mathcal{h}o m\biggl(K_{1}\mathop{\circ}_{\mu}K_{2}[-\operatorname{d i m}_{\mathbb{C}}Y],\Omega_{X\times Z/Z}\biggr)\Biggr),}\\ \end{aligned}
\]

其中  $ p_{ij}^{a} $  是投影  $ p_{ij} $  和第 j 个空间上的对映图的合成。

\statementlabel{证明}(a) 由定理 7.5.11 得出（参见练习 VII.5）。

(b) 根据命题 7.3.1，我们可以假设  $ (K_1, K_2) $  满足 (7.3.3) 和 (7.3.4)，以及  $ K_1 \circ_\mu K_2 = \lim_{\substack{V \\ V}} (K_1)_{X \times V} \circ K_2 $  ，其中  $ V $  范围在  $ \pi(p_Y) $  的开放邻域系统上。

根据引理 11.4.3，我们有态射

\[
\begin{aligned}{R p_{13}^{a}}&{{}!(p_{12}^{a-1}\mu\mathcal{h o m}(K_{1},\varOmega_{X\times Y/Y})\otimes p_{23}^{a-1}\mu\mathcal{h o m}(K_{2},\varOmega_{Y\times Z/Z}))}\\ {}&{{}\to\mu\mathcal{h o m}(K_{1}\circ K_{2}[-{\operatorname{d i m}}_{\mathbb{C}}Y],\varOmega_{X\times Z/Z})}\\ {}&{{}\to\mu\mathcal{h o m}\bigg(K_{1\atop\mu}\circ K_{2}[-{\operatorname{d i m}}_{\mathbb{C}}Y],\varOmega_{X\times Z/Z}\bigg).}\\ \end{aligned}
\]

现在应用相同的引理，即  $X, Y, Y, Z$  和  $K_1, \mathbb{C}_{\Delta_Y}$  和  $K_2$  。那么 $\mu \text{hom}(K_1, \Omega_{X \times Y/Y}) \otimes \mathcal{E}_Y^R \to \mu \text{hom}(K_1, \Omega_{X \times Y/Y})$ 和 $\mathcal{E}_Y^R \otimes \mu \text{hom}(K_2, \Omega_{Y \times Z/Z}) \to \mu \text{hom}(K_2, \Omega_{Y \times Z/Z})$ 得到的从 $R p_{13}^a(p_{12}^{a-1} \mu \text{hom}(K_1, \Omega_{X \times Y/Y}) \otimes \mathcal{E}_Y^R \otimes p_{23}^{a-1}(K_2, \Omega_{Y \times Z/Z}))$ 到 $\mu \text{hom}(K_1 \circ_\mu K_2 [-\dim_C Y], \Omega_{X \times Z/Z})$ 的两个态射 $\alpha$ 和 $\beta$ 是相等的。因此取上同调群，我们得到结果。   $\square$ 

在命题11.4.6的帮助下，我们现在准备对 $ \mathbf{D}^b(X; p) $ 中的 $ \mathcal{O}_X $ 图像进行复杂的接触变换。

令 $Y$  为与 $X$  具有相同维数的复流形，并令 $U_X$  和 $U_Y$  分别为 $T^*X$  和 $T^*Y$  的开子集。令  $A \subset U_X \times U_Y^a$  为闭合复拉格朗日子流形并假设：

\[
p_{1}|_{\mathcal{A}}\colon\varLambda\to U_{X}\quad\mathrm{a n d}\quad p_{2}^{a}|_{\mathcal{A}}\colon\varLambda\to U_{Y}\qquad\mathrm{a r e~i s o m o r p h i s m s}.
\]

用  $\chi$  表示与  $A$  关联的复杂接触变换，即  $\chi = (p_1|_{\mathcal{A}}) \circ (p_2^a|_{\mathcal{A}})^{-1}$  。让 $K \in \mathrm{Ob}(\mathbf{D}^b(X \times Y))$ 满意：

\[
\left\{\begin{aligned}{}&{{}K\mathrm{~i s~c o h o m o l o g i c a l l y~c o n s t r u c t i b l e},}\\ {}&{{}(p_{1}^{-1}(U_{X})\cup p_{2}^{a-1}(U_{Y}))\cap\mathrm{S S}(K)\subset\varLambda}\\ {}&{{}K\mathrm{~i s~s i m p l e~w i t h~s h i f t~0~a l o n g~}\varLambda.}\\ \end{aligned}\right.,
\]

请注意，给定满足（10.4.7）的 $A$ ，本地总是存在满足（10.4.8）的 $K$ ：这遵循与推论7.2.2的证明相同的论点。

根据定理 7.2.1，从  $ \mathbf{D}^b(Y; U_Y) $  到  $ \mathbf{D}^b(X; U_X) $  的函子  $ \Phi_K: G \to R q_{1!}(K \otimes q_2^{-1}G) $  是明确定义的，并且是范畴的等价。让 $ p = (p_X, p_Y^a) \in A $ 并让：

\[
s\in H^{0}(\mu\mathrm{h o m}(K,\Omega_{X\times Y/X}))_{p}.
\]

根据定理 6.1.2，s 定义态射：

\[
K\to\Omega_{X\times Y/X}\qquad\mathrm{i n}\qquad\mathbb{D}^{b}(X\times Y;p)~.
\]

我们得到  $ \mathbf{D}^{b}(X; p_{x}) $  中的态射链：

\[
\begin{aligned}{\varPhi_{K[n]}(\mathcal{O}_{Y})}&{{}=R q_{1!}(K\otimes q_{2}^{-1}\mathcal{O}_{Y})[n]}\\ {}&{{}\to R q_{1!}(\varOmega_{X\times Y/X}\otimes q_{2}^{-1}\mathcal{O}_{Y}[n])}\\ {}&{{}\to R q_{1!}\varOmega_{X\times Y/X}[n]}\\ {}&{{}\to\mathcal{O}_{X}.}\\ \end{aligned}
\]

因此，对于 (11.4.9) 中的 s，我们关联了一个态射：

\[
\alpha(s):\varPhi_{K[n]}(\mathcal{O}_{Y})\to\mathcal{O}_{X}\qquad\mathrm{~i n~}\qquad\mathbf{D}^{b}(X;p_{X})~.
\]

让我们解释一下如何构造这样的态射。

令 $Z$  为相同维度的第三流形， $U_Z$  为 $T^*Z$  的开子集， $A' \subset U_Y \times U_Z^a$  为拉格朗日流形， $K' \in \mathrm{Ob}(\mathbb{D}^b(Y \times Z))$  并假设 $(U_Y, U_Z, A', K')$  满足与(11.4.8)相同的假设。让 $p' = (p_Y, p_Z^a) \in A$  并让 $s' \in H^0(\mu \mathrm{hom}(K', \Omega_{Y \times Z/Y}))_{p'}$  。设置 $A'' = A \circ A', K'' = K \circ K' [n]$ 。那么  $K''$  是一个简单的层，沿着  $A''$  移动零并且

\[
\varPhi_{K[n]}\circ\varPhi_{K^{\prime}[n]}=\varPhi_{K^{\prime\prime}[n]}\;.
\]

我们得到  $ \mathbf{D}^{b}(X; p) $  中的态射：

\[
\varPhi_{K^{\prime\prime}[n]}(\mathcal{O}_{Z})\xrightarrow[\varPhi_{K[n]}(\alpha(s^{\prime}))]{}\varPhi_{K[n]}(\mathcal{O}_{Y})\xrightarrow[\alpha(s)]{}\mathcal{O}_{X}\enspace.
\]

\statementlabel{命题 11.4.7}其中一个有：

\[
\alpha(s)\circ\varPhi_{{\bf K}[n]}(\alpha(s^{\prime}))=\alpha(s\circ s^{\prime})~,
\]

其中 $s \circ s'$ 是 $s \otimes s'$ 通过命题11.4.6的态射的图像。

对于证明，我们参考 Kashiwara-Schapira [3，引理 11.1.4]。

\statementlabel{推论 11.4.8}让 $ p \in T^*X $  。环存在自然态射：

\[
\mathcal{E}_{X,p}^{\bf R}\to\mathrm{H o m}_{\mathsf{D}^{b}(X;p)}(\mathcal{O}_{X},\mathcal{O}_{X})\enspace.
\]

\statementlabel{证明}这是命题 11.4.7 当  $ A = A' = T_\delta^*(X \times X) $  和  $ K = K' = \mathbb{C}_\delta[-n] $  时的特殊情况。   $ \Box $ 

\statementlabel{定理 11.4.9}(i) 假设 (11.4.7), (11.4.8) 并设  $ p = (p_X, p_Y^a) \in \mathcal{A} $  。那么存在 $ s \in H^0(\mu \text{hom}(K, \Omega_{X \times Y/X}))_p $ 使得

\[
\mathcal{g}_{\mathbf{X},p_{X}}^{\mathbf{R}}\ni P\mapsto P s\quad{~a n d~}\quad\mathcal{E}_{Y,p_{Y}}^{\mathbf{R}}\ni Q\mapsto s Q\quad{~g i v e~i s o m o r p h i s n}
\]

(ii) 对于这样的 s，态射  $\alpha(s):\Phi_{\mathbb{K}[n]}(\mathcal{O}_{Y})\to\mathcal{O}_{X}$  是  $\mathbf{D}^{b}(X;p_{X})$  中的同构，并且  $\alpha(s)$  与  $\mathcal{E}_{X,p_{X}}^{\mathbb{R}}$  对  $\mathbf{D}^{b}(X;p_{X})$  中的  $\mathcal{O}_{X}$  和  $\mathcal{E}_{Y,p_{Y}}^{\mathbb{R}}$  对  $\mathbf{D}^{b}(Y;p_{Y})$  中的  $\mathcal{O}_{Y}$  的作用兼容。

\statementlabel{证明}(i) 我们在此承认 $s$  的存在满足(11.4.11)。这种  $s$  的构造是在 Kashiwara [5] 中执行的（也参见 Schapira [2, Chapter I, §5]）。 （在（loc. cit.）中，构造是在环  $\mathcal{E}_X$  和  $\mathcal{E}_Y$  的框架中完成的，但参数仍然适用于  $\mathcal{E}_X^R$  和  $\mathcal{E}_Y^R$  。）

(ii) 设  $ K^* = r_* K $  ，其中  $ r $  是规范映射  $ X \times Y \to Y \times X $  （参见 VII §2），并设  $ s' \in H^0(\mu \text{hom}(K^*, \Omega_{Y \times X/Y}))_p' $  和  $ p' = (p_Y, p_X^a) $  、 $ s' $  满足条件 (11.4.11)。然后  $ s'' = s \circ s' $  定义了  $ \mathcal{E}_X, p_X $  的自同构。由于 $ K \circ K^* \simeq \mathbb{C}_d $ （参见定理7.2.1的证明）， $ \mu \text{hom}(K \circ K^*[-n], \Omega_{X \times X/X}) $ 与 $ \mathcal{E}_X^R $ 同构， $ s \circ s' $ 是可逆的。根据命题 11.4.6，这意味着 $ \alpha(s) $  有一个右逆矩阵。类似地证明 $ \alpha(s) $  有左逆。因此 $ \alpha(s) $  是同构。

最后，再次根据命题 11.4.7，我们有 if  $ P \in \mathcal{E}_{X,p_X}^{\mathbb{R}}, Q \in \mathcal{E}_{Y,p_Y}^{\mathbb{R}} $  ：

\[
\alpha(s)\circ\varPhi_{K[n]}(Q)=\alpha(s\circ Q)=\alpha(P\circ s)=P\circ\alpha(s)\enspace.\quad\Box
\]

\statementlabel{定义 11.4.10}定理11.4.9的同构 $ \alpha(s) $ 在 $ \chi $ 之上被称为“量化接触变换”。

任何接触变换都可以被局部量化，但是“量化”不是唯一的（有许多部分满足（11.4.11））。

\statementlabel{推论 11.4.11}在定理11.4.9的情况下，令 $s$ 满足(11.4.11)。那么对于任何  $G \in \mathrm{Ob}(\mathbb{D}^b(Y))$  ，  $\alpha(s)$  定义了  $p_X: \chi_* \mu_{\mathrm{hom}}(G, \mathcal{O}_Y) \simeq \mu_{\mathrm{hom}}(\Phi_{K[n]}(G), \mathcal{O}_X)$  邻域中的同构。

回想一下， $ \chi $  是与 A 相关的接触变换。

\statementlabel{证明}应用定理 7.2.1 和 11.4.9。 ⑨

这个推论有很多应用，特别是当为  $Y$  的真实子流形  $N$  选择  $G = \mathbb{C}_N$  时，但我们不会在这里开发它（参见 Kashiwara-Schapira [3]）。

\statementlabel{备注 11.4.12}对于  $F \in \mathrm{Ob}(\mathbf{D}^b(X))$  ，  $H^j(\mu_{\mathrm{hom}}(F, \mathcal{O}_X))$  具有  $\mathcal{E}_X^\mathrm{R}$  模的结构，如命题 11.4.3 所示。但我们不知道  $\mu_{\mathrm{hom}}(F, \mathcal{O}_X)$  是否可以定义为  $\mathbf{D}(\mathcal{E}_X^\mathrm{R})$  的对象。

\subsection{微函数}
令  $M$  为维度  $n$  的实解析流形， $X$  为  $M$  的复化。

\statementlabel{定义 11.5.1}一套：

\[
\mathcal{C}_{M}=H^{n}(\mu_{M}(\mathcal{O}_{X})\otimes\operatorname{o r}_{M/X})\enspace,
\]

\[
\mathcal{B}_{M}=\mathcal{C}_{M}|_{M}=H_{M}^{n}(\mathcal{O}_{X})\otimes\operatorname{o r}_{M/X}.
\]

人们称 $\mathcal{C}_{M}$ （或 $\mathcal{P}_{M}$ ）为 $M$  上的佐藤微函数（或超函数）。

像往常一样，在  $\mathcal{C}_{M}$  的定义中，我们简写为  $o_{t_{M/X}}$  而不是  $\pi^{-1}or_{M/X}$  。

\statementlabel{命题 11.5.2}(i) 复数 $ \mu_{M}(\mathcal{O}_{X}) $  集中在度 $ n $  上。

(ii) 束 $ \mathcal{C}_M|\dot{\tau}_M^*X $  是圆锥形松弛的（即它在 $ \dot{\mathcal{T}}_M^*X/\mathbb{R}^+ $  上的直像是松弛的）。

(iii) 层 $ \mathscr{B}_{M} $  松弛。

(iv) M 上有一个自然精确的层序列：

\[
0\to\mathcal{A}_{M}\to\mathcal{B}_{M}\to\dot{\pi}_{*}\mathcal{C}_{M}\to0~.
\]

回想一下，  $ \mathcal{A}_M = \mathcal{O}_X|_{M} $  是  $ M $  上的实数分析函数。

\statementlabel{证明}(i) 由断言 (2.9.14) 和定理 11.1.3 得出。

(ii) 使用推论 11.4.11，可以证明  $ \mathcal{C}_M|_{\dot{T}_M X} $  与层  $ H^1(\mu_N(\mathcal{O}_X)) $  局部同构，其中  $ N $  是  $ X $  的严格伪凸开子集  $ \Omega $  的平滑边界。如果  $ \gamma $  表示映射  $ \dot{T}_M^*X \to \dot{T}_M^*X / \mathbb{R}^+ $  ，这意味着  $ \gamma_* \mathcal{C}_M $  与  $ H_{(X \setminus \Omega)}^1(\mathcal{O}_X|_{\bar{\Omega}}) $  局部同构，并且该束在 (2.9.16) 处是松弛的。

(iii) 令  $Z$  为  $M$  的局部闭子集。然后根据定理 11.1.3， $R\Gamma_Z\mathcal{B}_M = R\Gamma_Z(\mathcal{O}_X) \otimes \sigma_{M/X}[n]$  集中在零度。

(iv) 是所区分的三角形(4.3.1) 的一个特例。 □

请注意，从推论 11.4.11 可以得出，我们可以在层  $ \mathcal{C}_M $  上进行本地操作的真实解析接触变换。

用 sp 表示同构：

\[
\mathrm{sp}:\mathcal{B}_{M}\xrightarrow{\sim}\pi_{*}\mathcal{C}_{M}\enspace.
\]

\statementlabel{定义 11.5.3}令u 为M 上的超函数。 $ T_M^*X $  中sp(u) 的支持称为u 的奇异支持，记为SS(u)。

因此 SS(u) 是  $ T_M^*X $  和 SS(u)  $ \subset M \times_X T_X^*X $  的闭圆锥子集当且仅当 u 是实解析的。

让我们解释一下全纯函数的边界值的概念（参见 Schapira [3]）。令  $\Omega$  为  $X$  的开子集，满足：

\[
\begin{cases}{\bar{\Omega}\supset M{~a n d~t h e~i n c l u s i o n~}j\colon\Omega\hookrightarrow X{~i s~h o m e o m o r p h i c~t o~}}\\ {{t h e~i n c l u s i o n~o f~a n~o p e n~c o n v e x~s u b s e t~i n t o~}\mathbb{R}^{2n},{~l o c a l l y~o n~X~}.}\\ \end{cases}
\]

根据这个假设， $ R \mathcal{H}om(\mathbb{C}_{\bar{O}}, \mathbb{C}_X) \simeq \mathbb{C}_{\Omega} $  。将函子  $ R \mathcal{H}om(\cdot, \mathbb{C}_X) $  应用于态射  $ \mathbb{C}_{\bar{O}} \to \mathbb{C}_M $  我们得到  $ \mathbb{D}^b(X) $  中的态射：

\[
\omega_{M/X}\to\mathbb{C}_{\Omega}\;.
\]

（回想一下 $ \omega_{M/X} \simeq R\mathcal{H}_{om}(\mathbb{C}_{M},\mathbb{C}_{X}) \simeq or_{M/X}[-n] $ 。）

将函子  $ \mu_{\mathrm{hom}}(\cdot, \mathcal{O}_X) $  应用于 (11.5.4)，我们得到一个仍用 b 表示的态射：

\[
\mathfrak{b}:\mu\mathcal{h o m}(\mathbb{C}_{\Omega},\mathcal{O}_{X})\to\mu_{M}(\mathcal{O}_{X})\otimes\operatorname{o r}_{M/X}[n]
\]

特别是：

\[
\mathbf{b}:j_{*}j^{-1}\mathcal{O}_{X}\to\mathcal{B}_{M}\enspace.
\]

这最后一个态射通常被称为“边界值态射”。

如果  $ f \in \Gamma(\Omega; \mathcal{O}_X) $  ，  $ f $  定义  $ H^0(T^*X; \mu h o m(\mathbb{C}_\Omega, \mathcal{O}_X)) $  的元素。由于  $ \mu h o m(\mathbb{C}_\Omega, \mathcal{O}_X) $  由  $ \mathrm{SS}(\mathbb{C}_\Omega) $  支持，我们得到 (11.5.5)：

\[
\begin{array}{r}{\mathrm{S S}(\mathtt{b}(f))\subset T_{M}^{\ast}X\cap\mathrm{S S}(\mathbb{C}_{\Omega}),}\end{array}
\]

（参见 Delort-Lebeau [1] 以了解进一步的发展）。

注意，任何超函数都可以作为全纯函数的边界值之和来获得。当人们用边界值来“翻译”束 $ \mathcal{C}_M $ 的松弛度时，人们就得到了著名的楔形边缘定理。我们参考马蒂诺[2]。

现在让我们研究 $ \mathcal{D}_X $ 模的微功能解决方案。事实上，我们应该声明的大多数结果都适用于  $ \mathcal{E}_X $  甚至  $ \mathcal{E}_X^\mathbb{R} $  模的更一般框架（回想一下，根据命题 11.4.4， $ \mathcal{C}_M $  是  $ \mathcal{E}_X^\mathbb{R} $  模），但为了简单起见，我们更愿意将自己限制在  $ \mathcal{D}_X $  模。

\statementlabel{命题 11.5.4}设 M 是一个连贯的  $ \mathcal{D}_X $  模。然后：

\[
\begin{array}{r}{\mathrm{S S}(R\mathcal{H o m}_{\mathcal{D}_{X}}(\mathcal{M},\mathcal{C}_{M}))\subset C_{T_{M}^{*}X}(\mathrm{c h a r}(\mathcal{M}))~.}\end{array}
\]

特别是：

\[
\operatorname{s u p p}(R\mathcal{H}o m_{\mathcal{D}_{X}}(\mathcal{M},\mathcal{C}_{M}))\subset T_{M}^{\ast}X\cap\operatorname{c h a r}(\mathcal{M})\enspace.
\]

\statementlabel{证明}将定理 6.4.1 应用于  $F = \mathcal{R}\mathcal{H}om_{\mathcal{D}_X}(\mathcal{M},\mathcal{O}_X)$  和定理 11.3.3。 □

如果  $P$  是微分算子，我们可以通过这个结果看到  $P$  定义了集合  $T_M^*X \setminus \{\sigma(P)^{-1}(0)\}$  上的束  $\mathcal{C}_M$  的同构，并且如果  $u$  是超函数，则  $\mathrm{SS}(u) \subset \mathrm{SS}(Pu) \cup \{\sigma(P)^{-1}(0)\}$  。

\statementlabel{定义 11.5.5}(i) 如果  $ \dot{T}_M^* X \cap \text{char}(\mathcal{M}) = \emptyset $  ，则连贯的  $ \mathcal{D}_X $  -模  $ \mathcal{M} $  被称为椭圆。

(ii) 让 $ \theta \in T_{T_M^*X} T^*X $  。有人说，如果  $ \theta \notin C_{T_M^*X}(\text{char}(\mathcal{M})) $  ，则  $ \theta $  对于  $ \mathcal{M} $  是微双曲的。

如果  $ \theta \in T^*M $  （回想一下嵌入  $ T^*M \hookrightarrow T_{T_M^*X} T^*X $  ，VI §2），我们可以简单地说  $ \theta $  对于  $ \mathcal{M} $  是双曲的。

根据命题 11.5.4，如果  $\mathcal{M}$  是椭圆形，则态射：

\[
\mathrm{R}\mathcal{H}\mathrm{o}m_{\mathfrak{D}_{X}}(\mathcal{M},\mathcal{A}_{M})\to\mathrm{R}\mathcal{H}\mathrm{o}m_{\mathfrak{D}_{X}}(\mathcal{M},\mathcal{B}_{M})
\]

是同构（将  $ R_{\mathcal{H}om_{\mathfrak{D}}(\mathcal{M},\cdot)} $  应用于精确序列 11.5.1）。

如果  $\theta$  是微双曲线，则  $\theta \notin SS(R\mathcal{H}om_{\mathcal{D}_x}(\mathcal{M},\mathcal{C}_M))$  。如果  $\theta$  是双曲线，则  $\theta \notin SS(R\mathcal{H}om_{\mathcal{D}_x}(\mathcal{M},\mathcal{B}_M))$  。

偏微分方程理论中的一个经典问题是奇点传播问题。

令 $U$  为 $T^*X$  的开子集， $f$  为 $U$  上的全纯函数。假设  $\operatorname{Im}f$  在  $T_M^*X$  上消失。然后，对于每个 $p\in U\cap T_M^*X$ ，穿过 $(T^*X)^\mathbb{R}$ 的对合子流形 $\{ \operatorname{Im}f=0 \}$ 的 $p$ 的双特征曲线 $b_p$ 包含在 $p$ 邻域中的 $T_M^*X$ 中（参见附录）。

\statementlabel{命题 11.5.6}令  $\mathcal{M}$  为连贯的  $\mathcal{D}_X$  模，并让  $f$  为  $U \subset T^*X$  上的全纯函数。假设：

\[
\mathrm{I m}f|_{T_{M}^{*}X\cap U}=0,\qquad\mathrm{c h a r}(\mathcal{M})\cap U\subset\{f=0\}.
\]

那么对于流形  $\{\operatorname{Im}f=0\}$  的任何双特征  $b$  ，层  $\mathcal{E}x t_{\mathcal{D}_X}^{j}(\mathcal{M},\mathcal{C}_M)|_{b}$  是局部常数  $(j\in\mathbb{N})$  。

\statementlabel{证明}根据命题 11.5.4， $ \mathrm{SS}(R\mathcal{H}om_{\mathcal{D}_X}(\mathcal{M},\mathcal{C}_M)) $  包含在  $ C_{T_M^*X}(\{f^{-1}(0)\}) = \{\theta \in T^*T_M^*X; \langle \theta, H_f \rangle = 0\} $  中。因此应用命题 5.4.5 (ii) 就足够了。   $ \square $ 

示例 11.5.7。 (i) 令 $X$  为复流形。在 $X^R$  上， $\mathcal{D}_{X \times \bar{X}}$  模 $\mathcal{D}_X \boxtimes \mathcal{O}_{\bar{X}}$  是椭圆形的。这就是柯西-黎曼系统。请注意：

\[
\begin{array}{r}{\mathcal{O}_{X}\simeq R\mathcal{H}\mathcal{H}o m_{\mathcal{D}_{\bar{X}}}(\mathcal{O}_{\bar{X}},\mathcal{B}_{X^{\mathtt{R}}})\simeq R\mathcal{H}\mathcal{H}o m_{\mathcal{D}_{X\times\bar{X}}}(\mathcal{D}_{X}\boxtimes\mathcal{O}_{\bar{X}},\mathcal{B}_{X^{\mathtt{R}}}).}\end{array}
\]

(ii) 令 $ (z) = (z_1, \ldots, z_n) $  为 $ X $  上的全纯坐标系， $ (z;\zeta) $  为 $ T^*X $  上的相关坐标，以及 $ z = x + \sqrt{-1}y $  、 $ \zeta = \xi + \sqrt{-1}\eta $  ，并让 $ M = \{z \in X; y = 0\} $  使得 $ T_M^*X = \{y = \xi = 0\} $  。

如果关联系统  $\mathcal{M} = \mathcal{D}_X / \mathcal{D}_X P$  是椭圆形或双曲形等，则可以说运算符  $P$  是椭圆形或双曲形等。如果  $\sigma(P)(x; i\eta) \neq 0$  对于  $\eta \neq 0$  则  $P$  是椭圆形。特别是拉普拉斯算子 $\varDelta = \sum_{j=1}^n D_j^2$  是椭圆形的。

(iii) 令 $(z;\zeta)$  如上，并令 $p = (x_{\circ};i\eta_{\circ}) \in T_{M}^{*}X, \theta_{\circ} \in (T_{T_{M}^{*}X}T^{*}X)_{p}$  。那么 $\theta_{\circ}$ 对于 $P$ 在 $p$ 来说是微双曲如果

\[
\left\{\begin{aligned}{}&{{}\sigma(P)((x;i\eta)+\varepsilon\theta)\neq0}\\ {}&{{}\operatorname*{f o r\quad}\quad|x-x_{\circ}|+|\eta-\eta_{\circ}|+|\theta-\theta_{\circ}|+\varepsilon\ll1~,}&{\varepsilon>0~.}\\ \end{aligned}\right.
\]

事实上，使用博赫纳管定理的本地版本（参见 Komatsu [2]）足以检查 (11.5.9) 中的  $ \theta = \theta_o $  。波算子  $ \Box = D_1^2 - \sum_{j=2}^n D_j^2 $  满足 (11.5.9)  $ \eta_o = 0 $  、  $ \theta = dx_1 $  。

(iv) 假设 $ \sigma(P) $  在 $ T_M^*X $  上是真实的。如果 $ u $  是方程 $ Pu = 0 $  的微函数解， $ \text{supp}(u) $  将是哈密顿场 $ H_{\text{Im}\sigma(P)}^R $  的积分曲线的并集。回想一下，对于  $ T^*X $  上的实函数  $ h $  ，  $ H_h^R $  由下式给出（参见 11.1.4）：

\[
H_{h}^{\mathbb{R}}=\frac{1}{2}\sum_{j=1}^{n}\left(\frac{\partial h}{\partial\xi_{j}}\frac{\partial}{\partial x_{j}}-\frac{\partial h}{\partial x_{j}}\frac{\partial}{\partial\xi_{j}}+\frac{\partial h}{\partial y_{j}}\frac{\partial}{\partial\eta_{j}}-\frac{\partial h}{\partial\eta_{j}}\frac{\partial}{\partial y_{j}}\right).
\]

因此：

\[
\left.H_{\mathrm{l m}\sigma(P)}^{\mathrm{R}}\right|_{T_{M}^{*}X}=\frac{1}{2}\sum_{j=1}^{n}\left(\frac{\partial\operatorname{R e}\sigma(P)}{\partial x_{j}}\frac{\partial}{\partial\eta_{j}}-\frac{\partial\operatorname{R e}\sigma(P)}{\partial\eta_{j}}\frac{\partial}{\partial x_{j}}\right).
\]

现在我们将描述微函数的主要操作。令 N 为另一个实解析流形，Y 为 N 的复化。

自然态射（命题 4.3.6）：

\[
\mu_{M}(\mathcal{O}_{X})\boxtimes\mu_{N}(\mathcal{O}_{Y})\to\mu_{M\times N}(\mathcal{O}_{X}\boxtimes\mathcal{O}_{Y})
\]

与态射  $ \mathcal{O}_{X} \boxtimes \mathcal{O}_{Y} \to \mathcal{O}_{X \times Y} $  结合。定义态射：

\[
\mathcal{C}_{M}\boxtimes\mathcal{C}_{N}\to\mathcal{C}_{M\times N}\enspace.
\]

现在让  $f: N \to M$  成为一个实解析图，并且仍然用  $f$  表示复化  $f: Y \to X$  。考虑映射（参见 IV §3）：

\[
T_{N}^{*}Y\xleftarrow{\,f_{N}^{\prime}\,}N\underset{M}{\times}T_{M}^{*}X\xrightarrow{\,f_{N\pi}\,}T_{M}^{*}X\enspace.
\]

根据命题 4.3.5，存在自然态射：

\[
R^{\iota}f_{N!}^{\prime}(\omega_{N/M}\otimes f_{N\pi}^{-1}\mu_{M}(\mathcal{O}_{X}))\to\mu_{N}(f^{-1}\mathcal{O}_{X}\otimes\omega_{Y/X})\enspace.
\]

结合态射 $ f^{-1}\mathcal{O}_X \to \mathcal{O}_Y $ 我们得到态射：

\[
R^{\mathfrak{t}}f_{N}^{\prime}\underset{}{f_{N\pi}^{-1}}{\mathcal{C}}_{M}\to\mathcal{C}_{N}\enspace.
\]

这意味着，如果  $u$  是定义在  $T_M^*X$  的开子集  $U$  上的微函数，并且如果  $W$  是  $T_N^*Y$  的开子集，则：  $t'f_N': f_{N\pi}^{-1}(\text{supp}(u)) \cap (t'f_N')^{-1}W \to W$  是

正确的，那么我们可以通过  $f$  来定义  $u$  的逆像  $f^*u$  ，作为  $W$  上的微函数。特别是，如果  $u$  是  $M$  上的超函数，并且  $f_N'(f_N)$  是  $f_{N\pi}^{-1}(\mathrm{SS}(u))$  上的真函数，则可以定义超函数  $f^*u$  且  $\mathrm{SS}(f^*u)$  包含在  $f_N'(f_{N\pi}^{-1}(\mathrm{SS}(u)))$  中。当 $f:N \to M$ 是一种沉浸时，人们通常会写 $u|_{N}$ 而不是 $f^*u$ 

要定义直像，请在  $ M $  上引入分析密度的束  $ \mathcal{V}_M = \mathcal{A}_M^{(n)} \otimes \sigma t_M $  （ $ \mathcal{V}_N $  也类似）。根据命题 4.3.4，存在自然态射：

\[
R f_{N\pi}{}^{t}f_{N}^{\prime-1}\mu_{N}(\varOmega_{Y})\to\mu_{M}(R f_{!}\varOmega_{Y})\enspace.
\]

结合积分态射（定理11.1.4）我们得到态射：

\[
R f_{N\pi!}f_{N}^{\prime-1}\left(\mathcal{C}_{N}\mathop{{\otimes}}_{\mathcal{A}_{N}}\mathcal{V}_{N}\right)\to\mathcal{C}_{M}\mathop{{\otimes}}_{\mathcal{A}_{M}}\mathcal{V}_{M}\enspace.
\]

特别是，如果 $v$ 是 $N$ 上的“超函数密度”，即 $v \in \Gamma(N; \mathcal{B}_N \otimes_{\mathcal{A}_N} \mathcal{V}_N)$ 和 $f$ 在 $\operatorname{supp}(v)$ 上是真映射，那么可以通过 $f$ 定义 $v$ 的 $\operatorname{direct}$ 图像 $f_* v$ ， $M$ 上的超函数密度 $\operatorname{SS}(f_* v)$  包含在  $f_{N\pi} \iota' f_N'^{-1}(\operatorname{SS}(v))$  中。请注意，如果  $f: N \to M$  是平滑的，则  $N \times_M T_M^* X$  是  $T_N^* Y$  的闭合子流形  $V$ ，并且  $\operatorname{SS}(f_* v)$  由  $\operatorname{SS}(v) \cap V$  的  $f_\pi$  包含在图像中。人们经常写 $\int_f v$ 而不是 $f_* v$ 。

通过考虑微微分系统的微函数解的运算，可以推广前述结构。然而，我们将限制自己为 $ \mathcal{D}_x $  -模，并且我们将只考虑逆像。

令 $f: N \to M$  如前所述。利用(11.3.9)的态射和命题4.3.5，我们得到自然态射：

\[
R^{\prime}f_{N!}^{\prime}f_{N\pi}^{-1}R\mathcal{H}o m_{\mathcal{D}_{X}}(\mathcal{M},\mathcal{C}_{M})\to R\mathcal{H}o m_{\mathcal{D}_{Y}}(\underline{{f}}^{-1}\mathcal{M},\mathcal{C}_{N})\enspace.
\]

特别是假设  $f$  对于  $\mathcal{M}$  来说是非特征的。然后  $f^{-1}\mathcal{M}$  集中在零度，而且  $f_N'$  在  $f_{N\pi}^{-1}(\text{supp}(R\mathcal{H}om_{\mathcal{D}_x}(\mathcal{M},\mathcal{C}_M)))$  上是有限的，根据命题 11.5.4，因为  $f'$  在  $f_\pi^{-1}(\text{char}(\mathcal{M}))$  上是有限的。因此，在这种情况下， (11.5.14) 对每个  $j\in\mathbb{N}$  导出态射：

\[
\mathfrak{f}_{N}^{\prime}\mathfrak{f}_{N\pi}^{-1}\mathcal{E}x t_{\mathcal{D}_{X}}^{\mathfrak{j}}(\mathcal{M},\mathcal{C}_{M})\to\mathcal{E}x t_{\mathcal{D}_{Y}}^{\mathfrak{j}}(\underline{{\mathfrak{f}}}^{-1}\mathcal{M},\mathcal{C}_{N})\enspace.
\]

换句话说，系统 $ \mathcal{M} $ 的微函数 $ u $ 解的逆像 $ f^* u $ 是系统 $ f^{-1} \mathcal{M} $ 的微函数解。

将定理 6.7.1 应用于  $F = \mathcal{R}_{\mathcal{Hom}_{\mathcal{D}_X}}(\mathcal{M}, \mathcal{O}_X)$  并使用定理 11.3.5，我们可以求解微双曲系统微函数解的柯西问题：

\statementlabel{命题 11.5.8}令 V 为  $ T_{N}^{*} $  Y 的开子集。假设：

(i)  $f$  对于  $\mathcal{M}$  来说是非特征的（即：  $T_{Y}^{*}X\cap f_{\pi}^{-1}(\mathrm{char}(\mathcal{M}))\subset Y\times_{X}T_{Y}^{*}X)$  ，

(ii)  $ f_{N\pi}|_{t f_{N}^{\prime}-1}(V):t f_{N}^{\prime-1}(V) \to T_{M}^{*}X $  对于  $ C_{T_{M}^{*}X}(\mathrm{char}(\mathcal{M})) $  来说是非特征的，

（三） $ f^{\prime-1}(V) \cap f_{\pi}^{-1}(\text{char}(\mathcal{M})) \subset Y \times_X T_M^* X $ 。

那么态射 (11.5.14) 是  $V$  上的同构。

示例 11.5.9。让  $(z;\zeta)$  表示  $T^*X$  上的坐标，如例 11.5.7 (ii) 中所示，并设置  $\zeta=(\zeta_1,\zeta')$  。令  $N$  为  $M$  的超曲面  $\{x_1=0\}$  ，并令  $P$  为  $m$  阶微分算子，在  $\pm dx_1$  方向上为双曲线，即满足：

\[
\sigma(P)(x;i\eta+t\theta)\neq0\qquad\mathrm{f o r~a n y~}t\in\mathbb{R}\backslash\{0\}~,~\quad\mathrm{a n y~}(x;\eta)\in\mathbb{R}^{n}\times\mathbb{R}^{n}
\]

与 $ \theta = (1, 0, \ldots, 0) $  。

考虑柯西问题：

\[
\left\{\begin{aligned}{}&{{}P u=0~,}&{\gamma(u)=(w)}\\ {}&{{}\mathbf{w i t h}}&{\gamma(u)=(u|_{N},\ldots,\partial^{m-1}u/\partial^{m-1}x_{1}|_{N})~.}\\ \end{aligned}\right.
\]

应用命题 11.5.8，我们得到对于任何部分  $ (w) \in \mathcal{B}_N^n $  都存在 (11.5.16) 的唯一部分  $ u \in \mathcal{B}_M|_N $  解。此外 $ P: \mathcal{B}_M|_N \to \mathcal{B}_M|_N $  是满射的。

\subsection{习题}
\statementlabel{练习 XI.1}(i) 令 M 为左相干  $ \mathcal{D}_X $  模，并令  $ M^* $  如 (11.2.19) 中定义。证明同构：

\[
\Omega_{X}\overset{L}{\otimes}_{\mathcal{D}_{X}}\mathcal{M}[-n]\simeq R\mathcal{H}{}_{{O m}_{\mathcal{D}_{X}}}(\mathcal{M}^{\ast},\mathcal{O}_{X})
\]

其中 $ n = \dim_{\mathbb{C}} X $ 

(ii) 一组  $ \mathrm{DR}(\mathcal{M}) = \Omega_X \otimes_{\mathcal{D}_X}^L \mathcal{M} $  。证明如果  $ \mathcal{M} $  是完整的，则  $ \mathrm{DR}(\mathcal{M}) $  是  $ X $  上的反常层。

\statementlabel{练习 XI.2}(i) 如果假设 M 或 N 完整，证明态射 (11.3.8) 是同构。

(ii) 类似地，让 $\mathcal{M}$ （分别为 $\mathcal{N}$ ）成为连贯的 $\mathcal{D}_{X}$ 模（分别为 $\mathcal{D}_{Y}$ 模）。假设 $\mathcal{M}$  或 $\mathcal{N}$  是完整的。证明同构：

\[
\left(\begin{matrix}{{\Omega}_{X}^{{{\tiny~{L}~}}}{{\otimes}}\;{{\mathcal M}}}\\ {{\mathfrak D}_{X}^{{{\tiny~{L}~}}}}\\ \end{matrix}\right){{\boxtimes}}\left(\begin{matrix}{{\Omega}_{Y}^{{{\tiny~{L}~}}}{{\otimes}}\;{{\mathcal N}}}\\ {{\mathfrak D}_{Y}^{{{\tiny~{L}~}}}}\\ \end{matrix}\right)\simeq{\Omega}_{X\times Y}^{{{\tiny~{L}~}}}{{\otimes}}_{{\mathfrak D}_{X\times Y}}\left({{\mathcal M}}{{{{\boxtimes}}}}\;{{\mathcal N}}\right)\;.
\]

（提示：使用定理 11.3.7 和一些泛函分析。）

\statementlabel{练习 XI.3}令  $M$  为维度  $n$  的实解析流形， $X$  为  $M$  的复化。让 $M$  成为一个完整的 $\mathcal{D}_X$  模。证明：

\[
{\mathcal{E}}x t{\mathcal{I}}_{{\mathcal{D}}_{X}}({\mathcal{M}},{\mathcal{C}}_{M})=0\qquad{\mathrm{f o r}}\qquad j\geqslant n~.
\]

（提示：使用练习 X.8。）

\statementlabel{练习 XI.4}令 $M$  为实解析流形， $X$  为 $M$  的复化， $K$  为 $M$  的紧凑子集。证明可以自然地识别 $\Gamma_K(M; \mathcal{B}_M \otimes_{\mathcal{A}_M} \mathcal{V}_M)$  到 $(\mathcal{O}_X(K))',$  拓扑向量空间 $\mathcal{O}_X(K) = \Gamma(K; \mathcal{O}_X)$  的对偶空间。 （特别是， $\Gamma_K(M; \mathcal{B}_M)$  自然地被赋予了  $FS$  空间的拓扑。）

参见佐藤[1]、马蒂诺[1]、夏皮拉[1]。

\statementlabel{练习 XI.5}令 $M$  为实解析流形， $X$  为 $M$  的复化，让 $M$  为相干的 $\mathcal{D}_X$  模。假设：

(11.E.4) 在  $X$  上全局， $\mathcal{M}$  承认有限自由分辨率（如 (11.2.16)），

(11.E.5) M 是椭圆形。

令 $ \Omega \subset \mathcal{M} $  为具有平滑边界的相对紧凑的开子集。假设：

(11.E.6)  $ \partial\Omega $  对于  $ \mathcal{M} $  是双曲的（即：  $ \forall\theta\in\dot{T}_{\partial\Omega}^{*}M $  、  $ \theta\notin C_{T_{M}^{*}X}(\mathrm{char}(\mathcal{M})) $  ）。

证明空间 $ H^j(R\Gamma(\Omega, R\mathcal{H}_{om_{\mathcal{D}_Y}}(\mathcal{M},\mathcal{B}_M))) $  是有限维的。

（提示：让  $\mathcal{M}$  为  $\mathcal{M}$  的有限自由分辨率。证明同构  $RHom_{\mathfrak{D}_X}(\mathcal{M}, \mathcal{A}_M(\bar{\mathcal{O}})) \simeq RHom_{\mathfrak{D}_X}(\mathcal{M}, \mathcal{B}_M(\Omega))$  。然后使用  $\mathcal{B}_M(\Omega) \simeq (\Gamma_{\partial\Omega}(M;\mathcal{B}_M) \to \Gamma_{\bar{\Omega}}(M;\mathcal{B}_M)) [1]$  是 FS 空间的复形， $\mathcal{A}_M(\bar{\Omega})$  是 DFS 空间并通过泛函分析得出结论。）

有关进一步的发展，请参阅 Schapira-Schneiders [1] 或 Schapira [4]。

\statementlabel{练习 XI.7}令 $X$  为复流形， $\mathcal{M}$  为连贯的 $\mathcal{D}_X$  模。证明如果  $\theta \in T^*X$  对于  $\mathcal{M}$  来说是非特征的（即  $\theta \notin \mathrm{char}(\mathcal{M})$  ），则  $\theta$  被视为  $T^*X^\mathrm{R}$  的向量，对于  $\mathcal{M} \boxtimes \mathcal{O}_{\bar{X}}$  来说是双曲的（即：  $\theta \notin C_{T^{\#\mathrm{R}}(X \times \bar{X})}$  (char(  $\mathcal{M}$  )  $\times T_{\bar{X}}^* \bar{X}$  ））。

\statementlabel{练习 XI.8}令 $X$  为复流形， $\mathcal{M}$  为满足(11.E.4) 的连贯 $\mathcal{D}_X$  模，并令 $\Omega \subset \subset X$  为具有平滑边界的开子集。假设 $\dot{T}_{\partial\Omega}^*X \cap \mathrm{char}(\mathcal{M}) = \emptyset$ 。证明空间 $H^j(R\Gamma(\Omega; R\mathcal{H}om_{\mathcal{D}_X}(\mathcal{M},\mathcal{O}_X)))$  是有限维的。

（提示：使用练习 XI.6 和 XI.7。）

有关特定案例，请参阅 Bony-Schapira [1]

\statementlabel{练习 XI.9}令 $f: Y \to X$  为复流形的态射。

(i) 证明 $ \mathbf{D}^{b}(\mathcal{D}_{X}^{op}) $  中存在自然态射：

\[
R f_{!}\left(\begin{matrix}{L}\\ {\varOmega_{Y}\bigotimes_{\mathcal{D}_{Y}}\mathcal{D}_{Y\to X}}\\ \end{matrix}\right)[\mathsf{d i m}_{\mathbb{C}}Y]\to\varOmega_{X}[\mathsf{d i m}_{\mathbb{C}}X]~.
\]

（这是定理 11.1.4 的改进，参见 Schneiders [1]。）

(ii) 设 $ \mathcal{N} $  为右相干 $ \mathcal{D}_Y $  模。构造态射：

\[
R f_{!}(R\mathcal{H o m}_{\mathcal{D}_{Y}}(\mathcal{N},\Omega_{Y}[\dim_{\mathbb{C}}Y]))\to R\mathcal{H o m}_{\mathcal{D}_{X}}(\underline{{f}}_{*}\mathcal{N},\Omega_{X}[\dim_{\mathbb{C}}X])
\]

\[
R f_{*}(R\mathcal{H}o m_{\mathcal{D}_{Y}}(\mathcal{N},\Omega_{Y}[\dim_{\mathbb{C}}Y]))\to R\mathcal{H}o m_{\mathcal{D}_{X}}(\underline{{f}}_{!}\mathcal{N},\Omega_{X}[\dim_{\mathbb{C}}X])~.
\]

(iii) 设  $f: N \to M$  为实解析流形的态射，并且仍用  $f$  表示  $f$  的复化  $f: Y \to X$  。让 $\mathcal{N}$  成为一个连贯的右 $\mathcal{D}_{\Upsilon}$  模。构造自然态射：

\[
R f_{N\pi!}{}^{t}f_{N}^{t-1}R\mathcal{H}{o m}_{\mathcal{D}_{Y}}\bigg(\mathcal{N},\mathcal{C}_{N}\otimes\mathcal{V}_{N}\bigg)\to R\mathcal{H}{o m}_{\mathcal{D}_{X}}\bigg(\underline{{f}}_{*}\mathcal{N},\mathcal{C}_{M}\otimes\mathcal{V}_{M}\bigg)\;.
\]

(iv) 在 (iii) 的情况下，假设  $f$  优于  $N \cap \operatorname{supp}(\mathcal{N})$  。构造自然态射：

(11.E.9)

\[
\begin{array}{r l}&{\mathrm{E.9})}\\ &{\Gamma\bigg(N;\mathcal{H o m}_{\mathfrak{D}_{Y}}\bigg(\mathcal{N},\mathcal{B}_{N}\otimes\mathcal{V}_{N}\bigg)\bigg)\to\Gamma\bigg(M;\mathcal{H o m}_{\mathfrak{D}_{X}}\bigg(H^{0}(\underline{{\mathfrak{f}}}_{\ast}\mathcal{N}),\mathcal{B}_{M}\otimes\mathcal{V}_{M}\bigg)\bigg)}\end{array}
\]

（即：如果  $v$  是系统  $\mathcal{N}$  的  $N$  解上的超函数密度，则  $\int_{f} v$  是系统  $H^0(\underline{f}_*\mathcal{N})$  的解。

\statementlabel{练习 XI.10}设 X 为复流形，N 为余维 d 的实解析子流形，满足：

\[
T N+\sqrt{-1}T N=N\times T X_{x}.
\]

令 Y 为 N 的复化并考虑映射的交换图，其中  $ \delta $  是对角线嵌入：

\[
\begin{array}{l l l}{Y\subset\underbrace{\longrightarrow}_{g}X\times\bar{X}}&{\quad}&{Y\subset\underbrace{\longrightarrow}_{g}X\times\bar{X}}\\ {\quad\quad\quad\quad\quad\quad\quad\quad\quad\quad\quad\quad\quad\quad}\\ {\quad\quad\quad\quad\quad\quad\quad\quad\quad\quad\quad\quad\quad\quad}\\ {\quad\quad\quad\quad\quad\quad\quad\quad\quad\quad\quad\quad\quad}\\ {\quad\quad\quad\quad\quad\quad\quad\quad\quad\quad\quad\quad}\\ {\quad\quad\quad\quad\quad\quad\quad\quad\quad\quad\quad}\\ {N\subset\underbrace{\longrightarrow}_{i}X}&{\quad}&{X}\end{array}.
\]

设置 $ \mathcal{M} = g^{-1}(\mathcal{D}_X \boxtimes \mathcal{O}_{\bar{X}}) $ 。

(i) 证明： $ \overline{Y} \times_{X} T^{*} X = \text{char}(\mathcal{M}) $  。 （请注意 f 是平滑的。）

(ii) 证明同构：

\[
f^{-1}\mathcal{O}_{X}\simeq\mathrm{R H o m}_{\mathcal{D}_{Y}}(\mathcal{M},\mathcal{O}_{Y})\enspace.
\]

(iii) 证明同构：

\[
\mu_{N}(\mathcal{O}_{X})\otimes\operatorname{o r}_{N/X}[d]\simeq R\mathcal{H o m}_{\mathcal{D}_{Y}}(\mathcal{M},\mathcal{C}_{N})\enspace.
\]

（提示：使用推论 6.7.4。）

参见柏原河合 [1]。

注： $ \mathcal{D}_{Y} $  -模 M 称为诱导柯西-黎曼系统。

\statementlabel{练习 XI.11}令 $M$  为实解析流形， $N$  为子流形， $X$  为 $M$  的复化， $Y$  为 $X$  中 $N$  的复化。令  $\mathcal{M}$  成为一个连贯的  $\mathcal{D}_X$  模，使得  $Y$  对于  $\mathcal{M}$  来说是非特征的。

(i) 证明 Y 上的同构：

\[
R\mathcal{H}o m_{\mathfrak{D}_{Y}}(\mathcal{M}_{Y},\mathcal{O}_{Y})\simeq R\mathcal{H}o m_{\mathfrak{D}_{X}}(\mathcal{M},R\varGamma_{Y}(\mathcal{O}_{X}))\otimes\omega_{Y/X}^{\otimes-1}
\]

（使用推论 5.4.11）。

(ii) 证明 N 上的同构：

\[
R\mathcal{H}{o m}_{\mathfrak{D}_{Y}}(\mathcal{M}_{Y},\mathcal{B}_{N})\simeq R\mathcal{H}{o m}_{\mathfrak{D}_{X}}(\mathcal{M},\varGamma_{N}(\mathcal{B}_{M}))\otimes\omega_{N/M}^{\otimes-1}.
\]

(iii) 令 $\Omega$  为 $M$  的开子集，使得 $\bar{\Omega} \supset N$  和包含 $\Omega \subsetneq M$  局部同胚于将凸开子集包含到 $\mathbb{R}^n$  中。构造边值态射：

\[
\mathsf{b}:R\mathcal{H}\mathcal{o}m_{\mathcal{D}_{X}}(\mathcal{M},\varGamma_{\Omega}(\mathcal{B}_{M}))|_{N}\to R\mathcal{H}\mathcal{o}m_{\mathcal{D}_{Y}}(\mathcal{M}_{Y},\mathcal{B}_{N})\enspace.
\]

（提示：使用 (11.5.4)。参见 Schapira [3]。）

\statementlabel{练习 XI.12}令 $X$  为复流形， $U$  为 $T^*X$  、 $p \in U$  的开子集，并令 $\mathcal{M}$  为 $U$  上的有限自由左 $\mathcal{E}_X^\mathsf{R}$  -模的有界复数：

\[
\mathcal{M}^{\ast}:0\to(\mathcal{E}_{X}^{\mathbb{R}})^{N_{d}}\longrightarrow\cdots\xrightarrow[P_{0}]{\phantom{x}}\left(\mathcal{E}_{X}^{\mathbb{R}}\right)^{N_{0}}\longrightarrow0\quad.
\]

（  $P_j$  是  $U$  上  $\mathcal{E}_X^{\mathrm{R}}$  部分的矩阵，作用于右侧）

(i) 使用推论 11.4.8，表明复合体：

\[
0\longrightarrow\mathcal{O}_{X}^{N_{0}}\xrightarrow[P_{0}]{}\cdots\longrightarrow\mathcal{O}_{X}^{N_{d}}\longrightarrow0
\]

在  $ \mathbf{D}^{b}(X; p) $  中有明确定义。我们用 $ \mathrm{Sol}_{p}(\mathcal{M}) $  表示它。

(ii) 证明  $ \mathrm{SS}(\mathrm{Sol}_{p}(\mathcal{M})) \subset \mathrm{supp}(\mathcal{M}) $  在  $ p $  的邻域内。 （回想一下 $ \mathrm{supp}(\mathcal{M}) = \bigcup_{j} \mathrm{supp}(H^{j}(\mathcal{M})) $ 。）

参见 Kashiwara-Schapira [3，第 11 章，§4]。

\subsection{注记}
详细讨论连贯 $\mathcal{O}_{x}$ 模理论的历史超出了本书的范围。我们只强调 Oka [1] 的名字，

Cartan [2] 和 Serre [1] 附在其后，并请参阅 Hörmander [1] 进行评论。

同样，我们不回顾 $ \mathcal{D}_x $ 模的理论，而是参考Schapira [2]中的注释，简单地回顾一下这个理论出现在70年代，与Kashiwara [1]和Bernstein [1,2]一起。

Leray [3] 对经典的 Cauchy-Kowalewski 定理进行了改进，得到了定理 11.3.1。然后它在两个不同的方向上扩展到系统。其中之一是由 Kashiwara [1] 提出的，涉及“柯西问题”（定理 11.3.5）。另一条与“传播”相关：这是定理11.3.3，是由Zerner [1]、Bony-Schapira [1]和Kashiwara [5]的部分结果得出的Kashiwara-Schapira [2,3]。定理 11.3.7 已在 Kashiwara 中得到[3]。这是许多重要发展的起点，我们提到 Björk [2]。

微局部算子的环  $ \delta_{x} $ （命题 11.4.4）以及自然附加到它的各种模是由 Sato-Kawai-Kashiwara [1] 构建的，他发展了 Sato [2] 的微局部化和微函数的思想。 Sato [1] 于 1959 年将超函数引入为全纯函数的边界值之和，并用局部上同调来解释它们。在研究这些边界值时，特别是在尝试理解 Martineau [2] 的“楔形边缘定理”时，微函数的概念现在显得相当自然。它是研究线性偏微分方程的有效工具，如第 5 节所示。

实量子化接触变换的理论可以追溯到 Maslov [1]，并由 Hörmander [2] 在  $ C^{\infty} $  范畴中以及由 Sato-Kawai-Kashiwara [1] 在分析案例中系统地发展。这些作者证明了微函数和微局部算子的层通过这种变换被局部保留，这是推论 11.4.11 的一个特例。让我们强调这样一个事实，即 Kashiwara-Schapira [3] 提出的定理 11.4.9 更加强大，因为它涉及层  $ \mathcal{O}_{X} $  本身，而不是它的微局部化。推论 11.4.11 对于  $ \mathcal{O}_{X} $  沿实子流形微局域化的上同调消失，或具有简单特征的微微分方程组的上同调消失有很好的应用。为了简洁起见，我们没有将其包含在这里，我们指的是 Kashiwara-Schapira [3,4]。

有大量涉及解析偏微分方程的文献，§5 只是该主题的概览，其出发点是 Sato [2] 的命题 11.5.4 (ii)。典型的特征是 Sato-Kawai-Kashiwara [1] 的命题 11.5.6，该命题现已被许多作者大大改进（参见 Tose [1] 的评论）和由 Bony-Schapira [2]（处理了一名算子的情况）发起的双曲方程，然后由获得命题 11.5.4 (i) 和 11.5.8 的 Kashiwara-Schapira [1] 开发。现在，同样的技术也被用于边值问题的研究，这里没有讨论这些问题（除了公式（11.5.5）和（11.5.7），由 Schapira [3] 提出）。让我们强调一下，函子  $ \mu h o m $  允许我们定义新的微函数束和新的波前组，非常适合边值问题（参见 Schapira (loc. cit.)、Kataoka [1] 和 Uchida [1]）。

最后，让我们提一下，在处理分布的分析奇异性时，大多数人使用 Sjöstrand [1] 的“Fourier-Bros-Iagolnitzer 变换”（例如，参见 Delort-Lebeau [1]）。但现在可以通过使用 Kashiwara [6] 的函子 TH（“温和同态”）及其微局部化（Andronikof [1] 的函子 Tμhom）来扩展第 4 节和第 5 节的技术，以处理具有适度增长（即分布）的全纯函数的边界值。

\appendix
\ctexset{
    section = {
        name = {附录~,},
        number = \Alph{section}
    }
}
\section{辛几何}
\subsection{概要}
我们在这里收集了整本书中使用的辛几何的基本工具。在第 1 节和第 2 节中，我们讨论有关辛向量空间和齐次辛流形的一些基本概念。所有结果都是众所周知的并且或多或少是初级的。因此，我们将省略一些证明，并参考 Arnold [2]、Duistermaat [1]、Abraham-Mardsen [1]，尤其是 Hörmander [4 Chapter XXI]。

在第 3 节中，我们介绍了拉格朗日平面三元组的惯性指数。我们的演示接近 Lion-Vergne [1] 的演示，为了读者的方便，我们给出了所有的证明。我们还收集了第 7 章中需要的它的一些属性作为练习。

\subsection{辛向量空间}
在本节和下一节中，我们将处理实际情况中的辛空间，但在复杂情况下理论是相同的。

令 $E$  为实数有限维向量空间。  $E$  上的辛形式  $\sigma$  是  $E$  上的非简并交替双线性形式。赋予辛形式 $\sigma$ 的向量空间 $E$ 称为辛向量空间。如果 $E$  是辛，则 $\dim E$  是偶数。

如果 $ (E_1,\sigma_1) $  和 $ (E_2,\sigma_2) $  是两个辛向量空间，则线性映射 $ u:E_1\to E_2 $  称为辛如果 $ u^*\sigma_2=\sigma_1 $  。如果 u 是辛的，则 u 必然是内射的。

其中一个用  $Sp(E)$  表示辛向量空间  $E$  的辛自同构群。它是  $E$  线性自同构群  $\mathrm{GL}(E)$  的闭子群。

示例 A.1.1。令 $V$  为实数有限维向量空间， $V^*$  为其对偶空间。通过在 $V \oplus V^*$ 中设置 $(x; \xi)$ 和 $(x'; \xi')$ ，空间 $E = V \oplus V^*$ 自然地被赋予了辛结构。

\[
\sigma((x;\xi),(x^{\prime};\xi^{\prime}))=\left\langle x^{\prime},\xi\right\rangle-\left\langle x,\xi^{\prime}\right\rangle.
\]

我们称这种形式为 E 上的自然辛形式。

如果  $E = V \oplus V^*$  和  $u$  是  $V$  的线性同构，则映射  $\begin{pmatrix} u & 0 \\ 0 & ^{t} u^{-1} \end{pmatrix}$  是  $E$  的辛自同构。

令 $(E,\sigma)$  为辛向量空间。由于 $\sigma$ 是非简并的，它通过以下公式定义了从 $E^*$ 到 $E$ 的线性同构 $H$ ：

\[
\langle\theta,v\rangle=\sigma(v,H(\theta))~,\qquad v\in E,\quad\theta\in E^{\star}~.
\]

这种同构 $H$ 称为哈密顿同构。如果  $\theta \in E^*$  有时会写  $H_\theta$  而不是  $H(\theta)$  并将  $H_\theta$  称为  $\theta$  的哈密顿量向量。

由于  $H$  是同构，因此  $E^*: (u,v) \mapsto \sigma(H_u, H_v)$  上的斜双线性形式是  $E^*$  上的辛形式。它称为泊松括号，表示为  $\{u, v\}$  。因此：

\[
\{u,v\}=\sigma(H_{u},H_{v})=\langle v,H_{u}\rangle.
\]

令  $ \rho $  为 E 的线性子空间。一组：

\[
\rho^{\perp}=\left\{x\in E;\sigma(x,\rho)=0\right\}\;.
\]

然后：

\[
\rho^{\perp\perp}=\rho\quad,\qquad(\rho_{1}+\rho_{2})^{\perp}=\rho_{1}^{\perp}\cap\rho_{2}^{\perp}\quad,\qquad(\rho_{1}\cap\rho_{2})^{\perp}=\rho_{1}^{\perp}+\rho_{2}^{\perp}\quad.
\]

空间  $ \rho^{\perp} $  称为  $ \rho $  的正交空间。

\statementlabel{定义 A.1.2}如果  $ \rho \subset \rho^\perp $  （分别为  $ \rho = \rho^\perp $  和  $ \rho \supset \rho^\perp $  ），则 E 的线性子空间  $ \rho $  被称为各向同性（分别为拉格朗日和合合）。

一些作者使用“共各向同性”而不是“渐合”。

请注意，如果  $\rho$  是各向同性的（分别是拉格朗日、对合），则  $\dim \rho \leq n$  （分别是  $n$  和  $\geq n$  ），其中  $n = \frac{1}{2} \dim E$  。

线（或超平面）始终是各向同性（或对合）。如果 $ \dim \rho = n $  且 $ \rho $  是各向同性或对合的，则 $ \rho $  是拉格朗日。

假设  $\rho$  是各向同性的。然后通过设置 $\sigma(\dot{x},\dot{y})=\sigma(x,y)$ ，空间 $\rho^{\perp}/\rho$ 自然地被赋予辛结构，其中 $\dot{x}$ （或 $\dot{y}$ ）是 $\rho^{\perp}/\rho$ 中 $x$ （或 $y$ ）的图像。 （我们仍然用  $\sigma$  表示  $\rho^{\perp}/\rho$  上的辛形式。）

如果 $\lambda$ 是 $E$ 一组的线性子空间：

\[
\lambda^{\rho}=((\lambda\cap\rho^{\perp})+\rho)/\rho.
\]

特别是 $ E^\rho = \rho^\perp/\rho $  。那么很容易检查到：

\[
(\lambda^{\perp})^{\rho}=(\lambda^{\rho})^{\perp}.
\]

特别是，如果  $\lambda$  是  $E$  中的拉格朗日，则  $\lambda^{\rho}$  是  $E^{\rho}$  中的拉格朗日。

相反， $i$  表示嵌入  $\rho^\perp \hookrightarrow E$  ， $j$  表示投影  $\rho^\perp \to \rho^\perp/\rho$  。那么如果  $\mu$  是  $\rho^\perp/\rho$  的拉格朗日子空间，则  $ij^{-1}(\mu)$  是  $E$  的拉格朗日子空间，包含  $\rho$  。

例 A.1.1 中描述的辛空间并没有那么特殊。事实上，我们有

\statementlabel{命题 A.1.3}令  $\lambda_0$  为  $E$  的拉格朗日子空间。那么存在  $E$  的拉格朗日子空间  $\lambda_1$  使得  $E = \lambda_0 \oplus \lambda_1$  。此外，对于这样的拉格朗日空间 $\lambda_1$ ，由 $u(x, y) = x - H(y)$ 定义的映射 $u: \lambda_0 \oplus \lambda_0^* \to E$ 是辛同构。这里  $H(y)$  由  $\lambda_0^* \simeq (E/\lambda_1)^* \to E^* \xrightarrow{H} E$  给出。

\statementlabel{证明}(i) 令  $\rho$  为具有  $\rho \cap \lambda_0 = \{0\}$  的各向同性空间。如果  $\rho \neq \rho^\perp$  ，则  $\rho^\perp \notin \lambda_0 + \rho$  ，否则  $\rho \supset (\rho^\perp \cap \lambda_0)$  ，因此  $\rho^\perp \cap \lambda_0 = \{0\}$  ，这与  $\dim \rho^\perp > n$  相矛盾。选择 $e \in \rho^\perp \setminus (\lambda_0 + \rho)$ 。那么  $\rho + \mathbb{R}e$  是各向同性的， $(\rho + \mathbb{R}e) \cap \lambda_0 = \{0\}$  。然后我们通过对  $\dim \rho$  进行归纳论证得到  $\lambda_1$  。

(ii) 令 x 和  $ x' $  属于  $ \lambda_{0} $  且 y 和  $ y' $  属于  $ \lambda_{0}^{*} $  。然后：

\[
\sigma(x-H(y),x^{\prime}-H(y^{\prime}))=-\sigma(x,H(y^{\prime}))+\sigma(x^{\prime},H(y))=-\left\langle y^{\prime},x\right\rangle+\left\langle y,x^{\prime}\right\rangle.
\]

因此映射 $u$  是辛的。由于  $\dim(\lambda_0 \oplus \lambda_0^*) = \dim E$  ，  $u$  是同构。 □

令 $(E,\sigma)$  为辛向量空间。我们用  $E^{a}$  表示具有辛形式  $-\sigma$  的空间  $E$  ，即  $E^{a}=(E,-\sigma)$  。

对于两个辛向量空间  $ (E_1, \sigma_1) $  和  $ (E_2, \sigma_2) $  ，  $ E_1 \oplus E_2 $  也具有  $ \sigma_1 \oplus \sigma_2 $  的辛空间结构。

令  $E_{i}$  、  $(i=1,2,3)$  为三个辛向量空间，并用  $p_{ij}$  表示  $E_{1}\times E_{2}\times E_{3}$  上定义的  $(i,j)$  次投影（例如  $p_{13}$  是  $E_{1}\times E_{3}$  上的投影）。

\statementlabel{命题 A.1.4}令  $\lambda$  和  $\mu$  分别为  $E_1 \oplus E_2^a$  和  $E_2 \oplus E_3^a$  的两个拉格朗日子空间。设置 $\lambda \circ \mu = p_{13}(p_{12}^{-1}\lambda \cap p_{23}^{-1}\mu)$ 。那么  $\lambda \circ \mu$  是  $E_1 \oplus E_3^a$  的拉格朗日子空间。

\statementlabel{证明}$E_2^a \oplus E_2$  的对角线 $\mathcal{A}$  是拉格朗日子空间。因此  $\rho = \{0\} \times \mathcal{A} \times \{0\}$  是  $E_1 \oplus E_2^a \oplus E_2 \oplus E_3$  的各向同性子空间。从  $E_1 \oplus E_3 = (E_1 \oplus E_2^a \oplus E_2 \oplus E_3)^\rho$  和  $\lambda \circ \mu = (\lambda \oplus \mu)^\rho$  开始，我们得到了结果。   $\square$ 

现在我们将研究辛基。令  $(E,\sigma)$  为辛向量空间，维数为  $2n$  。

如果表示  $ (e_1^*, \ldots, e_n^*; f_1^*, \ldots, f_n^*) $  上  $ E^* $  的对偶基，则基  $ (e_1, \ldots, e_n; f_1, \ldots, f_n) $  称为辛基，我们有：

\[
\sigma=\sum_{j=1}^{n}f_{j}{}^{*}\wedge e_{j}^{*}\enspace.
\]

当然，(A.1.7)等价于：

\[
\left\{\begin{aligned}{\sigma(e_{j},e_{k})}&{{}=\sigma(f_{j},f_{k})=0~,}\\ {\sigma(e_{j},f_{k})}&{{}=-\sigma(f_{k},e_{j})=-\delta_{j,k}\qquad1\leqslant j,k\leqslant n~,}\\ \end{aligned}\right.
\]

其中  $ \delta_{jk} $  表示克罗内克符号（如果 j=k，则为  $ \delta_{jk}=1 $ ，否则为 0）。

对于这样的辛基，(A.1.2)中定义的辛同构 $H$ 满足：

\[
H(e_{j}^{*})=-f_{j}~,\qquad H(f_{j}^{*})=e_{j}~,\qquad1\leqslant j\leqslant n~.
\]

令 $J$  和 $K$  为 $\{1, \ldots, n\}$  的两个子集，并令 $\left((e_j)_{j \in J}, (f_k)_{k \in K}\right)$  为满足关系式(A.1.8) 的线性独立族。那么很容易证明这个族可以完备化为辛基。特别是，如果  $\rho$  是各向同性（分别为拉格朗日、分别为对合）子空间，则存在一个辛基，使得对于某些  $j$  和  $k$  ，  $\rho$  由  $(e_1, \ldots, e_j)$  （分别为  $(e_1, \ldots, e_n)$  和  $(e_1, \ldots, e_n, f_1, \ldots, f_k)$  ）生成。

用 $G(E,n)$  表示 $n$  的 $E$  维线性子空间的格拉斯曼流形（回想一下 $\dim E = 2n$  ）。这是一个紧凑流形（参见例如 Griffith-Harris [1]）。其中 $\mathcal{A}(E)$  表示由拉格朗日子空间组成的 $G(E,n)$  的子集。这是一个闭合（光滑）子流形，称为拉格朗日格拉斯曼流形。

让 $ \mu \in \mathcal{A}(E) $  。我们设定：

\[
\varLambda_{\mu}(E)=\left\{\lambda\in\varLambda(E);\lambda\cap\mu=\{0\}\right\}.
\]

假设  $E$  具有辛基，并让  $(x; \xi)$  表示相关联的线性坐标（即每个  $p \in E$  写作  $p = \sum_{j=1}^{n} (x_j e_j + \xi_j f_j)$  ）。对于任何  $\lambda \in G(E, n)$  ，都存在  $n \times n$  矩阵  $A$  和  $B$  使得：

(A.1.11) 矩阵  $ (A,B) $  的秩为 n， $ \lambda=\{(x;\xi);B\xi=Ax\} $  。

那么 $\lambda \in A(E)$  当且仅当 $A^t B$  是对称的。

请注意，在这种情况下， $ B\xi = Ax $  iff  $ x = 'Bz $  ， $ \xi = 'Az $  对于某些 $ z \in \mathbb{R}^n $  。

如果  $\mu$  是拉格朗日子空间  $\{x = 0\}$  ，则  $\lambda \in \mathcal{A}_\mu(E)$  如果  $\lambda = \{(x; \xi); \xi = Ax\}$  对于某些对称矩阵  $A$  。因此  $\mathcal{A}_\mu(E)$  在  $A(E)$  中是开放且密集的，并且与  $\mathbb{R}^{n(n+1)/2}$  同构。

我们有时会说，如果存在  $\mathcal{A}(E)$  的开稠子集  $\Omega$ ，使得属性“ $P$ ”适用于  $\lambda \in \Omega$ ，则属性“ $P$ ”适用于泛型  $\lambda$ （ $\mathcal{A}(E)$  中的  $\lambda$ ）。

我们还会遇到以下情况：  $\lambda$  是  $E$  中的拉格朗日并包含一行  $\rho$  。我们正在寻找  $\mu \in \mathcal{A}(E)$  和  $\mu \cap \lambda = \rho$  。考虑映射  $i: \rho^\perp \to E$  和  $j: \rho^\perp \to \rho^\perp/\rho$  。然后选择  $\mu' \in \mathcal{A}(E^\rho)$  和  $\mu' \cap \lambda^\rho = 0$  ，并设置  $\mu = i(j^{-1}(\mu'))$  就足够了。通过滥用语言，我们可以说，对于泛型  $\mu$  和  $\mu \supset \rho$  ，我们有  $\lambda \cap \mu = \rho$  。

\subsection{齐性辛流形}
这里考虑的所有流形和流形态射都是实数，属于  $ C^{\infty} $  类或实解析类。除非另有说明，流形上的实函数应该属于  $ C^{\infty} $  类或实解析函数。然而，我们将陈述的大多数结果在对复杂解析流形进行适当修改后仍然成立。

令 $X$  为流形。我们用 $\tau: TX \to X$  表示其切丛，用 $\pi: T^*X \to X$  表示其余切丛。我们用  $T X$  和  $T^*X$  表示束  $TX$  和  $T^*X$  ，并删除了零部分，并且我们用  $t$  和  $\pi$  表示映射  $\tau$  和  $\pi$  分别限制为  $T X$  和  $T^*X$  。让我们回想一下，如果  $M$  是  $X$  的子流形，则  $X$  中的法束  $T_M X$  和共形束  $T_M^*X$  到  $M$  由  $M$  上向量束的精确序列定义：

\[
\left\{\begin{aligned}{0}&{{}\to T M\to M\times\mathop{{T X}}_{X}\to T_{M}X\to0,}\\ {0}&{{}\to T_{M}^{*}X\to M\times\mathop{{T}^{*}}_{X}X\to T^{*}M\to0.}\\ \end{aligned}\right.
\]

令 $f: Y \to X$  为流形的态射。与  $f$  关联的是态射：

\[
\left\{\begin{array}{ccccc}TY&\xrightarrow{f'}&Y\times_X TX&\xrightarrow{f_t}&TX\\ T^*Y&\xleftarrow{{}^t f'}&Y\times_X T^*X&\xrightarrow{{}^t f_t}&T^*X\end{array}\right.,
\]

特别是，如果考虑投影 $ \pi: T^*X \to X $ ，我们会得到映射 $ t\pi' $ ： $ T^*X \times_X T^*X \to T^*T^*X $ 。

如果我们将此映射限制为  $ T^*X \times_X T^*X $  的对角线，我们会得到一个映射  $ T^*X \to T^*T^*X $  ，它是丛  $ T^*T^*X \to T^*X $  的一部分，即 1 次微分形式（有人说：1-形式）。  $ T^*X $  上的这个 1-形式称为规范 1-形式，并表示为  $ \alpha_v $  ，或者在不存在混淆风险的情况下简称为  $ \alpha $  。

令  $(x_1, \ldots, x_n)$  为  $X$  上的局部坐标系。那么  $x_j$  是在某个开放子集  $U$  上定义的实函数，满足  $U$  上的  $\mathrm{d}x_1 \wedge \cdots \wedge \mathrm{d}x_n \neq 0$  。在每个  $x \in U$  处，  $(\mathrm{d}x_1, \ldots, \mathrm{d}x_n)$  定义了向量空间  $T_x^* X$  的基础，并且向量  $\xi \in T_x^* X$  唯一地写为  $\xi = \sum_{j=1}^{n} \xi_j \, \mathrm{d}x_j$  。系统  $(x_1, \ldots, x_n; \xi_1, \ldots, \xi_n)$  称为与坐标系  $(x_1, \ldots, x_n)$  关联的  $T^* X$  上的坐标系。很容易检查到规范的 1-形式  $\alpha_X$  只是形式  $\sum_{j=1}^{n} \xi_j \, \mathrm{d}x_j$  。让 $\sigma = \mathrm{d}\alpha$  。因此  $\sigma = \sum_{j=1}^{n} \mathrm{d}\xi_j \wedge \mathrm{d}x_j$  是  $T^* X$  上的辛形式（即：在每个  $p \in T^* X$  处， $\sigma$  在向量空间  $T_p T^* X$  上引入辛结构）。换句话说，流形  $T^* X$  自然地被  $\mathrm{d}\alpha$  赋予了辛结构。然后我们可以将§1中引入的一些概念扩展到 $T^* X$ 。

如果在每个  $p\in V$  处，切线空间  $T_pV$  具有  $T_pT^*X$  中的相应属性，则  $T^*X$  的子流形  $V$  称为各向同性（分别为拉格朗日和合合）。如果  $f$  是在  $T^*X$  的某个开子集  $U$  上定义的实数函数，则

 $ f $  的哈密顿量场 $ H_f $  是 $ U $  上的矢量场， $ df $  是哈密顿量同构 $ H: T^*T^*X \simeq TT^*X $  的图像。

两个函数 f 和 g 的泊松括号定义为：

\[
\{f,g\}=H_{f}(g)=\mathrm{d}\alpha(H_{f},H_{\theta})\enspace.
\]

检查一下关系：

\[
\left\{\begin{aligned}{}&{{}\{f,g\}=-\{g,f\},}\\ {}&{{}\{f,h g\}=h\{f,g\}+g\{f,h\},}\\ {}&{{}\{\{f,g\},h\}+\{\{g,h\},f\}+\{\{h,f\},g\}=0.}\\ \end{aligned}\right.
\]

特别是 $ [H_{f}, H_{g}] = H_{\{f, g\}} $ ，其中 $ [u, v] = uv - vu $  是向量场u 和v 的交换子。

如果  $(x_1, \ldots, x_n)$  是  $X$  上的局部坐标系， $(x; \xi)$  上  $T^*X$  上的关联坐标和  $f$  是  $U \subset T^*X$  上的实函数，我们有：

\[
H_{f}=\sum_{j=1}^{n}\left(\frac{\partial f}{\partial\xi_{j}}\frac{\partial}{\partial x_{j}}-\frac{\partial f}{\partial x_{j}}\frac{\partial}{\partial\xi_{j}}\right).
\]

 $T^*X$  的子流形  $V$  是对合的，当且仅当泊松括号  $\{f, g\}$  在  $V$  上消失，对于任何在  $V$  上消失的函数  $f$  和  $g$  。事实上，向量丛  $(TV)^\perp$  是由向量场  $H_f$  和  $f|_{V}=0$  生成的。因此，对于任何  $f$  、  $g$  和  $f|_{V}=0$  、  $g|_{V}=0$  ，  $(TV)^\perp\subset TV$  等价于  $H_f(g)=0$  。此外，根据（A.2.4），我们发现如果 $V$ 是对合的，则 $TV$ 的子丛 $(TV)^\perp$ 满足Frobenius可积条件（即： $(TV)^\perp$ 的节束在括号 $\left[\cdot,\cdot\right]$ 下闭合）。根据 Frobenius 定理（参见 Hörmander [4，附录 C]），对合流形  $V$  承认一个叶状结构，并且该叶状结构的叶子被称为  $V$  的双特征叶子。请注意，叶子的尺寸是  $V$  的余尺寸。特别是，如果  $V$  是拉格朗日，则叶子在  $V$  中是开放的。

示例 A.2.1。令  $Z$  为  $X$  的子流形。流形 $Z \times_X T^*X$  是对合流形，流形 $T_Z^*X$  是拉格朗日流形。请注意  $Z = X$  的极端情况：我们得到拉格朗日流形  $T_X^*X$  ，  $T^*X$  的零部分。

 $ T^*X $  上的 1-形式  $ \alpha $  引发了比辛流形更丰富的结构，我们将在  $ T^*X $  上描述这种齐次辛结构。

令 $H(\alpha)$  为 $\alpha$  的哈密顿同构的图像。如果我们选择了上面的坐标 $(x;\xi)$ ，那么：

\[
\alpha=\sum_{j}\xi_{j}\mathrm{d}x_{j},\qquad H(\alpha)=-\sum_{j}\xi_{j}\frac{\partial}{\partial\xi_{j}}.
\]

因此  $-H(\alpha)$  只是向量束  $T^*X$  上的径向向量场（即  $\mathbb{R}^+$  对  $T^*X$  作用的无穷小生成元）。该矢量场也称为欧拉矢量场。

如果  $T^*X$  的子集  $S$  在  $\mathbb{R}^+$  的作用下保持不变（或局部不变），则我们说  $T^*X$  的子集  $S$  是圆锥曲线（或局部圆锥曲线）。因此，如果 $S$  与 $\mathbb{R}^+$  的任何轨道的交点在该轨道上是开放的，则 $S$  是局部圆锥曲线。如果  $f$  对于某些  $k \in \mathbb{C}$  满足微分方程  $H(\alpha)f = kf$  ，则在  $T^*X$  的开子集  $U$  上定义的函数  $f$  被称为齐次函数。请注意，子流形  $V$  是局部圆锥曲线，当且仅当  $H(\alpha)$  与  $V$  相切，或者等效地当且仅当  $V$  是由齐次方程局部定义的。

局部圆锥子流形  $V$  是各向同性当且仅当  $\alpha|_{\nu} \equiv 0$  ，因为  $\langle \alpha, v \rangle = \mathrm{d}\alpha(v, H(\alpha))$  ，对于  $v \in TT^*X$  。

有人说，如果  $\alpha|_{V}$  处处不为零，则局部圆锥对合子流形  $V$  是正则的。这相当于齐次函数  $f_{1}, \ldots, f_{r}$  在  $V$  和  $r = \mathrm{codim}V$  上消失的局部存在，这样：

\[
\left\{\begin{aligned}{}&{{}\left\{f_{i},f_{j}\right\}=0\quad\mathrm{o n~}V\mathrm{f o r~a n y}i,j\in\left\{1,\ldots,r\right\},}\\ {}&{{}\mathsf{d}f_{1}\wedge\cdots\wedge\mathsf{d}f_{r}\wedge\alpha\neq0\quad\mathrm{o n~}V.}\\ \end{aligned}\right.
\]

这又相当于说欧拉向量场在任何点都不与任何双特征叶子相切。

示例 A.2.2。令  $Z$  为  $X$  的子流形。那么  $Z \times_X T^*X$  是  $T_Z^*X$  之外的常规对合。

公约 A.2.3。在本附录中，除非另有说明， $ T^{*}X $  的所有子流形都是局部圆锥曲线。

令 $\rho(p)$  表示由 $p$  处的欧拉向量场生成的 $T_pT^*X$  的线性子空间。如果  $p\in T_X^*X$  、  $\rho(p)=\{0\}$  ，否则  $\rho(p)$  是一行。

如果  $V$  是一个（局部圆锥曲线）子流形，则在每个  $p \in V$  处，  $T_p V$  包含  $\rho(p)$  。让 $p \in T^* X$  。一套：

\[
\lambda_{0}(p)=T_{p}\pi^{-1}\pi(p).
\]

这是  $T_{p}T^{*}X$  的拉格朗日线性子空间。

令 $A$  为拉格朗日子流形。根据定义，投影  $\pi|_{\mathcal{A}}: A \to X$  的 corank 是空间  $T_p A \cap \lambda_0(p)$  的维度。在  $T^* X$  上，该软木塞至少为 1，因为  $T_p A$  和  $\lambda_0(p)$  都包含  $\rho(p)$  。如果这个 corank 是常数，比如  $d$  ，那么局部在  $A$  上，  $\pi(A)$  是 codimension  $d$  和  $A = T_M^* X$  的  $X$  的平滑子流形  $M$  。特别是，如果此软木塞在某个  $p$  处是一个，则  $A$  是  $p$  邻域中超曲面的共形丛。

现在让  $X$  和  $Y$  为相同维度的两个流形，  $U_X$  （分别为  $U_Y$  ）为  $T^*X$  （分别为  $T^*Y$  ）的开子集。令  $\chi$  成为从  $U_X$  到  $U_Y$  的微分同胚。如果  $\chi^*(\mathrm{d}\alpha_Y) = \mathrm{d}\alpha_X$  ，则表示  $\chi$  是辛同构。此外，如果  $\chi$  是齐次的（即：  $\chi$  通过  $\mathbb{R}^+$  的作用进行交换），那么  $\chi^*(\alpha_Y) = \alpha_X$  和我们应该说  $\chi$  是接触变换，尽管这并不真正正确，因为接触结构是在商空间  $T^*X/\mathbb{R}^+$  上获得的结构。

令  $\chi$  为齐次微分同胚  $U_{X} \simeq U_{Y}$  ，  $A_{x}$  其图在  $U_{X} \times U_{Y}$  中。  $U_{Y}$  上 1 型  $\beta$  的逆像  $\chi^*(\beta)$  的特征在于以下条件

 $ (\chi^*(\beta) - \beta)|_{\mathcal{A}_x} \equiv 0 $ 。让 $ \mathcal{A}_x^a $  表示 $ \mathcal{A}_x $  通过 $ T^*Y $  上的对映图的图像。然后  $ \chi^*(\alpha_Y) = \alpha_X $  iff  $ (\alpha_X + \alpha_Y)|_{\mathcal{A}_x^a} \equiv 0 $  ，即 iff  $ \mathcal{A}_x^a $  是各向同性的，因此 iff 它是拉格朗日。换句话说，  $ \chi $  是接触变换当且仅当  $ \mathcal{A}_x^a $  是  $ T^*(X \times Y) $  的（局部圆锥）拉格朗日子流形。

我们将与接触变换图 $ \chi $ 相关的拉格朗日流形称为 $ A_{\chi}^{a} $ 。

令 $(y)$  为 $Y$  上的局部坐标系，并令 $(y;\eta)$  表示 $T^*Y$  上的关联坐标。然后齐次映射 $\chi: U_X \to U_Y$ 由两组函数定义， $f_j$ 同次度 $0$ ， $g_k$ 同次度 $1$ ， $(1 \leq j, k \leq n)$ ，以及 $y_j = f_j$ ， $\eta_k = g_k$ 。映射 $\chi$  是一个接触变换，当且仅当 $A_x^a$  是对合的，即当且仅当：

\[
\{f_{j},f_{k}\}=0~,\qquad\{g_{j},g_{k}\}=0~,\qquad\{f_{j},g_{k}\}=-\delta_{j,k}~.
\]

示例 A.2.4。令  $(x; \xi)$  表示  $T^*\mathbb{R}^n$  上的坐标，并令  $\varphi(\xi)$  为在  $((\mathbb{R}^n)^*$  的某个开子集  $U$  上定义的一次齐次函数（例如： $\mathbb{R}^n \setminus \{0\}$  上的  $\varphi(\xi) = (\sum_j \xi_j^2)^{1/2}$  ）。那么映射 $\chi : (x; \xi) \mapsto (x + \varphi'(\xi); \xi)$ 就是一个接触变换。

下一个结果是构建接触转换的有用工具。

\statementlabel{命题 A.2.5}令  $V$  为  $T^*X$  、  $p\in V\cap\dot{T}^*X$  的正则对合子流形。令  $\lambda$  为  $T_pT^*X$  的拉格朗日线性子空间，使得  $\rho(p)\subset\lambda\subset T_pV$  。 （回想一下  $\rho(p)$  是由欧拉向量场生成的线。）然后存在拉格朗日流形  $A\subset T^*X$  使得  $A\subset V$  和  $T_pA=\lambda$  。

\statementlabel{证明}让 $n = \dim X, r = \operatorname{codim} V$  。令 $(f_1, \ldots, f_r)$  为在 $V$  上消失且满足(A.2.7) 的齐次函数系统。如果  $r = n - 1$  我们设置  $e = \alpha(p)$  。否则，我们选择 $v \in \lambda^\perp \setminus (T_p V + \mathbb{R} H(\alpha))$ 并设置 $e = H^{-1}(v)|_V$ 。根据经典微分方程理论，我们可以在 $V$ 上找到一个函数 $g$ ，使得：

\[
\left\{\begin{aligned}{}&{{}H(\alpha)|_{V}(g)=0~,\qquad H_{f_{j}}|_{V}(g)=0~,\qquad(j\leqslant r)}\\ {}&{{}\mathbf{d}g(p)=e~,\qquad g(p)=0~,\qquad}\\ \end{aligned}\right.
\]

因为 e 与经过 p 的双特征叶子不相切。

设置 $V_1 = \{q \in V; g(q) = 0\}$ 。那么 $V_1$  是一个圆锥流形，如果 $r = n - 1$  满足要求的性质，否则 $V_1$  是正则对合的。在这种情况下，我们通过  $r$  归纳进行论证。   $\square$ 

令  $X$  、  $Y$  、  $Z$  为三个流形。其中  $p_1$  和  $p_2$  表示在  $T^*X \times T^*Y$  或  $T^*Y \times T^*Z$  上定义的第一和第二投影， $p_{ij}$  表示在  $T^*X \times T^*Y \times T^*Z$  上定义的第  $(i, j)$  投影。我们设置 $p_2^a = a \circ p_2$ ，其中“ $a$ ”是对映图。令 $A_1 \subset T^*(X \times Y)$  、 $A_2 \subset T^*(Y \times Z)$  为两个拉格朗日流形。让 $(p_X, p_Y^a) \in A_1$ ， $(p_Y, p_Z^a) \in A_2$ ，并假设：

\[
\begin{cases}{{t h e~m a p s~}p_{2}^{a}|_{\varLambda_{1}}\colon\varLambda_{1}\to T^{*}Y{~a n d~}p_{1}|_{\varLambda_{2}}\colon\varLambda_{2}\to T^{*}Y}\\ {{a r e~t r a n s v e r s a l~a t~p_{Y}~.}}\\ \end{cases}
\]

然后用 $A_1 \cap U$ 和 $A_2 \cap V$ 替换 $A_1$ 和 $A_2$ ，其中 $U$ 和 $V$ 分别是 $(p_X, p_Y^a)$ 和 $(p_Y, p_Z^a)$ 的足够小的开邻域，映射 $p_{13}$ 诱导 $A_1 \times_{T^*Y} A_2$ 与拉格朗日流形 $A$ 的同构 $T^*(X \times Z)$  ，一组：

\[
\varLambda=\varLambda_{1}\circ\varLambda_{2}~.
\]

（参见引理 7.4.4 和定义 7.4.5。）

\statementlabel{命题 A.2.6}令  $ A_1 $  为  $ T^*(X \times Y) $  、  $ (p_X, p_Y^a) \in A_1 $  和  $ p_Y \notin T_Y^*Y $  的拉格朗日子流形。假设映射 $ p_1|_{\mathcal{A}_1} : \mathcal{A}_1 \to T^*X $  是光滑的。然后存在与  $ Y $  具有相同维度的流形  $ Z $  ，以及在  $ (p_Y, p_Z^a) $  邻域中定义的拉格朗日流形  $ A_2 \subset T^*(Y \times Z) $  ，使得：

(i)  $ A_{2} = T_{S}^{*}(Y \times Z) $  ，其中 S 是  $ Y \times Z $  的超曲面，

(ii)  $A_{2}$  与接触变换图相关联（即： $p_{1}|_{A_{2}}$  和  $p_{2}^{a}|_{A_{2}}$  是局部同构），

(iii)  $ A_1 \circ A_2 = T_S^*(X \times Z) $  ，其中  $ S' $  是  $ X \times Z $  的超曲面。

\statementlabel{证明}设 $E_X = T_{p_X} T^* X$  、 $E_Y = T_{p_Y} T^* Y$  并令 $E_Y^a$  表示具有相反辛结构的空间 $E_Y$ 。然后  $(p_1, p_2^a)$  定义了一个辛同构  $T_{(p_X, p_Y^a)} T^*(X \times Y) \simeq E_X \times E_Y^a$  ，我们将在  $T_{(p_X, p_Y^a)} T^*(X \times Y)$  中识别拉格朗日空间及其在  $E_X \times E_Y^a$  中的图像。

令 $\rho_X$  表示 $E_X$  中 $p_X$  处的欧拉向量场生成的线性空间，并类似地定义 $E_Y$  中的 $\rho_Y$ 、 $E_X \times E_Y^a$  中的 $\rho_{XY}$ 、 $E_Y \times E_Y^a$  中的 $\rho_{YY}$  。设定：

\[
\lambda_{1}=T_{(p_{X},p_{Y}^{\tt a})}\varLambda_{1},\qquad\lambda_{0X}=T_{p_{X}}\pi^{-1}\pi(p_{X}),\qquad\lambda_{0Y}=T_{p_{Y}}\pi^{-1}\pi(p_{Y}),
\]

并分别将  $\lambda_{1}$  和  $\lambda_{0x}$  识别为  $E_{x} \times E_{y}^{a}$  和  $E_{x}$  的拉格朗日子空间。

通过假设  $ p_1|_{\lambda_1} : \lambda_1 \to E_X $  是满射的，我们得到  $ p_2|_{\lambda_1} : \lambda_1 \to E_Y $  是单射的（参见练习 A.4）。此外 $ p_2^{-1}(\rho_Y) \cap \lambda_1 = \rho_{XY} $ 。由于  $ p_2(p_1^{-1}(\lambda_{0X}) \cap \lambda_1) $  是  $ E_Y $  中的拉格朗日（命题 A.1.4），因此我们有一个通用的拉格朗日空间  $ \lambda \subset E_Y $  和  $ \rho_Y \subset \lambda $  ：

\[
\lambda\cap p_{2}(p_{1}^{-1}(\lambda_{0X})\cap\lambda_{1})=\rho_{Y}.
\]

这意味着  $ \rho_{XY} = p_2^{-1}(\lambda) \cap p_1^{-1}(\lambda_{0X}) \cap \lambda_1 $  ，因此：

 $ \rho_{XY} = (\lambda_{0X} \times \lambda) \cap \lambda_1 $  ，对于通用拉格朗日空间

\[
\lambda\subset E_{Y}^{a}\mathrm{s u c h~t h a t}\rho_{Y}\subset\lambda.
\]

然后对于通用拉格朗日空间  $ \mu \subset E_Y \times E_Y^a $  和  $ \rho_{YY} \subset \mu $  ，我们有：

\[
p_{1}|_{\mu}\colon\mu\to E_{Y}\mathrm{a n d}p_{2}|_{\mu}\colon\mu\to E_{Y}^{a}\mathrm{a r e i s o m o r p h i s m s},
\]

\[
(\lambda_{0Y}\times\lambda_{0Y})\cap\mu=\rho_{Y Y}.
\]

由于 $ \lambda = p_1(\mu \cap p_2^{-1}(\lambda_{0Y})) $  是通用的，我们可以进一步假设 $ \lambda $  满足条件(A.2.12)。然后我们有：

(A.2.15)

\[
\left(\lambda_{1}\circ\mu\right)\cap\left(\lambda_{0X}\times\lambda_{0Y}\right)=\rho_{X Y}.
\]

现在将 $Y, p_Y$ 和 $\lambda_{0Y}$ 分别作为 $Z, p_Z$ 和 $\lambda_{0Z}$ 。根据命题 A.2.5，我们可以找到一个拉格朗日流形  $A_2 \subset T^*(Y \times Z)$  使得  $T_{(p_Y, p_Z^2)} A_2 = \mu$  。那么  $A_2$  将满足 (A.2.13)、(A.2.14) 和 (A.2.15) 所要求的所有条件。

\statementlabel{推论 A.2.7}令 $ A \subset T^*X $  为拉格朗日流形， $ p \in A $  。然后存在在 $ p $ 的邻域中定义的接触变换 $ \chi $ ，使得 $ \chi(A) $ 是超曲面的共形丛，而且与 $ \chi $ 的图关联的拉格朗日流形是超曲面的共形丛。

\statementlabel{证明}这只不过是  $ X = \{\mathrm{pt}\} $  时的命题 A.2.5。 □

\statementlabel{推论 A.2.8}令  $\varLambda \subset \tilde{T}^*(X \times Y)$  为与接触变换相关的拉格朗日流形（即： $p_1|_{\varLambda}$  和  $p_2^a|_{\varLambda}$  是局部同构）。然后，在  $\varLambda$  上的局部，存在与  $X$  相同维度的流形  $Z$  和两个拉格朗日流形  $\varLambda_1 \subset T^*(X \times Z)$  、  $\varLambda_2 \subset T^*(Z \times Y)$  ，使得  $\varLambda_1$  和  $\varLambda_2$  与接触变换相关联， $\varLambda_1$  和  $\varLambda_2$  是分别是 $X \times Z$  和 $Z \times Y$  以及 $\varLambda = \varLambda_1 \circ \varLambda_2$  。

\statementlabel{证明}根据命题 A.2.6，存在接触变换 $ \chi $ ，使得如果 $ A_2 $  是与 $ \chi $  图关联的拉格朗日流形，则 $ A_2 $  和 $ A \circ A_2 $  是超曲面的共形丛。然后  $ A = (A \circ A_2) \circ A_3 $  ，其中  $ A_3 $  是与  $ \chi^{-1} $  关联的拉格朗日流形。   $ \square $ 

为了结束本节，让我们回顾一下以下众所周知的结果。

\statementlabel{命题 A.2.9}令  $A$  为  $\tilde{T}^*X$  、  $p \in A$  的圆锥子流形。假设 $A$  是各向同性的（分别是拉格朗日，和正则对合）。然后存在在 $p$ 邻域中定义的接触变换 $\chi$ ，使得 $\chi(p) = (0; \mathrm{d}x_n) \in T^*\mathbb{R}^n$ 和 $A = \{(x; \xi); x = 0, \xi_1 = \cdots = \xi_r = 0 (r < n)\}$ （分别为 $A = \{(x; \xi); x = 0\}$ ，分别为 $A = \{(x; \xi); \xi_1 = \cdots = \xi_p = 0, (p < n)\}$ ）。

证明留作练习。

\subsection{惯性指数}
令 $(E,\sigma)$  为维数 $2n$  的实辛向量空间，并令 $\lambda_{1},\lambda_{2},\lambda_{3}$  为三个拉格朗日子空间。 （本节中，除非另有说明，子空间均指线性子空间。）

\statementlabel{定义 A.3.1}三元组  $ (\lambda_1, \lambda_2, \lambda_3) $  的惯性指数，表示为  $ \tau_E(\lambda_1, \lambda_2, \lambda_3) $  （或简称  $ \tau(\lambda_1, \lambda_2, \lambda_3) $  ）是在 3n 上定义的二次形式 q 的签名

维向量空间 $ \lambda_{1}\oplus\lambda_{2}\oplus\lambda_{3} $ ：

\[
q(x_{1},x_{2},x_{3})=\sigma(x_{1},x_{2})+\sigma(x_{2},x_{3})+\sigma(x_{3},x_{1}).
\]

该指数有时也称为“马斯洛夫指数”。

在  $\lambda_1 \oplus \lambda_2 \oplus \lambda_3$  的合适基础中，可以用一个对角矩阵来表示  $q$ ，该对角矩阵的对角元素由  $p_+$  - +1 的个数、 $p_-$  -1 的个数和  $3n - p_+ - p_-$  - 0 的个数组成。然后  $q$  的签名（表示为  $\mathrm{sgn}(q)$  ）等于  $p_+ - p_-$  。

索引 $ \tau $  具有以下属性。

\statementlabel{定理 A.3.2}(i)  $\tau(\lambda_{1},\lambda_{2},\lambda_{3})$  相对于  $\lambda_{j}$  的排列是交替的，即：

\[
\tau(\lambda_{1},\lambda_{2},\lambda_{3})=-\tau(\lambda_{2},\lambda_{1},\lambda_{3})=-\tau(\lambda_{1},\lambda_{3},\lambda_{2})~.
\]

(ii)  $ \tau $  满足“共循环条件”：对于拉格朗日空间的任何四元组  $ \lambda_{1}, \lambda_{2}, \lambda_{3}, \lambda_{4} $ ，我们有：

\[
\tau(\lambda_{1},\lambda_{2},\lambda_{3})-\tau(\lambda_{1},\lambda_{2},\lambda_{4})+\tau(\lambda_{1},\lambda_{3},\lambda_{4})-\tau(\lambda_{2},\lambda_{3},\lambda_{4})=0.
\]

(iii) 当拉格朗日子空间 $\lambda_{1},\lambda_{2},\lambda_{3}$ 在拉格朗日格拉斯曼 $\Lambda(E)$ 中以 $\dim(\lambda_{1}\cap\lambda_{2})$ 、 $\dim(\lambda_{2}\cap\lambda_{3})$ 和 $\dim(\lambda_{3}\cap\lambda_{1})$ 保持恒定的方式连续移动时， $\tau(\lambda_{1},\lambda_{2},\lambda_{3})$ 保持恒定。

\[
\begin{array}{r}{\mathrm{(i v)}\tau(\lambda_{1},\lambda_{2},\lambda_{3})\equiv n+\mathrm{d i m}(\lambda_{1}\cap\lambda_{2})+\mathrm{d i m}(\lambda_{2}\cap\lambda_{3})+\mathrm{d i m}(\lambda_{3}\cap\lambda_{1})\bmod2\mathbb{Z}.}\end{array}
\]

(v) 如果 $\rho$  是包含在 $(\lambda_1 \cap \lambda_2) + (\lambda_2 \cap \lambda_3) + (\lambda_3 \cap \lambda_1)$  中的各向同性空间，则：

\[
\tau_{E}(\lambda_{1},\lambda_{2},\lambda_{3})=\tau_{E^{\rho}}(\lambda_{1}^{\rho},\lambda_{2}^{\rho},\lambda_{3}^{\rho}).
\]

(vi) 令 $ E^{a} = (E, -\sigma) $  为具有辛形式 $ -\sigma $  的向量空间E。然后：

\[
\tau_{E^{a}}(\lambda_{1},\lambda_{2},\lambda_{3})=-\tau_{E}(\lambda_{1},\lambda_{2},\lambda_{3}).
\]

(vii) 令 $ (E_{1}, \sigma_{1}) $  和 $ (E_{2}, \sigma_{2}) $  为两个辛向量空间，并令 $ \lambda_{1}, \lambda_{2}, \lambda_{3} $  （分别为 $ \mu_{1}, \mu_{2}, \mu_{3} $  ）为 $ E_{1} $  （分别为 $ E_{2} $  ）的拉格朗日子空间的三元组。然后：

\[
\tau_{E_{1}\oplus E_{2}}(\lambda_{1}\oplus\mu_{2},\lambda_{2}\oplus\mu_{2},\lambda_{3}\oplus\mu_{3})=\tau_{E_{1}}(\lambda_{1},\lambda_{2},\lambda_{3})+\tau_{E_{2}}(\mu_{1},\mu_{2},\mu_{3})~.
\]

（这里 $ E_1 \oplus E_2 $ 被赋予辛形式 $ \sigma_1 \oplus \sigma_2 $ 。）

\statementlabel{证明}(i) 很清楚，因为二次形式  $q(x_{1}, x_{2}, x_{3})$  相对于  $\lambda_{j}$  的排列是交替的。

(iii) 对于 $ \lambda_1 \oplus \lambda_2 \oplus \lambda_3 $  中的 $ x = (x_1, x_2, x_3) $  和 $ y = (y_1, y_2, y_3) $  进行设置：

\[
B(x,y)=\sigma(x_{1},y_{2})+\sigma(x_{2},y_{3})+\sigma(x_{3},y_{1})+\sigma(y_{1},x_{2})+\sigma(y_{2},x_{3})+\sigma(y_{3},x_{1})
\]

q 的完全各向同性空间 I 的定义为：

\[
I=\left\{x\in\lambda_{1}\oplus\lambda_{2}\oplus\lambda_{3};B(x,y)=0\mathrm{f o r a n y}y\right\}.
\]

从 $ B(x,y)=\sigma(y_{1},x_{2}-x_{3})+\sigma(y_{2},x_{3}-x_{1})+\sigma(y_{3},x_{1}-x_{2}) $ 开始，我们有：

\[
I=\left\{(x_{1},x_{2},x_{3})\in\lambda_{1}\oplus\lambda_{2}\oplus\lambda_{3};x_{2}-x_{3}\in\lambda_{1},x_{3}-x_{1}\in\lambda_{2},x_{1}-x_{2}\in\lambda_{3}\right\}.
\]

设置 $ y_{1}=x_{2}+x_{3}-x_{1}, y_{2}=x_{3}+x_{1}-x_{2}, y_{3}=x_{1}+x_{2}-x_{3} $ 。然后：

\[
y_{1}\in\lambda_{2}\cap\lambda_{3},\qquad y_{2}\in\lambda_{3}\cap\lambda_{1},\qquad y_{3}\in\lambda_{1}\cap\lambda_{2}.
\]

而且

\[
2x_{1}=y_{2}+y_{3}~,~2x_{2}=y_{3}+y_{1}~,~2x_{3}=y_{1}+y_{2}~.
\]

因此，通过线性变换  $ (x_1, x_2, x_3) \mapsto (y_1, y_2, y_3) $  ，  $ I $  与  $ (\lambda_1 \cap \lambda_2) \oplus (\lambda_2 \cap \lambda_3) \oplus (\lambda_3 \cap \lambda_1) $  同构。由于  $ q $  （表示为  $ \mathrm{rk}(q) $  ）的等级是  $ 3n - \dim I $  ，我们得到：

\[
\mathrm{r k}(q)=3n-\operatorname{d i m}(\lambda_{1}\cap\lambda_{2})-\operatorname{d i m}(\lambda_{2}\cap\lambda_{3})-\operatorname{d i m}(\lambda_{3}\cap\lambda_{1})~.
\]

根据(iii)的假设，我们得到 $ \mathrm{rk}(q) $ 是常数。由于当  $ \lambda_{j} $  连续移动时 q 也连续移动，因此这意味着 (iii)。

（四）我们有：

\[
\operatorname{s g n}(q)\equiv\operatorname{r k}(q)\bmod{2\mathbb{Z}}.
\]

因此 (iv) 由 (A.3.2) 得出。

(ii) 暂时假设 $ \lambda_{1} $  和 $ \lambda_{2} $  横向相交并用 $ p_{1} $  和 $ p_{2} $  表示投影：

\[
p_{i}\colon E=\lambda_{1}\oplus\lambda_{2}\to\lambda_{i}\qquad(i=1,2).
\]

对于 E 中的 x 和 y，我们有：

\[
\begin{aligned}\sigma(p_{1}(x),y)&=\sigma(p_{1}(x),p_{1}(y)+p_{2}(y))\\&=\sigma(p_{1}(x),p_{2}(y))\\&=\sigma(x,p_{2}(y)).\end{aligned}
\]

\statementlabel{引理 A.3.3}假设 $\lambda_{1}$  和 $\lambda_{2}$  是横向的。那么  $\tau(\lambda_{1},\lambda_{2},\lambda_{3})$  等于  $\lambda_{3}$  上二次形式  $q_{3}$  的签名，定义如下：

\[
q_{3}(x_{3})=-\sigma(p_{1}(x_{3}),p_{2}(x_{3}))=-\sigma(p_{1}(x_{3}),x_{3}).
\]

引理 A.3.3 的证明。让 $ x = (x_1, x_2, x_3) \in \lambda_1 \oplus \lambda_2 \oplus \lambda_3 $  。然后：

\[
\begin{aligned}q(x)&=\sigma(x_{1},x_{2})+\sigma(x_{2},x_{3})+\sigma(x_{3},x_{1})\\&=\sigma(x_{1},x_{2})-\sigma(p_{1}(x_{3}),x_{2})-\sigma(x_{1},p_{2}(x_{3}))\\&=\sigma(x_{1}-p_{1}(x_{3}),x_{2}-p_{2}(x_{3}))-\sigma(p_{1}(x_{3}),p_{2}(x_{3}))~.\end{aligned}
\]

通过线性变换 (x\_1, x\_2, x\_3) \textbackslash{}mapsto (x\_1 - p\_1(x\_3), x\_2 - p\_2(x\_3), x\_3)，二次形式 (x\_1, x\_2, x\_3) \textbackslash{}mapsto \textbackslash{}sigma(x\_1 - p\_1(x\_3), x\_2 - p\_2(x\_3)) 等价于二次形式(x\_1, x\_2, x\_3) \textbackslash{}mapsto \textbackslash{}sigma(x\_1, x\_2)。因此它的签名为零，这证明了引理。   $ \square $ 

\statementlabel{引理 A.3.4}令  $ \lambda_j $  (  $ j = 1, 2, 3 $  ) 和  $ \mu $  为四个拉格朗日子空间，例如  $ \lambda_j \cap \mu = \{0\} $  、 (  $ j = 1, 2, 3 $  )。然后：

\[
\tau(\lambda_{1},\lambda_{2},\lambda_{3})=\tau(\lambda_{1},\lambda_{2},\mu)+\tau(\lambda_{2},\lambda_{3},\mu)+\tau(\lambda_{3},\lambda_{1},\mu).
\]

引理 A.3.4 的证明。根据引理 (A.3.3)，(A.3.4) 的右侧是  $ \lambda_1 \oplus \lambda_2 \oplus \lambda_3 $  上二次形式的签名：

\[
q^{\prime}(y_{1},y_{2},y_{3})=\sigma(p_{1}(y_{2}),y_{2})+\sigma(p_{2}(y_{3}),y_{3})+\sigma(p_{3}(y_{1}),y_{1})~,
\]

现在  $p_{j}$  表示投影：

\[
p_{j}\colon\lambda_{j}\oplus\mu\to\lambda_{j}\qquad(j=1,2,3).
\]

考虑  $ \lambda_1 \oplus \lambda_2 \oplus \lambda_3 $  的线性自同构，定义为：

\[
x_{1}=y_{1}+p_{1}(y_{2}),\quad x_{2}=y_{2}+p_{2}(y_{3}),\quad x_{3}=y_{3}+p_{3}(y_{1})
\]

\[
y_{1}=(x_{1}-p_{1}(x_{2})+p_{1}(x_{3}))/2,\quad y_{2}=(x_{2}-p_{2}(x_{3})+p_{2}(x_{1}))/2,
\]

\[
y_{3}=(x_{3}-p_{3}(x_{1})+p_{3}(x_{2}))/2.
\]

一个简单的计算表明：

\[
q(x_{1},x_{2},x_{3})=q^{\prime}(y_{1},y_{2},y_{3}),
\]

这证明了引理。 ⑨

定理A.3.2的证明结束。选择一个拉格朗日子空间  $ \mu $  横向于所有  $ \lambda_j $  (  $ j = 1, 2, 3, 4 $  ) 并应用 (A.3.4)。鉴于 (i)，则 (ii) 成立。

(v) 我们将把证明分解为几个步骤。

(a) 首先假设  $ \rho \subset \lambda_1 \cap \lambda_2 \cap \lambda_3 $  。那么  $ \lambda_1 \oplus \lambda_2 \oplus \lambda_3 $  上的二次形式  $ q $  是  $ \lambda_1^p \oplus \lambda_2^p \oplus \lambda_3^p $  上相应二次形式的拉回，通过满射映射  $ \lambda_1 \oplus \lambda_2 \oplus \lambda_3 \to \lambda_1^p \oplus \lambda_2^p \oplus \lambda_3^p $  。在这种情况下，断言如下。

(b) 现在假设  $ \rho \subset \lambda_2 \cap \lambda_3 $  和  $ \lambda_3^\rho = \lambda_1^\rho $  。考虑  $ \lambda_1 \oplus \lambda_2 \oplus (\lambda_1 \cap \rho^\perp) \oplus \rho $  上的二次形式  $ q'' $ ，定义如下：

\[
q^{\prime \prime}(x_{1},x_{2},u,v)=\sigma(x_{1},x_{2})+\sigma(x_{2},u+v)+\sigma(u+v,x_{1}).
\]

然后 $q''(x_1,x_2,u,v) = \sigma(x_1,x_2) + \sigma(x_2,u) + \sigma(v,x_1) = \sigma(x_1 - u,x_2 - v)$ 。因此 $q''$  的签名为零。根据假设， $\lambda_3 = (\lambda_1 \cap \rho^\perp) + \rho$  。这意味着  $q$  对  $\lambda_1 \oplus \lambda_2 \oplus \lambda_3$  的签名是  $q'$  的签名。因此在这种情况下 $\tau = 0$ 。

(c) 设置 $ \tilde{\lambda}_i = (\lambda_i \cap \rho^\perp) + \rho = (\lambda_i + \rho) \cap \rho^\perp $ 、 $ i = 1, 2, 3 $ 。我们有：

\[
\lambda_{1}^{\tilde{\lambda}_{1}\cap\lambda_{2}}=\tilde{\lambda}_{1}^{\tilde{\lambda}_{1}\cap\lambda_{2}}.
\]

事实上 $\rho \subset \lambda_{1} + \lambda_{2} \cap \lambda_{3} \subset \lambda_{1} + \lambda_{2} \cap \rho^{\perp}$ 意味着

\[
\rho\subset\lambda_{1}+(\lambda_{1}+\rho)\cap\lambda_{2}\cap\rho^{\perp}\subset\lambda_{1}+\tilde{\lambda}_{1}\cap\lambda_{2}.
\]

我们得到  $ \tilde{\lambda}_1 \subset \lambda_1 + (\tilde{\lambda}_1 \cap \lambda_2) $  因此  $ \tilde{\lambda}_1 \subset [\lambda_1 + (\tilde{\lambda}_1 \cap \lambda_2)] \cap (\tilde{\lambda}_1 + \lambda_2) $  给出：

\[
\tilde{\lambda}_{1}^{\tilde{\lambda}_{1}\cap\lambda_{2}}=\lambda_{1}^{\tilde{\lambda}_{1}\cap\lambda_{2}}.
\]

由于 (A.3.7) 中的两个空间都是拉格朗日，因此 (A.3.6) 如下。同样我们有：

\[
\lambda_{1}^{\tilde{\lambda}_{1}\cap\lambda_{3}}=\tilde{\lambda}_{1}^{\tilde{\lambda}_{1}\cap\lambda_{3}}.
\]

因此我们得到(b)：

\[
\tau(\lambda_{1},\tilde{\lambda}_{1},\lambda_{j})=0\qquad\mathrm{f o r}\qquad j=2,3.
\]

现在我们有：

\[
\begin{aligned}\tau(\lambda_{1},\lambda_{2},\lambda_{3})&=\tau(\lambda_{1},\lambda_{2},\tilde{\lambda}_{1})+\tau(\lambda_{2},\lambda_{3},\tilde{\lambda}_{1})+\tau(\lambda_{3},\lambda_{1},\tilde{\lambda}_{1})\\&=\tau(\tilde{\lambda}_{1},\lambda_{2},\lambda_{3}).\end{aligned}
\]

重复这个论证，我们得到

\[
\tau(\lambda_{1},\lambda_{2},\lambda_{3})=\tau(\tilde{\lambda}_{1},\tilde{\lambda}_{2},\tilde{\lambda}_{3})
\]

右边的项 eq.als  $ \tau(\lambda_{1}^{\rho},\lambda_{2}^{\rho},\lambda_{3}^{\rho}) $  (a)。

(vi)和(vii)是显而易见的。 □

\statementlabel{备注 A.3.5}定理 A.3.2 (ii) 中所写的“共循环条件”可以通过图 A.3.1 来形象化。

\begin{center}
\includegraphics[width=0.59\textwidth]{流行上的层_Chn-assets/img_in_image_box_224_1256_927_1667.png}
\end{center}

\begin{center}图。 A.3.1\end{center}

我们将在特殊情况下明确计算马斯洛夫指数。假设 $E$ 具有辛基，并用 $(x;\xi)$ 表示相关的线性坐标。令 $A$ 和 $B$ 为两个满足假设(A.1.11)且 $A^{t}B$ 对称的 $(n\times n)$ 矩阵，并定义：

\[
\lambda_{1}=\left\{x=0\right\},\quad\lambda_{2}=\left\{\xi=0\right\},\quad\lambda_{3}=\left\{(x;\xi);Ax=B\xi\right\}.
\]

\statementlabel{命题 A.3.6}其中一个有：

\[
\tau(\lambda_{1},\lambda_{2},\lambda_{3})=-\operatorname{sgn}(A^{t}B).
\]

\statementlabel{证明}分别用 $p_1$ 和 $p_2$ 表示从 $E = \lambda_1 \oplus \lambda_2$ 到 $\lambda_1$ 和 $\lambda_2$ 的投影。根据引理 A.3.3， $-\tau(\lambda_1, \lambda_2, \lambda_3)$  是二次形式  $q_3$  在  $\lambda_3$  上的签名，定义为：

\[
\begin{aligned}{q_{3}((x;\xi))}&{{}=\sigma(p_{1}(x;\xi),p_{2}(x;\xi))}\\ {}&{{}=\sigma(\xi,x)}\\ {}&{{}=\langle\xi,x\rangle.}\\ \end{aligned}
\]

 $z \mapsto (^{\prime}Bz, ^{\prime}Az)$  给出的从 $\mathbb{R}^n$  到 $\lambda_3$  的映射是线性同构。因此  $\langle \xi, x \rangle = \langle z, A^{\prime}Bz \rangle$  ，结果如下。   $\square$ 

令 $\lambda_{1}, \ldots, \lambda_{N}$  为 $N \geqslant 3$  的拉格朗日子空间，并令 $\mu$  为另一个拉格朗日子空间。根据定理 A.3.2 (ii)，我们有恒等式：

\[
\left\{\begin{aligned}{\tau}&{{}(\lambda_{1},\lambda_{2},\lambda_{3})+\tau(\lambda_{1},\lambda_{3},\lambda_{4})+\cdots+\tau(\lambda_{1},\lambda_{N-1},\lambda_{N})}\\ {}&{{}=\tau(\lambda_{1},\lambda_{2},\mu)+\tau(\lambda_{2},\lambda_{3},\mu)+\cdots+\tau(\lambda_{N-1},\lambda_{N},\mu)+\tau(\lambda_{N},\lambda_{1},\mu)~.}\\ \end{aligned}\right.
\]

这可以通过图 A.3.2 直观地看出（其中 N = 5）。

\begin{center}
\includegraphics[width=0.42\textwidth]{流行上的层_Chn-assets/img_in_image_box_342_1275_843_1659.png}
\end{center}

\begin{center}图。 A.3.2\end{center}

\statementlabel{定义 A.3.7}令 $ \lambda_1, \ldots, \lambda_N $  为拉格朗日子空间， $ N \geqslant 3 $  。将索引  $ \tau(\lambda_1, \ldots, \lambda_N) $  定义为 (A.3.10) 的左侧。

\statementlabel{命题 A.3.8}

\[
\tau(\lambda_{1},\lambda_{2},\ldots,\lambda_{N})=\tau(\lambda_{2},\lambda_{3},\ldots,\lambda_{N},\lambda_{1})=-\tau(\lambda_{N},\lambda_{N-1},\ldots,\lambda_{1}).
\]

(ii) 假设  $N \geqslant 4$  并设  $j \in \{3, \ldots, N-1\}$  。然后 $\tau(\lambda_1, \ldots, \lambda_N) = \tau(\lambda_1, \ldots, \lambda_j) + \tau(\lambda_1, \lambda_j, \lambda_{j+1}, \ldots, \lambda_N)$ 。

(iii) 当 $\lambda_j$ 以 $\dim(\lambda_1 \cap \lambda_2)$ 、 $\dim(\lambda_2 \cap \lambda_3)$ 、...、 $\dim(\lambda_{N-1} \cap \lambda_N)$ 和 $\dim(\lambda_N \cap \lambda_1)$ 保持恒定的方式连续移动时， $\tau(\lambda_1, \ldots, \lambda_N)$ 保持恒定。

\[
\begin{array}{r}{(\mathrm{i v})\tau(\lambda_{1},\ldots,\lambda_{N})\equiv n N+\dim(\lambda_{1}\cap\lambda_{2})+\cdots+\dim(\lambda_{N}\cap\lambda_{1})\bmod2\mathbb{Z}.}\end{array}
\]

\statementlabel{证明}根据定义和定理 A.3.2，(i) 和 (ii) 是显而易见的。当选择横向于所有 $ \lambda_j $  的拉格朗日空间 $ \mu $  时，(iii) 和(iv) 立即由定理A.3.2 和(A.3.10) 得出。   $ \Box $ 

如前所述，用  $E^{a}$  表示赋予辛形式  $-\sigma$  的空间  $E$  ，并用“  $a$  ”表示恒等式  $E\to E^{a}$  。如果  $\lambda$  是  $E$  的子空间，我们用  $\lambda^{a}$  表示其图像  $a$  。然后，我们有 $\tau_{E^{a}}(\lambda_{1}^{a},\lambda_{2}^{a},\lambda_{3}^{a})=-\tau_{E}(\lambda_{1},\lambda_{2},\lambda_{3})$ 。

\statementlabel{命题 A.3.9}令  $\lambda_1, \lambda_2, \mu_1, \mu_2$  为  $E$  的四个拉格朗日子空间，并用  $\Delta$  表示  $E^a \oplus E$  （拉格朗日）的对角线。然后：

\[
\tau_{E^{a}\oplus E}(\lambda_{1}^{a}\oplus\lambda_{2},\mu_{1}^{a}\oplus\mu_{2},\varDelta)=\tau_{E}(\lambda_{1},\lambda_{2},\mu_{2},\mu_{1})~.
\]

\statementlabel{证明}用 $ \tau $  表示(A.3.11) 的左侧。然后  $ \tau = \tau_1 + \tau_2 + \tau_3 $  ，其中：

\[
\tau_{1}=\tau(\lambda_{1}^{a}\oplus\lambda_{2},\mu_{1}^{a}\oplus\mu_{2},\lambda_{1}^{a}\oplus\mu_{2})~,
\]

\[
\tau_{2}=\tau(\mu_{1}^{a}\oplus\mu_{2},\varDelta,\lambda_{1}^{a}\oplus\mu_{2})~,
\]

\[
\tau_{3}=\tau(\varDelta,\lambda_{1}^{a}\oplus\lambda_{2},\lambda_{1}^{a}\oplus\mu_{2}).
\]

然后：

\[
\tau_{1}=\tau(\lambda_{1}^{a},\mu_{1}^{a},\lambda_{1}^{a})+\tau(\lambda_{2},\mu_{2},\mu_{2})=0,
\]

\[
\tau_{2}=\tau(\mu_{1}^{a},\mu_{2}^{a},\lambda_{1}^{a})~,
\]

\[
\tau_{3}=\tau(\lambda_{1},\lambda_{2},\mu_{2}).
\]

（为了计算  $\tau_2$  和  $\tau_3$ ，我们分别应用定理 A.3.2 (v) 和  $\rho = \{0\} \oplus \mu_2$  和  $\rho = \lambda_1^a \oplus \{0\}$  。）因此  $\tau = \tau(\lambda_1, \lambda_2, \mu_2) - \tau(\mu_1, \mu_2, \lambda_1) = \tau(\lambda_1, \lambda_2, \mu_2, \mu_1)$  。 □

为了结束本节，让我们描述  $Sp(E)$  在  $E$  拉格朗日子空间三元组的空间  $\mathcal{A}^{3}(E)$  上的作用。

对于 $ r = (r_0, r_1, r_2, r_3, d) \in \mathbb{N}^4 \times \mathbb{Z} $ ，我们设置：

\[
\begin{aligned}W_{r}&=\left\{(\lambda_{1},\lambda_{2},\lambda_{3})\in\varLambda^{3}(E)\;;\quad\operatorname{d i m}(\lambda_{1}\cap\lambda_{2}\cap\lambda_{3})=r_{0}\right.,\\&\quad\operatorname{d i m}(\lambda_{1}\cap\lambda_{2})=r_{3}\quad.\quad\operatorname{d i m}(\lambda_{2}\cap\lambda_{3})=r_{1}\quad,\\&\quad\operatorname{d i m}(\lambda_{3}\cap\lambda_{1})=r_{2}\quad,\quad\tau(\lambda_{1},\lambda_{2},\lambda_{3})=d\big\}.\end{aligned}
\]

然后  $Sp(E)$  自然地作用于  $A^{3}(E)$  ，并且  $W,^{\prime}$  s 对于  $Sp(E)$  是不变的。考虑条件：

\[
\begin{cases}0\leqslant r_{0}\leqslant r_{1},r_{2},r_{3}\leqslant n\ ,\quad r_{1}+r_{2}+r_{3}\leqslant n+2r_{0}\ ,\\|d|\leqslant n+2r_{0}-(r_{1}+r_{2}+r_{3})\ ,\quad d\equiv n+r_{1}+r_{2}+r_{3}\end{cases}
\]

可以很容易地证明，(A.3.12) 是  $W_{r}$  非空的充分必要条件。  $Sp(E)$  对  $A^{3}(E)$  的作用只有有限多个轨道，这些轨道是  $W_{r}$  ，其中  $r$  满足 (A.3.12)。由于我们不使用这个结果，因此我们将证明作为练习。

\subsection{习题}
\statementlabel{练习 A.1}令 $ (x) = (x_1, \ldots, x_n) $  为 $ X $  上的局部坐标系， $ (x; \xi) $  为 $ T^*X $  上的关联坐标。令 $ \varphi = (\varphi_1, \ldots, \varphi_p) $  :  $ X \to \mathbb{R}^p $  为平滑映射，并令 $ S = \{x; \varphi(x) = 0\} $  ，（因此 $ \mathrm{d}\varphi_1 \wedge \cdots \wedge \mathrm{d}\varphi_p \neq 0 $  和 $ S $  是平滑的）。让 $ x_0 \in S $ ， $ p = (x_0; \sum_i a_i \mathrm{d}\varphi_i(x_0)) \in T_S^*X $ 。证明：

\[
\begin{aligned}T_{S}^{*}X&=\left\{(x;\xi);\varphi(x)=0,\xi_{j}=\sum_{i=1}^{p}\lambda_{i}\frac{\partial\varphi_{i}}{\partial x_{j}},\lambda_{i}\in\mathbb{R}\right\},\\T_{p}T_{S}^{*}X&=\left\{(x;\overset{\circ}{\xi});\sum_{j=1}^{n}\frac{\partial\varphi_{i}}{\partial x_{j}}(x_{\circ})x_{j}=0,i=1,\ldots,p,\right.\\&\left.\xi_{j}=\sum_{i=1}^{p}\frac{\partial\varphi_{i}}{\partial x_{j}}(x_{\circ})\lambda_{i}+\sum_{k=1}^{n}\sum_{i=1}^{p}a_{i}\frac{\partial^{2}\varphi_{i}(x_{\circ})}{\partial x_{k}\partial x_{j}}x_{k},\lambda_{i}\in\mathbb{R}\right\}.\end{aligned}
\]

\statementlabel{练习 A.2}令  $A$  为  $T^*X$  的闭合圆锥拉格朗日子流形。证明存在  $X$  的子流形  $M$  使得  $A = T_M^*X$  。

\statementlabel{练习 A.3}令 $X$  和 $Y$  为两个相同维数的流形， $f$  为 $X \times Y$  上的实函数， $df \neq 0$  为 $S = \{f = 0\}$  上的实函数。设置 $A = \dot{T}_S^*(X \times Y)$ 。

当且仅当行列式  $\left(\begin{array}{cc}0&d_y f\\ d_x f&d_{xy}^2f\end{array}\right)$  不在  $S$  上消失时，证明  $p_1|_{\mathcal{A}}$  和  $p_2^a|_{\mathcal{A}}$  是局部同构（分别从  $\mathcal{A}$  到  $T^*X$  和  $\mathcal{A}$  到  $T^*Y$  ）。

\statementlabel{练习 A.4}令 $(E_1, \sigma_1)$  和 $(E_2, \sigma_2)$  为两个辛向量空间，并令 $\lambda$  为 $E_1 \oplus E_2$  的拉格朗日线性子空间。证明 $p_1|_{\lambda}$  是单射的当且仅当 $p_2|_{\lambda}$  是满射的。

\statementlabel{练习 A.5}令  $f: Y \to X$  为流形的态射，并令  $p \in Y \times_X T^*X$  。

(a) 令 $\mathcal{A}_{Y}$  为 $T^{*}$  Y 的拉格朗日子流形。假设 $t'f'$  相对于 $p_{Y}$  处的 $\mathcal{A}_{Y}$  是干净的。证明对于  $U$  的  $p$  足够小的邻域， $f_{\pi}(U\cap t'f'^{-1}(\mathcal{A}_{Y}))$  是光滑的拉格朗日流形。

(b) 令  $A_X$  为  $T^*X$  的拉格朗日子流形。假设  $f_\pi$  相对于  $p_X$  处的  $A_X$  是干净的。证明对于  $U$  的  $p$  足够小的邻域， $t f'(U \cap f_\pi^{-1}(A_X))$  是光滑的拉格朗日流形。

\statementlabel{练习 A.6}令 $ E_1 $  和 $ E_2 $  为两个辛向量空间， $ \nu $  为 $ E_1 \oplus E_2^a $  的拉格朗日子空间，并令 $ \lambda_i $  和 $ \mu_i $  为 $ E_i $  (  $ i = 1, 2 $  ) 的两个拉格朗日子空间。证明：

\[
\begin{aligned}\tau_{E_{1}\oplus E_{2}^{a}}(\lambda_{1}\oplus\lambda_{2}^{a},\mu_{1}\oplus\mu_{2}^{a},\nu)&=\tau_{E_{1}}(\lambda_{1},\mu_{1},\nu\circ\mu_{2})-\tau_{E_{2}}(\lambda_{2},\mu_{2},\lambda_{1}\circ\nu^{a})\\&=\tau_{E_{1}}(\lambda_{1},\mu_{1},\nu\circ\lambda_{2})-\tau_{E_{2}}(\lambda_{2},\mu_{2},\mu_{1}\circ\nu^{a})~.\end{aligned}
\]

\statementlabel{练习 A.7}令  $(E,\sigma)$  为复辛向量空间，并赋予实数底层向量空间  $E^{R}$  实数辛形式  $2\operatorname{Re}\sigma$  。令  $\lambda_{1},\lambda_{2},\lambda_{3}$  为  $E$  的三个复拉格朗日子空间。证明 $\tau_{E^{R}}(\lambda_{1},\lambda_{2},\lambda_{3})=0$ 。

\statementlabel{练习 A.8}令  $ E_i $  (  $ i = 1, 2, 3, 4 $  ) 为辛向量空间，并令  $ \lambda_i \subset E_i \oplus E_{i+1}^a $  (  $ i = 1, 2, 3 $  ) 为拉格朗日平面。

证明 $ (\lambda_{1} \circ \lambda_{2}) \circ \lambda_{3} = \lambda_{1} \circ (\lambda_{2} \circ \lambda_{3}) $ 。

\statementlabel{练习 A.9}令 $ (E_{i},\mu_{i}) $  为一对辛向量空间 $ E_{i} $  和 $ E_{i} $  (i=1,2,3,4) 的拉格朗日平面 $ \mu_{i} $ 。

然后对于拉格朗日平面  $ \lambda_1 \subset E_1 \oplus E_2^a $  和  $ \lambda_2 \subset E_2 \oplus E_3^a $  ，我们定义（参见 (7.5.10)）：

\[
\tau(\lambda_{1}:\lambda_{2})=\tau_{E_{2}}(\mu_{2},\lambda_{2}\circ\mu_{3},\mu_{1}\circ\lambda_{1}^{a})~.
\]

现在，令  $ \lambda_{i} $  为  $ E_{i}\oplus E_{i+1}^{a} $  (i = 1, 2, 3) 的拉格朗日平面。证明等式：

\[
\tau(\lambda_{1}:\lambda_{2}\circ\lambda_{3})+\tau(\lambda_{2}:\lambda_{3})=\tau(\lambda_{1}\circ\lambda_{2}:\lambda_{3})+\tau(\lambda_{1}:\lambda_{2})~.
\]

（提示：使用练习 A.6，证明两边都等于：

\[
\tau_{E_{2}\oplus E_{3}^{a}}(\mu_{2}\oplus\mu_{3}^{a},\mu_{1}\circ\lambda_{1}^{a}\oplus\lambda_{3}^{a}\circ\mu_{4}^{a},\lambda_{2}).)
\]

\statementlabel{练习 A.10}令 $E_i$  为辛向量空间， $\lambda_i$  和 $\mu_i$  为 $E_i$  (  $i=1,2,3$  ) 的两个拉格朗日子空间。令  $v$  和  $v'$  分别为  $E_1\oplus E_2^a$  和  $E_2\oplus E_3^a$  的拉格朗日子空间。证明：

\[
\begin{aligned}\tau_{E_{1}\oplus E_{2}^{a}}&(\lambda_{1}\oplus\lambda_{2}^{a},\nu,\mu_{1}\oplus\mu_{2}^{a})+\tau_{E_{2}\oplus E_{3}^{a}}(\lambda_{2}\oplus\lambda_{3}^{a},\nu^{\prime},\mu_{2}\oplus\mu_{3}^{a})\\&-\tau_{E_{1}\oplus E_{3}^{a}}(\lambda_{1}\oplus\lambda_{3}^{a},\nu\circ\nu^{\prime},\mu_{1}\oplus\mu_{3}^{a})\\&=\tau_{E_{2}}(\lambda_{2},\nu^{\prime}\circ\lambda_{3},\lambda_{1}\circ\nu^{a})-\tau_{E_{2}}(\mu_{2},\nu^{\prime}\circ\mu_{3},\mu_{1}\circ\nu^{a}).\end{aligned}
\]

（提示：使用练习 A.6。）

\subsection{注记}
辛几何和接触几何都是经典学科，可以追溯到汉密尔顿和雅可比，我们在这里不做回顾。正如本附录开头所指出的，§1 和 §2 的结果是众所周知的，例如可以在 Hömander [4] 中找到。

1965 年，为了计算“在焦散附近”的渐近展开式（即当平滑拉格朗日流形的投影没有恒定等级时），Maslov [1]（也参见 Keller [1]）引入了辛空间拉格朗日子流形中闭合曲线的指数。他的理论由 Arnold [1] 澄清和重新表述，然后由 Hörmander [2] 和 Leray [4] 澄清和重新表述，他们定义了三个横向相交的拉格朗日平面的索引，直到 Kashiwara（参见 Lion-Vergne [1]）通过我们在这里给出的简单方法在一般情况下定义了索引  $ \tau $ 。请注意，在 Lion-Vergne（前述）中，索引  $ \tau $  被推广到局部字段情况。

Abraham, R. 和 Mardsen, J.E. [1]：力学基础。卡明斯出版。 (1978)

Andronikof, E. [1]：分布和全息函数 I、II 的微局部化温度。 C.R.阿卡德。科学。 303, 347–350 (1986) 和 304, 511–514 (1987)

Arnold, V. [1]：关于进入量化条件的特征类。功能。肛门。应用。 1, 1–13 (1967)

—[2]：经典力学的数学方法。施普林格，纽约柏林海德堡 (1978)

Banica, C. 和 Stanasila, O. [1]：全球空间复合体理论中的代数方法。卷。一-二。高蒂尔-维拉斯-博达斯 (1977)

Beilinson, A.A.、Bernstein, J. 和 Deligne, P. [1]：Faisceaux pervers。星号 100 (1982)

Bengel, G. 和 Schapira, P. [1]：分布微区域分解分析。安.研究所。傅里叶格勒诺布尔 29, 101–124 (1979)

Bernstein, J. [1]：微分算子环上的模。研究具有常系数的方程的基本解。功能。肛门。应用。 5, 89–101 (1971)

— [2]：广义函数关于参数的解析延拓。功能。肛门。应用。 6, 272–285 (1972)

Berthelot, P.、Breen, L. 和 Messing, W. [1]：Théorie de Dieudonné cristalline II。莱克特。数学笔记。 930. 施普林格，柏林 海德堡 纽约 (1982)

Bierstone, E. 和 Milman, P.D. [1]：半解析集和亚解析集。出版。数学。 I.H.E.S. 67, 5–42 (1988)

比约克，J-E。 [1]：微分算子环。北荷兰数学。库。 (1979)

— [2]：分析 D 模。克鲁维尔出版社。出现 (1991)

Bloom, T. 和 Herrera, M. [1]：解析空间的 De Rham 上同调。发明。数学。 7、275–296（1969）

博尼，J-M。和 Schapira, P. [1]：粒子微分方程全态解的存在与延长。发明。数学。 17, 95–105 (1972)

— [2]：柯西问题的超函数解决方案。见：超函数和伪微分方程，Komatsu H.（编辑），Proceedings Katata 1971。Lect。数学笔记。 287, 82–98。施普林格，柏林海德堡纽约 (1973)

— [3]：粒子方程解的奇点分析传播。安.研究所。傅里叶格勒诺布尔 26, 81–140 (1976)

Borel, A. [1]：广义流形中的庞加莱对偶性。密歇根数学。 J. 4, 227–239 (1957)

Borel, A. 等人 [1]：交集上同调。数学进步。 50，伯克豪瑟，波士顿 (1984)

Borel, A. 和 Haeflinger, A. [1]：La classe d'homologie foldamentale d'un espace analytique。公牛。苏克。数学。法国 89, 461–513 (1961)

Borel, A. 和 Moore, J.C. [1]：局部紧空间的同调理论。密歇根数学。 J. 7, 137–159 (1960)

Bott, R. [1]：莫尔斯理论新旧讲座。公牛。阿米尔。数学。苏克。 7, 331–358 (1982)

Bourbaki, N. [1]：阿尔盖布尔，第 10 章。数学原理。马森 (1980)

Boutet de Monvel, L. 和 Malgrange, B. [1]：Le théorème de l'indice relatif。安.科学。例如。规范。高级。 23、151-192（1990）

布雷登，G.E. [1]：层理论。麦格劳-希尔 (1967)

本纳，A.V.和 Shubin, M.A. [1]：带边界流形的 Atiyah-Bott-Lefschetz 定理。功能。肛门。应用。 15, 286–287 (1981)

布林斯基，J-L。 [1]：(Co-)homologie d'intersection et faisceaux pervers。塞姆。布尔巴基 585 (1981-82)

— [2]：经典变换、二元射影、莱夫谢茨理论。星号 140/141, 3-134 (1986)

Brylinski, J-L.、Malgrange, B. 和 Verdier, J-L.。 [1]：傅里叶几何变换 I. C.R. Acad。科学。 297, 55–58 (1983)

Brylinski, J.-L.、Dubson, A. 和 Kashiwara, M. [1]：完整模指数公式和欧拉区域设置障碍。 C.R.阿卡德。科学。 293, 573–576 (1981)

嘉当，H. [1]：Sém。例如。规范。高级。 （1950-51），本杰明（1967）。

— [2]：Sém。例如。规范。高级。 （1951–52、1953–54、1960–61），本杰明（1967）

Cartan, H. 和 Chevalley, C. [1]：Sém。例如。规范。高级。 (1955–56)，本杰明 (1967)

Cartan, H. 和 Eilenberg, S. [1]：同调代数。普林斯顿大学出版社 (1956)

D'Agnolo, A. 和 Schapira, P. [1]：Un théorème d'image inverse pour les faisceaux。柯西问题的应用。 C.R.阿卡德。科学。 311, 23–26 (1990)

Deligne, P. [1]：Cohomologie à support propre。在[SGA 4]中，揭露XVII

— [2]：循环消失的形式主义。在[SGA 7]中，曝光XIII

德洛特，J-M。 [1]：Deuxième microlocalization simultanée et front d'onde de produits。安.科学。例如。规范。高级。 23, 257–310 (1990)

德洛特，J-M。和 Lebeau, G. [1]：拉格朗日微函数。 J.马斯。纯应用程序。 67, 39–84 (1988)

Denkowska, Z.、Lojasiewicz, S. 和 Stasica, J. [1]：分析中整体的某些属性。公牛。阿卡德。波隆.科学。数学。 27, 529–536 (1979)

Dubson, A. [1]：可构造复杂性指数和 D 模完整值的公式。 C.R.阿卡德。科学。 298, 113–164 (1984)

杜斯特马特，J.J. [1]：傅里叶积分算子。莱克特。库朗研究所注释纽约，（1973）

Freyd, P. [1]：阿贝尔范畴，函子理论简介。哈珀和罗，纽约 (1964)

Fulton, W. [1]：相交理论。施普林格，柏林 海德堡 纽约 东京 (1984)

Gabber, O. [1]：特征簇的可积性。阿米尔。 J.马斯。 103, 445–468 (1981)

Gabriel, P. 和 Zisman, M. [1]：分数微积分和同伦理论。施普林格，柏林海德堡纽约 (1967)

加布里埃洛夫，A.M. [1]：半解析集的投影。功能。肛门。应用。 2, 282–291 (1968)

格尔凡德 (Gelfand)，S. 和马宁 (Manin)，Yu。 [1]：同调代数方法I：上同调理论及其导出范畴简介。施普林格柏林海德堡纽约东京 (1991)

Ginsburg, V. [1]：关于微分系统指数和奇点流形几何的定理。苏联数学。多克尔。 31, 309–313 (1985)

— [2]：特征周期和消失循环。发明。数学。 84, 327–402 (1986)

戈洛文，V.D. [1]：解析层的上同调和对偶定理。诺卡，莫斯科 (1986)（俄语）

Godement, R. [1]：拓扑代数和 faisceaux 理论。赫尔曼·帕里斯 (1958)

Goresky, M. 和 MacPherson, R. [1]：交叉同源性 II。发明。数学。 71, 77–129 (1983)

— [2]：分层莫尔斯理论。施普林格，柏林海德堡 (1988)

— [3]：莫尔斯理论和交叉同源理论。在：奇异空间的分析和拓扑。星号 101/102, 135–192 (1983)

Grauert, H. [1]：关于李维问题和实解析流形的嵌入。安.数学。 68、460–472（1958）

Griffiths, P. 和 Harris, J. [1]：代数几何原理。约翰·威利父子 (1978)

Grothendieck, A. [1]：关于同调代数的某些点。东北数学. J. 9, 119–221 (1957)

— [2]：代数几何中的下降技术和存在定理 II。形式模理论中的存在定理。 SEM。布尔巴基 195 (1959–60)

— [3]：代数几何原理 III．出版。数学。 I.H.E.S. 11 (1961) 和 17 (1963)

— [4]：残差和对偶性。 “Hartshorne 研讨会”的前注，手稿（1963 年）

— [5]：关于图式理论的十次演讲。北荷兰省，阿姆斯特丹 (1968)

Guillemin, V.、Quillen, D. 和 Sternberg, S. [1]：特征的可积性。通讯。纯应用。数学。 23, 39–77 (1970)

Guillemin, V. 和 Pollack, A. [1] 微分拓扑。普伦蒂斯·霍尔 (1974)

Guillemin, V 和 Sternberg S. [1] 物理学中的辛技术。剑桥大学出版社 (1984)

哈特，R.M. [1]：真实分析映射和图像的分层。发明。数学。 28, 193–208 (1975)

— [2]：子分析集的三角测量和适当的光子分析图。发明。数学。 38, 207–217 (1977)

Hartshorne, R. [1]：残差和对偶性。莱克特。数学笔记。 20.施普林格，柏林海德堡纽约（1966）

Henry, J-P.、Merle, M. 和 Sabbah, C. [1]：关于复分析态射的严格 Thom 条件。安.科学。例如。规范。高级。 17, 227–268 (1984)

Herrera, M. [1]：半解析集上的积分。公牛。苏克。数学。法国 94, 141–180 (1966)

Hilton P.J. 和 Stammbach, U. [1]：同调代数课程。施普林格，纽约柏林海德堡 (1970)

Hironaka, H. [1]：亚解析集。见：数论、代数几何和交换代数（纪念 Y. Akizuki）。纪伊国屋，东京 453–493 (1973)

— [2]：实解析集和实解析图简介。比萨“L.托内利”学院 (1973)

— [3]：分层和平坦度。实几何和复几何，奥斯陆 1976，Sythoff \& Noordhoff 199–265 (1977)

Hörmander, L. [1]：多个变量的复杂分析简介。范·诺斯特兰德，普林斯顿 (1966)

— [2]：傅立叶积分算子 I. Acta Math。 127, 79–183 (1971)

—[3]：线性偏微分算子的分析 I. Springer，柏林 海德堡 纽约 东京（1983）

— [4]：线性偏微分算子的分析III-IV。施普林格，柏林 海德堡 纽约 东京 (1985)

Hotta, R. 和 Kashiwara, M. [1]：半简单李代数上的不变完整系统。发明。数学。 75, 327–358 (1984)

Houzel, C. 和 Schapira, P. [1]：差分模的直像。 C.R.阿卡德。科学。 298, 461–464 (1984)

Illusie, L. [1]：Deligne 的 l 进傅立叶变换。过程。症状。纯数学。 46, 151–163 (1987)

Iversen, B. [1]：层的上同调。施普林格，柏林 海德堡 纽约 东京 (1987)

—[2]：柯西残基和德拉姆同源性。军旗。数学。 35, 1–17 (1989)

Kashiwara, M. [1]：偏微分方程组的代数研究。论文，大学。东京 (1970)

— [2]：线性微分方程最大超定系统的指数定理。过程。日本科学院. 49, 803–804 (1973)

— [3]：关于线性微分方程的最大超定系统 I.Publ。 R.I.M.S.京都大学10, 563–579 (1975)

— [4]：b 函数和完整系统。发明。数学。 38, 33–53 (1976)

— [5]：微微分方程组。数学进步。 34. 伯克豪瑟，波士顿 (1983)

— [6]：完整系统的黎曼-希尔伯特问题。出版。 R.I.M.S.京都大学20, 319–365 (1984)（或：具有正则奇点的偏微分方程的可构造层和完整系统。Sém Eq. Dér. Part. Publ. Ec. Polyt. (1979/1980)）

— [7]：可构造层的指数定理。见：微分系统和奇点，A. Galligo、M. Maisonobe、Ph. Granger（主编）。星号 130, 193–209 (1985)

— [8]：字符、字符循环、不动点定理和群表示。副词。螺柱。纯数学。 14, 369–378 (1988)

Kashiwara, M. 和 Kawai, T. [1]：线性微分方程椭圆系统的边值问题。过程。日本科学院. 48, 712–715 (1971) 和 49, 164–168 (1972)

— [2]：第二次微局部化和渐近展开。见：复分析、微局域微积分和相对论量子理论，D. Iagolnitzer（主编）。莱克特。注释 物理。 126、21-76。施普林格，柏林 海德堡 纽约 (1980)

— [3]：关于微微分方程的完整系统，III。出版。 RIMS，京都大学。 17, 813–979 (1981)

柏原，M. 和蒙泰罗-费尔南德斯，T. [1]。 Involutvité des variétés microcaractéristiques。公牛。苏克。数学。法国 114, 393–402 (1986)

Kashiwara, M. 和 Schapira, P. [1]：微双曲系统。数学学报。 142, 1–55 (1979)

— [2]：微支撑 des faisceaux。见：Journées Complexes Nancy, Mai 1982（Publ. Inst. Elie Cartan），或：C.R. Acad。科学。 295, 487–490 (1982)

— [3]：层的微局部研究。星号 128 (1985)

— [4]：一类具有简单特征的系统的消失定理。发明。数学。 82, 579–592 (1985)

Kataoka, K. [1]：边值问题的微局域理论 I-II。 J.Fac。科学。大学。东京 27, 355–399 (1980) 和 28, 31–56 (1981)

Katz, N. 和 Laumon, G. [1]：傅立叶变换和指数指数变换。出版。数学。 I.H.E.S. 62, 146–202 (1985)

Keller, J.B. [1]：修正了不可分离系统的玻尔-索末菲量子条件。安.物理。 4、180–188（1958）

Komatsu, H. [1]：通过层超函数求解常系数微分方程的解。数学。安. 176, 77–86 (1968)

— [2]：博赫纳管定理的局部版本。 J.Fac。科学。大学。东京 19, 201–214 (1972)

Kuo, J.C. [1]：分析 Whitney 分层的比率检验。在：过程。利物浦奇点研讨会。莱克特。数学笔记。 192. 施普林格，柏林 海德堡 纽约 (1971)

Laumon, G. [1]：傅立叶变换、函数常数方程和韦尔猜想。出版。数学。 I.H.E.S. 65, 131–210 (1987)

Lazzeri, F. [1]：奇异空间莫尔斯理论。星号 7/8, 263–268 (1973)

Lebeau, G. [1]：半直线方程 II。非线性奇异点和腐蚀点控制。发明。数学。 95, 277–323 (1988)

Leray, J. [1]：L'anneau d'homologie d'une representation。 C.R.阿卡德。科学。 222、1367–1368。结构同源表示。同上，1419–1422 (1946)

— [2]：L'anneau 光谱和 l'anneau filtré d'un 空间定位紧凑和应用程序继续。 J.马斯。纯应用程序。 29, 1-139 (1950)

— [3]：Probleme de Cauchy I. Bull。苏克。数学。法国 85, 389–430 (1957)

— [4]：拉格朗日分析。麻省理工学院出版社，剑桥马萨诸塞州伦敦 (1981)，或：Cours Collège de France (1976/77)

Lion, G. 和 Vergne, M. [1]：Weil 表示、马斯洛夫指数和 theta 级数。数学进步。 6、博克豪瑟，波士顿 (1980)

Lojasiewicz, S. [1]：系综半分析。研究所。高级研究科学。伊维特河畔布尔 (1964)

— [2]：半解析集的三角剖分。安.斯库拉规范。高级。比萨 18, 449–474 (1964)

MacLane, S. [1]：同源理论。学术出版社（1963）

MacPherson, R. [1]：单变体的 Chern 类。安.数学。 100, 423–432 (1974)

MacPherson, R. 和 Vilonen, K. [1]：反常层的基本构造。库存。数学。 84, 403–435 (1986)

Malgrange, B. [1]：Faisceaux sur les variétés analytiques réelles。公牛。苏克。数学。法国 85, 231–237, (1957)

— [2]：傅立叶几何变换。塞姆。布尔巴基，692（1987–88）

Martineau, A. [1]：M. Sato 的超级功能。塞姆。布尔巴基 214 (1960–61)

— [2]：“楔形边缘”类型延长分析的理论。 SEM。布尔巴基 340 (1967/68)

马斯洛夫，副总裁[1]：摄动理论和渐近方法。莫斯科大学出版社 (1965)（俄语）

Mebkhout, Z. [1]：Une équivalence de catégories - Une autre équivalence de catégories。比较。数学。 51、55–62 和 63–64 (1984)

Milnor, J.M. [1]：莫尔斯理论。安.数学。研究 51，普林斯顿大学。按（1963）。

—[2]：复杂超曲面的奇异点。安.数学。研究 61，普林斯顿大学。出版社 (1968)

Mitchell, B. [1]：范畴论。纯应用程序。数学。 17、学术出版社（1965）

诺斯科特，D.G. [1]：同调代数简介。剑桥大学出版社 (1960)

Oka, K. [1]：关于优点变量复合体的功能分析。东京岩波商店 (1961)

Pignoni, N. [1]：分层空间上莫尔斯函数的密度和稳定性。安.科学。规范。高级。比萨，593–608 (1979)

保利，J-B。 [1]：残渣与分析链交集公式。这些大学。普瓦捷 (1974)

拉米斯，J-P。 [1]：Additif II à“GAGA 主题变奏”。莱克特。数学笔记。 694. 施普林格，柏林 海德堡 纽约 (1978)

de Rham, G. [1] Variétés différentiables。 Hermann, Paris (1955) 或：可微流形。施普林格，柏林 海德堡 纽约 东京 (1984)

Remmert, R. 和 Stein, K.[1]：我们的奇异性分析。数学。安. 126, 263–306 (1953)

Sabbah, C. [1]：关于几何空间的评论。见：系统差异和奇异点，A. Galligo、M. Maisonobe、Ph. Granger（主编）。星号 130, 161–192 (1985)

Sato, M. [1]：机能亢进理论 I-II。 J.Fac。科学。大学。东京 8, 139–193 和 387–436 (1959–1960)

Sato, M. [2]：超函数和偏微分方程。在：过程。国际。会议。东京大学泛函分析及相关主题。出版社，东京，91–94 (1969)

Sato, M.、Kawai, T. 和 Kashiwara, M. [1]：超函数和伪微分方程。见：超函数和伪微分方程，Komatsu H.（编辑），Proceedings Katata 1971。Lect。数学笔记。 287, 265–529。施普林格，柏林海德堡纽约（1973）。

Schapira, P. [1]：超函数理论。莱克特。数学笔记。 126.施普林格，柏林海德堡纽约（1970）

— [2]：复杂域中的微微分系统。施普林格，柏林 海德堡 纽约 东京 (1985)

— [3]：边值问题的微函数。见：代数分析，809–819（献给 M. Sato 的论文）。 M. Kashiwara，T. Kawai（编）。学术出版社（1988）

— [4]：偏微分方程的层理论。过程。国际米兰。丛。数学。 (1990) 出现

Schapira, P. 和 Tose, N. [1]：R 可构造层的莫尔斯不等式。副词。在数学中。出现

Schapira, P. 和 Schneiders, J-P。 [1]：Paires elliptiques I.有限性和二元性。 C.R.阿卡德。科学。 311, 83–86 (1990) 或：Finitude etclasses caractéristiques pour les D-模 et les faisceaux constructibles。塞姆。等式。德尔。部分。出版。例如。波利特。 (1989/90)

施奈德斯，J-P。 [1]：关于模差异的相对二元理论。 C.R.阿卡德。科学。 303, 235–238 (1986)

Schwartz, L. [1]：同态和应用的完善仍在继续。 C.R.阿卡德。科学。 236, 2472–2473 (1953)

— [2]：分布理论。赫尔曼，巴黎 (1966)

施瓦茨，M-H。 [1]：类特征定义了复杂分析的分层。 C.R.阿卡德。科学。 260、3262–3264 和 3535–3537 (1965)

[SHS]：Sém。海德堡-斯特拉斯堡（1966-67）。庞加莱对偶。出版。 I.R.M.A. 3、斯特拉斯堡（1969）

[SGA 2]：Sém。 géométrie algébrique (1962)，Grothendieck, A. Cohomologie locale des faisceaux cohérents et théorèmes de Lefschetz locaux et globaux。北荷兰省，阿姆斯特丹 (1968)

[SGA 4]：Sém。 géométrie algébrique (1963–64)，作者：Artin, M.、Grothendieck, A. 和 Verdier, J-L。拓扑理论和模式的同系学。莱克特。数学笔记。 269, 270, 305.施普林格，柏林海德堡纽约（1972-73）

[SGA 4½]：SEM。 géométrie algébrique，作者：Deligne, P. Cohomologie étale。莱克特。数学笔记。 569. 施普林格，柏林 海德堡 纽约 (1977)

[SGA 5]：Sém。 géométrie algébrique (1965–66) 作者：Grothendieck, A. Cohomologie l-adique et fonctions L. Lect。数学笔记。 589. 施普林格，柏林 海德堡 纽约 (1977)

[SGA 6]：Sém。 géométrie algébrique (1966-67)，作者：Berthelot, P. Illusie, L. 和 Grothendieck, A. Théorie des junctions et théorème de Riemann-Roch。莱克特。数学笔记。 225. 施普林格，柏林 海德堡 纽约 (1971)

[SGA 7]：Sém。几何阿尔格布里克（1967-69）。阿尔及几何的单一组。第一部分，格洛腾迪克，A. Lect。数学笔记。 288.施普林格，柏林海德堡纽约，(1972)。第二部分，作者：Deligne P. 和 Katz, N. Lect。数学笔记。 340. 施普林格，柏林 海德堡 纽约 (1973)

塞尔，J-P。 [1]：Faisceaux algébriques cohérents。安.数学。 61, 197–278 (1955)

Sjöstrand, J. [1]：微区域奇异性分析。星号 95 (1982)

Sullivan, D. [1]：分析空间的组合不变量。利物浦奇点研讨会论文集 I. Lect。数学笔记。 192、165–168。施普林格，柏林海德堡纽约 (1971)

Tamm, M. [1]：变分演算中的亚解析集。数学学报。 146, 167–199 (1981)。

Thom, R. [1]：系综和态射层。公牛。 A.M.S. 75, 240–284 (1969)

Teissier, B. [1]：关于态射分析的三角剖分。出版。数学。 I.H.E.S. 70 (1989)

Tose, N. [1]：层传播定理及其在微微分系统中的应用。 J.马斯。纯应用程序。 68, 137–151 (1989)

Trotman, D. [1]：比较分层的规律性条件。过程。症状。纯数学。 40, 575–585 (1983)

— [2]：Verdier 条件 (w) 的微区域版本。安.研究所。傅立叶格勒诺布尔 39, 825–829 (1989)

Uchida, M. [1]：障碍物角部衍射的微局域分析。出现，或：Sém。等式。德尔。部分。出版。例如。波利特。 (1989/90)

韦尔迪埃，J-L. [1]：空间位置契约中的双重性。塞姆。布尔巴基 300 (1965–66)

—[2]：导出范畴，状态 0。 4½]

— [3]：Whitney 分层和 Bertini-Sard 定理。发明。数学。 36, 295–312 (1976)

— [4]：与循环相关的同源类。在塞姆。吉姆.肛门。 A.杜阿迪，J-L。韦尔迪埃（编）。星号 36–37, 101–151 (1976)

— [5]：分册的专业化和中等的单一性。在：奇异空间的分析和拓扑。星号 101-102, 332-364 (1981)

威尔斯，R.O. Jr [1]：复流形的微分分析。施普林格，纽约柏林海德堡 (1980)

Whitney, H. [1]：分析簇的局部属性。在：微分和组合拓扑（纪念马斯顿·莫尔斯的研讨会）。普林斯顿大学出版社，205–244 (1965)

— [2]：解析簇的切线。安.数学。 81, 496–549 (1965)

— [3]：复解析簇。艾迪生韦斯利，《阅读弥撒》（1972）

Zerner, M. [1]：满足偏微分方程的函数的全纯域。 C.R.阿卡德。科学。 272, 1646–1648 (1971)

\subsection{记号与约定表}
\subsection{一般记号}
N：非负整数集 Z：整数环 Q：有理数域 R：实数域 C：复数域 R⁺：正实数乘法群 R⁻：负实数集  $ \mathbb{R}_{\geq 0} $ （分别为 $ \mathbb{R}_{\leq 0} $ ）： $ \{c \in \mathbb{R}; c \geq 0, (\text{resp.} c \leq 0)\} $  C×：非零复数乘法群

\subsection{集合与基本记号}
A \textbackslash{} B：A 中 B 的补集  $ \delta_{i,j} $  ：克罗内克符号， $ i \neq j $  的  $ \delta_{i,j} = 0 $  和  $ i = j $  的  $ \delta_{i,j} = 1  \bar{S} $  ：子集 S 的闭包  $ \partial S $  ： $ \bar{S} \setminus S $  cf. (9.2.4) X × S Y：超过 S 的纤维乘积 参见。符号 2.3.13  $ \mathbb{R}^n $ ：欧几里得 n 空间  $ \lim $ ：归纳极限  $ \lim $ ：射影极限 “ $ \lim $ ”：ind-object cf。 I §11“ $ \lim $ ”：pro-对象参见。 I §11  $ \{x_n\}_{n \in I} $ ：由  $ I $  索引的序列；   $ (I = \mathbb{N} \text{ or } \mathbb{Z})  x_n \xrightarrow{n} x $  ：序列  $ \{x_n\}_{n} $  收敛于 x  $ \{pt\} $  ：由单个元素组成的集合 □：表示平方是笛卡尔矩阵

\subsection{流形}
X, Y, ...：实流形或复流形  $ \delta $  或  $ \delta_X: X \hookrightarrow X \times X $ ：对角线嵌入  $ \tau: TX \to X $ ：A.2.1 的切向量束

 $ T_{M}^{*}X $ ：X 的子流形 M 的共形向量束，参见A.2.1  $ T_{X}^{*}X \simeq X $ ：零部分，标识为 X  $ E_{\gamma} $ ：赋予  $ \gamma $  拓扑的向量空间 E 参见III §5  $ X_{\gamma} $ ：具有  $ \gamma $  拓扑的空间 X 参见。 III §5  $ \varphi_{\gamma} $ ：连续映射 $ X \to X_{\gamma} $  参见。 III §5  $ f: Y \to X $ ：流形的态射  $ f|_{N}\colon N \to M $ ： $ f $ 、 $ (N \subset Y, M \subset X)  Tf $  引发的态射  $ TY \to T \times_{X} TX \to TX $  参见。 （4.1.8） $ T_{N}f $ ： $ T_{N}Y \xrightarrow{f_{N}} N \times_{M} T_{M}X \xrightarrow{f_{N_t}} T_{M}X $  参见。 (4.1.9)  $ T^{*}Y \xleftarrow{t_{f'}} Y \times_{X} T^{*}X \xrightarrow{f_{n}} T^{*}X $  参见。 (4.3.2)  $ T_{N}^{*}Y \xleftarrow{t_{f_N}} N \times_{M} T_{M}^{*}X \xrightarrow{f_{N_n}} T_{M}^{*}X $  参见。 （4.3.3） $ T_{Y}^{*}X $ ： $ \ker(t^{*}f' : Y \times_{X} T^{*}X \to T^{*}Y) $  参见。 (4.3.4)  $ a_{X} $  ：映射  $ X \to \{\text{pt}\}  \overline{X} $  ：与复流形相关的复共轭流形 X §1  $ X^{\mathbb{R}} $ ：复流形 X §1  $ \operatorname{dim} Y/X $  的实底层流形，  $ \operatorname{codim}_{X} Y $ ：相对维数和余维数 cf.符号 3.3.8  $ \operatorname{dim} X $  、 $ \operatorname{dim}_{\mathbb{R}} X $ ：X 的维数参见。 II §9 和符号 3.3.8（XI §3 除外）  $ \operatorname{dim}_{\mathbb{C}} X $ ：X 的复数维 cf.第二节 §9

\subsection{向量丛}
 $\tau: E \to Z$  ：  $Z$  上的向量丛  $\pi: E^{\ast} \to Z$  ： 对偶向量丛  $Z$  被识别为 $Z$  上向量丛的零部分  $\dot{E} = E \backslash Z, \dot{E}^{\ast} = E^{\ast} \backslash Z\dot{\tau} = \tau|_{\dot{E}}, \dot{\pi} = \pi|_{\dot{E}^{\ast}}a$  ：  $E$  上的对映图  $S^{a}$  ：  $S$  的图像 $a, S \subset ES^{\circ}$ ：极坐标设置为 $S$  cf。 (3.7.6)  $e$ ：欧拉向量场 cf. V §5  $A + B$ ：向量丛中的总和 cf. (5.4.5)  $S = \dot{E}/\mathbb{R}^{+}$ ：与向量丛  $E$  相关的球丛 参见(3.6.3)  $T^{*}(E/Z)$ ：相对余切丛 cf. (5.5.3)

\subsection{法锥}
 $ C_M(S) $ ：S 沿 M 的法向锥体 cf.定义4.1.1

 $ C(S_1, S_2) $ ： $ S_1 \times S_2 $  沿对角线的法锥体 cf.定义4.1.1

\[
C_{\mu}(A,B),\, f^{\#}(A,B),\, f_{\infty}^{\#}(A,B),\, f^{\#}(A),\, f_{\infty}^{\#}(A),\, A\, \widehat{+}\, B, \, A\, \widehat{+}_{\infty}\, B: \, \mathrm{cf.}\, \mathrm{Definition}\, 6.2.3
\]

\[
N(S) = TX\backslash C(X\backslash S,S): \, \mathrm{cf.}\, \mathrm{Definition}\, 5.3.6
\]

\[
N^{*}(S) = N(S)^{\circ} \colon \, \mathrm{cf.}\, \mathrm{Definition}\, 5.3.6
\]

\[
\widetilde{X}_{M}:\qquad\qquad\qquad\text{normal deformation of } M\text{ in } X\text{ cf. IV §1}
\]

\subsection{辛几何}
\begin{longtable}{@{}p{0.28\textwidth}p{0.66\textwidth}@{}}
$\sigma$ & 辛形式，参见 A §1\\
$\alpha_X$ 或 $\alpha$ & $T^*X$ 上的规范 $1$-形式，参见 A §2\\
$H$ & Hamilton 同构，参见 A §1\\
$H_f$ & Hamilton 向量场，参见 A §2\\
$\{\cdot,\cdot\}$ & Poisson 括号，参见 A §2\\
$\tau(\cdot,\cdot,\cdot)$ & 惯性指数，参见 A §3\\
$E^\rho$ & $\rho^\perp/\rho$，其中 $\rho$ 是 $E$ 中的迷向子空间，参见 A §1\\
$\lambda^\rho$ & $(\lambda\cap\rho^\perp+\rho)/\rho$，其中 $\lambda$ 为 Lagrange 子空间、$\rho$ 为迷向子空间，参见 A §1\\
$\lambda_1\circ\lambda_2$ & Lagrange 平面的复合，参见 A §1\\
$A_1\circ A_2$ & Lagrange 流形的复合，参见定义 7.4.5\\
$A_\varphi$ & 与 $\varphi$ 相伴的 Lagrange 流形，参见 (7.5.1)\\
$\lambda_\circ(p)=T_p\pi^{-1}\pi(p)$ & 参见 (7.5.2)\\
$\lambda_A(p)=T_pA$ & 切空间记号\\
$\tau_\varphi$ & 参见 (7.5.3)\\
$\tau(\lambda_1:\lambda_2)$ & 参见 (7.5.10)
\end{longtable}

\subsection{代数}
A：环（所有环都是单一的） A 模：左 A 模（所有模都是单一的）  $A^{\circ p}$  ： $A$  的相反环（ $A^{\circ p}$  -模是右 $A$  -模）  $\mathcal{C}$  ：一个范畴， $\mathrm{Ob}(\mathcal{C})$  、 $\mathrm{Hom}_{\mathcal{C}}(\cdot,\cdot)$  参见。 I §1  $\mathcal{C}^{\circ}$  ：相反的范畴参见。 I §1 id：恒等态射  $\mathcal{C}^{\vee}$  ：从  $\mathcal{C}^{\circ}$  到  $\mathcal{Set}\mathcal{C}^{\wedge}$  ： $\mathcal{C}^{\circ\vee\circ}$  Ker、Coker、Im、Coim 的函子范畴：参见I §2  $\mathbf{C}^{*}(\mathcal{C})$  :  $*= \emptyset, +, -, b$  :  $\mathcal{C}$  复合体的范畴 cf. I §3  $\mathbf{K}^{*}(\mathcal{C})$  :  $*= \emptyset, +, -, b$  : 参见。定义 1.3.4 Ht  $(X,Y)$  ：从  $X$  到  $Y$  同伦到零的态射群 cf. I §3  $\mathbf{K}(\mathcal{C})$  ：范畴  $\mathcal{C}$  的格洛腾迪克群 cf.练习 I.27  $X[k]$  ：翻译后的复数参见。定义 1.3.2  $\tau^{\leq}\tau^{\geq}$  ：截断函子 cf. （1.3.10）（1.3.11）。 X 81

 $ \tau^{\leq}, \tau^{\geq} $  ：截断函子参见(1.3.10), (1.3.11), X§1

 $ H^{k}(X) $  ：复数 X 的第 k 个上同调对象 cf.定义1.3.5

 $ M(f) $ ：f cf 的映射锥。 I §4  $ X \longrightarrow Y \longrightarrow Z \xrightarrow{+1} : \text{triangle cf. Notation 1.5.8}  \mathcal{C}_S $  ：由 S cf 进行的 C 本地化。定义 1.6.2  $ \mathcal{C}/N $ ：C 由 N 进行本地化 cf.符号 1.6.8  $ \mathbf{D}^*(\mathcal{C}), * = \emptyset, +, -, b $ ：导出范畴 cf.定义 1.7.1 RF：F 的右导函子 参见。定义 1.8.1 LF：F 的左导出函子 参见。 I §8  $ \mathbf{D}_c^*(\mathcal{C}) $  ： $ \mathbf{D}^*(\mathcal{C}) $  的满子范畴，由其上同调对象属于 C' 的复形组成。 I §7  $ H_I(X), H_{II}(X), s(X) $  ：与双复合体相关的复合体 cf. I §9  $ N \otimes_A M $  或  $ N \otimes M $  ：张量积（在环 A 上）  $ \operatorname{Hom}_A(N, M) $  或  $ \operatorname{Hom}(N, M) $  ：同态群（在环 A 上）  $ \operatorname{Tor}_n^A(N, M) = H^{-n}(N \otimes_A^L M) $  参见。示例 1.10.12  $ \operatorname{Ext}_A^n(N, M) = H^n(\operatorname{RHom}_A(N, M)) $  M-L：Mittag-Leffler 参见I §12  $ \oplus $  ：直和参见。 I §2  $ X \times_Z Y $  ：Z 上的乘积 cf.练习 I.6  $ X \oplus_Z Y $ ：Z 上的直和 cf.练习 I.6  $ \operatorname{Ext}^j(X, Y) = \operatorname{Hom}_{\mathbf{D}(\mathcal{C})}(X, Y[j]) $  参见练习 I.17  $ \operatorname{hd}(\mathcal{C}) $ ：C 的同调维数 cf.练习 I.17  $ \operatorname{gld}(A) $ ：A 的全局同调维数 cf.练习 I.28  $ \operatorname{wgld}(A) $ ：A 的弱全局同调维数 cf.练习 I.29  $ \operatorname{tr} $ ：trace cf.练习 I.32  $ \chi $ ：欧拉-庞加莱指数 cf.练习 I.32  $ \mathbf{b}_j(V) = \operatorname{dim} H^j(V) $  参见(5.4.17)  $ \mathbf{b}_j^*(V) = (-1)^j \sum_{k \leq j} (-1)^k \mathbf{b}_k(V) $  参见。练习 I.34 和 (5.4.18)

\subsection{层}
F, G, H, ... 层，R 空间 X Hom\_\{\textbackslash{}mathcal\{R\}\}(F, G) 或  $ \mathcal{Hom}(F, G) $  上的环层： $ \mathcal{R} $  层 - F 在 G 中的同态 参见定义 2.2.7 Hom(F, G) =  $ \Gamma(X; \mathcal{Hom}(F, G))  \mathcal{R}^{\mathrm{op}} $  ：对向环  $ F \otimes_{\mathcal{R}} G $  或  $ F \otimes G $  ：F 和 G 的张量积束（在  $ \mathcal{R} $  上） cf.定义 2.2.8  $ F|_{z} $  ：F 在 Z 上的逆像  $ F_{x} = F|_{\{x\}} $  ：F 在 x 处的茎  $ s|_{z}, s_{x} $  ：节 s 到 Z 的限制以及 s 在 x 处的萌芽 support(s)：节 s 的支持  $ \Gamma(X; F) $  ：F 在 X 上的全局节  $ \Gamma(Z; F) = \Gamma(Z; F|_{z})  f^{-1}F $  ：逆像层 F 参见定义 2.3.1  $ f_{*}F $  ：捆 F 的直像 参见。定义2.3.1

 $ f_{1}F $  ：具有适当支撑的直像参见。 (2.5.1)

 $f^*$  ：参见定义 2.7.4  $F_z$  ： $X$  上的束使得  $F_z|_{z}=F|_{z}, F_z|_{x\setminus z}=0$  、 $Z$  局部闭合 cf. II §3  $\Gamma_z(F)$  ： $F$  的子层，由 $Z$  支持的部分组成，参见。 II §3  $M_x=a_x^{-1}M$  ：  $X$  上的常值层，带茎  $M$  （M：  $A$  -模）  $M_z=(M_x)_z$  、  $Z\subset X$  、  $Z$  局部闭合  $F\boxtimes_S G$  ： 外部张量积（ $X\times_S Y$  上的束） cf.符号 2.3.12  $\Gamma_c(X;F)=a_{x!}F$  ：具有紧凑支持的全局部分 cf. (2.5.2)  $Rf_*$  、  $R\Gamma_z$  、  $R\Gamma(X,\cdot)$  、  $Rf_!$  、  $\otimes^L$  、  $\boxtimes^L$  、 RHom：前面的导出函子 cf. II §6  $H_z^j(F), H_z^j(X;F), H_c^j(X;F), \mathcal{E}xt_R^j(F,G), Ext_R^j(F,G), \mathcal{T}or_R^j(F,G)$  ：参见。符号2.6.8 wgld( $\mathcal{R}$ )：参见。定义 2.6.2 support(  $F$  ) =  $\bigcup_j$  support  $H^j(F)$  的闭包 参见(2.6.34) 和 II §2  $\mathcal{C}(\mathcal{U};F)$  :  $\mathcal{C}$  每个与开子集族相关的复合体参见。 II §8  $f^1$  ： $Rf_!$  的右伴随 cf. III §1  $D_xF = \mathbb{R}\mathcal{H}om(F,\omega_x)$  参见。定义 3.1.16  $D_xF = \mathbb{R}\mathcal{H}om(F,A_x)$  参见定义 3.1.16  $\int_X\colon$  态射  $H_c^n(X;or_x)\to A$  参见(3.3.15)  $\widehat{F}$：Fourier--Sato 变换，参见定义 3.7.8；$F^\vee$：逆 Fourier--Sato 变换，参见定义 3.7.8；$\Phi_K,\Psi_K$：与核 $K$ 相伴的函子，参见定义 3.6.1 和 7.1.3；$K_1\circ K_2$：核的复合，参见 (3.6.2) 和命题 7.1.2；$K_1\circ_\mu K_2$：核的微局部复合，参见定义 7.3.2；$\nu_M(F)$：$F$ 沿 $M$ 的特化，参见定义 4.2.2；$\mu_M(F)$：$F$ 沿 $M$ 的微局部化，参见定义 4.3.1；$f_\mu^{-1},f_\mu^!,f_\mu^*,f_!^\mu$：微局部操作，参见 VI §1；$\mu hom(G\to F),\mu hom(F\leftarrow G),\mu hom(F,G)$：微局部化函子，参见定义 4.4.1；$\mathrm{SS}(F)$：$F$ 的微支撑，参见 V §1；$\mathbf{D}^*(X)=\mathbf{D}^*(A_X)$（$*=\emptyset,+,b,-$）：$A_X$-模层范畴的导出范畴，参见 II §6；$N(X,Y;\Omega_X,\Omega_Y)$：核的范畴，参见定义 7.1.1；$\mathbf{N}(X,Y;p_X,p_Y)$：核的范畴，参见定义 7.3.7；$\phi_f$：消失循环函子，参见 VIII §6；$\Psi_f$：邻近循环函子，参见 VIII §6

\[
^{p}\dot{H}^{k}
\]

\subsection{特殊层}
 $ or_{x} $  ：定向层参见。定义3.3.3

 $ or_{Y/X} $  ：相对方向束，参见。 (3.3.3)

 $ \omega_{x} $  ：对偶复数参见。定义3.1.16

 $ \omega_{Y/X} $  ：相对对偶复数 cf.定义3.1.16

 $\mathcal{D}\mathcal{V}_{M}$  ：一堆分布参见。 II §9.6  $\mathcal{B}_{M}$  ：一堆功能亢进参见。 II §9.6 和定义 10.5.1  $\mathcal{V}_{M}$ ：流形上的密度束 cf. II §9.5  $\mathcal{C}_{M}$  ：一堆微函数 cf.定义11.5.1

\subsection{复流形上的层}
 $\mathcal{O}_{X}$  ：  $X$  上的全纯函数束  $\mathcal{O}_{X}^{(p)}$  ： 全纯  $p$  -形式  $\Omega_{X}$  ：  $\mathcal{O}_{X}^{(n)}\otimes\operatorname{or}_{X}(n=\operatorname{dim}_{\mathbb{C}}X)\mathcal{D}_{X}$  ： 有限阶全纯微分算子环层  $\mathcal{D}_{Y\to X}$  ： 从  $Y$  到  $Y$  的微分算子双模 $X$  参见。定义 11.2.8 Icar( $\mathcal{M}$ )：参见(11.2.7) char(  $\mathcal{M}$  )： $\mathcal{D}_{X}$  -模  $\mathcal{M}$  的特征变体 参见(11.2.14) $\mathcal{K}_{X}=\mathcal{D}_{X}\otimes_{\mathcal{O}_{X}}\Omega_{X}^{\otimes-1}[\operatorname{dim}_{\mathbb{C}}X]$  参见。 (11.2.18)  $\mathcal{N}^{*}=\mathrm{R}\mathcal{H}\mathcal{O}\mathcal{H}\mathcal{O}_{X}(\mathcal{N},\mathcal{H}_{X})$  参见。 (11.2.19)  $\underline{\mathfrak{d}}$  ,： $\mathcal{D}_{X}$  范畴中的外部张量积-模 cf. (11.2.21)  $\underline{f}^{-1}\mathcal{M}$  ： $\mathcal{D}_{X}$  范畴中的逆像-模 cf。定义 11.2.10  $\underline{f}_{\star}\mathcal{N},\underline{f}_{\star}\mathcal{N}$  ：  $\mathcal{D}_{X}$  范畴中的直接映像-模 cf.定义 11.2.10  $\mathcal{E}_{X}^{\mathbf{R}},\mathcal{E}_{Y\to X}^{\mathbf{R}},\mathcal{E}_{X\to Y}^{\mathbf{R}}$  ：全纯微局部算子的环和双模 cf.定义 11.4.2 和命题 11.4.3  $\mathcal{C}_{S|X}^{\mathbf{R}}$ ：复子流形上的微函数束  $S$  参见定义11.4.2

\subsection{范畴}
Set：集合的范畴 Ab：阿贝尔群的范畴 Mod(A)：左 A 模的阿贝尔范畴 Mod\textasciicircum{}\{f\}(A)：有限生成的左 A 模的范畴 Mod(R)：R 模的层范畴（R：X 上的环层） Mod(A\_\{x\})：X 上的 A 模的层范畴 Sh(X) = Mod(Z\_\{x\})：阿贝尔的层范畴X 上的群 D(A) = D(Mod(A)) cf.符号 1.7.14 和 2.6.1 D(X) = D(A\_\{x\}) = D(Mod(A\_\{x\})) 参见符号 2.6.11 D\textasciicircum{}\{b\}(X; \textbackslash{}Omega)：D\textasciicircum{}\{b\}(X) 在 \textbackslash{}Omega \textbackslash{}子集 T\textasciicircum{}\{*\}X 上的定位 参见VI §1 D\textasciicircum{}\{b\}(X; p) = D\textasciicircum{}\{b\}(X; \textbackslash{}\{p\textbackslash{}\})：D\textasciicircum{}\{b\}(X) 在 p 处的定位 参见VI §1 D\_\{R\}\textasciicircum{}\{+\}(E)：由圆锥形物体组成的 D\textasciicircum{}\{+\}(E) 的子范畴 cf.定义 3.7.1 Cons(S)、w-Cons(S)：S 上可构造和弱可构造层的范畴 参见。第八章 1

 $\mathbf{D}_{\mathbf{S}-c}^{b}(\mathbf{S})$  和  $\mathbf{D}_{w-\mathbf{S}-c}^{b}(\mathbf{S})$  ： $\mathbf{D}^{b}(|\mathbf{S}|)$  的子范畴，由具有可构造和弱可构造上同调的对象组成，参见第八章 1

 $\mathbb{R}-\mathbb{C}\mathrm{Cons}(X)$  和 $w\mathbb{-R}\mathbb{C}\mathrm{Cons}(X)$  ： $X$  上 $\mathbb{R}$  可构造和弱 $\mathbb{R}$  可构造层的范畴，参见 $X$ 。定义8.4.3

 $\mathbf{D}_{\mathbf{R}-c}^{b}(X)$  和  $\mathbf{D}_{w-\mathbf{R}-c}^{b}(X)$  ： $\mathbf{D}^{b}(X)$  的子范畴，由具有  $\mathbb{R}$  可构造和弱  $\mathbb{R}$  可构造上同调的对象组成，参见第八章 4

 $\mathbf{D}_{\mathbf{C}-c}^{b}(X)$  和 $\mathbf{D}_{w-\mathbf{C}-c}^{b}(X)$  ：参见。第八章 5

 $K_{\mathbf{R}-c}(X)$  ：  $\mathbf{D}_{\mathbf{R}-c}^{b}(X)$  的格洛腾迪克小组 参见。九§7

 $^{p}\mathbf{D}_{w-\mathbf{R}-c}^{≤0}(X)$  、 $^{p}\mathbf{D}_{w-\mathbf{R}-c}^{≥0}(X)$  等： $t$  -与反常性 $p$  相关的结构参见。 X §2 和 §3

 $\mu\mathbf{D}_{w-\mathbf{R}-c}^{≤0}(X)$  、 $\mu\mathbf{D}_{w-\mathbf{R}-c}^{≥0}(X)$  等： $t$  -微局部定义的结构参见。 X§3

\subsection{圈、迹与构造函数}
\[
\chi(F)(x),\chi_{c}(F)(x),\chi(X;F),\chi_{c}(X;F)\colon\text{Euler-Poincaré indices cf. IX } \S1
\]

\[
\text{tr}_{X}\colon\qquad\qquad\qquad\text{trace morphism cf. } (9.1.3)
\]

\[
\int_{X}\colon\qquad\qquad\qquad H_{c}^{0}(X,\omega_{X})\to k\text{ cf. III } \S3\text{ and IX } \S1
\]

\[
\mathcal{C}\mathcal{S}_{p}(F)\colon\qquad\qquad\qquad\text{取值于 }F\text{ 的次解析 }p\text{-链层，参见定义 }9.2.1
\]

\[
\mathcal{C}\mathcal{S}_{p}^{X}=\mathcal{C}\mathcal{S}_{p}(A_{X})=\mathcal{C}\mathcal{S}_{p}
\]

\[
\partial_{p}:\mathcal{C}\mathcal{S}_{p}\to\mathcal{C}\mathcal{S}_{p-1}\colon\text{边界算子，参见 }(9.2.10)
\]

\[
\mathcal{L}\mathcal{S}_{p}^{X}\colon\qquad\qquad\qquad\text{次解析 }p\text{-圈层，参见定义 }9.2.5
\]

\[
C_{1}\cap C_{2}\colon\qquad\text{两个圈的交，参见定义 }9.2.12
\]

\[
\#(C_{1}\cap C_{2})\colon\text{两个圈的交数，参见定义 } 9.2.12
\]

\[
\mathcal{L}_{X}\colon\qquad\qquad\qquad\text{Lagrange 圈层，参见定义 } 9.3.1
\]

\[
f^{*},f_{*}\colon\qquad\qquad\qquad\text{Lagrange 圈的逆像与直像，参见定义 }9.3.3
\]

\[
\boxtimes\colon\qquad\qquad\qquad\text{圈的外积，参见 }(9.3.2)
\]

\[
[T_{Y}^{*}X]\colon\qquad\qquad\qquad\text{与 }Y\hookrightarrow X\text{ 相伴的 Lagrange 圈，参见例 }9.3.4
\]

\[
\left[\sigma_{0}\right]\colon\qquad\qquad\qquad\text{与零截面相伴的圈，参见定义 }9.3.5
\]

\subsection{索引}
加法（范畴、函子）26

伴随函子（右、左）69

解析集 (C-) 344

反杰出三角形 40

对映图 169

关联坐标系（余切丛上）481

- 捆（到预层）86

- 简单复杂（到双重复杂）55

双特征（曲线、叶子）271、482

双圆锥组 242

- 捆 168

双功能56

边界算子 369

- 值 467

有界复数（从下到上）31

罐头352

规范 1-形式 481

笛卡尔平方 106

第23类

柯西-科瓦列夫斯基定理 453

柯西问题 454

柯西留数公式 183

柯西-黎曼系统 468

C*-圆锥曲线 344

特征等级（捆，φ） 362, 390

- 周期 377

-（Ω模的）簇449

切赫上同调 125

干净（映射）190

循环条件 487

共决赛81

相干 443

上同调维数（左精确函子）75

- 函子 39

上同调可构造复形 158

上同调 32

共图 28

焦仁28

复合体 30

圆锥曲线集（向量丛） 169, 483, 344

– 捆 167

共正规束 481

常值层90

可构造函数 398

– 束（R-、C-、S-）339、347、322

接触转化 483

凸集（向量丛中）169

卷积（层、可构造函数） 135, 409

逆变 25

余切丛 481

协变25

c-软尺寸 133

C-软轴104

杯产品 134

曲线选择引理 327

衍生范畴 45

– 函子（右、左）50, 52

去奇异化定理 328

（复数的）微分 31

微分算子（环） 448

有向有序集 63

直像（2-模）453

–（微功能）470

– （一个层）90

– 有适当的支撑 103

– 总和 27

– 被加 70

杰出的三角形35

多尔博复合体 128

双复合体54

双（层复合体）148

对偶复合体（相对）148

椭圆系统 468

足够的单射（投射）48

外态 24

范畴 25 的等效性

欧拉179级

欧拉态射 406

欧拉-庞加莱指数 240, 361

欧拉向量场 241, 482

精确（左、右）函子 29, 30

- 序列 29

扩展空间395

扩展接触变换 292

外张量积 97

F-无环 75

F-内射 50

F-投影 52

过滤环（和模） 445、446

过滤剂范畴62

最终对象24

更精细（覆盖）334

五个引理 72

松弛尺寸98

- 捆 98

扁平模77

- 捆 101

傅里叶佐藤变换172

f-软束 140

满子范畴 24

完全忠诚 25

函子 24

基础课445

γ-封闭 161

γ-开放 161

γ-拓扑 161

广义特征空间 395

通用 480

（部分的）胚芽 84

全局同源维度 77

良好的过滤446

分级环447

格洛腾迪克组 77, 399

吉辛映射 179

半双特征曲线271

哈密顿同构 258, 478

哈密顿向量场 482、469

心411

希尔伯特对称定理 443

全纯函数 128

全纯微分算子 449

完整（D 模）450

齐次辛 482

（范畴的）同源维度 75

同伦（循环）371

-（映射）119

-（态射）31

–（层）246

霍普夫指数定理 408

双曲系468

图 28

索引 361

工业对象 62

诱导柯西-黎曼系统 474

归纳系统 61

惯性指数 486

无限阶微分算子（环） 462

初始对象 24

内射对象 30

单射子范畴（关于函子） 50

逆傅立叶-佐藤变换 172

–（D 模的）图像 453

–（微功能的）图像470

–（一个层的）图像 91

对合 271, 478, 481

对合性定理 272

（对象、函子的）同构 24

–（在余切丛的子集上）221

各向同性 478, 481

– 集合（亚解析，C 解析）331, 344

内核 27, 164

昆尼斯公式 135

拉格朗日 478, 481

拉格朗日链380

拉格朗日循环373

拉格朗日格拉斯曼流形 480

拉格朗日集（亚解析、C 解析）331、344

拉普拉斯算子 469

莱布尼兹公式 449

Lefschetz 定点公式 392

莱夫谢茨不动点定理 390

勒让德变换（部分）318

勒雷无环定理 125

（范畴的）本地化 41

局部上同调平凡 178

局部圆锥 (R⁺, C⁺) 483, 344

– 常值层90

测绘锥 34

马斯洛夫指数 487

迈尔-维托里斯序列 115

微微分算子 462

微微分算子 462

微双曲系统468

微局部贝尔蒂尼-萨德定理 332

内核的组成 294

- 截止引理（对偶，精化）225、252、253、254

-（正确的）直像 255

- 逆像255

微局域莫尔斯引理 239

微局部运算符（环） 462

微局部化 198

微支撑221

中变态 426

米塔格-莱弗勒条件 64

模（在环层上）87

单态性 24

单一性 351

态射23

莫尔斯函数（相对于各向同性子集）388

莫尔斯不等式239

 $ \mu $  -条件 334

μ过滤337

微米202

 $ \mu $  -分层 334

乘法系41

邻近循环函子 350

诺特束 443

Noetherian过滤环446

非典型 235, 262

非特征变形引理 117

非特征（D 模）453

V 262 上 A 的非特征

普通捆绑包 185, 481

- 变形186

零系统 43

对象 23

相反范畴 24

- 环 29

顺序（在过滤环中）447

定向层（相对） 126, 153

超紧空间 102

完美（A 模复合体）78

反常层 427

反常性 419

庞加莱-维迪埃对偶性 140

泊松括号 478、482、447

极集 170

呈示（s-、有限自由、自由...）443

预层83

主符号 449

射影对象30

–（系统，限制）61

pro-对象62

奇点的传播 468

真圆锥（在向量丛中）169

- 映射103

真同伦映射 119

\par\noindent 量子化接触变换 465
拟逆 25

准同构40

常规对合 483

相对余切丛 238, 241

可表函子 25

代表对象 25

残基态射 182

德拉姆综合体 127

环（束）87

S-无环（束）324

佐藤功能亢进 127, 130, 466

佐藤微功能 466

施瓦茨分布 127

（预层的）第 84 部分

捆 85

转变（功能）31

缩小空间395

签名（二次形式）487

单捆312

单纯复形 321

单纯形 321

单一支持 467

软捆 132

专业化 191

分割（序列） 70

堆栈 424

（预层的）茎84

分层334

第334层

严格精确 $ \quad 446 $ 

亚分析链366

- 周期 369

- 设置 327

子范畴 24

软束 132

(一个层、一段的)支撑 85, 116

辛形式 477

- 向量空间 477

切丛 481 顶点 321 厚子范畴 49 Victoris-Begle 定理 121 Thom 类 179 拓扑淹没 151 迹 79, 361 横截（图） 190 三角形 35 三角形不等式 222 三角范畴 38 三角剖分定理 328 截断复数 33 T 结构 411消失环函子 351 var 352 波算子 469 弱全局维数（一个环或环层） 78, 110 弱可构造（R-, C-, S-） 339, 347, 322 魏尔斯特拉斯除法定理 458 惠特尼条件（(a) 和 (b)） 357 米田扩展 81 扎里什过滤环第446章
\end{document}
